\documentclass[11pt,a4paper]{article}
\usepackage[T1]{fontenc}
\usepackage[utf8]{inputenc}
\usepackage{lmodern}
\usepackage[margin=26mm]{geometry}
\usepackage{amsmath,amssymb,amsthm,mathtools}
\usepackage{mathrsfs}
\usepackage{array,longtable,booktabs}
\usepackage{microtype}
\usepackage{xcolor}
\usepackage{hyperref}
\hypersetup{colorlinks=true,linkcolor=blue!45!black,urlcolor=blue!55!black,
 pdftitle={MAS ARTIN: geometric traces, the S5 gate, and PF26},
 pdfauthor={MAS ARTIN research audit}}
\setlength{\emergencystretch}{3em}
\setlength{\parskip}{4pt}
\setlength{\parindent}{0pt}
\makeatletter
\renewcommand*\l@section[2]{\ifnum\c@tocdepth>\z@\addpenalty\@secpenalty\addvspace{1.0em \@plus\p@}\setlength\@tempdima{3.0em}\begingroup\parindent\z@\rightskip\@pnumwidth\parfillskip-\@pnumwidth\leavevmode\bfseries\advance\leftskip\@tempdima\hskip-\leftskip#1\nobreak\hfil\nobreak\hb@xt@\@pnumwidth{\hss#2\kern-\p@}\par\endgroup\fi}
\renewcommand*\l@subsection{\@dottedtocline{2}{1.5em}{3.6em}}
\makeatother
\newcommand{\RunStatus}{Active research checkpoint; open gates explicitly retained}
\title{MAS\_ARTIN\\[5pt]\large Masullo's central trace identity over
$\mathbb Q$, the $S_5$ Artin gate, and the analytic premise of PF26}
\author{Research audit and new deductions\\GPT-6 Astra Pro, MAS\_ARTIN run}
\date{22 November 2026\\\small\RunStatus}
\begin{document}
\maketitle
\begin{abstract}
We audit the specified source arguments and carry the two predecessor
runs forward by statement ID. For the central trace method we locate
a missing finite-prime half-weight, an invalid cutoff codomain, and
a path-dependent annular local limit, and prove local repairs. We
construct the actual spherical theta--M\"obius intertwiner, prove
escape of its finite Euler tails in the weighted Hilbert norm, and
show that its inverse has no bounded extension between matching
narrow-strip norms. The raw Euler-conjugated cutoff is also unbounded on every
narrow one-sided zeta, Dirichlet or nontrivial imaginary-quadratic
Hecke line; the proof uses
nonzero universal values and actual core norm estimates.
An exact polar comparison then gives genuine rational smooth
theta inputs with both relevant moments zero and unbounded
literal geometric diagonal outputs, all in the smooth space.
The original rational smaller common-sum space also fails
raw invariance. On its exact pole-free domain an actual
smoothed trace-class map exists in two stated analytic regimes,
but its closed domain has codimension two and every nonzero
smoothing exits that domain with rank two. The required
orthogonal correction changes the action. No nonzero
translation-invariant subspace lies in that domain, so it
cannot be the target of a nonzero equivariant repair.
For nontrivial finite-order character sectors, an unconditional
smooth-ideal division theorem instead identifies the actual
common-sum action on the whole smooth theta ideal. Its smoothing
extends to a trace-class endomorphism of the closed theta ideal
in the same stated regimes; its trace retains the actual ideal
projection. Actual cross-weight maps, finite theta Gram data and
a Blaschke kernel formula describe that projected trace. A
stronger nuclear return to the actual smooth local source then
evaluates it by the local Fourier multiplier, without dropping
the ideal projector. The two one-sided traces equal that value,
so the explicit Blaschke correction vanishes; omitting the
projector would instead double a generally nonzero trace.
It is not the full quotient trace. In the
rational principal sector a different completion retains both
actual polar source sections and gives an equivariant exact
sequence and a trace-class positive-source column. Its possible
outgoing quotient defect is explicitly finite dimensional; the
local source endomorphism and the full raw matrix remain distinct.
Under an explicitly used zero-free right half-plane, the full
rational matrix extends to a trace-class endomorphism of this
enlarged common completion. Its positive source space is not
invariant: every nonzero compact test has outgoing rank two.
An actual invariant relative operator instead gives trace and
Fredholm determinant comparisons. For value-zero tests the raw
matrix trace is precisely the omitted finite-prime sum, with
both uses of zero location in that comparison explicit.
The polar obstruction itself is subsequently made
unconditional using a published positive-proportion theorem
for known critical-line zeros. The actual full character
matrix and its outside-$S$ prime trace are also constructed
under their explicitly stated dual half-plane premise.
Conversely, any strictly exterior character zero prevents a
bounded full common-space action after every nonzero compact
smoothing, as proved on actual translated smooth sources.
An unretained boundary multiplicity gives the same obstruction,
including the equality threshold in the polynomial weight.
The corresponding necessity theorem on the rational augmented
completion retains both actual polar source sections.
Full-matrix sufficiency also holds for finitely many fully retained
boundary zeros when the remaining zeros have a fixed interior margin;
the proof uses the actual one-sided boundary distribution of the
wide inverse. Separately, a polynomially controlled contour strictly
inside the strip removes the Gevrey and generic-width restrictions
from the completed-ideal nuclear source return. That return and its
local trace now hold at every positive narrow width for every compact
smooth test, without a zero-location premise. This restricted-source
result does not extend the unrestricted one-sided inverse or prove
full-matrix sufficiency without a margin.
The full matrix does extend when total multiplicity in fixed-size
rectangles near the edge is uniformly bounded and the polynomial
weight exceeds the explicit whole-cluster threshold. A separate
criterion identifies polynomial reciprocal growth as exactly the
condition for unrestricted one-sided bounded smoothing throughout
a stated class of compact smooth cutoff perturbations. With a
predetermined physical cutoff, the same criterion is equivalent
to bounded smoothing for all tests on one fixed compact support. Its failure
would give actual normalized core countervectors, although the
finite-frequency traces still converge. The arithmetic reciprocal
condition is not decided independently of zero location; $A$
supplies the rational positive branch.
For the actual full common source, boundedness throughout a fixed
compact-test family forces a reciprocal estimate uniform in the
inward distance to the edge. Its failure would give fixed compact
source witnesses and residual bad tests; the rational version
retains both polar coordinates. A polynomial budget on total
cluster multiplicity and positive inward depth supplies an actual
full trace-class matrix, even for growing local counts. Separate
upper potential and lower whole-cluster estimates control the
original canonical quotient norms. Their arithmetic geometric
premises are not asserted. An exact localized negative-Sobolev
boundary condition then characterizes the universal bounded full
source family and automatically supplies trace class for all
compact tests under the containing-strip premise. Its arithmetic
polynomial norm bound remains unresolved.
A separate argument repairs trace class for
the actual convolution-after-projection operator under the precise
interior-strip hypothesis of MT-M26. An explicit one-sided Euler
commutator carries the omitted-prime terms in Meyer's actual common
space. An order-zero nuclear smoothing map then transports the
Hilbert quotient trace to that geometry under the same interior-strip
hypothesis. The independent narrow argument remains open at that
zero-location premise. A Hardy model first computes the partial
quotient at polynomial weight zero. Bounded convolution by the
smooth block idempotents then identifies the actual quotient and
its ambient cutoff trace at every polynomial weight, including all
surviving boundary jets, without a containing-strip hypothesis.
A closed-multiplier regularization gives an actual narrow
Euler commutator on generic lines with Gevrey tests, retaining
the crossed-zero and pole contributions. An actual partial nuclear
return to Meyer now transports every retained block at generic
widths; Gaussian regularization extends the trace transport to
all compact smooth tests, with a constructed discarded-divisor
correction. At arbitrary widths an explicit local eligibility condition
and actual finite-block costs characterize continuous returns; adaptive
regularizers give nuclear returns for every eligible test. An exact
polynomial bound on the canonical whole-quotient block costs,
together with absence of proper boundary truncation, characterizes
unrestricted returns for every smooth test. These arithmetic norm
conditions are not assumed to hold at an arbitrary width. An exact
positive resolvent defect also characterizes trace-norm convergence
of the actual ambient PW smoothings by the same arithmetic deficiency. An explicit local quotient Fourier transform carries
the finite prime terms, including character, conductor, and
different factors. The finite global Paley--Wiener projection
admits an exact Gram formula and a regularized cyclic trace
transport with a finite rational-dilation kernel. At polynomial
weight zero, a fixed trace-class Hankel tail proves the actual
slow-band joint asymptotic, including the character sectors. Its
difference from the local geometry is the upper exterior divisor.
A separate causal-range and meromorphic-resolvent proof establishes
the actual finite-band quotient at every polynomial weight, all
boundary jets, ambient trace class, and both iterated limits.
Nested actual PW ranges then prove a uniform causal slow-band
tail at every polynomial weight. An actual weighted adjoint defect
computes the negative-time limit, and a uniform resolvent estimate
handles the complete norm-one-vanishing test class. An actual dual-sector Fourier covariance argument extends the local
limit to arbitrary thin source profiles, and one regularizer works
for every compact smooth test along each prescribed cutoff path.
For every fixed number field and structural set of places, an
actual flat-source section, nuclear return and two weighted
commutator estimates now give the unchanged physical annular
cutoff rates, including all class groups and unit ranks.
Actual source counterexamples also refute the canonical theta-map
intertwining of the physical Fourier cutoff with the global
orthogonal PW projection. They cover the primitive odd character
modulo four and the principal rational sector, in the specified
corrected metric and in the exact printed twisted realization.
The principal witnesses are flat and preserve both polar conditions
before and after the cutoff. The non-induced cubic sector over
$\mathbb Q(\zeta_3)$ is included with the different prime retained
and both finite Fourier shells in the physical cutoff. These
counterexamples use no zero-location information
and do not refute the different Gram-corrected maps already constructed.
The vanishing of the discarded correction remains the exact zero-location gate. An actual joint-cutoff path
satisfying the printed thin-smoothing conditions has a nonzero
Gram remainder; it uses a retained critical zero, with that
classical zero-location input recorded. The positivity
audit separates ordinary positive cutoffs from faithful positive
full-jet actions. In Calegari's family, under the stipulated
$(\mathcal R_B)$, we prove the central equality $\beta=b+n$.
The proof extracts an integer positive-definite function from the
signed discrepancy divisor, then uses the coset-ring theorem and
prime-frequency independence to exclude the positive-excess branch.
Without assuming that real-zero hypothesis, failure of comparison
forces a shared real zero above the center, with normal-closure
multiplicity at least nine and an explicit discriminant-dependent
endpoint exclusion, sharpened by the reflected real zero, all gamma
terms, and a discriminant-aware lower contribution from every prime.
A further finite envelope of entire relative Dedekind zeta
functions strictly reduces that interval. Its exact completed
logarithmic derivatives remove the base-field poles and zero
divisors while retaining the actual comparison multiplicity.
The new minima may occur at one; explicit finite Euler tails
provide certified upper bounds without a zero-location premise.
This reduction does not exclude the remaining shared zero.
Boundary positivity alone does not supply
the ordinary Fourier identity needed to exclude that branch.
The general Hecke zero-location input remains necessary for the
stated all-factor GRH consequence. PF26's invoked algebraic
proofs pass a full reading; the bad-prime reduction is complete
relative to that input, and its residual analytic family is reduced
to explicit single Kummer Hecke functions. The four-way input register
covers 563 inherited findings with all secondary premises retained.
No unresolved assertion is promoted to a theorem in this paper.
\end{abstract}
\begingroup
\small
\setlength{\parskip}{0pt}
\tableofcontents
\endgroup
\newpage

\section*{Reading guide, conventions, and provenance}
\addcontentsline{toc}{section}{Reading guide, conventions, and provenance}
The user-specified theorem $A$ is Alp\"oge's Dirichlet GRH theorem;
its exact public step and predecessor-inspected Lean exports are
recorded in MA-IN-007. The symbol $M$ is Masullo's conditional
claim for Hecke functions over all number fields, draft
\S\S1.2--1.2.3, pp.~12--16. The empty Theorem 13.3.1.2 on p.~522
is not a theorem input. In the $S_5$ section only, $A(s)$ is also
the comparison $L$-function; there the Dirichlet theorem is
written $\mathsf A$.

Source page numbers for Masullo refer to the 940-page public PDF;
its printed and PDF page numbers agree. The retrieval hash and
source URL are recorded in the evidence receipt. No copy of that
PDF, extracted text, or rendered source page is included here.
The predecessor packages are GRH\_NIGHT checkpoint~07 and
ALPOGE\_BRIDGES checkpoint~09, as attached to the request.
Their established results are cited by ID and not silently
reclassified as new findings. A proof reading is not a formal
verification; a finite calculation verifies exactly its stated
certificate, not the surrounding conjecture.

Every new finding has a recorded whole-proof reading or an
explicit source-audit check in \texttt{CHECK\_EVENTS.tsv}.
The full table preserves who performed each event, the method,
the date, and the exact source locator. The model is GPT-6 Astra
Pro; the inherited reasoning-effort setting is not separately
exposed, so no unsupported effort label is supplied. Primary
bibliographic data and exact locators are given in the relevant
sections. Archive scan locators carried from the attached runs
are expressly distinguished from fresh source inspection.

\part{Checked findings and fully stated conditional deductions}
\section{The corrected identity and a prime-carrying smooth transport}
\label{sec:trace-global}

\subsection{The assertion, the actual objects, and the zero-location ledger}

Throughout this section, a statement attributed to a predecessor is used
with its recorded hypotheses. In particular, a theorem about convolution
on a quotient is not silently changed into a theorem about an arbitrary
lift of that convolution to its ambient Hilbert space.

Write $C_{\mathbb Q}=\widehat{\mathbb Z}^{\times}\times\mathbb R_{>0}$,
with probability Haar measure on the first factor and $dt/t$ on the
second. Put $x=\log t$, and let $\mathcal E$ be the space of smooth
functions on $\mathbb R$ all of whose derivatives decay faster than
$e^{-N|x|}$ for every $N$. Its Mellin transform is
\[
 \mathcal Mf(z)=\int_{\mathbb R}f(x)e^{zx}\,dx.
\]
For the trivial character let $\mathcal I$ be the closed convolution
ideal of GRH\_NIGHT MT-M06/MT-M18. MT-M24 identifies it, without
zero-location input, with the actual theta-summation image. The
centering operator is $T_{1/2}f(x)=e^{x/2}f(x)$. The corrected strip
space is
\[
 H_{\alpha,s}=L^2\bigl(\mathbb R,
      e^{2\alpha|x|}(1+x^2)^{s/2}\,dx\bigr),\qquad
 J_{\alpha,s}=\overline{T_{1/2}\mathcal I}^{\,H_{\alpha,s}},
 \qquad Z_{\alpha,s}=H_{\alpha,s}/J_{\alpha,s}.
 \tag{MA-TR-1}
\]
Here $\alpha>0$; a fixed $s\ge3$ is compatible with the source
geometric-kernel convention. Let $P_G$ and $P_Z$ be the orthogonal
projections onto $J_{\alpha,s}$ and its orthogonal complement. These
are Masullo's Definition 10.4.1.1, p.~456, on the corrected closed
subspace. For $h\in\mathcal D^0:=\{h\in C_c^\infty(\mathbb R_{>0}):
h(1)=0\}$, set
\[
 \widehat h(z)=\int_0^\infty h(t)t^z\frac{dt}{t},\qquad
 C_hf(x)=\int_{\mathbb R}h(e^y)f(x-y)\,dy.
\]
The source's convolution-after-projection operators are
$S_h=C_hP_Z$ and $G_{\rho,h}=M_\rho C_hP_G$; see Definition
10.4.1.5, p.~458, and (10.4.3.45), (10.4.3.48), p.~473.
The other order, $P_ZC_h$, is explicitly distinguished in
(10.4.1.10)--(10.4.1.11), pp.~457--458.

The corrected indentation may be taken to be the indicator of
$\mathbb R_{>0}\times\widehat{\mathbb Z}^{\times}$ inside the ideles,
as in ALPOGE\_BRIDGES F-SM-20. Its compact compression is valid, but
does not identify the compressed global operator with a separately
constructed local one. The closure and cutoff corrections are precisely
GRH\_NIGHT MT-M31--MT-M35; no injectivity of an indicator multiplier
or interchange of a cutoff image with a cofinal intersection is used.

In these normalizations put
\[
\begin{split}
 P(h)&=\widehat h(1/2)+\widehat h(-1/2),\\
 A_0(h)&=\int_0^\infty
   \frac{e^{-x/2}}{1-e^{-2x}}
       \bigl(h(e^x)+h(e^{-x})\bigr)\,dx,\\
 B(h)&=\sum_p\sum_{k\ge1}\frac{\log p}{p^{k/2}}
                   \bigl(h(p^k)+h(p^{-k})\bigr),\\
 \tau_\alpha(h)&=\sum_{|\Re\rho-1/2|<\alpha}
                         m_\rho\widehat h(\rho-1/2).
\end{split}
\tag{MA-TR-2}
\]
All the displayed integrals and sums converge absolutely. The prime
sum is finite for each compact test. Smoothness and $h(1)=0$ remove
the archimedean singularity. The spectral convergence uses only the
classical zero count and rapid vertical decay of $\widehat h$.

There is an essential distinction between the open-strip distribution
$\tau_\alpha$ just defined and an already established trace on
$Z_{\alpha,s}$. MT-M26 identifies the actual quotient trace when
\emph{all} zeros lie in a closed substrip of its interior. This holds
unconditionally for $\alpha>1/2$, and for every $\alpha>0$ under
$A$ in the present rational case. It does not, as stated, prove an
unconditional trace-class theorem for a partial strip. That additional theorem is proved later in MA-WP-001--003,
with its full boundary multiplicity convention; it is not
imported from MT-M26. At a boundary zero the polynomial
weight also retains some continuous jet functionals, as MA-TR-103
below makes precise. Thus a literal identification with an open-strip
sum would require a boundary convention or a boundary-free width.

The requested central geometric identity is
\[
 \boxed{\displaystyle
  \lim_{\rho\to1}\operatorname{Tr}G_{\rho,h}
     = A_0(h)+B(h)-P(h)=-\tau_\alpha(h)
       \quad(\alpha>0,\ h\in\mathcal D^0).}
 \tag{MA-TR-3}
\]
The first equality is the geometric transport assertion. The last is
the narrow spectral identity, with the source's sign convention.
Where the actual Hilbert trace is invoked, the preceding scope and
boundary qualification are part of the assertion. Proving it without
zero-location input includes, rather than presupposes, those missing
operator and comparison statements.
The local theorem alone gives $A_0(h)$ for $S=\{\infty\}$, before
the two removed moments, and gives $A_0(h)$ plus the primes in $S$
for larger $S$. It does not give the missing identification of global
operators. The exact final source location is the use of Theorem
12.4.1.1 in the proof on p.~520, followed by (12.5.1.4). Its proposed
bridge is Section~12.3, p.~514. The latter contains no proof.

The classical full explicit formula, in the conventions already
checked by F-SM-13 and MT-M30, is
\[
 A_0(h)+B(h)=P(h)-\tau_{\rm full}(h).
 \tag{MA-TR-4}
\]
Consequently the exact unconditional discrepancy between the numerical
geometric expression and the proposed restricted spectral expression is
\[
 A_0(h)+B(h)-P(h)+\tau_\alpha(h)
   =-\!\sum_{|\Re\rho-1/2|\ge\alpha}
                    m_\rho\widehat h(\rho-1/2).
 \tag{MA-TR-5}
\]
The equality criterion for these restricted tests is already proved in
MT-M30, (M30.14); it is not reproved or claimed as new here. Theorem
$A$ makes (MA-TR-5) zero for the rational and Dirichlet cases. Using
that fact in a proposed independent proof of (MA-TR-3) would be
circular with respect to the method being tested, even though it is a
valid application of an external theorem.

\begin{center}
\begin{tabular}{p{0.51\linewidth}p{0.40\linewidth}}
\hline
Step & Zero-location input \\
\hline
Theta image, closed ideal, compact averaging and valid indentation &
None; MT-M18/M24 and F-SM-20.\\
Local Euler factors, the corrected Poisson calculation, and the local
Fourier cutoff calculation & None.\\
Full spectral trace on a strip containing $[0,1]$ in its interior &
Classical containing strip only; MT-M26/M27.\\
Faithful spectral trace on every narrow strip & $A$ in the Dirichlet
case; $M$ in the general Hecke case, by MT-M26/M28.\\
Replacing the full zero sum by every narrow sum & Exactly the
zero-location assertion; circular in an independent proof.\\
Global-to-local transport of the specified cutoff kernels &
Not supplied by any of the foregoing.\\
\hline
\end{tabular}
\end{center}

\paragraph{MA-TR-103 (the endpoint jet convention).}
For the exact weight in (MA-TR-1), the functional
$f\mapsto\int x^j f(x)e^{\lambda x}\,dx$, initially on
$C_c^\infty(\mathbb R)$, extends continuously to $H_{\alpha,s}$
if and only if
\[
 |\Re\lambda|<\alpha,
 \quad\hbox{or}\quad
 |\Re\lambda|=\alpha\ \hbox{ and }\ 2j<s-1.
 \tag{MA-TR-12}
\]
It never extends when $|\Re\lambda|>\alpha$.
\emph{Proof.} By the weighted $L^2$ duality theorem, extension
is equivalent to finiteness of
\[
 \int_{\mathbb R}|x|^{2j}
        e^{2\Re\lambda x-2\alpha|x|}(1+x^2)^{-s/2}\,dx.
\]
The interior and exterior cases follow from exponential decay
and growth, respectively. At an endpoint one tail is comparable
to $\int_1^\infty x^{2j-s}\,dx$, which is finite precisely
under the displayed strict inequality. Density of compact smooth
functions proves necessity as well as sufficiency.
\hfill$\square$
This is the two-sided strip version of the dual-weight test in
GRH\_NIGHT MT-M25, not a proof of a complete partial-strip
spectral synthesis theorem. In particular, for $s\ge3$ every
boundary zero still imposes at least its value constraint on
$J_{\alpha,s}$.

\subsection{An explicit transport before the Hilbert completion}

\paragraph{MA-TR-101 (theta--M\"obius transport and the escape of its Euler tails).}
Let
\[
 \mathscr S_{0,+}=\{f\in\mathscr S(\mathbb R):
          f(-t)=f(t),\ f(0)=0,\ \int_{\mathbb R}f(t)\,dt=0\}.
\]
Associate to $f$ the radial function $F(t)=2\sqrt t\,f(t)$ for $t>0$.
The map
\[
 (\mathcal ZF)(t)=\sum_{n\ge1}n^{-1/2}F(nt)
                 =2\sqrt t\sum_{n\ge1}f(nt)
 \tag{MA-TR-6}
\]
is the centered global theta map in the trivial compact-character
sector. It is a continuous linear isomorphism from the displayed
Schwartz space onto $T_{1/2}\mathcal I$, with its smooth topology.
Its inverse on that image is explicitly
\[
 (\mathcal Jg)(t)=\sum_{n\ge1}\mu(n)n^{-1/2}g(nt).
 \tag{MA-TR-7}
\]
Both formulas converge with all derivatives locally uniformly for
$t>0$. They commute with norm translations on the $F$ variable
(the corresponding action on $f$ is
$f(t)\mapsto e^{-y/2}f(e^{-y}t)$). Thus there is a genuine
prime-carrying smooth transport; it has not yet been transported to
the source's narrow Hilbert kernels.

More precisely, for a finite set $S$ of rational primes define
\[
 P_S=\prod_{p\in S}(1-p^{-1/2}U_{-\log p}),\qquad
 (U_y g)(t)=g(e^{-y}t).
\]
Then
\[
 P_S\mathcal ZF(t)=
       \sum_{\substack{n\ge1\\(n,\prod_{p\in S}p)=1}}
                    n^{-1/2}F(nt).
 \tag{MA-TR-8}
\]
Suppose $F\in C_c^\infty([a,b])$, $0<a<b$, and
$\widehat F(1/2)=0$. Let $\rho$ be any nontrivial zeta zero with
$\lambda=\rho-1/2$ and $|\Re\lambda|<\alpha$. One may choose such
an $F$ with $\widehat F(\lambda)=1$. For this choice, every finite
$S$ satisfies the strict, $S$-independent bound
\[
 \left\|P_S\mathcal ZF-F\right\|_{H_{\alpha,s}}
 \ \ge\
 \left(\int_{\mathbb R}
     e^{2\Re\lambda x-2\alpha|x|}(1+x^2)^{-s/2}\,dx\right)^{-1/2}.
 \tag{MA-TR-9}
\]
If $S$ exhausts the primes and $q_S$ is the least omitted prime,
the difference in (MA-TR-9) is supported in $0<t\le b/q_S$ and
tends to zero in $C^\infty$ on every compact subset of
$\mathbb R_{>0}$. Thus local convergence of the prime-removal
operation conceals a nonvanishing Hilbert-space remainder escaping
to norm zero.

\emph{Proof.}
For fixed $t>0$, every derivative of the theta series converges
absolutely by the rapid decay of $f$ at infinity. Uniform convergence
holds when $t$ remains in a compact subinterval of $(0,\infty)$.
For large $t$ this also proves rapid decay of all logarithmic
derivatives. Poisson summation for the even Schwartz function gives
\[
 \sum_{n\in\mathbb Z}f(nt)=t^{-1}
                \sum_{n\in\mathbb Z}\widehat f(n/t).
\]
The two zero terms vanish by the conditions defining
$\mathscr S_{0,+}$. The right side and every logarithmic derivative
therefore decrease faster than every power of $t$ as $t\downarrow0$.
Multiplication by $\sqrt t$ preserves this conclusion. This proves
that $\mathcal ZF$ is exponential Schwartz in $\log t$, and the same
estimates, in terms of finitely many Schwartz seminorms for each
target seminorm, prove continuity.

For completeness, its image is the actual entire trivial-character
theta image, not just a collection of functions with the same
zeros. Average an adelic Schwartz representative over
$\widehat{\mathbb Z}^{\times}$. At a finite place, a locally constant
compactly supported unit-invariant function is a finite linear
combination of indicators of $p^j\mathbb Z_p$, $j\in\mathbb Z$.
Indeed its nonzero valuation shells are finite in number away from
a sufficiently small ball, on which local constancy makes it
constant; differences of adjacent ball indicators give the shells.
A restricted tensor product is therefore a finite linear combination
of rational dilates of $\prod_p1_{\mathbb Z_p}$. Absorb each rational
dilate into the real Schwartz function. A ball $q\widehat{\mathbb Z}$
contributes precisely the real function $f(qt)$ to the theta sum;
its finite additive volume is $q^{-1}$ for $q>0$. Thus both the
value at zero and the integral condition are preserved under this
absorption. Replacing the resulting real function by its even part
does not change its theta sum and preserves those conditions.
This proves surjectivity onto the actual compact-trivial image.
MT-M24 identifies that image with the closed ideal used here.

For $g=\mathcal ZF$ and $t>0$, the double series obtained from
(MA-TR-7) is absolutely convergent. For example, with $t$ in a fixed
compact interval and a sufficiently large $N$ its absolute value
and each logarithmic derivative are bounded by a constant times
\[
 \sum_{m,n\ge1}(mn)^{-1/2}(1+mnt)^{-N}<\infty.
\]
The coefficient of $k^{-1/2}F(kt)$ after regrouping is
$\sum_{n\mid k}\mu(n)$, which is $1$ at $k=1$ and $0$ otherwise.
This proves $\mathcal J\mathcal ZF=F$ and injectivity. The source
and target are Fr\'echet spaces: the source is a closed subspace
of Schwartz space, and the target is the closed ideal of MT-M24.
The open mapping theorem now proves continuity of the inverse in
these specified smooth topologies. Commutation with dilations is
immediate term by term. No zero-location result has been used.

Expanding the finite Euler product and regrouping the absolutely
convergent series gives (MA-TR-8) by the same divisor identity,
now restricted to square-free divisors composed of primes in $S$.
If $F$ is compactly supported, its corresponding even $f$ is
smooth and zero near the origin. Moreover
\[
 \int_{\mathbb R}f(t)\,dt
    =\int_0^\infty F(t)t^{-1/2}\,dt=\widehat F(1/2),
\]
so it belongs to $\mathscr S_{0,+}$. On any prescribed open
subinterval of $(a,b)$ the two functionals
$\widehat F(1/2)$ and $\widehat F(\lambda)$ are linearly independent:
otherwise the two smooth densities $t^{-1/2}$ and $t^{\lambda-1}$
would be proportional throughout that interval, forcing
$\lambda=1/2$, contrary to $\rho\ne1$. A linear map to
$\mathbb C^2$ with independent coordinate functionals is onto.
Choose $F$ with values $(0,1)$ under these two functionals.

Every $P_S\mathcal ZF$ belongs to the centered theta ideal, so its
Mellin value at $\lambda$ is zero by the full-jet description in
MT-M18/M24. This can also be seen first in $\Re z>1/2$ from
\[
 \widehat{\mathcal ZF}(z)=\zeta(z+1/2)\widehat F(z)
\]
and then by meromorphic continuation: the potential pole at $z=1/2$
is removed by the chosen moment. Cauchy--Schwarz on
$H_{\alpha,s}$ makes evaluation at $\lambda$ continuous with norm
at most the square root of the integral in (MA-TR-9). Applying it
to the difference proves that lower bound. Finally every $n>1$
in (MA-TR-8) is at least $q_S$; a nonzero summand requires $nt\le b$.
This proves the support assertion and exact eventual vanishing on
each compact interval.\hfill$\square$

\paragraph{Scope and input of MA-TR-101.}
The smooth isomorphism and Euler-tail identity are unconditional and
defined by the theta and M\"obius sums, without a spectral choice.
The lower bound is a conditional statement for each zero actually
retained by a chosen strip. A wide strip $\alpha>1/2$ retains every
nontrivial zero by the classical critical-strip theorem. For an
arbitrarily narrow strip, the existence of a critical-line zero already
suffices to instantiate the obstruction; this is classical, and is
strictly less than $A$. Such use of a known zero is recorded as a
zero-location input to the obstruction, never used to prove the
desired geometric equality. Alternatively the bound may simply be
left quantified over retained zeros, with no existence assertion.
The theorem does not prove that every other global-to-local
transport is impossible. It does rule out identifying compact
convergence of the finite Euler removals with convergence in the
specified Hilbert norm.

\subsection{An additional operator-order issue that is not formal}

\paragraph{MA-TR-102 (quotient trace class does not imply trace class after an orthogonal lift).}
There is a uniformly bounded strongly continuous representation
$V_t$ of $\mathbb R$ on a Hilbert space $H$, a closed invariant
subspace $J$, and a test $f\in C_c^\infty((0,\infty))$ such that
the integrated representation on $H/J$ is trace class, whereas
$C_fP_{J^\perp}$ on $H$ is not trace class. This remains true with
a simple purely imaginary quotient spectrum. Thus MT-M26/M28,
as quotient theorems, do not alone prove Masullo's separate
Theorem 10.4.2.8 about the opposite operator order.

\emph{Proof.}
Let $H=\ell^2(\mathbb N)\oplus\ell^2(\mathbb N)$ and
$J=\ell^2(\mathbb N)\oplus0$. Put
\[
 R_t=\operatorname{diag}(e^{int})_{n\ge1},\qquad
 K=\operatorname{diag}((in)^{-1})_{n\ge1},\qquad
 V_t=\begin{pmatrix}I&K(R_t-I)\\0&R_t\end{pmatrix}.
 \tag{MA-TR-10}
\]
The identity $K(R_{t+u}-I)=K(R_u-I)+K(R_t-I)R_u$ proves the
group law. Strong continuity follows from that of $R_t$, and
$\|K(R_t-I)\|\le2$ proves uniform boundedness. The quotient
action is $R_t$. For $F(z)=\int f(t)e^{zt}\,dt$ the integrated
matrix is
\[
 C_f=\begin{pmatrix}
 F(0)I&\operatorname{diag}((F(in)-F(0))/(in))\\
 0&\operatorname{diag}(F(in))
 \end{pmatrix}.
 \tag{MA-TR-11}
\]
Choose $f$ smooth, compactly supported away from zero, with
$F(0)\ne0$. Integration by parts gives rapid decay of $F(in)$,
so the quotient operator and $P_{J^\perp}C_f$ are trace class.
The upper-right corner of $C_fP_{J^\perp}$ has singular values
$|F(in)-F(0)|/n$, whose sum diverges. A corner of a trace-class
operator is trace class, so $C_fP_{J^\perp}$ is not.
The quotient eigenvalues $in$ are simple.\hfill$\square$

This is an abstract counterexample to an inference, not an arithmetic
counterexample to Theorem 10.4.2.8. The latter has its own unfinished
estimate on p.~464, (10.4.2.25)--(10.4.2.34). For the actual
arithmetic operators that additional estimate, or a valid replacement,
must be checked. F-SM-18 and F-SM-21 appropriately stated their
diagonal conclusions conditional on that actual lifted trace-class
input. The distinction is retained here.

\subsection{What the smooth transport does and does not accomplish}

The smooth map (MA-TR-7) explicitly contains the finite-prime
corrections. Its Mellin multiplier, wherever the Euler product
converges, is $1/\zeta(z+1/2)$, whose logarithmic derivative is
$-\zeta'/\zeta(z+1/2)$. With our translation convention the
corresponding distribution is
\[
 K(y)=\sum_{p,k\ge1}(\log p)p^{-k/2}\delta_{-k\log p}(y),
 \qquad
 (C_hK)(0)=\sum_{p,k\ge1}(\log p)p^{-k/2}h(p^k).
\]
Thus the centered contribution of the prime $p$ in this evaluation is
\[
 (\log p)\sum_{k\ge1}p^{-k/2}h(p^k)
\]
on tests supported above norm one. The reciprocal contribution is
obtained from the dual functional equation. This explains the
$\log2/\sqrt2$ correction in F-SM-21 using the actual arithmetic
map. It does not replace the missing trace estimate by a scalar
definition of a rank-one operator with the required trace.

For the source's central method, the outstanding assertion is that
the conjugated cutoff operators associated with $\mathcal J$ have
the correct common-space realization and a controlled remainder,
including the escape to norm zero in (MA-TR-9). The locally uniform
M\"obius limit is insufficient. An isomorphism of the smooth
geometric spaces supplies neither boundedness of its inverse in a
specified narrow Hilbert norm nor the convergence of nuclear traces.
The bounded Euler-module obstruction MT-M33 and the compatible
indentation obstruction MT-M37 continue to apply. They do not rule
out every non-equivariant family of cutoff operators.

In Grothendieck's terminology this is a question about transporting
an adelic trace on $C_{\mathbb Q}$ through a nonproper limiting
operation. Its Fredholm-determinant counterpart is the continuity
needed before taking $\det(1-wB)$ and its logarithmic derivative.
The exact archive comparison remains cote~1, scans~079--080 and
111--112, via ALPOGE\_BRIDGES BR-TR-01/02 and F-SM-11/12;
the archive pages themselves are not asserted freshly inspected.
At a finite spectral stage exterior powers and determinants live
in the usual $K_0$ and $\lambda$-ring formalism, but an analytic
limit is additional data. No theorem in EGA, FGA or SGA identifies
the above weighted Hilbert closure with a proper pushforward of
an algebraic sheaf; no such identification is claimed.

\section{New local checks in Masullo's Chapters 11--12}
\label{sec:local-source-audit}

\subsection{Scope and source normalization}

This section checks the public 940-page draft of Masullo, using printed
page numbers, which agree with its PDF page numbers.  The statements and
proofs on pp.~491--496 and 514--519 were read, together with the relevant
definitions on pp.~243, 259--260 and 427--428, Theorem~7.4.2.1 on
p.~367, and Lemma~C.3.2.2 on pp.~695--696.  The displayed formulas on
pp.~496, 514 and 517 were also checked in rendered page images.  No
copyrighted source PDF or page image is included in this deliverable.
These checks do not assume either $A$ or $M$, or any assertion concerning
the location of nontrivial zeta zeros.  They are additional to the
previously identified indentation and global-to-local problems.

The conclusions have different scopes.  MA-LA-001 is a literal
normalization inconsistency in a Poisson-summation step.  MA-LA-002 is
a false codomain for the printed cutoff.  MA-LA-003--004 establish a
path dependence for the natural annular repair, and give a sufficient
rate that repairs this \emph{local} trace limit.  MA-LA-005 identifies
two false intermediate kernel equalities but also explains their
cancellation in the real, even calculation; it is not an additional
counterexample to the repaired limiting trace formula.  MA-LA-006
separates the valid cyclic nuclear-trace lemma from an unverified
factorization hypothesis in its application.

\subsection{The missing half-weight in the orbital Poisson sum}

\paragraph{MA-LA-001 (the centered and uncentered orbital sums differ).}
For $h\in C_c^\infty(\mathbb R_{>0})$, set
\[
 \mathcal Mh(s)=\int_0^\infty h(t)t^s\,\frac{dt}{t}.
\]
For every rational prime $p$,
\[
 \sum_{n\in\mathbb Z}
 \mathcal Mh\!\left(\frac{2\pi i n}{\log p}-\frac12\right)
 =\log p\sum_{k\in\mathbb Z}p^{-k/2}h(p^k).
 \tag{MA-LA.1}
\]
Both sides are absolutely convergent.  In particular, if $h$ is
supported in $(3/2,5/2)$ and $h(2)=1$, then at $p=2$ the value is
$\log 2/\sqrt2$, whereas the unweighted sum is $\log 2$.

\emph{Proof.}
Put $L=\log p$ and $H(u)=e^{-u/2}h(e^u)$.  This is a compactly
supported smooth function on $\mathbb R$.  For the Fourier convention
$\widehat H(\xi)=\int H(u)e^{-2\pi i u\xi}\,du$, the summand on
the left is $\widehat H(-n/L)$.  Poisson summation gives
$\sum_n\widehat H(n/L)=L\sum_k H(kL)$, proving (MA-LA.1).
Rapid decay of $\widehat H$ proves absolute convergence; the other
sum is finite.  The support condition leaves only $2^1$.
\hfill$\square$

Here is the precise comparison with the draft.  In Theorem~7.4.2.1,
equation~(7.4.2.1), p.~367, $\varepsilon_p$ is defined by the zeros of
the \emph{ordinary} reciprocal local Euler factor, and the integrand
has $|g|^{\varepsilon_p}$.  For $F=\mathbb Q$ and $\chi=1$ these
zeros are $2\pi i n/\log p$, each of multiplicity one.  In the overview
on p.~494 and equation~(11.3.2.5), p.~496, the same test function and
the same local-zero parameters occur with
$|g|^{\varepsilon_p-1/2}$.  The following sentence identifies this
with the earlier right side, although its exponent differs by
$-1/2$.  Equation~(11.3.2.1), p.~495, and the Poisson conclusion on
p.~494 give an unweighted prime-power sum.

Take $\varphi(g)=h(|g|)$, normalize the compact norm-one fiber to
volume one, and use the standard unramified measure
$\mu_p(\mathbb Z_p^\times)=1$.  Then the displayed centered sum in
(11.3.2.5) is exactly the left side of (MA-LA.1).  A local uniformizer
has norm $p^{-1}$, so the printed unweighted sum
$\log p\sum_k\varphi(\pi_p^k)$ is
$\log p\sum_k h(p^k)$ after reindexing.  The corrected centered sum
in uniformizer notation is
\[
 \log p\sum_{k\in\mathbb Z}
 |\pi_p^k|^{-1/2}\varphi(\pi_p^k)
 =\log p\sum_{k\in\mathbb Z}p^{-k/2}h(p^k).
 \tag{MA-LA.2}
\]
Thus no implicit twist found in these source definitions removes the
discrepancy: twisting the test changes its value on the printed right
side as well.  Deleting the centering or inserting the half-weight are
different edits.  Inserting the weight repairs this Poisson step; it
does not repair the indentation or establish a global-to-local
intertwiner.  The example uses only zeros of $1-p^{-s}$ and Poisson
summation, not zeros of a global zeta function.

\subsection{The cutoff at the additive origin}

\paragraph{MA-LA-002 (an outer cutoff does not have the stated codomain).}
The source defines
\[
 \mathcal S_{00}(\mathbb R)
 =\{f\in\mathcal S(\mathbb R):f^{(j)}(0)=0\text{ for all }j\ge0\}
\]
in Theorem~5.3.2.1(1), p.~243; Definition~5.3.4.1, p.~259,
incorporates this factor at real places.  A smooth multiplier $a$ with
the usual Schwartz multiplier bounds satisfies
\[
 M_a\mathcal S(\mathbb R)\subseteq\mathcal S_{00}(\mathbb R)
 \quad\Longleftrightarrow\quad
 a^{(j)}(0)=0\quad(j\ge0).
 \tag{MA-LA.3}
\]
Consequently an outer cutoff equal to one near zero does not define
the map $P_{\rho,\lambda}:\mathcal S\to\mathcal S_{00}$ asserted
in Step II on p.~516.

\emph{Proof.}
Necessity follows by applying the multiplier to a test function equal
to one in a neighborhood of zero.  Sufficiency follows from the
Leibniz rule.  For the printed cutoff $a(0)=1$, even the zeroth
condition fails.\hfill$\square$

The preceding Fourier cutoff does not remove this issue, even on the
initial polar-free, zero-integral domain.  Let $b$ be a nonzero, even,
nonnegative compactly supported smooth frequency cutoff, and put
$q=\mathcal F^{-1}b$.  The function $q$ is nonconstant and even.
Choose two even smooth bumps supported away from zero, each of
integral one, whose integrals against $q$ differ; this is possible
by concentrating them near two positive points where $q$ differs.
Their difference $f$ belongs to $\mathcal S_{00}$ and has integral
zero, but
\[
 (\mathcal F^{-1}M_b\mathcal Ff)(0)=\int q(-y)f(y)\,dy\ne0.
\]
If $h\ge0$ is nonzero and supported in $(1,\infty)$, subsequent
averaging by $t^{-1/2}f(x/t)$ multiplies this value at zero by the
positive number $\int h(t)t^{-1/2}\,dt/t$.  An outer spatial cutoff
preserves the nonzero value.  Thus merely starting the composition
on $\mathcal S_{00}$ does not justify its claimed final codomain.

The footnote on p.~514 explicitly proposes replacing the spatial
cutoff by a smooth annulus $[1/\lambda'',\lambda]$, with
$\lambda''$ independent of $\lambda$, and asks for a convergence
study.  Vanishing on a neighborhood of zero does repair (MA-LA.3).
The next two results supply the required convergence study in the
real, trivial-isotype case and show why the relative rate matters.

\subsection{Exact real-place trace and the annular limit}

We use the Fourier normalization
\[
 (\mathcal F f)(\xi)=\int_{\mathbb R}f(x)e^{-2\pi i x\xi}\,dx,
 \qquad \check B=\mathcal F^{-1}B.
\]
Let $h\in C_c^\infty((1,\infty))$.  For $t>0$ set
\[
 D_tf(x)=t^{-1/2}f(x/t),\qquad
 A_h=\int_0^\infty h(t)D_t\,\frac{dt}{t}.
\]
Let $A,B\in C_c^\infty(\mathbb R)$ be even, satisfy $0\le A,B\le1$,
and equal one on a neighborhood of zero.  For $X,Y,\epsilon>0$, put
\[
 a_{X,\epsilon}(x)=A(x/X)\bigl(1-A(x/\epsilon)\bigr),
 \qquad b_Y(\xi)=B(\xi/Y),
 \qquad T_{X,Y,\epsilon}=M_{a_{X,\epsilon}}A_h
       \mathcal F^{-1}M_{b_Y}\mathcal F.
 \tag{MA-LA.4}
\]

\paragraph{MA-LA-003 (exact trace and two distinct annular limits).}
The operator in (MA-LA.4) is trace class, commutes with reflection,
and on $L^2_{\mathrm{even}}(\mathbb R,dx)$ has trace
\begin{align*}
 \operatorname{Tr}_{\mathrm{even}}T_{X,Y,\epsilon}
  =\frac12\int h(t)t^{-1/2}
   \int a_{X,\epsilon}(x)Y
   \left\{\check B\bigl(Y(t^{-1}-1)x\bigr)
          +\check B\bigl(Y(t^{-1}+1)x\bigr)\right\}
      dx\,\frac{dt}{t}.
 \tag{MA-LA.5}
\end{align*}
Write $L=XY$ and $r=Y\epsilon$.  If $L\to\infty$, then
\[
 \operatorname{Tr}_{\mathrm{even}}T_{X,Y,\epsilon}
 \longrightarrow
 \begin{cases}
 A_0(h):=\displaystyle\int_1^\infty
             h(t)\frac{\sqrt t}{t^2-1}\,dt,&r\to0,\\[6pt]
 0,&r\to\infty.
 \end{cases}
 \tag{MA-LA.6}
\]
For fixed $A,B,h$, the first error is
$O_N(L^{-N})+O(r)$, and the second is $O_N(r^{-N})$, for every
fixed $N$, after increasing the constants.  In particular the first
limit is strictly positive for nonnegative nonzero $h$.

\emph{Proof.}
The kernel, with respect to additive Lebesgue measure, is
\[
 K_T(x,y)=a_{X,\epsilon}(x)
 \int h(t)t^{-1/2}Y\check B\bigl(Y(x/t-y)\bigr)\,\frac{dt}{t}.
 \tag{MA-LA.7}
\]
It is a Schwartz function of $(x,y)$: $x$ has compact support and,
uniformly for $x$ there and $t$ in the fixed compact support of $h$,
all derivatives decay rapidly as $|y|\to\infty$.  A Schwartz kernel
is trace class and has trace its diagonal integral.  For completeness,
expand it in Hermite functions.  Applying arbitrary powers of the
harmonic oscillator in each variable gives coefficients
$c_{mn}=O_N((1+m)^{-N}(1+n)^{-N})$.  Their absolute sum is finite,
so the rank-one expansion converges in trace norm.  Moreover
$\int|e_m(x)e_n(x)|\,dx\le1$, which justifies integrating the
diagonal expansion and identifies its trace.

Let $Jf(x)=f(-x)$.  Evenness of both cutoffs makes $T$ commute with
$J$.  The even-sector trace is
$\tfrac12(\operatorname{Tr}T+\operatorname{Tr}(TJ))$; the two
diagonals are $K_T(x,x)$ and $K_T(x,-x)$.  This proves (MA-LA.5).

Substitute $z=Yx$.  Its inner integral becomes
\[
 \int A(z/L)(1-A(z/r))
 \{\check B(c_-(t)z)+\check B(c_+(t)z)\}\,dz,
 \qquad c_\pm(t)=t^{-1}\pm1.
 \tag{MA-LA.8}
\]
Because the support of $h$ is a compact subset of $(1,\infty)$,
both $|c_-|$ and $|c_+|$ have a positive lower bound there.  Thus
the Schwartz functions in (MA-LA.8) have an integrable dominating
function uniform in $t$.  If $r\to0$, the multiplier tends to one
for $z\ne0$; if $r\to\infty$, it tends to zero.  Dominated
convergence applies.  In the first case Fourier inversion and
$B(0)=1$ give
\[
 \int\check B(c_\pm(t)z)\,dz=\frac1{|c_\pm(t)|}.
\]
Consequently the limiting trace is
\[
 \frac12\int h(t)t^{-1/2}
 \left(\frac{t}{t-1}+\frac{t}{t+1}\right)\frac{dt}{t}
 = A_0(h).
\]
The outer-cutoff error is bounded by a Schwartz tail beyond a
constant times $L$.  For small $r$, the removed inner part is
supported on an interval of length $O(r)$ and its integrand is
uniformly bounded.  For large $r$, the remaining integrand vanishes
on an interval of radius a constant times $r$ and is bounded by a
Schwartz tail.  These prove the error estimates.\hfill$\square$

This is the normalized kernel that the draft itself writes in
(12.4.1.5) and (12.4.1.7), p.~516, followed by the three operators
listed in Step II and composed on p.~517.  Indeed, for
$F=\mathbb Q$, $S=\{\infty\}$, and the trivial compact isotype,
$\mathcal O_{\mathbb Q,S}^\times=\{\pm1\}$ and its invariant
sector is the even sector.  The normalized dilation in those
equations is precisely $t^{-1/2}f(x/t)$.  One may choose the lift
$\widetilde\varphi(t)=h(t)$ for $t>0$ and zero for $t<0$;
periodization by $\{\pm1\}$ gives $\varphi(g)=h(|g|)$.
Summing the two diagonals over $x>0$ equals the half of the full
real integral in (MA-LA.5).  The right side of the source's
(12.4.1.1), specialized to this case, is
\[
 \frac12\int_{\mathbb R^\times}
 h(|g|)\frac{|g|^{1/2}}{|1-g|}\,\frac{dg}{|g|}=A_0(h).
 \tag{MA-LA.9}
\]
There is no principal-value issue or logarithmic term for this $h$,
since its support is separated from $1$.  Thus the calculation tests
the actual displayed local kernel, without assuming the earlier
identification of a global geometric operator with that kernel.

\paragraph{MA-LA-004 (the path dependence persists for the source's
thin boundary smoothings).}
The fixed profiles in (MA-LA.4) are not needed for the obstruction.
Let $a_{X,\epsilon,\delta}$ be an even smooth approximation to
$1_{\{\epsilon\le|x|\le X\}}$, and $b_{Y,\delta}$ an even smooth
approximation to $1_{[-Y,Y]}$, with values in $[0,1]$, differing
from the indicators only in boundary intervals of total length
$O(\delta)$.  They may be chosen supported in the indicated
annulus and interval and equal to one outside boundary layers of
width $O(\delta)$.  Suppose
\[
 \delta<\tfrac14\epsilon,\qquad
 \delta(X+Y)\longrightarrow0.
 \tag{MA-LA.10}
\]
Then the trace of the same composition with these cutoffs differs
by $O_h(\delta(X+Y))$ from
\[
 \frac12\int h(t)t^{-1/2}
 \sum_{c\in\{c_-(t),c_+(t)\}}
 \frac{2}{\pi|c|}
 \left\{\operatorname{Si}(2\pi|c|XY)
       -\operatorname{Si}(2\pi|c|Y\epsilon)\right\}\,\frac{dt}{t},
 \tag{MA-LA.11}
\]
where $\operatorname{Si}(u)=\int_0^u(\sin v)/v\,dv$.
In particular, both paths
\[
 X=Y=N,\quad\delta=N^{-4},\qquad
 \epsilon=N^{-2}\quad\hbox{or}\quad\epsilon=N^{-1/2}
 \tag{MA-LA.12}
\]
satisfy the thin-smoothing requirements, but their traces tend to
$A_0(h)$ and $0$, respectively.

\emph{Proof.}
The inverse Fourier transform of the sharp frequency interval is
\[
 q_Y(v)=\frac{\sin(2\pi Yv)}{\pi v},\qquad q_Y(0)=2Y.
\]
For $c\ne0$, its integral over the sharp spatial annulus is
\[
 \int_{\epsilon\le|x|\le X}q_Y(cx)\,dx
 =\frac{2}{\pi|c|}
  \{\operatorname{Si}(2\pi|c|YX)
        -\operatorname{Si}(2\pi|c|Y\epsilon)\}.
\]
We use this merely as an explicit comparison integral; no trace
class assertion for discontinuous cutoff operators is required.
The frequency smoothing changes its inverse Fourier transform
uniformly by at most
$\|b_{Y,\delta}-1_{[-Y,Y]}\|_1=O(\delta)$, and integration over
the spatial support costs $O(X)$.
The spatial smoothing changes its $L^1$ norm by $O(\delta)$,
while $\|\mathcal F^{-1}b_{Y,\delta}\|_\infty\le
\|b_{Y,\delta}\|_1=O(Y)$.  Apply these estimates to each summand
in (MA-LA.5); the $t$ factors are bounded on the support of $h$.
This proves the stated error.

The classical improper integral
$\operatorname{Si}(u)\to\pi/2$ as $u\to\infty$, with error
$O(u^{-1})$ by integration by parts.  Also
$\operatorname{Si}(u)=O(u)$ as $u\to0$.  The constants are uniform
for the present $c_\pm(t)$.  Therefore (MA-LA.11) tends to
$A_0(h)$ when $XY\to\infty$, $Y\epsilon\to0$, and to zero
when $XY\to\infty$, $Y\epsilon\to\infty$.
All smoothing errors in (MA-LA.12) tend to zero.\hfill$\square$

Writing $X=\lambda$, $Y=\lambda'$ and
$\epsilon=1/\lambda''$, the essential ratio is
\[
 r=\frac{\lambda'}{\lambda''}.
 \tag{MA-LA.13}
\]
Both paths in (MA-LA.12) have
$\lambda,\lambda',\lambda''\to\infty$ and
$\delta\lambda,\delta\lambda'\to0$, the latter matching the
thin-smoothing regimes stated at the end of the source calculation
on p.~519.  Thus simultaneous divergence, even with those
smoothing conditions, is insufficient for the annular repair.
The sufficient rate $\lambda'/\lambda''\to0$ does restore the
local formula for tests supported away from $1$.  For the thin
boundary smoothing above the error is explicitly
\[
 O_h((\lambda\lambda')^{-1})+
 O_h(\lambda'/\lambda'')+
 O_h(\delta(\lambda+\lambda'))
 \tag{MA-LA.14}
\]
in this regime.  This result does not assert necessity of that rate
for each individual sign-changing test, and it does not settle
tests meeting $1$, the subtraction there, general $S$, or the
global intertwining step.

There is also a direct reading of the source's proposed volume
correction.  With a sharp annular spatial cutoff, the allowed
positive coordinate after $w=xu$ is
\[
 \max\{|w|/\lambda',1/\lambda''\}\le x\le\lambda.
\]
Its multiplicative volume is
\[
 \left[\log\lambda-
   \log\max\{|w|/\lambda',1/\lambda''\}\right]_+
 =\left[\log(\lambda\lambda')-
   \log\max\{|w|,r\}\right]_+.
\]
This is the correction suggested in the footnote on p.~514 to
(12.4.1.23), p.~518.  It displays the same residual parameter $r$;
one cannot discard it merely because both cutoffs diverge.

\subsection{The transposed kernel and why its errors can cancel}

\paragraph{MA-LA-005 (false intermediate equalities, with a limited
cancellation).}
The kernel convention in (12.4.1.5)--(12.4.1.7) is
\[
 K(x,y)=\int h(t)t^{-1/2}\delta(y-t^{-1}x)\,\frac{dt}{t}.
 \tag{MA-LA.15}
\]
For fixed $x$ and $\xi\in\{\pm1\}$, its substitution into
$K(x+z,\xi x)$ gives
\[
 \delta(\xi x-t^{-1}(x+z))
   =t\,\delta(z-(t\xi-1)x).
 \tag{MA-LA.16}
\]
Equation~(12.4.1.13), p.~517, instead substitutes
$\delta(x+z-t^{-1}\xi x)$, the kernel with the two arguments
exchanged.  Its factor is $t^{-1/2}$, whereas the actual integral
of the left kernel against a Fourier character has factor
$t^{1/2}$ and phase $(t\xi-1)xu$.

\emph{Verification.}
Equation (MA-LA.16) is the one-dimensional scaling rule for the
Dirac distribution, $\delta(t^{-1}v)=t\delta(v)$ for $t>0$.
Thus, with the negative Fourier character used in the source,
\[
 \int K(x+z,\xi x)e^{-2\pi i zu}\,dz
  =\int h(t)t^{1/2}
          e^{-2\pi i(t\xi-1)xu}\,\frac{dt}{t}.
 \tag{MA-LA.17}
\]
This is not its printed right side.\hfill$\square$

There is a preceding transposition as well.  The correct kernel
composition on p.~517 is
$a(x)\int K(x,z)q(z-y)\,dz$.  Its diagonal at $y=\xi x$
therefore uses $K(x,\xi x+z)$ against $q(z)$, not
$K(x+z,\xi x)$ against $q(z)$.  Parseval in its bilinear form
requires opposite Fourier signs.  The source's two characters in
(12.4.1.12) both have negative signs; that sign issue disappears
here because $b$, and hence $q$, is even.  Evenness does not by
itself justify interchanging the arguments of $K$.

It is necessary to record the cancellation, rather than count two
typographical slips as a further failure of the final real local
formula.  After the erroneous delta substitution, the source uses
the phase $(t^{-1}\xi-1)xu$ with the correct coefficient
$t^{-1/2}$.  For $\xi=1$ this is precisely the correct diagonal
phase.  For $\xi=-1$ its sign is the negative of the correct phase,
which again has no effect against even $b$.  Thus replacing the
two false intermediate lines by the direct kernel calculation
(MA-LA.7) yields (MA-LA.5).  The later variable change $t\mapsto
t^{-1}$ is legitimate with multiplicative Haar measure, and the
limiting density $|t|^{1/2}/|1-t|$ is inversion invariant.  These
facts explain why the source can arrive at the correct local
density despite the intermediate transpositions.  They do not
remove the annular path dependence of MA-LA-003--004.

The first erroneous transposition is not an exact finite-cutoff
identity that happens to be disguised.  For example, under the
fixed-profile scaling of MA-LA-003, with $L\to\infty$ and
$r\downarrow0$, its even trace would replace $t^{-1/2}$ and
$t^{-1}\pm1$ by $t^{1/2}$ and $t\pm1$.  Both expressions have
the same limit at $r=0$, but the coefficients of the removed inner
interval, to first order in $r$, are respectively
\[
 -r\check B(0)\!\left(\int A\right)
       \int h(t)t^{-1/2}\,\frac{dt}{t},
 \qquad
 -r\check B(0)\!\left(\int A\right)
       \int h(t)t^{1/2}\,\frac{dt}{t}.
\]
They differ for $h\ge0$, $h\ne0$ supported above $1$.
Indeed, substitute $z=rv$ into the removed term and use continuity
of $\check B$ at zero; $\check B(0)=\int B>0$.  Hence the
intermediate equality fails already for smooth local cutoffs.

\subsection{What is required for the trace factorization on p.~515}

\paragraph{MA-LA-006 (the cyclic trace lemma is valid; range inclusion
alone is insufficient).}
Lemma~C.3.2.2, pp.~695--696, requires a \emph{nuclear} lift
$A_0:E_2\to E_1$ through a continuous injection $T:E_1\to E_2$,
not only a nuclear endomorphism $A=TA_0$ of $E_2$ whose range
lies in $T(E_1)$.  Under its hypotheses, the claimed equality
\[
 \operatorname{Tr}(A_0T)=\operatorname{Tr}(TA_0)
 \tag{MA-LA.18}
\]
has a direct proof.

\emph{Proof.}
Choose a nuclear presentation
$A_0x=\sum_n\lambda_n\phi_n(x)y_n$ with
$\sum|\lambda_n|<\infty$, $(\phi_n)$ equicontinuous and
$(y_n)$ bounded.  The two compositions have presentations
\[
 TA_0x=\sum_n\lambda_n\phi_n(x)Ty_n,
 \qquad
 A_0Tz=\sum_n\lambda_n(\phi_n\circ T)(z)y_n.
\]
Their traces, defined independently of presentation using the
approximation properties assumed in the lemma, are the identical
absolutely convergent series
$\sum_n\lambda_n\phi_n(Ty_n)$.
\hfill$\square$

The finite-rank approximation explanation on p.~696 invokes
convergence uniform on precompact sets.  This topology by itself
does not make trace continuous: on $\ell^2$,
$n^{-1}P_{\{1,\ldots,n\}}\to0$ even in operator norm, but each
trace is one.  The nuclear presentation above supplies the needed
stronger control and proves the lemma without that inference.
Accordingly this observation is a repair of a proof detail, not a
counterexample to the lemma.

Nor does nuclearity automatically transfer to a stronger range
topology.  Put
\[
 E_2=\ell^2,\qquad
 E_1=\{y:\sum n^4|y_n|^2<\infty\},\qquad
 A x=(n^{-2}x_n)_n.
\]
The inclusion $T:E_1\to E_2$ is continuous, $A$ is trace class on
$E_2$, and its range is exactly $E_1$.  Its unique lift
$A_0:E_2\to E_1$ is an isometry onto $E_1$, hence is not compact
and is not nuclear.  This verifies the logical distinction needed
in the application on p.~515.

For the finite, smooth, annular real-place operator studied above,
the missing control \emph{can} be provided: its Schwartz kernel is
flat in the output coordinate near zero.  The Hermite expansion
in the input coordinate has coefficient functions in
$\mathcal S_{00}$ rapidly decreasing in every Schwartz seminorm.
Writing the expansion with scalar weights $(1+n)^{-2}$ gives a
nuclear map $L^2\to\mathcal S_{00}$ (the rescaled coefficient
functions are bounded in that Fr\'echet space and the Hermite
coefficient functionals have norm one on $L^2$).  Thus the local
regularization has a genuine nuclear lift.  This is a proof for
the explicitly specified real operator.  It does not verify an
unspecified lift from the global weighted Hilbert spaces, nor
establish the global-to-$S$-local operator identity required in
the missing Section~12.3.

\paragraph{Remaining gate.}
The proved local repairs are the centered prime weight and the
annular rate $\lambda'/\lambda''\to0$, with sufficiently thin
boundary smoothing.  An arithmetic transport still has to carry
the global prime contribution into the local comparison and to
control the relevant Hilbert or nuclear trace norms.  None of
MA-LA-001--006 claims that this is impossible.  Conversely, a
correct local cutoff limit alone does not provide that transport.

\section{The corrected annular limit on the full norm-one-vanishing test class}
\label{sec:local-extension}

The condition in task~1 is $h(1)=0$, not that the support of
$h$ be separated from $1$. The following extension closes that
local distinction. It does not treat nonzero $h(1)$, for which
a logarithmic subtraction is required, and it does not replace
the global arithmetic transport.

\paragraph{MA-LA-007 (all compact tests vanishing at norm one).}
In MA-LA-003 let instead $h\in C_c^\infty(\mathbb R_{>0})$
with $h(1)=0$. The exact even-sector trace (MA-LA.5) still holds.
Writing $L=XY$ and $r=Y\epsilon$, as there, one has
\[
 \operatorname{Tr}_{\rm even}T_{X,Y,\epsilon}
 \longrightarrow A_0(h)
 \quad\hbox{if }L\to\infty,\ r\to0,
 \qquad
 \operatorname{Tr}_{\rm even}T_{X,Y,\epsilon}
 \longrightarrow0
 \quad\hbox{if }r\to\infty,
 \tag{MA-LA.19}
\]
where
\[
 \begin{split}
 A_0(h)
 &=\frac12\int_0^\infty h(t)t^{-1/2}
       \left(\frac1{|t^{-1}-1|}+\frac1{t^{-1}+1}\right)\frac{dt}{t}\\
 &=\int_0^\infty\frac{e^{-x/2}}{1-e^{-2x}}
                 \bigl(h(e^x)+h(e^{-x})\bigr)\,dx.
 \end{split}
 \tag{MA-LA.20}
\]
The integrals are absolutely convergent. For fixed smooth
profiles $A,B$, the first error is $O_h(L^{-1}+r)$ for
$L\ge1$, $0<r\le1$, and the absolute trace in the second
regime is $O_h(r^{-1})$ for $r\ge1$.

For the thin boundary smoothings of MA-LA-004, with
$\delta<\epsilon/4$, the sufficient-rate error becomes
\[
 O_h\!\left(\frac{1+\log L}{L}+r+\delta(X+Y)\right).
 \tag{MA-LA.21}
\]
When $L,r\to\infty$ their absolute trace is bounded by
\[
 O_h\!\left(\frac{1+\log L}{L}
             +\frac{1+\log r}{r}+\delta(X+Y)\right).
 \tag{MA-LA.22}
\]
Thus both limiting paths in (MA-LA.12) remain valid on the full
requested test class, with the same two different limits whenever
$A_0(h)\ne0$.

\begin{proof}
The kernel in (MA-LA.7) is Schwartz for every fixed compact
support of $h$ inside $(0,\infty)$, whether or not that support
meets $1$. The trace-class and reflection arguments of
MA-LA-003 therefore give the same exact formula. Put
$c_-(t)=t^{-1}-1$ and $c_+(t)=t^{-1}+1$. The support is
contained in $[a,b]$ for some $0<a<b<\infty$.
For $t\ne1$, the absolute value of the inner integral
in (MA-LA.8) is at most
\[
 \|\check B\|_{L^1}
      \left(\frac1{|c_-(t)|}+\frac1{c_+(t)}\right),
\]
because the multiplier has absolute value at most one.
The factor $|h(t)|/|c_-(t)|$ is bounded near $1$:
$h(t)=O(|t-1|)$, while
$|c_-(t)|=|t-1|/t$. The other denominator stays bounded
away from zero throughout the compact support. Including
the bounded $t^{-1/2}$ and $dt/t$ factors, this is an
integrable dominating function in $t$. The point $t=1$
has measure zero and $h(1)=0$.

For each fixed $t\ne1$ the limiting inner integral has already
been computed in MA-LA-003. Dominated convergence now gives
both limits. For $t>1$, the first density in (MA-LA.20) is
$\sqrt t/(t^2-1)\,dt$; for $t<1$ it is
$t^{-1/2}/(1-t^2)\,dt$. Substituting $t=e^x$ and
$t=e^{-x}$, respectively, gives the second expression and
its absolute convergence.

For the quantitative bounds, the part removed by the outer
profile is at most
\[
 C_N|c|^{-1}(1+|c|L)^{-N}
\]
for either nonzero $c=c_\pm(t)$, by the Schwartz tail of
$\check B$. In a fixed neighborhood of $1$, multiply by
$|h(t)|$, cancel the factor $|c_-|^{-1}$, and use that
$dc_-/dt=-t^{-2}$ is bounded away from zero. For $N>1$,
\[
 \int_{|u|\le C}(1+|u|L)^{-N}\,du=O(L^{-1}).
\]
Away from that neighborhood the bound is $O_N(L^{-N})$;
the plus sign is harmless everywhere. The inner part removed
when $r\le1$ is supported on $|z|\le Cr$ and bounded by
$C\|\check B\|_\infty$, so its integral is $O(r)$,
uniformly in $t$. This proves the first error estimate.
When $r\ge1$, the remaining multiplier vanishes on
$|z|\le cr$. Replace $L$ by $r$ in the Schwartz-tail
argument just given to obtain $O(r^{-1})$.

For thin smoothings the error comparison with the sharp
annulus and sharp frequency interval is still
$O_h(\delta(X+Y))$: its proof used only bounded $t$ factors
and compact support, never separation of $c_-$ from zero.
For $u>0$ the sine integral satisfies
\[
 \left|\operatorname{Si}(u)-\frac\pi2\right|
       \le C\min(1,u^{-1}),\qquad
 |\operatorname{Si}(u)|\le C\min(u,1).
\]
The first inequality follows from integration by parts when
$u\ge1$, with boundedness for smaller $u$; the second is
immediate near zero and follows from boundedness otherwise.
In (MA-LA.11), the outer-end error is thus bounded by a
constant times
\[
 \int_{|v|\le C}\min(1,(|v|L)^{-1})\,dv
       =O((1+\log L)/L),
\]
using the same cancellation $|h|/|c_-|=O(1)$ and change
of variable. Away from $1$ the error is smaller. If $r\le1$,
the inner sine integral gives $|\operatorname{Si}(2\pi|c|r)|
\le C|c|r$, which yields $O_h(r)$. If $r\ge1$, apply
the outer-end estimate also with $L$ replaced by $r$.
Together with the smoothing comparison, these prove
(MA-LA.21)--(MA-LA.22).
\end{proof}

\paragraph{Scope, inputs, and programme placement.}
Every step uses the actual real dilation kernel and ordinary
Fourier analysis, with no zero-location premise. The sufficient
rate remains $\lambda'/\lambda''\to0$ in the notation of the
footnote on Masullo p.~514 and the proposed correction to
(12.4.1.23), p.~518. For the requested norm-one-vanishing
tests, this now proves the local annular repair even when their
support meets norm one. It closes the specific local question
Q-MA-LA-ENDPOINT; the nonzero-value logarithmic subtraction
and general-number-field cutoff estimates are different
questions and are not claimed.

In the adelic-trace interpretation this is the real-place
distribution on the full test class of GRH\_NIGHT MT-M30.
It refines the source-kernel side of ALPOGE\_BRIDGES
F-SM-20/21 without supplying the global prime transport.
The corresponding cote~1 issue (scans~079--080 and 111--112,
via BR-TR-01/02) is continuity of regularized trace under the
specified limiting process. No algebraic correspondence,
motivic weight statement, or new EGA/SGA identification is
asserted by this analytic estimate.

\section{The printed local operator chain and an explicit quotient repair}
\label{sec:source-operator-audit}

This section separates three spaces which the local calculation must
distinguish: the additive Schwartz space on $\mathbb A_{F,S}$,
its $S$-unit coinvariants, and the weighted global Hilbert space
with its orthogonal Paley--Wiener projections.  The first three
assertions below are source-level checks.  The subsequent construction
repairs a particular rational local quotient, with actual operators
and nuclear traces; it does not identify that quotient with the
global weighted projection.

\paragraph{MA-LA-008 (the centering change on pp.~515--516).}
The operator displayed near the bottom of Masullo p.~515 has,
before taking its output coinvariant class, the value
\[
 \begin{aligned}
 A_{\widetilde\varphi}f(x)
  &=\int \widetilde\varphi(g)|g^{-1}x|_S^{1/2}
                     f(g^{-1}x)\,d^\times g
  =|x|_S^{1/2}(Kf)(x),\\
 (Kf)(x)&=\int |g|_S^{-1/2}\widetilde\varphi(g)
                     f(g^{-1}x)\,d^\times g .
 \end{aligned}
 \tag{MA-LA.39}
\]
The second expression, without the output factor $|x|_S^{1/2}$,
is exactly (12.4.1.5) and (12.4.1.7), p.~516.  The paragraph
following (12.4.1.7) calls it the kernel of the same operator.
In coordinates the correct relation is $A=UK$, followed by
the output quotient, where $U=M_{|x|_S^{1/2}}$.
Equivalently $K$ represents $U^{-1}A$ after this explicit
change of target coordinates.  For $F=\mathbb Q$, $S=\{\infty\}$,
and even test functions the finite sign quotient does not remove
the factor: at $x=4$ the two nonzero displayed values differ
by a factor of two.

This is a repairable coordinate omission.  If the centering map
is used to identify operators, all other maps in that diagram
must be conjugated consistently; for example an additive Fourier
cutoff becomes $U^{-1}\mathcal F^{-1}M_b\mathcal F U$ in the
other coordinates.  The two formulas alone do not justify dropping
$U$ while retaining the other unaltered cutoff maps.
Both original pages were inspected as rendered pages, so this
finding is not an extraction ambiguity.

\subsection{A literal failure of the raw coinvariant lift}

Write $S=\{\infty,p\}$, $\mathbb A_S=\mathbb R\times\mathbb Q_p$,
$U_p^0=\mathbb Z_p^\times$, and
\[
 (L_t f)(x)=f(t^{-1}x),\qquad t\in\mathbb Q^\times.
 \tag{MA-LA.40}
\]
For $t=p$ the additive determinant is $|p|_S=1$, so $L_p$
is unitary on additive $L^2(\mathbb A_S)$.  For clarity the
polar-free test space used here is the concrete LF subspace
\[
 \mathcal V=\mathcal S_{\mathrm{flat}}(\mathbb R)
                 \otimes_\iota C_c^\infty(\mathbb Q_p^\times)
       \ \subset\ \mathcal S(\mathbb A_S),
 \qquad
 \mathcal S_{\mathrm{flat}}(\mathbb R)
       =\{u\in\mathcal S(\mathbb R):u^{(j)}(0)=0\ (j\ge0)\}.
 \tag{MA-LA.41}
\]
This spells out the local coordinate-hyperplane conditions used
in Definition~5.3.4.1, p.~259, and its proof, pp.~260--261.
It is stronger than vanishing only at the simultaneous adelic
origin.  Proposition~2.1.4.1(1), pp.~71--72, gives the actual
additive finite tensor structure: every function has a common
finite $p$-adic support and a common additive local-constancy
level.  The additional condition $f(x_\infty,0)=0$ consequently
gives a common neighborhood of $0$ in the $p$-adic variable on
which it vanishes.  Thus each LF stage of (MA-LA.41) has only
finitely many $p$-adic valuations.

There is a notation issue in the source which this explicit
definition avoids.  Remark~2.1.2.1, p.~63, uses
$\mathcal S(\mathbb Q_p^\times)$ for the larger Schwartz space
of the multiplicative group, including infinite rapidly decreasing
valuation sequences.  The proof on pp.~260--261 instead uses the
same notation for compactly supported, origin-vanishing functions
inside additive $\mathcal S(\mathbb Q_p)$, with its finest LF
topology.  No finite-valuation assertion about the larger
multiplicative Schwartz space is used here.

\paragraph{MA-LA-009 (the displayed raw kernel does not kill coinvariant relations).}
Let $q:\mathcal V\to\mathcal V_{\mathbb Z[1/p]^\times}$ be
the separated coinvariant map.  The tensor containment printed
after (12.4.1.7), p.~516, would imply that the raw-output operator
$K:\mathcal V\to\mathcal V$ factors through $q$ on its input.
This implication is false, even for a valid compact trivial-isotype
test supported away from norm one.

\begin{proof}
Choose $0\le h\in C_c^\infty(\mathbb R_{>0})$, nonzero,
supported in an interval disjoint from $1$, and put
\[
 \widetilde\varphi(g_\infty,g_p)
       =\tfrac12 h(|g_\infty|)\mathbf1_{U_p^0}(g_p).
 \tag{MA-LA.42}
\]
For any $g_p\ne0$ exactly one exponent $n$ makes
$p^n g_p$ a unit.  The two signs cancel the factor $1/2$,
so the $S$-unit periodization is exactly
\[
       \sum_{\xi\in\{\pm p^n:n\in\mathbb Z\}}
                \widetilde\varphi(\xi g)
                  =h(|g_\infty||g_p|_p).
 \tag{MA-LA.43}
\]
It is therefore a legitimate lift in (12.4.1.4), p.~515.
Use multiplicative measure with unit mass one at $p$;
other positive Haar normalizations only multiply the next
formula by a nonzero constant.  If $u$ is even and
$f=u\otimes\mathbf1_{U_p^0}$, then
\[
       Kf(x_\infty,x_p)
        =\mathbf1_{U_p^0}(x_p)
             \int_0^\infty h(t)t^{-1/2}u(x_\infty/t)\,\frac{dt}{t}.
 \tag{MA-LA.44}
\]
Take $u\ge0$ to be a nonzero smooth compact annular bump.
The right side is nonzero.  In contrast, $L_pf$ is supported
on the $p$-adic shell $pU_p^0$, and $K L_pf$ is supported on
the same shell, since the lift (MA-LA.42) has unit $p$-adic
component.  Thus $K(f-L_pf)$ is nonzero on $x_p\in U_p^0$,
although $q(f-L_pf)=0$.  The example lies in (MA-LA.41),
and in the source's intended polar-free subspace.
\end{proof}

The valid assertion is that $K$ commutes with the $S$-unit
action and hence induces $qK$ on the quotient.  Equivariance
does not make a raw-output lift constant on input coinvariant
classes.  Taking the output quotient, or specifying a section,
repairs this algebraic point but changes the map whose kernel
and trace have to be justified.

\paragraph{MA-LA-010 (product-norm cutoffs do not preserve additive Schwartz space).}
The Step~II multiplier on p.~516 is not generally an
endomorphism of $\mathcal S(\mathbb A_S)$ when $S$ has a
finite place.  This remains an issue after the proposed
annular position-cutoff repair.

\begin{proof}
Let $b:[0,\infty)\to\mathbb R$ be smooth, equal to one
near zero and zero for large arguments.  Apply the displayed
product-norm multiplier to
\[
       f(x_\infty,x_p)=e^{-x_\infty^2}\mathbf1_{\mathbb Z_p}(x_p).
 \tag{MA-LA.45}
\]
If $b(|x_\infty||x_p|_p)f(x_\infty,x_p)$ belonged to
additive Schwartz space, Proposition~2.1.4.1 would give an
integer $N\ge0$ such that its values at $x_p=0$ and
$x_p=p^N$ agreed for every $x_\infty$.
Choose $x_\infty=r p^N$, with $b(r)=0$.
The first value is $e^{-(r p^N)^2}$ and the second is zero.
This contradicts the required uniform local-constancy level.

For a nonzero smooth annular function $a$, with $a(0)=0$,
choose $r>0$ with $a(r)\ne0$ and use the same test.
The value at $x_p=0$ is zero and the value at $x_p=p^N$
is $a(r)e^{-(r p^N)^2}\ne0$.  Thus the annular multiplier
does not even map all additive Schwartz functions into
additive Schwartz space, and in particular does not have
the printed $\mathcal S\to\mathcal S_{00}$ codomain.
\end{proof}

These functions are perfectly good bounded multipliers on
additive $L^2$.  The failure is their asserted Schwartz-space
domain and codomain, not measurability or boundedness on $L^2$.
It also explains why correcting the real-place inner radius
alone does not complete the mixed-place operator argument.

\subsection{A continuous section and a unitary quotient Fourier transform}

\paragraph{MA-LA-011 (an explicit rational local repair).}
On the actual LF space (MA-LA.41), define
\[
 E_- =\tfrac12(1+L_{-1}),\qquad
 P_0 f=\mathbf1_{U_p^0}(x_p)f,\qquad
 B=P_0\sum_{n\in\mathbb Z}L_{p^n}E_- .
 \tag{MA-LA.46}
\]
This is a continuous projection, its kernel is precisely the
space of $S$-unit coinvariant relations, and $qB=q$.
Its range consists of the functions supported on the unit
valuation shell and invariant under simultaneous negation.
It consequently gives a continuous section of the separated
coinvariant quotient.

\begin{proof}
On a fixed LF stage the valuation support is a finite set.
After $P_0$ only those finitely many $n$ for which $-n$
belongs to that support survive.  Each resulting term is
a continuous dilation and a continuous shell projection.
Thus (MA-LA.46) is continuous on every stage and hence on
the strict LF space.  Reindexing this finite sum shows
$BL_p=B$ and $BL_{-1}=B$.

If $P_j$ denotes the projection onto the shell $p^jU_p^0$,
then, for each $f$, the following sums are finite:
\[
       f=\sum_jP_jf,\qquad
       Bf=\sum_jL_{p^{-j}}P_jE_-f.
 \tag{MA-LA.47}
\]
It follows that $Bf-f$ is a finite sum of unit relations
and hence $qBf=qf$.  On the stated range $B$ is the identity,
so $B^2=B$.  If $Bf=0$, (MA-LA.47) expresses $f$ as a
finite sum of unit relations.  Thus this relation space is
already closed, being $\ker B$, and the induced section is
continuous for the quotient topology.
\end{proof}

The finite-sum assertion must not be reused after additive
Fourier transformation.  With self-dual measures,
\[
 \mathcal F_p\mathbf1_{U_p^0}
       =\mathbf1_{\mathbb Z_p}
                   -p^{-1}\mathbf1_{p^{-1}\mathbb Z_p}.
 \tag{MA-LA.48}
\]
It is nonzero near zero.  Nevertheless its section-periodization
can be calculated and completed on a Hilbert space, rather than
asserted to lie in the original polar-free LF space.

Restrict to the compact trivial isotype, so section vectors are
$u\otimes\mathbf1_{U_p^0}$ with $u$ even.  Let
$\mathcal F_e$ be the usual unitary Fourier transform on
even $L^2(\mathbb R,dx)$ and set
\[
 \begin{aligned}
 (V_pu)(x)&=p^{1/2}u(px),& q_p&=p^{-1/2},\\
 D_p&=(1-q_pV_p^*)(1-q_pV_p)^{-1},&
 H_p&=D_p\mathcal F_e .
 \end{aligned}
 \tag{MA-LA.49}
\]
Then $D_p$ is unitary, $H_p$ is a unitary involution, and
for $u\in\mathcal S_{\mathrm{flat}}(\mathbb R)$ the actual
periodization of $\mathcal F_S(u\otimes\mathbf1_{U_p^0})$
on the unit shell is
\[
 H_pu(x)=(1-p^{-1})\sum_{k\ge0}(\mathcal F_eu)(p^kx)
                         -p^{-1}(\mathcal F_eu)(x/p).
 \tag{MA-LA.50}
\]
The sum converges in $L^2$ and normally with all derivatives
on compact annuli.  It need not be flat at zero.

\begin{proof}
At $x_p\in U_p^0$, (MA-LA.48) gives the coefficient
$1-p^{-1}$ for $n\le0$, the coefficient $-p^{-1}$ for
$n=1$, and zero for $n\ge2$ in the periodization
$\sum_nL_{p^n}\mathcal F_S f$.  This proves the displayed
formula pointwise on annuli.  Rapid decay of $\mathcal F_eu$
proves normal convergence there.  Its $k$th dilation has
$L^2$ norm $p^{-k/2}\|u\|_2$, so the series also converges
absolutely in operator norm.
In terms of $V_p$ its multiplier is
\[
 (1-q_p^2)\sum_{k\ge0}q_p^kV_p^k-q_pV_p^*
       =\frac{1-q_pV_p^*}{1-q_pV_p}=D_p.
 \tag{MA-LA.51}
\]
Since $V_p$ is unitary, numerator and denominator commute
and are adjoints, and the denominator is boundedly invertible.
Thus $D_p^*D_p=1$.
The Fourier scaling identity gives
$\mathcal F_e V_p\mathcal F_e=V_p^*$ and
$\mathcal F_e D_p\mathcal F_e=D_p^{-1}$.
Together with $\mathcal F_e^2=1$, this proves $H_p^2=1$.

For an explicit warning even with zero integral, choose an
even nonnegative annular bump $\psi$ and let
$u(x)=\psi(x)-p^{-1}\psi(x/p)$.
Then $u$ is flat at zero and $\int u=0$.
The series in (MA-LA.50) telescopes to
\[
                H_pu(x)=(\mathcal F_e\psi)(x)
                                -p^{-1}(\mathcal F_e\psi)(x/p),
 \tag{MA-LA.52}
\]
whose value at zero is
$(1-p^{-1})(\mathcal F_e\psi)(0)\ne0$.
Here (MA-LA.52) is the identity on nonzero real annuli,
hence also an $L^2$ identity; its right side gives the
smooth continuation to zero.  Termwise evaluation of
the original periodization at the single point zero
need not equal that continuation.
Thus the intermediate Fourier-periodized function is not
in the flat source space, although it belongs to the stated
Hilbert completion.
\end{proof}

There is now an actual local nuclear construction.  For an even
smooth annular multiplier $a$, a bounded even multiplier $b$,
and $h\in C_c^\infty(\mathbb R_{>0})$, put
\[
 (A_hu)(x)=\int_0^\infty h(t)t^{-1/2}u(x/t)\,\frac{dt}{t},
 \qquad
 T^{\mathrm{quot}}_{a,b,h}=M_a A_h H_pM_bH_p.
 \tag{MA-LA.53}
\]
This is trace class on even $L^2(\mathbb R)$ and factors
nuclearly through the flat unit-shell section.
Indeed the unitary logarithmic map
$u\mapsto\sqrt2e^{y/2}u(e^y)$ converts $M_aA_h$ into
the integral operator with kernel
\[
                  a(e^y)h(e^{y-z}).
 \tag{MA-LA.54}
\]
It is smooth and supported in a compact rectangle.
A Fourier-series expansion on a larger rectangle has coefficients
decreasing faster than every power in both indices.  It gives
an absolutely summable rank-one presentation and, with all output
derivatives, a nuclear map into smooth functions on the compact
output interval.  Multiply each output Fourier vector by a fixed
smooth cutoff supported inside that larger interval and equal to
one on the kernel's output support.  This leaves the kernel
unchanged and makes every output vector extend smoothly by zero.
For the latter nuclear assertion one may choose a
diagonal, faster-than-polynomial scaling of the output Fourier
vectors: the rapid coefficient bounds permit summable functional
norms while the rescaled vectors are bounded in every smooth
seminorm.  Returning from logarithmic coordinates gives a nuclear
map into $\mathcal S_{\mathrm{flat}}(\mathbb R)$.
The remaining factor $H_pM_bH_p$ is bounded, proving the claim.
Insertion into the unit shell and its continuous quotient map
give the trace-compatible factorization of Lemma~C.3.2.2;
the common nuclear rank-one series gives both traces.

This construction is on the quotient Hilbert completion just
specified.  It does not assert that every intermediate operation
preserves additive Schwartz space.  Nor does it assert equality
with the global weighted orthogonal Paley--Wiener projection.
The coupled annular cutoff limit requires its own proof;
MA-LA-014 supplies it below for this rational local sector.
Finite-cutoff nuclearity alone does not fix that limit.

\subsection{The full additive cover cannot supply a nonzero compact trace operator}

\paragraph{MA-LA-012 (unit symmetry prevents compactness on the cover).}
Let $a,b$ be bounded functions of the product norm and let
$K$ be the compactly averaged unitary dilation operator in
(MA-LA.39).  On ordinary additive $L^2(\mathbb A_S)$ the operator
\[
             T=M_{a(|x|_S)}K\mathcal F_S^{-1}
                         M_{b(|\xi|_S)}\mathcal F_S
 \tag{MA-LA.55}
\]
is bounded and commutes with $L_p$.
If it is nonzero it is not compact, hence not trace class.
There are nonzero examples with the valid test (MA-LA.42),
with the printed outer cutoff and also with every nonzero
smooth annular position cutoff.

\begin{proof}
The normalized dilation $|g|_S^{-1/2}f(g^{-1}x)$ is unitary,
so $\|K\|\le\|\widetilde\varphi\|_{L^1(d^\times g)}$.
Both multipliers are bounded and the additive Fourier transform
is unitary.  Norm-dependent multipliers commute with $L_p$,
and multiplicative convolution does also.  Since $|p|_S=1$,
\[
           \mathcal F_S L_p=L_{p^{-1}}\mathcal F_S.
 \tag{MA-LA.56}
\]
This shows that the middle Fourier multiplier commutes with $L_p$.

The powers $L_p^n$ tend weakly to zero on additive $L^2$.
To see this first take two functions with $p$-adic supports
in finite unions of valuation shells; their translated supports
are disjoint for large $n$.  Such functions are dense in $L^2$,
since $x_p=0$ has measure zero, and unitarity gives the general
claim by approximation.  If $T$ were compact, then
$TL_p^nf\to0$ in norm for every $f$.  Commutation instead gives
\[
                     \|TL_p^nf\|_2=\|Tf\|_2.
 \tag{MA-LA.57}
\]
Thus every compact commuting $T$ is zero.

For nonzeroness take $\widetilde\varphi$ as in (MA-LA.42),
and $0\le b\le1$ equal to one near zero and with compact
norm support.  Use the self-dual positive test
$f=e^{-\pi x_\infty^2}\mathbf1_{\mathbb Z_p}(x_p)$.
The function $b\mathcal F_Sf$ is integrable and nonnegative,
with positive integral, so
$v=\mathcal F_S^{-1}(b\mathcal F_S f)$ is continuous and
$v(0)>0$.  Compactness of the dilation parameters and
nonnegativity of the test imply $Kv>0$ in a neighborhood
of the additive origin.  This already proves nonzeroness
for outer position cutoffs, and for sufficiently small
annular inner radii.

For any prescribed nonzero smooth annular $a$, the same
example works.  Gaussian domination shows that $v$, and
then $Kv$, are entire in $x_\infty$ for each fixed $x_p$.
Choose a sufficiently small nonzero $x_p$ so that
$Kv(0,x_p)>0$.  This entire function cannot vanish on an
open real interval.  The set of positive $x_\infty$ for
which $a(|x_\infty||x_p|_p)\ne0$ contains such an interval,
so $M_aKv$ is nonzero.  This completes the example.
\end{proof}

The section-periodization itself has a related continuity
boundary.  If $f_0$ is a nonzero unit-shell vector fixed by
simultaneous negation, then
\[
 f_N=N^{-1/2}\sum_{n=1}^NL_{p^n}f_0,
 \qquad \|f_N\|_2=\|f_0\|_2,
 \qquad Bf_N=\sqrt N f_0.
 \tag{MA-LA.58}
\]
The supports are disjoint valuation shells.  The same argument
works for any norm-dependent additive $L^2$ weight, since
$|p|_S=1$.  Hence the LF section-periodization has no bounded
extension from that full additive Hilbert space.  This does
not contradict the bounded compressed Fourier transform
$H_p$ in MA-LA-011, whose input is the chosen quotient section.

The attribution boundary matters.  The explicit nuclearity
claim on p.~515 concerns a quotient Paley--Wiener level.
On p.~517 Step~III invokes nuclearity after assigning the same
$k^{\mathrm{PW}}$ notation to the raw additive kernel, and
then integrates over a fundamental domain and sums over units.
MA-LA-012 rules out interpreting this as trace class on the
full additive cover.  It does not rule out a quotient trace:
MA-LA-011 constructs one.  A compatible quotient or nuclear
factorization must be supplied before that raw kernel can
be used for the claimed quotient trace.

\subsection{The prime term carried by the repaired quotient transform}

\paragraph{MA-LA-013 (the local all-pass factor has the prime trace).}
In logarithmic coordinates let $\ell=\log p$,
$(V_pf)(x)=f(x+\ell)$, and retain $D_p$ from (MA-LA.49).
Let $g\in C_c^\infty(\mathbb R)$ and let $\phi$ be smooth,
with $\phi(-\infty)=1$, $\phi(+\infty)=0$ and compactly
supported derivative.  Then the actual operator
\[
                  C_g(D_p^*M_\phi D_p-M_\phi)
 \tag{MA-LA.59}
\]
is trace class on $L^2(\mathbb R)$ and
\[
 \operatorname{Tr}\bigl(C_g(D_p^*M_\phi D_p-M_\phi)\bigr)
      = (\log p)\sum_{k\ge1}p^{-k/2}
                    \bigl(g(k\log p)+g(-k\log p)\bigr).
 \tag{MA-LA.60}
\]
For a finite set of primes, $D_S=\prod_{p\in S_f}D_p$
gives the sum of these local prime traces.

\begin{proof}
Write $q=p^{-1/2}$ and the absolutely convergent Laurent series
\[
 D_p=\sum_{n\in\mathbb Z}d_nV_p^n,\qquad
 d_{-1}=-q,\quad d_n=(1-q^2)q^n\ (n\ge0),\quad
 d_n=0\ (n<-1).
 \tag{MA-LA.61}
\]
Unitarity gives $\sum_m\overline d_m d_{m+r}=\delta_{r0}$.
Consequently the difference in (MA-LA.59), before convolution,
is the double sum with coefficients $\overline d_m d_n$ of
\[
    M_{\phi(x-m\ell)-\phi(x)}V_p^{n-m}.
 \tag{MA-LA.62}
\]
Each multiplier difference is compactly supported in an interval
of length $O(1+|m|)$, with all positive-order derivatives
uniformly bounded.  A fixed smooth partition into unit intervals
and the compact smooth kernel of $C_g$ on each piece give
\[
 \bigl\|C_gM_{\phi(x-m\ell)-\phi(x)}V_p^{n-m}\bigr\|_1
                         \le C_{g,\phi,p}(1+|m|).
 \tag{MA-LA.63}
\]
Here translation of a piece is unitary and the final shift
does not change trace norm.  The coefficient series weighted
by $1+|m|$ is summable, proving an actual trace-norm double
series before any traces are taken.

The diagonal kernel of each term gives
$g((n-m)\ell)\int(\phi(x-m\ell)-\phi(x))\,dx$.
The integral equals $m\ell$, by integrating $\phi'$ and
using its total integral $-1$.  Thus the coefficient of
$\ell g(r\ell)$ is
\[
                  a_r=\sum_m m\overline d_m d_{m+r}.
 \tag{MA-LA.64}
\]
For $r\ge1$ its $m=-1$ term is $(1-q^2)q^r$ and its
$m\ge0$ terms sum to $q^{r+2}$, so $a_r=q^r$.
For $r=-k$, $k\ge1$, the $m=k-1$ term is
$-(k-1)(1-q^2)q^k$ and the $m\ge k$ terms sum to
$q^k(q^2+k(1-q^2))$, again giving $a_{-k}=q^k$.
For $r=0$ the two contributions $-q^2$ and $q^2$ cancel.
This proves (MA-LA.60), including both directions and its sign.

All the $D_p$ commute with one another and with $C_g$.
For two factors, expand
\[
 (D_1D_2)^*M_\phi(D_1D_2)-M_\phi
 =D_2^*(D_1^*M_\phi D_1-M_\phi)D_2
                  +(D_2^*M_\phi D_2-M_\phi).
\]
The preceding trace-class estimate and bounded cyclicity make
the trace of the first summand equal to the first local trace.
Induction proves the finite-set assertion.  The same finite
product of (MA-LA.50) computes the compact trivial quotient
Fourier transform for $S=\{\infty\}\cup S_f$.
\end{proof}

Thus the repaired quotient Fourier transform carries the local
prime terms through an explicit trace-class commutator.  This
is an operator statement, not a trace assigned from an explicit
formula.  It does not replace the separate relative-rate analysis
for two compact annular/Fourier cutoffs by a one-sided cutoff.

The same operator calculation has a precise character version.
For $|c|=1$ put
\[
 D_{p,c}=(1-\overline c\,p^{-1/2}V_p^*)
                         (1-cp^{-1/2}V_p)^{-1}.
 \tag{MA-LA.65}
\]
It is unitary.  Its Laurent coefficients are
$d_{-1}=-\overline c\,q$ and
$d_n=(1-q^2)c^nq^n$ for $n\ge0$.
The absolute trace-norm estimates above are unchanged, while
(MA-LA.64) becomes $a_k=c^kq^k$ and
$a_{-k}=\overline c^{\,k}q^k$ for $k\ge1$, with $a_0=0$.
Consequently the trace is
\[
 (\log p)\sum_{k\ge1}p^{-k/2}
       \bigl(c^kg(k\log p)+\overline c^{\,k}g(-k\log p)\bigr).
 \tag{MA-LA.66}
\]
This proves the dual coefficient orientation at the operator
level, in agreement with MA-CC-001 and MA-RS-005.
It is an algebraic generalization of the local unitary factor;
no unproved full character-sector section is inferred from it.

\subsection{The exact remaining source bridge}

Definition~11.2.1.3(2), (11.2.1.8), p.~489, defines
$Q^{\mathrm{PW}}_{S,\Lambda_1,a}$ to be an orthogonal
projection in the weighted global Hilbert space.  Theorem~11.2.1.4
uses it in $P^{\mathrm{geom}}_{\rho,a}A(\varphi)Q^{\mathrm{PW}}$.
The local calculation instead uses a smooth additive Fourier
multiplier.  Literal equality of those two candidates fails
already because a nontrivial smooth multiplier $b$ is not
idempotent, whereas an orthogonal projection is.  This observation
refutes that literal candidate identification; the source does
not actually print such an equality, so it is not attributed to it.

The p.~515 assertion of a linear topological copy also cites
Proposition~5.3.4.4, pp.~263--264.  That proposition states an
equivalence between an entire local zeta transform, existence
of a polar-free representative with the same transform, and
rapid decay of its theta image.  Its displayed statement does
not itself specify the continuous inverse or a nuclear lift
used on p.~515.  The explicit section above supplies such
a map in its stated rational local sector; it is not a claim
that the cited proposition already supplied the full global
operator diagram.

Section~12.3 on p.~514 has no argument, only its heading and
the footnote ``Include section.''  The p.~520 proof then uses
the local theorem to identify the global $k_S$ diagonal integral,
with the two polar terms subtracted.  Proposition~12.2.2.2,
pp.~508--513, is a diagonal integration identity for a
\emph{specified} lifted quotient kernel, under its stated
regularity and absolute-integrability hypotheses.  It supplies
neither the global-to-local map nor the equality of these kernels.
The shell overview on pp.~491--494 also first defines its
valuation cutoffs upstairs on ideles, rather than on idele
classes; a descent or indentation is needed before using them
as quotient kernels.  For example, the condition that a $p$-adic
valuation be zero takes different values on the idele $1$ and
the principal idele $p$, although they represent the same class.

The narrow missing operator statement is therefore a compatible
factorization of the actual condensed global operator through
the local quotient, with the second cyclic composition equal
to the claimed local cutoff operator, together with control
of the limiting remainder and the polar finite-dimensional
part.  Range inclusion alone is insufficient for this use of
Lemma~C.3.2.2.  MA-LA-011 supplies a genuine local nuclear
diagram, and MA-LA-013 supplies its prime-bearing unitary
factor, but neither identifies that diagram with the global
weighted orthogonal projection.  The precise actual global
cutoff value and its lost-divisor criterion are already given
by MA-WP and MA-EG, without adding this missing identification.
The bare inverse obstruction MA-MO and the smoothed
commutators MA-RS retain their separate scopes.

All source locators in this section refer to the inspected
public Masullo draft; no source PDF, text extraction, or
rendered source page is included in the deliverable.

\section{The corrected coupled cutoff limit on the rational local quotient}
\label{sec:local-quotient-limit}

\paragraph{MA-LA-014 (the actual quotient limit includes the finite-prime terms).}
Let $S_f$ be a finite set of rational primes and let
$D_S=\prod_{p\in S_f}D_p$, with $D_p$ as in MA-LA-011.
All these factors are commuting unitaries on even
$L^2(\mathbb R,dx)$.  The quotient Fourier transform
\[
                      H_S=D_S\mathcal F_e
 \tag{MA-LA.67}
\]
is a unitary involution.  It is the finite-place product
of the actual section-periodization in (MA-LA.50).

Choose even smooth profiles $A,B$, between zero and one,
equal to one on $[-1,1]$ and supported on $[-2,2]$.
For $X\ge2\epsilon>0$ and $Y>0$, set
\[
 \begin{aligned}
 a_{\epsilon,X}(x)
    &=A(x/X)-A(x/\epsilon)
      =A(x/X)(1-A(x/\epsilon)),\\
 Q_Y&=\mathcal F_e M_{B(\xi/Y)}\mathcal F_e,\\
 T^S_{X,Y,\epsilon}(h)
    &=M_{a_{\epsilon,X}}A_hH_SM_{B(\xi/Y)}H_S
      =M_{a_{\epsilon,X}}A_hD_SQ_YD_S^* .
 \end{aligned}
 \tag{MA-LA.68}
\]
For $h\in C_c^\infty(\mathbb R_{>0})$ this is an actual
trace-class operator on the stated quotient completion,
with a nuclear lift to its flat unit-shell section.
If $h(1)=0$, then
\[
 \begin{aligned}
 \operatorname{Tr}T^S_{X,Y,\epsilon}(h)
       &\longrightarrow A_0(h)+B_S(h)
           &&\text{if }XY\to\infty,\quad\epsilon Y\to0,\\
 \operatorname{Tr}T^S_{X,Y,\epsilon}(h)
       &\longrightarrow0
           &&\text{if }\epsilon Y\to\infty,
 \end{aligned}
 \tag{MA-LA.69}
\]
where $A_0$ is (MA-LA.20) and
\[
 B_S(h)=\sum_{p\in S_f}\sum_{k\ge1}
       (\log p)p^{-k/2}\bigl(h(p^k)+h(p^{-k})\bigr).
 \tag{MA-LA.70}
\]
The second hypothesis implies $XY\to\infty$ because
$X\ge2\epsilon$.
No zero-location statement is used.

\begin{proof}
The operators $D_p$ are functions of commuting multiplicative
dilations.  Thus their finite product is unitary, and
$\mathcal F_eD_S\mathcal F_e=D_S^*$ proves
$H_S^2=1$.  At each finite place the Fourier transform
of the unit-shell section is (MA-LA.48).  Its folded
series converges absolutely in operator norm by
(MA-LA.50)--(MA-LA.51).  The finite product of these
series is therefore absolutely convergent as well and
is precisely $D_S\mathcal F_e$.  This proves the asserted
local quotient identification, including its topology.
Intermediate vectors are allowed to lie in the Hilbert
completion; they are not asserted to remain polar-free.

The multiplier identity for $a_{\epsilon,X}$ holds because
$A(x/\epsilon)\ne0$ implies $|x|\le2\epsilon\le X$,
where $A(x/X)=1$.
Finite-cutoff trace class and the nuclear lift follow
from (MA-LA.53)--(MA-LA.54), since
$H_SM_{B(\xi/Y)}H_S$ is bounded.

We now compare the trace with the real-place operator
\[
                    T^\varnothing_{X,Y,\epsilon}(h)
                          =M_{a_{\epsilon,X}}A_hQ_Y.
 \tag{MA-LA.71}
\]
Pass to logarithmic coordinates by the unitary map
$u\mapsto\sqrt2e^{y/2}u(e^y)$, put $g(y)=h(e^y)$,
and write $\Phi(y)=A(e^y)$.  Then $A_h=C_g$ in these
coordinates; $\Phi(-\infty)=1$, $\Phi(+\infty)=0$,
and $\Phi'$ has compact support.
Define the operator with convolution in the indicated order:
\[
                E_S(g)=(D_S^*M_\Phi D_S-M_\Phi)C_g.
 \tag{MA-LA.72}
\]
It is trace class and has trace $B_S(h)$.
For completeness this order follows directly from
MA-LA-013: the difference in parentheses is self-adjoint,
so (MA-LA.72) is the adjoint of
$C_{g^*}(D_S^*M_\Phi D_S-M_\Phi)$, where
$g^*(y)=\overline{g(-y)}$.
The latter is trace class and has trace $B_S(g^*)$;
complex conjugation gives $B_S(g)$ because the local sum
contains both directions with real coefficients.

Let $(W_Rf)(y)=f(y+R)$ be the unitary log translation.
Every $D_p$ and $C_g$ commutes with $W_R$.
Writing $\Phi_R(y)=\Phi(y-R)$, we have
\[
 (D_S^*M_{\Phi_R}D_S-M_{\Phi_R})C_g
                          =W_{-R}E_S(g)W_R.
 \tag{MA-LA.73}
\]
In additive coordinates $W_R$ is the normalized dilation
$u(x)\mapsto e^{R/2}u(e^Rx)$.  Fourier scaling gives the
exact bounded-operator identity
\[
                       W_RQ_YW_{-R}=Q_{e^RY}.
 \tag{MA-LA.74}
\]
All products in the following trace calculation are trace
class: the annular $M_aC_g$ is trace class and the separated
endpoint operators (MA-LA.73) are trace class.  Bounded
cyclicity, and commutation of $C_g$ with $D_S$, therefore give
\[
 \begin{aligned}
 \operatorname{Tr}T^S_{X,Y,\epsilon}(h)
       -\operatorname{Tr}T^\varnothing_{X,Y,\epsilon}(h)
   &=\operatorname{Tr}\bigl((D_S^*M_aD_S-M_a)C_gQ_Y\bigr)\\
   &=\operatorname{Tr}\bigl(E_S(g)Q_{XY}\bigr)
                   -\operatorname{Tr}\bigl(E_S(g)Q_{\epsilon Y}\bigr).
 \end{aligned}
 \tag{MA-LA.75}
\]
Here $M_a=M_{\Phi_{\log X}}-M_{\Phi_{\log\epsilon}}$
in logarithmic coordinates.  This proof does not commute
$C_g$ past $Q_Y$.

The uniformly bounded family $Q_L$ tends strongly to the
identity as $L\to\infty$ and to zero as $L\to0$.
Indeed its Fourier multiplier $B(\xi/L)$ tends respectively
to one and to zero almost everywhere, and dominated
convergence proves the strong limits.  Pairing a uniformly
bounded strongly convergent sequence with a trace-class
operator is continuous for the trace; this follows first
for finite-rank operators and then by trace-norm approximation.
Consequently the right side of (MA-LA.75) tends to
$\operatorname{Tr}E_S(g)=B_S(h)$ when $XY\to\infty$
and $\epsilon Y\to0$.  If $\epsilon Y\to\infty$, both
terms tend to the same trace, so their difference tends to zero.
The actual real-place limits in MA-LA-007 now give
(MA-LA.69), on the full class $h(1)=0$.
\end{proof}

\paragraph{The repaired diagram and the remaining global issue.}
The section on the flat additive LF space, its Hilbert
Fourier extension $H_S$, and the nuclear annular smoothing
are actual maps.  Their common nuclear presentation justifies
the quotient trace.  Equations (MA-LA.72)--(MA-LA.75) show
exactly how the prime terms survive the coupled cutoff limit.
The same proof works algebraically for the unitary phase
factors (MA-LA.65), replacing the prime sum by (MA-LA.66)
and using $D_SQ_YD_S^*$; this does not assert an additional
character-sector section that has not been constructed.

The quotient construction does not make the full-cover
operator compact: MA-LA-012 still applies upstairs.
It also does not identify the global weighted orthogonal
Paley--Wiener condensation with this local quotient operator.
That source bridge remains the one isolated after
MA-LA-013 and measured exactly by MA-WP/MA-EG.
The present theorem closes the coupled local rational
trivial-sector limit, including its finite-prime terms,
without importing that missing global identification.

\section{Trace class for the actual convolution-after-projection operator}
\label{sec:actual-lift-trace}

This section repairs a separate operator-order issue in the real
weighted model. It uses the actual quotient theorem MT-M26, with
its uniform interior containing-strip hypothesis. The new ingredient
is the differential form of translation on the specified weighted
$L^2$ space, including the finite boundary contribution from the
cusp of the weight at zero.

Fix $\alpha>0$, $s\ge0$, and write
\[
 w(x)=e^{2\alpha|x|}(1+x^2)^{s/2},\qquad
 H=L^2(\mathbb R,w(x)\,dx),\qquad (U_tf)(x)=f(x-t).
 \tag{MA-LIFT.1}
\]
Let $J$ be the Hilbert closure of the centered theta ideal, let
$Q=H/J$, and identify $Q$ isometrically with $Z=J^\perp$ by its
orthogonal representative. Write $P_Z$ and $P_G=1-P_Z$ for the
orthogonal projections. For $g\in C_c^\infty(\mathbb R)$, put
$C_gf=\int g(t)U_tf\,dt$. A multiplicative test $h$ corresponds
to $g(t)=h(e^t)$. All results below are for one fixed character
sector; no summation over varying conductors is inferred.

\paragraph{Input and zero-location scope.}
Assume precisely the hypothesis of MT-M26: a fixed closed strip
$[c_0,c_1]$ contains the real parts of all zeros of the relevant
completed function, and
\[
 [c_0-1/2,c_1-1/2]\subset(-\alpha,\alpha).
 \tag{MA-LIFT.2}
\]
For rational zeta one can take $[c_0,c_1]=[0,1]$ unconditionally,
so $\alpha>1/2$ suffices. Under the adopted Dirichlet input $A$,
one can take $c_0=c_1=1/2$ in those sectors, so every $\alpha>0$
suffices. For general Hecke characters this latter instantiation
uses the conditional input $M$. No unconditional partial-strip
version of MT-M26 is assumed.

\paragraph{MA-LIFT-001 (a trace-class resolvent power on the quotient).}
Under (MA-LIFT.2), let $A_Q$ generate the induced translation group
on $Q$. For a sufficiently large real $\kappa$ and an integer $k\ge2$,
\[
                  (\kappa-A_Q)^{-k}\in\mathcal S_1(Q).
 \tag{MA-LIFT.3}
\]
One may choose $k>B+2$, where $B$ is the polynomial exponent in
MT-M26's fixed-strip lifting bound (M26.9).

\emph{Proof.}
MT-M26 supplies finite block algebras $J_n$, uniformly bounded
maps $r_n:Q\to J_n$, polynomially bounded maps $s_n:J_n\to Q$, and
\[
 r_m s_n=\delta_{mn}\operatorname{id},\qquad
 d_n=\dim J_n=O(\log(2+|n|)),\qquad
 \|s_n\|\le C(2+|n|)^B.
 \tag{MA-LIFT.4}
\]
Their images span a dense subspace of $Q$: the smooth quotient has
absolute block synthesis, maps densely into $Q$, and its block
series converges in $Q$ by the polynomial lift bound. On
$s_n(J_n)$, translation is multiplication by the block class of
$e^{(z-1/2)t}$, by the smooth action and the exact zero-jet kernel
statement (M26.3).

Translation satisfies
\[
 \|U_t\|\le 2^{s/4}e^{\alpha|t|}(1+t^2)^{s/4}.
 \tag{MA-LIFT.5}
\]
The same holds on $Q$. Choose $\kappa>\alpha+1$; its excess
exponential decay absorbs the polynomial factor. The resolvent
power is the operator-norm convergent Laplace integral
\[
 (\kappa-A_Q)^{-k}=\frac1{(k-1)!}
 \int_0^\infty e^{-\kappa t}t^{k-1}U_t^Q\,dt.
 \tag{MA-LIFT.6}
\]
On block $n$ this is multiplication by
$F_k(z)=(\kappa-(z-1/2))^{-k}$. The contours used in MT-M26
lie in a fixed vertical strip, at heights comparable to $n$.
The quotient algebra norm thus gives
$\|M_{F_k,n}\|\le C_k(2+|n|)^{-k}$, and
\[
 \sum_n\|s_nM_{F_k,n}r_n\|_1
 \le C_k\sum_n d_n(2+|n|)^{B-k}<\infty\qquad(k>B+2).
 \tag{MA-LIFT.7}
\]
This trace-class sum agrees with (MA-LIFT.6) on each finite
block, hence on their dense span and then on $Q$.
\hfill$\square$

\paragraph{MA-LIFT-002 (the specified orthogonal lift is trace class).}
Under (MA-LIFT.2), the actual operator $S_g=C_gP_Z:H\to H$
is trace class for every $g\in C_c^\infty(\mathbb R)$, and
\[
 \operatorname{Tr}_H(C_gP_Z)
 =\operatorname{Tr}_Q(C_g^Q)
 =\sum_\rho m_\rho\,\mathcal Mg(\rho-1/2).
 \tag{MA-LIFT.8}
\]
The sum is absolutely convergent. On every fixed compact
test-support space, $g\mapsto S_g$ is continuous into trace norm.

\emph{Proof.}
The adjoint translation group is explicitly
\[
 (U_t^*u)(x)=\frac{w(x+t)}{w(x)}u(x+t),\qquad
 Lu=w^{-1}\partial_x(wu).
 \tag{MA-LIFT.9}
\]
It is strongly continuous by (MA-LIFT.5) and compactly supported
smooth approximation; $L$ is its generator. Derivatives in this
formula are weak derivatives. Multiplication by $w$ is a unitary
map from $H$ onto $L^2(w^{-1}dx)$, intertwining $U_t^*$ with
$v(x)\mapsto v(x+t)$. Consequently $u\in\mathcal D(L^k)$ exactly
when $v=wu$ has weak derivatives through order $k$ in
$L^2(w^{-1}dx)$, and $L^ju=w^{-1}\partial^jv$.
For the first derivative this follows from
$v(\cdot+t)-v=\int_0^t v'(\cdot+r)\,dr$ and strong continuity
of translation; the converse follows by testing a generator
difference quotient against compact smooth functions. Iteration
gives the asserted power domains.

Since $J$ is invariant under every $U_t$, $Z$ is invariant under
every $U_t^*$. This restricted group is the Hilbert adjoint of
the quotient group. MA-LIFT-001 gives
\[
 R_Z^k:=(\kappa-L|_Z)^{-k}\in\mathcal S_1(Z).
 \tag{MA-LIFT.10}
\]
Here $L|_Z$ means the generator of the restricted group. Its
resolvent is the restriction of $R=(\kappa-L)^{-1}$, by the
Laplace integral. The extended map
$\mathcal R_k=R^kP_Z:H\to H$ is therefore trace class and takes
values in $\mathcal D(L^k)$. It is bounded in the graph norm:
$L^jR^k$ is bounded for $j\le k$, using $LR=\kappa R-1$.

We now use the specified weight. On each open half-line it is
smooth, and
\[
 a_j(y):=w(y)\partial_y^j(w(y)^{-1})\quad(y\ne0)
       \ \in L^\infty(\mathbb R)\qquad(j\ge0).
 \tag{MA-LIFT.11}
\]
Indeed, $\log w=\pm2\alpha y+(s/2)\log(1+y^2)$ there; all
positive-order derivatives of this expression are bounded,
and each $a_j$ is a polynomial in them. All one-sided
derivatives of $w^{-1}$ at zero also exist.

For $u\in\mathcal D(L^k)$ put $v=wu$. Then
\[
 C_g(\kappa-L)^ku(x)=
 \int_{\mathbb R}\frac{g(x-y)}{w(y)}
                    (\kappa-\partial_y)^kv(y)\,dy.
 \tag{MA-LIFT.12}
\]
Integrate by parts separately on $y<0$ and $y>0$. Terms at
infinity vanish because $g(x-y)$ has compact support in $y$.
The bulk term is
\[
 B_{g,k}u(x)=\int_{y\ne0}
 w(y)(\kappa+\partial_y)^k
           \left(\frac{g(x-y)}{w(y)}\right)u(y)\,dy.
 \tag{MA-LIFT.13}
\]
Expansion makes this a finite sum of $C_{g^{(i)}}M_{a_j}$
with coefficients depending on $\kappa,k$. Each summand is
bounded on $H$ by (MA-LIFT.5), the weighted $L^1$ norm of
$g^{(i)}$, and (MA-LIFT.11). Hence $B_{g,k}$ is bounded.

The remaining terms have the form
\[
 F_{g,k}u=\sum_{m=0}^{k-1}\phi_{g,k,m}\,\ell_m(u),
 \qquad \ell_m(u)=\partial^m(wu)(0).
 \tag{MA-LIFT.14}
\]
Each $\phi_{g,k,m}$ is a finite linear combination of
$g,g',\ldots,g^{(k-1)}$ with coefficients given by jumps of
one-sided derivatives of $w^{-1}$ at zero. In particular
$\phi_{g,k,m}\in C_c^\infty(\mathbb R)\subset H$.
Explicitly the boundary terms of the derivative of order $j$
are linear combinations of
\[
 \bigl[\partial_y^l(g(x-y)/w(y))\bigr]_{0-}^{0+}
                     v^{(j-1-l)}(0),\qquad 0\le l<j\le k.
\]
These common derivatives of $v$ exist since $v\in H^k_{\rm loc}$.
This is the ordinary Sobolev integration-by-parts formula;
the cusp $|x|$ has not been treated as globally smooth.

The one-dimensional Sobolev trace inequality gives
\[
 |\ell_m(u)|\le C_k\|wu\|_{H^k([-1,1])}
 \le C_{k,\alpha,s}\sum_{j=0}^k\|L^ju\|_H
 \qquad(0\le m\le k-1).
 \tag{MA-LIFT.15}
\]
For the last inequality use boundedness above and below of $w$
on this interval. Thus $\ell_mR^kP_Z$ is bounded on $H$.
The trace inequality needed here is elementary: the fundamental
theorem of calculus, followed by averaging and Cauchy--Schwarz,
bounds $|q(0)|$ by $C(\|q\|_2+\|q'\|_2)$ on $[-1,1]$.
Apply it to $q=v^{(m)}$, with $m+1\le k$.

Apply (MA-LIFT.12)--(MA-LIFT.14) to $u=\mathcal R_kf$.
Since $(\kappa-L)^ku=P_Zf$, the exact identity is
\[
 C_gP_Z=B_{g,k}\mathcal R_k+
     \sum_{m=0}^{k-1}\phi_{g,k,m}\otimes(\ell_m\mathcal R_k).
 \tag{MA-LIFT.16}
\]
The first term is trace class and the rest have finite rank.
More quantitatively,
\[
 \|C_gP_Z\|_1\le \|B_{g,k}\|\,\|R_Z^k\|_1+
   \sum_{m=0}^{k-1}\|\phi_{g,k,m}\|_H
                    \|\ell_mR^kP_Z\|_{H'}.
 \tag{MA-LIFT.17}
\]
All dependence on $g$ is controlled by finitely many derivatives
on its fixed compact support, proving the stated continuity.

Finally, trace cyclicity with the bounded projection gives
$\operatorname{Tr}(C_gP_Z)=\operatorname{Tr}(P_ZC_gP_Z)$.
The latter is the orthogonal representative of quotient
convolution, with zero action on $J$. MT-M26 computes its
trace as (MA-LIFT.8), with absolute convergence.
\hfill$\square$

\paragraph{MA-LIFT-003 (actual cutoff spectral and geometric traces).}
Under the same hypotheses, let $\rho_\nu\in C_c^\infty(\mathbb R)$
have uniformly bounded suprema and tend to $1$ almost everywhere.
Then
\[
       \|(M_{\rho_\nu}-1)C_gP_Z\|_1\longrightarrow0.
 \tag{MA-LIFT.18}
\]
For every compactly supported smooth $\rho$, $M_\rho C_g$
is trace class on $H$ and
\[
 \operatorname{Tr}_H(M_\rho C_g)=g(0)\int_{\mathbb R}\rho(x)\,dx.
 \tag{MA-LIFT.19}
\]
In particular, if $g(0)=0$, the actual geometric-order operator
$G_{\rho,g}=M_\rho C_gP_G$ is trace class, and
\[
 \lim_\nu\operatorname{Tr}_H G_{\rho_\nu,g}
       =-\sum_\rho m_\rho\,\mathcal Mg(\rho-1/2).
 \tag{MA-LIFT.20}
\]
Here $\rho$ in the sum denotes zeros, independently of the
cutoff notation.

\emph{Proof.}
Dominated convergence gives $M_{\rho_\nu}\to1$ strongly on $H$,
with uniformly bounded operator norms. Strong convergence of a
uniformly bounded family, followed by a fixed trace-class
operator, is convergence in trace norm: first verify it for
finite-rank operators and then approximate in trace norm.
MA-LIFT-002 therefore gives (MA-LIFT.18).

On ordinary $L^2(\mathbb R)$ the kernel $\rho(x)g(x-y)$ is
smooth and compactly supported in both variables, so it is
trace class with trace $g(0)\int\rho$. An absolutely summable
Hermite expansion proves this as in MA-LA-003.
Conjugate to $H$ by multiplication by $\sqrt w$. Choose a smooth
compact $\eta$ equal to one on both coordinate projections of
the kernel support. The conjugated operator is
\[
 M_{\eta\sqrt w}\,(M_\rho C_g)_{L^2}\,M_{\eta/\sqrt w}.
\]
Both outside factors are bounded. Trace cyclicity replaces their
product by $M_{\eta^2}$, which leaves the middle kernel unchanged.
This proves (MA-LIFT.19), including trace class, without requiring
smoothness of the weight at zero.

Now $G_{\rho,g}=M_\rho C_g-M_\rho C_gP_Z$ is a difference of
trace-class operators. If $g(0)=0$, its trace equals
$-\operatorname{Tr}(M_\rho C_gP_Z)$ by (MA-LIFT.19).
Use (MA-LIFT.18) and (MA-LIFT.8) to pass to the limit.
\hfill$\square$

\paragraph{What this closes, and what it does not.}
MA-LIFT-002 replaces the unfinished lifted trace-class estimate
in Masullo, p.~464, equations~(10.4.2.25)--(10.4.2.34), for this
specified fixed-sector weighted model under (MA-LIFT.2).
MA-LIFT-003 supplies cutoff convergence in the operator order
of Theorem~10.4.2.8. It discharges the actual lifted trace-class
premise of F-SM-18/F-SM-21 in the rational wide-strip case, and
in every rational narrow strip when $A$ is explicitly used.
It does not contradict the abstract obstruction MA-TR-102.

The arithmetic quotient is used only through the verified
finite block data and the uniform interior strip hypothesis.
The rest is a calculation with the specified weight and actual
translations. Equation (MA-LIFT.20) is a trace theorem on that
actual global Hilbert space, rather than a scalar assigned from
the explicit formula. Nevertheless its computation is spectral.
Identifying it with an independently transported local geometric
kernel carrying the prime terms remains the global-to-local
question. Using $A$ or $M$ to supply (MA-LIFT.2) in every narrow
strip and then presenting this as an independent proof of their
zero-location conclusion would be circular.


\section{The canonical prime-carrying inverse has no bounded narrow-strip extension}
\label{sec:mobius-unbounded}

The finite-Euler remainder of MA-TR-101 rules out one proposed
limit argument. A different route gives a stronger obstruction
for the actual smooth theta--M\"obius inverse itself. It uses
small values of zeta on a fixed vertical line, not the location
of any zero. The obstruction therefore persists even when the
Dirichlet theorem $A$ is available.

\paragraph{MA-MO-001 (unboundedness of the canonical inverse).}
Let $0<\alpha<1/2$ and $s\ge0$, and give both the radial
half-density variable $F$ and the centered theta variable
$g=\mathcal ZF$ the norm
\[
 \|u\|_{\alpha,s}^2=
 \int_{\mathbb R}|u(x)|^2
             e^{2\alpha|x|}(1+x^2)^{s/2}\,dx,
 \qquad x=\log t.
 \tag{MA-MO-1}
\]
The inverse $\mathcal J$ in (MA-TR-7), on the actual smooth
theta image, has no bounded extension from
$J_{\alpha,s}$ to $H_{\alpha,s}$ agreeing with that inverse.
Equivalently, no constant $C$ satisfies
\[
        \|F\|_{\alpha,s}\le C\|\mathcal ZF\|_{\alpha,s}
 \tag{MA-MO-2}
\]
even on the smaller domain
$\{F\in\mathcal E:\mathcal MF(1/2)=0\}$.
Consequently there are smooth pairs $g_j=\mathcal ZF_j$ with
$\|F_j\|_{\alpha,s}=1$ and $\|g_j\|_{\alpha,s}\to0$.

\begin{proof}
First this smaller domain is contained in the one used in
MA-TR-101. If $F$ is exponential Schwartz as a function of
$x=\log t$, set $f(t)=F(\log t)/(2\sqrt t)$ for $t>0$,
extend it evenly, and set $f(0)=0$. Every derivative tends
to zero at the origin and decays faster than every power at
infinity; this follows by differentiating in $t$, which only
introduces fixed powers of $t^{-1}$ and logarithmic derivatives
of $F$. Thus $f$ is even Schwartz and flat at zero. Moreover
\[
       \int_{\mathbb R}f(t)\,dt=\mathcal MF(1/2)=0.
\]
MA-TR-101 therefore gives $g=\mathcal ZF\in\mathcal E$ and
the exact identity
\[
 \mathcal Mg(z)=\zeta(z+1/2)\mathcal MF(z).
 \tag{MA-MO-3}
\]
Initially this follows by absolute convergence for $\Re z>1/2$;
the zero moment removes the only zeta pole, so analytic
continuation gives it everywhere. No zero-free region is used.

Suppose (MA-MO-2) holds. Write $U_XF(x)=F(x-X)$. It preserves
the domain because its Mellin transform is
$e^{Xz}\mathcal MF(z)$, and $\mathcal ZU_X=U_X\mathcal Z$.
For each fixed $F\in\mathcal E$,
\[
 \frac{\|U_XF\|_{\alpha,s}^2}
      {e^{2\alpha X}(1+X^2)^{s/2}}
 \longrightarrow
       \int_{\mathbb R}|F(u)|^2e^{2\alpha u}\,du
       \qquad(X\to+\infty).
 \tag{MA-MO-4}
\]
Indeed the ratio of weights tends pointwise to $e^{2\alpha u}$.
For $X\ge1$ it is bounded by
$2^{s/2}e^{2\alpha|u|}(1+u^2)^{s/2}$, using
$|X+u|-X\le|u|$ and
$1+(X+u)^2\le2(1+X^2)(1+u^2)$.
This is an integrable dominating weight against $|F|^2$.
The same argument applies to $g\in\mathcal E$.
Apply (MA-MO-2) to the translated pair and pass to the limit:
\[
 \int|F(x)|^2e^{2\alpha x}\,dx
 \le C^2\int|g(x)|^2e^{2\alpha x}\,dx.
\]
Plancherel on the vertical Mellin line and (MA-MO-3) give
\[
 \int_{\mathbb R}|\mathcal MF(\alpha+it)|^2\,dt
 \le C^2\int_{\mathbb R}
     |\zeta(1/2+\alpha+it)|^2
                  |\mathcal MF(\alpha+it)|^2\,dt.
 \tag{MA-MO-5}
\]
The common Fourier normalization constant cancels.

Fix any $t_0\in\mathbb R$ and put, for $\delta>0$,
\[
 \phi_\delta(x)=e^{-\delta x^2-\alpha x-it_0x},
 \qquad F_\delta=(\partial_x+1/2)\phi_\delta.
\]
These are in $\mathcal E$, and integration by parts gives
$\mathcal MF_\delta(z)=(1/2-z)\mathcal M\phi_\delta(z)$.
In particular their moment at $1/2$ vanishes. On the required
line, their squared Mellin transform is a positive constant
times
\[
 \bigl((1/2-\alpha)^2+t^2\bigr)
                     e^{-(t-t_0)^2/(2\delta)}.
\]
After dividing by its integral, this is an approximate point
mass at $t_0$, because $1/2-\alpha>0$.
The classical fixed-strip polynomial growth bound for zeta,
together with the Gaussian tail, justifies this limit also
against the factor $|\zeta(1/2+\alpha+it)|^2$.
Thus (MA-MO-5) implies
\[
          |\zeta(1/2+\alpha+it_0)|\ge C^{-1}
                     \qquad\hbox{for every }t_0\in\mathbb R.
 \tag{MA-MO-6}
\]

This contradicts the classical universality theorem of Voronin
(1975), the unnumbered theorem in \S I, pp.~443--444.
For clarity its consequence needed here is only
\[
       \inf_{t\in\mathbb R}|\zeta(\sigma+it)|=0
                    \qquad(1/2<\sigma<1).
\]
It follows directly from the theorem as printed, without using
a later stronger formulation: choose
$|\sigma-3/4|<r<1/4$, approximate the nonzero constant
$\varepsilon$ on the disc centered at $3/4$ with radius $r$
by a vertical translate of zeta to error $\varepsilon$, and
evaluate at its real point $\sigma$.
This proves the contradiction. Finally the failure of every
bound (MA-MO-2) supplies pairs with
$\|F_j\|>j\|g_j\|$; rescale to obtain the asserted sequence.
\end{proof}

\paragraph{Corollary: the original Schwartz target does not admit a continuous extension either.}
The inverse, now read as $g\mapsto f$ in
$\mathscr S_{0,+}\subset\mathscr S(\mathbb R)$, has no continuous
extension $J_{\alpha,s}\to\mathscr S_{0,+}$ for
$0<\alpha<1/2$. Indeed the map
$f\mapsto F(t)=2\sqrt t\,f(t)$ is continuous from this Schwartz
subspace into $H_{\alpha,s}$. Near zero, evenness and $f(0)=0$
give $|f(t)|\le(t^2/2)\sup_{|u|\le1}|f''(u)|$; at infinity
any sufficiently high Schwartz seminorm suffices. Explicitly,
\[
 \|F\|_{\alpha,s}^2
 =4\int_0^\infty |f(t)|^2
       \max(t^{2\alpha},t^{-2\alpha})
                    (1+\log^2t)^{s/2}\,dt,
\]
and the displayed bounds make this integral a continuous
Schwartz seminorm estimate (the near-zero estimate even allows
$\alpha<5/2$). Composing the proposed continuous extension
with this embedding contradicts MA-MO-001. This does not
exclude individually regularized maps obtained by additional
spatial and Fourier cutoffs.

\paragraph{Precise consequence for the proposed transport.}
The theorem rules out a bounded Hilbert-space realization of
the \emph{specified} inverse with these two natural radial
weights. Since a bounded extension is determined by its values
on the dense smooth geometric image, there is no different
bounded extension that agrees with that same arithmetic map
there. This is not an assertion that every prime-corrected
intertwiner between every possible pair of completions is
impossible. In particular, it leaves open an unbounded closed
operator, a target norm transported through $\mathcal Z$, a
different scale of spaces, and a directly controlled nuclear
regularization. Nor does the theorem contradict MA-LIFT-002,
which proves trace class for the actual convolution lift
without introducing this inverse as a bounded operator.

The source's missing step is still the passage in Section~12.3,
p.~514, used in the proof on p.~520, between the global weighted
geometric operator and the specified local one. MA-TR-101
constructs the smooth map with the needed arithmetic content;
MA-MO-001 shows that declaring this map bounded on the narrow
spaces does not fill that step. The remaining problem is now
to construct and estimate an \emph{unbounded or differently
completed} trace transport, including its prime correction,
or prove that no such transport realizes the source's kernels.
The universality input concerns small nonzero values. It uses
neither $A$, $M$, nor a hypothesized zero outside the critical
line, and hence is not a circular zero-location argument.

\paragraph{Grothendieck placement.}
This is an analytic completion obstruction for the adelic
theta map on $C_{\mathbb Q}$, in the precise smooth-quotient
identification of GRH\_NIGHT MT-M24. It refines
Q-MAS-GEOMETRIC-NARROW and ALPOGE\_BRIDGES R-MAS-KERNEL by
excluding the canonical bounded Hilbert transport.
In cote~1, scans~079--080 and 111--112, as located by
BR-TR-01/02 and F-SM-11/12, Fredholm trace and determinant
operations require specified topologies and nuclearity.
A smooth topological isomorphism does not supply those
properties in a different Hilbert completion. The theorem
does not construct a sheaf, a motive, or an algebraic
correspondence, and there is no additional EGA/SGA algebraic
theorem identifying the two analytic norms. The corresponding
$K_0$ and $\lambda$-ring identities do not remove the missing
norm estimate.

\paragraph{Relation to the spectral-operator literature.}
The one-sided small-value obstruction is established background.
Herichi and Lapidus (2014), Theorem~3.11, p.~641, identify the
spectrum of $\zeta(\partial_c)$ with the closure of
$\zeta(c+i\mathbb R)$; their \S6, p.~652, records that this
closure is all of $\mathbb C$ for $1/2<c<1$ by universality.
Thus its bounded invertibility fails even when a vertical line
contains no zero. Their Theorem~3.12 concerns a different,
truncated quasi-invertibility condition, which is equivalent to RH.
MA-MO-001's specific deduction is the translation-limit transfer
of the small-value obstruction to the two-sided polynomially
weighted theta-image norm and its original Schwartz target.
No novelty is claimed for the underlying one-sided spectral
obstruction.

\paragraph{Primary reference and check boundary.}
Voronin, S. M. (1975). Theorem on the ``universality'' of the
Riemann zeta-function. \emph{Mathematics of the USSR-Izvestiya,
9}(3), 443--453.
\url{https://doi.org/10.1070/IM1975v009n03ABEH001485}.
The exact theorem statement and its two printed pages were
freshly inspected; the classical universality proof is cited
as background, not asserted to have been independently
re-proved here. The deduction above is new work of this run.

Herichi, H., \& Lapidus, M. L. (2014). Truncated infinitesimal
shifts, spectral operators and quantized universality of the
Riemann zeta function. \emph{Annales de la Facult\'e des Sciences
de Toulouse: Math\'ematiques, 23}(3), 621--664.
\url{https://doi.org/10.5802/afst.1419}.
The cited statements and the concluding universality application
were freshly inspected; their full proofs are carried background.

\section{The raw Euler-conjugated cutoff is unbounded on a narrow line}
\label{sec:euler-cutoff-obstruction}

Unboundedness of the canonical inverse in MA-MO-001 does not
by itself prove unboundedness inside a cutoff conjugation.
We now prove that stronger assertion for the actual one-sided
Euler multiplier model. The argument uses nonzero universal
values, not an assumed zero on the line.

Fix \(1/2<\sigma<1\), put \(a=\sigma-1/2\), and use
\[
 \mathcal H_a=L^2(\mathbb R,e^{2ax}\,dx),\qquad
 V_a=\mathcal F_+M_{e^{ax}},\qquad
 \mathcal F_+f(t)=(2\pi)^{-1/2}\int e^{itx}f(x)\,dx.
 \tag{MA-MO.20}
\]
Let \(L_0(s)\) be zeta or a fixed primitive Dirichlet
\(L\)-function. For a finite set \(S\) of finite primes put
\[
 \begin{aligned}
 A_0(t)&=L_0(\sigma+it),&
 E_S(t)&=\prod_{p\in S}(1-\chi(p)p^{-\sigma-it}),\\
 A_S(t)&=E_S(t)A_0(t),&
 j_S(t)&=A_S(t)^{-1}.
 \end{aligned}
 \tag{MA-MO.21}
\]
For zeta, \(\chi(p)=1\); a ramified Dirichlet value zero
gives the factor one. The closed operators \(Z_{a,S}\)
and \(J_{a,S}\) are the maximal multipliers \(M_{A_S}\)
and \(M_{j_S}\), conjugated by \(V_a\), as in MA-RS.9
and its character version MA-RS-005.

Take any nonconstant bounded \(\phi\in C^\infty(\mathbb R)\)
with compactly supported derivative. Its limiting constants
may agree; nonconstant compact position cutoffs are allowed.
Put
\[
 T_\phi=\mathcal F_+M_\phi\mathcal F_+^{-1},\qquad
 p(v)=\int e^{ivx}\phi'(x)\,dx,\qquad
 q(v)=\frac{i\,p(v)}{2\pi v}\quad(v\ne0),
 \tag{MA-MO.22}
\]
and define the dense Fourier core
\[
 \mathscr C_\sigma=C_c^\infty
    \bigl(\mathbb R\setminus\{t:L_0(\sigma+it)=0\}\bigr).
 \tag{MA-MO.23}
\]
No assertion that this line is zero free is needed.

\paragraph{MA-MO-002 (no bounded raw cutoff conjugation).}
The actual composition
\[
                    \mathcal Q_{\phi,S}
                         =Z_{a,S}M_\phi J_{a,S}
 \tag{MA-MO.24}
\]
is well defined on \(V_a^{-1}\mathscr C_\sigma\).
Neither it nor \(M_\phi-\mathcal Q_{\phi,S}\) has a
bounded extension to \(\mathcal H_a\).
In the Fourier picture there are unit vectors
\(v_M\in\mathscr C_\sigma\) with
\[
 \|M_{A_0}T_\phi M_{A_0^{-1}}v_M\|_2
                         \ge c\,e^{Md}/M\longrightarrow\infty
 \tag{MA-MO.25}
\]
for fixed \(c,d>0\). Finite Euler deletion changes only
a fixed positive comparison constant. The result remains
true on a line which is known to have no zero.

\begin{proof}
\emph{Domains before kernels.}
The line-zero set is discrete and has measure zero.
Removing small neighborhoods of the finitely many zeros
in each compact interval proves density of
\(\mathscr C_\sigma\) in \(L^2\).
The finite Euler factors do not vanish when \(\sigma>0\).
For \(v\in\mathscr C_\sigma\), \(j_Sv\in C_c^\infty\).
All derivatives of \(\phi\) are bounded, so \(M_\phi\)
preserves Schwartz space and \(T_\phi(j_Sv)\) is Schwartz.
The unconditional fixed-line polynomial growth of the
fixed zeta or Dirichlet function, already used in
MA-GW-001/004 and MA-RS, makes
\(A_ST_\phi(j_Sv)\) square integrable.
This proves every domain inclusion needed for (MA-MO.24).
No domain assertion for \(J_{a,S}M_\phi\) is being
silently inferred.

\emph{The actual cross-window kernel.}
The distribution kernel of \(T_\phi\) is
\((2\pi)^{-1}\widehat\phi_+(t-u)\).
Since \(\widehat{\phi'}_+(v)=-iv\widehat\phi_+(v)\),
the kernel off the diagonal is \(q(t-u)\).
On disjoint compact input and output windows the
already defined composition for \(S=\varnothing\)
therefore has the ordinary kernel
\[
                      \frac{A_0(t)}{A_0(u)}q(t-u).
 \tag{MA-MO.26}
\]
There is no delta or principal-value term on these
disjoint windows. Smooth core pairings can be computed
from this kernel.

\emph{The exact universality statement used.}
Voronin (1975), the unnumbered theorem in \S I,
printed pp.~443--444, gives uniform approximation of
every nonvanishing analytic target on \(|z|\le r<1/4\)
by \(z\mapsto\zeta(z+3/4+iT)\) for a real \(T\).
The following Remark on p.~444 gives the same assertion
for Dirichlet \(L\)-functions.
We use precisely those printed forms. No density
theorem or assertion that the approximating shifts
can be chosen arbitrarily large is required.
The classical universality proofs are published
background; the operator deduction is proved here.

Set \(h=\sigma-3/4\), and choose
\[
                      |h|<r<1/4 .
 \tag{MA-MO.27}
\]
The transform \(p\) is entire because \(\phi'\) has
compact support. It is not identically zero, by Fourier
injectivity and nonconstancy of \(\phi\).
The identity theorem therefore permits a fixed
sufficiently small \(d>0\) with
\[
                         h^2+d^2<r^2,\qquad q(d)\ne0.
 \tag{MA-MO.28}
\]
This also handles \(p(0)=0\) when the limiting values
of the cutoff agree.

Fix \(0<\epsilon<1/100\). For each positive integer \(M\),
the entire nonvanishing target
\[
                         f_M(z)=e^{-iM(z-h)}
 \tag{MA-MO.29}
\]
has minimum modulus \(e^{-Mr}\) on the disk.
Apply the printed theorem with absolute tolerance
\(\epsilon e^{-Mr}\). It gives a real \(T_M\) such that
\[
 L_0(z+3/4+iT_M)=f_M(z)(1+\eta_M(z)),\qquad
        |\eta_M(z)|<\epsilon\quad(|z|\le r).
 \tag{MA-MO.30}
\]
The relative error follows by division by the
displayed minimum modulus; it is not a stronger
universality input. On \(z=h+iv\) the target is \(e^{Mv}\).
In particular the windows below contain no line zero.

\emph{A lower bound for the actual operator norm.}
Choose a fixed nonnegative real
\(b\in C_c^\infty((-1,1))\), with \(\|b\|_2=1\), and put
\[
 \begin{aligned}
 v_M(u)&=M^{1/2}b(M(u-T_M)),\\
 w_M(t)&=M^{1/2}b(M(t-T_M-d)).
 \end{aligned}
 \tag{MA-MO.31}
\]
These are unit vectors. For large \(M\) their supports
are disjoint and lie in the two vertical windows strictly
inside the fixed disk. In particular \(v_M\) is in the
declared core.
Substitute \(u=T_M+\xi/M\) and
\(t=T_M+d+\eta/M\) in (MA-MO.26). With the inner
product linear in its first argument, the pairing is
\[
 \begin{split}
 &\langle M_{A_0}T_\phi M_{A_0^{-1}}v_M,w_M\rangle\\
 &\quad=\frac{e^{Md}}{M}\iint b(\xi)b(\eta)e^{\eta-\xi}
       q\left(d+\frac{\eta-\xi}{M}\right)
       \frac{1+\eta_M(h+i(d+\eta/M))}
            {1+\eta_M(h+i\xi/M)}\,d\xi\,d\eta .
 \end{split}
 \tag{MA-MO.32}
\]
The error ratio differs from one by at most
\(2\epsilon/(1-\epsilon)\).
Continuity of \(q\) and \(q(d)\ne0\) give, for large \(M\),
\[
 \sup_{|\xi|,|\eta|\le1}
 \left|q\left(d+\frac{\eta-\xi}{M}\right)/q(d)-1\right|<1/100.
\]
The product of these two normalized factors is within
\(1/2\) of one. The remaining weight is nonnegative,
with strictly positive integral
\[
 I_b=\left(\int b(\eta)e^\eta d\eta\right)
                   \left(\int b(\xi)e^{-\xi}d\xi\right)>0.
 \tag{MA-MO.33}
\]
Subtracting the main value \(q(d)I_b\) inside the
double integral proves
\[
 |\langle M_{A_0}T_\phi M_{A_0^{-1}}v_M,w_M\rangle|
               \ge\tfrac12|q(d)|I_b\,e^{Md}/M.
 \tag{MA-MO.34}
\]
The unit norm of \(w_M\) proves (MA-MO.25).
This argument proves a norm lower bound on actual
core vectors; it does not infer it from a large
kernel value at a single point.

\emph{Finite Euler factors.}
There are constants \(0<m_S\le M_S<\infty\) such that
\(m_S\le|E_S(t)|\le M_S\) for all real \(t\).
Use the factors \(1-|\chi(p)|p^{-\sigma}>0\)
and \(1+|\chi(p)|p^{-\sigma}\).
Both \(E_S\) and its inverse preserve the core, and
\[
 M_{A_S}T_\phi M_{j_S}
       =M_{E_S}
          (M_{A_0}T_\phi M_{A_0^{-1}})
                                M_{E_S}^{-1}.
 \tag{MA-MO.35}
\]
This is an exact bounded similarity on that core.
The normalized inputs \(E_Sv_M/\|E_Sv_M\|_2\)
give at least \(m_S/M_S\) times the lower bound
(MA-MO.25). Thus no finite \(S\) has a bounded
extension. Since \(T_\phi\) is bounded, subtracting
it cannot turn this unbounded operator into a
bounded one. Conjugation by \(V_a\) finishes the proof.
\end{proof}

\paragraph{MA-MO-003 (the rational smooth-image graph transfer).}
In the rational sector the obstruction already holds
on actual smooth theta-image vectors. More precisely,
the graph closure of the restriction of
\(\mathcal Q_{\phi,S}\) to those vectors on which its
closed multiplier composition is defined contains
its graph on \(V_a^{-1}\mathscr C_\sigma\).
Consequently no bounded estimate on that smooth image
can extend the declared composition to \(\mathcal H_a\).
This remains a statement about the specified multiplier
continuation of its output.

\begin{proof}
Fix \(v\in\mathscr C_\sigma\) and let \(q_0=j_Sv\),
a smooth compact function. Choose \(h_0\in\mathscr S(\mathbb R)\)
with
\[
          \mathcal F_+h_0(t)=q_0(t)/(1-\sigma-it).
 \tag{MA-MO.36}
\]
There is no real zero of this denominator.
For \(\epsilon>0\) put
\[
 H_\epsilon(x)=e^{-\epsilon x^2-ax}h_0(x),\qquad
 F_\epsilon=(\partial_x+\tfrac12)H_\epsilon .
 \tag{MA-MO.37}
\]
These functions belong to \(\mathcal E\).
Integration by parts gives
\(\mathcal M F_\epsilon(1/2)=0\).
Moreover
\[
 q_\epsilon:=V_aF_\epsilon
       =(1-\sigma-it)\mathcal F_+(e^{-\epsilon x^2}h_0)
                           \longrightarrow q_0
                               \quad\hbox{in Schwartz space}.
 \tag{MA-MO.38}
\]
For completeness, Gaussian multiplication tends to the
identity in every Schwartz seminorm: split into a fixed
compact interval and its complement; derivatives converge
uniformly on the compact part, while the product rule,
the Schwartz tails of \(h_0\), and uniformly polynomially
bounded derivatives of \(e^{-\epsilon x^2}\) for
\(0<\epsilon\le1\) control the complement.
Fourier transformation and the linear multiplier preserve
this convergence.

To see the actual theta-image assertion, define the even
archimedean function, first for \(t>0\), by
\[
                   f_\epsilon(t)=
                    F_\epsilon(\log t)/(2\sqrt t).
 \tag{MA-MO.39}
\]
Superexponential decay with all derivatives makes it
Schwartz and flat at zero. Its full real integral is
\(\mathcal M F_\epsilon(1/2)=0\).
Thus both theta polar conditions hold, and the actual
theta formula of MA-TR-101/MA-GP.1 gives
\[
 (\Theta f_\epsilon)(x)
       =\sum_{n\ge1}n^{-1/2}F_\epsilon(x+\log n)
                       \in T_{1/2}\mathcal I .
 \tag{MA-MO.40}
\]
On a right Mellin line its transform is the product
of zeta and the transform of \(F_\epsilon\).
The moment condition removes the sole zeta pole.
The actual smooth theta identity therefore continues
to the line \(\sigma\).
Finite Euler deletion is a finite combination of
translations and preserves this smooth ideal.
Consequently the actual smooth-image vector
\(g_\epsilon=Z_{a,S}F_\epsilon\) has
\(V_ag_\epsilon=A_Sq_\epsilon\).
It is in \(\mathcal D(J_{a,S})\), and
\(J_{a,S}g_\epsilon=F_\epsilon\), since multiplication
by \(j_S A_S\) is one almost everywhere.

The convergence (MA-MO.38), polynomial growth of \(A_S\),
and continuity of \(T_\phi\) on Schwartz space now give
\[
 \begin{aligned}
 V_ag_\epsilon&=A_Sq_\epsilon\longrightarrow A_Sq_0=v
                                                   &&\text{in }L^2,\\
 V_a\mathcal Q_{\phi,S}g_\epsilon
       &=A_ST_\phi q_\epsilon
                        \longrightarrow A_ST_\phi q_0
                                                   &&\text{in }L^2.
 \end{aligned}
 \tag{MA-MO.41}
\]
In particular each \(g_\epsilon\) lies in the composition
domain and its output lies in \(\mathcal H_a\), regardless
of whether that output is in the smooth ideal.
This is exactly graph approximation.
A bounded estimate on these smooth-image vectors would
pass to every core vector in (MA-MO.41), contradicting
MA-MO-002. The same conclusion holds for the difference
with \(M_\phi\), since that multiplier is bounded.
\end{proof}

\paragraph{The exact common-space consequence and its limits.}
For Meyer's source cutoff, the first diagonal entry of
the centered difference \(P_+^S-P_-\) is
\[
 D=Z_SM_{1-\phi}J_S-M_{1-\phi}
                   =M_\phi-Z_SM_\phi J_S
 \tag{MA-MO.42}
\]
on its source domain, as derived from the definitions
in MA-MEY.13--.14. The exact locator is Meyer (2005),
Lemmas~5.3--5.4 and the following projection displays,
then \texttt{eq:global\_perturbation} in the proof of
Lemma~5.6, author-manuscript pp.~39--42.
The already recorded sign correction is retained.

A bounded operator on
\(\mathcal H_a\oplus\mathcal H_{-a}\), \(0<a<1/2\),
which agrees with this multiplier continuation on its
declared core, would give a bounded \(D\) by coordinate
compression. MA-MO-002 rules that out. MA-MO-003 also
rules out obtaining the bound merely by restricting
this composition to the rational smooth theta image.
This settles that specific bounded narrow matrix route.
It does not identify the multiplier-continued output,
with its possible pole subtraction, with every
uncontinued geometric inverse composition.
The subsequent MA-MO-004 identifies the exact rational source
polar correction, and MA-MO-005 removes it on actual smooth
theta-image witnesses with unbounded literal raw outputs.
Those strengthened statements retain the distinction between
the full direct product and the smaller common-sum subspace.

Meyer's wide-weight nuclear theorem uses \(a>1/2\)
after centering and is unaffected.
On the generic lines and Gevrey cutoffs where MA-RS
checks the literal commutator domains,
\(D=Z_{a,S}[J_{a,S},M_\phi]\) on the stated core;
the present result proves this unsmoothed operator
unbounded. Its product with \(C_g\) can still be
trace class, as MA-RS-002 proves.
No unclosability assertion or impossibility of
every nonlocal transport is made.
The one-sided \(\mathcal H_a\) is not substituted
for the two-sided weighted PW space or its Gram factors.
The separate MA-MO-009 supplies the required primary
Hecke universality input and proves the non-induced
imaginary-quadratic extension, including literal
smooth ideal-sum inputs and outputs.

\paragraph{Trace-framework placement and source boundary.}
This is an analytic obstruction to bounded entries
of the specified Hilbert matrix, not a new algebraic
obstruction to the finite-length modules or their
\(K_0\) and exterior-power identities.
In the inherited cote~1 nuclear-trace framework,
scans~079--080 and~111--112 through
ALPOGE\_BRIDGES BR-TR-01/02, the relevant topology
must be justified before cyclicity is used.
The smoothed constructions do so on their stated
domains; the raw narrow matrix has no such bounded
extension. No motivic positivity or weight theorem
is inferred, and no zero-location theorem is an input.

Voronin, S. M. (1975). Theorem on the universality
of the Riemann zeta-function.
\emph{Mathematics of the USSR-Izvestiya, 9}(3), 443--453.
\url{https://doi.org/10.1070/IM1975v009n03ABEH001485}.
The exact disk theorem in \S I, pp.~443--444, and
the Dirichlet Remark on p.~444 were freshly read
in retrieved primary text. This is a statement-scope
check; the full classical universality proof remains
published background. The new deductions above
receive separate whole-proof checks.

\section{The exact polar correction and a pole-free raw diagonal obstruction}
\label{sec:euler-polar-identification}

The multiplier continuation in MA-MO-003 can be compared with the
original summation operator, including its polar term. In the rational
sector that comparison also produces actual smooth inputs on which the
uncontinued first diagonal is unbounded. All assertions below concern
the specified operators; they do not exclude cancellation after
restriction to a smaller common-sum space.

Fix \(0<a<1/2\), \(\sigma=a+1/2\), and a finite set \(S\) of rational
primes. Keep \(\mathcal H_a,V_a\) from (MA-MO.20), and set
\[
 \begin{aligned}
 \mathcal E&=\{U\in C^\infty(\mathbb R):e^{bx}U(x)\in
                    \mathscr S(\mathbb R)\ \hbox{for every }b\in\mathbb R\},\\
 \mathcal E_0&=\{F\in\mathcal E:\mathcal MF(1/2)=0\},\\
 \mathcal MU(z)&=\int_{\mathbb R}e^{zx}U(x)\,dx,\\
 E_S(s)&=\prod_{p\in S}(1-p^{-s}),\qquad
 \kappa_S=E_S(1)>0,\qquad p_0(x)=e^{-x/2}.
 \end{aligned}
 \tag{MA-MO.43}
\]
Here \(E_S(s)\) is the same finite Euler multiplier as in MA-MO.21,
now written as a function of its complex argument. Let \(\mathcal N_S\)
be the positive integers prime to every prime in \(S\). Write
\[
 Z_S^{\rm geom}U(x)=\sum_{n\in\mathcal N_S}n^{-1/2}U(x+\log n),
 \qquad T_lU(x)=U(x+l).
 \tag{MA-MO.44}
\]
The superscript means the actual convergent summation, with no Mellin
continuation or pole removal. The operators \(Z_{a,S},J_{a,S}\) remain
the maximal closed multipliers with symbols
\(A_S(t)=E_S(\sigma+it)\zeta(\sigma+it)\) and \(A_S(t)^{-1}\).
The geometric inverse \(J_S^{\rm geom}\) is the original Euler inverse
on Meyer's first smooth component, obtained on every wider
\(\mathcal H_b\), \(b>1/2\), by its absolutely convergent operator sum.

Use precisely the source orientation of the cutoff:
\(\phi\in C^\infty(\mathbb R)\) is one sufficiently far to the left
and zero sufficiently far to the right. In particular its derivative
is compactly supported and \(\phi\) is nonconstant. Define
\[
                 \ell_\phi(F)=\mathcal M(\phi F)(1/2).
 \tag{MA-MO.45}
\]
Symmetry \(\phi(x)+\phi(-x)=1\) will only be needed for the reflected
second-diagonal observation at the end.

\paragraph{MA-MO-004 (exact geometric and continued comparison).}
For every \(U\in\mathcal E\), the following equality holds as an
equality of smooth functions, and its right side belongs to
\(\mathcal H_a\):
\[
 Z_{a,S}U=Z_S^{\rm geom}U-\kappa_S\mathcal MU(1/2)p_0.
 \tag{MA-MO.46}
\]
If \(F\in\mathcal E_0\) and \(y=Z_S^{\rm geom}F\), then \(y\in
\mathcal E\) is an actual smooth theta-image vector and
\[
                       J_S^{\rm geom}y=J_{a,S}y=F.
 \tag{MA-MO.47}
\]
All the following compositions are defined on these vectors:
\[
 \begin{aligned}
 Q_\phi^{\rm geom}y&=Z_S^{\rm geom}(\phi F),\\
 Q_\phi^{\rm cont}y&=Z_{a,S}(\phi F),\\
 D^{\rm geom}y&=Z_S^{\rm geom}((1-\phi)F)-(1-\phi)y\\
              &=\phi y-Q_\phi^{\rm geom}y,\\
 D^{\rm cont}y&=\phi y-Q_\phi^{\rm cont}y.
 \end{aligned}
 \tag{MA-MO.48}
\]
Their exact comparison is
\[
 \begin{aligned}
 Q_\phi^{\rm cont}y
       &=Q_\phi^{\rm geom}y-\kappa_S\ell_\phi(F)p_0,\\
 D^{\rm cont}y
       &=D^{\rm geom}y+\kappa_S\ell_\phi(F)p_0.
 \end{aligned}
 \tag{MA-MO.49}
\]
Moreover
\[
 Q_\phi^{\rm geom}y\in\mathcal H_a
 \quad\Longleftrightarrow\quad
 D^{\rm geom}y\in\mathcal H_a
 \quad\Longleftrightarrow\quad \ell_\phi(F)=0.
 \tag{MA-MO.50}
\]
When these conditions hold, all four outputs in (MA-MO.48) belong
to \(\mathcal E\), and the continued and geometric outputs agree
exactly. The correction in (MA-MO.49) has rank one in the ambient
smooth-function space. It is not a bounded rank-one operator on
\(\mathcal H_a\), since \(p_0\notin\mathcal H_a\).

\begin{proof}
\emph{The original sum and its polar asymptotic.}
For \(U\in\mathcal E\), the sum (MA-MO.44) and all its derivatives
converge locally uniformly: on a fixed compact interval in \(x\),
the summands are bounded by \(C_N n^{-N-1/2}\) for arbitrary
\(N\). The same estimate with its \(x\)-dependence retained proves
rapid exponential decay, with every derivative, at the right end.
For \(S=\varnothing\), define the even function
\[
              u(t)=\frac{U(\log t)}{2\sqrt t}\quad(t>0),
              \qquad u(0)=0.
 \tag{MA-MO.51}
\]
It is Schwartz and flat at zero. Indeed, each derivative away from
zero is a finite sum of powers of \(t^{-1}\) times derivatives
of \(U(\log t)\), which decay faster than every power at both
ends. Also \(\int_{\mathbb R}u(t)\,dt=\mathcal MU(1/2)\).
The self-dual additive Fourier transform, with kernel
\(e^{-2\pi it\xi}\), and Poisson summation give
\[
 Z_\varnothing^{\rm geom}U(x)
   =\mathcal MU(1/2)e^{-x/2}
       +2e^{-x/2}\sum_{n\ge1}\widehat u(ne^{-x}).
 \tag{MA-MO.52}
\]
The last term and every logarithmic derivative decay faster than
every exponential as \(x\to-\infty\). This follows from the
Schwartz estimates for \(\widehat u\), after choosing their order
arbitrarily large and summing \(n^{-N}\). The condition \(u(0)=0\)
has removed the other Poisson polar term.

Finite deletion is the exact identity
\[
 Z_S^{\rm geom}
    =\prod_{p\in S}(1-p^{-1/2}T_{\log p})
                         Z_\varnothing^{\rm geom}.
 \tag{MA-MO.53}
\]
Absolute local convergence justifies reindexing. Since
\(T_{\log p}p_0=p^{-1/2}p_0\), the coefficient of the left polar
term becomes \(\kappa_S\mathcal MU(1/2)\). Thus the actual sum
belongs to \(\mathcal E\) exactly when this moment vanishes.
Without that cancellation it still lies in the first smooth
component of Meyer: every weight \(b>1/2\) makes its left polar
tail rapidly decreasing, and its right tail is already rapid.

\emph{Identification of the continued multiplier.}
For \(b>1/2\), the original Euler sum is bounded on
\(\mathcal H_b\), by the absolutely convergent norm sum
\(\sum_{n\in\mathcal N_S}n^{-b-1/2}\). Its Mellin transform there is
\[
 E_S(z+1/2)\zeta(z+1/2)\mathcal MU(z),\qquad \Re z=b.
 \tag{MA-MO.54}
\]
Move the inverse Mellin line to \(\Re z=a\). The Mellin transform
of \(U\) decreases faster than every power on this fixed strip,
whereas zeta has polynomial growth away from its sole pole.
Horizontal integrals therefore tend to zero. The only pole crossed
is \(z=1/2\), with residue
\(\kappa_S\mathcal MU(1/2)e^{-x/2}\) after multiplication by
the inverse Mellin kernel. The right-line integral is the
left-line integral plus this residue. This proves (MA-MO.46),
including its sign. Zeros of zeta are not poles in this step,
so no generic-line or zero-free hypothesis is used.

The maximal operator domain also follows directly:
\(V_aU\) is Schwartz and \(A_S\) has polynomial growth, so
\(A_SV_aU\in L^2\). The pole-subtracted function in (MA-MO.46)
decays rapidly at the left; at the right it is a rapidly
decreasing function minus a multiple of \(e^{-x/2}\), which is
square integrable with weight \(e^{2ax}\) because \(a<1/2\).

\emph{The two actual inverse domains.}
If \(F\in\mathcal E_0\), (MA-MO.52)--(MA-MO.53) show that
\(y\in\mathcal E\). Formula (MA-MO.51), now with \(U=F\),
has both value at zero and integral zero, so its theta sum is
the actual smooth rational theta image, as in MA-TR-101 and
MA-MO.39--.40. Finite Euler deletion is a finite linear
combination of translations and preserves that image.
On each \(\mathcal H_b\), \(b>1/2\), the original Euler inverse
is bounded and inverse to the original sum. Consequently it
returns \(F\) on \(y\); this is the source smooth inverse,
not an independently chosen inverse.

On the narrow line, (MA-MO.46) gives
\(V_ay=A_SV_aF\). The line zeros of \(A_S\) are discrete
and have measure zero, while the finite Euler factors never
vanish for \(\sigma>0\). Thus \(A_S^{-1}V_ay=V_aF\) almost
everywhere, and this function is in \(L^2\). This proves
membership in the maximal inverse domain and (MA-MO.47).
Both \(\phi F\) and \((1-\phi)F\) are in \(\mathcal E\).
Their original sums and their continued multipliers are
therefore defined. Linearity of the actual sum proves
(MA-MO.48), and (MA-MO.46) proves (MA-MO.49).

Finally \(p_0\notin\mathcal H_a\), since
\(e^{2ax}|p_0(x)|^2=e^{(2a-1)x}\) is not integrable at
the left end. The continued outputs are in \(\mathcal H_a\)
and \(\phi y\in\mathcal E\). Thus (MA-MO.49) proves the
equivalences (MA-MO.50). If \(\ell_\phi(F)=0\), both
\(\phi F\) and \((1-\phi)F\) have moment zero, so their
actual sums are in \(\mathcal E\) by the proved Poisson
estimate. All remaining assertions follow.
\end{proof}

\paragraph{MA-MO-005 (unboundedness with both polar moments killed).}
For the source cutoff and every finite \(S\), there exist
\(F_M\in\mathcal E\) and actual smooth theta-image vectors
\(y_M=Z_S^{\rm geom}F_M\in\mathcal E\) such that
\[
 \begin{gathered}
 \mathcal MF_M(1/2)=\mathcal M(\phi F_M)(1/2)=0,
       \qquad \sup_M\|y_M\|_{\mathcal H_a}<\infty,\\
 \|Q_\phi^{\rm geom}y_M\|_{\mathcal H_a}\longrightarrow\infty,
 \qquad
 \|D^{\rm geom}y_M\|_{\mathcal H_a}\longrightarrow\infty.
 \end{gathered}
 \tag{MA-MO.55}
\]
Every input, inverse, and final geometric output in this assertion
is a genuine smooth function in \(\mathcal E\). In particular
neither original raw operator has a bounded extension to
\(\mathcal H_a\) agreeing on this pole-free smooth subspace.
The witnesses require neither A nor M, nor any zero on the
chosen line. The published universality input is exactly that
already specified in MA-MO-002.

\begin{proof}
\emph{A bounded moment functional and a fixed correction.}
The orientation of the source cutoff is essential here. Since
\(\phi=0\) sufficiently far right and \(1-2a>0\),
\[
 L_{\phi,a}^2=\int_{\mathbb R}|\phi(x)|^2e^{(1-2a)x}\,dx<\infty,
 \qquad |\ell_\phi(F)|\le L_{\phi,a}\|F\|_{\mathcal H_a}.
 \tag{MA-MO.56}
\]
Choose nonnegative smooth functions \(\eta_L,\eta_R\), each
of integral one, compactly supported respectively in regions
where \(\phi=1\) and \(\phi=0\). Put
\[
 F_*(x)=e^{-x/2}(\eta_L(x)-\eta_R(x)),\qquad
 y_*=Z_S^{\rm geom}F_*.
 \tag{MA-MO.57}
\]
Then \(F_*\in C_c^\infty\), \(\mathcal MF_*(1/2)=0\),
\(\ell_\phi(F_*)=1\), and \(y_*\in\mathcal E\).
The continued output \(Q_\phi^{\rm cont}y_*\) is a fixed
vector of \(\mathcal H_a\). We do not assert that the
separate raw output \(Q_\phi^{\rm geom}y_*\) is in that
Hilbert space; in fact (MA-MO.50) shows it is not.

\emph{Uniformly bounded inverses of the universality inputs.}
Use the rational inputs \(v_M\) of MA-MO.31. On their support,
write \(t=T_M+\xi/M\), \(|\xi|\le1\). By MA-MO.30,
\[
 A_0(t)=e^\xi(1+\eta_M(t)),\qquad
      |\eta_M(t)|<\varepsilon_0<1/100,
      \qquad A_0(t)=\zeta(\sigma+it).
 \tag{MA-MO.58}
\]
The fixed Euler product has bounds
\(0<m_S\le|E_S(\sigma+it)|\le M_S<\infty\).
Normalize the finite-deletion inputs by
\[
 \begin{aligned}
 v_M^S&=\frac{E_S(\sigma+i\,\cdot)\,v_M}
                    {\|E_S(\sigma+i\,\cdot)\,v_M\|_2},
 &u_M&=V_a^{-1}v_M^S,\\
 q_{0,M}&=A_S^{-1}v_M^S
   =\frac{v_M}{A_0\|E_S(\sigma+i\,\cdot)\,v_M\|_2}.
 \end{aligned}
 \tag{MA-MO.59}
\]
They belong to the exact Fourier core of MA-MO-002. Its
norm lower bound, together with (MA-MO.58), gives fixed
constants \(K,c_S,d>0\) such that
\[
 \begin{gathered}
 \|u_M\|_{\mathcal H_a}=1,\qquad
 \|q_{0,M}\|_2\le K:=\frac{e}{m_S(1-\varepsilon_0)},\\
 \|Q_\phi^{\rm cont}u_M\|_{\mathcal H_a}
             \ge R_M:=c_Se^{Md}/M\longrightarrow\infty.
 \end{gathered}
 \tag{MA-MO.60}
\]
The inverse estimate is uniform in the approximating shifts
\(T_M\). It is not an estimate for the inverse on the entire
line or for arbitrary core inputs.

\emph{Graph approximation with inverse control.}
Apply the explicit Gaussian construction MA-MO.36--.41
separately to each \(v_M^S\). It supplies
\(F_{M,\epsilon}\in\mathcal E_0\) and actual smooth images
\(y_{M,\epsilon}=Z_S^{\rm geom}F_{M,\epsilon}\) satisfying
\[
 \begin{aligned}
 V_aF_{M,\epsilon}&\longrightarrow q_{0,M}&&\text{in }L^2,\\
 y_{M,\epsilon}&\longrightarrow u_M&&\text{in }\mathcal H_a,\\
 Q_\phi^{\rm cont}y_{M,\epsilon}
      &\longrightarrow Q_\phi^{\rm cont}u_M&&\text{in }\mathcal H_a.
 \end{aligned}
 \tag{MA-MO.61}
\]
The first convergence follows from the Schwartz convergence
MA-MO.38; it is part of the explicit construction, not a
consequence of a bounded inverse. For each \(M\), choose
\(\epsilon_M>0\) so that all three errors are at most one.
Abbreviate the resulting vectors by \(F'_M,y'_M\). Then
\[
 \|F'_M\|_{\mathcal H_a}\le K+1,\qquad
 \|y'_M\|_{\mathcal H_a}\le2,\qquad
 \|Q_\phi^{\rm cont}y'_M\|_{\mathcal H_a}\ge R_M-1.
 \tag{MA-MO.62}
\]

\emph{An exact bounded correction, before any raw norm comparison.}
Set
\[
 \begin{aligned}
 c_M&=\ell_\phi(F'_M),&|c_M|&\le C_*:=L_{\phi,a}(K+1),\\
 F_M&=F'_M-c_MF_*,&y_M&=y'_M-c_My_*=Z_S^{\rm geom}F_M.
 \end{aligned}
 \tag{MA-MO.63}
\]
Both required moments of \(F_M\) vanish exactly. Moreover
\(\|y_M\|_{\mathcal H_a}\le C_Y:=2+C_*\|y_*\|_{\mathcal H_a}\).
All terms of the following estimate are continued outputs
in the actual Hilbert space, so the triangle inequality is valid:
\[
 \|Q_\phi^{\rm cont}y_M\|_{\mathcal H_a}
 \ge R_M-1-C_*\|Q_\phi^{\rm cont}y_*\|_{\mathcal H_a}
                              \longrightarrow\infty.
 \tag{MA-MO.64}
\]
Now apply MA-MO-004 to the corrected \(F_M\), whose two
moments are zero. Its continued output equals its literal
geometric output, and that output belongs to \(\mathcal E\).
Hence (MA-MO.64) is the desired raw norm divergence. Also
\[
 \|D^{\rm geom}y_M\|_{\mathcal H_a}
    \ge\|Q_\phi^{\rm geom}y_M\|_{\mathcal H_a}
                         -\|\phi\|_\infty C_Y
                                \longrightarrow\infty.
 \tag{MA-MO.65}
\]
The inverse in the original composition is exactly \(F_M\)
by (MA-MO.47). Since
\(\mathcal M((1-\phi)F_M)(1/2)=0\) too, its original
summation output is in \(\mathcal E\). Thus both literal
source expressions in (MA-MO.48) have the asserted domains
and smooth outputs. No polar subtraction remains in the
final witnesses. This proves (MA-MO.55) and rules out a
bounded extension agreeing with either raw operator there.
\end{proof}

\paragraph{Exact original-source domain and the reserved smaller space.}
Meyer (2005), \S5.2, Lemmas~5.3--5.4, and the following
projection definitions, author-manuscript pp.~39--40,
define the full smooth product
\(\mathcal S(\mathcal C)_{>}\oplus\mathcal S(\mathcal C)_{<}\)
and \(p_+^S\) on it. The primary labels are
\texttt{sec:summation}, \texttt{lem:summation}, and
\texttt{lem:Hilm\_p\_embed}. Centering turns the first
component into the space of functions whose products with
\(e^{bx}\) are Schwartz for every \(b>1/2\). Thus the
vectors \((y_M,0)\), with \(y_M\in\mathcal E\), are in
that genuine full smooth product. The actual inverse
\(\mathcal L_S^{-1}\) there returns \(F_M\), and the
first diagonal derived from the defining maps is
\[
       Z_S^{\rm geom}M_{1-\phi}J_S^{\rm geom}-M_{1-\phi}
                         =D^{\rm geom}.
 \tag{MA-MO.66}
\]
It is the first diagonal of \(P_+^S-P_-\). The printed
\texttt{eq:global\_perturbation} has the reversed label on
its left side; the defining-map sign already audited in
MA-MEY-003 is retained here.

A bounded operator on \(\mathcal H_a\oplus\mathcal H_{-a}\)
agreeing with this original matrix on the coordinate core
containing \((y_M,0)\) would have a bounded first diagonal
compression. MA-MO-005 rules out that extension. The source
also uses a different, smaller sum space: by
\texttt{lem:Poisson\_sumpm}, in \S5.3, its centered form is
\[
 \{(y_0,y_1):y_0-y_1=c_0e^{x/2}-c_1e^{-x/2}\}.
 \tag{MA-MO.67}
\]
A general \((y_M,0)\) does not lie in that smaller space.
Consequently the diagonal argument alone does not exclude
cancellation between entries on that subspace with a separately
specified norm, nor does it settle every possible common-space
transport. This distinction is part of the theorem's scope.

If the source cutoff is symmetric, reflection \(Ru(x)=u(-x)\)
is unitary \(\mathcal H_a\to\mathcal H_{-a}\), and the other
diagonal is \(RD^{\rm geom}R\). Reflecting (MA-MO.49) replaces
\(p_0\) by \(e^{x/2}\), and reflecting the two-moment-zero
witnesses proves the corresponding second-diagonal statement.
Again no assertion about the off-diagonal domains is required.

\paragraph{Source and Grothendieck placement.}
The cutoff orientation is stated explicitly in the supplied
Meyer primary text at lines~1403--1407 and~1703--1709.
The fresh source check here covers lines~3790--4088, including
\texttt{eq:Poisson\_summation} and the full-product versus
smaller-sum distinction. Lemma~5.6, pp.~41--42, assumes raw
weight \(\alpha>1\), hence centered \(a>1/2\). Its smoothed
nuclear result is unaffected: MA-MO-005 rules out the stated
raw extension to \(0<a<1/2\), not the source lemma.
The RS smoothed trace-class results and PR actual smoothed
quotient returns retain their proved domains and hypotheses.

The polar correction lies in the same two-dimensional
representation as the source singular term, whose centered
generators are
\[
                            e^{x/2},\quad e^{-x/2}.
 \tag{MA-MO.68}
\]
This is an explicit finite-rank comparison in the ambient
smooth-function category, not a splitting of a Hilbert
extension or of the finite-length spectral jet modules.
The inherited cote~1 nuclear-trace locators, scans~079--080
and~111--112 through ALPOGE\_BRIDGES BR-TR-01/02, require
actual nuclear maps before cyclicity is used. They do not
turn the non-Hilbert polar term into a bounded finite-rank
perturbation. No motivic positivity, new zero-location input,
or impossibility of all alternative transports is asserted.

\section{A line zero prevents closability of the specified smoothed Euler composition}
\label{sec:boundary-euler-nonclosability}

This is conditional on a zero on the chosen Mellin line. It
asserts neither the existence of an off-critical zero nor
failure of every possible transport. The sequential operator
domain is part of the statement.

Use (MA-MO.20), with \(0<\sigma<1\) and \(a=\sigma-1/2\).
Let \(A(t)=L(\sigma+it)E_S(\sigma+it)\) be the zeta or
primitive Dirichlet multiplier of MA-MO.21. More generally,
the argument applies to a function analytic near the real line,
not identically zero, whose derivatives of every fixed order
have polynomial growth there. Assume that \(A\) has a real
zero \(u_0\) of order \(m\ge1\). Let \(\phi\) be smooth,
one sufficiently far left and zero sufficiently far right,
as in the original source. Set
\[
 T_\phi=\mathcal F_+M_\phi\mathcal F_+^{-1},\qquad
 \mathscr C_A=C_c^\infty(\mathbb R\setminus\{A=0\}).
 \tag{MA-MO.69}
\]
The zero set is locally finite, so this core is dense in
ordinary \(L^2(\mathbb R,dt)\). For a nonzero
\(g\in C_c^\infty(\mathbb R)\), put
\[
 G(t)=\mathcal Mg(a+it),\qquad
 Q=M_AT_\phi M_{1/A},\qquad Q_g=M_GM_AT_\phi M_{1/A}
                  \quad\hbox{on }\mathscr C_A.
 \tag{MA-MO.70}
\]
Every displayed composition is sequentially defined on this
core: \(v/A\) is smooth and compactly supported,
\(T_\phi(v/A)\) is Schwartz, and multiplying by \(A\)
gives an \(L^2\) function. These are restrictions of the
literal maximal multiplier compositions.

\paragraph{MA-MO-006 (conditional nonclosability, including left smoothing).}
Under this line-zero hypothesis, neither \(Q\) nor \(Q_g\)
is closable on \(\mathscr C_A\). Neither is
\[
                 M_G(T_\phi-M_AT_\phi M_{1/A}).
 \tag{MA-MO.71}
\]
None therefore admits a closable extension agreeing on that
core. Under \(V_a\), (MA-MO.71) is the specified sequential
difference \(C_g(M_\phi-Z_{a,S}M_\phi J_{a,S})\).

\begin{proof}
Define
\[
 P(v)=\int e^{ivx}\phi'(x)\,dx,
             \qquad q(v)=\frac{iP(v)}{2\pi v}\quad(v\ne0).
 \tag{MA-MO.72}
\]
Then \(P\) is entire and Schwartz on the real line, and
\(P(0)=-1\). This is the off-diagonal kernel of MA-MO.22.
Distributional differentiation gives the full convolution
kernel \(\operatorname{pv}q+c\delta_0\) for some constant
\(c\): multiplying it by \(v\) gives \(iP(v)/(2\pi)\),
and the kernel of multiplication by \(v\) on distributions
is the span of \(\delta_0\). Its constant is immaterial.

For a smooth polynomially growing function \(B\), the
commutator on a compact smooth test has the ordinary kernel
\[
 \begin{aligned}
 [M_B,T_\phi]h(t)&=\int r_B(t,u)h(u)\,du,\\
 r_B(t,u)&=(B(t)-B(u))q(t-u)\quad(t\ne u),\\
 r_B(u,u)&=\frac{iP(0)}{2\pi}B'(u).
 \end{aligned}
 \tag{MA-MO.73}
\]
The delta term cancels and the difference removes the
principal-value singularity. For \(B=A\), both terms of
the commutator are in \(L^2\), by the domain observation
above.

The map \(u\mapsto r_A(\cdot,u)\) is continuous into
\(L^2(dt)\) for \(u\) near \(u_0\). To verify this,
restrict \(u\) to a fixed compact neighborhood. When
\(|t-u|\le1\), the difference quotient of \(A\) is
bounded by \(\sup|A'|\) on one fixed compact interval.
When \(|t-u|\ge1\), rapid decay of \(q(t-u)\) beats
the polynomial growth of \(A(t)-A(u)\), uniformly in
that neighborhood. These estimates give an integrable
square majorant. Pointwise continuity, including the
removable diagonal value, and dominated convergence give
the claimed \(L^2\) continuity.

Choose nonnegative \(h\in C_c^\infty((1,2))\) of integral
one and set
\[
 h_\epsilon(u)=\epsilon^{-1}h((u-u_0)/\epsilon),
                  \qquad v_\epsilon(u)=A(u)h_\epsilon(u).
 \tag{MA-MO.74}
\]
For sufficiently small positive \(\epsilon\), this support
avoids every real zero of \(A\), including \(u_0\).
Thus \(v_\epsilon\in\mathscr C_A\) and
\(M_{1/A}v_\epsilon=h_\epsilon\). Taylor expansion gives
\[
            \|v_\epsilon\|_2=O(\epsilon^{m-1/2})\to0.
 \tag{MA-MO.75}
\]
The exact commutator identity now yields
\[
 \begin{aligned}
 Qv_\epsilon
 &=T_\phi(Ah_\epsilon)+[M_A,T_\phi]h_\epsilon\\
 &\longrightarrow r_A(\cdot,u_0)=A(t)q(t-u_0)
                                      \quad\hbox{in }L^2.
 \end{aligned}
 \tag{MA-MO.76}
\]
The first term tends to zero by boundedness of \(T_\phi\).
The second is the approximate-identity average of the
proved \(L^2\)-continuous map. The last equality uses
\(A(u_0)=0\), with its removable value at \(u_0\).
Near that point, write \(A(t)=(t-u_0)^m a_0(t)\),
where \(a_0(u_0)\ne0\). Since \(P(0)=-1\), the limit
has a nonzero leading term of order \(m-1\). It is an
\(L^2\) function by the estimates above. The graph closure
therefore contains a nonzero vector above zero input,
which proves nonclosability of \(Q\).

The function \(G\) is bounded and Schwartz on the real
line and is the restriction of an entire function. It is
not identically zero, since Fourier uniqueness would
otherwise give \(e^{ax}g(x)=0\). Multiplying (MA-MO.76)
by this bounded function gives
\[
 Q_gv_\epsilon\longrightarrow G(t)A(t)q(t-u_0)\ne0
                                      \quad\hbox{in }L^2.
 \tag{MA-MO.77}
\]
If \(G(u_0)=0\), that zero has finite order, so the same
local expansion proves nonvanishing. Finally \(M_GT_\phi\)
is bounded and sends \(v_\epsilon\) to zero in norm.
Thus (MA-MO.71) has the negative nonzero limit and is
nonclosable too. Any extension containing the core has
the same graph sequence, proving the extension assertion.
\end{proof}

\paragraph{Domains and remaining scope.}
The inverse is applied only to tests supported away from
zeros, with every following domain verified. A factored
commutator involving \(J_{a,S}M_\phi\) would require an
additional inverse domain; no such equality of unbounded
compositions is used. The rational geometric comparison
MA-MO-004/005 has its own two-moment-zero smooth domain.
The present compact-frequency sequence is not asserted
to lie in that domain, so its multiplier nonclosability
is not imported into every original common-space model.
The theorem does not address the PR finite-jet quotient
returns or Gram-normalized maps with other domains.
The RS results on lines avoiding zeros retain their
hypotheses. On a line containing a zero, left multiplication
by nonzero compact smooth convolution does not make the
particular composition (MA-MO.70) closable.

\paragraph{Finite-length and nuclear-framework placement.}
The order \(m\) is the length of the local zero-divisor
module. Sequential domains and closability are additional
analytic structure, not consequences of that finite length.
The inherited cote~1 nuclear-trace locators, scans~079--080
and~111--112 through ALPOGE\_BRIDGES BR-TR-01/02, require
the relevant maps and their nuclear hypotheses before
traces or determinants are formed; a proposed closed
realization must also have a valid operator domain.
Neither finite-length \(K_0\) data nor arithmetic weights
repair the nonclosable composition proved here. No new
primary-source reading is claimed by this placement.

\section{The actual narrow common-sum domain is not preserved by the raw matrix}
\label{sec:smaller-sum-noninvariance}

MA-MO-005 concerns a coordinate core in the full smooth product.
The following result instead uses a genuine vector of Meyer's smaller
common-sum space. It determines its intersection with the narrow
Hilbert product and exhibits a literal raw output outside that product.
The operator is the original approximate projection; it is not replaced
by an orthogonal projection or an exact idempotent.

Fix \(0<a<1/2\), and put
\[
 \begin{aligned}
 H_{\Sigma,a}&=L^2\bigl(\mathbb R,
                         (e^{2ax}+e^{-2ax})\,dx\bigr),\\
 \mathcal S_>&=\{y:e^{bx}y\in\mathscr S(\mathbb R)
                                      \text{ for every }b>1/2\},\\
 \mathcal S_<&=\{y:e^{bx}y\in\mathscr S(\mathbb R)
                                      \text{ for every }b<-1/2\},\\
 \mathscr S_{\rm sum}
   &=\{(y_0,y_1)\in\mathcal S_>\oplus\mathcal S_<:\\
   &\hspace{2em}y_0-y_1=c_+e^{x/2}-c_-e^{-x/2}
                     \text{ for some }c_+,c_-\in\mathbb C\}.
 \end{aligned}
 \tag{MA-MO.78}
\]
These are the centered rational trivial-compact-isotype spaces of
Meyer (2005), \S\S5.2--5.3. The equation for the difference can also
be imposed in distributions on any locally square-integrable pair.

\paragraph{MA-MO-007 (exact intersection and original raw noninvariance).}
The pairs in \(\mathcal H_a\oplus\mathcal H_{-a}\) satisfying the
polar-difference equation in (MA-MO.78) are exactly
\[
 \{(y,y):y\in H_{\Sigma,a}\},\qquad
 \|(y,y)\|_{\mathcal H_a\oplus\mathcal H_{-a}}^2
                         =\|y\|_{H_{\Sigma,a}}^2.
 \tag{MA-MO.79}
\]
At the original smooth level,
\[
 \mathscr S_{\rm sum}\cap(\mathcal H_a\oplus\mathcal H_{-a})
                          =\{(y,y):y\in\mathcal E\}.
 \tag{MA-MO.80}
\]
Its Hilbert closure is (MA-MO.79).

Already for the source place set \(\{\infty\}\), the actual raw
operator \(P_+-P_-\) fails to preserve this narrow domain. There is
an even \(f_*\in C_c^\infty(\mathbb R\setminus\{0\})\) with
\(f_*(0)=\int f_*=0\), whose actual centered theta image
\(y_*\in\mathcal E\) satisfies
\[
 i_+f_*=(y_*,y_*),\qquad
 \bigl((P_+-P_-)(y_*,y_*)\bigr)_0\notin\mathcal H_a.
 \tag{MA-MO.81}
\]
The source compositions in this formula are all defined on their
original smooth domains. Consequently no Hilbert-valued operator
whose domain contains the entire original smooth diagonal can agree
with this raw matrix there. This is a failure of invariance of the
specified domain, not an assertion of nonclosability of an already
Hilbert-valued restriction.

\begin{proof}
\emph{The exact narrow intersection.}
If \(y_0\in\mathcal H_a\) and \(y_1\in\mathcal H_{-a}\), then
both \(e^{-ax}y_0\) and \(e^{-ax}y_1\) belong to
\(L^2((0,\infty))\): for the first use
\(e^{-ax}y_0=e^{-2ax}(e^{ax}y_0)\). Thus on this half-line
their difference belongs to \(L^2\). In the polar-difference
equation it equals
\(c_+e^{(1/2-a)x}-c_-e^{-(1/2+a)x}\). The second term
is square integrable, while the first is not unless \(c_+=0\).
On the negative half-line both \(e^{ax}y_j\) are square
integrable, since \(e^{ax}y_1=e^{2ax}(e^{-ax}y_1)\) there.
With \(c_+=0\), the remaining term
\(-c_-e^{(a-1/2)x}\) is not square integrable at the left
end unless \(c_-=0\). Equality in distributions gives the
same almost-everywhere equality, so this argument also covers
the stated distributional formulation. It proves (MA-MO.79).

The intersection \(\mathcal S_>\cap\mathcal S_<\) is exactly
\(\mathcal E\). To check the nontrivial inclusion, any desired
real exponential weight lies strictly between one permitted
left weight and one permitted right weight. On each half-line
its weighted derivatives are bounded by a decaying exponential
times the corresponding Schwartz estimates; a compact partition
joins the two bounds. This holds for every Schwartz seminorm.
Thus (MA-MO.80) follows from (MA-MO.79). The smooth compact
functions belong to \(\mathcal E\) and are dense in
\(H_{\Sigma,a}\), proving the closure assertion. The norm in
(MA-MO.79) is exactly the inherited norm; no equivalent
one-sided norm is substituted. It is also equivalent to the
usual weight \(e^{2a|x|}\), since the sum of the two
exponentials lies between that weight and twice that weight.

\emph{The actual archimedean source maps.}
Use self-dual additive Haar measure on \(\mathbb R\), and let
\(\mathfrak F\) have kernel \(e^{-2\pi it\xi}\). Work in the
even sector, where \(\mathfrak F^*=\mathfrak F\) and
\(\mathfrak F^2=1\). This is the source \(S=\{\infty\}\)
local model: the rational \(S\)-units are \(\{1,-1\}\).
Write the original cutoff in additive coordinates as
\[
 v(t)=\phi(\log|t|)\quad(t\ne0),\qquad v(0)=1,
                         \qquad q(t)=1-v(t).
 \tag{MA-MO.82}
\]
Because \(\phi\) equals one sufficiently far left and zero
sufficiently far right, \(v\) is an even compact smooth
function, equal to one near zero. Thus \(q=0\) near zero
and multiplication by \(q\) preserves Schwartz space.

The centered actual theta map is
\[
 \Theta f(x)=e^{x/2}\sum_{n\in\mathbb Z\setminus\{0\}}f(ne^x).
 \tag{MA-MO.83}
\]
Poisson summation gives the exact source identity
\[
 R\Theta\mathfrak F f
      =\Theta f+f(0)e^{x/2}
                     -\left(\int_{\mathbb R}f\right)e^{-x/2},
                          \qquad Ry(x)=y(-x).
 \tag{MA-MO.84}
\]
Hence \(i_+f=(\Theta f,R\Theta\mathfrak F^*f)\) is diagonal
when \(f(0)=\int f=0\). From the defining source maps,
not from an assumed idempotence, one has
\[
 p_+i_+f=qf+\mathfrak F(q\mathfrak F^*f),
                \qquad P_+=i_+p_+,\qquad
 P_-(y,y)=(y,y).
 \tag{MA-MO.85}
\]
The formula for \(p_+i_+\) is an approximate-section formula.
It is not simplified to the identity. Both summands on its
right are actual even Schwartz functions.

\emph{A genuine diagonal input with a polar output.}
Choose nonnegative \(\eta_L,\eta_R\in C_c^\infty(\mathbb R)\)
of integral one, supported respectively where \(\phi=1\)
and where \(\phi=0\). Set
\[
 \begin{aligned}
 F_*(x)&=e^{-x/2}(\eta_L(x)-\eta_R(x)),\\
 f_*(t)&=\frac{F_*(\log|t|)}{2\sqrt{|t|}}\\
       &=\frac{\eta_L(\log|t|)-\eta_R(\log|t|)}{2|t|}
                                     \quad(t\ne0),\qquad f_*(0)=0.
 \end{aligned}
 \tag{MA-MO.86}
\]
This is even, compactly supported away from zero, and smooth.
Changing variables \(t=e^x\) on the positive half-line shows
\(\int f_*=0\) and \(\int vf_*=1\). Both original theta
polar conditions therefore hold. The right tail of
\(y_*=\Theta f_*\) is rapidly decreasing by the actual sum,
and the left tail is rapidly decreasing by (MA-MO.84) and
Schwartz estimates for \(\mathfrak F f_*\). The same
estimates apply to all derivatives, so \(y_*\in\mathcal E\).
In particular \((y_*,y_*)\) lies in the exact domain
(MA-MO.80), not merely in the full smooth product.

Let \(h=p_+i_+f_*\), with its literal value from (MA-MO.85).
Its additive integral is
\[
 \int_{\mathbb R}h
   =\int qf_*+(q\mathfrak F^*f_*)(0)
   =-1.
 \tag{MA-MO.87}
\]
Here the second term is zero because \(q=0\) near zero;
the first is \(\int f_*-\int vf_*=-1\). Fourier inversion
at zero justifies the first equality for these Schwartz
functions. In particular the image has a genuine nonzero
source polar term.

For any even Schwartz \(h\), the full Poisson formula gives
\(\Theta h(x)=(\int h)e^{-x/2}-h(0)e^{x/2}\) plus a
remainder rapidly decreasing with all derivatives as
\(x\to-\infty\). Apply this to the present \(h\). The
first component of \((P_+-P_-)(y_*,y_*)\) is
\(\Theta h-y_*\), whose left asymptotic is
\(-e^{-x/2}-h(0)e^{x/2}\) plus a rapid remainder. Its
leading coefficient is nonzero, so it does not belong to
\(\mathcal H_a\) for \(0<a<1/2\). This proves (MA-MO.81)
on the actual source domain and completes the proof.
\end{proof}

\paragraph{Compact smoothing and the value-zero test condition.}
The same example also gives an exact scope statement for smoothing.
If \(g\in C_c^\infty(\mathbb R)\), compact convolution preserves
the displayed asymptotics and acts on each polar exponential by
its Mellin value. Consequently
\[
 \bigl(C_g(P_+-P_-)(y_*,y_*)\bigr)_0
   =-\mathcal Mg(1/2)e^{-x/2}
     -h(0)\mathcal Mg(-1/2)e^{x/2}
       +\text{a rapid left remainder}.
 \tag{MA-MO.88}
\]
If \(\mathcal Mg(1/2)\ne0\), this raw smoothed output is
still outside \(\mathcal H_a\). The condition \(g(0)=0\)
does not remove that term: a nonzero nonnegative test supported
in \((1,2)\) satisfies \(g(0)=0\) but has positive
\(\mathcal Mg(1/2)\). No conflict with MA-RS is implied,
since its continued operators and explicitly subtracted polar
terms have different, specified domains.

\paragraph{Further pole-free restriction and source boundary.}
The result does not decide the norm of the raw action after
restricting its inputs further so that the output also has
both polar coefficients zero. Nor does it identify the closure
of such a further domain, or infer nonclosability from an
output that is not in the proposed Hilbert target. These are
separate questions. In particular MA-MO-007 is a genuine
common-sum domain test, not a reuse of the coordinate-axis
witness from MA-MO-005.

The exact primary locators are Meyer (2005),
\texttt{sec:summation}, \texttt{lem:Hilm\_p\_embed}, the
display for \(p_+^S i_+\), the defining matrices \(P_+^S,P_-\),
\texttt{eq:Poisson\_summation}, and \texttt{lem:Poisson\_sumpm};
author-manuscript pp.~39--41, supplied \texttt{main.tex}
lines~3808--4043. The cutoff orientation is also explicit
at lines~1403--1407 and~1703--1709. The sign is the one
derived from the defining matrices in MA-MEY-003.
Meyer's Lemma~5.6 concerns centered weights \(a>1/2\),
where these source polar tails have the allowed weighted
behavior, and is unaffected by this narrow-domain result.

The two-dimensional singular representation is part of the
actual source exact-sequence geometry. Its polar generators
are absent from the inherited narrow Hilbert intersection,
as (MA-MO.79) proves. Adding them by a new norm or removing
them by a chosen finite-rank correction changes that analytic
realization and requires its own map. Neither the finite-length
spectral \(K_0\) class nor the inherited cote~1 nuclear-trace
locators, scans~079--080 and~111--112 through
ALPOGE\_BRIDGES BR-TR-01/02, supply that analytic change.
No A, M, or other zero-location premise is used here.

\section{Killing the two output poles does not give an invariant raw domain}
\label{sec:pole-free-noninvariance}

The noninvariance in MA-MO-007 has a precise further refinement.
One may require the first raw output to have both polar coefficients
zero. That restriction is a closed subspace of finite codimension in
the additive Schwartz source, but it is not preserved by the raw
action. We prove this with genuine compact additive inputs and two
successive original source compositions. No estimate for the norm of
the once-defined restricted operator is inferred from this result.

Keep the rational source place set \(\{\infty\}\), the real even
cutoff \(v(t)=\phi(\log|t|)\), \(v(0)=1\), and the actual maps
of MA-MO-007. Thus \(v\in C_c^\infty(\mathbb R)\) is one near
zero. Write \(\mathscr S_{\rm ev}\) for the even Schwartz functions.
On that space \(\mathfrak F^*=\mathfrak F\) and
\(\mathfrak F^2=1\). Define
\[
 \begin{aligned}
 L&=1-M_v-\mathfrak F M_v\mathfrak F,\\
 \ell_0(f)&=\int_{\mathbb R}v(t)f(t)\,dt,\\
 \ell_1(f)&=\int_{\mathbb R}v(t)\mathfrak Ff(t)\,dt
             =\int_{\mathbb R}\mathfrak Fv(t)f(t)\,dt,\\
 \mathscr U_0&=\{f\in\mathscr S_{\rm ev}:f(0)=\int f=0\},\\
 \mathscr U_{\rm pf}&=\{f\in\mathscr U_0:
                                      \ell_0(f)=\ell_1(f)=0\}.
 \end{aligned}
 \tag{MA-MO.89}
\]
The equality of the two integrals for \(\ell_1\) is ordinary
Fourier Fubini on Schwartz functions. All functionals in this display
are continuous in the Schwartz topology. Let
\(D=P_+-P_-\) denote the actual source matrix.

\paragraph{MA-MO-008 (exact pole-free source and failure at the second iterate).}
The subspace \(\mathscr U_{\rm pf}\) is closed of codimension four
in \(\mathscr S_{\rm ev}\), and of codimension two in
\(\mathscr U_0\). For \(f\in\mathscr U_0\), the actual raw output
\(D i_+f\) belongs to the inherited narrow Hilbert product if and
only if \(f\in\mathscr U_{\rm pf}\). In that case it is a diagonal
pair in \(\mathcal E\oplus\mathcal E\).

Nevertheless there is a real even
\(f_0\in C_c^\infty(\mathbb R\setminus\{0\})\cap
\mathscr U_{\rm pf}\) such that, with \(f_1=Lf_0\) and
\(f_2=L^2f_0\),
\[
 \begin{gathered}
 f_1\in\mathscr U_0,\qquad \ell_0(f_1)=1,
                          \qquad \int f_2=-1,\\
 D i_+f_0=i_+f_1=(\Theta f_1,\Theta f_1)
                                     \in\mathcal E\oplus\mathcal E,\\
 D^2 i_+f_0=i_+f_2,\qquad
                         \Theta f_2\notin\mathcal H_a
                                     \quad(0<a<1/2).
 \end{gathered}
 \tag{MA-MO.90}
\]
Both iterates in this formula are defined in the original smooth
source product. Thus restricting once to pole-free outputs does not
produce an invariant raw domain.

\begin{proof}
\emph{The exact raw action and its polar functionals.}
The operator \(L\) preserves even Schwartz space. As an operator on
the ordinary additive even \(L^2\) space, it is bounded and
self-adjoint, since \(v\) is real and bounded and the Fourier
transform is unitary and involutive there. It commutes with
\(\mathfrak F\). Direct Fourier evaluation gives, for every
\(f\in\mathscr S_{\rm ev}\),
\[
                 (Lf)(0)=-\ell_1(f),\qquad
                        \int Lf=-\ell_0(f).
 \tag{MA-MO.91}
\]
Indeed \(v(0)=1\) cancels the first two values in \(Lf(0)\).
For the integral, \(\int\mathfrak F(v\mathfrak Ff)
=(v\mathfrak Ff)(0)=\int f\), leaving \(-\ell_0(f)\).

Formula (MA-MO.85) is exactly
\(p_+i_+=1+L\). For \(f\in\mathscr U_0\), the source vector
\(i_+f\) is diagonal and \(P_-i_+f=i_+f\). Therefore
\[
                       D i_+f=i_+Lf
                              \qquad(f\in\mathscr U_0).
 \tag{MA-MO.92}
\]
This identity uses the actual source maps, with no idempotence
assumption. By MA-MO-007, a source common-sum vector belongs to
the narrow Hilbert product exactly when its polar difference
vanishes. The two additive polar moments of \(Lf\) are
\(-\ell_1(f)\) and \(-\ell_0(f)\). By (MA-MO.84), the
difference of the first and second components of \(i_+Lf\)
is therefore \(\ell_1(f)e^{x/2}-\ell_0(f)e^{-x/2}\).
Thus the claimed condition is necessary.
It is sufficient because both polar moments of \(Lf\) then
vanish, and full Poisson summation makes
\(\Theta Lf\in\mathcal E\). Here \(Lf\) is an additive
Schwartz function; it need not be flat at zero, and its
centered preimage is not asserted to belong to \(\mathcal E\).

\emph{The stated finite codimensions.}
The three smooth functions
\(1,v,\mathfrak Fv\) are linearly independent on
\(\mathbb R\setminus\{0\}\). To see this, an identity
\(b_0+b_1v+b_2\mathfrak Fv=0\) outside \(\operatorname{supp}v\)
first gives \(b_0=0\) by Schwartz decay. If \(b_2\ne0\),
the entire function \(\mathfrak Fv\) vanishes on an open
interval, hence vanishes identically, contradicting \(v=1\)
near zero. Thus \(b_2=0\), and then \(b_1=0\).
The transform is entire because \(v\) has compact support;
differentiating its integral proves this directly.

These three functions represent the integral, \(\ell_0\),
and \(\ell_1\). Their independence as functionals on even
compact tests away from zero follows by testing with even
bumps, since all three are even and continuous. Evaluation
at zero is independent of them on \(\mathscr S_{\rm ev}\):
a distributional relation would express a multiple of
\(\delta_0\) as a smooth function; testing away from zero
and using the independence just proved forces the smooth
part to vanish, and then the multiple of \(\delta_0\) is
zero too. Thus the four continuous functionals are independent.
Their joint map has rank four, and its kernel is
\(\mathscr U_{\rm pf}\). The first two have rank two,
so the asserted codimensions and closedness follow.

\emph{The additional functional is not a consequence of the first three.}
Put
\[
                      W=Lv=v-v^2-\mathfrak F(v\mathfrak Fv).
 \tag{MA-MO.93}
\]
This is a real even Schwartz function. We claim
\[
                  W\notin\operatorname{span}\{1,v,\mathfrak Fv\}
                  \quad\hbox{on }\mathbb R\setminus\{0\}.
 \tag{MA-MO.94}
\]
Suppose otherwise. Outside \(\operatorname{supp}v\), the
purported identity is
\(-\mathfrak F(v\mathfrak Fv)=b_0+b_2\mathfrak Fv\).
Both Fourier transforms here are Schwartz, since their
preimages are compact smooth functions. Sending \(t\) to
positive infinity gives \(b_0=0\). Both transforms are also
entire, so equality on an exterior interval implies equality
everywhere. Fourier inversion in the even sector yields
\[
                         v\mathfrak Fv=-b_2v.
 \tag{MA-MO.95}
\]
On a neighborhood of zero where \(v=1\), this says
\(\mathfrak Fv=-b_2\). The entire function \(\mathfrak Fv\)
would therefore be constant everywhere, and its Schwartz
decay would force it to be zero. Fourier injectivity then
forces \(v=0\), a contradiction. This proves (MA-MO.94).

\emph{An actual finite-bump construction.}
For clarity, the required annihilating test can be obtained
from a finite linear system. The four real even functions
\[
                   k_1=1,\quad k_2=v,\quad
                   k_3=\mathfrak Fv,\quad k_4=W
 \tag{MA-MO.96}
\]
are linearly independent on \((0,\infty)\), by the preceding
arguments and evenness. Hence there are distinct positive
points \(t_1,\ldots,t_4\) for which the evaluation matrix
\((k_i(t_j))_{i,j=1}^4\) is invertible. Otherwise all evaluation
vectors would span a proper subspace of \(\mathbb R^4\),
and a nonzero linear functional vanishing on their span
would give a forbidden relation among the four functions.

Choose an even nonnegative \(\eta\in C_c^\infty((-1,1))\)
of integral one, and define, for sufficiently small \(\epsilon>0\),
\[
 \begin{aligned}
 \beta_j^\epsilon(t)
   &=\frac{1}{2\epsilon}
       \left[\eta\left(\frac{t-t_j}{\epsilon}\right)
              +\eta\left(\frac{t+t_j}{\epsilon}\right)\right],\\
 M_{ij}(\epsilon)&=\int_{\mathbb R}k_i(t)\beta_j^\epsilon(t)\,dt.
 \end{aligned}
 \tag{MA-MO.97}
\]
These are even compact smooth bumps of mass one with disjoint
supports away from zero. The matrix tends to the invertible
evaluation matrix, so is invertible for all sufficiently
small \(\epsilon\). Fix one such value. Solve
\(M(\epsilon)c=(0,0,0,1)^{\mathsf T}\), and set
\(f_0=\sum_{j=1}^4c_j\beta_j^\epsilon\). The coefficients
and the function are real. By construction
\[
 f_0(0)=\int f_0=\ell_0(f_0)=\ell_1(f_0)=0,
                         \qquad \int Wf_0=1.
 \tag{MA-MO.98}
\]
This is an actual compact additive test, not a distribution
or an assumed vector of a completed quotient.

\emph{Two original compositions and the exit at the second one.}
By (MA-MO.91), \(f_1=Lf_0\) is even Schwartz with
\(f_1(0)=\int f_1=0\). Since the data are real, the
self-adjointness of \(L\), or direct Fourier Fubini, gives
\[
 \ell_0(f_1)=\int vLf_0=\int (Lv)f_0=1,
               \qquad \int f_2=-\ell_0(f_1)=-1.
 \tag{MA-MO.99}
\]
The first source action is (MA-MO.92), and its output is the
genuine diagonal \(i_+f_1\). Applying (MA-MO.92) again is
legitimate because \(f_1\in\mathscr U_0\). It gives the
second identity in (MA-MO.90). The first component
\(\Theta f_2\) now has left polar term \(-e^{-x/2}\),
by the full Poisson formula and \(\int f_2=-1\). Its other
polar term is a multiple of \(e^{x/2}\); the remainder is
rapid at the left. Hence \(\Theta f_2\notin\mathcal H_a\)
for \(0<a<1/2\). This proves every assertion.
\end{proof}

\paragraph{What the result does and does not repair.}
The pole-free source conditions are exact and finite in number,
but applying them once is insufficient for a raw invariant
realization: the explicitly constructed first output is a valid
diagonal smooth vector whose next original output leaves the
Hilbert target. A larger family of conditions, or a chosen
polar correction after each action, defines a different domain
or operator and needs a separate proof.

This is not a sequence with unbounded one-step Hilbert norm,
and it is not a nonclosability proof for the restriction from
\(\mathscr U_{\rm pf}\). In particular boundedness of \(L\)
on ordinary additive \(L^2\) does not establish boundedness
of its theta-conjugated action for the inherited
\(H_{\Sigma,a}\) norm. No density or closure description
for the pole-free theta-image domain in that norm is asserted.
The PR quotient returns and the RS continued operators retain
their distinct, already proved scopes.

The source identities used here are exactly the approximate-section
and Poisson maps of MA-MO.84--.85, with Meyer (2005),
\texttt{lem:Hilm\_p\_embed}, the \(p_+^S i_+\) display,
the defining matrices, and \texttt{eq:Poisson\_summation},
author-manuscript pp.~39--41. The inherited finite polar
representation and finite-length spectral \(K_0\) data do
not make an arbitrary kernel of polar functionals invariant.
The cote~1 nuclear-trace placement, scans~079--080 and
111--112 through ALPOGE\_BRIDGES BR-TR-01/02, still requires
actual maps on the claimed domains before iteration or trace
identities. No new primary-source audit is claimed beyond
the precise MO004--007 checks, and no zero-location premise
enters the present argument.

\section{The non-induced Hecke cutoff obstruction and its literal smooth domain}
\label{sec:hecke-cutoff-obstruction}

The universality input missing from the scope of MA-MO-002 is
available for an imaginary quadratic field. A second change is
favorable: a nontrivial finite-order Hecke function is entire, so
the original ideal sum has no polar discrepancy on the exponential
Schwartz space. We prove both statements with the actual domains
specified. The compact-character identification is the one already
proved in MA-CC-001 and MT-M24; no radial section for an arbitrary
ideal class is assumed.

Let \(K\) be imaginary quadratic and let \(\eta\) be a fixed
nontrivial primitive finite-order Hecke character. The character
need not be induced from \(\mathbb Q\). Extend its coefficients
by zero at conductor ideals as in (MA-RS.40). For a finite set
\(S\) of prime ideals put
\[
 \begin{aligned}
 A(t)&=E_{\eta,S}(\sigma+it)L_\eta(\sigma+it),
       &j(t)&=A(t)^{-1},\\
 E_{\eta,S}(s)&=\prod_{\mathfrak p\in S}
             (1-\eta(\mathfrak p)(N\mathfrak p)^{-s}),
       &a&=\sigma-\tfrac12,\qquad \tfrac12<\sigma<1.
 \end{aligned}
 \tag{MA-MO.100}
\]
Use \(\mathcal H_a,V_a,T_\phi\) from (MA-MO.20)--(MA-MO.22),
where \(\phi\) is any nonconstant bounded smooth function with
compactly supported derivative. The closed operators
\(Z_{a,\eta,S}\) and \(J_{a,\eta,S}\) are exactly (MA-RS.41).
Their compact-frequency core for the present composition is
\[
 \mathscr C_{\eta,\sigma}
 =C_c^\infty\bigl(\mathbb R\setminus
                  \{t:L_\eta(\sigma+it)=0\}\bigr).
 \tag{MA-MO.101}
\]
Write \(\mathcal E\) for the exponential Schwartz space and
\(\mathcal I_\eta\) for the uncentered closed ideal of MT-M18.

\paragraph{MA-MO-009 (unbounded Hecke cutoff on actual smooth ideal sums).}
The composition
\[
 Q_{\eta,\phi,S}=Z_{a,\eta,S}M_\phi J_{a,\eta,S}
 \tag{MA-MO.102}
\]
is defined on \(V_a^{-1}\mathscr C_{\eta,\sigma}\).
Neither it nor \(M_\phi-Q_{\eta,\phi,S}\) has a bounded
extension to \(\mathcal H_a\). This assertion holds for every
line \(1/2<\sigma<1\), including a line with zeros or a line
proved zero free.

More strongly, for every \(F\in\mathcal E\) the literal sum
\[
 (Z_{\eta,S}^{\rm geom}F)(x)
   =\sum_{\substack{\mathfrak a\text{ integral}\\
                         (\mathfrak a,S)=1}}
       \eta(\mathfrak a)(N\mathfrak a)^{-1/2}
                              F(x+\log N\mathfrak a)
 \tag{MA-MO.103}
\]
converges with all derivatives locally uniformly and belongs to
the actual centered smooth theta ideal \(T_{1/2}\mathcal I_\eta\).
For \(y=Z_{\eta,S}^{\rm geom}F\), the original wide-weight
inverse \(J_{\eta,S}^{\rm geom}\) and the narrow maximal inverse
both return \(F\), and
\[
 \begin{aligned}
 Z_{a,\eta,S}F&=Z_{\eta,S}^{\rm geom}F=y,\\
 J_{a,\eta,S}y&=J_{\eta,S}^{\rm geom}y=F,\\
 Q_{\eta,\phi,S}y
   &=Z_{\eta,S}^{\rm geom}(\phi F)
                    \in T_{1/2}\mathcal I_\eta,\\
 M_\phi y-Q_{\eta,\phi,S}y&\in\mathcal E.
 \end{aligned}
 \tag{MA-MO.104}
\]
There is a sequence of such genuine smooth vectors \(y_M\)
with bounded \(\mathcal H_a\) norm for which the norms of
both last two outputs in (MA-MO.104) tend to infinity.
Thus the same unboundedness holds for the literal ideal-sum
conjugation on this specified smooth domain, without a polar
subtraction or an output outside the Hilbert space.

The conclusion applies in particular to the non-induced cubic
character of \(\mathbb Q(\zeta_3)\) in MA-CH-003, and to its
conjugate. The smooth-domain argument also applies to every
nontrivial primitive Dirichlet character, using MA-MO-002 for
its universality step. No use of \(A\) or \(M\) enters these
assertions.

\begin{proof}
\emph{The precise additional published input.}
Mishou (2003), the unnumbered main Theorem on printed p.~46,
proves universality on compact sets with connected complement
in \(\sigma_K<\operatorname{Re}s<1\), where
\(\sigma_K=1-1/[K:\mathbb Q]\) for \(K\ne\mathbb Q\).
It applies to a Gr\"ossencharacter as defined in \S2,
printed pp.~47--49. Our finite-order character has trivial
archimedean type and satisfies that definition. Since
\([K:\mathbb Q]=2\), its permitted strip is exactly
\(1/2<\operatorname{Re}s<1\). The target must be continuous
and nonzero on the compact set and holomorphic inside. The
theorem gives positive lower density of approximating shifts;
only existence is used here. There is no zero-location or
extra zero-density hypothesis.

\emph{Actual composition and the norm obstruction.}
The zero set removed in (MA-MO.101) is discrete, since the
analytic function is not identically zero. Its complement
supports a dense compact smooth core in \(L^2\). On this
core \(jv\) is compact smooth. Multiplication by \(\phi\)
preserves Schwartz space, so \(T_\phi(jv)\) is Schwartz.
The fixed-strip polynomial bounds for \(L_\eta\) and all
its derivatives, already supplied by NT-Z10/MT-M18 and used
in MA-RS-005, imply \(AT_\phi(jv)\in L^2\). Thus every
sequential domain needed in (MA-MO.102) is verified.

Here are the full changes in the norm argument of MA-MO-002.
Set \(h=\sigma-3/4\) and choose \(|h|<r<1/4\).
The compact disk \(C=\{3/4+z:|z|\le r\}\) lies in the
permitted strip and has connected complement. With
\(q(v)=i\widehat{\phi'}_+(v)/(2\pi v)\), choose a fixed
\(d>0\) so small that \(h^2+d^2<r^2\) and \(q(d)\ne0\).
Such a choice follows from the identity theorem for the
nonzero entire transform of \(\phi'\).

For each positive integer \(M\), the target
\(f_M(3/4+z)=e^{-iM(z-h)}\) is entire and nonvanishing,
with minimum modulus \(e^{-Mr}\) on \(C\).
Apply Mishou's theorem with tolerance
\(\epsilon e^{-Mr}\), where \(0<\epsilon<1/100\).
It gives \(T_M\) with
\[
 L_\eta(3/4+z+iT_M)
   =e^{-iM(z-h)}(1+\delta_M(z)),\qquad
      |\delta_M(z)|<\epsilon\quad(|z|\le r).
 \tag{MA-MO.105}
\]
Choose a nonnegative real \(b\in C_c^\infty((-1,1))\)
of \(L^2\) norm one, and use the two unit vectors
\[
 v_M(u)=M^{1/2}b(M(u-T_M)),\qquad
 w_M(t)=M^{1/2}b(M(t-T_M-d)).
 \tag{MA-MO.106}
\]
For large \(M\) their disjoint windows lie inside the
disk and contain no zero, by (MA-MO.105). Off the diagonal
the kernel of the already defined operator with \(S\) empty
is exactly
\(L_\eta(\sigma+it)q(t-u)/L_\eta(\sigma+iu)\).
After \(u=T_M+\xi/M\), \(t=T_M+d+\nu/M\), its pairing
with these vectors is
\[
 \frac{e^{Md}}{M}\iint b(\xi)b(\nu)e^{\nu-\xi}
 q\!\left(d+\frac{\nu-\xi}{M}\right)
 \frac{1+\delta_M(h+i(d+\nu/M))}
      {1+\delta_M(h+i\xi/M)}\,d\xi\,d\nu.
 \tag{MA-MO.107}
\]
The relative error ratio differs from one by at most
\(2\epsilon/(1-\epsilon)\). Uniform continuity of \(q\)
on these windows makes its ratio to \(q(d)\) differ from
one by less than \(1/100\) for large \(M\). The combined
relative error is less than \(1/2\). The remaining weight
is nonnegative and has the strictly positive integral
\(I_b=(\int b(\nu)e^\nu d\nu)(\int b(\xi)e^{-\xi}d\xi)\).
Consequently the pairing has modulus at least
\(\tfrac12|q(d)|I_b e^{Md}/M\). This is a norm bound
on actual vectors, not an inference from a kernel value.

Each finite prime-ideal factor and its inverse is bounded
on this line: use \(1-(N\mathfrak p)^{-\sigma}>0\)
and \(1+(N\mathfrak p)^{-\sigma}\). The exact identity
\[
 M_A T_\phi M_j
  =M_{E_{\eta,S}}
       (M_{L_\eta}T_\phi M_{L_\eta^{-1}})
                               M_{E_{\eta,S}}^{-1}
 \tag{MA-MO.108}
\]
therefore preserves unboundedness, using normalized inputs
\(E_{\eta,S}v_M/\|E_{\eta,S}v_M\|_2\). Subtraction of
the bounded \(T_\phi\), and conjugation by \(V_a\), prove
the first assertion.

\emph{Literal ideal sums and the absence of a polar term.}
For \(F\in\mathcal E\), every derivative of the series
(MA-MO.103) converges absolutely and uniformly for \(x\)
in a compact interval. Indeed exponential decay at the
right bounds its terms by a constant times
\((N\mathfrak a)^{-c-1/2}\), for any fixed \(c>1/2\),
and that ideal series converges. On \(\mathcal H_b\),
\(b>1/2\), it is exactly the absolutely norm-convergent
operator of (MA-CC.1), with Mellin multiplier
\(E_{\eta,S}(z+1/2)L_\eta(z+1/2)\) on
\(\operatorname{Re}z=b\).

The function
\[
 H_F(z)=E_{\eta,S}(z+1/2)L_\eta(z+1/2)\mathcal MF(z)
 \tag{MA-MO.109}
\]
is entire. This is precisely where nontriviality of
\(\eta\) is used: the finite Hecke function has no pole.
On every fixed closed vertical strip, \(\mathcal MF\)
and its derivatives decay faster than every power of
\(|\operatorname{Im}z|\), whereas the other factors and
their derivatives grow at most polynomially. Thus the
same rapid bounds hold for \(H_F\) and all its derivatives.
These transform facts follow directly by integration by
parts for \(F\in\mathcal E\), and by the fixed-strip
Hecke bounds just stated.

Inverse Mellin integration along any real line defines
the same smooth function: shift a finite-height rectangle
between two lines and let its height tend to infinity.
The horizontal integrals tend to zero by these uniform
rapid bounds, and no residue is crossed. On each chosen
line its exponentially weighted inverse is Schwartz,
by Fourier inversion. Since the line is arbitrary, that
inverse belongs to \(\mathcal E\). On a line \(b>1/2\)
it is the literal sum just proved, so the two are equal
everywhere by smoothness. Hence
\[
 \mathcal M(Z_{\eta,S}^{\rm geom}F)(z)=H_F(z)
                         \quad\hbox{for every }z\in\mathbb C.
 \tag{MA-MO.110}
\]
This proves actual equality on the narrow line, not merely
a definition of its output by analytic continuation.

At every completed nontrivial zero \(\rho\), the finite
\(L_\eta\) has at least the completed multiplicity: the
gamma factors have no zeros. Thus (MA-MO.110) vanishes to
full required order at every \(\rho-1/2\). The exact
kernel equality (M18.4), with the explicit center
\(T_{1/2}F=e^{x/2}F\), gives
\(Z_{\eta,S}^{\rm geom}F\in T_{1/2}\mathcal I_\eta\).
The actual theta-image equality (M24.9), together with
(M24.3), identifies this with the original smooth
compact-character theta image. This argument uses the
proved image equality and does not invent an adelic
radial representative for each ideal class.

The original wide inverse of (MA-CC.1) returns \(F\),
since its absolutely convergent ideal M\"obius series is
inverse to \(Z\) on \(\mathcal H_b\), \(b>1/2\).
On the narrow line, (MA-MO.110) gives
\(V_ay=A V_aF\). Since \(F\in\mathcal H_a\) and
\(jA=1\) almost everywhere, \(y\) is in the maximal
inverse domain and its inverse is again \(F\).
The same argument applies to \(\phi F\in\mathcal E\).
It proves every equality and domain assertion in
(MA-MO.104), with no moment conditions.

\emph{Graph approximation by the actual smooth ideal.}
Fix \(v\in\mathscr C_{\eta,\sigma}\), put \(q_0=jv\),
and choose \(h_0=\mathcal F_+^{-1}q_0\in\mathscr S\).
For \(\epsilon>0\) set
\[
 F_\epsilon(x)=e^{-\epsilon x^2-ax}h_0(x),\qquad
 q_\epsilon=V_aF_\epsilon
             =\mathcal F_+(e^{-\epsilon x^2}h_0).
 \tag{MA-MO.111}
\]
Every \(F_\epsilon\) belongs to \(\mathcal E\), and
\(q_\epsilon\to q_0\) in Schwartz space. The latter
follows by the product rule: Gaussian derivatives are
uniformly polynomially bounded for \(0<\epsilon\le1\),
their derivatives converge on compact sets, and the
Schwartz tails of \(h_0\) control the complementary
region in every seminorm. Fourier transformation
preserves that convergence.

Let \(y_\epsilon=Z_{\eta,S}^{\rm geom}F_\epsilon\).
The equalities already proved and polynomial bounds
on \(A\) give
\[
 \begin{aligned}
 V_ay_\epsilon&=Aq_\epsilon\longrightarrow v,\\
 V_aQ_{\eta,\phi,S}y_\epsilon
       &=AT_\phi q_\epsilon
                    \longrightarrow AT_\phi q_0
                         \qquad\text{in }L^2.
 \end{aligned}
 \tag{MA-MO.112}
\]
Continuity of \(T_\phi\) on Schwartz space justifies
the second limit. Every vector and every cutoff output
in this display is the literal smooth ideal sum of
(MA-MO.104). Thus the graph closure of this restriction
contains the entire compact-frequency core graph.

For an explicit unbounded sequence, use the normalized
core vectors from (MA-MO.108), and for each \(M\) choose
\(\epsilon_M\) so that the input and output errors in
(MA-MO.112) are at most one. Then
\(\|y_{\epsilon_M}\|_{\mathcal H_a}\le2\), whereas
\(\|Q_{\eta,\phi,S}y_{\epsilon_M}\|_{\mathcal H_a}
\ge c_S e^{Md}/M-1\to\infty\).
The bounded multiplier \(M_\phi\) gives the same
divergence for the difference. This proves the claimed
sequence and the literal geometric unboundedness.

All preceding smooth-domain steps use only that the
fixed nontrivial Hecke function is entire with the
standard strip bounds. For a primitive nontrivial
Dirichlet function these hypotheses hold as well,
and MA-MO-002 supplies the first norm obstruction.
For the conjugate imaginary-quadratic character,
Mishou's theorem and each argument apply unchanged.
The reflected source component uses that conjugate
character exactly as in (MA-CC.9) and (MA-RS.44).
\end{proof}

\paragraph{Original source scope and zero-location bookkeeping.}
For the source cutoff and a structural place set \(S\) as in MA-CC-001,
(MA-CC.8) identifies the literal ideal operator with
the actual compact-character restriction of Meyer's
Euler summation. Hence the source first diagonal
\(M_\phi-Z_{\eta,S}^{\rm geom}M_\phi
J_{\eta,S}^{\rm geom}\) has the stated obstruction
on genuine smooth ideal inputs. This refutes a bounded
extension of that diagonal, and of a full product
matrix which agrees on its coordinate core.
It does not rule out cancellations on a separately
restricted common-sum domain or a smoothed operator.
It does not change the discarded-divisor gate in
MA-FT-002 or prove Hecke GRH for the cubic example.

If a zero is assumed to lie on the chosen line,
the analytic multiplier \(A\) satisfies every
hypothesis of MA-MO-006: it is analytic near the
line zero, its finite factors there are nonzero,
and it and its derivatives have polynomial growth.
That already proved theorem therefore also gives
conditional nonclosability on the compact-frequency
core, even after any nonzero compact left smoothing.
No existence of such a zero is asserted, and this
core conclusion is not transferred without proof
to the distinct geometric restriction above.

In Grothendieck's inherited trace setting, the
actual topological operator remains the required
input before a nuclear trace or Fredholm determinant
can be used: cote~1, scans~079--080 and~111--112,
through ALPOGE\_BRIDGES BR-TR-01/02. The character
is the rank-one finite-order local system already
identified in MA-CH/MA-CQ; prime-ideal coefficients
replace rational primes in its Euler realization.
No new weight, positivity, or motivic assertion is
being substituted for the analytic domain check.

Mishou, H. (2003). The universality theorem for Hecke
\(L\)-functions. \emph{Acta Arithmetica, 110}(1), 45--71.
\url{https://doi.org/10.4064/aa110-1-4}.
Exact locators: main unnumbered Theorem, p.~46;
Gr\"ossencharacter definition, \S2, pp.~47--49.
The primary statement and its applicability were freshly
read. Its complete classical proof is carried as
published background, not recorded as a new full proof
audit. The new operator deduction above receives its
own complete reading events.

\section{An actual smoothed map on the pole-free common-sum domain}
\label{sec:pole-free-smoothed-transport}

The raw iteration obstruction in MA-MO-008 does not preclude a
bounded once-defined smoothed map. We now construct such a map on
the exact pole-free source domain, with its inherited two-sided
Hilbert norm. The result is a trace-class map between two specified
Hilbert spaces. It does not assert an endomorphism of the smaller
domain or assign it an endomorphism trace.

Keep the rational source place set \(\{\infty\}\), the real source
cutoff \(\phi\), its additive version \(v\), and
\(L=1-M_v-\mathfrak F M_v\mathfrak F\) from MA-MO-008.
Thus \(\phi=1\) sufficiently far left and \(\phi=0\)
sufficiently far right. For \(0<a<1/2\) and \(s\ge0\), define
\[
 \begin{aligned}
 p_s(x)&=(1+x^2)^{s/2},\\
 H_{b,s}&=L^2(\mathbb R,e^{2bx}p_s(x)\,dx)\quad(b=\pm a),\\
 H_{\Sigma,a,s}
   &=L^2\bigl(\mathbb R,(e^{2ax}+e^{-2ax})p_s(x)\,dx\bigr),\\
 \mathscr Y_\phi&=\Theta\mathscr U_{\rm pf},\qquad
 K_{a,s,\phi}=\overline{\mathscr Y_\phi}^{\,H_{\Sigma,a,s}},\qquad
 J_{a,s}=\overline{\Theta\mathscr U_0}^{\,H_{\Sigma,a,s}}.
 \end{aligned}
 \tag{MA-MO.120}
\]
The source spaces \(\mathscr U_0,\mathscr U_{\rm pf}\) are
exactly (MA-MO.89), including both cutoff moment conditions.
The theta map is injective on \(\mathscr U_0\), by the
actual theta--M\"obius identity MA-TR-101. Consequently the
following action on \(\mathscr Y_\phi\) is well defined:
\[
       B_{\phi,g}^{\rm raw}(\Theta f)=C_g\Theta Lf,
                   \qquad f\in\mathscr U_{\rm pf}.
 \tag{MA-MO.121}
\]
It is the first component of the original
\(C_g(P_+-P_-)i_+f\); its second component is the same.

There are two analytic regimes in the theorem:
\begin{enumerate}
\item \emph{Fixed horizontal gap.} For some \(\delta>0\),
\[
 \left|\tfrac12+a-\operatorname{Re}\rho\right|\ge\delta
                  \quad\hbox{for every nontrivial zeta zero }\rho.
 \tag{MA-MO.122}
\]
Here \(g\in C_c^\infty(\mathbb R)\) is arbitrary and
\(\phi'\in C_c^\infty(\mathbb R)\). Under the stipulated
input A for rational zeta, this regime holds at every fixed
\(0<a<1/2\), with \(\delta=a\).
\item \emph{Unconditional generic width.} The number
\(\tfrac12+a\) belongs to the full-measure set
\(\mathscr G\) of MA-RS-001. Both \(g\) and \(\phi'\)
are compactly supported of some finite Gevrey order \(q>1\).
No zero-location hypothesis is used in this regime.
\end{enumerate}

\paragraph{MA-MO-010 (actual pole-free smoothing in every polynomial weight).}
In either stated regime, (MA-MO.121) extends uniquely to a
trace-class map
\[
             B_{\phi,g}:K_{a,s,\phi}\longrightarrow J_{a,s}
                    \subseteq H_{\Sigma,a,s}.
 \tag{MA-MO.123}
\]
Thus the original smoothed matrix, restricted to its genuine
pole-free smooth diagonal, has a bounded extension into the
same ambient narrow Hilbert product, for every \(s\ge0\).
Its two components are obtained from compatible one-sided
extensions, not by replacing the sum norm by a one-sided norm.

For fixed \(a,s,\phi\), the map \(g\mapsto B_{\phi,g}\)
is continuous into trace class on each fixed compact smooth
test-support space in the gap regime. In the generic regime
it is uniformly bounded for families with fixed support and
common finite-order Gevrey bounds. No invariance of
\(K_{a,s,\phi}\) under this map is asserted.

\begin{proof}
\emph{The estimates on both lines.}
Put
\[
 A_b(t)=\zeta(\tfrac12+b+it),\qquad j_b(t)=A_b(t)^{-1},
 \qquad G_b(t)=\mathcal Mg(b+it),\qquad b=\pm a.
 \tag{MA-MO.124}
\]
The nontrivial zero set is invariant under
\(\rho\mapsto1-\overline\rho\). Consequently (MA-MO.122)
gives the same fixed horizontal gap for \(\tfrac12-a\).
Likewise (MA-RS.2) for \(\tfrac12+a\) gives precisely the
same bound for \(\tfrac12-a\), since that reflection keeps
the height \(\operatorname{Im}\rho\) unchanged.

The minimum-modulus proofs of MA-RS-001 and MA-RS-004 apply
on both lines. To check their scope, their disk is centered
at \(2+it\), has radius between 12 and 13, and contains both
fixed real parts in its inner quarter. At large height it
contains no zeta pole or trivial zero. The same local zero
count, polynomial upper bound, and Harnack estimate apply.
The reflected distance bound supplies the same lower bound
for each finite Blaschke factor. At bounded height the line
has no zero, and its real part lies strictly between 0 and 1.
Thus, for every fixed derivative order, \(j_b\) has
polynomial growth in the gap regime, and is bounded by
\(\exp(C\log^2(2+|t|))\) in the generic regime. All
derivatives of \(A_b\) have polynomial growth in both cases.
This explicitly verifies the left-line estimate; it is not
deduced merely from a zero-free line without a rate.

\emph{The two one-sided smoothing kernels.}
Let \(V_b=\mathcal F_+M_{e^{bx}}\), as in MA-RS.8, and
let \(Z_b,J_b\) denote the maximal multiplier operators
with symbols \(A_b,j_b\). Initially on compact-frequency
tests define the sequential difference
\(D_b=M_\phi-Z_bM_\phi J_b\). The RS distributional
calculation gives for \(C_gD_b\) the kernel
\[
 \begin{aligned}
 k_{b,g}(t,u)
    &=\frac{i}{2\pi}G_b(t)A_b(t)
             \frac{j_b(t)-j_b(u)}{t-u}\,P(t-u),\\
 P(w)&=\int e^{iwx}\phi'(x)\,dx.
 \end{aligned}
 \tag{MA-MO.125}
\]
It is a Schwartz kernel. In the gap regime this follows
from the polynomial derivative estimates, the integral
formula for the divided difference, and the Schwartz
functions \(G_b(t)\) and \(P(t-u)\), exactly as in
MA-RS.37. In the generic regime their Gevrey Fourier decay
beats \(\exp(C\log^2(2+|t|+|u|))\), as explicitly proved
in MA-RS.22--.23. These arguments apply to every mixed
derivative. Denote the resulting operator by \(K_b(g)\).

Here is why the same kernel gives trace class in every
polynomial weight, not only the exponential weight.
Let \(r=s/2\) and
\(\Lambda^r=(1-\partial_t^2)^{r/2}\). Plancherel gives
\[
   \|u\|_{H_{b,s}}
      =\|\Lambda^r V_bu\|_2.
 \tag{MA-MO.126}
\]
Thus \(V_b\) identifies this space with the Sobolev space
\(H^r(\mathbb R)\). Conjugating a Schwartz-kernel operator
by \(\Lambda^r\) gives the kernel
\(\Lambda_t^r\Lambda_u^{-r}k_{b,g}(t,u)\). It remains
Schwartz: after Fourier transformation in the two variables,
the multipliers are \((1+\xi^2)^{r/2}\) and
\((1+\eta^2)^{-r/2}\), whose derivatives have polynomial
growth. The Hermite expansion argument of MA-RS-002 now
proves trace class on ordinary \(L^2\), hence on
\(H_{b,s}\) before conjugation. The estimates use finitely
many Schwartz seminorms and give the claimed test-family
bounds. This also proves the corresponding statement for
any other Schwartz kernel used below.

\emph{Agreement on the actual theta domain.}
The additive source need not be flat at zero. If
\(f\in\mathscr U_0\), define
\[
 F(x)=2e^{x/2}f(e^x),\qquad
 \widetilde F(x)=2e^{x/2}\mathfrak Ff(e^x),\qquad
                         y=\Theta f.
 \tag{MA-MO.127}
\]
Both weighted functions \(e^{bx}F\) and
\(e^{bx}\widetilde F\) are Schwartz for \(b=\pm a\).
Indeed evenness and the vanishing values of \(f\) and
\(\mathfrak Ff\) at zero give \(O(t^2)\) there,
with the corresponding bounds for all logarithmic
derivatives. The radial functions therefore decrease like
\(e^{5x/2}\) at the left, while their right tails decrease
faster than every exponential. This suffices for all the
weights presently used; it does not assert \(F\in\mathcal E\).

The actual theta sums are
\(Z^{\rm geom}F=y\) and
\(Z^{\rm geom}\widetilde F=Ry\), by Poisson summation.
Their moments at \(1/2\) are zero. The contour argument
of MA-MO-004 applies here on the strip between a fixed
right line \(b'>1/2\) and either \(b=\pm a\): the
Mellin transforms of these radial functions are holomorphic
for \(\Re z>-5/2\), with rapid vertical decay on each
smaller closed strip. No nonzero pole residue remains.
Consequently
\[
        V_by=A_bV_bF,\qquad
        V_bRy=A_bV_b\widetilde F,
        \qquad J_by=F,\quad J_bRy=\widetilde F.
 \tag{MA-MO.128}
\]
These are identities in the actual maximal inverse domains.
The same reasoning applies to \(\phi F\) and
\(\phi\widetilde F\). For \(f\in\mathscr U_{\rm pf}\)
their two moments vanish exactly, so
\[
   Z_b(\phi F)=\Theta(vf),\qquad
   Z_b(\phi\widetilde F)=\Theta(v\mathfrak Ff).
 \tag{MA-MO.129}
\]
Every output in this display is in \(\mathcal E\).

We also verify that the bounded kernel extension agrees
with these sequential outputs. Put \(u=V_by\) and
\(q=j_bu=V_bF\), both Schwartz. Let \(\chi_N\) be
standard smooth frequency cutoffs tending to one, and
put \(u_N=\chi_Nu\). Then
\(j_bu_N=\chi_Nq\to q\) in Schwartz space, and
\(u_N\to u\) there. Multiplication by \(\phi\) in
physical space and by \(A_b\) in frequency space preserve
Schwartz space. Hence the sequential outputs converge
in Schwartz topology, thus in every Sobolev norm.
The kernel extensions are continuous in those norms and
agree on \(u_N\), so they agree on \(u\). This argument
also works for \(Ry\). In particular it does not require
\(j_b\) to preserve every Schwartz function in the
generic regime; it uses the exact cancellation
\(j_bu_N=\chi_Nq\).

\emph{The original two-component action and the compact residual.}
For \(f\in\mathscr U_{\rm pf}\), both \(vf\) and
\(v\mathfrak Ff\) have value and integral zero.
Poisson summation therefore gives
\[
 \begin{aligned}
 \Theta Lf
 &=y-\Theta(vf)-R\Theta(v\mathfrak Ff)\\
 &=\psi y+D_by+RD_{-b}Ry,\qquad
                   \psi(x)=1-\phi(x)-\phi(-x).
 \end{aligned}
 \tag{MA-MO.130}
\]
The function \(\psi\) is compactly supported and smooth.
It is identically zero for a symmetric source cutoff.
Reflection is unitary from \(H_{b,s}\) to \(H_{-b,s}\),
because \(p_s\) is even, and
\(C_gR=RC_{g^\vee}\), where \(g^\vee(x)=g(-x)\).
It follows that the actual output in (MA-MO.121) agrees,
on the entire stated smooth core, with the trace-class
operator on \(H_{b,s}\)
\[
 T_b(g)=C_gM_\psi+K_b(g)+RK_{-b}(g^\vee)R.
 \tag{MA-MO.131}
\]
For completeness, the weighted Fourier kernel of the first
term is
\(G_b(t)(2\pi)^{-1}\int e^{i(t-u)x}\psi(x)\,dx\).
It is Schwartz, since both factors have rapid decay in
the independent variables \(t,t-u\). Thus the polynomial
weight proof above applies to it too. No nonsymmetric
cutoff term has been omitted.

\emph{Compatibility in the inherited sum norm.}
Set \(\mathbb H=H_{a,s}\oplus H_{-a,s}\) and define
the isometric embedding
\[
 \iota:H_{\Sigma,a,s}\longrightarrow\mathbb H,
             \qquad \iota y=(y,y),\qquad
 \mathbb T_g=\operatorname{diag}(T_a(g),T_{-a}(g)).
 \tag{MA-MO.132}
\]
The range of \(\iota\) is closed. The operator
\(\mathbb T_g\) is trace class, and by the actual-domain
identity just proved,
\[
       \mathbb T_g\iota y=\iota B_{\phi,g}^{\rm raw}y
                       \qquad(y\in\mathscr Y_\phi).
 \tag{MA-MO.133}
\]
Thus, if such \(y_n\) converge to \(y\in K_{a,s,\phi}\),
their output pairs converge in \(\mathbb H\) and their
limit is still in the closed diagonal range. Define
\[
 B_{\phi,g}=\left.\iota^*\mathbb T_g\iota\right|_{K_{a,s,\phi}}.
 \tag{MA-MO.134}
\]
On that domain this is precisely the inverse diagonal
identification of the output pair, by (MA-MO.133) and
closedness. It does not average two incompatible outputs.
Composition with bounded maps proves it is trace class,
with
\[
       \|B_{\phi,g}\|_1
           \le\|T_a(g)\|_{1,H_{a,s}}
                    +\|T_{-a}(g)\|_{1,H_{-a,s}}.
 \tag{MA-MO.135}
\]
Uniqueness follows from density of \(\mathscr Y_\phi\)
in its defined closure.

The range is in \(J_{a,s}\), not just in the ambient
space. On the smooth core \(Lf\in\mathscr U_0\), and
\[
 C_g\Theta Lf=\Theta h_g,\qquad
 h_g(t)=\int g(z)e^{-z/2}(Lf)(e^{-z}t)\,dz.
 \tag{MA-MO.136}
\]
Compactness of the dilation parameter gives an even
Schwartz function; its value and integral are zero.
Taking closure proves the range assertion. Restriction
of the codomain to this closed subspace preserves the
singular values. All trace-norm continuity and uniformity
assertions follow from the earlier kernel bounds and
the bounded compositions in (MA-MO.134).
\end{proof}

\paragraph{What has and has not been transported.}
The action on the smooth core is the original smoothed
common-sum matrix, with its two additional polar conditions
verified before any Hilbert norm is taken. In particular
this is an actual source identification, beyond a formal
comparison of scalar traces or of continued multipliers.
Both components use the same inherited sum norm, including
every polynomial exponent \(s\ge0\).

The closed domain \(K_{a,s,\phi}\) has been defined as a
closure, not identified with \(J_{a,s}\) or with an
arbitrarily chosen finite-codimensional Hilbert subspace.
No claim that \(B_{\phi,g}(K_{a,s,\phi})\subseteq
K_{a,s,\phi}\) is made. Trace class between two Hilbert
spaces does not itself supply an endomorphism trace,
a determinant, or an iterated original source action.
MA-MO-008's raw second-iteration obstruction remains valid.
The result also does not identify an orthogonal global
Paley--Wiener projection with a local Fourier cutoff.
No continuous lift from the Hilbert completion into the
original all-weight smooth space is inferred merely from
these fixed-weight Schwartz kernels.

The generic-width regime uses the proved RS distance
estimate and Gevrey tails, and covers a full-measure
set of widths, not every width. The arbitrary-smooth
regime uses the explicit fixed horizontal gap; invoking
A is one sufficient source of that gap. The arithmetic
polynomial cost in MA-PR-010/011 is not proved by either
regime. The inherited cote~1 locators, scans~079--080 and
111--112 via BR-TR-01/02, apply to actual nuclear maps
with their stated domains; finite-length spectral data
do not turn the map (MA-MO.123) into an endomorphism.

\section{The exact codimension-two closure and the required rank-two correction}
\label{sec:pole-free-closure}

The two cutoff moments are continuous for the inherited Hilbert
norms in the regimes of MA-MO-010. Their common kernel is therefore
a proper closed subspace of the usual theta ideal. We identify it,
including its actual Riesz Gram matrix, and prove that every nonzero
admissible smoothing still has a rank-two outgoing polar-moment
defect. Orthogonal removal of that defect produces a defined
endomorphism, but changes the original source action.

Retain all hypotheses, spaces and notation of MA-MO-010. In
particular \(0<a<1/2\), \(s\ge0\), and either the fixed-gap
smooth regime or the unconditional generic-width Gevrey regime
holds. Write \(J=J_{a,s}\), \(K=K_{a,s,\phi}\), and
\[
 w(x)=(e^{2ax}+e^{-2ax})(1+x^2)^{s/2},\qquad
 H_\Sigma=L^2(\mathbb R,w(x)\,dx).
 \tag{MA-MO.137}
\]
Inner products in this section are linear in the first variable.

\paragraph{MA-MO-011 (exact closure, noninvariance for every nonzero test,
and canonical correction).}
There are continuous linear functionals \(\lambda_0,\lambda_1\)
on \(H_\Sigma\) whose restrictions to the actual theta source are
\[
 \lambda_0(\Theta f)=\ell_0(f),\qquad
 \lambda_1(\Theta f)=\ell_1(f)
                         \qquad(f\in\mathscr U_0).
 \tag{MA-MO.138}
\]
On \(J\) they are independent. If
\(\Lambda=(\lambda_0,\lambda_1):J\to\mathbb C^2\), then
\[
                  K=J\cap\ker\lambda_0\cap\ker\lambda_1
                    =\ker\Lambda,
                    \qquad \operatorname{codim}_J K=2.
 \tag{MA-MO.139}
\]
Let \(r_0,r_1\in J\) be their Riesz vectors and
\(M=\Lambda\Lambda^*\). This positive definite two-by-two
matrix, with entries \(M_{ij}=\lambda_i(r_j)\), gives
\[
        P_K=1_J-\Lambda^*M^{-1}\Lambda.
 \tag{MA-MO.140}
\]

For every nonzero test \(g\) admitted in MA-MO-010, the actual
trace-class map \(B_{\phi,g}:K\to J\) satisfies
\[
             \Lambda B_{\phi,g}:K\longrightarrow\mathbb C^2
                       \quad\hbox{surjective}.
 \tag{MA-MO.141}
\]
In particular its range is not contained in \(K\). The corrected
map and its defect are
\[
 \begin{aligned}
 B_{\phi,g}^{\rm corr}&=P_KB_{\phi,g}:K\longrightarrow K,\\
 C_{\phi,g}&=\Lambda^*M^{-1}\Lambda B_{\phi,g}:K\longrightarrow J,
              &\operatorname{rank}C_{\phi,g}&=2.
 \end{aligned}
 \tag{MA-MO.142}
\]
The first map is trace class. The correction has the minimum
possible rank among linear corrections whose subtraction makes
the range lie in \(K\). Its construction uses the actual inherited
Hilbert Gram matrix, not the additive \(L^2\) norm.

\begin{proof}
\emph{Continuous moment functionals in the actual norm.}
Put \(c=1/2-a>0\) and
\[
 p_a(x)=e^{cx}\phi(x),\qquad
 \widehat p_a(t)=\mathcal F_+p_a(t),\qquad
 \beta_a(t)=j_a(t)\widehat p_a(-t),
             \qquad j_a(t)=\zeta(\tfrac12+a+it)^{-1}.
 \tag{MA-MO.143}
\]
The function \(p_a\) is Schwartz: it decays exponentially
at the left and vanishes sufficiently far right, with the
same properties for its derivatives. In the fixed-gap regime,
the polynomial reciprocal estimate therefore gives
\(\beta_a\in L^2\).

In the generic regime a sharper estimate is needed. Integration
by parts, with vanishing boundary terms at both ends, gives
\[
 \widehat p_a(t)=
       -\frac{\mathcal F_+(e^{c\,\cdot}\phi')(t)}{c+it}.
 \tag{MA-MO.144}
\]
The numerator is the transform of a compact Gevrey function
of the same finite order: the product rule bounds its
\(n\)-th derivative by \(C A^n(n!)^q\), since the exponential
factor is taken on a fixed compact interval. Repeated
integration by parts and optimization in \(n\), as in
MA-RS.19, yield a bound \(C\exp(-d|t|^{1/q})\).
The denominator has modulus at least \(c\). This decay beats
the proved bound \(\exp(C'\log^2(2+|t|))\) for \(j_a\).
Thus \(\beta_a\in L^2\) in this regime too. Formula
(MA-MO.144) verifies the noncompact weighted-step estimate
without treating \(p_a\) itself as compactly supported.

For \(y\in H_{a,0}\), define
\[
        \lambda_0(y)=\int_{\mathbb R}(V_ay)(t)\beta_a(t)\,dt,
                 \qquad \lambda_1(y)=\lambda_0(Ry)
                       \quad(y\in H_{-a,0}).
 \tag{MA-MO.145}
\]
Each integral is absolutely controlled by Cauchy--Schwarz.
Since \(s\ge0\), the inherited sum norm dominates both
component norms, so these functionals are continuous on
\(H_\Sigma\), with
\(\|\lambda_i\|\le\|\beta_a\|_2\).

For an actual source vector \(y=\Theta f\), write
\(F(x)=2e^{x/2}f(e^x)\) as in MA-MO.127.
The checked domain identity gives \(j_aV_ay=V_aF\).
The bilinear Plancherel identity, with the unitary Fourier
normalization used throughout, therefore gives
\[
 \begin{aligned}
 \lambda_0(y)
 &=\int (V_aF)(t)\widehat p_a(-t)\,dt\\
 &=\int e^{ax}F(x)p_a(x)\,dx
   =\int e^{x/2}\phi(x)F(x)\,dx
   =\int_{\mathbb R}v(t)f(t)\,dt.
 \end{aligned}
 \tag{MA-MO.146}
\]
Poisson gives \(Ry=\Theta\mathfrak Ff\), proving the
second identity in (MA-MO.138). In particular these are
functionals for the actual theta norm, not an assumed
bounded inverse into an additive source norm.

\emph{The exact kernel closure and its actual Gram matrix.}
Independence of \(\ell_0,\ell_1\) on \(\mathscr U_0\)
was proved in MA-MO-008. Thus \(\Lambda\) has rank two
already on \(\Theta\mathscr U_0\), and hence on \(J\).
Choose \(y_0,y_1\) in that smooth theta image with
\(\Lambda y_i=e_i\), the two standard basis vectors.
Continuity immediately gives \(K\subseteq\ker\Lambda\).
Conversely, if \(y\in J\) and \(\Lambda y=0\), choose
\(z_n\in\Theta\mathscr U_0\) tending to \(y\) in the
inherited norm and set
\[
             z_n'=z_n-\lambda_0(z_n)y_0-\lambda_1(z_n)y_1.
 \tag{MA-MO.147}
\]
Then \(z_n'\in\Theta\mathscr U_{\rm pf}\) and
\(z_n'\to y\). This proves (MA-MO.139) and its exact
codimension, without assuming that finite codimension
survives completion automatically.

Here is an explicit construction of the Riesz vectors in
the declared metric. Set
\[
 h_a=V_a^{-1}\overline{\beta_a},\qquad
 d_a(x)=\frac{e^{2ax}}{w(x)}h_a(x),\qquad
 r_0=P_Jd_a,\qquad r_1=Rr_0.
 \tag{MA-MO.148}
\]
The vector \(h_a\) belongs to \(H_{a,0}\), and
\[
 \|d_a\|_{H_\Sigma}^2
  =\int\frac{e^{4ax}}{w(x)}|h_a(x)|^2\,dx
  \le\int e^{2ax}|h_a(x)|^2\,dx.
\]
Direct substitution shows that \(d_a\) represents
\(\lambda_0\) on \(H_\Sigma\). Orthogonal projection
therefore gives its representative on \(J\). Reflection
is unitary in this even weight, preserves \(J\) by the
theta--Fourier identity, and commutes with \(P_J\).
Thus \(r_1=Rr_0\) represents \(\lambda_1\) on \(J\).
Their independence makes \(M=\Lambda\Lambda^*\)
positive definite, and the standard orthogonal decomposition
\(J=\ker\Lambda\oplus\operatorname{ran}\Lambda^*\)
proves (MA-MO.140). For example,
\(\Lambda\Lambda^*M^{-1}\Lambda y=\Lambda y\), so the
displayed complementary projection removes exactly both
moments, and its range is orthogonal to their common kernel.

\emph{The exact outgoing moments of the smoothed source action.}
For the fixed nonzero test \(g\), define even compact
smooth additive functions
\[
 W_g(t)=\int g(z)e^{z/2}v(e^zt)\,dz,
             \qquad U=W_{g^\vee},\qquad W=W_g,
             \qquad g^\vee(z)=g(-z).
 \tag{MA-MO.149}
\]
Compactness and smoothness follow from the bounded dilation
parameter and compact smooth \(v\); near zero the first
function is the constant \(\mathcal Mg(1/2)\).
Both \(W\) and \(U\) are nonzero. Indeed, if \(W_g=0\),
differentiate its restriction to positive arguments in
logarithmic coordinates and take the ordinary Fourier
transform. It gives
\[
      \mathcal Mg(1/2-it)P(t)=0\quad(t\in\mathbb R),
                  \qquad P(t)=\int e^{itx}\phi'(x)\,dx.
 \tag{MA-MO.150}
\]
Here \(P(0)=-1\). The entire function
\(\mathcal Mg(1/2-it)\) must vanish on an interval, hence
identically; Fourier uniqueness implies \(g=0\), a
contradiction. The same argument applies to \(g^\vee\).

Let \(\mathcal D_g h(t)=\int g(z)e^{-z/2}h(e^{-z}t)\,dz\),
the actual additive dilation operator in (MA-MO.136).
Changing variables in its bilinear pairing gives
\(\int v\mathcal D_g h=\int W_g h\).
Also \(\mathfrak F\mathcal D_g
=\mathcal D_{g^\vee}\mathfrak F\), directly by Fourier
scaling. The operator \(L\) is its own transpose for
this bilinear pairing and commutes with \(\mathfrak F\).
For every \(f\in\mathscr U_{\rm pf}\), therefore,
\[
 \begin{aligned}
 \lambda_0(B_{\phi,g}\Theta f)&=\int f(t)(LW)(t)\,dt,\\
 \lambda_1(B_{\phi,g}\Theta f)&=\int f(t)(\mathfrak F LU)(t)\,dt.
 \end{aligned}
 \tag{MA-MO.151}
\]
No adjoint conjugation has been silently introduced: all
these source moment identities are complex bilinear, and
Fourier Fubini supplies the stated transpose identities.

\emph{Independence of the two outgoing moments modulo all input constraints.}
We show that the five even smooth functions
\[
                 1,\quad v,\quad\mathfrak Fv,
                       \quad LW,\quad\mathfrak F LU
 \tag{MA-MO.152}
\]
are linearly independent on \(\mathbb R\setminus\{0\}\).
The first three are independent by MA-MO-008. Suppose
\[
        A LW+B\mathfrak F LU=c_0+c_1v+c_2\mathfrak Fv.
\]
All terms except \(c_0\) are Schwartz, so \(c_0=0\).
Write \(q=1-v\). Outside a sufficiently large compact
interval the remaining identity is
\[
           -A\mathfrak F(v\mathfrak FW)+B\mathfrak F(qU)
                       =c_2\mathfrak Fv.
\]
The preimages of all three transforms are compact smooth
functions. Their transforms are entire, so equality on
that exterior interval holds everywhere. Fourier inversion
in the even sector gives
\[
                  -A v\mathfrak FW+BqU=c_2v.
 \tag{MA-MO.153}
\]
Near zero \(v=1\) and \(q=0\). If \(A\ne0\), this
forces the entire function \(\mathfrak FW\) to be
constant on an interval, hence constant everywhere.
Its Schwartz decay then forces \(W=0\), contrary to
(MA-MO.150). Thus \(A=0\), and (MA-MO.153) near zero
also gives \(c_2=0\).

If \(B\ne0\), (MA-MO.153) now gives \(qU=0\).
Since
\(\mathfrak F LU=\mathfrak F(qU)-v\mathfrak FU\),
the original identity becomes
\(-Bv\mathfrak FU=c_1v\). On the same neighborhood
of zero, \(\mathfrak FU\) is constant. The preceding
entire-function argument forces \(U=0\), again a
contradiction. Hence \(B=0\), and the remaining coefficients
vanish by the independence of the first three functions.
This proves (MA-MO.152).

The independence is realized by actual source tests. Choose
five distinct positive points with invertible evaluation
matrix for (MA-MO.152); they exist because the evaluation
vectors span \(\mathbb C^5\). Replace each point and its
negative by an even compact smooth bump of mass one,
supported away from zero. For sufficiently narrow bumps
their moment matrix remains invertible, as in MA-MO.97.
Solving that matrix with right side
\((0,0,0,z_0,z_1)^{\mathsf T}\) gives an even compact
\(f\), supported away from zero, with
\(f\in\mathscr U_{\rm pf}\) and outgoing moments
\((z_0,z_1)\) in (MA-MO.151). Thus surjectivity in
(MA-MO.141) already holds on genuine smooth source
vectors, and passes to the stated Hilbert domain.

\emph{The correction and its necessary rank.}
Formula (MA-MO.140) gives (MA-MO.142), and composition
of the trace-class map \(B_{\phi,g}\) with this bounded
projection is trace class. Its removed part has range
\(\operatorname{ran}\Lambda^*\), which is two dimensional.
Surjectivity of \(\Lambda B_{\phi,g}\) makes that range
exactly two dimensional. If any linear correction \(C\)
has \((B_{\phi,g}-C)(K)\subseteq K\), then
\(\Lambda C=\Lambda B_{\phi,g}\), so \(\operatorname{rank}C\ge2\).
The orthogonal correction attains this lower bound.
\end{proof}

\paragraph{Trace and source scope.}
The corrected endomorphism has a legitimate ordinary nuclear
trace and Fredholm determinant. If desired, bounded cyclicity
between the two Hilbert spaces gives the exact identities
\[
 \begin{aligned}
 \operatorname{Tr}_K(P_KB_{\phi,g})
      &=\operatorname{Tr}_J(B_{\phi,g}P_K),\\
 \det_K(1-zP_KB_{\phi,g})
      &=\det_J(1-zB_{\phi,g}P_K).
 \end{aligned}
 \tag{MA-MO.154}
\]
These use the actual trace-class maps just proved, as in
the inherited nuclear-pairing/Sylvester argument of MA-GP-008.
They are not an evaluation by an arithmetic explicit formula.
The operator \(P_KB_{\phi,g}\) is a changed source action,
depending on the cutoff and the inherited weight through
the Gram matrix. It does not agree with the uncorrected
source action on all of \(K\); their defect has rank two
for every nonzero admissible \(g\), including nonzero
admissible tests satisfying \(g(0)=0\).

The proper closure (MA-MO.139) resolves the closure question
left open in MA-MO-010 in its two stated regimes. Nothing
here asserts the same functional bounds on every exceptional
width with arbitrary smooth cutoff. No new arithmetic
zero-location assertion or polynomial interpolation cost
bound is obtained. The finite polar representation and
the inherited cote~1 nuclear-trace locators, scans~079--080
and~111--112 via BR-TR-01/02, distinguish the finite-rank
correction from the analytic construction of the maps;
finite-length \(K_0\) data alone neither choose its Gram
metric nor identify its trace with the original formula.

\section{An actual nuclear endomorphism on the nontrivial character theta ideal}
\label{sec:character-common-sum-transport}

In a nontrivial finite-order compact character sector the source polar
representation is absent. This removes the two output conditions of
MA-MO-010, but it does not by itself prove the required inverse domains.
We first prove an unconditional division statement for every actual
smooth theta-ideal vector. We then identify the complete original
common-sum action on that core and extend its smoothing to its actual
Hilbert closure. This is an endomorphism of the completed theta ideal;
it is not a construction of the full smooth zero quotient or its trace.

Let \(K\) be a fixed number field and let \(\eta\) be a nontrivial
primitive finite-order Hecke character. This includes a nontrivial
primitive Dirichlet character and the non-induced cubic character
\(\psi_{269,3}\) already constructed over \(\mathbb Q(\zeta_3)\).
Choose a structural finite place set \(S\) as in MA-CC-001,
containing the conductor and the required additive-character data.
Its finite part is understood in Euler products. Use the same set
for \(\overline\eta\). Put
\[
 \begin{aligned}
 A_{\omega,S}(s)&=E_{\omega,S}(s)L_\omega(s),\\
 E_{\omega,S}(s)&=\prod_{\mathfrak p\in S}
          (1-\omega(\mathfrak p)(N\mathfrak p)^{-s}),\\
 \mathscr I^c_\omega&=T_{1/2}\mathcal I_\omega\\
       &=\{y\in\mathcal E:(\mathcal My)^{(j)}(\rho-1/2)=0\\
       &\hspace{5em}\text{for }\rho\in Z_\omega,
                                      \ 0\le j<m_\rho\},\\
 H_{\Sigma,a,s}&=L^2\bigl(\mathbb R,
         (e^{2ax}+e^{-2ax})(1+x^2)^{s/2}\,dx\bigr),\\
 J_{\omega,a,s}&=\overline{\mathscr I^c_\omega}^{\,H_{\Sigma,a,s}},\\
 &0<a<1/2,\quad s\ge0,\quad\omega\in\{\eta,\overline\eta\}.
 \end{aligned}
 \tag{MA-MO.170}
\]
Here \(Z_\omega\) is the full nontrivial completed zero divisor,
with multiplicities. The equality defining \(\mathscr I^c_\omega\) is the exact
MT-M18/MT-M24 image theorem, not a newly defined Euler ideal.
All constants concern the fixed field, character, place set and width.
Use the probability compact Haar measure and norm splitting of MT-M24
and MA-CC.8. Thus \(C_g\) is the actual centered scalar convolution
associated with the compact-character test in MA-CC.5.

Let \(\phi\) be the original real source cutoff, equal to one
sufficiently far left and zero sufficiently far right. It need not
be symmetric. We use either of the following two regimes:
\[
 \begin{array}{ll}
 \text{gap/smooth:}&
  |1/2+a-\operatorname{Re}\rho|\ge\delta>0
     \quad(\rho\in Z_\eta\cup Z_{\overline\eta}),\\
 &g,\phi'\in C_c^\infty;\\[2mm]
 \text{generic/Gevrey:}&
  1/2+a\in\mathscr G_{\eta,\overline\eta},\\
 &g,\phi'\text{ are compact Gevrey of a common finite order }q>1.
 \end{array}
 \tag{MA-MO.171}
\]
Here \(\mathscr G_{\eta,\overline\eta}\) is the simultaneous
full-measure set of MA-RS-005. The gap is an explicit hypothesis.
The generic regime has no
zero-location premise.

\paragraph{MA-MO-012 (actual character common-sum smoothing on its ideal closure).}
For every \(y\in\mathscr I^c_\eta\), the complete original
source action of \(C_g(P_+^S-P_-)\) on \((y,y)\) is a diagonal
pair \((B^{\rm raw}_{\eta,g}y,B^{\rm raw}_{\eta,g}y)\), with
\(B^{\rm raw}_{\eta,g}y\in\mathscr I^c_\eta\).
In either regime (MA-MO.171), this action extends uniquely to
an actual trace-class endomorphism
\[
                   B_{\eta,g}:J_{\eta,a,s}\longrightarrow J_{\eta,a,s}
                                      \qquad(s\ge0).
 \tag{MA-MO.172}
\]
The map is continuous in trace norm on every fixed compact smooth
test-support space in the gap regime, and uniformly bounded on
families with fixed support and common finite-order Gevrey bounds
in the generic regime. No additional moment conditions on \(g\)
or on the source vector are imposed.

The proof includes the following domain statement without either
analytic regime: for every \(y\in\mathscr I^c_\omega\), its actual
original wide inverse \(F=J^{\rm geom}_{\omega,S}y\) satisfies
\[
 \begin{gathered}
 e^{bx}F(x)\in\mathscr S(\mathbb R)\quad\hbox{for every }b>-1/2,
       \qquad (1-\phi)F\in\mathcal E,\\
 J_{b,\omega,S}y=F\quad(b=\pm a),\\
 Z_{b,\omega,S}(\phi F)=Z^{\rm geom}_{\omega,S}(\phi F)
                                         \in\mathscr I^c_\omega.
 \end{gathered}
 \tag{MA-MO.173}
\]
The narrow symbols in this formula denote the maximal multipliers
on the indicated lines. The vector \(F\) need not belong to
\(\mathcal E\); the weighted statement is the one proved and used.

\begin{proof}
\emph{Unconditional division on the positive finite-\(L\) half-plane.}
For \(y\in\mathscr I^c_\omega\), set
\[
             H_y(z)=\mathcal My(z-1/2),\qquad
                         Q_y(z)=H_y(z)/L_\omega(z).
 \tag{MA-MO.174}
\]
The numerator is entire and rapidly decreasing, with all
derivatives, on every fixed closed vertical strip. All nontrivial
completed zero jets vanish to their full order. Since the gamma
factors have no zeros, the possible poles of \(Q_y\) that these
jets do not remove occur only at the trivial finite-\(L\) zeros.
For a finite-order character those occur among the nonpositive
real gamma-pole locations. There are only finitely many in any
fixed vertical strip. This also covers a possible overlap at
an endpoint: its remaining pole order is at most the gamma-pole
order, after the full completed-zero multiplicity is canceled.

Fix a central closed strip \(\epsilon\le\Re z\le R\), with
\(\epsilon>0\). Choose a wider strip \(-A\le\Re z\le B\),
where \(A>2\), \(B>\max(2,R)\), and the left edge avoids the
finitely located trivial-zero real parts. Let \(P(z)\) be a
finite polynomial clearing all the possible trivial-zero poles
of \(Q_y\) in that wider strip; its zeros can all be taken in
\(\Re z\le0\). The orders can be fixed from the gamma factors,
independently of \(y\). Then
\[
                         P(z)Q_y(z)
       \quad\hbox{is holomorphic on the wider closed strip.}
 \tag{MA-MO.175}
\]

The selected heights of MT-M18, (M18.7)--(M18.9), give polynomial
lower bounds for the finite \(L_\omega\), uniformly on each
horizontal segment of this wider strip. Indeed the exponential
Stirling factor in the lower bound for the completion cancels
against the same factor in \(\Xi_\omega/L_\omega\).
On the two chosen vertical edges the finite function likewise
has a polynomial lower bound: at the right edge use the Euler
product; at the left use the completed functional equation,
Stirling, and the right-half-plane Euler lower bound. The left
edge was chosen to avoid the finite low-height gamma exceptions.
All these statements are unconditional for the fixed character.

For each sufficiently large positive or negative height block,
the rapid bounds for \(H_y\), multiplied by \(P\) and these
polynomial reciprocal bounds, therefore give, for every \(N\),
\[
       |P(z)Q_y(z)|\le C_{N,y}(2+|n|)^{-N}
 \quad\hbox{on the boundary of }
                  [-A,B]+i[h_n,h_{n+1}].
 \tag{MA-MO.176}
\]
The maximum principle gives the same estimate inside each
rectangle. The heights have gaps between \(1/2\) and \(3/2\),
so these estimates cover both infinite ends of the central strip.
On that strip \(P\) is bounded away from zero at bounded height
and has polynomial size at large height. The finitely many
remaining portions of the central strip are compact and \(Q_y\)
is holomorphic there. Thus \(Q_y\) is rapidly decreasing on that strip.
Cauchy estimates on a slightly wider positive strip give the
same conclusion for every derivative.

Only now divide by the finite deleted Euler product. On
\(\Re z\ge\epsilon>0\), every nontrivial factor obeys
\(|1-\omega(\mathfrak p)(N\mathfrak p)^{-z}|
\ge1-(N\mathfrak p)^{-\epsilon}>0\), and its inverse and
all derivatives are uniformly bounded on each such strip.
We have proved
\[
 \widetilde Q_y(z)=\frac{H_y(z)}{E_{\omega,S}(z)L_\omega(z)},
 \qquad \widetilde Q_y\in\operatorname{Hol}(\{\Re z>0\}).
 \tag{MA-MO.177}
\]
This function is rapid with all derivatives on every positive closed
strip. No lower bound for \(E_{\omega,S}\) was used on a strip crossing
\(\Re z=0\), where its zeros need not be finite in number.

\emph{The actual inverse and the cutoff domain.}
Inverse Mellin integration of
\(\widetilde Q_y(z+1/2)\) along any \(\Re z=b>-1/2\)
defines the same smooth function \(F\). The contours may be
shifted between any two such lines by (MA-MO.177), with no
poles and vanishing horizontal integrals. Fourier inversion
gives
\[
 \begin{gathered}
 \mathcal MF(z)=\frac{\mathcal My(z)}
                   {E_{\omega,S}(z+1/2)L_\omega(z+1/2)}
                  \quad(\Re z>-1/2),\\
             e^{bx}F\in\mathscr S\quad(b>-1/2).
 \end{gathered}
 \tag{MA-MO.178}
\]
The Mellin identity is understood through the holomorphic
quotient, including its canceled zero values. On every wider
line \(b>1/2\), it is exactly the original absolutely
convergent Euler inverse of MA-CC.1. Hence this \(F\) is
\(J^{\rm geom}_{\omega,S}y\), the source inverse, rather
than a chosen lift of the Hilbert vector. The source
identification is valid for the structural set \(S\) by
MA-CC.8 and Meyer's Lemmas~5.3--5.4.
More explicitly, MT-M24 gives an actual positive-image class for
\(y\). Its compact character is invariant outside this structural
set \(S\), so Meyer's local identification supplies a vector of
the corresponding \(\mathcal T(\mathcal C_S)\) sector whose
summation is \(y\). The wide Euler inverse is unique, and
therefore that vector is exactly \(F\). This verifies its local
source membership before any source Fourier map is applied.

The weights with arbitrarily large positive \(b\) show
rapid decay of \(F\) and all its derivatives at the right
end. The cutoff \(q=1-\phi\) vanishes at the left end.
Consequently
\[
                           qF\in\mathcal E.
 \tag{MA-MO.179}
\]
For any \(U\in\mathcal E\), the literal ideal sum
\(Z^{\rm geom}_{\omega,S}U\) belongs to
\(\mathscr I^c_\omega\). This uses only the nontrivial
finite character, not universality: on a wide line its
Mellin transform is \(A_{\omega,S}(z+1/2)\mathcal MU(z)\).
It is entire with rapid decay on every fixed strip, because
\(L_\omega\) is entire and has the standard polynomial
strip bounds. Shifting to every line gives an element of
\(\mathcal E\); its full completed-zero jets vanish, so
MT-M24 identifies it with the actual theta image. This is
the residue-free part of MA-MO-009, which applies here
without the degree-two universality hypothesis.

Apply it to \(U=qF\). Linearity of the actual wide Euler
sum gives the literal identity
\[
       Q^{\rm geom}_\omega y
          :=Z^{\rm geom}_{\omega,S}(\phi F)
           =y-Z^{\rm geom}_{\omega,S}(qF)
                                   \in\mathscr I^c_\omega.
 \tag{MA-MO.180}
\]
All these original sums exist on the wide source domain.
For \(b=\pm a\), write
\(A_{\omega,b}(t)=A_{\omega,S}(1/2+b+it)\).
Formula (MA-MO.178) gives
\(V_by=A_{\omega,b}V_bF\), with \(V_bF\) Schwartz.
Multiplication by \(A_{\omega,b}^{-1}\) therefore returns
\(V_bF\) almost everywhere, even if the line has zeros.
This verifies the maximal inverse domain. Both \(F\) and
\(qF\) are in the maximal \(Z_{b,\omega,S}\) domain, so
their difference \(\phi F\) is too. Its output is precisely
(MA-MO.180). This proves all of (MA-MO.173) without a gap
or a generic-width assumption.

\emph{The actual complete common-sum formula and the dual sector.}
A finite-order character is trivial on the connected norm
splitting. Since \(\eta\ne1\), its compact restriction is
nontrivial. Both source polar characters are trivial on that
compact group, so their averages in this sector vanish.
The source Poisson identity therefore makes the entire
positive image diagonal; by MT-M24 it is exactly
\(\mathscr I^c_\eta\). These are statements about the
specified compact average and image, not a radial model
assumed for a single ideal class.

For \(y\in\mathscr I^c_\eta\), let \(F\) be its actual
local source inverse and put
\(\widetilde F=J^{\rm geom}_{\overline\eta,S}Ry\).
The actual source embeddings and approximate section give
\[
 \begin{aligned}
 i_+F&=(y,y),\qquad
            \widetilde F=\mathfrak F^*F,\\
 \bigl(P_+^S(y,y)\bigr)_0
       &=Z^{\rm geom}_{\eta,S}(qF)
          +Z^{\rm geom}_{\eta,S}\mathfrak F(q\widetilde F),\\
 Z^{\rm geom}_{\eta,S}\mathfrak F(q\widetilde F)
       &=RZ^{\rm geom}_{\overline\eta,S}(q\widetilde F).
 \end{aligned}
 \tag{MA-MO.181}
\]
Here \(\mathfrak F,\mathfrak F^*\) are Meyer's actual
additive Fourier maps on the local coinvariant spaces,
including the finite-place phases. They are not replaced
by a bare real or complex archimedean Fourier transform.
The last identity is the source Poisson formula, whose
singular term is zero in this nontrivial compact sector.
The cutoff vectors \(qF,q\widetilde F\) belong to
\(\mathcal E\) by the division statement and are valid
local source vectors. These assertions also respect the
odd Dirichlet Fourier-square phases already checked in
MA-CQ; no separate assumption of an archimedean involution
is used.

Inversion in MA-CC.9 exchanges \(\eta\) and
\(\overline\eta\) and becomes the linear reflection \(R\).
In particular
\[
                 R\mathscr I^c_\eta=\mathscr I^c_{\overline\eta}.
 \tag{MA-MO.182}
\]
The primitive finite local factors and the conductor remain
those of \(\eta\); reflection conjugates their coefficients.
Using \(Z^{\rm geom}_{\eta,S}F=y\) and
\(Z^{\rm geom}_{\overline\eta,S}\widetilde F=Ry\) in
(MA-MO.181), and subtracting \(P_-(y,y)=(y,y)\), gives
\[
 \begin{aligned}
 (P_+^S-P_-)(y,y)&=(\Delta_\eta y,\Delta_\eta y),\\
 \Delta_\eta y&=y-Q^{\rm geom}_\eta y
                         -R Q^{\rm geom}_{\overline\eta}Ry
                                      \in\mathscr I^c_\eta.
 \end{aligned}
 \tag{MA-MO.183}
\]
The second component agrees by the zero-singular-term
common-sum relation. Every term here is an actual source
function; the preceding argument checked all its inverse
and cutoff domains. Convolution preserves the actual
theta image by equivariance, or directly by multiplication
of all its Mellin jets. Thus
\(B^{\rm raw}_{\eta,g}y=C_g\Delta_\eta y\) belongs to
\(\mathscr I^c_\eta\).

\emph{Trace-class kernels on both lines and both character sectors.}
The completed functional equation sends a zero \(\rho\) of
\(\eta\) to \(1-\rho\) of \(\overline\eta\), with the
same multiplicity. Hence the union of the two divisors is
stable under this reflection. The gap or generic-distance
bound at \(1/2+a\) supplies the same bound at \(1/2-a\).
The disk minimum-modulus proofs of MA-RS-005 apply to
both lines: their positive real parts lie inside the same
fixed inner disk at large height, and there are no trivial
zeros there. The reciprocal derivatives are polynomial
in the gap regime and bounded by
\(\exp(C_j\log^2(2+|t|))\) in the generic regime.
Finite Euler inverses have bounded derivatives on these
positive closed strips. No estimate for a reciprocal on
the negative half-plane is invoked here.

For \(\omega=\eta,\overline\eta\) and \(b=\pm a\), put
\(j_{\omega,b}=A_{\omega,b}^{-1}\),
\(G_b(t)=\mathcal Mg(b+it)\), and
\(P(w)=\int e^{iwx}\phi'(x)\,dx\). The sequential
core difference \(D_{\omega,b}=M_\phi-
Z_{b,\omega,S}M_\phi J_{b,\omega,S}\), after smoothing,
has the exact Fourier kernel
\[
 k_{\omega,b,g}(t,u)
   =\frac{i}{2\pi}G_b(t)A_{\omega,b}(t)
            \frac{j_{\omega,b}(t)-j_{\omega,b}(u)}{t-u}P(t-u).
 \tag{MA-MO.184}
\]
The canceled diagonal has its ordinary limiting value.
The RS proofs show that this is a Schwartz kernel: in the
gap regime use polynomial reciprocal derivatives and
Schwartz decay; in the generic regime the two Gevrey
Fourier tails dominate the logarithmic-square growth.
The argument controls every mixed derivative and applies
to all four line/character choices.

Write \(K_{\omega,b}(g)\) for its extension. It is trace
class on every \(H_{b,s}=L^2(e^{2bx}(1+x^2)^{s/2}dx)\).
For completeness, with \(r=s/2\) and
\(\Lambda^r=(1-\partial_t^2)^{r/2}\),
\[
 \begin{gathered}
       \|u\|_{H_{b,s}}=\|\Lambda^rV_bu\|_2,\\
 k^{(r)}_{\omega,b,g}
                  =\Lambda_t^r\Lambda_u^{-r}k_{\omega,b,g}.
 \end{gathered}
 \tag{MA-MO.185}
\]
The second line is the kernel of the operator conjugated by
\(\Lambda^r\) in the \(V_b\) Fourier picture. It remains
Schwartz, since the two Bessel multipliers and their
derivatives have polynomial growth. Its Hermite expansion
is absolutely summable in rank-one trace norms. This proves
trace class in the actual weighted space, as in MA-MO-010.

\emph{Agreement of these kernels on every smooth ideal vector.}
The unconditional division lemma supplies
\(u=V_by\in\mathscr S\) and
\(q=j_{\omega,b}u=V_bF\in\mathscr S\) for every
\(y\in\mathscr I^c_\omega\). In either regime the line
has no zeros. Choose compact smooth frequency cutoffs
\(\chi_N\to1\), and set \(u_N=\chi_Nu\). Then
\[
       u_N\longrightarrow u,\qquad
       j_{\omega,b}u_N=\chi_Nq\longrightarrow q
                              \quad\hbox{in Schwartz topology}.
 \tag{MA-MO.186}
\]
Multiplication by \(\phi\) in physical coordinates and
by \(A_{\omega,b}\) in Fourier coordinates preserves
Schwartz topology. Thus the sequential outputs converge
there, and in every Sobolev norm. The bounded kernel
extensions agree on \(u_N\) and hence on \(u\).
This proves agreement with the actual geometric outputs
(MA-MO.180). It does not assume that the reciprocal
multiplier sends every Schwartz function to Schwartz
in the generic regime. The same argument applies to
\(Ry\) in the dual sector.

\emph{Assembly in the inherited sum norm.}
Let \(\psi(x)=1-\phi(x)-\phi(-x)\), a compact smooth
function, and \(g^\vee(x)=g(-x)\). For each \(b=\pm a\)
the actual smoothing in (MA-MO.183) agrees on its entire
smooth ideal core with
\[
        T_{\eta,b}(g)=C_gM_\psi+K_{\eta,b}(g)
                                +RK_{\overline\eta,-b}(g^\vee)R.
 \tag{MA-MO.187}
\]
Indeed \(D_{\eta,b}+RD_{\overline\eta,-b}R\) has
the multiplication part \(\phi+\phi(-\cdot)\), while
(MA-MO.183) has the multiplication part one.
Reflection is unitary \(H_{b,s}\to H_{-b,s}\) and
\(C_gR=RC_{g^\vee}\). The remaining term \(C_gM_\psi\)
has the Schwartz Fourier kernel
\(G_b(t)\widehat\psi_+(t-u)/(2\pi)\), so is also
trace class in every polynomial weight. Thus (MA-MO.187)
is trace class on the specified one-sided space.

Set \(\mathbb H=H_{a,s}\oplus H_{-a,s}\), and define
\[
 \iota:H_{\Sigma,a,s}\longrightarrow\mathbb H,
       \qquad\iota y=(y,y),\qquad
 \mathbb T_{\eta,g}
       =\operatorname{diag}(T_{\eta,a}(g),T_{\eta,-a}(g)).
 \tag{MA-MO.188}
\]
The map \(\iota\) is isometric with closed range.
On every \(y\in\mathscr I^c_\eta\), the two components
of \(\mathbb T_{\eta,g}\iota y\) agree with the same
actual source vector \(B^{\rm raw}_{\eta,g}y\), which
belongs to \(\mathscr I^c_\eta\). Passing to the defined
closure therefore gives a diagonal output in
\(\iota J_{\eta,a,s}\) for every input from that subspace.
Consequently
\[
 \begin{aligned}
 B_{\eta,g}&=\left.\iota^*\mathbb T_{\eta,g}\iota
                                           \right|_{J_{\eta,a,s}},\\
 \|B_{\eta,g}\|_1
       &\le\|T_{\eta,a}(g)\|_{1,H_{a,s}}
                         +\|T_{\eta,-a}(g)\|_{1,H_{-a,s}}.
 \end{aligned}
 \tag{MA-MO.189}
\]
On this domain \(\iota^*\) is the inverse diagonal
identification of the already compatible outputs; it
does not average inconsistent extensions. Composition
of trace-class and bounded maps gives (MA-MO.172).
The range lies in the same closed ideal, so this time
the map is an endomorphism. Density gives uniqueness.
The finite Schwartz-seminorm estimates in the kernel
proof give all the stated test-family bounds.
\end{proof}

\paragraph{The legitimate trace and the projection that cannot be omitted.}
Let \(\iota_J=\iota|_{J_{\eta,a,s}}\) and
\(\Pi_J=\iota_J\iota_J^*\), the orthogonal projection
onto its actual diagonal ideal in \(\mathbb H\).
Bounded trace cyclicity gives
\[
       \operatorname{Tr}_{J_{\eta,a,s}}B_{\eta,g}
          =\operatorname{Tr}_{\mathbb H}(\mathbb T_{\eta,g}\Pi_J).
 \tag{MA-MO.190}
\]
For example, one absolutely summable rank-one representation
of \(\mathbb T_{\eta,g}\) gives both sides as the same
series, with the same \(\Pi_J\) pairing. This is an
actual operator identity, not a prescription by zero sums.
It does not permit dropping \(\Pi_J\), and hence does
not identify this trace with the sum of the two unrestricted
one-sided traces or with Meyer's full common-space trace.
The endomorphism has its ordinary Fredholm determinant,
but no arithmetic evaluation of it is asserted here.

\paragraph{Requested characters, source placement, and the remaining gate.}
For a nontrivial primitive Dirichlet character, the stipulated
input A supplies the gap in (MA-MO.171) at every fixed width.
For a non-induced imaginary-quadratic character, its own
Hecke RH, or the stipulated input M, supplies that gap;
the generic/Gevrey assertion for almost every width is
unconditional. In particular the result applies to both
\(\psi_{269,3}\) and its conjugate with their primitive
ramified factors and structural place set from MA-CH/MA-CQ.
The principal character is excluded: its polar terms lead
to the distinct domain and correction in MA-MO-007--011.

The exact inherited sources are MT-M18 (M18.4) and the
complete selected-height and minimum-modulus proof
(M18.7)--(M18.9); MT-M24 (M24.3), (M24.6)--(M24.9) for
the actual averaged theta image; Meyer (2005),
\texttt{lem:summation}, \texttt{lem:Hilm\_p\_embed}, the
defining \(P_+^S,P_-\) maps and
\texttt{eq:Poisson\_summation}, pp.~39--42; and MA-CC.8--.9
for the centered compact-character and dual-sector factors.
The inverse proof uses those full multiplicity and actual
image statements, not a guessed adelic representative.

The finite-order character is the previously specified
rank-one finite-monodromy local system; its Euler factor
is the corresponding Frobenius determinant. The analytic
norm-half centering here is not an integral Tate twist.
In the inherited cote~1 trace framework, scans~079--080
and~111--112 through BR-TR-01/02, (MA-MO.189) supplies an
actual nuclear endomorphism and (MA-MO.190) its valid
pairing. No weight-zero positivity or global zero bound
is inferred from that placement.

Finally the domain is the completed \emph{theta ideal},
not all of \(\mathcal E\) completed in the narrow norm.
The proof does not establish the original action on every
smooth zero-quotient class, identify the full Meyer trace
with a narrow quotient trace, remove the projector in
(MA-MO.190), or erase the discarded-divisor gate of
MA-FT-002. Those assertions require additional maps or
estimates. The improvement here is the actual original
common-sum endomorphism on every smooth ideal input and
its trace-class extension in the explicitly stated regimes.

\section{The pole-free closure is not a convolution submodule}
\label{sec:pole-free-module-obstruction}

The correction in MA-MO-011 gives an endomorphism on a genuine
Hilbert space, but the corrected domain is not preserved by the
translation representation. The failure holds for every nonzero
compact smooth convolution, including tests vanishing at logarithmic
zero. This supplies a precise obstruction to an equivariant
retraction onto that domain, without ruling out a different transport.

Keep \(J,K,\Lambda,v,\phi\) and their exact inherited norms
from MA-MO-011. Thus its fixed-gap smooth or generic-width Gevrey
regime is assumed for the cutoff and the two moment functionals.
In the present result \(g\in C_c^\infty(\mathbb R)\) may be
arbitrary in either regime: only boundedness of convolution is
needed, not the trace-class assertion for the raw smoothed defect.
Let
\[
 (U_\tau y)(x)=y(x-\tau),\qquad
 C_gy=\int g(z)U_zy\,dz.
 \tag{MA-MO.200}
\]

\paragraph{MA-MO-013 (no equivariant retraction onto the pole-free closure).}
Both \(U_\tau\) and \(C_g\) are bounded on the inherited
sum Hilbert space and preserve its closed theta ideal \(J\).
Nevertheless
\[
 \begin{aligned}
 U_\tau(K)&\not\subseteq K&&\text{for every }\tau\ne0,\\
 \Lambda C_g:K&\longrightarrow\mathbb C^2
     \text{ is onto}&&\text{for every }g\ne0.
 \end{aligned}
 \tag{MA-MO.201}
\]
The second assertion holds already on genuine compact additive
theta inputs in \(\mathscr U_{\rm pf}\). It is therefore
stronger than failure of invariance in the Hilbert closure.

No bounded retraction \(P:J\to J\) with
\(P^2=P\) and \(\operatorname{ran}P=K\) can commute with
any nontrivial \(U_\tau\), or with any nonzero \(C_g\)
of this compact smooth class. In particular the orthogonal
projection (MA-MO.140) is not a convolution-module retraction.
In fact the only linear subspace of \(K\) invariant under
all \(U_\tau\), \(\tau\in\mathbb R\), is \(\{0\}\).
Consequently any intertwiner from a representation of this
translation group into \(J\) whose range is contained in
\(K\) is zero.

\begin{proof}
\emph{Actual bounded representation and smooth source.}
For the weight \(w\) of (MA-MO.137), the ratio
\(w(x+\tau)/w(x)\) is bounded uniformly in \(x\) for
each fixed \(\tau\), locally uniformly in \(\tau\).
Indeed the exponential part is bounded by
\(e^{2a|\tau|}\), and
\(1+(x+\tau)^2\le2(1+x^2)(1+\tau^2)\)
bounds the polynomial factor. Thus translation and its
inverse are bounded, and integration over the compact
support of \(g\), followed by the triangle inequality
in the Hilbert norm, proves boundedness of \(C_g\).

On the even additive Schwartz source these operations are
\[
 \begin{aligned}
 \rho_\tau f(t)&=e^{-\tau/2}f(e^{-\tau}t),\\
 \mathcal D_g f(t)&=\int g(z)e^{-z/2}f(e^{-z}t)\,dz,\\
 U_\tau\Theta f&=\Theta\rho_\tau f,\qquad
 C_g\Theta f=\Theta\mathcal D_gf.
 \end{aligned}
 \tag{MA-MO.202}
\]
The sums and integrals interchange absolutely, with all
derivatives on compact norm intervals. These formulas also
follow directly from the actual theta sum in MA-MO.83.
Dilation and its compact integral preserve even Schwartz
space, the value zero at the additive origin and the zero
integral. Hence they preserve \(\mathscr U_0\), and
boundedness gives preservation of \(J\) by taking closure.

\emph{A compact smooth convolution cannot reproduce the cutoff.}
Use the actual functions \(W_g\) of (MA-MO.149). We claim
\[
        W_g\notin\mathbb C v
                      \quad(g\in C_c^\infty,\ g\ne0).
 \tag{MA-MO.203}
\]
If \(W_g=c v\), differentiate at \(t=e^x\) and take
the Fourier transform in \(x\), just as in (MA-MO.150).
It gives
\[
       \bigl(\mathcal Mg(1/2-it)-c\bigr)P(t)=0,
               \qquad P(t)=\widehat{\phi'}_+(t),\quad P(0)=-1.
 \tag{MA-MO.204}
\]
The entire function \(\mathcal Mg(1/2-it)\) therefore
equals \(c\) on an interval, hence everywhere. Along
the real \(t\) axis it is a Schwartz Fourier transform,
so it tends to zero at infinity; consequently \(c=0\).
Fourier injectivity then gives \(g=0\), a contradiction.
This proves (MA-MO.203), including the possibility \(c=0\),
and applies as well to \(g^\vee\).

\emph{Both outgoing convolution moments have full rank.}
Let \(W=W_g\) and \(U=W_{g^\vee}\). The bilinear
scaling and Fourier identities proved in MA-MO-011 give,
without the source operator \(L\),
\[
 \lambda_0(C_g\Theta f)=\int fW,\qquad
 \lambda_1(C_g\Theta f)=\int f\mathfrak FU
                        \quad(f\in\mathscr U_{\rm pf}).
 \tag{MA-MO.205}
\]
All functions here are ordinary smooth functions and all
pairings are absolutely convergent. We prove independence
of the five even functions
\[
                     1,\quad v,\quad\mathfrak Fv,
                              \quad W,\quad\mathfrak FU
                  \quad\text{on }\mathbb R\setminus\{0\}.
 \tag{MA-MO.206}
\]
Suppose \(A W+B\mathfrak FU=c_0+c_1v+c_2\mathfrak Fv\).
All terms other than \(c_0\) are Schwartz, so decay gives
\(c_0=0\). Outside the compact supports of \(W\) and
\(v\), the identity becomes
\(B\mathfrak FU=c_2\mathfrak Fv\). Both transforms
are entire, because \(U\) and \(v\) are compact smooth.
It follows by the identity theorem and Fourier inversion
that \(BU=c_2v\). If \(B\ne0\), this contradicts
(MA-MO.203) for \(g^\vee\). Thus \(B=0\) and
\(c_2=0\). The remaining identity \(AW=c_1v\)
contradicts (MA-MO.203) if \(A\ne0\). Hence \(A=0\)
and \(c_1=0\), proving the claimed independence.

As in MA-MO.97, choose five positive evaluation points
with an invertible matrix for these five functions,
replace each point and its negative by an even bump
of mass one supported away from zero, and take the
bumps narrow enough that their moment matrix remains
invertible. Given arbitrary \((z_0,z_1)\in\mathbb C^2\),
solve the resulting system with right side
\((0,0,0,z_0,z_1)^{\mathsf T}\). Its actual compact
even solution \(f\) has
\[
        f(0)=\int f=\ell_0(f)=\ell_1(f)=0,
                \qquad \int fW=z_0,\quad\int f\mathfrak FU=z_1.
 \tag{MA-MO.207}
\]
Equations (MA-MO.139) and (MA-MO.205) prove the second
assertion of (MA-MO.201), on the claimed smooth domain.
The invertible matrix exists because otherwise the span
of all evaluation vectors would have a nonzero annihilator,
contradicting (MA-MO.206); its bump approximation follows
from continuity of these functions. Thus the construction
does not rely on a distributional input.

\emph{Every nontrivial translation also fails.}
The first moment of \(U_\tau\Theta f\) pairs \(f\)
with \(v_\tau(t)=e^{\tau/2}v(e^\tau t)\).
If \(v_\tau\) belonged to
\(\operatorname{span}\{1,v,\mathfrak Fv\}\) away zero,
the exterior identity and decay would remove the constant
term, and entire continuation would remove the
\(\mathfrak Fv\) term. Thus \(v_\tau=c v\).
Near zero both undilated cutoff values are one, so
\(c=e^{\tau/2}\). This would imply
\(v(e^\tau t)=v(t)\) everywhere. Iterating the dilation
or its inverse takes a fixed positive argument outside
the compact support, contradicting \(v=1\) near zero.
Therefore the four functions \(1,v,\mathfrak Fv,v_\tau\)
are independent for \(\tau\ne0\). The same four-bump
construction gives \(f\in\mathscr U_{\rm pf}\) with
\(\lambda_0(U_\tau\Theta f)=1\). This proves the first
assertion of (MA-MO.201).

\emph{The retraction obstruction.}
If \(P\) is such a retraction and commutes with a
bounded operator \(T\) on \(J\), then for every
\(y\in K\) one has \(Ty=TPy=PTy\in K\).
Apply this to either class of operators in (MA-MO.201)
to obtain a contradiction. This proves the stated
impossibility for every proposed bounded retraction,
not only for the orthogonal one.

\emph{No nonzero invariant subspace remains inside the kernel.}
There is also a direct stronger argument, using the actual
functional constructed in MA-MO-011. Let \(y\in H_\Sigma\),
put \(u=V_ay\in L^2\), and retain \(\beta_a\) of
(MA-MO.143). Fourier scaling gives
\[
 \lambda_0(U_\tau y)
   =e^{a\tau}\int_{\mathbb R}u(t)\beta_a(t)e^{it\tau}\,dt.
 \tag{MA-MO.208}
\]
The product is integrable by Cauchy--Schwarz, since
\(\beta_a\in L^2\). Moreover \(\beta_a\ne0\) almost
everywhere. Indeed \(j_a\) is finite and nonzero on the
chosen line in both stated regimes. By (MA-MO.144), the
remaining factor is the nonzero entire transform of
\(e^{cx}\phi'(x)\), with \(c=1/2-a\), divided by a nonvanishing linear
factor on the real axis. The numerator is not identically
zero by Fourier uniqueness, and its real zeros are discrete.

If every translate of \(y\) belongs to \(K\), the left
side of (MA-MO.208) vanishes for every real \(\tau\).
Uniqueness of the Fourier transform of an \(L^1\) function
gives \(u\beta_a=0\) almost everywhere, hence \(u=0\)
and \(y=0\). Every element of a translation-invariant
linear subspace of \(K\) has that property, proving the
stronger assertion without a closedness assumption on
that subspace. If an intertwiner has range in \(K\),
the orbit of each of its values is also in \(K\), so
each value is zero. No claim about a different target
space or an unequivariant map is implied.
\end{proof}

\paragraph{Scope for an intertwiner and a trace.}
The original representation on \(J\) is well defined,
and MA-MO-011 still gives legitimate corrected endomorphisms
and determinants on \(K\). What fails is using a retraction
onto \(K\) as a translation- or convolution-equivariant
domain repair. In particular the rank-two correction is
not an intertwining argument for the original source
representation. This holds also for every nonzero compact
test with \(g(0)=0\); that pointwise condition does not
remove either independent outgoing moment.

The result concerns this precisely defined rational
pole-free closure in the two analytic regimes already
stated. It neither rules out all global-to-local transports
nor changes the exact arithmetic deficiency condition in
MA-FT-002. In the inherited cote~1 framework, scans~079--080
and~111--112 through BR-TR-01/02, existence of an ordinary
nuclear endomorphism trace is separate from compatibility
with a representation or its convolution algebra. The
finite-length \(K_0\) bookkeeping does not supply that
missing equivariance.

\section{Evaluating the actual projected ideal trace by finite Gram data}
\label{sec:projected-ideal-gram}

The trace in MA-MO.190 contains the orthogonal projection onto the
actual diagonal ideal.  This section retains that projection.  It
gives a finite theta Gram evaluation, an explicit error certificate,
and a comparison between polynomial weights by actual nuclear maps.
It does not identify this trace with the full zero-quotient trace.

Fix the character, place set, width and one of the two regimes of
MA-MO-012.  For a polynomial exponent \(s\ge0\), write
\[
 \begin{gathered}
 w_s(x)=(e^{2ax}+e^{-2ax})(1+x^2)^{s/2},\qquad
 H_s=L^2(w_s\,dx),\qquad J_s=\overline{\mathscr I^c_\eta}^{H_s},\\
 P_s=P_{J_s},\qquad
 \mathcal A_{g,s}=\iota_s^*\mathbb T_{\eta,g}\iota_s
       \in\mathcal S_1(H_s),\qquad
 B_{g,s}=\mathcal A_{g,s}|_{J_s}.
 \end{gathered}
 \tag{MA-MO.210}
\]
The operator on the right is the original source smoothing of
MA-MO-012 on the actual completed ideal.  In particular
\(\mathcal A_{g,s}P_s\) has its range in \(J_s\), and
\(\operatorname{Tr}_{J_s}B_{g,s}=\operatorname{Tr}_{H_s}
(\mathcal A_{g,s}P_s)\).  No projection has been deleted.

\paragraph{MA-MO-014 (cross-weight comparison and a certified Gram evaluation).}
The actual ideal trace and its ordinary Fredholm determinant are
independent of \(s\ge0\).  More precisely, there are trace-class maps
\[
 B_g^{s,t}:J_s\longrightarrow J_t\quad(s,t\ge0)
 \tag{MA-MO.211}
\]
agreeing on every smooth ideal vector with the same original source
action.  Their compositions with the natural inclusions give the
endomorphisms in (MA-MO.210), and therefore
\[
 \operatorname{Tr}_{J_s}B_{g,s}=\operatorname{Tr}_{J_t}B_{g,t},
 \qquad
 \det_{J_s}(1-zB_{g,s})=\det_{J_t}(1-zB_{g,t}).
 \tag{MA-MO.212}
\]

For the two requested character families, namely nonprincipal
primitive Dirichlet characters and the characters over
\(\mathbb Q(\zeta_3)\) covered by MA-CGP-001, the same trace admits
the exact Gram-resolvent evaluation below in the metric \(w_s\).
Its finite-matrix approximation has the error bound
(MA-MO.226).  This is an evaluation of the completed-ideal trace,
with no zero-location premise added to the specified regime.

\begin{proof}
\emph{Actual maps between polynomial weights.}
Let \(k_b(t,u)\), \(b=\pm a\), be the full Fourier kernel of
\(T_{\eta,b}(g)\) in (MA-MO.187), including its compact residual.
It is Schwartz.  Use the unitary Fourier transform
\[
 V_b f(t)=(2\pi)^{-1/2}\int_{\mathbb R}e^{(b+it)x}f(x)\,dx,
 \qquad \Lambda^r=(1-\partial_t^2)^{r/2}.
\]
The map from \(H_{b,s}\) to \(H_{b,t}\) has, after these
unitary identifications, the Schwartz kernel
\[
                  \Lambda_\xi^{t/2}\Lambda_\upsilon^{-s/2}
                                           k_b(\xi,\upsilon).
 \tag{MA-MO.213}
\]
Here \(t\) is the target polynomial exponent, while
\(\xi,\upsilon\) are kernel coordinates.  Bessel multipliers of either real order
preserve Schwartz space.  The Hermite expansion thus proves
trace class between these two Hilbert spaces.

Restrict the two input components to \(J_s\).  Approximate each
input in its \(H_s\) norm by vectors of \(\mathscr I^c_\eta\).
On those vectors the two output components are the same actual
source vector, and lie in \(\mathscr I^c_\eta\).  The two
trace-class extensions converge in the respective target norms.
The closed diagonal in \(H_{a,t}\oplus H_{-a,t}\) therefore
contains the limit, which belongs to its closed ideal \(J_t\).
This proves (MA-MO.211), including its actual range assertion.
For \(t\ge s\), let \(j_{t,s}:J_t\to J_s\) be inclusion;
its norm is at most one.  Density on the common core gives
\[
 B_{g,s}=j_{t,s}B_g^{s,t},\qquad
 B_{g,t}=B_g^{s,t}j_{t,s}.
 \tag{MA-MO.214}
\]
Bounded/nuclear trace cyclicity gives the first equality in
(MA-MO.212).  The determinant equality follows first for
finite-rank approximations to \(B_g^{s,t}\) and then by
trace-norm continuity of the Fredholm determinant, exactly as
in MA-GP.75.  This argument compares actual operators; it does
not identify the orthogonal projections in different metrics.

\emph{The actual finite-band source and its exact metric.}
For the remainder of this proof fix \(s\), and write
\(H=H_s\), \(J=J_s\), \(P=P_s\), and
\(\mathcal A=\mathcal A_{g,s}\).  Use the notation of MA-CGP:
\[
 \begin{gathered}
 Q=|D_K|N\mathfrak f_\eta,\qquad
 a_n=\sum_{N\mathfrak a=n}\overline{\eta(\mathfrak a)},\qquad
 \gamma_K=2\ \text{or}\ 2\pi,\\
 \mathcal P_\Lambda=L^2((0,\Lambda),\gamma_K\,dy),\qquad
 (T_\Lambda U)(t)=
   \frac{w_K\epsilon_\eta^\theta}{\sqrt{Qt}}
                        \sum_{n\ge1}a_nU(n/(Qt)).
 \end{gathered}
 \tag{MA-MO.215}
\]
The initial domain is the actual smooth archimedean frequency
core of the required parity or radial type, with support in the
open band.  It includes the origin whenever that type permits;
no extra moment condition is imposed.  The phase and the scale
are exactly MA-CGP.6 and .16.

The target norm is now exactly
\(\int_0^\infty |F(t)|^2w_s(\log t)\,dt/t\).
The ratio of this weight to the maximum weight in MA-CGP is
between one and two.  Consequently closability and the graph
closure of \(T_\Lambda\) persist, as does its actual range
closure.  We denote the closed map by \(\overline T_\Lambda\),
its Gram operator by \(G_\Lambda=\overline T_\Lambda^*
\overline T_\Lambda\), and its range projection by \(Q_\Lambda\).
The Gram operator is formed in this exact sum metric, not copied
from an equivalent norm.

The increasing ranges have dense union in \(J\).  To see this
in the stated two-sided norm, start with any actual smooth ideal
vector.  The topological isomorphism MA-CGP.3 writes it as
\(\Theta_\eta f\).  Smooth compact frequency cutoffs preserving
the parity or radial type approximate \(f\) in Schwartz topology.
Their theta images converge in \(\mathcal E\), by the actual
isomorphism, and hence in every \(H_s\).  This proves density
without any finite-level inverse bound.  Thus
\[
 Q_\Lambda\longrightarrow P\ \text{strongly},\qquad
 \|\mathcal A Q_\Lambda-\mathcal A P\|_1\longrightarrow0.
 \tag{MA-MO.216}
\]
The second assertion follows by finite-rank approximation of
the fixed trace-class operator \(\mathcal A\).  It is not a
trace-norm assertion about the moving quotient convolutions
of MA-TN-001--002.

\emph{Normal-equation transport with all domains retained.}
Let \(\overline T_\Lambda=V_\Lambda G_\Lambda^{1/2}\) be its
polar decomposition.  For \(\varepsilon>0\), put
\[
 \begin{aligned}
 R_{\Lambda,\varepsilon}&=(G_\Lambda+\varepsilon)^{-1/2},\\
 F_{\Lambda,\varepsilon}&=\overline T_\Lambda
             R_{\Lambda,\varepsilon}
   =V_\Lambda\sqrt{G_\Lambda/(G_\Lambda+\varepsilon)}.
 \end{aligned}
 \tag{MA-MO.217}
\]
These are bounded maps on their entire source;
\(\|F_{\Lambda,\varepsilon}\|\le1\).  Their products
\(F_{\Lambda,\varepsilon}F_{\Lambda,\varepsilon}^*\)
converge strongly to \(Q_\Lambda\).  Their use does not
assert boundedness of the unregularized map \(T_\Lambda\).

Choose real compact smooth cutoffs \(0\le\chi_R\le1\),
increasing to one in the norm variable, and put
\[
 T_{\Lambda,R}=M_{\chi_R}\overline T_\Lambda,
 \qquad K_{\Lambda,R}=T_{\Lambda,R}^*\mathcal A T_{\Lambda,R}.
 \tag{MA-MO.218}
\]
The first map extends boundedly to \(\mathcal P_\Lambda\)
by the finite Poisson sum in (MA-MO.215); its values there
agree with the displayed graph-domain composition.
Consequently \(K_{\Lambda,R}\) is trace class.  Without
ever asserting that \(\overline T_\Lambda^*\mathcal A
\overline T_\Lambda\) is bounded, we have
\[
 \begin{aligned}
 R_{\Lambda,\varepsilon}K_{\Lambda,R}R_{\Lambda,\varepsilon}
 &=F_{\Lambda,\varepsilon}^*M_{\chi_R}\mathcal A
                    M_{\chi_R}F_{\Lambda,\varepsilon},\\
 \operatorname{Tr}_{J}B_{g,s}
 &=\lim_{\Lambda\to\infty}\lim_{\varepsilon\downarrow0}
       \lim_{R\to\infty}
       \operatorname{Tr}_{\mathcal P_\Lambda}
       (R_{\Lambda,\varepsilon}K_{\Lambda,R}
                                      R_{\Lambda,\varepsilon}).
 \end{aligned}
 \tag{MA-MO.219}
\]
Indeed the first line follows on the whole source because
\(R_{\Lambda,\varepsilon}\) takes its values in the graph
domain.  The error at the innermost limit is bounded by
\(\|\mathcal A-M_{\chi_R}\mathcal A M_{\chi_R}\|_1\),
uniformly in \(\Lambda,\varepsilon\).  At fixed \(\Lambda\),
strong convergence of the regularized polar maps and their
adjoints gives trace-norm convergence to
\(V_\Lambda^*\mathcal A V_\Lambda\).  Cyclicity gives its
trace as \(\operatorname{Tr}(\mathcal A Q_\Lambda)\), and
(MA-MO.216) proves the outer limit.  No unrestricted joint
choice of \(\varepsilon\) and \(\Lambda\) is implicit.

\emph{An explicit finite local kernel.}
Write \(\mathfrak a(t,z)\) for the kernel of \(\mathcal A\)
relative to \(dz/z\).  It is smooth on compact positive
rectangles.  Explicitly, if \(\ell_b(x,y)\) is the inverse
double Fourier transform of \(k_b\), then
\[
 \begin{aligned}
 \mathfrak a(e^x,e^y)
   &=\sum_{b=\pm a}\theta_b(x)e^{-b(x-y)}\ell_b(x,y),\\
 \theta_b(x)&=\frac{e^{2bx}}{e^{2ax}+e^{-2ax}},\qquad
 \ell_b(x,y)=\frac1{2\pi}\iint
                 e^{-itx}k_b(t,u)e^{iuy}\,dt\,du .
 \end{aligned}
 \tag{MA-MO.220}
\]
The adjoint of the inclusion \(H_s\to H_{b,s}\) is
precisely multiplication by \(\theta_b\).  This proves
(MA-MO.220) in the exact metric, including its normalization.

Substitute (MA-MO.215) twice into (MA-MO.218), and set
\(y=m/(Qt)\), \(v=n/(Qz)\).  The resulting action, with
integration measure \(dv\), is
\[
 \begin{aligned}
 (K_{\Lambda,R}U)(y)
  ={}&\frac{w_K^2}{\gamma_K}
       \sum_{m,n\ge1}\frac{\overline{a_m}a_n}{\sqrt{mn}}
       \int_0^\Lambda
       \frac{w_s(\log(m/(Qy)))}{\sqrt{yv}}\\
     &\quad\cdot\chi_R(m/(Qy))\chi_R(n/(Qv))
       \mathfrak a(m/(Qy),n/(Qv))U(v)\,dv .
 \end{aligned}
 \tag{MA-MO.221}
\]
If the cutoff is supported in \([r_-,r_+]\subset(0,\infty)\),
only \(m,n\le Q\Lambda r_+\) occur, and every nonzero
term is away from \(y=0\) and \(v=0\).  This is a finite
smooth-kernel formula on the actual auxiliary band space.
Its trace-class property also follows directly by extending
each kernel smoothly to a larger square and using an
absolutely summable double Fourier series.  The scalar theta
phase cancels; the primitive scale \(Q\), the conjugate
coefficients, and the exact target weight do not.

For reference the localized positive Gram form uses the
finite rational-dilation operator
\[
 (G_{\Lambda,R}U)(y)=\frac{w_K^2}{\gamma_K}
       \sum_{m,n\ge1}\frac{\overline{a_m}a_n}{m}
       \chi_R(m/(Qy))^2w_s(\log(m/(Qy)))U(ny/m),
 \tag{MA-MO.222}
\]
with extension by zero off the band.  Its forms increase
on the graph core to \(\|\overline T_\Lambda U\|_H^2\).
This supplies the actual normal equation in (MA-MO.217);
an uncut double sum is not interchanged with an integral.

\emph{A finite-matrix evaluation with a stopping certificate.}
Choose a countable sequence of actual smooth frequency tests,
each lying in some finite band, whose theta images \(v_j\)
have dense span in \(J\).  Such a sequence follows from the
frequency approximation just proved and separability of the
Schwartz source; delete linear dependencies.  Let
\(V_N:\mathbb C^N\to H\) have columns \(v_1,\ldots,v_N\),
and define
\[
 \begin{gathered}
 C_N=V_N^*V_N>0,\qquad D_N=V_N^*\mathcal A V_N,\qquad
 P_N=V_NC_N^{-1}V_N^*,\\
 \tau_N=\operatorname{tr}(C_N^{-1}D_N)
                   =\operatorname{Tr}_H(\mathcal A P_N).
 \end{gathered}
 \tag{MA-MO.223}
\]
Every matrix entry is an actual theta pairing.  It can be
computed by compact truncation using (MA-MO.221)--(MA-MO.222),
with the Schwartz tails of these finitely many smooth vectors
retained in the truncation error.

The complementary data are equally explicit.  For any finite
set of retained zero jets from MA-WP.3, take their Riesz vectors
\[
 h_{\rho,j}(x)=\frac{x^j e^{\overline{(\rho-1/2)}x}}{w_s(x)}.
 \tag{MA-MO.224}
\]
The inner product is linear in its first variable, so
\(\langle f,h_{\rho,j}\rangle=(\mathcal Mf)^{(j)}(\rho-1/2)\).
These vectors belong to \(J^\perp\) in the exact metric.
Their Gram projection \(E_M=W_M(W_M^*W_M)^{-1}W_M^*\)
is defined by their column map.  Distinct retained jets are
linearly independent, since exponential polynomials are
linearly independent.  If there are no retained jets use
\(E_M=0\).  The full exhaustive family spans \(J^\perp\)
by MA-WP.4, which applies to the sum norm as stated there.
In particular \(P_NE_M=0\).

Here is a completely specified nuclear approximation of
\(\mathcal A\).  Put \(W_{b,s}=\Lambda^{s/2}V_b\), let
\(e_m\) be the normalized Hermite basis of \(L^2(\mathbb R)\),
and let \(c^b_{mn}\) be the Hermite matrix of the Schwartz
kernel \(\Lambda_t^{s/2}\Lambda_u^{-s/2}k_b\).
If \(i_b:H\to H_{b,s}\) is inclusion, put
\(u^b_m=i_b^*W_{b,s}^{-1}e_m\).  Then
\[
 \begin{gathered}
 \mathcal A=\sum_{b=\pm a}\sum_{m,n\ge0}
      c^b_{mn}\,u^b_m\otimes(u^b_n)^*,\qquad \|u^b_m\|\le1,\\
 \delta_{N,M}(u)=\|u\|^2-\|P_Nu\|^2-\|E_Mu\|^2\ge0 .
 \end{gathered}
 \tag{MA-MO.225}
\]
All sums converge in trace norm.  If
\(e_L=\sum_b\sum_{\max(m,n)>L}|c^b_{mn}|\), the error bound is
\[
 \left|\operatorname{Tr}_{J}B_{g,s}-\tau_N\right|
 \le e_L+
 \sum_{b=\pm a}\sum_{m,n\le L}|c^b_{mn}|
       \sqrt{\delta_{N,M}(u^b_m)\delta_{N,M}(u^b_n)} .
 \tag{MA-MO.226}
\]
To prove it, \(D=P-P_N\) is an orthogonal projection and
\(0\le D\le1-P_N-E_M\).  Hence
\( |\langle Du,v\rangle|
 \le\sqrt{\delta_{N,M}(u)\delta_{N,M}(v)}\).
Apply this to the finite part of (MA-MO.225); the remaining
trace is bounded by its nuclear tail.  Every quantity on
the right uses finite Gram matrices and explicitly specified
kernel coefficients.  An incomplete finite set of certified
zero jets still gives a valid upper bound; it is not assumed
to list all zeros.

The nuclear tail has a concrete seminorm majorant.  With
\(\mathcal H=-\partial_t^2+t^2\), for any integer \(q>1\)
let
\[
 C_{b,q}=\| (\mathcal H_t+1)^q(\mathcal H_u+1)^q
           (\Lambda_t^{s/2}\Lambda_u^{-s/2}k_b)\|_{L^2(\mathbb R^2)}.
\]
Integration against Hermite eigenfunctions gives
\[
 \begin{aligned}
 |c^b_{mn}|&\le C_{b,q}(2m+2)^{-q}(2n+2)^{-q},\\
 e_L&\le\sum_b C_{b,q}\left[
       \left(\sum_{m\ge0}(2m+2)^{-q}\right)^2
       -\left(\sum_{m=0}^L(2m+2)^{-q}\right)^2\right].
 \end{aligned}
 \tag{MA-MO.227}
\]
Thus the bound does not hide an unknown projector norm.
For fixed \(L\), \(P_N\to P\) and \(E_M\to1-P\) strongly
make its finite part tend to zero.  Then \(L\to\infty\)
makes its tail zero.  This proves a convergent finite-matrix
evaluation of the specified trace.
\end{proof}

\paragraph{What is and is not evaluated.}
Equations (MA-MO.219) and (MA-MO.223)--(MA-MO.227) evaluate
\(\operatorname{Tr}_{J_s}B_{g,s}\), retaining the theta Gram
normalization and the actual ideal projection.  The bound is
an a posteriori analytic certificate when its finite pairings,
certified zero jets and kernel-seminorm bounds are supplied.
The full-measure generic-width theorem by itself does not
provide an effective zero-distance constant for an arbitrarily
chosen decimal width, and no uniform numerical convergence
rate is claimed here.

The source inputs are MA-CGP-001--003, especially .3 and
.15--.19; MA-GP-005/008 for the closed normal equation and
bounded regularization; MA-WP-001 for the exact complementary
jet span; and the actual kernels and source agreement of
MA-MO-012.  The finite matrices use the inherited two-sided
metric throughout.  They are Hilbert Gram matrices, not
categorical identifications of different completions.  Their
trace limits use the actual nuclear pairing of the inherited
cote~1 framework.  The retained jet spaces are finite-length
primary modules, but their algebraic classes alone do not
determine this Gram correction.

No equality with \(\operatorname{Tr}C_g\) on the full smooth
zero quotient follows from this evaluation.  Such an equality
requires an additional actual comparison or return map.
In particular no part of (MA-MO.212) or (MA-MO.226) removes
the omitted-divisor condition of MA-FT-002.

\section{The exact Blaschke correction in the projected ideal trace}
\label{sec:projected-ideal-blaschke}

At polynomial weight zero the actual ideal projector has an
explicit Hardy-space form.  This makes its contribution to the
trace a concrete, absolutely convergent kernel integral.  The
cross-weight nuclear comparison in MA-MO-014 then gives the
same value at every polynomial exponent.  The correction below
is part of the projected ideal trace; it is not silently set
to zero or identified with the full zero-quotient trace.

Keep all hypotheses of MA-MO-012.  This section allows its
general fixed number field and nontrivial primitive finite-order
character.  Use \(s=0\) initially, put
\(\Pi_a=\{z:-a<\operatorname{Re}z<a\}\), and write
\[
 \begin{gathered}
 \varphi_a(z)=\tan\frac{\pi z}{4a},\qquad
 \mathcal B(z)=B(\varphi_a(z)),\\
 Z(B)=\{\varphi_a(\rho-1/2):
                    |\operatorname{Re}(\rho-1/2)|<a\},\\
 b_\pm(t)=\mathcal B(\pm a+it).
 \end{gathered}
 \tag{MA-MO.230}
\]
Multiplicities are retained, \(B\) is the disk Blaschke product
of MA-HY-001, and \(B=1\) when the interior divisor is empty.
Its irrelevant unimodular normalization will cancel.  The
two regimes of MA-MO-012 have no zeros on either edge.

Let \(\mathcal K_b\) denote the full operator
\(V_bT_{\eta,b}(g)V_b^{-1}\), with its Schwartz kernel
\(k_b(t,u)\), for \(b=\pm a\).  This includes all terms of
(MA-MO.187), in particular the compact residual when the
original cutoff is not symmetric.  Define
\[
 l_{\pm a}(t,u)=\overline{b_\pm(t)}\,k_{\pm a}(t,u)b_\pm(u).
 \tag{MA-MO.231}
\]

\paragraph{MA-MO-015 (an explicit projected-trace correction).}
The kernels in (MA-MO.231) are Schwartz.  For every \(s\ge0\),
the actual completed-ideal endomorphism has trace
\[
 \begin{aligned}
 \operatorname{Tr}_{J_s}B_{g,s}
  ={}&\frac12\sum_{\epsilon=\pm1}
              \int_{\mathbb R} k_{\epsilon a}(t,t)\,dt
                         +\mathfrak C_{\eta,a}(g,\phi),\\
 \mathfrak C_{\eta,a}(g,\phi)
  ={}&\frac{i}{8a}\sum_{\epsilon=\pm1}\epsilon
       \int_{\mathbb R}\int_0^\infty
       \frac{l_{\epsilon a}(t,t+v)-l_{\epsilon a}(t,t-v)}
                         {\sinh(\pi v/(4a))}\,dv\,dt .
 \end{aligned}
 \tag{MA-MO.232}
\]
Both displayed integrals are absolutely convergent.  The first
term is one half of the sum of the two unrestricted one-sided
traces.  The second is the exact correction required by the
inherited diagonal-ideal projection.  No vanishing assertion
for that correction is part of this theorem.

\begin{proof}
\emph{The actual strip ideal and its boundary unitary.}
Take the centering parameter of MA-HY to be \(1/2\), its
two strip edges to be \(-a,a\), and its polynomial exponent
to be zero.  Its onto Mellin theorem and its half-derivative
conformal unitary give
\[
 \mathcal M H_0=\mathscr H^2(\Pi_a),\qquad
 \mathcal M J_0=\mathcal B\mathscr H^2(\Pi_a).
 \tag{MA-MO.233}
\]
The notation \(J_0\) in this section means the exponent-zero
ideal, not the inversion operator of the original source.
Equation (MA-MO.233) is the actual closure theorem, not a
new choice of a spectral kernel.  The half-derivative factor
commutes with multiplication as in MA-HY.7, so it introduces
no further factor into \(\mathcal B\).

Boundary values satisfy \(|b_\pm(t)|=1\) almost everywhere.
Thus \(U_{\mathcal B}=\operatorname{diag}(M_{b_+},M_{b_-})\)
is unitary on the two boundary \(L^2\) spaces.  Let
\(\widehat\Pi_D\) be the projection onto the Mellin boundary
image of the entire diagonal \(H_0\), and let
\(\widehat\Pi_J\) project onto its actual ideal image.
Equation (MA-MO.233) implies the equality of orthogonal
projections
\[
          \widehat\Pi_J
            =U_{\mathcal B}\widehat\Pi_D U_{\mathcal B}^*.
 \tag{MA-MO.234}
\]
This retains the whole ideal projector.

\emph{Boundary regularity with the exact generic scope.}
We give the needed derivative bound explicitly.  Write an
interior centered zero as \(\lambda=\beta+i\gamma\), put
\(\kappa=\pi/(2a)\), and map the strip to the upper half-plane
by \(w(z)=\exp(i\kappa(z+a))\).  Up to a constant phase,
the elementary Blaschke factor is
\((w(z)-w(\lambda))/(w(z)-\overline{w(\lambda)})\).
The magnitude of its boundary logarithmic derivative is
\[
 \begin{aligned}
 p_{\lambda,+}(t)
    &=\frac{\kappa\sin(\kappa(\beta+a))}
        {\cosh(\kappa(t-\gamma))+
                              \cos(\kappa(\beta+a))},\\
 p_{\lambda,-}(t)
    &=\frac{\kappa\sin(\kappa(\beta+a))}
        {\cosh(\kappa(t-\gamma))-
                              \cos(\kappa(\beta+a))}.
 \end{aligned}
 \tag{MA-MO.235}
\]
These follow by differentiating the elementary factor on
the two real boundary rays in the \(w\)-plane.  The signs
of the logarithmic derivatives are immaterial for the
following absolute estimates.

For \(|t-\gamma|\le1\), their \(j\)th derivatives are
bounded by \(C_{a,j}\delta_{\lambda,\pm}^{-j-1}\), where
\(\delta_{\lambda,+}=a-\beta\) and
\(\delta_{\lambda,-}=a+\beta\).  Indeed the denominator
is comparable near the relevant edge to
\((t-\gamma)^2+\delta_{\lambda,\pm}^2\); away from that
edge it is bounded below.  For \(|t-\gamma|\ge1\), all
these derivatives are bounded by
\(C_{a,j}e^{-\kappa|t-\gamma|}\), uniformly in \(\beta\).
The local completed-zero count, with multiplicity, is
\(O(\log(2+|T|))\), as in the analytic inputs and
(M18.2) of MT-M18.

In the generic regime, the distance condition and functional
reflection used in MA-MO-012 give
\(\delta_{\lambda,\pm}\ge c(1+|\gamma|)^{-2}\).
In the gap regime they give a fixed positive lower bound.
Summing (MA-MO.235) and its differentiated series therefore
gives polynomial bounds in \(t\) for every derivative of
the boundary logarithmic derivative.  The far-height sums
converge locally uniformly by their exponential tails.
Finite Blaschke products and their derivatives converge
across each finite boundary segment: no zero accumulates
there, and the same logarithmic-derivative estimates apply.
Consequently the asserted boundary values are smooth, and
\[
        |\partial_t^j b_\pm(t)|
           \le C_j(1+|t|)^{N_j}\qquad(j\ge0).
 \tag{MA-MO.236}
\]
One can obtain this last bound recursively from
\(b_\pm'=b_\pm(\log b_\pm)'\) and \(|b_\pm|=1\).
This proves that multiplying the Schwartz kernel \(k_b\)
by the two factors in (MA-MO.231) preserves Schwartz
space.  It does not assume an edge separation at widths
outside the two regimes under consideration.

\emph{The inherited diagonal projection.}
Before Fourier transform, the unweighted two-component
embedding sends \(f\) to \((e^{ax}f,e^{-ax}f)\).
Its pointwise orthogonal projection is
\[
 \Pi_D(x)=\frac1{e^{2ax}+e^{-2ax}}
       \begin{pmatrix}e^{2ax}&1\\1&e^{-2ax}\end{pmatrix},
 \qquad
 \theta_\pm(x)=\frac{e^{\pm2ax}}{e^{2ax}+e^{-2ax}}.
 \tag{MA-MO.237}
\]
Thus the two diagonal blocks of \(\widehat\Pi_D\) are
\(\mathcal F_+M_{\theta_\pm}\mathcal F_+^{-1}\), where
\(\mathcal F_+\) is unitary Fourier transform with the
positive exponential.  Using (MA-MO.190), (MA-MO.234)
and bounded trace cyclicity gives the exact identity
\[
 \operatorname{Tr}_{J_0}B_{g,0}
   =\sum_{\epsilon=\pm1}
       \operatorname{Tr}_{L^2}\!\left(
          \mathcal L_{\epsilon a}\,
          \mathcal F_+M_{\theta_\epsilon}\mathcal F_+^{-1}
                                \right),
 \tag{MA-MO.238}
\]
where \(\mathcal L_b\) has kernel \(l_b\).  Off-diagonal
blocks of \(\widehat\Pi_D\) have not been declared zero;
they simply do not enter the trace of a product whose
other factor is block diagonal.

\emph{Fourier evaluation and the sign.}
For this calculation only write
\(\widehat f_+(v)=\int e^{ivx}f(x)\,dx\), without
the unitary scalar.  Since
\(\theta_+'(x)=a\operatorname{sech}^2(2ax)\), elementary
contour integration, or the beta integral after
\(y=e^{4ax}\), gives
\[
 \begin{aligned}
 \widehat{\theta_+'}_+(v)
       &=\frac{\pi v}{4a\sinh(\pi v/(4a))},\\
 \widehat{\theta_\pm}_+(v)
       &=\pi\delta_0(v)
           \ \pm\frac{i\pi}{4a}\operatorname{PV}
                       \frac1{\sinh(\pi v/(4a))}.
 \end{aligned}
 \tag{MA-MO.239}
\]
Indeed the substitution in the first line gives
\(\int_0^\infty y^{iv/(4a)}(1+y)^{-2}\,dy
=\Gamma(1+iv/(4a))\Gamma(1-iv/(4a))\), and Euler's
reflection identity gives the displayed value, including its
continuous value at zero.  The derivative convention is
\(\widehat{f'}_+=-iv\widehat f_+\); the delta coefficient
is fixed by \(\theta_+(x)+\theta_+(-x)=1\).

For a Schwartz kernel \(l\), the physical kernel obtained
by inverse double Fourier transform is Schwartz.  Hence
the trace against the bounded multiplier \(\theta_\pm\)
is its absolutely integrable diagonal pairing.  Equivalently
it is the distributional pairing
\[
 \operatorname{Tr}\bigl(\mathcal L\mathcal F_+
                    M_{\theta_\pm}\mathcal F_+^{-1}\bigr)
     =\frac1{2\pi}\int\!\!\int
             l(t,u)\widehat{\theta_\pm}_+(u-t)\,dt\,du.
 \tag{MA-MO.240}
\]
The order \(u-t\), rather than \(t-u\), comes from the
second kernel evaluated at \((u,t)\) in the trace product.
Substituting (MA-MO.239), and pairing positive and negative
\(v=u-t\), yields
\[
 \frac12\int l(t,t)\,dt
       \ \pm\frac{i}{8a}\int_{\mathbb R}\int_0^\infty
          \frac{l(t,t+v)-l(t,t-v)}{\sinh(\pi v/(4a))}
                                                  \,dv\,dt.
 \tag{MA-MO.241}
\]
Near zero, the numerator is \(O(v)\), with an integrable
Schwartz majorant in \(t\), by the mean value formula.
For \(v\ge1\), the denominator gives exponential decay,
and the Schwartz bounds supply an integrable majorant.
The second integral is therefore absolutely convergent.
Since \(l_b(t,t)=k_b(t,t)\), (MA-MO.238)--(MA-MO.241)
prove (MA-MO.232) at \(s=0\).  The actual cross-weight
cyclicity (MA-MO.212) proves it at every \(s\ge0\).
\end{proof}

\paragraph{Scope of the explicit correction.}
The kernels \(k_b\) are the explicit original-source
Schwartz kernels of MA-MO-012.  The product \(\mathcal B\)
encodes precisely the retained interior divisor through the
already proved actual closure theorem MA-HY-001.  Thus
(MA-MO.232) gives an explicit projected-trace correction,
not a replacement of the projector by an assigned zero sum.
Its finite theta Gram evaluation and independent matrix
error certificate are MA-MO-014 for the two requested
character families.

The proof uses the exact Mellin and conformal normalization
MA-HY.3--.9, including the independent reproducing-kernel
verification MA-HY.6a--.6b, and the full kernel construction
and genuine ideal range of MA-MO-012.  The local zero count
is the inherited MT-M18 analytic input, not a newly assumed
spacing hypothesis.  The Fourier constant and the diagonal
sum-norm projector were derived explicitly above.

The statement does not identify the result with the
unrestricted sum of the one-sided traces, the original
positive-source trace, or the full zero-quotient trace.
A further genuine nuclear comparison can establish such
an identification if its domains and common trace pairing
are proved; density of the smooth ideal alone cannot.
No arithmetic assertion that an exterior divisor is absent,
and no additional motivic or weight-theoretic positivity,
is inferred from the finite-length divisor or this Hilbert
trace calculation.

\section{A polar-retaining completion and the actual positive source column}
\label{sec:polar-source-completion}

The narrow inherited intersection in MA-MO-007 loses both source
polar modes. We now retain those modes by completing actual source
sections together with the inherited diagonal norm. This produces
a bounded representation and an exact quotient map. The original
smoothed positive source column extends to a trace-class map into
the completed common space. The target of that map is kept explicit:
it is not automatically an endomorphism of the completed positive
source image.

Keep the rational place set \(S=\{\infty\}\), the centered theta
map (MA-MO.83), and the original smooth common-sum space
\(\mathscr S_{\rm sum}\) in (MA-MO.78). Put
\[
 \begin{aligned}
 H&=H_{\Sigma,a,s},&0<a&<1/2,&s&\ge0,\\
 I&=\Theta\mathscr U_0,&J&=\overline I^{\,H},
       &K&=K_{a,s,\phi},\\
 \mathscr H_+&=i_+\mathscr S_{\rm even}(\mathbb R),
       &\mathscr H_-&=i_-\mathcal E.
 \end{aligned}
 \tag{MA-MO.250}
\]
Thus \(i_+f=(\Theta f,R\Theta\mathfrak Ff)\),
\(i_-u=(u,u)\), and \(\mathscr U_0\) imposes
\(f(0)=\int f=0\). Let
\[
 \begin{aligned}
 f_0(t)&=(1-2\pi t^2)e^{-\pi t^2},\\
 f_1(t)&=2\pi t^2e^{-\pi t^2},&s_j&=i_+f_j.
 \end{aligned}
 \tag{MA-MO.251}
\]
These are actual even Schwartz sources, with
\((f_j(0),\int f_j)=e_j\). No bare polar exponential is
being declared an element of a narrow weighted \(L^2\) space.

\paragraph{MA-MO-016 (equivariant polar completion and actual source-column
transport).}
The map
\[
 \begin{aligned}
 \mathfrak s(u,c)&=i_-u+c_0s_0+c_1s_1,\\
 \mathcal W&=H\oplus\mathbb C^2,&
 \mathcal W_+&=J\oplus\mathbb C^2,&
 \mathcal W_-&=H\oplus\{0\}
 \end{aligned}
 \tag{MA-MO.252}
\]
identifies \(\mathcal W\) with a Hilbert completion of
\(\mathscr S_{\rm sum}\), and \(\mathcal W_+\) and
\(\mathcal W_-\) with its closed positive and negative source
image completions. The embedding \(\mathfrak s\) into pairs
of locally square-integrable functions is continuous and injective.
The diagonal norm is precisely the inherited norm \(H\).

Simultaneous translations of the actual pair functions extend
to a strongly continuous representation on \(\mathcal W\).
Both \(\mathcal W_+\) and \(\mathcal W_-\) are invariant.
There are equivariant Hilbert exact sequences
\[
 \begin{gathered}
 0\longrightarrow J\longrightarrow\mathcal W_+
       \xrightarrow{\,\operatorname{pol}\,}\mathbb C^2
       \longrightarrow0,\\
 0\longrightarrow\mathcal W_+\longrightarrow\mathcal W
       \xrightarrow{\,\mathfrak q\,}H/J\longrightarrow0,
       \qquad \mathfrak q(u,c)=u+J.
 \end{gathered}
 \tag{MA-MO.253}
\]
The last quotient identification is isometric in the displayed
direct-sum norm, and its representation is the actual quotient
translation representation. Also
\(\mathcal W_+\cap\mathcal W_-=J\) and
\(\mathcal W_++\mathcal W_-=\mathcal W\).
Changing the two actual source sections gives a boundedly
isomorphic completion, by an isomorphism equal to the identity
on the original common-sum space.

Assume either analytic regime of MA-MO-010, with its stated test
class, and write \(D=P_+-P_-\) for the original source matrix.
Then the actual smooth map \(C_gD|_{\mathscr H_+}\) extends
uniquely to
\[
       \mathcal B_g:\mathcal W_+\longrightarrow\mathcal W,
                      \qquad \mathcal B_g\in\mathcal S_1.
 \tag{MA-MO.254}
\]
Its restriction to \(J\) takes values in \(\mathcal W_+\).
The possible outgoing quotient defect has rank at most two.
More precisely, define actual smooth functions
\[
       r_j=(1-\phi)(s_j)_0+\phi(s_j)_1\in\mathcal E.
 \tag{MA-MO.255}
\]
Then
\[
       \mathfrak q\mathcal B_g(y,c)
             =-\sum_{j=0}^1c_j\,[C_gr_j]\quad\hbox{in }H/J.
 \tag{MA-MO.256}
\]
In particular \(\mathcal B_g\) preserves \(\mathcal W_+\)
if and only if both \(C_gr_j\) belong to \(J\). The theorem
does not assert that condition. Nor does it extend the entire
original matrix to \(\mathcal W\), assign an endomorphism
trace to (MA-MO.254), or identify a quotient trace with a local
explicit formula.

\begin{proof}
\emph{Actual sections and the exact algebraic decomposition.}
The Gaussian integral and its derivative give
\(\int e^{-\pi t^2}dt=1\) and
\(\int2\pi t^2e^{-\pi t^2}dt=1\), proving the stated
moment normalization in (MA-MO.251). In fact
\(\mathfrak Ff_0=f_1\) and \(\mathfrak Ff_1=f_0\),
though only the moment normalization is needed below.
Poisson summation (MA-MO.84) gives
\[
 (i_+f)_0-(i_+f)_1
      =-f(0)e^{x/2}+\left(\int f\right)e^{-x/2}.
 \tag{MA-MO.257}
\]
Use the two coefficients \((f(0),\int f)\) as the polar
coordinates. The corresponding two exponentials are linearly
independent in distributions. Every pair in
\(\mathscr S_{\rm sum}\) has unique such coordinates;
subtracting the matching combination of \(s_0,s_1\)
makes its components equal. The equal component lies in
\(\mathcal S_>\cap\mathcal S_<=\mathcal E\), as proved
in MA-MO-007. Conversely every pair in (MA-MO.252) with
\(u\in\mathcal E\) belongs to the original common-sum
space. Thus
\[
 \mathscr S_{\rm sum}=\mathfrak s(\mathcal E\oplus\mathbb C^2),
 \qquad
 \mathscr H_+=\mathfrak s(I\oplus\mathbb C^2),
 \qquad
 \mathscr H_-=\mathfrak s(\mathcal E\oplus\{0\}).
 \tag{MA-MO.258}
\]
For the middle equality, write any even Schwartz source as
\(f=h+f(0)f_0+(\int f)f_1\), with \(h\in\mathscr U_0\).
Then \(i_+h=i_-\Theta h\). This also proves the reverse
inclusion and uses no hypothesized theta surjectivity beyond
the definition of \(I\).

On each compact logarithmic interval the weight of \(H\)
is bounded below by a positive constant. Thus the diagonal
\(H\) component maps continuously to local \(L^2\), and
the fixed sections do also. Injectivity of \(\mathfrak s\)
follows first from (MA-MO.257), which forces both coefficients
to vanish, and then from its diagonal component. This proves
the Hilbert realization and the asserted closures in
(MA-MO.252). If other normalized actual source sections are
chosen, their differences from \(s_j\) are diagonal vectors
in \(I\); the coordinate change is
\((u,c)\mapsto(u+Tc,c)\), with \(T:\mathbb C^2\to I\)
bounded. Its inverse has the opposite finite column.
Unnormalized independent moment sections differ additionally
by an invertible two-by-two matrix. This proves the claimed
section independence up to the specified bounded isomorphism.

\emph{Translations and their actual theta cocycle.}
Set
\[
 \begin{aligned}
 U_\tau u(x)&=u(x-\tau),&
 \rho_\tau f(t)&=e^{-\tau/2}f(e^{-\tau}t),\\
 d_0(\tau)&=e^{-\tau/2},&d_1(\tau)&=e^{\tau/2},\\
 k_j(\tau)&=\Theta\bigl(\rho_\tau f_j-d_j(\tau)f_j\bigr).
 \end{aligned}
 \tag{MA-MO.259}
\]
Direct substitution in the theta sum gives
\(\Theta\rho_\tau=U_\tau\Theta\), and Fourier scaling
gives \(\mathfrak F\rho_\tau=\rho_{-\tau}\mathfrak F\).
Hence \(i_+\rho_\tau=U_\tau i_+\), with simultaneous
translation on the pair. The two source moments of
\(\rho_\tau f\) scale by \(d_0(\tau),d_1(\tau)\).
Consequently every \(k_j(\tau)\) lies in \(I\), and the
actual coordinate formula is
\[
 \pi_\tau(u,c)=
 \left(U_\tau u+\sum_{j=0}^1c_jk_j(\tau),
                 (d_0(\tau)c_0,d_1(\tau)c_1)\right).
 \tag{MA-MO.260}
\]
This is a representation because it agrees with the actual
translations on a dense source domain; alternatively its
cocycle identity follows directly from (MA-MO.259).

For completeness the required continuity is in the declared
norms. The weight \(w\) of \(H\) satisfies
\[
       \frac{w(x+\tau)}{w(x)}
       \le e^{2a|\tau|}2^{s/2}(1+\tau^2)^{s/2}.
 \tag{MA-MO.261}
\]
Thus \(U_\tau\) is bounded, locally uniformly in \(\tau\),
and is strongly continuous by approximation with compact
smooth functions. The map \(\Theta:\mathscr U_0\to\mathcal E\)
is continuous: at the right end, estimates for the sum and
each logarithmic derivative follow from finitely many
Schwartz seminorms of the source; at the left end use
\(\Theta f=R\Theta\mathfrak Ff\) and the same estimates
for \(\mathfrak Ff\). Arbitrarily large additive decay
powers give every prescribed exponential logarithmic weight.
Since dilation is smooth on Schwartz space, \(k_j\) is
continuous in \(\mathcal E\), hence in \(H\). Formula
(MA-MO.260) therefore defines a bounded strongly continuous
representation. It preserves \(J\) by the density of its
translation-invariant subspace \(I\), and the cocycle lies
in \(J\). This proves invariance of both completed images.

The kernel and surjectivity assertions in (MA-MO.253) now
follow from the direct-sum coordinates. In the second quotient
the \(\mathbb C^2\) component may be removed at no remaining
cost, and the infimum over the first component is precisely
the quotient norm of \(H/J\). The cocycle dies in that
quotient, so its action is \(u+J\mapsto U_\tau u+J\).
The polar quotient in the first sequence has the two
characters \(d_0,d_1\). None of these statements asserts
that the displayed source-section splitting is equivariant.

\emph{The exact negative projection and finite polar columns.}
The original definition of \(p_-\) gives, in these coordinates,
\[
       p_-\mathfrak s(u,c)=u+\sum_jc_jr_j.
 \tag{MA-MO.262}
\]
Each \(r_j\) belongs to \(\mathcal E\): sufficiently far
left only \((s_j)_1\) remains, and sufficiently far right
only \((s_j)_0\) remains, with the rapid weighted tails
specified by \(\mathcal S_<\) and \(\mathcal S_>\).
The intervening region is compact and smooth. Consequently
\(p_-:\mathcal W\to H\) and
\(P_-=i_-p_-:\mathcal W\to\mathcal W_-\) extend
boundedly, and \(P_-\) remains a projection since
\(p_-i_-=1_H\). This does not assert that \(p_+\) has
a comparable bounded extension.

Compact convolution \(C_g=\int g(\tau)\pi_\tau\,d\tau\)
is bounded on \(\mathcal W\), by the local bounds just
proved. On source vectors it equals the actual simultaneous
convolution, and
\[
 C_gi_+f=i_+\mathcal D_gf,\qquad
 \mathcal D_gf(t)=\int g(\tau)e^{-\tau/2}f(e^{-\tau}t)\,d\tau.
 \tag{MA-MO.263}
\]
This is an even Schwartz function. Thus convolution preserves
\(\mathscr H_+\), \(\mathscr H_-\), and their completions.
Both \(P_+s_j\) and \(P_-s_j\) are genuine original source
outputs: the first is \(i_+p_+s_j\) with even Schwartz
\(p_+s_j\), and the second is \(i_-r_j\).
Therefore \(C_gDs_j\in\mathscr S_{\rm sum}\subset\mathcal W\)
without any narrow multiplier continuation.

\emph{The actual trace-class map on the ideal column.}
The bounded functionals of MA-MO-011 give
\(\Lambda=(\lambda_0,\lambda_1):J\to\mathbb C^2\),
with \(K=\ker\Lambda\). Choose fixed
\(h_0,h_1\in\mathscr U_0\), put \(y_j=\Theta h_j\),
and normalize \(\Lambda y_j=e_j\), as in (MA-MO.147).
The bounded map
\[
          R_\phi y=y-\sum_j\lambda_j(y)y_j:J\to K
 \tag{MA-MO.264}
\]
preserves the actual smooth domains:
\(R_\phi I=\Theta\mathscr U_{\rm pf}\).
For \(f\in\mathscr U_0\), the defining source equations
(MA-MO.85) give \(D i_+f=i_+Lf\), exactly as in
MA-MO-008. Hence the two fixed vectors
\(b_j=C_gD i_-y_j=i_+\mathcal D_gLh_j\) belong to
\(\mathscr H_+\subset\mathcal W_+\).
Let \(B_{\phi,g}:K\to J\) be the actual map constructed
in MA-MO-010. The formula
\[
 \mathcal B_g(y,c)=
   i_-B_{\phi,g}R_\phi y
       +\sum_j\lambda_j(y)b_j
       +\sum_jc_j C_gDs_j
 \tag{MA-MO.265}
\]
now defines a trace-class map \(\mathcal W_+\to\mathcal W\).
Indeed the first term is a trace-class map composed with
bounded maps; the other four columns have finite rank.
It agrees with the original operator on
\(I\oplus\mathbb C^2\) by the exact smooth decomposition
above, so density proves uniqueness. This also proves
independence from the auxiliary cutoff-moment correctors.
The trace-norm continuity in smooth fixed-support test spaces,
or boundedness for common Gevrey test bounds, follows from
MA-MO-010 for the first term and ordinary Schwartz continuity
of the finitely many actual source columns for the others.

The first two terms in (MA-MO.265) lie in \(\mathcal W_+\).
On a polar section \(\mathfrak q C_gP_+s_j=0\), since
\(P_+s_j\in\mathscr H_+\) and convolution preserves that
space. The negative term is
\(\mathfrak q C_gP_-s_j=[C_gr_j]\). This proves the
exact formula (MA-MO.256) and its stated if-and-only-if
criterion. It also explains why an endomorphism target
cannot be obtained merely by completing the two additional
source vectors.
\end{proof}

\paragraph{Source scope and the remaining matrix problem.}
The completion and equivariant exact sequences use only Poisson
summation and the declared weights; their construction has no
zero-location premise. The analytic regimes enter precisely in
the bounded moment functionals and the trace-class ideal-column
estimate inherited from MA-MO-010--011. In particular the result
does not assert unconditional smooth-test transport at every
narrow width. It changes the completion of the actual common-sum
source by retaining its polar coordinates, while retaining the
exact inherited norm on its diagonal part.

The original smooth common-space matrix is still defined on
\(\mathfrak s(\mathcal E\oplus\mathbb C^2)\). The theorem
does not provide its bounded action on the complementary
off-ideal diagonal directions in \(H\). That is the remaining
negative-column problem. A trace of a completed full matrix,
a relative trace, or a quotient trace would require those
additional maps and the corresponding trace-class comparisons.
The finite outgoing formula (MA-MO.256) is a separate condition
even on the positive source column.

The exact primary locators are Meyer (2005),
\texttt{sec:summation}, \texttt{lem:Hilm\_p\_embed}, the
definitions of \(p_-,P_-\), \texttt{eq:Poisson\_summation},
and \texttt{lem:Poisson\_sumpm}; author-manuscript
pp.~39--41, supplied \texttt{main.tex} lines~3808--4043.
Their complete displayed maps and Poisson-sum argument were
rechecked here. The centering and original operator signs are
those independently derived in MA-MO-007--008.
The two polar characters are genuine finite-dimensional
representation data. In the cote~1 framework, scans~079--080
and~111--112 through ALPOGE\_BRIDGES BR-TR-01/02, the
equivariant exact sequences belong to the additive
representation bookkeeping; a trace or determinant still
requires a defined endomorphism and its analytic trace ideal.
No such trace is assigned to the two-space map (MA-MO.254).

\section{An actual nuclear return to the local source and its trace}
\label{sec:local-source-return}

The ideal projector in MA-MO-012 cannot be omitted by formal
cyclicity. There is, however, a different possible comparison:
return the completed ideal operator to the actual smooth local
source before taking its trace. We prove that comparison here.
The resulting trace is a local source trace. It is not the trace
of the full narrow zero quotient or of an unconstructed ambient
common-space matrix.

Fix a nontrivial primitive finite-order character \(\eta\), the
structural finite place set \(S\), the actual norm splitting and
local source Fourier maps of MA-CC/MA-CQ/MA-MO-012. Write
\(\mathcal T_\eta\) for the corresponding compact-character
sector of Meyer's \(\mathcal T(C_S)\), with its actual local
coinvariant topology. Use the fixed probability Haar measure
on the compact norm-one factor and center all norm coordinates,
as in MA-CC/MA-CQ. Its global
summation embedding is
\[
 \vartheta_\eta:\mathcal T_\eta\longrightarrow
 I_\eta^c\subset\mathcal E,\qquad
 J_{\eta,a,s}=\overline{I_\eta^c}^{\,H_{\Sigma,a,s}}.
 \tag{MA-MO.280}
\]
This is the original source embedding, not a chosen radial
representative for an ideal class. On the fixed compact sector
the source is a nuclear Fr\'echet space. In the centered norm
coordinate its topology is measured by the Schwartz seminorms
of \(e^{cx}F\) and of the actual dual Fourier vector
\(e^{cx}\mathfrak S_\eta F\), for every \(c>-1/2\).
The dual vector belongs to the conjugate compact sector; finite
Fourier phases and conductor normalizations are retained in
\(\mathfrak S_\eta\).

Let \(R F(x)=F(-x)\). The equivariant local Fourier operator
\(\mathfrak U_\eta=\mathfrak S_{\bar\eta}R\), with sectors
chosen so that it acts on \(\mathcal T_\eta\), has scalar
Mellin multiplier \(m_\eta(z)\). Equivalently this notation
defines \(m_\eta\) by the actual local map; no conductor phase
is removed. The elementary local functional equations give:
\[
 \begin{gathered}
 m_\eta\text{ is holomorphic for }\Re z>-1/2,\\
 m_\eta\text{ and all its derivatives have polynomial growth}\\
 \text{on every fixed closed vertical substrip there},\\
 m_\eta(b+it)^{-1}\text{ and all its derivatives have polynomial growth}\\
 \text{for }-1/2<b<1/2.
 \end{gathered}
 \tag{MA-MO.281}
\]
Here and below polynomial exponents may depend on the fixed
strip and derivative order. These assertions concern the local
multiplier only. For the rational even archimedean sector, for
example, with \(\Gamma_{\mathbb R}(w)=\pi^{-w/2}\Gamma(w/2)\),
\[
       m(z)=\frac{\Gamma_{\mathbb R}(1/2+z)}
                       {\Gamma_{\mathbb R}(1/2-z)}.
 \tag{MA-MO.282}
\]
At a real place the parity shifts both gamma arguments by zero
or one; a complex finite-order sector has its corresponding
complex gamma quotient. Up to the nonzero exponential monomial
from the different and the chosen additive character, an
unramified finite factor is a ratio
\((1-\omega q^{z-1/2})/(1-\omega' q^{-z-1/2})\), with unit
modulus constants in the numerator and denominator; ramified
factors are nonzero exponential monomials times unit constants.
Thus their possible poles lie on \(\Re z=-1/2\), and their
possible positive-strip zeros are on \(\Re z=1/2\).
Stirling's formula proves the gamma bounds, including derivatives;
the finite factors and their derivatives are bounded on each
fixed admissible input line. This verifies (MA-MO.281) without
using any global zero location.

Let \(\phi=1\) sufficiently far left and \(\phi=0\)
sufficiently far right, as before, and put
\[
 \begin{aligned}
 \phi^\vee(x)&=\phi(-x),&
 \psi(x)&=1-\phi(x)-\phi(-x),\\
 T_{\eta,\phi}&=M_\phi-
       \mathfrak U_\eta M_\phi\mathfrak U_\eta^{-1},&
 \mathcal L_\eta&=-T_{\eta,\phi}
        +\mathfrak U_\eta M_\psi\mathfrak U_\eta^{-1}.
 \end{aligned}
 \tag{MA-MO.283}
\]
The last expression is exactly
\(1-M_\phi-\mathfrak S_{\bar\eta}M_\phi\mathfrak S_\eta\)
on the original source, by inversion of the cutoff. In the
rational even additive coordinates it is the \(L\) of
MA-MO-008. It equals \(-T_{\eta,\phi}\) for a symmetric
cutoff; otherwise the compact residual \(\psi\) is retained.
The actual two Fourier maps conjugate \(\mathcal L_\eta\)
to \(\mathcal L_{\bar\eta}\). No printed Fourier
anticommutation identity is needed.

\paragraph{MA-MO-017 (nuclear local-source return and evaluated ideal trace).}
Assume either analytic regime of MA-MO-012: generic widths with
compact Gevrey \(g,\phi'\), or the stated fixed horizontal zero
gap with arbitrary compact smooth \(g,\phi'\). For every
\(s\ge0\) there is a continuous nuclear map
\[
 \begin{aligned}
 A_{\eta,g}:J_{\eta,a,s}&\longrightarrow\mathcal T_\eta,\\
 A_{\eta,g}\vartheta_\eta F&=C_g\mathcal L_\eta F,\\
 \vartheta_\eta A_{\eta,g}&=B_{\eta,g},
 \end{aligned}
 \tag{MA-MO.284}
\]
where \(B_{\eta,g}\) is the actual completed-ideal endomorphism
of MA-MO-012. Thus its trace has the exact evaluation
\[
 \boxed{\displaystyle
 \operatorname{Tr}_{J_{\eta,a,s}}B_{\eta,g}
   =\frac1{2\pi}\int_{\mathbb R}
          \mathcal Mg(it)\frac{m_\eta'(it)}{m_\eta(it)}\,dt
       +g(0)\int_{\mathbb R}\psi(x)\,dx .}
 \tag{MA-MO.285}
\]
The integral is absolutely convergent. In particular the extra
cutoff term vanishes for the requested \(g(0)=0\) class.
The answer is independent of \(a,s\) within the explicit
regimes. In the notation of the actual local Weil trace,
the first term is \(-W^{\rm loc}_{\eta,S}(g)\), with
all finite-place and different/conductor factors in \(m_\eta\).

The projected trace in (MA-MO.190) therefore equals the right
side of (MA-MO.285). This evaluates that projected trace by a
genuine nuclear factorization; it does not assert
\(\Pi_J=1\) or equality with the unrestricted ambient trace.

There is also a rational principal-sector statement. Let
\(\mathcal W_+\) be the polar-retaining source completion of
MA-MO-016 and \(\jmath_+:\mathscr S_{\rm even}(\mathbb R)
\to\mathcal W_+\) its original theta-pair embedding. The
actual source endomorphism \(C_gL\) has a continuous nuclear
return
\[
 A_{+,g}:\mathcal W_+\longrightarrow
                    \mathscr S_{\rm even}(\mathbb R),\qquad
 A_{+,g}\jmath_+f=\mathcal D_gLf.
 \tag{MA-MO.286}
\]
Its resulting endomorphism \(\jmath_+A_{+,g}\) has trace
(MA-MO.285) with (MA-MO.282). This is distinct from the full
source column \(C_g(P_+-P_-)|_{\mathcal W_+}\), whose two
outgoing polar columns are explicitly retained in MA-MO-016.

\begin{proof}
\emph{An exact kernel between different exponential weights.}
Write \(V_b=\mathcal F_+M_{e^{bx}}\), with unitary Fourier
normalization, and set
\[
 \begin{gathered}
 z=c+it,\quad w=b+iu,\quad G(z)=\mathcal Mg(z),\\
 P(z-w)=\int e^{(z-w)x}\phi'(x)\,dx,\qquad
 \Psi(z-w)=\int e^{(z-w)x}\psi(x)\,dx.
 \end{gathered}
 \tag{MA-MO.287}
\]
For every \(c>-1/2\) and fixed \(-1/2<b<1/2\), the
kernel of \(V_c C_g\mathcal L_\eta V_b^{-1}\) on its
original smooth source is
\[
 \ell_{c,b}(t,u)=\frac{G(z)}{2\pi}
 \left[
 -\frac{P(z-w)}{z-w}
              \left(\frac{m_\eta(z)}{m_\eta(w)}-1\right)
       +\frac{m_\eta(z)}{m_\eta(w)}\Psi(z-w)
 \right].
 \tag{MA-MO.288}
\]
At \(z=w\) the divided difference has its removable value.

Here is a direct derivation, so that no inverse local multiplier
is moved across one of its positive zeros. If \(h\) and its
two exponential weights are Schwartz, the kernel of
\([M_\phi,C_h]\), between these weighted lines, has double
Fourier transform
\[
 -\frac{P(z-w)}{2\pi(z-w)}
              \bigl(\mathcal Mh(w)-\mathcal Mh(z)\bigr).
 \tag{MA-MO.289}
\]
Indeed, putting \(r=x-y\), integration in \(y\) uses the
compactly supported difference \(\phi(y+r)-\phi(y)\).
An integration by parts gives
\[
 \int e^{(z-w)y}\{\phi(y+r)-\phi(y)\}\,dy
       =-\frac{e^{-(z-w)r}-1}{z-w}P(z-w).
 \tag{MA-MO.290}
\]
This identity is valid for either sign of \(c-b\); it is not
an unjustified transform of the nonintegrable step itself.
At equality use its continuous limit. The decay hypotheses on
\(h\) justify the remaining integration and differentiation.

By (MA-MO.281), the Mellin transform of
\(\mathfrak U_\eta g\) is \(m_\eta(z)G(z)\), rapidly
decreasing on every fixed positive finite-\(L\) half-plane
line. Contour shifting, or the actual local Fourier estimate,
places this function in all the required weighted Schwartz
spaces. Now use the exact operator identity
\[
 C_gT_{\eta,\phi}
     =-[M_\phi,C_g]
       +[M_\phi,C_{\mathfrak U_\eta g}]
                    \mathfrak U_\eta^{-1}.
 \tag{MA-MO.291}
\]
Substitution of (MA-MO.289) cancels the two terms containing
\(G(w)\). It gives the negative of the first summand in
(MA-MO.288). The compact multiplier term in (MA-MO.283)
has kernel \(G(z)m_\eta(z)\Psi(z-w)/(2\pi m_\eta(w))\).
This proves (MA-MO.288), with the inverse evaluated only on
the fixed input line \(b\).

\emph{Composition with the genuine global inverse.}
Use the actual finite Euler multiplier of MA-MO-012:
\[
 Z_\eta(z)=E_{\eta,S}(1/2+z)L_\eta(1/2+z),
       \qquad j_{\eta,b}(u)=Z_\eta(b+iu)^{-1}.
 \tag{MA-MO.292}
\]
For every smooth ideal vector \(y=\vartheta_\eta F\),
the unconditional division result (MA-MO.173--178) identifies
its original inverse with \(V_bF=j_{\eta,b}V_by\).
It also supplies the corresponding identity for the actual
dual Fourier vector and \(Ry\). Thus the kernel giving the
first source component of the proposed return is exactly
\[
                  \ell_{c,b}(t,u)j_{\eta,b}(u).
 \tag{MA-MO.293}
\]
There is no assertion that \(j_{\eta,b}\) separately acts
on every Schwartz vector.

For clarity, (MA-MO.293) is a Schwartz function of \((t,u)\)
for every \(c>-1/2\). In the gap regime all derivatives of
\(j_{\eta,b}\) are polynomially bounded; (MA-MO.281), the
rapid transforms \(G\), \(P\), \(\Psi\), and the divided
difference prove the assertion. In the generic regime their
derivatives are bounded by \(\exp(C\log^2(2+|u|))\).
The compact residual \(\psi\) has the same Gevrey order as
\(\phi'\), since its derivatives are sums of derivatives of
that fixed step and its reflection. Compact Gevrey data give decay
\[
 C_N(1+|t|+|t-u|)^{N}
       \exp\{-d(|t|^{1/q}+|t-u|^{1/q})\}
 \tag{MA-MO.294}
\]
for suitable constants in each derivative estimate. The
exponential defeats that logarithmic-square growth since
\(|u|\le |t|+|t-u|\). The divided difference creates no
singularity: write
\((m(z)-m(w))/(z-w)=\int_0^1m'(w+v(z-w))\,dv\).
The segment stays in a fixed positive finite-\(L\) half-plane
strip. All remaining factors have the polynomial bounds in
(MA-MO.281). These estimates hold in every derivative order
and locally uniformly in \(c\).

Every such kernel maps \(L^2(du)\) continuously into
\(\mathscr S(dt)\), by Cauchy--Schwarz applied to its
weighted derivatives. The input map \(y\mapsto V_by\)
is bounded from \(H_{\Sigma,a,s}\) to \(L^2\), for
\(b=a\) or \(-a\), since \(s\ge0\). Consequently all
weighted source seminorms of the first output are bounded by
the declared Hilbert norm of \(y\).

Apply exactly the same construction in the dual compact
character sector to \(Ry\), with \(g^\vee(x)=g(-x)\).
The actual source Fourier identities and (MA-MO.283) give
\[
 \mathfrak S_\eta C_g\mathcal L_\eta F
       =C_{g^\vee}\mathcal L_{\bar\eta}
                            \mathfrak S_\eta F.
 \tag{MA-MO.295}
\]
Thus the second family of kernels controls every seminorm of
the dual Fourier output as well. This step retains the dual
character and the actual finite Fourier transform.

Take a sequence from the smooth ideal tending to \(y\in J\).
The two outputs converge in their weighted Schwartz families.
At the central local \(L^2\) weight the source Fourier map is
unitary, so their limits remain Fourier transforms of each
other. The source topology described above therefore puts the
limit in \(\mathcal T_\eta\), continuously. It is independent
of the approximating sequence. This constructs \(A_{\eta,g}\)
and proves its first identity on the complete smooth source.

\emph{Nuclearity into the actual source topology.}
We spell out the functional-analytic point used here. The fixed
compact-character source is nuclear Fr\'echet, with a cofinal
increasing family of Hilbertian norms \(q_n\), each dominating
the injective central local \(L^2\) norm, whose successive
completion maps may be chosen trace class. Thus none of the
finite-dimensional source images below has a seminorm-null
vector. A
continuous map from a Hilbert space into this source is trace
class in each \(q_n\): factor it through a stronger seminorm
and the corresponding trace-class completion map. Choose
finite-rank orthogonal domain projections \(P_N\) increasing
strongly to the identity. For each \(n\), the \(q_n\)
trace norm of \(A(1-P_N)\) tends to zero. Choose an increasing
subsequence \(N_k\) making these tails at most \(2^{-k}\)
for every \(n\le k\). Decompose each finite-rank block
\(A(P_{N_{k+1}}-P_{N_k})\) by its singular-value expansion
in \(q_k\). Its vectors lie in the actual source, its
coefficient functionals have norm one, and the sum of its
singular values is at most \(2^{-k}\). For each fixed
\(q_n\), the resulting normalized source vectors are bounded
once \(k\ge n\); the finitely many earlier blocks are bounded
as well. Together with the first finite block this is a
summable nuclear series with a bounded set of actual source
vectors. This proves the required nuclearity, rather than
only separate one-sided Hilbert trace classes.

The source's nuclear Fr\'echet property and Hilbertian
description are the standard fixed-sector consequence of
Meyer's split extension (24), Proposition 4.4 and the
weighted Schwartz topology; the same bounded-map argument is
used in Proposition 4.6. Alternatively, in the rational
archimedean sector one may use the even Hermite basis of
\(\mathscr S(\mathbb R)\) directly. Its coefficient norms
decay faster than every power by the source seminorm estimates
just proved and give the same nuclear series.

\emph{Actual factorization and the trace.}
The global source embedding (MA-MO.280) is continuous into
\(J_{\eta,a,s}\): its two Poisson-related tails give every
exponential seminorm for a nontrivial character, as in
MA-MO-012. The first source identity in (MA-MO.284) and the
literal common-sum formula imply
\(\vartheta_\eta A_{\eta,g}=B_{\eta,g}\) on a dense smooth
ideal. Both sides are continuous, so the identity holds on
the Hilbert completion. The reverse composition is the
original local operator \(C_g\mathcal L_\eta\) on
\(\mathcal T_\eta\). Nuclear trace cyclicity for these
actual maps therefore gives
\[
 \operatorname{Tr}_J(\vartheta_\eta A_{\eta,g})
       =\operatorname{Tr}_{\mathcal T_\eta}
                    (A_{\eta,g}\vartheta_\eta)
       =\operatorname{Tr}_{\mathcal T_\eta}
                              C_g\mathcal L_\eta.
 \tag{MA-MO.296}
\]
Concretely both traces are the same absolutely summable
pairing series from the preceding nuclear expansion. The
smooth source has the approximation property from its split
weighted-Schwartz model, so this is its canonical nuclear
trace. Meyer's Proposition 4.6 and Corollary 4.7 identify it
with the local central \(L^2\) trace; the compact residual
term has the same comparison by the elementary kernel lemma.
Equivalently, repeating (MA-MO.288--295) with the global
inverse factor replaced by one proves the local central
\(L^2\)-to-source nuclear factorization directly. This
also avoids relying on the printed Fourier-sign shorthand
in the source proof of Proposition 4.6.

Its normalization may also be read directly from (MA-MO.288)
without the global inverse, at \(c=b=0\). That kernel is
Schwartz and hence trace class. Since \(P(0)=-1\), its
diagonal is
\[
 \ell_{0,0}(t,t)=\frac{G(it)}{2\pi}
       \left(\frac{m_\eta'(it)}{m_\eta(it)}
                                      +\int\psi\right).
 \tag{MA-MO.297}
\]
The Schwartz-kernel trace equals its diagonal integral, and
\((2\pi)^{-1}\int G(it)\,dt=g(0)\). Formula (MA-MO.285)
follows, including its sign and compact-residual term.
The logarithmic derivative of the local gamma/finite factors
has at most polynomial growth, so the displayed integral is
absolutely convergent. This is also the local Weil trace of
Meyer's equation (20), with the sign from (MA-MO.283).

\emph{The rational polar source.}
On \(I=\Theta\mathscr U_0\), the same construction uses
the actual inverse identities in MA-MO-010. The paired
Fourier estimates put its output in the actual even additive
Schwartz space. This can be verified without relying on a
radial description alone: taking \(c=-1/2+\varepsilon\)
gives local integrability at additive zero, all positive
\(c\) give rapid decay at infinity, and the paired Fourier
estimates give integrability of all Fourier moments. Fourier
inversion then gives smooth bounded derivatives at zero and
the ordinary Schwartz estimates. This verifies the actual
additive target.

For \((y,c_0,c_1)\in J\oplus\mathbb C^2=\mathcal W_+\),
using the two actual source sections of MA-MO-016, define
\[
 A_{+,g}(y,c_0,c_1)
       =A_{0,g}y+\sum_{j=0}^1c_j\mathcal D_gLf_j.
 \tag{MA-MO.298}
\]
The extra columns are actual Schwartz vectors and have finite
rank. The result is continuous and nuclear, and on the dense
original positive source it is exactly \(\mathcal D_gL\).
The theta-pair inclusion into \(\mathcal W_+\) is continuous
by the seminorm estimates in MA-MO-016. Applying the same
nuclear pairing argument proves (MA-MO.286) and its trace.
On the polar sections \(i_+f_j\), however,
\[
 (P_+-P_-)i_+f_j
       =i_+Lf_j+(i_+f_j-P_-i_+f_j).
 \tag{MA-MO.299}
\]
The latter finite columns are precisely why this source
endomorphism is not silently identified with the full raw
matrix defect. This completes the proof.
\end{proof}

\paragraph{Primary locators and the remaining arithmetic gate.}
The supplied Meyer (2005) source is used at Lemma 4.3
(author-manuscript statement p.~25, proof p.~26),
Proposition 4.4 (pp.~26--27),
the source topology and split extension (24) (pp.~27--28),
Proposition 4.6 and Corollary 4.7 (p.~28), and the
local trace equation (20) (p.~22; (19) is its normalized-character
version). Their exact TeX labels are
\texttt{lem:mother\_of\_nuclearity},
\texttt{pro:Fourier\_estimate},
\texttt{eq:Twist\_convariants},
\texttt{pro:T\_on\_Twist}, \texttt{cor:T\_on\_Twist}, and
\texttt{eq:local\_trace\_formula\_psi}.
The newly needed assertion is the return after the possibly
unbounded global inverse; it is (MA-MO.288--295), not a direct
application of the local theorem to that inverse.

The elementary local factors in (MA-MO.281--282) are Tate
(1950), Theorem 2.4.1, thesis p.~(2.11), with the explicit
tables in \S2.5: real p.~(2.14), complex p.~(2.16), and
nonarchimedean p.~(2.19). These are PDF pages 18, 21, 23,
and 26 of the original-thesis scan. Those tables are used for
the gamma/Euler forms and their analytic domains; the actual
Fourier convention, different, conductor and phases here are
fixed by the source definition of \(m_\eta\), not transferred
from an inconsistent Fourier-sign convention.

\smallskip
\noindent Tate, J. (1950). \emph{Fourier analysis in number
fields and Hecke's zeta-functions} [Doctoral dissertation,
Princeton University]. Original-thesis scan,
\url{https://sites.math.rutgers.edu/~alexk/2023S572/Tate1950.pdf}.

No zero-location premise enters the local multiplier or
cross-weight construction. Global reciprocal bounds use the
explicit generic/Gevrey or gap/smooth hypotheses. For rational
and Dirichlet sectors A supplies the gap at every positive
narrow width; using it is recorded as zero-location input.
For the non-induced character, M or its own RH supplies that
gap, while the generic/Gevrey alternative remains unconditional.
The full quotient trace and the FT002 discarded-divisor
condition are unchanged. The present theorem closes the
specific evaluation of the completed positive-ideal source
trace, not the independent full narrow identity.

In the inherited cote~1 framework, scans~079--080 and~111--112
through ALPOGE\_BRIDGES BR-TR-01/02, (MA-MO.284) is an actual
nuclear factorization with compatible pairings. It therefore
justifies a trace comparison and its associated Fredholm
endomorphism determinant. The finite-order character remains
the specified rank-one finite-monodromy local system; no
positive faithful full-jet realization or new arithmetic weight
is inferred from this nuclear comparison.

\section{The full polar-retaining rational matrix under a right zero-free half-plane}
\label{sec:polar-common-matrix}

The positive source-column result MA-MO-016 does not by itself
extend the original matrix on the complementary diagonal source.
Here that second column is constructed under an explicit
zero-location premise. Both original inverses are evaluated on
the positive narrow line, and both polar moments are retained.
In particular a zero on the critical line does not obstruct this
construction when the stipulated premise A is available.

Fix
\[
 \begin{gathered}
 0\le a_*<a<1/2,\qquad s\ge0,\\
 \zeta(\sigma+it)\ne0
          \quad\hbox{whenever }\sigma>1/2+a_*.
 \end{gathered}
 \tag{MA-MO.320}
\]
The inequality refers to zeros, not to the pole at one. It is
an additional arithmetic hypothesis in this section; A gives
it with \(a_*=0\). Keep the rational place set
\(S=\{\infty\}\), the actual cutoff \(\phi\), and the
Gaussian sections of MA-MO-016. The cutoff need not be symmetric.
Put
\[
 \begin{aligned}
 H&=H_{\Sigma,a,s},&H_a&=H_{a,s},&
 \mathcal W&=H\oplus\mathbb C^2,\\
 q(x)&=1-\phi(x),&p_-(x)&=e^{-x/2},&
 D&=P_+-P_- .
 \end{aligned}
 \tag{MA-MO.321}
\]
In this section \(p_-(x)\) denotes the displayed scalar
polar function; the source retraction remains denoted
\(p_-^{\rm src}\) when it is needed. The operator \(P_-\)
is always the original source projection.

\paragraph{MA-MO-018 (actual full common-space trace-class extension).}
For every \(g\in C_c^\infty(\mathbb R)\), the original
smooth operator
\[
       C_gD:\mathscr S_{\rm sum}\longrightarrow
                                     \mathscr S_{\rm sum}
\]
extends uniquely to a trace-class endomorphism
\[
       \mathscr D_g:\mathcal W\longrightarrow\mathcal W.
 \tag{MA-MO.322}
\]
Its restriction to the positive completion is the actual
source-column map of MA-MO-016. On each fixed compact test
support, \(g\mapsto\mathscr D_g\) is continuous into trace
class. This is an extension of the original wide-inverse
source action, with no substitution of an orthogonal projector
or of an unrelated continued inverse.

The trace and Fredholm determinant of (MA-MO.322) are therefore
defined. The theorem does not compute that trace, assert that
\(\mathcal W_+\) is invariant under \(\mathscr D_g\), or
define its action on \(\mathcal W/\mathcal W_+\). The exact
outgoing condition of (MA-MO.256) remains in force. In
particular this is not a zero-location-free proof of a full
narrow identity.

\begin{proof}
\emph{The original inverse on every negative smooth input.}
For \(y\in\mathcal E\), let
\[
 F_y(x)=\sum_{n\ge1}\mu(n)n^{-1/2}y(x+\log n).
 \tag{MA-MO.323}
\]
This is the original wide Euler inverse. The series and its
derivatives converge locally absolutely, by the rapid right
decay of \(y\). On a line \(b>1/2\), it converges in the
weighted Hilbert operator norm and has Mellin transform
\(\mathcal My(z)/\zeta(1/2+z)\), as in MA-MEY-001.

Hypothesis (MA-MO.320) gives a fixed horizontal margin
between every line \(\Re z\ge a\) and every nontrivial
zero. The reciprocal minimum-modulus proof of MA-RS-004
therefore gives polynomial bounds for the reciprocal and
all its derivatives, uniformly on each fixed closed vertical
strip in \(\Re z\ge a\). The inverse is holomorphic there:
the zeta pole at \(z=1/2\) gives a zero of the inverse,
not a pole. Thus contour shifting in the rapidly decreasing
Mellin quotient is legitimate. It proves
\[
 \begin{gathered}
 e^{bx}F_y\in\mathscr S(\mathbb R)\quad\hbox{for every }b\ge a,
 \qquad
 V_aF_y=j_aV_ay,\\
 A(z)=\zeta(1/2+z),\qquad j_a(t)=A(a+it)^{-1},\qquad
 \mathcal MF_y(1/2)=0.
 \end{gathered}
 \tag{MA-MO.324}
\]
Here \(V_a=\mathcal F_+M_{e^{ax}}\) is the unitary
Fourier normalization. The last Mellin integral converges
absolutely: use a line \(a<1/2\) at the left and a line
\(b>1/2\) at the right. Its value is zero because the
holomorphic quotient vanishes at \(z=1/2\). The equality
with the original function (MA-MO.323) follows by starting
on the absolutely convergent wide line before shifting.
This argument is for every \(y\in\mathcal E\), with no
ideal membership assumption.

Write \(J_a\) for the maximal inverse multiplier on
\(H_{a,0}\). Then (MA-MO.324) says precisely
\(F_y=J_ay\) on this actual source core. The same is true
for \(Ry\). No contour is moved to the negative narrow
line in making either inverse identification.

\emph{The two actual source functions and their polar moments.}
Define an even additive function by
\[
 u_y(t)=\frac{q(\log|t|)F_y(\log|t|)}{2\sqrt{|t|}}
        \quad(t\ne0),\qquad u_y(0)=0.
 \tag{MA-MO.325}
\]
It is Schwartz: it vanishes near additive zero, and all
its derivatives decrease rapidly at infinity by the
arbitrarily large right weights in (MA-MO.324).
The exact source formula for \(p_+\) on a diagonal vector
is consequently
\[
        p_+i_-y=h_y:=u_y+\mathfrak F u_{Ry},
             \qquad h_{Ry}=\mathfrak F h_y.
 \tag{MA-MO.326}
\]
This is the original map of Meyer's
\texttt{lem:Hilm\_p\_embed}, specialized to the rational
even sector. It uses the wide inverse (MA-MO.323), now
identified with its positive narrow multiplier value.

Let \(\lambda_0,\lambda_1\) be the bounded functionals
of MA-MO-011 on \(H\). For arbitrary \(y\in\mathcal E\),
not only ideal vectors, the same bilinear Fourier computation
and (MA-MO.324) give
\[
 \lambda_0(y)=\int e^{x/2}\phi(x)F_y(x)\,dx,
                  \qquad\lambda_1(y)=\lambda_0(Ry).
 \tag{MA-MO.327}
\]
The functionals are bounded on \(H\), since the local
weighted step is Schwartz and \(j_a\) has polynomial
growth on this line. The zero total moment in (MA-MO.324)
now proves both source moment identities
\[
       h_y(0)=-\lambda_1(y),\qquad
       \int h_y=-\lambda_0(y).
 \tag{MA-MO.328}
\]
Indeed \(\int u_y=\int e^{x/2}qF_y=-\lambda_0(y)\),
while \(u_y(0)=0\); Fourier evaluation and integration
at zero give the other term. Thus neither source polar
coefficient has been discarded.

Let \(G(z)=\mathcal Mg(z)\). Under the actual additive
dilation operator \(\mathcal D_g\) from (MA-MO.263),
the two moments are multiplied by \(G(-1/2)\) and
\(G(1/2)\), respectively. Hence the polar coordinates
of \(C_gDi_-y\) are
\[
 \nu_g(y)=
       \bigl(-G(-1/2)\lambda_1(y),
             -G(1/2)\lambda_0(y)\bigr).
 \tag{MA-MO.329}
\]
This is already a bounded finite-rank map \(H\to\mathbb C^2\).

\emph{An explicit positive-line off-diagonal kernel.}
Let \(\mathcal H\) be the radial form of the actual even
additive Fourier transform, so that
\[
 \mathcal M\mathcal H F(z)=m(z)\mathcal MF(-z),\qquad
 m(z)=\frac{\Gamma_{\mathbb R}(1/2+z)}
                    {\Gamma_{\mathbb R}(1/2-z)}.
 \tag{MA-MO.330}
\]
Only \(z=a+it\) is used in this step. Put
\[
       P(\xi)=\int e^{\xi x}\phi'(x)\,dx,
                  \qquad z=a+it,\quad w=a+iu.
\]
For \(\Re\xi<0\), integration by parts, with zero
boundary terms because \(q=0\) far left, gives
\(\int e^{\xi x}q(x)\,dx=P(\xi)/\xi\).
It follows directly by Fourier inversion that the source
operator \(C_gZ_a\mathcal H M_qJ_a\) has kernel
\[
       o_{a,g}(t,u)=\frac{G(z)A(z)m(z)j_a(u)}{2\pi}
                          \frac{P(-z-w)}{-z-w}
 \tag{MA-MO.331}
\]
in the \(V_a\) coordinates. Here \(Z_a\) is the maximal
multiplier \(A(a+it)\). On the source core its sequential
domain is valid: \(qF_y\) comes from an additive Schwartz
function, its local Fourier transform has weighted
Schwartz estimates at \(a>-1/2\), and the final zeta
multiplier has polynomial growth. The preceding inversion
can be justified absolutely after inserting the Schwartz
input transform; the separated real part
\(\Re(-z-w)=-2a\) is crucial.

All derivatives of (MA-MO.331) are Schwartz in \((t,u)\).
The factors \(A,m,j_a\) and their derivatives have
polynomial growth on their fixed lines. The two rapid
factors are \(G(a+it)\) and \(P(-2a-i(t+u))\), and
\(|u|\le |t|+|t+u|\). The denominator stays a distance
\(2a\) from zero. Thus (MA-MO.331) extends to a
trace-class operator
\(O_a(g):H_{a,s}\to H_{a,s}\) for every \(s\ge0\),
by the Bessel/Hermite argument of MA-MO.126. This is an
extension of the actual sequential operator on the specified
smooth core, not a claim that all four unbounded factors
are separately defined on the whole Hilbert space.

\emph{The remaining half-line multiplier uses the sum norm.}
Let \(K_a(g)\) be the trace-class operator of MA-MO.125,
which agrees on \(\mathcal E\) with
\(C_g(M_\phi-Z_aM_\phi J_a)\). The agreement here holds
on all of \(\mathcal E\): (MA-MO.324) makes the inverse
transform Schwartz, and compact frequency cutoff with
\(j_a\chi_NV_ay=\chi_NV_aF_y\) gives the same graph
approximation as in MA-MO.128--129.

The additional map
\(N_a(g)=C_gM_\phi:H\to H_{a,s}\) is trace class as
well. After the two weight isometries, its physical kernel is
\[
 \frac{e^{ax}(1+x^2)^{s/4}}
      {(e^{2ay}+e^{-2ay})^{1/2}(1+y^2)^{s/4}}
                         g(x-y)\phi(y).
 \tag{MA-MO.332}
\]
The difference \(x-y\) lies in a fixed compact set. At
\(y\to-\infty\) the exponential ratio decreases like
\(e^{2ay}\), up to a bounded factor depending on that
compact set; at the other end \(\phi(y)\) vanishes.
The same estimates hold after every derivative, and the
polynomial weight ratio is harmless on a bounded difference.
Thus (MA-MO.332) is Schwartz in \((x,y)\), and its Hermite
expansion proves the assertion. No trace-class claim for
\(C_gM_\phi\) on a single unchanged one-sided space is
being made.

\emph{Removing the actual polar sections before combining the rows.}
Write \(f_0,f_1\) for (MA-MO.251), and set
\[
       e_0=\Theta f_0,\qquad
       e_1=\Theta f_1-p_-,\qquad Re_0=e_1.
 \tag{MA-MO.333}
\]
Both vectors belong to \(H\). The theta estimate at the
right and Poisson summation at the left show that \(e_0\)
has left tail \(-e^{x/2}\), while \(e_1\) has right tail
\(-e^{-x/2}\), with rapid opposite tails. These are
integrable in the sum weight since \(a<1/2\). The reflection
identity follows from \(\mathfrak Ff_0=f_1\) and (MA-MO.84).

For any even Schwartz source \(h\), the exact centered
pole subtraction gives, in locally integrable functions,
\[
       \Theta h=Z_aF_h+\left(\int h\right)p_-,
                    \qquad F_h(x)=2e^{x/2}h(e^x).
 \tag{MA-MO.334}
\]
This is the contour identity of MA-MO-004: shifting the
wide direct theta transform to \(a<1/2\) crosses only
the zeta pole, with the displayed positive residue. Zeta
zeros are not poles in this direct summation step.

Subtract the two actual sections from \(C_gDi_-y\)
using its moments (MA-MO.329). The first component of the
remaining diagonal vector is exactly
\[
 \begin{aligned}
 T_{a,g}y={}&K_a(g)y-N_a(g)y+O_a(g)Ry\\
       &+G(-1/2)\lambda_1(y)e_0
         +G(1/2)\lambda_0(y)e_1.
 \end{aligned}
 \tag{MA-MO.335}
\]
To check the formula, the radial vector of \(h_y\) is
\(qF_y+\mathcal H(qF_{Ry})\). Apply (MA-MO.334) and
(MA-MO.328), and use
\(Z_aM_qJ_a-1=(M_\phi-Z_aM_\phi J_a)-M_\phi\).
The subtracted source moments turn \(\Theta f_1\) into
\(\Theta f_1-p_-\), with exactly the signs in
(MA-MO.335). Every term in that formula is now a
trace-class map \(H\to H_{a,s}\): use the bounded input
inclusion for \(K_a\), (MA-MO.332) for \(N_a\), reflection
followed by that inclusion for \(O_a\), and the two
bounded functionals for the finite columns.

The second component of the same diagonal remainder is
\(RT_{a,g^\vee}Ry\), where \(g^\vee(x)=g(-x)\).
Indeed (MA-MO.326) gives \(h_{Ry}=\mathfrak Fh_y\),
Fourier scaling gives
\(\mathfrak F\mathcal D_g=\mathcal D_{g^\vee}\mathfrak F\),
and the two Gaussian sections are interchanged by Fourier.
The two components therefore agree on every
\(y\in\mathcal E\), by the original Poisson identity.
Each defines a trace-class map from \(H\) to its respective
weighted component. The diagonal isometry
\(\iota:H\to H_{a,s}\oplus H_{-a,s}\) has closed range.
Since \(\mathcal E\) is dense in \(H\), the paired map
has range in this diagonal on the entire completion.
Consequently there is a unique trace-class map
\[
       T_g:H\to H,\qquad
       \iota T_g y=(T_{a,g}y,RT_{a,g^\vee}Ry).
 \tag{MA-MO.336}
\]
This uses the actual sum norm and source compatibility;
it does not replace that norm by either one-sided norm.

\emph{The complete original matrix.}
The negative source column has now extended as
\(y\mapsto(T_gy,\nu_g(y))\in\mathcal W\), a trace-class
map. The two polar source columns
\(C_gDs_j\) are actual fixed vectors of
\(\mathscr S_{\rm sum}\subset\mathcal W\), as proved in
MA-MO-016. Thus in the source-section coordinates the full
extension is
\[
       \mathscr D_g(y,c)=
             (T_gy,\nu_g(y))
                       +\sum_{j=0}^1c_j\,C_gDs_j.
 \tag{MA-MO.337}
\]
It is trace class and agrees with the original operator on
the dense smooth domain
\(\mathfrak s(\mathcal E\oplus\mathbb C^2)\).
This proves uniqueness and its agreement with MA-MO-016
on \(\mathcal W_+\). All kernel estimates use finitely
many seminorms of \(g\) on a fixed compact support, and
the finitely many actual source columns depend continuously
on those tests. This proves the stated trace-norm continuity.
\end{proof}

\paragraph{Exact scope of the repaired matrix.}
The two positive-line inverse identities (MA-MO.324) are
where the zero-free right half-plane is used. The argument
does not infer them from a generic zero-free line alone, nor
does it transport an original inverse across unaccounted
right-exterior zeros. The two cutoff moments and actual
polar sections are required by (MA-MO.328)--(MA-MO.335).
They are not zero for a general negative source input.

The matrix (MA-MO.322) is a defined trace-class endomorphism
of the specified enlarged common completion. The positive
source endomorphism \(\jmath_+A_{+,g}\) evaluated in
MA-MO-017 is a different operator on a different space;
its finite-column difference is (MA-MO.299). Likewise
the raw full matrix need not preserve \(\mathcal W_+\),
so its quotient trace cannot be assigned by deleting that
distinction. No new arithmetic zero-location assertion,
or unconditional narrow trace identity, follows here.

The source maps are the freshly checked Meyer (2005)
\texttt{lem:Hilm\_p\_embed}, defining \(P_+,P_-\),
\texttt{eq:Poisson\_summation} and
\texttt{lem:Poisson\_sumpm}, author-manuscript pp.~39--41,
supplied \texttt{main.tex} lines~3808--4043.
The new analytic comparison is the explicit source inverse
shift and the two-row construction above. In the inherited
cote~1 nuclear framework, scans~079--080 and~111--112
through ALPOGE\_BRIDGES BR-TR-01/02, this now supplies a
defined endomorphism in the trace ideal, rather than only
a formal finite-length or exact-sequence class. It does not
by itself evaluate that endomorphism trace.

\section{The evaluated one-sided traces and cancellation of the explicit correction}
\label{sec:projected-trace-cancellation}

MA-MO-015 computed the ideal projection without asserting that
its explicit correction vanished.  The actual local-source
return of MA-MO-017 now supplies the additional comparison
needed to evaluate that correction.  The argument below keeps
the source Fourier multiplier and every projection in place.

Retain the nontrivial character and either regime of MA-MO-012.
Use the centered incomplete global multiplier
\[
 A_\eta(z)=E_{\eta,S}(1/2+z)L_\eta(1/2+z),\qquad
 A_{\bar\eta}(z)=E_{\bar\eta,S}(1/2+z)L_{\bar\eta}(1/2+z),
 \tag{MA-MO.360}
\]
and the actual local equivariant-Fourier multiplier
\(m_\eta(z)\) of MA-MO.281.  In particular its finite
conductor, different and Fourier phases are retained.  Put
\[
 \tau_{\eta,S}(g,\phi)=\frac1{2\pi}\int_{\mathbb R}
        \mathcal Mg(it)\frac{m_\eta'(it)}{m_\eta(it)}\,dt
                    +g(0)\int\psi,
       \qquad \psi=1-\phi-\phi^\vee.
 \tag{MA-MO.361}
\]

\paragraph{MA-MO-019 (actual one-sided cancellation and the projected trace).}
For every polynomial exponent \(s\ge0\),
\[
 \begin{aligned}
 \operatorname{Tr}_{J_{\eta,a,s}}B_{\eta,g}
 &=\operatorname{Tr}_{H_{a,s}}T_{\eta,a}(g)\\
 &=\operatorname{Tr}_{H_{-a,s}}T_{\eta,-a}(g)
  =\tau_{\eta,S}(g,\phi).
 \end{aligned}
 \tag{MA-MO.362}
\]
Consequently the correction explicitly defined in (MA-MO.232)
satisfies
\[
                    \mathfrak C_{\eta,a}(g,\phi)=0.
 \tag{MA-MO.363}
\]
This is a consequence of proved source and kernel comparisons.
It is not obtained by omitting the ideal projector.  Indeed,
for the assembled operator \(\mathbb T\) of (MA-MO.188),
\[
 \operatorname{Tr}_{\mathbb H}(\mathbb T\Pi_J)=\tau_{\eta,S},
 \qquad
 \operatorname{Tr}_{\mathbb H}\mathbb T=2\tau_{\eta,S}.
 \tag{MA-MO.364}
\]
The two expressions need not be equal.

\begin{proof}
\emph{The source functional equation with its actual multiplier.}
The nontrivial-sector Poisson equation has zero singular term.
On the original local source it is the identity relating
\(RZ^{\rm geom}_{\bar\eta,S}\mathfrak S_\eta\) to
\(Z^{\rm geom}_{\eta,S}\), as used in MA-MO.181 and
MA-MO.295.  The two actual additive Fourier maps are inverse
on these coinvariants; the global element \(-1\) accounts
for their Fourier-square signs, including the finite factors
in the odd Dirichlet sector.  Since
\(\mathfrak U_\eta=\mathfrak S_{\bar\eta}R\), their
scalar functional equation is exactly
\[
                     A_{\bar\eta}(-z)=A_\eta(z)m_\eta(z).
 \tag{MA-MO.365}
\]
One may obtain it first by Mellin transformation of the actual
Poisson identity on smooth local vectors and then meromorphic
continuation.  Transforms of compact annular source tests
separate each scalar multiplier, so no factor is lost in this
identification.  Equivalently it is the elementary local and
global functional equation with the actual Fourier convention
already fixed in MA-CC/MA-CQ/MA-MO-017.  Replacing the local
map by a bare archimedean transform would not justify (MA-MO.365).

\emph{An exact kernel identity on the one-sided space.}
Fix \(b=a\) or \(-a\).  Write \(z=b+it\), \(w=b+iu\),
and use \(\ell_{b,b}\) from (MA-MO.288).  The full Fourier
kernel \(k_b\) of the assembled operator (MA-MO.187) is
\[
                k_b(t,u)=A_\eta(z)\ell_{b,b}(t,u)
                                             A_\eta(w)^{-1}.
 \tag{MA-MO.366}
\]
Here is a domain-valid verification.  On the compact-frequency
core \(V_b^{-1}C_c^\infty\), the global inverse multiplier
is smooth on the compact input support.  Multiplication by
\(\phi\) preserves weighted Schwartz space, the global
forward multiplier has polynomial growth, and both local
multipliers \(m_\eta^{\pm1}\) have polynomial derivatives
on the input line.  Thus all sequential expressions below
exist on that core.

In the Fourier picture the reflected second source term has
multiplier \(A_{\bar\eta}(-z)\) and the cutoff
\(\phi^\vee\).  Formula (MA-MO.365) therefore gives the
literal algebra
\[
 \begin{aligned}
 T_{\eta,b}(g)
 &=C_g\{1-A_\eta M_\phi A_\eta^{-1}
          -A_{\bar\eta}(-\cdot)M_{\phi^\vee}
                                  A_{\bar\eta}(-\cdot)^{-1}\}\\
 &=A_\eta C_g\{1-M_\phi-m_\eta M_{\phi^\vee}
                         m_\eta^{-1}\}A_\eta^{-1}
   =A_\eta C_g\mathcal L_\eta A_\eta^{-1}.
 \end{aligned}
 \tag{MA-MO.367}
\]
Each Fourier multiplier in this display is taken on the fixed
line.  The local bracket is precisely (MA-MO.283), including
its nonsymmetric residual.  Substitution of its proved kernel
gives (MA-MO.366) on the core.

The right side of (MA-MO.366) is Schwartz: in the gap regime
all reciprocal derivatives are polynomial, while in the
generic regime the Gevrey estimates (MA-MO.294) dominate
the logarithmic-square growth of \(A_\eta(w)^{-1}\).
Multiplication by \(A_\eta(z)\) costs only polynomial
growth.  Hence it defines a bounded trace-class operator.
The compact-frequency core is dense in \(H_{b,0}\), and
both sides agree there.  This proves equality of the actual
bounded operators and of their kernels, without an assertion
that the global inverse is bounded on the whole space.

On the kernel diagonal the actual factors cancel:
\(A_\eta(b+it)A_\eta(b+it)^{-1}=1\).
Equation (MA-MO.297), applied on this line, consequently gives
\[
 \operatorname{Tr}_{H_{b,0}}T_{\eta,b}(g)
   =\frac1{2\pi}\int_{\mathbb R}
      \mathcal Mg(b+it)\frac{m_\eta'(b+it)}{m_\eta(b+it)}\,dt
                              +g(0)\int\psi.
 \tag{MA-MO.368}
\]
The local multiplier is holomorphic and nonzero throughout
\(|\operatorname{Re}z|<1/2\).  Its logarithmic derivative
has polynomial growth on every closed substrip.  The compact
smooth transform \(\mathcal Mg\) decreases faster than
every power there.  Ordinary rectangles therefore move this
integral from \(b\) to zero without residues; their horizontal
integrals tend to zero.  This gives \(\tau_{\eta,S}\).
No global reciprocal is used on the central line.

The same Schwartz kernels give trace-class maps between any
two polynomially weighted versions of \(H_b\).  Composition
with the bounded inclusion from the larger exponent to the
smaller, as in (MA-MO.213)--(MA-MO.214), proves that their
traces are equal.  Thus (MA-MO.368) holds with exponent \(s\).
The actual ideal trace is \(\tau_{\eta,S}\) by the nuclear
source return (MA-MO.284)--(MA-MO.297), already proved on
that completed ideal.  This proves (MA-MO.362).

Now substitute the two evaluated one-sided traces into the
independently proved projector formula (MA-MO.232).  Its first
term is \(\tau_{\eta,S}\), and its left side is also that
value.  Hence (MA-MO.363) follows.  Finally the trace of a
block diagonal operator is the sum of its two block traces,
whereas the projected trace is (MA-MO.190).  This proves
(MA-MO.364).
\end{proof}

\paragraph{A nonzero test, so that the projector still matters.}
Choose \(\delta>0\) smaller than every positive prime logarithm
occurring in the finitely many unramified local factors of
\(m_\eta\); if there are none, choose any fixed positive
\(\delta\).  Let \(g\) be a nonnegative, nonzero compact
Gevrey function supported in \((0,\delta)\).  It satisfies
\(g(0)=0\).  Finite unramified local factors have their
nonconstant Fourier contributions at
\(\pm k\log N\mathfrak p\), and the ramified/conductor
monomials contribute only a multiple of \(g(0)\).  All
such terms therefore vanish on this test.  The real gamma
factor of parity \(e_v\in\{0,1\}\) and the complex gamma
factor give, respectively, the elementary digamma integrals
already derived in MA-CH-001:
\[
 \begin{aligned}
 \tau_{\eta,S}(g,\phi)
  ={}&-\sum_{v\ {
m real}}\int_0^\infty
       \frac{e^{-(e_v+1/2)x}}{1-e^{-2x}}g(x)\,dx\\
     &-\sum_{v\ {
m complex}}\int_0^\infty
       \frac{e^{-x/2}}{1-e^{-x}}g(x)\,dx<0.
 \end{aligned}
 \tag{MA-MO.369}
\]
There is at least one archimedean place.  This proves that
dropping \(\Pi_J\) in (MA-MO.364) would in general double a
nonzero trace, even for the requested test class.  The
example is available unconditionally at every width in the
specified generic/Gevrey set.  It asserts no global
zero-location result.

The operator \(\mathbb T\) in (MA-MO.364) is the assembled
block diagonal operator of MA-MO-012, not a newly constructed
extension of every entry of the original common-space matrix.
The result evaluates the positive-ideal source trace.  The
full narrow zero quotient and the omitted-divisor condition
of MA-FT-002 remain separate.  Its source locators are the
actual Poisson and Fourier equations used in MA-MO-012/017,
the complete source-kernel return (MA-MO.288)--(MA-MO.297),
and the previously derived gamma integrals MA-CH.1 and .6.
All trace equalities use actual operators with compatible
nuclear pairings; no categorical weight or algebraic divisor
class substitutes for that analytic comparison.

\section{The polar source column has an exact rank-two quotient defect}
\label{sec:polar-quotient-rank}

Under the half-plane hypothesis of MA-MO-018, all nontrivial
zeros lie strictly inside the completed strip. This lets the
finite outgoing condition of MA-MO-016 be decided exactly.
The two polar columns remain independent after passage to the
actual zero quotient, for every nonzero compact test. The
argument retains the full multiplicity of every zero.

Keep \(0\le a_*<a<1/2\), \(s\ge0\), and the explicit
zero-free half-plane (MA-MO.320). Write
\[
 \begin{gathered}
 H=H_{\Sigma,a,s},\qquad J=\overline{\Theta\mathscr U_0}^{\,H},
 \qquad\mathcal W=H\oplus\mathbb C^2,
 \qquad\mathcal W_+=J\oplus\mathbb C^2,\\
 \mathfrak q:\mathcal W\longrightarrow H/J,
                    \qquad\mathfrak q(y,c)=y+J.
 \end{gathered}
 \tag{MA-MO.400}
\]
The source sections are the actual Gaussian vectors
\(s_j=i_+f_j\) of (MA-MO.251), and \(r_j\) are the
actual functions in (MA-MO.255).

\paragraph{MA-MO-020 (exact outgoing rank and noninvariance).}
For every nonzero \(g\in C_c^\infty(\mathbb R)\), the
map
\[
 \begin{aligned}
 \mathbb C^2&\longrightarrow H/J,\\
 (c_0,c_1)&\longmapsto[C_g(c_0r_0+c_1r_1)]
 \end{aligned}
 \tag{MA-MO.401}
\]
is injective. Consequently
\[
        \operatorname{rank}
         (\mathfrak q\mathscr D_g|_{\mathcal W_+})=2,
                \qquad
        \mathscr D_g(\mathcal W_+)\not\subseteq\mathcal W_+.
 \tag{MA-MO.402}
\]
Here \(\mathscr D_g\) is the actual full trace-class matrix
of MA-MO-018, and its positive source restriction is the map
of MA-MO-016. A correction making that source restriction
take values in \(\mathcal W_+\) must have rank at least two.
The orthogonal correction
\[
       (1-P_{\mathcal W_+})
                        \mathscr D_g|_{\mathcal W_+}
 \tag{MA-MO.403}
\]
has rank exactly two and attains that minimum. It changes
the original source action.

The conclusions still hold when \(g(0)=0\), and even when
\(\mathcal Mg(1/2)=\mathcal Mg(-1/2)=0\), provided
\(g\ne0\). They concern this particular completed source
image and the original operator. They are not an obstruction
to every possible return map or relative-trace construction.

\begin{proof}
\emph{The exact two source Mellin residues.}
Poisson summation and the Gaussian Fourier interchange give
\[
       r_0=\Theta f_0+\phi e^{x/2},\qquad
       r_1=\Theta f_1-\phi e^{-x/2}.
 \tag{MA-MO.404}
\]
For example \(R\Theta f_1=\Theta f_0+e^{x/2}\),
whereas \(R\Theta f_0=\Theta f_1-e^{-x/2}\).
Substitution in (MA-MO.255) proves both signs in
(MA-MO.404). Each \(r_j\) lies in \(\mathcal E\):
the displayed term cancels its left Poisson tail, and the
cutoff vanishes at the right.

Let
\[
 \begin{aligned}
 F_j(x)&=2e^{x/2}f_j(e^x),&Z(z)&=\zeta(1/2+z),\\
 P(z)&=\int e^{zx}\phi'(x)\,dx,&P(0)&=-1.
 \end{aligned}
 \tag{MA-MO.405}
\]
The Gaussian integral and gamma recurrence give, first for
\(\Re z>-1/2\),
\[
 \mathcal MF_0(z)=(1/2-z)\Gamma_{\mathbb R}(1/2+z),
 \qquad
 \mathcal MF_1(z)=(1/2+z)\Gamma_{\mathbb R}(1/2+z).
 \tag{MA-MO.406}
\]
In particular both are holomorphic in the open strip
\(|\Re z|<1/2\). The second has value one at \(z=1/2\),
and the first has value zero there, in agreement with the
additive moments of the sections.

On \(\Re z>1/2\), absolute theta summation gives
\(\mathcal M\Theta f_j(z)=Z(z)\mathcal MF_j(z)\).
Integration by parts gives
\(\int e^{\xi x}\phi(x)\,dx=-P(\xi)/\xi\)
when \(\Re\xi>0\). Thus (MA-MO.404) gives the exact
continued identities
\[
 \begin{aligned}
 \mathcal Mr_0(z)
     &=Z(z)\mathcal MF_0(z)-\frac{P(z+1/2)}{z+1/2},\\
 \mathcal Mr_1(z)
     &=Z(z)\mathcal MF_1(z)+\frac{P(z-1/2)}{z-1/2}.
 \end{aligned}
 \tag{MA-MO.407}
\]
They hold throughout the critical open strip by meromorphic
continuation. The left sides are entire. At \(z=1/2\)
the pole in the second formula cancels because its zeta
residue is one and \(P(0)=-1\). At \(z=-1/2\) the
apparent pole in the first formula cancels its Gaussian
term: \(\Gamma_{\mathbb R}(w)=2/w+O(1)\) and
\(\zeta(0)=-1/2\) give residue \(-1\) in its first
term and residue \(+1\) in its second. The value
\(\zeta(0)=-1/2\) also follows directly from the reflection
formula and the residue one at \(s=1\). Neither endpoint
is a nontrivial zero in the argument
below, so no endpoint division is needed.

If \(\lambda=\rho-1/2\) is a nontrivial zero of
multiplicity \(m_\rho\), the holomorphic first terms in
(MA-MO.407) vanish to order at least \(m_\rho\) there.
For \(c=(c_0,c_1)\), put
\[
 \begin{aligned}
 E_c(z)&=-c_0(z-1/2)P(z+1/2)
                       +c_1(z+1/2)P(z-1/2),\\
 E_c(1/2)&=-c_1,&E_c(-1/2)&=-c_0.
 \end{aligned}
 \tag{MA-MO.408}
\]
Thus \(E_c\) is a nonzero entire function whenever
\(c\ne0\). These two evaluations are legitimate values
of the entire numerator; they are not evaluations of the
singular fractions in (MA-MO.407).

\emph{Full retained jets, not only point evaluations.}
The functional equation reflects every nontrivial zero
\(\rho\) to \(1-\overline\rho\). Hypothesis (MA-MO.320)
therefore implies
\[
          |\Re(\rho-1/2)|\le a_*<a.
 \tag{MA-MO.409}
\]
By the exact sum-weight closure theorem MA-WP-001, every
zero is strictly interior and all of its
\(m_\rho\) jets are retained, for every \(s\ge0\).
Suppose \(C_g(c_0r_0+c_1r_1)\in J\), and write
\(G(z)=\mathcal Mg(z)\). Its Mellin transform is
\(G(z)\mathcal M(c_0r_0+c_1r_1)(z)\).
It must vanish to order at least \(m_\rho\) at every
\(\lambda=\rho-1/2\). Equations (MA-MO.407)--(MA-MO.408)
then show that
\[
                G(z)E_c(z)
   \text{ vanishes to order at least }m_\rho
                     \text{ at every }z=\rho-1/2.
 \tag{MA-MO.410}
\]
Indeed multiplication by \(z^2-1/4\) is invertible in
each of these local jet rings. No simplicity assumption
has entered this implication.

\emph{The exponential-type zero count.}
If \(c\ne0\) and \(g\ne0\), the product \(GE_c\)
is not identically zero, by Fourier uniqueness for \(g\)
and (MA-MO.408). Compactness of the supports of \(g\)
and \(\phi'\) gives constants \(C,R\) such that
\[
             |G(z)E_c(z)|\le C(1+|z|)e^{R|z|}
                                     \quad(z\in\mathbb C).
 \tag{MA-MO.411}
\]
Here the factors \(e^{\pm x/2}\) on the fixed support
of \(\phi'\) only change \(C\).

To spell out the zero-count consequence, choose a point
\(z_0\) where \(G(z_0)E_c(z_0)\ne0\). Jensen's formula
on the disk of radius \(2r\) about \(z_0\) gives
\[
 n(r)\log2
 \le\frac1{2\pi}\int_0^{2\pi}
       \log|G(z_0+2re^{it})E_c(z_0+2re^{it})|\,dt
           -\log|G(z_0)E_c(z_0)|=O(r).
 \tag{MA-MO.412}
\]
Here \(n(r)\) counts zeros in the smaller disk, with
multiplicity. Radii meeting a zero are handled by a limit
or a slightly larger radius. Thus (MA-MO.410) would imply
that the full nontrivial zeta divisor has only \(O(T)\)
zeros of height at most \(T\), counted with multiplicity.

\emph{Why the full zeta divisor cannot have that count.}
We prove the needed incompatibility, so a distinct-zero
count or an unexamined asymptotic is not being substituted.
Let
\[
       \xi(s)=\tfrac12s(s-1)\pi^{-s/2}\Gamma(s/2)\zeta(s).
 \tag{MA-MO.413}
\]
It is entire of order at most one, with nonzero values
at zero and one and with precisely the nontrivial zeta
divisor. The order bound follows from the Euler bound on
\(\Re s\ge2\), the functional equation on
\(\Re s\le-1\), and the fixed-strip polynomial zeta
estimate and gamma growth between them, exactly as in the
Hadamard derivation (MA-RB.8)--(MA-RB.10).

Suppose its zero count, including multiplicities for
\(|\Im\rho|\le T\), were \(O(T)\). Hadamard
factorization and logarithmic differentiation give
\(\xi'/\xi(s)=B+\sum_\rho((s-\rho)^{-1}+\rho^{-1})\),
where \(B=\xi'(0)/\xi(0)\) is real. All sums of real
parts below converge absolutely: \(0<\Re\rho<1\)
and the assumed count imply convergence of
\(\sum_{|\rho|>1}|\rho|^{-2}\).
Put \(S=\sum_\rho\Re(1/\rho)\). The functional equation
at zero and one, and symmetry under \(\rho\mapsto1-\rho\),
give \(-B=B+2S\). Therefore for real \(\sigma>1\),
\[
       \frac{\xi'}{\xi}(\sigma)
          =\sum_\rho\frac{\sigma-\Re\rho}
                    {(\sigma-\Re\rho)^2+(\Im\rho)^2}.
 \tag{MA-MO.414}
\]
For \(\sigma\ge2\), each summand is at most
\(\sigma/(\sigma^2/4+(\Im\rho)^2)\). The zeros with
\(|\Im\rho|\le\sigma\) contribute \(O(1)\), using
their assumed \(O(\sigma)\) count. The dyadic shell
\(2^k\sigma<|\Im\rho|\le2^{k+1}\sigma\) contributes
\(O(2^{-k})\). Consequently the whole right side of
(MA-MO.414) is \(O(1)\), uniformly as \(\sigma\to\infty\).

On the other hand, direct differentiation of (MA-MO.413)
gives
\[
 \frac{\xi'}{\xi}(\sigma)
   =\frac1\sigma+\frac1{\sigma-1}-\frac12\log\pi
       +\frac12\psi(\sigma/2)
       +\frac{\zeta'}{\zeta}(\sigma).
 \tag{MA-MO.415}
\]
The final term tends to zero by its absolutely convergent
Euler series, dominated for \(\sigma\ge2\). For \(x>1\),
let \(U\) have density \(u^{x-1}e^{-u}/\Gamma(x)\).
Then \(\psi(x)=\mathbb E\log U\) and
\(\mathbb E(U^{-1})=1/(x-1)\). Jensen applied to
\(\log(U^{-1})\) proves
\(\psi(x)\ge\log(x-1)\). Thus (MA-MO.415) tends to
\(+\infty\), contradicting the bounded value of
(MA-MO.414). This proves the required failure of an
\(O(T)\) full zero count, including every multiplicity.
The contradiction with (MA-MO.412) proves that membership
in \(J\) was impossible for \(c\ne0\).

\emph{The exact outgoing rank and correction.}
We have proved injectivity of (MA-MO.401). Formula
(MA-MO.256) says that the outgoing map of the positive
source column is its negative, applied to the two polar
coordinates, and vanishes on its \(J\) summand. Its
rank is consequently exactly two. If a linear correction
\(C:\mathcal W_+\to\mathcal W\) makes
\(\mathscr D_g-C\) take values in \(\mathcal W_+\),
then \(\mathfrak qC=\mathfrak q\mathscr D_g\) has rank
two, so \(\operatorname{rank}C\ge2\). Orthogonal
projection onto the complement of \(\mathcal W_+\)
identifies that complement with \(H/J\); hence (MA-MO.403)
has rank exactly two. This proves all the assertions.
\end{proof}

\paragraph{Counting input and framework scope.}
The only zero-location use is (MA-MO.409), which makes the
entire nontrivial divisor retained in the actual Hilbert
quotient. The counting contradiction itself is unconditional.
For comparison, the stronger classical count is explicitly
quantified by Hasanalizade, Shen, and Wong (2022),
Corollary~1.2, equation~(1.5), p.~221; its completed
normalization and argument-principle setup are in \S2,
equations~(2.1)--(2.3), pp.~222--223. The argument above
does not require their numerical estimate.

The primary reference is: Hasanalizade, E., Shen, Q., \&
Wong, P.-J. (2022). Counting zeros of the Riemann zeta
function. \emph{Journal of Number Theory, 235}, 219--241.
\url{https://doi.org/10.1016/j.jnt.2021.06.032}.
The completed normalization and reflection identity also
agree with NIST DLMF, equations~25.4.3--25.4.4.
Only the stated source locators and counting setup were
checked here; the numerical estimate is corroborative,
not a new unproved premise of MA-MO-020.

At each retained zero the argument uses the full local
finite-length factor \(\mathbb C[z]/(z^{m_\rho})\),
with all its nilpotent jets. The inherited cote~1 bookkeeping
therefore retains precisely the multiplicities needed in
Jensen and Hadamard. The rank-two orthogonal correction
is a concrete changed analytic map. Neither a finite-length
\(K_0\) identity nor a weight assertion makes the original
source image invariant or assigns a quotient trace to the
uncorrected matrix.

\section{The actual rational common-space trace and its prime contribution}
\label{sec:polar-relative-trace}

The full matrix of MA-MO-018 has a trace on its specified
completion. Its failure to preserve the positive source space
does not prevent a relative trace comparison: a different,
explicit trace-class operator does preserve that space. We
construct that operator before invoking trace additivity.

Retain exactly the rational setting and zero-free half-plane
hypothesis (MA-MO.320). Write
\[
 \begin{gathered}
 Q=H/J=\mathcal W/\mathcal W_+,\qquad
 G(z)=\int g(x)e^{zx}\,dx,\\
 p(g)=G(-1/2)+G(1/2),\\
 \tau(g,\phi)=\frac1{2\pi}\int G(it)\frac{m'(it)}{m(it)}\,dt
                  +g(0)\int\psi .
 \end{gathered}
 \tag{MA-MO.440}
\]
Here \(g\in C_c^\infty(\mathbb R)\), \(s\ge0\),
\(\psi=1-\phi-\phi^\vee\), and \(m\) is the actual
rational multiplier (MA-MO.282). Let
\(L_g=\jmath_+A_{+,g}\) be the trace-class endomorphism of
\(\mathcal W_+\) in MA-MO-017. The notation \(C_g\) on
each space denotes its actual integrated translation action.

\paragraph{MA-MO-021 (relative trace with the actual prime-bearing matrix).}
The operator initially defined on the original smooth common
space by \(C_g(1-P_+)\) extends uniquely to
\[
 \mathscr E_g=C_g(1-P_-)-\mathscr D_g
                         \in\mathcal S_1(\mathcal W).
 \tag{MA-MO.441}
\]
It preserves \(\mathcal W_+\), its restriction is \(-L_g\),
and its induced operator on \(Q\) is exactly \(C_g\).
Consequently
\[
 \begin{aligned}
 \operatorname{Tr}_{\mathcal W}\mathscr E_g
       &=-\tau(g,\phi)+\operatorname{Tr}_{Q}C_g,\\
 \operatorname{Tr}_{\mathcal W}\mathscr D_g
       &=p(g)+\tau(g,\phi)-\operatorname{Tr}_{Q}C_g.
 \end{aligned}
 \tag{MA-MO.442}
\]
For every \(z\in\mathbb C\), there is also the actual
Fredholm determinant identity
\[
 \det_{\mathcal W}(1-z\mathscr E_g)
   =\det_{\mathcal W_+}(1+zL_g)\,
                       \det_Q(1-zC_g).
 \tag{MA-MO.443}
\]
This is a factorization for \(\mathscr E_g\), not for the
raw matrix \(\mathscr D_g\).

For the requested tests \(g(0)=0\), the full matrix trace is
the finite prime distribution
\[
 \boxed{\displaystyle
 \operatorname{Tr}_{\mathcal W}\mathscr D_g
   =\sum_p\sum_{k\ge1}\frac{\log p}{p^{k/2}}
                       \bigl(g(k\log p)+g(-k\log p)\bigr).}
 \tag{MA-MO.444}
\]
The sum is finite. Formula (MA-MO.442) applies without the
value-zero restriction; its nonsymmetric-cutoff term
\(g(0)\int\psi\) must then be retained.

\begin{proof}
\emph{The finite-rank term and its two actual diagonal entries.}
In the source-section coordinates of MA-MO-016, the original
bounded projection is
\[
 \begin{gathered}
 P_-(u,c)=(u+c_0r_0+c_1r_1,0),\\
 (1-P_-)(u,c)=\sum_{j=0}^1c_jv_j,
 \quad v_j=(-r_j,e_j).
 \end{gathered}
 \tag{MA-MO.445}
\]
The actual functions \(r_j\in\mathcal E\subset H\) are
(MA-MO.255); the \(e_j\) here are the standard basis of
\(\mathbb C^2\). Thus \(1-P_-\) is a bounded rank-two
operator on \(\mathcal W\). The integrated representation
\(C_g\) is bounded by the strong continuity and compact-support
estimates of MA-MO-016. This proves that
\(F_g=C_g(1-P_-)\) has rank at most two.

Let \(\pi_j\) be its two polar coordinate functionals.
The rank-one trace formula gives
\(\operatorname{Tr}F_g=\sum_j\pi_j(C_gv_j)\).
Translation acts on the polar quotient by
\(\operatorname{diag}(e^{-t/2},e^{t/2})\), including on
these actual vectors; its diagonal polar coordinates after
integration are \(G(-1/2)\) and \(G(1/2)\). Therefore
\[
                         \operatorname{Tr}F_g=p(g).
 \tag{MA-MO.446}
\]
No invariance of the range of \(1-P_-\) is needed in this
rank-one calculation. By MA-MO-018, \(\mathscr D_g\) is
trace class, so (MA-MO.441) is trace class as well. It agrees
with the original \(C_g(1-P_+)\) on the dense smooth
common space, proving its stated extension and uniqueness.

\emph{The actual invariant subspace and quotient.}
On the original smooth space, \(P_+=i_+p_+\) takes values
in the positive source image. That image is preserved by
translation. Thus \(C_g(1-P_+)\) takes a positive source
vector into the same image. Since that image is dense in the
closed subspace \(\mathcal W_+\), continuity proves that
\(\mathscr E_g\mathcal W_+\subseteq\mathcal W_+\).

The exact rational source formula, before any completion, is
\[
 \begin{aligned}
 p_+i_+&=M_{1-\phi}+\mathfrak F M_{1-\phi}\mathfrak F
        =1+L,\\
 L&=1-M_\phi-\mathfrak F M_\phi\mathfrak F.
 \end{aligned}
 \tag{MA-MO.447}
\]
Here the additive source cutoff means
\(\phi(\log|t|)\), extended by one at \(t=0\), and
\(\mathfrak F^2=1\) on the rational even additive
source. This is precisely the local operator used in
MA-MO.286; it does not assert that \(P_+\) is a projection.
Equivariance of \(i_+\) now gives
\[
 \mathscr E_gi_+f=-i_+\mathcal D_gLf=-L_gi_+f.
 \tag{MA-MO.448}
\]
Both sides are bounded on \(\mathcal W_+\), and the
displayed source is dense there. Thus the restriction equals
\(-L_g\). Its trace is \(-\tau(g,\phi)\) by the actual
nuclear return and trace comparison of MA-MO-017.

Let \(\mathfrak q:\mathcal W\to Q\) be (MA-MO.253).
For a smooth common-space vector \(w\), the positive range
of \(C_gP_+\) gives
\(\mathfrak q\mathscr E_gw=\mathfrak q C_gw\).
Both sides are continuous on \(\mathcal W\). The identity
therefore holds everywhere and, by equivariance of
\(\mathfrak q\), identifies the induced quotient operator
with the actual \(C_g\) on \(H/J\).

\emph{Trace and determinant additivity on these spaces.}
Use the orthogonal decomposition
\(\mathcal W=\mathcal W_+\oplus\mathcal W_+^\perp\).
Invariance makes \(\mathscr E_g\) upper triangular.
Its diagonal blocks are \(-L_g\) and the operator unitarily
identified with \(C_g\) on the Hilbert quotient. Both blocks
are trace class by compression of \(\mathscr E_g\).
Summing the diagonal against the union of orthonormal bases
of the two summands, with absolute convergence, proves the
first equation in (MA-MO.442). Subtract it from (MA-MO.446)
to obtain the second.

For completeness, choose increasing finite-dimensional
orthogonal projections in each summand, and take their direct
sum. Compression converges in trace norm for each of the
three trace-class operators. The finite-dimensional matrices
are upper triangular in the same block decomposition, so
their determinants multiply. Continuity of the Fredholm
determinant in trace norm proves (MA-MO.443) for every \(z\).
This establishes the exact-sequence formula on actual
operators, rather than invoking a formal quotient class.

\emph{Prime evaluation and the precise use of zero location.}
Hypothesis (MA-MO.320), combined with the functional equation
of \(\zeta\), puts every nontrivial zero in
\(|\Re\rho-1/2|\le a_*<a\). All its multiplicities are
therefore interior retained jets. The already proved actual
all-weight quotient theorem MA-WP-001--003 gives
\[
 \operatorname{Tr}_Q C_g
    =\sum_\rho m_\rho G(\rho-1/2).
 \tag{MA-MO.449}
\]
The change from the maximum exponential weight there to the
sum weight here is a bounded isomorphism of the same
completed quotient, since their ratio lies between one and
two; its trace is unchanged. This step uses the stated
zero-location premise to retain the entire divisor.

If \(g(0)=0\), the explicit local gamma computation in
MA-MO-017 gives \(\tau(g,\phi)=-A_0(h)\), with
\(h(e^x)=g(x)\) and \(A_0\) from (MA-TR-2).
The unconditional full explicit formula (MA-TR-4), inherited
from the checked F-SM-13/MT-M30 conventions and Meyer's
Theorems 5.8/5.11, reads
\(\sum_\rho m_\rho G(\rho-1/2)=p(g)-A_0(h)-B(h)\).
Substitution into (MA-MO.442) proves (MA-MO.444).
The full explicit formula itself uses no zero location;
the narrow inverse extension and (MA-MO.449) do.
\end{proof}

\paragraph{Scope, source locators, and the original trace gap.}
The actual source formula (MA-MO.447) is Meyer (2005),
immediately following \texttt{lem:Hilm\_p\_embed}, supplied
\texttt{main.tex} lines 3937--3966. The relative trace argument
parallels his \texttt{the:global\_trace\_formula}, lines
4204--4266, but here every map is first proved on the narrow
Hilbert completion in use. The set of places remains
\(S=\{\infty\}\). Meyer's later choice of \(S\) large
enough to contain every contributing prime is not imposed.
Accordingly the common matrix need not have trace zero:
its trace is exactly the omitted finite-prime distribution
on the requested value-zero tests. For a test whose support
meets \(\{\pm k\log p\}\) only at \(\log2\), with
\(g(\log2)=1\), that trace is \(\log2/\sqrt2\), the finite-prime
size in ALPOGE\_BRIDGES F-SM-21.

Input A supplies (MA-MO.320) at every positive narrow width,
so this is a valid conditional method comparison over
\(\mathbb Q\). Applying A at the inverse-identification
step or to (MA-MO.449) is circular for the requested
zero-location-free test of the method; these uses have not
been hidden inside a trace convention. The theorem does not
identify \(\mathscr D_g\) with the unchanged local cutoff,
make it preserve \(\mathcal W_+\), or prove FT002's
arithmetic deficiency empty without the premise.

In the cote 1 framework, scans 079--080 and 111--112 through
ALPOGE\_BRIDGES BR-TR-01/02, (MA-MO.443) is an actual
nuclear Fredholm identity along an invariant Hilbert exact
sequence. The two polar characters and the full retained
zero-jet quotient remain distinct terms. It creates no
faithful positive full-jet realization and assigns no
additional motivic or algebraic weight structure.

\section{The actual full character matrix under an explicit half-plane hypothesis}
\label{sec:character-common-matrix}

For a nontrivial character the original common-sum sector is
diagonal and has no polar modes. We construct its entire
smoothed matrix on the inherited narrow Hilbert completion,
using the actual local source and both dual inverses. The
construction has an explicit zero-free half-plane premise.
Its positive-ideal restriction is the operator already
evaluated in MA-MO-017, while its quotient action retains
the negative convolution contribution.

Let \(K\) be a fixed number field, and let \(\eta\) be a
nontrivial primitive finite-order Hecke character. Choose the
actual coherent additive character, self-dual measures, and
norm splitting of MA-NR/MA-CC/MA-CQ. The finite place set
\(S\) contains the archimedean places, all ramification of
\(\eta\), and all places required by the fixed additive
character. In addition it is \emph{structural} in Meyer's
precise sense:
\[
      C_S\simeq C_K/\mathcal O_{\overline S}^{\times}.
 \tag{MA-MO.470}
\]
Here \(\mathcal O_{\overline S}^{\times}\) is the product
of the local unit groups outside \(S\). Thus the chosen
global compact character sector is the actual corresponding
sector of \(\mathcal T(C_S)\), not an arbitrary radial
representative. Such sets exist by finiteness of the ideal
class group. This requirement is stronger than merely
containing the conductor primes.

Use probability Haar on the compact norm-one quotient and
the same centered conventions as MA-MO-012. Put
\[
 \begin{gathered}
 H=H_{\Sigma,a,s},\quad 0\le a_*<a<1/2,\quad s\ge0,\\
 I_\eta=I_\eta^c=\vartheta_\eta\mathcal T_\eta\subset\mathcal E,
 \qquad J_\eta=\overline{I_\eta}^{\,H},\qquad Q_\eta=H/J_\eta.
 \end{gathered}
 \tag{MA-MO.471}
\]
The same norm spaces in the dual compact sector are understood
when \(\bar\eta\) appears. Assume the explicit premise
\[
 L_K(\sigma+it,\omega)\ne0
       \quad\text{for }\sigma>1/2+a_*,
             \qquad \omega\in\{\eta,\bar\eta\}.
 \tag{MA-MO.472}
\]
The finite functions are entire for these nontrivial characters.
The premise concerns their positive-half-plane zeros, which
there coincide with the completed nontrivial zeros.
Let \(\phi\) be the real smooth source cutoff, equal to one
sufficiently far left and zero sufficiently far right, with
\(\phi'\in C_c^\infty\); no symmetry is imposed.

\paragraph{MA-MO-022 (full common-space character operator and its trace).}
The original smoothed common-space matrix
\(C_g(P_+-P_-)\), on its actual nontrivial smooth sector,
extends uniquely to a trace-class endomorphism
\[
          \mathscr D_{\eta,g}:H\longrightarrow H
                    \qquad(g\in C_c^\infty(\mathbb R)).
 \tag{MA-MO.473}
\]
It preserves \(J_\eta\), its restriction there is the
actual \(B_{\eta,g}\) of MA-MO-012/017, and its induced
operator on \(Q_\eta\) is exactly \(-C_g\). In particular
\[
 \begin{aligned}
 \operatorname{Tr}_H\mathscr D_{\eta,g}
   &=\tau_{\eta,S}(g,\phi)
                         -\operatorname{Tr}_{Q_\eta}C_g,\\
 \det_H(1-z\mathscr D_{\eta,g})
   &=\det_{J_\eta}(1-zB_{\eta,g})
                         \det_{Q_\eta}(1+zC_g),
                   \qquad z\in\mathbb C.
 \end{aligned}
 \tag{MA-MO.474}
\]
The local trace \(\tau_{\eta,S}\) is exactly (MA-MO.361),
including the actual finite Fourier, conductor and different
factors in \(m_\eta\), and the nonsymmetric-cutoff residual.

For the requested \(g(0)=0\) tests the full matrix trace is
\[
 \boxed{\begin{aligned}
 \operatorname{Tr}_H\mathscr D_{\eta,g}
  &=\sum_{\substack{\mathfrak p\notin S\\k\ge1}}
       \frac{\log N\mathfrak p}{(N\mathfrak p)^{k/2}}
       \bigl[c_{\mathfrak p}^{\,k}g(k\log N\mathfrak p)\\
  &\hspace{10em}
       +\overline{c_{\mathfrak p}}^{\,k}g(-k\log N\mathfrak p)\bigr].
 \end{aligned}}
 \tag{MA-MO.475}
\]
where \(c_{\mathfrak p}=\eta(\mathfrak p)\) at unramified
places. The sum is finite. No assumption that \(S\) contains
all primes contributing to the test has been made.

\begin{proof}
\emph{The actual structural sector and its smooth common space.}
Condition (MA-MO.470) is Meyer's sufficiently-large condition,
supplied \texttt{main.tex} lines~3592--3602. His actual
global/local coinvariant theorem, lines~3671--3707, identifies
the specified sector with \(\mathcal T_\eta\); it applies
without a class-number-one or unit-rank restriction. The
norm-one quotient is compact, so the norm map is proper.
On a fixed compact character, weighted Schwartz functions
are therefore measured by the one real norm coordinate.
The actual inclusion
\[
        \mathcal E_\omega
          =\mathscr S(C_S)_{\cup,\omega}
                          \longrightarrow\mathcal T_\omega,
                   \qquad\omega=\eta,\bar\eta,
 \tag{MA-MO.476}
\]
is the inclusion following the split sequence (24),
\texttt{eq:Twist\_convariants}. It is an inclusion into
the local coinvariant source topology, not a choice of
additive representative.

The singular two-dimensional Poisson representation has
trivial compact type. A nontrivial finite-order global
character is trivial on the connected norm splitting and
nontrivial on the compact quotient. By (MA-MO.470) the
specified local sector is that same nontrivial compact type.
Consequently its singular Poisson term is zero. In this
sector the original common-sum space is exactly
\(i_-\mathcal E=\{(y,y):y\in\mathcal E\}\), and
\(P_-=1\) on it. Its inherited completion is the diagonal
sum-norm space \(H\) in (MA-MO.471).

\emph{Original wide inverses on every smooth diagonal input.}
For \(\omega=\eta,\bar\eta\), write
\[
 \begin{aligned}
 A_\omega(z)&=E_{\omega,S}(1/2+z)L_K(1/2+z,\omega),\\
 E_{\omega,S}(s)&=\prod_{\substack{\mathfrak p\in S\ \mathrm{finite}\\
                                  \omega\ \mathrm{unramified}\ }}
                         (1-c_{\omega,\mathfrak p}(N\mathfrak p)^{-s}),\\
 F_{\omega,y}(x)&=
   \sum_{\mathfrak a\ \mathrm{prime\ to}\ S}
        \mu_K(\mathfrak a)c_{\omega,\mathfrak a}
        (N\mathfrak a)^{-1/2}y(x+\log N\mathfrak a).
 \end{aligned}
 \tag{MA-MO.477}
\]
These are the exact original Euler inverse coefficients of
MA-CC.1. At a wide weight \(b>1/2\) the series converges
absolutely in operator norm, dominated by
\(\zeta_K(b+1/2)\). On \(y\in\mathcal E\) the series
and every derivative converge locally absolutely, using
arbitrarily rapid right decay. Its Mellin transform on
the wide line is \(\mathcal My(z)/A_\omega(z)\).

All zeros of the finite Euler factors in (MA-MO.477) have
\(\Re z=-1/2\), so they do not affect the positive
inverse lines. Hypothesis (MA-MO.472) gives a fixed margin
to every global zero on every fixed strip \(\Re z\ge a\).
The reciprocal proof of MA-RS-004, in its fixed-Hecke version
used in MA-MO-012, gives polynomial bounds for
\(A_\omega^{-1}\) and all derivatives there. There is
no pole of a nontrivial finite \(L\)-function at one.
The entire numerator \(\mathcal My\) is rapidly decreasing
on each fixed strip. Shifting from the original wide line
to any \(b\ge a\) therefore proves
\[
 \begin{gathered}
 e^{bx}F_{\omega,y}\in\mathscr S\quad(b\ge a),\qquad
 F_{\omega,y}=J_{\omega,a}y,\\
 V_aF_{\omega,y}=j_{\omega,a}V_ay,\qquad
 j_{\omega,a}(u)=A_\omega(a+iu)^{-1}.
 \end{gathered}
 \tag{MA-MO.478}
\]
This holds for \emph{every} \(y\in\mathcal E\), not
merely ideal vectors. Equality with the original series is
proved on the wide line before the shift. No inverse is
moved through zeros in the interior narrow strip.

Put \(q=1-\phi\). The function \(qF_{\omega,y}\) belongs
to \(\mathcal E_\omega\): it is zero sufficiently far
left and has all the right weighted derivative estimates
in (MA-MO.478). By (MA-MO.476) it is consequently an actual
local source vector. This is the step that justifies the
next Fourier transform on its original domain.

\emph{Both Fourier sectors and the original positive output.}
Let \(\mathfrak S_\eta:\mathcal T_\eta\to
\mathcal T_{\bar\eta}\) and
\(\mathfrak S_{\bar\eta}:\mathcal T_{\bar\eta}\to
\mathcal T_\eta\) be the actual mutually inverse source
Fourier maps of MA-MO-017/019. They include the finite
phases, conductor and different; in particular no false
archimedean odd-parity involution is used. The original
map \(p_+\) on the smooth diagonal is
\[
 \begin{aligned}
 H_{\eta,y}&=qF_{\eta,y}
                 +\mathfrak S_{\bar\eta}(qF_{\bar\eta,Ry})
                                      \in\mathcal T_\eta,\\
 \mathfrak S_\eta H_{\eta,y}&=H_{\bar\eta,Ry},\\
 (P_+-P_-)i_-y
    &=i_-\bigl(\vartheta_\eta H_{\eta,y}-y\bigr).
 \end{aligned}
 \tag{MA-MO.479}
\]
The first line is the defining source formula with its two
wide inverses; (MA-MO.478) has identified their values.
The second line follows by applying the actual inverse
Fourier pair. Nontrivial-sector Poisson summation makes
the two global output components equal, proving the last
line. In particular the positive output is in the actual
ideal \(I_\eta\).

For any \(F\in\mathcal T_\omega\), the first source
family has weighted Schwartz estimates for every
\(c>-1/2\). Its Mellin transform is holomorphic there.
The direct theta transform is \(A_\omega\mathcal MF\)
on a wide line; shifting it to \(a\) crosses no pole,
because \(A_\omega\) is entire. Thus
\[
                  \vartheta_\omega F=Z_{\omega,a}F
                     \quad\text{on the actual local source}.
 \tag{MA-MO.480}
\]
The source output lies in \(\mathcal E\) by its two
Poisson-related tails. This direct summation step uses
no zero-free assertion: zeros are not poles of the forward
multiplier. There is no rational principal polar subtraction
or hidden moment condition in (MA-MO.480).

\emph{The positive-line crossed kernel.}
Use the actual equivariant multiplier
\(\mathfrak U_\eta=\mathfrak S_{\bar\eta}R\) and
\(m_\eta\) from MA-MO.281. Thus
\(\mathcal M\mathfrak S_{\bar\eta}F(z)
       =m_\eta(z)\mathcal MF(-z)\).
Set
\[
       z=a+it,\qquad w=a+iu,\qquad
       G(z)=\mathcal Mg(z),\qquad
       P(\xi)=\int e^{\xi x}\phi'(x)\,dx.
\]
The operator
\(C_g Z_{\eta,a}\mathfrak S_{\bar\eta}M_qJ_{\bar\eta,a}\)
has, on its original smooth core, the kernel
\[
       o_{\eta,a,g}(t,u)
          =\frac{G(z)A_\eta(z)m_\eta(z)j_{\bar\eta,a}(u)}{2\pi}
                           \frac{P(-z-w)}{-z-w}
 \tag{MA-MO.481}
\]
between the positive-line dual input and positive-line
original output. Indeed Fourier inversion on the input
line and \(\int q(x)e^{\xi x}dx=P(\xi)/\xi\) for
\(\Re\xi<0\) give exactly this formula. Absolute
justification on the core follows from its Schwartz
input transform and \(\Re(-z-w)=-2a\). The source
cutoff vector and its actual Fourier transform have the
required domains by (MA-MO.476)--(MA-MO.480).

All derivatives of (MA-MO.481) are Schwartz: the local
factor \(m_\eta\), the forward multiplier and the
dual inverse have polynomial derivative bounds on these
fixed lines; \(G(a+it)\) and \(P(-2a-i(t+u))\) give
the two rapid directions, and the denominator is separated
from zero. The Bessel/Hermite proof therefore gives an
actual extension
\(O_{\eta,a}(g):H_{a,s}^{\bar\eta}\to H_{a,s}^\eta\)
in trace class, for every \(s\ge0\). Interchanging the
characters gives the other crossed kernel. The input
factor in (MA-MO.481) is the \emph{dual} inverse; replacing
it by the same-character inverse is not permissible.

\emph{Compatible rows in the exact sum norm.}
Let \(K_{\eta,a}(g)\) be the one-sided commutator kernel
(MA-MO.184). Under (MA-MO.472) it agrees on all of
\(\mathcal E\) with
\(C_g(M_\phi-Z_{\eta,a}M_\phi J_{\eta,a})\), by
(MA-MO.478) and the graph cutoff identity
\(j_{\eta,a}\chi_NV_ay=\chi_NV_aF_{\eta,y}\).
It is trace class on \(H_{a,s}\). Also
\(N_a(g)=C_gM_\phi:H\to H_{a,s}\) is trace class:
its normalized physical kernel is precisely (MA-MO.332),
whose compact difference and exponential left decay
were proved there. These norm-coordinate statements do
not depend on the number of places of the fixed field.

The first actual smoothed row of (MA-MO.479) is now
\[
        \mathscr R_{\eta,a}(g)y
          =K_{\eta,a}(g)y-N_a(g)y+O_{\eta,a}(g)Ry.
 \tag{MA-MO.482}
\]
To verify it, use (MA-MO.480) on \(H_{\eta,y}\) and
\(Z_{\eta,a}M_qJ_{\eta,a}-1
   =(M_\phi-Z_{\eta,a}M_\phi J_{\eta,a})-M_\phi\).
Every term is trace class from \(H\) to \(H_{a,s}\).

The second row is
\(R\mathscr R_{\bar\eta,a}(g^\vee)R\), with
\(g^\vee(x)=g(-x)\). The actual Fourier relation in
(MA-MO.479), its scaling by \(g^\vee\), and zero-singular
Poisson summation show that the two rows agree as functions
for every \(y\in\mathcal E\). The paired extension has
values in \(H_{a,s}\oplus H_{-a,s}\), and is trace class
there. Its range is in the closed diagonal, by density of
\(\mathcal E\) in the sum norm. Applying the inverse of
that diagonal isometry gives exactly (MA-MO.473), agreeing
with the original smooth matrix on its whole common sector.
This also proves uniqueness. No unbounded individual
factor has been declared defined on every Hilbert input.

\emph{The actual ideal restriction and quotient action.}
On the dense ideal source \(I_\eta\), the map just
constructed agrees with the original action whose
completion is \(B_{\eta,g}\) in MA-MO-012. Both maps
are bounded, so their restrictions on \(J_\eta\) agree.
In particular \(J_\eta\) is invariant, and its trace is
\(\tau_{\eta,S}(g,\phi)\) by MA-MO-017.

The bounded map \(\mathscr D_{\eta,g}+C_g\) agrees
on \(\mathcal E\) with \(C_gP_+\), whose range is
in \(I_\eta\). Closure therefore puts its entire range
in \(J_\eta\). Since \(C_g\) preserves \(J_\eta\),
the induced quotient map of \(\mathscr D_{\eta,g}\)
is exactly \(-C_g\). In the orthogonal decomposition
\(H=J_\eta\oplus J_\eta^\perp\), the operator is
upper triangular with diagonal blocks \(B_{\eta,g}\)
and the quotient operator \(-C_g\). Compression makes
both blocks trace class. Absolute summation against the
two orthonormal bases gives the trace equality in
(MA-MO.474). Finite-dimensional block compressions converge
in trace norm; their triangular determinant identity and
Fredholm continuity prove the determinant equality as well.

\emph{Full retained zeros and the finite prime evaluation.}
The functional equation between \(\eta\) and \(\bar\eta\),
together with (MA-MO.472) for both, places every completed
zero in \(|\Re\rho-1/2|\le a_*<a\). Thus MA-WP-001--003
retains all multiplicities, with no boundary truncation,
and gives the actual quotient trace
\[
        \operatorname{Tr}_{Q_\eta}C_g
                  =\sum_\rho m_\eta(\rho)G(\rho-1/2).
 \tag{MA-MO.483}
\]
The series is absolutely convergent by the fixed-character
local zero count and rapid Mellin decay. It is the trace
of the actual completed quotient, not an assigned spectral
functional.

For \(g(0)=0\), the full primitive explicit formula in
MA-CH/MT-M30 and Meyer's Theorems~5.8/5.11 reads
\[
       \sum_\rho m_\eta(\rho)G(\rho-1/2)
                   =-\mathcal A_{\infty,\eta}(g)-B_\eta(g).
 \tag{MA-MO.484}
\]
Here \(\mathcal A_{\infty,\eta}\) is the sum of the real
parity and complex gamma integrals, and \(B_\eta\) uses
all primitive finite Euler factors with conjugate coefficients
on the negative argument, as in MA-CH.2 and .6. There is
no principal polar term. The conductor logarithm pairs
with \(g(0)\), hence is zero in this displayed formula.

The actual multiplier \(m_\eta\) contains the local
factors for \(S\). Each unramified finite factor has the
form, with \(q_v=N\mathfrak p_v\),
\[
       \frac{1-\overline{c_{\mathfrak p_v}}q_v^{z-1/2}}
            {1-c_{\mathfrak p_v}q_v^{-z-1/2}}.
 \tag{MA-MO.485}
\]
Its logarithmic derivative on \(\Re z=0\) has the
absolutely convergent expansion
\[
 -\sum_{k\ge1}(\log q_v)q_v^{-k/2}
 \bigl(c_{\mathfrak p_v}^{\,k}e^{-kz\log q_v}
  +\overline{c_{\mathfrak p_v}}^{\,k}e^{kz\log q_v}\bigr).
\]
Mellin inversion therefore contributes the negative of
that place's primitive prime distribution. Ramified finite
factors and different/conductor dilations are nonzero
exponential monomials times phases; their logarithmic
derivatives are constants and pair with \(g(0)\).
The gamma calculation of MA-MO.369, with both signs of
the test argument retained, consequently gives
\[
        \tau_{\eta,S}(g,\phi)
          =-\mathcal A_{\infty,\eta}(g)-B_{\eta,S}(g)
                            \qquad(g(0)=0).
 \tag{MA-MO.486}
\]
Here \(B_{\eta,S}\) sums the finite places in \(S\).
Ramification of \(K/\mathbb Q\) alone does not delete
an unramified character factor. The nonsymmetric residual
also vanishes because it is \(g(0)\int\psi\).
Subtracting (MA-MO.484) from (MA-MO.486) in the actual
trace identity (MA-MO.474) proves (MA-MO.475).
\end{proof}

\paragraph{Conditional scope and actual operators distinguished.}
The half-plane premise is used twice: to identify the
original wide inverse on every smooth common input with
the positive narrow-line inverse, and to retain the entire
completed divisor in (MA-MO.483). The direct forward
theta shift, the local finite phases, and the full explicit
formula themselves require no zero-location assertion.
For a nonprincipal primitive Dirichlet character, A supplies
(MA-MO.472) with \(a_*=0\). For the non-induced cubic
imaginary-quadratic character of MA-CH-003, its own Hecke RH
or the stipulated M supplies that premise; Dirichlet A
alone does not. No claim about its truth is added here.

The full operator \(\mathscr D_{\eta,g}\) is not the
unrestricted assembled block diagonal operator
\(\mathbb T\) of MA-MO-012/019. They agree on the actual
ideal only after the stated identifications. The former has
the negative quotient convolution in (MA-MO.474); the latter
has the separately computed unrestricted trace
\(2\tau_{\eta,S}\). No ideal projector has been dropped
in passing between these statements.

The structural-source locators freshly checked here are
Meyer (2005), supplied \texttt{main.tex} lines~3570--3725
for the sufficiently-large condition and full coinvariant
proof; lines~2630--2675 for the split sequence and
Schwartz-cup inclusion; and \texttt{lem:Hilm\_p\_embed},
\texttt{eq:Poisson\_summation}, and
\texttt{lem:Poisson\_sumpm}, lines~3808--4043, for the
actual matrix. The two-sided local topology and nuclear
return are the independently checked MA-MO-017.
In the inherited cote~1 framework, scans~079--080 and
111--112 through ALPOGE\_BRIDGES BR-TR-01/02, (MA-MO.474)
is an actual nuclear determinant formula along an invariant
Hilbert subspace and quotient. It does not create a positive
full-jet realization, a new motivic weight, or an unconditional
answer to the remaining FT002 arithmetic gate.

\section{The polar obstruction without a zero-free half-plane premise}
\label{sec:polar-critical-line-obstruction}

The rank-two obstruction can be strengthened independently of
the full-matrix extension. The known critical-line zeros alone
already make the two polar quotient classes independent.
This uses a proved partial zero-location theorem, not A, M,
or a supposition that the whole divisor lies in the strip.

Fix any \(0<a<1/2\), \(s\ge0\), and the actual rational
completion \(\mathcal W=H_{\Sigma,a,s}\oplus\mathbb C^2\)
and \(\mathcal W_+=J\oplus\mathbb C^2\) of MA-MO-016.
The sections, cutoff and functions \(r_0,r_1\) are unchanged;
write \(G=\mathcal Mg\) and \(P=\mathcal M\phi'\).

\paragraph{MA-MO-023 (unconditional polar quotient independence).}
For every nonzero \(g\in C_c^\infty(\mathbb R)\),
\[
 \mathbb C^2\longrightarrow H_{\Sigma,a,s}/J,
       \qquad c\longmapsto[C_g(c_0r_0+c_1r_1)]
 \tag{MA-MO.500}
\]
is injective, without (MA-MO.320). Consequently the original
smooth positive source column of \(C_g(P_+-P_-)\) has an
outgoing quotient map of rank exactly two and does not take
all its source vectors into \(\mathcal W_+\).
Whenever that column has the bounded trace-class extension
of MA-MO-016, its outgoing rank is also exactly two and the
minimum correction rank of MA-MO.403 is two. This includes
the unconditional generic-width Gevrey regime. The same
statements hold when \(g(0)=G(1/2)=G(-1/2)=0\), as long
as \(g\ne0\). No unrestricted full-matrix extension is
asserted here.

\begin{proof}
For every real \(\gamma\), the centered Mellin evaluation
\[
              f\longmapsto\int f(x)e^{i\gamma x}\,dx
 \tag{MA-MO.501}
\]
is a continuous functional on \(H_{\Sigma,a,s}\).
Indeed its squared dual norm is
\(\int[(e^{2ax}+e^{-2ax})(1+x^2)^{s/2}]^{-1}dx<\infty\).
If \(\rho=1/2+i\gamma\) is a nontrivial zeta zero, this
functional annihilates the smooth theta ideal and therefore
its closure \(J\). This follows from the actual Mellin
theta identity, already used in MA-MO.407, with the two
source moments zero. Only a value constraint is needed below.

The exact residue identities (MA-MO.404)--(MA-MO.408)
are unconditional identities of the original smooth source.
For \(c\ne0\), their entire numerator
\[
 E_c(z)=-c_0(z-1/2)P(z+1/2)+c_1(z+1/2)P(z-1/2)
 \tag{MA-MO.502}
\]
is nonzero because \(E_c(1/2)=-c_1\) and
\(E_c(-1/2)=-c_0\). If the class in (MA-MO.500)
were zero, (MA-MO.501) and these identities would force
\[
                    G(i\gamma)E_c(i\gamma)=0
 \tag{MA-MO.503}
\]
at every critical-line zero. The denominators
\(z\pm1/2\) are nonzero at these points, so this
implication entails no endpoint division.

For completeness, specify precisely the known counting
input. Conrey (1989), Theorem 1, printed p.~4, with the
counts in equations (8), (11), and (12), printed p.~3,
implies that there are constants \(c>0\) and \(T_0\)
such that the number of simple critical-line zeros with
\(0<\gamma<T\) is at least \(cT\log T\) for all
\(T\ge T_0\). Only this weaker consequence is used.
They are distinct points \(i\gamma\), so (MA-MO.503)
forces that many distinct zeros of \(GE_c\).

But \(g\ne0\) makes \(G\) nonzero, and compact support
of \(g\) and \(\phi'\) gives the exponential-type bound
(MA-MO.411). The Jensen argument (MA-MO.412), which does
not use the half-plane premise, therefore bounds its number
of zeros in \(|z|\le T\), even with multiplicity, by
\(O(T)\). This contradicts the preceding lower bound.
Thus (MA-MO.500) is injective.

On the original positive source image, the algebraic
outgoing formula (MA-MO.256) holds before any analytic
extension is invoked. Every polar coordinate in
\(\mathbb C^2\) occurs there, by the actual two Gaussian
sections. The just-proved injectivity therefore gives rank
two already on this smooth image. If the bounded extension
exists, the same outgoing formula and density preserve that
rank. Factoring a correction through the quotient gives
the lower bound two on its rank; the orthogonal complement
realizes it exactly, as in the final paragraph of MA-MO-020.
None of these steps uses the value \(g(0)\) or either
polar Mellin value \(G(1/2),G(-1/2)\), so the stated
subclasses follow.
\end{proof}

\paragraph{Source and zero-location scope.}
Conrey's theorem is a known published input, cited rather
than reproved. Its statement, counting conventions and
locators were read directly in the original article; this
is not a new audit of its complete mollifier proof. The
unconditional lower bound concerns an established subset
of zeros. It does not say where the other zeros lie and
does not imply their simplicity. Thus MA-MO-023 removes
the containing-half-plane premise from the obstruction,
not from the extension or prime-trace evaluation in
MA-MO-018/021.

Conrey, J. B. (1989). More than two fifths of the zeros
of the Riemann zeta function are on the critical line.
\emph{Journal f\"ur die reine und angewandte Mathematik,
399}, 1--26. \url{https://doi.org/10.1515/crll.1989.399.1}.

In the inherited cote 1 framework, scans 079--080 and
111--112 through BR-TR-01/02, the obstruction already
appears in the value factors of known retained local
zero modules. Allowing their full nilpotent jet structure
cannot remove it. The correction still changes the
actual analytic source map; no positivity or motivic
weight assertion is supplied by that change.

\section{Strict exterior zeros obstruct the actual full character matrix}
\label{sec:character-exterior-obstruction}

The sufficient half-plane premise in MA-MO-022 is not a
disposable device for its inverse estimate. A strictly exterior
zero obstructs boundedness of the original full common-sum
action itself, even after any nonzero compact smoothing.
The proof uses genuine translated smooth inputs and an actual
bounded output functional. It does not replace the original
inverse by a narrow multiplier continuation.

Fix a number field \(K\), a nontrivial primitive finite-order
Hecke character \(\eta\), the structural place set and actual
source conventions of (MA-MO.470), and
\[
 \begin{gathered}
 H=H_{\Sigma,a,s},\qquad 0<a<1/2,\qquad s\ge0,\\
 C_gf(x)=\int g(t)f(x-t)\,dt,\qquad
 U_tf(x)=f(x-t),\qquad Rf(x)=f(-x),\\
 \phi=1\ \hbox{sufficiently far left},\qquad
 \phi=0\ \hbox{sufficiently far right},\qquad
 \phi'\in C_c^\infty,\qquad q=1-\phi .
 \end{gathered}
 \tag{MA-MO.550}
\]
The cutoff is real and smooth; no monotonicity or symmetry
is assumed. In this section there is \emph{no} zero-free
half-plane hypothesis. Put
\[
 \begin{aligned}
 A_\omega(z)&=E_{\omega,S}(1/2+z)L_K(1/2+z,\omega),\\
 F_{\omega,y}(x)&=
   \sum_{\mathfrak a\ {\rm prime\ to}\ S}
     \mu_K(\mathfrak a)c_{\omega,\mathfrak a}
       (N\mathfrak a)^{-1/2}y(x+\log N\mathfrak a),
       \qquad \omega=\eta,\bar\eta .
 \end{aligned}
 \tag{MA-MO.551}
\]
These are the original wide Euler inverses of MA-CC.1.
Write \(\mathfrak S_\eta,\mathfrak S_{\bar\eta}\) for the
actual inverse Fourier pair on the two local source sectors.
The original smooth common-space action is
\[
 \begin{aligned}
 \mathscr D^{\,\rm raw}_{\eta,g}y
   &=C_g\{\vartheta_\eta H_{\eta,y}-y\},\\
 H_{\eta,y}
   &=qF_{\eta,y}
       +\mathfrak S_{\bar\eta}(qF_{\bar\eta,Ry}),
                         \qquad y\in\mathcal E .
 \end{aligned}
 \tag{MA-MO.552}
\]
It is \(C_g(P_+-P_-)\) on the actual nontrivial common
sector, where \(P_-=1\). All finite Fourier phases are
those of the original source.

\paragraph{MA-MO-024 (actual-source strict-exterior obstruction).}
If a completed nontrivial zero \(\rho\) of \(L_K(s,\eta)\)
satisfies
\[
                  |\Re\rho-1/2|>a,
 \tag{MA-MO.553}
\]
then for every nonzero \(g\in C_c^\infty(\mathbb R)\)
the map (MA-MO.552) has no bounded extension \(H\to H\).
The conclusion holds for every polynomial exponent \(s\ge0\).
For a right-exterior zero there are a fixed
\(h\in C_c^\infty(\mathbb R)\) and \(t_n\to+\infty\)
such that the genuine smooth inputs have
\[
       \frac{\|\mathscr D^{\,\rm raw}_{\eta,g}U_{t_n}h\|_H}
            {\|U_{t_n}h\|_H}\longrightarrow\infty .
 \tag{MA-MO.554}
\]
For a left-exterior zero the corresponding reflected
translation sequence gives the conclusion. Both inputs
and outputs in these sequences belong to the original
smooth common sector.

In particular, existence of a bounded extension for even
one nonzero test implies that the entire nontrivial divisor
lies in the \emph{closed} strip
\[
                   |\Re\rho-1/2|\le a .
 \tag{MA-MO.555}
\]
No conclusion about zeros on its boundary is asserted.
There is no exception when \(\mathcal Mg(\rho-1/2)=0\).
The theorem also applies to nonzero tests with
\(g(0)=\mathcal Mg(1/2)=\mathcal Mg(-1/2)=0\).

\begin{proof}
\emph{The source domain before any narrow estimate.}
For \(y\in\mathcal E\), the series (MA-MO.551) and
every derivative converge locally absolutely. Its right
tail and all derivatives decrease with every exponential
weight, by the wide Euler estimates at arbitrarily large
weights. Multiplication by \(q\) kills the other tail.
Thus \(qF_{\omega,y}\in\mathcal E_\omega\), without
moving any inverse contour into the critical strip.
The actual Schwartz-cup inclusion (MA-MO.476) puts it
in \(\mathcal T_\omega\). The two Fourier maps in
(MA-MO.552) consequently act on their original domains.
Nontrivial-sector Poisson summation puts
\(\vartheta_\eta H_{\eta,y}\) in \(I_\eta\subset\mathcal E\).
Therefore (MA-MO.552) is a genuine map
\(\mathcal E\to\mathcal E\) independently of the bounded
extension under investigation.

\emph{Support and elementary wide estimates.}
For now suppose
\(\lambda=\rho-1/2\) satisfies \(\Re\lambda>a\).
Use bilinear, complex-linear output functionals
\[
                  \ell_v(f)=\int f(x)v(x)\,dx,
                         \qquad v\in C_c^\infty(\mathbb R).
 \tag{MA-MO.556}
\]
Each is continuous on \(H\) by weighted Cauchy--Schwarz.
There is no complex conjugation of \(g\) in its transpose:
\[
 \ell_v(C_gf)=\ell_{v_g}(f),\qquad
 v_g=C_{g^\vee}v,\qquad g^\vee(t)=g(-t).
 \tag{MA-MO.557}
\]
The choices of \(h\) and \(v\) that detect \(\lambda\)
will be made below; first take any compact smooth \(h,v\).
Write
\[
 \begin{gathered}
 \operatorname{supp}h\subset[\alpha,\beta],\qquad
 \operatorname{supp}v_g\subset[c,d],\qquad
 \phi(x)=1\quad(x\le\ell),\\
 F=F_{\eta,h},\quad
 H_h(w)=\int h(x)e^{wx}\,dx,\quad
 V_g(w)=\int v_g(x)e^{wx}\,dx,\\
 W(y)=\sum_{\mathfrak a\ {\rm prime\ to}\ S}
       c_{\eta,\mathfrak a}(N\mathfrak a)^{-1/2}
                            v_g(y-\log N\mathfrak a).
 \end{gathered}
 \tag{MA-MO.558}
\]
Both sums defining \(F\) and \(W\) are locally finite
because the two tests are compactly supported and there
are only finitely many integral ideals of bounded norm.
In particular
\[
           F(x)=0\ (x>\beta),\qquad
           W(y)=0\ (y<c).
 \tag{MA-MO.559}
\]
For every fixed \(b>1/2\), absolute ideal Euler convergence
also gives
\[
             |F(x)|\le C_b e^{-bx},\qquad
             |W(y)|\le C_b' e^{by}.
 \tag{MA-MO.560}
\]
For example, factor \(e^{-bx}(N\mathfrak a)^{-b}\)
out of \(h(x+\log N\mathfrak a)\); the remaining sum
is bounded by \(\zeta_K(b+1/2)\). The other estimate
uses \(e^{-b(y-\log N\mathfrak a)}v_g(y-\log N\mathfrak a)\).
The same argument works for every derivative.

Choose \(t_0>0\) large enough that
\[
                 -t-\alpha<\ell,\qquad t+\alpha>d
                              \quad(t\ge t_0).
 \tag{MA-MO.561}
\]
The support of \(RU_th\) is
\([-t-\beta,-t-\alpha]\). For \(x>\ell\), every
argument \(x+\log N\mathfrak a\) is greater than that
support. Hence the literal inverse series proves
\[
             qF_{\bar\eta,RU_th}=0\qquad(t\ge t_0).
 \tag{MA-MO.562}
\]
This is exact vanishing of the dual source term before
Fourier transform, not an asymptotic estimate. The second
condition in (MA-MO.561) gives
\(\ell_{v_g}(U_th)=0\). Translation commutes with the
original inverse series, so
\(F_{\eta,U_th}(y)=F(y-t)\). Pairing the remaining
actual theta sum against \(v_g\) now gives
\[
 \begin{aligned}
 f(t)&=\ell_v(\mathscr D^{\,\rm raw}_{\eta,g}U_th)\\
     &=\int q(y)W(y)F(y-t)\,dy,\qquad t\ge t_0 .
 \end{aligned}
 \tag{MA-MO.563}
\]
For each such \(t\), the cutoff inverse is supported
between \(\ell\) and \(t+\beta\). Only finitely many
ideal shifts can meet the compact output test, so the
change of variables in this pairing is justified by
finite sums. It is the actual theta action of MA-CC.8.
The formula also shows that \(f\) is continuous on
every compact \(t\)-interval.

\emph{The actual paired Laplace transform.}
For \(\Re w>1/2\), define
\[
 \begin{aligned}
 \mathfrak F(w)&=\int_{t_0}^\infty e^{-wt}f(t)\,dt,\\
 B_v(w)&=\int \phi(y)W(y)e^{-wy}\,dy,\\
 \mathcal R_{h,v}(w)
    &=\int q(y)W(y)e^{-wy}
          \int_{y-t_0}^\infty e^{wu}F(u)\,du\,dy .
 \end{aligned}
 \tag{MA-MO.564}
\]
The first integral converges absolutely on that half-plane.
Indeed choose \(1/2<b<\Re w\) in (MA-MO.560). The
integrand in (MA-MO.563) is bounded by a constant times
\(e^{bt}\), and its \(y\)-interval has length at most
\(t+C\). This gives an integrable bound
\((t+C)e^{-(\Re w-b)t}\) for the double integral and
justifies Fubini.

Both \(B_v\) and \(\mathcal R_{h,v}\) are entire.
The function \(\phi W\) is compactly supported: \(W\)
is zero to the left of \(c\), and \(\phi\) is zero
sufficiently far right. For the second assertion put
\(y_0=\max(\ell,c)\). The nonzero part of the remainder
lies in
\[
       y_0\le y\le\beta+t_0,\qquad
                    y-t_0\le u\le\beta .
 \tag{MA-MO.565}
\]
Thus it is a compact double integral of a smooth
function times \(e^{w(u-y)}\), which proves entire
dependence with all derivatives. An empty interval
contributes zero.

Absolute Euler convergence on \(\Re w>1/2\) gives
\[
       \int F(u)e^{wu}\,du=\frac{H_h(w)}{A_\eta(w)},
       \qquad
       \int W(y)e^{-wy}\,dy=A_\eta(w)V_g(-w).
 \tag{MA-MO.566}
\]
The coefficients in the first equality are the literal
M\"obius coefficients; those in the second are the
literal forward theta coefficients. From
\[
 \int_{t_0}^\infty e^{-wt}F(y-t)\,dt
       =e^{-wy}\int_{-\infty}^{y-t_0}e^{wu}F(u)\,du
\]
and \(q=1-\phi\), equations (MA-MO.563)--(MA-MO.566)
therefore give the exact identity
\[
 \mathfrak F(w)
   =H_h(w)V_g(-w)
       -\frac{H_h(w)B_v(w)}{A_\eta(w)}
       -\mathcal R_{h,v}(w),
                    \qquad\Re w>1/2 .
 \tag{MA-MO.567}
\]
No narrow reciprocal multiplier is used to obtain it.

\emph{A cutoff coefficient survives every nonzero smoothing.}
At the fixed exterior point define
\[
 b_\lambda(x)=e^{-\lambda x}
       \sum_{\mathfrak a\ {\rm prime\ to}\ S}
         c_{\eta,\mathfrak a}(N\mathfrak a)^{-1/2-\lambda}
                                \phi(x+\log N\mathfrak a).
 \tag{MA-MO.568}
\]
The sum is locally finite. Let
\(r=\sup\{x:\phi(x)\ne0\}<\infty\). The function is
zero for \(x\ge r\). If \(r-\log2<x<r\), every
nonunit ideal has \(x+\log N\mathfrak a>r\); its
contribution is zero. Thus
\[
              b_\lambda(x)=e^{-\lambda x}\phi(x)
                              \quad(r-\log2<x<r).
 \tag{MA-MO.569}
\]
By the definition of \(r\), \(\phi\) is not identically
zero on that interval. Consequently \(b_\lambda\ne0\).

There is a direct exponential-order bound on its left
tail. Choose
\(\sigma_0>\max(1,1/2+\Re\lambda)\). For \(x\le r\),
only \(N\mathfrak a\le e^{r-x}\) contribute, so
\[
 \begin{aligned}
 |b_\lambda(x)|
 &\le C e^{-\Re\lambda x}
   e^{(r-x)(\sigma_0-1/2-\Re\lambda)}
                         \zeta_K(\sigma_0)\\
 &\le C_\lambda e^{-(\sigma_0-1/2)x}.
 \end{aligned}
 \tag{MA-MO.570}
\]
Its one-sided Mellin transform is therefore holomorphic
on \(\Re z>\sigma_0-1/2\). If \(C_gb_\lambda\) were
zero, absolute convolution on that half-plane would give
\[
                  \mathcal Mg(z)\,\mathcal Mb_\lambda(z)=0.
 \tag{MA-MO.571}
\]
The first factor is a nonzero entire function, by
Fourier uniqueness for the nonzero compact test \(g\).
The second would then vanish identically on that
half-plane by the identity theorem. On a fixed vertical
line there, Fourier uniqueness for the integrable function
\(e^{(\Re z)x}b_\lambda(x)\) would imply
\(b_\lambda=0\), a contradiction. Hence
\(C_gb_\lambda\ne0\), for arbitrary complex \(g\ne0\).

Choose \(v\in C_c^\infty\) with
\(\int v(x)C_gb_\lambda(x)\,dx\ne0\). Such a test
exists because the convolution is a nonzero smooth
function; one may multiply its complex conjugate by
a nonnegative compact bump in a region where it is
nonzero. Finite summation on the compact support of
\(v_g\) and the bilinear transpose identity give
\[
            B_v(\lambda)
               =\int v_g(x)b_\lambda(x)\,dx
               =\int v(x)C_gb_\lambda(x)\,dx\ne0 .
 \tag{MA-MO.572}
\]
This is why a possible zero of \(\mathcal Mg(\lambda)\)
does not remove the obstruction: the cutoff function
\(b_\lambda\), with its last support edge, is not a
pure exponential.

Finally choose a nonnegative nonzero
\(h_0\in C_c^\infty\) and \(h(x)=e^{-\lambda x}h_0(x)\).
Then \(H_h(\lambda)=\int h_0>0\). The preceding
identities hold for these \(h,v\), with \(t_0\)
increased to satisfy (MA-MO.561) if necessary.

\emph{A bounded translated orbit contradicts the actual pole.}
The elementary translation estimate in the exact sum norm is
\[
 \|U_th\|_H
    \le 2^{s/4}(1+t^2)^{s/4}e^{at}\|h\|_H
                                  \qquad(t\ge0).
 \tag{MA-MO.573}
\]
If (MA-MO.552) had a bounded \(H\)-extension, continuity
of \(\ell_v\) would imply
\(|f(t)|\le C e^{at}(1+t)^{s/2}\). Dominated
differentiation would then make \(\mathfrak F(w)\)
holomorphic throughout \(\Re w>a\). All other
nonreciprocal factors in (MA-MO.567) are entire.
Multiplying that identity by \(A_\eta(w)\) and using
the identity theorem on the connected half-plane
\(\Re w>a\) gives
\[
 A_\eta(w)\bigl[
      H_h(w)V_g(-w)-\mathcal R_{h,v}(w)-\mathfrak F(w)
            \bigr]=H_h(w)B_v(w)
                      \qquad(\Re w>a).
 \tag{MA-MO.574}
\]
At \(w=\lambda\), its left side is zero and its right
side is nonzero by construction. The finite Euler
deletions have no zeros on this half-plane, so a
completed zero of any multiplicity gives an actual
zero of \(A_\eta\) here. Equivalently the quotient
in (MA-MO.567) has a pole of that full multiplicity;
the numerator has been chosen nonzero. Thus no
higher-order cancellation has been assumed away.

The same argument shows more directly that
\[
       \sup_{t\ge t_0}
          \frac{|f(t)|}{e^{at}(1+t)^{s/2}}=\infty .
 \tag{MA-MO.575}
\]
Otherwise its finite bound alone would justify the same
holomorphic continuation and contradiction. Since \(f\)
is continuous on compact intervals, one can choose
\(t_n\to\infty\) for which that ratio tends to infinity.
Weighted Cauchy--Schwarz for \(\ell_v\) and
(MA-MO.573) imply (MA-MO.554).

\emph{The left side and the exact remaining boundary scope.}
The actual inverse Fourier pair and zero-singular
Poisson identity in (MA-MO.552) give, on the original
smooth sources,
\[
 \mathfrak S_\eta H_{\eta,y}=H_{\bar\eta,Ry},
 \qquad
 R\vartheta_\eta=\vartheta_{\bar\eta}\mathfrak S_\eta .
\]
Together with \(RC_g=C_{g^\vee}R\), these identities give
\[
           R\mathscr D^{\,\rm raw}_{\eta,g}R
                 =\mathscr D^{\,\rm raw}_{\bar\eta,g^\vee}.
 \tag{MA-MO.576}
\]
This uses both source sectors, not complex conjugation
of a vector. Reflection is an isometry of the sum
norm. A left-exterior completed zero of \(\eta\)
gives a right-exterior completed zero of \(\bar\eta\)
by the functional equation. Apply the proved case to
\(\bar\eta,g^\vee\) and then reflect the genuine inputs.
This proves (MA-MO.553)--(MA-MO.555).

Every part of the nonannihilation argument required
only \(g\ne0\). None required a nonzero value of its
transform at \(\lambda\), at either polar exponent,
or a nonzero \(g(0)\). The claimed test subclasses
therefore follow.
\end{proof}

\paragraph{An every-width operator criterion, with its arithmetic input exposed.}
For this fixed nontrivial character the following statements
are equivalent:
\[
 \begin{array}{ll}
 {\rm (i)}&
   \Re\rho=1/2\quad\text{for every nontrivial completed zero};\\
 {\rm (ii)}&
   \forall\,0<a<1/2,\ \exists s\ge0,\
                          \exists\,0\ne g\in C_c^\infty:\\
 &\quad\mathscr D^{\,\rm raw}_{\eta,g}
       \text{ has a bounded }H_{\Sigma,a,s}\text{ extension};\\
 {\rm (iii)}&
   \forall\,0<a<1/2,\ \forall s\ge0,\
                              \forall g\in C_c^\infty:\\
 &\quad\mathscr D^{\,\rm raw}_{\eta,g}
       \text{ has an }\mathcal S_1(H_{\Sigma,a,s})\text{ extension}.
 \end{array}
 \tag{MA-MO.577}
\]
The first statement and conjugation give the premise of
MA-MO-022 for both characters with \(a_*=0\), proving
the third. The third implies the second. The second
and (MA-MO.555), for arbitrarily small \(a\), force
the first. This is an actual original-operator criterion;
it asserts none of these arithmetic antecedents anew.
The second statement can also restrict its nonzero tests
by \(g(0)=\mathcal Mg(\pm1/2)=0\). That class is
nonempty: for nonzero \(h_0\in C_c^\infty((1,2))\),
\(g=(\partial_x^2-1/4)h_0\) has these properties and
cannot be zero by the elementary constant-coefficient ODE.

At a fixed width, containment of the entire divisor in
a strictly smaller closed strip is sufficient by
MA-MO-022, and (MA-MO.555) is necessary.
The present Laplace argument is holomorphic only on
\(\Re w>a\). It does not settle zeros on the boundary,
their polynomial-weight thresholds, or sufficiency when
zeros approach an edge without a uniform margin.

\paragraph{Original-source and framework placement.}
The only source identities used are the literal ideal
Euler sums and inverse coefficients (MA-CC.1 and .8),
the structural coinvariant inclusion (MA-MO.470,.476),
the original \(p_+\) formula, and nontrivial-sector
Poisson summation. Their primary locations are Meyer
(2005), supplied \texttt{main.tex} lines 3592--3707,
2630--2675, and 3808--4043. The proof uses the
unconditional source identities from MA-MO-022; it does
not import that theorem's conditional narrow inverse
identification. The finite ideal sums and the cleared
Laplace identity above provide the new comparison.

No existence of a strictly exterior arithmetic zero is
asserted. The theorem proves what such a zero would do
to these actual source maps, rather than only to an
uncontinued scalar trace or a bare reciprocal multiplier.
The source restriction and quotient-return maps in
MA-MO-012/017 and MA-PR remain different domains.
In the inherited cote 1 framework, scans 079--080 and
111--112 through BR-TR-01/02, an exterior divisor factor
here prevents a bounded realization of this particular
full source action. Its finite-length algebraic class
does not repair the missing analytic norm estimate.
No positivity, additional motivic weight, or independent
FT002 deficiency vanishing follows from (MA-MO.577).

\section{Boundary multiplicity obstructs the actual character source}
\label{sec:character-boundary-obstruction}

The strict-exterior result MA-MO-024 leaves open what happens
on the edge itself. Integrating its actual translated source
orbit gives a stronger estimate than testing one translation
at a time. At the equality threshold the decisive convergence
is weak convergence in the original Hilbert norm.

Retain the nontrivial character, structural place set, actual
Fourier pair, real cutoff, and original smooth source action
of (MA-MO.550)--(MA-MO.552). In particular
\[
 \begin{gathered}
 H=H_{\Sigma,a,s}=L^2(\mathbb R,w_{a,s}(x)\,dx),
       \qquad 0<a<1/2,\quad s\ge0,\\
 w_{a,s}(x)=(e^{2ax}+e^{-2ax})(1+x^2)^{s/2},\qquad
 U_th(x)=h(x-t),\\
 D_g^{\,0}=\mathscr D^{\,\rm raw}_{\eta,g}:
                         \mathcal E\longrightarrow\mathcal E .
 \end{gathered}
 \tag{MA-MO.600}
\]
The last map is defined on the actual smooth common sector
by the original wide inverse and the original source Fourier
maps. No bounded narrow extension is assumed in its definition.

\paragraph{MA-MO-025 (the exact deficient-boundary obstruction).}
Let \(\rho\) be a completed nontrivial zero of \(L_K(s,\eta)\),
of multiplicity \(m\), such that
\[
             |\Re\rho-1/2|=a,\qquad s\le 2m-1 .
 \tag{MA-MO.601}
\]
For every nonzero complex \(g\in C_c^\infty(\mathbb R)\),
the actual map \(D_g^{\,0}\) has no bounded extension
\(H\to H\). The conclusion includes the equality
\(s=2m-1\), and has no exception when the Mellin transform
of \(g\) vanishes at the zero. It also holds within the
nonzero test class \(g(0)=\mathcal Mg(\pm1/2)=0\).

Consequently, boundedness of this actual full source action
for even one nonzero compact test requires full retention
of every nontrivial zero, with its complete multiplicity,
in the weighted partial-strip quotient. This is a necessary
condition. No fixed-width sufficiency without an interior
margin is asserted.

\begin{proof}
\emph{A boundary coefficient that survives the smoothing.}
First suppose \(\lambda=\rho-1/2=a+i\gamma\).
Use the notation \(A_\eta,H_h,V_g,B_v,\mathcal R_{h,v}\)
from (MA-MO.551) and (MA-MO.558)--(MA-MO.567).
The nonannihilation proof (MA-MO.568)--(MA-MO.572)
applies at this boundary point as well: none of its steps
requires the strict inequality \(\Re\lambda>a\).
For clarity, the function used there is
\[
 b_\lambda(x)=e^{-\lambda x}
     \sum_{\mathfrak a\ {\rm prime\ to}\ S}
       c_{\eta,\mathfrak a}(N\mathfrak a)^{-1/2-\lambda}
                              \phi(x+\log N\mathfrak a).
 \tag{MA-MO.602}
\]
It vanishes to the right of the last support edge \(r\)
of \(\phi\). In the interval \(r-\log2<x<r\), the
unit ideal alone contributes, so
\(b_\lambda=e^{-\lambda x}\phi(x)\) is not identically
zero. For any
\(\sigma_0>\max(1,1/2+\Re\lambda)\), the finite ideal
sum gives the left bound
\(\lvert b_\lambda(x)\rvert\le
C_\lambda e^{-(\sigma_0-1/2)x}\).
Its one-sided Mellin transform therefore exists on a
nonempty right half-plane. If \(C_gb_\lambda=0\), the
product of that transform and the nonzero entire function
\(\mathcal Mg\) would vanish there. Holomorphic uniqueness
and then Fourier uniqueness would give \(b_\lambda=0\),
a contradiction. This argument permits arbitrary complex
\(g\ne0\).

Choose \(v\in C_c^\infty\) detecting \(C_gb_\lambda\),
and choose \(h=e^{-\lambda x}h_0(x)\), where
\(0\ne h_0\in C_c^\infty\) is nonnegative. The bilinear
transpose \(v_g=C_{g^\vee}v\), without complex conjugation
of \(g\), gives
\[
 B_v(\lambda)=\int v(x)C_gb_\lambda(x)\,dx\ne0,\qquad
 H_h(\lambda)=\int h_0(x)\,dx>0 .
 \tag{MA-MO.603}
\]
Fix \(\operatorname{supp}h\subset[\alpha,\beta]\).
Increase \(t_0>0\) so that all support conditions
(MA-MO.561) hold and \(\beta+t_0\ge1\).
The actual dual source term then vanishes identically
for \(U_th\), \(t\ge t_0\), before its Fourier map is
applied. With the continuous output functional
\(\ell_v(y)=\int y(x)v(x)\,dx\), put
\[
 f(t)=\ell_v(D_g^{\,0}U_th),\qquad
 \mathfrak F(w)=\int_{t_0}^{\infty}e^{-wt}f(t)\,dt
                                      \quad(\Re w>1/2).
 \tag{MA-MO.604}
\]
The original source sums, the finite theta pairing, and
absolute wide Euler convergence give (MA-MO.567).
In particular its remainder \(\mathcal R_{h,v}\) and
both numerator factors in (MA-MO.603) are entire.

\emph{What a bounded extension would imply.}
Suppose that a bounded \(D:H\to H\) agrees with
\(D_g^{\,0}\) on \(\mathcal E\).
The exact weighted translation bound is
\[
 \|U_th\|_H\le
       2^{s/4}(1+t^2)^{s/4}e^{at}\|h\|_H
                                             \qquad(t\ge0).
 \tag{MA-MO.605}
\]
It follows that \(f(t)=O(e^{at}(1+t)^{s/2})\), and
dominated differentiation makes the actual integral
\(\mathfrak F\) holomorphic on \(\Re w>a\).
The cleared identity (MA-MO.574), initially proved by
wide Euler sums, extends by holomorphic uniqueness:
\[
 A_\eta(w)\bigl[
 H_h(w)V_g(-w)-\mathcal R_{h,v}(w)-\mathfrak F(w)
                 \bigr]=H_h(w)B_v(w),\qquad \Re w>a .
 \tag{MA-MO.606}
\]
There is no use here of a reciprocal multiplier on the
boundary line.

The finite Euler deletions have no zeros when
\(\Re(1/2+w)>0\); the gamma factors are finite and
nonzero at the present completed zero. Thus \(A_\eta\)
has exactly the same multiplicity \(m\) at \(\lambda\).
Write
\[
 A_\eta(w)=a_m(w-\lambda)^m+
                      O((w-\lambda)^{m+1}),\qquad
 a_m=\frac{A_\eta^{(m)}(\lambda)}{m!}\ne0 .
 \tag{MA-MO.607}
\]
In a sufficiently small disk \(\lambda\) is its only
zero. Evaluating (MA-MO.606) at
\(w_\delta=\lambda+\delta\), \(0<\delta<\delta_0\),
and using the nonzero numerator in (MA-MO.603), gives
\[
 \lim_{\delta\downarrow0}
       \delta^m\mathfrak F(\lambda+\delta)
       =-\frac{H_h(\lambda)B_v(\lambda)}{a_m}\ne0 .
 \tag{MA-MO.608}
\]
The other two terms in the meromorphic expression
(MA-MO.567) are holomorphic near \(\lambda\), so they
do not affect this limit. This retains the full
multiplicity of the zero.

\emph{An actual Hilbert orbit integral and its original-core
approximants.}
For each fixed \(\delta>0\), define
\[
 \begin{aligned}
 u_{\delta,T}
     &=\int_{t_0}^{T}e^{-w_\delta t}U_th\,dt,
                         &&T>t_0,\\
 u_\delta
     &=\int_{t_0}^{\infty}e^{-w_\delta t}U_th\,dt .
 \end{aligned}
 \tag{MA-MO.609}
\]
These are Bochner integrals in \(H\).
Translations are strongly continuous in this exact norm:
first verify this on compact smooth functions, and then
use their density and the locally uniform operator bound
(MA-MO.605). The infinite integral converges because
its norm integrand is bounded by
\(C e^{-\delta t}(1+t)^{s/2}\).
For finite \(T\), differentiation under the compact
parameter integral shows
\[
 u_{\delta,T}\in C_c^\infty
                 ([\alpha+t_0,\beta+T])\subset\mathcal E,
 \qquad
 u_{\delta,T}\longrightarrow u_\delta\quad\hbox{in }H .
 \tag{MA-MO.610}
\]
Thus the approximating inputs lie in the original
source domain, with no inverse continuation.
Since \(D\) is assumed bounded, it commutes with these
Bochner integrals and agrees with the actual source map
on every integrand and every finite approximant.
Applying the fixed continuous functional yields
\[
 \begin{aligned}
 \ell_v(Du_{\delta,T})
       &=\int_{t_0}^{T}e^{-w_\delta t}f(t)\,dt,\\
 \ell_v(Du_\delta)&=\mathfrak F(w_\delta).
 \end{aligned}
 \tag{MA-MO.611}
\]
We do not assert that the uncontinued raw source formula
is independently defined on \(u_\delta\). Only the
hypothetical bounded extension is applied to that limit.

\emph{The exact tail in the inherited sum norm.}
Changing variables in the orbit integral gives
\[
 u_\delta(x)=e^{-w_\delta x}
       \int_{-\infty}^{\,x-t_0}e^{w_\delta y}h(y)\,dy .
 \tag{MA-MO.612}
\]
This pointwise identity also follows from the finite
integrals and their local uniform convergence.
In particular
\[
 \begin{cases}
 u_\delta(x)=0,&x\le\alpha+t_0,\\
 u_\delta(x)=H_h(w_\delta)e^{-w_\delta x},
                                      &x\ge\beta+t_0 .
 \end{cases}
 \tag{MA-MO.613}
\]
On the fixed interval between these two endpoints,
\(u_\delta\) is uniformly bounded for
\(0<\delta\le\delta_0\); the same is true on every
fixed compact set. Put \(R_0=\beta+t_0\ge1\).
Since \(H_h(w_\delta)\) is bounded for small \(\delta\),
the exact right-tail weight gives
\[
 \begin{aligned}
 \|u_\delta\|_H^2
 &\le C+
 C\int_{R_0}^{\infty}e^{-2\delta x}
             (1+e^{-4ax})(1+x^2)^{s/2}\,dx\\
 &\le C_s\delta^{-(s+1)},\qquad 0<\delta\le\delta_0\le1 .
 \end{aligned}
 \tag{MA-MO.614}
\]
For the last bound, use
\((1+x^2)^{s/2}\le 2^{s/2}x^s\) for \(x\ge1\),
and integrate \(e^{-2\delta x}x^s\) over the positive
half-line. This computation retains both exponential
terms of the sum norm.

If \(s<2m-1\), (MA-MO.614) implies
\[
       \|\delta^m u_\delta\|_H
          \le C\delta^{\,m-(s+1)/2}\longrightarrow0 .
 \tag{MA-MO.615}
\]
Boundedness of \(D\) and continuity of \(\ell_v\) force
\(\delta^m\mathfrak F(w_\delta)\to0\) by
(MA-MO.611), contradicting (MA-MO.608).

\emph{The equality threshold.}
If \(s=2m-1\), the vectors
\(z_\delta=\delta^m u_\delta\) are uniformly bounded
in \(H\), by (MA-MO.614).
They tend uniformly to zero on each fixed compact
interval, by (MA-MO.612)--(MA-MO.613), since \(m\ge1\).
For a compactly supported \(k\in H\), weighted
Cauchy--Schwarz on its support therefore gives
\(\langle z_\delta,k\rangle_H\to0\).
For arbitrary \(k\in H\), approximate \(k\) in \(H\)
by compactly supported vectors and use the uniform
bound on \(\|z_\delta\|_H\). Consequently
\[
                 \delta^m u_\delta
                       \rightharpoonup0\quad\hbox{in }H .
 \tag{MA-MO.616}
\]
The fixed functional \(\ell_v\circ D\) is continuous,
so its values on these vectors must tend to zero.
Equations (MA-MO.608) and (MA-MO.611) contradict this.
The equality case does not rely on norm convergence.

\emph{The other boundary.}
If \(\Re(\rho-1/2)=-a\), the functional equation gives
a zero of the same multiplicity for the dual character
at \(1-\rho\), whose centered real part is \(a\).
The exact source identity (MA-MO.576) is
\[
       R D_{\eta,g}^{\,0}R=D_{\bar\eta,g^\vee}^{\,0}.
 \tag{MA-MO.617}
\]
It retains the actual finite Fourier phases, and \(R\)
is an isometry of the norm (MA-MO.600).
A bounded extension on the left would therefore give
one for the right-boundary problem already contradicted.
Since \(g^\vee\ne0\), this proves the full statement.
The proof required only \(g\ne0\), so the stated moment
conditions on the test do not change the conclusion.
\end{proof}

\paragraph{The exact necessary divisor condition.}
In the centered sum norm, MA-WP-001 gives retained
multiplicity
\[
 \begin{aligned}
 b_s&=\#\{j\in\mathbb Z_{\ge0}:2j<s-1\},\\
 r_\rho&=
 \begin{cases}
 m_\rho,&|\Re\rho-1/2|<a,\\
 \min(m_\rho,b_s),&|\Re\rho-1/2|=a,\\
 0,&|\Re\rho-1/2|>a .
 \end{cases}
 \end{aligned}
 \tag{MA-MO.618}
\]
At an edge, full retention is equivalent to
\(2(m_\rho-1)<s-1\), or \(s>2m_\rho-1\).
Combining MA-MO-024 and MA-MO-025 therefore gives
\[
 \begin{gathered}
 \exists\,0\ne g\in C_c^\infty:
 D_g^{\,0}\text{ has a bounded extension on }H_{\Sigma,a,s}\\
 \Longrightarrow\quad
 r_\rho=m_\rho\quad\text{for every nontrivial completed zero}.
 \end{gathered}
 \tag{MA-MO.619}
\]
For \(0\le s\le1\), this excludes every boundary zero.
For general \(s\), a zero at the boundary is excluded
exactly when its local algebra is truncated by this
Hilbert norm. The word \emph{exactly} here describes the
proved obstruction threshold, not a sufficiency theorem.
Even when every \(r_\rho=m_\rho\), no boundedness or
trace-class estimate for the original source action is
deduced from (MA-MO.619).
The sufficient hypothesis in MA-MO-022 remains the
explicit zero-free half-plane with a strict margin.
The existing quotient traces and quotient-return maps
do not by themselves supply that missing full-source
operator estimate.

\paragraph{Source scope and finite-length placement.}
All arithmetic source identities are inherited from
the completely written MA-MO-024 proof: literal wide
ideal sums, actual structural-\(S\) source inclusion,
the original \(p_+\) map, and nontrivial-sector Poisson
summation. Their exact primary locators are Meyer
(2005), supplied \texttt{main.tex} lines 3592--3707,
2630--2675, and 3808--4043, through
(MA-MO.551)--(MA-MO.567) and (MA-MO.576).
The new argument is the actual Hilbert orbit integral,
its compact original-core approximants, and the
strong/weak limiting contradiction.
No additional zero-distribution theorem, nonvanishing
claim, or arithmetic occurrence of a boundary zero is
used.

In the inherited cote 1 trace framework, scans 079--080
and 111--112 through BR-TR-01/02, the local divisor has
finite length \(m_\rho\), while the completed quotient
retains length \(r_\rho\). These are classes in the
category of finite-length \(\mathbb C[z]\)-modules;
their underlying finite-dimensional vector spaces
carry the determinant constructions already used in
MA-WP. Equality of those lengths is necessary here
for a bounded realization of the original full source
action. It is not an analytic extension theorem, a
positivity statement, or a consequence of a motivic
weight. The newly separated boundary question is
decided at this stated threshold; the arithmetic
deficiency condition of FT002 is not established anew.

\section{The actual rational full matrix: exterior and boundary necessity}
\label{sec:polar-exterior-obstruction}

The rational completion in MA-MO-016 carries actual polar source
sections. A compactly supported functional of either represented
component is therefore bounded on that completion, although a bare
polar exponential does not lie in its diagonal Hilbert summand.
This observation transfers the actual-source obstruction to the
principal character without discarding either polar coordinate.

Keep \(S=\{\infty\}\), the arbitrary smooth real cutoff \(\phi\)
of MA-MO-018, and the original common source maps. Put
\[
 \mathcal W_{a,s}=H_{\Sigma,a,s}\oplus\mathbb C^2,
 \qquad 0<a<1/2,\quad s\geq0,\qquad
 \mathscr D_g^{\rm raw}=C_g(P_+-P_-).
 \tag{MA-MO.650}
\]
The raw map here is defined on \(\mathscr S_{\rm sum}\), with
its actual pair realization \(\mathfrak s\) from (MA-MO.252).
No zero-free half-plane is assumed in this section.

\paragraph{MA-MO-026 (necessary full-divisor retention in the polar completion).}
Let \(\rho\) be a nontrivial zeta zero of multiplicity \(m\).
For every nonzero \(g\in C_c^\infty(\mathbb R)\), the original
map \(\mathscr D_g^{\rm raw}\) has no bounded endomorphism
extension on \(\mathcal W_{a,s}\) if either
\[
 \begin{array}{ll}
 \text{(a)}&|\Re\rho-1/2|>a,\\
 \text{(b)}&|\Re\rho-1/2|=a\quad\text{and}\quad s\leq2m-1.
 \end{array}
 \tag{MA-MO.651}
\]
This includes tests satisfying
\(g(0)=\mathcal Mg(1/2)=\mathcal Mg(-1/2)=0\), and tests
whose Mellin transform vanishes at \(\rho-1/2\). In case (a)
there is a fixed compact smooth negative source \(i_-h\) and
\(t_n\to+\infty\) such that
\[
 \frac{\|\mathscr D_g^{\rm raw}i_-U_{t_n}h\|_{\mathcal W_{a,s}}}
      {\|i_-U_{t_n}h\|_{\mathcal W_{a,s}}}\longrightarrow\infty.
 \tag{MA-MO.652}
\]
Every displayed input and output belongs to the original smooth
common source. In case (b) the contradiction uses convergent
integrals of these same actual inputs, and compact smooth
approximations to those integrals.

\begin{proof}
\emph{The bounded output functional and the literal rational source.}
For \(v\in C_c^\infty(\mathbb R)\) define
\[
 \Lambda_v(u,c)=\int v(x)
       \{u(x)+c_0\Theta f_0(x)+c_1\Theta f_1(x)\}\,dx.
 \tag{MA-MO.653}
\]
This is the bilinear integral of \(v\) against the first actual
pair component of \(\mathfrak s(u,c)\). It is bounded on
\(\mathcal W_{a,s}\): the diagonal term is bounded by weighted
Cauchy--Schwarz, and the two fixed smooth section functions are
integrable on the compact support of \(v\). The polar terms
are included in this functional, not subtracted from the source.
Also \(i_-y=(y,0)\) is an isometry from \(H_{\Sigma,a,s}\)
into \(\mathcal W_{a,s}\). These statements follow directly
from the actual completion, without its conditional matrix theorem.

For a compact smooth \(h\), let
\[
 F(x)=\sum_{n\geq1}\mu(n)n^{-1/2}h(x+\log n),
 \qquad A(w)=\zeta(1/2+w),\qquad q=1-\phi.
 \tag{MA-MO.654}
\]
This is the original wide inverse. It is locally a finite sum,
has a finite right support endpoint, and obeys the absolute
wide bounds (MA-MO.560), with \(\zeta(b+1/2)\) in place of
the ideal majorant. The actual additive source for \(i_-y\)
is (MA-MO.325--326), namely
\[
 u_y(t)=\frac{q(\log|t|)F_y(\log|t|)}{2\sqrt{|t|}},
 \qquad h_y=u_y+\mathfrak F u_{Ry}.
 \tag{MA-MO.655}
\]
For compact \(y\), \(qF_y\) has compact logarithmic support,
so \(u_y\) is an actual compact annular even Schwartz function.
The source formula (MA-MO.655) consequently holds before any
narrow contour shift. The conditional estimates in MA-MO-018
are not being imported here. For general \(y\in\mathcal E\),
the arbitrarily rapid right tail and the literal wide estimates
still make \(u_y\) Schwartz and zero near additive zero;
compact right support is not needed. Thus the original maps
act on their stated smooth common source without a narrow
boundedness premise.

The normalization of the actual theta map gives exactly
\[
 \Theta u_y(x)=\sum_{n\geq1}n^{-1/2}
                   q(x+\log n)F_y(x+\log n).
 \tag{MA-MO.656}
\]
If \(y=U_th\) and \(t\) is sufficiently large, the support
argument (MA-MO.561--562) gives \(qF_{RU_th}=0\).
Hence \(h_{U_th}=u_{U_th}\) exactly, before Fourier
transformation. Since \(P_-i_-y=i_-y\), the first component
of the actual raw output is \(C_g(\Theta u_{U_th}-U_th)\).
With \(v_g=C_{g^\vee}v\) and
\[
 W(y)=\sum_{n\geq1}n^{-1/2}v_g(y-\log n),
 \qquad f(t)=\Lambda_v(\mathscr D_g^{\rm raw}i_-U_th),
\]
the literal finite-support calculation therefore gives
\[
                 f(t)=\int q(y)W(y)F(y-t)\,dy
                         \qquad(t\geq t_0).
 \tag{MA-MO.657}
\]
The subtracted input pairs to zero after increasing \(t_0\).
This establishes the exact hypotheses of the paired calculation
in MA-MO-024 on the actual polar completion.

\emph{The pole-cleared Laplace identity.}
Use the definitions \(H_h,V_g,B_v,\mathcal R_{h,v}\) of
(MA-MO.558,.564), now with the ordinary positive integers.
The proof (MA-MO.559--567) applies verbatim to these literal
sums: both \(B_v\) and \(\mathcal R_{h,v}\) are entire,
the latter is a compact double integral, and for \(\Re w>1/2\)
the absolutely convergent actual paired Laplace transform obeys
\[
 \mathfrak F(w)=H_h(w)V_g(-w)
             -\frac{H_h(w)B_v(w)}{A(w)}-\mathcal R_{h,v}(w).
 \tag{MA-MO.658}
\]
The only change from the nontrivial-character case is the pole
of \(A\) at \(w=1/2\). Clear it by the entire function
\(\widetilde A(w)=(w-1/2)A(w)\). Thus a bounded extension
of the raw matrix, which would force \(\mathfrak F\) to be
holomorphic on \(\Re w>a\), would imply
\[
 \widetilde A(w)
  [H_h(w)V_g(-w)-\mathcal R_{h,v}(w)-\mathfrak F(w)]
       =(w-1/2)H_h(w)B_v(w),\qquad \Re w>a.
 \tag{MA-MO.659}
\]
This follows by the identity theorem from the original wide
half-plane. Clearing the zeta pole does not change the order
of any nontrivial zero \(\lambda=\rho-1/2\), since
\(\lambda\ne1/2\).

At any such \(\lambda\) the nonannihilation construction
(MA-MO.568--572) also applies without change: its coefficients
are now \(n^{-1/2-\lambda}\), the term \(n=1\) is the sole
term near the last support edge, and the majorant is the
convergent ordinary zeta series. For every complex nonzero
\(g\) it therefore supplies a compact smooth \(v\) with
\(B_v(\lambda)\ne0\). Choose
\(h(x)=e^{-\lambda x}h_0(x)\) for a nonnegative nonzero
compact smooth \(h_0\), so \(H_h(\lambda)>0\).
No nonzero value of \(\mathcal Mg(\lambda)\) is required.

If \(\Re\lambda>a\), (MA-MO.659) is already a contradiction
at \(w=\lambda\). More precisely the argument of
(MA-MO.573--575), using the bounded functional (MA-MO.653)
and the isometric negative embedding, proves (MA-MO.652).
Every raw output is still a smooth common-source vector by
(MA-MO.655). This proves case (a) for a right exterior zero.

\emph{A boundary zero and its exact polynomial threshold.}
Suppose \(\lambda=a+i\gamma\) has multiplicity \(m\), and
assume a bounded extension exists. For \(\delta>0\), the
Bochner integral
\[
 y_\delta=\int_{t_0}^{\infty}
                   e^{-(\lambda+\delta)t}U_th\,dt
                    \quad\text{in }H_{\Sigma,a,s}
 \tag{MA-MO.660}
\]
converges because the translation norm is at most a constant
times \(e^{at}(1+t)^{s/2}\). Truncating the integral at a
finite upper endpoint gives a compact smooth original source;
the truncations converge in \(H\), and their negative embeddings
converge in \(\mathcal W\). Boundedness then justifies
\[
 \Lambda_v(\mathscr D_g i_-y_\delta)
                        =\mathfrak F(\lambda+\delta).
 \tag{MA-MO.661}
\]
If \(\operatorname{supp}h\subset[\alpha,\beta]\), direct
substitution in (MA-MO.660) gives
\[
 y_\delta(x)=e^{-(\lambda+\delta)x}
       \int_{-\infty}^{x-t_0}e^{(\lambda+\delta)u}h(u)\,du.
 \tag{MA-MO.662}
\]
It vanishes for \(x<\alpha+t_0\), is uniformly bounded with
all derivatives on each compact set as \(0<\delta\leq1\),
and for \(x\geq\beta+t_0\) equals
\(H_h(\lambda+\delta)e^{-(\lambda+\delta)x}\).
Since the exact sum weight is
\((e^{2ax}+e^{-2ax})(1+x^2)^{s/2}\), integration of this
tail gives
\[
                  \|y_\delta\|_H
                         =O(\delta^{-(s+1)/2}).
 \tag{MA-MO.663}
\]
Increase \(t_0\), if necessary, so that \(\beta+t_0\geq1\).
The contribution on the remaining fixed compact interval is
uniformly bounded. The tail bound follows, for example, by
\(u=2\delta x\) in the integral of
\(e^{-2\delta x}(1+x)^s\), with a harmless fixed lower
endpoint. The other exponential in the sum weight only improves
integrability.

If \(s<2m-1\), it follows that \(\delta^m y_\delta\to0\)
strongly in \(H\). If \(s=2m-1\), this family is bounded
and tends to zero in \(L^2\) on every compact set. Pairing
first with compactly supported \(H\)-vectors and then using
their density and the uniform bound shows
\(\delta^m y_\delta\rightharpoonup0\) weakly in \(H\).
In either case a fixed bounded functional composed with the
assumed bounded matrix gives
\[
       \delta^m\Lambda_v(\mathscr D_g i_-y_\delta)\longrightarrow0.
 \tag{MA-MO.664}
\]
But near \(\lambda\) the right side of (MA-MO.658), justified
there by (MA-MO.659), has a pole of order exactly \(m\): its
numerator \(H_h(\lambda)B_v(\lambda)\) is nonzero, and
\(\mathcal R_{h,v},H_h,V_g\) are entire. If
\(A(w)=(w-\lambda)^m a_0(w)\) with \(a_0(\lambda)\ne0\),
then (MA-MO.661) instead gives
\[
 \delta^m\Lambda_v(\mathscr D_g i_-y_\delta)
       \longrightarrow
       -\frac{H_h(\lambda)B_v(\lambda)}{a_0(\lambda)}\ne0.
 \tag{MA-MO.665}
\]
This contradicts (MA-MO.664), proving the boundary assertion.

Finally the zeta functional equation carries a left exterior
or left boundary zero to a right one with the same multiplicity.
Applying the proved right-side argument to that actual zero
proves all of (MA-MO.651). This use of the functional equation
asserts no new location of zeros. The construction of \(v\)
only required \(g\ne0\), which proves the stated test scopes.
\end{proof}

\paragraph{The exact arithmetic meaning of all widths.}
For this actual completion, the following are equivalent:
zeta RH; existence, at each \(0<a<1/2\), of some \(s\geq0\)
and some nonzero compact smooth test for which the full original
matrix extends boundedly; and trace-class extension for every
such width, weight and test. Necessity is (MA-MO.651)(a) at
arbitrarily small widths. Sufficiency is precisely MA-MO-018
with \(a_*=0\); it explicitly uses zeta RH. The test may also
be required to satisfy \(g(0)=\mathcal Mg(\pm1/2)=0\), since
the nonzero class constructed after (MA-MO.577) has those
properties. Under the user-specified theorem A, this arithmetic
premise is supplied. It remains a named zero-location input in
an audit of the method.

At a fixed width, bounded full-source action therefore requires
no exterior zero and retention of every boundary multiplicity
under the sharp MA-WP condition \(s>2m-1\). Sufficiency without
a strictly smaller containing strip is not proved here. The
result concerns the actual polar completion, not the full
coordinate product already excluded by MA-MO-005, nor the
restricted positive source column of MA-MO-016/023.

The source locators are Meyer (2005),
\texttt{lem:Hilm\_p\_embed} and its original \(p_\pm\) formulas,
supplied \texttt{main.tex} lines 3808--4043, and the actual
completion and source calculations (MA-MO.250--263,.323--326).
In Grothendieck's cote~1 framework, scans 079--080 and 111--112
through BR-TR-01/02, a trace or Fredholm determinant requires a
defined bounded nuclear realization. Adjoining the two actual
polar sections does not remove the analytic obstruction from an
unretained zero or jet. The finite-length \(K_0\) class records
that multiplicity; it does not supply the missing operator norm.
No faithful full-jet positivity, motivic purity, or independent
arithmetic vanishing theorem is asserted.

\section{Full common-space extension with finitely many retained boundary zeros}
\label{sec:finite-boundary-full-matrix}

The necessary boundary condition in MA-MO-025/026 has a matching
sufficient case. A finite boundary divisor can be treated by its
actual right boundary distributions, provided the remaining zeros
have a uniform interior margin. The polynomial weight then absorbs
the boundary poles at exactly the previously established jet
threshold. The construction extends the original wide-inverse
source matrix; it does not replace that inverse by an arbitrary
principal value.

Use either of the following already fixed source settings:
the nontrivial finite-order character sector of MA-MO-022, with
its number field, coherent Fourier normalizations and structural
place set (MA-MO.470); or the rational principal sector
\(K=\mathbb Q,S=\{\infty\}\) and the actual polar completion
of MA-MO-016/018. In the first case let
\(\omega\in\{\eta,\bar\eta\}\); in the second case there is
one self-dual multiplier. Write
\[
 \begin{gathered}
 0<a<1/2,\qquad s\ge0,\qquad r=s/2,\\
 H_{\pm a,s}
   =L^2\bigl(e^{\pm2ax}(1+x^2)^{s/2}\,dx\bigr),\\
 H=H_{\Sigma,a,s}=H_{a,s}\cap H_{-a,s},\qquad
 \|y\|_H^2=\|y\|_{H_{a,s}}^2+\|y\|_{H_{-a,s}}^2 .
 \end{gathered}
 \tag{MA-MO.700}
\]
The function \(\phi\) is the original real smooth step,
equal to one far left and zero far right, with
\(\phi'\in C_c^\infty\). Put \(q=1-\phi\). No symmetry
or Gevrey condition is imposed.

The new arithmetic hypothesis is the following finite-boundary
condition on each completed centered divisor:
\[
 \begin{gathered}
 |\Re\lambda|\le a\quad\hbox{for every completed zero }\lambda,\\
 \{\lambda:|\Re\lambda|=a\}\ \hbox{is finite},\\
 \text{for some }\delta>0,\qquad
 |\Re\lambda|<a\ \Longrightarrow\
                       |\Re\lambda|\le a-\delta .
 \end{gathered}
 \tag{MA-MO.701}
\]
Here \(\lambda=\rho-1/2\), with its full multiplicity
\(m_\omega(\lambda)\). The condition is required for both
dual characters in the nontrivial case. Shrinking \(\delta\)
allows \(0<\delta<a\). No claim that (MA-MO.701) holds
arithmetically is made.

\paragraph{MA-MO-027 (finite-boundary sufficiency in the actual norm).}
Assume (MA-MO.701) and
\[
        s>2m_\omega(\lambda)-1
               \quad\hbox{at every boundary zero }|\Re\lambda|=a .
 \tag{MA-MO.702}
\]
The condition is vacuous if the boundary divisor is empty.
For every \(g\in C_c^\infty(\mathbb R)\), the original
smoothed matrix \(C_g(P_+-P_-)\) has a unique trace-class
extension on the actual completion:
\[
 \begin{array}{ll}
 \mathscr D_{\eta,g}:H\longrightarrow H,
       &\text{nontrivial character sector},\\[2pt]
 \mathscr D_g:\mathcal W_{a,s}\longrightarrow\mathcal W_{a,s},
       \quad\mathcal W_{a,s}=H\oplus\mathbb C^2,
       &\text{rational principal sector}.
 \end{array}
 \tag{MA-MO.703}
\]
On each fixed compact test support the extensions depend
continuously on \(g\) in trace norm. The two polar coordinates
in the second line are the actual source-section coordinates,
not extra bare exponentials inserted into \(H\).

In the nontrivial case \(J_\eta=\overline{I_\eta}^{\,H}\)
is invariant, and the induced quotient action is \(-C_g\).
Writing \(B_{\eta,g}^{\partial}\) for the actual restriction
to this closed ideal, one has
\[
 \begin{aligned}
 \operatorname{Tr}_H\mathscr D_{\eta,g}
   &=\operatorname{Tr}_{J_\eta}B_{\eta,g}^{\partial}
                  -\sum_\rho m_\eta(\rho)G(\rho-1/2),\\
 \det_H(1-z\mathscr D_{\eta,g})
   &=\det_{J_\eta}(1-zB_{\eta,g}^{\partial})
                         \det_{H/J_\eta}(1+zC_g),
               \qquad G=\mathcal Mg .
 \end{aligned}
 \tag{MA-MO.704}
\]
This statement supplies the actual boundary extension and
its invariant-subspace identities. It does not identify
the first trace in (MA-MO.704) with a local scalar formula
by silently reusing a no-boundary inverse estimate.

Within the class (MA-MO.701), condition (MA-MO.702) is
also necessary for any nonzero compact smooth test to have
a bounded original full-matrix extension: this is exactly
MA-MO-025 for the nontrivial sector and MA-MO-026 for the
rational polar completion. Thus it characterizes bounded
full-matrix existence in this finite-boundary, uniform-margin
case. It is not a sufficiency theorem for an infinite divisor
approaching an edge without a uniform margin.

\begin{proof}
\emph{Boundary reciprocals and the exact weight threshold.}
Use the original forward multipliers
\[
 A_\omega(z)=
 \begin{cases}
 E_{\omega,S}(1/2+z)L_K(1/2+z,\omega),
                           &\text{nontrivial case},\\
 \zeta(1/2+z),              &\text{rational case},
 \end{cases}
 \qquad j_{\omega,\epsilon}(u)=A_\omega(a+\epsilon+iu)^{-1}.
 \tag{MA-MO.705}
\]
The finite Euler factors have their zeros on \(\Re z=-1/2\).
Choose \(0<\epsilon_0<1/2-a\) sufficiently small. A narrow
closed strip about the line \(\Re z=a\) contains no zeros
except the finitely many right boundary zeros
\(\lambda_j=a+i\gamma_j\). It also avoids the rational
pole at \(z=1/2\).

In a disk around \(\lambda_j\), write the genuine Laurent
principal part
\[
 A_\omega(z)^{-1}
       =\sum_{k=1}^{m_j}c_{j,k}(z-\lambda_j)^{-k}
                               +\text{holomorphic function},
               \qquad c_{j,m_j}\ne0 .
 \tag{MA-MO.706}
\]
Choose compact smooth functions \(\chi_j(u)\), equal to
one near \(\gamma_j\), with sufficiently small disjoint
supports. Subtracting these localized principal parts gives
\[
 j_{\omega,\epsilon}(u)=r_{\omega,\epsilon}(u)
    +\sum_j\sum_{k=1}^{m_j}
       c_{j,k}\chi_j(u)(\epsilon+i(u-\gamma_j))^{-k},
                   \qquad 0<\epsilon\le\epsilon_0 .
 \tag{MA-MO.707}
\]
Every derivative of \(r_{\omega,\epsilon}\) has a polynomial
bound in \(u\), uniform for these \(\epsilon\), and these
functions converge to a smooth \(r_{\omega,0}\) with the
same weighted derivative bounds. More precisely, for each
derivative there is a fixed inverse polynomial weight in
which this convergence is uniform; a further derivative
in the real part and the mean value theorem give this
assertion. Here is the exact
minimum-modulus input. Factoring the finitely many zeros
in (MA-MO.706) leaves a multiplier without zeros in a
fixed neighborhood of the positive line, with a fixed
horizontal margin to every remaining zero. The
fixed-gap proof of MA-RS-004, already used for each fixed
Hecke character in MA-MO-012/022, gives polynomial
reciprocal derivative bounds outside the finitely many
disks. The factors removed are polynomials, so preserve
such bounds. Inside the disks the Laurent remainders
are holomorphic on a slightly larger compact disk;
Cauchy estimates are uniform. The transition regions
of the \(\chi_j\) stay a positive distance from the
poles. These three regions prove the asserted estimates
and convergence. In the rational case the distant pole
of \(A\) is a zero of its reciprocal; it creates no
additional inverse singularity.

For each positive integer \(k\), the elementary Laplace
identity is
\[
 (\epsilon+iv)^{-k}
     =\frac1{(k-1)!}\int_0^\infty
                 t^{k-1}e^{-\epsilon t}e^{-ivt}\,dt .
 \tag{MA-MO.708}
\]
With the unitary plus Fourier transform, its transform
is \(\sqrt{2\pi}\,t^{k-1}e^{-\epsilon t}
\mathbf1_{t>0}/(k-1)!\). Consequently it converges as
\(\epsilon\downarrow0\) in \(H^{-r}(\mathbb R)\) if
and only if
\[
       \int_0^\infty t^{2k-2}(1+t^2)^{-r}\,dt<\infty,
               \qquad\text{equivalently }r>k-\tfrac12 .
 \tag{MA-MO.709}
\]
Necessity follows from the same nonnegative integral,
and sufficiency and convergence follow by dominated
convergence. Translation and multiplication by a fixed
compact smooth function are bounded on \(H^{-r}\).
Thus (MA-MO.702) gives the actual boundary distribution
\[
 j_{\omega,+}
    =r_{\omega,0}+
       \sum_{j,k}c_{j,k}\chi_j
             (0+i(\,\cdot-\gamma_j))^{-k}.
 \tag{MA-MO.710}
\]
Each singular summand has compact support and belongs to
\(H^{-r}\). Its sign and delta-derivative contributions
are fixed by the limit from \(\Re z>a\). No alternative
principal-value prescription has been chosen.

\emph{Identification with the original inverse before using a kernel.}
For \(y\in\mathcal E\), let \(F_{\omega,y}\) be the
original wide inverse series (MA-MO.477), or (MA-MO.323)
in the rational case. That series is locally absolutely
convergent with all derivatives. On every line \(b>a\)
its Mellin transform is the original wide quotient shifted
without crossing any zero:
\[
 \begin{gathered}
 e^{bx}F_{\omega,y}\in\mathscr S(\mathbb R),\qquad
 V_bF_{\omega,y}=j_{\omega,b-a}V_by,\qquad
 V_b=\mathcal F_+M_{e^{bx}},\\
                         a<b<a+\epsilon_0 .
 \end{gathered}
 \tag{MA-MO.711}
\]
Indeed such a line has a positive gap from every zero;
the already proved wide-inverse contour argument of
MA-MO-018/022 applies, with \(a_*=a<b\).
The same argument supplies every fixed right weight
\(b>a\), including arbitrarily large ones.

The transforms \(V_by\) converge in Schwartz space to
\(V_ay\). Equations (MA-MO.707)--(MA-MO.710) therefore
give convergence in tempered distributions, and
Fourier inversion gives
\[
        V_aF_{\omega,y}=j_{\omega,+}V_ay
                    \quad\hbox{as tempered distributions}.
 \tag{MA-MO.712}
\]
This is an identity with the original function: on every
compact \(x\)-interval,
\(e^{bx}F_{\omega,y}\to e^{ax}F_{\omega,y}\) smoothly,
so the distributional limit just obtained has that
local value. The compact singular distributions in
(MA-MO.710), after multiplication by \(V_ay\), have
finite order. Their inverse Fourier transforms, and
all derivatives, have at most polynomial growth.
The regular part times \(V_ay\) is Schwartz. This
also proves the needed polynomial growth of
\(e^{ax}F_{\omega,y}\), rather than assuming it
from a formal boundary integral.

Multiplying (MA-MO.707) by \(A_\omega(a+\epsilon+iu)\)
and taking the same distributional limit proves
\[
                    A_\omega(a+iu)j_{\omega,+}(u)=1 .
 \tag{MA-MO.713}
\]
The multiplication is legitimate because the forward
factor and all its derivatives have polynomial bounds;
on each compact singular support the factors converge
in every smooth norm. This identity is established
before any diagonal commutator cancellation.

Finally \(qF_{\omega,y}\in\mathcal E_\omega\):
it vanishes far left and has all right weighted
Schwartz estimates from (MA-MO.711).
For a nontrivial sector, the structural-\(S\)
inclusion (MA-MO.476) puts it in the actual local
\(\mathcal T_\omega\), with its source topology.
For the rational sector,
\[
       u_y(t)=\frac{q(\log|t|)F_y(\log|t|)}
                         {2\sqrt{|t|}},\qquad u_y(0)=0
 \tag{MA-MO.714}
\]
is the actual even additive Schwartz function,
zero near additive zero. Thus both subsequent
source Fourier transforms are on their original domains.

\emph{A trace-class lemma for the singular input distributions.}
We give the norm argument needed at the boundary.
If \(d\in H^{-r}(\mathbb R)\) has compact support
and \(L\in\mathscr S(\mathbb R^2)\), then
\[
       T_{d,L}f(t)=\langle d(u),L(t,u)f(u)\rangle
          \quad\hbox{defines }\quad
       T_{d,L}\in\mathcal S_1(H^r,H^R)
                    \quad\hbox{for every fixed }R\ge0 .
 \tag{MA-MO.715}
\]
The pairing is bilinear distributional pairing;
it is not a pointwise multiplication of two
undefined distributions. To prove the assertion,
expand in the normalized Hermite functions of
the output variable:
\(L(t,u)=\sum_{n\ge0}h_n(t)L_n(u)\).
Repeated use of the harmonic oscillator in \(t\)
shows that every Schwartz seminorm of \(L_n\)
decreases faster than every power of \(n+1\).
For a smooth multiplier \(v\), the Fourier
convolution formula and
\((1+|\xi+\eta|^2)^{r/2}
 \le C_r(1+|\xi|^2)^{r/2}(1+|\eta|^2)^{r/2}\)
give
\[
 \|vf\|_{H^r}
   \le C_r\bigl\|(1+|\xi|^2)^{r/2}\widehat v(\xi)
                      \bigr\|_{L^1}\|f\|_{H^r}.
 \tag{MA-MO.716}
\]
Hence the norms of
\(\ell_n(f)=\langle d,L_nf\rangle\) decrease rapidly.
The \(H^R\)-norms of \(h_n\) grow at most polynomially:
integer derivative norms follow from the two Hermite
ladder identities, and a larger integer derivative
norm bounds the fractional one. It follows that
\(\sum_n\|h_n\|_{H^R}\|\ell_n\|<\infty\).
This is an explicit nuclear series between Hilbert
spaces and proves (MA-MO.715).

The proof also proves continuity in \(d\) with its
\(H^{-r}\) norm and in finitely many Schwartz
seminorms of \(L\). In particular \(d_\epsilon\to d\)
in \(H^{-r}\) and \(L_\epsilon\to L\) in Schwartz
space imply convergence in trace norm. A smooth
polynomially bounded input factor times \(L\) is
an ordinary Schwartz kernel and has the same
conclusion. Finally \(V_a\) identifies the actual
weighted norm with \(H^r\):
\[
           \|y\|_{H_{a,s}}
                =\|(1-\partial_t^2)^{r/2}V_ay\|_2,
                          \qquad r=s/2 .
 \tag{MA-MO.717}
\]
Thus the lemma concerns exactly the specified
polynomial weight.

\emph{Cancellation before the boundary pairing.}
Put \(P(\xi)=\int e^{\xi x}\phi'(x)\,dx\).
For \(b=a+\epsilon\), let \(z=b+it,w=b+iu\).
The ordinary commutator kernel (MA-MO.184) can
be written, with its sign unchanged, as
\[
 \begin{aligned}
 k_{\omega,b,g}(t,u)&=L_{\omega,b,g}(t,u)
                                      j_{\omega,\epsilon}(u),\\
 L_{\omega,b,g}(t,u)
   &=\frac{G(z)P(z-w)}{2\pi}
                          \frac{A_\omega(z)-A_\omega(w)}{z-w}.
 \end{aligned}
 \tag{MA-MO.718}
\]
The divided difference has its removable ordinary
value at \(z=w\), before the singular reciprocal
is paired with anything. The equality follows from
\(A(z)[j(z)-j(w)]=(A(w)-A(z))j(w)\) and
\(z-w=i(t-u)\). Equivalently it uses (MA-MO.713)
in the limit. There is no product of a diagonal
principal-value kernel with an independent pole
distribution.

The functions \(L_{\omega,b,g}\) converge in
\(\mathscr S(\mathbb R^2)\) as \(b\downarrow a\).
To check all derivatives, the forward multiplier
has polynomial derivative bounds on this closed
strip. Its divided difference has polynomial
derivative bounds too, by integrating its derivative
along the vertical segment from \(w\) to \(z\).
The factors \(G(b+it)\) and \(P(i(t-u))\) supply
rapid decrease in the independent directions
\(t\) and \(t-u\). Equations (MA-MO.707),
(MA-MO.710) and (MA-MO.715) therefore give a
trace-class limit \(K_{\omega,a}^{\partial}(g)\)
on \(H_{a,s}\).

The crossed Fourier kernel has the same type of
controlled singularity, now with the \emph{dual}
reciprocal:
\[
 \begin{aligned}
 o_{\eta,b,g}(t,u)
   &=L^{\mathrm{cross}}_{\eta,b,g}(t,u)
                                      j_{\bar\eta,\epsilon}(u),\\
 L^{\mathrm{cross}}_{\eta,b,g}(t,u)
   &=\frac{G(z)A_\eta(z)m_\eta(z)}{2\pi}
                              \frac{P(-z-w)}{-z-w}.
 \end{aligned}
 \tag{MA-MO.719}
\]
In the rational case omit the subscripts on the
self-dual factors. This is exactly (MA-MO.481)
or (MA-MO.331) on the line \(b>a\), not a
new Fourier normalization. The real part of
\(-z-w\) stays at most \(-2a<0\).
The actual local multiplier \(m_\eta\), including
the finite phases, conductor and different, is
holomorphic with polynomial derivative bounds
on the compact positive strip under consideration.
The rapid directions are \(t\) and \(t+u\).
Thus \(L^{\mathrm{cross}}_{\eta,b,g}\) also
converges in Schwartz space, and (MA-MO.715)
gives an actual trace-class limit
\(O_{\eta,a}^{\partial}(g):H_{a,s}^{\bar\eta}
\to H_{a,s}^{\eta}\).

\emph{Actual compatible character rows.}
The third term needed in the first row is
\(N_a(g)=C_gM_\phi:H\to H_{a,s}\).
It is trace class by the exact physical kernel
(MA-MO.332): \(x-y\) has compact support,
the normalized weight ratio decreases exponentially
as \(y\to-\infty\), and \(\phi(y)=0\) far right.
Every derivative obeys the same estimates.
This argument uses the sum norm on the input;
it makes no trace-class assertion for an unchanged
one-sided half-line multiplication.

Define
\[
 \mathscr R_{\eta,a}^{\partial}(g)y
       =K_{\eta,a}^{\partial}(g)y
                     -N_a(g)y+O_{\eta,a}^{\partial}(g)Ry .
 \tag{MA-MO.720}
\]
It is trace class from \(H\) to \(H_{a,s}\).
For \(y\in\mathcal E\), it is the first row of
the original source matrix. To prove this exact
agreement, first use the already proved source
identity (MA-MO.479)--(MA-MO.482) at every
\(b>a\). Both inverse functions there are the
fixed original functions identified in
(MA-MO.711). Their cutoff vectors are actual
local source vectors by (MA-MO.714) and the
structural inclusion, and their theta outputs
are in \(\mathcal E\). In the Fourier picture
the input \(V_by\) converges in \(H^r\), the
two kernel operators converge in trace norm by
the preceding construction, and the transform
of that fixed actual output converges to its
\(a\)-line transform. The half-line term has
the same direct limit. Thus the identity passes
to (MA-MO.720). This argument takes limits
on the genuine smooth core; it does not assert
that identity maps between all varying-width
Hilbert spaces are bounded.

The second row is
\(R\mathscr R_{\bar\eta,a}^{\partial}(g^\vee)R\).
On \(\mathcal E\) the two functions agree by
the actual dual Fourier relation and nonpolar
Poisson summation (MA-MO.479). The paired map
is trace class into \(H_{a,s}\oplus H_{-a,s}\).
The diagonal image of \(H\) is closed and
\(\mathcal E\) is dense in \(H\), so the range
of the paired map is in that diagonal on the
whole completion. Its inverse diagonal isometry
gives the first operator in (MA-MO.703),
with agreement on the entire original smooth
nontrivial common sector.

\emph{The rational cutoff moments and polar rows.}
The rational case requires two further checks.
For \(y\in\mathcal E\), the original \(F_y\)
has polynomially weighted \(a\)-line growth
at the left by (MA-MO.712), and arbitrary
right weighted decrease. Therefore its Mellin
integral at \(1/2\) is absolutely convergent.
The wide inverse quotient continues through
that point and has a zero there, so
\[
          \mathcal MF_y(1/2)=0 .
 \tag{MA-MO.721}
\]
This is the original zero total moment, not
a moment imposed on the input.

Set \(p_a(x)=\phi(x)e^{(1/2-a)x}\).
It is Schwartz, since \(a<1/2\) and
\(\phi\) is zero far right. With the unitary
plus Fourier normalization define
\[
 \begin{aligned}
 \lambda_0(y)
    &=\big\langle j_+(u),
                     \widehat p_a(-u)(V_ay)(u)\big\rangle,\\
 \lambda_1(y)&=\lambda_0(Ry).
 \end{aligned}
 \tag{MA-MO.722}
\]
The regular summand in (MA-MO.710) times
\(\widehat p_a(-u)\) is Schwartz.
Each compact singular summand is in \(H^{-r}\),
and multiplication by \(\widehat p_a(-u)\)
is bounded on \(H^r\), by (MA-MO.716).
Thus \(\lambda_0\) is bounded on \(H_{a,s}\),
\(\lambda_1\) is bounded on \(H_{-a,s}\), and
both are bounded on the actual sum space.

For \(y\in\mathcal E\), bilinear Parseval
on the line \(b>a\), followed by the same
boundary limit, gives exactly
\[
 \lambda_0(y)=\int\phi(x)e^{x/2}F_y(x)\,dx,
              \qquad \lambda_1(y)=\lambda_0(Ry).
 \tag{MA-MO.723}
\]
There is no extra Fourier constant: both
transforms in (MA-MO.722) are unitary.
The physical integral converges by the
left estimate proved before (MA-MO.721);
on lines \(b>a\) its weight is
\(\phi(x)e^{(1/2-b)x}\), whose Fourier
transform converges in Schwartz space to
\(\widehat p_a\).

The actual source \(h_y=u_y+\mathfrak F u_{Ry}\)
from (MA-MO.326) consequently has
\[
       h_y(0)=-\lambda_1(y),\qquad
       \int h_y=-\lambda_0(y).
 \tag{MA-MO.724}
\]
Indeed \(u_y(0)=0\), while its integral is
\(\int e^{x/2}qF_y=-\lambda_0(y)\), by
(MA-MO.721); Fourier interchanges evaluation
and integration. These statements use the
actual additive Schwartz functions (MA-MO.714).

Retain the fixed Gaussian sections \(f_0,f_1\)
and vectors \(e_0=\Theta f_0\),
\(e_1=\Theta f_1-e^{-x/2}\) of
(MA-MO.333), with \(Re_0=e_1\).
They belong to \(H\) since \(a<1/2\).
After subtracting the actual polar sections,
the first row and the polar coordinates are
\[
 \begin{aligned}
 T_{a,g}^{\partial}y
   ={}&K_a^{\partial}(g)y-N_a(g)y
                                 +O_a^{\partial}(g)Ry\\
      &+G(-1/2)\lambda_1(y)e_0
                     +G(1/2)\lambda_0(y)e_1,\\
 \nu_g(y)
   ={}&\bigl(-G(-1/2)\lambda_1(y),
                         -G(1/2)\lambda_0(y)\bigr).
 \end{aligned}
 \tag{MA-MO.725}
\]
The signs are those of the actual centered residue
\(\Theta h=Z_bF_h+(\int h)e^{-x/2}\)
in (MA-MO.334). That direct theta contour
shift crosses only the pole of zeta, not a pole
at any of its zeros. Alternatively, all
identities (MA-MO.328)--(MA-MO.335) hold at
\(b>a\); the now proved moment and kernel
limits give exactly (MA-MO.725).
Every term in its first line is trace class
\(H\to H_{a,s}\), and \(\nu_g\) is bounded
finite rank.

The reflected second row is
\(RT_{a,g^\vee}^{\partial}R\).
It agrees on \(\mathcal E\) by
\(h_{Ry}=\mathfrak Fh_y\) and the interchange
of the two Gaussian sections. As in the
character case, the closed diagonal gives
a unique trace-class \(T_g^\partial:H\to H\).
The negative source column is therefore
\(y\mapsto(T_g^\partial y,\nu_g(y))\).
The two remaining columns \(C_gDs_j\) are
fixed actual smooth common-sum vectors,
already constructed without a zero-location
assumption in MA-MO-016. Appending those
finite columns gives
\[
 \mathscr D_g(y,c)
       =(T_g^\partial y,\nu_g(y))
                           +\sum_{j=0}^1c_jC_gDs_j .
 \tag{MA-MO.726}
\]
This is trace class on \(\mathcal W_{a,s}\)
and agrees with the original matrix on the
dense source \(\mathfrak s(\mathcal E\oplus\mathbb C^2)\).
This proves the rational part of (MA-MO.703).

\emph{The quotient statements and the scope of the result.}
In the nontrivial case the exact smooth formula
has \(P_+\mathcal E\subset I_\eta\).
It also preserves \(I_\eta\) after subtracting
the identity and convolving. The bounded
extension therefore preserves \(J_\eta\), and
\(\mathscr D_{\eta,g}+C_g\) has range in
\(J_\eta\) by density. Hence its quotient is
\(-C_g\). Condition (MA-MO.702) retains all
boundary jets, by the exact inequality
\(2j<s-1\) of MA-WP-001--003: for a
multiplicity \(m\), every \(0\le j<m\)
is retained precisely when \(s>2m-1\).
All interior zeros are retained as well.
The actual quotient trace is consequently
the full absolutely convergent sum in
(MA-MO.704). Triangular compression on
\(J_\eta\oplus J_\eta^\perp\), followed by
finite-rank trace-norm approximation, proves
both the trace and Fredholm determinant
identities there.

The source core is dense in each claimed
completion, proving uniqueness. Each
Schwartz kernel estimate and each finite
distribution pairing uses only finitely
many derivatives of \(g\) on a fixed
compact support. The finite source columns
have the same continuity. Thus the maps
are continuous into trace class.
Finally MA-MO-025/026 already refutes
boundedness for every nonzero compact
test when a boundary multiplicity violates
(MA-MO.702). This proves the stated exact
characterization within (MA-MO.701).
\end{proof}

\paragraph{Remaining scope and framework placement.}
The proof uses the finite edge divisor twice:
to have only finitely many \(H^{-s/2}\) input
singularities, and to retain uniform polynomial
estimates for the remaining reciprocal.
It gives no summability estimate for infinitely
many approaching interior poles, and does not
replace their actual lifting costs by their
formal multiplicities. The broad full-retention
sufficiency question of MA-MO-025 remains
separate.

The rational operator is an endomorphism of
the polar completion \(\mathcal W\), not
automatically of \(\mathcal W_+\); the
rank-two outgoing obstruction of MA-MO-023
is unchanged. In the nontrivial sector the
ideal restriction in (MA-MO.704) is actual
and trace class, but a local evaluation of
that restriction is not assumed merely
from the boundary construction.

The primary source maps and structural
interfaces are the already checked Meyer
locations in MA-MO-018/022; the present
proof introduces the boundary-distribution
estimate, not a new reading of those
sources. In the inherited cote~1 framework,
scans~079--080 and~111--112 through
ALPOGE\_BRIDGES BR-TR-01/02, a boundary
multiplicity is a finite-length datum,
whereas (MA-MO.709), (MA-MO.715) and
the actual source comparison supply the
additional analytic continuity and
trace-ideal structure. No arithmetic
positivity or new zero-location theorem
is inferred from the finite-length class.

\section{An interior contour removes the ideal return's Gevrey restriction}
\label{sec:interior-contour-source-return}

The reciprocal on a fixed boundary line can have the
logarithmic-square loss of MA-RS-001. The actual theta ideal
has additional structure: its Mellin numerator cancels the
entire nontrivial divisor. An input contour can therefore
move inside the open strip of its inherited two-sided norm.
We construct such a contour with polynomial reciprocal bounds
and use it in the original local-source formula. This
removes both the Gevrey and generic-width assumptions for
that restricted source return. It does not remove them by
an estimate on the unrestricted one-sided inverse.

Retain the nontrivial primitive finite-order character
\(\eta\), structural \(S\), actual compact-sector source
\(\mathcal T_\eta\), Fourier pair, and theta embedding of
MA-MO-012 and MA-MO-017. Put
\[
 \begin{gathered}
 0<a<1/2,\quad s\ge0,\quad
 H=H_{\Sigma,a,s},\quad
 J=\overline{I_\eta^c}^{\,H},\\
 A_\eta(z)=E_{\eta,S}(1/2+z)L_\eta(1/2+z),\\
 \phi=1\text{ sufficiently far left},\quad
 \phi=0\text{ sufficiently far right},\quad
 \phi'\in C_c^\infty,\quad
 \psi=1-\phi-\phi^\vee .
 \end{gathered}
 \tag{MA-MO.800}
\]
The cutoff is real and smooth. Its derivative and the
test \(g\in C_c^\infty(\mathbb R)\) are arbitrary compact
smooth functions; no Gevrey bounds are assumed.
The local multiplier \(m_\eta\), source endomorphism
\(\mathcal L_\eta\), and exact normalizations are those
of (MA-MO.281)--(MA-MO.283).

\paragraph{MA-MO-028 (all-width, all-smooth actual ideal return).}
For every width and polynomial weight in (MA-MO.800)
there is a unique continuous nuclear map
\[
 \begin{aligned}
 \mathcal A_{\eta,g}:J&\longrightarrow\mathcal T_\eta,\\
 \mathcal A_{\eta,g}\vartheta_\eta F
       &=C_g\mathcal L_\eta F
                         \qquad(F\in\mathcal T_\eta).
 \end{aligned}
 \tag{MA-MO.801}
\]
Its theta composition is the actual smoothed common-sum
action on the entire smooth theta ideal and is trace class:
\[
 \mathcal B_{\eta,g}=\vartheta_\eta\mathcal A_{\eta,g}
       \in\mathcal S_1(J),\qquad
 \mathcal B_{\eta,g}|_{I_\eta^c}
                    =B^{\rm raw}_{\eta,g}.
 \tag{MA-MO.802}
\]
The return in fact takes values in the original smooth
source, and its theta composition takes values in
\(I_\eta^c\) with its \(\mathcal E\) topology.
On each fixed compact test-support space,
\(g\mapsto\mathcal B_{\eta,g}\) is continuous into trace
class. The construction and its actual trace comparison
give
\[
 \operatorname{Tr}_{J}\mathcal B_{\eta,g}
 =\frac1{2\pi}\int_{\mathbb R}
      \mathcal Mg(it)\frac{m_\eta'(it)}{m_\eta(it)}\,dt
                  +g(0)\int_{\mathbb R}\psi(x)\,dx .
 \tag{MA-MO.803}
\]
There is no zero-location or simplicity hypothesis,
and a width whose boundary contains zeros is allowed.
This is an assertion on \(J\), not the full Hilbert
completion of the original common-sum source.

\begin{proof}
\emph{Polynomial good paths strictly inside the input strip.}
We first construct a path depending only on the fixed
function and width, not on \(g\) or the source vector.
Use the horizontal heights of MT-M18,
\[
 n+\tfrac14<h_n<n+\tfrac34,\qquad
             \tfrac12<h_{n+1}-h_n<\tfrac32 .
 \tag{MA-MO.804}
\]
On each horizontal segment
\(\{c+ih_n:|c|\le a/4\}\), (M18.8), after canceling
the common Stirling factor, gives
\[
             |A_\eta(c+ih_n)|^{-1}
                     \le C(2+|n|)^D .
 \tag{MA-MO.805}
\]
The finite deleted Euler factors are bounded away from
zero on this fixed positive finite-\(L\) strip.
The height choice avoids every nontrivial zero ordinate.
There are no trivial finite-\(L\) zeros in this strip.
Thus the finitely many bounded-height segments are also
covered after increasing \(C\).

Here is the corresponding vertical choice. For each
sufficiently large \(|n|\), center a disk at
\(2+i(h_n+h_{n+1})/2\), in the uncentered finite-\(L\)
variable, and choose a zero-free boundary radius
\(12<R_n<13\). At these heights its only zeros are
nontrivial zeros, and its total multiplicity \(M_n\)
is \(O(\log(2+|n|))\). For \(c\in[-a/4,a/4]\), put
\[
 V_n(c)=
 \sum_{\rho\text{ in this disk}}
       m_\rho\log^+\frac1{|1/2+c-\Re\rho|}.
 \tag{MA-MO.806}
\]
Each summand has integral at most
\(\int_{-1}^{1}-\log|u|\,du=2\) over this interval.
Since the interval has length \(a/2\), averaging
chooses \(c_n\in[-a/4,a/4]\), avoiding the finitely
many zero real parts in the sum, with
\[
 \begin{aligned}
 V_n(c_n)&\le (4/a+1)M_n,\\
 \prod_{\rho\text{ in the disk}}
       \min\{1,|1/2+c_n-\Re\rho|\}^{m_\rho}
           &\ge e^{-(4/a+1)M_n}
            \ge C^{-1}(2+|n|)^{-D}.
 \end{aligned}
 \tag{MA-MO.807}
\]
When \(M_n=0\), take the midpoint.

We spell out why this product controls the whole
vertical segment, rather than one selected point.
Write \(L_\eta=H_nB_n\) on the disk, where \(B_n\)
is its finite disk Blaschke product, with multiplicities.
The polynomial strip upper bound gives
\(|H_n|\le C(2+|n|)^C\) on the disk. At its center,
the Euler-product lower bound and \(|B_n|\le1\)
give \(|H_n|\ge c>0\).
Harnack's inequality applied to
\(\log(C(2+|n|)^C)-\log|H_n|\) therefore gives
\[
                     |H_n(w)|\ge C^{-1}(2+|n|)^{-D}
 \tag{MA-MO.808}
\]
on the inner half disk.
Every point \(w=1/2+c_n+it\),
\(h_n\le t\le h_{n+1}\), lies in its inner quarter.
Each normalized Blaschke factor there has modulus
at least
\[
       \frac{\min\{1,|1/2+c_n-\Re\rho|\}}{2R_n}.
\]
Multiplying these bounds, using (MA-MO.807) and
\((2R_n)^{M_n}\le C(2+|n|)^D\), proves
\[
 |A_\eta(c_n+it)|^{-1}\le C(2+|n|)^D
                      \qquad(h_n\le t\le h_{n+1}).
 \tag{MA-MO.809}
\]
This is the same finite-disk argument as MT-M18,
but with averaging in the real coordinate instead
of the imaginary coordinate. For the finitely many
remaining blocks, choose \(c_n\) avoiding all zero
real parts meeting the compact vertical rectangle.
Compactness then supplies (MA-MO.809) with an enlarged
constant. No zero-separation hypothesis was used.

Let \(\Gamma\) be the upward-oriented polygonal path
consisting of the vertical segment at \(c_n\) from
height \(h_n\) to \(h_{n+1}\), followed by the
horizontal segment to \(c_{n+1}+ih_{n+1}\).
Its horizontal segments may point in either direction.
Equations (MA-MO.805) and (MA-MO.809) give
\[
 \begin{gathered}
 |\Re w|\le a/4\quad(w\in\Gamma),\\
 \operatorname{length}\bigl(\Gamma
          \cap\{h_n\le\Im w\le h_{n+1}\}\bigr)
                    \le 3/2+a,\\
                  |A_\eta(w)|^{-1}
                       \le C(2+|\Im w|)^D\quad(w\in\Gamma).
 \end{gathered}
 \tag{MA-MO.810}
\]
The length bound permits either convention for assigning
the horizontal endpoints to their adjacent blocks.
The path may cross the central vertical line; this
does not affect any contour identity below.

\emph{The Hilbert norm controls interior evaluations.}
For every \(y\in H\), its Mellin transform is holomorphic
on \(|\Re w|<a\). For \(|\Re w|\le a/4\),
weighted Cauchy--Schwarz gives, for every integer \(j\ge0\),
\[
 \begin{aligned}
 |(\mathcal My)^{(j)}(w)|
 &\le \|y\|_H
 \left(\int_{\mathbb R}
       \frac{|x|^{2j}e^{2\Re(w)x}}
       {(e^{2ax}+e^{-2ax})(1+x^2)^{s/2}}\,dx\right)^{1/2}\\
 &\le C_{a,s,j}\|y\|_H .
 \end{aligned}
 \tag{MA-MO.811}
\]
The last integrand has a fixed exponential decrease
at both ends, uniformly in that interior interval.
Holomorphy and its differentiated formula follow by
the same domination on every smaller open strip.
In particular (MA-MO.810)--(MA-MO.811) imply
\[
       \left|\frac{\mathcal My(w)}{A_\eta(w)}\right|
           \le C(2+|\Im w|)^D\|y\|_H
                                      \qquad(w\in\Gamma).
 \tag{MA-MO.812}
\]
This estimate is for evaluations on the chosen path.
It makes no assertion that the reciprocal is a
Schwartz multiplier on either boundary line.

\emph{The original local kernel as a holomorphic contour kernel.}
Set
\[
 \begin{gathered}
 G(z)=\mathcal Mg(z),\qquad
 P(\xi)=\int e^{\xi x}\phi'(x)\,dx,\qquad
 \Psi(\xi)=\int e^{\xi x}\psi(x)\,dx,\\
 \mathcal K_\eta(z,w)=G(z)\left[
   -\frac{P(z-w)}{z-w}
          \left(\frac{m_\eta(z)}{m_\eta(w)}-1\right)
       +\frac{m_\eta(z)}{m_\eta(w)}\Psi(z-w)\right].
 \end{gathered}
 \tag{MA-MO.813}
\]
For fixed \(\Re z>-1/2\), this is holomorphic in
\(|\Re w|<a\). The apparent singularity \(w=z\)
is removable, since \(m_\eta\) is holomorphic and
nonzero on the possible input strip.
The local multiplier and its derivatives have
polynomial bounds on fixed strips by (MA-MO.281).
Its inverse has polynomial bounds uniformly on
\(|\Re w|\le a/4\), which is strictly inside
\((-1/2,1/2)\).

For every fixed output line \(z=c+it\), \(c>-1/2\),
every derivative order \(j\), and every \(N\), these
bounds and integration by parts in the three compact
smooth transforms give a bound of the form
\[
 \left|\partial_t^j\mathcal K_\eta(c+it,w)\right|
 \le C_{c,j,N}
   \sum_{\nu\le j}|G^{(\nu)}(c+it)|
       (1+|t|+|\Im w|)^{D_{c,j}}
       (1+|t-\Im w|)^{-N},
       \qquad w\in\Gamma.
 \tag{MA-MO.814}
\]
For \(|z-w|\ge1\), this follows directly from
(MA-MO.813). For \(|z-w|\le1\), write
\((m_\eta(z)-m_\eta(w))/(z-w)\) as the integral of
\(m_\eta'\) along the segment. The segment stays in a
fixed closed substrip of \(\Re z>-1/2\). The same
argument applies after every output derivative.
Thus no singular denominator is hidden in (MA-MO.814).
All constants are uniform when \(c\) lies in a fixed
compact interval in this half-plane.

For an actual smooth ideal vector \(y=\vartheta_\eta F\),
the unconditional division statement (MA-MO.173)--(MA-MO.178)
gives
\[
 Q_y(w)=\frac{\mathcal My(w)}{A_\eta(w)}
             =\mathcal MF(w)
       \quad\text{holomorphic for }\Re w>-1/2,
 \tag{MA-MO.815}
\]
with rapid decrease, including derivatives, on every
fixed closed substrip. This is its original source
inverse, not an arbitrary section of the Hilbert vector.
The local formula (MA-MO.288), whose proof uses only
the polynomial local multiplier estimates and compact
smoothness, yields
\[
 \mathcal M(C_g\mathcal L_\eta F)(z)
       =\frac1{2\pi i}\int_{i\mathbb R}
                     \mathcal K_\eta(z,w)Q_y(w)\,dw .
 \tag{MA-MO.816}
\]
The factor \(1/(2\pi i)\) corresponds to the
unnormalized Mellin transforms; on \(w=iu\), this
is exactly the \(1/(2\pi)\) Fourier kernel in
(MA-MO.288).

\emph{Deformation on the actual ideal, before taking its closure.}
Truncate the central contour at two selected heights
and the polygonal contour at the same heights, joining
their endpoints horizontally inside \(|\Re w|\le a/4\).
For the actual smooth ideal input, the integrand in
(MA-MO.816) is holomorphic throughout this strip by
(MA-MO.815), including all its canceled global zeros.
Cauchy's theorem applies to the finite polygon. If
the polygon crosses the central line, apply it to
the successive rectangles; the intermediate sides
cancel with their orientations.
The two final horizontal connectors tend to zero:
their lengths are uniformly bounded, \(Q_y\) is
rapid there, and the local kernel and its derivatives
are polynomially bounded for fixed \(z\).
Equations (MA-MO.810) and (MA-MO.814) give absolute
convergence of the infinite path integral. Consequently
\[
 \mathcal M(C_g\mathcal L_\eta F)(z)
       =\frac1{2\pi i}\int_{\Gamma}
          \mathcal K_\eta(z,w)
               \frac{\mathcal My(w)}{A_\eta(w)}\,dw .
 \tag{MA-MO.817}
\]
There is no residue term: this deformation is performed
only after the actual ideal's full zero jets have
canceled every pole in the quotient.

\emph{Uniform source estimates from the inherited Hilbert norm.}
Combine (MA-MO.810), (MA-MO.812) and
(MA-MO.814). The length per height block is bounded.
For \(N\) sufficiently large compared with the fixed
polynomial exponents,
\[
 \sum_{n\in\mathbb Z}
  (1+|t|+|n|)^{D_1}(1+|n|)^{D_2}
        (1+|t-n|)^{-N}
                \le C(1+|t|)^{D_3}.
 \tag{MA-MO.818}
\]
For example, bound \(1+|n|\) by
\((1+|t|)(1+|t-n|)\), do the same for
\(1+|t|+|n|\), and sum a remaining power
less than \(-2\).
The rapid decrease of every \(G^{(\nu)}(c+it)\)
now gives, for all integers \(M,j\ge0\),
\[
 \sup_{t\in\mathbb R}(1+|t|)^M
  \left|\partial_t^j
        \mathcal M(C_g\mathcal L_\eta F)(c+it)\right|
          \le C_{c,M,j,g,\phi}\|y\|_H .
 \tag{MA-MO.819}
\]
On a fixed compact support for \(g\), each constant
is bounded by finitely many ordinary derivative
seminorms of \(g\). The same statement holds locally
uniformly in \(c>-1/2\).
Fourier inversion and the usual Schwartz estimates
therefore control every seminorm of
\(e^{cx}C_g\mathcal L_\eta F\) by \(\|y\|_H\).
No Gevrey estimate occurs in this calculation.

The actual dual Fourier identity is (MA-MO.295):
\[
 \mathfrak S_\eta C_g\mathcal L_\eta F
      =C_{g^\vee}\mathcal L_{\bar\eta}\mathfrak S_\eta F,
 \qquad
       \vartheta_{\bar\eta}\mathfrak S_\eta F=Ry .
 \tag{MA-MO.820}
\]
Construct an interior path for the dual character
by the same procedure, and apply (MA-MO.819) to
\(Ry\). Reflection is isometric on the exact sum
norm. Hence every weighted Schwartz seminorm of
the actual dual output is also bounded by \(\|y\|_H\).
All conductor, different and compact-character phases
remain those of (MA-MO.820); no unshifted archimedean
Fourier map is substituted.

\emph{Extension to the whole completed ideal and the actual target.}
For a Cauchy sequence from \(I_\eta^c\) in \(H\),
the two output families just estimated converge in
all their weighted Schwartz seminorms. On the central
local \(L^2\) line the actual source Fourier map is
unitary, so their limits remain the correct Fourier
pair. Meyer's fixed-sector source topology is precisely
this paired weighted-Schwartz topology. Thus the
limit belongs to \(\mathcal T_\eta\), continuously
and independently of the approximating sequence.
This defines \(\mathcal A_{\eta,g}\) on all of \(J\)
and proves (MA-MO.801), rather than only a formula
for individually smooth \(y\).

The scalar path integral can be formed on arbitrary
vectors of \(H\), since (MA-MO.811) holds there.
We do not assert the requisite dual-source compatibility
or the original-source identification on those vectors.
The preceding extension uses the actual ideal and
its density. This domain distinction is essential.

\emph{Nuclearity with one bounded family of genuine source vectors.}
We include the functional-analytic step so that the
all-weight return is stronger than separate Hilbert
kernel estimates. Choose increasing Hilbertian norms
\(q_n\) defining the nuclear Fr\'echet topology of
\(\mathcal T_\eta\), each dominating its injective
central local \(L^2\) norm, and sufficiently far apart
that their successive completion maps are trace class.
The source property is the one used and explicitly
located in MA-MO-017. Estimate (MA-MO.819) and its
dual give a bounded map \(J\to(\mathcal T_\eta,q_n)\)
for every \(n\).
Factoring through a stronger norm shows that
\(\mathcal A_{\eta,g}\) is trace class into each
of these Hilbert completions.

Let \(P_N\) be increasing finite-rank orthogonal
projections on \(J\), tending strongly to one.
For each \(n\), the \(q_n\) trace norm of
\(\mathcal A_{\eta,g}(1-P_N)\) tends to zero.
Select an increasing sequence \(N_k\) so that this
tail is at most \(2^{-k}\) for all \(n\le k\).
The singular-value decomposition of
\(\mathcal A_{\eta,g}(P_{N_{k+1}}-P_{N_k})\),
measured in \(q_k\), has coefficient sum at most
\(2^{-k}\), norm-one domain functionals, and
normalized output vectors in the actual source.
They lie in that source because each finite-dimensional
image is spanned by images of vectors of \(J\).
For each fixed \(q_n\), these normalized vectors
are bounded for \(k\ge n\), and the finitely many
earlier blocks are bounded as well. Including the
first finite block gives
\[
 \mathcal A_{\eta,g}y
     =\sum_\nu \lambda_\nu\,\ell_\nu(y)F_\nu,
 \quad
 \sum_\nu|\lambda_\nu|<\infty,\quad
 \sup_\nu\|\ell_\nu\|_{J^*}<\infty,\quad
 \{F_\nu\}\text{ bounded in }\mathcal T_\eta .
 \tag{MA-MO.821}
\]
This is a single genuine nuclear representation.
It is not a collection of unrelated one-sided traces.

The original theta embedding is continuous
\(\mathcal T_\eta\to I_\eta^c\subset\mathcal E\)
by the structural source comparison and nontrivial-sector
Poisson summation in MA-MO-012. Thus its composition
with (MA-MO.821) is trace class on \(J\) and actually
smooths into \(I_\eta^c\). On the dense smooth ideal,
the unconditional source identity (MA-MO.183) gives
\(\vartheta_\eta C_g\mathcal L_\eta F
=B^{\rm raw}_{\eta,g}\vartheta_\eta F\).
Continuity proves (MA-MO.802) everywhere.
To verify test continuity, choose a source Hilbert
norm that controls the theta map into \(H\), then
a stronger norm with trace-class inclusion. Its
operator bound in (MA-MO.819) uses only finitely
many derivative seminorms of \(g\). This bounds
\(\|\mathcal B_{\eta,g}\|_1\) by those same finite
seminorms on each fixed compact test support.

\emph{The actual trace comparison.}
The reverse composition on the original source is
\[
 \mathcal A_{\eta,g}\vartheta_\eta
                       =C_g\mathcal L_\eta .
 \tag{MA-MO.822}
\]
The source has the approximation property from its
split weighted-Schwartz model. Applying the same
absolutely summable pairings in (MA-MO.821) to
the two compositions gives
\[
 \operatorname{Tr}_{J}
          (\vartheta_\eta\mathcal A_{\eta,g})
     =\operatorname{Tr}_{\mathcal T_\eta}
          (\mathcal A_{\eta,g}\vartheta_\eta)
     =\operatorname{Tr}_{\mathcal T_\eta}C_g\mathcal L_\eta .
 \tag{MA-MO.823}
\]
The latter is the local source trace already computed
without a global inverse. Indeed (MA-MO.813) on the
central local line, without \(A_\eta^{-1}\), is a
Schwartz kernel divided by \(2\pi\). Its diagonal is
\[
       \frac{G(it)}{2\pi}
          \left(\frac{m_\eta'(it)}{m_\eta(it)}
                         +\int\psi\right),
 \tag{MA-MO.824}
\]
because \(P(0)=\int\phi'=-1\).
The local-to-source nuclear factorization with global
inverse replaced by one is exactly the unconditional
local argument in MA-MO-017; its paired weighted
estimates use polynomial local multipliers and compact
smooth data only. It equates the actual local \(L^2\)
trace to the source trace. Integrating (MA-MO.824)
and using \((2\pi)^{-1}\int G(it)\,dt=g(0)\)
proves (MA-MO.803).
This concludes the all-width, all-smooth theorem.
\end{proof}

\paragraph{The rational positive source and its two actual moments.}
There are corresponding unconditional, all-width conclusions
for the rational source with \(S=\{\infty\}\).
Use its exact spaces \(\mathscr U_0,\mathscr U_{\rm pf}\),
\(J_{a,s},K_{a,s,\phi}\), and \(\mathcal W_+\) from
MA-MO-010/011/016. The original additive cutoff is
\(v(t)=\phi(\log|t|)\), extended by one at zero.
For \(y=\Theta f\), \(f\in\mathscr U_0\), put
\[
 F(x)=2e^{x/2}f(e^x),\qquad
 \widetilde F(x)=2e^{x/2}\mathfrak Ff(e^x),\qquad
                A(z)=\zeta(1/2+z).
 \tag{MA-MO.825}
\]
The unconditional actual source identities
(MA-MO.127)--(MA-MO.128) identify their Mellin transforms
with \(\mathcal My/A\) and \(\mathcal M(Ry)/A\)
throughout the interior contour strip. Their zero
singularities are canceled, and their weighted transforms
are rapid there. The pole of \(A\) at \(z=1/2\) lies
outside that strip and creates no contour term.
The same path construction therefore applies to both
actual radial sources.

We explicitly check the source moments as continuous
functionals on the completed ideal, rather than importing
an unstated inverse-domain bound. Define
\[
 \begin{aligned}
 C_\phi(w)
   &=\int_{\mathbb R}\phi(x)e^{(1/2-w)x}\,dx
       =-\frac{P(1/2-w)}{1/2-w},\qquad \Re w<1/2,\\
 \lambda_0(y)&=\int_{\mathbb R}v(t)f(t)\,dt,\\
 \lambda_1(y)&=\int_{\mathbb R}v(t)\mathfrak Ff(t)\,dt .
 \end{aligned}
 \tag{MA-MO.826}
\]
The integration by parts for \(C_\phi\) has no endpoint
term: at the left the exponential decreases, and at
the right \(\phi\) is zero. This transform is rapidly
decreasing on every closed vertical substrip of
\(\Re w<1/2\).
Changing the additive variable \(t=e^x\), using
evenness and the factor two in (MA-MO.825), gives
\[
 \begin{aligned}
 \lambda_0(y)
   &=\int_{\mathbb R}\phi(x)e^{x/2}F(x)\,dx\\
   &=\frac1{2\pi i}\int_{\Gamma}
           C_\phi(w)\frac{\mathcal My(w)}{A(w)}\,dw,\\
 \lambda_1(y)
   &=\frac1{2\pi i}\int_{\Gamma}
           C_\phi(w)\frac{\mathcal M(Ry)(w)}{A(w)}\,dw .
 \end{aligned}
 \tag{MA-MO.827}
\]
First derive these identities on the central line by
absolute Fourier inversion and then deform exactly
as in (MA-MO.817). On the path the rapid factor
\(C_\phi\) defeats the polynomial bound (MA-MO.812).
Thus
\[
       |\lambda_0(y)|+|\lambda_1(y)|
                         \le C_{a,s,\phi}\|y\|_{H_{\Sigma,a,s}} .
 \tag{MA-MO.828}
\]
They extend uniquely to all of \(J_{a,s}\).
On overlap with MA-MO-011's regimes, these are the
same actual moment extensions, since their values
agree on the dense smooth theta ideal. Its closure
and codimension proof uses only moment continuity
and two fixed smooth correcting vectors. The
correction (MA-MO.147) therefore repeats at every
width here, and gives
\(K_{a,s,\phi}=J_{a,s}\cap\ker\lambda_0\cap\ker\lambda_1\).
The boundary-line formula involving \(\beta_a\)
in (MA-MO.148) is not asserted when \(j_a\)
is singular. The newly continuous functionals
on \(J_{a,s}\) nevertheless have their canonical
Riesz vectors and the same rank-two Gram construction.

The local kernel (MA-MO.813), with the rational
multiplier (MA-MO.282), and the two radial sources
in (MA-MO.825) now give a continuous nuclear map
\(\mathcal A_{0,g}:J_{a,s}\to\mathscr S_{\rm even}(\mathbb R)\)
agreeing with the actual \(\mathcal D_gLf\) on
\(y=\Theta f\).
Here \(\mathcal D_g\) is the original additive dilation
operator corresponding to centered \(C_g\).
To verify the target, the bounds for all \(c>-1/2\)
give integrability at additive zero and rapid decrease
at infinity; the paired Fourier bounds give integrability
of every Fourier moment. Fourier inversion then gives
all ordinary Schwartz derivatives, including at zero,
exactly as in the rational target verification of
MA-MO-017. The common nuclear-series proof applies to
the actual even Schwartz space.

Add the two actual polar source sections \(f_0,f_1\)
of MA-MO-016:
\[
 \begin{aligned}
 \mathcal A_{+,g}: \mathcal W_+=J_{a,s}\oplus\mathbb C^2
                      &\longrightarrow\mathscr S_{\rm even}(\mathbb R),\\
 \mathcal A_{+,g}(y,c_0,c_1)
       &=\mathcal A_{0,g}y+
                 \sum_{j=0}^1 c_j\mathcal D_gLf_j .
 \end{aligned}
 \tag{MA-MO.829}
\]
The added columns are genuine fixed Schwartz vectors.
Hence this is a nuclear return for every compact smooth
test and every width, with its prescribed source-section
values. The endomorphism
\(\jmath_+\mathcal A_{+,g}\) of \(\mathcal W_+\)
has the local trace (MA-MO.803) with the rational
multiplier, by the same nuclear pairing as above.
It remains distinct from the full original positive
column, whose outgoing polar terms are those in
(MA-MO.299).

Finally \(Lf(0)=-\lambda_1(\Theta f)\) and
\(\int Lf=-\lambda_0(\Theta f)\) on \(\mathscr U_0\),
by the actual source algebra MA-MO-008.
After compact dilation these values are multiplied
by \(G(-1/2)\) and \(G(1/2)\), respectively.
They therefore vanish on \(\mathscr U_{\rm pf}\),
and the continuous extensions (MA-MO.828) show the
same for every input from \(K_{a,s,\phi}\).
The zero-moment source \(\mathscr U_0\) is a closed
finite-codimensional Schwartz subspace. The same
Hilbertian source-norm argument therefore makes this
restricted map nuclear with target \(\mathscr U_0\);
equivalently one may apply a continuous finite-rank
moment correction to its source nuclear series.
Restricting the return and then applying
the actual theta embedding proves, without a gap
or generic-width condition,
\[
       B_{\phi,g}:K_{a,s,\phi}\longrightarrow J_{a,s}
                      \quad\text{is trace class}
 \tag{MA-MO.830}
\]
and agrees with the original pole-free source action.
There is still no \(K\)-invariance claim or
endomorphism trace for this last map. The previously
proved domain and outgoing-moment obstructions are
unchanged.

\paragraph{What this decides and what it does not decide.}
The theorem applies directly to the requested primitive
nontrivial Dirichlet characters and to both
\(\psi_{269,3}\) sectors with their primitive factors
and structural place set. No use of A, M, or a
generic-width assertion occurs in this return.
It includes widths whose edge contains a zero,
because the input contour has a fixed strict
interior margin and the actual ideal cancels every
zero met during deformation.

The extension is an endomorphism of the completed
theta ideal and has its genuine Fredholm determinant.
It does not construct the full raw common-sum matrix
on \(H\), an inverse on either one-sided Hilbert space,
or a return of discarded quotient factors.
For a general source numerator, the quotient in
(MA-MO.815) has poles; the deformation then has
residues and the original wide inverse need not
be identified by the interior contour. The exterior
and deficient-boundary obstructions MA-MO-024/025
are therefore unaffected.

In particular this result does not show that the
finite-frequency operators of MA-RS-007 converge
on their unrestricted one-sided spaces for every
compact smooth test. Nor does it justify removing
the ideal projection from a common-space trace.
The trace in (MA-MO.803) is compared by the actual
return (MA-MO.801), not assigned from a scalar
zero sum. The full quotient and FT002 arithmetic
deficiency gates remain distinct.

\paragraph{Exact source and framework placement.}
The new path construction uses only the unconditional
analytic inputs and fully proved minimum-modulus
argument in the supplied GRH\_NIGHT MT-M18:
\texttt{work/root/block\_arithmetic\_synthesis.tex},
(M18.5)--(M18.9), together with NT-Z10.
The underlying fixed-Hecke facts have the inherited
primary locators Ramakrishnan and Valenza (1999),
\emph{Fourier Analysis on Number Fields}, \S7.4,
Theorem 7-19, p.~279, and \S7.7, Theorem 7-28,
p.~289. The real-coordinate averaging and the
interior contour estimates are proved here.

The original source inclusion and inverse identification
are MA-MO.173--183, through MT-M18.4 and MT-M24.6--9,
and Meyer (2005), supplied \texttt{main.tex}
lines 3592--3707 and 3808--4043.
The local multiplier, paired source topology and
nuclear comparison have the precise primary
locators recorded in MA-MO-017: Meyer, Lemma 4.3,
Propositions 4.4 and 4.6, Corollary 4.7, split
extension (24), and local trace equation (20);
Tate (1950), Theorem 2.4.1 and \S2.5, thesis
pp.~(2.11), (2.14), (2.16), (2.19).
All finite Fourier phases and conductor/different
monomials are fixed by those actual source maps.

In the inherited cote 1 framework, scans 079--080
and 111--112 through BR-TR-01/02, (MA-MO.821)--(MA-MO.823)
give the required genuine nuclear maps and common
absolutely summable trace pairing. The finite-order
character remains the specified rank-one finite-monodromy
local system, with its Frobenius-determinant Euler
factors. Analytic norm-half centering is not an
integral Tate twist. No new positivity or global
zero bound is inferred from this source return.

\section{Evaluated full-matrix traces in the finite-boundary case}
\label{sec:finite-boundary-trace-corollary}

This corollary combines two separately checked operator
constructions. MA-MO-027 extends the original full matrix
under its finite-boundary and uniform-interior-gap hypotheses.
MA-MO-028 evaluates the actual ideal or positive-source
endomorphism at every width. Their combination removes
the reserved ideal-trace evaluation in (MA-MO.704).
It does not weaken the full-source hypotheses of MA-MO-027.

Retain precisely (MA-MO.700)--(MA-MO.702), including
the strict inequality \(s>2m-1\) at every boundary
zero of multiplicity \(m\). Use the structural nontrivial
character setting of MA-MO-022, or the rational principal
setting \(S=\{\infty\}\) with the actual polar completion.
The cutoff is the same arbitrary real smooth step.
For a compact smooth test put
\[
 \begin{aligned}
 G(z)&=\mathcal Mg(z),&
 \psi&=1-\phi-\phi^\vee,\\
 Z_\omega(g)&=\sum_\rho m_\omega(\rho)G(\rho-1/2),&
 p(g)&=G(-1/2)+G(1/2),\\
 \tau_{\omega,S}(g,\phi)
   &=\frac1{2\pi}\int_{\mathbb R}
             G(it)\frac{m_\omega'(it)}{m_\omega(it)}\,dt
                         +g(0)\int_{\mathbb R}\psi(x)\,dx .
 \end{aligned}
 \tag{MA-MO.850}
\]
The local multiplier is the actual one, including all
finite Fourier phases and conductor/different monomials.
For the rational case write \(Z_0,\tau_0\).
The completed-zero sum is absolutely convergent.

\paragraph{MA-MO-029 (finite-boundary full traces and genuine determinants).}
In the nontrivial character setting, the original full
extension from MA-MO-027 satisfies
\[
 \operatorname{Tr}_H\mathscr D_{\eta,g}
                 =\tau_{\eta,S}(g,\phi)-Z_\eta(g).
 \tag{MA-MO.851}
\]
Its actual ideal restriction is the all-width operator
\(\mathcal B_{\eta,g}\) of MA-MO-028, and its determinant is
\[
 \begin{aligned}
 \det_H(1-z\mathscr D_{\eta,g})
   &=\det_{J_\eta}(1-z\mathcal B_{\eta,g})
                                \det_{Q_\eta}(1+zC_g)\\
   &=\det_{J_\eta}(1-z\mathcal B_{\eta,g})
       \prod_\rho\bigl(1+zG(\rho-1/2)\bigr)^{m_\eta(\rho)} .
 \end{aligned}
 \tag{MA-MO.852}
\]
Both equalities hold for every complex \(z\).
For \(g(0)=0\) the trace has the exact finite-prime value
\[
 \begin{aligned}
 \operatorname{Tr}_H\mathscr D_{\eta,g}
  =\sum_{\substack{\mathfrak p\notin S\\k\ge1}}
    &\frac{\log N\mathfrak p}{(N\mathfrak p)^{k/2}}
       \bigl[
        c_{\mathfrak p}^{\,k}g(k\log N\mathfrak p)\\
    &\hspace{9em}
       +\overline{c_{\mathfrak p}}^{\,k}g(-k\log N\mathfrak p)
       \bigr].
 \end{aligned}
 \tag{MA-MO.853}
\]
Here \(c_{\mathfrak p}\) is the primitive unramified
character value. No contributing prime is implicitly
added to \(S\).

In the rational setting put
\[
 \begin{gathered}
 \mathcal W=H\oplus\mathbb C^2,\qquad
 \mathcal W_+=J\oplus\mathbb C^2,\qquad Q=H/J,\\
 L_g=\jmath_+\mathcal A_{+,g}\in\mathcal S_1(\mathcal W_+),\\
 F_g=C_g(1-P_-),\qquad
 \mathscr E_g=F_g-\mathscr D_g .
 \end{gathered}
 \tag{MA-MO.854}
\]
The return \(\mathcal A_{+,g}\) is now the all-width
one in (MA-MO.829). The operator \(\mathscr E_g\)
extends the original \(C_g(1-P_+)\), preserves
\(\mathcal W_+\), restricts there to \(-L_g\),
and induces \(C_g\) on \(Q\). In particular
\[
 \begin{aligned}
 \operatorname{Tr}_{\mathcal W}\mathscr E_g
                       &=-\tau_0(g,\phi)+Z_0(g),\\
 \operatorname{Tr}_{\mathcal W}\mathscr D_g
                       &=p(g)+\tau_0(g,\phi)-Z_0(g),\\
 \det_{\mathcal W}(1-z\mathscr E_g)
                       &=\det_{\mathcal W_+}(1+zL_g)
                                          \det_Q(1-zC_g).
 \end{aligned}
 \tag{MA-MO.855}
\]
The quotient determinant is the entire product
\(\prod_\rho(1-zG(\rho-1/2))^{m_\rho}\).
This triangular factorization is for \(\mathscr E_g\);
it does not assert that the raw \(\mathscr D_g\)
preserves \(\mathcal W_+\).

There is also an exact finite-rank correction formula
for the determinant of that raw matrix. In the actual
source-section coordinates let
\[
 \begin{gathered}
 V:\mathcal W\longrightarrow\mathbb C^2,\qquad V(u,c)=c,\\
 v_j=(-r_j,e_j),\qquad
 U_g:\mathbb C^2\longrightarrow\mathcal W,\qquad
 U_ge_j=C_gv_j ,
 \end{gathered}
 \tag{MA-MO.856}
\]
where \(r_j\) are (MA-MO.255) and \(e_j\) are the
standard polar coordinate vectors. Whenever
\(1+z\mathscr E_g\) is invertible,
\[
 \begin{aligned}
 \det_{\mathcal W}(1-z\mathscr D_g)
    ={}&\det_{\mathcal W_+}(1-zL_g)\det_Q(1+zC_g)\\
      &{}\cdot
       \det_{\mathbb C^2}\!
        \left(1-zV(1+z\mathscr E_g)^{-1}U_g\right).
 \end{aligned}
 \tag{MA-MO.857}
\]
The equality extends across the exceptional \(z\)'s
after removing the singularities of the product on
the right. No value of the two-by-two resolvent factor
alone is assigned at its poles.
For \(g(0)=0\), the raw trace is
\[
 \operatorname{Tr}_{\mathcal W}\mathscr D_g
       =\sum_p\sum_{k\ge1}
             \frac{\log p}{p^{k/2}}
                 \bigl[g(k\log p)+g(-k\log p)\bigr].
 \tag{MA-MO.858}
\]
The sums in (MA-MO.853) and (MA-MO.858) are finite.
For arbitrary \(g\), the trace statements are
(MA-MO.851) and (MA-MO.855), with the displayed
polar and nonsymmetric-cutoff terms retained.

\begin{proof}
\emph{Actual restrictions and the complete retained divisor.}
On the dense smooth ideal, the operator constructed in
MA-MO-027 and the theta composition in MA-MO-028
are the same original smoothed source action.
Both are bounded on \(J_\eta\). Their restrictions
therefore agree on all of \(J_\eta\), and
(MA-MO.803) gives its trace
\(\tau_{\eta,S}(g,\phi)\). Substitute this into
the already established actual triangular identity
(MA-MO.704) to prove (MA-MO.851) and the first
line of (MA-MO.852).

Every interior zero and every boundary jet is retained
under the explicit premises of this corollary.
MA-WP-001--003 therefore gives the full quotient
trace \(Z_\omega(g)\) on the actual Hilbert quotient.
On each finite zero block the diagonal value of
convolution is \(G(\rho-1/2)\), repeated according
to the full multiplicity. Nilpotent jet entries
do not alter a determinant. The trace-norm block
convergence proved in MA-WP, followed by Fredholm
determinant continuity, gives the stated infinite
products. They also converge locally uniformly
because \(\sum_\rho m_\rho|G(\rho-1/2)|<\infty\).
Thus neither a formally assigned divisor trace
nor a claim of simple zeros is used.

\emph{The rational relative operator on its correct invariant space.}
The bounded original negative projection is
\[
       P_-(u,c)=(u+c_0r_0+c_1r_1,0).
 \tag{MA-MO.859}
\]
Consequently \(F_g=U_gV\) has rank at most two.
The rank-one trace formula and the two actual
translation characters on the polar coordinates
give
\(\operatorname{Tr}F_g=G(-1/2)+G(1/2)=p(g)\).
This argument does not require the range of
\(1-P_-\) to be invariant under translation.
Since \(\mathscr D_g\) is trace class by
MA-MO-027, \(\mathscr E_g=F_g-\mathscr D_g\)
is trace class and agrees with the original
\(C_g(1-P_+)\) on the dense common source.

The actual source identity (MA-MO.447) is
\(p_+i_+=1+L\). Hence
\[
       C_g(1-P_+)i_+f=-i_+\mathcal D_gLf .
 \tag{MA-MO.860}
\]
The all-width return (MA-MO.829) supplies the
right side as \(-L_g\) on the completed
\(\mathcal W_+\). Density proves invariance
and the exact restriction, with trace
\(-\tau_0(g,\phi)\). The original \(C_gP_+\)
has range in the positive source image.
Passing to the actual quotient therefore
identifies the quotient action of
\(\mathscr E_g\) with \(C_g\). These are
exactly the source-level identities in
MA-MO-021; their narrow continuity is now
provided by MA-MO-027/028 instead of the
earlier strict half-plane margin.

In the orthogonal decomposition
\(\mathcal W=\mathcal W_+\oplus\mathcal W_+^\perp\),
\(\mathscr E_g\) is upper triangular with
diagonal blocks \(-L_g\) and the actual
quotient \(C_g\). Summing their absolutely
convergent diagonal traces and then taking
finite-rank compressions proves the trace
and determinant formulas for \(\mathscr E_g\).
Subtract its trace from \(\operatorname{Tr}F_g=p(g)\)
to obtain the raw trace in (MA-MO.855).

For the additional raw determinant formula,
\(\mathscr D_g=U_gV-\mathscr E_g\) gives
\[
  1-z\mathscr D_g
     =(1+z\mathscr E_g)
       \left(1-z(1+z\mathscr E_g)^{-1}U_gV\right).
 \tag{MA-MO.861}
\]
Fredholm determinant multiplicativity and the
finite-rank identity \(\det(1-AB)=\det(1-BA)\)
give a two-by-two determinant for the second
factor. Apply (MA-MO.855) with parameter
\(-z\) to the first factor. This proves
(MA-MO.857) on its stated resolvent set,
which contains a neighborhood of zero.
The finite-rank formula is meromorphic away
from the discrete exceptional set and agrees
there with the entire determinant on the
left; all singularities of its total product
are consequently removable. Indeed near an
exceptional point, the finite-dimensional
generalized eigenspace of the compact
\(\mathscr E_g\) gives the finite Laurent
principal part of its resolvent, while the
complementary resolvent is holomorphic.

\emph{Primitive prime evaluation and the value-zero restriction.}
For \(g(0)=0\), the unconditional primitive
explicit formulas and the actual local multiplier
calculation, in the already fixed conventions
of (MA-MO.484)--(MA-MO.486), are
\[
 \begin{aligned}
 Z_\eta(g)
    &=-\mathcal A_{\infty,\eta}(g)-B_\eta(g),\\
 \tau_{\eta,S}(g,\phi)
    &=-\mathcal A_{\infty,\eta}(g)-B_{\eta,S}(g),\\
 Z_0(g)&=p(g)-A_0(g)-B_0(g),\\
 \tau_0(g,\phi)&=-A_0(g).
 \end{aligned}
 \tag{MA-MO.862}
\]
Here \(A_0(g)\) abbreviates the same archimedean
term \(A_0(h)\) with \(h(e^x)=g(x)\) used
in MA-MO-021. The \(B\)'s are the full
primitive finite-prime distributions, and
\(B_{\eta,S}\) is their finite-place part
inside \(S\). In particular ramification of
the field alone does not delete an unramified
character factor. The negative test argument
has the conjugate character coefficient,
as written in (MA-MO.853).

The conductor and different monomials have
constant logarithmic derivatives and hence
pair with \(g(0)\); the compact residual is
also \(g(0)\int\psi\). Their disappearance
in (MA-MO.862) is exactly the stated
value-zero restriction. Their normalizations
remain present in (MA-MO.850) for general
tests. Substituting (MA-MO.862) into
(MA-MO.851) and (MA-MO.855) proves
(MA-MO.853) and (MA-MO.858), with their
positive prime signs.
\end{proof}

\paragraph{Scope and framework placement.}
The evaluated ideal trace is unconditional
in MA-MO-028. The full-matrix existence and
full quotient retention used here still
require precisely the divisor hypotheses
of MA-MO-027. In particular no new assertion
about the arithmetic zeros, and no
arbitrary no-margin sufficiency statement,
is a consequence of this corollary.

The rational Fredholm factorization for
the raw matrix retains the actual finite
rank-two factor in (MA-MO.857). The
rank-two outgoing obstruction is not
removed by relabeling the invariant
relative operator. In the inherited cote~1
framework, scans~079--080 and~111--112
through BR-TR-01/02, these are identities
of defined nuclear operators and their
actual invariant exact sequences. The
finite-length zero factors retain their
multiplicities without a diagonalizability
assumption; no new positivity or algebraic
weight claim is made.

\section{Full-source extension with uniformly bounded local multiplicity near an edge}
\label{sec:bounded-boundary-clusters}

The finite boundary divisor in MA-MO-027 can be replaced
by a uniform bound on the total multiplicity in short
rectangles near the edge. Infinitely many boundary zeros,
and infinitely many interior zeros approaching the edge,
are then allowed. The proof keeps each nearby group of
poles in one product; unstable partial-fraction
coefficients are never estimated separately.

Use the actual source settings, forward multipliers
\(A_\omega\), Fourier normalizations, step cutoff and
norms of (MA-MO.700), (MA-MO.705). Thus the setting is
either a nontrivial finite-order Hecke character with
its dual and a structural \(S\), or the rational
principal sector \(S=\{\infty\}\) with its actual
polar completion. Fix \(0<a<1/2\) and \(s\ge0\).
The new divisor hypotheses are
\[
 \begin{gathered}
 |\Re\lambda|\le a
       \quad\hbox{for every completed centered zero of each sector},\\
 \text{there exist }\delta>0,\ \ell>0,\ M\in\mathbb Z_{\ge0}
       \text{ such that for every }T\in\mathbb R,\\
 \sum_{\substack{\lambda:\ a-\delta\le\Re\lambda\le a\\
                         |\Im\lambda-T|\le2\ell}}
                         m_\omega(\lambda)\le M .
 \end{gathered}
 \tag{MA-MO.900}
\]
The bound is imposed for both \(\omega=\eta,\bar\eta\)
in the nontrivial case. Multiplicities are included.
There is no separation condition on distinct zeros.
Shrinking the parameters if necessary, assume
\[
 0<\delta<\tfrac18\min\{a,1/2-a\},
                  \qquad 0<\ell\le1,\qquad
                  0<\epsilon_0<\delta/8 .
 \tag{MA-MO.901}
\]
This shrinking cannot invalidate (MA-MO.900).
All constants concern the fixed field, characters,
place set, \(a,\delta,\ell,M,s\) and cutoff.

\paragraph{MA-MO-030 (bounded local multiplicity suffices for the full matrix).}
Assume (MA-MO.900) and
\[
                         s>2M-1 .
 \tag{MA-MO.902}
\]
When \(M=0\), this inequality holds for every \(s\ge0\).
For every compact smooth test \(g\), the original
\(C_g(P_+-P_-)\) has a unique trace-class extension
on \(H_{\Sigma,a,s}\) in the nontrivial sector and
on \(\mathcal W_{a,s}=H_{\Sigma,a,s}\oplus\mathbb C^2\)
in the rational principal sector. The extensions
are continuous in trace norm on each fixed compact
test-support space. They agree with the original
wide-inverse source formulas.

All completed zero multiplicities are retained.
The evaluated trace and determinant identities
(MA-MO.851)--(MA-MO.858) hold on these actual
operators, including the raw rational determinant's
finite-rank correction and the \(g(0)=0\) condition
for the stated prime formulas.

This is a sufficient condition with a threshold on
the \emph{total local multiplicity} \(M\), not a
claim that the individual boundary-jet condition
alone suffices. The ordinary local zero count
\(O(\log(2+|T|))\) does not imply the fixed bound
in (MA-MO.900). No arithmetic assertion that the
new condition holds is made.

\begin{proof}
\emph{Uniform holomorphic factors after removing nearby zeros.}
Write \(t_n=n\ell\), \(\nu_n=2+|n|\), and let
\(Z_{\omega,n}\) be the multiset of the zeros in
the rectangle appearing in (MA-MO.900) with \(T=t_n\).
Its cardinality with multiplicity \(k_{\omega,n}\)
is at most \(M\). Define
\[
 \begin{aligned}
 P_{\omega,n}(z)&=\prod_{\lambda\in Z_{\omega,n}}(z-\lambda),\\
 b_{\omega,n}(z)&=\frac{P_{\omega,n}(z)}{A_\omega(z)},\\
 \mathcal R_n&=\{z:\ a-\delta/2<\Re z<a+\delta/2,\
                         |\Im z-t_n|<3\ell/2\}.
 \end{aligned}
 \tag{MA-MO.903}
\]
An empty product is one. Every zero of \(A_\omega\)
in \(\mathcal R_n\) is included in \(Z_{\omega,n}\).
The other finite-function zeros, and the finite Euler
zeros on \(\Re z=-1/2\), are outside this rectangle.
The rational pole is outside it by (MA-MO.901).
Thus \(b_{\omega,n}\) is holomorphic in \(\mathcal R_n\).

There are constants \(C_j,D_j\), independent of \(n\),
such that
\[
 \sup_{\substack{a-\delta/4\le\Re z\le a+\delta/4\\
                       |\Im z-t_n|\le5\ell/4}}
                   |b_{\omega,n}^{(j)}(z)|
                         \le C_j\nu_n^{D_j}
                         \qquad(j\ge0).
 \tag{MA-MO.904}
\]
We verify this estimate without assuming separation
inside \(Z_{\omega,n}\). For sufficiently large
\(|t_n|\), use the uncentered finite-\(L\) variable
and a disk centered at \(w_n=2+it_n\), with a
zero-free boundary radius \(12<R_n<13\).
Its total zero multiplicity is \(O(\log\nu_n)\).
All its zeros are nontrivial at these heights.
The entire rectangle \(\mathcal R_n+1/2\)
lies in the disk's inner quarter.

Write \(L_\omega=H_nB_n\), where \(B_n\) is the
finite disk Blaschke product with multiplicities.
The finite-\(L\) polynomial strip bound gives
\(|H_n|\le C\nu_n^D\) on the disk by the maximum
principle. The Euler-product lower bound at
\(2+it_n\), and \(|B_n|\le1\), give
\(|H_n(w_n)|\ge c>0\). Applying Harnack's
inequality to
\(\log(C\nu_n^D)-\log|H_n|\) gives
\[
                        |H_n(w)|^{-1}\le C\nu_n^D
             \quad\hbox{on the inner half disk}.
 \tag{MA-MO.905}
\]
This is the complete finite-disk estimate used
in MT-M18 and (MA-MO.808).

Each normalized Blaschke factor has, up to a phase,
the form
\[
       B_{\rho,n}(w)=
       \frac{R_n(w-\rho)}
        {R_n^2-\overline{(\rho-w_n)}(w-w_n)} .
 \tag{MA-MO.906}
\]
For a zero in \(Z_{\omega,n}\), its numerator
\(w-\rho=z-\lambda\) is canceled exactly by
\(P_{\omega,n}\). The remaining factor
\((w-\rho)/B_{\rho,n}(w)\) is bounded by a
fixed constant on the inner half disk.
There are at most \(M\) such factors.

Every other disk zero is uniformly separated
from \(\mathcal R_n+1/2\). Indeed, either its
centered real part is less than \(a-\delta\),
giving distance at least \(\delta/2\), or
its ordinate differs from \(t_n\) by more
than \(2\ell\), giving distance at least
\(\ell/2\). There are no zeros strictly
to the right of \(a\). Thus the inverses
of all these remaining Blaschke factors
are bounded by a fixed constant depending
only on \(\delta,\ell\). Multiplying
\(O(\log\nu_n)\) such constants costs
only a fixed power of \(\nu_n\).
The deleted Euler factors and their
reciprocals are uniformly bounded on
this closed positive strip. Together
with (MA-MO.905), this proves the
polynomial bound for \(b_{\omega,n}\)
on \(\mathcal R_n\).

For the finitely many bounded-height \(n\)'s,
the canceled quotient is holomorphic on the
same rectangle and bounded on a slightly
smaller fixed rectangle. Its finite bounds
can be incorporated into the constants.
Cauchy estimates on a fixed-radius disk,
using the gap between the two displayed
rectangles, prove (MA-MO.904) for every
derivative. In particular no Laurent
coefficient at an individually isolated
zero has entered the estimate.

\emph{A uniform Sobolev bound for the whole pole product.}
Set \(r=s/2\). Enumerate \(Z_{\omega,n}\),
with repetitions for multiplicity, by
\(\lambda_j=a-d_j+i(t_n+\gamma_j)\),
where \(0\le d_j\le\delta\) and
\(|\gamma_j|\le2\ell\).
For \(k=k_{\omega,n}\ge1\), put
\[
 R_{n,\epsilon}(v)
       =\prod_{j=1}^k
                    (\epsilon+d_j+i(v-\gamma_j))^{-1},
                  \qquad 0<\epsilon\le\epsilon_0 .
 \tag{MA-MO.907}
\]
The character index is suppressed in this local calculation.
For \(k=0\), put \(R_{n,\epsilon}=1\).
For \(k\ge1\), each factor has the elementary
integral representation
\(\int_0^\infty e^{-(\epsilon+d_j)t}
                  e^{i\gamma_jt}e^{-ivt}\,dt\).
Multiplying the Laplace integrals shows
\[
 \begin{gathered}
 R_{n,\epsilon}(v)
       =\int_0^\infty e^{-ivt}e^{-\epsilon t}h_n(t)\,dt,\\
 |h_n(t)|\le\frac{t^{k-1}}{(k-1)!}\quad(t>0).
 \end{gathered}
 \tag{MA-MO.908}
\]
Here \(h_n\) is the convolution of the \(k\)
functions \(\mathbf1_{t>0}e^{-d_jt}e^{i\gamma_jt}\).
The bound is the volume of the \((k-1)\)-simplex,
and is uniform in all \(d_j,\gamma_j\), including
coinciding zeros and zeros on the edge.

The unitary plus Fourier transform of
\(R_{n,\epsilon}\) is
\(\sqrt{2\pi}e^{-\epsilon t}h_n(t)\).
Since \(r>M-\tfrac12\), (MA-MO.908) gives
\[
 \begin{gathered}
 \sup_{\substack{n,\epsilon\\k_{\omega,n}\ge1}}
                    \|R_{n,\epsilon}\|_{H^{-r}}\le C,\\
 R_{n,\epsilon}\longrightarrow R_{n,0}
             \quad\hbox{in }H^{-r}
             \text{ uniformly over the nonempty products}.
 \end{gathered}
 \tag{MA-MO.909}
\]
For the second assertion, the difference has
Fourier transform
\(\sqrt{2\pi}(e^{-\epsilon t}-1)h_n(t)\).
The square of its weighted norm is dominated
by a constant times
\[
 \int_0^\infty
    (1+t^{2M-2})(1+t^2)^{-r}
                         |e^{-\epsilon t}-1|^2\,dt
\]
when \(M\ge1\), with the bounded \(t\le1\)
part interpreted separately. This integrable
majorant and dominated convergence are
independent of \(n\) and of the pole locations.
For \(k=0\) there is no pole product; its
compactly localized contribution below is
handled directly as a smooth function.
Thus \(M=0\) also causes no exception.

\emph{Polynomially controlled localized boundary distributions.}
Choose a smooth partition of unity
\(\sum_{n\in\mathbb Z}\chi_n(u)=1\), with
\(\operatorname{supp}\chi_n\subset
[t_n-\ell,t_n+\ell]\), obtained by translating
and scaling one fixed compact profile.
For \(0<\epsilon\le\epsilon_0\), the exact
localized reciprocal is
\[
 \begin{aligned}
 d_{\omega,n,\epsilon}(u)
   &=\frac{\chi_n(u)}{A_\omega(a+\epsilon+iu)}\\
   &=\chi_n(u)b_{\omega,n}(a+\epsilon+iu)
                                  R_{n,\epsilon}(u-t_n).
 \end{aligned}
 \tag{MA-MO.910}
\]
The smooth compact multiplier
\(\chi_n b_{\omega,n}(a+\epsilon+i\,\cdot)\)
has polynomially bounded derivatives of
every fixed order, uniformly in \(\epsilon\),
by (MA-MO.904). The Sobolev multiplier
estimate (MA-MO.716), applied by duality
to \(H^{-r}\), therefore gives
\[
 \begin{gathered}
 \|d_{\omega,n,\epsilon}\|_{H^{-r}}
                                \le C\nu_n^D,\\
 \|d_{\omega,n,\epsilon}-d_{\omega,n,0}\|_{H^{-r}}
                    \le C\nu_n^D\omega(\epsilon),\qquad
                         \omega(\epsilon)\longrightarrow0 .
 \end{gathered}
 \tag{MA-MO.911}
\]
To see the second bound explicitly, split the
difference between the two smooth multipliers,
which is \(O(\epsilon\nu_n^D)\) in the
required finite derivative norms by Cauchy
estimates, and the pole-product difference
in (MA-MO.909). Both contributions have
the displayed bound after increasing \(D\).
When \(k=0\), the first contribution alone
proves both assertions. All translated
Sobolev norms are the same, so there is
no hidden dependence on the height.

The locally finite sum
\[
                         j_{\omega,+}=\sum_n d_{\omega,n,0}
 \tag{MA-MO.912}
\]
defines a tempered distribution. Indeed,
when paired with a Schwartz function,
insert a fixed translate of a compact
cutoff equal to one on
\(\operatorname{supp}\chi_n\). Its product
with the Schwartz test has \(H^r\) norm
decreasing faster than every power of
\(\nu_n\). This defeats (MA-MO.911)
and proves absolute convergence of the
pairing. The same argument proves
\[
 A_\omega(a+\epsilon+i\,\cdot)^{-1}
       \longrightarrow j_{\omega,+}
                  \quad\hbox{in tempered distributions},\qquad
 A_\omega(a+i\,\cdot)j_{\omega,+}=1 .
 \tag{MA-MO.913}
\]
For the last identity multiply before
taking the limit. The forward multiplier
and its derivatives have uniform polynomial
bounds on these strips, so its product
with a Schwartz test converges in Schwartz
space. Thus (MA-MO.913) is a legitimate
distributional identity.

\emph{Schwartz smoothing sums the entire infinite family in trace norm.}
If \(L(t,u)\in\mathscr S(\mathbb R^2)\), define
\[
       T_{\omega,\epsilon,L}f(t)
            =\sum_n
                 \langle d_{\omega,n,\epsilon}(u),
                                      L(t,u)f(u)\rangle .
 \tag{MA-MO.914}
\]
This defines a trace-class map \(H^r\to H^R\)
for each fixed \(R\ge0\). It converges in
trace norm as \(\epsilon\downarrow0\);
the same conclusion holds if
\(L_\epsilon\to L_0\) in Schwartz space.
Here are the summability details. Insert a
cutoff equal to one on \(\operatorname{supp}\chi_n\)
in the \(u\) variable and translate \(u=t_n+v\).
Translation is isometric on \(H^{\pm r}\).
Every fixed Schwartz seminorm of
\(L(t,t_n+v)\) on the fixed compact \(v\)
support is \(O_N(\nu_n^{-N})\), for every
\(N\), including the necessary polynomial
weights and derivatives in \(t\).
The explicit Hermite-series lemma
(MA-MO.715) and (MA-MO.911) consequently give
\[
  \bigl\|T_{d_{\omega,n,\epsilon},L}\bigr\|_{\mathcal S_1(H^r,H^R)}
                        \le C_{L,N}\nu_n^{D-N}.
 \tag{MA-MO.915}
\]
Choose \(N>D+2\) and sum. The identical
estimate for the differences, using
(MA-MO.911) and the Schwartz convergence
of \(L_\epsilon\), proves the asserted
trace-norm convergence. Only finitely
many Schwartz seminorms are required
for each specified pair of norms.

This constructs the \emph{smoothed} operator.
It does not assert that multiplication by
\(j_{\omega,+}\) is separately defined
as a bounded operator on every \(H^r\) input.

\emph{The boundary value is still the original wide inverse.}
For \(y\in\mathcal E\), let \(F_{\omega,y}\)
be its original series (MA-MO.477) or
(MA-MO.323). Since there are no zeros
strictly to the right of \(a\), the
original wide contour can be shifted
to every \(b>a\); on each such fixed
line it has a positive horizontal gap.
Thus the exact identity (MA-MO.711) holds.

For each Schwartz function \(f\),
(MA-MO.911) also proves that
\(j_{\omega,\epsilon}f\) converges in
\(H^{-r}\), not just in distributions:
the Sobolev multiplier norms of the
localized \(f\)'s decrease rapidly, and
the sum of their product norms is finite.
The same holds for \(u^jf(u)\) for every
integer \(j\ge0\).
Apply this to the uniformly Schwartz
family \(V_by\), with \(b=a+\epsilon\).
Taking Fourier inverses and using local
smooth convergence of \(e^{bx}F_{\omega,y}\)
proves
\[
 \begin{gathered}
      V_aF_{\omega,y}=j_{\omega,+}V_ay,\\
 D_x^j(e^{ax}F_{\omega,y})
             \in L^2((1+x^2)^{-r}\,dx)
                         \qquad(j\ge0).
 \end{gathered}
 \tag{MA-MO.916}
\]
This identifies the boundary prescription
with the original function, without crossing
an unrecorded pole.

For completeness, (MA-MO.916) gives
pointwise polynomial bounds. The elementary
one-dimensional \(H^1\) estimate on
\([x-1,x+1]\), and comparability of
\((1+t^2)^r\) on that interval, imply
\[
 |D_x^j(e^{ax}F_{\omega,y})(x)|
                   \le C_{j,y}(1+x^2)^{r/2}.
 \tag{MA-MO.917}
\]
The original function is smooth, so this
applies to its actual derivatives. At the
right it still has every stronger
exponential weight from the original
wide contour shift. Hence
\(qF_{\omega,y}\in\mathcal E_\omega\),
and is in the actual local source by
the structural-\(S\) inclusion. In the
rational case (MA-MO.714) is again an
actual even additive Schwartz function
vanishing near zero.

\emph{The actual two kernels and their compatible rows.}
Use exactly the removable commutator and
separated crossed kernels (MA-MO.718)--(MA-MO.719).
Their nonsingular factors
\(L_{\omega,b,g}\) and
\(L^{\rm cross}_{\eta,b,g}\)
converge in Schwartz space as \(b\downarrow a\).
The forward multipliers have polynomial
derivative bounds, \(G\) and \(P\) supply
the two rapid directions, and
\(\Re(-z-w)=-2b\le-2a\).
Their proofs do not use a bound for the
reciprocal. The division by \(z-w\)
in the commutator remains removable
before the distribution pairing.
Equation (MA-MO.913) supplies the same
identity \(A j_+=1\).
Applying (MA-MO.914)--(MA-MO.915) gives
trace-class limits
\[
     K_{\omega,a}^{\partial}(g),\qquad
     O_{\eta,a}^{\partial}(g):
               H_{a,s}\longrightarrow H_{a,s},
 \tag{MA-MO.918}
\]
with the appropriate dual input sector
for \(O_{\eta,a}^{\partial}\).
The input reciprocal in that crossed
kernel is \(j_{\bar\eta,+}\), as in
(MA-MO.719), not \(j_{\eta,+}\).

The independent half-line term
\(N_a(g)=C_gM_\phi:H_{\Sigma,a,s}\to H_{a,s}\)
has the trace-class physical kernel
(MA-MO.332). Therefore the first
character row is the actual
\[
       K_{\eta,a}^{\partial}(g)y
              -N_a(g)y+O_{\eta,a}^{\partial}(g)Ry .
 \tag{MA-MO.919}
\]
For every \(y\in\mathcal E\), the original
identity is first valid on every \(b>a\),
with its two fixed original wide inverses
and their actual cutoff source vectors.
The kernel convergence just proved and
\(V_by\to V_ay\) pass that identity
to (MA-MO.919). The actual smooth
theta output is fixed during this
passage. No bounded identity between
whole varying-width Hilbert spaces
is inferred.

The reflected dual row is
\(R\mathscr R_{\bar\eta,a}^{\partial}(g^\vee)R\).
It agrees with the first on the dense
smooth common sector by the actual
Fourier pair and nonpolar Poisson
identity. The paired trace-class
map has range in the closed diagonal
of \(H_{a,s}\oplus H_{-a,s}\), by
density. This proves the claimed
actual nontrivial endomorphism in
the inherited sum norm.

\emph{The rational moments and complete polar matrix.}
The only extra input for the rational
rows is boundedness of the two actual
cutoff moments. Let
\(p_a(x)=\phi(x)e^{(1/2-a)x}\), which
is Schwartz. The functional
\[
 \lambda_0(y)=
       \sum_n\big\langle d_{n,0}(u),
              \widehat p_a(-u)(V_ay)(u)\big\rangle,
                   \qquad\lambda_1(y)=\lambda_0(Ry)
 \tag{MA-MO.920}
\]
is bounded on \(H_{a,s}\), respectively
\(H_{-a,s}\). Indeed the local multiplier
norms of \(\widehat p_a(-u)\) decrease
faster than every power of \(\nu_n\),
and (MA-MO.911) makes this sum of
functional norms absolutely convergent.
For \(y\in\mathcal E\), passage from
the lines \(b>a\) proves precisely
the physical moments (MA-MO.723).
The polynomial left bound
(MA-MO.917) and the stronger right
weights give absolute convergence
at Mellin parameter \(1/2\).
The original zeta quotient has a
zero there, so
\(\mathcal MF_y(1/2)=0\).
Consequently the actual source
\(h_y=u_y+\mathfrak F u_{Ry}\)
has \(h_y(0)=-\lambda_1(y)\) and
\(\int h_y=-\lambda_0(y)\).

All ingredients of the exact
polar-subtracted formula (MA-MO.725)
are now bounded in the required
norms. Its two finite terms remain
\[
 \begin{gathered}
 G(-1/2)\lambda_1(y)e_0+
                   G(1/2)\lambda_0(y)e_1,\\
 \nu_g(y)=
   \bigl(-G(-1/2)\lambda_1(y),
         -G(1/2)\lambda_0(y)\bigr).
 \end{gathered}
 \tag{MA-MO.921}
\]
These signs come from the same
actual theta residue and the two
Gaussian source sections, without
a changed pole subtraction.
The first row is the sum of
(MA-MO.919), in the self-dual case,
and the displayed finite terms.
Its reflected row agrees on the
smooth core. The closed-diagonal
argument gives \(T_g^\partial:H\to H\)
in trace class. Appending the two
fixed genuine source columns gives
exactly the actual full matrix
\[
 \mathscr D_g(y,c)
       =(T_g^\partial y,\nu_g(y))
                             +\sum_{j=0}^1 c_jC_gDs_j .
 \tag{MA-MO.922}
\]
This proves the rational conclusion.

\emph{Continuity, full retention and evaluated traces.}
The original cores are dense, giving
uniqueness. The trace-norm summation
(MA-MO.915) uses a fixed finite
number of Schwartz seminorms for
each target norm. The kernel
seminorms are controlled by finitely
many derivatives of \(g\) on each
fixed compact support. This proves
the asserted trace-norm continuity.

Each right boundary multiplicity \(m\)
is at most \(M\). The completed
functional equation sends a left
boundary zero to a right boundary
zero of the dual sector, with the
same multiplicity; in the rational
case the function is self-dual.
Thus the same bound holds on the
left edge. Consequently (MA-MO.902)
implies \(s>2m-1\). Every boundary
jet is retained by MA-WP; every
interior zero is retained as well.
The actual nontrivial ideal is
invariant and its quotient action
is \(-C_g\), by the same smooth
source formula and density used
in (MA-MO.704). Its restriction
agrees with the unconditional
all-width return of MA-MO-028.
In the rational case the actual
relative operator
\(\mathscr E_g=C_g(1-P_-)-\mathscr D_g\)
has the invariant restriction
\(-L_g\) and quotient \(C_g\).
These are precisely the three
inputs used in the proof of
MA-MO-029: existence of the actual
full matrix, full quotient retention,
and the evaluated positive-source
trace. Its trace, finite-rank
determinant and primitive-prime
calculations therefore give
(MA-MO.851)--(MA-MO.858) here too.
No finite-divisor hypothesis is
used in those algebraic calculations.
\end{proof}

\paragraph{Comparison with the finite-boundary theorem and remaining scope.}
MA-MO-027 is contained in this result.
Choose \(\delta\) smaller than its
interior margin, so only the finitely
many edge zeros remain in the band.
Choose \(\ell\) so small that a
window of length \(4\ell\) contains
at most one distinct edge ordinate,
for either sector. Then \(M\) is
the maximum edge multiplicity, and
(MA-MO.902) is exactly the condition
there. If the edge divisor is empty,
take \(M=0\).

In contrast, (MA-MO.900) permits
infinitely many zeros in the band,
arbitrarily close horizontal
approach to the edge, and arbitrarily
close zero pairs within a window,
provided their \emph{total}
multiplicity stays bounded by \(M\).
The proof's Sobolev exponent pays
for that whole product. It asserts
no necessity of this cluster bound
for a particular arithmetic function.
If the near-edge counts grow with
height, the uniform estimate
(MA-MO.909) at a fixed \(s\) is not
available. The general no-margin
full-retention problem is still open.

The unconditional analytic ingredients
are the polynomial finite-\(L\)
strip bound, Euler-product lower
bound and local \(O(\log|T|)\)
zero count already proved and
located in MT-M18.5--9 and
MA-MO-028. The new assumption is
explicitly the stronger near-edge
count (MA-MO.900), not a conclusion
of that ordinary local count.
The source interfaces and phases
are the already checked
MA-MO-018/022/027/028.

In the inherited cote~1 framework,
scans~079--080 and~111--112 through
BR-TR-01/02, the local multiplicities
remain finite-length divisor data.
The convolution estimate and the
summed actual nuclear operators
supply the additional analytic
structure needed for traces and
determinants. No new arithmetic
positivity or zero-location claim
is inferred from those lengths.

\section{An exact reciprocal-growth criterion for the unrestricted one-sided operator}
\label{sec:reciprocal-growth-criterion}

This section concerns the actual one-sided Hilbert space and
compact-frequency core of MA-RS-001--007.  Its universal quantifier
is over smooth physical cutoffs near one fixed cutoff.  It does
not decide a prescribed individual cutoff, and it does not
assert that the arithmetic reciprocal has nonpolynomial growth.
The completed-ideal return of MA-MO-028 has a different domain
and remains unconditional.

Fix a finite set \(S\) of rational primes and
\(\sigma\in\mathscr G\) as in (MA-RS.2).  Retain
\[
\begin{gathered}
 a=\sigma-\tfrac12,\qquad
 \mathcal H_a=L^2(\mathbb R,e^{2ax}\,dx),\qquad
 V_a=\mathcal F_+M_{e^{ax}},\\
 \mathcal F_+f(t)=(2\pi)^{-1/2}\int e^{itx}f(x)\,dx,\qquad
 \mathscr C=C_c^\infty(\mathbb R),\qquad
 \mathscr D_a=V_a^{-1}\mathscr C,\\
 A(t)=\zeta_S(\sigma+it),\qquad j(t)=A(t)^{-1}.
\end{gathered}
\tag{MA-MO.1000}
\]
All derivatives of \(A\) have polynomial growth.  The
already proved generic-line estimates give
\[
 |j^{(k)}(t)|\le
       \exp(C_k\log^2(2+|t|)),\qquad
 \left|\frac{j'(t)}{j(t)}\right|
                       \le C(2+|t|)^D .
\tag{MA-MO.1001}
\]
The second assertion is (MA-RS.64), with the derivative
converted to the real frequency variable.  In particular
\(A\) and \(j\) are smooth and nonzero on the real line.
The maximal domain of \(M_j\) is
\(\{v\in L^2:jv\in L^2\}\).

Fix a nonzero, possibly complex, test
\(g\in C_c^\infty(\mathbb R)\), and write
\[
 G(z)=\int e^{zx}g(x)\,dx,\qquad m(t)=G(a+it).
\tag{MA-MO.1002}
\]
There is no restriction on \(g(0)\); in particular the
theorem applies to every nonzero test with \(g(0)=0\).
Choose a real smooth step \(\phi_0\), equal to one
far left and zero far right, such that
\(0\le\phi_0\le1\) and \(\phi_0'\) has compact support
and finite Gevrey order.  Choose a closed bounded
interval \(I\) with nonempty interior such that
\[
 0<\delta\le\phi_0(x)\le1-\delta\quad(x\in I).
\]
Let \(\mathscr D_I^\mathbb R\) be the real Frechet
space of real smooth functions supported in \(I\), with
seminorms \(\max_{0\le k\le N}\|\psi^{(k)}\|_\infty\).
Fix \(0<\epsilon<\delta\) and put
\[
 \mathscr U=\{\psi\in\mathscr D_I^\mathbb R:
                         \|\psi\|_\infty<\epsilon\},\qquad
 \phi_\psi=\phi_0+\psi,\qquad
 T_\phi=\mathcal F_+M_\phi\mathcal F_+^{-1}.
\tag{MA-MO.1003}
\]
Every \(\phi_\psi\) is a real smooth step with
\(0\le\phi_\psi\le1\) and compactly supported derivative.
Monotonicity and reflection symmetry are not imposed.

Since \(g\ne0\), the entire transform \(G\) is not
identically zero.  Choose a bounded open frequency
interval \(\Omega\) on which \(mA\ne0\), and fix once
and for all
\[
 0\ne q\in C_c^\infty(\Omega),\qquad
 w=\frac{q}{\overline{mA}}\in C_c^\infty(\Omega),\qquad
 f=\mathcal F_+^{-1}q .
\tag{MA-MO.1004}
\]
The vector \(q\) and the output-testing vector \(w\)
do not vary with the cutoff perturbation.  The function
\(f\) is entire and is not identically zero.

Throughout this section a step \(\phi\) has the normalized
endpoints: it is identically one far left and identically
zero far right.  For every such smooth bounded step with
\(\phi'\in C_c^\infty\), define on the same fixed core
\(\mathscr C\) the actual operator
\[
 K^\circ_{g,\phi}v
      =M_mT_\phi v-M_{mA}T_\phi M_jv,\qquad v\in\mathscr C .
\tag{MA-MO.1005}
\]
Both terms in this formula are defined and are Schwartz
functions: \(jv\) is compact smooth, \(T_\phi\) maps
Schwartz functions to Schwartz functions, and \(A\)
has polynomially bounded derivatives.  Thus (MA-MO.1005)
does not use a hypothetical bounded inverse or an assigned
trace.  On any vector where the sequential commutator
is defined, it agrees with
\(V_aC_gZ_{a,S}[J_{a,S},M_\phi]V_a^{-1}\).
It need not imply that this latter sequential expression
is defined on every vector of \(\mathscr C\).

\paragraph{MA-MO-031 (polynomial growth, universal cutoff scope, and actual failure witnesses).}
In the fixed setting above, the following conditions
are equivalent:
\[
\begin{array}{ll}
\mathrm{(P)}&
 |j(t)|\le C(2+|t|)^N
          \text{ for some }C,N\text{ and every }t;\\[2pt]
\mathrm{(D)}&
 T_{\phi_\psi}q\in\mathcal D(M_j)
                           \text{ for every }\psi\in\mathscr U;\\[2pt]
\mathrm{(B)}&
 K^\circ_{g,\phi_\psi}\text{ has a bounded extension to }L^2
                           \text{ for every }\psi\in\mathscr U;\\[2pt]
\mathrm{(T)}&
 K^\circ_{g,\phi_\psi}\text{ has a trace-class extension to }L^2
                           \text{ for every }\psi\in\mathscr U .
\end{array}
\tag{MA-MO.1006}
\]
In fact (P) implies a stronger positive conclusion:
for every compact smooth test \(h\) and every smooth
bounded step \(\phi\) with compactly supported derivative,
the original sequential commutator is defined on
\(\mathscr D_a\), has the Schwartz kernel (MA-RS.16),
and extends to a trace-class operator.  The actual
finite-frequency operators (MA-RS.54) converge to it
in trace norm.  Its trace and residue correction are
exactly (MA-RS.56).

If (P) fails, the following set is a dense
\(G_\delta\) subset of \(\mathscr D_I^\mathbb R\):
\[
 \{\psi:
          jT_{\phi_0+\psi}q\notin L^2\}.
\tag{MA-MO.1007}
\]
Consequently every neighborhood of zero contains a
real perturbation \(\psi\in\mathscr U\) for which
the original sequential commutator fails its inverse
domain already on the fixed core vector \(q\).
For the same cutoff the always-defined operator
(MA-MO.1005) has no bounded extension.  There are
explicit normalized vectors \(v_R\in\mathscr C\)
with
\[
       \|K^\circ_{g,\phi_\psi}v_R\|_2\longrightarrow\infty.
\tag{MA-MO.1008}
\]
For the actual operators \(K_N(g)\) of (MA-RS.54)
with this fixed cutoff,
\[
 \sup_N\|K_N(g)\|_{\mathcal H_a\to\mathcal H_a}=\infty .
\tag{MA-MO.1009}
\]
Their traces nevertheless converge to the value
in (MA-RS.56).  In particular there is no bounded
strong operator limit and no trace-norm limit in
this conditional failure case.

\begin{proof}
\emph{1. Polynomial reciprocal growth gives all smooth
domains and the actual kernel.}
Assume (P).  Repeatedly differentiate \(Aj=1\).
Every derivative \(j^{(k)}\) is a finite sum of products
of powers of \(j\) and derivatives of \(A\).
It therefore has polynomial growth, with an exponent
allowed to depend on \(k\).  Thus both \(A\) and \(j\)
are continuous multipliers of the Schwartz space.
For a core vector \(v\), the functions
\(T_\phi v\) and \(T_\phi jv\) are Schwartz.
It follows that \(M_jT_\phi v\) is in the maximal
domain of \(M_A\), and that \(T_\phi jv\) is also in
that domain.  Moreover
\(M_AM_jT_\phi v=T_\phi v\).
This checks the two actual sequential compositions
and gives (MA-MO.1005) as their difference.

For \(h\in C_c^\infty\), put
\(m_h(t)=\mathcal Mh(a+it)\), and put
\(p(\xi)=\int e^{i\xi x}\phi'(x)\,dx\).
Both are Schwartz functions.  The exact distribution
calculation (MA-RS.59)--(MA-RS.60) gives the kernel
\[
 k_{h,\phi}(t,u)
  =\frac{i}{2\pi}m_h(t)A(t)
       \left(\int_0^1j'(u+\theta(t-u))\,d\theta\right)
                                      p(t-u).
\tag{MA-MO.1010}
\]
Every derivative of the parenthesized function has
polynomial growth in \((t,u)\).  The variables \(t\)
and \(t-u\) form invertible linear coordinates on
\(\mathbb R^2\).  Schwartz decay of \(m_h(t)\) and
\(p(t-u)\), with all their derivatives, consequently
proves \(k_{h,\phi}\in\mathcal S(\mathbb R^2)\).

Expansion in the two-variable Hermite basis has
coefficients decaying faster than every power of
both indices.  The resulting absolutely summable
rank-one series proves trace class, and its trace
is the integral of the diagonal.  Equivalently this
is the Schwartz-kernel argument already used in
MA-RS-002/004.  This implication is continuous
from Schwartz kernels to trace class, since finitely
many oscillator seminorms bound the sum of the
absolute coefficients.

Let \(\chi_N(t)=\chi(t/N)\) be exactly the cutoffs
in MA-RS-007.  Formula (MA-RS.54) agrees with
the kernel \(\chi_N(t)k_{h,\phi}(t,u)\chi_N(u)\).
Since \(\chi=1\) near zero, this converges to
\(k_{h,\phi}\) in every Schwartz seminorm.
The preceding nuclear bound proves trace-norm
convergence.  The diagonal, its sign and its
residue evaluation are therefore exactly
(MA-RS.56), now as a limit of operators as well
as of traces.  This proves (P) implies (D) and (T);
(T) implies (B) directly.

\emph{2. A fixed output functional forces the fixed
inverse-domain condition.}
The step is real, so \(T_\phi\) is self-adjoint.
Use the Hilbert inner product linear in its first
variable.  For \(v\in\mathscr C\), the actual two
compositions give
\[
 \begin{split}
 \langle K^\circ_{g,\phi}v,w\rangle
 &=\left\langle v,\,
       T_\phi(\overline m\,w)
            -\overline j\,T_\phi q\right\rangle,\\
 H_\phi&=T_\phi(\overline m\,w)
            -\overline j\,T_\phi q .
 \end{split}
\tag{MA-MO.1011}
\]
The pairing on the right is well defined as an
ordinary integral over the compact support of \(v\).
Indeed \(H_\phi\) is a smooth function, though it
need not be in \(L^2\).  The identity uses precisely
\(\overline{mA}\,w=q\), with no restriction on the
phase of the test \(g\).

If \(K^\circ_{g,\phi}\) has a bounded extension,
the left side is a bounded functional of \(v\)
in \(L^2\).  The Riesz representation and equality
on all compact smooth vectors imply that the
smooth function \(H_\phi\) is in \(L^2\).
Its first term is already in \(L^2\), since
\(T_\phi\) is bounded.  Hence
\(\overline j\,T_\phi q\in L^2\), or equivalently
\(jT_\phi q\in L^2\).  This proves (B) implies (D).

\emph{3. Universal cutoff domains force a polynomial
local reciprocal bound.}
The fixed Gevrey step satisfies
\[
                 jT_{\phi_0}q\in L^2 .
\tag{MA-MO.1012}
\]
For completeness, outside a fixed interval containing
the support of \(q\), the transform \(T_{\phi_0}q\)
is the ordinary integral (MA-RS.20).  The Fourier
transform of \(\phi_0'\) has
\(\exp(-c|t|^\theta)\) decay for some \(\theta>0\).
This beats the \(\exp(C\log^2(2+|t|))\) growth
of \(j\), and proves (MA-MO.1012).

Assume (D).  Subtract (MA-MO.1012) and scale each
real perturbation by a sufficiently small nonzero
constant to put it in \(\mathscr U\).  One obtains
\[
 \mathcal L\psi
       :=j\,\mathcal F_+(\psi f)\in L^2
                  \quad(\psi\in\mathscr D_I^\mathbb R).
\tag{MA-MO.1013}
\]
Taking real and imaginary parts extends this
assertion to the complex Frechet space \(\mathscr D_I\).
The linear map \(\mathcal L:\mathscr D_I\to L^2\)
is continuous.  Here is an explicit Baire argument.
For each positive integer \(n\), the set
\(\{\psi:\|\mathcal L\psi\|_2\le n\}\) is closed:
convergence in \(\mathscr D_I\) gives pointwise
convergence of the Fourier transforms, and Fatou's
lemma gives the norm inequality.  These closed
sets cover \(\mathscr D_I\).  One has interior,
so taking differences gives a neighborhood of
zero on which the norm is bounded.  Homogeneity
therefore yields finite constants \(C,N\) with
\[
 \|j\mathcal F_+(\psi f)\|_2
       \le C\max_{0\le k\le N}\|\psi^{(k)}\|_\infty
                         \quad(\psi\in\mathscr D_I).
\tag{MA-MO.1014}
\]

The nonzero entire function \(f\) has only isolated
real zeros.  Choose a closed subinterval \(J\)
in the interior of \(I\) on which \(f\ne0\).
Take a nonnegative nonzero \(b\in C_c^\infty(J)\).
Its normalized Fourier transform is nonzero at
zero; for some \(\delta_1,c_1>0\),
\[
 |\mathcal F_+b(t)|\ge c_1\quad(|t|\le\delta_1).
\]
For each real \(T\), the compact smooth function
\[
          \psi_T(x)=\frac{b(x)}{f(x)}e^{-iTx}
\tag{MA-MO.1015}
\]
belongs to \(\mathscr D_I\), and its derivative
seminorm through order \(N\) is
\(O((2+|T|)^N)\).  Substitution in
(MA-MO.1014) gives
\[
 \int_{T-\delta_1}^{T+\delta_1}|j(t)|^2\,dt
                        \le C'(2+|T|)^{2N}.
\tag{MA-MO.1016}
\]
The use of complex \(\psi_T\) here is legitimate
because the real domain assertion was first extended
complex linearly; the failure cutoff constructed
below remains real.

\emph{4. Local \(L^2\) control implies pointwise
polynomial growth on the generic line.}
By (MA-MO.1001),
\[
 \left|\frac{d}{dt}\log|j(t)|\right|
                        \le C(2+|t|)^D.
\]
Choose
\[
 r_T=\min\{\delta_1/2,\ 1/2,\
                       [4C(3+|T|)^D]^{-1}\}.
\tag{MA-MO.1017}
\]
Enlarging \(C\) if necessary, integration on
\([T-r_T,T+r_T]\) proves
\(|j(t)|\ge e^{-1/2}|j(T)|\) there.
Combining with (MA-MO.1016) yields
\[
 2e^{-1}r_T|j(T)|^2
       \le C'(2+|T|)^{2N}.
\tag{MA-MO.1018}
\]
Since \(r_T\ge c(2+|T|)^{-D}\), this is a
polynomial pointwise bound for \(j(T)\).
Thus (D) implies (P), proving all equivalences.
The polynomial logarithmic-derivative estimate is
used exactly here; a smooth zero-free line without
that estimate is not substituted for the generic
hypothesis.

\emph{5. Real smooth failure cutoffs are residual.}
Suppose (P) fails.  On the real space
\(\mathscr D_I^\mathbb R\), define
\[
 E_n=\{\psi:
       \|j\mathcal F_+(\psi f)\|_2\le n\},
\tag{MA-MO.1019}
\]
where a nonintegrable square norm is assigned
\(+\infty\).  Each \(E_n\) is closed by the
same pointwise convergence and Fatou argument.
If any \(E_n\) had interior, its difference set
would imply that every sufficiently small real
perturbation has finite square norm.  Scaling
would then give that property for every real
perturbation, and complex linearity and
steps3--4 would prove (P), a contradiction.
Thus each \(E_n\) has empty interior.
Their union is exactly the inverse-domain set
for the perturbation part.  Its complement is
a dense \(G_\delta\), by Baire's theorem.
Adding the already integrable base term
(MA-MO.1012) leaves membership unchanged.
This proves (MA-MO.1007).

Choose a bad \(\psi\) in \(\mathscr U\).
The real cutoff retains all plateau and range
conditions in (MA-MO.1003), but \(M_jT_\phi q\)
does not exist in \(L^2\).  Thus the original
sequential expression is not defined on the
entire fixed core.  At the same time
the actual function \(H_\phi\) in
(MA-MO.1011) is not in \(L^2\).
Choose real \(\eta\in C_c^\infty\) with
\(0\le\eta\le1\), equal to one near zero,
and put \(\eta_R(t)=\eta(t/R)\).  For all
sufficiently large \(R\), define
\[
 v_R=\frac{\eta_R H_\phi}{\|\eta_R H_\phi\|_2}
                                      \in\mathscr C.
\tag{MA-MO.1020}
\]
The denominator tends to infinity: for each
fixed compact interval the cutoff is eventually
one, and \(H_\phi\notin L^2\).
By (MA-MO.1011) and \(0\le\eta_R\le1\),
\[
 \begin{split}
 |\langle K^\circ_{g,\phi}v_R,w\rangle|
 &=\frac{\int\eta_R|H_\phi|^2}
              {\|\eta_R H_\phi\|_2}\\
 &\ge\|\eta_R H_\phi\|_2\longrightarrow\infty .
 \end{split}
\tag{MA-MO.1021}
\]
This proves the actual normalized-core divergence,
including for arbitrary complex nonzero \(g\).
The physical vectors \(V_a^{-1}v_R\) have norm
one in \(\mathcal H_a\).  No assertion that
they belong to the two-sided completed ideal
is made.

\emph{6. The finite operators cannot converge,
although their traces do.}
Every finite operator (MA-RS.54) exists and is
trace class for this same smooth cutoff.  For
a fixed \(v\in\mathscr C\), \(\chi_Nv=v\)
for all sufficiently large \(N\); the same
is true of \(jv\).  Consequently the exact
bounded-composition formula gives
\[
 V_aK_N(g)V_a^{-1}v
        =\chi_NK^\circ_{g,\phi}v
                 \longrightarrow K^\circ_{g,\phi}v
                                      \quad\hbox{in }L^2 .
\tag{MA-MO.1022}
\]
If the operator norms of \(K_N(g)\) were
uniformly bounded, this identity on the
dense fixed core would give a bounded
extension of \(K^\circ_{g,\phi}\), contrary
to (MA-MO.1021).  Hence (MA-MO.1009) holds.
The uniform boundedness principle also rules
out a bounded strong operator limit.
Norm and trace-norm convergence are excluded
in particular.  The trace calculation of
MA-RS-007 uses only the polynomial logarithmic
derivative, not a polynomial reciprocal.
It therefore still gives the convergent
traces and their exact residue correction.
\end{proof}

\paragraph{A concrete sufficient zero configuration for the conditional failure branch.}
The analytic hypothesis that (P) fails can be implied
by the following explicit, but not asserted, arithmetic
configuration.  Fix \(R=1/4\).  Suppose there are
actual nontrivial zeros
\(\rho_n=\beta_n+i\gamma_n\), of multiplicity \(m_n\),
with \(|\gamma_n|\to\infty\), such that
\[
 \begin{gathered}
 0<d_n=|\sigma-\beta_n|<R/2,\\
 \frac{m_n\log(R/d_n)}
                   {\log(2+|\gamma_n|)}\longrightarrow\infty .
 \end{gathered}
\tag{MA-MO.1023}
\]
Then (P) fails and the real-cutoff/core/operator
conclusions above apply.

Indeed the fixed-strip upper bound gives
\(|\zeta(s)|\le C(2+|\gamma_n|)^C\) on the
circle \(|s-\rho_n|=R\), for sufficiently large
\(n\).  The disk contains no pole of zeta.
The function \(\zeta(s)/(s-\rho_n)^{m_n}\)
is holomorphic throughout it, regardless of
other zeros in the disk.  The maximum principle
therefore gives
\[
 |\zeta(\sigma+i\gamma_n)|
        \le C(2+|\gamma_n|)^C(d_n/R)^{m_n}.
\tag{MA-MO.1024}
\]
The finitely many Euler factors \(E_S\) are
uniformly bounded and bounded away from zero
on the fixed line \(\sigma>0\).  Thus
\[
 |j(\gamma_n)|
       \ge c(2+|\gamma_n|)^{-C}(R/d_n)^{m_n},
\tag{MA-MO.1025}
\]
which exceeds every fixed polynomial along
this sequence under (MA-MO.1023).
This is a conditional implication, not evidence
that such zeros or multiplicities occur.
The generic-distance bound (MA-RS.2) alone
does not establish (MA-MO.1023), and no result
in this section asserts nonpolynomial growth
of an actual arithmetic reciprocal.

\paragraph{Source, completion and remaining scope.}
The source for the actual closed multipliers,
sequential domains and Fourier sign is
MA-RS.8--9 and MA-RS.15--21; the finite
operators and polynomial logarithmic-derivative
input are MA-RS.52--64.  Their finite-disk
and selected-height input is the supplied
MT-M18 proof, (M18.7)--(M18.9), already
audited in those statements.  The positive
Schwartz-kernel argument also recovers
MA-RS-004 when its fixed zero gap supplies
(P).  No new universal zero estimate is
imported here.

In the rational setting of this section, the
user-stipulated input \(A\) supplies the fixed
gap \(\sigma-1/2>0\), and MA-RS-004 therefore
proves (P) and all the positive operator
conclusions under \(A\).  The unresolved
arithmetic issue below is an independent
proof without that zero-location input.
Using this \(A\)-supported implication to
justify an independent narrow trace argument
would retain the same arithmetic premise.

The theorem sharpens the unrestricted operator
gate in \texttt{MA-Q1-RF002}/\texttt{MA-Q1-RS003}
to a precise analytic condition for the stated
universal cutoff class.  It does not independently
decide that arithmetic condition without
zero-location input, an individual
predetermined cutoff, or singular boundary
lines.  Its normalized countervectors are
vectors of the one-sided Hilbert core; they
are not asserted to lie in the actual
two-sided completed theta ideal.  In
MA-MO-028 that ideal cancels the full divisor
before contour deformation, so its
unconditional return is compatible with
every assertion here.  The polynomial
partial-return costs of MA-PR and the
full-source divisor premises of
MA-MO-027/030 remain different questions.

In the inherited Grothendieck trace framework,
the relevant additional structure is the
actual continuous nuclear map and its
domain, not a scalar assignment of its
trace.  The arithmetic Euler determinants
and finite-length divisor factors are
unchanged.  The trace/determinant programme
recorded in cote~1 BR-TR-01/02, scans~079--080
and~111--112, supplies no implication from
convergent scalar traces to bounded
operator convergence; (MA-MO.1009) and
(MA-RS.56) explicitly separate those
properties in the conditional failure
case.  These archive locators are inherited,
not claimed freshly read for this theorem.

\section{The reciprocal criterion with the physical cutoff fixed}
\label{sec:fixed-cutoff-reciprocal-criterion}

MA-MO-031 varies the cutoff while keeping the test fixed.
Here the cutoff is fixed in advance.  The test varies in
one specified compact-support Frechet space, and the
operator remains the actual fixed-core composition.
No nonpolynomial growth of an arithmetic reciprocal
is asserted.

Retain the generic rational line, finite \(S\),
\(\mathcal H_a,V_a,A,j,\mathscr C\) of (MA-MO.1000).
In particular
\[
 |A^{(k)}(t)|\le C_k(2+|t|)^{D_k},\qquad
 \left|\frac{j'(t)}{j(t)}\right|
                         \le C(2+|t|)^D,\qquad
 |j(t)|\le\exp(C\log^2(2+|t|)).
\tag{MA-MO.1100}
\]
Fix once and for all a real smooth bounded cutoff
\(\phi\), identically one far left and identically
zero far right, with \(\phi'\in C_c^\infty\).
The proof permits \(0\le\phi\le1\), including any
predetermined monotone cutoff of that kind; no
cutoff perturbation will be made.  Put
\[
 T_\phi=\mathcal F_+M_\phi\mathcal F_+^{-1},\qquad
 p(v)=\int e^{ivx}\phi'(x)\,dx,\qquad
 p(0)=-1 .
\tag{MA-MO.1101}
\]
Let \(\mathscr D_K\) be the complex Frechet space
of smooth functions supported in \(K=[1,2]\),
with its usual derivative seminorms, and let
\(\mathscr D_K^\mathbb R\) be its real subspace.
All these tests satisfy \(g(0)=0\).
For \(g\in\mathscr D_K\), write
\[
 m_g(t)=\mathcal Mg(a+it),\qquad
 K^\circ_gv
   =M_{m_g}T_\phi v-M_{m_gA}T_\phi M_jv,
                           \quad v\in\mathscr C .
\tag{MA-MO.1102}
\]
Both terms are literal \(L^2\) functions on this
fixed core, as verified in (MA-MO.1005).

\paragraph{MA-MO-032 (a predetermined cutoff and a fixed test-support space).}
The following are equivalent:
\[
\begin{array}{ll}
\mathrm{(P)}&
 |j(t)|\le C(2+|t|)^N\text{ for some }C,N;\\[2pt]
\mathrm{(B_{\mathbb C})}&
 K^\circ_g\text{ extends boundedly to }L^2
                         \text{ for every }g\in\mathscr D_K;\\[2pt]
\mathrm{(B_{\mathbb R})}&
 K^\circ_g\text{ extends boundedly to }L^2
                         \text{ for every }g\in\mathscr D_K^\mathbb R;\\[2pt]
\mathrm{(T)}&
 K^\circ_g\text{ extends to a trace-class operator}
                         \text{ for every }g\in\mathscr D_K.
\end{array}
\tag{MA-MO.1103}
\]
When these conditions hold, the original sequential
commutator is defined on the whole fixed core for
every compact smooth test, including tests with
other compact supports.  It has the Schwartz kernel
of MA-RS-002, and the actual finite-frequency
operators of MA-RS-007 converge in trace norm.
The operator depends continuously in trace norm
on the test on each fixed compact support.

If (P) fails, the set of \(g\) for which
\(K^\circ_g\) has no bounded extension is a dense
\(G_\delta\) in each of \(\mathscr D_K\) and
\(\mathscr D_K^\mathbb R\).  For every such test
there are unit vectors \(v_n\in\mathscr C\)
with \(\|K^\circ_gv_n\|_2\to\infty\), and the
finite operators for this same fixed \(g,\phi\)
satisfy
\[
                   \sup_N\|K_N(g)\|=\infty ,
\tag{MA-MO.1104}
\]
although their scalar traces converge by MA-RS-007.

If the predetermined \(\phi'\) has finite Gevrey
order, the original sequential composition
\[
                C_gZ_{a,S}[J_{a,S},M_\phi]
\tag{MA-MO.1105}
\]
is defined on \(\mathscr D_a\) for every
\(g\in C_c^\infty\), without assuming (P).
Thus in the conditional failure branch the
unbounded operators are genuine sequential
commutators on their entire fixed core, not
merely alternate formulas used after a failed
inverse-domain condition.

\begin{proof}
\emph{1. A polynomial lower value at a nearby
ordinate, derived from the finite disk.}
Choose \(c>0\), with \(c<1/8\), so small that
\[
 |p(v)|\ge\tfrac12
                      \quad(c/2\le v\le3c).
\tag{MA-MO.1106}
\]
This follows from \(p(0)=-1\) and continuity.
For every sufficiently large \(|T|\), there
is \(v_T\in(c,2c)\) for which, putting
\(t_T=T+v_T\),
\[
                    |A(t_T)|\ge c_0(2+|T|)^{-D_0},
\tag{MA-MO.1107}
\]
with constants independent of \(T\).
We give the full lower-bound argument.

In the raw zeta variable use a disk centered
at \(w_T=2+iT\), with a zero-free boundary
radius \(12<R_T<13\).  Its zero count with
multiplicities is \(M_T=O(\log(2+|T|))\).
For large \(|T|\) it contains neither the
pole nor trivial zeros.  Write
\(\zeta=H_TB_T\), where \(B_T\) is the finite
disk Blaschke product with multiplicities.
The fixed-strip polynomial upper bound and
\(|B_T|=1\) on the boundary give a polynomial
upper bound on \(|H_T|\) throughout the disk.
At its center,
\[
 |H_T(w_T)|\ge|\zeta(2+iT)|\ge\zeta(2)^{-1}.
\]
Harnack's inequality applied to a positive
constant minus \(\log|H_T|\) therefore gives
\[
 |H_T(w)|^{-1}\le C(2+|T|)^C
                       \quad\hbox{on the inner half disk}.
\tag{MA-MO.1108}
\]
This is the actual finite-disk argument of
MA-RS-001, not an assertion of a polynomial
reciprocal along the entire original line.

The points \(w=\sigma+i(T+v)\) with
\(v\in[c,2c]\) lie in the inner quarter.
For a disk zero \(\rho=\beta+i\gamma\),
the normalized factor satisfies
\[
 |B_{\rho,T}(w)|\ge\frac{|w-\rho|}{2R_T}
                 \ge\frac{|T+v-\gamma|}{2R_T}.
\]
Put
\[
 V_T(v)=\sum_{\rho\ {\rm in\ the\ disk}}
                 m_\rho\log^+\frac1{|T+v-\gamma|}.
\tag{MA-MO.1109}
\]
For every real \(\gamma-T\),
\[
 \int_c^{2c}\log^+\frac1{|v-(\gamma-T)|}\,dv
       \le\int_{-1}^1\log(1/|u|)\,du=2 .
\]
Hence \(\int_c^{2c}V_T(v)\,dv\le2M_T\).
There is a point \(v_T\) outside the finitely
many exact ordinates with
\(V_T(v_T)\le(2/c+1)M_T\).
Multiplying the factor estimates yields
\[
 |B_T(\sigma+it_T)|
       \ge(2R_T)^{-M_T}e^{-V_T(v_T)}
                       \ge c(2+|T|)^{-C}.
\tag{MA-MO.1110}
\]
Together with (MA-MO.1108) this gives the
same lower bound for \(|\zeta(\sigma+it_T)|\).
The finitely many factors of \(E_S\) are
bounded and bounded away from zero on the
fixed line \(\sigma>0\), proving
(MA-MO.1107).  No separation among disk
zeros, or simplicity assumption, was used.

\emph{2. Actual separated packets and quantitative
noncancellation.}
Fix a nonnegative \(b\in C_c^\infty((1,2))\)
with \(\int b=1\), and put
\[
 \begin{gathered}
 B(v)=\int e^{ivx}b(x)\,dx,\qquad B(0)=1,\\
 g_T(x)=e^{-ax}e^{-it_Tx}b(x),\qquad
 m_{g_T}(t)=B(t-t_T).
 \end{gathered}
\tag{MA-MO.1111}
\]
These are tests in the one fixed space
\(\mathscr D_K\).  For every integer \(N\),
\[
 \max_{k\le N}\|g_T^{(k)}\|_\infty
                           \le C_N(2+|T|)^N .
\tag{MA-MO.1112}
\]
The off-diagonal distribution kernel of
\(T_\phi\) is the ordinary smooth function
\[
                    \kappa(v)=\frac{i\,p(v)}{2\pi v}
                                      \quad(v\ne0).
\tag{MA-MO.1113}
\]
By (MA-MO.1106), \(\kappa\) is bounded away
from zero, with bounded derivatives, on
\([c/2,3c]\).  Its possible delta component
at zero never enters the separated supports
used below.

There are fixed \(h_0,D_1>0\) such that
\[
                        h_T=h_0(2+|T|)^{-D_1}
\tag{MA-MO.1114}
\]
satisfies all the following for large \(|T|\).
If \(|u-T|\le h_T\) and
\(|t-t_T|\le h_T\), then \(t-u\in[c/2,3c]\)
and each of the four ratios
\[
 B(t-t_T),\qquad
 \frac{A(t)}{A(t_T)},\qquad
 \frac{j(u)}{j(T)},\qquad
 \frac{\kappa(t-u)}{\kappa(v_T)}
\tag{MA-MO.1115}
\]
differs from one by at most \(1/100\).
Here is a quantitative choice argument.
The first and fourth ratios obey the
required estimate on sufficiently small
fixed neighborhoods, by their bounded
derivatives and nonzero lower bounds.
For the second, the polynomial bound for
\(A'\) and (MA-MO.1107) give
\[
 \left|\frac{A(t)-A(t_T)}{A(t_T)}\right|
                         \le C(2+|T|)^{D_A+D_0}h_T .
\]
For the third, integrate \(j'/j\) along
the real segment; since \(j\ne0\),
\[
 \left|\frac{j(u)}{j(T)}-1\right|
       \le \exp(C(2+|T|)^D h_T)-1 .
\]
Choose \(D_1\ge\max\{D_A+D_0,D\}+1\),
and then \(h_0\) small enough, also
less than \(c/8\).  This proves every
claim in (MA-MO.1115).  The product of
the four ratios differs from one by
less than \(1/4\), so its real part
is at least \(3/4\).

Choose a nonnegative \(\alpha\in C_c^\infty((-1,1))\)
with \(\|\alpha\|_2=1\).  Use the Hilbert inner
product linear in its first variable, and set
\[
 \begin{gathered}
 u_T(u)=h_T^{-1/2}\alpha((u-T)/h_T),\\
 C_T=A(t_T)j(T)\kappa(t_T-T),\qquad
 w_T(t)=\frac{C_T}{|C_T|}
               h_T^{-1/2}\alpha((t-t_T)/h_T).
 \end{gathered}
\tag{MA-MO.1116}
\]
All factors in \(C_T\) are nonzero.
The packets are compact smooth, have
norm one, and have separated supports.

For clarity let
\(\mathcal S_T=M_{m_{g_T}A}T_\phi M_j\)
on the core.  The actual separated-support
kernel, the phase in \(w_T\), and
(MA-MO.1115) give
\[
 \begin{split}
 \operatorname{Re}\langle\mathcal S_Tu_T,w_T\rangle
 &\ge\frac34|A(t_T)j(T)\kappa(t_T-T)|\,h_T
                                  \left(\int\alpha\right)^2\\
 &\ge c_1(2+|T|)^{-D_0-D_1}|j(T)|.
 \end{split}
\tag{MA-MO.1117}
\]
This calculation is an ordinary double
integral on two compact separated intervals.
All compositions are defined before that
calculation.  It proves noncancellation by
a common phase and positive packet profiles;
it does not estimate an unsigned kernel
and assume that its operator norm is large.

The other term in (MA-MO.1102) has norm
at most \(\|B\|_\infty\|\phi\|_\infty
\le\|\phi\|_\infty\), independent of \(T\).
Thus, whenever \(K^\circ_{g_T}\) has a
bounded extension, and also as a bound
on its core operator norm if not,
\[
 \|K^\circ_{g_T}\|
       \ge c_1(2+|T|)^{-D_0-D_1}|j(T)|
                              -\|\phi\|_\infty .
\tag{MA-MO.1118}
\]

\emph{3. The universal fixed-support condition
has a finite derivative estimate.}
Assume \(\mathrm{(B_{\mathbb C})}\), and denote
the unique bounded extensions by \(\mathcal K(g)\).
The map
\(\mathcal K:\mathscr D_K\to\mathcal B(L^2)\)
is linear and has closed graph.
Indeed, suppose \(g_n\to g\) in \(\mathscr D_K\)
and \(\mathcal K(g_n)\to T\) in operator norm.
The Mellin transforms converge in every
Schwartz seminorm.  For each core vector
\(v\), both \(T_\phi v\) and \(T_\phi jv\)
are Schwartz, and \(A\) is a polynomial
Schwartz multiplier.  Consequently
\[
           K^\circ_{g_n}v\longrightarrow K^\circ_gv
                                      \quad\hbox{in }L^2.
\tag{MA-MO.1119}
\]
The operator-norm convergence gives
\(Tv=K^\circ_gv\).  Density and uniqueness
of the bounded extension imply
\(T=\mathcal K(g)\), proving closedness.

The finite-seminorm consequence can also
be proved directly, to make its uniform
content explicit.  The sets
\[
 F_n=\{g\in\mathscr D_K:
        \|K^\circ_gv\|_2\le n\|v\|_2
                              \text{ for every }v\in\mathscr C\}
\tag{MA-MO.1120}
\]
are closed, by (MA-MO.1119), and cover
\(\mathscr D_K\) under the assumption.
Baire's theorem gives interior to one
\(F_n\).  Differences give a neighborhood
of zero with operator norm at most \(2n\).
Homogeneity then yields \(C,N\) with
\[
              \|\mathcal K(g)\|
                  \le C\max_{k\le N}\|g^{(k)}\|_\infty
                                      \quad(g\in\mathscr D_K).
\tag{MA-MO.1121}
\]
Thus no uniform estimate over tests has
been assumed without proof.

Apply (MA-MO.1121) to \(g_T\) and use
(MA-MO.1112), (MA-MO.1118).  One obtains
\[
                    |j(T)|\le C'(2+|T|)^{N+D_0+D_1}
\tag{MA-MO.1122}
\]
for large \(|T|\).  The compact range is
bounded by smooth nonvanishing on the line.
This proves (P).

If only \(\mathrm{(B_{\mathbb R})}\) is
assumed, write each complex test as the
sum of its real part and \(i\) times its
imaginary part.  Core linearity gives a
bounded extension for the sum.  Hence
\(\mathrm{(B_{\mathbb C})}\) follows, and
the complex modulated tests in
(MA-MO.1111) are now legitimate.  This
does not silently assume complex test
bounds from a real assertion.

Conversely, (P) gives polynomial growth
of every derivative of \(j\), as proved
in MA-MO-031.  The actual sequential
domain and Schwartz-kernel proof there
applies to this fixed cutoff and every
compact smooth test.  It proves (T),
the trace-norm finite-frequency convergence
and test continuity.  Since (T) implies
\(\mathrm{(B_{\mathbb C})}\), all
conditions in (MA-MO.1103) are equivalent.

\emph{4. Residual genuine test failures
when the reciprocal bound fails.}
Suppose (P) fails.  Define the closed
sets \(F_n\) in (MA-MO.1120) without
assuming that they cover the test space.
Their union is exactly the set of tests
whose core operator has a bounded
extension: the core is dense, and a
uniform bound on it gives that unique
extension.  No \(F_n\) has interior.
Otherwise differences and scaling would
make every test have a bounded extension,
contrary to step3.  Thus their complement
is a dense \(G_\delta\).  The identical
argument on the real Frechet space works,
because the resulting universal real
bound would complexify as above.

For every test in this complement,
unboundedness on the fixed core supplies
unit \(v_n\in\mathscr C\) with
\(\|K^\circ_gv_n\|_2\ge n\).
These are actual core vectors and actual
outputs of (MA-MO.1102).  The choice of
the fixed bad test is by Baire's theorem;
no explicit numerical arithmetic example
is claimed.  Each test is supported in
\([1,2]\), so its value at zero vanishes.

For this same fixed \(g,\phi\), the two
bounded compositions of MA-RS-007 obey
\[
 V_aK_N(g)V_a^{-1}v
              =\chi_NK^\circ_gv
                       \longrightarrow K^\circ_gv
                       \quad(v\in\mathscr C)
\tag{MA-MO.1123}
\]
once \(N\) is large enough for its input
cutoff to equal one on the compact support.
Uniformly bounded operator norms would
therefore bound the core operator, a
contradiction.  This proves (MA-MO.1104)
and excludes strong or trace-norm
convergence to a bounded operator.
The polynomial logarithmic derivative
still gives the actual convergent trace
formula (MA-RS.56), with its residue
correction unchanged.

\emph{5. A fixed Gevrey cutoff keeps the
original sequential domain for every smooth test.}
This last assertion does not use (P).
For \(v\in\mathscr C\), both \(v\) and
\(jv\) are compact smooth.  Outside their
fixed frequency supports, the formula
(MA-RS.20) expresses \(T_\phi v\) and
\(T_\phi jv\) using \(p(t-u)/(t-u)\).
Gevrey decay of \(p\) dominates the
\(\exp(C\log^2(2+|t|))\) growth of \(j\).
Thus \(M_jT_\phi v\in L^2\), and
\(M_AM_jT_\phi v=T_\phi v\in L^2\).
Also \(T_\phi jv\) is in the maximal
domain of \(M_A\).  This checks both
sequential compositions on the entire
core.  The left convolution \(C_g\)
is bounded on \(\mathcal H_a\) for
every compact smooth \(g\), with norm
at most \(\int|g(x)|e^{ax}\,dx\).
No Gevrey estimate on \(g\) enters these
domain checks.  Consequently (MA-MO.1105)
is genuinely defined and equals
\(V_a^{-1}K^\circ_gV_a\) on \(\mathscr D_a\).
Its conditional unboundedness follows
from the preceding actual core witnesses.
\end{proof}

\paragraph{Arithmetic inputs and exact remaining question.}
The stipulated rational input \(A\) places
all nontrivial zeros on the critical line.
For every \(\sigma\in(1/2,1)\) considered
here, it therefore supplies a fixed
horizontal gap and proves (P) by
MA-RS-004.  Under \(A\), every positive
conclusion above is available.  The
failure branch is an explicit conditional
operator theorem; it neither asserts
nonpolynomial reciprocal growth nor a
zero configuration causing it.  The
sufficient multiplicity/distance condition
of (MA-MO.1023) remains conditional too.

Compared with MA-MO-031, the new result
settles the quantifier in which the
physical cutoff is prescribed and the
compact test varies.  It still does
not decide boundedness for one fixed
individual pair \(g,\phi\), or prove
the reciprocal bound independently of
zero-location input.  It is not an
unconditional arithmetic resolution
of RF002/RS003 or FT002.  Singular
boundary distributions and the actual
two-sided ideal return have their
separately proved domains; no packet
above is asserted to belong to that
completed ideal.

The finite-disk argument was derived
in (MA-MO.1108)--(MA-MO.1110) from the
actual polynomial strip bounds, Euler
lower bound and local zero count used
in the supplied MT-M18.7--9 and
MA-RS-001.  The normalized off-diagonal
kernel is (MA-RS.20)/(MA-RS.53).
The core and finite operators are
exactly (MA-RS.8--9)/(MA-RS.54);
the polynomial logarithmic derivative
is (MA-RS.64).  No whole-line
polynomial reciprocal estimate is
hidden in the selected nearby ordinate.

In the inherited cote~1 BR-TR-01/02
framework, scans~079--080 and~111--112,
the distinction remains one between
actual defined continuous/nuclear
maps and a scalar trace distribution.
The new fixed-cutoff criterion concerns
that analytic structure.  It changes
no finite-monodromy representation,
Euler determinant or motivic weight,
and supplies no positivity assertion
from convergence of scalar traces.
These archive locators are inherited.

\section{An actual-source reciprocal bound and a conditional interior-cluster obstruction}
\label{sec:full-source-cluster-obstruction}

\paragraph{Status.}
The complete author proof reading and two independent whole-proof
readings have passed.  Their exact source and final-scope check
records are kept separately.

This result concerns the original full nontrivial common-sector
map, not the completed ideal return or the whole-quotient return.
It extracts a quantitative necessary condition from a family of
genuine compact inputs.  The cluster conclusion below is conditional
on the stated zero configuration; no occurrence of that configuration
for a Hecke function is asserted.

Keep the number field, nontrivial primitive finite-order character
$\eta$, dual character, structural $S$, and actual source maps of
(MA-MO.550)--(MA-MO.552).  In particular $\phi$ is a fixed real
smooth step, equal to one sufficiently far left and zero sufficiently
far right, and $q=1-\phi$.  Write
\[
\begin{gathered}
 A_\omega(z)=E_{\omega,S}(1/2+z)L_K(1/2+z,\omega),
       \qquad \omega\in\{\eta,\bar\eta\},\\
 H=H_{\Sigma,a,s},\qquad
 w_{a,s}(x)=(e^{2ax}+e^{-2ax})(1+x^2)^{s/2},\\
 0<a<1/2,\qquad s\ge0,\qquad
 p=\frac{s+1}{2},\qquad
 0<\delta_0<\min\{1,(1/2-a)/2\},\\
 D_g^0=C_g(P_+-P_-):\mathcal E\longrightarrow\mathcal E,
       \qquad \mathcal D=C_c^\infty([1,2]).
\end{gathered}
\tag{MA-MO.1200}
\]
The meaning of $\mathcal D$ is the closed Fr\'echet space of
smooth functions on $\mathbb R$ supported in $[1,2]$.
Its real subspace will also be used.  Every test in this space
has $g(0)=0$.

Assume, for the comparison below, that every completed nontrivial
zero of both sectors lies in $|\Re z|\le a$ in centered coordinates.
The additional full-retention condition means that every zero on
an edge has multiplicity $m$ satisfying $m<p$.  The latter condition
is necessary by MA-MO-025, but is not used to assume any inverse
estimate near an edge.  Put $\nu(T)=2+|T|$ and define the condition
\[
 \begin{gathered}
 (\mathrm{BR}_{a,s,\eta}):\quad
 \exists C,N<\infty\quad
 \delta^p|A_\eta(a+\delta+iT)|^{-1}
       \le C\nu(T)^N\\
 \hfill(T\in\mathbb R,\quad0<\delta\le\delta_0).
 \end{gathered}
 \tag{MA-MO.1201}
\]
This is a two-variable estimate from the open right half-plane.
It does not assign a boundary reciprocal multiplier or a boundary
distribution to arbitrary Hilbert data.

\paragraph{MA-MO-033 (a full-source family requires boundary-uniform reciprocal control).}
Under the closed-strip premise above, if $D_g^0$ has a bounded
extension on $H$ for every complex $g\in\mathcal D$, then
$(\mathrm{BR}_{a,s,\eta})$ holds.  The same conclusion follows
if boundedness is assumed only for every real $g\in\mathcal D$.
The positive-support family controls the right edge of $\eta$.
By reflection, boundedness for the negative-support $\eta$ family
controls the right edge of $\bar\eta$, equivalently the left
edge of $\eta$.  No dual bound is inferred from the positive
family alone.

If (MA-MO.1201) fails, there is a fixed sequence
\[
 e_n\in C_c^\infty(\mathbb R)\subset\mathcal E,
           \qquad\|e_n\|_H=1,
 \tag{MA-MO.1202}
\]
such that the set of $g\in\mathcal D$ for which
$\sup_n\|D_g^0e_n\|_H=\infty$ is a dense $G_\delta$.
This holds separately in the real and complex test spaces.
Thus these fixed tests admit no bounded extension of their actual
source action.  The inputs in (MA-MO.1202) are normalized finite
translation-orbit integrals, specified in the proof.

A sufficient, explicitly hypothetical configuration for this failure
is a sequence $|T_n|\to\infty$, $d_n\downarrow0$, with at least
$k_n$ centered zeros, counted with multiplicity, in
\[
 \left\{z:\left|z-(a-d_n+iT_n)\right|\le d_n/2\right\},
 \qquad
 \frac{(k_n-p)\log(1/d_n)}{\log\nu(T_n)}\longrightarrow\infty.
 \tag{MA-MO.1203}
\]
All zeros in each displayed cluster are strictly inside the strip.
Here and in the proof discard a finite prefix so that
$d_n<\min\{\delta_0,4a/3\}$; this makes both edge inequalities
strict and permits the choice $\delta=d_n$.
In particular, a fixed integer $k>p$ and
$\log(1/d_n)/\log\nu(T_n)\to\infty$ suffice.
Unbounded $k_n$ with suitably small polynomial $d_n$ also suffice.
The configuration may consist of distinct simple zeros.  Neither
these possibilities nor the full-retention premise are asserted
to occur arithmetically.

\begin{proof}
\emph{1. Fixed compact supports and uniform noncancellation.}
Let $r=\sup\{x:\phi(x)\ne0\}$.  Choose a compact interval
$J$ in $(r-\log2,r)$ on which $\phi$ has a constant nonzero
sign and is bounded away from zero.  Such an interval exists
by the definition of $r$ and continuity.  Choose nonnegative,
nonzero $g_0\in C_c^\infty((1,2))$ and
$v_0\in C_c^\infty(\mathbb R)$, with sufficiently small
supports, so that
\[
 k_0=C_{g_0^\vee}v_0\ne0,\qquad
             k_0\ge0,\qquad \operatorname{supp}k_0\subset J.
 \tag{MA-MO.1204}
\]
For example, place the support of $g_0$ near $t_*$ in $(1,2)$
and that of $v_0$ near $x_*+t_*$, with $x_*$ in the interior
of $J$.  Their support difference then lies in $J$.
Fix also $0\ne h_0\ge0$ compact and smooth, with
$\operatorname{supp}h_0\subset[\alpha,\beta]$.
For real $T$ set
\[
 g_T(t)=e^{-iTt}g_0(t),\qquad
 v_T(x)=e^{iTx}v_0(x),\qquad
 h_T(x)=e^{-iTx}h_0(x).
 \tag{MA-MO.1205}
\]
The bilinear transpose in (MA-MO.557) gives exactly
$C_{g_T^\vee}v_T=e^{iTx}k_0(x)$.  The norm of the output
functional $\ell_{v_T}(y)=\int yv_T$ on $H$ is independent
of $T$.  Also
\[
 \|g_T\|_{C^j}\le C_j\nu(T)^j
                   \qquad(j\ge0).
 \tag{MA-MO.1206}
\]

Use the literal $F,W,H_h,V_g,B_v,\mathcal R_{h,v}$ of
(MA-MO.558)--(MA-MO.567), with these modulated tests.
Every nonunit ideal has norm at least two.  Since
$\operatorname{supp}k_0\subset(r-\log2,r)$, the only term
of $W$ meeting $\operatorname{supp}\phi$ is the unit ideal.
At $w=a+\delta+iT$ this proves
\[
 \begin{aligned}
 H_{h_T}(w)&=\int h_0(x)e^{(a+\delta)x}\,dx,\\
 B_{v_T}(w)&=\int_J\phi(x)k_0(x)e^{-(a+\delta)x}\,dx,\\
 V_{g_T}(-w)&=\int_Jk_0(x)e^{-(a+\delta)x}\,dx.
 \end{aligned}
 \tag{MA-MO.1207}
\]
Thus, for $0<\delta\le\delta_0$, both of the first two
absolute values are bounded below by fixed positive constants,
and all three are bounded above.  This is exact unit-ideal
noncancellation after modulation, including for complex $g_T$;
no estimate of a global inverse has been used.

Choose $t_0$ as in (MA-MO.561), using these fixed supports,
and enlarge it so that $\beta+t_0\ge1$.
The dual source term on $U_th_T$ is identically zero for
$t\ge t_0$, before applying the actual source Fourier map.
The direct convolution pairing also vanishes there.
The compact region (MA-MO.565) is independent of $T$.
On that region both ideal sums are finite, with uniformly
bounded absolute values: modulation only changes their terms
by phases.  The exponential factor in the remainder has
bounded modulus for $a\le\Re w\le a+\delta_0$.
Consequently
\[
       |\mathcal R_{h_T,v_T}(a+\delta+iT)|\le C
                    \quad(T\in\mathbb R,\ 0<\delta\le\delta_0).
 \tag{MA-MO.1208}
\]
This bound uses the compact double integral, not a bound on
individual partial-fraction coefficients.

\emph{2. The orbit identity on its actual source domain.}
Put $w=a+\delta+iT$ and, for $L>t_0$, define
\[
 \begin{aligned}
 u_{\delta,T,L}
      &=\int_{t_0}^{L}e^{-wt}U_th_T\,dt,\\
 u_{\delta,T}(x)
      &=e^{-wx}\int_{-\infty}^{x-t_0}
                         e^{(a+\delta)y}h_0(y)\,dy .
 \end{aligned}
 \tag{MA-MO.1209}
\]
The first vector is compact and smooth, supported in
$[\alpha+t_0,\beta+L]$.  It belongs to the original smooth
common sector, and (MA-MO.552) acts on it using the literal
wide inverse and actual local Fourier pair.
Changing variables gives its more explicit expression
\[
 u_{\delta,T,L}(x)=e^{-wx}
       \int_{x-L}^{x-t_0}e^{(a+\delta)y}h_0(y)\,dy.
 \tag{MA-MO.1210}
\]
Positivity of $h_0$ shows that its absolute value is bounded
by that of $u_{\delta,T}$.  The latter vanishes left of
$\alpha+t_0$ and has the exact tail
$H_{h_T}(w)e^{-wx}$ right of $\beta+t_0$.
The compact middle interval is uniformly bounded.  Retaining
both exponentials in the sum norm, exactly as in (MA-MO.614),
gives
\[
 \|u_{\delta,T,L}\|_H
       \le\|u_{\delta,T}\|_H\le C\delta^{-p}.
 \tag{MA-MO.1211}
\]
All constants here are independent of $T,\delta,L$.
The finite vectors converge to $u_{\delta,T}$ in $H$ by
dominated convergence.

There is also an independent justification for the scalar
output limit which does not assume the desired extension on
$H$.  For fixed $\delta>0$, take $b=a+\delta/2$.
The closed-strip hypothesis gives the strict right-half-plane
margin required by MA-MO-022 at width $b$, for both sectors.
That theorem supplies the actual source extension on
$H_{\Sigma,b,s}$, with constants allowed to depend on $b$.
The vectors in (MA-MO.1209) converge in this wider norm too,
since $\Re w>b$.  Equivalently, the actual paired function
$f(t)=\ell_{v_T}(D_{g_T}^0U_th_T)$ has an exponential bound
with every fixed exponent strictly above $a$.  Hence its
Laplace integral converges whenever $\Re w>a$.

The wide identity (MA-MO.567), initially valid for
$\Re w>1/2$, therefore continues holomorphically to
$\Re w>a$; $A_\eta$ has no zero there.  At the specified
$w$ it yields
\[
 \begin{aligned}
 \lim_{L\to\infty}\ell_{v_T}(D_{g_T}^0u_{\delta,T,L})
   &=\mathfrak F_T(w),\\
 \mathfrak F_T(w)
   &=H_{h_T}(w)V_{g_T}(-w)
       -\frac{H_{h_T}(w)B_{v_T}(w)}{A_\eta(w)}
       -\mathcal R_{h_T,v_T}(w),\\
 |\mathfrak F_T(w)|&\ge c|A_\eta(w)|^{-1}-C.
 \end{aligned}
 \tag{MA-MO.1212}
\]
Finite-interval integration is on the actual compact source.
The wider-space result is used only to justify this scalar
limit for each fixed $\delta$; no uniform $b\downarrow a$
operator estimate is imported.

\emph{3. Baire gives the required uniform test estimate.}
Suppose first that every complex test in $\mathcal D$ has
a bounded extension on $H$.  Uniqueness follows from density
of compact smooth inputs.  For each fixed compact input $u$,
the map $g\mapsto D_g^0u$ is continuous into $H$:
$D_g^0u=C_gD^0u$, with $D^0u\in\mathcal E$, and the
weighted translation bound controls convolution on a fixed
test-support interval.  Therefore the sets
\[
 \mathcal F_m=\{g\in\mathcal D:
      \|D_g^0u\|_H\le m\|u\|_H
                 \text{ for every }u\in C_c^\infty\}
 \tag{MA-MO.1213}
\]
are closed.  They cover $\mathcal D$ under the hypothesis.
Baire's theorem, subtraction of a point in an interior ball,
linearity, and scaling give some finite $N,C$ such that
\[
                \|D_g\|_{H\to H}\le C\|g\|_{C^N}
                              \qquad(g\in\mathcal D).
 \tag{MA-MO.1214}
\]
The identical argument in the real Fr\'echet space gives this
bound first for real tests; writing a complex test as its real
and imaginary parts extends it with at most a fixed factor.
Thus the real-only hypothesis suffices as well.

Apply the bound to (MA-MO.1212), then use (MA-MO.1211), the
constant norm of $\ell_{v_T}$, and (MA-MO.1206).  It follows that
\[
 c|A_\eta(a+\delta+iT)|^{-1}-C
      \le C'\nu(T)^N\delta^{-p}.
 \tag{MA-MO.1215}
\]
Multiplication by $\delta^p$, with $\delta\le\delta_0$,
proves (MA-MO.1201).

\emph{4. A fixed genuine-core sequence and residual bad tests.}
Assume that (MA-MO.1201) fails.  Choose pairs $(T_n,\delta_n)$
so that $\delta_n^p|A_\eta(a+\delta_n+iT_n)|^{-1}$ exceeds
any prescribed constant times $n\nu(T_n)^n$.
By (MA-MO.1212), choose a finite $L_n>t_0$ such that its
actual scalar output differs from $\mathfrak F_{T_n}$ by
at most one, and put
\[
 e_n=\frac{u_{\delta_n,T_n,L_n}}
              {\|u_{\delta_n,T_n,L_n}\|_H}.
 \tag{MA-MO.1216}
\]
The denominator is nonzero by (MA-MO.1210).  Increasing the
chosen reciprocal lower bounds if necessary, equations
(MA-MO.1211)--(MA-MO.1212) and the uniform functional norm give
\[
                  \|D_{g_{T_n}}^0e_n\|_H
                                    \ge n\nu(T_n)^n.
 \tag{MA-MO.1217}
\]
This construction uses only finite original-source inputs.
It may choose $L_n$ adaptively; no uniform truncation rate is claimed.

For these now fixed inputs define the closed sets
\[
 \mathcal B_m=\{g\in\mathcal D:
                    \sup_n\|D_g^0e_n\|_H\le m\}.
 \tag{MA-MO.1218}
\]
If one had interior, the same subtraction and scaling argument
would give $\sup_n\|D_g^0e_n\|_H\le C\|g\|_{C^N}$ for all
tests.  Applied to $g_{T_n}$, this contradicts
(MA-MO.1206) and (MA-MO.1217) for $n>N$.
Every $\mathcal B_m$ is therefore nowhere dense.
Their union is exactly the set of tests with bounded outputs
on the fixed sequence.  The complement is a dense $G_\delta$.
For real tests the same interior argument complexifies before
using $g_{T_n}$, so the real statement follows too.

If full individual retention is additionally assumed, reciprocal
growth failure cannot be confined to bounded heights: near a
boundary zero of order $m<p$, $\delta^p|A(a+\delta+iT)|^{-1}$
is bounded locally, and away from finitely many such zeros
compactness applies.  Thus in that setting the pairs may be
chosen with $|T_n|\to\infty$.  This observation is not needed
for the genuine-core construction itself.

\emph{5. A conditional cluster forces failure.}
We verify the finite-disk estimate needed here directly.
At large $|T|$, center the uncentered finite-$L$ variable at
$2+iT$ and choose a zero-free boundary radius between $12$
and $13$.  Divide $L_K$ by the finite disk Blaschke product
of its zeros, with multiplicity.  The resulting holomorphic
factor $H_T$ has polynomial upper bound $C\nu(T)^D$ on
the disk, by the finite-$L$ polynomial strip bound and the
maximum principle; each Blaschke factor has modulus one
on its boundary.  On the inner half disk each normalized
factor belonging to a zero $\rho$ satisfies
\[
      |B_{\rho,T}(w)|\le C|w-\rho|,
                       \qquad |B_{\rho,T}(w)|\le1.
 \tag{MA-MO.1219}
\]
The first estimate follows by bounding the denominator in
the explicit disk factor (MA-MO.906) below by a fixed positive
constant.  Its constant is independent of the zero and height.

For a cluster in (MA-MO.1203), evaluate at
$z_n=a+d_n+iT_n$, or $w_n=1/2+z_n$ in the uncentered
finite-$L$ disk.  Its distance to every counted centered zero is
at most $5d_n/2$.  Use the first bound for those factors and
the second for the rest.  The deleted finite Euler factor
is uniformly bounded on this positive central strip.  Hence
\[
 |A_\eta(z_n)|\le C\nu(T_n)^D(C_1d_n)^{k_n}.
 \tag{MA-MO.1220}
\]
The usual fixed-window zero bound is $k_n=O(\log\nu(T_n))$,
as in MT-M18 and MA-MO-030.  Thus $C_1^{k_n}$ can be absorbed
in a fixed power of $\nu(T_n)$, and
\[
 d_n^p|A_\eta(a+d_n+iT_n)|^{-1}
           \ge C^{-1}\nu(T_n)^{-D'}d_n^{\,p-k_n}.
 \tag{MA-MO.1221}
\]
Condition (MA-MO.1203) makes the right side exceed every
polynomial in $\nu(T_n)$.  It therefore contradicts
(MA-MO.1201), and the already constructed residual obstruction
applies.  No existence statement about such zeros was used.

Finally the exact identity
$R D_{\eta,g}^0R=D_{\bar\eta,g^\vee}^0$ and the isometry of
$R$ in the sum norm give the reflected version.  The original
finite Fourier phases and primitive Euler conventions remain
those in (MA-MO.550)--(MA-MO.576).
\end{proof}

\paragraph{What this adds, and what it does not decide.}
The new necessity records quantitative information beyond the
individual local lengths in MA-MO.618.  It is a uniform estimate forced by a
whole fixed-support test family on the actual full source.
The cluster implication explains a possible additional analytic
failure mechanism even with only strictly interior, hence fully
retained, local factors.  It does not refute full-retention
sufficiency for actual Hecke arithmetic without an arithmetic
configuration satisfying (MA-MO.1203).  Nor is the reciprocal
estimate asserted sufficient for a trace-class matrix; an
appropriate boundary-distribution estimate is still needed.
MA-MO-030 provides one sufficient class and remains unchanged.

The stipulated input A gives the positive fixed-gap situation for
the Dirichlet family, and M gives it for the stated Hecke family.
Then MA-MO-022 already supplies the full source operators and
(MA-MO.1201) follows.  Invoking that premise does not provide an
independent estimate without zero-location input.  The completed
ideal result MA-MO-028 and the whole-quotient return criteria
MA-PR-009--011 have different domains and are not contradicted.

\paragraph{Source and finite-length placement.}
The original inverse, source membership, dual-term support,
bilinear transpose and compact remainder are those completely
proved in MA-MO-024/025.  Their primary interfaces are Meyer
(2005), supplied \texttt{main.tex} lines 2630--2675,
3592--3707 and 3808--4043, as recorded there.
The fixed-disk analytic ingredients are the finite-$L$ polynomial
strip bound and local zero count used in MT-M18.7--9 and
MA-MO-030; the upper factor estimate (MA-MO.1219)--(MA-MO.1221)
has been derived above.  This module does not claim a new full
reading of those primary proofs.

In the inherited cote~1 trace framework, BR-TR-01/02,
scans 079--080 and 111--112, equality of retained local lengths
is a statement in the category of finite-length
$\mathbb C[z]$-modules.  The height-uniform original-source
bound proved necessary here is additional analytic structure.
Neither that $K_0$ equality, the exterior-power determinants
of the finite underlying vector spaces, nor a motivic weight
supplies it.  No arithmetic progress of FT002 is claimed.

\section{The rational polar completion and the boundary-uniform reciprocal necessity}
\label{sec:polar-cluster-obstruction}

\paragraph{Status.}
The complete author and independent root whole-proof readings
have passed; their exact check scopes are recorded separately.
The two actual polar source coordinates are retained throughout;
the statement does not replace the common source by a bare
diagonal component.

Keep $S=\{\infty\}$, the real smooth normalized step $\phi$,
and the actual Gaussian sections of MA-MO-016.  In its coordinates
write
\[
\begin{gathered}
 H=H_{\Sigma,a,s},\qquad
 \mathcal W=H\oplus\mathbb C^2,\qquad 0<a<1/2,\ s\ge0,\\
 \mathfrak s(u,c)=i_-u+c_0s_0+c_1s_1,\qquad
 s_j=i_+f_j,\qquad i_-u=(u,0)\text{ in these coordinates},\\
 D=P_+-P_-,\qquad
 T_g^0=\mathscr D_g^{\rm raw}i_-=C_gD i_-:
                                  \mathcal E\longrightarrow\mathcal W,\\
 A(w)=\zeta(1/2+w),\qquad p=(s+1)/2,\qquad
 0<\delta_0<\min\{1,(1/2-a)/2\},\qquad
 \mathcal D=C_c^\infty([1,2]).
\end{gathered}
\tag{MA-MO.1300}
\]
The norm of $H$ is exactly the inherited sum norm with weight
$(e^{2ax}+e^{-2ax})(1+x^2)^{s/2}$, and $i_-$ is isometric.
The smooth map $T_g^0$ is the original negative source column,
before any bounded extension is assumed.

Assume for the conditional comparison that all nontrivial zeta
zeros satisfy $|\Re\rho-1/2|\le a$.  Define
\[
 \begin{gathered}
 (\mathrm{BR}_{a,s,\zeta}):\quad
 \exists C,N<\infty\quad
 \delta^p|\zeta(1/2+a+\delta+iT)|^{-1}
                    \le C(2+|T|)^N\\
 \hfill(T\in\mathbb R,\quad0<\delta\le\delta_0).
 \end{gathered}
 \tag{MA-MO.1301}
\]

\paragraph{MA-MO-034 (the actual polar negative column requires reciprocal control).}
If every complex $g\in\mathcal D$ gives a bounded extension
$T_g:H\to\mathcal W$, then (MA-MO.1301) holds.  The same is
true if boundedness is assumed only for every real test in
$\mathcal D$.  In particular the conclusion is necessary for
boundedness of the full original matrix on $\mathcal W$ for
that test family.

If (MA-MO.1301) fails, there is a fixed sequence of compact smooth
inputs $e_n$ of $H$-norm one such that
\[
 \left\{g\in\mathcal D:
       \sup_n\|\mathscr D_g^{\rm raw}i_-e_n\|_{\mathcal W}=\infty\right\}
       \quad\text{is a dense }G_\delta.
 \tag{MA-MO.1302}
\]
This holds separately for the real and complex test spaces.
Each input $i_-e_n$ and every displayed output lies in the
original smooth common source.  Both polar output coordinates
are retained in its norm.

The explicitly hypothetical cluster condition
\[
 \begin{gathered}
 |T_n|\to\infty,\quad d_n\downarrow0,\quad
 0<d_n<\min\{\delta_0,4a/3\},\\
 \#\left\{\rho:\left|\rho-1/2-(a-d_n+iT_n)\right|
                           \le d_n/2\right\}\ge k_n,\\
 \frac{(k_n-p)\log(1/d_n)}{\log(2+|T_n|)}\longrightarrow\infty
 \end{gathered}
 \tag{MA-MO.1303}
\]
implies failure of (MA-MO.1301); the count includes multiplicities.
Every clustered zero is strictly interior.  No such zero
configuration, or independent arithmetic failure of the bound,
is asserted.  The zeta functional equation transfers a cluster
near the left edge to a right-edge cluster of the same form.

\begin{proof}
\emph{1. The actual completion, source and output functional.}
For a compact smooth $v$ use the bounded functional
\[
 \Lambda_v(u,c)=\int v(x)
       \{u(x)+c_0\Theta f_0(x)+c_1\Theta f_1(x)\}\,dx.
 \tag{MA-MO.1304}
\]
It integrates the first actual pair component of $\mathfrak s$.
Weighted Cauchy--Schwarz handles $u$; the fixed section functions
are smooth on the compact support of $v$.
In particular, if $v_T=e^{iTx}v_0$ for fixed compact $v_0$, then
\[
 \sup_T\|\Lambda_{v_T}\|_{\mathcal W^*}<\infty.
 \tag{MA-MO.1305}
\]
For the diagonal part the norm is independent of $T$; for each
polar part the absolute value of its compact integral is bounded
by $\int|v_0\Theta f_j|$.  No polar exponential is put in $H$.

For a compact smooth $y$, the original inverse and its actual
additive source are
\[
 \begin{aligned}
 F_y(x)&=\sum_{n\ge1}\mu(n)n^{-1/2}y(x+\log n),\\
 u_y(t)&=\frac{q(\log|t|)F_y(\log|t|)}{2\sqrt{|t|}},
          \qquad q=1-\phi,\\
 h_y&=u_y+\mathfrak F u_{Ry}.
 \end{aligned}
 \tag{MA-MO.1306}
\]
The inverse sum is locally finite and has a finite right support
endpoint.  The cutoff kills its left end, so $u_y$ is an actual
compact annular even Schwartz function.  Thus $i_+h_y-i_-y$
belongs to the original common-sum space.  Its coordinates in
$\mathcal W$ are the unique actual source-section coordinates
of MA-MO-016, not a discarded-pole representative.
The exact radial normalization is
\[
 \Theta u_y(x)=\sum_{n\ge1}n^{-1/2}
                     q(x+\log n)F_y(x+\log n).
 \tag{MA-MO.1307}
\]
For $y\in\mathcal E$ the same construction is valid using
the rapid right tail, as established in MA-MO-026.

Fix nonzero nonnegative compact smooth functions $h_0,g_0,v_0$ as in
(MA-MO.1204): $g_0$ is supported in $(1,2)$, and
$k_0=C_{g_0^\vee}v_0$ is nonzero, nonnegative, and supported
in an interval $J\subset(r-\log2,r)$ on which $\phi$ has
one nonzero sign, where $r=\sup\operatorname{supp}\phi$.
Set $h_T=e^{-iTx}h_0$, $g_T=e^{-iTx}g_0$, $v_T=e^{iTx}v_0$.
The source dictionary is then
\[
 \begin{gathered}
 v_{g_T}=C_{g_T^\vee}v_T=e^{iTx}k_0(x),\\
 W_T(y)=\sum_{n\ge1}n^{-1/2}v_{g_T}(y-\log n),\qquad
 H_T(w)=\int e^{wx}h_T(x)\,dx,\\
 V_T(w)=\int e^{wx}v_{g_T}(x)\,dx,
 \qquad B_T(w)=\int\phi(y)W_T(y)e^{-wy}\,dy.
 \end{gathered}
 \tag{MA-MO.1308}
\]
The transpose is bilinear and introduces no conjugate of $g_T$.
Every nonunit integer shifts $J$ past $r$, so only $n=1$
contributes to $B_T$.  At $w=a+\delta+iT$,
\[
 \begin{aligned}
 H_T(w)&=\int h_0(x)e^{(a+\delta)x}\,dx,\\
 B_T(w)&=\int_J\phi(x)k_0(x)e^{-(a+\delta)x}\,dx,\\
 V_T(-w)&=\int_Jk_0(x)e^{-(a+\delta)x}\,dx.
 \end{aligned}
 \tag{MA-MO.1309}
\]
The first two have absolute values bounded below uniformly
for $0<\delta\le\delta_0$, and all three are bounded above.

Choose $t_0>0$ by the fixed support conditions (MA-MO.561).
Then $qF_{RU_th_T}=0$ for $t\ge t_0$, before applying
Fourier transformation.  Also the subtracted input has zero
pairing with $v_{g_T}$.  Equations (MA-MO.1306)--(MA-MO.1308)
therefore give, on the actual source,
\[
 \begin{aligned}
 f_T(t)&=\Lambda_{v_T}(T_{g_T}^0U_th_T)\\
       &=\int q(y)W_T(y)F_{h_T}(y-t)\,dy
                                      \quad(t\ge t_0).
 \end{aligned}
 \tag{MA-MO.1310}
\]
This evaluates the represented first component and retains its
two source polar coordinates through $\Lambda_{v_T}$.

\emph{2. The actual orbit limit and the pole-cleared identity.}
The literal wide Euler calculation of (MA-MO.658) yields
\[
 \mathfrak F_T(w)=H_T(w)V_T(-w)
                  -\frac{H_T(w)B_T(w)}{A(w)}-\mathcal R_T(w),
                    \qquad \Re w>1/2,
 \tag{MA-MO.1311}
\]
where $\mathfrak F_T$ is the Laplace integral of (MA-MO.1310)
from $t_0$ to infinity.  The remainder is the compact double
integral (MA-MO.564)--(MA-MO.565) with positive integers.
Its domain is independent of the modulation, and its finite sums
have uniformly bounded moduli.  Therefore
\[
       |\mathcal R_T(a+\delta+iT)|\le C
                         \qquad(T\in\mathbb R,\ 0<\delta\le\delta_0).
 \tag{MA-MO.1312}
\]
It is entire in $w$.  Clearing the zeta pole with
$\widetilde A(w)=(w-1/2)A(w)$ gives the pole-cleared identity
on the wide half-plane
\[
 \widetilde A(w)[H_T(w)V_T(-w)-\mathcal R_T(w)-\mathfrak F_T(w)]
                       =(w-1/2)H_T(w)B_T(w).
 \tag{MA-MO.1313}
\]
The pole of $A$ supplies a removable zero of its reciprocal.

For $w=a+\delta+iT$ define the finite original negative inputs
\[
 \begin{aligned}
 u_{\delta,T,L}&=\int_{t_0}^L e^{-wt}U_th_T\,dt,\\
 u_{\delta,T,L}(x)
      &=e^{-wx}\int_{x-L}^{x-t_0}e^{(a+\delta)y}h_0(y)\,dy.
 \end{aligned}
 \tag{MA-MO.1314}
\]
They are compact and smooth.  The second formula shows that
their moduli are bounded by the infinite-orbit modulus; that
vector has a uniformly bounded compact middle part and exact
right tail $H_T(w)e^{-wx}$.  After increasing $t_0$ so that
the tail starts beyond one, the actual sum weight gives
\[
 \|i_-u_{\delta,T,L}\|_{\mathcal W}
       =\|u_{\delta,T,L}\|_H\le C\delta^{-p}.
 \tag{MA-MO.1315}
\]
Indeed its tail square integral is bounded by a constant times
$\int_1^\infty e^{-2\delta x}(1+e^{-4ax})(1+x^2)^{s/2}\,dx$.
The bound is independent of $T,L$, and uses the explicit tail,
not the weaker triangle estimate on the orbit integral.

For each fixed $\delta$, put $b=a+\delta/2<1/2$.
The closed-strip premise implies the strict margin required by
MA-MO-018 on $\mathcal W_{b,s}$, with $a_*=a$.
The finite orbit inputs converge there, and the actual source
matrix has its proved bounded extension there.  The functionals
$\Lambda_{v_T}$ represent the same compact first-component
integral in every such completion.  Consequently their actual
scalar outputs converge.  The same fixed-width bounds give
holomorphy of the paired Laplace integral on $\Re w>a$.
The identity (MA-MO.1313) continues there from the wide half-plane;
equivalently (MA-MO.1311) continues with the reciprocal's
removable value at $w=1/2$.
No uniform estimate as $b\downarrow a$ is being assumed.
At the displayed $w$, it follows that
\[
 \begin{gathered}
 \Lambda_{v_T}(T_{g_T}^0u_{\delta,T,L})\longrightarrow
                  \mathfrak F_T(a+\delta+iT),\\
 |\mathfrak F_T(a+\delta+iT)|
                  \ge c|A(a+\delta+iT)|^{-1}-C.
 \end{gathered}
 \tag{MA-MO.1316}
\]
Every finite input/output in this limit is a genuine smooth
common-source vector.  Its nonpolar and polar coordinates are
not estimated separately or silently canceled from the norm.

\emph{3. The fixed-support family and Baire's theorem.}
For a fixed compact smooth $u$, the dependence
$g\mapsto T_g^0u$ is continuous into $\mathcal W$ without
assuming a bounded matrix on arbitrary inputs.
Indeed $D i_-u$ is one fixed actual smooth common-source vector,
and $C_g$ is its integrated representation.
The exact representation (MA-MO.259)--(MA-MO.261) is strongly
continuous and locally bounded on $\mathcal W$, including
its theta cocycle and both polar characters.  Thus on a fixed
compact support its integral has norm bounded by a constant
times $\|g\|_{L^1}$ times that vector's norm.

If every complex test has a bounded column, the closed sets
\[
 \mathcal F_m=\{g\in\mathcal D:
       \|T_g^0u\|_{\mathcal W}\le m\|u\|_H
                            \text{ for every }u\in C_c^\infty\}
 \tag{MA-MO.1317}
\]
cover $\mathcal D$.  Baire's theorem, subtraction and scaling
give finite $N,C$ with
$\|T_g\|_{H\to\mathcal W}\le C\|g\|_{C^N}$.
If the hypothesis concerns only real tests, first apply Baire
over the real Fr\'echet space and then complexify this bound.
Since $\|g_T\|_{C^N}\le C_N(2+|T|)^N$, equations
(MA-MO.1305), (MA-MO.1315) and (MA-MO.1316) give
\[
 c|A(a+\delta+iT)|^{-1}-C
          \le C'(2+|T|)^N\delta^{-p}.
 \tag{MA-MO.1318}
\]
This proves the required estimate (MA-MO.1301).

\emph{4. The fixed genuine-input sequence under failure.}
If (MA-MO.1301) fails, choose pairs $(T_n,\delta_n)$ for which
$\delta_n^p|A(a+\delta_n+iT_n)|^{-1}$ exceeds an arbitrarily
large fixed multiple of $n(2+|T_n|)^n$.
For each pair choose a finite $L_n$ so that the scalar error
in (MA-MO.1316) is at most one.  Normalize
$e_n=u_{\delta_n,T_n,L_n}/\|u_{\delta_n,T_n,L_n}\|_H$.
The denominator is nonzero, and the preceding bounds imply
\[
           \|T_{g_{T_n}}^0e_n\|_{\mathcal W}
                               \ge n(2+|T_n|)^n.
 \tag{MA-MO.1319}
\]
The input has zero source polar coordinates; its output retains
both coordinates.  The fixed sequence is chosen through finite
actual orbit integrals, not through a proposed completed inverse.

For these fixed $e_n$, each set
$\{g:\sup_n\|T_g^0e_n\|_{\mathcal W}\le m\}$ is closed.
If it had interior, the same subtraction/scaling argument would
give a finite $C^N$ seminorm bound, contradicting (MA-MO.1319).
For the real space the bound complexifies before evaluating the
complex modulations.  These sets are nowhere dense, so their
complement is the asserted dense $G_\delta$ in either space.
Unboundedness of this column excludes a bounded full matrix.

\emph{5. The conditional zeta cluster.}
At large height the finite disk centered at $2+iT_n$ with radius
between $12$ and $13$ contains no zeta pole.  Divide zeta by its
finite disk Blaschke product.  The polynomial strip upper bound
and maximum principle bound the remaining holomorphic factor
by $C(2+|T_n|)^D$.  At the raw evaluation point
$1/2+a+d_n+iT_n$, every clustered factor has modulus at most
$C_1d_n$, and every other factor has modulus at most one.
Thus, exactly as in (MA-MO.1219)--(MA-MO.1221),
\[
 d_n^p|\zeta(1/2+a+d_n+iT_n)|^{-1}
       \ge C^{-1}(2+|T_n|)^{-D'}d_n^{p-k_n}.
 \tag{MA-MO.1320}
\]
The ordinary fixed-window count $k_n=O(\log(2+|T_n|))$
absorbs the constant $C_1^{k_n}$ in a power of the height.
Condition (MA-MO.1303) now forces failure of (MA-MO.1301).
All source conclusions were proved before making this conditional
specialization.  No existence of such a cluster is inferred.
\end{proof}

\paragraph{Exact relation to the arithmetic and operator gates.}
The theorem concerns the actual negative column into the full
polar completion, and so supplies a necessary condition for the
original full matrix.  It is not a statement about the bounded
positive source return, the quotient-return norms, or an
orthogonal-PW identification.  The pole at one is retained through
the actual source sections and the removable inverse value; it
does not exempt the full matrix from the reciprocal estimate.
Neither this estimate nor full individual jet retention is
asserted sufficient at arbitrary widths.

The user-stipulated A supplies zeta RH.  Then MA-MO-018 gives
the positive fixed-gap branch for every $0<a<1/2$, $s\ge0$
and every compact smooth test, and (MA-MO.1301) follows.
This explicitly named zero-location input is not replaced by
an independent arithmetic reciprocal estimate in this theorem.
No new arithmetic movement of FT002 is claimed.

\paragraph{Source and finite-length placement.}
The actual polar-section completion, local representation and
negative embedding are (MA-MO.250)--(MA-MO.261).  The source
dictionary, bounded compact output functional and pole-cleared
Laplace identity are (MA-MO.653)--(MA-MO.659), proved from
Meyer's original common-space maps, supplied \texttt{main.tex}
lines 3808--4043.  The wider-width source agreement is MA-MO-018,
and the exact norm and cluster analysis have been written above.
No new full primary-source proof reading is claimed here.
In the inherited cote~1 framework, BR-TR-01/02, scans 079--080
and 111--112, the two actual polar sections and finite-length
divisor modules remain separate data.  Their $K_0$ classes and
finite-dimensional exterior-power determinants do not supply the
additional uniform source norm required for a nuclear trace.

\section{A depth budget for the actual full source matrix}
\label{sec:weighted-cluster-depth-sufficiency}

\paragraph{Status.}
The complete author and two independent whole-proof readings
have passed; their exact check scopes are recorded separately.
No arithmetic validity of the stated cluster geometry is asserted.

The whole-product method of MA-MO-030 permits a more precise
estimate than a fixed bound on the total local multiplicity.
The polynomial weight pays for part of the pole order; positive
inward depth controls the rest.  Actual boundary factors require
a separate integrability condition, retained explicitly below.

Use the two source settings of (MA-MO.700): either a nontrivial
finite-order Hecke character with its dual and structural $S$,
or the rational principal sector with its actual polar completion.
Keep the original source Fourier maps, real smooth step $\phi$,
primitive finite functions $A_\omega$, and inherited norms.
Fix $0<a<1/2$, $s\ge0$, and put
\[
 r=s/2,\qquad p=r+1/2=(s+1)/2,\qquad
 0<\delta<\tfrac18\min\{a,1/2-a\},\qquad0<\ell\le1.
 \tag{MA-MO.1400}
\]
Fix also $0<\epsilon_0<\delta/8$ for the right-boundary limits.
Assume that every completed centered zero lies in $|\Re z|\le a$,
for both dual sectors when applicable.  At $t_n=n\ell$, let
$Z_{\omega,n}$ be the full multiset in
\[
 a-\delta\le\Re\lambda\le a,
                  \qquad |\Im\lambda-t_n|\le2\ell.
 \tag{MA-MO.1401}
\]
Write $k_{\omega,n}$ for its total multiplicity and
$m_{\omega,n}$ for the total multiplicity of its factors on
$\Re\lambda=a$.  Impose
\[
                         m_{\omega,n}<p
                          \qquad(\omega,n).
 \tag{MA-MO.1402}
\]
This concerns the total boundary multiplicity in the rectangle,
even when the boundary zeros have distinct ordinates.
It is stronger than merely retaining each of those factors
individually.  When positive-depth factors exist, define
\[
 d_{\omega,n}
   =\min_{\substack{\lambda\in Z_{\omega,n}\\\Re\lambda<a}}
                               (a-\Re\lambda)\in(0,\delta].
 \tag{MA-MO.1403}
\]
If no such factor exists, put $d_{\omega,n}=1$.
In that case $k_{\omega,n}=m_{\omega,n}<p$.
Define the explicit cost
\[
 B_{\omega,n}=
 \begin{cases}
 1,&k_{\omega,n}<p,\\
 \bigl(1+\log(1/d_{\omega,n})\bigr)^{1/2},&k_{\omega,n}=p,\\
 d_{\omega,n}^{-(k_{\omega,n}-p)},&k_{\omega,n}>p.
 \end{cases}
 \tag{MA-MO.1404}
\]
In the last two cases (MA-MO.1402) guarantees the presence of
a positive-depth factor.  Suppose that, for fixed $C,D$,
\[
                  B_{\omega,n}\le C(2+|n|)^D
                              \qquad(\omega,n).
 \tag{MA-MO.1405}
\]
All multiplicities in these definitions are the actual ones.

\paragraph{MA-MO-035 (a polynomial depth budget suffices).}
Under (MA-MO.1401)--(MA-MO.1405) and the closed-strip premise,
the original $C_g(P_+-P_-)$ extends uniquely in trace class
on $H_{\Sigma,a,s}$ in the nontrivial sector, and on
$\mathcal W_{a,s}=H_{\Sigma,a,s}\oplus\mathbb C^2$ in the
rational principal sector, for every compact smooth $g$.
The extensions agree with the original wide-inverse source maps
and depend continuously on $g$ in trace norm on each fixed
compact support.

All completed jets are retained.  The actual relative-trace,
Fredholm-determinant and prime formulas of MA-MO-029 apply,
with their original signs, finite polar correction and condition
$g(0)=0$ for the stated prime formulas.  The assertion is a
sufficient condition; neither (MA-MO.1402) nor (MA-MO.1405)
is claimed necessary.

\begin{proof}
\emph{1. A uniform pole-product estimate.}
Consider one rectangle and suppress its indices.  Enumerate
its $k$ factors by
$\lambda_j=a-d_j+i(t_n+\gamma_j)$, with
$0\le d_j\le\delta$ and $|\gamma_j|\le2\ell$.
For $\epsilon>0$ let
\[
 R_\epsilon(v)=\prod_{j=1}^k
                  (\epsilon+d_j+i(v-\gamma_j))^{-1}.
 \tag{MA-MO.1406}
\]
If $k\ge1$, the actual Laplace formula is
\[
 R_\epsilon(v)=\int_0^\infty e^{-ivt}e^{-\epsilon t}h(t)\,dt,
 \qquad
 h=\mathop{*}_{j=1}^k
       \bigl(\mathbf1_{t>0}e^{-d_jt}e^{i\gamma_jt}\bigr).
 \tag{MA-MO.1407}
\]
Consequently
$|h(t)|\le t^{k-1}/(k-1)!$.  All estimates below are
independent of the ordinates and their separations.
With the unitary plus Fourier transform,
\[
 \|R_\epsilon\|_{H^{-r}}^2
       =2\pi\int_0^\infty
           (1+t^2)^{-r}e^{-2\epsilon t}|h(t)|^2\,dt.
 \tag{MA-MO.1408}
\]

If $1\le k<p$, the simplex bound is integrable in (MA-MO.1408),
uniformly even when every depth is zero.  There are only finitely
many such integers $k$, so their bounds are controlled by a
constant depending on $s$ alone.

Suppose first that every depth is positive, with least depth $d$.
Then $|h(t)|\le e^{-dt}t^{k-1}/(k-1)!$.
If $k>p$, split the integral at one and use
$(1+t^2)^{-r}\le t^{-2r}$ for $t\ge1$.  With
$q_k=2(k-p)>0$, the result is bounded by
\[
 C+\frac{2\pi\Gamma(q_k)}{((k-1)!)^2}
                                      (2d)^{-q_k}.
 \tag{MA-MO.1409}
\]
Here the constant from $t\le1$ is uniform in $k$.
The gamma ratio has a uniform estimate, not a separate
constant for each growing multiplicity.  If $q_k\ge2$,
log-convexity of the gamma integral, between $2$ and
$2k-1$, gives
\[
 \Gamma(q_k)\le\Gamma(2k-1)=(2k-2)!,\qquad
 \frac{(2k-2)!}{((k-1)!)^2}\le4^{k-1}.
 \tag{MA-MO.1410}
\]
For $0<q_k<2$ there are only finitely many eligible integers
$k$ for this fixed $s$, and a constant depending on $s$ covers
them.  Since $q_k=2(k-p)$, the power of two in (MA-MO.1409)
cancels $4^{k-1}$ up to the fixed factor $4^{p-1}$.
It follows that $\|R_\epsilon\|_{H^{-r}}\le C_s d^{-(k-p)}$.
If $k=p$, which can occur only for this fixed integer $p$,
split instead at $1$ and $1/d$.  The middle integral is
bounded by $\log(1/d)$, and the tail by
$\int_1^\infty e^{-2u}\,du/u$.  This gives the critical
square-root logarithm in (MA-MO.1404).

Now allow $m\ge1$ boundary factors, where $m<p$, and put
$l=k-m$ for the number of positive-depth factors.
If $l=0$, then $k<p$ and the first case already applies.
Otherwise let $d$ be the least of their positive depths.
Convolving the separate simplex bounds gives both estimates
\[
 |h(t)|\le
 \min\left\{\frac{t^{k-1}}{(k-1)!},
             \frac{t^{m-1}d^{-l}}{(m-1)!}\right\}.
 \tag{MA-MO.1411}
\]
For the second, write the convolution parameter of the
positive-depth factors as $u$, bound $(t-u)^{m-1}$ by
$t^{m-1}$, and integrate
$e^{-du}u^{l-1}/(l-1)!$ over $u>0$ to get $d^{-l}$.
No exponential damping is assigned to the boundary factors.

Use the first bound up to $t=1/d$ and the second thereafter.
For $k>p$ the two large-$t$ integrals are bounded, apart from
their factorials, by
\[
 \int_1^{1/d}t^{2(k-p)-1}\,dt
       \le\frac{d^{-2(k-p)}}{2(k-p)},\qquad
 d^{-2l}\int_{1/d}^\infty t^{2(m-p)-1}\,dt
       =\frac{d^{-2(k-p)}}{2(p-m)}.
 \tag{MA-MO.1412}
\]
The first positive denominator is bounded away from zero
as $k$ varies through integers greater than the fixed $p$.
The second is also bounded away from zero, because $m$ ranges
over the finite set of positive integers less than $p$.
The factorial denominators only improve the bound.
For $k=p$, the first integral is $\log(1/d)$ and the second
is a fixed constant.  The case $k<p$ was already handled.
Thus in every nonempty case,
\[
 \sup_{0<\epsilon\le\epsilon_0}\|R_\epsilon\|_{H^{-r}}
          \le C_s B,\qquad
 R_\epsilon\longrightarrow R_0\quad\hbox{in }H^{-r}
                              \quad(\epsilon\downarrow0).
 \tag{MA-MO.1413}
\]
Convergence follows by dominated convergence in (MA-MO.1408)
for this fixed multiset.  A common modulus across all rectangles
is neither asserted nor needed.
For $k=0$ the unlocalized constant $R_\epsilon=1$ is not
put in $H^{-r}$; its compactly localized smooth term below
is handled separately.

\emph{2. The canceled holomorphic factors remain polynomial.}
Define $P_{\omega,n}$, $b_{\omega,n}=P_{\omega,n}/A_\omega$
and the nested rectangles exactly as in (MA-MO.903).
The ordinary local zero bound gives
\[
                    k_{\omega,n}=O(\log(2+|n|)).
 \tag{MA-MO.1414}
\]
The proof of (MA-MO.904) still gives a polynomial bound
for every fixed derivative of $b_{\omega,n}$, uniformly on
the smaller rectangle.  To check the dependence on $k$, remove
all finite-function zeros in the fixed disk centered at $2+it_n$.
The zero-free factor has polynomial inverse by the Euler lower
bound and Harnack.  For every zero in $Z_{\omega,n}$, the
numerator of its Blaschke factor is canceled by $P_{\omega,n}$;
the remaining factor costs a fixed constant.  All other disk
zeros are a fixed distance from the smaller rectangle, so each
also costs a fixed constant.  There are $O(\log(2+|n|))$
factors altogether.  Thus even a factor $C^{k_{\omega,n}}$
costs only a polynomial.  Fixed-radius Cauchy estimates give
the same conclusion for each fixed derivative.  The finite
Euler factors are invertible on this positive strip, and the
rational pole is outside the rectangle.

Use the fixed smooth partition $\chi_n$ of (MA-MO.910), and set
\[
 d_{\omega,n,\epsilon}(u)
   =\chi_n(u)b_{\omega,n}(a+\epsilon+iu)
                              R_{n,\epsilon}(u-t_n).
 \tag{MA-MO.1415}
\]
For $\epsilon>0$ this is precisely the localized actual
reciprocal.  Sobolev multiplication, (MA-MO.1405) and
(MA-MO.1413) yield, for some fixed $C,D_1$,
\[
 \|d_{\omega,n,\epsilon}\|_{H^{-r}}\le C(2+|n|)^{D_1},\qquad
 d_{\omega,n,\epsilon}\to d_{\omega,n,0}\quad\hbox{in }H^{-r}
                         \quad\text{for each fixed }n.
 \tag{MA-MO.1416}
\]
The estimates are uniform for $0\le\epsilon\le\epsilon_0$.
When $k=0$ the compactly supported smooth factor has this
bound directly.  This is the only use of the empty product.

The locally finite sum $j_{\omega,+}=\sum_n d_{\omega,n,0}$
is a tempered distribution.  Polynomial bounds and the rapid
decrease of localized Schwartz tests give
\[
 \frac1{A_\omega(a+\epsilon+i\,\cdot)}\to j_{\omega,+}
          \quad\hbox{in distributions},\qquad
                 A_\omega(a+i\,\cdot)j_{\omega,+}=1.
 \tag{MA-MO.1417}
\]
The second identity is taken after multiplying at $\epsilon>0$;
the forward multiplier and its derivatives have uniform polynomial
bounds.  For each fixed Schwartz $f$, the products
$j_{\omega,\epsilon}f$ also converge in $H^{-r}$.
These conclusions use dominated summation of (MA-MO.1416),
not a uniform-in-$n$ boundary-convergence modulus.

\emph{3. Actual source identification and trace-norm summation.}
For a two-variable Schwartz kernel $L$, the compact-distribution
Hermite lemma (MA-MO.715), translated to the $n$th window, gives
\[
 \|T_{d_{\omega,n,\epsilon},L}\|_{\mathcal S_1(H^r,H^R)}
     \le C_{L,N}(2+|n|)^{D_1-N}\qquad(R\ge0)
 \tag{MA-MO.1418}
\]
for every $N$, with finitely many Schwartz seminorms for fixed
$r,R,N$.  Choose $N>D_1+2$.  The series is uniformly summable
in trace norm.  Each individual term converges by
(MA-MO.1416), so dominated summation proves trace-norm
convergence of the series.  The same proof works when
$L_\epsilon\to L_0$ in Schwartz space.  This is the required
replacement for the stronger common-modulus estimate in MO030.

For $y\in\mathcal E$, the original wide inverse $F_{\omega,y}$
can first be shifted to every fixed $b>a$, because that line
has a fixed margin from all zeros.  Hence its actual identity is
$V_bF_{\omega,y}=A_\omega(b+i\,\cdot)^{-1}V_by$.
For fixed $y\in\mathcal E$, $V_by\to V_ay$ in Schwartz space
as $b\downarrow a$ on the indicated compact range of real parts.
The just-proved $H^{-r}$ convergence, also after multiplication
by every polynomial in frequency, identifies its boundary value:
\[
 V_aF_{\omega,y}=j_{\omega,+}V_ay,\qquad
 D_x^j(e^{ax}F_{\omega,y})\in
                   L^2((1+x^2)^{-r}\,dx)\quad(j\ge0).
 \tag{MA-MO.1419}
\]
The local $H^1$ estimate therefore gives polynomial left bounds
for this actual function and its derivatives.  Its original
right tail has every stronger exponential weight.  Thus
$qF_{\omega,y}$ is in the original local Schwartz-cup source,
and the boundary prescription agrees with the original maps.
There is no replacement of an original inverse by a chosen
unrelated boundary distribution.

The nonsingular factors in the actual removable commutator and
crossed kernels (MA-MO.718)--(MA-MO.719) converge in Schwartz
space as $b\downarrow a$.  This uses only forward-multiplier
bounds, the compact smooth test and cutoff derivative, and the
separated denominator $\Re(-z-w)=-2b\le-2a$.
Apply (MA-MO.1418) to their input reciprocals.  In the crossed
term the reciprocal is that of the dual sector, unchanged from
the original source formula.  The first row is exactly
\[
 \mathscr R_{\eta,a}^{\partial}(g)y
      =K_{\eta,a}^{\partial}(g)y-N_a(g)y
                         +O_{\eta,a}^{\partial}(g)Ry,
 \qquad N_a(g)=C_gM_\phi.
 \tag{MA-MO.1420}
\]
The half-line map $N_a(g):H_{\Sigma,a,s}\to H_{a,s}$ is
trace class by its physical kernel (MA-MO.332).
The other terms are trace class by (MA-MO.1418) and the
exact Fourier identification $H_{a,s}\simeq H^r$.
On each smooth input, the row is obtained by passing the
original $b>a$ formula to the limit, keeping that input and
its actual smooth theta output fixed.  No bounded embedding
between whole varying-width Hilbert spaces is assumed.
The reflected dual row agrees on the dense smooth common sector.
Its paired map therefore has range in the closed diagonal of
$H_{a,s}\oplus H_{-a,s}$, giving the trace-class endomorphism
in the exact inherited sum norm.

\emph{4. The rational polar completion and evaluated identities.}
The two rational cutoff moments are also controlled by
(MA-MO.1416).  With $p_a(x)=\phi(x)e^{(1/2-a)x}$,
\[
 \lambda_0(y)=\sum_n\left\langle d_{n,0}(u),
                 \widehat p_a(-u)(V_ay)(u)\right\rangle,
                \qquad\lambda_1(y)=\lambda_0(Ry).
 \tag{MA-MO.1421}
\]
The localized multiplier norms of the Schwartz function
$\widehat p_a(-u)$ decay rapidly; the sum of functional norms
is finite.  Passage from the original lines $b>a$, justified
above, identifies these with the two actual physical moments.
The left bound in (MA-MO.1419) and the wide right weights
make the full Mellin integral at $1/2$ converge.  The original
zeta inverse has a zero there, so $\mathcal MF_y(1/2)=0$.
Thus the actual source vector has the same two values
$h_y(0)=-\lambda_1(y)$ and $\int h_y=-\lambda_0(y)$.

Consequently the polar-subtracted diagonal rows and original
source columns assemble exactly as in (MA-MO.921)--(MA-MO.922):
\[
 \begin{aligned}
 \mathscr D_g(y,c)
     &=(T_g^\partial y,\nu_g(y))+\sum_{j=0}^1c_jC_gDs_j,\\
 \nu_g(y)&=(-G(-1/2)\lambda_1(y),-G(1/2)\lambda_0(y)).
 \end{aligned}
 \tag{MA-MO.1422}
\]
The first component includes the unchanged finite corrections
$G(-1/2)\lambda_1(y)e_0+G(1/2)\lambda_0(y)e_1$.
All terms are trace class or finite rank in the specified
source-section completion.  This proves the rational statement.

Every boundary multiplicity is at most the corresponding
$m_{\omega,n}<p$, and reflection supplies the other edge.
Hence $s>2m_\rho-1$ for every boundary zero.  Interior
multiplicities are automatically fully retained.  The actual
full operator, this full retention, and the ideal trace of
MA-MO-028 are exactly the hypotheses used in MA-MO-029.
Its relative-trace and determinant algebra therefore applies
without any finite-divisor assumption.  The source and primitive
Euler factors, both polar signs, and the $g(0)=0$ prime-formula
condition remain unchanged.
Finally (MA-MO.1418) involves only finitely many kernel seminorms,
which are controlled by finitely many derivatives of $g$ on a
fixed compact support.  This proves the asserted trace-norm
continuity.  Density of the original cores gives uniqueness.
\end{proof}

\paragraph{Concrete sufficient cases and exact limitations.}
MA-MO-030 is contained in this theorem: if $k_{\omega,n}\le M<p$,
then (MA-MO.1402) holds and every $B_{\omega,n}=1$, regardless
of inward distances or boundary collisions.  Beyond that class,
suppose there are no boundary zeros, $k_{\omega,n}\le M$,
and $d_{\omega,n}\ge c(2+|n|)^{-D_0}$.  Then (MA-MO.1405)
holds for every fixed $s\ge0$, without requiring $s>2M-1$.
Unbounded $k_{\omega,n}$ are also permitted whenever their
displayed depth budget is polynomial.  In the supercritical
case its logarithmic form is
\[
       (k_{\omega,n}-p)\log(1/d_{\omega,n})=O(\log(2+|n|)).
 \tag{MA-MO.1423}
\]
At the critical integer $k=p$ the separate square-root logarithm
in (MA-MO.1404) is essential.

The estimate uses a nonoscillatory bound for the whole boundary
subproduct.  If its total multiplicity is at least $p$, even at
distinct ordinates with each individual order below $p$, the
argument above does not apply.  This is a limitation of this
sufficient condition, not a proved operator obstruction in that
case.  Likewise failure of the positive-depth budget is not
asserted sufficient for failure of the actual operator: ordinate
separation and cancellations may improve the bound.
MA-MO-033 gives a distinct necessary estimate and a conditional
obstruction for a specified tight cluster configuration.
Neither result proves that any proposed configuration occurs.
The general no-margin full-retention arithmetic question remains
open; A or M supplies the previously separated fixed-gap branch.

\paragraph{Sources and trace-framework placement.}
The fixed-strip polynomial bound, Euler lower bound and ordinary
local zero count are the unconditional inputs in MT-M18.5--9.
The actual inverse, local source, Fourier phases, kernel lemma and
polar sections are the already proved MA-MO-018/022/027/028/030
interfaces.  The new estimates are (MA-MO.1409)--(MA-MO.1413)
and their blockwise dominated trace-norm assembly.  The source
boundary distribution is always tied back to the original wide
inverse before it is used in a completed operator.
In the inherited cote~1 framework, BR-TR-01/02, scans 079--080
and 111--112, finite-length $\mathbb C[z]$-module data give the
local multiplicities and exterior-power determinants.  The new
analytic budget and actual nuclear construction provide additional
structure; no positivity or arithmetic zero bound follows from
the finite-length $K_0$ classes alone.

\section{An exact local boundary-distribution criterion for the full source}
\label{sec:local-boundary-distribution-criterion}

\paragraph{Status.}
The complete author and two independent whole-proof readings
have passed.  Their exact scopes are recorded separately.
No arithmetic validity of the equivalent conditions is asserted.

Use the original source settings of (MA-MO.700): a nontrivial
primitive finite-order Hecke character $\eta$ with its dual,
structural $S$, and actual Fourier phases; or the rational principal
sector with $S=\{\infty\}$ and its two actual polar sections.
Keep the original real smooth step $\phi$, equal to one far left
and zero far right, with compactly supported derivative; put
$q=1-\phi$.
Write
\[
\begin{gathered}
 H=H_{\Sigma,a,s},\qquad
 w(x)=(e^{2ax}+e^{-2ax})(1+x^2)^r,\qquad r=s/2,\\
 0<a<1/2,\qquad s\ge0,\qquad
 0<\epsilon_0<\tfrac12(1/2-a),\\
 A_\omega(z)=E_{\omega,S}(1/2+z)L_K(1/2+z,\omega)
       \quad(\omega=\eta,\bar\eta),\\
 A(z)=\zeta(1/2+z)\quad\hbox{in the rational case},\qquad
 \mathcal W=H\oplus\mathbb C^2 .
\end{gathered}
\tag{MA-MO.1500}
\]
The norm of $H$ is exactly $L^2(w\,dx)$.  The two polar
coordinates of $\mathcal W$ are the actual section coordinates
of MA-MO-016; no bare polar function is put in $H$.

Throughout this theorem assume the closed-strip premise
\[
 |\Re(\rho-1/2)|\le a
 \quad\hbox{for every completed nontrivial zero, in both
 dual sectors when applicable}.
\tag{MA-MO.1501}
\]
No individual boundary-retention hypothesis is added.
The finite Euler factors have no zeros on the positive
central strip used below.  The rational pole at $z=1/2$
is outside $a\le\Re z\le a+\epsilon_0$.

Choose a smooth uniform partition of unity on the frequency
line: for some $\ell>0$, $t_n=n\ell$ and
$\chi_n(u)=\chi(u-t_n)$, where $\chi\in C_c^\infty$,
the supports have uniformly bounded overlap, and
$\sum_{n\in\mathbb Z}\chi_n=1$.  The choice can, for example,
have $\operatorname{supp}\chi\subset[-2\ell,2\ell]$.
Define the following condition, denoted $\mathrm{LB}_{a,s}$:
for every indicated $\omega$ there are
$d_{\omega,n,0}\in H^{-r}(\mathbb R)$ and constants $C,D$
such that
\[
\begin{aligned}
 d_{\omega,n,\epsilon}(u)
   &:=\frac{\chi_n(u)}{A_\omega(a+\epsilon+iu)}
       \longrightarrow d_{\omega,n,0}(u)
                &&\text{strongly in }H^{-r}\text{ for each }n,\\
 \|d_{\omega,n,\epsilon}\|_{H^{-r}}
   &\le C(2+|n|)^D
                &&(n\in\mathbb Z,\ 0<\epsilon\le\epsilon_0).
\end{aligned}
\tag{MA-MO.1502}
\]
The rational condition has the one self-dual multiplier $A$.
The Sobolev norm uses the unitary Fourier transform in $u$.
There is no uniform convergence rate in $n$ in (MA-MO.1502).
Existence for one such partition is equivalent to existence
for every fixed smooth uniform partition, as proved below.

\paragraph{MA-MO-036 (the actual universal bounded family has an exact boundary criterion).}
Under (MA-MO.1501), the following three conditions are equivalent.
In the nontrivial sector they are:
\[
\begin{array}{ll}
{\rm (B)}&
 \mathscr D_{\eta,g}^{\,\mathrm{raw}}
   =C_g(P_+-P_-):\mathcal E\to\mathcal E
 \text{ has a bounded extension on }H\\
&\hfill
 \text{for every }g\in C_c^\infty([-2,-1]\cup[1,2]);\\
{\rm (L)}&\mathrm{LB}_{a,s}\text{ holds for both }\eta,\bar\eta;\\
{\rm (S)}&
 \mathscr D_{\eta,g}^{\,\mathrm{raw}}\text{ has a unique }
 \mathcal S_1(H)\text{ extension for every }g\in C_c^\infty,\\
&\hfill\text{continuously in trace norm on each fixed support}.
\end{array}
\tag{MA-MO.1503}
\]
In the rational principal sector replace (B) by the weaker
negative-column assertion
\[
\begin{gathered}
 \mathscr D_g^{\,\mathrm{raw}}i_-:\mathcal E\to\mathcal W
 \text{ has a bounded extension }H\to\mathcal W\\
       \text{for every }g\in C_c^\infty([1,2]),
\end{gathered}
\tag{MA-MO.1504}
\]
and replace (S) by the full original matrix extension
$\mathscr D_g\in\mathcal S_1(\mathcal W)$, with the same
continuity.  In both settings (B) can be required only for
real tests, without changing the equivalence.

Whenever these conditions hold, every boundary zero of
multiplicity $m$ satisfies
\[
                         s>2m-1 .
\tag{MA-MO.1505}
\]
Thus the actual full-retention, relative-trace and determinant
conclusions of MA-MO-029 apply, with their unchanged primitive
Euler factors, polar correction and $g(0)=0$ restriction for
the stated prime formulas.  This is an analytic equivalence
under (MA-MO.1501), not a proof of either premise for an
arithmetic divisor.

\begin{proof}
\emph{1. Universal boundedness gives a finite test seminorm.}
Use $\mathcal D_+=C_c^\infty([1,2])$ with its closed
Fr\'echet topology.  In the nontrivial case let
$T_g=\mathscr D_{\eta,g}$; in the rational case let
$T_g=\mathscr D_g i_-$, with output space $\mathcal W$.
For a fixed compact smooth input $u$, $g\mapsto T_g^0u$
is continuous in the indicated output norm.  In the first
case this follows from $T_g^0u=C_gD^0u$ and the weighted
translation bound, since $D^0u\in\mathcal E$.  In the
second case it follows from the actual strongly continuous,
locally bounded representation on $\mathcal W$,
(MA-MO.259)--(MA-MO.261), including its cocycle and both
polar characters.  No bounded full matrix is assumed here.

The closed sets
$\{g:\|T_g^0u\|\le j\|u\|_H\text{ for all compact smooth }u\}$
cover $\mathcal D_+$ under (B).  Baire's theorem, subtraction,
linearity and scaling therefore give
\[
              \|T_g\|\le C\|g\|_{C^N}
                         \qquad(g\in\mathcal D_+)
\tag{MA-MO.1506}
\]
for some finite $N$.  The original compact core is dense,
so each extension is unique.  If (B) is assumed only for
real tests, apply Baire in the real Fr\'echet space first
and then write a complex test as its real and imaginary
parts.  This proves (MA-MO.1506) for complex tests too.

\emph{2. The actual scalar functional yields a strong Sobolev boundary value.}
Use the fixed nonnegative compact $g_0,h_0,v_0$ of
(MA-MO.1204)--(MA-MO.1205).  Thus $g_0$ is supported
strictly inside $(1,2)$ and
\[
 k_0=C_{g_0^\vee}v_0\ge0,\qquad k_0\ne0,\qquad
 \operatorname{supp}k_0\subset J\subset(r_\phi-\log2,r_\phi),
\tag{MA-MO.1507}
\]
where $\phi$ has one nonzero sign on $J$ and
$r_\phi=\sup\operatorname{supp}\phi$.
For $T\in\mathbb R$ set
$g_T=e^{-iT\,\cdot}g_0$, $h_T=e^{-iT\,\cdot}h_0$ and
$v_T=e^{iT\,\cdot}v_0$.
Let $\ell_T$ be $\int v_Ty$ in the nontrivial case and
the actual functional $\Lambda_{v_T}$ of (MA-MO.1304)
in the rational case.  In either case
$\sup_T\|\ell_T\|<\infty$.  In the latter case this
includes the two compact polar-section integrals.

The bounded complex-linear functional $\ell_T T_{g_T}$ on
$L^2(w\,dx)$ has a bilinear representation
\[
       \ell_T(T_{g_T}y)=\int y(x)b_T(x)\,dx,\qquad
       b_T\in L^2(w^{-1}dx),\qquad
       \|b_T\|_{L^2(w^{-1})}\le C\|T_{g_T}\|.
\tag{MA-MO.1508}
\]
Complex conjugation is absorbed in $b_T$ by the Riesz
representation; the displayed pairing remains bilinear.
With $U_th(x)=h(x-t)$, the actual translated-source response is
\[
\begin{aligned}
 f_T(t)&=\ell_T(T_{g_T}^0U_th_T)
        =\int h_T(x-t)b_T(x)\,dx\\
       &=C_{h_T^\vee}b_T(t).
\end{aligned}
\tag{MA-MO.1509}
\]
The last equality has the indicated reflected convolution
kernel, without a conjugate.  For shifts in the fixed compact
support of $h_0^\vee$, $w(t-y)^{-1}/w(t)^{-1}$ is uniformly
bounded.  Minkowski's inequality therefore gives the weighted
Young estimate
\[
       \|f_T\|_{L^2(w^{-1})}
          \le C\|b_T\|_{L^2(w^{-1})}
          \le C'\|T_{g_T}\|.
\tag{MA-MO.1510}
\]
All constants are independent of $T$; the modulation changes
only the absolute phase of the convolution kernel.

Choose one $t_0>0$ satisfying the fixed support conditions
in (MA-MO.561).  Put
\[
       \eta_T(t)=\mathbf1_{[t_0,\infty)}(t)
                      e^{-(a+iT)t}f_T(t).
\tag{MA-MO.1511}
\]
Since
$w(t)^{-1}=e^{-2at}(1+e^{-4at})^{-1}(1+t^2)^{-r}$ for $t>0$,
\[
 \|\eta_T\|_{L^2((1+t^2)^{-r}dt)}
           \le\sqrt2\,\|f_T\|_{L^2(w^{-1})}
           \le C\|T_{g_T}\|.
\tag{MA-MO.1512}
\]
Consequently the actual half-line Laplace integral satisfies
\[
\begin{aligned}
 \mathfrak F_T(a+\epsilon+i(T+v))
   &=\int_{t_0}^{\infty}
             e^{-[a+\epsilon+i(T+v)]t}f_T(t)\,dt\\
   &=\sqrt{2\pi}\,\mathcal F_-
                      (e^{-\epsilon t}\eta_T)(v),\\
 \mathfrak F_T(a+\epsilon+i(T+\cdot))
   &\longrightarrow \sqrt{2\pi}\,\mathcal F_-\eta_T
                           \quad\text{strongly in }H^{-r}.
\end{aligned}
\tag{MA-MO.1513}
\]
Here $\mathcal F_-$ is the unitary minus Fourier transform.
The Sobolev norm on the second line is exactly the physical
$L^2((1+t^2)^{-r}dt)$ norm, up to $\sqrt{2\pi}$.
Dominated convergence proves the last line and the uniform
bound $C\|T_{g_T}\|$ for $0<\epsilon\le\epsilon_0$.
Absolute Laplace convergence and its holomorphy for
$\Re w>a$ follow either by weighted Cauchy--Schwarz or
the compact-source translation estimate.  This uses the
assumed bounded actual column, not a proposed inverse
multiplier on arbitrary data.

\emph{3. Exact noncancellation recovers the localized reciprocal.}
For $w=a+\epsilon+i(T+v)$ the original wide Euler calculation,
continued from $\Re w>1/2$ using (MA-MO.1501), is
\[
 \mathfrak F_T(w)
    =H_T(w)V_T(-w)
           -\frac{H_T(w)B_T(w)}{A_\eta(w)}-\mathcal R_T(w).
\tag{MA-MO.1514}
\]
The functions and actual compact remainder are exactly those
in (MA-MO.558)--(MA-MO.567).  In the rational case use
(MA-MO.1311)--(MA-MO.1313): first clear the zeta pole,
continue that identity, and remove the reciprocal's harmless
zero at $w=1/2$.  No output polar coordinate is discarded.
The Laplace function in (MA-MO.1514) is (MA-MO.1513), because
the source identity is valid on every compact translated
input and the wide Laplace integral is absolutely convergent.

Every nonunit ideal, or integer, shifts $J$ past $r_\phi$.
Thus only its unit term contributes to $B_T$.  Modulation
cancels exactly and gives
\[
\begin{aligned}
 H_T(w)&=H_0(a+\epsilon+iv)
             :=\int h_0(x)e^{(a+\epsilon+iv)x}\,dx,\\
 B_T(w)&=B_0(a+\epsilon+iv)
             :=\int_J\phi(x)k_0(x)e^{-(a+\epsilon+iv)x}\,dx,\\
 V_T(-w)&=V_0(-a-\epsilon-iv)
             :=\int_Jk_0(x)e^{-(a+\epsilon+iv)x}\,dx.
\end{aligned}
\tag{MA-MO.1515}
\]
At $v=0$ the first two have fixed nonzero signs and absolute
values uniformly bounded below for $0\le\epsilon\le\epsilon_0$.
Compactness therefore gives $v_*>0$ with
\[
 |U_\epsilon(v)|\ge c>0,\qquad
 U_\epsilon(v)=H_0(a+\epsilon+iv)B_0(a+\epsilon+iv),
 \qquad |v|\le v_* .
\tag{MA-MO.1516}
\]
All derivatives of $U_\epsilon^{-1}$ on a smaller fixed
interval are uniformly bounded and converge as $\epsilon\downarrow0$.
On the compact double-integration region defining $\mathcal R_T$,
the ideal sums are finite with uniformly bounded absolute
values.  Every fixed $v$ or $\epsilon$ derivative only inserts
a bounded power of a compact integration variable.  Hence
$\mathcal R_T(a+\epsilon+i(T+v))$ has uniform bounds in every
fixed smooth seminorm on this rectangle, and converges there
smoothly for each $T$ as $\epsilon\downarrow0$.

For any fixed $\chi_*\in C_c^\infty((-v_*,v_*))$, (MA-MO.1514)
now gives the actual identity
\[
\begin{aligned}
 \frac{\chi_*(v)}{A_\eta(a+\epsilon+i(T+v))}
 ={}&\frac{\chi_*(v)}{U_\epsilon(v)}
  \bigl[H_0(a+\epsilon+iv)V_0(-a-\epsilon-iv)\\
 &\quad-\mathcal R_T(a+\epsilon+i(T+v))
            -\mathfrak F_T(a+\epsilon+i(T+v))\bigr].
\end{aligned}
\tag{MA-MO.1517}
\]
Multiplication by a compact smooth function
is bounded on $H^{-r}$; its operator norm is controlled by
finitely many derivatives, by duality from (MA-MO.716).
The multipliers here converge in those norms.  Equations
(MA-MO.1513)--(MA-MO.1517) prove a strong $H^{-r}$ limit and
\[
 \sup_{0<\epsilon\le\epsilon_0}
 \left\|\frac{\chi_*(\,\cdot\,)}
 {A_\eta(a+\epsilon+i(T+\,\cdot\,))}\right\|_{H^{-r}}
       \le C(1+\|T_{g_T}\|)
       \le C'(2+|T|)^N .
\tag{MA-MO.1518}
\]
The final inequality is (MA-MO.1506), since
$\|g_T\|_{C^N}\le C_N(2+|T|)^N$.
Choose a uniform partition with sufficiently short supports
inside these windows and take $T=t_n$.  Frequency translation
is an isometry on $H^{-r}$, so this is (MA-MO.1502).

For the nontrivial sector, the positive tests have so far
given only the $\eta$ right boundary.  The actual source
reflection identity is
\[
 R\mathscr D_{\eta,g}^{\,\mathrm{raw}}R
       =\mathscr D_{\bar\eta,g^\vee}^{\,\mathrm{raw}},
                         \qquad Ry(x)=y(-x).
\tag{MA-MO.1519}
\]
It is (MA-MO.576), including the original Fourier/dual phases;
$R$ is an isometry of the sum norm.  The negative-support
part of (B) therefore supplies the positive-support dual
family, and the same argument supplies its right-boundary
condition.  No dual estimate is inferred from the positive
$\eta$ family alone.  The rational multiplier is self-dual,
so its original negative-column hypothesis already suffices.
This proves (B)$\Rightarrow$(L).

\emph{4. The boundary condition supplies the actual source inverse.}
First note that the partition convention is immaterial.
Given a second uniform smooth partition, multiply its $m$th
term by $\sum_n\chi_n=1$.  Only a bounded number of comparable
windows occur.  Their smooth multiplier norms are uniform,
and their indices have comparable height.  Thus both the
polynomial estimate and each strong limit in (MA-MO.1502)
pass to the second partition.  Refinement of either partition
if necessary gives the same argument for different fixed
window lengths.

Put $j_{\omega,+}=\sum_n d_{\omega,n,0}$.  This is a tempered
distribution: localized Schwartz tests decrease more rapidly
than the polynomial norms in (MA-MO.1502).  The same dominated
sum gives, for each $f\in\mathscr S$,
\[
\begin{gathered}
 A_\omega(a+\epsilon+i\,\cdot\,)^{-1}f
       \longrightarrow j_{\omega,+}f
                             \quad\text{strongly in }H^{-r},\\
 A_\omega(a+i\,\cdot\,)j_{\omega,+}=1 .
\end{gathered}
\tag{MA-MO.1520}
\]
The second equality follows from the identity at $\epsilon>0$
and the polynomial derivative bounds for the forward factor,
before taking the boundary limit.  It is not an arbitrary
principal-value prescription.  The first assertion also holds
when $f_\epsilon\to f$ in Schwartz space, and after any fixed
polynomial multiplication of $f$.

For $y\in\mathcal E$, let $F_{\omega,y}$ be the original
wide inverse.  At every fixed $b>a$ sufficiently close to $a$
the closed-strip premise supplies a positive margin, so the
already proved wide-source shift identifies
$V_bF_{\omega,y}=A_\omega(b+i\,\cdot\,)^{-1}V_by$.
For this fixed smooth input, $V_by\to V_ay$ in Schwartz space.
Equation (MA-MO.1520) therefore identifies the boundary of
that same original function, and all its fixed derivatives:
\[
 V_aF_{\omega,y}=j_{\omega,+}V_ay,\qquad
 D_x^j(e^{ax}F_{\omega,y})
                   \in L^2((1+x^2)^{-r}dx)\quad(j\ge0).
\tag{MA-MO.1521}
\]
The boundary convergence can also be read distributionally
in $x$: the original function is fixed on every compact
interval, so the limit cannot choose a different inverse.
Local Sobolev bounds give a polynomial bound for
$e^{ax}F_{\omega,y}$ and its derivatives on the left.
On the right it has the original rapid weighted bounds
at every wider line.  Thus $qF_{\omega,y}$ belongs to the
original local source, including the structural finite
vectors.  In the rational case it is precisely the actual
even additive Schwartz vector (MA-MO.714).

\emph{5. Dominated nuclear assembly retains the original matrix.}
We recall the needed estimate rather than infer boundedness
from a scalar boundary value.  If $d\in H^{-r}$ has compact
support and $L\in\mathscr S(\mathbb R^2)$, then
\[
 T_{d,L}f(t)=\langle d(u),L(t,u)f(u)\rangle
       \in\mathcal S_1(H^r,H^R)\quad(R\ge0).
\tag{MA-MO.1522}
\]
Expand $L$ in output Hermite functions.  Their $H^R$ norms
grow polynomially, while all multiplier seminorms of their
input coefficients decrease faster than any power.  The
$H^{-r}$--$H^r$ pairing and (MA-MO.716) give an absolutely
summable rank-one series.  This is the proof of (MA-MO.715)
and shows continuity in $d$ and finitely many Schwartz
seminorms of $L$.  Applied to the $n$th localized distribution,
Schwartz decrease in its support gives, for every $M$,
\[
 \|T_{d_{\omega,n,\epsilon},L}\|_{\mathcal S_1(H^r,H^R)}
          \le C_{L,M}(2+|n|)^{D-M}.
\tag{MA-MO.1523}
\]
Choose $M>D+2$.  Each summand converges in trace norm by
the strong limit in (MA-MO.1502); dominated summation proves
convergence of the entire series.  This remains true when
$L_\epsilon\to L_0$ in Schwartz space.  No uniform local
boundary-convergence modulus is used.

For completeness the two nonsingular source kernels are,
with $z=b+it,w=b+iu$, $G=\mathcal Mg$ and
$P(\xi)=\int e^{\xi x}\phi'(x)\,dx$,
\[
\begin{aligned}
 L_{\omega,b,g}(t,u)
   &=\frac{G(z)P(z-w)}{2\pi}
                         \frac{A_\omega(z)-A_\omega(w)}{z-w},\\
 L^{\mathrm{cross}}_{\eta,b,g}(t,u)
   &=\frac{G(z)A_\eta(z)m_\eta(z)}{2\pi}
                                      \frac{P(-z-w)}{-z-w}.
\end{aligned}
\tag{MA-MO.1524}
\]
These are exactly (MA-MO.718)--(MA-MO.719).  The first
is paired with the $\omega$ reciprocal and the crossed
one with the dual reciprocal.  The divided difference
is removed before distributional pairing; in the second
line $\Re(-z-w)\le-2a$.  Forward polynomial derivative
bounds, rapid decrease of $G$ and $P$, and the actual
local multiplier $m_\eta$ with its conductor, different
and finite phases show Schwartz convergence as $b\downarrow a$.
No reciprocal estimate is hidden in these nonsingular factors.

Apply (MA-MO.1523) and the exact identification
$H_{a,s}\simeq H^r$ of (MA-MO.717).  With the resulting
operators $K_{\eta,a}^\partial(g),O_{\eta,a}^\partial(g)$,
the first actual row is
\[
 \mathscr R_{\eta,a}^\partial(g)y
       =K_{\eta,a}^\partial(g)y-C_gM_\phi y
                              +O_{\eta,a}^\partial(g)Ry .
\tag{MA-MO.1525}
\]
The middle map is trace class from the sum space to
$H_{a,s}$ by its physical half-line kernel (MA-MO.332).
The other two are trace class by (MA-MO.1523).
For $y\in\mathcal E$, pass the exact original $b>a$ source
identity to the boundary using (MA-MO.1521), the strong
input convergence, and trace-norm kernel convergence.
The original smooth theta output is fixed during this
limit.  Thus (MA-MO.1525) agrees with the original row,
without an assumed bounded embedding between whole
varying-width spaces.

In the nontrivial sector the reflected dual row agrees
on $\mathcal E$ by actual source Poisson summation.
The paired map into $H_{a,s}\oplus H_{-a,s}$ is trace
class.  Its range lies in the closed diagonal copy of
$H$, first on the dense original core and hence everywhere.
The inverse diagonal isometry gives the required actual
trace-class endomorphism on the inherited sum norm.

\emph{6. The two rational coordinates and the remaining implications.}
In the rational case put $p_a(x)=\phi(x)e^{(1/2-a)x}$.
It is Schwartz.  Both actual cutoff moments extend boundedly:
\[
\begin{aligned}
 \lambda_0(y)&=\sum_n\langle d_{n,0}(u),
                         \widehat p_a(-u)(V_ay)(u)\rangle,\\
 \lambda_1(y)&=\lambda_0(Ry).
\end{aligned}
\tag{MA-MO.1526}
\]
The localized multiplier norms of $\widehat p_a(-u)$ decrease
rapidly, so (MA-MO.1502) makes this a norm-convergent sum
of bounded functionals.  Bilinear unitary Parseval on
$b>a$, followed by (MA-MO.1520), identifies them with the
original physical moments, with no extra Fourier constant.
The left bounds in (MA-MO.1521), $a<1/2$, and the original
right bounds make $\mathcal MF_y(1/2)$ absolutely convergent.
The original zeta inverse has a zero there, so it is zero.
Accordingly the actual source has
$h_y(0)=-\lambda_1(y)$ and $\int h_y=-\lambda_0(y)$.

Retain the fixed actual sections and $e_0,e_1$ of
(MA-MO.725).  The original diagonal and polar coordinates are
\[
\begin{aligned}
 T_g^\partial y={}&K_a^\partial(g)y-C_gM_\phi y
                       +O_a^\partial(g)Ry\\
 &+G(-1/2)\lambda_1(y)e_0+G(1/2)\lambda_0(y)e_1,\\
 \nu_g(y)={}&(-G(-1/2)\lambda_1(y),-G(1/2)\lambda_0(y)),\\
 \mathscr D_g(y,c)={}&(T_g^\partial y,\nu_g(y))
                                  +\sum_{j=0}^1c_jC_gDs_j .
\end{aligned}
\tag{MA-MO.1527}
\]
These are precisely the original polar signs and finite
columns.  Every term is trace class or finite rank; the
compatible diagonal argument is unchanged after the
displayed polar subtraction.  This proves (S) on $\mathcal W$.

All estimates in (MA-MO.1523)--(MA-MO.1527) involve only
finitely many derivatives of $g$ on a fixed support.
Thus the claimed trace-norm continuity holds.  The original
smooth cores are dense, proving uniqueness.  We have proved
(L)$\Rightarrow$(S), and (S)$\Rightarrow$(B) is immediate,
including restriction to the isometric rational negative
column.  Real and complex versions are equivalent by
Step 1.

Finally choose a nonzero compact test.  The necessary
actual boundary obstructions MA-MO-025 and MA-MO-026
applied to the just-constructed full matrix force
$s>2m-1$ for every boundary zero.  Therefore full individual
retention is a conclusion, not an unproved input to the
sufficiency argument.  MA-MO-028 supplies the actual ideal
trace; with full retention, the relative-trace and
determinant argument of MA-MO-029 applies verbatim.
\end{proof}

\paragraph{Meaning and limits of the criterion.}
The condition (L) contains no cutoff $\phi$.  Therefore universal
boundedness for one normalized smooth step implies the full
trace-class conclusion for every such step, with constants
allowed to depend on the step.  This follows by applying the
proved necessity for the first step and sufficiency for the
second; it asserts no cutoff independence for one isolated test.

This theorem removes the gap between the particular sufficient
cluster budgets and a necessary analytic boundary condition.
MA-MO-030 and MA-MO-035 prove (MA-MO.1502) for their stated
geometries.  The present equivalence does not show that full
individual jet retention alone implies it.  Its local
$H^{-s/2}$ norm permits oscillation and ordinate separation;
it is not replaced by the stronger nonoscillatory minimum-depth
budget.  Neither a hypothetical bad cluster nor failure of
that sufficient budget is asserted to occur.

The conclusion is for a universal fixed-support test family.
It makes no equivalence claim for one predetermined pair
$(g,\phi)$, for an unrestricted one-sided inverse composition,
or for the distinct whole-quotient return cost.  In particular
it does not decide RF002/RS003 or the PR cost gate merely by
rewriting them.  Under A the rational and Dirichlet fixed-gap
branch supplies the criterion; M supplies the indicated
Hecke branch.  Those zero-location inputs remain explicit.
No independent arithmetic deficiency vanishing in FT002
has been proved here.

\paragraph{Source and framework placement.}
The genuine source definitions, dual reflection and wide Euler
identity are MA-MO-024/026, with primary Meyer (2005) locations
in the supplied \texttt{main.tex}, lines 2630--2675,
3592--3707 and 3808--4043.  The exact polar completion is
MA-MO-016; the bilinear functional and universal-test
interfaces are MA-MO-033/034.  The distributional kernel
lemma, source comparison and compatible rows are
MA-MO-027/030/035.  The new necessity is the actual
weighted Young and strong boundary argument
(MA-MO.1508)--(MA-MO.1518), not a scalar trace assignment.
The unconditional fixed-strip forward estimates are
MT-M18.5--9.

In the inherited cote 1 placement, BR-TR-01/02,
scans 079--080 and 111--112, the finite-length
$\mathbb C[z]$-modules encode local divisor multiplicities
and their exterior-power determinants.  The strong
boundary topology and the actual nuclear series above
are additional analytic structure.  Finite-length
$K_0$ classes alone do not supply their norms, their
arithmetic validity, or a new positivity statement.

\section{An actual Euler commutator in Meyer's common-space trace}
\label{sec:meyer-transport}

This section supplies a prime-carrying operator calculation on the
one-sided completions used by Meyer.  It neither assumes a zero-location
theorem nor extends the canonical inverse boundedly to the matching
narrow two-sided norms refuted in MA-MO-001.  The prime contribution is
the trace of the indicated operator, obtained from its kernel.  No
operator is defined by prescribing the value of its trace.

We use $x=\log t$, convolution
$C_g u(x)=\int_{\mathbb R}g(y)u(x-y)\,dy$, and
$T_\ell u(x)=u(x+\ell)$.  Let $S$ be a finite set of rational primes,
let $\mathcal N_S$ be the positive integers whose prime factors are
outside $S$, and set
\[
 \mathcal H_a=L^2(\mathbb R,e^{2ax}\,dx),\qquad a>1/2,
 \qquad
 Z_S=\sum_{n\in\mathcal N_S}n^{-1/2}T_{\log n},\quad
 J_S=\sum_{n\in\mathcal N_S}\mu(n)n^{-1/2}T_{\log n}.
 \tag{MA-MEY.1}
\]
Here $J_S$ denotes the M\"obius inverse.  Meyer's letter $J$ denotes
an inversion operator; we denote that different operator by $J_0$
below.  The notation is deliberately kept separate.

\paragraph{MA-MEY-001 (a smoothed Euler commutator and its actual trace).}
Both series in (MA-MEY.1) converge absolutely in operator norm,
and $Z_SJ_S=J_SZ_S=1$.  Let $g\in C_c^\infty(\mathbb R)$ and let
$\rho\in C^\infty(\mathbb R)$ have compactly supported derivative,
with limiting constants $\rho_-$ at $-\infty$ and $\rho_+$ at
$+\infty$.  Then
\[
 K_{S,\rho}(g)=C_gZ_S[J_S,M_\rho]
             =C_g(M_\rho-Z_SM_\rho J_S)
 \tag{MA-MEY.2}
\]
is trace class on $\mathcal H_a$, and
\[
 \boxed{\operatorname{Tr}_{\mathcal H_a}K_{S,\rho}(g)
 = (\rho_- -\rho_+)
   \sum_{\substack{p\notin S\\ k\ge1}}
     (\log p)p^{-k/2}g(k\log p).}
 \tag{MA-MEY.3}
\]
The prime sum is finite.  The trace-norm bound is uniform for $g$
in a bounded subset of $C_c^\infty$ with fixed compact support.
There is also a trace-class formula before multiplication by $Z_S$:
\[
 \operatorname{Tr}_{\mathcal H_a}C_g[J_S,M_\rho]
 = (\rho_+-\rho_-)
   \sum_{n\in\mathcal N_S}\mu(n)(\log n)n^{-1/2}g(\log n).
 \tag{MA-MEY.4}
\]

\begin{proof}
Since $\|T_{\log n}\|_{\mathcal H_a}=n^{-a}$, both series are
absolutely convergent for $a>1/2$.  Their product is an absolutely
convergent double series.  Its coefficient at $T_{\log r}$ is
$r^{-1/2}\sum_{d\mid r}\mu(d)$ for $r\in\mathcal N_S$.
This is one at $r=1$ and zero otherwise, proving the inverse
identities.  Compact smooth convolution is bounded: conjugating by
the unitary $W_a u(x)=e^{ax}u(x)$ gives
$W_aC_gW_a^{-1}=C_{g_a}$, where $g_a(y)=e^{ay}g(y)$, and
Young's inequality applies.

We give a trace-norm estimate, not merely a formal diagonal
calculation.  If $q\in C_c^\infty(\mathbb R)$ and
$b\in C_c^\infty(\mathbb R)$ has support in an interval of length
$L$, then
\[
 \|C_qM_b\|_1\le
 c_q(1+L)\max_{0\le j\le2}\|b^{(j)}\|_\infty.
 \tag{MA-MEY.5}
\]
To prove this, choose a fixed smooth partition of unity
$\sum_{r\in\mathbb Z}\eta(x-r)=1$ with compactly supported
$\eta$.  Only $O(1+L)$ of the functions
$b_r(x)=\eta(x-r)b(x)$ are nonzero.  Simultaneously translating
both variables reduces the kernel
$q(x-z)b_r(z)$ to a smooth kernel in a fixed compact rectangle,
whose size depends only on $\operatorname{supp}q$ and $\eta$.
Choose a larger fixed square with the rectangle in its interior
and expand the kernel in the orthonormal Fourier basis on that
square, extended by zero to $\mathbb R$.  Integrating twice in
each variable bounds its coefficients by
\[
 c_q\max_{0\le j\le2}\|b^{(j)}\|_\infty
 (1+m^2)^{-1}(1+n^2)^{-1}.
\]
Only derivatives of $q$ through order four and of $b_r$ through
order two occur.  The absolute sum of these coefficients bounds
the trace norm of the corresponding rank-one series.  Summing
over $r$ proves (MA-MEY.5).  The same absolute expansion, whose
diagonal terms are integrable by Cauchy--Schwarz, proves that the
trace of each of these compact smooth kernels is its diagonal
integral.

For $\ell\ge0$, put
$D_\ell(x)=\rho(x+\ell)-\rho(x)$.  If
$\operatorname{supp}\rho'\subseteq[A,B]$, then
$\operatorname{supp}D_\ell\subseteq[A-\ell,B]$;
all its derivatives through order two are bounded uniformly in
$\ell$.  The identity
\[
 [T_\ell,M_\rho]=M_{D_\ell}T_\ell
\]
and (MA-MEY.5), after conjugation by $W_a$, imply
\[
 \|C_g[T_\ell,M_\rho]\|_{1,\mathcal H_a}
       \le c_{g,a,\rho}e^{-a\ell}(1+\ell).
 \tag{MA-MEY.6}
\]
Here the final translation is unitary on ordinary $L^2$, and
$W_aT_\ell W_a^{-1}=e^{-a\ell}T_\ell$.
Thus the series for $C_g[J_S,M_\rho]$ converges absolutely in
trace norm.  More precisely, the double series
\[
 K_{S,\rho}(g)
  =\sum_{m,n\in\mathcal N_S}
       \mu(n)(mn)^{-1/2}
       C_gT_{\log m}[T_{\log n},M_\rho]
 \tag{MA-MEY.7}
\]
does so, since the trace norm of the term without its scalar
coefficient is at most
$c_{g,a,\rho}(mn)^{-a}(1+\log n)$.  This follows from
(MA-MEY.5) also when $D_{\log n}$ is translated by $\log m$.
The convergent majorant is
\[
 c_{g,a,\rho}
 \sum_{m\ge1}m^{-a-1/2}
 \sum_{n\ge1}n^{-a-1/2}(1+\log n).
\]
The constants use only a fixed finite list of seminorms of $g$
on its support.  This proves the stated uniformity.

Write $v=\log m$ and $w=\log n$.  The kernel with respect to
ordinary $dz$ of the term $C_gT_v[T_w,M_\rho]$ is
\[
     g(x+v+w-z)\{\rho(z)-\rho(z-w)\}.
 \tag{MA-MEY.8}
\]
Its conjugate by $W_a$ has the additional factor $e^{a(x-z)}$,
which is one on the diagonal.  Equivalently, its diagonal
factor is $e^{-a(v+w)}g_a(v+w)=g(v+w)$.  The compact-kernel
trace calculation just proved therefore gives
\[
 \operatorname{Tr} C_gT_v[T_w,M_\rho]
   =g(v+w)\int_{\mathbb R}\{\rho(x)-\rho(x-w)\}\,dx
   =g(v+w)w(\rho_+-\rho_-).
 \tag{MA-MEY.9}
\]
For the last equality, write the difference as
$\int_{x-w}^x\rho'(u)\,du$ and use Fubini; the integrand is
absolutely integrable since $\rho'$ has compact support.

Taking traces in the absolutely trace-norm convergent series
(MA-MEY.7), and regrouping by $r=mn$, gives
\[
 \operatorname{Tr}K_{S,\rho}(g)
 = (\rho_+-\rho_-)
   \sum_{r\in\mathcal N_S}r^{-1/2}g(\log r)
              \sum_{n\mid r}\mu(n)\log n.
\]
The inner sum is $-\Lambda(r)$.  Indeed, for a prime power the
only squarefree divisors are $1$ and $p$, giving $-\log p$;
if $r$ has at least two distinct prime factors, the coefficient
of each $\log p$ cancels in the product expansion of
$\prod_{q\mid r}(1-1)$.  It is zero also at $r=1$.
This proves (MA-MEY.3).  Taking only $m=1$ in the kernel
calculation proves (MA-MEY.4).  Compactness of
$\operatorname{supp}g$ makes both final arithmetic sums finite.
\end{proof}

The multiplication by $Z_S$ is essential to the interpretation:
(MA-MEY.4) differentiates the reciprocal Euler multiplier,
whereas (MA-MEY.3) produces its logarithmic derivative.  With
the Mellin convention of MA-TR-101, their multipliers are
$J_S(z)=1/\zeta_{\overline S}(z+1/2)$ and
$J_S'(z)/J_S(z)=-\zeta_{\overline S}'/
\zeta_{\overline S}(z+1/2)$ in $\Re z>1/2$.
The signs in (MA-MEY.3)--(MA-MEY.4) are fixed by the displayed
operator order and the limits of $\rho$.

\paragraph{MA-MEY-002 (a boundary contribution invisible to compact cutoffs).}
If $\rho$ has compact support, then
$\operatorname{Tr}K_{S,\rho}(g)=0$.  If
$\phi\in C^\infty(\mathbb R)$ satisfies $\phi=1$ near
$-\infty$ and $\phi=0$ near $+\infty$, then
\[
 \operatorname{Tr}K_{S,\phi}(g)=P_{\overline S}^+(g),\qquad
 P_{\overline S}^+(g)=
 \sum_{\substack{p\notin S\\k\ge1}}
       (\log p)p^{-k/2}g(k\log p).
 \tag{MA-MEY.10}
\]
There are uniformly bounded $\rho_R\in C_c^\infty$ with
$\rho_R\to\phi$ pointwise, and for every such family
$K_{S,\rho_R}(g)\to K_{S,\phi}(g)$ strongly on $\mathcal H_a$.
Whenever $P_{\overline S}^+(g)\ne0$, this convergence is not in
trace norm: for every $R$,
\[
 \|K_{S,\rho_R}(g)-K_{S,\phi}(g)\|_1
       \ge |P_{\overline S}^+(g)|.
 \tag{MA-MEY.11}
\]

\begin{proof}
The first two claims follow from (MA-MEY.3), since the two
limits of a compact cutoff are zero.  Multiplying $\phi$ by
a smooth function that is zero on $(-\infty,-R-1]$ and one on
$[-R,\infty)$ supplies the stated family.  Uniform boundedness
and dominated convergence give $M_{\rho_R}\to M_\phi$ strongly.
All other operators in (MA-MEY.2) are fixed bounded operators,
so the same is true of $K$.  Each difference is trace class;
the trace inequality $|\operatorname{Tr}T|\le\|T\|_1$ gives
(MA-MEY.11).  For a concrete nonzero case, take $S=\varnothing$
and $g\ge0$ supported sufficiently close to $\log2$ that it
contains no other positive logarithm of a prime power, with
$g(\log2)=1$.  Then
$P_{\overline S}^+(g)=\log2/\sqrt2$.
\end{proof}

This is a precise obstruction to replacing Meyer's one-sided
cutoff by compact cutoffs in this particular trace-class
commutator.  It does not exclude other regularizations.
Nor is there a contradiction with cyclicity of the trace:
$C_gM_\phi$ and $C_gZ_SM_\phi J_S$ have not separately been
shown trace class.  Their difference has the trace just computed.
For compact $\rho$, both separated terms are trace class and
cyclicity indeed gives zero.

\paragraph{MA-MEY-003 (the fixed-$S$ defect in Meyer's actual common space).}
Restrict Meyer's construction over $\mathbb Q$ to the trivial
compact norm-one isotype, with that compact group of volume one.
Use the finite place set $S\cup\{\infty\}$ and a symmetric
cutoff $\phi(x)+\phi(-x)=1$ in logarithmic coordinates, with
the endpoint values in MA-MEY-002.  Let $P_+^S=i_+p_+^S$
and $P_-=i_-p_-$ be Meyer's actual approximate projection and
projection.  On the centered common completion
$\mathcal H_a\oplus\mathcal H_{-a}$, $a>1/2$, their smoothed
difference is trace class, and
\[
 \operatorname{Tr}_{\mathcal H_a\oplus\mathcal H_{-a}}
        C_g(P_- -P_+^S)
 =-B_{\overline S}(g),\qquad
 B_{\overline S}(g)=P_{\overline S}^+(g)
                         +P_{\overline S}^+(g^\vee),
 \quad g^\vee(x)=g(-x).
 \tag{MA-MEY.12}
\]
Here $C_g$ acts diagonally on the common space.  Thus the
omitted-prime correction in Meyer's trace additivity is an
actual trace of the common-space defect, even for a fixed $S$
which has not absorbed every prime meeting the support of $g$.

\begin{proof}
We first record the exact source maps to fix their spaces and
normalizations.  In Meyer's notation let $\mathfrak F$ be the
additive Fourier transform, $J_0 f(t)=t^{-1}f(t^{-1})$ the
inversion, and $\mathsf F=\mathfrak FJ_0$ the equivariant
Fourier transform.  The common smooth space is
$\mathcal S(\mathcal C)_{>}\oplus
 \mathcal S(\mathcal C)_{<}$.
The actual maps are
\begin{align*}
 i_+v&=(\Sigma v,J_0\Sigma\mathfrak F^*v),\\
 p_+^S(v_0,v_1)
   &=M_{1-\phi}\mathcal L_S^{-1}v_0
       +\mathfrak F M_{1-\phi}\mathcal L_S^{-1}J_0v_1,\\
 i_-v&=(v,v),\qquad
 p_-(v_0,v_1)=M_{1-\phi}v_0+M_\phi v_1.
 \tag{MA-MEY.13}
\end{align*}
Here $\Sigma$ is arithmetic summation and
$\mathcal L_S=\prod_{p\notin S}(1-\lambda_p^{-1})^{-1}$
is its realization on the $S$-local space $\mathcal T$.
These statements are Meyer (2005), Lemmas~5.3--5.4,
pp.~39--40, and the displays immediately following Lemma~5.4.
The maps are defined on the local and global smooth spaces
there; the proof of the nuclear defect supplies the indicated
Hilbert extensions, rather than assuming that all projection
entries individually act on that Hilbert direct sum.
For $\mathbb Q$, every set $S\cup\{\infty\}$ here is
``sufficiently large'' in Meyer's structural sense: unique
factorization identifies
$\mathbb Q^\times/\mathbb Z[S^{-1}]^\times$ with the free
valuation group on primes outside $S$. This is the definition
in his \S5.1, before the summation construction; it does not
require enlarging $S$ to fit the test support.

Center functions by $v(t)\mapsto u(x)=e^{x/2}v(e^x)$.
The raw Hilbert weight $\alpha>1$ becomes
$a=\alpha-1/2>1/2$, and its paired weight $1-\alpha$
becomes $-a$.  Inversion becomes the unitary reflection
$Ru(x)=u(-x)$ from $\mathcal H_a$ to $\mathcal H_{-a}$.
The operator $\mathcal L_S$ becomes $Z_S$, its inverse
becomes $J_S$, and raw convolution by
$f(t)=t^{-1/2}g(\log t)$ becomes $C_g$.  These statements
follow directly from $\lambda_t v(u)=v(t^{-1}u)$; in
particular the half-density factor in $\lambda_{1/n}$ is
$n^{-1/2}$, exactly as in (MA-MEY.1).

Meyer's Lemma~5.6, pp.~41--42, proves that the smoothed
difference $C_g(P_- -P_+^S)$ is nuclear from the raw weighted
Hilbert direct sum into the common smooth space.  Composing
with the continuous inclusion into the same Hilbert sum makes
it trace class there. The nuclear trace on the common smooth
space equals this Hilbert trace: one nuclear representation
$\sum_j\lambda_j v_j\otimes\ell_j$ of the Hilbert-to-smooth
map gives both traces as the same absolutely convergent series
$\sum_j\lambda_j\ell_j(v_j)$, with the inclusion understood.
This is precisely the compatible-pairing argument in Meyer's
Corollary \texttt{cor:T\_on\_Twist}; it uses the actual nuclear
lift, not just a range inclusion for a trace-class endomorphism.
This applies in particular on the
invariant compact isotype considered here.  No condition
$q_p>\sup|\operatorname{supp}f|$ is required for that lemma;
the support-adapted enlargement of $S$ is used later to make
the trace vanish.

From the definitions (MA-MEY.13), the two diagonal entries
of the centered difference $P_+^S-P_-$ are $D$ and $RDR$,
where
\[
 D=Z_SM_{1-\phi}J_S-M_{1-\phi}
   =M_\phi-Z_SM_\phi J_S=Z_S[J_S,M_\phi].
 \tag{MA-MEY.14}
\]
For the second diagonal one uses
$RM_{1-\phi}R=M_\phi$ and $Z_SJ_S=1$.
The trace of a trace-class operator on a Hilbert direct sum
is the sum of the traces of its diagonal entries.  Since
$RC_gR=C_{g^\vee}$, (MA-MEY.3) gives
\[
 \operatorname{Tr}C_g(P_+^S-P_-)
 =\operatorname{Tr}_{\mathcal H_a}C_gD
  +\operatorname{Tr}_{\mathcal H_a}C_{g^\vee}D
 =B_{\overline S}(g).
\]
This proves (MA-MEY.12), including its sign.
\end{proof}

There is a printed sign inconsistency that matters for this
nonzero refinement.  In the proof of Meyer's Lemma~5.6,
the display labelled \texttt{eq:global\_perturbation} in the
attached primary \LaTeX\ labels its left side $P_- -P_+^S$
but gives the entries of $P_+^S-P_-$.  The definitions of
$P_+^S$ and $P_-$ just before \S5.3 determine the sign in
(MA-MEY.12).  Meyer's later vanishing argument is unchanged
by this printed sign; no failure of the global trace theorem
is inferred from it.

For completeness, the off-diagonal part of that source proof
also has signs that should be read with the defining Poisson
formula. Here is the short corrected reduction. In raw
coordinates put $q=1-\phi$, $L=\mathcal L_S$ and
$\mathcal C_u=\lambda(u)$ for convolution. Write
$s=\operatorname{Sing}(J_0f)$, so Poisson gives
$LJ_0f=h-\phi s$ with $h\in\mathcal S(\mathcal C_S)_\cup$.
The relevant off-diagonal operator is
\[
 \begin{split}
 O_f&=\mathcal C_f
          (L\mathfrak F M_qL^{-1}-M_\phi J_0)\\
 &=-[\mathcal C_f,M_\phi]J_0
   -J_0\bigl([M_q,\mathcal C_{LJ_0f}]
                         -\mathcal C_sM_q\bigr)L^{-1},\\
 [M_q,\mathcal C_{LJ_0f}]-\mathcal C_sM_q
   &=[M_q,\mathcal C_h]
                  -M_q\mathcal C_{\phi s}-\mathcal C_{qs}M_q.
 \end{split}
\]
These identities follow from $J_0L\mathfrak F=L+\operatorname{Sing}$
and the convolution intertwining relation. The last two kernels
are separately nuclear from the raw $L^2_\alpha$, $\alpha>1$,
to $\mathcal S(\mathcal C_S)_{<}$: $s$ is a linear combination
of $1$ and $t^{-1}$, and their supports respectively force
$1\ll x\ll y$ and $1\ll y\ll x$, up to fixed constants.
After weighting by $x^\beta y^{-\alpha}$, $\beta<0$, both
kernels and every logarithmic derivative decay rapidly.
The other term is the source's nuclear commutator lemma.
Reflection gives the required $>$ target. Thus the corrected
signs preserve the full nuclear conclusion of Lemma~5.6.

\subsection{The completed geometric piece and the remaining transport}

For precision, the trace additivity used in the proof of
Meyer's Theorem~5.8, pp.~43--44, is
\begin{align*}
 \chi_{\pi_+}(f)-\chi_{\pi_-}(f)
 &=\operatorname{Tr}_{\mathcal T}
       \lambda(f)(1-p_+^Si_+)
       -\operatorname{Tr}_{\mathrm{common}}
          \lambda(f)(P_- -P_+^S),\\
 1-p_+^Si_+&=M_\phi-\mathsf F M_\phi\mathsf F^{-1}.
 \tag{MA-MEY.15}
\end{align*}
The quotient $\pi_+$ is the two-dimensional polar quotient;
$\pi_-$ is the negative smooth quotient.  The maps and the
additivity are taken in the actual common space, using Poisson
summation and nuclearity.  The first trace is the actual local
Fourier-cutoff trace.  Meyer proves that its computation on
$\mathcal T$ agrees with its computation on
$L^2(\mathcal C_S)_{1/2}$ in Proposition
\texttt{pro:T\_on\_Twist} and Corollary
\texttt{cor:T\_on\_Twist}; the local value is calculated in
Lemma~3.3 and the following local principal-value calculation.
Combining (MA-MEY.12) and (MA-MEY.15) over $\mathbb Q$ gives
\[
 \chi_{\pi_+}(f)-\chi_{\pi_-}(f)
   =W_S(f)+B_{\overline S}(g),
 \qquad f(t)=t^{-1/2}g(\log t),
 \tag{MA-MEY.16}
\]
where $W_S$ is this independently computed local trace, with
the source's fixed Haar and principal-value conventions.
The discriminant term is zero for $\mathbb Q$.  This proof
uses the source's trace additivity before its final choice of
$S$, and (MA-MEY.3) computes the term that that choice makes
zero.  It does not begin with the numerical explicit formula.

In particular, Meyer takes $S$ large enough that every omitted
prime exceeds the support bound for $f$ and its reciprocal.
Then every summand in $B_{\overline S}(g)$ vanishes.  Formula
(MA-MEY.16) retains the same common-space construction at a
fixed $S$ and exposes the actual prime-carrying remainder.
For $S=\varnothing$ and a test concentrated at $\log2$, the
remainder is exactly $\log2/\sqrt2$, with its reciprocal term
zero.  Thus a prime contribution can be recovered at the
operator level without a bounded narrow-strip M\"obius inverse.

This closes the explicit-commutator and fixed-$S$ common-space
piece.  It does not identify every individual orthogonal
cutoff $M_\rho C_gP_G$ of the two-sided weighted Hilbert model
with an additive $S$-local cutoff kernel.  The equality of
smooth quotient spaces is MT-M24; transporting their
convolution traces to a specified Hilbert quotient requires
the corresponding actual maps and nuclear estimates, under
the precise MT-M26 interior-strip hypothesis.  Likewise,
MA-MEY-002 shows why the compact-cutoff limit cannot simply
be substituted for Meyer's one-sided defect in trace norm.
The source's local additive kernel and the equivariant-Fourier
cutoff in (MA-MEY.15) are actual operators on related local
spaces; equality of their limiting scalar local traces alone
is not an identification of their finite cutoffs.

The geometric interpretation here is the trace of an extension
and its boundary defect.  Meyer explicitly identifies the
cutoff commutator in Lemma~3.2 with the cyclic $1$-cocycle of
a $1$-summable odd Fredholm module in the paragraph following
that lemma, pp.~19--20; the Euler commutator above evaluates
the same one-sided boundary mechanism on the arithmetic
summation operator.  The inherited Grothendieck comparison
with determinants and logarithmic derivatives remains
ALPOGE\_BRIDGES BR-TR-01/02 and F-SM-11/12, cote~1,
scans~079--080 and 111--112.  These archive scans were not
freshly reread in this check.  No identification of this
nonproper analytic limiting operation with a proper algebraic
pushforward in EGA, FGA or SGA is asserted.

\paragraph{Source and check record.}
Meyer, R. (2005). \emph{On a representation of the idele class
group related to primes and zeros of $L$-functions}
(arXiv:math/0311468, version~3, 16 March 2005).
\url{https://arxiv.org/abs/math/0311468}.
The journal version is \emph{Duke Mathematical Journal},
127, 519--595, \url{https://doi.org/10.1215/S0012-7094-04-12734-4};
all page locators here refer to the 53-page arXiv version.
The attached primary \LaTeX\ was used for the complete proofs
of Lemmas~3.2--3.3, the following local principal-value lemmas,
the nuclear-kernel lemma \texttt{lem:mother\_of\_nuclearity},
Proposition/Corollary \texttt{T\_on\_Twist}, and
Lemmas~5.3--5.6 and Theorem~5.8 as used here.  The primary PDF
was consulted for version and page/number locators.  This is
a targeted reading of the complete invoked local/global
proofs, not a claim to have freshly re-proved all upstream
bornological analysis in the paper.  MA-MEY-001--003 have
complete proofs above and individual post-writing check
events in the run record.

\section{Smoothing from the actual Hilbert quotient to Meyer's quotient}
\label{sec:hilbert-meyer-bridge}

The comparison in this section is an operator factorization.  It does
not identify two traces by first assigning both the same zero sum.
It gives a nuclear smoothing map from the specified Hilbert quotient
to the actual smooth quotient of MT-M24, and then uses trace cyclicity.
The uniform-interior hypothesis of MT-M26 is retained throughout.
Consequently the comparison is unconditional on a wide strip and
applies to every positive narrow Dirichlet strip using $A$.  It is
not an independent proof that all zeros belong to those narrow strips.

\subsection{The spaces and the one hypothesis about the zero divisor}

Fix one primitive unitary Hecke function and the completed nontrivial
zero divisor of $\Xi_\eta$.  Let $\mathcal E$ and its closed ideal
$\mathcal I_\eta$ be those of GRH\_NIGHT MT-M18, and put
$Q=\mathcal E/\mathcal I_\eta$, with quotient map $q$.
For $a\in\mathbb R$, $s\ge0$ and $I=[u,v]$, $u<v$, use the actual
spaces of MT-M26:
\[
 \begin{aligned}
 H&=L^2\!\left(\mathbb R,
     \max(e^{2ux},e^{2vx})(1+x^2)^{s/2}\,dx\right),\\
 J&=\overline{T_a\mathcal I_\eta}^{\,H},
       &T_af(x)&=e^{ax}f(x),\\
 Z&=H/J,\qquad
       i_a:Q\longrightarrow Z,\quad [f]\longmapsto[T_af].
 \end{aligned}
 \tag{MA-HM.1}
\]
The hypothesis is that a closed strip contains the real parts of
every nontrivial completed zero and has a positive margin inside
the available Mellin strip:
\[
       \operatorname{Re}\rho\in[c_0,c_1]\quad\hbox{for every }\rho,
       \qquad [c_0-a,c_1-a]\subset(u,v).
 \tag{MA-HM.2}
\]
This is exactly MT-M26 (M26.2), not a condition on only the zeros
already selected inside $I$.  The closed interval $[0,1]$ works
unconditionally.  At $a=1/2$ and $I=[-\alpha,\alpha]$, it works
unconditionally for $\alpha>1/2$, and for every $\alpha>0$ under
the RH of the particular function.  In the Dirichlet case that
last input is $A$; the general sufficient input is $M$.

Choose the adapted rectangles $\Omega_n$ strictly between the
edges $a+u,a+v$ as in the proof of MT-M26, and put $\nu_n=2+|n|$.
The exact inherited data used below are
\[
 \begin{aligned}
 d_n&=\dim J_n=O(\log\nu_n),\\
 R_m\mathcal T_n&=\delta_{mn}\operatorname{id}_{J_n},
       &\mathcal T_n&:J_n\longrightarrow\mathcal E,\\
 p_k(\mathcal T_n z)&\le C_k\nu_n^{B_k}\|z\|_n,\\
 r_n&:Z\longrightarrow J_n,\qquad
       \sup_n\|r_n\|\le C,\qquad r_ni_a=R_n .
 \end{aligned}
 \tag{MA-HM.3}
\]
Here $p_k$ is any enumeration of the defining exponential Schwartz
seminorms.  The $J_n$ are the quotient holomorphic supremum-norm
algebras with their actual full zero jets; $R_n$ is the block map
on $Q$.  MT-M18 identifies $Q$ topologically with the rapid block
space $s(J)$ and supplies the lifts.  MT-M26 (M26.5)--(M26.9)
supplies the same data in the narrower adapted rectangles and
the uniform Hilbert restriction bound.  In particular no
unstable basis of individual zero evaluations is being substituted
for these stable finite algebras.

For $g\in\mathcal E$, let $M_{g,n}$ be multiplication in $J_n$
by the class of $w\mapsto\mathcal Mg(w-a)$.  Set
\[
 b_n(g)=\sup_{w\in\partial\Omega_n}|\mathcal Mg(w-a)|.
 \qquad
 \|M_{g,n}\|\le b_n(g),\qquad
 b_n(g)\le C_A(g)\nu_n^{-A}\quad(A>0).
 \tag{MA-HM.4}
\]
Every $C_A(g)$ can be bounded by finitely many seminorms of $g$,
and is bounded on bounded subsets of $\mathcal E$, by the
integration-by-parts estimate in MT-M26 (M26.12).

\subsection{A lift into the smooth test space, not just into its completion}

\paragraph{MA-HM-001 (order-zero nuclear smoothing through the actual quotient).}
Under (MA-HM.2), the series
\[
 \widetilde S_g z
       =\sum_{n\in\mathbb Z}\mathcal T_nM_{g,n}r_nz
       \quad(z\in Z),\qquad
 S_g=q\widetilde S_g:Z\longrightarrow Q
 \tag{MA-HM.5}
\]
converges absolutely in every seminorm of $\mathcal E$.
It defines a continuous linear map into $\mathcal E$, and is
of order zero: it admits a nuclear rank-one representation with
coefficients in $\ell^p$ for every $p>0$, bounded target vectors
in $\mathcal E$, and bounded functionals on $Z$.
The maps $g\mapsto\widetilde S_g$ are continuous for uniform
convergence on the unit ball of $Z$ in every seminorm of
$\mathcal E$.  Bounded test families give bounded order-zero
nuclear families.  These assertions also hold for $S_g$.

Writing $k_g(x)=e^{-ax}g(x)$, one has identities of actual operators
\[
             S_gi_a=C_{k_g}^{Q},\qquad
             i_aS_g=C_g^{Z}.
 \tag{MA-HM.6}
\]
Thus $C_g^Z$ factors through the stronger smooth quotient topology,
although the unsmoothed dense injection $i_a$ is not a topological
isomorphism onto $Z$.  The quotient-valued map $S_g$ is independent
of the chosen block lifts.  It is a genuine smoothed norm
intertwiner: with $U_yf(x)=f(x-y)$ on both radial models,
\[
             S_gU_y^Z=e^{-ay}U_y^QS_g .
 \tag{MA-HM.6a}
\]
Only the chosen representative lift $\widetilde S_g$ into
$\mathcal E$ need depend on the interpolation section.

\emph{Proof.}
For every seminorm $p_k$ and every $z\in Z$, (MA-HM.3)--(MA-HM.4)
give
\[
 \sum_n p_k(\mathcal T_nM_{g,n}r_nz)
      \le CC_k\|z\|_Z\sum_n\nu_n^{B_k}b_n(g)<\infty.
 \tag{MA-HM.7}
\]
The completeness of $\mathcal E$ gives a unique sum.  Choosing
$A>B_k+2$ in (MA-HM.4) bounds its $p_k$ seminorm by
a constant times $C_A(g)\|z\|_Z$.  This proves continuity and
the asserted continuous dependence on $g$.  It also justifies
applying every continuous block restriction to the sum.

Per-seminorm convergence alone would not be an adequate
substitute for a nuclear tensor representation, so we give that
representation explicitly.  Choose an Auerbach basis
$e_{n,j},e_{n,j}^*$ of each nonzero $J_n$,
$1\le j\le d_n$, with both norms equal to one.  Put
\[
     \ell_{n,j}=e_{n,j}^*r_n\in Z',\qquad
           \|\ell_{n,j}\|\le C.
 \tag{MA-HM.8}
\]
Then the finite-dimensional identity in each block yields
\[
       \widetilde S_g
           =\sum_{n,j}\mathcal T_nM_{g,n}e_{n,j}
                                      \otimes\ell_{n,j}.
 \tag{MA-HM.9}
\]
If $b_n(g)>0$, write its summand as
$\lambda_{n,j}v_{n,j}\otimes\ell_{n,j}$ with
\[
   \lambda_{n,j}=\sqrt{b_n(g)},\qquad
   v_{n,j}=b_n(g)^{-1/2}\mathcal T_nM_{g,n}e_{n,j}.
 \tag{MA-HM.10}
\]
If $b_n(g)=0$, the block multiplier is zero and its summands
are omitted.  For each fixed $k$,
\[
       p_k(v_{n,j})\le C_k\nu_n^{B_k}\sqrt{b_n(g)}.
 \tag{MA-HM.11}
\]
Thus the entire family of vectors is bounded in $\mathcal E$.
For every $p>0$,
\[
       \sum_{n,j}|\lambda_{n,j}|^p
          =\sum_nd_n b_n(g)^{p/2}<\infty.
 \tag{MA-HM.12}
\]
These are precisely the bounded-vector and equibounded-functional
conditions for an order-zero nuclear map.  If a Banach-disk
description is desired, take the closed absolutely convex hull
of the bounded target vectors: completeness of the Fr\'echet
space makes it a complete bounded disk.  Formula (MA-HM.9)
then factors through that Banach disk and the indicated
$\ell^p$ coefficients.  No common polynomial exponent for
all Schwartz seminorms is required.

For a bounded test set $\mathcal B\subset\mathcal E$, replace
$b_n(g)$ in the normalizing square root by
$b_n^{\mathcal B}=\sup_{g\in\mathcal B}b_n(g)$.
This is again a rapidly decreasing sequence.  The target
vectors in (MA-HM.10), now with $g$ varying over $\mathcal B$,
belong to one bounded disk by (MA-HM.11); the functionals
in (MA-HM.8) are fixed and uniformly bounded; and
(MA-HM.12) gives one summable coefficient sequence for each
$p>0$.  This proves the asserted boundedness in the
order-zero operator bornology.  Applying the continuous
quotient map $q$ proves the same assertions for $S_g$.

To prove the first operator identity in (MA-HM.6), let $[f]\in Q$.
By (MA-HM.3), the $n$th block of $S_gi_a[f]$ is
$M_{g,n}R_n[f]$.  On the other hand
$\mathcal M k_g(w)=\mathcal Mg(w-a)$, so the same is true
of $C_{k_g}^Q[f]$.  The block isomorphism $Q\simeq s(J)$
makes the two classes equal.  This proves an equality in
the actual quotient, not merely equality of a list of scalar
eigenvalues.

For the second identity, convolution is bounded on $H$ by
the translation estimate (M26.10).  It preserves $J$ since
\[
       C_gT_af=T_aC_{k_g}f
       \quad(f\in\mathcal E),
 \tag{MA-HM.13}
\]
and $\mathcal I_\eta$ is a convolution ideal.  Consequently
$C_g^Z$ is the bounded operator induced from the literal
ambient convolution.  The same equation proves
$C_g^Zi_a=i_aC_{k_g}^Q$.  Composing the first identity in
(MA-HM.6) with $i_a$ therefore shows that $i_aS_g$ and
$C_g^Z$ agree on $i_aQ$.  That image is dense, because
$T_a\mathcal E$ contains $C_c^\infty(\mathbb R)$, dense in
the weighted $L^2$ space.  Both operators on $Z$ are
bounded, so they agree on all of $Z$.

If two continuous maps $Z\to Q$ have the first identity in
(MA-HM.6), they agree on the dense subspace $i_aQ$, hence
everywhere because $Q$ is Hausdorff.  This proves the
independence of $S_g$.  The same argument proves (MA-HM.6a):
$U_y^Zi_a=e^{-ay}i_aU_y^Q$ by the definition of $T_a$,
and smooth convolution commutes with $U_y^Q$.  The two
sides of (MA-HM.6a) are continuous maps $Z\to Q$ and
therefore agree everywhere once they agree on $i_aQ$.

Only the available data (MA-HM.3) used (MA-HM.2).
No scalar trace formula, simplicity assumption, or positivity
was used in this proof.  The construction proves the
factorization even if one uses only the continuity and
density of $i_a$ before invoking the injectivity supplied
by the rest of MT-M26. \hfill$\square$

\subsection{Trace cyclicity and the actual Meyer common-space expression}

\paragraph{MA-HM-002 (the Hilbert trace is the actual Meyer quotient trace).}
Under the same hypothesis, one has
\[
 \operatorname{Tr}_{Z}C_g^Z
       =\operatorname{tr}_{Q}C_{k_g}^Q
       \qquad(g\in\mathcal E),
 \tag{MA-HM.14}
\]
where the left side is the Hilbert trace and the right side
is the canonical nuclear trace on the smooth quotient.
For $g\in C_c^\infty(\mathbb R)$ let $\beta$ be the idele
character in the inverse-Euler convention of MT-M22 and set
\[
       f_{g,a}(b)=\beta(b)^{-1}|b|^{-a}g(\log|b|)
                \quad(b\in C_K).
 \tag{MA-HM.15}
\]
Probability Haar on $C_K^1$ is understood.  Then
\[
       \operatorname{Tr}_{Z}C_g^Z
                 =\chi_{\pi_-}(f_{g,a}),
 \tag{MA-HM.16}
\]
where $\pi_-$ is Meyer's actual quotient
$H_-/(H_+\cap H_-)$, restricted to the indicated compact
character component.  In particular, this is an operator
trace comparison transported by MT-M24, not an inference
from equal spectra.

Let $\mathcal A_S(f)$ and $\mathcal D_S(f)$ denote the actual
traces in (MA-MEY.15):
\[
 \begin{aligned}
 \mathcal A_S(f)
   &=\operatorname{tr}_{\mathcal T}
           \lambda(f)(1-p_+^Si_+),\\
 \mathcal D_S(f)
   &=\operatorname{tr}_{\mathrm{common}}
           \lambda(f)(P_- -P_+^S).
 \end{aligned}
 \tag{MA-HM.17}
\]
The first is the local Fourier-cutoff trace and the second
is Meyer's common-space defect.  With all source domains
and nuclear extensions as in MA-MEY-003, the exact identity is
\[
       \boxed{\displaystyle
       \operatorname{Tr}_{Z}C_g^Z
          =\chi_{\pi_+}(f_{g,a})
                 -\mathcal A_S(f_{g,a})
                 +\mathcal D_S(f_{g,a}).}
 \tag{MA-HM.18}
\]
Over $\mathbb Q$, on the trivial compact character component
and with $a=1/2$, (MA-MEY.12) computes the defect as the
negative omitted-prime contribution.  Therefore
\[
       \operatorname{Tr}_{Z}C_g^Z
           =P(g)-W_S(f_{g,1/2})-B_{\overline S}(g).
 \tag{MA-HM.19}
\]
All three expressions on the right come from the actual
Meyer maps.  The omitted-prime term is the trace-class
Euler commutator of MA-MEY-001, with the sign fixed from
the definitions of $P_+^S$ and $P_-$.

\emph{Proof.}
Apply $q$ to the nuclear expansion in the proof of
MA-HM-001 and denote its bounded vectors in $Q$ by
$u_{n,j}=qv_{n,j}$.  The factorization gives
\[
 \begin{aligned}
 i_aS_g&=\sum_{n,j}\lambda_{n,j}
                     (i_au_{n,j})\otimes\ell_{n,j},\\
 S_gi_a&=\sum_{n,j}\lambda_{n,j}
                      u_{n,j}\otimes(\ell_{n,j}i_a).
 \end{aligned}
 \tag{MA-HM.20}
\]
The first series converges in Hilbert trace norm: $i_a$
maps the bounded vector family into a bounded subset of
$Z$, the functionals have uniformly bounded norm, and the
coefficients are in $\ell^1$.  The second is an order-zero
nuclear representation on $Q$, because the composed
functionals are equicontinuous and the target vectors are
bounded.  Each scalar trace series converges absolutely,
and both traces are
\[
          \sum_{n,j}\lambda_{n,j}
                            \ell_{n,j}(i_au_{n,j}).
 \tag{MA-HM.21}
\]
This proves cyclicity here directly from one common
rank-one representation.  It does not appeal to a trace
identity for arbitrary bounded lifts on $H$.

For completeness, the trace on $Q$ is unambiguous.
The order-zero framework is Meyer (2005), \S2.1, in
particular Theorem~2.1 and the trace pairing preceding it.
One can also see the needed approximation property
directly from MT-M18.  On $s(J)$ let $E_N$ keep the
blocks with $|n|\le N$.  For a bounded subset
$\mathcal U\subset s(J)$,
\[
    \sup_{z\in\mathcal U}q_k((1-E_N)z)
       \le (N+2)^{-1}\sup_{z\in\mathcal U}q_{k+1}(z)
                        \longrightarrow0.
 \tag{MA-HM.22}
\]
Thus finite-rank maps approximate the identity uniformly
on bounded sets.  The canonical nuclear trace agrees
with the trace obtained from these finite-rank
approximations.  Equations (MA-HM.6) and (MA-HM.21)
give (MA-HM.14).

The topological and bounded-bornological isomorphism in
MT-M22/MT-M24 sends $[F]\in Q$ to the class of
$b\mapsto\beta(b)^{-1}F(\log|b|)$ in the actual Meyer
negative quotient.  Its translation formula is
(M22.4).  Integrating that formula against (MA-HM.15)
cancels the character $\beta$ and leaves convolution by
$e^{-ax}g(x)=k_g(x)$.  The compact group has volume one.
Thus it carries $C_{k_g}^Q$ to the actual integrated
operator $\pi_-(f_{g,a})$.
It also carries the bounded vectors and equicontinuous
functionals in (MA-HM.20) to corresponding ones, preserving
each rank-one trace pairing.  This proves (MA-HM.16)
without using the zero-divisor trace formula.

Finally the actual trace additivity in the proof of
Meyer Theorem~5.8, author pp.~43--44, is exactly
\[
       \chi_{\pi_+}(f)-\chi_{\pi_-}(f)
                         =\mathcal A_S(f)-\mathcal D_S(f).
 \tag{MA-HM.23}
\]
Its spaces, maps and fixed-$S$ correction are recorded
and checked in MA-MEY.13--MA-MEY.16.  Rearrangement with
(MA-HM.16) proves (MA-HM.18), and (MA-MEY.12) proves
(MA-HM.19).  The polar trace for the trivial rational
character is
$P(g)=\int_{\mathbb R}g(x)(e^{-x/2}+e^{x/2})\,dx$.
For a nontrivial finite-order compact character the
polar quotient is killed by character projection.
\hfill$\square$

\subsection{A structural obstruction on a genuinely partial strip}

\paragraph{MA-HM-003 (exterior jet blocks are lost before any smoothing).}
Retain the spaces in (MA-HM.1), but do not assume (MA-HM.2).
The natural map $i_a:Q\to Z$ is still continuous with dense
image.  For each nontrivial completed zero $\rho$, let
$Q_\rho\subset Q$ be its finite-dimensional local jet block,
of dimension $m_\rho$, supplied by MT-M18 and the Chinese
remainder decomposition of its finite block algebra.
If
\[
              \operatorname{Re}(\rho-a)\notin[u,v],
 \tag{MA-HM.24}
\]
then
\[
                         i_a(Q_\rho)=0.
 \tag{MA-HM.25}
\]
This statement concerns zeros strictly outside the closed
strip.  It makes no assertion that boundary jets are lost.

Consequently, if a linear map $T_g:Z\to Q$, even without a
continuity assumption, is to satisfy
\[
                    T_gi_a=C_{e^{-a\cdot}g}^{Q},
 \tag{MA-HM.26}
\]
then at every zero in (MA-HM.24) it is necessary that
\[
               (\mathcal Mg)^{(j)}(\rho-a)=0
                       \quad(0\le j<m_\rho).
 \tag{MA-HM.27}
\]
For any one such zero there exists
$g\in C_c^\infty(\mathbb R)$, identically zero near $0$,
with $\mathcal Mg(\rho-a)=1$; for that allowed test
no map satisfying (MA-HM.26) exists.

In particular, for $a=1/2$ the following property is
equivalent to RH for this one function: for every
$\alpha>0$ and every compact smooth $g$ with $g(0)=0$,
the actual quotient $Z_{1/2,[-\alpha,\alpha],s}$ admits
the two smoothing factorizations (MA-HM.6) through
the full smooth quotient $Q$.  Under RH those maps
can be taken of order zero by MA-HM-001; the reverse
implication already follows from (MA-HM.26).

\emph{Proof.}
Continuity and density of $i_a$ use only the definitions:
$T_a:\mathcal E\to H$ is continuous, its image of
$\mathcal I_\eta$ lies in $J$, and its range contains
the compact smooth functions.  The ideal and its
Hilbert closure are invariant under all translations.
For the actual weight in (MA-HM.1) the sharper form
of the elementary translation estimate is
\[
 \|U_y\|_{H\to H}
 \le 2^{s/4}\max(e^{uy},e^{vy})(1+y^2)^{s/4}.
 \tag{MA-HM.28}
\]
Indeed the exponential part of $w(x+y)/w(x)$ is at
most $\max(e^{2uy},e^{2vy})$, and
$1+(x+y)^2\le2(1+x^2)(1+y^2)$ controls the polynomial
part.  The same bound holds for the induced group
on $Z$.  No zero-location assertion enters this bound.

On $Q_\rho$, $U_y^Q$ is multiplication by the class
of $e^{wy}$ in the local algebra at $w=\rho$.
After the intertwining identity
$U_y^Zi_a=e^{-ay}i_aU_y^Q$, its image
$Y_\rho=i_a(Q_\rho)$ is a finite-dimensional invariant
space.  If it is nonzero, the sole eigenvalue of
$U_1^Z$ on it is $e^{\rho-a}$: it is a nonzero
quotient of the generalized eigenspace with precisely
that eigenvalue.  Choose an eigenvector in $Y_\rho$.
For positive integers $n$, (MA-HM.28) then implies
\[
       e^{n\operatorname{Re}(\rho-a)}
          \le 2^{s/4}e^{vn}(1+n^2)^{s/4}.
 \tag{MA-HM.29}
\]
This is impossible when
$\operatorname{Re}(\rho-a)>v$.
Applying the same argument to $U_{-1}$ makes it
impossible when $\operatorname{Re}(\rho-a)<u$.
Hence $Y_\rho=0$, proving (MA-HM.25) for the entire
local jet algebra, not only its reduced eigenline.

Equation (MA-HM.26) would force
$C_{e^{-a\cdot}g}^Q$ to vanish on each such $Q_\rho$.
There it is multiplication by the class of
$\mathcal Mg(w-a)$ in
$\mathbb C[w]/(w-\rho)^{m_\rho}$.
Multiplication is the zero operator exactly when
this class is zero, which is (MA-HM.27).
Choose $b\in C_c^\infty((1,2))$ with $\int b=1$
and set $g(x)=e^{-(\rho-a)x}b(x)$.
Then $g$ vanishes near $0$ and
$\mathcal Mg(\rho-a)=1$, giving the stated obstruction.

If a zero is off the critical line, choose
$0<\alpha<|\operatorname{Re}\rho-1/2|$ and apply
this test to its exterior block.  Thus the
every-positive-width factorization property forces
RH.  Conversely RH supplies (MA-HM.2) with
$c_0=c_1=1/2$ for every such width, and MA-HM-001
constructs the required maps.
\hfill$\square$

The obstruction is loss of an invariant finite-length
spectral block under the actual completion map.
It is not an obstruction inferred from equal or
unequal scalar traces.  It also does not assert that
no useful map can factor through a quotient retaining
only the interior blocks.  Such a quotient has
discarded the exterior information in the full Meyer
quotient; composing it with smoothing or the Euler
commutator cannot recover that information by a
return map satisfying (MA-HM.26).
Boundary behavior remains separate, as in MA-TR-103.

\subsection{The actual cutoff limit after the stated strip premise}

\paragraph{MA-HM-004 (completed operator transport under the containing-strip premise).}
Specialize to the trivial rational character, $a=1/2$,
$I=[-\alpha,\alpha]$, and assume (MA-HM.2).
Use the literal orthogonal projections $P_Z$ and
$P_G=1-P_Z$ of the corrected Hilbert model.
Let $g\in C_c^\infty(\mathbb R)$ satisfy $g(0)=0$,
put $h(t)=g(\log t)$, and let
$\rho_\nu\in C_c^\infty(\mathbb R)$ be uniformly
bounded and tend to $1$ almost everywhere.
Then the actual geometric-order operators are
trace class and
\[
 \begin{aligned}
 \lim_\nu\operatorname{Tr}_{H}
                  (M_{\rho_\nu}C_gP_G)
   &=\mathcal A_S(f_{g,1/2})
          -\mathcal D_S(f_{g,1/2})
          -\chi_{\pi_+}(f_{g,1/2})\\
   &=W_S(f_{g,1/2})+B_{\overline S}(g)-P(g)\\
   &=A_0(h)+B(h)-P(h).
 \end{aligned}
 \tag{MA-HM.30}
\]
Every equality is obtained from traces of the actual
operators described in MA-LIFT-001--003,
MA-HM-001--002, and MA-MEY-001--003.
The proof does not substitute an equality of scalar
zero sums.  It gives a trace-limit comparison;
it does not identify each individual finite
orthogonal cutoff operator with a local Fourier
cutoff operator.

\emph{Proof.}
MA-LIFT-002 proves that the literal ambient operator
$C_gP_Z$ is trace class and that its trace equals
the quotient convolution trace.  We use that
operator assertion and its proof by trace cyclicity,
not the subsequent evaluation of the quotient
trace as a zero sum.  MA-LIFT-003 proves
\[
 \|(M_{\rho_\nu}-1)C_gP_Z\|_1\longrightarrow0,
 \qquad
 \operatorname{Tr}_{H}(M_\rho C_g)
                         =g(0)\int_{\mathbb R}\rho(x)\,dx.
 \tag{MA-HM.31}
\]
Thus, since $g(0)=0$,
\[
 \lim_\nu\operatorname{Tr}_{H}
                  (M_{\rho_\nu}C_gP_G)
       =-\operatorname{Tr}_{H}(C_gP_Z)
       =-\operatorname{Tr}_{Z}(C_g^Z).
 \tag{MA-HM.32}
\]
These are limits of traces of the actual trace-class
operators in the order specified by the draft.
Only the spectral product in (MA-HM.31) is asserted
to converge in trace norm.  The cusp of the
weight and its finite boundary contribution were
included in MA-LIFT-002.

Now apply the operator trace comparison
(MA-HM.18).  Its negative is the first line of
(MA-HM.30).  The common-space defect is
$\mathcal D_S=-B_{\overline S}$ by MA-MEY-003,
with the omitted-prime contribution computed
from the trace-class commutator in MA-MEY-001.
The polar quotient has trace $P(g)$.
Finally the independently computed local
Fourier-cutoff trace is
$W_S=A_0(h)+B_S(h)$ when $h(1)=0$.
Combining its finite prime terms with the
omitted-prime trace gives $B(h)$, proving the
last line.  This local evaluation is the one
used in MA-MEY.15--MA-MEY.16; it requires no
zero-location input.
\hfill$\square$

For every rational or Dirichlet narrow strip,
the operator part of this transport can therefore
be instantiated using $A$.  On unconditional
wide strips it needs no such input.  The
operator bridge is complete under its stated
premise; obtaining that premise on arbitrarily
narrow strips without using zero location is
the remaining independent-method question.

\subsection{What this bridge settles and what it does not}

The smoothed Hilbert operator now has an explicit route to
the genuine common-space geometry:
$C_g^Z=i_aS_g$, $S_gi_a=C_{k_g}^Q$, the MT-M24 map identifies
that quotient operator with $\pi_-(f_{g,a})$, and
(MA-HM.18) expresses its trace by the actual local and
arithmetic-defect operators.  The route carries prime
contributions through MA-MEY-001/003 and makes no demand
that a narrow-strip M\"obius inverse be bounded.

The only zero-location premise in this route is
(MA-HM.2), used in obtaining the uniformly bounded
Hilbert restrictions and the interior adapted blocks.
It is harmless on wide strips.  Using $A$ to verify it
on every Dirichlet narrow strip is a valid conditional
application of the already granted theorem, but is
circular in a proposed zero-location-free proof of the
method under test.  For the imaginary quadratic
character of MA-CH-003, the corresponding invocation
needs its own RH, sufficiently $M$.

MA-LIFT-001--003 and MA-HM-004 also give the
specific ambient cutoff trace limit under that same
premise.  This discharges the ambient comparison
for the corrected model at the level asserted in
(MA-HM.30).  No equality of individual finite
cutoffs is inferred.  In particular, a compact
cutoff cannot silently replace Meyer's one-sided
commutator: MA-MEY-002 gives an explicit failure
of trace-norm convergence for that substitution.
MA-HM-003 explains structurally why a completion
that genuinely drops exterior zero blocks cannot
be used to recover the full smooth quotient by
the same smoothing factorization.

The functional-analytic location is Grothendieck's
nuclear tensor product and trace pairing: the common
rank-one expansion (MA-HM.20) is the concrete object
that makes the trace comparison valid.  It fits
the inherited cote~1 determinant/trace programme,
scans~079--080 and~111--112, through ALPOGE\_BRIDGES
BR-TR-01/02 and F-SM-11/12.  Those archive placements
are carried from the supplied runs; no new archive
scan reading is claimed here.  The quotient still
records the full finite-length spectral jet algebras.
Neither this nuclear comparison nor its Hilbert
factorization supplies faithful positivity, replaces
those jets by a reduced divisor, or realizes them
as the cohomology of a proper scheme.

\paragraph{Exact source scope.}
The inherited analytic data are GRH\_NIGHT MT-M18
(M18.1)--(M18.4), (M18.13), and MT-M26
(M26.2), (M26.5)--(M26.12); the actual smooth-quotient
comparison is MT-M22 (M22.2)--(M22.4), (M22.10) and
MT-M24 (M24.3), (M24.6)--(M24.9).
Meyer's published source is
\emph{On a representation of the idele class group related
to primes and zeros of $L$-functions},
\emph{Duke Mathematical Journal, 127} (2005), 519--595,
arXiv:math/0311468v3: \S2.1 for the nuclear trace
pairing and Theorem~2.1; Lemmas~5.4--5.6 for the
common-space maps and nuclear defect; Theorem~5.8,
author pp.~43--44, for trace additivity.
Its Theorem~5.11 or the resulting scalar zero sum
is not needed to establish the operator trace
identity (MA-HM.18).

\section{The actual partial-strip quotient at polynomial weight zero}
\label{sec:hardy-partial-strip}

This section computes the actual Hilbert closure and its convolution
trace when the polynomial weight parameter is exactly $s=0$.
The strip need not contain every completed zero.  No uniform distance
of its interior zeros from its edges is assumed.  The argument uses
the canonical generator of MT-M18, a Hardy-space closure calculation,
and actual compressed multiplication operators.  Its conclusion is
strictly about the partial quotient; it supplies no return map to the
full smooth quotient when an exterior block is lost, as MA-HM-003
already proves.

\subsection{The norm and the actual closed ideal}

Fix the primitive unitary Hecke function, its entire pole-removed
completion $\Xi=\Xi_\eta$, and
\[
 \mathcal I_\eta=\overline{\psi_\eta*\mathcal E}^{\,\mathcal E},
 \qquad \mathcal M\psi_\eta=\Xi_\eta
\]
of MT-M18.  Let $a\in\mathbb R$, $u<v$, $L=v-u$, and $c=(u+v)/2$.
Throughout this section the polynomial weight parameter is $s=0$.
Write
\[
 \begin{aligned}
 H_{\max}&=L^2(\mathbb R,\max(e^{2ux},e^{2vx})\,dx),\\
 H_\Sigma&=L^2(\mathbb R,(e^{2ux}+e^{2vx})\,dx),\\
 J&=\overline{T_a\mathcal I_\eta}^{\,H_{\max}}
   =\overline{T_a\mathcal I_\eta}^{\,H_\Sigma},
       &T_af(x)&=e^{ax}f(x),\\
 Z_{\max}&=H_{\max}/J,\qquad Z_\Sigma=H_\Sigma/J,\qquad
 \Pi=\{z:u<\operatorname{Re}z<v\}.
 \end{aligned}
 \tag{MA-HY.1}
\]
The same set $J$ occurs because
\[
       \|f\|_{\max}\le\|f\|_\Sigma\le\sqrt2\,\|f\|_{\max}.
 \tag{MA-HY.2}
\]
The identity on representatives consequently induces a bounded
invertible map between the two actual quotients.  It intertwines
translation and convolution.  In particular it preserves trace
class and trace by bounded similarity; it need not identify the
two orthogonal projections onto the respective complements of $J$.

Let $\mathscr H^2(\Pi)$ denote the translation-invariant Hardy space
with norm
\[
 \|F\|_{\mathscr H^2}^2
 =\sup_{0<\delta<L/2}\frac1{2\pi}\int_{\mathbb R}
       \bigl(|F(u+\delta+it)|^2+
             |F(v-\delta+it)|^2\bigr)\,dt .
 \tag{MA-HY.3}
\]
This is the unweighted boundary-norm Hardy space, rather than the
space obtained by composition alone with a disk conformal map.
The Mellin transform is a unitary map
\[
             \mathcal M:H_\Sigma\longrightarrow\mathscr H^2(\Pi).
 \tag{MA-HY.4}
\]
Indeed, Plancherel on each vertical line identifies the integral
in (MA-HY.3) with
\[
 \int_{\mathbb R}|f(x)|^2
       \bigl(e^{2(u+\delta)x}+e^{2(v-\delta)x}\bigr)\,dx .
\]
Its integrand increases as $\delta$ decreases to zero, since the
parenthesis equals $2e^{2cx}\cosh((L-2\delta)x)$.
Monotone convergence proves the norm identity.  Conversely, if
$F$ satisfies (MA-HY.3), take inverse Fourier transforms $f_\sigma$
of $t\mapsto F(\sigma+it)$ on the interior lines.  The
Cauchy--Riemann equation gives, in distributions,
$\partial_\sigma f_\sigma=x f_\sigma$.  Testing first on compact
$x$ intervals and solving this scalar equation gives
$f_\sigma=e^{\sigma x}f$ for one distribution $f$, which is a
locally square-integrable function because any one $f_\sigma$
is in $L^2$.  The uniform bound and Fatou's lemma as
$\sigma\to u,v$ give $f\in H_\Sigma$.  The interior Mellin
integral is absolutely convergent by Cauchy--Schwarz against
$(e^{2ux}+e^{2vx})^{-1}$, and equals $F$ by Fourier uniqueness.
This proves surjectivity, and also the $L^2$ boundary limits.

For later use choose the explicit conformal maps
\[
 q(w)=c+\frac{2L}{\pi}\arctan w,\qquad
 \varphi(z)=\tan\frac{\pi(z-c)}{2L},\qquad
       q'(w)=\frac{2L}{\pi(1+w^2)} .
 \tag{MA-HY.5}
\]
Thus $q:\mathbb D\to\Pi$, and the two infinite ends correspond to
$w=i$ and $w=-i$.  With any analytic square root of $q'$,
\[
     \mathcal U F(w)=q'(w)^{1/2}F(q(w))
        :\mathscr H^2(\Pi)\longrightarrow H^2(\mathbb D)
 \tag{MA-HY.6}
\]
is unitary.  The normalization follows from the boundary change
of variables $|q'(e^{i\theta})|\,d\theta=dt$, with both edges
included.  The onto statement, including the Hardy-space
condition at the two ends, is the rotated and rescaled
$p=2$ case of Bakan--Kaijser (2007), Theorems~2.1--2.2 and
Corollary~2.1, pp.~349--350, with proofs in \S4, pp.~352--355.
In particular the half-derivative factor in (MA-HY.6) is required.

There is also a direct verification of the precise onto statement,
including its end condition.  The reproducing kernel of the
Mellin image in (MA-HY.4) is
\[
\begin{aligned}
 K_\Pi(z,\zeta)
 &=\int_{\mathbb R}
     \frac{e^{(z+\overline\zeta)x}}{e^{2ux}+e^{2vx}}\,dx\\
 &=\frac{\pi}{2L}
     \sec\frac{\pi(z+\overline\zeta-2c)}{2L}.
\end{aligned}
\tag{MA-HY.6a}
\]
The integral converges absolutely for $z,\zeta\in\Pi$.
Write its denominator as $2e^{2cx}\cosh(Lx)$ and substitute
$y=e^{2Lx}$.  The resulting beta integral is
$\int_0^\infty y^{r-1}(1+y)^{-1}\,dy=\pi/\sin(\pi r)$,
where $0<\operatorname{Re}r<1$, proving the displayed value.
Cauchy--Schwarz against the same weight proves boundedness
of every point evaluation.  Choose the analytic square root
of $q'$ positive at zero.  The cosine addition identity gives
\[
 q'(w)^{1/2}\overline{q'(\omega)^{1/2}}\,
 K_\Pi(q(w),q(\omega))
             =\frac1{1-w\overline\omega}.
\tag{MA-HY.6b}
\]
Indeed the secant argument becomes
$\arctan w+\arctan\overline\omega$, whose cosine is
\[
 (1-w\overline\omega)/
 \sqrt{(1+w^2)(1+\overline\omega^2)}.
\]
The right side of (MA-HY.6b) is the disk $H^2$ kernel.
The weighted composition map in (MA-HY.6), with its transported
norm, therefore has exactly the disk kernel.  Equality of the
inner products on the dense spans of kernel functions proves
that its range completes to all of $H^2(\mathbb D)$.
Point evaluations identify this unitary with the stated
weighted composition map, proving both normalization and
surjectivity directly.

It commutes with multiplication in the exact sense
\[
        \mathcal U M_F\mathcal U^{-1}=M_{F\circ q}
                    \qquad(F\in H^\infty(\Pi)).
 \tag{MA-HY.7}
\]

\paragraph{MA-HY-001 (the actual closure is the interior Blaschke ideal).}
Let $B$ be the disk Blaschke product whose zeros, with their
actual multiplicities, are
\[
       \varphi(\rho-a),\qquad
       \Xi(\rho)=0,\quad u<\operatorname{Re}(\rho-a)<v.
 \tag{MA-HY.8}
\]
If there are no such zeros, take $B=1$.
Then
\[
 \mathcal U\mathcal M J=B H^2(\mathbb D),\qquad
       Z_\Sigma\simeq K_B:=H^2(\mathbb D)\ominus B H^2(\mathbb D).
 \tag{MA-HY.9}
\]
The first equality concerns the actual ideal closure in (MA-HY.1).
Equivalently it is the simultaneous vanishing, to the prescribed
multiplicities, of the Mellin jets at all zeros strictly inside
$a+\Pi$.  Zeros on the edges and outside the closed strip impose
no further conditions in this $s=0$ completion.

\begin{proof}
Put $\Psi(z)=\Xi(a+z)$.  MT-M18 (M18.6) shows that $\Psi$ is
bounded on the closed strip, and $\Psi$ is entire.  Multiplication
by $\Psi$ is therefore bounded on $\mathscr H^2(\Pi)$.
The map $T_a$ is an automorphism of $\mathcal E$, and
\[
       T_a(\psi_\eta*f)=(T_a\psi_\eta)*(T_af).
\]
The inclusion $\mathcal E\to H_\Sigma$ is continuous and has
dense range, because it contains the compact smooth functions.
Taking the closure first in $\mathcal E$ and then in $H_\Sigma$
therefore gives exactly
\[
 \mathcal M J
     =\overline{\Psi\,\mathcal M(\mathcal E)}^{\,\mathscr H^2}
     =\overline{\Psi\,\mathscr H^2(\Pi)}^{\,\mathscr H^2}.
 \tag{MA-HY.10}
\]
Thus there is no replacement of the specified theta ideal by
an independently chosen spectral kernel.

Set $F_0=\Psi\circ q\in H^\infty(\mathbb D)$.
Its canonical factorization is
\[
       F_0=B S_\mu O,\qquad
 |S_\mu(w)|=\exp\left(-\int_{\mathbb T}
              \frac{1-|w|^2}{|\zeta-w|^2}\,d\mu(\zeta)\right),
 \tag{MA-HY.11}
\]
where $\mu$ is a finite positive measure singular with respect
to arc length, and $O$ is outer.  The unimodular constants are
irrelevant.  The existence and this precise form are Beurling's
Factorization Lemma, \S7, pp.~245--246; his Theorem~I, \S9,
pp.~247--249, says that
$\overline{F_0H^2}=B S_\mu H^2$.  Here the zeros of $B$
are exactly (MA-HY.8), so their Blaschke summability is also
part of the factorization theorem.  It remains to prove
that $\mu=0$; this is a substantive step at the two ends.

Write $M=\|F_0\|_\infty>0$.  The outer-function formula gives
$|O|\le M$, whence
\[
           \log|F_0(w)|\le\log M-P\mu(w).
 \tag{MA-HY.12}
\]
On any closed arc avoiding $\{\pm i\}$ and the boundary zeros
of $F_0$, the function $F_0$ extends holomorphically across the
arc and is bounded away from zero in a neighborhood of it.
Equation (MA-HY.12) bounds $P\mu$ in the adjacent disk sector.
This forces $\mu$ to be absolutely continuous on the arc:
for every sufficiently short subarc $I$ there, the Poisson
kernel at radius $1-|I|$ over its midpoint is at least
$C/|I|$ on $I$, so $\mu(I)\le C'|I|$.  Since $\mu$ is
singular, its restriction to that arc is zero.

The finite boundary zeros of $F_0$ are isolated away from
$\{\pm i\}$, and form a countable set.  Hence $\mu$ is
supported on that countable set together with the two ends,
and must be purely atomic.  At a finite boundary zero
$\zeta$, analytic continuation gives
$|F_0(r\zeta)|\ge C(1-r)^m$ for $r$ sufficiently close to one,
where $m$ is the finite order of the zero.  An atom
$\mu(\{\zeta\})=\tau>0$ would instead give, by (MA-HY.12),
\[
       |F_0(r\zeta)|\le M
                 \exp\left(-\tau\frac{1+r}{1-r}\right),
\]
which is incompatible with the preceding lower bound.
All such atoms vanish.

Finally MT-M18 (M18.8) supplies ordinates $h_n\to+\infty$
and $h_n\to-\infty$ such that, at the fixed real part $a+c$,
\[
       \log|\Xi(a+c+ih_n)|
                 \ge -C(1+|h_n|)-C\log(2+|h_n|).
 \tag{MA-HY.13}
\]
These are unconditional selected-height minimum-modulus bounds.
By (MA-HY.5),
\[
       \varphi(c+it)=i\tanh\frac{\pi t}{2L},\qquad
       1-\left|\varphi(c+it)\right|
                    \asymp e^{-\pi|t|/L}.
\]
An atom at $i$ or at $-i$ would therefore make the upper
bound in (MA-HY.12), along the corresponding sequence,
at most $\log M-C'\exp(\pi|h_n|/L)$.  This contradicts
(MA-HY.13).  Thus $\mu=0$.  Equations (MA-HY.7),
(MA-HY.10) and Beurling's closure theorem prove (MA-HY.9).

Divisibility by $B$ in $H^2$ is equivalent to vanishing
at its zeros with their multiplicities, by the same
canonical factorization.  On a boundary line, in contrast,
the formal evaluation of a Mellin derivative of order $j$
would have squared dual norm
\[
    \int_{\mathbb R}
       \frac{|x|^{2j}e^{2ux}}{e^{2ux}+e^{2vx}}\,dx=\infty
\]
at the left edge, and the analogous integral is infinite
at the right edge.  Riesz representation, tested on compact
smooth functions, shows that none of these boundary jets
is continuous.  The closure calculation already proves that
no additional boundary restriction survives.
\end{proof}

\subsection{Compressed multipliers and the partial trace}

For $F\in H^\infty(\Pi)$ write
\[
       T_F=P_{K_B}M_{F\circ q}\big|_{K_B}.
 \tag{MA-HY.14}
\]
This is the quotient multiplier, represented on the orthogonal
model space.  In general $M_{F\circ q}$ itself does not leave
$K_B$ invariant.  Because $BH^2$ is invariant under every
bounded analytic multiplier, the compressed operators satisfy
\[
               T_{F_1}T_{F_2}=T_{F_1F_2}.
 \tag{MA-HY.15}
\]
Indeed the omitted term in inserting $1=P_{K_B}+P_{BH^2}$
has range in $BH^2$, and its final projection to $K_B$ is zero.

\paragraph{MA-HY-002 (trace class without a uniform edge margin).}
For every $g\in C_c^\infty(\mathbb R)$, the actual quotient
convolution $C_g^{Z_{\max}}$ is trace class and
\[
       \operatorname{Tr}_{Z_{\max}}C_g^{Z_{\max}}
       =\sum_{\substack{\Xi(\rho)=0\\
                  u<\operatorname{Re}(\rho-a)<v}}
                     m_\rho\,\mathcal Mg(\rho-a).
 \tag{MA-HY.16}
\]
The series is absolutely convergent.  The result also holds
for $g\in\mathcal E$.  For every real $\kappa>v$ and integer
$k\ge2$, the actual quotient translation generator $A_Z$
satisfies
\[
                    (\kappa-A_Z)^{-k}\in\mathcal S_1(Z_{\max}).
 \tag{MA-HY.17}
\]
All these assertions are unconditional for the one fixed
Hecke function.

\begin{proof}
List the zeros $\lambda_j=\rho_j-a$ strictly inside $\Pi$,
with repetitions for multiplicity, and put $b_j=\varphi(\lambda_j)$.
Let $B_n$ be the product of the first $n$ normalized Blaschke
factors.  The functions
\[
       e_j(w)=B_{j-1}(w)
               \frac{\sqrt{1-|b_j|^2}}{1-\overline b_j w}
 \tag{MA-HY.18}
\]
form an orthonormal basis of $K_B$.  To verify this without
any separation condition, use
$K_{UV}=K_U\oplus U K_V$ for inner functions $U,V$.
A single Blaschke factor has one-dimensional model space
spanned by the displayed normalized reproducing kernel.
This proves the finite-product assertion inductively.
Moreover $\bigcap_nB_nH^2=BH^2$: division by each $B_n$
preserves the boundary norm, and a weakly convergent
subsequence of these bounded quotients, or canonical
factorization, gives division by $B$ in $H^2$.
Taking orthogonal complements proves completeness in the
infinite case.  Repeated zeros cause no change in the argument.

Since $|B_{j-1}|=1$ almost everywhere on the circle,
the boundary density $|e_j|^2\,d\theta/(2\pi)$ is exactly
harmonic measure at $b_j$.  Transferring it to the strip,
for $\lambda_j=u+x_j+it_j$ with $0<x_j<L$, its two
densities on the edges are
\[
 \begin{aligned}
 P_{x_j,0}(t-t_j)&=
  \frac{\sin(\pi x_j/L)}
       {2L\{\cosh(\pi(t-t_j)/L)-\cos(\pi x_j/L)\}},\\
 P_{x_j,L}(t-t_j)&=
  \frac{\sin(\pi x_j/L)}
       {2L\{\cosh(\pi(t-t_j)/L)+\cos(\pi x_j/L)\}}.
 \end{aligned}
 \tag{MA-HY.19}
\]
These formulas follow by the disk Poisson kernel and
(MA-HY.5); their combined integral is one.  In particular,
for $|t-t_j|\ge1$ both densities are bounded by
$C_L e^{-\pi|t-t_j|/L}$, uniformly for every
$0<x_j<L$, including arbitrarily close to either edge.

Suppose $F\in H^\infty(\Pi)$ has boundary bounds
$|F(u+it)|+|F(v+it)|\le C_N(1+|t|)^{-N}$.
For $|t_j|\ge2$, divide the integral in (MA-HY.19) into
$|t-t_j|\le |t_j|/2$ and its complement.  On the first
set use $|t|\ge|t_j|/2$ and total harmonic mass at most
one.  On the complement use the exponential tail and
$\|F\|_\infty$.  Enlarging the constant handles
$|t_j|<2$ and gives
\[
 \|M_{F\circ q}e_j\|_{H^2}^2
       =\int_{\partial\Pi}|F|^2\,d\omega_{\lambda_j}
       \le C_{F,N,L}(1+|t_j|)^{-2N}.
 \tag{MA-HY.20}
\]
The constant has no inverse edge-distance factor.
MT-M18's unconditional local zero count is
$O(\log(2+|T|))$ in each fixed-height window, including
multiplicity.  Consequently (MA-HY.20) is summable
over $j$ for any $N>1/2$.  Thus
$M_{F\circ q}:K_B\to H^2$ is Hilbert--Schmidt, and so
is its compression $T_F$.

Put $G=\mathcal Mg$.  Its values on the closed strip
decay faster than every power, by integration by parts.
Choose a real $d>\max(|u|,|v|)+1$ and factor
\[
       G(z)=F_1(z)F_2(z),\qquad
       F_1(z)=\frac1{d^2-z^2},\quad
       F_2(z)=(d^2-z^2)G(z).
 \tag{MA-HY.21}
\]
Both factors are bounded analytic on $\Pi$, with
boundary decay sufficient for (MA-HY.20).
They give two Hilbert--Schmidt compressed multipliers.
Equation (MA-HY.15) then gives
$T_G=T_{F_1}T_{F_2}\in\mathcal S_1(K_B)$.
On the actual Mellin quotient convolution is
multiplication by $G$, by Fubini on the dense smooth
subspace and bounded extension.  Equations (MA-HY.2),
(MA-HY.4), (MA-HY.7) and (MA-HY.9) therefore identify
this trace-class operator with the specified quotient
convolution, up to bounded similarity.

Its trace can now be computed in the actual orthonormal
basis (MA-HY.18), since trace class has already been proved.
The analytic Poisson formula gives
\[
 \langle T_Ge_j,e_j\rangle
      =\int_{\mathbb T}(G\circ q)|e_j|^2\,\frac{d\theta}{2\pi}
      =(G\circ q)(b_j)=G(\lambda_j).
 \tag{MA-HY.22}
\]
The local zero count and the rapid vertical decay also
give absolute convergence of these diagonal values.
Summing proves (MA-HY.16), with all multiplicities.

Finally the actual translation group on $H_{\max}$ has
$\|U_t\|\le e^{vt}$ for $t\ge0$.  Its quotient has the same
bound, so for $\kappa>v$ its resolvent is the operator-norm
convergent integral
$\int_0^\infty e^{-\kappa t}U_t\,dt$.
Its Mellin multiplier is $(\kappa-z)^{-1}$.
This multiplier has the $N=1$ decay required above,
so its compression is Hilbert--Schmidt.
Its square is trace class by (MA-HY.15); higher powers
are trace class by multiplication with a bounded
resolvent.  Bounded similarity gives (MA-HY.17).
\end{proof}

\subsection{The actual ambient lift and its partial geometric limit}

\paragraph{MA-HY-003 (the geometric-order cutoff trace at $s=0$).}
Take $a=1/2$, $u=-\alpha$, $v=\alpha$, with any $\alpha>0$,
and use the actual weight $w=e^{2\alpha|x|}$.
Let $P_Z$ and $P_G=1-P_Z$ be the orthogonal projections in
$H_{\max}$ onto $J^\perp$ and $J$, respectively.
For $g\in C_c^\infty(\mathbb R)$,
$C_gP_Z$ is trace class and
\[
 \operatorname{Tr}_{H_{\max}}(C_gP_Z)
  =\sum_{|\operatorname{Re}\rho-1/2|<\alpha}
                    m_\rho\,\mathcal Mg(\rho-1/2).
 \tag{MA-HY.23}
\]
Let $\chi_R\in C_c^\infty(\mathbb R)$ have uniformly bounded
suprema and converge to one almost everywhere.
Then
\[
       \|(M_{\chi_R}-1)C_gP_Z\|_1\longrightarrow0.
 \tag{MA-HY.24}
\]
If $g(0)=0$, every $M_{\chi_R}C_gP_G$ is trace class and
\[
 \lim_R\operatorname{Tr}_{H_{\max}}(M_{\chi_R}C_gP_G)
   =-\sum_{|\operatorname{Re}\rho-1/2|<\alpha}
                    m_\rho\,\mathcal Mg(\rho-1/2).
 \tag{MA-HY.25}
\]
This is a limit of traces of the actual operators in the
specified order, without any zero-location premise.

\begin{proof}
The operator part of MA-LIFT-002 uses arithmetic only to
obtain a trace-class power of the actual quotient resolvent.
MA-HY-002 provides precisely this premise here, already
with power two.  To make the replacement explicit, let
$L=w^{-1}\partial_x(w\,\cdot)$ generate the actual adjoint
translation group, put $R=(\kappa-L)^{-1}$ for
$\kappa>\alpha+1$, and extend its restriction to
$J^\perp$ by $P_Z$.  Equation (MA-HY.17) gives
$R^2P_Z\in\mathcal S_1(H_{\max})$.
It maps boundedly into the graph domain of $L^2$.
The integration by parts on the two half-lines in
MA-LIFT.12--MA-LIFT.16, with $k=2$ and $s=0$, is the
exact identity
\[
        C_gP_Z=B_{g,2}R^2P_Z+
              \sum_{m=0}^{1}
                \phi_{g,2,m}\otimes(\ell_mR^2P_Z).
 \tag{MA-HY.26}
\]
Here $B_{g,2}$ is bounded, $\phi_{g,2,m}$ is compact
smooth, and the graph-domain trace functionals
$\ell_m(u)=\partial^m(wu)(0)$ are bounded after
$R^2P_Z$.  The finite boundary contribution at the
cusp of $w$ is retained.  Thus (MA-HY.26) proves trace
class of the actual lift.  Trace cyclicity with $P_Z$
identifies its trace with the actual quotient trace,
which MA-HY-002 computes as (MA-HY.23).

The proof of MA-LIFT-003 is purely an operator calculation
once this trace-class lift is available.
Uniform boundedness and strong convergence of the
$M_{\chi_R}$ imply (MA-HY.24) by finite-rank approximation.
The smooth compact kernel of $M_{\chi_R}C_g$ has trace
$g(0)\int\chi_R$, also for the weighted space by bounded
local conjugation and cyclicity.  Subtracting
$M_{\chi_R}C_gP_Z$ proves trace class of
$M_{\chi_R}C_gP_G$ and, for $g(0)=0$, (MA-HY.25).
No convergence of $M_{\chi_R}C_gP_G$ itself in trace norm
is asserted.
\end{proof}

\paragraph{Scope and the remaining boundary problem.}
This is an unconditional partial-strip theorem.  It does
not identify (MA-HY.25) with the full Meyer geometric trace:
the excluded zeros, including any on the two edges, are
absent from this actual $s=0$ quotient.  The difference from
the full spectral trace is exactly their absolutely
convergent Mellin sum.  MA-HM-003 prevents a full return
intertwiner for arbitrary tests if a strictly exterior
zero exists.  Under the single-character RH, every positive
centered strip contains the full divisor and the earlier
actual Meyer transport applies.

The condition $s=0$ is essential to the proof given here.
For positive polynomial weights the Mellin boundary space
has Sobolev regularity; at the thresholds established in
MA-TR-103, boundary jets are continuous and can survive in
the actual quotient.  Multiplication by the disk
half-derivative is then not the plain Hardy-space model
used above.  No assertion about the full Hardy--Sobolev
closed ideal or its trace is obtained by replacing $s=0$
by $s\ge3$ in these statements.

\paragraph{Placement in the requested geometric programme.}
The object completed here is the actual spectral jet quotient
of MT-M18/MT-M24.  The inner function $B$ encodes the divisor
that remains visible to this Hilbert norm, with all its interior
multiplicities; its Hardy-space model supplies an analytic
realization of that quotient and its trace.  This is compatible
with the nuclear trace and determinant perspective attached in
BR-TR-01/02 and F-SM-11/12 to cote~1, scans~079--080 and
111--112.  Those archival locators are inherited framework
placements, not a new independent rereading of the scans in
this section.  The Hardy model supplies neither a return to
discarded zero blocks nor an order or positivity on a
Grothendieck group.  In particular the inner factor is not an
assertion that the entire divisor is critical-line supported,
and the quotient construction does not remove the
zero-location premise from the full Meyer transport.

\paragraph{Primary literature and inherited interfaces.}
Beurling, A. (1949).
\emph{On two problems concerning linear transformations in
Hilbert space}. Acta Mathematica, 81, 239--255.
doi:10.1007/BF02395019.
The invoked locators are the Factorization Lemma in \S7,
pp.~245--246, and Theorem~I with proof in \S9,
pp.~247--249.
Bakan, A., \& Kaijser, S. (2007).
\emph{Hardy spaces for the strip}.
Journal of Mathematical Analysis and Applications, 333(1),
347--364. doi:10.1016/j.jmaa.2006.10.088.
The invoked strip-space locators are equations~(2.1)--(2.3),
Theorems~2.1--2.2 and Corollary~2.1, pp.~349--350;
the conformal/Poisson discussion is \S3 and the proofs of
the two isomorphism theorems are \S4, pp.~352--355.
The arithmetic premises are MT-M18 (M18.6), (M18.8) and
its unconditional local zero count, together with MT-M24's
identification of the actual theta ideal.
No unpublished outside result or zero-location input is
used in MA-HY-001--003.

\section{The actual partial-strip quotient at every polynomial weight}
\label{sec:weighted-partial-strip}

The arithmetic block synthesis of MT-M18 gives a direct extension
of the partial-strip result to the draft's polynomial weights.
Its block restrictions need not extend boundedly to the Hilbert
space. Instead, convolution by the smooth block idempotents does
extend, and supplies the finite-rank projections that the argument
requires. This distinction also retains the boundary jets at their
exact continuity thresholds.

Fix one primitive unitary Hecke function and the entire
pole-removed completion $\Xi$ of MT-M18. Let $\psi\in\mathcal E$
be its inverse Mellin transform and
$\mathcal I=\overline{\psi*\mathcal E}^{\,\mathcal E}$.
Fix $a\in\mathbb R$, $u<v$, and $s\ge0$, and put
\[
 \begin{aligned}
 w_{u,v,s}(x)&=\max(e^{2ux},e^{2vx})(1+x^2)^{s/2},\\
 H&=L^2(\mathbb R,w_{u,v,s}(x)\,dx),\qquad
 I_a=T_a\mathcal I,\qquad T_af(x)=e^{ax}f(x),\\
 J&=\overline{I_a}^{\,H},\qquad Q=H/J,\qquad
 q:H\longrightarrow Q,\qquad U_tf(x)=f(x-t).
 \end{aligned}
 \tag{MA-WP.1}
\]
All constants concern this fixed function and these fixed parameters.
There is no hypothesis that the strip contains every zero, and no
uniform separation of interior zeros from its edges.

For $\lambda=\rho-a$, where $\Xi(\rho)=0$ has multiplicity
$m_\rho$, define
\[
 b_s=\#\{j\in\mathbb Z_{\ge0}:2j<s-1\}
     =\max\left(0,\left\lceil\frac{s-1}{2}\right\rceil\right),
 \tag{MA-WP.2}
\]
and the retained multiplicity
\[
 r_\rho=
 \begin{cases}
 m_\rho,&u<\operatorname{Re}(\rho-a)<v,\\
 \min(m_\rho,b_s),&\operatorname{Re}(\rho-a)\in\{u,v\},\\
 0,&\operatorname{Re}(\rho-a)\notin[u,v].
 \end{cases}
 \tag{MA-WP.3}
\]
Thus $s=3$ retains a boundary value, but not a boundary first
derivative. At $s=4$ it retains derivatives of orders zero and one,
up to the actual zero multiplicity.

\paragraph{MA-WP-001 (the literal weighted closure and its boundary blocks).}
The actual closed ideal in (MA-WP.1) is
\[
 J=\left\{f\in H:
    (\mathcal Mf)^{(j)}(\rho-a)=0
       \text{ for every }\rho\text{ and }0\le j<r_\rho\right\}.
 \tag{MA-WP.4}
\]
Every displayed functional is continuous on $H$; at an edge it
is the continuous Mellin boundary derivative. The actual quotient
has finite-dimensional, mutually annihilating bounded block
projections with densely spanning ranges. In a block its algebraic
module is exactly
\[
             \prod_{\rho\ \text{in the block},\ r_\rho>0}
                         \mathbb C[z]/(z^{r_\rho}),
 \tag{MA-WP.5}
\]
with translation acting by multiplication by
$e^{(\rho-a+z)t}$ on the corresponding factor.

Equivalently, replace the maximum in (MA-WP.1) by
$e^{2ux}+e^{2vx}$. Mellin transform then identifies the space
with the analytic Hardy--Sobolev space on $u<\operatorname{Re}z<v$
whose two boundary functions have Sobolev order $s/2$.
In that literal boundary norm, the closed principal ideal generated
by $\Xi(a+z)$ is precisely the simultaneous vanishing ideal in
(MA-WP.4), with all the stated boundary truncations.

\begin{proof}
We give the block construction first; it will also prove the
operator estimates needed for the trace.

\emph{Bounded convolution and the ideal action.}
The weight gives the explicit translation bound
\[
 \|U_t\|_{H\to H}
 \le 2^{s/4}(1+t^2)^{s/4}\max(e^{ut},e^{vt}).
 \tag{MA-WP.6}
\]
Indeed $\max(e^{2u(x+t)},e^{2v(x+t)})$ is at most
$\max(e^{2ut},e^{2vt})\max(e^{2ux},e^{2vx})$, and
$1+(x+t)^2\le2(1+x^2)(1+t^2)$.
The translations are strongly continuous, first on compact smooth
functions and then by density and the local uniform bound.
Consequently, for $h\in\mathcal E$,
\[
 \|C_h\|_{H\to H}
 \le 2^{s/4}\int |h(t)|(1+t^2)^{s/4}
                         \max(e^{ut},e^{vt})\,dt
 \le C_N p_{N,0}(h)
 \tag{MA-WP.7}
\]
for a fixed sufficiently large integer $N$.

The space $I_a$ is a closed convolution ideal of $\mathcal E$,
since $T_a$ is a continuous convolution-algebra automorphism.
If $h\in I_a$, then $C_h(\mathcal E)\subset I_a$.
The inclusion $\mathcal E\subset H$ is continuous and dense,
so (MA-WP.7) implies $C_h(H)\subset J$. Hence every
$h\in I_a$ acts as zero on the actual quotient $Q$.
All convolutions by $\mathcal E$ preserve $J$ and descend to $Q$.

\emph{The precise MT-M18 data used.}
Retain the finite block algebras $(J_n,\|\cdot\|_n)$, block
restrictions $R_n$, and entire interpolation lifts $\mathcal T_n$
of MT-M18, (M18.1)--(M18.4) and (M18.10)--(M18.13).
Their defining facts are
\[
 \begin{gathered}
 d_n=\dim J_n=O(\log(2+|n|)),\qquad \|1\|_n\le1,\\
 R_m\mathcal T_n=\delta_{mn}\operatorname{id},\qquad
 p_{N,j}(\mathcal T_n b)
       \le C_{N,j}(2+|n|)^{D_{N,j}}\|b\|_n,\\
 f-\sum_n\mathcal T_n(R_nf)\in\mathcal I
             \quad(f\in\mathcal E),
 \end{gathered}
 \tag{MA-WP.8}
\]
where the series converges absolutely in every seminorm of
$\mathcal E$. Empty blocks have zero lift. The norm is the
holomorphic quotient norm of MT-M18, not the norm of individual
Hermite coordinates. Thus the displayed constants have no
inverse zero-separation or edge-distance factor.

Center the lifts and restrictions by
\[
 S_n=T_a\mathcal T_n,
 \qquad r_ng=R_n(T_{-a}g),
 \qquad e_n=S_n(1),
 \qquad h_{n,g}=S_n(r_ng).
 \tag{MA-WP.9}
\]
The symbol $r_n$ is used only on $\mathcal E$; no bounded
extension $H\to J_n$ is assumed. For every $g\in\mathcal E$,
\[
 g-\sum_n h_{n,g}\in I_a,
 \qquad \sum_n h_{n,g}\text{ converges in }\mathcal E.
 \tag{MA-WP.10}
\]

\emph{Actual finite-rank projections.}
Let $P_n=C_{e_n}^{Q}$. The full smooth jet description of
MT-M18 implies
\[
 e_n*e_n-e_n\in I_a,\qquad
 e_n*e_m\in I_a\ (n\ne m).
\]
The preceding ideal-action argument therefore gives
$P_n^2=P_n$ and $P_nP_m=0$ for $n\ne m$ on $Q$.
For $f\in\mathcal E$, its block jets give the exact identity
\[
                 P_nq(f)=q(S_n(r_nf)).
 \tag{MA-WP.11}
\]
Thus $P_nq(\mathcal E)$ lies in the finite-dimensional space
$V_n=q(S_nJ_n)$, of dimension at most $d_n$. This space is
closed. Density of $q(\mathcal E)$ and boundedness of $P_n$
show that $P_n(Q)\subset V_n$. Conversely $P_n$ is the
identity on $V_n$, again by its smooth jets. Hence
$P_n(Q)=V_n$ and $P_n$ is finite rank. Equations
(MA-WP.7)--(MA-WP.9) also give
\[
                    \|P_n\|\le C(2+|n|)^D.
 \tag{MA-WP.12}
\]
The block synthesis in (MA-WP.8), followed by $q$, shows
that the algebraic span of the $V_n$ is dense in $Q$.

For later use let $A_{n,g}=C_{h_{n,g}}^Q$.
On the smooth quotient $h_{n,g}$ is supported in block $n$;
therefore the same bounded extension argument gives
\[
 A_{n,g}=P_nA_{n,g}=A_{n,g}P_n,
 \qquad \operatorname{rank} A_{n,g}\le d_n.
 \tag{MA-WP.13}
\]
For the single $N$ in (MA-WP.7), there is an exponent $D$
such that
\[
 \|A_{n,g}\|_1
 \le d_n\|A_{n,g}\|
 \le C d_n(2+|n|)^D\|r_ng\|_n.
 \tag{MA-WP.14}
\]
The last norm decreases faster than every power for
$g\in\mathcal E$, by MT-M18's restriction estimate.
Thus (MA-WP.10) gives the actual operator identity
\[
                  C_g^Q=\sum_n A_{n,g},
 \qquad\text{with convergence in trace norm.}
 \tag{MA-WP.15}
\]
In particular $g\mapsto C_g^Q$ is continuous
$\mathcal E\to\mathcal S_1(Q)$.

This identity removes a possible completion ambiguity.
If $P_n v=0$ for every $n$, (MA-WP.13)--(MA-WP.15)
give $C_g^Qv=0$ for every $g\in\mathcal E$.
Choose a compact smooth approximate identity of integral
one. Strong continuity in (MA-WP.6) gives $C_g^Qv\to v$,
so $v=0$. Therefore
\[
                         \bigcap_n\ker P_n=\{0\}.
 \tag{MA-WP.16}
\]
Dense block ranges alone would not justify this last assertion;
the approximate-identity argument is required.

\emph{Which finite jet combinations are continuous.}
On compact smooth functions write
\[
        \ell_{\lambda,j}(f)
             =\int x^j e^{\lambda x}f(x)\,dx.
\]
Weighted $L^2$ duality gives
\[
 \ell_{\lambda,j}\in H^*
 \quad\Longleftrightarrow\quad
 \int\frac{|x|^{2j}e^{2\operatorname{Re}\lambda x}}
                    {w_{u,v,s}(x)}\,dx<\infty.
 \tag{MA-WP.17}
\]
This holds for every $j$ in the open strip, for
$2j<s-1$ at an edge, and for no $j$ outside the closed
strip.

There is also no cancellation that allows an excluded
higher jet to survive in a finite linear combination.
To see this explicitly, take a finite sum of such jets
that is continuous on $H$. On its finite-dimensional
formal jet space, precomposition with $U_t$ has, at
$\lambda$, the matrix
\[
 \ell_{\lambda,j}(U_tf)
       =e^{\lambda t}\sum_{i=0}^{j}
                       \binom{j}{i}t^{j-i}\ell_{\lambda,i}(f).
 \tag{MA-WP.18}
\]
Choose one real $t$ such that the finitely many numbers
$e^{\lambda t}$ for distinct $\lambda$ are distinct.
A finite Hermite interpolation polynomial in this bounded
dual translation isolates each generalized eigenspace.
Each resulting one-$\lambda$ component is therefore
continuous on $H$. Its integral kernel is
$p(x)e^{\lambda x}$ for a polynomial $p$.
If $p\ne0$ has degree $j$, then $|p(x)|$ is bounded above
and below by positive constants times $|x|^j$ on each
sufficiently distant real tail. Its dual integrability
is exactly (MA-WP.17) for that highest degree.
Hence every nonzero coefficient in a continuous finite
jet combination is among the permitted jets in
(MA-WP.3). No uniform bound for these finite interpolation
polynomials is needed.

\emph{The actual finite blocks and the full closed kernel.}
Let $L\in V_n^*$. The functional $L P_nq$ is continuous on
$H$. Restricted to $\mathcal E$, (MA-WP.11) expresses it
as a finite combination of the full jets in block $n$.
By the preceding argument it is a combination of precisely
the permitted jets. Conversely every permitted jet in that
block is continuous on $H$, vanishes on $J$, and satisfies
$\ell_{\lambda,j}P_n=\ell_{\lambda,j}$ on the quotient.
The finite Hermite interpolation in $J_n$ shows that these
functionals are linearly independent. Thus $V_n^*$ is
exactly their span. The kernel of $J_n\to V_n$ is
therefore the ideal that discards the other coordinates
and the higher boundary powers. This proves (MA-WP.5)
and its translation action.

Every element of $J$ annihilates the permitted jets.
Conversely, if $f\in H$ annihilates them, the description
of $V_n^*$ gives $P_nq(f)=0$ for every $n$.
Equation (MA-WP.16) gives $q(f)=0$, proving (MA-WP.4).

\emph{The literal Hardy--Sobolev interpretation.}
Let $H_\Sigma$ use the weight
$(e^{2ux}+e^{2vx})(1+x^2)^{s/2}$.
The two norms differ by factors between one and $\sqrt2$,
so they have the same closed ideal $J$ and boundedly
similar quotient convolutions. Put $r=s/2$.
The Mellin boundary functions
$F_\sigma(t)=\mathcal Mf(\sigma+it)$, $\sigma=u,v$,
satisfy, by Plancherel,
\[
 \|f\|_{H_\Sigma}^2
  =\frac1{2\pi}\sum_{\sigma\in\{u,v\}}
       \|(1-\partial_t^2)^{r/2}F_\sigma\|_{L^2(dt)}^2.
 \tag{MA-WP.19}
\]
Conversely the $s=0$ strip transform theorem
(MA-HY.3)--(MA-HY.4) gives a
unique underlying $f$ for an analytic Hardy function;
Sobolev regularity of its two boundary values, by the
same Fourier identity, is exactly $f\in H_\Sigma$.
This proves the asserted Hardy--Sobolev identification.
The continuous edge jets are actual derivatives of the
boundary values, also given by the absolutely convergent
integrals in (MA-WP.17).

Finally $\overline{C_{T_a\psi}H_\Sigma}=J$.
Indeed $\mathcal E$ is dense in $H_\Sigma$, convolution
is bounded, and $I_a$ is its $\mathcal E$-closure on the
dense smooth subspace. Under Mellin transform this is
the closure of multiplication by $\Xi(a+z)$.
Thus (MA-WP.4) is the stated closed principal ideal in
the actual Sobolev norm, not an independently imposed
spectral kernel.
\end{proof}

\paragraph{MA-WP-002 (partial trace with exactly the allowed edge multiplicities).}
For every $g\in\mathcal E$, including every compact smooth
test, the actual quotient convolution is trace class and
\[
        \operatorname{Tr}_Q C_g^Q
          =\sum_{\Xi(\rho)=0}r_\rho\,
                                  \mathcal Mg(\rho-a).
 \tag{MA-WP.20}
\]
The sum is absolutely convergent and the map into trace
class is continuous in the topology of $\mathcal E$.
For the actual quotient translation generator $A_Q$,
a sufficiently large real $\kappa$ and an integer $k$
satisfy
\[
                      (\kappa-A_Q)^{-k}\in\mathcal S_1(Q).
 \tag{MA-WP.21}
\]
One may take any $k>D+2$ in terms of the fixed exponent
in (MA-WP.14), and
$\kappa>\max(v,2-a)+1$.

\begin{proof}
Trace class and continuity were proved in (MA-WP.14)--(MA-WP.15)
before any trace was assigned. On $V_n$, the operator
$A_{n,g}$ is multiplication by the class of $\mathcal Mg$.
In the factor $\mathbb C[z]/(z^{r_\rho})$, its diagonal
entries are all $\mathcal Mg(\rho-a)$, so its trace is
$r_\rho\mathcal Mg(\rho-a)$. It vanishes on $\ker P_n$.
Trace-norm convergence in (MA-WP.15) therefore proves
(MA-WP.20). The unconditional local zero count and rapid
vertical decay also prove the displayed scalar series
absolutely convergent directly.

For the resolvent put $F_k(z)=(\kappa-z)^{-k}$.
It is analytic on the fixed strip containing all the
centered MT-M18 block rectangles. Write $[F_k]_n$ for
the class of $w\mapsto F_k(w-a)$ in the original
block algebra $J_n$. These classes satisfy
\[
             \|[F_k]_n\|_n\le C_k(2+|n|)^{-k}.
\]
Apply the centered interpolation lifts to these classes.
The resulting quotient convolutions have rank at most
$d_n$ and, by (MA-WP.7)--(MA-WP.8), trace norms bounded
by $C_kd_n(2+|n|)^{D-k}$. Their sum is trace class for
$k>D+2$.

On the other hand (MA-WP.6) gives the actual resolvent
power as the operator-norm convergent integral
\[
 (\kappa-A_Q)^{-k}
   =\frac1{(k-1)!}\int_0^\infty
              e^{-\kappa t}t^{k-1}U_t^Q\,dt.
 \tag{MA-WP.22}
\]
On each $V_n$ the finite-dimensional translation formula
in (MA-WP.5) identifies this integral with multiplication
by $F_k$. It therefore agrees there with the trace-class
sum just constructed. The block ranges span densely and
both operators are bounded, so they agree on $Q$.
This proves (MA-WP.21).
\end{proof}

\paragraph{MA-WP-003 (the actual ambient lift and cutoff limit for $s\ge3$).}
In particular take $a=1/2$, $u=-\alpha$, $v=\alpha$,
with arbitrary $\alpha>0$ and $s\ge0$, and let
$P_Z$ be the orthogonal projection onto $J^\perp$ in
the exact space $H_{\alpha,s}$ of the draft.
For $g\in C_c^\infty(\mathbb R)$,
\[
 C_gP_Z\in\mathcal S_1(H_{\alpha,s}),\qquad
 \operatorname{Tr}(C_gP_Z)
       =\sum_\rho r_\rho\,\mathcal Mg(\rho-1/2).
 \tag{MA-WP.23}
\]
If $\chi_R$ are compact smooth cutoffs, uniformly bounded
in supremum norm and tending to one almost everywhere,
then
\[
           \|(M_{\chi_R}-1)C_gP_Z\|_1\longrightarrow0.
 \tag{MA-WP.24}
\]
For $g(0)=0$, the actual geometric-order operators satisfy
\[
 \lim_R\operatorname{Tr}(M_{\chi_R}C_gP_G)
       =-\sum_\rho r_\rho\,\mathcal Mg(\rho-1/2),
 \qquad P_G=1-P_Z.
 \tag{MA-WP.25}
\]
These assertions include the requested $s\ge3$ weights.

\begin{proof}
The operator proof of MA-LIFT-002 needs a trace-class
power of the actual quotient resolvent and the explicit
weight. MA-WP-002 supplies the first premise here without
an interior-margin assumption. The weight is unchanged,
so its adjoint generator is still
$L=w^{-1}\partial_x(w\,\cdot)$, and the integration by
parts on the two half-lines, including the finite cusp
jets in (MA-LIFT.12)--(MA-LIFT.16), gives the same bounded
operator plus finite-rank factorization through
$(\kappa-L)^{-k}P_Z$. That operator is trace class by
(MA-WP.21) and Hilbert adjunction. Thus $C_gP_Z$ is trace
class in the specified order. Trace cyclicity with $P_Z$
and (MA-WP.20) give (MA-WP.23).

Uniform boundedness and strong convergence of the
cutoff multipliers imply (MA-WP.24) by finite-rank
approximation in trace norm. The compactly localized
convolution has its usual diagonal trace
$\operatorname{Tr}(M_{\chi_R}C_g)=g(0)\int\chi_R$;
MA-LIFT-003 proves this for the actual weighted space.
Subtracting $M_{\chi_R}C_gP_Z$ gives (MA-WP.25).
This is a trace limit, not a trace-norm convergence
claim for $M_{\chi_R}C_gP_G$ itself.
\end{proof}

\paragraph{Placement in the requested Grothendieck framework.}
The object in (MA-WP.5) is an actual finite-length local
algebra $\mathbb C[z]/(z^{r_\rho})$ in the completed
spectral quotient. The polynomial weight changes its
visible jet length at an edge; it does not change the
arithmetic multiplicity $m_\rho$ of the original entire
function. In the Grothendieck group of the abelian category
of finite-length $\mathbb C[z]$-modules its class has
length $r_\rho$ at its support point, equal to the
class of $r_\rho$ copies of its residue field. Likewise
multiplication by a test transform $G$ has trace
$r_\rho G(\rho-a)$ and determinant $G(\rho-a)^{r_\rho}$.
The determinant here is formed from exterior powers over
$\mathbb C$ of the underlying finite-dimensional vector
space with its endomorphism.
These scalar invariants retain the multiplicity but
do not distinguish the nilpotent Jordan structure
from its semisimplification. The actual truncated
algebra is retained in the operator construction.
This is the nuclear-trace and determinant placement
of BR-TR-01/02 and F-SM-11/12, with inherited cote~1
scans~079--080 and 111--112. Those scans are framework
locators inherited from the predecessor work, not
newly reread source pages. The argument supplies
neither an arithmetic positivity theorem nor an
$\ell$-adic weight realization of this analytic
quotient.

\paragraph{Exact scope and source boundary.}
The new ingredients are the use of bounded convolution
by the MT-M18 block idempotents, the approximate-identity
argument excluding simultaneous-kernel vectors, and the
finite-block boundary classification. The polynomial
lift bound is exactly (M18.13); its proof uses the
Gaussian interpolation kernel (M18.10) and the matching
Stirling and selected-height bounds (M18.6)--(M18.9).
That complete argument was reread for this extension.
No bounded Hilbert block restriction, individual zero
separation estimate, or independent Hardy--Sobolev
cyclicity theorem is imported.

The conclusion is an unconditional computation of the
actual \emph{partial} quotient and its geometric-order
cutoff trace for the fixed Hecke function. It includes
only the multiplicities in (MA-WP.3). It does not turn
that partial trace into the full Meyer trace if omitted
zero jets contribute. MA-HM-003's obstruction to a full
return intertwiner remains, and the full all-zero
transport retains its stated zero-location premise.

\section{A regularized one-sided Euler commutator inside the critical strip}
\label{sec:regularized-strip}

This section constructs actual closed multiplier operators and an
actual trace-class commutator on almost every one-sided weight inside
the critical strip.  The test and the derivative of the cutoff have
compact support and finite Gevrey order.  The trace is obtained from
the operator kernel before a contour shift is made.  The shift then
exhibits the pole and the zeros crossed by this regularization.
No bounded M\"obius inverse on a matching narrow Hilbert space is
asserted, and the result is not stated for every width or for every
compact smooth test.

For a finite set $S$ of rational primes write
\[
 E_S(s)=\prod_{p\in S}(1-p^{-s}),\qquad
 \zeta_S(s)=E_S(s)\zeta(s).
 \tag{MA-RS.1}
\]
The subscript removes the primes in $S$ from the Euler product.
Let $\rho=\beta+i\gamma$ run through the distinct nontrivial
zeros of $\zeta$, with multiplicities $m_\rho$.

\paragraph{MA-RS-001 (generic vertical lines and the inverse multiplier).}
There is a full-measure subset
$\mathscr G\subset(1/2,1)$ such that for every
$\sigma\in\mathscr G$ there is $c_\sigma>0$ with
\[
       |\sigma-\beta|
          \ge c_\sigma(2+|\gamma|)^{-2}
                    \quad\hbox{for every nontrivial zero }\rho.
 \tag{MA-RS.2}
\]
For this $\sigma$ and every integer $r\ge0$ there is
$C_{\sigma,S,r}>0$ such that
\[
 \left|\frac{d^r}{dt^r}\frac1{\zeta_S(\sigma+it)}\right|
       \le \exp\!\left(C_{\sigma,S,r}
                            \log^2(2+|t|)\right)
                    \qquad(t\in\mathbb R).
 \tag{MA-RS.3}
\]
The functions $\zeta_S(\sigma+it)$ and all of their
$t$ derivatives have polynomial growth.  These assertions
are unconditional.

\begin{proof}
The unconditional local zero count used in NT-Z10 and
MT-M18 is $O(\log(2+|T|))$ in each fixed-height window,
with multiplicities.  Consequently
\[
                  \sum_\rho (2+|\gamma|)^{-2}<\infty.
 \tag{MA-RS.4}
\]
The intervals centered at $\beta$ of radius
$(2+|\gamma|)^{-2}$ have summable lengths.
The first Borel--Cantelli lemma says that almost every
$\sigma$ belongs to only finitely many of them.
Remove also the countable set of exact real parts $\beta$.
For each remaining $\sigma$, the finitely many exceptional
distances have a positive minimum after division by their
respective radii.  Decreasing that minimum to at most one
proves (MA-RS.2).

We give the needed minimum-modulus argument, because a
zero-free vertical line alone would not imply (MA-RS.3).
At large $|t|$, center a disk at $2+it$ and choose a
zero-free boundary radius $R$ in a fixed sufficiently large
interval, for example $12<R<13$.
The disk contains $M=O(\log(2+|t|))$ zeros and no pole or
trivial zero of $\zeta$.  Divide them out with the finite
disk Blaschke product $B$, with multiplicities, and write
$\zeta=G B$.  On the boundary $|B|=1$.
The fixed-strip polynomial upper bound gives
$|G|\le C(2+|t|)^C$ inside the disk.  At its center,
$|G(2+it)|\ge|\zeta(2+it)|\ge\zeta(2)^{-1}$.
Harnack's inequality, applied to
$\log(C(2+|t|)^C)-\log|G|$, therefore gives
\[
          |G(w)|\ge c(2+|t|)^{-C}
 \tag{MA-RS.5}
\]
on the inner half disk.  This is the same finite-disk
argument used in the proof of MT-M18 (M18.8).

Take $w=\sigma+it$.  In the normalized disk coordinate
it lies in the inner quarter disk, so each Blaschke
factor corresponding to $\rho$ is at least
$|w-\rho|/(2R)$.  For every zero in the disk,
(MA-RS.2) and $|\gamma-t|<13$ give
\[
       |w-\rho|\ge c'_\sigma(2+|t|)^{-2}.
 \tag{MA-RS.6}
\]
Multiplying at most $C\log(2+|t|)$ such lower bounds and
then using (MA-RS.5) yields
\[
        |\zeta(\sigma+it)|^{-1}
              \le \exp(C_\sigma\log^2(2+|t|)).
 \tag{MA-RS.7}
\]
The same argument works uniformly when $w$ ranges over
a disk of radius $c''_\sigma(2+|t|)^{-2}$ about
$\sigma+it$, after halving the constant in (MA-RS.6).
That disk is zero free.  Cauchy's estimate for $1/\zeta$
on it proves (MA-RS.3) for $S=\varnothing$: the extra
factor $r!c_\sigma^{-r}(2+|t|)^{2r}$ is absorbed into
the right side.  A compact range of $t$ is covered by
enlarging the constants.

On a small closed strip about $\sigma>1/2$, all factors
$1-p^{-s}$, $p\in S$, are bounded away from zero, and
their inverses and derivatives are bounded.
This proves (MA-RS.3) for finite $S$.
The fixed-strip upper bound for $\zeta$, with Cauchy
estimates for its derivatives, proves the remaining
polynomial-growth assertion.  The pole at $1$ is
avoided by the fixed line and causes no issue on the
compact low-height portion.
\end{proof}

\subsection{The closed operators and their dense common core}

Fix $\sigma\in\mathscr G$ and put $a=\sigma-1/2$.
Use
\[
 \begin{aligned}
 \mathcal H_a&=L^2(\mathbb R,e^{2ax}\,dx),&
 W_au(x)&=e^{ax}u(x),\\
 \mathcal F_+v(t)&=(2\pi)^{-1/2}
                       \int_{\mathbb R}e^{itx}v(x)\,dx,&
 V_a&=\mathcal F_+W_a,\\
 \mathscr D_a&=V_a^{-1}C_c^\infty(\mathbb R),&
 A(t)&=\zeta_S(\sigma+it),\qquad j(t)=A(t)^{-1}.
 \end{aligned}
 \tag{MA-RS.8}
\]
Thus $V_a$ is unitary and $\mathscr D_a$ is the dense
core with compact support in the weighted Fourier variable.
Define
\[
 \begin{aligned}
 Z_{a,S}&=V_a^{-1}M_A V_a,&
 \mathcal D(M_A)&=\{v\in L^2:Av\in L^2\},\\
 J_{a,S}&=V_a^{-1}M_j V_a,&
 \mathcal D(M_j)&=\{v\in L^2:jv\in L^2\}.
 \end{aligned}
 \tag{MA-RS.9}
\]
These are closed, densely defined inverse operators, with
their maximal multiplier domains.  The equality
$M_A M_j=1$ holds on $\mathcal D(M_j)$, and
$M_j M_A=1$ on $\mathcal D(M_A)$.
Both have $C_c^\infty$ as a graph core in the Fourier
picture, by cutoff and smooth approximation in their
weighted graph norms.  Neither inverse identity says
that both operators are bounded on all of $\mathcal H_a$.

This continuation of $Z$ is also an explicit regularization
of its Euler sums.  For every positive integer $N$ and
$\operatorname{Re}s>0$, $s\ne1$, Abel summation gives
\[
 \zeta(s)=\sum_{n\le N}n^{-s}+\frac{N^{1-s}}{s-1}
           -s\int_N^\infty\{x\}x^{-s-1}\,dx .
 \tag{MA-RS.10}
\]
The remainder has modulus at most
$|s|N^{-\operatorname{Re}s}/\operatorname{Re}s$.
In particular the bounded multipliers
\[
        E_S(s)\left(\sum_{n\le N}n^{-s}
                              +\frac{N^{1-s}}{s-1}\right),
                  \qquad s=\sigma+it,
 \tag{MA-RS.11}
\]
converge to $A(t)$ on each compact frequency interval,
and their operators converge to $Z_{a,S}$ on
$\mathscr D_a$.  The closed graph closure gives (MA-RS.9).
Formula (MA-RS.10) follows first in $\operatorname{Re}s>1$
by writing $\lfloor x\rfloor=x-\{x\}$ in the tail
Stieltjes integral, and then continues to
$\operatorname{Re}s>0$ because the displayed integral
converges there.  Its exact published reference is
DLMF~25.2.8, equivalently Apostol (1976),
Theorem~12.21 and equation~(25), p.~269.
For $\operatorname{Re}s>1$ the counterterm tends to zero,
and this agrees with the operator-norm Euler construction
of MA-MEY-001.  Inside the critical strip the counterterm
is part of the definition, not an omitted convergence claim.

For completeness the inverse can be approximated on its
actual domain by the bounded multipliers
\[
                j_\varepsilon(t)
                  =\frac{\overline{A(t)}}{|A(t)|^2+\varepsilon}.
 \tag{MA-RS.12}
\]
They converge to $j$ in $L^2$ after multiplication by any
vector in $\mathcal D(M_j)$, since $|j_\varepsilon|\le|j|$.
This is a domain approximation; no uniform bound for their
operator norms as $\varepsilon\downarrow0$ is asserted.

\subsection{Gevrey decay and an actual trace-class kernel}

A compactly supported function has Gevrey order $q>1$
if its derivatives satisfy
$\|f^{(N)}\|_\infty\le C R^N(N!)^q$ for every $N\ge0$.
Choose $g$ of compact support and finite Gevrey order,
and a bounded smooth cutoff $\phi$ such that
\[
 \phi(-\infty)=1,\qquad \phi(+\infty)=0,\qquad
                 \phi'\text{ has compact support and finite Gevrey order}.
 \tag{MA-RS.13}
\]
Increase $q>1$ so that one common order works for $g$ and
$\phi'$.  Compact Gevrey bumps yield such cutoffs by
integration; a symmetric bump gives
$\phi(x)+\phi(-x)=1$ if that additional convention is desired.
Put
\[
 G(z)=\mathcal Mg(z)=\int e^{zx}g(x)\,dx,\qquad
 m(t)=G(a+it),\qquad
 p(\xi)=\int e^{i\xi x}\phi'(x)\,dx .
 \tag{MA-RS.14}
\]
The transforms here, unlike $\mathcal F_+$, are not
normalized by $(2\pi)^{-1/2}$.

\paragraph{MA-RS-002 (the regularized commutator is defined on its core).}
On $\mathscr D_a$ the composition
\[
           K_{a,S,\phi}(g)
              =C_gZ_{a,S}[J_{a,S},M_\phi]
 \tag{MA-RS.15}
\]
is well defined.  It extends uniquely to a trace-class
operator on $\mathcal H_a$.  In the unitary $V_a$ picture
its kernel with respect to $du$ is the Schwartz function
\[
 \mathcal K(t,u)=\frac{i}{2\pi}\,
          m(t)A(t)\frac{j(t)-j(u)}{t-u}\,p(t-u),
 \tag{MA-RS.16}
\]
with its smooth divided-difference value at $t=u$.
Its actual operator trace is
\[
       \operatorname{Tr}K_{a,S,\phi}(g)
        =-\frac1{2\pi}\int_{\mathbb R}
               G(a+it)\frac{\zeta_S'}{\zeta_S}(\sigma+it)\,dt .
 \tag{MA-RS.17}
\]
The integral is absolutely convergent.  The trace-norm
estimates are uniform when test supports and Gevrey
derivative constants are fixed.

\begin{proof}
We first give the decay estimate used both for domains
and for the kernel.  Repeated integration by parts gives
for a compact Gevrey function $f$
\[
       |\widehat f_+(t)|
          \le C R^N(N!)^q|t|^{-N}\qquad(N\ge0).
 \tag{MA-RS.18}
\]
Choose $N$ to be a fixed positive multiple of
$|t|^{1/q}$, rounded down, so that
$R N^q/|t|\le1/2$, and use $N!\le N^N$.
This proves $|\widehat f_+(t)|\le C'e^{-c|t|^{1/q}}$.
Multiplying $f$ by any fixed power of $x$ or by $e^{ax}$
preserves its Gevrey order on the fixed compact support,
by Leibniz's formula.  Hence for every fixed $r\ge0$,
\[
          |m^{(r)}(t)|+|p^{(r)}(t)|
                 \le C_r e^{-c_r|t|^{1/q}}.
 \tag{MA-RS.19}
\]
These estimates are uniform for the bounded Gevrey
families specified in the statement.

Convolution satisfies $V_a C_gV_a^{-1}=M_m$.
Let $v\in C_c^\infty$ be a weighted Fourier core vector,
and put $f=V_a^{-1}v\in\mathscr D_a$.
Then $jv\in C_c^\infty$.
For a bounded smooth $\phi$ with compactly supported
derivative, $\phi\,\mathcal F_+^{-1}v$ is Schwartz,
because $\mathcal F_+^{-1}v$ is Schwartz and all
derivatives of $\phi$ are bounded.
Outside a sufficiently large interval containing
$\operatorname{supp}v$, its Fourier transform is the
ordinary integral
\[
    \frac{i}{2\pi}\int
                 \frac{p(t-u)}{t-u}\,v(u)\,du .
 \tag{MA-RS.20}
\]
Indeed the distribution identity
$\widehat{\phi'}_+(\xi)=-i\xi\widehat\phi_+(\xi)$
gives $\widehat\phi_+(\xi)=i p(\xi)/\xi$
away from zero.  The same formula and its derivatives
show Gevrey exponential decay in that exterior region,
by (MA-RS.19).  The assertion applies equally with
$v$ replaced by $jv$.

It follows from (MA-RS.3) that $M_\phi$ applied to a
core vector lies in $\mathcal D(J_{a,S})$:
the tails from (MA-RS.20) dominate
$\exp(C\log^2(2+|t|))$.
Moreover $J_{a,S}M_\phi f$ lies in
$\mathcal D(Z_{a,S})$, and $Z_{a,S}J_{a,S}M_\phi f=M_\phi f$.
The second composition $Z_{a,S}M_\phi J_{a,S}f$
is also defined, because the Fourier transform after
$M_\phi J_{a,S}$ has those same tails and $A$
has polynomial growth.  This verifies all domains in
(MA-RS.15), including the application of $Z$ to the
commutator.  On this core it also proves
\[
       K_{a,S,\phi}(g)
          =C_g(M_\phi-Z_{a,S}M_\phi J_{a,S}).
 \tag{MA-RS.21}
\]

In Fourier variables multiplication by $\phi$ has
distribution kernel $(2\pi)^{-1}\widehat\phi_+(t-u)$.
The commutator multiplies this by $j(t)-j(u)$.
Since this difference vanishes at $t=u$, the possible
delta term in $\widehat\phi_+$ disappears, and the
identity $\xi\widehat\phi_+=i p$ gives precisely
(MA-RS.16).  This use of a distribution kernel is
initially on compact-frequency core vectors, where
all its products are defined; the resulting kernel
will now be seen to extend to all of $L^2$.

Write
\[
       \frac{j(t)-j(u)}{t-u}
           =\int_0^1 j'(u+\theta(t-u))\,d\theta .
 \tag{MA-RS.22}
\]
Every derivative of this function is bounded by
$\exp(C\log^2(2+|t|+|u|))$, with a derivative-dependent
constant, by (MA-RS.3).
All derivatives of $A$ have polynomial growth.
Combining (MA-RS.19) with the product rule gives,
for each fixed pair of derivative orders, an upper
bound of the form
\[
 C\exp\!\left(
     -c|t|^{1/q}-c|t-u|^{1/q}
             +C'\log^2(2+|t|+|u|)\right).
 \tag{MA-RS.23}
\]
The logarithmic term is absorbed by half the two
negative Gevrey powers; use
$|u|\le|t|+|t-u|$.
The same bound therefore dominates every polynomial
weight in $(t,u)$.  This proves
$\mathcal K\in\mathcal S(\mathbb R^2)$.

A Schwartz kernel defines a trace-class operator,
and its trace is its diagonal integral.  One explicit
justification is to expand in normalized Hermite
functions in the two variables.  Applying an arbitrary
power of the harmonic oscillator in each variable
bounds the coefficients by
$C_N(1+r)^{-N}(1+s)^{-N}$.
For $N>1$ their absolute sum is finite, so the
rank-one series converges in trace norm.
Higher powers give convergence of the kernel series
in Schwartz seminorms; integration of the diagonal
then gives the same trace.  This also proves the
claimed uniform trace-norm bounds from finitely many
Schwartz seminorms of (MA-RS.16).
Agreement on the dense core proves uniqueness of
the bounded extension.

Finally $p(0)=\int\phi'=-1$, and
\[
             j'(t)=-i\frac{\zeta_S'(\sigma+it)}
                                  {\zeta_S(\sigma+it)^2}.
 \tag{MA-RS.24}
\]
Thus the diagonal of (MA-RS.16) is
$-(2\pi)^{-1}G(a+it)(\zeta_S'/\zeta_S)(\sigma+it)$.
It is integrable by (MA-RS.3) and (MA-RS.19).
This proves (MA-RS.17), including its sign.
\end{proof}

\subsection{The residues carried by this operator regularization}

For compact $g$ put
\[
 \begin{aligned}
 P_{\overline S}^{+}(g)&=
       \sum_{\substack{p\notin S\\k\ge1}}
          (\log p)p^{-k/2}g(k\log p),\\
 B_{\overline S}(g)&=
          P_{\overline S}^{+}(g)+P_{\overline S}^{+}(g^\vee),
       &g^\vee(x)&=g(-x),\\
 P(g)&=G(1/2)+G(-1/2).
 \end{aligned}
 \tag{MA-RS.25}
\]
Both prime sums are finite.

\paragraph{MA-RS-003 (crossed zeros are an actual commutator correction).}
For the operators and tests just constructed,
\[
 \operatorname{Tr}K_{a,S,\phi}(g)
       =P_{\overline S}^{+}(g)-G(1/2)
                +\sum_{\operatorname{Re}\rho>\sigma}
                             m_\rho G(\rho-1/2).
 \tag{MA-RS.26}
\]
The zero sum is absolutely convergent.
Let $R:\mathcal H_a\to\mathcal H_{-a}$ be reflection,
$Ru(x)=u(-x)$.  The paired diagonal regularization
on $\mathcal H_a\oplus\mathcal H_{-a}$ is the actual
trace-class operator
\[
        \mathbb K_{a,S,\phi}(g)
             =K_{a,S,\phi}(g)\oplus
                       R K_{a,S,\phi}(g^\vee)R .
 \tag{MA-RS.27}
\]
It satisfies
\[
 \boxed{\ 
 \operatorname{Tr}\mathbb K_{a,S,\phi}(g)
    =B_{\overline S}(g)-P(g)
       +\sum_{|\operatorname{Re}\rho-1/2|>a}
                         m_\rho G(\rho-1/2).\ }
 \tag{MA-RS.28}
\]
There are no boundary zeros at this generic width.
Thus the regularized diagonal expression retains precisely
the contributions crossed in moving the defining Mellin
lines; its trace is not the unchanged wide-strip prime sum.

\begin{proof}
Rewrite (MA-RS.17) as
\[
    -\frac1{2\pi i}\int_{\sigma-i\infty}^{\sigma+i\infty}
           G(s-1/2)\frac{\zeta_S'}{\zeta_S}(s)\,ds .
 \tag{MA-RS.29}
\]
Fix $c>1$.  On the line $\operatorname{Re}s=c$ the
absolutely convergent Euler logarithmic derivative is
\[
          -\frac{\zeta_S'}{\zeta_S}(s)
           =\sum_{\substack{p\notin S\\k\ge1}}
                              (\log p)p^{-ks}.
 \tag{MA-RS.30}
\]
Interchange this absolute series with the integrable
Mellin test.  Fourier inversion on the vertical line
gives
\[
 -\frac1{2\pi i}\int_{c-i\infty}^{c+i\infty}
        G(s-1/2)\frac{\zeta_S'}{\zeta_S}(s)\,ds
                      =P_{\overline S}^{+}(g).
 \tag{MA-RS.31}
\]
This also follows from the already established actual
wide-strip commutator trace MA-MEY-001.

Use the selected upper and lower horizontal ordinates
of MT-M18 (M18.7)--(M18.8).
Its minimum-modulus proof, together with Stirling's
upper and lower bounds for $\Xi/\zeta$, makes
$|\zeta'/\zeta|$ polynomially bounded on the horizontal
segments $\sigma\le\operatorname{Re}s\le c$ at these
heights.  The factors $E_S$ and their inverses are
uniformly bounded there since $\sigma>1/2$.
The Gevrey estimate for $G(s-1/2)$ is uniform on this
fixed real interval.  Therefore both horizontal
integrals tend to zero.  The vertical integrals
converge absolutely by the same estimates as in
MA-RS-002.

The residue of the integrand in (MA-RS.29) before
its outside minus sign is $m_\rho G(\rho-1/2)$
at a nontrivial zero and $-G(1/2)$ at the pole $s=1$.
The extra zeros of $E_S$ have real part zero and
are not crossed.  The counterclockwise rectangle
has the $c$ edge oriented upwards and the $\sigma$
edge downwards.  Thus the upward $\sigma$ integral
equals the upward $c$ integral minus $2\pi i$
times the residue sum.  Multiplying by
$-(2\pi i)^{-1}$ gives (MA-RS.26).
The local zero count and uniform vertical decay
of $G$ prove absolute convergence of the resulting
zero sum, so the chosen horizontal subsequence
does not change its value.

Reflection is unitary between the two weighted
spaces, and $RC_gR=C_{g^\vee}$.
Thus (MA-RS.27) is the smoothed pair of diagonal
expressions $D$ and $RDR$, where
$D=Z_{a,S}[J_{a,S},M_\phi]$ is initially defined
on its core.  It is not a definition by its trace.
Apply (MA-RS.26) to $g$ and $g^\vee$ and add.
Their Mellin transforms satisfy
$\mathcal M(g^\vee)(z)=G(-z)$.
The nontrivial zeta divisor is invariant, with
multiplicity, under $\rho\mapsto1-\rho$.
Consequently the second right-hand zero sum becomes
the sum over $\operatorname{Re}\rho<1-\sigma$.
These two disjoint sets are exactly
$|\operatorname{Re}\rho-1/2|>a$.
The two pole terms sum to $-P(g)$.
Equation (MA-RS.2), and its reflected version from
the same zero symmetry, exclude zeros on either
boundary.  This proves (MA-RS.28).
\end{proof}

\paragraph{Compact cutoffs and the limits of the result.}
If $\phi$ in MA-RS-002 is replaced by a compactly supported
Gevrey function, the same kernel and domain proof applies,
but $p(0)=0$, so the commutator trace is zero.
Whenever (MA-RS.26) is nonzero, no family of these compact
cutoff commutators can converge to the one-sided operator
in trace norm.  This conclusion uses only the trace
inequality and does not assert strong convergence for
the unbounded inverse compositions.

The theorem covers $\sigma\in\mathscr G$, a full-measure
set, and compact tests and cutoff derivatives of finite
Gevrey order.  A zero-free line not satisfying the
quantitative estimate has not been covered.  Arbitrary
$C_c^\infty$ Fourier decay need not dominate the inverse
growth estimate proved here, so passage to all compact
smooth tests has not been justified by this argument.
At a line containing a zero, the multiplier inverse
and boundary-domain issue must be treated separately.

There is a distinction between extending the scalar functional
and extending the operator.  The disk argument in MA-RS-001
also makes $\zeta_S'/\zeta_S$ polynomially bounded on the
chosen line.  Indeed $G'/G$ is polynomially bounded by
Cauchy's estimate and (MA-RS.5), while the logarithmic
derivative of the finite Blaschke product is bounded by
$C(2+|t|)^2\log(2+|t|)$ using (MA-RS.6).
Thus the integral in (MA-RS.17), considered only as a
scalar distribution, extends continuously to every
$C_c^\infty$ test and has the same residue identity.
This does not improve the available bounds for the
two-variable divided difference of $1/\zeta_S$, and
does not prove a trace-class operator extension for
those arbitrary smooth tests.

The paired operator in (MA-RS.27) realizes the regularized
Euler diagonal expression.  It is not a proof that all
of Meyer's common-space projection maps, with their
off-diagonal terms, extend to this critical-strip
completion.  Such an identification would require its
own domain and intertwining proof.  The residue correction
in (MA-RS.28) is already forced for this concrete
regularization; deleting it would change the trace of
the actual constructed operator.

In the requested geometric interpretation, the crossed
divisor is the spectral boundary contribution to a change
of completion.  Its occurrence is compatible with the
actual partial quotients of MA-HY and MA-WP and with
MA-HM-003's obstruction to recovering discarded blocks.
The nuclear trace viewpoint has the inherited
BR-TR-01/02 and F-SM-11/12 placement at cote~1,
scans~079--080 and 111--112.  This section does not infer
a positive Grothendieck-group realization or a
zero-location theorem from that interpretation.

\paragraph{Exact source and proof scope.}
The arithmetic inputs are the unconditional fixed-strip
upper bound and local zero count of NT-Z10, together with
the finite-disk Harnack argument and selected-height
minimum modulus in MT-M18 (M18.5)--(M18.9).
MA-RS-001 supplies the new generic-line version in full.
The explicit regularization formula is DLMF~25.2.8,
\url{https://dlmf.nist.gov/25.2.E8}, sourced there to
Apostol, T. M. (1976), \emph{Introduction to analytic
number theory}, Theorem~12.21 and equation~(25), p.~269.
The Fourier sign, core-domain argument, Gevrey estimate,
Schwartz-kernel estimate and contour correction are all
proved in this section.  MA-MEY-001--003 supplies the
already checked wide-strip operator and source
diagonal conventions; its zero-location-free source is
Meyer (2005), Lemmas~5.3--5.6 and Theorem~5.8, author
pp.~39--44.  No A, M, or new outside preprint is used.

\section{Smooth tests for the regularized commutator under a fixed zero gap}
\label{sec:regularized-smooth}

\paragraph{MA-RS-004 (arbitrary compact smooth tests under a fixed horizontal gap).}
Retain the actual closed multiplier operators of MA-RS-001--003.
Suppose $1/2<\sigma<1$ and, for some $\delta>0$,
\[
       |\sigma-\operatorname{Re}\rho|\ge\delta
          \quad\text{for every nontrivial zero of }\zeta.
 \tag{MA-RS.32}
\]
Put $a=\sigma-1/2$. Then the conclusions of MA-RS-002--003
hold for every $g\in C_c^\infty(\mathbb R)$ and every bounded
smooth $\phi$ with
\[
       \phi(-\infty)=1,\qquad \phi(+\infty)=0,
                         \qquad\phi'\in C_c^\infty(\mathbb R).
 \tag{MA-RS.33}
\]
In particular the actual composition on the compact-frequency core
\[
             C_gZ_{a,S}[J_{a,S},M_\phi]
 \tag{MA-RS.34}
\]
extends uniquely to a trace-class operator on
$\mathcal H_a=L^2(e^{2ax}dx)$, with the exact kernel and trace
in (MA-RS.16)--(MA-RS.17). Its residue formula is still
(MA-RS.26), and its reflected pair still has (MA-RS.28).
The map $g\mapsto K_{a,S,\phi}(g)$ is continuous into trace
class on each fixed compact test-support space.

Under the stipulated Dirichlet input $A$, applied to rational
zeta, (MA-RS.32) holds for every $1/2<\sigma<1$.
For these lines the zero correction is absent, and the formulas
become
\[
 \begin{aligned}
 \operatorname{Tr}K_{a,S,\phi}(g)
     &=P_{\overline S}^{+}(g)-\mathcal Mg(1/2),\\
 \operatorname{Tr}\mathbb K_{a,S,\phi}(g)
     &=B_{\overline S}(g)
          -\mathcal Mg(1/2)-\mathcal Mg(-1/2).
 \end{aligned}
 \tag{MA-RS.35}
\]
Thus a smoothed commutator can be well defined and trace class
even though the bare M\"obius inverse has no bounded extension
between matching narrow weighted norms, as MA-MO-001 proves.

\begin{proof}
We first improve the inverse-multiplier estimate, keeping its
zero-location premise explicit. Use the finite-disk proof in
MA-RS-001, centered at $2+it$ with a zero-free boundary radius
$12<R<13$ for large $|t|$. Its factorization $\zeta=G B$
has $M=O(\log(2+|t|))$ zeros in the disk. Harnack's inequality
and the fixed-strip polynomial upper bound give
$|G(w)|\ge c(2+|t|)^{-C}$ on the inner half disk.

At $w=\sigma+it$ each finite Blaschke factor has modulus at
least $\delta/(2R)$, by (MA-RS.32). Replace this lower bound
by its minimum with $1/2$ if necessary. Its product over the
$M$ factors is at least a fixed negative power of $2+|t|$.
The same estimate is uniform on a fixed disk of radius
$\epsilon>0$ about $\sigma+it$, where $\epsilon$ is smaller
than $\delta/2$ and than the distance of the line from the
pole at $1$. This disk contains no zero or pole.
Cauchy's estimate for the reciprocal on a slightly smaller
disk consequently proves, for every integer $r\ge0$,
\[
 \left|\frac{d^r}{dt^r}\frac1{\zeta_S(\sigma+it)}\right|
       \le C_r(2+|t|)^{D_r}.
 \tag{MA-RS.36}
\]
The finitely many removed Euler factors and their inverses
are bounded with all derivatives on this fixed narrow strip.
The compact range of $t$ follows by continuity and the absence
of zeros on the line. All derivatives of
$A(t)=\zeta_S(\sigma+it)$ have polynomial growth as before.
This uses the actual gap (MA-RS.32), not merely a zero-free
line without quantitative control.

For compact smooth $g$ and $\phi'$, the functions
\[
     m(t)=\mathcal Mg(a+it),\qquad
     p(t)=\int e^{itx}\phi'(x)\,dx
\]
belong to $\mathcal S(\mathbb R)$, with continuous Schwartz
seminorms on each fixed compact support space. In the unitary
weighted Fourier picture, multiplication by $\phi$ preserves
the Schwartz space: $\phi$ is bounded and every positive-order
derivative is compactly supported and smooth. The multipliers
$A$ and $j=A^{-1}$ preserve Schwartz functions by their
polynomial derivative bounds. Thus, for every compact-frequency
core vector, both $J_{a,S}M_\phi$ and
$Z_{a,S}M_\phi J_{a,S}$ are defined, and their required
images lie in the domain of the next operator. This checks
the actual composition (MA-RS.34), rather than defining
it by a trace.

The exact distributional commutator calculation of MA-RS-002
then gives the kernel
\[
       \frac{i}{2\pi}m(t)A(t)
                   \frac{j(t)-j(u)}{t-u}\,p(t-u).
 \tag{MA-RS.37}
\]
Every mixed derivative of the divided difference has polynomial
growth in $(t,u)$: differentiate
$\int_0^1j'(u+\theta(t-u))\,d\theta$ and use
(MA-RS.36). The linear change of variables $(t,u)\mapsto
(t,t-u)$ shows that the product of any such polynomial with
any derivatives of $m(t)p(t-u)$ decreases faster than every
power in both variables. Hence (MA-RS.37) is a Schwartz
kernel. The Hermite-expansion proof in MA-RS-002 gives trace
class, the diagonal-integral formula, and continuity in trace
norm from finitely many Schwartz seminorms.

The contour argument of MA-RS-003 also applies unchanged.
At its selected horizontal heights the logarithmic derivative
has polynomial growth, while $\mathcal Mg$ decreases faster
than every power on the fixed strip. The horizontal integrals
therefore tend to zero, and the same residues give
(MA-RS.26) and (MA-RS.28). No Gevrey estimate is needed once
(MA-RS.36) is available.

Under $A$ all the nontrivial zeta zeros have real part $1/2$.
For each fixed $\sigma>1/2$ use $\delta=\sigma-1/2$.
There are no crossed nontrivial zeros on the right of the
line, nor outside its reflected pair. The pole at $1$
is still crossed and gives exactly (MA-RS.35).
\end{proof}

\paragraph{A concrete cutoff distinction, still present under $A$.}
Choose $g\ge0$, $g\ne0$, with compact support in
$(0,\log2)$. Both prime sums in (MA-RS.35) vanish,
regardless of the finite set $S$. Nevertheless
\[
 \operatorname{Tr}K_{a,S,\phi}(g)=-\mathcal Mg(1/2)<0,
 \qquad
 \operatorname{Tr}\mathbb K_{a,S,\phi}(g)
        =-\mathcal Mg(1/2)-\mathcal Mg(-1/2)<0.
 \tag{MA-RS.38}
\]
If the cutoff itself is compactly supported, the same
Schwartz-kernel proof applies but $p(0)=\int\phi'=0$,
so its commutator trace is zero. Such compact-cutoff
commutators cannot converge in trace norm to either
nonzero-trace operator in (MA-RS.38). This is an actual
operator distinction and a necessary pole correction;
it is not a refutation of a formula that already includes
the two polar terms.

\paragraph{Scope of this improvement.}
The unconditional generic-line result MA-RS-001--003 retains
its stated Gevrey scope. MA-RS-004 uses the additional fixed
horizontal gap and covers all compact smooth tests. For the
rational zeta sectors that gap follows from $A$; this use of
zero location is explicit. The construction concerns the actual
one-sided closed multipliers and their smoothed commutator.
It does not by itself identify them with the orthogonal
Paley--Wiener projections of Masullo's global weighted space.

\section{The regularized operator in Dirichlet and Hecke character sectors}
\label{sec:regularized-characters}

This extension concerns actual one-sided closed multipliers and their
smoothed commutators. It carries the construction through the
non-induced imaginary-quadratic example, with the conjugate character
in the reflected component. The actual weighted quotient and cutoff
trace in these sectors have already been supplied by MA-WP.

Fix a number field $K$, a primitive finite-order Hecke character
$\eta$, and a finite set $S$ of finite primes. Extend its ideal
coefficients by zero on ideals not coprime to its conductor, and set
\[
 \begin{aligned}
 L_\eta(s)&=\sum_{\mathfrak a}\eta(\mathfrak a)
                                       (N\mathfrak a)^{-s},\\
 E_{\eta,S}(s)&=\prod_{\mathfrak p\in S}
                       (1-\eta(\mathfrak p)(N\mathfrak p)^{-s}),\\
 A_\eta(s)&=E_{\eta,S}(s)L_\eta(s).
 \end{aligned}
 \tag{MA-RS.40}
\]
The first series is used in $\operatorname{Re}s>1$ and continued
thereafter by the classical Hecke function. Write
$\delta_\eta=1$ for the principal character and zero otherwise,
and $c_\eta=\operatorname{Res}_{s=1}L_\eta(s)$, taking
$c_\eta=0$ when the function is entire. Let $Z_\eta$ be the
nontrivial completed zero divisor, with multiplicities.

\paragraph{MA-RS-005 (actual character commutators and their crossed divisors).}
There is a full-measure subset of $\sigma\in(1/2,1)$ for which
the following assertions hold simultaneously for $\eta$ and
$\overline\eta$. Put $a=\sigma-1/2$ and use the same unitary
$V_a:\mathcal H_a\to L^2(\mathbb R)$ as in (MA-RS.8).
Define the maximal closed multipliers
\[
 Z_{a,\eta,S}=V_a^{-1}M_{A_\eta(\sigma+it)}V_a,
 \qquad
 J_{a,\eta,S}=V_a^{-1}M_{A_\eta(\sigma+it)^{-1}}V_a.
 \tag{MA-RS.41}
\]
For a compact Gevrey test $g$ and a one-sided cutoff $\phi$
with compact Gevrey derivative and endpoints $1,0$,
the actual core composition
\[
        K_{a,\eta,S,\phi}(g)
              =C_g Z_{a,\eta,S}[J_{a,\eta,S},M_\phi]
 \tag{MA-RS.42}
\]
extends to a trace-class operator. With $G=\mathcal Mg$,
\[
\begin{aligned}
 \operatorname{Tr}K_{a,\eta,S,\phi}(g)
  &=-\frac1{2\pi}\int_{\mathbb R}G(a+it)
             \frac{A_\eta'}{A_\eta}(\sigma+it)\,dt\\
  &=\sum_{\substack{\mathfrak p\notin S\\k\ge1}}
       \eta(\mathfrak p)^k\frac{\log N\mathfrak p}
                                    {(N\mathfrak p)^{k/2}}
                         g(k\log N\mathfrak p)
       -\delta_\eta G(1/2)
       +\sum_{\substack{\rho\in Z_\eta\\\operatorname{Re}\rho>\sigma}}
                         m_\rho G(\rho-1/2).
\end{aligned}
\tag{MA-RS.43}
\]
Write the finite prime sum as $P^+_{\eta,\overline S}(g)$ and
put
\[
 B_{\eta,\overline S}(g)
   =P^+_{\eta,\overline S}(g)
                  +P^+_{\overline\eta,\overline S}(g^\vee).
\]
The actual reflected pair on $\mathcal H_a\oplus\mathcal H_{-a}$ is
\[
 \mathbb K_{a,\eta,S,\phi}(g)
   =K_{a,\eta,S,\phi}(g)
                  \oplus R K_{a,\overline\eta,S,\phi}(g^\vee)R.
 \tag{MA-RS.44}
\]
Its trace is
\[
 \operatorname{Tr}\mathbb K_{a,\eta,S,\phi}(g)
  =B_{\eta,\overline S}(g)-\delta_\eta\bigl(G(1/2)+G(-1/2)\bigr)
    +\sum_{\substack{\rho\in Z_\eta\\
                |\operatorname{Re}\rho-1/2|>a}}
                         m_\rho G(\rho-1/2).
 \tag{MA-RS.45}
\]
Every infinite zero sum here is absolutely convergent.

If both divisors instead have a fixed positive horizontal distance
from the chosen line, all these assertions hold for arbitrary
$g\in C_c^\infty$ and $\phi'\in C_c^\infty$, with continuity into
trace norm on fixed compact test supports. Under their Hecke GRH,
this applies at every $1/2<\sigma<1$ and the last sum in (MA-RS.45)
is empty. For Dirichlet characters this uses the stipulated theorem
$A$. For a general non-induced imaginary-quadratic character it uses
its own Hecke zero-location statement; $M$ is sufficient.

\begin{proof}
\emph{Generic inverse estimates and actual domains.}
NT-Z10 and MT-M18 give, for this fixed Hecke function, polynomial
upper bounds on fixed vertical strips, the local zero count
$O(\log(2+|T|))$, and the Euler lower bound
$|L_\eta(2+it)|\ge\zeta_K(2)^{-1}$.
Apply the Borel--Cantelli construction in MA-RS-001 to the union
of the two zero divisors. For almost every $\sigma$, both satisfy
\[
 |\sigma-\operatorname{Re}\rho|
                 \ge c_\sigma(2+|\operatorname{Im}\rho|)^{-2}.
 \tag{MA-RS.46}
\]
At large height the fixed disk centered at $2+it$ contains no
trivial zero or pole of $L_\eta$, and contains only
$O(\log(2+|t|))$ nontrivial zeros. Dividing them out with the
finite Blaschke product and applying Harnack as in (MA-RS.5)
gives the same polynomial lower bound for the zero-free factor.
The product of the bounds in (MA-RS.46) then gives
\[
 |\partial_t^j A_\eta(\sigma+it)^{-1}|
       \le \exp(C_j\log^2(2+|t|))\quad(j\ge0).
 \tag{MA-RS.47}
\]
Cauchy estimates on disks of radius a fixed multiple of
$(2+|t|)^{-2}$ prove the derivative assertion exactly as in
MA-RS-001. All derivatives of $A_\eta(\sigma+it)$ have
polynomial growth. The finite factors $E_{\eta,S}$ and their
inverses have bounded derivatives on a small closed strip about
$\sigma>1/2$, since $|\eta(\mathfrak p)|$ is zero or one.
The compact height interval is covered by nonvanishing on the
line. All statements apply also to $\overline\eta$.

The maximal multiplier operators in (MA-RS.41) are therefore
closed and have the same dense compact-frequency graph core as
in MA-RS.9. Multiplication by $\phi$ sends such a core vector
and its image under $J$ to Fourier tails bounded by
$\exp(-c|t|^{1/q})$, where $q>1$ is a common Gevrey order.
Formula (MA-RS.20) proves this directly, independently of the
arithmetic multiplier. These tails dominate (MA-RS.47), so both
$J M_\phi$ and $Z M_\phi J$ are defined on that core and have
the requisite domain for the next operation.
In weighted Fourier variables the resulting kernel is exactly
\[
 \frac{i}{2\pi}G(a+it)A_\eta(\sigma+it)
  \frac{A_\eta(\sigma+it)^{-1}-A_\eta(\sigma+iu)^{-1}}{t-u}
                 \widehat{\phi'}_+(t-u).
 \tag{MA-RS.48}
\]
The divided difference has subexponential derivative bounds
(MA-RS.47). The two Gevrey factors dominate them in both
variables, as in (MA-RS.22)--(MA-RS.23). Thus this is a Schwartz
kernel, and its Hermite expansion converges in trace norm.
Its diagonal integral, using $\int\phi'=-1$, is the first line
of (MA-RS.43). This proves trace class for the actual composition
before its arithmetic trace is evaluated.

\emph{An explicit ideal-sum regularization.}
The closed operator $Z$ also arises from bounded, explicitly
regularized ideal sums; analytic continuation is not being
substituted for convergence of the unaltered Euler sums.
For $\varepsilon>0$ set
\[
 \begin{aligned}
 L_{\eta,\varepsilon}(s)
     &=\sum_{\mathfrak a}\eta(\mathfrak a)(N\mathfrak a)^{-s}
                                           e^{-\varepsilon N\mathfrak a},\\
 A_{\eta,\varepsilon}(s)
     &=E_{\eta,S}(s)\left(
        L_{\eta,\varepsilon}(s)
             -c_\eta\Gamma(1-s)\varepsilon^{s-1}\right).
 \end{aligned}
 \tag{MA-RS.49}
\]
For fixed $\varepsilon$, these are bounded functions on the
line $\operatorname{Re}s=\sigma$: the ideal series is absolutely
summable there, and $\Gamma(1-\sigma-it)$ is bounded on the
real $t$ line. The finite Euler factor is bounded as well.

Mellin inversion of $e^{-x}$ gives, for $b>\max(0,1-\sigma)$,
\[
 L_{\eta,\varepsilon}(s)
   =\frac1{2\pi i}\int_{b-i\infty}^{b+i\infty}
                  \Gamma(z)L_\eta(s+z)\varepsilon^{-z}\,dz.
 \tag{MA-RS.50}
\]
The interchange with the ideal series is absolute because
$\sigma+b>1$ and $\Gamma(b+iv)$ is integrable.
Fix $0<\theta<\min(1,\sigma)$. Move this contour to
$\operatorname{Re}z=-\theta$.
The crossed poles are $z=0$, with residue $L_\eta(s)$,
and, only for the principal character, $z=1-s$, with residue
$c_\eta\Gamma(1-s)\varepsilon^{s-1}$.
There are no other poles between these lines.
Stirling's exponential vertical decay for $\Gamma$, together
with the fixed-strip polynomial bound for $L_\eta$, makes
the horizontal integrals vanish. The integral on the new
line is $O_T(\varepsilon^\theta)$ uniformly for
$s=\sigma+it$, $|t|\le T$.
Consequently $A_{\eta,\varepsilon}(\sigma+it)$ converges to
$A_\eta(\sigma+it)$ uniformly on every compact frequency
interval. The corresponding bounded multipliers converge on
the compact-frequency core, and the graph closure of the
limit is exactly the maximal $Z$ in (MA-RS.41).
This is an actual regularization with its pole counterterm;
no operator-norm convergence on the whole Hilbert space is claimed.

\emph{The prime and reflected terms.}
On $\operatorname{Re}s=c>1$, the Euler logarithmic derivative is
\[
 -\frac{A_\eta'}{A_\eta}(s)
   =\sum_{\substack{\mathfrak p\notin S\\k\ge1}}
              \eta(\mathfrak p)^k(\log N\mathfrak p)
                                     (N\mathfrak p)^{-ks}.
\]
Absolute convergence and Fourier inversion give
$P^+_{\eta,\overline S}(g)$ for the line integral there.
Shift the line to $\sigma$, using the selected horizontal
heights of MT-M18 for this function. Its minimum-modulus
proof and Stirling bounds give polynomial logarithmic-derivative
bounds on those horizontal segments. The Gevrey Mellin factor
makes their integrals tend to zero.
Each crossed zero has residue $m_\rho$ in $A_\eta'/A_\eta$,
and the pole at one has residue $-\delta_\eta$.
The finite removed Euler factors have their zeros on
$\operatorname{Re}s=0$ and add no crossed term.
The contour orientation is that of MA-RS-003; it gives exactly
(MA-RS.43). The zero count and Mellin decay give absolute
convergence of its residue sum.

The raw character-sector reflection in MA-CC-001 exchanges
$\eta$ and $\overline\eta$. The completed functional equation
identifies the zero divisor of $\overline\eta$ with
$1-Z_\eta$, including multiplicity. Apply (MA-RS.43) to
$\overline\eta$ and $g^\vee$; since
$\mathcal M(g^\vee)(z)=G(-z)$, its right-hand zeros become
exactly the left-hand zeros of $\eta$. Adding gives (MA-RS.45).
A root number is a nonzero constant and introduces no extra
logarithmic-derivative or divisor term.

Finally, with a fixed horizontal gap, every finite Blaschke
factor has a fixed positive lower bound. Since the number of
factors is $O(\log(2+|t|))$, the reciprocal bounds and all their
derivatives are polynomial. The proof of MA-RS-004 then applies
to (MA-RS.48): arbitrary compact smooth tests give a Schwartz
kernel and a continuous trace-class family. Under the stated
Hecke zero-location premise the gap is $\sigma-1/2$ and both
outside sums are empty. This proves all claimed extensions.
\end{proof}

\paragraph{MA-RS-006 (the non-induced cubic sector in the narrow regularization).}
Take the verified character of MA-CH-003:
\[
 K=\mathbb Q(\zeta_3),\qquad
 K(269^{1/3})/K,\qquad \eta=\psi,\qquad
 S=\{\mathfrak p\mid3\cdot269\}.
\]
The test of MA-CH-004 can be chosen of finite Gevrey order.
For that choice it has $g(0)=0$, its local $S$ contribution
is zero, and its full unramified prime contribution is $\log2$.
At every generic width covered by MA-RS-005,
\[
 \operatorname{Tr}\mathbb K_{a,\psi,S,\phi}(g)
     =\log2+
       \sum_{\substack{\rho\in Z_\psi\\
                    |\operatorname{Re}\rho-1/2|>a}}
                       m_\rho\mathcal Mg(\rho-1/2).
 \tag{MA-RS.51}
\]
Under $M$, the same construction works at every $0<a<1/2$
with arbitrary compact smooth versions of this test, and its
trace is exactly $\log2$.

\begin{proof}
The construction in MA-CH-004 uses one bump near norm four,
normalized to have value one there, and subtracts a scalar
multiple of a second bump supported in the interval $(1.4,1.6)$
to cancel the complex archimedean integral. Both bumps can be
chosen compact Gevrey; taking the same scalar combination
preserves that class. The second archimedean integral is nonzero
because its kernel is strictly positive and the bump is chosen
nonnegative and nonzero. All support, value and cancellation
requirements therefore remain exactly those of MA-CH-004.

The sole prime-power contribution in those supports is the
unique prime above two, of norm four and character value one.
Its coefficient is $(\log4)/\sqrt4=\log2$.
The reflected positive-prime sum is zero because $g$ is
supported on positive logarithmic norms. The character is
nonprincipal, so $\delta_\psi=0$. Substitution in (MA-RS.45)
gives (MA-RS.51). Under $M$, all zeros are central, the fixed
gap hypothesis of the smooth extension holds, and the last sum
is empty. This proves the assertion for the actual operator.
It does not use the Dirichlet theorem $A$ to supply a Hecke
premise that $A$ does not contain.
\end{proof}

\paragraph{Character geometry and source locators.}
The raw compact character convention and its inversion under
reflection are those checked in MA-CH-001 and MA-CC-001.
The prime terms are Frobenius traces of the corresponding
rank-one finite-monodromy local system on the open arithmetic
curve where the character is unramified. In the cubic example,
induction to $\mathbb Q$ gives the verified odd rank-two
$S_3$ Artin representation of MA-CH-003.
This is the SGA~1, Expos\'e~V, \S7 finite-\'etale/monodromy
dictionary, and the Artin-motive convention of Deligne,
\emph{Valeurs de fonctions L}, Definition~6.1 and \S\S6.2--6.3.
The actual nuclear trace remains an analytic realization on a
chosen completion, with the cote~1 comparison at inherited
scans~079--080 and~111--112 through BR-TR-01/02.
No positive arithmetic cohomology or new $\ell$-adic weight
statement is inferred from the regularization.

The unconditional Hecke analytic inputs are those of NT-Z10
and MT-M18: Ramakrishnan and Valenza, \emph{Fourier Analysis on
Number Fields}, \S7.4, Theorem~7-19, p.~279, and \S7.7,
Theorem~7-28, p.~289, with the precise strip/count consequences
proved in the inherited modules. The inverse bounds and actual
operator arguments are given above, using the already checked
MA-RS-001--004 kernel proof. The ideal-sum regularization
(MA-RS.49)--(MA-RS.50) uses Mellin inversion of $e^{-x}$:
NIST Digital Library of Mathematical Functions,
(2.5.1)--(2.5.2) and (5.9.1), with $\mu=z=1$ in the latter.
The official formulas were checked on 21 November 2026 at
\url{https://dlmf.nist.gov/2.5.E2} and
\url{https://dlmf.nist.gov/5.9.E1}. The residue computation and
uniform compact-frequency convergence required here are proved
in full, rather than left as an appeal to analytic continuation.

\section{Frequency cutoffs for all compact smooth tests on generic lines}
\label{sec:regularized-frequency}

The generic-line logarithmic derivative in MA-RS-001 is better
controlled than the reciprocal multiplier itself. This distinction
allows a further explicit regularization: each finite-frequency
expression is an actual trace-class operator, and its trace has a
limit for every compact smooth test. The construction below uses two
bounded compositions before any kernel or trace is evaluated.
Its limit is a limit of traces. It does not assert convergence of
the operators on a merely generic line.

Use the notation of (MA-RS.8), with
$\sigma\in\mathscr G$, $a=\sigma-1/2$, and
\[
 A(t)=\zeta_S(\sigma+it),\qquad
 j(t)=A(t)^{-1},\qquad
 G(z)=\mathcal Mg(z),\qquad m(t)=G(a+it).
 \tag{MA-RS.52}
\]
Here $g\in C_c^\infty(\mathbb R)$ is arbitrary.
Let $\phi$ be bounded and smooth, with
$\phi(-\infty)=1$, $\phi(+\infty)=0$, and
$\phi'\in C_c^\infty(\mathbb R)$, and put
\[
 T_\phi=V_aM_\phi V_a^{-1}
        =\mathcal F_+M_\phi\mathcal F_+^{-1},
 \qquad
 p(\xi)=\int_{\mathbb R}e^{i\xi x}\phi'(x)\,dx.
 \tag{MA-RS.53}
\]
In particular $T_\phi$ is bounded on $L^2(\mathbb R)$; it
is the Fourier conjugate of multiplication by the physical cutoff.
It is not multiplication by $\phi$ in the frequency variable.
Fix a real $\chi\in C_c^\infty(\mathbb R)$ equal to one in a
neighborhood of zero, and set $\chi_N(t)=\chi(t/N)$ for integers
$N\ge1$.

\paragraph{MA-RS-007 (actual finite-frequency traces and their smooth limit).}
Define the operator $K_N(g)$ on $\mathcal H_a$ by the literal
bounded-composition formula
\[
 \begin{split}
 \widetilde K_N(g)&:=V_aK_N(g)V_a^{-1}\\
  &=M_{m\chi_N}\,T_\phi\,M_{\chi_N}
       -M_{mA\chi_N}\,T_\phi\,M_{j\chi_N}.
 \end{split}
 \tag{MA-RS.54}
\]
Both displayed compositions are bounded. Their difference is trace
class, with smooth compactly supported kernel
\[
 \mathcal K_N(t,u)
   =\frac{i}{2\pi}\chi_N(t)\chi_N(u)m(t)A(t)
          \frac{j(t)-j(u)}{t-u}\,p(t-u).
 \tag{MA-RS.55}
\]
The divided difference is assigned its smooth value $j'(t)$
on the diagonal. The trace and its limit are
\[
 \begin{aligned}
 \operatorname{Tr}K_N(g)
   &=-\frac1{2\pi}\int_{\mathbb R}\chi_N(t)^2G(a+it)
                    \frac{\zeta_S'}{\zeta_S}(\sigma+it)\,dt,\\
 \lim_{N\to\infty}\operatorname{Tr}K_N(g)
   &=-\frac1{2\pi}\int_{\mathbb R}G(a+it)
                    \frac{\zeta_S'}{\zeta_S}(\sigma+it)\,dt\\
   &=P_{\overline S}^{+}(g)-G(1/2)
       +\sum_{\operatorname{Re}\rho>\sigma}
                      m_\rho G(\rho-1/2).
 \end{aligned}
 \tag{MA-RS.56}
\]
The limiting integral and zero sum are absolutely convergent.
The limit is a continuous distribution in $g$ and is independent
of $\chi$ and of the shape of $\phi$ with its fixed endpoints.

Let $R:\mathcal H_a\longrightarrow\mathcal H_{-a}$ denote
physical reflection as in MA-RS-003. Define actual trace-class
operators on the direct sum by
\[
 \mathbb K_N(g)=K_N(g)\oplus R K_N(g^\vee)R.
 \tag{MA-RS.57}
\]
Then
\[
 \lim_{N\to\infty}\operatorname{Tr}\mathbb K_N(g)
   =B_{\overline S}(g)-G(1/2)-G(-1/2)
       +\sum_{|\operatorname{Re}\rho-1/2|>a}
                         m_\rho G(\rho-1/2).
 \tag{MA-RS.58}
\]
There are no boundary zeros at this generic width. These assertions
use neither $A$ nor $M$.

\begin{proof}
\emph{The two operators are defined before the trace.}
The generic-line estimate (MA-RS.2) excludes zeros of $\zeta$
on the line. The pole at one is off it, and the finitely many
removed Euler factors have no zeros there. Thus $A$ and $j$
are smooth on the real frequency axis. All four multipliers
in (MA-RS.54) are smooth and compactly supported, hence bounded.
Since $T_\phi$ is bounded, (MA-RS.54) defines two bounded
operators on all of $L^2$, without any domain assertion for
the uncut expression $J_{a,S}M_\phi$.

The distribution kernel of $T_\phi$ is
$(2\pi)^{-1}\widehat\phi_+(t-u)$, where the hat here is the
unnormalized transform used for $p$ in (MA-RS.53).
On compactly supported smooth frequency vectors, the difference
of the two kernels in (MA-RS.54) is this distribution multiplied
by
\[
 \chi_N(t)\chi_N(u)m(t)\bigl(1-A(t)j(u)\bigr)
  =\chi_N(t)\chi_N(u)m(t)A(t)\bigl(j(t)-j(u)\bigr).
 \tag{MA-RS.59}
\]
Since $\widehat{\phi'}_+(\xi)=-i\xi\widehat\phi_+(\xi)$,
one has $\xi\widehat\phi_+(\xi)=ip(\xi)$ as distributions.
The identity
\[
 j(t)-j(u)=(t-u)\int_0^1j'(u+\theta(t-u))\,d\theta
 \tag{MA-RS.60}
\]
therefore converts (MA-RS.59) exactly into (MA-RS.55).
It also cancels any delta term of $\widehat\phi_+$ at zero.
The integral in (MA-RS.60) is smooth in both variables.
Multiplication by the two compact cutoffs consequently gives
$\mathcal K_N\in C_c^\infty(\mathbb R^2)$.
The resulting integral operator agrees on a dense set with
the already defined bounded difference in (MA-RS.54), hence
agrees everywhere.

\emph{Trace class and the diagonal formula.}
Here is a direct trace-class argument that does not assume a
trace from the diagonal alone. Choose $L$ so that the support
of $\mathcal K_N$ lies in the interior of $[-L,L]^2$.
Use the normalized Fourier basis
$e_r(t)=(2L)^{-1/2}e^{\pi irt/L}$, $r\in\mathbb Z$,
on $[-L,L]$, extended by zero outside that interval.
The smooth kernel is zero near the boundary of the square, so
its periodic Fourier expansion has the form
\[
 \mathcal K_N(t,u)=\sum_{r,s\in\mathbb Z}
              c_{r,s}e_r(t)\overline{e_s(u)},\qquad
 |c_{r,s}|\le
       C_N(1+r^2)^{-1}(1+s^2)^{-1}.
 \tag{MA-RS.61}
\]
The estimate follows by two integrations by parts in each
variable, or by applying $1-\partial_t^2$ and
$1-\partial_u^2$ before taking the Fourier coefficients.
The coefficients are absolutely summable. Accordingly the
associated rank-one operator series converges in trace norm,
and the kernel series converges uniformly on the square.
Its trace is $\sum_r c_{r,r}$, which, by the same absolute
summability and the orthonormality of this basis, equals
$\int\mathcal K_N(t,t)\,dt$.
This proves trace class and the diagonal formula for the
actual operator. No bound uniform in $N$ is needed here.

Finally $p(0)=\int\phi'=-1$ and
$j'(t)=-i\zeta_S'(\sigma+it)/\zeta_S(\sigma+it)^2$.
The diagonal of (MA-RS.55) is therefore
\[
 \mathcal K_N(t,t)=
    -\frac1{2\pi}\chi_N(t)^2G(a+it)
                         \frac{\zeta_S'}{\zeta_S}(\sigma+it).
 \tag{MA-RS.62}
\]
This gives the first line of (MA-RS.56), including its sign
and its Fourier normalization.

\emph{A polynomial bound for the logarithmic derivative.}
We give the needed estimate explicitly; the superpolynomial
upper bound for $1/\zeta_S$ in (MA-RS.3) by itself would not
give it. At large $|t|$ use the finite disk in the proof of
MA-RS-001, centered at $c_t=2+it$ and with a zero-free boundary
of radius $12<R_t<13$. Write $\zeta=H_t B_t$ on that disk,
where $B_t$ is the finite disk Blaschke product over its
$O(\log(2+|t|))$ zeros, with multiplicities.
That proof gives a polynomial upper bound for $|H_t|$
throughout the disk and the polynomial lower bound (MA-RS.5)
on its inner half. At $w=\sigma+it$, Cauchy's estimate on a
fixed-radius circle inside this inner half disk gives a
polynomial bound for $H_t'(w)$. Therefore $H_t'(w)/H_t(w)$
is polynomially bounded.

For a disk Blaschke factor with zero $\rho$, its logarithmic
derivative at $w$ is
\[
 \frac1{w-\rho}+
     \frac{\overline{\rho-c_t}}
          {R_t^2-\overline{\rho-c_t}(w-c_t)}.
 \tag{MA-RS.63}
\]
The unimodular normalizing constant has no derivative.
At $w=\sigma+it$ the second summand is bounded by a fixed
constant, since $|w-c_t|<R_t/4$. The first summand is
$O_\sigma((2+|t|)^2)$ by (MA-RS.2) and the bounded height
range of this disk. Summing over its $O(\log(2+|t|))$
zeros proves a polynomial bound for $B_t'(w)/B_t(w)$.
Adding the two logarithmic derivatives and the bounded
logarithmic derivative of $E_S$ gives constants $C,D$ with
\[
 \left|\frac{\zeta_S'}{\zeta_S}(\sigma+it)\right|
              \le C(2+|t|)^D\qquad(t\in\mathbb R).
 \tag{MA-RS.64}
\]
The compact height range follows from continuity and
nonvanishing on the line.

For arbitrary $g\in C_c^\infty$, repeated integration by
parts shows that $m(t)=G(a+it)$ is Schwartz. Thus its product
with (MA-RS.64) is integrable. Moreover
$|\chi_N(t)|\le\|\chi\|_\infty$ and $\chi_N(t)\to1$
pointwise. Dominated convergence in (MA-RS.62) proves the
second line of (MA-RS.56). The same bound and integration
by parts bound this limiting functional by finitely many
derivative seminorms of $g$ on every fixed compact support.
This proves its continuity, and proves independence of
the frequency cutoff and of the physical cutoff's shape.

\emph{Arithmetic evaluation of the established trace limit.}
Only now rewrite the limiting integral as the contour
integral (MA-RS.29). On a line $\operatorname{Re}s=c>1$
the Euler logarithmic derivative converges absolutely,
and Fourier inversion gives $P_{\overline S}^{+}(g)$
exactly as in (MA-RS.30)--(MA-RS.31).
On the fixed strip $\sigma\le\operatorname{Re}s\le c$,
integration by parts gives
\[
 |G(s-1/2)|\le C_{g,c,\sigma,M}
                    (1+|\operatorname{Im}s|)^{-M}
                  \quad\text{for every integer }M\ge0.
 \tag{MA-RS.65}
\]
Use the selected upper and lower heights of MT-M18
(M18.7)--(M18.8). Its finite-disk minimum-modulus proof
and the fixed-strip estimates give a polynomial bound
for $\zeta_S'/\zeta_S$ on the horizontal segments at
these heights, as checked in MA-RS-003.
Choosing $M$ larger than this polynomial exponent makes
the horizontal integrals tend to zero. On the generic
vertical line absolute convergence follows from
(MA-RS.64); on the line $c$ it follows from the Euler
series and (MA-RS.65).

The counterclockwise rectangle has its right edge directed
upwards and its left edge downwards. The residue of
$G(s-1/2)\zeta_S'(s)/\zeta_S(s)$ is
$m_\rho G(\rho-1/2)$ at a crossed nontrivial zero and
$-G(1/2)$ at the pole at one. Zeros of the removed Euler
factors have real part zero and are not crossed.
Multiplying the resulting left-line integral by
$-(2\pi i)^{-1}$ gives the final line of (MA-RS.56).
The local zero count $O(\log(2+|T|))$ in fixed-height
windows, together with (MA-RS.65), proves absolute
convergence of the zero sum.

Physical reflection is unitary from $\mathcal H_a$ onto
$\mathcal H_{-a}$, so (MA-RS.57) is an actual trace-class
direct sum for each $N$. Apply (MA-RS.56) to $g$ and
$g^\vee$ and add their limits. The identity
$\mathcal M(g^\vee)(z)=G(-z)$ and the multiplicity-preserving
symmetry $\rho\mapsto1-\rho$ turn the second right-hand
zero sum into the sum over $\operatorname{Re}\rho<1-\sigma$.
This proves (MA-RS.58). The symmetry also excludes zeros
on the reflected boundary. No parity of $\chi$ is required:
reflection is applied to the already defined operator
$K_N(g^\vee)$.
\end{proof}

\paragraph{The same finite cutoffs in character sectors.}
For the primitive finite-order Hecke character $\eta$ of
MA-RS-005, choose $\sigma$ in the simultaneous generic set
for $\eta$ and $\overline\eta$. Replace $A(t)$ in
(MA-RS.52)--(MA-RS.55) by
$A_\eta(\sigma+it)$ and $j(t)$ by its reciprocal, and denote
the resulting actual operators by $K_{\eta,N}(g)$.
For arbitrary compact smooth $g$ and $\phi'$ they satisfy
\[
 \begin{aligned}
 \lim_{N\to\infty}\operatorname{Tr}K_{\eta,N}(g)
  &=P^+_{\eta,\overline S}(g)-\delta_\eta G(1/2)
       +\sum_{\substack{\rho\in Z_\eta\\
                         \operatorname{Re}\rho>\sigma}}
                   m_\rho G(\rho-1/2),\\
 \lim_{N\to\infty}\operatorname{Tr}
   \bigl(K_{\eta,N}(g)\oplus
                R K_{\overline\eta,N}(g^\vee)R\bigr)
  &=B_{\eta,\overline S}(g)
         -\delta_\eta\bigl(G(1/2)+G(-1/2)\bigr)\\
  &\qquad +\sum_{\substack{\rho\in Z_\eta\\
                    |\operatorname{Re}\rho-1/2|>a}}
                       m_\rho G(\rho-1/2).
 \end{aligned}
 \tag{MA-RS.66}
\]
Indeed the finite-cutoff boundedness, distributional kernel,
Fourier-series trace-class proof and diagonal evaluation
above use only smooth nonvanishing of $A_\eta$ on the line,
which MA-RS-005 supplies. That proof also supplies its
finite-disk Blaschke factorization, polynomial upper and
zero-free-factor lower bounds, and the distance (MA-RS.46).
Applying exactly (MA-RS.63) to those disk factors gives
a polynomial bound for $A_\eta'/A_\eta$ on the line,
including its bounded finite Euler-factor contribution.
Hence dominated convergence works for every compact smooth
test. The selected-height bounds, Euler derivative and
residues in the last part of MA-RS-005 then give the first
identity of (MA-RS.66), using (MA-RS.65) in place of
Gevrey decay. Its pole has order $\delta_\eta$.
Finally its completed functional equation identifies
$Z_{\overline\eta}$ with $1-Z_\eta$, with multiplicity.
This is precisely the conjugate reflection needed for the
second identity. The root number is a constant and adds
no term. This proves the character extension without
assuming its Hecke GRH.

\paragraph{Compatibility and the precise remaining limit issue.}
When the uncut operator $K_{a,S,\phi}(g)$ of MA-RS-002
is known to be trace class, put
$Q_N=V_a^{-1}M_{\chi_N}V_a$. Its kernel formula and
(MA-RS.55) give the actual equality
\[
           K_N(g)=Q_NK_{a,S,\phi}(g)Q_N.
 \tag{MA-RS.67}
\]
Thus $K_N(g)\to K_{a,S,\phi}(g)$ in trace norm in that
already proved case. To justify this last statement,
$Q_N$ are uniformly bounded self-adjoint operators and
converge strongly to one. The conclusion first holds for
each rank-one operator, since both its defining vectors
converge in norm after applying $Q_N$, then for finite-rank
operators by addition, and then for all trace-class
operators by finite-rank approximation and the uniform
bound on $\|Q_N\|$. It applies to generic-line Gevrey
data by MA-RS-002, or to arbitrary compact smooth data
under the fixed gap of MA-RS-004. The same reasoning
applies in the character cases already covered by
MA-RS-005.

For arbitrary compact smooth data on a merely generic line,
this section proves only the existence and the stated
limit of the actual finite-frequency traces.
It asserts neither trace-norm nor operator convergence
of $K_N(g)$ in that larger scope. The regularization has
changed by inserting frequency cutoffs into the two
bounded compositions. It does not identify their limit
with Masullo's physical cutoff, an orthogonal quotient
projection, or all off-diagonal maps in Meyer's common
space. The polar and exterior-divisor corrections in
(MA-RS.58) and (MA-RS.66) remain part of the formula.

\paragraph{Sources and geometric placement.}
The generic set, disk bounds and selected heights are
the unconditional inputs of MA-RS-001--003, with their
exact inherited source NT-Z10 and MT-M18
(M18.5)--(M18.9). The character extension uses the
Hecke inputs and exact conventions of MA-RS-005,
including Ramakrishnan and Valenza (1999),
\emph{Fourier Analysis on Number Fields}, \S7.4,
Theorem~7-19, p.~279, and \S7.7, Theorem~7-28, p.~289,
through the already checked inherited consequences.
The bounded-composition, compact-kernel, diagonal and
limit arguments are proved above. Meyer's original
wide-strip diagonal trace is located at Meyer (2005),
Lemmas~5.3--5.6 and Theorem~5.8, author pp.~39--44,
as checked in MA-MEY-001--003; no new narrow-width
claim is attributed to that source.
In a character sector, the prime coefficients retain
the Frobenius-trace interpretation of MA-RS-005 and
MA-CH-001: SGA~1, Expos\'e~V, \S7, and Deligne,
\emph{Valeurs de fonctions L}, Definition~6.1 and
\S\S6.2--6.3. The trace is the trace of the actual
regularized analytic operators constructed here;
this interpretation does not supply positivity or
zero location.

\section{An actual partial return to the smooth Meyer quotient}
\label{sec:partial-return}

The obstruction in MA-HM-003 concerns a return that recovers
every smooth spectral block.  There is a different return to
the closed smooth subrepresentation supported on the zeros that
the Hilbert completion retains.  At generic widths, compact
Gevrey smoothing makes this return nuclear.  The construction
below uses actual Mellin restrictions and entire interpolation
maps; comparison of scalar zero sums is not used to establish
its operator identities.

Fix the primitive unitary Hecke function and notation of
MT-M18 and MA-WP.  To distinguish the two quotients, write
\[
\begin{aligned}
 Q_\infty&=\mathcal E/\mathcal I,\qquad q_\infty:\mathcal E\to Q_\infty,
       &\mathcal I&=\overline{\psi*\mathcal E}^{\,\mathcal E},\\
 H&=L^2(\mathbb R,w_{u,v,s}(x)\,dx),&
 w_{u,v,s}(x)&=\max(e^{2ux},e^{2vx})(1+x^2)^{s/2},\\
 J_H&=\overline{T_a\mathcal I}^{\,H},&
 Q_H&=H/J_H,\qquad \pi:H\to Q_H,\\
 i:Q_\infty&\longrightarrow Q_H,&
 i[q_\infty f]&=\pi(T_af),\qquad T_af(x)=e^{ax}f(x).
\end{aligned}
\tag{MA-PR.1}
\]
Here $u<v$, $s\ge0$, and all constants concern these fixed
parameters and this fixed Hecke function.  The map $i$ is
continuous and has dense image, but it need not be injective.
Let $\rho=\beta+i\gamma$ be the distinct zeros of the entire
pole-removed completion $\Xi$, with multiplicities $m_\rho$.

Assume the following quantitative avoidance of both edges:
\[
 \min\{|\beta-(a+u)|,\ |\beta-(a+v)|\}
       \ge c_0(2+|\gamma|)^{-2}
       \qquad(\Xi(\rho)=0)
\tag{MA-PR.2}
\]
for one $c_0>0$.  This excludes boundary zeros.
For the fixed function, it holds for almost every pair
$(u,v)$ with $u<v$.  For centered intervals
$a=1/2$, $u=-\alpha$, $v=\alpha$, it holds for almost every
$\alpha>0$.  Indeed the local zero count makes
$\sum_\rho(2+|\gamma|)^{-2}<\infty$; apply the
Borel--Cantelli argument of MA-RS-001 to both edge
parameters and remove the countable set of exact equalities.
For centered intervals the two families of forbidden intervals
are centered at $\beta-1/2$ and $1/2-\beta$.
No zero-location assertion is required.

Set
\[
\begin{aligned}
 Z_{\rm in}&=\{\rho:a+u<\operatorname{Re}\rho<a+v\},\\
 Z_{\rm out}&=\{\rho:\operatorname{Re}\rho<a+u
                         \text{ or }\operatorname{Re}\rho>a+v\},\\
 Q_{\rm in}
   &=\{[f]\in Q_\infty:
       (\mathcal Mf)^{(j)}(\rho)=0
       \text{ for }\rho\in Z_{\rm out},\ 0\le j<m_\rho\}.
\end{aligned}
\tag{MA-PR.3}
\]
This is a closed, translation-invariant subspace of the actual
smooth quotient.  Define $Q_{\rm out}$ by the same formula with
$Z_{\rm in}$ and $Z_{\rm out}$ interchanged.  These definitions
do not assert a continuous unsmoothed projection onto either
subspace.

\paragraph{MA-PR-001 (inner contour restrictions and partial entire lifts).}
There are finite-dimensional Banach algebras $J_n^{\rm in}$,
continuous maps
\[
 r_n:Q_H\longrightarrow J_n^{\rm in},
 \qquad \mathcal T_n^{\rm in}:J_n^{\rm in}\longrightarrow\mathcal E,
\tag{MA-PR.4}
\]
and constants, with $\nu_n=2+|n|$, such that
\[
\begin{gathered}
 d_n^{\rm in}:=\dim J_n^{\rm in}=O(\log\nu_n),\qquad
             \|r_n\|\le C\nu_n,\\
 p_{N,j}(\mathcal T_n^{\rm in}b)
        \le C_{N,j}\exp(C_{N,j}\log^2\nu_n)\|b\|_n,\\
 r_m\,i\,q_\infty\mathcal T_n^{\rm in}
        =\delta_{mn}\operatorname{id}_{J_n^{\rm in}}.
\end{gathered}
\tag{MA-PR.5}
\]
The smooth class of every lift has exactly its prescribed
full jets in the interior zeros of block $n$, and zero full
jets at all other zeros of $\Xi$.
The restriction $i_{\rm in}:Q_{\rm in}\to Q_H$ is continuous,
injective and has dense image.

\begin{proof}
Use the horizontal heights $h_n$ of MT-M18, with
$n+1/4<h_n<n+3/4$.  Choose a sufficiently small fixed
$\delta_0>0$ and put
\[
 \delta_n=\delta_0\nu_n^{-2},\qquad
 \Omega_n^{\rm in}
   =\{a+u+\delta_n<\operatorname{Re}z<a+v-\delta_n,\
                          h_n<\operatorname{Im}z<h_{n+1}\}.
\tag{MA-PR.6}
\]
Choose $\delta_0<(v-u)/4$ and $\delta_0\le1$.
In each fixed-height window about block $n$, (MA-PR.2)
bounds the distance of every zero from either original
vertical edge below by $c\nu_n^{-2}$, with $c$ independent
of $n$.  Make $\delta_0$ smaller than one quarter of this
constant.  The two inner edges then stay at distance at
least $c'\nu_n^{-2}$ from every zero in that window.
The rectangle contains exactly the zeros of $Z_{\rm in}$
whose ordinates are between $h_n$ and $h_{n+1}$.
Its boundary contains no zero.  Its width stays bounded
below by $(v-u)/2$, its height is between $1/2$ and $3/2$,
and its perimeter is uniformly bounded.

Let $J_n^{\rm in}$ be the algebra $A(\Omega_n^{\rm in})$
modulo its ideal of full jets at those interior zeros,
with the quotient of the boundary supremum norm.
This is the same finite holomorphic quotient norm as
(M18.1), for the new rectangle.  Polynomial Hermite
interpolation proves the finite dimension, and the local
zero count gives $d_n^{\rm in}=O(\log\nu_n)$.
Its norm is submultiplicative.  An empty block is the zero
algebra.

For $h\in H$, the function
$w\mapsto\mathcal Mh(w-a)$ is holomorphic on
$a+u<\operatorname{Re}w<a+v$ and continuous on the
closed inner rectangle.  Cauchy--Schwarz gives, for
$u+\delta_n\le\operatorname{Re}z\le v-\delta_n$,
\[
\begin{aligned}
 |\mathcal Mh(z)|
 &\le\|h\|_H
  \left(\int_{\mathbb R}
       \frac{e^{2\operatorname{Re}z\,x}}
            {\max(e^{2ux},e^{2vx})(1+x^2)^{s/2}}\,dx\right)^{1/2}\\
 &\le C\delta_n^{-1/2}\|h\|_H
 \le C'\nu_n\|h\|_H .
\end{aligned}
\tag{MA-PR.7}
\]
The elementary integral bound uses $s\ge0$ and
$e^{-2\delta_n|x|}$ on the two tails.  Restricting this
representative to the quotient algebra defines a bounded
map $R_n^H:H\to J_n^{\rm in}$ of norm at most $C'\nu_n$.
It vanishes on $T_a\mathcal I$, hence on $J_H$, and
therefore descends to $r_n$ in (MA-PR.4).
For $h=T_af$ it gives the ordinary block jets of
$\mathcal Mf$.  No individual-jet condition number enters
this restriction estimate.

We next establish the exact lower bound needed for the
entire lift.  On the horizontal edges, (M18.8) gives its
polynomial minimum modulus on the fixed real interval
containing all the new rectangles.  On a vertical inner
edge use the finite-disk Blaschke argument of MA-RS-001,
or the proof of (M18.8), centered at $2+it$.
Choose its fixed radius large enough that this entire
real interval lies in the inner quarter disk.
There are $O(\log\nu_n)$ zeros in that disk.
Condition (MA-PR.2) and the choice of $\delta_0$ bound
each normalized Blaschke factor below by
$c\nu_n^{-2}$.  The zero-free factor is bounded below
by a fixed inverse power of $\nu_n$, by Harnack's
inequality and the Euler lower bound at $2+it$.
The upper and lower Stirling bounds for $\Xi/L$ then give
\[
 \inf_{w\in\partial\Omega_n^{\rm in}}|\Xi(w)|
       \ge c\exp(-\kappa|n|-C\log^2\nu_n),
       \qquad \kappa=\pi[K:\mathbb Q]/4.
\tag{MA-PR.8}
\]
Only the fixed-height windows used in the chosen disk
matter in selecting $\delta_0$ above.  At the finitely
many low heights, all the contours avoid zeros and
compactness supplies the bound after decreasing $c$.
Finite gamma exceptions and possible finite-$L$ poles
are absent from the large-height disk argument;
$\Xi$ itself is entire on every contour.

Fix $\tau>0$.  For a representative
$F\in A(\Omega_n^{\rm in})$, define
\[
 H_n^{\rm in}F(z)
   =\frac1{2\pi i}\int_{\partial\Omega_n^{\rm in}}
      F(w)\frac{\Xi(z)-\Xi(w)}{\Xi(w)(z-w)}
                         e^{\tau(z-w)^2}\,dw .
\tag{MA-PR.9}
\]
The divided difference is entire in $z$.
The Cauchy argument of (M18.10)--(M18.11) gives
$H_n^{\rm in}F=F+\Xi A$ inside the rectangle,
where $A$ is holomorphic there, and divisibility by
$\Xi$ in a neighborhood of every exterior zero.
If the interior jets of $F$ vanish, $F/\Xi$ is
holomorphic on the rectangle.  The two exterior
Cauchy integrals then vanish; entireness implies
$H_n^{\rm in}F=0$ everywhere.
Thus (MA-PR.9) depends only on the class $b\in J_n^{\rm in}$
and has precisely the stated full jet values.

For each fixed real strip $|\operatorname{Re}z|\le A$,
the proof of (M18.12), with (MA-PR.8) in place of
(M18.9), gives
\[
 |H_n^{\rm in}b(\sigma+it)|
    \le C_A\exp(C_A\log^2\nu_n)\|b\|_n
                                   e^{-c_\tau(t-n)^2}.
\tag{MA-PR.10}
\]
To spell out its estimates, when $|z-w|\ge1$, the
ratio of the upper bound for $\Xi(z)$ to (MA-PR.8)
is bounded by $\exp(C_A\log^2\nu_n)$ times a fixed
polynomial in $2+|t|$ and $e^{\kappa|t-n|}$.
The factor $e^{\tau(z-w)^2}$ absorbs that polynomial
and linear exponential into a smaller Gaussian.
The second term of the divided difference has the
same bound.  When $|z-w|<1$, integrate $\Xi'$ along
the segment from $w$ to $z$ and use its fixed-strip
Stirling upper bound; division by (MA-PR.8) again
costs only $\exp(C_A\log^2\nu_n)$.
The perimeter is bounded, and taking the infimum over
representatives gives the displayed quotient-norm estimate.
Inverse Mellin transformation, shifting to the two
lines $\operatorname{Re}z=\pm(N+1)$ as in MT-M18,
now defines $\mathcal T_n^{\rm in}b\in\mathcal E$.
The integrals of $(1+|t|)^j$ against that Gaussian
prove the second estimate in (MA-PR.5).
The jet identities prove its last equality.

The restriction of $i$ to $Q_{\rm in}$ is injective:
MA-WP-001 says that its kernel has all interior jets
zero, whereas membership in $Q_{\rm in}$ says all
exterior jets are zero.  MT-M18's full smooth kernel
then makes the smooth class zero.  Its image is dense.
Indeed MA-WP's bounded projections $P_n$ have densely
spanning finite-dimensional ranges, each with exactly
the interior local factors because there are no
boundary zeros.  The lifts in (MA-PR.9) realize every
one of those finite factors in $i(Q_{\rm in})$.
This proves the density assertion.
\end{proof}

\paragraph{MA-PR-002 (nuclear partial smoothing with actual factorization).}
Let $g$ have compact support and finite Gevrey order $q>1$,
and put $k_g(x)=e^{-ax}g(x)$.
There is a canonical continuous map
$S_g:Q_H\to Q_{\rm in}\subset Q_\infty$, nuclear of every
positive order, such that
\[
             i_{\rm in}S_g=C_g^{Q_H},\qquad
             S_g i_{\rm in}=C_{k_g}^{Q_{\rm in}}.
\tag{MA-PR.11}
\]
It has an actual continuous lift to $\mathcal E$.
The nuclear bounds are uniform for test families with
fixed compact support, fixed Gevrey order and fixed
Gevrey derivative constants.

On the full smooth quotient define
\[
       B_g=S_g i,\qquad E_g=C_{k_g}^{Q_\infty}-B_g .
\tag{MA-PR.12}
\]
Both are nuclear of every positive order.
$B_g$ has exactly the smoothed interior full jets and
zero exterior full jets; $E_g$ has exactly the smoothed
exterior full jets and zero interior full jets.
Their ranges are contained in $Q_{\rm in}$ and
$Q_{\rm out}$, respectively.  In particular
\[
\begin{gathered}
 B_gB_h=B_{g*h},\qquad E_gE_h=E_{g*h},\\
 B_gE_h=E_gB_h=0
\end{gathered}
\tag{MA-PR.13}
\]
for compact finite-order Gevrey $g,h$.
No unsmoothed projection is asserted.

\begin{proof}
In $J_n^{\rm in}$ let $M_{g,n}$ multiply by the class
of $w\mapsto\mathcal Mg(w-a)$.
The elementary Gevrey Fourier estimate (MA-RS.18)--(MA-RS.19),
uniform on the fixed real interval of the contours, gives
\[
 \|M_{g,n}\|\le
   \sup_{\partial\Omega_n^{\rm in}}|\mathcal Mg(w-a)|
       \le b_n:=C e^{-c\nu_n^{1/q}} .
\tag{MA-PR.14}
\]
We use a positive upper bound $b_n$, including when the
actual multiplier is zero.  The constants $C,c>0$ can
be chosen uniformly for the stated bounded Gevrey families.
Define
\[
 \widetilde S_g
    =\sum_n\mathcal T_n^{\rm in}M_{g,n}r_n:
                        Q_H\longrightarrow\mathcal E,
       \qquad S_g=q_\infty\widetilde S_g .
\tag{MA-PR.15}
\]
Equations (MA-PR.5) and (MA-PR.14) show that the series
converges absolutely in every $\mathcal E$ seminorm,
uniformly on the Hilbert unit ball:
\[
 p_{N,j}(\mathcal T_n^{\rm in}M_{g,n}r_n v)
 \le C_{N,j}\nu_n
       e^{C_{N,j}\log^2\nu_n-c\nu_n^{1/q}}\|v\|_{Q_H}.
\tag{MA-PR.16}
\]
This proves continuity and gives the claimed actual lift.
Its exterior full jets vanish termwise and by continuity,
so $S_g$ has range in the closed subspace $Q_{\rm in}$.

Here is a common nuclear series, to control the topology
as well as the scalar trace.  Choose an Auerbach basis
$e_{n,j}$ of $J_n^{\rm in}$ with dual basis $\ell_{n,j}$;
both vectors and functionals have norm one.  Take
$R_n=C_1\nu_n$ above $\|r_n\|$.  Then
\[
\begin{aligned}
 u_{n,j}&=b_n^{1/2}\mathcal T_n^{\rm in}e_{n,j},\\
 \alpha_{n,j}(v)&=
       \frac{\ell_{n,j}(M_{g,n}r_nv)}{b_nR_n},\\
 \lambda_{n,j}&=R_n b_n^{1/2},\qquad
 \widetilde S_gv=\sum_{n,j}
                       \lambda_{n,j}u_{n,j}\alpha_{n,j}(v).
\end{aligned}
\tag{MA-PR.17}
\]
The $\alpha_{n,j}$ have norm at most one on $Q_H$.
The family $(u_{n,j})$ is bounded in every seminorm of
$\mathcal E$, because its negative Gevrey exponential
dominates each $\exp(C_{N,j}\log^2\nu_n)$.
The coefficients, repeated $d_n^{\rm in}=O(\log\nu_n)$
times, belong to $\ell^p$ for every $p>0$.
This proves nuclearity of every positive order, with
one common bounded target disk for a bounded Gevrey
family.  The images of this disk under $q_\infty$ lie
in $Q_{\rm in}$ and are bounded for its subspace
topology.  Thus the same assertion holds with that
space as target.

For every $m$, the jet identities of the lifts give
\[
       r_m i S_g=M_{g,m}r_m.
\tag{MA-PR.18}
\]
Convolution on the actual quotient also satisfies
$r_m C_g^{Q_H}=M_{g,m}r_m$.
This follows first on the dense smooth image, from
the Mellin convolution identity, then on $Q_H$ by
boundedness of both maps.
By MA-WP-001, the intersection of the kernels of
all the $r_m$ is zero: together they are exactly all
the retained interior jets.
Consequently (MA-PR.18) proves
$i S_g=C_g^{Q_H}$ as an equality of actual bounded
operators on all of $Q_H$.
No comparison of traces has occurred.

For a class in $Q_{\rm in}$, $S_g i$ has the same
interior full jets as $C_{k_g}$, and both have zero
exterior full jets.  The full smooth jet kernel of
MT-M18 proves the second identity in (MA-PR.11).
For an arbitrary class in $Q_\infty$, the same argument
proves the asserted jet description of $B_g$.
It also shows that $S_g$ is independent, as a map to
the smooth quotient, of the allowed choices of
contours, representatives and lifts.
Alternatively its first identity uniquely determines
it among maps into $Q_{\rm in}$, because $i_{\rm in}$
is injective.

The series (MA-PR.17) composed with the continuous $i$
proves nuclearity of every positive order of $B_g$.
The full smooth convolution $C_{k_g}^{Q_\infty}$
has that property by the arithmetic block expansion
of MT-M18--M19.  Their difference $E_g$ therefore
has it too.  Its jet description follows from
subtraction, and gives its stated range.
All identities in (MA-PR.13) follow by multiplying
the full local Taylor classes, using
$\mathcal M(g*h)=\mathcal Mg\,\mathcal Mh$ and
$k_g*k_h=k_{g*h}$.  Equality of all full jets is
equality in the actual smooth quotient.
These identities are for the constructed operators;
they do not presuppose a direct-sum topology
$Q_\infty=Q_{\rm in}\oplus Q_{\rm out}$.
\end{proof}

The construction also gives the genuine covariance
$S_gU_t^{Q_H}=e^{-at}U_t^{Q_\infty}S_g$.
On every retained local algebra both sides multiply
by $e^{(\rho-a+z)t}\mathcal Mg(\rho-a+z)$, with the
full Taylor class understood.  The exterior jets
are zero.  This is another equality of actual maps,
by the full jet kernel.

\paragraph{MA-PR-003 (partial Meyer transport with an operator correction).}
The traces of the actual operators satisfy
\[
 \operatorname{Tr}_{Q_H}C_g
    =\operatorname{Tr}_{Q_\infty}B_g
    =\operatorname{Tr}_{Q_\infty}C_{k_g}
                         -\operatorname{Tr}_{Q_\infty}E_g .
\tag{MA-PR.19}
\]
Under the fixed-sector isomorphism of MT-M22/MT-M24,
the middle and last terms are traces of the corresponding
actual operators on Meyer's negative smooth quotient.
Thus the Hilbert trace has been transported to a
geometric smooth representation, with the discarded
part represented by the actual nuclear operator $E_g$.

In the rational trivial sector, specialize to $a=1/2$,
$u=-\alpha$, $v=\alpha$ satisfying (MA-PR.2).
Let $W_S(g)$ be the actual local trace in (MA-MEY.15),
and retain $B_{\overline S}(g)$ and $P(g)$ from MA-RS.
Then
\[
 \operatorname{Tr}_{Q_H}C_g
      =P(g)-W_S(g)-B_{\overline S}(g)
                            -\operatorname{Tr}_{Q_\infty}E_g .
\tag{MA-PR.20}
\]
For $g(0)=0$ the literal cutoff limit of MA-WP-003 is
\[
\begin{aligned}
 \lim_R\operatorname{Tr}(M_{\chi_R}C_gP_G)
    &=W_S(g)+B_{\overline S}(g)-P(g)
                           +\operatorname{Tr}_{Q_\infty}E_g\\
    &=A_0(g)+B(g)-P(g)
                           +\operatorname{Tr}_{Q_\infty}E_g .
\end{aligned}
\tag{MA-PR.21}
\]
The transport to (MA-PR.21) uses actual operator
factorization, trace cyclicity and Meyer's actual
common-space geometry.  After this transport has
been established, the finite-block trace calculation
identifies its correction as
\[
 \operatorname{Tr}_{Q_\infty}E_g
       =\sum_{\rho\in Z_{\rm out}}
                         m_\rho\mathcal Mg(\rho-a).
\tag{MA-PR.22}
\]
This last scalar description is not used in the
construction or in the trace cyclicity step.

\begin{proof}
Pass the target vectors of (MA-PR.17) to $Q_\infty$.
Composing that series with the continuous $i$ gives
a trace-class rank-one series on the Hilbert space.
Composing the source functionals with $i$ gives a
nuclear rank-one series for $B_g$ on $Q_\infty$.
The two rank-one traces are termwise equal:
\[
 \lambda_{n,j}\alpha_{n,j}
                       (i q_\infty u_{n,j}).
\tag{MA-PR.23}
\]
Their absolute summability follows from boundedness
of the target disk, continuity of $i$, the unit
functional bound and $\sum|\lambda_{n,j}|<\infty$.
The smooth quotient has the approximation property
directly from the absolute finite-block synthesis
of MT-M18.  Equivalently the order-zero nuclear
trace is the same common series, as in MA-HM-002.
This proves the first equality of (MA-PR.19) with
the actual identity $iS_g=C_g^{Q_H}$.
Linearity of the nuclear trace and (MA-PR.12)
prove its second equality.

The MT-M22/MT-M24 map is an actual continuous
equivariant isomorphism to the fixed compact-character
sector of Meyer's negative quotient.
The normalization check in MA-HM-002, especially
(MA-HM.15) and its proof, carries $C_{k_g}$ to the
integrated Meyer operator with test
$f(ca(x))=\beta(c)^{-1}e^{-ax}g(x)$, with compact
group Haar measure of mass one.
It also carries $B_g$ and $E_g$ to their actual
conjugate operators and preserves every rank-one
pairing in (MA-PR.23).
For the rational specialization, Meyer's source
trace additivity (MA-MEY.15)--(MA-MEY.16) expresses
the full smooth trace as $P-W_S-B_{\overline S}$.
Substitution into the operator identity (MA-PR.19),
not substitution of a zero sum, proves (MA-PR.20).

MA-WP-003 proves trace class of $C_gP_Z$ on the
specified weighted Hilbert space and, in the stated
cutoff order, identifies the cutoff limit with the
negative of the actual quotient trace when $g(0)=0$.
Applying that result to (MA-PR.20) proves the first
line of (MA-PR.21).
The checked rational local-place normalization
$W_S+B_{\overline S}=A_0+B$ gives its second line.
Only the spectral product against a cutoff is
asserted to converge in trace norm; the geometric
cutoff operators themselves need not do so.

Finally $E_g$ is zero on the retained local factors
and is multiplication by the full Taylor class of
$\mathcal Mg(w-a)$ on each exterior factor.
The actual finite-block trace is therefore
$m_\rho\mathcal Mg(\rho-a)$ on an exterior factor
and zero on an interior factor.
The absolute block nuclear series just used, or
the full and partial series of MA-WP, permits
summation of these finite-dimensional traces.
The local zero count and vertical Gevrey decay
also prove absolute convergence directly.
This proves (MA-PR.22).
\end{proof}

\paragraph{MA-PR-004 (sharp obstruction at a truncated boundary block).}
For this paragraph do not assume (MA-PR.2).
Suppose a zero $\rho$ is on an edge of the Hilbert strip and
MA-WP retains a proper nonzero part of its multiplicity:
\[
       m=m_\rho,\qquad r=r_\rho,\qquad 0<r<m,\qquad
                         \lambda=\rho-a .
\tag{MA-PR.24}
\]
The actual smooth and Hilbert primary factors at this point
are, with their translation actions,
\[
 A=\mathbb C[z]/(z^m),\qquad B=\mathbb C[z]/(z^r),
 \qquad \pi_{m,r}:A\longrightarrow B .
\tag{MA-PR.25}
\]
The restriction of $i$ is exactly this quotient map.
For a test $g$ write $G(z)=\mathcal Mg(z)$.
A translation-equivariant return of this primary factor
which factors the actual convolution must satisfy
\[
 L:B\longrightarrow A,\qquad
       \pi_{m,r}L=M_{G(\lambda+z)}^B .
\tag{MA-PR.26}
\]
Here equivariance includes the centering factor
$e^{-at}$, as for $S_g$ above.
Such a linear equivariant $L$ exists if and only if
\[
     G^{(j)}(\lambda)=0
           \quad(0\le j<\min(m-r,r)).
\tag{MA-PR.27}
\]
This is a sharp finite-block condition.  In particular,
if $G(\lambda)\ne0$, no such return exists, even before
continuity is required.

Every prescribed boundary point admits a compact
finite-order Gevrey test with $g(0)=0$ and
$G(\lambda)=1$.  Thus no translation-equivariant
return can factor convolution for all these tests
at a proper truncated boundary factor.  In particular
this factor cannot be realized as a smooth invariant
subrepresentation on which $i$ is an isomorphism.
The assertion is conditional on such a truncated
boundary factor actually occurring; it does not
assert a multiple zero of any particular $L$-function.

\begin{proof}
MA-WP-001 identifies the actual Hilbert factor with
$B$, with translation $e^{(\lambda+z)t}$.
The full smooth factor has translation
$e^{(\rho+z)t}$ by MT-M18.
After the centering factor is removed, equivariance
is
\[
                      L e^{tz}=e^{tz}L\quad(t\in\mathbb R).
\tag{MA-PR.28}
\]
These are finite polynomial matrices in $t$, so
differentiation at zero gives $Lz=zL$.
Equivalently $L$ is a $\mathbb C[z]$-module map.
It is determined by $c(z)=L(1)\in A$, and its
well-definedness on $B$ is exactly $z^r c=0$ in $A$.
Consequently
\[
\begin{aligned}
 \operatorname{Hom}_{\mathbb C[z]}(B,A)
       &\simeq \operatorname{Ann}_A(z^r)=z^{m-r}A,\\
 \pi_{m,r}\operatorname{Ann}_A(z^r)
       &=z^{m-r}B .
\end{aligned}
\tag{MA-PR.29}
\]
If $m-r\ge r$, the ideal on the second line is zero.
Otherwise it consists precisely of the classes
whose coefficients of degrees $0,\ldots,m-r-1$
vanish.  Equation (MA-PR.26), applied to $1$,
therefore has a solution exactly when
$G(\lambda+z)\bmod z^r$ belongs to that ideal.
This is (MA-PR.27).

For the converse, when $m-r<r$ take as $c(z)$ the
degree-less-than-$r$ Taylor polynomial of
$G(\lambda+z)$.  Under (MA-PR.27) it is divisible
by $z^{m-r}$, so $z^r c=0$ in $A$ and
$L(b)=bc$ gives the required module map.
When $m-r\ge r$, the condition makes the multiplier
on $B$ zero and $L=0$ suffices.
More generally all choices are
\[
              c=c_0+z^{\max(r,m-r)}A,
\tag{MA-PR.30}
\]
where $c_0$ is any one solution, because the difference
must lie in both $z^{m-r}A$ and $\ker\pi_{m,r}=z^r A$.
This proves the sharpness assertion without a
semisimplicity assumption.

To see why the finite-block condition is necessary
for a proposed global return, embed the actual
Hilbert factor $B$ in its finite block using the
primary idempotent of MA-WP, and then project the
smooth output to its full $\rho$-jet factor.
The finite-dimensional jet maps are continuous.
A globally equivariant return yields (MA-PR.28)
on these factors, and the actual identity
$iS=C_g^{Q_H}$ yields (MA-PR.26).
Thus no other block can remove this obstruction.
Conversely the finite solution just constructed
can be represented in the smooth quotient by
MT-M18's finite-block interpolation.  This converse
is local to the one primary factor; it does not
assert estimates for a return on an infinite
collection of blocks.

Choose a compact Gevrey bump $\eta$ supported in
$(1,2)$ with $\int\eta=1$, and put
$g(x)=e^{-\lambda x}\eta(x)$.
Multiplication by this analytic exponential
preserves its Gevrey order on the compact support.
Then $g(0)=0$ and $\mathcal Mg(\lambda)=1$.
The condition (MA-PR.27) fails.
Finally an invariant smooth realization of $B$
with $i$ an isomorphism would give a module section
of $\pi_{m,r}$, corresponding to the unit multiplier
in (MA-PR.26), which has just been excluded.
\end{proof}

\paragraph{MA-PR-005 (Gaussian-regularized actual returns for all smooth tests).}
Return to the generic-edge hypothesis (MA-PR.2).
For every $g\in\mathcal E$, put
\[
 \gamma_\varepsilon(x)
       =(4\pi\varepsilon)^{-1/2}e^{-x^2/(4\varepsilon)},
 \qquad g_\varepsilon=g*\gamma_\varepsilon,\qquad
                \mathcal M\gamma_\varepsilon(z)=e^{\varepsilon z^2}.
\tag{MA-PR.31}
\]
One may take $h_\varepsilon=g_\varepsilon$.
If $g(0)=0$, fix once and for all a compact finite-order
Gevrey function $b$ with $b(0)=1$, and instead take
\[
       h_\varepsilon
            =g_\varepsilon-g_\varepsilon(0)b .
\tag{MA-PR.32}
\]
Then $h_\varepsilon(0)=0$ and in either case
$h_\varepsilon\to g$ in $\mathcal E$.

For every fixed $\varepsilon>0$, (MA-PR.15) defines
an actual continuous return
$S_\varepsilon:Q_H\to Q_{\rm in}$, nuclear of every
positive order, with
\[
\begin{aligned}
 iS_\varepsilon&=C_{h_\varepsilon}^{Q_H},\\
 B_\varepsilon&=S_\varepsilon i,\qquad
 E_\varepsilon
   =C_{e^{-a\cdot}h_\varepsilon}^{Q_\infty}-B_\varepsilon .
\end{aligned}
\tag{MA-PR.33}
\]
The actual Hilbert operators satisfy
\[
       C_{h_\varepsilon}^{Q_H}
                \longrightarrow C_g^{Q_H}
                           \quad\hbox{in trace norm}.
\tag{MA-PR.34}
\]
The full smooth trace and the discarded correction traces
have the separate limits
\[
\begin{aligned}
 \operatorname{Tr}_{Q_\infty}
       C_{e^{-a\cdot}h_\varepsilon}
     &\longrightarrow
       \operatorname{Tr}_{Q_\infty}C_{e^{-a\cdot}g},\\
 \operatorname{Tr}_{Q_\infty}E_\varepsilon
     &\longrightarrow
        \operatorname{Tr}_{Q_\infty}C_{e^{-a\cdot}g}
                     -\operatorname{Tr}_{Q_H}C_g .
\end{aligned}
\tag{MA-PR.35}
\]
No convergence of $S_\varepsilon$, $B_\varepsilon$
or $E_\varepsilon$ as maps into $Q_\infty$ is
asserted for general $g$.

In particular, for every compact smooth $g$ with $g(0)=0$,
the rational centered geometric cutoff limit at this
generic width has the actual-operator transport
\[
 \lim_R\operatorname{Tr}(M_{\chi_R}C_gP_G)
     =A_0(g)+B(g)-P(g)
             +\lim_{\varepsilon\downarrow0}
                         \operatorname{Tr}_{Q_\infty}E_\varepsilon .
\tag{MA-PR.36}
\]
Its correction is the trace limit of the constructed
nuclear discarded operators.  As an additional
evaluation of that limit, for every $g\in\mathcal E$,
\[
 \lim_{\varepsilon\downarrow0}
             \operatorname{Tr}_{Q_\infty}E_\varepsilon
       =\sum_{\rho\in Z_{\rm out}}
                              m_\rho\mathcal Mg(\rho-a).
\tag{MA-PR.37}
\]

\begin{proof}
The Gaussian belongs to $\mathcal E$ for each
$\varepsilon>0$, and its Mellin transform in
(MA-PR.31) follows by the Gaussian integral and
analytic continuation in $z$.
The convolution estimates of MT-M19 give
$g_\varepsilon\in\mathcal E$ and
\[
 \mathcal Mg_\varepsilon(z)
       =\mathcal Mg(z)e^{\varepsilon z^2}.
\tag{MA-PR.38}
\]
Here the actual multiplier in the uncentered
$w$-block is
$\mathcal Mg(w-a)e^{\varepsilon(w-a)^2}$.
This fixes the centering in the regularization.

For completeness the approximation takes place in
every defining seminorm, not merely in $L^2$.
The fundamental theorem of calculus gives
\[
 p_{N,j}(U_y g-g)
       \le |y|e^{N|y|}p_{N,j+1}(g).
\tag{MA-PR.39}
\]
Integrate against $\gamma_\varepsilon(y)\,dy$.
After $y=\sqrt{\varepsilon}\,t$, its right-hand
moment is at most $C_N\sqrt{\varepsilon}$ for
$0<\varepsilon\le1$.
Thus
\[
 p_{N,j}(g_\varepsilon-g)
       \le C_N\sqrt{\varepsilon}\,p_{N,j+1}(g).
\tag{MA-PR.40}
\]
When $g(0)=0$, this also proves
$g_\varepsilon(0)\to0$, so subtraction of the
fixed $b$ in (MA-PR.32) preserves convergence
in every $\mathcal E$ seminorm.

The only use of compact Gevrey regularity in the
proof of MA-PR-002 was the multiplier estimate
(MA-PR.14).  The remainder of that proof applies
to any $h\in\mathcal E$ whose contour multipliers
are bounded by $C e^{-c\nu_n^\theta}$ for some
$\theta>0$: (MA-PR.16) is summable in every
seminorm, the square-root construction (MA-PR.17)
gives the common nuclear series, and the actual
jet identities give (MA-PR.18).
For $g_\varepsilon$, (MA-PR.38) on the fixed
real strip gives precisely such a bound:
\[
 \sup_{\partial\Omega_n^{\rm in}}
       |\mathcal Mg_\varepsilon(w-a)|
            \le C_\varepsilon e^{-c_\varepsilon\nu_n^2}.
\tag{MA-PR.41}
\]
The unregularized transform $\mathcal Mg$ is
bounded on that fixed strip, and the real part
of $\varepsilon(w-a)^2$ is
$\varepsilon((\operatorname{Re}w-a)^2-(\operatorname{Im}w)^2)$.
This proves (MA-PR.41) with adjusted constants
for the finitely many low blocks.
If (MA-PR.32) is used, the additional compact
Gevrey transform has a bound
$C e^{-c\nu_n^{1/q}}$.
The sum still has the required positive-power
exponential decay for each fixed $\varepsilon$.
No uniform estimate as $\varepsilon\downarrow0$
is needed or claimed for these return maps.
The same construction therefore proves
(MA-PR.33), together with its actual full-jet
descriptions and nuclear trace cyclicity.

We now pass limits in two different spaces.
MA-WP.14--MA-WP.15 gives a continuous map
$\mathcal E\to\mathcal S_1(Q_H)$.
More explicitly, summing its rank-times-operator-norm
bound after a sufficiently strong vertical-decay
estimate bounds
$\|C_f^{Q_H}\|_1$ by finitely many
$\mathcal E$ seminorms of $f$.
Apply this to $f=h_\varepsilon-g$ and (MA-PR.40).
This proves the Hilbert trace-norm convergence
(MA-PR.34).

For the full smooth term use the fixed block
series of MT-M19, (M19.3)--(M19.4), on
$Q_\infty\simeq s(J)$.
Its Auerbach coefficient functionals are
equicontinuous, its coefficients
$(2+|n|)^{-2}$ are summable with the block
dimensions, and the target vectors depend
continuously on the test in every bounded-family
seminorm.
For the difference $f=e^{-a\cdot}(h_\varepsilon-g)$,
the sum of absolute rank-one trace pairings is
at most
\[
       \sum_n d_n\,
          \sup_{w\in\partial\Omega_n}|\mathcal Mf(w)|
                \longrightarrow0 .
\tag{MA-PR.42}
\]
Indeed integration by parts on the fixed full
strip bounds the supremum by $\nu_n^{-2}$
times finitely many seminorms of $f$, and
$d_n=O(\log\nu_n)$.
This proves the first limit in (MA-PR.35)
from the actual smooth nuclear series.

For each fixed $\varepsilon$, the common series
of the already constructed return gives
\[
 \operatorname{Tr}_{Q_H}C_{h_\varepsilon}
    =\operatorname{Tr}_{Q_\infty}
            C_{e^{-a\cdot}h_\varepsilon}
                   -\operatorname{Tr}_{Q_\infty}E_\varepsilon .
\tag{MA-PR.43}
\]
The first trace converges by (MA-PR.34), and
the full smooth trace converges by (MA-PR.42).
Their difference proves the second limit of
(MA-PR.35).  It is a statement about correction
traces, without an operator convergence claim
for $E_\varepsilon$ on the smooth quotient.

For the rational centered compact test $g$,
apply Meyer's common-space formula to the
unregularized full smooth operator
$C_{e^{-x/2}g}$, whose trace is the limit of the
full smooth terms just proved.
It gives $P-W_S-B_{\overline S}$ as before.
Apply the actual MA-WP-003 ambient cutoff
theorem to $g$ itself, giving the negative
of the Hilbert trace.
Equations (MA-PR.35) and these two actual
identifications prove (MA-PR.36).
This argument does not require a new
common-space geometric formula or a compact
support assertion for the Gaussian-regularized
tests.

Lastly the constructed $E_\varepsilon$ has
zero retained jets and the full exterior
Taylor multiplier of $h_\varepsilon$.
The finite-block calculation of MA-PR-003
therefore evaluates its trace as the
exterior divisor sum for $h_\varepsilon$.
Convergence $h_\varepsilon\to g$ in
$\mathcal E$, the local zero count and
rapid vertical Mellin decay give convergence
of these sums absolutely.
This proves (MA-PR.37) after the actual
operator and trace limits have been established.
\end{proof}

\paragraph{Scope, boundary and the remaining independent-method step.}
The scope below records MA-PR-001--005.
The later MA-PR-006--008 globalize the local boundary
converse: every locally eligible test admits actual
adaptive regularized returns at every width, while
an unregularized continuous return is characterized
by the rapid decay of its attained finite-block costs.
Those later results retain the separate original
common-space off-diagonal domain question.

The maps (MA-PR.11)--(MA-PR.12) are unconditional under
the explicit generic-edge condition and Gevrey test
hypothesis.  They apply at every fixed polynomial
weight $s\ge0$, including the requested $s\ge3$.
No boundary jets occur in those unregularized maps
because (MA-PR.2) excludes boundary zeros.
The exact boundary truncations at other widths
remain those of MA-WP.  MA-PR-004 proves a sharp
local obstruction to an equivariant return when
a nonzero proper part of a boundary block survives;
it gives the local converse for tests satisfying
the necessary vanishing condition, without
asserting a global return at exceptional widths.
MA-PR-005 gives a family of actual regularized
returns for arbitrary $\mathcal E$ tests at
generic widths, with Hilbert trace-norm convergence
and correction-trace convergence.  It does not
remove the Gevrey hypothesis for an unregularized
map $S_g$ into $Q_\infty$.

If exterior zeros occur, $S_gi$ is the retained
operator $B_g$, not the full convolution $C_{k_g}$.
MA-HM-003 excludes the latter identity for a test
that detects a discarded block.
Nothing above proves an unsmoothed continuous
spectral selector, nor a full return on the same
partial completion.  The identity (MA-PR.21)
therefore retains its actual discarded-divisor
operator.  Removing that term for every test is
the exact criterion of MA-EG-001.

The paired regularized Euler diagonals of MA-RS
carry the same discarded-divisor trace in their
contour correction.  That agreement does not
identify the unbounded inverse composition there
with every off-diagonal map in Meyer's original
common-space projection.  The factorization
proved here is the explicit map (MA-PR.15) through
the actual retained smooth subrepresentation;
these are distinct domain statements.

The Grothendieck placement concerns actual finite-length
$\mathbb C[z]$-modules at the retained and discarded
spectral points.  Their local lengths and underlying
finite-dimensional complex endomorphisms enter the
nuclear traces in (MA-PR.19)--(MA-PR.23).
This is the same inherited cote~1, scans~079--080
and~111--112, placement as MA-HM and MA-WP.
The analytic factorization does not construct an
arithmetic positivity or weight realization.

\paragraph{Exact inputs and proof scope.}
The arithmetic estimates are the unconditional local
zero count and fixed-strip finite-$L$ bounds of NT-Z10,
and the entire lift, selected heights, Stirling bounds
and full smooth zero-jet kernel of MT-M18,
(M18.5)--(M18.13).  The new inner-contour lower bound
and all restriction and nuclear estimates are proved
above.  MA-WP-001 supplies the already proved literal
Hilbert closure and its finite retained blocks;
MA-WP-003 supplies the actual ambient cutoff limit.
The Meyer identifications and common-space signs
are exactly MT-M22/MT-M24, MA-HM-002 and
MA-MEY.15--MA-MEY.16.  No A, M, or new outside
preprint is an input to this generic partial return.
The boundary obstruction uses only the exact
local modules of MA-WP and elementary finite
$\mathbb C[z]$-module algebra.  The Gaussian
extension uses the explicit approximate-identity
estimate (MA-PR.40), the literal Hilbert
trace-norm continuity in MA-WP, and the full
smooth rank-one continuity in MT-M19.

\section{Exceptional-width returns with the actual finite-block costs}
\label{sec:exceptional-partial-return}

The local converse of MA-PR-004 can be assembled globally after
an adaptive regularization.  For an unregularized test there is
an exact additional topological criterion.  Both statements use
the original adapted block norms and the actual Hilbert quotient.
No polynomial bound for a boundary selector, a reciprocal edge
distance, or an individual interpolation matrix is assumed.

Retain (MA-PR.1), with arbitrary \(u<v\), \(a\in\mathbb R\)
and \(s\ge0\), but do not assume (MA-PR.2).
The multiplicities \(r_\rho\) are exactly (MA-WP.3).
Write \(\nu_n=2+|n|\).  The full smooth blocks \(J_n\),
their holomorphic quotient norms \(\|\cdot\|_n\), the rectangles
\(\Omega_n\), and the entire lifts \(\mathcal T_n\) are those
of MT-M18, without changing their norms.
In particular
\[
 \begin{aligned}
 Q_\infty&\simeq s(J),&
 q_k((b_n))&=\sup_n\nu_n^k\|b_n\|_n,\\
 p_{N,j}(\mathcal T_n b)&\le C_{N,j}\nu_n^{D_{N,j}}\|b\|_n .
 \end{aligned}
 \tag{MA-PR.44}
\]
Let \(\mathfrak R_n:Q_\infty\to J_n\) be the full block
restriction.  The actual Hilbert block and its algebraic
coordinate space are
\[
 V_n=P_nQ_H,\qquad
 B_n=\prod_{\substack{\rho\in\Omega_n\\r_\rho>0}}
                      \mathbb C[z]/(z^{r_\rho}),\qquad
 \pi_n:J_n\longrightarrow B_n .
 \tag{MA-PR.45}
\]
Here \(P_n\) is the bounded convolution projection of
(MA-WP.11)--(MA-WP.12).  It need not be orthogonal.
The finite algebra \(B_n\) in (MA-PR.45) is not assigned an
independent Hilbert norm in what follows.

The continuous retained-jet map on the \emph{whole} Hilbert
quotient is denoted by \(\mathscr R_n:Q_H\to B_n\).
For \(v=\pi h\), its coordinates are
\[
 (\mathscr R_nv)_{\rho,j}
    =\frac1{j!}\int_{\mathbb R}x^j e^{(\rho-a)x}h(x)\,dx,
       \qquad 0\le j<r_\rho .
 \tag{MA-PR.46}
\]
The integral is absolutely convergent by (MA-WP.17).
These functionals annihilate \(J_H\), so (MA-PR.46) is
independent of the representative \(h\).
The exact block description in MA-WP gives
\[
 \mathscr R_nP_n=\mathscr R_n,\qquad
 \mathscr R_n|_{V_n}:V_n\xrightarrow{\ \sim\ }B_n,\qquad
 \mathscr R_ni=\pi_n\mathfrak R_n .
 \tag{MA-PR.47}
\]
All three are identities of actual continuous maps.

For \(g\in\mathcal E\), put \(G=\mathcal Mg\) and
\[
 b_n(g)=R_n(T_{-a}g)\in J_n,\qquad
 b_n(g)(w)=G(w-a)\quad\hbox{as a full Taylor class}.
 \tag{MA-PR.48}
\]
Call \(g\) \emph{eligible} if at every proper nonzero
boundary truncation it satisfies
\[
 G^{(j)}(\rho-a)=0
       \quad\bigl(0\le j<\min(m_\rho-r_\rho,r_\rho)\bigr),
       \qquad 0<r_\rho<m_\rho .
 \tag{MA-PR.49}
\]
This defines a closed linear subspace \(\mathcal E_{\rm el}\)
of \(\mathcal E\).  No condition is imposed at a factor
with \(r_\rho=0\), or with \(r_\rho=m_\rho\).
Whether any particular boundary or exterior factor occurs
is not asserted.

\paragraph{MA-PR-006 (an exact global lifting criterion).}
For eligible \(g\), let
\[
 \begin{aligned}
 \mathcal F_n(g)
   =\{L\in\operatorname{Hom}_{\mathbb C[w]}(B_n,J_n):
                        \pi_nL=M_{\pi_nb_n(g)}\},\\
 c_n(g)=\min_{L\in\mathcal F_n(g)}
                \|L\mathscr R_n\|_{\mathcal L(Q_H,J_n)} .
 \end{aligned}
 \tag{MA-PR.50}
\]
For \(B_n=0\), set \(c_n(g)=0\).
Every nonempty minimum in (MA-PR.50) is attained.
These are finite-dimensional quantities with the actual source
norm; they do not discard the norm of an oblique projection.

There is a continuous linear map \(S_g:Q_H\to Q_\infty\)
with
\[
 iS_g=C_g^{Q_H},\qquad
 S_gU_t^{Q_H}=e^{-at}U_t^{Q_\infty}S_g
                  \quad(t\in\mathbb R)
 \tag{MA-PR.51}
\]
if and only if \(g\) is eligible and
\[
                  \sup_n\nu_n^k c_n(g)<\infty
                         \quad\hbox{for every integer }k\ge0 .
 \tag{MA-PR.52}
\]
Every continuous map satisfying (MA-PR.51) is nuclear of
every positive order and has a continuous linear lift
from \(Q_H\) to \(\mathcal E\).
There is no claim of uniqueness at a truncated boundary factor.

The costs in (MA-PR.50) admit an explicit Gram description.
Use the Taylor coordinates of (MA-PR.46) and their ordinary
Euclidean scalar product solely for this formula.  Then
\[
 \begin{aligned}
 K_n&=\mathscr R_n\mathscr R_n^*>0,\\
 (K_n)_{(\rho,j),(\sigma,k)}
   &=\int_{\mathbb R}
       \frac{x^{j+k}e^{(\rho-a+\overline{\sigma}-a)x}}
              {j!\,k!\,w_{u,v,s}(x)}\,dx,\\
 c_n(g)&=\min_{L\in\mathcal F_n(g)}
             \|LK_n^{1/2}\|_{\mathcal L(\ell^2,J_n)} .
 \end{aligned}
 \tag{MA-PR.53}
\]
The target in the last line retains its original
holomorphic quotient norm.

\begin{proof}
\emph{The finite module problem and the actual norms.}
At a primary factor, write
\(A=\mathbb C[z]/(z^m)\) and \(B=\mathbb C[z]/(z^r)\).
A module map \(L:B\to A\) is determined by \(c=L(1)\)
and is well-defined exactly when
\[
        c\in\operatorname{Ann}_A(z^r)=z^{m-r}A .
 \tag{MA-PR.54}
\]
Its reduction is multiplication by \(G(\rho-a+z)\)
if and only if \(c\equiv G(\rho-a+z)\pmod {z^r}\).
This is solvable exactly under (MA-PR.49), as follows
by reducing \(z^{m-r}A\) modulo \(z^r\).
For \(r=m\) the map is forced to be multiplication by
the full class; for \(r=0\) there is no source factor.
Module maps between primary factors at distinct \(\rho\)
are zero, by the coprimality of their annihilating
polynomials.  Thus (MA-PR.49) is precisely the
nonemptiness of every \(\mathcal F_n(g)\).

The maps \(\mathscr R_n\) are surjective because their
restriction to \(V_n\) is an isomorphism.  Consequently
\(K_n\) is positive definite when \(B_n\ne0\).
The same Gram matrix can be computed on \(H\):
the quotient map \(\pi:H\to Q_H\) is a coisometry and
the retained functionals annihilate its kernel.
Weighted Hilbert duality then gives the integral in
(MA-PR.53).  Absolute convergence of its entries follows
from Cauchy--Schwarz applied to the two continuous jets.
The polar decomposition
\(\mathscr R_n=K_n^{1/2}W_n\) has \(W_nW_n^*=1\).
Its coisometry carries the unit ball onto the Euclidean
unit ball, so
\[
                  \|L\mathscr R_n\|=\|LK_n^{1/2}\|.
 \tag{MA-PR.55}
\]
This remains true for the Banach target \(J_n\).
In particular (MA-PR.53) computes the norm on all of
\(Q_H\), not just the inherited norm on \(V_n\).

The feasible set is a closed affine subspace of the
finite-dimensional space of module maps.
Since \(K_n^{1/2}\) is invertible, the right-hand side
of (MA-PR.55) is a genuine norm on that space.
A minimizing sequence is therefore bounded, has a
convergent subsequence, and its limit is feasible.
This proves attainment and the exact Gram formula.

\emph{Necessity of the global condition.}
Suppose that \(S_g\) is continuous and satisfies
(MA-PR.51).  Put \(F_n=\mathfrak R_nS_g\).
For every \(k\), continuity into (MA-PR.44) gives
\[
             \sup_n\nu_n^k\|F_n\|_{\mathcal L(Q_H,J_n)}
                           <\infty .
 \tag{MA-PR.56}
\]
We must verify that these maps factor through the
actual \(P_n\); orthogonality is not being assumed.
On each finite-dimensional \(V_m\), equivariance
intertwines the matrices \(e^{t(w-a)}\).
Differentiating these finite matrix identities shows
that \(F_n|_{V_m}\) is a module map.
For \(m\ne n\), the two annihilating polynomials have
disjoint primary support, so this map is zero.
For \(m=n\), use the isomorphism in (MA-PR.47)
to write \(F_n|_{V_n}=L_n\mathscr R_n|_{V_n}\).
The algebraic span of the \(V_m\) is dense by
MA-WP-001.  Boundedness of \(F_n\), \(P_n\), and
\(\mathscr R_n\) therefore proves on the whole quotient
\[
                  F_n=F_nP_n=L_n\mathscr R_n .
 \tag{MA-PR.57}
\]
Applying the retained jets to \(iS_g=C_g^{Q_H}\)
and using (MA-PR.47) gives
\(\pi_nL_n=M_{\pi_nb_n(g)}\).
The finite module calculation proves eligibility,
and \(c_n(g)\le\|F_n\|\) proves (MA-PR.52).
This proves necessity without replacing \(P_n\)
by an orthogonal projection.

\emph{Sufficiency and the actual lift.}
Now assume eligibility and (MA-PR.52).
Choose a minimizer \(L_n\) in each block and set
\(F_n=L_n\mathscr R_n=F_nP_n\).
Define
\[
       \widetilde S_gv=\sum_n\mathcal T_nF_nv
              \ \in\mathcal E,\qquad
       S_g=q_\infty\widetilde S_g .
 \tag{MA-PR.58}
\]
The polynomial bounds in (MA-PR.44) and the rapid
decay (MA-PR.52) make this series absolutely convergent
in every defining seminorm, uniformly for \(\|v\|\le1\).
It is a continuous linear lift.
The full jet of its output in block \(n\) is \(F_nv\).
Consequently its retained jets are
\(\pi_nF_nv=M_{\pi_nb_n(g)}\mathscr R_nv\),
which are exactly the retained jets of \(C_g^{Q_H}v\).
The common retained-jet kernel on \(Q_H\) is zero
by (MA-WP.16) and its complete block description.
Hence \(iS_g=C_g^{Q_H}\) as an actual bounded-operator
identity.  The module identities for each \(L_n\)
and the full smooth jet kernel prove the equivariance
in (MA-PR.51).

For the topology of nuclearity, choose an Auerbach
basis \(e_{n,j}\), \(\ell_{n,j}\) of \(J_n\).
Omit blocks for which \(F_n=0\), and put
\[
 \begin{aligned}
 u_{n,j}&=c_n(g)^{1/2}\mathcal T_ne_{n,j},&
 \lambda_{n,j}&=c_n(g)^{1/2},\\
 \alpha_{n,j}(v)&=\ell_{n,j}(F_nv)/c_n(g),&
 \widetilde S_gv&=\sum_{n,j}
                 \lambda_{n,j}u_{n,j}\alpha_{n,j}(v).
 \end{aligned}
 \tag{MA-PR.59}
\]
The functionals have norm at most one.
The target vectors form one bounded set in \(\mathcal E\),
by (MA-PR.44) and rapid decay of \(c_n(g)^{1/2}\).
The coefficients, repeated \(d_n=O(\log\nu_n)\) times,
belong to every \(\ell^p\), \(p>0\).
This proves nuclearity of every positive order with
a common bounded target disk.
For an arbitrary continuous return, use its actual
norms \(\|F_n\|\) from (MA-PR.56) in the same argument.
The resulting lifted series has the same full jets
as that return, so gives that return after \(q_\infty\).
This proves the last assertion for every continuous
solution, not just the chosen minimizers.
\end{proof}

\paragraph{MA-PR-007 (adaptive returns at every width).}
For these fixed Hilbert parameters there are continuous
linear regularizers
\(\mathcal D_N:\mathcal E\to\mathcal E\), chosen
independently of the test, with the following properties.
They preserve \(\mathcal E_{\rm el}\), preserve \(g(0)\),
and satisfy \(\mathcal D_Ng\to g\) in \(\mathcal E\),
uniformly on bounded subsets.
For every eligible \(g\), the test
\(h_{N,g}=\mathcal D_Ng\) has an actual return
\[
 \begin{aligned}
 S_{N,g}:Q_H&\longrightarrow Q_\infty,&
 iS_{N,g}&=C_{h_{N,g}}^{Q_H},\\
 B_{N,g}&=S_{N,g}i,&
 E_{N,g}&=C_{T_{-a}h_{N,g}}^{Q_\infty}-B_{N,g}.
 \end{aligned}
 \tag{MA-PR.60}
\]
Each return has an actual \(\mathcal E\)-valued lift
and is nuclear of every positive order.
For fixed \(N\) these nuclear bounds are uniform for
eligible \(g\) in a bounded subset of \(\mathcal E\).
The construction is linear in \(g\).
The tests \(h_{N,g}\) need not be compactly supported.
No bound uniform in \(N\) for the return maps is claimed.
Eligibility is essential even for approximation: since
\(\mathcal E_{\rm el}\) is closed, no ineligible \(g\)
can be an \(\mathcal E\)-limit of tests each admitting
a continuous equivariant return.  MA-PR-006 makes
every such test eligible.

On the actual Hilbert quotient and on the full smooth
quotient, respectively,
\[
 \begin{aligned}
 C_{h_{N,g}}^{Q_H}&\longrightarrow C_g^{Q_H}
                                      &&\text{in trace norm},\\
 \operatorname{Tr}_{Q_\infty}E_{N,g}
       &\longrightarrow
        \operatorname{Tr}_{Q_\infty}C_{T_{-a}g}
                          -\operatorname{Tr}_{Q_H}C_g .
 \end{aligned}
 \tag{MA-PR.61}
\]
These are operator constructions and trace limits even
when (MA-PR.52) fails for \(g\).
No convergence of \(S_{N,g}\), \(B_{N,g}\), or \(E_{N,g}\)
as maps into \(Q_\infty\) is inferred.

In the rational centered case, for every eligible
compact smooth \(g\) with \(g(0)=0\), the actual
ambient cutoff limit has the transport
\[
 \lim_R\operatorname{Tr}(M_{\chi_R}C_gP_G)
       =A_0(g)+B(g)-P(g)
                       +\lim_{N\to\infty}
                              \operatorname{Tr}_{Q_\infty}E_{N,g}.
 \tag{MA-PR.62}
\]
This now holds at arbitrary widths with the actual
edge multiplicities.  It does not eliminate its correction.

\begin{proof}
\emph{A fixed linear selector in each finite block.}
Let \(\mathcal A_n\subset J_n\) be the subspace satisfying
the local conditions (MA-PR.49).
Choose the following explicit linear solution of the
finite module problem for \(b\in\mathcal A_n\).
At a full retained factor take multiplication by \(b\).
At a proper factor with \(m-r<r\), take
\(c=b\bmod z^r\), represented by its polynomial of
degree less than \(r\), and set \(L(y)=cy\) in \(A\).
Eligibility makes \(c\) divisible by \(z^{m-r}\).
If \(m-r\ge r\), eligibility makes the multiplier
on \(B\) zero and take \(L=0\).
Take zero at discarded factors.
Denote the resulting linear selector by
\(\mathcal L_n(b):B_n\to J_n\).
All choices here are finite-dimensional and satisfy
\[
       \pi_n\mathcal L_n(b)=M_{\pi_nb}.
 \tag{MA-PR.63}
\]
Measure its exact cost by
\[
 A_n=\max\left\{1,\
   \sup_{\substack{b\in\mathcal A_n\\\|b\|_n\le1}}
       \|\mathcal L_n(b)\mathscr R_n\|_{\mathcal L(Q_H,J_n)}
                                    \right\}<\infty .
 \tag{MA-PR.64}
\]
The supremum is finite by finite dimensionality.
It uses the full source norm in (MA-PR.53), hence
retains all oblique-projection and near-edge costs.
No estimate for the growth of \(A_n\) is required.

Set
\[
 \theta_{N,n}=
 \begin{cases}
 1,& |n|\le N,\\
 e^{-\nu_n}/A_n,& |n|>N .
 \end{cases}
 \qquad 0<\theta_{N,n}\le1 .
 \tag{MA-PR.65}
\]
Using the centered lifts of MA-WP, define first
\[
 D_Ng=g-\sum_n(1-\theta_{N,n})
                         T_a\mathcal T_n b_n(g).
 \tag{MA-PR.66}
\]
The series converges absolutely in every
\(\mathcal E\) seminorm by (MA-WP.8)--(MA-WP.10).
Its full block class is
\(\theta_{N,n}b_n(g)\).
Thus \(D_N\) is linear, continuous, and preserves
eligibility.  Its positive scalars do not delete a
nonzero block class.

More explicitly, for every defining seminorm \(p\)
there are \(C,D\), independent of \(N\) and \(g\), with
\[
 p(D_Ng-g)
       \le C\sum_{|n|>N}\nu_n^D\|b_n(g)\|_n
                              \longrightarrow0 .
 \tag{MA-PR.67}
\]
The block restrictions of a bounded subset of
\(\mathcal E\) decrease rapidly with uniform constants,
so the displayed convergence is uniform on such a set.
This estimate concerns the regularizers in
\(\mathcal E\), not the return norms.

\emph{Preserving the value-zero test class without changing jets.}
There is a fixed \(b\in I_a\) with \(b(0)=1\).
Indeed let \(\varphi=T_a\psi\ne0\) and
\(\varphi^\dagger(x)=\overline{\varphi(-x)}\).
The full jet kernel of MT-M18 gives
\(\varphi\in I_a\), and the convolution ideal property gives
\[
          b=\frac{\varphi*\varphi^\dagger}
                         {\int_{\mathbb R}|\varphi(x)|^2\,dx}
                   \in I_a,\qquad b(0)=1 .
 \tag{MA-PR.68}
\]
No stability of the ideal under \(\dagger\) is used:
only the first convolution factor must belong to the ideal.
Define
\[
       \mathcal D_Ng
            =D_Ng-\bigl(D_Ng(0)-g(0)\bigr)b .
 \tag{MA-PR.69}
\]
The correction leaves every full spectral jet of \(D_Ng\)
unchanged and restores the value \(g(0)\).
Continuity of evaluation and (MA-PR.67) prove all the
stated \(\mathcal E\)-convergence properties.
In particular the norm-one class \(g(0)=0\) is preserved.
Neither \(b\) nor the resulting regularizers are asserted
to have compact support.

\emph{The actual nuclear maps.}
For eligible \(g\) set
\[
 F_{N,n,g}=\theta_{N,n}\mathcal L_n(b_n(g))\mathscr R_n,
 \qquad
 \widetilde S_{N,g}
       =\sum_n\mathcal T_nF_{N,n,g}.
 \tag{MA-PR.70}
\]
Every \(F_{N,n,g}=F_{N,n,g}P_n\) and is an equivariant
solution for the full class of \(h_{N,g}\).
Its actual source norm is bounded by
\[
 \|F_{N,n,g}\|
   \le
   \begin{cases}
       A_n\|b_n(g)\|_n,& |n|\le N,\\
       e^{-\nu_n}\|b_n(g)\|_n,& |n|>N .
   \end{cases}
 \tag{MA-PR.71}
\]
For fixed \(N\) this sequence decreases rapidly.
Consequently (MA-PR.58)--(MA-PR.59), now with these
actual norms or their positive upper bounds, prove
the convergence of (MA-PR.70), its all-order nuclearity,
and the actual factorization (MA-PR.60).
Linearity follows from the fixed linear selectors.

For completeness the uniformity in the statement
does not hide a uniform bound for \(A_n\).
There is one continuous seminorm \(p_0\) of
\(\mathcal E\) such that
\(\sup_n\|b_n(g)\|_n\le p_0(g)\), by the Mellin
bound on the fixed strip containing the block contours.
For fixed \(N\), use in (MA-PR.71) the positive scalar
sequence \(A_n\) on its finite head and \(e^{-\nu_n}\)
on its tail.  The square-root nuclear construction
then gives one bounded target disk and coefficients
in every \(\ell^p\), uniformly when \(p_0(g)\) is bounded.
The finite head, and hence these return bounds,
can grow without control as \(N\to\infty\).

\emph{Actual trace cyclicity and the limits.}
Compose the common nuclear series for \(S_{N,g}\)
first with \(i\) on the target and then with \(i\)
on the source.  Its two rank-one trace pairings agree
term by term and are absolutely summable.
The first composition is the actual Hilbert operator
\(C_{h_{N,g}}^{Q_H}\); the second is \(B_{N,g}\).
The approximation property and absolute smooth block
expansion in MT-M18 justify the same nuclear trace
as in MA-HM-002 and MA-PR-003.  Therefore
\[
 \operatorname{Tr}_{Q_H}C_{h_{N,g}}
    =\operatorname{Tr}_{Q_\infty}B_{N,g}
    =\operatorname{Tr}_{Q_\infty}C_{T_{-a}h_{N,g}}
                       -\operatorname{Tr}_{Q_\infty}E_{N,g}.
 \tag{MA-PR.72}
\]
These equalities precede any scalar evaluation of a trace.

MA-WP.14--MA-WP.15 is a continuous map
\(\mathcal E\to\mathcal S_1(Q_H)\).
Together with (MA-PR.67)--(MA-PR.69) it proves
the first limit in (MA-PR.61).
For the full smooth term, the absolute rank-one trace
bound from MT-M19 is
\[
 \left|\operatorname{Tr}_{Q_\infty}
                   C_{T_{-a}(h_{N,g}-g)}\right|
 \le\sum_n d_n\|b_n(h_{N,g}-g)\|_n\longrightarrow0 .
 \tag{MA-PR.73}
\]
Rapid restriction decay and continuity of the transform
justify the convergence.  Subtracting the two limits
in (MA-PR.72) proves the second limit in (MA-PR.61).

For the original compact \(g\) in the rational
specialization, apply the actual ambient theorem
MA-WP-003 to \(g\) itself and Meyer's common-space
identity MA-MEY.15--MA-MEY.16 to its full smooth
operator \(C_{T_{-1/2}g}\).
The first gives the negative of the Hilbert quotient
trace; the second gives \(P-A_0-B\).
Their actual operator identifications and (MA-PR.61)
give (MA-PR.62).
This reasoning asserts no ambient trace-class
extension for noncompact \(h_{N,g}\) and requires
no geometric formula for those regularized tests.

Only after this transport, the finite-dimensional
cyclic identity
\(\operatorname{tr}_{J_n}(L\pi_n)
 =\operatorname{tr}_{B_n}(\pi_nL)\)
evaluates the correction limit as
\[
 \lim_N\operatorname{Tr}_{Q_\infty}E_{N,g}
          =\sum_{\Xi(\rho)=0}
              (m_\rho-r_\rho)G(\rho-a).
 \tag{MA-PR.74}
\]
The local zero count and rapid Mellin decay give
absolute convergence.  At a proper nonzero
boundary truncation, eligibility forces \(G(\rho-a)=0\);
this does not make its actual nilpotent correction
operator zero.  Neither a direct-sum splitting
at that boundary nor a semisimple replacement is used.
\end{proof}

\paragraph{MA-PR-008 (a conditional close-pair topological obstruction).}
Suppose that a block contains a fully retained zero \(\rho\)
and a discarded zero \(\sigma\), so \(r_\rho=m_\rho>0\)
and \(r_\sigma=0\).  Set
\[
 \delta_n(\rho,\sigma)
  =\min\{1,\operatorname{dist}(\rho,\partial\Omega_n),
                  \operatorname{dist}(\sigma,\partial\Omega_n)\}.
 \tag{MA-PR.75}
\]
For every eligible test,
\[
 \begin{aligned}
 c_n(g)&\ge
 \frac{\delta_n(\rho,\sigma)}{2|\rho-\sigma|}
       |G(\rho-a)|\,\|\ell_{\rho-a,0}\|_{H^*},\\
 \ell_{\lambda,0}(h)&=\int e^{\lambda x}h(x)\,dx .
 \end{aligned}
 \tag{MA-PR.76}
\]
In particular, along any sequence of such pairs in blocks
whose heights tend to infinity, an unregularized continuous
return requires
\[
       \frac{|G(\rho_n-a)|}{|\rho_n-\sigma_n|}
                       \quad\hbox{to decrease rapidly in }\nu_n .
 \tag{MA-PR.77}
\]
If that necessary rapid decay fails, no continuous
equivariant return for \(g\) exists.
No such arithmetic pair sequence, or failure of
(MA-PR.77) for an arithmetic function, is asserted here.

\begin{proof}
Take any feasible \(L\) in (MA-PR.50) and let
\(F=L\mathscr R_n\).
At the full retained factor its value coordinate is forced:
\[
                  (Fv)(\rho)=G(\rho-a)\ell_{\rho-a,0}(v).
 \tag{MA-PR.78}
\]
At \(\sigma\) every full output jet is zero.
Indeed the source has no \(\sigma\)-primary factor,
and module maps from every other primary factor
to that factor are zero.
The line segment joining \(\rho\) to \(\sigma\)
lies inside the convex rectangle and has distance
at least \(\delta_n(\rho,\sigma)\) from its boundary.
For any holomorphic representative \(H\) of \(Fv\),
Cauchy's estimate on disks of radius \(\delta_n/2\)
and integration along that segment give
\[
 |H(\rho)-H(\sigma)|
       \le 2|\rho-\sigma|\delta_n^{-1}
                            \sup_{\partial\Omega_n}|H| .
 \tag{MA-PR.79}
\]
Take the infimum over representatives, then the
supremum over the unit ball of the \emph{whole}
\(Q_H\).  The continuous value functional annihilates
\(J_H\), so its quotient norm is exactly its \(H^*\)
norm.  Finally minimize over \(L\).
This proves (MA-PR.76).

The selected-height product bound (M18.7) gives
\(\delta_n(\rho,\sigma)\ge c\nu_n^{-D}\).
To see this directly, every horizontal distance
less than one is a factor in the product in (M18.7);
all factors are at most one.  Each such factor,
even with its multiplicity, is therefore bounded
below by that product.  The vertical edges
at real parts \(-1\) and \(2\) are at least one
from the critical strip.  This proves the asserted
polynomial lower bound for the distance to the
\emph{full} block boundary.  It is not a lower
bound for \(|\rho-\sigma|\).

Furthermore \(\|\ell_{\rho-a,0}\|_{H^*}\ge c'>0\),
uniformly for \(0\le\operatorname{Re}\rho\le1\):
restrict its squared dual integral to \([-1,1]\),
where the fixed weight is bounded above and
the real exponential is bounded below.
If a return exists, (MA-PR.52) and (MA-PR.76)
now imply (MA-PR.77), since multiplication by a
fixed power of \(\nu_n\) preserves rapid decay.
Its failure is thus a precise topological obstruction.
\end{proof}

\paragraph{Scope and placement.}
MA-PR-006 separates the finite algebraic condition
from the global continuity condition.  MA-PR-007
closes the exceptional-width regularized return problem
for every locally eligible test and every \(s\ge0\),
including the requested polynomial weights \(s\ge3\).
The remaining unregularized question is exactly
the actual cost condition (MA-PR.52); local eligibility
alone is not promoted to that bound.
The constructed maps land in the full smooth quotient.
At a proper truncated boundary factor there is no
invariant section on which \(i\) is an isomorphism,
as MA-PR-004 proves.
There is therefore no claim of a canonical retained
smooth direct summand, nor of the projector identities
from the generic-edge construction.

The finite objects belong to the abelian category of
finite-length \(\mathbb C[w]\)-modules.
Locally the exact sequence
\[
 0\longrightarrow z^r\mathbb C[z]/(z^m)
   \longrightarrow\mathbb C[z]/(z^m)
   \longrightarrow\mathbb C[z]/(z^r)
   \longrightarrow0
 \tag{MA-PR.80}
\]
has lengths \(m-r,m,r\).  Its relation in that category's
Grothendieck group records these lengths; it does not
split the sequence or erase its nilpotent structure.
Finite traces and determinants refer to endomorphisms
of the underlying finite-dimensional complex vector spaces,
with determinants formed by their top exterior powers.
The common nuclear series in (MA-PR.59) and (MA-PR.72)
then implements the trace passage on the actual quotients.
This is the inherited nuclear-trace placement of
cote~1, scans~079--080 and~111--112, as already specified
in MA-HM and MA-WP.  No new reading of those scans,
arithmetic positivity theorem, or weight realization
is claimed.

The exact analytic inputs are the full MT-M18
block theorem (M18.1)--(M18.4), its selected-height
bound (M18.7), and its entire interpolation estimates
(M18.10)--(M18.13), in the supplied
\texttt{block\_arithmetic\_synthesis.tex};
MT-M19 (M19.3)--(M19.4), in the supplied
\texttt{block\_trace\_hilbert.tex};
and the actual quotient, projections, kernel and
trace-norm synthesis of MA-WP.4--MA-WP.18.
The complete MT-M18 proof and the relevant MT-M19
nuclear-series proof were reread for this construction.
The Meyer transport uses the same actual fixed-sector
isomorphisms MT-M22/MT-M24 and the same common-space
trace signs MA-MEY.15--MA-MEY.16 as MA-PR-003.

The half-line finite-Paley--Wiener quotient in MA-GW
has its own causal domain and boundary blocks.
It has not been substituted for the two-sided
\(Q_H\) or for its canonical map \(i\) here.
Likewise none of the regularized returns identifies
the unbounded off-diagonal maps in Meyer's original
common-space projection on a new Hilbert domain.
No assumption A, M, generic-edge avoidance, simplicity,
or new outside preprint is used.  The actual arithmetic
deficiency and its possible absence remain unasserted.

\section{The sharp all-compact-test return criterion}
\label{sec:uniform-partial-return}

MA-PR-006 characterizes the return for one test by its
actual finite-block costs.  Uniform boundedness turns that
criterion into a precise statement for every compact test.
It also permits a single linear choice of the return maps.
The result below does not assert that its arithmetic cost
estimates hold, or that they fail.

Retain the fixed function, Hilbert parameters, actual
quotients, and block norms of MA-PR-006.
For a compact set \(K\subset\mathbb R\), put
\[
 \begin{aligned}
 \mathcal D_K
    &=\{g\in C^\infty(\mathbb R):\operatorname{supp}g\subset K\},\\
 p_j(g)&=\max_{0\le l\le j}\sup_{x\in K}|g^{(l)}(x)|.
 \end{aligned}
 \tag{MA-PR.81}
\]
Fix either of the following closed test spaces:
\[
 X_K=\mathcal D_K\cap\mathcal E_{\rm el},
 \qquad\hbox{or}\qquad
 X_K=\mathcal D_K\cap\mathcal E_{\rm el}
                              \cap\{g:g(0)=0\}.
 \tag{MA-PR.82}
\]
The choice may include the norm-one condition throughout.
All topological conclusions will be stated on these
fixed-support Fr\'echet spaces.
For \(g\in X_K\), define the possibly infinite seminorm
\[
              Q_k(g)=\sup_n\nu_n^k c_n(g),\qquad k\ge0,
 \tag{MA-PR.83}
\]
where \(c_n(g)\) is the attained minimum in (MA-PR.50),
measured on the whole \(Q_H\).

\paragraph{MA-PR-009 (uniform boundedness and a coherent linear return).}
For every fixed compact \(K\), the following are equivalent.
\begin{enumerate}
\item Every \(g\in X_K\) admits an actual continuous
equivariant return \(S_g:Q_H\to Q_\infty\) satisfying
(MA-PR.51).
\item For every integer \(k\ge0\), there are an integer
\(j=j(K,k)\) and a finite constant \(C=C(K,k)\) such that
\[
                   Q_k(g)\le C p_j(g)\qquad(g\in X_K).
 \tag{MA-PR.84}
\]
\item The returns can be chosen linearly in \(g\), with
actual linear lifts \(\widetilde S_g:Q_H\to\mathcal E\),
so that
\[
           (g,v)\longmapsto\widetilde S_gv
                     \quad\hbox{is jointly continuous on }
                     X_K\times Q_H .
 \tag{MA-PR.85}
\]
They satisfy the actual factorization and equivariance
of (MA-PR.51), and their common nuclear representations
have coefficients in every \(\ell^p\), \(p>0\), uniformly
when \(g\) ranges over a bounded subset of \(X_K\).
\end{enumerate}
The finite linear sections used for the third condition
may be chosen once for the fixed Hilbert parameters,
independently of \(K\).
Consequently, if these equivalent conditions hold on
every compact-support space, there is one linear return
assignment on all eligible compact tests in the chosen class, with the
joint continuity and nuclear bounds just stated on
each \(X_K\).
It suffices to test the intervals \(K=[-R,R]\),
\(R=1,2,\ldots\).

There is a sharp failure alternative.  Suppose that
for one fixed \(K\) and one fixed \(k\), no estimate
(MA-PR.84) holds for any finite \(j,C\).
Then
\[
       \{g\in X_K:Q_k(g)=\infty\}
          \quad\hbox{is a dense }G_\delta\hbox{ subset of }X_K .
 \tag{MA-PR.86}
\]
Every test in that set has no continuous equivariant
return.  This is a conditional topological alternative,
not a claim that the displayed failure premise occurs
for an arithmetic \(L\)-function.

\begin{proof}
\emph{The exact Fr\'echet space and its cost seminorms.}
The increasing seminorms in (MA-PR.81) make
\(\mathcal D_K\) a complete metrizable locally convex space.
Indeed uniform convergence of every derivative on a
fixed interval containing \(K\) gives a smooth limit,
by integrating the derivative limits; the limit vanishes
off \(K\).  All endpoint flatness conditions are retained.
Each Mellin jet in (MA-PR.49) is continuous on this
space, because its integral is over the fixed compact
support.  Evaluation at zero is continuous as well.
Thus either space in (MA-PR.82) is a closed Fr\'echet
subspace.

On \(X_K\), every \(c_n\) is a continuous seminorm.
Homogeneity follows by scaling feasible module maps.
For the triangle inequality, the sum of a feasible
map for \(g\) and one for \(h\) is feasible for \(g+h\),
so minimizing the triangle inequality for their actual
operator norms gives
\[
        c_n(g+h)\le c_n(g)+c_n(h),\qquad
        c_n(tg)=|t|c_n(g).
 \tag{MA-PR.87}
\]
For continuity, the fixed linear selector of
(MA-PR.63)--(MA-PR.64) gives
\(c_n(g)\le A_n\|b_n(g)\|_n\).
If \(K\subset[-R,R]\), the fixed full strip of the
M18 block contours gives
\[
 \|b_n(g)\|_n
       \le \sup_{w\in\partial\Omega_n}
                   \left|\int_K e^{(w-a)x}g(x)\,dx\right|
       \le 2R e^{(2+|a|)R}p_0(g).
 \tag{MA-PR.88}
\]
Take a positive \(R\) even for a degenerate compact set.
For each fixed \(n\), \(A_n\) is finite; no uniform
growth bound for these constants is used.
This proves continuity of \(c_n\).
The supremum \(Q_k\) is therefore a lower
semicontinuous extended seminorm.

\emph{Pointwise returns imply finite-seminorm bounds.}
Assume the first condition.
MA-PR-006 makes \(Q_k(g)<\infty\) for every \(g\in X_K\)
and every fixed \(k\).
For this \(k\), let
\[
             E_M=\{g\in X_K:Q_k(g)\le M\},
                         \qquad M=1,2,\ldots .
 \tag{MA-PR.89}
\]
These are closed, balanced and convex, and they cover \(X_K\).
The Baire property of the complete metrizable space
makes one \(E_M\) have nonempty interior.

For clarity, the complete-metric fact used here can
be seen directly.  If countably many closed sets all
had empty interior, start inside an arbitrary nonempty
open set and choose successively nonempty closed balls,
each contained in the interior of its predecessor,
avoiding the next closed set, with radii tending to zero.
Their centers form a Cauchy sequence and their common
limit, by completeness, avoids every closed set.
Thus such sets cannot cover the space.
The same construction starting in any open set proves
the density of the complement of their union.

Choose \(g_0\) in the interior of this \(E_M\).
For some neighborhood \(U\) of zero,
\(g_0+U\subset E_M\).  Since \(g_0\in E_M\),
(MA-PR.87) gives \(Q_k(u)\le2M\) for \(u\in U\).
The increasing Fr\'echet seminorms imply that \(U\)
contains \(\{g:p_j(g)<\epsilon\}\) for some \(j\)
and \(\epsilon>0\).
Scaling, and then letting the strict inequality
approach \(\epsilon\), gives
\[
                  Q_k(g)\le (2M/\epsilon)p_j(g).
 \tag{MA-PR.90}
\]
This is (MA-PR.84).  The index \(j\) and the constant
may depend on \(K\) and \(k\).

\emph{A finite linear section with the precise dimension loss.}
We next construct the coherent family required for
the third condition.  This avoids selecting minimizers
nonlinearly as the test varies.
Define the finite-dimensional normed space
\[
 \mathcal X_n=\operatorname{Hom}_{\mathbb C[w]}(B_n,J_n),
 \qquad
       \|L\|_{\mathcal X_n}=\|L\mathscr R_n\| .
 \tag{MA-PR.91}
\]
This is a genuine norm since \(\mathscr R_n\) is
surjective.  All source norms remain those on the
whole Hilbert quotient.
Let
\[
 \begin{aligned}
 \Delta_n:\mathcal X_n&\longrightarrow
                 \operatorname{End}_{\mathbb C[w]}(B_n),
             &\Delta_nL&=\pi_nL,\\
 \mathcal Y_n&=\operatorname{im}\Delta_n,
             &\|y\|_{\mathcal Y_n}
                 &=\min_{\Delta_nL=y}\|L\|_{\mathcal X_n}.
 \end{aligned}
 \tag{MA-PR.92}
\]
This is the actual quotient norm, with attained
minimum by the proof of (MA-PR.53).
An endomorphism of \(B_n\) is multiplication by
its value on the cyclic vector \(1\).
The primary calculation (MA-PR.54) therefore gives
\[
 e_n:=\dim\mathcal Y_n
       =\sum_{\rho\in\Omega_n}
                     \max(0,2r_\rho-m_\rho)
       \le\dim B_n\le d_n=O(\log\nu_n).
 \tag{MA-PR.93}
\]
At a full retained factor the summand is \(m_\rho\);
at a discarded factor it is zero.

Choose an Auerbach basis \(y_{n,l}\), \(y^*_{n,l}\)
of \(\mathcal Y_n\): all vectors and functionals
have norm one.  The elementary determinant-maximization
construction of MT-M18 supplies this basis.
For each \(y_{n,l}\), choose an attained minimum
preimage \(L_{n,l}\), of \(\mathcal X_n\)-norm one.
Then
\[
 \begin{aligned}
 \sigma_n(y)&=\sum_{l=1}^{e_n}y^*_{n,l}(y)L_{n,l},&
 \Delta_n\sigma_n&=1,\\
 \|\sigma_n(y)\|_{\mathcal X_n}
                  &\le e_n\|y\|_{\mathcal Y_n}.
 \end{aligned}
 \tag{MA-PR.94}
\]
If \(e_n=0\), use the zero map.
These are fixed complex-linear sections of the
finite spaces of feasible homomorphism data.
They are not module sections of
\(\pi_n:J_n\to B_n\), which a proper truncation
would forbid.

For an eligible test, set
\[
 \begin{aligned}
 y_n(g)&=M_{\pi_nb_n(g)}\in\mathcal Y_n,&
 \|y_n(g)\|_{\mathcal Y_n}&=c_n(g),\\
 \widehat F_{n,g}&=\sigma_n(y_n(g))\mathscr R_n .
 \end{aligned}
 \tag{MA-PR.95}
\]
These maps are linear in \(g\) and satisfy
\[
 \begin{aligned}
 \widehat F_{n,g}&=\widehat F_{n,g}P_n,&
 \|\widehat F_{n,g}\|&\le e_nc_n(g),\\
 \pi_n\sigma_n(y_n(g))&=M_{\pi_nb_n(g)} .
 \end{aligned}
 \tag{MA-PR.96}
\]
Thus the dimension loss is at most \(O(\log\nu_n)\),
with no hidden interpolation or edge-distance constant.

\emph{Uniform estimates give the actual continuous family.}
Assume (MA-PR.84) for every \(k\).
Define
\[
        \widetilde S_gv
              =\sum_n\mathcal T_n\widehat F_{n,g}v,
              \qquad S_g=q_\infty\widetilde S_g .
 \tag{MA-PR.97}
\]
For every defining seminorm \(p\) of \(\mathcal E\),
choose a sufficiently large integer \(k_p\) relative
to its fixed polynomial lift exponent in (MA-PR.44).
Equations (MA-PR.93) and (MA-PR.96) give
\[
 \begin{aligned}
 p(\widetilde S_gv)
   &\le C_p\sum_n\nu_n^{D_p}e_nc_n(g)\|v\|\\
   &\le C'_p Q_{k_p}(g)\|v\|
    \le C''_{K,p}p_{j(K,k_p)}(g)\|v\|.
 \end{aligned}
 \tag{MA-PR.98}
\]
For example any integer \(k_p>D_p+2\) makes the
sum converge.  The series is absolutely convergent
in every seminorm, uniformly on the indicated
bounded sets.  This proves its linearity in \(g\)
and the joint continuity in (MA-PR.85).
Its maps satisfy the actual factorization and
equivariance by the same full-jet and actual
common-kernel argument as (MA-PR.58).

Let \(\mathcal B\subset X_K\) be bounded.
The estimates (MA-PR.84) for every \(k\) show that
\[
       \beta_n=e^{-\nu_n}
                   +\sup_{g\in\mathcal B}\|\widehat F_{n,g}\|
                   \quad\hbox{decreases rapidly}.
 \tag{MA-PR.99}
\]
For an Auerbach basis \(e_{n,j},\ell_{n,j}\) of \(J_n\),
use target vectors
\(\beta_n^{1/2}\mathcal T_ne_{n,j}\),
coefficients \(\beta_n^{1/2}\), and source functionals
\(\ell_{n,j}\widehat F_{n,g}/\beta_n\).
The source norms are at most one for all
\(g\in\mathcal B\), the target vectors form one
bounded set in \(\mathcal E\), and the coefficients,
repeated \(d_n\) times, belong to every \(\ell^p\).
This proves exactly the asserted uniform
all-order nuclear representation.
There is no passage from pointwise bounds to
an unjustified uniform near-boundary estimate:
the only uniform step is the proved Baire implication
(MA-PR.90).

The sections \(\sigma_n\) depend only on the fixed
finite normed spaces, so the same formula (MA-PR.97)
works on every compact-support space for which
(MA-PR.84) holds.  This proves coherence across
supports.  Continuity is asserted on each \(X_K\);
no additional assertion about a different inherited
LF topology is needed.
The third condition plainly implies the first,
which completes the equivalence.

\emph{The residual failure alternative.}
Fix \(K,k\) and suppose that every estimate in
(MA-PR.84) fails.
The closed sets \(E_M\) of (MA-PR.89) then have
empty interior: otherwise the exact neighborhood
and scaling argument proving (MA-PR.90) would
give one of the forbidden finite-seminorm estimates.
Their open complements are dense.
The complete-metric argument given above shows
that
\[
       X_K\setminus\bigcup_{M\ge1}E_M
          =\bigcap_{M\ge1}(X_K\setminus E_M)
\]
is a dense \(G_\delta\).  This set is precisely
\(\{g:Q_k(g)=\infty\}\).
MA-PR-006 excludes a continuous equivariant return
for every one of its elements.  This proves
(MA-PR.86) without constructing or assuming
an arithmetic close-pair sequence.
\end{proof}

\paragraph{Exact remaining gate and source placement.}
The all-compact-test question is now exactly the
family of finite derivative estimates (MA-PR.84).
The theorem proves neither these estimates nor
their failure for the arithmetic function.
MA-PR-008 supplies one possible necessary quantitative
test of them, conditional on an actual retained/discarded
configuration; it supplies no such configuration.
The adaptive returns of MA-PR-007 remain available
for every eligible test without (MA-PR.84).

The spaces \(J_n\) and \(B_n\) remain actual
finite-length \(\mathbb C[w]\)-modules.
The spaces \(\mathcal X_n\) and \(\mathcal Y_n\)
in this argument organize feasible module maps
and carry their measured operator and quotient norms.
The Auerbach section (MA-PR.94) is only complex linear
in this space of maps.  It does not split the local
module extension (MA-PR.80), and no splitting is
inferred from its relation in the Grothendieck group
of finite-length modules.
The common nuclear series in (MA-PR.97)--(MA-PR.99)
is the actual analytic realization; finite traces
and exterior-power determinants retain the scope
already specified in MA-PR-006--008.
No new arithmetic positivity or zero-location claim
is obtained.

The exact inputs are MA-PR-006, including its
whole-quotient Gram norm, minimum attainment and
actual return criterion; the finite module calculation
MA-PR.54; the full smooth adapted topology
MT-M18 (M18.1)--(M18.3), its polynomial entire-lift
bound (M18.13), and its complete Auerbach-basis
construction.  The Baire implication and the
finite linear-section estimate used here are proved
above.  The inherited cote~1 nuclear-trace locators
remain scans~079--080 and~111--112, as in MA-WP
and MA-PR-006--008; no new reading of those scans
is claimed.  No A, M, generic-edge assumption,
or new outside preprint is an input.

\section{One polynomial block-cost criterion without proper boundary truncation}
\label{sec:polynomial-partial-return}

The all-compact-test criterion admits a further simplification
when the Hilbert quotient retains or discards each primary
factor in full.  In this case the finite return has no
algebraic freedom.  A single polynomial estimate for its
actual norm is equivalent to returns for every smooth test.
The estimate itself is not established here at a general width.

Retain all fixed data and the original adapted block norms
of MA-PR-006.  Assume explicitly that
\[
                 r_\rho\in\{0,m_\rho\}
                    \qquad\hbox{for every zero }\rho.
 \tag{MA-PR.100}
\]
This holds automatically when \(s\le1\), because then
no edge jet is continuous, and at widths with no boundary
zero for every \(s\ge0\).  It also permits fully retained
boundary factors.  It is not automatic at other widths
with \(s>1\).  Under (MA-PR.100), eligibility imposes
no condition: \(\mathcal E_{\rm el}=\mathcal E\).

The primary decomposition supplies the canonical module
inclusion \(\iota_n:B_n\to J_n\): it is the identity
on every fully retained factor and zero at every
discarded factor.  Define its measured whole-quotient cost by
\[
       \mathcal C_n
           =\|\iota_n\mathscr R_n\|_{\mathcal L(Q_H,J_n)}.
 \tag{MA-PR.101}
\]
Set \(\mathcal C_n=0\) if \(B_n=0\).
The target has its original holomorphic quotient norm,
and \(\mathscr R_n\) acts on the whole \(Q_H\).
Thus (MA-PR.101) includes every oblique-projection cost.

\paragraph{MA-PR-010 (a single polynomial gate for all tests).}
Under (MA-PR.100), the following conditions are equivalent.
\begin{enumerate}
\item There are fixed finite constants \(C>0\), \(D\ge0\) with
\[
                \mathcal C_n\le C\nu_n^D
                           \qquad(n\in\mathbb Z).
 \tag{MA-PR.102}
\]
\item Every \(g\in\mathcal E\) has a continuous
equivariant return satisfying (MA-PR.51).
\item Every \(g\in C_c^\infty(\mathbb R)\) with
\(g(0)=0\) has such a return.
\item Every \(g\in\mathcal D_{[1,2]}\) has such a return.
\end{enumerate}
The same equivalence holds if condition~3 is stated
for all compact smooth tests without the value-zero
restriction.  Condition~4 concerns the entire fixed-support
Fr\'echet space, not just a chosen sequence of tests.

Whenever these conditions hold, the return \(S_g\) is
unique as a map to \(Q_\infty\), is linear in \(g\),
and has an \(\mathcal E\)-valued lift such that
\[
       (g,v)\longmapsto\widetilde S_gv
             \quad\hbox{is jointly continuous on }
                         \mathcal E\times Q_H .
 \tag{MA-PR.103}
\]
Its common nuclear representations have coefficients
in every positive \(\ell^p\), uniformly for \(g\)
in bounded subsets of \(\mathcal E\).
The quotient-valued return is unique; its representative
lift into \(\mathcal E\) uses the chosen MT-M18 lifts.
All constants concern the fixed function, Hilbert
parameters and adapted blocks.  No uniformity as these
data vary is asserted.

If (MA-PR.102) fails, then
\[
 \{g\in\mathcal D_{[1,2]}:
                   \sup_n c_n(g)=\infty\}
       \quad\hbox{is a dense }G_\delta
                    \hbox{ in }\mathcal D_{[1,2]}.
 \tag{MA-PR.104}
\]
Every such compact test satisfies \(g(0)=0\) and
has no continuous equivariant return.  This alternative
does not assert the arithmetic failure of (MA-PR.102).

\begin{proof}
\emph{Uniqueness of the finite map and exact multiplication.}
At a fully retained factor the reduction is the identity
\(\mathbb C[z]/(z^m)\to\mathbb C[z]/(z^m)\).
A module map factoring convolution is therefore forced
to be multiplication by the full Taylor class.
At a discarded factor there is no source factor;
all maps from other primary supports to it vanish by
coprimality of their annihilating polynomials.
Thus, for any \(b\in J_n\), the sole feasible map
with reduction \(M_{\pi_nb}\) is
\[
       L_n(b)=M_b\iota_n=\iota_nM_{\pi_nb}.
 \tag{MA-PR.105}
\]
The equality follows factor by factor in the actual
finite modules; it does not use an orthogonal Hilbert
decomposition.  In particular
\[
 \begin{aligned}
 c_n(g)&=\|M_{b_n(g)}\iota_n\mathscr R_n\|,\\
 c_n(g)&\le\|b_n(g)\|_n\mathcal C_n.
 \end{aligned}
 \tag{MA-PR.106}
\]
The latter estimate is the submultiplicative original
quotient norm on \(J_n\).
Every continuous global return has exactly these
block maps by (MA-PR.57).  Since the full smooth
jet kernel is zero in \(Q_\infty\), they determine
the global return uniquely whenever it exists.

\emph{The polynomial estimate gives all smooth returns.}
Assume (MA-PR.102).  For every \(g\in\mathcal E\),
the sequence \(\|b_n(g)\|_n\) decreases faster than
every power.  Equation (MA-PR.106) therefore satisfies
the exact return criterion (MA-PR.52).
More explicitly, choose an integer \(N_0>3+|a|\).
Integration by parts on the fixed strip containing
all full block contours gives, for each integer \(M\ge0\),
\[
 \|b_n(g)\|_n\le C_{a,M}\nu_n^{-M}
                  \sum_{l=0}^{M}p_{N_0,l}(g).
 \tag{MA-PR.107}
\]
At large \(|n|\), apply \(M\) integrations by parts
to \(e^{(\operatorname{Re}w-a)x}g(x)\) in its Fourier
integral.  Its differentiated integrands are bounded
by the displayed exponential seminorms times an
integrable exponential in \(|x|\).
The finitely many small heights follow from the
undifferentiated integral after enlarging the constant.
This proves (MA-PR.107) for every derivative order;
it is an estimate in the original quotient norm.

Set
\[
 \begin{aligned}
 F_{n,g}&=M_{b_n(g)}\iota_n\mathscr R_n,\\
 \widetilde S_gv&=\sum_n\mathcal T_nF_{n,g}v,
             \qquad S_g=q_\infty\widetilde S_g .
 \end{aligned}
 \tag{MA-PR.108}
\]
The finite map and \(\mathcal T_n\) are fixed,
so the series is linear in \(g\) and in \(v\).
For a defining target seminorm \(p_{N,j}\), let
\(D_{N,j}\) be its polynomial lift exponent in
(MA-PR.44).  Choose an integer
\(M>D+D_{N,j}+2\).
Then (MA-PR.102), (MA-PR.107) and the lift bound give
\[
 p_{N,j}(\widetilde S_gv)
       \le C_{N,j,M}\left(\sum_{l=0}^{M}
                                   p_{N_0,l}(g)\right)\|v\|.
 \tag{MA-PR.109}
\]
The series converges absolutely in every seminorm,
uniformly on the corresponding bounded product sets.
This proves the claimed joint continuity and the
actual lift.  Its full block class is \(F_{n,g}v\).
Taking retained jets proves \(iS_g=C_g^{Q_H}\) by
the actual common-kernel theorem MA-WP-001.
The module identities in (MA-PR.105) prove the
centered equivariance in (MA-PR.51).

For a bounded test set \(\mathcal B\subset\mathcal E\),
the sequence
\(\beta_n=e^{-\nu_n}+\sup_{g\in\mathcal B}\|F_{n,g}\|\)
is positive and rapidly decreasing by (MA-PR.107).
Using an Auerbach basis of \(J_n\), take target
vectors \(\beta_n^{1/2}\mathcal T_ne_{n,l}\),
coefficients \(\beta_n^{1/2}\), and functionals
\(\ell_{n,l}F_{n,g}/\beta_n\).
The target vectors form one bounded family in
\(\mathcal E\), the source norms are at most one,
and the coefficients repeated \(d_n=O(\log\nu_n)\)
times belong to every positive \(\ell^p\).
This proves the uniform all-order nuclear assertion.
Condition~1 has now given condition~2 and the
stated coherent topology.  The implications
\(2\Rightarrow3\Rightarrow4\) follow from the
test-space inclusions.

\emph{One fixed positive bump with uniform inverse classes.}
For the converse fix \(x_0=3/2\) and the compact
rectangle
\[
 \mathcal R_a=[-1-a,2-a]
                           +i[1/4,7/4].
 \tag{MA-PR.110}
\]
Let \(M_a=\max_{z\in\mathcal R_a}|z|\).
Choose \(0<\epsilon<1/4\) so small that
\(e^{M_a\epsilon}-1\le1/2\).
Take a nonnegative \(\phi\in C_c^\infty
((x_0-\epsilon,x_0+\epsilon))\) with integral one.
For \(z\in\mathcal R_a\), the elementary exponential
estimate gives
\[
 \begin{split}
 \left|e^{-x_0z}\mathcal M\phi(z)-1\right|
   &=\left|\int\bigl(e^{z(x-x_0)}-1\bigr)\phi(x)\,dx\right|\\
   &\le e^{M_a\epsilon}-1\le\frac12 .
 \end{split}
 \tag{MA-PR.111}
\]
In particular the transform is bounded above and
away from zero on this whole rectangle, with
constants depending only on the fixed \(a\) and bump.

Put \(g_n(x)=e^{-inx}\phi(x)\).
It is supported in the same compact subset of
\((1,2)\) and has \(g_n(0)=0\).
For \(w\in\overline\Omega_n\), the selected heights
give \(w-a-in\in\mathcal R_a\), and
\[
 \begin{aligned}
 b_n(g_n)(w)&=\mathcal M\phi(w-a-in),\\
 \|b_n(g_n)\|_n+\|b_n(g_n)^{-1}\|_n&\le C_a .
 \end{aligned}
 \tag{MA-PR.112}
\]
For the inverse, use the actual holomorphic function
\(1/\mathcal M\phi(w-a-in)\) on a neighborhood of
the closed rectangle, and take its class in \(J_n\).
Equation (MA-PR.111) makes its supremum norm
uniform in \(n\).  There is no inverse Hermite
matrix or zero-separation factor in this argument.
Blocks with no retained factor require no inverse
estimate, since their cost is zero.

Every derivative of the modulated bump has the
explicit formula
\[
 \begin{aligned}
 g_n^{(l)}(x)
   &=e^{-inx}\sum_{q=0}^{l}\binom lq
                             (-in)^{l-q}\phi^{(q)}(x),\\
 p_j^{[1,2]}(g_n)&\le C_{\phi,j}\nu_n^j,
       \qquad
 p_{N,j}(g_n)\le C_{\phi,N,j}\nu_n^j .
 \end{aligned}
 \tag{MA-PR.113}
\]
The last bound also follows directly because
the common support is contained in \([1,2]\).
Constants may depend on the displayed derivative
orders, but never on \(n\).

\emph{All fixed-support returns force the polynomial cost.}
Assume condition~4.  There are no eligibility
constraints under (MA-PR.100), so MA-PR-009 with
\(X_K=\mathcal D_{[1,2]}\) and \(k=0\) gives
one finite \(j\) and one \(C\) such that
\[
                     \sup_m c_m(g)\le C p_j^{[1,2]}(g)
                          \qquad(g\in\mathcal D_{[1,2]}).
 \tag{MA-PR.114}
\]
This uniform conclusion uses the return hypothesis
for the entire Fr\'echet space.  Existence merely
for the sequence \(g_n\) would not justify it.

Use the unique finite map (MA-PR.105) and multiply
it by the inverse full class in (MA-PR.112).
Submultiplicativity and (MA-PR.113)--(MA-PR.114) give
\[
 \begin{aligned}
 \mathcal C_n
   &=\|M_{b_n(g_n)^{-1}}F_{n,g_n}\|\\
   &\le C_a c_n(g_n)
    \le C_a C p_j^{[1,2]}(g_n)
    \le C'\nu_n^j .
 \end{aligned}
 \tag{MA-PR.115}
\]
Thus condition~1 holds.  This completes all four
equivalences.  If condition~3 omits the value-zero
restriction, condition~2 still implies it, and it
still implies condition~4, so that version is
equivalent as well.

If the polynomial estimate fails, (MA-PR.115)
shows that no bound of the form (MA-PR.114) can
hold for any finite \(j,C\).
The \(k=0\) failure alternative of MA-PR-009
then gives precisely (MA-PR.104).
No arithmetic instance of that failure is needed
for this conditional implication.
\end{proof}

\paragraph{Trace transport and the remaining domains.}
Under these equivalent conditions, the common nuclear
series in (MA-PR.108) permits the actual trace passage
\[
 \begin{gathered}
 B_g=S_g i,\qquad E_g=C_{T_{-a}g}^{Q_\infty}-B_g,\\
 \begin{aligned}
 \operatorname{Tr}_{Q_H}C_g
   &=\operatorname{Tr}_{Q_\infty}B_g,\\
 \operatorname{Tr}_{Q_\infty}B_g
   &=\operatorname{Tr}_{Q_\infty}C_{T_{-a}g}
                         -\operatorname{Tr}_{Q_\infty}E_g .
 \end{aligned}
 \end{gathered}
 \tag{MA-PR.116}
\]
The rank-one pairings coincide and are absolutely
summable, exactly as proved in MA-PR.72.
In each full retained factor \(B_g\) is multiplication
by the full class; in each discarded factor it is zero.
Thus \(E_g\) is an actual nuclear correction on the
discarded smooth factors.  For an original compact
test with \(g(0)=0\) in the rational centered sector, the existing
Meyer identification and MA-WP-003 supply the same
geometric transport as MA-PR.62, now with
\(\operatorname{Tr}E_g\) itself instead of a
regularized correction limit.
The arithmetic disappearance of that correction
has not been proved.

These return maps act on the specified smooth and
Hilbert quotients.  They do not extend the original
off-diagonal maps of Meyer's common-space projection
to new Hilbert domains, or identify that projection
with a source PW cutoff.  No ambient trace-class
claim for arbitrary noncompact \(g\in\mathcal E\)
is made by (MA-PR.116).
The assumption \(s\le1\) removes proper edge
truncation; it does not by itself prove
(MA-PR.102).  Likewise the theorem does not supply
the polynomial cost bound merely from absence of
an exact boundary zero.

The exact inputs are the actual criterion and
whole-quotient norms MA-PR-006, the proved Baire
implication MA-PR-009, and the full adapted
MT-M18 block topology and entire-lift estimates.
The modulated-bump inverse estimate is proved here
in the original holomorphic norm.  No A, M,
simplicity, generic-edge distance estimate, or
new outside preprint is used.

The finite-length \(\mathbb C[w]\)-module and nuclear-trace
placement is the one specified at (MA-PR.80) and in
MA-PR-009.  Under (MA-PR.100), the primary decomposition
itself supplies the finite retained/discarded inclusion;
its global analytic cost is still the separate estimate
(MA-PR.102).  A Grothendieck-group relation alone supplies
no such estimate.  The inherited cote~1 scan locators
and exterior-power determinant conventions retain their
previously stated scope; no new reading of those scans
or arithmetic positivity conclusion is claimed.

\section{Actual quotient costs controlled by the edge zero potential}
\label{sec:edge-potential-quotient-cost}

This section estimates the canonical cost of MA-PR-010 in its
original whole-quotient norm.  It gives sufficient local zero
geometry for polynomial costs, and records exactly what the
standard count and generic distance estimates alone supply.
It asserts neither that the new sufficient geometry holds
for the fixed arithmetic function nor that a contrary
arithmetic zero configuration occurs.

Keep the fixed primitive unitary Hecke function, pole-removed
entire completion \(\Xi\), Hilbert space and actual quotients
of (MA-PR.1).  Write
\[
 e_-=a+u,\qquad e_+=a+v,\qquad \nu_n=2+|n|,
 \qquad
 \Gamma_n=\partial\bigl([-1,2]+i[h_n,h_{n+1}]\bigr),
 \tag{MA-PR.150}
\]
with the original selected heights and holomorphic quotient
norms \(J_n\) of MT-M18.  Assume the generic edge-distance
condition in exactly its proved form:
\[
 \min\{|\operatorname{Re}\rho-e_-|,
               |\operatorname{Re}\rho-e_+|\}
       \ge c_0(2+|\operatorname{Im}\rho|)^{-2}
                     \qquad(\Xi(\rho)=0).
 \tag{MA-PR.151}
\]
All zeros are counted with their actual multiplicities.
There are no boundary zeros, so every retained primary
factor is retained in full, for every fixed \(s\ge0\).
The canonical maps and costs are therefore defined:
\[
 \mathcal C_n
   =\|\iota_n\mathscr R_n\|_{\mathcal L(Q_H,J_n)}.
 \tag{MA-PR.152}
\]
The norm on the source is the norm on the whole \(Q_H\);
in particular no oblique finite-block projection is
discarded or replaced by an orthogonal one.

Let
\[
 \begin{split}
 \mathcal Z_n&=\{\rho:n-2\le\operatorname{Im}\rho\le n+3\},\\
 U_n&=\max_{e\in\{e_-,e_+\}}
       \sup_{h_n\le t\le h_{n+1}}
       \sum_{\rho\in\mathcal Z_n}m_\rho
             \log^+\frac1{|e+it-\rho|}.
 \end{split}
 \tag{MA-PR.153}
\]
This is a finite nonnegative real number.  Unlike a
bound using only real-part distances, it retains the
actual distances in the complex plane and the spread
of zero ordinates within a block.

\paragraph{MA-PR-012 (whole-quotient cost and a sufficient local count).}
Under (MA-PR.151), there are fixed finite constants
\(C,D\), depending on the fixed arithmetic function
and Hilbert parameters, such that
\[
                    \mathcal C_n\le C\nu_n^D e^{U_n}
                                   \qquad(n\in\mathbb Z).
 \tag{MA-PR.154}
\]
Thus the sufficient condition \(U_n=O(\log\nu_n)\)
gives all the actual returns, coherent lifts and
all-order nuclear representations of MA-PR-010
for every \(g\in\mathcal E\).
This conclusion holds for every fixed \(s\ge0\).

Here are two explicit sufficient conditions for that
potential bound.  Either may be imposed at both edges.
\begin{enumerate}
\item There are constants \(A_0,A_1\) such that, for every
block \(n\), \(t\in[h_n,h_{n+1}]\), and \(0<r\le1\),
\[
 N_e(t,r):=\sum_{|e+it-\rho|<r}m_\rho
                   \le A_0+A_1r\log\nu_n .
 \tag{MA-PR.155}
\]
\item There are fixed \(\eta>0\) and \(M<\infty\) such that
\[
 \sum_{\substack{\rho\in\mathcal Z_n\\
                     |\operatorname{Re}\rho-e|<\eta}}
                         m_\rho\le M
                         \qquad(n\in\mathbb Z).
 \tag{MA-PR.156}
\]
\end{enumerate}
The first condition permits a growing total block count
provided it has the stated distribution at small radii.
The second is a bounded count only in a fixed band about
the edges, rather than in the whole critical strip.
Neither condition is deduced here from the ordinary
\(O(\log\nu_n)\) count.

Without either additional condition, the standard count
and (MA-PR.151) yield only the established estimate
\[
                U_n=O(\log^2\nu_n),\qquad
                \mathcal C_n\le C\exp(C'\log^2\nu_n).
 \tag{MA-PR.157}
\]
This is an upper bound supplied by these estimates,
not a lower bound or an asserted arithmetic failure of
polynomial costs.  Absence of proper boundary truncation
does not improve this argument by itself.

\begin{proof}
\emph{1. The actual inner restriction and the full output class.}
Choose a fixed small \(\delta_0>0\) and set
\[
 \delta_n=\delta_0\nu_n^{-2},\qquad
 \Omega_n^{\rm in}
  =\{e_-+\delta_n<\operatorname{Re}w<e_+-\delta_n,\
                       h_n<\operatorname{Im}w<h_{n+1}\}.
 \tag{MA-PR.158}
\]
Take \(\delta_0<(v-u)/4\).  For \(\rho\in\mathcal Z_n\),
\(2+|\operatorname{Im}\rho|\) and \(\nu_n\) are comparable
with absolute constants.  Decrease \(\delta_0\) so that
\[
 \delta_n\le\tfrac14
           \min\{|\operatorname{Re}\rho-e_-|,
                          |\operatorname{Re}\rho-e_+|\}
                      \qquad(\rho\in\mathcal Z_n).
 \tag{MA-PR.159}
\]
The rectangle contains exactly the retained zeros of
block \(n\), with their full multiplicities.  Its height,
width and perimeter have the same uniform bounds as
in (MA-PR.6).  Empty retained blocks have cost zero
and may be omitted.

For \(h\in H\), the function
\(F_h(w)=\mathcal Mh(w-a)\) is holomorphic on the
open Hilbert strip and continuous on this closed
inner rectangle.  The actual weighted
Cauchy--Schwarz estimate gives
\[
 \sup_{\partial\Omega_n^{\rm in}}|F_h|
 \le
 \|h\|_H
 \sup_{u+\delta_n\le b\le v-\delta_n}
 \left(\int_{\mathbb R}
      \frac{e^{2bx}}{\max(e^{2ux},e^{2vx})(1+x^2)^{s/2}}
                  \,dx\right)^{1/2}
 \le C\nu_n\|h\|_H .
 \tag{MA-PR.160}
\]
Only \(s\ge0\) is used: discard the polynomial denominator
and integrate \(e^{-2\delta_n|x|}\).
Let \(J_n^{\rm in}\) be the quotient of
\(A(\Omega_n^{\rm in})\) by the full retained jets,
with its boundary supremum quotient norm.  The
restriction in (MA-PR.160), followed by this quotient,
annihilates \(T_a\mathcal I\), hence its closure \(J_H\).
It is consequently a continuous map from the whole
\(Q_H\), of norm at most \(C\nu_n\).

Fix \(\tau>0\).  For a representative \(F\) of a class
in \(J_n^{\rm in}\), use the actual entire interpolant
\[
 (\mathcal H_nF)(z)
   =\frac1{2\pi i}\int_{\partial\Omega_n^{\rm in}}
       F(w)\frac{\Xi(z)-\Xi(w)}{\Xi(w)(z-w)}
                        e^{\tau(z-w)^2}\,dw .
 \tag{MA-PR.161}
\]
The divided difference is entire in \(z\), including
\(z=w\).  The same Cauchy argument as (MA-PR.9)
gives \(\mathcal H_nF=F+\Xi A\) in the inner
rectangle, where \(A\) is holomorphic there, and
divisibility by \(\Xi\) near every exterior zero.
If the full retained jets of \(F\) vanish, then
\(F/\Xi\) is holomorphic inside; both exterior
Cauchy integrals vanish.  Entireness makes
\(\mathcal H_nF=0\) everywhere.  Thus (MA-PR.161)
depends on the actual quotient class.

For each fixed \(n\), its entire function decreases
rapidly on every vertical strip: the fixed-strip
bound for \(\Xi\), its derivative bound near the
removable diagonal, and the Gaussian factor give
this directly.  Mellin inversion as in MT-M18
therefore supplies an actual element of \(\mathcal E\).
It has the input's complete retained jets and zero
full jets at every discarded zero and every other
block.  In particular, on taking its class in the
\emph{original full} \(J_n\), one obtains exactly
\[
 \bigl[\mathcal H_nF_h\bigr]_{J_n}
                 =\iota_n\mathscr R_n(\pi h).
 \tag{MA-PR.162}
\]
This identity follows from all primary Taylor classes,
not from a trace comparison.  It proves which actual
whole-quotient operator the following norm estimate
measures.

\emph{2. The potential controls the genuine contour denominator.}
For \(e=e_-\) or \(e_+\), let \(e_n\) be its displaced
inner edge.  Equation (MA-PR.159) gives, for every
\(\rho\in\mathcal Z_n\),
\[
 |e_n+it-\rho|\ge\tfrac34|e+it-\rho|.
 \tag{MA-PR.163}
\]
Indeed the real component changes by at most one
quarter of its absolute value and the imaginary
component does not change.  Thus the sum of
logarithmic losses on either inner edge is at most
\(U_n+(\log(4/3))\sum_{\mathcal Z_n}m_\rho\).
Zeros outside \(\mathcal Z_n\) have imaginary
distance greater than one from every such edge
point: \(h_n>n+1/4\) and \(h_{n+1}<n+7/4\).
They contribute no positive logarithm of an inverse
complex distance.

At large height use the fixed finite disk centered
at \(2+it\), with radius chosen so that both fixed
edge real parts and their small displacements lie
in its inner quarter.  The radius ranges over a
fixed interval and its boundary avoids zeros.
There are \(O(\log\nu_n)\) zeros in the disk.
The Euler lower bound at its center, polynomial
finite-\(L\) upper bound, and Harnack argument of
MT-M18 give a polynomial inverse bound for the
zero-free factor.  A normalized Blaschke factor
satisfies
\[
 |B_\rho(w)|\ge (2R)^{-1}\min\{1,|w-\rho|\}.
 \tag{MA-PR.164}
\]
All factors with distance at least one cost only
\((2R)^{-1}\); their product is a fixed power of
\(\nu_n\).  The remaining factors are controlled
by (MA-PR.163) and \(U_n\).
Multiplying by the lower Stirling bound for the
completion therefore gives, on the two inner
vertical edges,
\[
 |\Xi(w)|\ge c\nu_n^{-D_0}
                         e^{-\kappa|n|}e^{-U_n},
 \qquad \kappa=\pi[K:\mathbb Q]/4.
 \tag{MA-PR.165}
\]
On the horizontal edges the selected-height
bound (M18.8) is polynomial with the same
completion exponential.  Enlarging \(D_0\)
gives (MA-PR.165) on the entire inner contour.
The finitely many small heights have nonzero
contour denominators, so adjustment of the
constant covers them too.

\emph{3. Estimate in the original full block norm.}
For \(z\in\Gamma_n\) and
\(w\in\partial\Omega_n^{\rm in}\), all real parts
range over one fixed interval, and their imaginary
parts differ by less than \(3/2\).  The derivative
bound following (M18.6), applied on the segment
between \(z\) and \(w\), gives
\[
 \left|\frac{\Xi(z)-\Xi(w)}{z-w}\right|
             \le C\nu_n^{D_1}e^{-\kappa|n|}.
 \tag{MA-PR.166}
\]
This includes the removable value at \(z=w\).
The Gaussian factor is bounded on these pairs,
and the inner perimeter is bounded.  Equations
(MA-PR.161), (MA-PR.165) and (MA-PR.166) yield
\[
 \sup_{z\in\Gamma_n}|\mathcal H_nF(z)|
       \le C\nu_n^{D_2}e^{U_n}
                        \sup_{\partial\Omega_n^{\rm in}}|F|.
 \tag{MA-PR.167}
\]
Taking the infimum over representatives gives
the same bound from \(J_n^{\rm in}\) to the
original \(J_n\).  Equivalently, (MA-PR.167)
directly applies to \(F_h\); (MA-PR.160) then
bounds the operator on \(H\) by
\(C\nu_n^{D_2+1}e^{U_n}\).
Its kernel contains \(J_H\), and (MA-PR.162)
identifies its quotient map.  Infimizing over
the representatives of \(\pi h\) proves (MA-PR.154).
This is a bound on \(\iota_n\mathscr R_n\) on
all of \(Q_H\).  No separate estimate on an
oblique \(P_n\), or scalar interpolation surrogate,
has replaced that norm.

\emph{4. The stated zero-count conditions and the exact loss.}
For \(t\in[h_n,h_{n+1}]\), every zero in
\(|e+it-\rho|<1\) belongs to \(\mathcal Z_n\).
The layer-cake identity for this finite multiset is
\[
 \sum_{\rho\in\mathcal Z_n}m_\rho
          \log^+\frac1{|e+it-\rho|}
                =\int_0^1N_e(t,r)\,\frac{dr}{r}.
 \tag{MA-PR.168}
\]
Equation (MA-PR.151) makes this count zero when
\(r<\varepsilon_n:=\min(1,c_*\nu_n^{-2})\), for a fixed
\(c_*>0\).  Under (MA-PR.155), the integral is at most
\[
 A_0\log(1/\varepsilon_n)+A_1(1-\varepsilon_n)\log\nu_n
                                      =O(\log\nu_n).
 \tag{MA-PR.169}
\]
Under (MA-PR.156), the at most \(M\) near-edge zeros
cost \(M\log(1/\varepsilon_n)\).  Every other zero in the
window costs at most \(\log^+(1/\eta)\); the ordinary
window count makes their sum \(O(\log\nu_n)\).
This again bounds \(U_n\) by \(O(\log\nu_n)\).

If only the ordinary bound
\(N_e(t,r)\le C\log\nu_n\) is available, (MA-PR.168)
instead gives
\[
       \int_{\varepsilon_n}^1 C\log\nu_n\,\frac{dr}{r}
                              =O(\log^2\nu_n).
 \tag{MA-PR.170}
\]
This is the explicit product-of-multiplicities loss;
the whole-cluster construction has not silently
replaced it by a fixed exponent.  Formula (MA-PR.154)
now gives (MA-PR.157).  If the stronger potential
bound holds, it gives the polynomial cost condition
of MA-PR-010, whose complete proof supplies all
the asserted actual returns and topology.
\end{proof}

\paragraph{A separate exact single-jet lower bound.}
The potential is sufficient, not asserted necessary.
There is nevertheless an intrinsic norm lower bound
which is useful when assessing possible improvements.
For this separate statement, (MA-PR.151) may be
replaced by the no-proper-truncation assumption
(MA-PR.100) on every primary factor, so that
(MA-PR.152) remains defined.
Let \(\rho=\beta+i\gamma\) be any fully retained zero
in block \(n\), and \(0\le j<m_\rho\).  Put
\[
 d_\Omega(\rho)
      =\min\{1,\operatorname{dist}(\rho,\Gamma_n)\}.
\]
Using the actual norms of (MA-PR.152), one has
\[
 \mathcal C_n\ge
 \frac{d_\Omega(\rho)^j}{j!}
 \left(\int_{\mathbb R}
       \frac{x^{2j}e^{2(\beta-a)x}}{w_{u,v,s}(x)}\,dx\right)^{1/2}.
 \tag{MA-PR.171}
\]
The hypothesis that this factor is fully retained
ensures that every displayed jet is continuous.
This statement also applies when a fully retained
boundary factor exists, although the main upper
bound above was proved with the generic hypothesis.

To prove (MA-PR.171), Cauchy's estimate on disks
of radius increasing to \(d_\Omega(\rho)\) gives
\[
        \left|\frac{F^{(j)}(\rho)}{j!}\right|
               \le d_\Omega(\rho)^{-j}\|[F]\|_n .
 \tag{MA-PR.172}
\]
This is well defined on the full holomorphic quotient.
Its composite with \(\iota_n\mathscr R_n\) is exactly
the retained Hilbert jet in (MA-PR.46).
That functional annihilates \(J_H\), so its norm
on \(Q_H\) equals its weighted Hilbert dual norm;
the square of that norm is the integral in
(MA-PR.171) divided by \(j!^2\).
Taking the supremum over the whole quotient unit
ball proves the lower bound.

The original selected-height product bound also gives
fixed \(c,D_3>0\) such that
\[
              d_\Omega(\rho)^j\ge c\nu_n^{-D_3}
                         \qquad(0\le j<m_\rho).
 \tag{MA-PR.173}
\]
Indeed each horizontal distance, truncated at one
and raised to \(m_\rho\), is at least the entire
product in (M18.7).  Both adjacent products are
bounded below by a fixed power of \(\nu_n\).
The vertical boundaries are at least one from
the critical strip.  Hence
\(d_\Omega(\rho)^{m_\rho}\ge c\nu_n^{-D_3}\);
since \(0<d_\Omega\le1\) and \(j<m_\rho\), this
implies (MA-PR.173).

For a strictly interior zero at right-edge distance
\(\delta=e_+-\beta>0\), and \(2j-s+1>0\), restricting
the integral to \(x\ge1\) gives the more explicit
bound
\[
 \mathcal C_n\ge
 \frac{2^{-s/4}d_\Omega(\rho)^j}{j!}\,
 (2\delta)^{-j+(s-1)/2}
                  \Gamma(2j-s+1,2\delta)^{1/2},
 \tag{MA-PR.174}
\]
where \(\Gamma(q,z)=\int_z^\infty t^{q-1}e^{-t}\,dt\).
This follows from
\((1+x^2)^{s/2}\le2^{s/2}x^s\) for \(x\ge1\).
The left edge has the identical bound with
\(\delta=\beta-e_-\), by using \(x\le-1\).
These inequalities do not assert any growing
multiplicity or very small distance for actual
arithmetic zeros.  They identify an actual
whole-quotient jet norm which cannot be deleted
merely by replacing separate interpolation
formulas by a whole-cluster formula.

\paragraph{Remaining arithmetic scope and source placement.}
The ordinary count and the generic distance are
proved arithmetic inputs; their consequence here
is (MA-PR.157).  The stronger local disk or edge-band
conditions are additional explicitly stated sufficient
conditions.  Their validity at a particular arithmetic
width, or polynomial costs without them, is not
established.  In particular this argument does not
close the unrestricted PR011 gate from absence of
proper truncation alone.  It does not prove a
full common-space matrix theorem, replace the
boundary distributions of MA-MO-027/030, or remove
the discarded-divisor trace correction.

The inputs are the complete MT-M18 block theorem,
selected-height and minimum-modulus arguments,
the actual inner restriction and contour construction
MA-PR-001, the whole-quotient Gram and module
identifications MA-PR-006, and the proved
polynomial-cost implication MA-PR-010.
The new estimates keep those actual norms throughout.
Their finite-length module and nuclear-trace meaning
is the one already specified in (MA-PR.80) and
MA-PR-009--011.  No new reading of the inherited
cote~1 scans, additional outside arithmetic premise,
or Grothendieck-group splitting is asserted.

\section{An actual whole-cluster lower bound for the quotient cost}
\label{sec:cluster-quotient-cost}

The single-jet estimate in MA-PR-012 can be replaced by
a functional on an entire cluster.  The functional has
polynomial norm in the original full block algebra,
irrespective of the separations inside that cluster.
Its composite with the canonical Hilbert return measures
an actual weighted exponential-polynomial norm.  This
gives a conditional obstruction to polynomial costs even
for clusters of distinct simple zeros.  No occurrence
of such clusters for the fixed arithmetic function is
asserted.

Retain the fixed data of (MA-PR.1), the original
MT-M18 rectangles and quotient norms, and the
no-proper-truncation condition (MA-PR.100).  Thus
\[
 \begin{gathered}
 e_-=a+u,\qquad e_+=a+v,\qquad p=\frac{s+1}{2},
 \qquad \nu_n=2+|n|,\\
 \mathcal C_n
   =\|\iota_n\mathscr R_n\|_{\mathcal L(Q_H,J_n)} .
 \end{gathered}
 \tag{MA-PR.175}
\]
The inclusion \(\iota_n\) is the identity at every
fully retained primary factor and zero at every
discarded factor.  The source norm in (MA-PR.175)
is the norm on the whole \(Q_H\).  No generic edge
distance, simplicity, or separation between zeros
is assumed below.

\paragraph{MA-PR-013 (cluster depth in the actual quotient norm).}
Let \(\rho_1,\ldots,\rho_k\), \(k\ge1\), be a multiset
of fully retained zeros in one original block
\(\Omega_n\).  A zero may be repeated no more than
its actual multiplicity.  Suppose that, for some
\(0<d\le1\) and \(T\in\mathbb R\),
\[
       |\rho_j-(e_+-d+iT)|\le d/2
                      \qquad(1\le j\le k).
 \tag{MA-PR.176}
\]
There are constants \(c_s>0\) and \(C,D>0\),
depending only on the fixed arithmetic function,
Hilbert parameters and selected block system, for which
\[
 \mathcal C_n\ge
       c_s C^{-1}\nu_n^{-D}
                   3^{-k}k^{-s/2}d^{\,p-k}.
 \tag{MA-PR.177}
\]
In particular, after changing the fixed constants,
\[
                  \mathcal C_n\ge
                         c\,\nu_n^{-D'}d^{\,p-k}.
 \tag{MA-PR.178}
\]
Both estimates hold with the disk in (MA-PR.176)
replaced by the disk centered at \(e_-+d+iT\)
of radius \(d/2\).  Constants are independent of
the cluster size, its internal separations, and
its depth \(d\).

Here is a consequence in which membership in a
single block need not be assumed in advance.
Suppose there is a sequence of multisets of actual
zeros in the disks of (MA-PR.176), with sizes \(k_j\),
centers \(T_j\), and depths \(d_j\), such that
\[
 \begin{gathered}
 |T_j|\longrightarrow\infty,\qquad
 d_j\longrightarrow0,\qquad
 \frac{(k_j-p)\log(1/d_j)}
                     {\log(2+|T_j|)}
                         \longrightarrow+\infty .
 \end{gathered}
 \tag{MA-PR.179}
\]
The same assertion is allowed at the left edge.
After a finite prefix, every such multiset lies
in one original block, consists of strictly
interior fully retained zeros, and forces failure
of every polynomial bound on \(\mathcal C_n\).
Consequently
\[
 \{g\in\mathcal D_{[1,2]}:
                    \sup_n c_n(g)=\infty\}
       \quad\hbox{is a dense }G_\delta
                       \hbox{ in }\mathcal D_{[1,2]} .
 \tag{MA-PR.180}
\]
Every test in that set has \(g(0)=0\) and has no
continuous equivariant return of (MA-PR.51).
This conclusion concerns the actual quotient-return
problem.  It does not identify it with the distinct
full common-space matrix problem of MA-MO-033.

\begin{proof}
\emph{1. A cluster functional on the original full block.}
Set \(P(z)=\prod_{j=1}^k(z-\rho_j)\), with all the
chosen repetitions retained.  For \(F\in A(\Omega_n)\)
define
\[
       \Delta_P[F]
          =\frac1{2\pi i}
                   \int_{\Gamma_n}\frac{F(z)}{P(z)}\,dz .
 \tag{MA-PR.181}
\]
If \(F\in I_n\), then it vanishes at each chosen
node to at least its chosen multiplicity.  Thus
\(F/P\) is holomorphic throughout the rectangle,
and (MA-PR.181) is zero.  The functional is therefore
well defined on the original \(J_n=A(\Omega_n)/I_n\).

On either vertical side of \(\Gamma_n\), every
factor \(|z-\rho_j|\) is at least one, since the
arithmetic zeros lie in \(0\le\operatorname{Re}z\le1\)
and the sides are at \(-1\) and \(2\).  On a
horizontal side at height \(h_l\), \(l=n,n+1\),
\[
 \begin{split}
 |P(z)|
 &\ge\prod_{j=1}^k
               \min\{1,|h_l-\operatorname{Im}\rho_j|\}\\
 &\ge\prod_{\rho:\,\Xi(\rho)=0}
               \min\{1,|h_l-\operatorname{Im}\rho|\}^{m_\rho}
       \ge c\,\nu_n^{-D}.
 \end{split}
 \tag{MA-PR.182}
\]
The middle inequality holds because omitted factors
are at most one and no chosen multiplicity exceeds
the actual multiplicity.  The final inequality is
exactly (M18.7), with the two adjacent height indices
comparable to \(n\).  The uniformly bounded perimeter
of \(\Gamma_n\) now gives
\(\|\Delta_P\|_{J_n^*}\le C\nu_n^D\), by taking the
infimum over all representatives \(F\) of a class.
This estimate has no inverse difference between
two cluster nodes.

\emph{2. Identify its composite on the whole Hilbert quotient.}
Define the entire exponential polynomial
\[
          K_P(x)=\frac1{2\pi i}
                    \int_{\Gamma_n}\frac{e^{zx}}{P(z)}\,dz .
 \tag{MA-PR.183}
\]
Its residue expansion uses derivatives at the chosen
zeros up to their chosen multiplicities minus one.
The identity of the canonical inclusion on those
full factors therefore implies
\[
 \Delta_P\iota_n\mathscr R_n(\pi h)
       =\int_{\mathbb R}h(x)e^{-ax}K_P(x)\,dx .
 \tag{MA-PR.184}
\]
One can first verify this for \(h\in\mathcal E\):
its Mellin transform is entire, and integration
over the fixed compact contour may be interchanged
with its absolutely convergent Mellin integral.
The residue formula gives exactly the chosen
finite linear combination of (MA-PR.46).
Every such jet is continuous on \(H\), since its
factor is fully retained.  Density of
\(\mathcal E\) in \(H\) then proves (MA-PR.184)
for every \(h\in H\).  In particular no Mellin
evaluation of a general \(h\in H\) on the outer
contour has been assumed.

The functional in (MA-PR.184) annihilates \(J_H\).
The dual of the Hilbert quotient is isometric to
that annihilator in \(H^*\), so its exact norm is
\[
 \|\Delta_P\iota_n\mathscr R_n\|_{Q_H^*}^{\,2}
     =\int_{\mathbb R}
          \frac{e^{-2ax}|K_P(x)|^2}
                {\max(e^{2ux},e^{2vx})(1+x^2)^{s/2}}\,dx .
 \tag{MA-PR.185}
\]
This equality uses the whole quotient norm, rather
than a norm imposed separately on the algebraic
coordinates \(B_n\).

\emph{3. A separation-free estimate of the positive half-line.}
Under (MA-PR.176), write
\(d_j=e_+-\operatorname{Re}\rho_j\) and
\(\epsilon_j=\operatorname{Im}\rho_j-T\).  Then
\(d/2\le d_j\le3d/2\) and \(|\epsilon_j|\le d/2\).
For \(x\ge0\) put
\[
 \begin{split}
 \kappa_P(x)&=e^{-(e_++iT)x}K_P(x),\\
 \kappa_P&=
       \mathop{*}_{j=1}^k
       \bigl(e^{-(d_j-i\epsilon_j)x}
                         {\bf1}_{[0,\infty)}(x)\bigr).
 \end{split}
 \tag{MA-PR.186}
\]
For completeness, the Laplace transform of
\(K_P|_{[0,\infty)}\), to the right of all its
exponents, is \(1/P(z)\).  To justify the initial
interchange on the whole outer contour, first take
\(\operatorname{Re}z>2\).  Integrating
\(e^{-zx}e^{wx}\) first in (MA-PR.183), the
resulting sum of residues of
\(1/((z-w)P(w))\) is \(1/P(z)\); holomorphy
continues this identity to the right of all exponents.
The convolution
of the \(k\) exponentials has the same Laplace
transform.  Uniqueness of the Fourier transform
after an exponentially decaying twist proves the
identity, including repeated nodes.  Each factor
in (MA-PR.186) is in \(L^1\cap L^2\); its
convolution also is.  With the unnormalized Fourier
transform \(\widehat f(\xi)=\int e^{-i\xi x}f(x)\,dx\),
\[
       \widehat\kappa_P(\xi)
            =\prod_{j=1}^k
                     \frac1{d_j+i(\xi-\epsilon_j)} .
 \tag{MA-PR.187}
\]
For \(|\xi|\le d\), each denominator has modulus
at most \(3d\).  Plancherel therefore gives
\[
     \|\kappa_P\|_2^2
        \ge\frac1\pi\,9^{-k}d^{\,1-2k}.
 \tag{MA-PR.188}
\]
This step handles arbitrary phases within the
disk.  No common-phase estimate of an oscillatory
simplex integral is needed.

The convolution simplex also gives the pointwise
upper bound
\[
       |\kappa_P(x)|
          \le e^{-dx/2}\frac{x^{k-1}}{(k-1)!}
                         \qquad(x\ge0).
 \tag{MA-PR.189}
\]
Let \(\Gamma(q,b)=\int_b^\infty t^{q-1}e^{-t}\,dt\).
The \(L^2\) mass after \(16k/d\) is consequently at most
\[
 \begin{split}
 \int_{16k/d}^\infty|\kappa_P(x)|^2\,dx
 &\le d^{\,1-2k}
                  \frac{\Gamma(2k-1,16k)}{((k-1)!)^2}\\
 &\le d^{\,1-2k}e^{-8k}
             2^{2k-1}\frac{(2k-2)!}{((k-1)!)^2}\\
 &\le \frac18(16e^{-8})^k d^{\,1-2k}
 \le \frac18\,9^{-k}d^{\,1-2k}.
 \end{split}
 \tag{MA-PR.190}
\]
The second inequality uses \(e^{-t}\le
e^{-8k}e^{-t/2}\) for \(t\ge16k\); the third uses
\(\binom{2k-2}{k-1}\le4^{k-1}\).
The last inequality follows from \(e^8>144\).
Combining (MA-PR.188)--(MA-PR.190), and using
\(\pi<4\), leaves at least
\(9^{-k}d^{1-2k}/8\) of the squared norm in
\([0,16k/d]\).

On \(x\ge0\) the maximum in the weight is
\(e^{2vx}\), so the positive half of
(MA-PR.185) is exactly the weighted squared
norm of \(\kappa_P\).  Since \(16k/d\ge1\),
\[
 \begin{split}
 \int_0^\infty
       |\kappa_P(x)|^2(1+x^2)^{-s/2}\,dx
 &\ge \frac{2^{-s/2}}{8\cdot16^s}\,
            9^{-k}k^{-s}d^{\,s+1-2k},\\
 \|\Delta_P\iota_n\mathscr R_n\|
 &\ge c_s\,3^{-k}k^{-s/2}d^{\,p-k}.
 \end{split}
 \tag{MA-PR.191}
\]
Together with the actual \(J_n^*\) bound following
(MA-PR.182), this proves (MA-PR.177).

The ordinary arithmetic block count (M18.2) gives
\(k\le C_0\log\nu_n\).  Thus \(3^{-k}\) is
bounded below by a fixed negative power of
\(\nu_n\), and, for each fixed \(s\), so is
\(k^{-s/2}\).  Absorbing these factors proves
(MA-PR.178) with constants independent of \(k\).

\emph{4. The left edge.}
If the cluster is centered at \(e_-+d+iT\), use
the negative half of (MA-PR.185), put \(x=-y\),
and remove the exponent \(e^{-(e_-+iT)y}\).
Changing the residue variable from \(z\) to \(-z\)
gives
\[
 e^{(e_-+iT)y}K_P(-y)
   =(-1)^{k-1}
      \mathop{*}_{j=1}^k
       \bigl(e^{-(\rho_j-e_--iT)y}
                         {\bf1}_{[0,\infty)}(y)\bigr).
 \tag{MA-PR.192}
\]
The real decay rates again lie between \(d/2\)
and \(3d/2\), and the imaginary deviations have
absolute value at most \(d/2\).  On the negative
half-line the maximum in the weight is
\(e^{2ux}\).  Hence it cancels the removed
exponent exactly as before.  All the estimates
(MA-PR.187)--(MA-PR.191) apply unchanged in
absolute value, proving the left-edge assertion.

\emph{5. An arbitrarily small bad cluster cannot cross a selected boundary.}
Consider (MA-PR.179), and delete a finite prefix
so that \(d_j<\min(1/4,2(v-u)/3)\).  The entire
disk then lies strictly between \(e_-\) and
\(e_+\).  Every zero in it is retained in full
by the actual Hilbert jet classification.
The chosen zero ordinates occupy an interval
of length at most \(d_j\).

If they do not all belong to one block, a
selected height \(h_l\) lies strictly between
two of their ordinates.  No ordinate equals
\(h_l\).  Every one of the \(k_j\) chosen
ordinates then has distance at most \(d_j\)
from \(h_l\), while \(2+|l|\) is comparable
to \(2+|T_j|\).  Applying (M18.7) to this
subset of factors gives
\[
       c(2+|T_j|)^{-D}
             \le d_j^{\,k_j},
 \qquad
       k_j\log(1/d_j)
             \le D\log(2+|T_j|)+O(1).
 \tag{MA-PR.193}
\]
This contradicts (MA-PR.179), since
\(k_j\ge k_j-p\) and \(p>0\).
Thus eventually every cluster lies in one
original block \(\Omega_{n_j}\), with
\(\nu_{n_j}\) comparable to \(2+|T_j|\).
The already proved estimate now gives
\[
 \frac{\log\mathcal C_{n_j}}
                   {\log(2+|T_j|)}
 \ge
 \frac{(k_j-p)\log(1/d_j)}
                   {\log(2+|T_j|)}
                  -D'+o(1)
       \longrightarrow+\infty .
 \tag{MA-PR.194}
\]
No polynomial cost bound is possible.  The
precisely quantified fixed-support Baire
alternative (MA-PR.104) proves (MA-PR.180).
The same boundary-crossing argument applies
to left-edge clusters.
\end{proof}

\paragraph{Exact scope and source placement.}
The implication just proved is conditional on
the actual cluster geometry (MA-PR.179).
Neither that geometry nor superpolynomial
canonical costs are asserted for the fixed
arithmetic zero divisor.  The ordinary local
zero count alone does not imply this geometry.
Nor does the theorem assert that the sufficient
potential of MA-PR-012 is necessary.

The only zero-count input in the lower bound is
the established block count and the original
selected-height product (M18.2), (M18.7).
The actual quotient and retained-jet interfaces
are (MA-PR.1), (MA-PR.45)--(MA-PR.47), and
the canonical no-proper-truncation inclusion
(MA-PR.100)--(MA-PR.106).  In particular the
argument has not changed block norms, chosen
separated scalar interpolation coordinates, or
deleted an oblique projection.

Repeated roots in \(P\) represent their actual
finite-length primary factors.  The residue
functional respects those full Taylor classes,
but its Hilbert dual norm contains additional
analytic information.  The finite-length,
Grothendieck-group and nuclear-trace scope is
the one already specified in (MA-PR.80) and
MA-PR-009--011; those algebraic descriptions
do not supply a polynomial norm estimate.
No fresh reading of the inherited cote~1
locators or new outside arithmetic premise
is claimed here.

\section{The unrestricted all-test return criterion at every width}
\label{sec:all-width-return-corollary}

The finite obstruction of MA-PR-004 and the polynomial
criterion of MA-PR-010 combine into an exact characterization
without an initial restriction on boundary multiplicities.
All quotients, norms, translations and centering conventions
remain those of (MA-PR.1), with arbitrary fixed \(u<v\),
\(a\in\mathbb R\) and \(s\ge0\).

\paragraph{MA-PR-011 (both necessary gates for unrestricted tests).}
The following conditions are equivalent.
\begin{enumerate}
\item Every \(g\in\mathcal E\) admits a continuous
equivariant return \(S_g:Q_H\to Q_\infty\) satisfying
\[
 \begin{aligned}
          iS_g&=C_g^{Q_H},\\
          S_gU_t^{Q_H}&=e^{-at}U_t^{Q_\infty}S_g
                         \quad(t\in\mathbb R).
 \end{aligned}
 \tag{MA-PR.117}
\]
\item Every compact smooth \(g\) with \(g(0)=0\)
admits a return satisfying (MA-PR.117).
\item Every \(g\in\mathcal D_{[1,2]}\) admits a return
satisfying (MA-PR.117).
\item There is no proper nonzero retained primary factor:
\[
                  r_\rho\in\{0,m_\rho\}
                             \quad\hbox{for every }\rho.
 \tag{MA-PR.118}
\]
Once this holds, form the canonical inclusions and their
actual whole-quotient costs as in (MA-PR.101).  They satisfy
\[
       \mathcal C_n=\|\iota_n\mathscr R_n\|_{Q_H\to J_n}
                    \le C\nu_n^D
                 \quad(n\in\mathbb Z)
 \tag{MA-PR.119}
\]
for some fixed finite \(C>0,D\ge0\).
\end{enumerate}
Condition~2 is equivalently stated for all compact smooth
tests, without the value-zero restriction.
The canonical inclusion in condition~4 is invoked only
after its preceding no-proper-truncation condition is satisfied.

When these conditions hold, the quotient-valued return
is unique.  Its linear choice, jointly continuous
\(\mathcal E\)-valued interpolation lift, and common
all-order nuclear representations are exactly those
proved in MA-PR-010.  If condition~4 fails, the tests
in \(\mathcal D_{[1,2]}\) admitting no such return
contain a dense \(G_\delta\) subset.
Neither validity nor failure of either arithmetic
condition is asserted at a particular width.

\begin{proof}
The test-space inclusions give \(1\Rightarrow2\Rightarrow3\).
Suppose condition~3 holds and that there were a primary
factor with \(0<r_\rho<m_\rho\).  Put
\(\lambda=\rho-a\), choose
\(\eta\in C_c^\infty((1,2))\) with integral one, and set
\[
              g(x)=e^{-\lambda x}\eta(x),
              \qquad g(0)=0,\qquad
              \mathcal Mg(\lambda)=1.
 \tag{MA-PR.120}
\]
This is an actual test in \(\mathcal D_{[1,2]}\).
The finite module necessity in MA-PR-004,
(MA-PR.25)--(MA-PR.29), requires
\(\mathcal Mg(\lambda)=0\) for every return at that
proper factor, since \(\min(m_\rho-r_\rho,r_\rho)\ge1\).
It applies to any global return by the actual finite
Hilbert block and full smooth jet restrictions proved
there.  Equation (MA-PR.120) contradicts that necessity.
Thus condition~3 forces (MA-PR.118).

There are now no eligibility restrictions.  Apply the
\(4\Rightarrow1\) implication of MA-PR-010 to the
entire fixed-support test space in the present condition~3.
In particular its proved Baire bound (MA-PR.114) and
uniform inverse-class argument (MA-PR.115) give
(MA-PR.119).  This proves \(3\Rightarrow4\).
Conversely, (MA-PR.118)--(MA-PR.119) are precisely the
hypotheses of the \(1\Rightarrow2\) implication of
MA-PR-010.  Its actual series (MA-PR.108)--(MA-PR.109)
constructs the return for every \(g\in\mathcal E\),
proving \(4\Rightarrow1\) and the stated uniqueness,
topology and nuclear properties.
The version of condition~2 without \(g(0)=0\) follows
from the same inclusions between conditions~1 and~3.

For the conditional failure statement, first suppose
there is a proper retained factor as above.  The functional
\(g\mapsto\mathcal Mg(\lambda)\) is continuous on
\(\mathcal D_{[1,2]}\), and (MA-PR.120) shows it is
nonzero.  Therefore
\[
       \{g\in\mathcal D_{[1,2]}:
                         \mathcal Mg(\lambda)\ne0\}
 \tag{MA-PR.121}
\]
is open and dense.  Density follows also by adding an
arbitrarily small nonzero multiple of the test in
(MA-PR.120) to any element of its kernel.
Every test in (MA-PR.121) fails the finite return
condition.  This open dense set is itself a
\(G_\delta\).  If instead (MA-PR.118) holds but the
polynomial estimate fails, the dense \(G_\delta\)
alternative (MA-PR.104) applies.  These cases exhaust
failure of condition~4 and prove the assertion.
\end{proof}

The first gate is algebraic at the actual boundary
jets; the second is a measured analytic growth condition
on the whole Hilbert quotient.  For \(s\le1\) only
the first gate is automatic.  At higher weights,
absence of boundary zeros also removes only that
first obstruction.  No polynomial estimate is inferred
from these facts alone.

This characterizes the actual quotient returns in
(MA-PR.117).  It does not erase the discarded nuclear
correction of (MA-PR.116), identify an original Meyer
common-space projection on a new Hilbert domain, or
prove an arithmetic zero-location or positivity theorem.
The exact proof inputs are the complete local module
obstruction MA-PR-004 and the complete norm/topology
equivalence MA-PR-010; no new outside source is used.

\section{The exact remaining geometric criterion and the actual determinant}
\label{sec:exact-geometric-gate}

The preceding partial-strip theorem allows the remaining question to
be stated for the literal cutoff operators at every polynomial weight.
It no longer has an unproved trace-class or boundary-closure premise.
This section makes that consequence explicit and records the actual
Fredholm determinant of the same quotient.

Use the rational trivial sector, put $a=1/2$, and fix
$\alpha>0$, $s\ge0$. Let $r_\rho=r_\rho(\alpha,s)$ be
(MA-WP.3) for $u=-\alpha$, $v=\alpha$.
For $g\in C_c^\infty(\mathbb R)$ with $g(0)=0$, write
$G(z)=\mathcal Mg(z)$ and
\[
\begin{aligned}
 A_0(g)&=\int_0^\infty
   \frac{e^{-x/2}}{1-e^{-2x}}\bigl(g(x)+g(-x)\bigr)\,dx,\\
 B(g)&=\sum_{p,k\ge1}(\log p)p^{-k/2}
                 \bigl(g(k\log p)+g(-k\log p)\bigr),\\
 P(g)&=G(1/2)+G(-1/2),\qquad
 \mathcal W(g)=A_0(g)+B(g)-P(g),\\
 \mathcal G_{\alpha,s}(g)&=
      \lim_R\operatorname{Tr}(M_{\chi_R}C_gP_G).
\end{aligned}
\tag{MA-EG.1}
\]
The cutoffs satisfy the exact hypotheses of MA-WP-003.
The limit exists unconditionally by that theorem, and does not
depend on their particular shape or exhaustion rate.
The prime sum in (MA-EG.1) is finite. The archimedean integral
is ordinary convergent at zero because $g(0)=0$.

\paragraph{MA-EG-001 (criterion for the actual unchanged geometric identity).}
For these actual operators,
\[
 \mathcal G_{\alpha,s}(g)-\mathcal W(g)
    =\sum_\rho (m_\rho-r_\rho)G(\rho-1/2).
 \tag{MA-EG.2}
\]
In particular, the following are equivalent:
\begin{enumerate}
\item $\mathcal G_{\alpha,s}(g)=\mathcal W(g)$ for every compact
smooth $g$ with $g(0)=0$;
\item $r_\rho=m_\rho$ for every nontrivial zero;
\item no zero lies strictly outside the closed centered strip,
and every zero on its boundary has $m_\rho\le b_s$,
where $b_s$ is (MA-WP.2).
\end{enumerate}
If a zero fails condition~2, a compact smooth witness with
$g(0)=0$ and nonzero difference in (MA-EG.2) exists.
Requiring condition~1 for every $\alpha>0$, with any fixed
$s\ge0$, is exactly the Riemann zero-location statement.
This is an equivalence with an unproved independent-method step,
not a new zero-location proof.

\begin{proof}
The classical full explicit formula in the normalization of
MA-TR-2--5, equivalently Meyer's Theorem~5.8 and its spectral
identification Theorem~5.11, gives
\[
             \mathcal W(g)=-\sum_\rho m_\rho G(\rho-1/2).
\]
MA-WP-003 gives the actual operator limit
$\mathcal G_{\alpha,s}(g)=-\sum_\rho r_\rho G(\rho-1/2)$.
Both sums converge absolutely. Subtraction proves (MA-EG.2).
The equivalence of conditions~2 and~3 is exactly (MA-WP.3),
and condition~2 implies condition~1.

For the converse, the two spectral functionals extend continuously
to $\mathcal E$ and annihilate the centered smooth ideal
$I_{1/2}$: the local zero count and rapid vertical Mellin decay
prove continuity, and every ideal element has all the required
zero jets. The restricted-trace criterion MT-M30, (M30.13)--(M30.14),
therefore applies with weights $r_\rho$.
Here is its explicit fixed-sector witness, to make clear what
happens at a lost boundary multiplicity. If
$m_{\rho_0}-r_{\rho_0}>0$, MT-M18 gives $f\in\mathcal E$
whose centered Mellin value is one at $\rho_0$ and zero at every
other zero, with the prescribed higher jets. Choose a nonzero
$\theta\in I_{1/2}$ and put
\[
 b=\frac{\theta*\theta^*}{\|\theta\|_2^2},\qquad
 \theta^*(x)=\overline{\theta(-x)}.
\]
The ideal property gives $b\in I_{1/2}$ and direct evaluation
gives $b(0)=1$. Thus $f_0=f-f(0)b$ vanishes at zero and has
exactly the same spectral values as $f$. Choose a fixed
$\kappa\in C_c^\infty$ equal to one near zero and put
$\kappa_R(x)=\kappa(x/R)$. Then
$\kappa_R f_0\to f_0$ in $\mathcal E$. Its spectral difference
tends to the nonzero number $m_{\rho_0}-r_{\rho_0}$.
For all sufficiently large $R$, $g=\kappa_R f_0$ is the required
compact witness. This proves the converse and the witness assertion.
The construction is conditional on an actually omitted zero;
it does not assert that such a zero exists for $\zeta$.

If all zeros are on the central line, every positive-width
strip contains them strictly, regardless of their multiplicities.
Conversely, any zero with $|\operatorname{Re}\rho-1/2|>0$
is exterior for a smaller positive width, contradicting
condition~2 for that width. This proves the final equivalence.
\end{proof}

The same proof works in each fixed Dirichlet or finite-order Hecke
character sector, with its own normalized full Weil distribution
from MA-CH-001 and its own pole term. The full-family norm-one
splitting is precisely MT-M30. Thus the actual cutoff identity in
all rational Dirichlet strips follows using the specified theorem
$A$. Using that theorem to remove the discarded-divisor term is
a zero-location input, hence circular for the requested independent
method test. For the non-induced imaginary-quadratic character of
MA-CH-003, the corresponding Hecke zero-location statement is still
the condition; $M$ supplies it, whereas $A$ does not.

\paragraph{MA-EG-002 (Fredholm determinant of the actual partial quotient).}
In the general fixed-character setting of MA-WP, for every
$g\in\mathcal E$ the actual trace-class quotient convolution has
\[
 \det\bigl(1-wC_g^Q\bigr)
       =\prod_{\Xi(\rho)=0}
           \bigl(1-w\mathcal Mg(\rho-a)\bigr)^{r_\rho},
 \qquad w\in\mathbb C.
 \tag{MA-EG.3}
\]
The product converges locally uniformly. Its derivative at zero
is the negative of the actual trace (MA-WP.20).
It retains every surviving finite-block length, including the exact
boundary truncations, without an assumption of simple zeros or of
a positive operator.

\begin{proof}
Let $A_{n,g}$ be the actual finite-rank operators (MA-WP.13).
Their finite sums $T_N=\sum_{|n|\le N}A_{n,g}$ converge to
$C_g^Q$ in trace norm by (MA-WP.15).
The mutually annihilating bounded projections $P_n$ give a finite
invariant direct sum $\bigoplus_{|n|\le N}V_n$, and $T_N$ is
zero on the kernel of the bounded projection
$\sum_{|n|\le N}P_n$. On a retained local factor
$\mathbb C[z]/(z^{r_\rho})$, multiplication by the Taylor class
of $\mathcal Mg$ is triangular with diagonal
$\mathcal Mg(\rho-a)$. Finite-dimensional determinant theory gives
\[
 \det(1-wT_N)=
   \prod_{\rho\text{ in the blocks }|n|\le N}
        \bigl(1-w\mathcal Mg(\rho-a)\bigr)^{r_\rho}.
\]
Trace-norm continuity of the Fredholm determinant now gives
convergence of the left side locally uniformly in $w$.
An exact reference, including the needed quantitative estimate,
is Kostenko, \emph{Trace Ideals with Applications},
Theorem~3.4.4 and (3.4.19), printed pp.~39--40:
\[
 |\det(1+A)-\det(1+B)|
       \le\|A-B\|_1\exp(1+\|A\|_1+\|B\|_1).
\]
The right-side products converge locally uniformly because
$\sum_\rho r_\rho|\mathcal Mg(\rho-a)|<\infty$, as proved
in MA-WP-002. Taking the limits proves (MA-EG.3).
Differentiation at zero is justified by local uniform convergence
and gives the displayed trace. No claim that the finite projections
are orthogonal has entered the argument.
\end{proof}

\paragraph{Grothendieck placement and the precise method boundary.}
The retained local algebra has length $r_\rho$ over
$\mathbb C$ and defines a finite-length module supported at the
spectral point $\rho-a$ on the complex affine line. Its class in
the Grothendieck group $K_0$ of the abelian category of finite-length
$\mathbb C[z]$-modules records this length at that point.
Exterior powers over $\mathbb C$ of the underlying finite-dimensional
vector spaces, with their induced endomorphisms, produce the finite
block determinants used above. Formula (MA-EG.3) is their nuclear
Fredholm limit in the actual completion. This is the cote~1
comparison at scans~079--080 and~111--112, via
ALPOGE\_BRIDGES BR-TR-01/02 and F-SM-11/12, with the actual
operator rather than only an associated graded used here.
The scan locators are inherited, not newly inspected in this run.
The determinant and trace do not recover the nilpotent extension
structure within $\mathbb C[z]/(z^{r_\rho})$.
No scheme over $\mathbb Z$, arithmetic cohomology, or
$\ell$-adic weight realization is constructed by this analytic
change of completion. In particular it supplies no faithful
full-jet positivity contrary to MT-M20/M23.

The source bridge remains the identification from the unfinished
\S12.3 on Masullo p.~514 to the global use of the local theorem
in (12.5.1.4), p.~520. MA-EG-001 identifies exactly what the
corrected actual objects require at that bridge. It does not
turn a missing proof into a mathematical counterexample under
$A$, and it does not infer nonexistence of every possible new
unbounded intertwiner. The direct bounded inverse obstruction,
the proved conditional full bridge, and the concrete regularized
commutator have their separate scopes in MA-MO, MA-HM and MA-RS.

\paragraph{Source details.}
Kostenko, A. (n.d.). \emph{Trace ideals with applications}
[Lecture notes], Theorem~3.4.4 and (3.4.19), pp.~39--40.
\url{https://users.fmf.uni-lj.si/kostenko/teach/IdealsNotes.pdf}.
The cited theorem and its complete proof were read in the
author-hosted version on 21 November 2026; no unverified publication
year is assigned. Meyer (2005), Theorem~5.8, author pp.~43--44,
and Theorem~5.11, p.~47, are used with the exact normalizations
already checked in MA-MEY. The inherited restricted-test theorem
is GRH\_NIGHT MT-M30, (M30.13)--(M30.14); the new input is
MA-WP's unconditional identification of the literal weighted
quotient and cutoff trace.

\section{The actual global Paley--Wiener condensation and its half-plane range}
\label{sec:global-condensation}

The local quotient calculation does not determine the global orthogonal
projection.  This section computes more of that actual projection.  The
first result extends its fixed-cutoff convergence to every corrected
strip, using MA-WP.  The next two results identify the exact smooth
Paley--Wiener range and its Hilbert closure at polynomial weight zero.
The latter is a weighted version of the theta-image Hardy-space method,
with the source's particular range and centering retained.

Write
\[
 \mathscr S_{0,+}=\{f\in\mathscr S(\mathbb R):f(-t)=f(t),
                         \ f(0)=\widehat f(0)=0\},\qquad
 (\Theta f)(x)=2e^{x/2}\sum_{n\ge1}f(ne^x).
 \tag{MA-GP.1}
\]
The Fourier transform is the unitary transform with kernel
$e^{-2\pi it\xi}$.  MA-TR-101 identifies $\Theta$ with an actual
topological isomorphism onto $T_{1/2}\mathcal I$.
For $Y>0$ define
\[
 \begin{aligned}
 \mathscr S_Y&=\{f\in\mathscr S_{0,+}:
                       \operatorname{supp}\widehat f\subseteq[-Y,Y]\},\\
 \mathscr P_Y&=\Theta\mathscr S_Y,\qquad
 J_Y=\overline{\mathscr P_Y}^{\,H_{\alpha,s}},\qquad
 Q_Y^{\rm glob}=P_{J_Y}.
 \end{aligned}
 \tag{MA-GP.2}
\]
For $S=\{\infty\}$ these are the rational trivial-sector
Paley--Wiener levels of Masullo, Definition~7.3.2.1,
(7.3.2.1)--(7.3.2.7), p.~351, and Definition~7.3.2.3,
(7.3.2.12)--(7.3.2.17), pp.~354--355: the absolute and
invariant support conditions coincide in this one-place case.
The Hilbert projections are Definitions~11.2.1.1 and~11.2.1.3,
pp.~487--489, on the corrected actual theta image.  The notation
$Q_Y^{\rm glob}$ is deliberately different from the local additive
Fourier multiplier $Q_Y$ in MA-LA-014.

\paragraph{MA-GP-001 (actual fixed-cutoff condensation on every strip).}
For every $\alpha>0$ and $s\ge0$,
\[
 Q_Y^{\rm glob}\longrightarrow P_G\quad\hbox{strongly as }Y\to\infty.
 \tag{MA-GP.3}
\]
For $\chi,g\in C_c^\infty(\mathbb R)$ the actual operators satisfy
\[
 \bigl\|M_\chi C_gQ_Y^{\rm glob}
                  -M_\chi C_gP_G\bigr\|_1\longrightarrow0.
 \tag{MA-GP.4}
\]
For the fixed exhaustion $\chi_R(x)=\chi_0(x/R)$ of MA-WP-003,
with $g(0)=0$, the iterated trace limit is therefore
\[
 \lim_{R\to\infty}\lim_{Y\to\infty}
        \operatorname{Tr}(M_{\chi_R}C_gQ_Y^{\rm glob})
       =-\sum_\rho r_\rho\,\mathcal Mg(\rho-1/2).
 \tag{MA-GP.5}
\]
The retained multiplicities $r_\rho$ are exactly those of MA-WP.
For each fixed test $g$ one can also choose a diagonal
$Y=Y(R)\to\infty$ with the same limit.  This assertion gives no
prescribed quantitative relationship between $Y$ and $R$, and
does not identify this diagonal with the local coupled cutoffs.
No zero-location input is used.

\begin{proof}
Here is a direct density argument respecting both polar conditions.
Choose an even $\kappa\in C_c^\infty$ equal to one near zero.
For $f\in\mathscr S_{0,+}$ set
$f_N=\mathcal F^{-1}(\kappa(\,\cdot/N)\widehat f)$.
Then $f_N\to f$ in Schwartz space, $f_N$ is even and
$\widehat f_N(0)=0$, but $f_N(0)$ need not be zero.
Choose an even $v$ whose Fourier transform is compactly supported
away from zero and has integral one.  Thus
$v(0)=1$ and $\widehat v(0)=0$.
The functions
\[
                         \widetilde f_N=f_N-f_N(0)v
 \tag{MA-GP.6}
\]
have both polar conditions, are band limited, and converge to $f$
in Schwartz space.  Continuity of the actual theta map into
$\mathcal E$, followed by its continuous inclusion into
$H_{\alpha,s}$, proves that $\bigcup_YJ_Y$ is dense in $J$.
The spaces $J_Y$ increase.  For $j\in J$, approximate it by
$j_0\in J_{Y_0}$; then
$\|(1-Q_Y^{\rm glob})j\|\le\|j-j_0\|$ for $Y\ge Y_0$.
On $J^\perp$ every $Q_Y^{\rm glob}$ is zero.  This proves
(MA-GP.3).

MA-WP-003 and the actual ambient lift of MA-LIFT give
$T=M_\chi C_gP_G\in\mathcal S_1(H_{\alpha,s})$.
Since $P_GQ_Y^{\rm glob}=Q_Y^{\rm glob}$ and $TP_G=T$,
the difference in (MA-GP.4) is $T(Q_Y^{\rm glob}-P_G)$.
For a rank-one operator $u\otimes\ell$, its trace norm after
this composition is $\|u\|\,\|(Q_Y^{\rm glob}-P_G)^*\ell\|$,
which tends to zero by self-adjoint strong convergence.
Finite-rank approximation of $T$, with the uniform bound
$\|Q_Y^{\rm glob}-P_G\|\le1$, proves (MA-GP.4).
This is the usual trace-ideal continuity argument, now applied
to the actual all-weight operators rather than to an unproved
ambient lift.

Take traces and then apply MA-WP-003 to obtain (MA-GP.5).
For each $R\ge1$, (MA-GP.4) permits a choice $Y(R)\ge R$
for which the trace-norm difference is at most $1/R$.
This proves the stated diagonal version without a uniform
rate or a change in the order of the two limits.
\end{proof}

\paragraph{MA-GP-002 (the exact smooth support description).}
For $c=-\log Y$,
\[
 \mathscr P_Y
   =\{g\in T_{1/2}\mathcal I:
                              \operatorname{supp}g\subseteq[c,\infty)\}.
 \tag{MA-GP.7}
\]
This is an equality of actual smooth theta-image spaces.  It
does not assert that taking this intersection commutes with
Hilbert closure.

\begin{proof}
Put $u=\widehat f$.  The two polar conditions and Poisson
summation give the exact identity
\[
 (\Theta f)(x)=2e^{-x/2}\sum_{n\ge1}u(ne^{-x})
                              =(\Theta u)(-x).
 \tag{MA-GP.8}
\]
If $\operatorname{supp}u\subseteq[-Y,Y]$, every term vanishes
when $x< -\log Y$.  This proves one inclusion.
Conversely let $g=\Theta f$ have the support in (MA-GP.7).
Equation (MA-GP.8) gives $\Theta u(\log t)=0$ for $t>Y$.
The actual smooth M\"obius inverse of MA-TR-101 gives
\[
 2\sqrt t\,u(t)
    =\sum_{n\ge1}\mu(n)n^{-1/2}(\Theta u)(\log(nt)).
 \tag{MA-GP.9}
\]
The series converges with all derivatives on compact
positive intervals.  If $t>Y$, all its arguments satisfy
$nt>Y$, so it is zero.  Evenness gives $u=0$ outside
$[-Y,Y]$.  Thus $f\in\mathscr S_Y$.
The band-support condition is closed in Schwartz space;
MA-TR-101 also shows that $\mathscr P_Y$ is closed in
the smooth theta-image topology, as required by the
source's smooth closure convention.
\end{proof}

\paragraph{MA-GP-003 (the actual finite-level Hardy range).}
Assume $0<\alpha<1/2$ and $s=0$.  On the half-line $[c,\infty)$
first use the equivalent norm
\[
 K_c=L^2([c,\infty),e^{2\alpha x}\,dx),\qquad
 (U_cf)(y)=e^{\alpha(c+y)}f(c+y),\quad y\ge0.
 \tag{MA-GP.10}
\]
Let $\mathcal L$ denote the Laplace isometry onto
$H^2(\operatorname{Re}z>0)$, with boundary norm
$(2\pi)^{-1}\int|F(it)|^2dt$.
There is a pure Blaschke product $B_\alpha$ with zeros
\[
 z_\rho=\alpha+\tfrac12-\rho,
 \qquad \operatorname{Re}\rho<\tfrac12+\alpha,
 \qquad\hbox{multiplicity }m_\rho,
 \tag{MA-GP.11}
\]
such that the actual closed range satisfies
\[
            \mathcal L U_cJ_Y=B_\alpha H^2(\operatorname{Re}z>0).
 \tag{MA-GP.12}
\]
The equality concerns sets with equivalent Hilbert topologies;
the actual symmetric-weight orthogonal projection still needs
the Gram operator computed in the next section.  In particular,
finite Paley--Wiener levels impose all zero jets strictly to
the left of the upper strip edge, including any zeros to the
left of the full strip.  No hypothesis about whether such
exterior zeros occur is used.

\begin{proof}
\emph{The actual closure and shift invariance.}
On $[c,\infty)$ the ratio of the original weight to the weight
in (MA-GP.10) is
\[
 R_c(x)=e^{2\alpha|x|-2\alpha x},\qquad
            1\le R_c(x)\le e^{4\alpha\max(-c,0)}.
 \tag{MA-GP.13}
\]
Thus the two closures of $\mathscr P_Y$ are the same set.
By (MA-GP.7), $\mathscr P_Y$ is invariant under right
translations $g(x)\mapsto g(x-t)$ for $t\ge0$.
The map $U_c$ converts these to $e^{\alpha t}$ times
the unilateral translations on $L^2(0,\infty)$.

For clarity the Laplace normalization and surjectivity can
be checked directly: its reproducing kernel is
$K(z,\zeta)=\int_0^\infty e^{-(z+\overline\zeta)y}dy
=(z+\overline\zeta)^{-1}$.
With $z=(1+w)/(1-w)$, the weighted composition
\[
 F(z)\longmapsto \frac{\sqrt2}{1-w}
                         F\!\left(\frac{1+w}{1-w}\right)
 \tag{MA-GP.14}
\]
has reproducing kernel $(1-w\overline\omega)^{-1}$.
Density of the spans of the evaluation kernels identifies
it isometrically onto disk $H^2$; Plancherel gives the
stated half-plane boundary norm.

Write $\mathcal M_Y$ for the Laplace image of the actual
closure.  Invariance under unilateral translations makes
it invariant under multiplication by $e^{-tz}$, $t\ge0$.
For each vector their strongly continuous Bochner average with
weight $e^{-t}$ converges in Hilbert norm and defines the bounded
operator of multiplication by $(1+z)^{-1}$.
It is consequently invariant under
$(z-1)/(z+1)=1-2/(1+z)$.
Under (MA-GP.14) this is disk multiplication by $w$.
Beurling's invariant-subspace theorem therefore gives
$\mathcal M_Y=\Theta_Y H^2$ for an inner function $\Theta_Y$.
The space is nonzero by any nonzero source test and the
injectivity of MA-TR-101.

\emph{The exact common zero divisor.}
The relation between Laplace and Mellin transforms is
\[
 (\mathcal L U_cg)(z)=e^{zc}\mathcal Mg(\alpha-z).
 \tag{MA-GP.15}
\]
All core functions therefore vanish with the multiplicities
in (MA-GP.11), by the full smooth ideal description.
These are all the common zeros, with no extra multiplicity.
To verify that assertion using actual source tests, take
an even $u\in C_c^\infty((-Y,Y)\setminus\{0\})$ with
$\int u=0$, and put $f=\mathcal F^{-1}u$.
It belongs to $\mathscr S_Y$.  For
$s_z=z+1/2-\alpha$, termwise Mellin integration first in
$\operatorname{Re}s_z>1$, followed by analytic continuation,
gives
\[
 (\mathcal L U_c\Theta f)(z)
       =2e^{zc}\zeta(s_z)
                    \int_0^Y u(t)t^{s_z-1}\,dt.
 \tag{MA-GP.16}
\]
Here $\operatorname{Re}s_z>1/2-\alpha>0$ throughout the
half-plane; the zero of the integral at $s_z=1$ removes
the simple pole of zeta.
For any prescribed $s_z\ne1$ the integral can be made
nonzero subject to $\int_0^Y u=0$: the densities
$1$ and $t^{s_z-1}$ are linearly independent on any
positive open interval.  At $s_z=1$ one can instead
make $\int u(t)\log t\,dt\ne0$, so the removable value
of the product is nonzero.
Thus its common zeros are exactly the nontrivial zeros
of $\zeta(s_z)$ in this half-plane, to their exact orders.
The functional equation identifies that multiset with
(MA-GP.11), through $\rho\mapsto1-\rho$.
This uses the functional equation and the classical
containing strip, not RH or any narrower zero location.

\emph{No singular inner factor or positive delay.}
Every $g\in\mathscr P_Y$ is smooth, flat at $c$, and
decreases faster than all exponentials as $x\to\infty$.
Consequently $v=U_cg$, extended by zero to the negative
half-line, is smooth, flat at zero, and has all derivatives
in $L^1(0,\infty)$.
Integration by parts gives, for all $N$,
\[
 |(\mathcal Lv)(z)|\le C_N(1+|z|)^{-N}
                       \quad(\operatorname{Re}z\ge0).
 \tag{MA-GP.17}
\]
Its Laplace transform is entire by the same exponential
decay.  Therefore its disk image in (MA-GP.14) is bounded
and extends holomorphically across every circle arc
except the point $w=1$ corresponding to infinity.

Fix one nonzero core function.  The singular measure of
the common inner factor is bounded above by that of this
bounded disk function.  The elementary Poisson-kernel
argument in MA-HY-001 shows that it is supported on the
isolated boundary zeros of that function and at $w=1$:
on an arc with a nonzero analytic continuation,
boundedness of its Poisson integral makes the singular
measure absolutely continuous and hence zero there.
At an isolated zero of finite order, a positive atom
would force decay $\exp(-C/(1-r))$ along a radius,
contradicting the finite-order analytic lower bound.
Thus only an atom at $w=1$ can remain.  In half-plane
coordinates its factor is $e^{-dz}$ for some $d\ge0$.

Such a common factor would force every $U_cg$ to vanish
on $[0,d)$.  Indeed multiplication by $e^{-dz}$ under
the Laplace isometry is precisely the unilateral shift
by $d$, as is seen first by direct integration and then
by $L^2$ continuity and uniqueness.
There is a core function with no such delay.
Choose $u$ even and smooth, supported in $[-Y,Y]$, positive
on $(Y-\delta,Y)$ and flat at $Y$, with a compensating
negative bump below $Y/2$ so that $\int u=0$.
It vanishes near zero.  For $x>c$ sufficiently close to
$c$, (MA-GP.8) has only its $n=1$ term, and
\[
                  (\Theta\mathcal F^{-1}u)(x)
                           =2e^{-x/2}u(e^{-x})>0.
 \tag{MA-GP.18}
\]
Hence $d=0$.
The common inner factor is the pure Blaschke product
with the zeros already identified.  Beurling's theorem
now proves (MA-GP.12), including Blaschke summability.
\end{proof}

\paragraph{What the finite-level result does and does not identify.}
The source's support condition has now become an actual Hardy
range, independently of zero location.  For a left-exterior
point $\lambda$ with $\operatorname{Re}\lambda< -\alpha$,
evaluation on the supported half-line is continuous, but its
squared dual norm contains
\[
 \int_c^0e^{2(\operatorname{Re}\lambda+\alpha)x}\,dx,
 \tag{MA-GP.19}
\]
which diverges as $c\to-\infty$.  At equality it is $-c$.
Thus the extra finite-level constraints can disappear when
the union is completed.  This is consistent with both
(MA-GP.3) and MA-WP's actual full-strip multiplicities.
It supplies no license to omit those constraints from
a finite-level kernel or to replace its weighted projection
by the local additive Fourier multiplier.

\paragraph{Primary-source and Grothendieck placement.}
The density argument in Masullo, Proposition~7.3.2.2(1),
pp.~352--353, and the condensation statements
Theorems~11.2.1.2/11.2.1.4, pp.~488--489, were read at
their stated scope.  MA-GP-001 applies the repaired all-weight
ambient trace theorem; it does not rely on the source's
unrepaired nuclear lift.
The Hardy input is Beurling (1949), Factorization Lemma,
\S7, pp.~245--246, and Theorem~I, \S9, pp.~247--249,
already checked in MA-HY.  A closely related unweighted
theta-image closure is Burnol (2001; author manuscript 2000), \S2, the proof preceding
Theorem~2.6, author pp.~9--10; its adelic version is
Theorem~3.2, author pp.~14--15.  The present proof supplies
the shifted norm, the actual source support condition,
and the finite-level projection interface needed here.
No priority claim is made for the Hardy-space method.

In Grothendieck's nuclear framework, (MA-GP.4) is an actual
trace-ideal approximation before the Fredholm trace is taken,
as in the cote~1 comparison, scans~079--080 and~111--112,
ALPOGE\_BRIDGES BR-TR-01/02.  Those archive locators remain
inherited.  The finite-length modules of MA-WP still give
the spectral multiplicities; passing through an equivalent
half-line norm does not construct a new arithmetic motive
or an $\ell$-adic weight realization.

Burnol, J.-F. (2001). An adelic causality problem related
to abelian $L$-functions. \emph{Journal of Number Theory, 87}(2),
253--269. \url{https://doi.org/10.1006/jnth.2000.2616}.
Author manuscript (2000), version~3:
\url{https://arxiv.org/abs/math/0001013}.

\section{The finite Paley--Wiener projection and its Gram operator}
\label{sec:global-gram}

The finite-level orthogonal projector is a concrete additional
operator.  The same range can be parametrized by a local theta
map, by a Hardy-space inner multiplier, or by frequency data,
but these parametrizations are not automatically isometries for
the specified global norm.  This section gives the actual Gram
factor, with its domain and inverse, and an explicit finite-level
theta Gram kernel.  It does not infer an unbounded finite-level
inverse from the uncut obstruction MA-MO.

\paragraph{MA-GP-004 (the exact bounded Gram correction on a half-line).}
Fix $\alpha>0$, $s\ge0$ and $c\in\mathbb R$.  On $[c,\infty)$
put
\[
 \begin{aligned}
 w_+(x)&=e^{2\alpha x}(1+x^2)^{s/2},&
 w(x)&=e^{2\alpha|x|}(1+x^2)^{s/2},\\
 \mathcal K_+&=L^2([c,\infty),w_+(x)\,dx),&
 R(x)&=w(x)/w_+(x)=e^{2\alpha|x|-2\alpha x}.
 \end{aligned}
 \tag{MA-GP.40}
\]
Let $\mathcal M\subset\mathcal K_+$ be the closed finite-level
subspace supplied by the actual smooth PW range, and let
$i:\mathcal M\hookrightarrow\mathcal K_+$ be its isometric
inclusion for the $w_+$ norm.  The two norms on this half-line
are equivalent, since
\[
            1\le R(x)\le C_c:=e^{4\alpha\max(-c,0)}.
 \tag{MA-GP.41}
\]
Consequently either norm gives the same range closure.  The
positive operator
\[
              G_c=i^*M_Ri:\mathcal M\to\mathcal M
 \tag{MA-GP.42}
\]
is boundedly invertible, with $1\le G_c\le C_c$ and
$C_c^{-1}\le G_c^{-1}\le1$.  The orthogonal projection for
the \emph{actual} $w$ norm is exactly
\[
                  Q_c^{(w)}=iG_c^{-1}i^*M_R.
 \tag{MA-GP.43}
\]
On the whole global space first restrict the input to
$[c,\infty)$ and extend the output by zero.  Formula
(MA-GP.43) then gives the actual global finite-level projector.

\begin{proof}
The weight ratio is (MA-GP.41), including for every polynomial
weight $s$.  For $u\in\mathcal M$,
$\langle G_cu,u\rangle_+=\|iu\|_w^2$, which lies between
$\|u\|_+^2$ and $C_c\|u\|_+^2$.  The spectral theorem
therefore gives the asserted bounded inverse.
For $h\in\mathcal K_+$ the vector in (MA-GP.43) belongs
to $\mathcal M$, and for every $v\in\mathcal M$,
\[
 \langle h-iG_c^{-1}i^*M_Rh,iv\rangle_w
   =\langle i^*M_Rh-G_cG_c^{-1}i^*M_Rh,v\rangle_+=0.
 \tag{MA-GP.44}
\]
This is the defining property of the weighted orthogonal
projection.  The two support half-lines are orthogonal in
the ambient weighted space, so restriction and zero extension
prove the final assertion.
\end{proof}

For $s=0$ the actual Hardy identification of MA-GP-003,
in its stated range $0<\alpha<1/2$ and $c=-\log Y$,
makes this formula explicit.  Write
\[
 (U_cf)(r)=e^{\alpha(c+r)}f(c+r),\quad r\ge0,
 \qquad
 (\mathcal Lv)(z)=\int_0^\infty e^{-zr}v(r)\,dr.
 \tag{MA-GP.45}
\]
The boundary norm $(2\pi)^{-1}\int|F(it)|^2\,dt$, as in
MA-GP-003, makes $\mathcal L$ unitary onto
$H^2(\operatorname{Re}z>0)$.  If $M_{B_\alpha}$ is the
isometric inner multiplier describing the actual finite PW
range, define
\[
 \mathcal R_c=\mathcal L U_cM_RU_c^{-1}\mathcal L^{-1},
 \qquad
 \mathcal G_c=M_{B_\alpha}^*\mathcal R_cM_{B_\alpha}.
 \tag{MA-GP.46}
\]
In this Hardy picture the projector is
\[
 M_{B_\alpha}\mathcal G_c^{-1}
                   M_{B_\alpha}^*\mathcal R_c.
 \tag{MA-GP.47}
\]
The zeros of $B_\alpha$ are those of MA-GP-003, namely
$z=\alpha+1/2-\rho$ in the right half-plane, with full
multiplicity.  Thus left-exterior global zeros can enter this
finite-level inner factor even when they do not survive in
the completed full-strip quotient of MA-WP.

There is an explicit kernel for the weight correction.
If $c<0$ and $L=-c$, then $\mathcal R_c=1+\mathcal K_c$,
where on Hardy boundary values
\[
 (\mathcal K_cF)(t)=\frac1{2\pi}\int_{\mathbb R}
                       \kappa_c(t-u)F(u)\,du,
 \qquad
 \kappa_c(v)=\int_0^L
            \bigl(e^{4\alpha(L-r)}-1\bigr)e^{-ivr}\,dr.
 \tag{MA-GP.48}
\]
The convolution is interpreted in the bounded $L^2$ sense,
as the Fourier transform of multiplication by the displayed
bounded compactly supported function.  Hence
\[
       \mathcal G_c=1+M_{B_\alpha}^*\mathcal K_cM_{B_\alpha},
       \qquad 1\le\mathcal G_c\le Y^{4\alpha}
                    \quad(Y>1).
 \tag{MA-GP.49}
\]
For $Y\le1$ the ratio $R$ is one and this particular weight
Gram is the identity.  The Hardy inner factor remains part
of the projection.  For $s>0$, (MA-GP.43) remains valid, but
no replacement of the actual Hardy--Sobolev range by the
$s=0$ inner-multiplier model is asserted.

\paragraph{The Gram factor changes the projection.}
Suppose $c<0$, $s=0$, and the inner factor has a zero and
has no positive delay, as in MA-GP-003.  Then the weighted
projection (MA-GP.47) differs from the ordinary Hardy
projection $M_{B_\alpha}M_{B_\alpha}^*$.

\begin{proof}
If the projections agreed, their range and ordinary orthogonal
complement would be orthogonal also for the $R$ inner product.
Thus multiplication by $R(c+r)$ would commute with the
ordinary projection onto the inverse-Laplace image of
$B_\alpha H^2$.  Its complement would reduce that multiplier.
On $0<r<L=-c$, $R(c+r)=e^{4\alpha(L-r)}$ is strictly
monotone.  Its spectral projections therefore include
multiplication by $\mathbf1_{(a,b)}$ for every
$0<a<b<L$.

If $B_\alpha(\mu)=0$ with $\operatorname{Re}\mu>0$,
the Hardy reproducing kernel at $\mu$ belongs to the
orthogonal complement of $B_\alpha H^2$.  Its inverse
Laplace transform is a nonzero multiple of
$e^{-\overline\mu r}$.
All the vectors $\mathbf1_{(a,b)}e^{-\overline\mu r}$
would consequently belong to that complement.  Their span
is dense in $L^2(0,L)$, since multiplication by this
exponential is boundedly invertible on the finite interval.
Every function in the $B_\alpha$ range would then vanish
on $(0,L)$, contrary to the no-delay conclusion of
MA-GP-003.  This proves the distinction.
\end{proof}

This argument is a precise obstruction to dropping the Gram
factor in the finite projection.  It is not an obstruction to
its existence: the bounded inverse was proved in (MA-GP.42).
Nor does it disprove a later asymptotic trace comparison with
all factors and remainders retained.

\subsection{An explicit theta Gram kernel before the Hardy identification}

\paragraph{MA-GP-005 (the actual frequency parametrization and normal equation).}
The Gram form can also be written directly from the source
PW definition, independently of the Hardy identification.
Take $F=\mathbb Q$, $S=\{\infty\}$ and a finite frequency
level $\Lambda>0$.  Let
\[
 \begin{aligned}
 \mathcal P_\Lambda^0
   &=\{u\in L^2_{\mathrm{even}}([-\Lambda,\Lambda]):
                                       \int_{-\Lambda}^{\Lambda}u=0\},\\
 \mathscr D_\Lambda
   &=\{u\in C_c^\infty((-\Lambda,\Lambda)):
                      u\text{ even},\ u(0)=0,\ \int u=0\}.
 \end{aligned}
 \tag{MA-GP.50}
\]
Use the ordinary even $L^2$ norm in this \emph{auxiliary}
source space.  It is not the weighted local norm in the draft.
For $u\in\mathscr D_\Lambda$, put $f=\mathcal F_eu$.
The two source conditions $f(0)=\int f=0$ hold.
Poisson summation gives the actual centered global theta map
\[
 (T_\Lambda u)(t)
       =2\sqrt t\sum_{n\ge1}f(nt)
       =2t^{-1/2}\sum_{n\ge1}u(n/t),\qquad t>0.
 \tag{MA-GP.51}
\]
Its target norm is
\[
 \|F\|_H^2=\int_0^\infty|F(t)|^2
           w_{\alpha,s}(\log t)\,\frac{dt}{t},
 \qquad w_{\alpha,s}(x)=e^{2\alpha|x|}(1+x^2)^{s/2}.
 \tag{MA-GP.52}
\]
Every vector (MA-GP.51) belongs to $H$, and the closure of
its range is the actual finite source PW range.
The map $T_\Lambda$ is densely defined and closable from
$\mathcal P_\Lambda^0$ into $H$.

Let $0\le\chi\in C_c^\infty(\mathbb R_{>0})$.
The localized map $M_{\sqrt\chi}T_\Lambda$ extends boundedly
to $\mathcal P_\Lambda^0$.  Its bounded positive Gram
operator is the restriction followed by zero-mean projection
of the explicit finite rational-dilation kernel
\[
 G_{\chi,\Lambda}u(y)=P_\Lambda^0\left[
 2\sum_{m,n\ge1}\frac1m\,
       \chi(m/y)w_{\alpha,s}(\log(m/y))u(ny/m)
                         \right],\qquad 0<y<\Lambda.
 \tag{MA-GP.53}
\]
Here $u$ is extended evenly by zero outside the band and
$P_\Lambda^0$ is the ordinary $L^2$ projection onto the
zero-integral even subspace.  All sums are finite for a
fixed compact output cutoff, uniformly in $y$ where they
are nonzero.

\begin{proof}
The test domain is dense in $\mathcal P_\Lambda^0$:
approximate by smooth even functions, cut away a shrinking
neighborhood of zero, and subtract their integral times
one fixed smooth even annular function of integral one.
The corrections tend to zero in $L^2$.
Poisson summation applies to $f$ and its Fourier transform
$u$, since both zeroth moments vanish, and gives (MA-GP.51).
The second expression is zero for $t<1/\Lambda$.
The first decreases faster than every power for $t\to\infty$,
uniformly after every derivative, by the Schwartz estimates
for $f$.  Thus it belongs to every target (MA-GP.52).

For this one-place source the absolute and invariant PW
support conditions coincide with the even Fourier band.
Definitions~7.3.2.1 and~7.3.2.3, pp.~351 and~355,
and Definition~11.2.1.3, pp.~488--489, give precisely the
range closure stated here.  Allowing smooth transforms
whose support meets the endpoints of the closed band
does not change it: a dilation into a slightly smaller
band converges in Schwartz topology and preserves the
two zero conditions.

If $\operatorname{supp}\chi\subset[a,b]$, the second
sum in (MA-GP.51) has at most $\lfloor\Lambda b\rfloor$
terms.  Cauchy--Schwarz and the change of variable $y=n/t$
give a finite bound by $\|u\|_2^2$ for
$\|M_{\sqrt\chi}T_\Lambda u\|_H^2$; the weight is bounded
on this compact interval.  Thus each output localization
is bounded on the auxiliary source.
If $u_j\to0$ in that source and $T_\Lambda u_j\to F$
in $H$, every such localization of $F$ is zero.
Hence $F=0$, proving closability.

For the Gram calculation, finite summation and the same
substitution yield
\[
 \begin{aligned}
 \langle \chi T_\Lambda u,T_\Lambda v\rangle_H
  &=4\sum_{m,n\ge1}\frac1n\int_0^\Lambda
       \chi(n/y)w_{\alpha,s}(\log(n/y))
                       u(y)\overline{v(my/n)}\,dy.
 \end{aligned}
 \tag{MA-GP.54}
\]
Changing the variable once more to the argument of $v$
and using the even source norm $2\int_0^\Lambda|u|^2$
gives (MA-GP.53).  The zero-mean projection appears because
the source is $\mathcal P_\Lambda^0$, rather than all even
band functions.  Both $m$ and $n$ are at most $\Lambda b$
in a nonzero term, proving uniform finiteness.
\end{proof}

Let $\overline T_\Lambda$ be this closed theta map and
$G_\Lambda=\overline T_\Lambda^{\,*}\overline T_\Lambda$
its positive self-adjoint Gram operator.  If
$\overline T_\Lambda=V G_\Lambda^{1/2}$ is its polar
decomposition, the actual global orthogonal PW projection is
\[
 Q_\Lambda
     =\operatorname*{s-lim}_{\varepsilon\downarrow0}
          V\,G_\Lambda(G_\Lambda+\varepsilon)^{-1}V^*.
 \tag{MA-GP.55}
\]
On the adjoint domain the bounded operator inside this limit
agrees with
$\overline T_\Lambda(G_\Lambda+\varepsilon)^{-1}
\overline T_\Lambda^{\,*}$.
Indeed polar decomposition gives this identity, with norm
at most one, and the spectral theorem gives the strong limit
$VV^*$ onto the range closure.  This is the domain-correct
version of a normal-equation inverse.  An unqualified formula
using $T(T^*T)^{-1}T^*$ would require additional hypotheses.
The localized forms (MA-GP.54) increase to the full target
norm on the core under a nonnegative compact exhaustion;
their uncut double expansion must not be interchanged with
the integral without a convergence proof.

\paragraph{MA-GP-006 (the frequency Gram is substantive, already at one place).}
The localized Gram kernel in (MA-GP.53) has genuine
off-diagonal rational-dilation terms, even on the two-moment
test domain.  Moreover, for every $\alpha>0$ and $s\ge0$,
$T_\Lambda$ is unbounded from the auxiliary ordinary
$L^2$ source in (MA-GP.50) into (MA-GP.52).
Neither assertion claims an unbounded inverse in a different
weighted local norm.

\begin{proof}
For the first assertion choose the output cutoff $\chi$
supported very close to $t_0=3/\Lambda$.
Take a positive input bump $u_+$ supported near $\Lambda/3$
and a positive bump $v_+$ near $2\Lambda/3$, aligned by
$v_+(2y)=u_+(y)$ on a smaller interval.
Correct the integral of $u_+$ by a bump near $\Lambda/6$,
and that of $v_+$ by a bump near $5\Lambda/6$.
Choose all four intervals sufficiently small and disjoint.
On the support of $\chi$ none of the correction intervals
is sampled by $n/t$ for an integer $n$; the only relevant
samples are near the multiples of $\Lambda/3$.
Thus the even zero-integral corrections $u,v$ are orthogonal
in the local source norm, while
\[
                    \langle\chi T_\Lambda u,T_\Lambda v\rangle_H>0.
 \tag{MA-GP.56}
\]
They vanish near the origin and hence belong to the test
domain.  This exhibits an actual off-diagonal Gram coupling.

For the second assertion choose an even nonnegative
$\psi\in C_c^\infty((-\Lambda,\Lambda))$, supported away
from zero, with nonzero integral.  Put $\Psi=\mathcal F_e\psi$.
Choose another fixed even annular test $\psi_0$ with
nonzero integral, and let
\[
 u_R(\xi)=\psi(\xi)\cos(2\pi R\xi)-c_R\psi_0(\xi),
 \qquad c_R=\frac{\int\psi(\xi)\cos(2\pi R\xi)\,d\xi}
                         {\int\psi_0(\xi)\,d\xi}.
 \tag{MA-GP.57}
\]
Then $u_R\in\mathscr D_\Lambda$, its source norm is
uniformly bounded, and $c_R=O_N(R^{-N})$ for every $N$.
Its Fourier transform is
\[
 f_R(t)=\tfrac12\bigl(\Psi(t-R)+\Psi(t+R)\bigr)
                                  -c_R(\mathcal F_e\psi_0)(t).
 \tag{MA-GP.58}
\]
Since $\Psi(0)>0$, choose $d>0$ with $\Psi(t)>0$
uniformly on $[-d,d]$.  For $t\in[R-d,R+d]$ and large
$R$, the $n=1$ summand of $\sum_{n\ge1}f_R(nt)$ is
bounded below by a fixed positive constant.  The remaining
summands are $O_N(R^{-N})$ after summing, by the Schwartz
decay at $nt-R\gg(n-1)R$ and at $nt+R$; the correction
has the same bound.  Consequently
\[
 \|T_\Lambda u_R\|_H^2
   \ge C\int_{R-d}^{R+d}w_{\alpha,s}(\log t)\,dt
   \ge C'R^{2\alpha}(1+\log^2R)^{s/2}.
 \tag{MA-GP.59}
\]
This diverges, proving the assertion with the specified
source norm and at the fixed frequency level $\Lambda$.
\end{proof}

These formulas isolate the actual finite-level normal equation.
The bounded half-line Gram inverse in MA-GP-004 and the
possibly unbounded auxiliary frequency Gram in MA-GP-005
are two different parametrizations of the same range, with
different source norms; they are not contradictory claims.
Neither may be silently replaced by the identity when
transporting a weighted orthogonal projection to an additive
Fourier cutoff.  The actual fixed-position condensation limit
is handled separately by MA-GP-001.  Any asymptotic local
trace comparison must retain the projection factors just
identified or prove that their contribution has the required
limit.

\paragraph{Grothendieck placement and source scope.}
The Gram inverse is an ordinary Hilbert-space normalization
of the actual range inclusion.  Its positivity comes from
the scalar norm weight $R$; it supplies no arithmetic
positive form on the full zero jets.  The finite-length
$\mathbb C[z]$-modules at those zeros retain the multiplicities
already identified in MA-WP, while a change of Hilbert metric
changes orthogonal representatives.  The actual trace-ideal
approximation in MA-GP-001 and the range-normalized projection
here fit the inherited cote~1 nuclear-trace comparison,
scans~079--080 and~111--112, through BR-TR-01/02 and
F-SM-11/12.  Those scan locators are not newly inspected here.
The finite-prime unitary factors in MA-LA-013 have their own
rank-one character interpretation when genuine Hecke coefficients
are supplied; the Gram calculation does not construct an
arithmetic correspondence, a new Artin motive, or an
$\ell$-adic weight statement.  The explicit Hardy range used
in (MA-GP.46)--(MA-GP.49) is exactly the theorem proved in
MA-GP-003, with the Burnol/Beurling primary-source scope stated
there; the normal-equation and frequency-kernel proofs above
are supplied in full.

\subsection{A sharp-band comparison and an actual regularized trace transport}

\paragraph{MA-GP-007 (the same-coordinate sharp-band candidate has a different range).}
The distinction from an additive Fourier cutoff persists even if its
frequency profile is a sharp characteristic function, so that it really
is an orthogonal projection.  This observation concerns a precisely
specified candidate identification; it is not attributed to a printed
operator equality in Masullo.

For the weight in (MA-GP.52) define the unitary radial map
\[
 (\mathcal U_w f)(t)
       =\sqrt{\frac{2t}{w_{\alpha,s}(\log t)}}\,f(t),\qquad t>0,
 \quad
 \mathcal U_w:L^2_{\mathrm{even}}(\mathbb R)\longrightarrow H.
 \tag{MA-GP.60}
\]
For any finite $L>0$ the same-coordinate additive candidate is
\[
 \Pi_L=\mathcal U_w\mathcal F_e
         M_{\mathbf1_{[-L,L]}}\mathcal F_e\mathcal U_w^{-1}.
 \tag{MA-GP.61}
\]
For every finite actual PW level $\Lambda>0$,
\[
               J_\Lambda\cap\operatorname{ran}\Pi_L=\{0\}.
 \tag{MA-GP.62}
\]
In particular $Q_\Lambda\ne\Pi_L$.  The same conclusion holds
after imposing the continuous zero-integral condition on the
auxiliary frequency band.  No choice of two finite positive
bandwidths changes this range obstruction.

\begin{proof}
The factor in (MA-GP.60) gives
$\|\mathcal U_w f\|_H^2=2\int_0^\infty|f(t)|^2dt$,
so it is unitary.  Every vector in $J_\Lambda$ is zero for
$0<t<1/\Lambda$, by (MA-GP.51) and closedness of this
support subspace in $H$.
If $F=\mathcal U_w f$ also belongs to $\operatorname{ran}\Pi_L$,
then $\widehat f$ is an $L^2$ function supported in $[-L,L]$.
It is in $L^1$ on that finite interval, and
\[
                 f(z)=\int_{-L}^{L}\widehat f(\xi)
                                         e^{2\pi iz\xi}\,d\xi
 \tag{MA-GP.63}
\]
is entire.  Since the radial factor is strictly positive for
$t>0$, $f$ vanishes almost everywhere on $(0,1/\Lambda)$.
Continuity and the identity theorem give $f=0$.
The actual range $J_\Lambda$ is nonzero, by a nonzero test
in (MA-GP.50) and the injectivity of the smooth theta map.
Thus the two projections cannot agree.  Passing to a smaller
moment-zero additive range preserves the zero intersection.
\end{proof}

The unweighted centering factor $\sqrt{2t}$ gives the same
support argument wherever both spaces are compared as actual
functions.  Formula (MA-GP.60) additionally puts the comparison
in one specified Hilbert space.  This argument rules out the
literal same-coordinate substitution (MA-GP.61); it neither
rules out a nonlocal transport nor contradicts an asymptotic
trace formula with its missing factors restored.

\paragraph{MA-GP-008 (an actual Gram-normalized local cyclic operator).}
Fix $\alpha>0$, $s\ge0$ and $\Lambda>0$, and retain the
auxiliary frequency space $\mathcal P_\Lambda^0$ and closed
theta map $\overline T=\overline T_\Lambda$ of MA-GP-005.
Write $G=\overline T^{\,*}\overline T$ and
$\overline T=VG^{1/2}$.
For $\chi,h\in C_c^\infty(\mathbb R_{>0})$ put
\[
 A=M_\chi C_h,\qquad
 (C_hF)(t)=\int_0^\infty h(r)F(t/r)\,\frac{dr}{r}.
 \tag{MA-GP.64}
\]
Thus $A\in\mathcal S_1(H)$, by the compact-kernel argument
of MA-LIFT-003.  There is a bounded trace-class operator
$K_{\chi,h}$ on $\mathcal P_\Lambda^0$ whose exact core form is
\[
       \langle A T_\Lambda u,T_\Lambda v\rangle_H
                         =\langle K_{\chi,h}u,v\rangle_{\mathcal P_\Lambda^0}.
 \tag{MA-GP.65}
\]
On positive band coordinates its explicit formula is
\[
 \begin{split}
 (K_{\chi,h}u)(y)&=P_\Lambda^0\biggl[
 2\sum_{m,n\ge1}\frac{\chi(m/y)w_{\alpha,s}(\log(m/y))}{m}\\
 &\hspace{18mm}\cdot\int_0^\infty h(r)r^{1/2}u(nry/m)\,\frac{dr}{r}
                                      \biggr].
 \end{split}
 \tag{MA-GP.66}
\]
Here $0<y<\Lambda$, and the output inside the projection is extended evenly.
If $\operatorname{supp}\chi\subset[a,b]$ and
$\operatorname{supp}h\subset[r_-,r_+]$ with $a,r_->0$,
only $m\le\Lambda b$ and $n\le\Lambda b/r_-$ contribute.

For every $\varepsilon>0$ define bounded operators
\[
 \begin{aligned}
 R_\varepsilon&=(G+\varepsilon)^{-1/2},&
 F_\varepsilon&=\overline T R_\varepsilon
                 =V\sqrt{G/(G+\varepsilon)},\\
 Q_{\Lambda,\varepsilon}&=F_\varepsilon F_\varepsilon^*,&
 L_{\Lambda,\varepsilon}(\chi,h)
       &=R_\varepsilon K_{\chi,h}R_\varepsilon
         =F_\varepsilon^*AF_\varepsilon.
 \end{aligned}
 \tag{MA-GP.67}
\]
Here $R_\varepsilon$ maps the entire auxiliary space into
$\operatorname{Dom}\overline T$, so the displayed maps
require no unverified adjoint domain.  They satisfy
\[
 \begin{gathered}
 \|F_\varepsilon\|\le1,\qquad
 \|L_{\Lambda,\varepsilon}(\chi,h)\|_1\le\|A\|_1,\\
 \operatorname{Tr}_H(AQ_{\Lambda,\varepsilon})
       =\operatorname{Tr}_{\mathcal P_\Lambda^0}
                           L_{\Lambda,\varepsilon}(\chi,h),\\
 \|AQ_{\Lambda,\varepsilon}-AQ_\Lambda\|_1
  +\|L_{\Lambda,\varepsilon}(\chi,h)-V^*AV\|_1
                                  \longrightarrow0.
 \end{gathered}
 \tag{MA-GP.68}
\]
In particular this is an actual local cyclic trace formula
for the specified global finite-level projection:
\[
 \operatorname{Tr}_H(M_\chi C_hQ_\Lambda)
    =\lim_{\varepsilon\downarrow0}
       \operatorname{Tr}_{\mathcal P_\Lambda^0}
       \bigl((G+\varepsilon)^{-1/2}
                K_{\chi,h}(G+\varepsilon)^{-1/2}\bigr).
 \tag{MA-GP.69}
\]

\begin{proof}
\emph{The explicit trace-class core form.}
In logarithmic coordinates the unweighted kernel of $A$ is
$\chi(e^x)h(e^{x-y})$, smooth with compact support in both
variables.  It is trace class.  Conjugating to the weighted
space introduces bounded multiplication factors on these
two compact coordinate projections; no differentiability of
$w_{\alpha,s}$ at zero is needed.  This proves the stated
trace-class property of $A$ directly.

For core vectors, insert (MA-GP.51) into (MA-GP.65).
The inner integral contributing to $AT_\Lambda u(t)$ is
\[
       2\chi(t)t^{-1/2}\sum_n
                   \int h(r)r^{1/2}u(nr/t)\,\frac{dr}{r}.
 \tag{MA-GP.70}
\]
Pair it with $2t^{-1/2}\sum_m v(m/t)$ in (MA-GP.52)
and put $y=m/t$.  The measure $t^{-2}dt$ becomes
$dy/m$, with the integration orientation reversed.
The factor four in the form becomes the factor two in
(MA-GP.66), because the source norm is
$2\int_0^\Lambda|u(y)|^2dy$.
The compact supports and the band bounds give precisely
the two finite summation bounds following (MA-GP.66).

For trace class, substitute $z=nry/m$ in that equation.
Before the bounded projection $P_\Lambda^0$, each term is
an integral operator with kernel, relative to $dz$,
\[
 \frac{2\,\chi(m/y)w_{\alpha,s}(\log(m/y))}
                        {\sqrt{mn}\sqrt{yz}}
                          h\!\left(\frac{mz}{ny}\right).
 \tag{MA-GP.71}
\]
On its support $y\ge1/b$ and $z\ge r_-/b$.
Insert smooth cutoffs equal to one on these portions of
$[0,\Lambda]$.  After removing the bounded factor
$\chi(m/y)w_{\alpha,s}(\log(m/y))$, the remaining kernel
extends smoothly to a neighborhood of the closed square
$[0,\Lambda]^2$.  Such a kernel is trace class: extend it
smoothly to a larger periodic square, expand in a double
Fourier series whose coefficients are absolutely summable
after repeated integration by parts, and restrict back to
the interval.  Bounded multiplication, the finite sum and
$P_\Lambda^0$ preserve this property.  This proves that
(MA-GP.66) defines a bounded $K_{\chi,h}\in\mathcal S_1$,
even though its weight factor can have a cusp.

\emph{Closure and domains.}
The original test domain is a graph core for $\overline T$
by its definition as the closure of $T_\Lambda$.
Approximating both vectors in graph norm therefore extends
(MA-GP.65) to all $u,v\in\operatorname{Dom}\overline T$:
the left side converges because $A$ is bounded and the right
side because $K_{\chi,h}$ is bounded.
In particular $A\overline T u\in\operatorname{Dom}\overline T^{\,*}$
for such $u$, and
$\overline T^{\,*}A\overline T u=K_{\chi,h}u$ there.
This adjoint-domain assertion is a consequence of the form
estimate, not an assumption in the initial calculation.

The spectral theorem gives
$\operatorname{Dom}\overline T=\operatorname{Dom}G^{1/2}$
and $\|G^{1/2}R_\varepsilon\|\le1$.
Thus $F_\varepsilon$ is a bounded contraction on its entire
source.  Substituting $R_\varepsilon u,R_\varepsilon v$
in the extended form proves the last identity of (MA-GP.67).
It proves it as an equality of bounded trace-class operators,
without writing an unrestricted product with $\overline T^{\,*}$.

In fact $\overline T$ is injective in this model.
The compact-output bounded maps of MA-GP-005 extend the
finite Poisson expression to every graph-domain vector.
If $\overline T u=0$, then
$S_u(y)=\sum_{m\ge1}u(my)=0$ almost everywhere on $y>0$.
For each $y>0$ the band support makes the sums finite, and
finite M\"obius inversion gives
\[
                     u(y)=\sum_{n\ge1}\mu(n)S_u(ny)=0
                                \quad\hbox{almost everywhere}.
 \tag{MA-GP.72}
\]
Dilations preserve null sets, so the almost-everywhere
statement survives the countably many possible indices.
Consequently $V$ is an isometry onto $J_\Lambda$.
The argument below would also work with its polar partial
isometry if a kernel were present.

\emph{Actual cyclic traces and their limit.}
Spectral calculus gives $F_\varepsilon\to V$ and
$F_\varepsilon^*\to V^*$ strongly, with uniform norm at
most one.  It also gives $Q_{\Lambda,\varepsilon}\to Q_\Lambda$
strongly; these are self-adjoint contractions.
For a fixed trace-class operator, uniformly bounded strong
convergence on the left, and adjoint strong convergence on
the right, imply trace-norm convergence.  This follows first
for rank-one operators and then by finite-rank approximation.
Apply it to $A$ to obtain both limits in (MA-GP.68) and
the uniform trace-norm bound there.
The cross-space cyclic identity for a trace-class factor
and a bounded factor gives
$\operatorname{Tr}(AF_\varepsilon F_\varepsilon^*)
 =\operatorname{Tr}(F_\varepsilon^*AF_\varepsilon)$.
Taking the proved trace-norm limits yields (MA-GP.69).
\end{proof}

\paragraph{The exact correction before taking any position or band limit.}
Let $E_\Lambda:\mathcal P_\Lambda^0\to L^2_{\mathrm{even}}(\mathbb R)$
be band inclusion and put
\[
 W_\Lambda=\mathcal U_w\mathcal F_eE_\Lambda,
 \qquad \Pi_\Lambda^0=W_\Lambda W_\Lambda^*.
 \tag{MA-GP.73}
\]
This is the sharp additive candidate with the same auxiliary
zero-integral source condition; $W_\Lambda$ is an isometry.
The actual finite comparison is the trace-class identity
\[
 \begin{split}
 \operatorname{Tr}_H\bigl(A(Q_{\Lambda,\varepsilon}-\Pi_\Lambda^0)\bigr)
   =\operatorname{Tr}_{\mathcal P_\Lambda^0}\bigl(
          R_\varepsilon K_{\chi,h}R_\varepsilon
                                  -W_\Lambda^*AW_\Lambda\bigr).
 \end{split}
 \tag{MA-GP.74}
\]
The first local term has the explicit rational-dilation kernel
(MA-GP.66), with its two normal-equation resolvents.
The second is the actual conjugated additive-band operator.
Thus this comparison retains both the theta summation and
the Gram normalization.  Neither can be removed by trace
cyclicity.  MA-GP-006 already proves that the theta map is
not an isometry for the auxiliary norm, and MA-GP-007
proves that the two finite ranges differ.
No nonzero value of the difference is asserted for every
particular test, and no asymptotic estimate for (MA-GP.74)
is inferred from the range distinction alone.

For completeness, at fixed $\chi,h,\Lambda,\varepsilon$
the ordinary Fredholm determinants satisfy
\[
 \det_{\mathcal P_\Lambda^0}(1-zL_{\Lambda,\varepsilon})
       =\det_H(1-zAQ_{\Lambda,\varepsilon}).
 \tag{MA-GP.75}
\]
This is Sylvester's identity, first for a finite-rank
approximation to $A$ and then by trace-norm continuity of
the Fredholm determinant.  The same reasoning and
(MA-GP.68) pass to the limit $\varepsilon\downarrow0$
at these fixed data, uniformly for $z$ in compact sets.
No determinant limit in an expanding position cutoff is
claimed.

The maps $F_\varepsilon$ provide a genuine Hilbert trace
transport.  They are not asserted to intertwine the source
and target representations: the weighted Gram need not
commute with the dilation action, and that action also
changes the fixed frequency level.  Likewise no continuity
of its resolvents on the polar-free LF source, or common
nuclear LF presentation, has been proved here.
The ordinary Hilbert trace-ideal and Fredholm-determinant
claims above have their own complete proofs.
The fixed-position condensation of MA-GP-001 can now be
expressed through (MA-GP.69); the remaining geometric
problem is an asymptotic estimate for the explicit
Gram-corrected cyclic operator, not the existence of such
an operator or an unverified use of $T^*$.

\section{The finite Paley--Wiener trace and the order of its limits}
\label{sec:global-pw-trace}

This section uses the actual range of MA-GP-003 and the specified
symmetric weight.  In particular it retains the bounded Gram change
of MA-GP-004.  It proves trace class of the actual convolution after
the orthogonal projection, before computing its trace.
All test functions in this section are additive tests
$g\in C_c^\infty((0,\infty))$, extended by zero to $\mathbb R$.
This support condition makes the half-line subspace and its
finite-level ideal invariant under convolution.  In multiplicative
coordinates the corresponding test is supported in $(1,\infty)$.

\paragraph{MA-GP-009 (the actual finite-level trace and two iterated limits).}
Let $0<\alpha<1/2$, $s=0$, and
$H=L^2(\mathbb R,e^{2\alpha|x|}\,dx)$.
For $Y>0$ set $c=-\log Y$ and
\[
 H_c=\{f\in H:\operatorname{supp}f\subseteq[c,\infty)\},
 \qquad \mathcal Q_c=H_c/J_Y,\qquad Z_c=H_c\ominus J_Y .
 \tag{MA-GP.80}
\]
Here $J_Y$ and $Q_Y^{\rm glob}=P_{J_Y}$ are the actual
Paley--Wiener closure and orthogonal projection of MA-GP-001--003.
Let $P_c^Z$ denote the orthogonal projection onto $Z_c$,
extended by zero on $H_c^\perp$.
Then the actual operators satisfy
\[
 \begin{gathered}
 C_gP_c^Z\in\mathcal S_1(H),\qquad
 C_g^{\mathcal Q_c}\in\mathcal S_1(\mathcal Q_c),\\
 \operatorname{Tr}_H(C_gP_c^Z)
 =\operatorname{Tr}_{\mathcal Q_c}(C_g^{\mathcal Q_c})
 =\sum_{\operatorname{Re}\rho<1/2+\alpha}
           m_\rho\,\mathcal Mg(\rho-1/2).
 \end{gathered}
 \tag{MA-GP.81}
\]
The sum is over nontrivial zeta zeros, with multiplicity, and
converges absolutely.  It is independent of $Y$.

Choose the uniformly bounded smooth compact exhaustion
$\chi_R(x)=\chi_0(x/R)$ with $\chi_0=1$ near zero.
Both iterated limits of the actual localized operators exist:
\[
 \begin{aligned}
 \lim_{Y\to\infty}\lim_{R\to\infty}
   \operatorname{Tr}(M_{\chi_R}C_gQ_Y^{\rm glob})
   &=-\sum_{\operatorname{Re}\rho<1/2+\alpha}
                  m_\rho\,\mathcal Mg(\rho-1/2),\\
 \lim_{R\to\infty}\lim_{Y\to\infty}
   \operatorname{Tr}(M_{\chi_R}C_gQ_Y^{\rm glob})
   &=-\sum_{1/2-\alpha<\operatorname{Re}\rho<1/2+\alpha}
                  m_\rho\,\mathcal Mg(\rho-1/2).
 \end{aligned}
 \tag{MA-GP.82}
\]
Their difference, in the displayed order, is
\[
          -\sum_{\operatorname{Re}\rho\le1/2-\alpha}
                         m_\rho\,\mathcal Mg(\rho-1/2).
 \tag{MA-GP.83}
\]
The formula includes the lower-edge zeros.  It does not assert
that such exterior or boundary zeros occur, or that the
difference is nonzero for any specified test.

\begin{proof}
\emph{The actual quotient action and its Hardy representative.}
Put
\[
 (U_cf)(y)=e^{\alpha(c+y)}f(c+y),\qquad
 R_c(y)=e^{\,2\alpha|c+y|-2\alpha(c+y)},\qquad y\ge0.
 \tag{MA-GP.84}
\]
The map $U_c$ is unitary from the actual $H_c$ onto
$H_{R_c}=L^2((0,\infty),R_c(y)\,dy)$.
For this fixed $c$ the norm is equivalent to ordinary $L^2$,
since $1\le R_c\le e^{4\alpha\max(-c,0)}$.
If $c\ge0$ it is identically one; if $c<0$ and $L=-c$,
\[
 R_c(y)=
 \begin{cases}
       e^{4\alpha(L-y)},&0\le y\le L,\\
       1,&y\ge L.
 \end{cases}
 \tag{MA-GP.85}
\]
MA-GP-003 says that, as an actual closed set,
the image of $J_Y$ under the Laplace map $\mathcal L U_c$
is $B_\alpha H^2(\operatorname{Re}z>0)$.
Here $B_\alpha$ is the pure Blaschke product with zeros
\[
                 z_j=\alpha+\tfrac12-\rho_j=a_j+ib_j,
          \qquad a_j>0,
 \tag{MA-GP.86}
\]
repeated according to multiplicity.
The norms induce a bounded invertible map between their
quotient Hilbert spaces.  Thus the actual quotient action is
boundedly similar to the standard model-space action on
$K_{B_\alpha}=H^2\ominus B_\alpha H^2$.
This statement concerns the action on the quotient; it does
not replace the actual weighted orthogonal projection by the
ordinary Hardy projection.

The unilateral shifts $S_t v(y)=\mathbf1_{y\ge t}v(y-t)$,
$t\ge0$, are contractions on $H_{R_c}$ because $R_c$ is
nonincreasing.  They preserve the image of $J_Y$.
For
\[
      k(t)=e^{\alpha t}g(t),\qquad
      G(z)=\int_0^\infty k(t)e^{-zt}\,dt
                           =\mathcal Mg(\alpha-z),
 \tag{MA-GP.87}
\]
the map $U_c C_g U_c^{-1}$ on $H_{R_c}$ is the actual
causal convolution $C_k^+=\int_0^\infty k(t)S_t\,dt$.
The vectorwise integrals define bounded operators.
The corresponding Hardy quotient action is
\[
                 T_G=P_{K_{B_\alpha}}M_G|_{K_{B_\alpha}}.
 \tag{MA-GP.88}
\]
Since $k$ is smooth and compactly supported inside $(0,\infty)$,
integration by parts gives, for every $N$,
$|G(z)|\le C_N(1+|z|)^{-N}$ on $\operatorname{Re}z\ge0$.
In particular all multipliers used below are bounded analytic
multipliers of $H^2$.

\emph{Uniform half-plane estimates at arbitrarily close edge zeros.}
Use the boundary norm $(2\pi)^{-1}\int|F(it)|^2dt$.
The normalized Malmquist basis of $K_{B_\alpha}$ is
\[
 e_j(z)=\frac{\sqrt{2a_j}}{z+\overline{z_j}}
                         \prod_{\ell<j} b_\ell(z),
 \qquad |b_\ell(it)|=1,\quad b_\ell(z_\ell)=0.
 \tag{MA-GP.89}
\]
Unimodular normalization constants in $b_\ell$ do not affect
the argument.  The elementary orthogonal decomposition
$K_{UV}=K_U\oplus U K_V$ for inner $U,V$ proves orthonormality,
including repeated zeros.  It also proves completeness.
Indeed the finite-product subspaces increase, and their
orthogonal complements decrease to $B_\alpha H^2$:
if $F$ is divisible in $H^2$ by every finite product $B_N$,
then $F/B_N$ has bounded $H^2$ norm.  A weakly convergent
subsequence, evaluated at each point in the half-plane,
has limit $F/B_\alpha$ after the standard normalizations
of $B_N$.  Hence $F/B_\alpha\in H^2$.
This proves the claimed intersection and the basis assertion.
The empty or finite product cases have the same interpretation.

For a bounded analytic multiplier $F$ its basis-vector norm is
\[
 \|M_F e_j\|_{H^2}^2
   =\frac1\pi\int_{\mathbb R}
         |F(it)|^2\frac{a_j}{(t-b_j)^2+a_j^2}\,dt .
 \tag{MA-GP.90}
\]
Suppose $|F(it)|\le C/(1+|t|)$.
There is a constant $C_F'$ independent of $j$ such that
\[
                   \|M_F e_j\|_{H^2}^2
                              \le\frac{C_F'}{1+b_j^2}.
 \tag{MA-GP.91}
\]
To prove it, the containing strip for zeta gives
$0<a_j\le\alpha+1/2<1$.
For $|b_j|\ge2$, split the integral into
$|t-b_j|\le|b_j|/2$ and its complement.
On the first part $|t|\ge|b_j|/2$, and the Poisson
kernel has total mass at most one.
On the second part its density is at most
$4(\alpha+1/2)/(\pi b_j^2)$; use
$\int|F(it)|^2dt<\infty$.
These give (MA-GP.91); bounded $|b_j|$ follows directly
from the total Poisson mass and boundedness of $F$.
The unconditional local zero count NT-Z10 yields
\[
                      \sum_j(1+b_j^2)^{-1}<\infty.
 \tag{MA-GP.92}
\]
No positive lower bound on $a_j$ has been used.
It follows from (MA-GP.89)--(MA-GP.92) that
$M_F|_{K_{B_\alpha}}$ and its compression $T_F$ are
Hilbert--Schmidt.

Fix $\kappa>0$.  Apply this result to
\[
                    r(z)=(\kappa+z)^{-1},
               \qquad F_2(z)=(\kappa+z)G(z).
 \tag{MA-GP.93}
\]
Both meet the bounded-analytic and boundary-decay requirements.
The quotient compressions form a homomorphism:
$T_FT_H=T_{FH}$, because $B_\alpha H^2$ is invariant under
every bounded analytic multiplier.
Consequently
\[
                          T_G=T_rT_{F_2}\in\mathcal S_1,
             \qquad T_r^2\in\mathcal S_1.
 \tag{MA-GP.94}
\]
This proves quotient trace class before a trace is assigned.
Trace in the actual orthonormal basis (MA-GP.89) and the
analytic Poisson formula now give
\[
 \begin{split}
 \operatorname{Tr}T_G
   &=\sum_j\frac1\pi\int_{\mathbb R}
                   G(it)\frac{a_j}{(t-b_j)^2+a_j^2}\,dt\\
   &=\sum_jG(z_j)
    =\sum_{\operatorname{Re}\rho<1/2+\alpha}
                      m_\rho\,\mathcal Mg(\rho-1/2).
 \end{split}
 \tag{MA-GP.95}
\]
The Poisson identity follows also by the Cauchy integral
formula on the right half-plane, using the rapid decay of
$G$.  Absolute convergence follows both from trace class
for the diagonal sum and directly from rapid vertical
decay of $\mathcal Mg$ on the fixed containing strip and
NT-Z10.  Bounded similarity proves the quotient trace
assertion of (MA-GP.81).

\emph{A trace-class resolvent square on the actual adjoint subspace.}
The generator $A_c$ of the quotient semigroup induced by $S_t$
has resolvent represented by $T_r$ in the Hardy quotient.
Indeed the vectorwise strong Laplace integral
$\int_0^\infty e^{-\kappa t}S_t\,dt$ is bounded, and
its multiplier is $(\kappa+z)^{-1}$.
The usual resolvent identity follows directly by changing
variables in its shift difference quotient:
$A_cR_\kappa=\kappa R_\kappa-1$; integration by parts
on the generator domain gives the inverse identity there.
Its tail has operator norm at most $e^{-\kappa T}/\kappa$,
although no operator-norm continuity of $t\mapsto S_t$
is required.  By (MA-GP.94) and bounded similarity,
\[
                        (\kappa-A_c)^{-2}\in\mathcal S_1.
 \tag{MA-GP.96}
\]

Write $R=R_c$ in the remainder of this calculation, and
let $Z_R=U_cZ_c$, with its actual $H_R$ orthogonal projection
$P_{Z_R}$.  The adjoint unilateral semigroup is
\[
  (S_t^{*,R}u)(y)=\frac{R(y+t)}{R(y)}u(y+t),
       \qquad \mathcal L_Ru=R^{-1}\partial_y(Ru).
 \tag{MA-GP.97}
\]
It preserves $Z_R$ because the ideal is invariant under $S_t$.
Multiplication by $R$ is unitary from $H_R$ to
$L^2((0,\infty),R^{-1}dy)$ and converts this adjoint
semigroup to left translation $v(y)\mapsto v(y+t)$.
Its generator has no boundary condition at zero.
Precisely, for $j\le2$ the graph domain consists of
$v=Ru$ with weak derivatives through order $j$ in this
weighted $L^2$, and $\mathcal L_R^j u=R^{-1}v^{(j)}$.
For the first derivative the direct assertion follows from
the fundamental theorem of calculus and strong translation
continuity; the converse follows by testing the generator
difference quotient against interior smooth tests.
Iteration proves the second-derivative statement.
The weight is bounded above and below, so $v\in H^2(0,\infty)$
when $u\in\mathcal D(\mathcal L_R^2)$.

The restriction of (MA-GP.97) to $Z_R$ is the Hilbert
adjoint of the actual quotient semigroup.  Equation
(MA-GP.96), adjunction, and the restricted resolvent integral
therefore give
\[
 \mathcal R_2=(\kappa-\mathcal L_R)^{-2}P_{Z_R}
                       \in\mathcal S_1(H_R).
 \tag{MA-GP.98}
\]
It maps boundedly into the second graph domain, since
$\mathcal L_R(\kappa-\mathcal L_R)^{-1}
=\kappa(\kappa-\mathcal L_R)^{-1}-1$.

\emph{The actual half-line lift, including its endpoint and cusp.}
Set $a(y)=R(y)^{-1}$.  For
$u\in\mathcal D(\mathcal L_R^2)$ and $v=Ru$,
\[
 C_k^+(\kappa-\mathcal L_R)^2u(x)
       =\int_0^\infty k(x-y)a(y)
                             (\kappa-\partial_y)^2v(y)\,dy .
 \tag{MA-GP.99}
\]
The test $k$ is extended by zero on the negative half-line.
Integrate twice by parts on $(0,L)$ and $(L,\infty)$
when $c<0$, or on $(0,\infty)$ when $c\ge0$.
All terms at infinity vanish because $k(x-y)$ has compact
support in $y$.  The bulk operator is
\[
 \mathcal B_k u(x)=
   \int_{y>0,\ y\ne L} R(y)
       (\kappa+\partial_y)^2
                 \bigl(k(x-y)a(y)\bigr)u(y)\,dy ,
 \tag{MA-GP.100}
\]
where the exclusion of $L$ is omitted when $c\ge0$.
After expansion its kernel is the sum of
\[
  [\kappa^2k-2\kappa k'+k''](x-y)
 +[2\kappa k-2k'](x-y)R(y)a'(y)
 +k(x-y)R(y)a''(y).
 \tag{MA-GP.101}
\]
The piecewise functions $Ra'$ and $Ra''$ are bounded;
they equal $4\alpha,16\alpha^2$ on $(0,L)$ and zero
outside when $c<0$.  Hence $\mathcal B_k$ is bounded
on $H_R$, as a finite sum of bounded causal convolutions
and bounded multiplications.

The exact endpoint contribution at zero is
\[
 \bigl[a(0)(2\kappa k-k')+a'(0+)k\bigr]v(0)
                                  -a(0)k\,v'(0).
 \tag{MA-GP.102}
\]
When $c<0$ the only interior contribution is
\[
          [a']_{L-}^{L+}\,k(\,\cdot-L)v(L)
                      =-4\alpha\,k(\,\cdot-L)v(L).
 \tag{MA-GP.103}
\]
To verify the signs, the boundary expression before summing
the intervals is
$qv'-(2\kappa q+q')v$, where $q(y)=k(x-y)a(y)$.
Subtracting its value at zero gives (MA-GP.102).
At $L$, the continuity of $a,v,v'$ cancels all terms
except the jump of $a'$ times $k(x-L)v(L)$.
For $R=1$ these formulas reduce to the explicit identity
\[
 C_k^+(\kappa-\partial)^2u
  =C_{\kappa^2k-2\kappa k'+k''}^+u
                 +(2\kappa k-k')u(0)-k\,u'(0).
 \tag{MA-GP.104}
\]
Thus the half-line endpoint has not been dropped.

The functions multiplying $v(0),v'(0),v(L)$ in
(MA-GP.102)--(MA-GP.103) belong to $H_R$.
The corresponding evaluation functionals are bounded
on the graph domain of $\mathcal L_R^2$.
This follows from the one-dimensional Sobolev trace
inequality on a fixed interval containing the relevant
endpoint, and norm equivalence there.  One can prove
the inequality by averaging the fundamental theorem
of calculus and applying Cauchy--Schwarz to $v$ and
$v'$, then to $v'$ and $v''$.
Consequently their compositions with $\mathcal R_2$
are bounded functionals on $H_R$.

Apply (MA-GP.99)--(MA-GP.103) to
$u=\mathcal R_2f$.  Its left side is $C_k^+P_{Z_R}f$.
The right side is the product of the bounded
$\mathcal B_k$ with the trace-class $\mathcal R_2$,
plus the displayed finite-rank boundary terms.
This proves
\[
                         C_k^+P_{Z_R}\in\mathcal S_1(H_R).
 \tag{MA-GP.105}
\]
Conjugating by the actual unitary $U_c$, and extending by
zero on $H_c^\perp$, proves $C_gP_c^Z\in\mathcal S_1(H)$.
The support of $g$ ensures that $C_g$ maps $H_c$ into itself.
Trace cyclicity with the actual orthogonal projection gives
\[
 \operatorname{Tr}(C_gP_c^Z)
       =\operatorname{Tr}(P_c^ZC_gP_c^Z)
       =\operatorname{Tr}_{\mathcal Q_c}(C_g^{\mathcal Q_c}).
 \tag{MA-GP.106}
\]
The middle compression is the orthogonal representative
of quotient convolution.  Together with (MA-GP.95),
this finishes (MA-GP.81) by an actual lift proof.

\emph{Position cutoff first and the comparison of orders.}
Let $P_c$ be multiplication by $\mathbf1_{[c,\infty)}$.
As actual orthogonal projections on $H$,
\[
                        Q_Y^{\rm glob}=P_c-P_c^Z.
 \tag{MA-GP.107}
\]
For compactly supported smooth $\chi$, the operator
$M_\chi C_gP_c$ is trace class and has trace zero.
Indeed $M_\chi C_g$ is trace class by the actual compact
kernel argument of MA-LIFT-003, and right multiplication
by $P_c$ preserves trace class.
Approximate $P_c$ strongly by uniformly bounded smooth
multiplication functions.  The corresponding products
converge in trace norm, by the finite-rank approximation
argument and adjoint strong convergence.
Each approximating trace is the ordinary diagonal integral
$g(0)\int\chi(x)\eta(x)\,dx=0$.
The weight is handled by the same two bounded compact
coordinate multiplication factors as in MA-LIFT-003;
no smoothness of the weight at zero is required.
Thus
\[
 \operatorname{Tr}(M_{\chi_R}C_gQ_Y^{\rm glob})
                   =-\operatorname{Tr}(M_{\chi_R}C_gP_c^Z).
 \tag{MA-GP.108}
\]
For fixed $Y$, the uniformly bounded multipliers
$M_{\chi_R}$ converge strongly to one.
Since $C_gP_c^Z$ is trace class, this proves
\[
          \|(M_{\chi_R}-1)C_gP_c^Z\|_1\longrightarrow0.
 \tag{MA-GP.109}
\]
Equations (MA-GP.108)--(MA-GP.109) and (MA-GP.81)
give the first limit in (MA-GP.82); its inner limit
is already independent of $Y$.

For the opposite order, MA-GP-001 gives the actual
fixed-cutoff trace-norm convergence as $Y\to\infty$,
and then the all-weight ambient theorem MA-WP-003
gives the full-strip quotient trace.
At $s=0$ no boundary jet is continuous, since the
edge condition is $2j<s-1=-1$.
Thus precisely the strict strip
$1/2-\alpha<\operatorname{Re}\rho<1/2+\alpha$
is retained.  This is the second limit in (MA-GP.82).
Subtracting the two absolutely convergent sums proves
(MA-GP.83), including the lower-edge convention.
\end{proof}

\paragraph{Scope, source placement, and the remaining transport problem.}
All estimates above hold for a fixed finite $Y$.
The bound relating the actual symmetric weight to the half-line
exponential weight grows as $Y^{4\alpha}$ for $Y>1$.
No uniform trace-norm bound for the lifted operators as
$Y\to\infty$, no joint-limit theorem, and no prescribed
diagonal rate are inferred.
The defect (MA-GP.83) is an exact statement about two orders
of limits of the actual global Paley--Wiener projections.
It is separate from the local additive Fourier-cutoff
calculation in MA-LA-014.

The proof relies on the actual range theorem MA-GP-003,
with its primary Beurling/Burnol locators, and on NT-Z10's
unconditional zero count.  The half-plane Malmquist calculation
is supplied explicitly; it is the half-plane version of the
compressed-multiplier argument in MA-HY-002.
The actual lift is a half-line version of MA-LIFT-002 with
the endpoint at zero and the weight cusp both retained.
These are proofs for the specified Hilbert model, not an
appeal to the source's uncorrected ambient nuclearity claim.
Neither $\mathsf A$ nor $M$ is used.

In the inherited Grothendieck trace framework, the finite
zero jets give the finite-length spectral blocks and the
bounded-to-trace-class factorization makes the ambient
trace legitimate before it is evaluated.  The applicable
cote~1 locators are the already recorded scans~079--080
and~111--112, ALPOGE\_BRIDGES BR-TR-01/02.
No additional archival reading or new arithmetic motive
realization is asserted by this operator calculation.

\section{Causal spectral synthesis for the actual weighted Paley--Wiener levels}
\label{sec:global-weighted-pw}

The polynomial weight cannot be removed by the equivalent-norm
argument of MA-GP-003.  We instead identify the actual theta range
as the closure of the range of a bounded causal convolution.
An explicit divided-difference identity gives the quotient
resolvent.  Its principal parts, polynomial bounds, and spectral
completeness are proved below.  This supplies the half-line
Hardy--Sobolev closure without invoking a weighted analogue
of Beurling's theorem.

Fix $0<\alpha<1/2$, $s\ge0$, $Y>0$, and set
\[
 c=-\log Y,\qquad d=\tfrac12-\alpha,\qquad
 u=\tfrac12+\alpha=1-d,\qquad
 \mathcal H_s=L^2((0,\infty),(1+y^2)^{s/2}\,dy).
 \tag{MA-GW.1}
\]
The actual global space is
$H=L^2(\mathbb R,e^{2\alpha|x|}(1+x^2)^{s/2}dx)$,
and $H_c$ is its supported half-line subspace.
The map $U_cf(y)=e^{\alpha(c+y)}f(c+y)$ carries its actual
norm to
\[
 \|U_cf\|_{W_c}^2,\qquad
 W_c(y)=e^{2\alpha|c+y|-2\alpha(c+y)}
                              (1+(c+y)^2)^{s/2}.
 \tag{MA-GW.2}
\]
For this fixed $c$ the $W_c$ norm and the $\mathcal H_s$
norm are equivalent.  The actual closed subspace $J_Y$
is always the theta-image closure from MA-GP-001--002;
it is not replaced by an unproved weighted inner-function
range.

\paragraph{MA-GW-001 (an actual bounded causal range).}
For a sufficiently large integer $N$ define
\[
               \Phi(z)=\frac{(z-u)\zeta(z+d)}{(z+2)^N}.
 \tag{MA-GW.3}
\]
The singularity at $z=u$ is removable and its value is nonzero.
There is a continuous causal kernel $\phi$, supported in
$[0,\infty)$, with $|\phi(t)|\le C e^{-\eta t}$ for some
$\eta>0$, whose Laplace transform is $\Phi$.
Consequently
\[
 T_\phi f=\int_0^\infty\phi(t)S_tf\,dt,\qquad
 (S_tf)(y)=\mathbf1_{y\ge t}f(y-t),
 \tag{MA-GW.4}
\]
is bounded on $\mathcal H_s$, and the actual range satisfies
\[
                  U_cJ_Y=\overline{\operatorname{ran}T_\phi}
                                           ^{\,\mathcal H_s}.
 \tag{MA-GW.5}
\]
The same equality holds as a set in the equivalent $W_c$ norm.

\begin{proof}
\emph{The kernel and the weighted semigroup.}
Take $0<\eta<d$, for example $\eta=d/2$.
The numerator of (MA-GW.3) is holomorphic on the whole plane,
including $z=u$, because the residue of zeta at $1$ is one.
Its denominator has no zero on $\operatorname{Re}z\ge-\eta$.
The unconditional fixed-strip polynomial bound for zeta
from NT-Z10, and the Euler bound on a right half-plane,
allow $N$ to be chosen so that
\[
 |\Phi(\sigma+it)|\le C(1+|t|)^{-2}
                       \quad(-\eta\le\sigma\le2),
 \tag{MA-GW.6}
\]
with the removable value included on bounded sets.
For $\sigma\ge2$, its Euler product also gives the uniform
bound $|\Phi(\sigma+it)|\le C|\sigma+it|^{1-N}$.
Increase $N$ if necessary so $N>3$.

Define the inverse Laplace integral on a vertical line with
positive real part.  It is absolutely convergent by
(MA-GW.6).  Cauchy's theorem, first on finite rectangles
and then by the uniform vertical decay, allows its line
to move to $\operatorname{Re}z=-\eta$.
This proves the continuous-kernel bound
$|\phi(t)|\le C e^{-\eta t}$ for $t\ge0$.
For $t<0$, move the line to $\operatorname{Re}z=M$ and
then let $M\to\infty$.
The absolute integral there is at most
$C e^{Mt}M^{2-N}$, so it tends to zero.
Thus the kernel is causal.  Fourier inversion on a
positive vertical line proves that its Laplace transform
is exactly $\Phi$; the identity then holds throughout
$\operatorname{Re}z>-\eta$ by analytic continuation.

The shifts are strongly continuous on $\mathcal H_s$ and obey
\[
 \|S_t\|\le2^{s/4}(1+t^2)^{s/4}\qquad(t\ge0).
 \tag{MA-GW.7}
\]
The inequality follows from
$1+(y+t)^2\le2(1+y^2)(1+t^2)$; continuity follows first
on compact smooth functions and then by density.
The exponential kernel bound therefore makes (MA-GW.4)
a bounded vectorwise integral.  Its Laplace transform
is multiplication by $\Phi$, by Fubini on a right
half-plane and then uniqueness.

\emph{The exact source primitive.}
Let $v$ be an even smooth frequency test supported in
$[-Y,Y]$ with $v(0)=0$ and $\int v=0$.
Write $f=\mathcal F_e v$ for the even Fourier inverse.
The actual Poisson theta formula gives, with
$z$ initially far enough to the right,
\[
 (\mathcal L U_c\Theta f)(z)
       =2e^{zc}\zeta(z+d)
                        \int_0^Yv(t)t^{z+d-1}\,dt .
 \tag{MA-GW.8}
\]
Both sides continue to the open right half-plane.
Make the change $t=Y e^{-q}$ and put
\[
 a(q)=e^{-dq}v(Y e^{-q}),\qquad
 b(q)=e^{uq}\int_0^q e^{-ur}a(r)\,dr.
 \tag{MA-GW.9}
\]
The moment condition gives
\[
       \int_0^\infty e^{-uq}a(q)\,dq
                         =Y^{-1}\int_0^Yv(t)\,dt=0.
 \tag{MA-GW.10}
\]
Thus $a=b'-ub$ and also
$b(q)=-e^{uq}\int_q^\infty e^{-ur}a(r)\,dr$.
At zero, $a$ and $b$ are flat because $v$ is smooth and
flat at the band endpoint $Y$.
At infinity, evenness and $v(0)=0$ give
$v(t)=O(t^2)$, including the corresponding logarithmic
derivative bounds.  Hence $a,b$ and every derivative
are $O(e^{-(d+2)q})$.
In particular $b$ and all its derivatives belong to
$\mathcal H_s$, with all zero boundary traces at zero.

If $B(z)=\mathcal Lb(z)$, (MA-GW.8)--(MA-GW.10) become
the exact identity
\[
             (\mathcal L U_c\Theta f)(z)
                 =2Y^d(z-u)\zeta(z+d)B(z)
                 =2Y^d\Phi(z)(z+2)^N B(z).
 \tag{MA-GW.11}
\]
The factor $2Y^d$ is fixed and nonzero.
Since $(\partial+2)^N b\in\mathcal H_s$ and all its
integration-by-parts boundary terms vanish, (MA-GW.11)
shows that every such actual theta core vector belongs
to $\operatorname{ran}T_\phi$.
The band-supported Schwartz tests used to define the
source range have exactly this smoothness and flatness;
equivalently one may use the inward-dilation approximation
already justified in MA-GP-005.

Conversely, every $b\in C_c^\infty((0,\infty))$ gives a
source test by setting $a=b'-ub$ and
$v(Y e^{-q})=e^{dq}a(q)$, then extending evenly.
The result is smooth, supported away from both zero and
the band endpoints, and has zero integral by (MA-GW.10).
Thus the vector with Laplace transform
$(z-u)\zeta(z+d)B(z)$ is in the actual theta core.

It remains to approximate all of $\operatorname{ran}T_\phi$.
Take $h\in C_c^\infty((0,\infty))$ and solve
$(\partial+2)^N b=h$ with zero initial traces by the
causal kernel $e^{-2q}q^{N-1}/(N-1)!$.
The resulting $b$ is flat at zero and has exponentially
decreasing derivatives at infinity.
Smooth compact truncations $b_M$ therefore satisfy
$(\partial+2)^N b_M\to h$ in $\mathcal H_s$.
Boundedness of $T_\phi$ and the preceding source construction
show that $T_\phi h$ is in $U_cJ_Y$.
Such $h$ are dense in $\mathcal H_s$.
This proves both inclusions in (MA-GW.5).
Equivalence of the two fixed-$c$ norms proves the final assertion.
\end{proof}

\paragraph{MA-GW-002 (weighted half-line spectral synthesis and trace).}
Let
\[
 b_s=\#\{j\in\mathbb Z_{\ge0}:2j<s-1\},\qquad
 r^+_\rho=
 \begin{cases}
 m_\rho,&\operatorname{Re}\rho<u,\\
 \min(m_\rho,b_s),&\operatorname{Re}\rho=u,\\
 0,&\operatorname{Re}\rho>u.
 \end{cases}
 \tag{MA-GW.12}
\]
The exact actual closure in $H_c$ is the common kernel
of the continuous Mellin jets of orders
$0\le j<r^+_\rho$ at $\rho-1/2$.
The quotient $\mathcal Q_c=H_c/J_Y$ has dense span of
its finite spectral blocks, with precisely these lengths.
For every additive $g\in C_c^\infty((0,\infty))$ its
actual quotient convolution is trace class and
\[
       \operatorname{Tr}_{\mathcal Q_c}C_g^{\mathcal Q_c}
           =\sum_\rho r^+_\rho\,\mathcal Mg(\rho-1/2).
 \tag{MA-GW.13}
\]
The sum converges absolutely and is independent of $Y$.
A sufficiently high power of the actual quotient
semigroup resolvent is trace class.
No zero separation or positive distance from the upper
edge is assumed.

\begin{proof}
\emph{The quotient and its meromorphic resolvent are actual operators.}
Work first on $\mathcal H_s$ and let
$\mathcal J=\overline{\operatorname{ran}T_\phi}$.
The shifts commute with $T_\phi$, so they preserve
$\mathcal J$ and induce a strongly continuous semigroup
on $\mathcal Q=\mathcal H_s/\mathcal J$ with the same
polynomial bound (MA-GW.7).
Let $A$ be its generator.
For $\operatorname{Re}\lambda>0$ the vectorwise
Laplace integral gives its bounded resolvent
\[
          R(\lambda)=(\lambda-A)^{-1}
                     =\int_0^\infty e^{-\lambda t}S_t^{\mathcal Q}\,dt .
 \tag{MA-GW.14}
\]
The direct difference-quotient and generator-domain
integration-by-parts verification is the one given in
MA-GP-009; (MA-GW.7) supplies every integrable majorant.

For $\operatorname{Re}\lambda<\eta$ define the scalar kernel
and actual bounded operator
\[
 \psi_\lambda(r)=\int_r^\infty\phi(t)e^{\lambda(t-r)}\,dt,
 \qquad
 T_\lambda=\int_0^\infty\psi_\lambda(r)S_r^{\mathcal Q}\,dr .
 \tag{MA-GW.15}
\]
The map $\lambda\mapsto\psi_\lambda$ is holomorphic in
the weighted $L^1$ kernel norm used with (MA-GW.7).
Differentiating its scalar exponential and using
$|\phi(t)|\le Ce^{-\eta t}$ proves this on each smaller
half-plane.  Hence $T_\lambda$ is operator-norm
holomorphic there.  Moreover
\[
 \sup_{\operatorname{Re}\lambda\le\epsilon}\|T_\lambda\|<\infty
                      \qquad(\epsilon<\eta).
 \tag{MA-GW.16}
\]
For example, bound its double integral by
$C_s\int_0^\infty|\phi(t)|e^{\max(\epsilon,0)t}
\int_0^t(1+r)^{s/2}\,dr\,dt$.
No bound here depends on the imaginary part or on how
far $\operatorname{Re}\lambda$ is negative.

The derivative $\psi_\lambda'=-\phi-\lambda\psi_\lambda$
and the value $\psi_\lambda(0)=\Phi(-\lambda)$ give
the bounded divided-difference identity
\[
 (\lambda-A)T_\lambda=\Phi(-\lambda)\,1,\qquad
 T_\lambda(\lambda-A)=\Phi(-\lambda)\,1
                                      \quad\hbox{on }\mathcal D(A).
 \tag{MA-GW.17}
\]
For clarity, before taking the quotient the right side is
$\Phi(-\lambda)1-T_\phi$.
Integrating the semigroup derivative against
$\psi_\lambda$ gives this formula, with boundary value
$\psi_\lambda(0)$; a shift difference quotient proves
that the resulting operator maps into the generator
domain.  On the quotient $T_\phi$ induces the zero
operator by (MA-GW.5), proving (MA-GW.17).
All these convolution operators preserve the range
closure because they commute with $T_\phi$.

Whenever $\Phi(-\lambda)\ne0$, the two identities show that
\[
                         R(\lambda)=T_\lambda/\Phi(-\lambda).
 \tag{MA-GW.18}
\]
This is a two-sided bounded inverse of the actual closed
generator, not just a formal meromorphic function.
It agrees with (MA-GW.14) on their common domain.
Together they give a meromorphic continuation of the
resolvent to the plane.  In the right half-plane any
apparent singularity of (MA-GW.18) is removable by
(MA-GW.14).  The remaining possible poles are
\[
                \lambda_\rho=\rho-u,\qquad
                \operatorname{Re}\rho\le u,
 \tag{MA-GW.19}
\]
of order at most $m_\rho$.
Indeed the zeros of $\Phi(-\lambda)$ there correspond
to $\zeta(d-\lambda)=0$, and the functional equation
pairs this multiset with (MA-GW.19).
The factor at $d-\lambda=1$ is removable and nonzero.
Trivial zeta zeros would give $\operatorname{Re}\lambda>0$
and contribute no singularity to the actual resolvent.
Thus all possible poles have real parts in $[-u,0]$.

\emph{Every principal part has finite rank, with the exact edge truncation.}
Let $\lambda_0$ be one possible pole, with zero multiplicity $m$.
The contour projection
$E_{\lambda_0}=(2\pi i)^{-1}\int_\gamma R(\lambda)\,d\lambda$
is a bounded projection.  Here $\gamma$ is a small
positively oriented circle containing no other possible pole.
These familiar resolvent facts follow directly in the present
setting: integrate
$R(z)R(w)=(R(w)-R(z))/(z-w)$ over nested circles to obtain
$E_{\lambda_0}^2=E_{\lambda_0}$ and commutation with the
resolvent.  The identity $AR(z)=zR(z)-1$ allows integration
in the graph norm and gives
$A^jE_{\lambda_0}=(2\pi i)^{-1}\int_\gamma z^jR(z)\,dz$.
Since the possible pole order is at most $m$,
\[
                  (A-\lambda_0)^mE_{\lambda_0}=0.
 \tag{MA-GW.20}
\]
More explicitly, the coefficient of
$(z-\lambda_0)^{-j}$ in the principal part is
$(A-\lambda_0)^{j-1}E_{\lambda_0}$, by the same
graph-norm contour identity.
The same statements hold on the continuous linear dual.

Identify that dual with functions $v$ such that
\[
       \ell_v(f)=\int_0^\infty f(y)v(y)\,dy,\qquad
       \int_0^\infty |v(y)|^2(1+y^2)^{-s/2}\,dy<\infty.
 \tag{MA-GW.21}
\]
The dual of the quotient is the annihilator of $\mathcal J$.
The transpose semigroup acts by $v(y)\mapsto v(y+t)$.
Therefore a vector in the dual range of $E_{\lambda_0}$
satisfies $(\partial-\lambda_0)^m v=0$ on $(0,\infty)$
in the distributional sense.
Its complete solution space consists of
$e^{\lambda_0y}p(y)$ with $\deg p<m$:
multiply by $e^{-\lambda_0y}$ and successively integrate
the equation $\partial^m q=0$ on the connected interval.
This conclusion requires no diagonalizability assumption.

If $\operatorname{Re}\lambda_0<0$, all these $m$
polynomials meet (MA-GW.21).
If $\operatorname{Re}\lambda_0=0$, the polynomial of
highest degree $j$ meets it exactly when $2j<s-1$.
This gives the upper bounds $m$ and $\min(m,b_s)$
for the projection rank.
They are also lower bounds.
The functionals
\[
                         \ell_j(f)=\int_0^\infty
                                      y^j e^{\lambda_0y}f(y)\,dy
 \tag{MA-GW.22}
\]
for precisely these allowed $j<m$ are continuous and
annihilate $\operatorname{ran}T_\phi$.
To see the latter, substitute the causal convolution
and expand $(y+t)^j$: every coefficient is a derivative
of $\Phi(-\lambda)$ at $\lambda=\lambda_0$ of order
at most $j<m$, and is zero.
Absolute convergence follows from the exponential
kernel bound and (MA-GW.21).
Their translates form the usual finite triangular
Jordan chain.  The resolvent integral (MA-GW.14),
applied to these functionals, is the finite geometric
resolvent of that chain; its small contour integral
acts as the identity.  Thus they belong to the dual
range of $E_{\lambda_0}$.
They are linearly independent by their polynomial
densities.  This proves equality of the rank bounds,
and hence finite rank of every principal part.
It also identifies its full dual range as (MA-GW.22).

Under $U_c^{-1}$ the functionals (MA-GW.22) become,
up to a nonzero scalar,
\[
              f\longmapsto\int_c^\infty
                     (x-c)^j e^{(\lambda_0+\alpha)x}f(x)\,dx .
 \tag{MA-GW.23}
\]
Since $\lambda_0+\alpha=\rho-1/2$, their spans are
the ordinary Mellin jets at that point, by the
invertible triangular change from $(x-c)^j$ to $x^j$.
This proves the local multiplicity assertion (MA-GW.12).

\emph{Uniform polynomial contour bounds.}
Choose $0<\epsilon<\eta$ and $C_0>u+2$.
On $\operatorname{Re}\lambda=\epsilon$,
(MA-GW.14) and (MA-GW.7) bound $R(\lambda)$ uniformly
in its imaginary part.
On the entire region $\operatorname{Re}\lambda\le-C_0$,
the Euler lower bound for $\zeta(d-\lambda)$ and
(MA-GW.16)--(MA-GW.18) give the uniform estimate
\[
                 \|R(\lambda)\|\le C(1+|\lambda|)^{N-1}.
 \tag{MA-GW.24}
\]
Here $|-\lambda-u|\ge c(1+|\lambda|)$ after increasing
$C_0$, and the Euler lower bound is uniform for
$\operatorname{Re}(d-\lambda)\ge d+C_0>1$.
In particular the constant in (MA-GW.24) does not grow
as the left line is moved farther left.

Use the good heights of MT-M18, (M18.7)--(M18.8), reflected
if necessary, to choose increasing $t_n=n+O(1)$ with
$1/2<t_{n+1}-t_n<3/2$.
That full minimum-modulus proof supplies a polynomial
lower bound for the finite zeta function on
$[-C_0,\epsilon]+it_n$ after the change $z=d-\lambda$.
It is independent of zero separation inside each block.
Combining it with (MA-GW.16) and the explicit polynomial
factors in $\Phi$ gives constants $C,B$ such that
\[
 \sup_{\lambda\in\Gamma_n}\|R(\lambda)\|
      \le C(2+|n|)^B,\qquad
 \Gamma_n=\partial\bigl([-C_0,\epsilon]
                              +i[t_n,t_{n+1}]\bigr).
 \tag{MA-GW.25}
\]
The bound on the right side uses the actual resolvent
integral, even if an apparent zero of the denominator
there has been canceled.
The horizontal sides avoid all zeros by construction.
All low-height contours can be adjusted individually.
Their lengths are uniformly bounded.
Their projections $E_n$ have finite rank
\[
                           \operatorname{rank}E_n
                                     =O(\log(2+|n|)),
 \tag{MA-GW.26}
\]
by the finite-rank proof and NT-Z10.
These estimates use one fixed central strip, together
with the globally uniform left-region bound
(MA-GW.24); no unstated uniform minimum modulus on
expanding strips is needed.

\emph{Causal smoothing removes every residual subspace.}
For $k\in C_c^\infty((0,\infty))$ put
$H_k(\lambda)=\int_0^\infty k(t)e^{\lambda t}\,dt$.
It decreases faster than every power of
$|\operatorname{Im}\lambda|$ on the fixed central strip.
Therefore the finite-rank operators
\[
       V_n(k)=\frac1{2\pi i}\int_{\Gamma_n}
                                  H_k(\lambda)R(\lambda)\,d\lambda
 \tag{MA-GW.27}
\]
have an absolutely convergent sum in trace norm:
\[
 \sum_n\|V_n(k)\|_1
 \le C_M\sum_n\log(2+|n|)(2+|n|)^{B-M}<\infty
                                             \quad(M>B+2).
 \tag{MA-GW.28}
\]
Their ranges are contained in $\operatorname{ran}E_n$,
by the finite-dimensional resolvent calculus already
proved from the contour identities.

To identify the sum as an actual operator, first insert
(MA-GW.14) on the line $\operatorname{Re}\lambda=\epsilon$.
Fourier inversion and Fubini, justified by the rapid
decay of $H_k$ and the exponentially damped polynomial
semigroup bound, give
\[
 C_k^{\mathcal Q}
     =\frac1{2\pi i}\int_{\epsilon-i\infty}^{\epsilon+i\infty}
                                    H_k(\lambda)R(\lambda)\,d\lambda .
 \tag{MA-GW.29}
\]
Summing finite rectangles and letting their good heights
tend to both infinities moves this integral to the line
$-C_0$ and picks up exactly (MA-GW.27).
The horizontal integrals tend to zero by (MA-GW.25).
One may then move the left line to $-M$ for any $M>C_0$:
there are no poles to its left, and (MA-GW.24) controls
the added horizontal segments.

Finally let $M\to\infty$.
If $\operatorname{supp}k\subset[a,b]$ with $a>0$,
integration by parts gives, for every integer $\ell$,
\[
 |H_k(-M+it)|
    \le C_{\ell,k}e^{-Ma}(1+M)^\ell(1+|t|)^{-\ell}.
 \tag{MA-GW.30}
\]
Combining (MA-GW.24) with (MA-GW.30), for $\ell>N+1$,
makes the norm of the remaining vertical integral tend
to zero, bounded by an exponentially decreasing factor
times a fixed power of $1+M$.
Consequently
\[
                    C_k^{\mathcal Q}=\sum_n V_n(k)
                              \quad\hbox{in trace norm}.
 \tag{MA-GW.31}
\]
This establishes an actual trace-class operator identity,
not just equality of its possible residues.

Choose a smooth positive-time approximate identity
$k_\varepsilon$ of integral one with support in
$[\varepsilon,2\varepsilon]$.
Strong continuity of the semigroup gives
$C_{k_\varepsilon}^{\mathcal Q}\to1$ strongly.
Equation (MA-GW.31) then proves density of the span
of the finite spectral blocks.
It also proves the exact closure assertion:
if every functional (MA-GW.22) vanishes on a quotient
vector, every $E_n$ kills that vector, hence
every $V_n(k)$ and every $C_k^{\mathcal Q}$ kill it,
and the approximate identity makes the vector zero.
There is no remaining subspace invisible to the
local principal parts.

On each finite block, the trace in (MA-GW.27) is its
length times $H_k(\lambda_\rho)$; the higher Jordan
terms have trace zero.
Trace-norm convergence permits summation of those
finite traces.
For the physical test $g$, take $k(t)=e^{\alpha t}g(t)$.
Then $H_k(\lambda_\rho)=\mathcal Mg(\rho-1/2)$.
Bounded similarity back to the actual $H_c/J_Y$
proves (MA-GW.13), including its absolute convergence.

\emph{A trace-class resolvent power.}
Fix $\nu>\epsilon$ and choose an integer
$q>\max\{B+2,N+1\}$.
Repeat the same contours with
$H(\lambda)=(\nu-\lambda)^{-q}$.
The right-line integral is
$(\nu-A)^{-q}$, by its vectorwise Laplace representation
with kernel $t^{q-1}e^{-\nu t}/(q-1)!$ and Fourier inversion.
The trace-norm block series converges by (MA-GW.25)--(MA-GW.26).
On the line $-M$, (MA-GW.24) bounds the remaining
integral by $C M^{N-q}$, which tends to zero.
The central horizontal integrals vanish by $q>B+2$,
and the farther-left ones by $q>N+1$.
Thus
\[
                              (\nu-A)^{-q}\in\mathcal S_1(\mathcal Q).
 \tag{MA-GW.32}
\]
Bounded similarity gives the same assertion for the
actual normalized quotient semigroup.
\end{proof}

\paragraph{MA-GW-003 (the actual weighted lift and the two orders of limits).}
Let $P_c^Z$ be the actual orthogonal projection onto
$Z_c=H_c\ominus J_Y$, extended by zero on $H_c^\perp$.
For every $g\in C_c^\infty((0,\infty))$,
\[
 C_gP_c^Z\in\mathcal S_1(H),\qquad
 \operatorname{Tr}_H(C_gP_c^Z)
        =\sum_\rho r^+_\rho\,\mathcal Mg(\rho-1/2).
 \tag{MA-GW.33}
\]
This is the actual convolution-after-projection operator.
For the fixed smooth exhaustion $\chi_R$ of MA-GP-009,
\[
 \begin{aligned}
 \lim_{Y\to\infty}\lim_{R\to\infty}
    \operatorname{Tr}(M_{\chi_R}C_gQ_Y^{\rm glob})
      &=-\sum_\rho r^+_\rho\,\mathcal Mg(\rho-1/2),\\
 \lim_{R\to\infty}\lim_{Y\to\infty}
    \operatorname{Tr}(M_{\chi_R}C_gQ_Y^{\rm glob})
      &=-\sum_\rho r^{\rm strip}_\rho\,\mathcal Mg(\rho-1/2),
 \end{aligned}
 \tag{MA-GW.34}
\]
where the actual full-strip multiplicities are
\[
 r^{\rm strip}_\rho=
 \begin{cases}
 m_\rho,&d<\operatorname{Re}\rho<u,\\
 \min(m_\rho,b_s),&\operatorname{Re}\rho=d
                                \ \hbox{or}\ \operatorname{Re}\rho=u,\\
 0,&\hbox{otherwise}.
 \end{cases}
 \tag{MA-GW.35}
\]
The first limit minus the second is exactly
\[
 -\sum_{\operatorname{Re}\rho<d}m_\rho\,\mathcal Mg(\rho-1/2)
 -\sum_{\operatorname{Re}\rho=d}
          \bigl(m_\rho-\min(m_\rho,b_s)\bigr)
                                    \mathcal Mg(\rho-1/2).
 \tag{MA-GW.36}
\]
Thus all $s\ge0$, including the draft's $s\ge3$, are covered.
No assertion that an omitted zero or omitted edge jet actually
exists is part of this formula.

\begin{proof}
\emph{The weighted adjoint domain on the actual half-line.}
Use the actual unitary $U_c:H_c\to H_{W_c}$ from
(MA-GW.2), and write $W=W_c$.
The normalized unilateral semigroup has adjoint
\[
 (S_t^{*,W}v)(y)=\frac{W(y+t)}{W(y)}v(y+t),
            \qquad L_Wv=W^{-1}\partial_y(Wv).
 \tag{MA-GW.37}
\]
The orthogonal complement of the actual ideal is invariant
under these adjoints.
Multiplication by $W$ is unitary onto $L^2(W^{-1}dy)$,
where the adjoints become left translations.
Their generator has no boundary condition at zero.
Exactly as in the explicit domain proof of MA-GP-009,
\[
 v\in\mathcal D(L_W^q)
 \quad\Longleftrightarrow\quad
 \partial^j(Wv)\in L^2(W^{-1}dy)\quad(0\le j\le q).
 \tag{MA-GW.38}
\]
The proof uses the fundamental theorem of calculus for
the direct difference quotient and interior test functions
for its converse, followed by induction.
The polynomial shift bound and fixed-$c$ norm equivalence
give the required translation bounds.
In particular $Wv\in H^q_{\rm loc}([0,\infty))$.

The trace-class quotient resolvent power of (MA-GW.32),
bounded similarity to the actual quotient, and Hilbert
adjunction give
\[
 \mathcal R_q=(\nu-L_W)^{-q}P_{Z_W}
                         \in\mathcal S_1(H_W),
                 \qquad Z_W=U_cZ_c .
 \tag{MA-GW.39}
\]
The full resolvent restricts to the adjoint invariant
subspace by its vectorwise Laplace integral.
The identities
$L_W(\nu-L_W)^{-1}=\nu(\nu-L_W)^{-1}-1$
show that $\mathcal R_q$ maps boundedly into the
$q$th graph domain.

\emph{All endpoint and cusp terms are finite rank.}
Let $k(t)=e^{\alpha t}g(t)$, extended by zero outside
$(0,\infty)$, and take $v\in\mathcal D(L_W^q)$.
Put $b=Wv$ and $a=W^{-1}$.
The exact integral before integration by parts is
\[
 C_k^+(\nu-L_W)^qv(x)
       =\int_0^\infty k(x-y)a(y)
                              (\nu-\partial_y)^q b(y)\,dy .
 \tag{MA-GW.40}
\]
When $c<0$, split this integral at $L=-c$; otherwise
there is no interior split.
On either smooth part all functions
$W\partial^j(W^{-1})$ are bounded.
Indeed
\[
 \log W(y)=2\alpha|c+y|-2\alpha(c+y)
                              +\frac{s}{2}\log(1+(c+y)^2),
 \tag{MA-GW.41}
\]
whose positive-order derivatives are bounded on each
part and have finite one-sided values at the cusp.
The bulk after $q$ integrations by parts is
\[
 \mathcal B_{k,q}v(x)
  =\int W(y)(\nu+\partial_y)^q
                         \bigl(k(x-y)/W(y)\bigr)v(y)\,dy .
 \tag{MA-GW.42}
\]
It is a finite sum of bounded causal convolutions
$C_{k^{(i)}}^+$ followed by bounded input multipliers.
Hence it is bounded on $H_W$.

For explicit control of the boundary signs and orders,
the term of derivative order $j\ge1$ on each interval
has boundary expression
\[
 \binom qj\nu^{q-j}
       \sum_{\ell=0}^{j-1}(-1)^{j+\ell}
        \left[\partial_y^\ell(k(x-y)a(y))\,
                    b^{(j-1-\ell)}(y)\right]_{\rm left}^{\rm right}.
 \tag{MA-GW.43}
\]
Terms at infinity vanish by the compact support of $k$.
At zero take the lower-endpoint term, and at $L$ take
the difference of the two one-sided expressions.
Because $b\in H^q_{\rm loc}$, its derivatives through
order $q-1$ have common traces at the cusp.
Every remaining output profile is a finite linear
combination of $k^{(i)}(x)$ or $k^{(i)}(x-L)$.
It is smooth and compactly supported in the half-line.
Every input functional is a trace $b^{(j)}(0)$ or
$b^{(j)}(L)$ with $j<q$.
The elementary one-dimensional Sobolev trace inequality
bounds these functionals by the graph norm in (MA-GW.38).
For $q=2$ and $s=0$, (MA-GW.43) specializes exactly
to the endpoint and $-4\alpha$ cusp terms of
(MA-GP.102)--(MA-GP.104).

Apply (MA-GW.40)--(MA-GW.43) to $v=\mathcal R_q f$.
Its left side is $C_k^+P_{Z_W}f$.
Its right side is the product
$\mathcal B_{k,q}\mathcal R_q$ plus finitely many
rank-one operators, since each graph trace composed
with $\mathcal R_q$ is bounded.
This proves actual ambient trace class.
Conjugating by $U_c$ and extending by zero gives
$C_gP_c^Z\in\mathcal S_1(H)$.
Causality ensures that $C_g$ preserves $H_c$.
Trace cyclicity with the actual orthogonal projection
therefore identifies its trace with quotient convolution,
already computed in (MA-GW.13).
This proves (MA-GW.33).

\emph{Both actual iterated limits.}
Let $P_c=\mathbf1_{[c,\infty)}$ on $H$.
The orthogonal decomposition is still
$Q_Y^{\rm glob}=P_c-P_c^Z$.
For compact smooth $\chi$, $M_\chi C_gP_c$ is trace
class with trace zero.
The proof is the one given in MA-GP-009: restrict a
smooth compact convolution kernel by bounded position
multipliers, conjugate by the two bounded compact
coordinate factors of the weight, and approximate
the half-line indicator strongly by smooth multipliers.
The diagonal is $g(0)\chi(x)=0$.
All these steps apply to the present polynomial weight
because it is bounded above and below on each compact
coordinate projection.
Thus
\[
 \operatorname{Tr}(M_{\chi_R}C_gQ_Y^{\rm glob})
                    =-\operatorname{Tr}(M_{\chi_R}C_gP_c^Z).
 \tag{MA-GW.44}
\]
For fixed $Y$, strong convergence of the uniformly
bounded $M_{\chi_R}$ and (MA-GW.33) give convergence
in trace norm of the right-hand product to $C_gP_c^Z$.
Its trace is independent of $Y$, proving the first
line of (MA-GW.34).

For the opposite order, MA-GP-001 gives the actual
fixed-position trace-norm condensation for every
$s\ge0$.  The subsequent actual global cutoff limit
is MA-WP-003, with its exact two boundary-jet counts.
That is precisely (MA-GW.35), proving the second line.
All sums are absolutely convergent by NT-Z10 and
rapid vertical decay of the test transform.
Their difference is (MA-GW.36); the upper-edge terms
cancel, while only the unretained lower-edge jets
remain there.
\end{proof}

\paragraph{Scope, primary inputs, and trace interpretation.}
The argument uses the exact source theta and Poisson formulas,
the classical containing strip and unconditional fixed-strip
growth for zeta, and the complete minimum-modulus construction
of MT-M18.  Its required lower bound is (M18.7)--(M18.8)
for the finite zeta factor on one fixed strip, with the
disk-Blaschke/Harnack proof given there.
The analytic antecedent and local multiplicity count are
NT-Z10, with the Ramakrishnan--Valenza locators recorded
in that module.  All new causal range, resolvent, finite-rank,
contour, and actual lift arguments are supplied above.
Neither $\mathsf A$ nor $M$ is used, and no stronger
zero-location assertion than the classical containing
strip is an input.

The result supplies the precise polynomial-weight extension
that an ordinary $s=0$ Hardy identification cannot supply.
At the critical equalities $2j=s-1$, the norm in
(MA-GW.21) diverges logarithmically, so those jets are
excluded without an unsupported density assertion at the
Sobolev threshold.
No uniform trace-norm bound as $Y\to\infty$, no unspecified
interchange of joint limits, and no equality with a local
additive Fourier projection is inferred.

In the inherited Grothendieck trace framework, the actual
finite-length modules arise as the finite-rank principal
parts of the quotient resolvent, not as a presupposed
spectral basis.  Their dense smoothed synthesis and the
actual ambient factorization justify the traces before
their finite-dimensional values are summed.
The available cote~1 locators remain scans~079--080 and
111--112, ALPOGE\_BRIDGES BR-TR-01/02; no new archive
reading or Artin-motive realization is claimed here.

\section{An actual joint-cutoff Gram remainder}
\label{sec:joint-cutoff-counterpath}

The source's two divergent cutoffs do not specify a relative rate.
The following construction keeps the actual global PW projection,
the actual Gram-normalized transport, and the repaired finite-place
local quotient as three explicitly defined operators.  It gives a
nonzero remainder along an admissible joint path.  The argument
does not evaluate the remainder by subtracting two zero sums.
It first constructs a nonzero trace on the actual quotient and
then uses trace-ideal convergence of the actual operators.

Throughout this section take the rational trivial sector,
$\alpha>0$, $s\ge0$, and
\[
 \begin{aligned}
 H&=L^2(\mathbb R,w(x)\,dx),&
 w(x)&=e^{2\alpha|x|}(1+x^2)^{s/2},\\
 J&=\overline{I_{1/2}}^{\,H},&
 Z&=J^\perp,&
 Q_Y^{\rm glob}&=P_{J_Y}.
 \end{aligned}
 \tag{MA-JP.1}
\]
The notation $I_{1/2}$ is the actual centered smooth theta ideal,
and $J_Y$ is the actual finite Fourier-band theta range of
MA-GP-001--008.  The quotient trace from MA-WP is denoted by
\[
 \tau(g)=\operatorname{Tr}_{H/J}C_g^{H/J}
        =\operatorname{Tr}_H(C_gP_Z),
 \qquad
 G(z)=\int_{\mathbb R}g(x)e^{zx}\,dx .
 \tag{MA-JP.2}
\]
The second equality is the actual ambient lift and trace cyclicity
of MA-WP/MA-LIFT, not a definition through a scalar divisor.

\paragraph{MA-JP-001 (the exact source rate conditions).}
Masullo, Definition~9.3.3.1, printed p.~427, specifies a
multiplicatively symmetric cutoff supported on $[1/X,X]$,
equal to one on $[1/X+\delta,X-\delta]$.
At the end of the local calculation, printed p.~519,
(12.4.1.25)--(12.4.1.28), the limits stated are
\[
       X,Y\longrightarrow\infty,\qquad
       \delta X\longrightarrow0,\qquad
       \delta Y\longrightarrow0 .
 \tag{MA-JP.3}
\]
Here $X=\lambda$, $Y=\lambda'$, and $\delta$ is the source's
boundary-smoothing parameter.  It is not the lower spatial
endpoint.  The footnote on p.~514 explicitly proposes the
annular replacement $[1/\lambda'',\lambda]$ with independent
$\lambda''$, and requests the convergence study of its volume
integral.  Neither that footnote nor the conditions on p.~519
prescribe $Y/\lambda''\to0$.  For the symmetric cutoff of
Definition~9.3.3.1, $\lambda''=X$, so the parameter of
MA-LA-014 is
\[
                     \epsilon Y=Y/X.
 \tag{MA-JP.4}
\]
The proof of Theorem~12.5.1.1, p.~520, invokes the local
calculation for the global geometric kernel but supplies no
additional relation between these two cutoffs or a quantitative
Gram estimate.  This is a precise reading of the cited passages,
not a claim that no possible additional scaling could work.

\paragraph{MA-JP-002 (a compact witness with both polar moments removed).}
Suppose the retained quotient divisor of MA-WP is nonempty:
$r_{\rho_0}>0$ for some nontrivial completed zero $\rho_0$.
Then there is a test $g\in C_c^\infty(\mathbb R)$ such that
\[
        g(0)=0,\qquad G(1/2)=G(-1/2)=0,\qquad \tau(g)=1 .
 \tag{MA-JP.5}
\]
The normalization of the last value is a harmless complex
scalar normalization of the test.  In particular its polar
term $P(g)=G(1/2)+G(-1/2)$ is zero.

\begin{proof}
Put $\lambda_0=\rho_0-1/2$ and
$q(z)=z^2-1/4$.  Because $\rho_0$ is a nontrivial completed
zero, $\lambda_0\ne\pm1/2$.
The actual full-jet interpolation of MT-M18 supplies
$f\in\mathcal E$ whose centered Mellin value at $\rho_0$
is $1/(r_{\rho_0}q(\lambda_0))$ and whose jets at all
other zeros are zero.  The finitely many jets in the
chosen block can be prescribed by the finite-dimensional
Hermite classes; no assertion about individual zero
separation is needed.
The differential operator
\[
                     \mathcal D=\frac{d^2}{dx^2}-\frac14
 \tag{MA-JP.6}
\]
multiplies the Mellin transform by $q(z)$.
Consequently $f_1=\mathcal D f$ has both polar moments
zero and $\tau(f_1)=1$.  Here the finite-block trace
$r_\rho G(\rho-1/2)$ is used only to choose a nonzero
vector in the already constructed actual trace functional.
Continuity of that functional on $\mathcal E$ follows
from the rapid block estimates of MA-WP, or directly
from the zero count and rapid Mellin decay.

There is $v\in I_{1/2}$ with both polar moments zero and
$v(0)=1$.  To see this, choose any nonzero
$\theta\in I_{1/2}$.  Then $\mathcal D\theta\ne0$:
the only smooth solutions of $\mathcal D\theta=0$
are linear combinations of $e^{x/2}$ and $e^{-x/2}$,
and neither nonzero combination belongs to $\mathcal E$.
The ideal is stable under differentiation and translation,
as is immediate from its full Mellin-jet description.
Translate $\mathcal D\theta$ to put a nonzero value at
zero, and divide by that value.
This produces $v$.  In particular $\tau(v)=0$.
Thus
\[
                        f_2=f_1-f_1(0)v
 \tag{MA-JP.7}
\]
still has the two zero polar moments and trace one,
and now vanishes at zero.

Choose $\kappa\in C_c^\infty(\mathbb R)$ equal to one
near zero and put $f_N=\kappa(x/N)f_2(x)$.
Then $f_N\to f_2$ in $\mathcal E$ and $f_N(0)=0$.
To correct both small polar errors exactly, choose a
nonnegative smooth bump $b$ supported near zero and
two distinct positive numbers $x_1,x_2$ so that
$b_j(x)=b(x-x_j)$ are supported away from zero.
The matrix
\[
 \begin{pmatrix}
 \mathcal Mb_1(1/2)&\mathcal Mb_2(1/2)\\
 \mathcal Mb_1(-1/2)&\mathcal Mb_2(-1/2)
 \end{pmatrix}
 =
 \begin{pmatrix}
 e^{x_1/2}\mathcal Mb(1/2)&e^{x_2/2}\mathcal Mb(1/2)\\
 e^{-x_1/2}\mathcal Mb(-1/2)&e^{-x_2/2}\mathcal Mb(-1/2)
 \end{pmatrix}
 \tag{MA-JP.8}
\]
is invertible.  Its determinant is a nonzero constant
times $2\sinh((x_1-x_2)/2)$.
Let $(c_{1,N},c_{2,N})$ solve this fixed system with
right side $(\mathcal Mf_N(1/2),\mathcal Mf_N(-1/2))$.
These coefficients tend to zero.  Hence
$u_N=f_N-c_{1,N}b_1-c_{2,N}b_2$ converges to $f_2$
in $\mathcal E$, vanishes at zero, and has both polar
moments exactly zero.  For all sufficiently large $N$,
$\tau(u_N)\ne0$ by continuity.  Set
$g=u_N/\tau(u_N)$.  This proves (MA-JP.5) with an
actual compact smooth test.
\end{proof}

\paragraph{The classical zeta instance and its exact input.}
For every $\alpha>0$ and every $s\ge0$, the hypothesis of
MA-JP-002 is satisfied for zeta by the classical theorem
of Hardy (1914), which gives infinitely many zeros on
$\operatorname{Re}s=1/2$.  Each such zero lies strictly
inside the corrected strip and retains its full multiplicity.
Only existence of one such zero is needed.
This invocation is unconditional arithmetic information
about a zero's location; it is not an absence of all
zero-location input.  It assumes neither the Riemann
hypothesis nor the theorem $\mathsf A$ nor $M$.
The primary statement is Hardy's opening paragraph,
printed p.~1012, with its argument on pp.~1012--1014.

\subsection{Two local comparisons with their Hilbert metrics stated}

Write $h(t)=g(\log t)$ and let $S$ be any fixed finite set
of rational primes, together with the real place.  It may
be chosen to include every prime power norm in the support
of $h$, as required when the source passes to p.~520.
On $\mathscr H_e=L^2_{\rm even}(\mathbb R,dt)$ use the
actual local operators of MA-LA-011--014,
\[
 \begin{aligned}
 (A_hu)(t)&=\int_0^\infty h(r)r^{-1/2}u(t/r)\,\frac{dr}{r},\\
 Q_Y&=\mathcal F_e^{-1}M_{\mathbf1_{[-Y,Y]}}\mathcal F_e,&
 D_S&=\prod_{p\in S_f}D_p .
 \end{aligned}
 \tag{MA-JP.9}
\]
The finite-place product $D_S$ is the actual unitary
quotient Fourier factor, not an assigned prime scalar.
All factors in this paragraph have fixed $S$.
The logarithmic unitary
\[
 (\Gamma_\infty u)(x)=\sqrt2\,e^{x/2}u(e^x)
 \tag{MA-JP.10}
\]
sends $A_h$ to $C_g$ on unweighted $L^2(\mathbb R,dx)$.
The global weighted unitary of MA-GP-007 is instead
\[
 (\mathcal U_wu)(t)
       =\sqrt{\frac{2t}{w(\log t)}}u(t),\qquad t>0 .
 \tag{MA-JP.11}
\]
It maps $\mathscr H_e$ to the radial form of $H$.
For clarity, transporting global convolution by this
second unitary gives the distinct operator
\[
 \begin{split}
 (\mathcal U_w^{-1}C_g\mathcal U_wu)(t)
 =\int_0^\infty h(r)r^{-1/2}
       \sqrt{\frac{w(\log t)}{w(\log t-\log r)}}
                  u(t/r)\,\frac{dr}{r}.
 \end{split}
 \tag{MA-JP.12}
\]
Thus the actual local convolution and the same-norm
additive comparison must not be silently equated.
Both comparisons below are carried out explicitly.

Let $E_Y:\mathcal P_Y^0\to\mathscr H_e$ be the ordinary
band inclusion of MA-GP-008 and set
\[
 W_Y^S=\mathcal U_wD_S\mathcal F_eE_Y,\qquad
 \Pi_Y^{S,0}=W_Y^S(W_Y^S)^* .
 \tag{MA-JP.13}
\]
This is an orthogonal projection in the same weighted
space $H$ as $Q_Y^{\rm glob}$.
For $S_f=\varnothing$ it is exactly the zero-integral
additive candidate in (MA-GP.73).
It is a precisely specified comparison, not an
operator identity asserted to be printed in the source.

The finite set $S$ may also be retained in the \emph{global}
PW family itself.  Let $\widehat J_{S,Y}$ be the $H$-closure
of the actual global theta images of polar-free
$S$-adelic Schwartz functions whose additive Fourier
support has product norm at most $Y$, extended by
$\mathbf1_{\mathbb Z_p}$ outside $S$ and projected to
the trivial compact type.  This is the rational invariant
PW level of Definition~7.3.2.3, with no additional
finite transition set.  Write
$\widehat Q_{S,Y}=P_{\widehat J_{S,Y}}$.
There are literal inclusions
\[
         J_Y\subseteq\widehat J_{S,Y}\subseteq J,
       \qquad \widehat Q_{S,Y}\longrightarrow P_G
                          \quad\hbox{strongly}.
 \tag{MA-JP.13a}
\]
For the first inclusion use the actual vector
$f_\infty\otimes\bigotimes_{p\in S_f}\mathbf1_{\mathbb Z_p}$.
Its finite Fourier factors are again
$\mathbf1_{\mathbb Z_p}$, so their norms are at most one,
and an archimedean band $Y$ meets the stated product
support condition.  After extending by the outside
integer balls, the global rational summation is exactly
the same $2\sqrt t\sum_{n\ge1}f_\infty(nt)$ as before.
Both global polar conditions follow from those of
$f_\infty$.  The second inclusion holds because these
are actual global polar-free theta images in the
centered smooth ideal, using its Meyer/MT-M24
identification.  The levels increase, and the union
of the smaller $J_Y$ is dense in $J$ by MA-GP-001.
The orthogonal projection argument of that theorem
proves the last assertion.  This proves only the
inclusions and strong limit needed here; it does not
identify the mixed-place finite range with the
one-place Hardy model or its frequency Gram.

\paragraph{MA-JP-003 (a nonzero actual Gram remainder along a source-admissible path).}
For $g$ as in (MA-JP.5) and fixed $S$, there are
$X_n=n+2$, integers $Y_n\ge X_n^2$, and positive
$\delta_n\to0$ with
\[
 \delta_nX_n\longrightarrow0,\qquad
 \delta_nY_n\longrightarrow0,\qquad Y_n/X_n\longrightarrow\infty,
 \tag{MA-JP.14}
\]
and smooth, multiplicatively symmetric spatial cutoffs
$\rho_n$ of Definition~9.3.3.1, of level $(X_n,\delta_n)$,
and even smooth frequency cutoffs $b_n$ at $Y_n$, such that
the following \emph{actual} traces have the limits
\[
 \begin{aligned}
 \operatorname{Tr}_H(M_{\rho_n}C_gQ_{Y_n}^{\rm glob})
       &\longrightarrow-1,\\
 \operatorname{Tr}_H(M_{\rho_n}C_g\widehat Q_{S,Y_n})
       &\longrightarrow-1,\\
 \operatorname{Tr}_H(M_{\rho_n}C_g\Pi_{Y_n}^{S,0})
       &\longrightarrow0,\\
 \operatorname{Tr}_{\mathscr H_e}
   \bigl(M_{\rho_n(|t|)}A_hD_S
           \mathcal F_e^{-1}M_{b_n}\mathcal F_eD_S^*\bigr)
       &\longrightarrow0 .
 \end{aligned}
 \tag{MA-JP.15}
\]
In the first three lines, $M_{\rho_n}$ means multiplication
by $\rho_n(e^x)$ in logarithmic coordinates.
The frequency cutoff satisfies $0\le b_n\le1$, is
supported on $[-Y_n,Y_n]$, and is one on
$[-Y_n+\delta_n,Y_n-\delta_n]$.
The last line is the actual repaired finite-$S$ local
quotient operator, with its nuclear section lift.
Its coordinate realization, measure, and finite-place
factors are precisely those of MA-LA-014.
Because $P(g)=0$, subtracting the polar term does not
alter any of these limits.

In particular each global trace minus either local
comparison trace tends to $-1$.  The one-place global
trace minus the same-norm additive trace also has the
literal Gram-resolvent expression of (MA-GP.74) with
$W_{Y_n}$ replaced by the stated isometry $W_{Y_n}^S$.

\begin{proof}
\emph{Compact-position trace ideals, including sharp cutoffs.}
For $X>1$ let
$a_X(t)=\mathbf1_{[1/X,X]}(t)$ and
$\chi_X(x)=a_X(e^x)$.
The operators
\[
 A_X=M_{\chi_X}C_g\quad\hbox{on }H,\qquad
 B_X=M_{a_X(|t|)}A_h\quad\hbox{on }\mathscr H_e
 \tag{MA-JP.16}
\]
are trace class, and both have trace zero.
Indeed insert a fixed smooth cutoff equal to one
on the compact output interval, and a second smooth
cutoff on its enlargement by $\operatorname{supp}g$.
The resulting convolution kernel is smooth on a compact
rectangle and hence trace class by its absolutely
summable Fourier expansion.  Multiplication by the
bounded characteristic function preserves trace class.
On the weighted space, the factors $w^{1/2}$ and
$w^{-1/2}$ are bounded on the two coordinate
projections.  They need no smoothness at the cusp.
They cancel on the diagonal, as may also be checked
by cyclicity after these compact cutoffs are inserted.
Smooth bounded approximations of $\chi_X$ converge
strongly; multiplication of the fixed trace-class
compact-kernel operator then converges in trace norm.
Its diagonal traces converge to
$g(0)\int\chi_X(x)\,dx=0$.
Formula (MA-JP.10) proves the identical assertion for
$B_X$.  This argument also proves trace-norm convergence
when either sharp spatial cutoff is smoothed at fixed $X$.

As $X\to\infty$, multiplication by $\chi_X$ tends
strongly to the identity, with norm one.  The actual
ambient lift $C_gP_Z$ is trace class by MA-WP.
Consequently
\[
 \begin{split}
 \operatorname{Tr}_H(A_XP_G)
 &=\operatorname{Tr}_H A_X-
            \operatorname{Tr}_H(M_{\chi_X}C_gP_Z)\\
 &\longrightarrow-\operatorname{Tr}_H(C_gP_Z)=-1.
 \end{split}
 \tag{MA-JP.17}
\]
This is a limit of traces of actual trace-class
operators.  No spectral scalar sum has been
substituted for either cutoff operator.

\emph{The fixed-position convergence statements.}
For every fixed $X$, MA-GP-001 and finite-rank
trace-ideal approximation give
\[
            \|A_X(Q_Y^{\rm glob}-P_G)\|_1\longrightarrow0 .
 \tag{MA-JP.18}
\]
The same argument applies to the sharp $A_X$ just
proved trace class, since the PW projections are
self-adjoint and converge strongly.
It also gives
$\|A_X(\widehat Q_{S,Y}-P_G)\|_1\to0$
by (MA-JP.13a).

Also $\Pi_Y^{S,0}\to1$ strongly on $H$.
On the additive frequency space the band projection
onto zero-integral even functions is
\[
 P_Y^0=M_{\mathbf1_{[-Y,Y]}}-e_Y\otimes e_Y,\qquad
              e_Y=(2Y)^{-1/2}\mathbf1_{[-Y,Y]} .
 \tag{MA-JP.19}
\]
The first term converges strongly to one.
The unit vectors $e_Y$ converge weakly to zero:
this follows first for compactly supported $L^2$
functions by their $L^1$ bound, and then by density.
Thus their rank-one projections converge strongly
to zero.  Conjugation by the fixed unitaries in
(MA-JP.13) proves the claim.
Finally $D_SQ_YD_S^*\to1$ strongly on $\mathscr H_e$.
All these operators are self-adjoint contractions.
Therefore
\[
 \begin{aligned}
 \|A_X(\Pi_Y^{S,0}-1)\|_1&\longrightarrow0,\\
 \|B_X(D_SQ_YD_S^*-1)\|_1&\longrightarrow0.
 \end{aligned}
 \tag{MA-JP.20}
\]
The elementary trace-ideal fact used here follows
for a rank-one $u\otimes\ell$ from strong convergence
of the adjoints acting on $\ell$, and then for every
trace-class operator by finite-rank approximation.

\emph{Selecting one path before smoothing.}
For $X_n=n+2$, choose an integer $Y_n\ge X_n^2$
so large that each norm in (MA-JP.18) and
(MA-JP.20), and the corresponding norm for
$\widehat Q_{S,Y}$, at $X=X_n$, is at most $1/n$.
All four requirements hold for every sufficiently
large $Y$, so they can be met simultaneously.
Together with $\operatorname{Tr}A_{X_n}
=\operatorname{Tr}B_{X_n}=0$ and (MA-JP.17),
this already proves all four limits (MA-JP.15)
for sharp cutoffs.
It is an actual diagonal choice specified by four
trace-norm inequalities.  No unproved rate such as
the sufficiency of the exact power $Y=X^2$ is asserted.

\emph{The source's smooth symmetric cutoffs.}
Choose a fixed smooth $b:\mathbb R\to[0,1]$ with
$b(u)=0$ for $u\le0$ and $b(u)=1$ for $u\ge1$.
For fixed $X\ge2$ and sufficiently small $\delta>0$
put
\[
 \begin{aligned}
 \rho_{X,\delta}(t)
  &=b((X-t)/\delta)\,b((X-t^{-1})/\delta),\qquad t>0,\\
 b_{Y,\delta}(\xi)
  &=b((Y-|\xi|)/\delta).
 \end{aligned}
 \tag{MA-JP.21}
\]
The second function is smooth at zero because it
is identically one in a neighborhood of zero.
The first is smooth, multiplicatively symmetric,
supported on $[1/X,X]$, and equal to one on
$[1/X+\delta,X-\delta]$ when $X(X-\delta)\ge1$.
Indeed $1/(X-\delta)-1/X
=\delta/(X(X-\delta))\le\delta$.
Thus it is a cutoff of the exact source definition.
As $\delta\downarrow0$ both cutoffs converge boundedly
and pointwise almost everywhere to their sharp versions.

At the already fixed $X_n,Y_n$, the spatial trace-norm
convergences proved after (MA-JP.16) hold.
Moreover
$\mathcal F_e^{-1}M_{b_{Y_n,\delta}}\mathcal F_e$
converges strongly, with its adjoint, to $Q_{Y_n}$.
Splitting the difference of the two local products
at the fixed sharp $B_{X_n}$ and using (MA-JP.20)'s
trace-ideal argument shows that the entire operator
in the last line of (MA-JP.15) converges in trace norm to
$B_{X_n}D_SQ_{Y_n}D_S^*$ as $\delta\downarrow0$.
Every factor has its stated actual Hilbert domain.

Choose $\delta_n>0$ after $X_n,Y_n$ so that all
four resulting changes of trace have absolute
value at most $1/n$ and
\[
             \delta_n<
        \min\{1/(4X_n),\,1/(n(X_n+Y_n))\}.
 \tag{MA-JP.22}
\]
This proves (MA-JP.14)--(MA-JP.15), with
$\rho_n=\rho_{X_n,\delta_n}$ and
$b_n=b_{Y_n,\delta_n}$.
The smoothing does not change the band defining
the \emph{global} orthogonal PW projection; it
smooths the actual local additive Fourier cutoff.
Those are two distinct operations in the source
comparison, as their definitions require.

\emph{The literal Gram remainder.}
For the smoothed data put
$A_n=M_{\rho_n}C_g$ and use the closed
$T_{Y_n}$, Gram $G_{Y_n}$ and core-form operator
$K_{\rho_n,h}$ of MA-GP-008.
At each fixed $n$ that theorem gives
\[
 \begin{split}
 &\operatorname{Tr}_H
      \bigl(A_n(Q_{Y_n}^{\rm glob}-\Pi_{Y_n}^{S,0})\bigr)\\
 &\quad=\lim_{\eta\downarrow0}
    \operatorname{Tr}_{\mathcal P_{Y_n}^0}
     \left((G_{Y_n}+\eta)^{-1/2}
             K_{\rho_n,h}(G_{Y_n}+\eta)^{-1/2}
                    -(W_{Y_n}^S)^*A_nW_{Y_n}^S\right).
 \end{split}
 \tag{MA-JP.23}
\]
Only the bounded isometry $W_{Y_n}^S$ has changed
from (MA-GP.74), so the same cross-space cyclicity
proves this identity.  If a single fully finite
sequence is desired, choose $\eta_n>0$ at each
fixed $n$ so that the Gram trace differs from
its limit by at most $1/n$.
Then the trace of the fully bounded, trace-class
operator inside (MA-JP.23), with $\eta=\eta_n$,
tends to $-1$.  This retains both Gram resolvents
and is an actual nonvanishing finite-operator
remainder.
\end{proof}

\paragraph{What this path does and does not decide.}
The result disproves a rate-independent identification
of the two actual traces under only the joint
divergence and thin-smoothing conditions printed on
p.~519, when the annular repair is made.
It also prevents dropping the exact Gram correction
uniformly on that allowed collection of paths.
The path has $Y_n/X_n\to\infty$.
It does not refute the separately imposed local
repair $Y/X\to0$, nor prove a quantitative Gram
asymptotic on that different regime.
It does not exchange the two iterated limits of
MA-GP-009.
The nonzero quotient trace is selected from a
retained block, so this argument is different
from MA-EG's criterion involving discarded blocks.
In particular it is already instantiated by a
critical-line zero, independently of whether any
exterior zero exists.

The finite-prime set can meet the source's support
condition exactly: after $g$ is fixed, choose
$T>\max\{e^x,e^{-x}:x\in\operatorname{supp}g\}$ and
include every rational prime at most $T$ in $S_f$.
Then every omitted prime has norm greater than $T$.
The test is in the trivial compact isotype, so the
outside-unit invariance condition also holds.
The proof used only that this set is fixed and
finite; its actual $D_S$ was retained throughout.
The second line of (MA-JP.15) also retains the
actual global invariant PW family for that same
set $S$.  The explicit Gram-resolvent kernel
(MA-JP.23) remains the one-place parametrization
for which MA-GP-008 proves it; no uncomputed
mixed-place Gram kernel is assigned to the
second line.

\paragraph{Grothendieck placement and source scope.}
The witness is chosen in one finite-length
$\mathbb C[z]$-module of the actual completed
spectral quotient.  Its trace normalization is
the finite-dimensional trace with its retained
length.  The subsequent limits use actual
Hilbert trace ideals and the nuclear section
for the local quotient.  This is an instance of
the inherited cote~1 trace comparison
(scans~079--080 and~111--112, BR-TR-01/02);
it is not a new arithmetic positivity form,
an Artin-motive zero bound, or an identification
of the local and global Hilbert metrics.

The source-rate statement was checked in the
supplied Masullo text at Definition~9.3.3.1,
p.~427, the footnote to Theorem~12.4.1.1,
p.~514, (12.4.1.25)--(12.4.1.28), p.~519,
and the complete proof of Theorem~12.5.1.1,
p.~520.  These are fresh text readings;
no new rendered-page inspection is claimed here.
Hardy's primary theorem statement was checked
in its reproduced text at printed p.~1012.
Its full 1914 proof is used as a classical
published theorem, not presented as a newly
verified proof or numerical zero certificate.

Hardy, G.~H. (1914). Sur les z\'eros de la fonction
$\zeta(s)$ de Riemann. \emph{Comptes rendus
hebdomadaires des s\'eances de l'Acad\'emie
des sciences, 158}, 1012--1014.
Primary-page reproduction:
\url{https://fr.wikisource.org/wiki/Page:Comptes_rendus_hebdomadaires_des_s\%C3\%A9ances_de_l\%E2\%80\%99Acad\%C3\%A9mie_des_sciences,_tome_158,_1914.djvu/1014}.

\section{The exact rational finite-place invariant PW range}
\label{sec:rational-mixed-pw}

The mixed-place global range in MA-JP can be identified exactly
in the rational trivial sector.  The Fourier support condition
here is the invariant product-norm condition, and no additional
unit-shell transition set is inserted.  This distinction keeps
finite Euler-deletion factors visible.

\paragraph{MA-GP-010 (the actual invariant global PW range is independent of the finite set).}
Let $S=\{\infty\}\cup S_f$ be any finite set of rational
places, and take $S'=S$ in Masullo Definition~7.3.2.3,
pp.~354--355.  For $Y>0$ let $\mathscr W_{S,Y}$ be the
smooth closed global image of the following actual data:
\[
 \begin{gathered}
 \varphi_S\in\mathcal S_0(\mathbb A_{\mathbb Q,S}),\qquad
 \widehat\varphi_S\text{ compactly supported},\\
 \operatorname{supp}\widehat\varphi_S
       \subseteq\{\xi:|\xi_\infty|\prod_{p\in S_f}|\xi_p|_p\le Y\},\\
 \varphi=\varphi_S\otimes
                  \bigotimes_{p\notin S_f}\mathbf1_{\mathbb Z_p},
 \end{gathered}
 \tag{MA-GP.110}
\]
followed by the trivial compact-character projection and the
norm-half centering of the actual global theta map.
The polar conditions on $\varphi_S$ mean its value and its
Fourier value at zero are both zero.
The closure is taken in the actual exponential Schwartz
image topology, as in MT-M24 and the source PW definition.
Then, for the one-place space of MA-GP-002,
\[
                    \mathscr W_{S,Y}=\mathscr P_Y
      =\{g\in I_{1/2}:\operatorname{supp}g
                                      \subseteq[-\log Y,\infty)\}.
 \tag{MA-GP.111}
\]
In the last expression $g$ is written in logarithmic
coordinates.  Consequently, at every $\alpha>0$ and $s\ge0$,
\[
            \widehat J_{S,Y}=J_Y,\qquad
            \widehat Q_{S,Y}=Q_Y^{\rm glob}.
 \tag{MA-GP.112}
\]
These are literal equalities of the actual global weighted
ranges and their orthogonal projectors.  They strengthen the
inclusions in (MA-JP.13a).

\begin{proof}
\emph{Normalization and the canonical inclusion.}
Use the coherent Fourier character of MA-NR-001: positive
at infinity, negative fractional part at finite places,
and conjugate Fourier kernel.  Every finite integer ball
$\mathbf1_{\mathbb Z_p}$ is self-dual, with additive volume
one.  Conjugating a finite additive sign reflects that
Fourier coordinate and preserves the product-norm bound
in (MA-GP.110), so the subspace described by the source
support condition is unaffected by this explicit repair.

For any even real Schwartz function $f$ with
$f(0)=\widehat f(0)=0$ and Fourier support in $[-Y,Y]$,
take
\[
           \varphi_S=f\otimes
                    \bigotimes_{p\in S_f}\mathbf1_{\mathbb Z_p}.
 \tag{MA-GP.113}
\]
It is polar free.  Its Fourier support is compact and has
product norm at most $Y$, because each finite Fourier
coordinate has norm at most one.
After the outside factors in (MA-GP.110) are attached, the
nonzero principal elements in the actual theta sum are
precisely the nonzero integers.  The centered sum on the
norm splitting is therefore exactly
\[
                       2\sqrt t\sum_{n\ge1}f(nt).
 \tag{MA-GP.114}
\]
The vector is invariant under every finite unit and under
the real sign, so the probability compact projection fixes
it.  This is the normalization of MA-GP-001, with no finite
Euler factors deleted and no additional volume scalar.
Thus $\mathscr P_Y\subseteq\mathscr W_{S,Y}$.

\emph{The exact reverse support bound.}
Take arbitrary data (MA-GP.110).  The global value and integral
at zero still vanish, since the outside factors have both
value and integral one.  Global Poisson consequently relates
its centered theta sum at norm $t$ to $t^{-1/2}$ times the
sum of its Fourier transform at the inverse idele.
All outside-$S$ Fourier factors remain
$\mathbf1_{\mathbb Z_p}$.

Every norm-one rational idele class has a representative
with real component in $\{\pm1\}$ and all finite components
in $\mathbb Z_p^\times$: remove its finite valuations by a
principal rational element and use the product formula.
Thus these representatives have product norm one in $S$
as well.  We may perform compact averaging on this compact
cover, modulo its simultaneous sign.
For any such representative $u$ and any nonzero rational
$a$ contributing to the reflected theta sum, the outside
conditions give
$|a|_p\le1$ for $p\notin S_f$.
Hence
\[
          |a|_S:=|a|_\infty\prod_{p\in S_f}|a|_p
                =\left(\prod_{p\notin S_f}|a|_p\right)^{-1}
                \ge1.
 \tag{MA-GP.115}
\]
At the inverse norm splitting, its $S$-argument has product
norm $|a|_S/t\ge1/t$, because $|u|_S=1$.
If $0<t<1/Y$, this exceeds $Y$, and the source Fourier
support condition makes every summand zero.
The conclusion holds for every compact representative,
so compact averaging preserves it.

The resulting centered function belongs to the actual smooth
ideal $I_{1/2}$ by MT-M24, and is zero for $t<1/Y$.
MA-GP-002 now places it in $\mathscr P_Y$.
That space is closed in the smooth image topology, by the
actual theta isomorphism and the closed Fourier-band condition.
Thus the source's smooth closure introduces no further vectors.
This proves the reverse inclusion and (MA-GP.111).
Taking the same weighted Hilbert closure on both sides proves
(MA-GP.112) without an interchange of a support intersection
and a Hilbert closure.
\end{proof}

\paragraph{The exact consequence for the Gram comparison.}
The one-place frequency parametrization $T_Y$ of MA-GP-005
therefore parametrizes this same actual finite-$S$ global
range.  Its closed theta Gram $G_Y$, the compact core kernel
$K_{\chi,h}$ and both Gram resolvents in MA-GP-008 are a
literal parametrization of $\widehat Q_{S,Y}$ as well as
of $Q_Y^{\rm glob}$.
In particular the two global traces in (MA-JP.15) are equal
at every finite $n$, not merely in their limit, and the
actual Gram remainder (MA-JP.23) applies to that mixed-place
source range itself.

This does not identify an auxiliary archimedean frequency
norm with the additive $S$-adelic cover norm or with the
local quotient norm.  It identifies the global image and
its projection in the one specified weighted global norm.
The repaired local quotient Fourier factor $D_S$ of
MA-LA-011--014 remains part of the distinct local operator;
it has not disappeared from a local Fourier calculation.

\paragraph{Why the absence of an extra transition set matters.}
At the level of actual global functions, insertion of
$\mathbf1_{\mathbb Z_p^\times}$ at an additional unramified
prime replaces the spherical factor by
$\mathbf1_{\mathbb Z_p}-\mathbf1_{p\mathbb Z_p}$.
For a theta image $F(t)$ this gives the exact operator
\[
                 F(t)\longmapsto F(t)-p^{-1/2}F(pt).
 \tag{MA-GP.116}
\]
Indeed the second theta sum runs over integers divisible
by $p$, and reindexing by $n=pm$ gives the factor
$p^{-1/2}$ after centering.  The same reindexing works with
arbitrary fixed $S$-local data by moving the principal
$p$ through the sum; its unit components at primes in
$S_f$ are absorbed by the trivial compact average.
For a fixed finite extra set $T$ the product of these
commuting Euler factors is retained.  It can shift the lower
support edge by $-\log\prod_{p\in T}p$, and the invariant
Fourier product bound changes correspondingly.
The theorem does not replace this Euler-deleted image by
the primitive image at the same level.
This is exactly the distinction between $S'=S$ and the
additional unit-shell transitions in (7.3.2.13).

\paragraph{Source and trace-framework scope.}
The exact data and compact Fourier support are
Definition~7.3.2.3, (7.3.2.12)--(7.3.2.17), pp.~354--355;
its support proof is Proposition~7.3.2.4, pp.~356--357.
The corresponding global orthogonal projection is
Definition~11.2.1.3, (11.2.1.7)--(11.2.1.8), pp.~488--489.
These passages were freshly read in full for the stated
specialization.  The theorem supplies the exact rational
support constant, using the actual compact representatives
and both directions of the already proved one-place support
identity; it does not import the source's claimed trace
factorization.

The underlying smooth quotient is the actual Tate--Poisson
image of MT-M24, and the passage to its Gram-normalized
trace is MA-GP-008.  In the inherited cote~1 framework
(BR-TR-01/02, scans~079--080 and~111--112) this identifies
which concrete range the trace-ideal maps act on before
cyclicity is used.  It adds no motive or weight claim and
uses no zero-location assumption.

\section{The actual slow-band Gram asymptotic}
\label{sec:slow-band-gram}

The preceding counterpath does not address the repaired regime
$Y/X\to0$.  This section proves the actual global asymptotic
in that regime at polynomial weight zero.  The proof retains
the weighted Gram projection and estimates its trace correction.
It requires no generic choice of strip edge and no separation
of zeros from that edge.

\subsection{An abstract trace-class Hankel tail}

\paragraph{MA-SB-001 (uniform rectangles from the actual model-space lift).}
Let $B$ be an inner function on the right half-plane.
Use the Laplace isometry
$\mathcal L:L^2(0,\infty)\to H^2(\operatorname{Re}z>0)$
with boundary norm $(2\pi)^{-1}\int|F(it)|^2dt$.
On $\mathcal K_+=L^2(0,\infty)$ put
\[
 V=\mathcal L^{-1}M_B\mathcal L,\qquad
 P_K=1-VV^* .
 \tag{MA-SB.1}
\]
Thus $V$ is an isometry and $P_K$ is the actual orthogonal
model-space projection in time coordinates.
Let $k\in C_c^\infty((0,\infty))$ and let
$C_k^+$ be its causal convolution on $\mathcal K_+$.
Assume the \emph{actual ambient lift}
\[
                         C_k^+P_K\in\mathcal S_1(\mathcal K_+).
 \tag{MA-SB.2}
\]
This is a hypothesis about an actual operator, stronger
than merely assigning a summable scalar spectrum.
For the zeta inner function used below, MA-GP-009 proves it
without an edge-separation assumption.

Let $\mathcal K=L^2(\mathbb R)$, and let $P_+$ and $P_-$
be its positive- and negative-half-line projections.
Multiplication by the boundary function $B(it)$ defines
a full-line unitary convolution operator $\mathcal V$.
Define
\[
 \begin{aligned}
 H_B&=P_+\mathcal V|_{\mathcal K_-},&
 V_-&=P_-\mathcal V|_{\mathcal K_-},\\
 H_k&=P_+C_k|_{\mathcal K_-},&
 H_{Bk}&=P_+C_k\mathcal V|_{\mathcal K_-}.
 \end{aligned}
 \tag{MA-SB.3}
\]
The symbol of the last operator is
$B(it)\int_0^\infty k(r)e^{-itr}dr$.
Then
\[
 \begin{gathered}
 H_BH_B^*=P_K,\qquad H_B=P_KH_B,\\
 H_{Bk}=C_k^+P_KH_B+H_kV_-\in
                  \mathcal S_1(\mathcal K_-,\mathcal K_+).
 \end{gathered}
 \tag{MA-SB.4}
\]
For $R>L\ge0$ the actual tail rectangle satisfies
\[
 \left\|P_{[R,\infty)}VC_k^+P_{[0,L]}\right\|_1
 \le
 \left\|P_{[R-L,\infty)}H_{Bk}\right\|_1.
 \tag{MA-SB.5}
\]
The right side tends to zero as $R-L\to\infty$,
uniformly in $L$.  The same bound holds if the output
projection on the left is replaced by any multiplication
operator of norm at most one supported in $[R,\infty)$.

\begin{proof}
The Hardy invariance corresponding to causality gives
$P_-\mathcal VP_+=0$.  Relative to
$\mathcal K_+\oplus\mathcal K_-$, its block form is
\[
                      \mathcal V=
                       \begin{pmatrix}V&H_B\\0&V_-\end{pmatrix}.
 \tag{MA-SB.6}
\]
The first diagonal block of $\mathcal V\mathcal V^*=1$
gives $VV^*+H_BH_B^*=1$, proving the first identity
in (MA-SB.4).
Since $P_K$ is an orthogonal projection,
$(1-P_K)H_BH_B^*(1-P_K)=0$, and hence
$(1-P_K)H_B=0$.
All the displayed blocks have norm at most one.

After identifying a negative variable with its opposite,
the kernel of $H_k$ is $k(x+y)$ for $x,y>0$.
It is the restriction of a smooth kernel supported on
a fixed compact square: choose smooth factors equal to
one on $[0,b]$, where $\operatorname{supp}k\subset(0,b)$.
Multiplying $k(x+y)$ by those factors does not change it
on the positive quadrant.
A double Fourier expansion on a larger square is
absolutely summable after repeated integration by parts.
It gives a rank-one trace-class representation; restricting
to the positive quadrant preserves trace class.
Thus $H_k\in\mathcal S_1$.

Inserting $P_++P_-$ between $C_k$ and $\mathcal V$ gives
\[
                      H_{Bk}=C_k^+H_B+H_kV_- .
 \tag{MA-SB.7}
\]
The preceding identity $H_B=P_KH_B$ and hypothesis
(MA-SB.2) prove the trace-class assertion in (MA-SB.4).
No differentiability of $B(it)$ is used.
In particular zeros may approach the boundary arbitrarily
closely at large heights.

The full-line convolution operators $C_k$ and
$\mathcal V$ commute.  On positive-half-line inputs their
product agrees with $VC_k^+$.
Extend the operator in (MA-SB.5) by zero to $\mathcal K$
and translate both input and output coordinates by $L$.
Translation commutes with this full-line convolution.
The resulting operator is exactly
\[
        P_{[R-L,\infty)}\,C_k\mathcal V\,P_{[-L,0]}
       =P_{[R-L,\infty)}\,H_{Bk}\,P_{[-L,0]} .
 \tag{MA-SB.8}
\]
Translations are unitary and compression is contractive
for the trace norm, proving (MA-SB.5).
Finally $P_{[r,\infty)}\to0$ strongly on $\mathcal K_+$
as $r\to\infty$.
Left multiplication of a fixed trace-class operator
by these uniformly bounded projections converges in
trace norm: this holds first for rank-one operators
and then by finite-rank approximation.
The multiplication-profile version follows by factoring
its norm-one multiplier through the same tail projection.
\end{proof}

\subsection{The weighted Gram correction on the positive tail}

Fix $0<\alpha<1/2$, polynomial weight $s=0$, and
$g\in C_c^\infty((0,\infty))$.
Put
\[
 \begin{gathered}
 H=L^2(\mathbb R,e^{2\alpha|x|}\,dx),\qquad
 k(t)=e^{\alpha t}g(t),\\
 \sigma_+=\tfrac12+\alpha,\qquad
 \tau_+(g)=\operatorname{Tr}_{\mathcal K_+}(C_k^+P_K),
 \end{gathered}
 \tag{MA-SB.9}
\]
where $B=B_\alpha$, $V$ and $P_K$ are the actual
half-plane objects of MA-GP-003/009.
The same fixed $B$ and $k$ are used for every band $Y$.
MA-GP-009, including its endpoint term at zero, gives
(MA-SB.2) and identifies this trace with the actual
finite-level quotient trace.

For $Y\ge1$ write $L=\log Y$ and use $r=x+L\ge0$ on
the actual support half-line $x\ge-\log Y$.
Its metric is
\[
 R_L(r)=
 \begin{cases}
 e^{4\alpha(L-r)},&0\le r\le L,\\
 1,&r\ge L.
 \end{cases}
 \tag{MA-SB.10}
\]
The map $f(x)\mapsto e^{\alpha x}f(x)$, followed by this
translation of coordinates, is unitary from that actual
weighted half-line into $L^2(R_L(r)\,dr)$.
Let
\[
 \mathfrak G_L=V^*M_{R_L}V,\qquad
 D_L=\mathfrak G_L^{-1}-1,\qquad
 Q_L^w=V\mathfrak G_L^{-1}V^*M_{R_L}.
 \tag{MA-SB.11}
\]
The last operator is the actual orthogonal projection
for the $R_L$ metric, as in MA-GP-004.
The identity map from $L^2(R_Ldr)$ to $\mathcal K_+$
is a bounded isomorphism for each fixed $L$.
It preserves traces under conjugation, but is not
being asserted uniformly isometric as $L$ grows.

\paragraph{MA-SB-002 (uniform actual slow-band trace at every strip edge).}
Define the two fixed trace-norm tails
\[
 \begin{aligned}
 \omega_k(r)&=\|P_{[r,\infty)}C_k^+P_K\|_1,\\
 \nu_k(r)&=\|P_{[r,\infty)}H_{Bk}\|_1 .
 \end{aligned}
 \tag{MA-SB.12}
\]
Both tend to zero as $r\to\infty$, by
MA-GP-009 and MA-SB-001.
Let $0\le\phi\le1$ be a compactly supported
multiplication profile on $[0,\infty)$ equal to one
on $[0,R]$, where $R>L$.
Then the actual weighted trace satisfies
\[
 \left|
 \operatorname{Tr}_{L^2(R_Ldr)}(M_\phi C_k^+Q_L^w)
                     +\tau_+(g)\right|
            \le\omega_k(R)+\nu_k(R-L).
 \tag{MA-SB.13}
\]
In particular let $\rho$ be any smooth source cutoff
of Definition~9.3.3.1, supported on $[1/X,X]$ and
equal to one on $[1/X+\delta,X-\delta]$.
If
\[
           Y\ge1,\qquad X\ge2Y,\qquad
           0<\delta\le1/(4X),
 \tag{MA-SB.14}
\]
then, with $Q_Y^{\rm glob}$ the actual global PW
projection,
\[
 \begin{aligned}
 &\left|\operatorname{Tr}_H(M_\rho C_gQ_Y^{\rm glob})
                    +\tau_+(g)\right|\\
 &\qquad\le
 \omega_k(\log Y+\log(X-\delta))
                 +\nu_k(\log(X-\delta))\\
 &\qquad\le
 \omega_k(\log(X/2))+\nu_k(\log(X/2)).
 \end{aligned}
 \tag{MA-SB.15}
\]
As before, $M_\rho$ denotes $\rho(e^x)$ in logarithmic
coordinates.
Thus along every path with $X,Y\to\infty$,
$Y/X\to0$ and the source's thin smoothing condition,
\[
           \operatorname{Tr}_H(M_\rho C_gQ_Y^{\rm glob})
                      \longrightarrow-\tau_+(g).
 \tag{MA-SB.16}
\]
This limit is uniform in $1\le Y\le X/2$ in the
explicit sense of (MA-SB.15).
No estimate of the rate at which the two fixed
tails themselves decay is assumed.
By MA-GP-010 the same conclusion holds for the
actual finite-$S$ invariant global PW projector
$\widehat Q_{S,Y}$, with $S'=S$.

\begin{proof}
\emph{The Gram correction has uniformly bounded support and norm.}
Put $P_L=P_{[0,L]}$ on $\mathcal K_+$.
Since $V$ is causal,
$P_LV=P_LVP_L$ and $V^*P_L=P_LV^*P_L$.
The difference $R_L-1$ is supported on $[0,L]$.
Consequently
\[
       \mathfrak G_L-1=P_L(\mathfrak G_L-1)P_L,\qquad
       D_L=P_LD_LP_L,\qquad -1\le D_L\le0 .
 \tag{MA-SB.17}
\]
Indeed $\mathfrak G_L\ge1$ and it is the identity
on $P_L^\perp$; its inverse has the same block
decomposition and lies between zero and one.
In particular $\|D_L\|\le1$, independently of $L$.
This is a support statement for the inverse Gram
correction, not a uniform bound for the whole
projection difference with its right metric factor.

\emph{All cyclic products are trace class before cancellation.}
In these half-line coordinates $M_\phi C_k^+$ is
trace class.  Its kernel is the compact output
cutoff times $k(r-s)$; causality and compact
support of $k$ also confine its input variable
to a fixed compact interval.
A smooth compact-square extension followed by
bounded coordinate multipliers proves trace
class, just as in MA-GP-008.
This proof applies also to arbitrary bounded compact profiles.
Its trace is $k(0)\int\phi=0$.
The diagonal assertion follows first for smooth
profiles by the same nuclear kernel expansion
and then for bounded compact profiles by bounded smooth
approximation, strong multiplier convergence and trace-norm convergence.
Similarly $P_LC_k^+$ is trace class.
Thus $D_LC_k^+$ is trace class by (MA-SB.17).

For this fixed $L$, all metric multiplication
operators are bounded and boundedly invertible.
Conjugating the actual weighted operator by
the identity isomorphism and using the displayed
trace-class factor gives
\[
 \begin{split}
 \operatorname{Tr}_{L^2(R_Ldr)}(M_\phi C_k^+Q_L^w)
 &=\operatorname{Tr}_{\mathcal K_+}
     \bigl(\mathfrak G_L^{-1}V^*M_{R_L}M_\phi VC_k^+\bigr).
 \end{split}
 \tag{MA-SB.18}
\]
Here $C_k^+V=VC_k^+$ because their Hardy
representatives are analytic multipliers.
No Fourier cutoff was moved past convolution.

Set $F_\phi=V^*M_\phi V$.
Since $\phi=1$ on the support of $R_L-1$,
\[
         R_L\phi=\phi+(R_L-1),\qquad
         V^*M_{R_L}M_\phi V
                           =F_\phi+\mathfrak G_L-1 .
 \tag{MA-SB.19}
\]
The operators $F_\phi C_k^+$ and
$(\mathfrak G_L-1)C_k^+$ are trace class by
the preceding compact-position argument.
Substitution in (MA-SB.18), followed only by
algebra of these trace-class products, gives
the exact trace correction
\[
 \begin{split}
 &\operatorname{Tr}_{L^2(R_Ldr)}(M_\phi C_k^+Q_L^w)
       -\operatorname{Tr}_{\mathcal K_+}
                                (M_\phi C_k^+VV^*)\\
 &\qquad=\operatorname{Tr}
             \bigl(D_L(F_\phi-1)C_k^+\bigr)
       =-\operatorname{Tr}
             \bigl(D_LV^*M_{1-\phi}VC_k^+\bigr).
 \end{split}
 \tag{MA-SB.20}
\]
In particular the right metric factor has not
been discarded from an uncompressed projection.
It was accounted for by (MA-SB.19) inside a
proved trace-class cyclic product.

\emph{Uniform control of the actual correction.}
The operator in the last trace of (MA-SB.20)
is trace class, as the difference of the
two preceding trace-class products.
Its range is in $P_L\mathcal K_+$.
Trace cyclicity with the bounded projection
$P_L$ therefore permits insertion of $P_L$
on its input:
\[
 \begin{split}
 \operatorname{Tr}
    (D_LV^*M_{1-\phi}VC_k^+)
 &=\operatorname{Tr}
    (D_LV^*M_{1-\phi}VC_k^+P_L),\\
 \left|\operatorname{Tr}
    (D_LV^*M_{1-\phi}VC_k^+)\right|
 &\le\|M_{1-\phi}VC_k^+P_L\|_1
 \le\nu_k(R-L).
 \end{split}
 \tag{MA-SB.21}
\]
The last inequality is exactly the translated
Hankel rectangle of MA-SB-001.
This estimates a scalar trace through an actual
trace-class factorization.  It does not assert
that an uncompressed weighted tail operator
has the same norm bound.

For the ordinary model term, $VV^*=1-P_K$,
and its identity contribution has trace zero.
The actual fixed lift $C_k^+P_K$ is trace class.
Hence
\[
 \begin{split}
 \operatorname{Tr}(M_\phi C_k^+VV^*)
 &=-\operatorname{Tr}(M_\phi C_k^+P_K),\\
 \left|\operatorname{Tr}(M_\phi C_k^+VV^*)+\tau_+(g)\right|
 &\le\omega_k(R).
 \end{split}
 \tag{MA-SB.22}
\]
Combining (MA-SB.20)--(MA-SB.22) proves
(MA-SB.13), with fixed tails independent of $L$.

\emph{Returning to the actual symmetric cutoff.}
The actual global PW range is supported on
$x\ge-\log Y$.  The convolution by $g$, supported
in $(0,\infty)$, preserves this half-line.
The global operator is zero on its orthogonal
complement and has range in it, so its trace
equals the trace of its restriction.
The actual unitary (MA-GP.84) converts that
restriction to (MA-SB.18), with
$\phi(r)=\rho(e^{r-L})$.

Under (MA-SB.14), every $r\ge0$ has
$e^{r-L}\ge1/Y\ge2/X>1/X+\delta$.
Thus the entire lower transition is absent
on this half-line.  Moreover $\phi(r)=1$
for
$0\le r\le L+\log(X-\delta)$.
Use this value of $R$ in (MA-SB.13).
It exceeds $L$, and $X-\delta\ge X/2$.
Both trace-norm tails are nonincreasing under
further left multiplication by tail projections,
so the second bound in (MA-SB.15) follows.
The eventual hypotheses in (MA-SB.14) hold
along every stated slow-band path.
Finally MA-GP-010 identifies the actual
finite-$S$ invariant global projector with
this same $Q_Y^{\rm glob}$, in the same norm.
\end{proof}

\subsection{The remaining correction is one-sided}

\paragraph{MA-SB-003 (actual slow transport and its precise geometric correction).}
The preceding theorem also applies to the common annular
profiles of MA-LA-014:
\[
       a_X(t)=A(t/X)-A(Xt),\qquad
       h(t)=g(\log t),
 \tag{MA-SB.23}
\]
where $A=1$ on $[0,1]$, $A=0$ on $[2,\infty)$,
and $0\le A\le1$ is smooth.
Indeed, for $X\ge2Y$ the lower term vanishes
on the actual global support half-line, and
the remaining profile is one through
$r=L+\log X$.
The same proof gives
\[
 \operatorname{Tr}_H(M_{a_X}C_gQ_Y^{\rm glob})
                         \longrightarrow-\tau_+(g)
       \qquad(X,Y\to\infty,\ Y/X\to0).
 \tag{MA-SB.24}
\]
All position multipliers in this display act through
$t=e^x$.

Choose $S$ to include every prime power norm in
the compact support of $h$; equivalently choose
it large enough for the support condition on
Masullo p.~520.  For the \emph{same} annular
position profile, the actual local quotient
operator of MA-LA-014 is
\[
 T^S_{X,Y,X^{-1}}(h)
       =M_{a_X}A_hD_SQ_Y^{\rm sm}D_S^*,
 \qquad
 Q_Y^{\rm sm}=\mathcal F_e^{-1}
                   M_{B_0(\,\cdot/Y)}\mathcal F_e ,
 \tag{MA-SB.25}
\]
with the fixed smooth even frequency profile
$B_0$ used there.
MA-LA-014 proves its actual trace limit
$A_0(g)+B_S(g)$ in this regime.
Since $B_S(g)=B(g)$ for the chosen support,
Meyer's actual common-space trace identity
MA-MEY identifies its polar-corrected limit
as $-\tau_{\rm all}(g)$, where
$\tau_{\rm all}$ is the trace of the full
smooth zero quotient.
Thus the limit of the difference of actual
traces is
\[
 \begin{split}
 &\lim_{\substack{X,Y\to\infty\\Y/X\to0}}
 \left\{\operatorname{Tr}_H(M_{a_X}C_gQ_Y^{\rm glob})
      -\left(\operatorname{Tr}T^S_{X,Y,X^{-1}}(h)-P(g)\right)\right\}\\
 &\hspace{35mm}=\tau_{\rm all}(g)-\tau_+(g).
 \end{split}
 \tag{MA-SB.26}
\]
The actual operators and their limits have been
constructed before this trace difference is
evaluated.  By the spectral identifications in
MA-GP-009 and Meyer, its value is
\[
                \sum_{\operatorname{Re}\rho\ge1/2+\alpha}
                        m_\rho\,\mathcal Mg(\rho-1/2).
 \tag{MA-SB.27}
\]
The upper-edge zeros are included, since the
half-plane Hardy quotient retains only the
strict inequality.
No such exterior or boundary zero is asserted
to exist.

For comparison, the actual completed full-strip
cutoff trace of MA-WP has retained zeros in
$1/2-\alpha<\operatorname{Re}\rho<1/2+\alpha$.
The slow-band limit minus that actual full-strip
trace is
\[
             -\sum_{\operatorname{Re}\rho\le1/2-\alpha}
                       m_\rho\,\mathcal Mg(\rho-1/2).
 \tag{MA-SB.28}
\]
This includes the lower edge and has the opposite
sign to the corresponding left-subset trace.
Both formulas concern an exact remaining
correction; neither establishes its nonvanishing.

The literal Gram-normalized finite cyclic operator
in MA-GP-008 supplies the same limit (MA-SB.24):
at each fixed $X,Y$ first use its proved
$\eta\downarrow0$ trace-norm limit, then the
uniform estimate (MA-SB.15).
Equivalently, along a specified slow-band sequence
choose $\eta_{X,Y}>0$ so that this fixed-data
Gram regularization changes the trace by at
most a prescribed error tending to zero.
This uses its two actual resolvents.
It asserts neither a uniform regularizer rate
nor that either Gram resolvent may be omitted.

\paragraph{Exact scope and trace-framework placement.}
The global conclusion (MA-SB.16) includes the source's
smooth symmetric thin-boundary cutoffs.
The paired local comparison (MA-SB.26) uses the
specified common fixed annular profiles for which
MA-LA-014 proves its local operator limit.
No unproved uniform estimate for an arbitrary
varying finite-place local smoothing is imported
by identifying those two profile conventions.

The polynomial weight is $s=0$ throughout the
application.  The abstract lemma MA-SB-001 applies
to any inner multiplier with the actual ambient
lift (MA-SB.2), but it does not alone identify a
positive-polynomial-weight Gram metric.
There is no generic-edge restriction, no
uniform-interior premise, and no use of
$\mathsf A$ or $M$.
The hypotheses of the actual half-plane range
theorem, including $0<\alpha<1/2$, remain explicit.

The finite-length zero modules determine the
quotient action, while the fixed trace-class
Hankel and model-space operators determine the
uniform tail control before scalar traces are
taken.  This is an actual trace-ideal comparison
in the inherited Grothendieck framework
(cote~1 scans~079--080 and~111--112,
BR-TR-01/02).  The new factorization is proved
above and uses no new archive citation.
It constructs no arithmetic positivity form,
Tate twist, or new zero-location theorem.
The source PW definitions and coherent global
Poisson normalization are precisely those
checked in MA-GP-003/010.

\section{One reference level controls every weighted slow-band tail}
\label{sec:weighted-slow-band}

The actual nested PW ranges give a shorter uniform tail
argument than a growing-metric Gram estimate.  It works
in the single physical Hilbert space, with its specified
weight and orthogonal projections.  It therefore retains
the boundary jets that occur at positive polynomial weight.

\paragraph{MA-SB-005 (the actual slow-band limit at every polynomial weight).}
Fix the rational trivial sector,
$0<\alpha<1/2$, $s\ge0$, and
\[
 \begin{aligned}
 H&=L^2(\mathbb R,w(x)\,dx),&
 w(x)&=e^{2\alpha|x|}(1+x^2)^{s/2},\\
 P_c&=M_{\mathbf1_{[c,\infty)}},& H_c&=P_cH,\\
 J_c&=J_{e^{-c}},& Z_c&=H_c\ominus J_c.
 \end{aligned}
 \tag{MA-SB.50}
\]
Here $J_Y$ is the actual source PW theta-image closure,
as in MA-GP-002/010 and MA-GW.  Write $P_c^Z$ for
its actual complement projection, extended by zero
on $H_c^\perp$, so that
$Q_Y^{\rm glob}=P_c-P_c^Z$ when $c=-\log Y$.
Let $g\in C_c^\infty((0,\infty))$ and choose
$0<a<b$ with $\operatorname{supp}g\subset[a,b]$.
The single fixed reference operator is
\[
 K_{0,g}=C_gP_0^Z\in\mathcal S_1(H),\qquad
 \Omega_g(R)=\|P_RK_{0,g}\|_1
                       \longrightarrow0\quad(R\to\infty).
 \tag{MA-SB.51}
\]
Its trace class is MA-GW-003 at source level $Y=1$.
For every $c\le0$ and $R>b$ one has the exact identity
and uniform bound
\[
 \begin{aligned}
 P_RC_gP_c^Z&=P_RC_gP_0^ZP_0P_c^Z,\\
 \|P_RC_gP_c^Z\|_1&\le\Omega_g(R).
 \end{aligned}
 \tag{MA-SB.52}
\]
The same estimate holds with any multiplication
profile of norm at most one supported on $[R,\infty)$
in place of the leftmost projection.

Put $d=1/2-\alpha$, $u=1/2+\alpha$, and let
\[
 \begin{aligned}
 b_s&=\#\{j\in\mathbb Z_{\ge0}:2j<s-1\},\\
 r_\rho^+&=
 \begin{cases}
 m_\rho,&\operatorname{Re}\rho<u,\\
 \min(m_\rho,b_s),&\operatorname{Re}\rho=u,\\
 0,&\operatorname{Re}\rho>u,
 \end{cases}\\
 \tau_{+,s}(g)&=\sum_\rho r_\rho^+G_0(\rho-1/2),\qquad
 G_0(z)=\int g(x)e^{zx}\,dx.
 \end{aligned}
 \tag{MA-SB.53}
\]
Take a source cutoff $0\le\rho_{X,\delta}\le1$
supported on $[1/X,X]$ and equal to one on
$[1/X+\delta,X-\delta]$.
For
$Y\ge1$, $X\ge2Y$, $0<\delta\le1/(4X)$, and
$\log(X-\delta)>b$, the actual trace satisfies
\[
 \begin{split}
 &\left|\operatorname{Tr}_H
       (M_{\rho_{X,\delta}}C_gQ_Y^{\rm glob})
                     +\tau_{+,s}(g)\right|\\
 &\hspace{20mm}\le\Omega_g(\log(X-\delta))
                         \le\Omega_g(\log(X/2)).
 \end{split}
 \tag{MA-SB.54}
\]
Thus every source-thin path with $X,Y\to\infty$ and
$Y/X\to0$ has actual limit $-\tau_{+,s}(g)$.
The estimate is uniform for $1\le Y\le X/2$.
It covers every $s\ge0$, in particular the draft's
$s\ge3$, with precisely the upper-edge lengths
in (MA-SB.53).

\begin{proof}
\emph{An exact orthogonality identity in the actual norm.}
The smooth source range at $Y=1$ is contained in
the smooth source range at every $Y\ge1$.
Equivalently MA-GP-002 gives the supported-ideal
inclusion at the two endpoints.  Taking closures
in this same $H$ gives
\[
                         J_0\subset J_c\qquad(c\le0).
 \tag{MA-SB.55}
\]
No equivalence of norms depending on $c$ is used.
If $z\in Z_c$ and $j\in J_0$, then $P_0j=j$ and
\[
              \langle P_0z,j\rangle_H
                 =\langle z,j\rangle_H=0.
\]
As $P_0z\in H_0$, it follows that $P_0z\in Z_0$.
Consequently the actual orthogonal projections obey
\[
                  P_0P_c^Z=P_0^ZP_0P_c^Z,
                     \qquad\|P_0P_c^Z\|\le1.
 \tag{MA-SB.56}
\]
The final norm bound is in the same physical weighted
space on both sides.

\emph{Causality puts every far right tail through the reference level.}
Convolution $C_g$ is bounded on $H$ by the translation
estimate of MA-WP.6--.7.  If an input is supported
on $(-\infty,0)$, its output under $C_g$ is supported
on $(-\infty,b]$.  This follows directly from
$(C_gf)(x)=\int g(t)f(x-t)dt$, first on compact
smooth inputs and then by boundedness and density.
Hence $P_RC_g=P_RC_gP_0$ for $R>b$.
Insert (MA-SB.56) to obtain
\[
 \begin{split}
 P_RC_gP_c^Z
   &=P_RC_gP_0P_c^Z\\
   &=P_RC_gP_0^ZP_0P_c^Z .
 \end{split}
 \tag{MA-SB.57}
\]
The reference factor $K_{0,g}$ is trace class by
MA-GW-003.  The trace-ideal inequality and the
contraction in (MA-SB.56) give (MA-SB.52).
The profile variant factors a norm-one multiplication
operator through $P_R$.
Strong convergence $P_R\to0$ as $R\to\infty$,
followed by finite-rank approximation of $K_{0,g}$,
proves (MA-SB.51).  Monotonicity of $\Omega_g$ follows
from $P_{R'}P_R=P_{R'}$ when $R'\ge R$.

This tail factorization in fact needs only the
reference-level ambient lift.  Each full
$C_gP_c^Z$ is then trace class: its right tail
is (MA-SB.57), and its output on $[c,R]$ is
a compact-position smooth convolution kernel,
with bounded compact weight factors and support
projections.  Such a kernel is trace class by
the compact Fourier-series argument already
used in MA-GP-008.  Causality gives zero output
below $c$.  This also proves the statement
without any uniform estimate of the full
trace norm of $C_gP_c^Z$.

\emph{The actual trace and the source cutoff.}
MA-GW-002--003 identify the actual quotient and
ambient trace as $\tau_{+,s}(g)$, independent of
$c$.  The retained boundary lengths are those
in (MA-SB.53); they have not been replaced by
the ordinary $s=0$ Hardy divisor.
For every fixed compact position profile $\rho$,
the actual decomposition
$Q_Y^{\rm glob}=P_c-P_c^Z$ gives
\[
 \begin{aligned}
 \operatorname{Tr}(M_\rho C_gQ_Y^{\rm glob})
      &=-\operatorname{Tr}(M_\rho C_gP_c^Z),\\
 \operatorname{Tr}(M_\rho C_gQ_Y^{\rm glob})
                       +\tau_{+,s}(g)
      &=\operatorname{Tr}((1-M_\rho)C_gP_c^Z).
 \end{aligned}
 \tag{MA-SB.58}
\]
All products are trace class before taking traces.
The first equality uses the compact-position
kernel trace $\operatorname{Tr}(M_\rho C_gP_c)=0$
proved in MA-GW.44; its diagonal is zero since
$g(0)=0$.

Now take the source profile and $c=-\log Y$.
Convolution is causal, so $C_gP_c^Z$ has range
in $H_c$.  On that half-line,
$e^x\ge1/Y\ge2/X>1/X+\delta$.
The entire lower transition is therefore absent.
The profile is one through
$R=\log(X-\delta)$.
Consequently the multiplier
$(1-\rho(e^x))\mathbf1_{[c,\infty)}(x)$ has norm
at most one and is supported on $[R,\infty)$.
Apply the profile version of (MA-SB.52) to the
second line of (MA-SB.58).  This proves the first
bound of (MA-SB.54).  The second follows from
$X-\delta\ge X/2$ and monotonicity of $\Omega_g$.
The hypotheses are eventually satisfied on every
stated slow-band path with the source's thin
boundary condition $\delta X\to0$.
\end{proof}

\paragraph{Actual local comparison and both boundary deficiencies.}
Use the common annular profile
$a_X(t)=A(t/X)-A(Xt)$ from MA-LA-014,
with $A=1$ on $[0,1]$, zero on $[2,\infty)$,
and values in $[0,1]$.
For $X\ge2Y$ its lower term vanishes on $H_c$,
and its upper term is one through $x=\log X$.
The same proof gives error at most
$\Omega_g(\log X)$ and the same global limit.
By MA-GP-010 this is also the actual finite-$S$
invariant global projection, with $S'=S$.

Choose $S$ to contain every prime-power contribution
to the fixed $h(t)=g(\log t)$.
The actual local quotient operator
$T^S_{X,Y,X^{-1}}(h)$ of MA-LA-014, with the
same annular profile, has limit $A_0(h)+B_S(h)$.
This local operator and its Hilbert norm are the
ones already constructed there; no polynomial
global weight is inserted into its definition.
The unconditional smooth Meyer trace identity
gives $A_0+B_S-P(g)=-\tau_{\rm all}(g)$, where
$P(g)=G_0(1/2)+G_0(-1/2)$.
The two actual limits thus give
\[
 \begin{split}
 &\lim_{\substack{X,Y\to\infty\\Y/X\to0}}
 \left\{\operatorname{Tr}(M_{a_X}C_gQ_Y^{\rm glob})
       -\bigl(\operatorname{Tr}T^S_{X,Y,X^{-1}}(h)-P(g)\bigr)\right\}\\
 &\quad=
   \sum_{\operatorname{Re}\rho>u}m_\rho G_0(\rho-1/2)\\
 &\qquad+\sum_{\operatorname{Re}\rho=u}
      \bigl(m_\rho-\min(m_\rho,b_s)\bigr)G_0(\rho-1/2).
 \end{split}
 \tag{MA-SB.59}
\]
Thus at the upper edge only the unretained
length remains, rather than the full multiplicity
that occurs at $s=0$.
Relative to the actual completed full-strip
trace of MA-WP, the difference is
\[
 \begin{split}
 &\lim_{\substack{X,Y\to\infty\\Y/X\to0}}
       \operatorname{Tr}(M_{a_X}C_gQ_Y^{\rm glob})
                       -\bigl(-\tau_{{\rm strip},s}(g)\bigr)\\
 &\quad=-\sum_{\operatorname{Re}\rho<d}
                         m_\rho G_0(\rho-1/2)\\
 &\qquad-\sum_{\operatorname{Re}\rho=d}
      \bigl(m_\rho-\min(m_\rho,b_s)\bigr)G_0(\rho-1/2).
 \end{split}
 \tag{MA-SB.60}
\]
These series are absolutely convergent.
They express the differences of already constructed
operator traces.  No omitted zero or omitted edge
jet is asserted to exist.
The positive-length divisor of each nonzero deficiency
also falls under the exact separation lemma MA-JP-004.

\paragraph{Abstract character scope and the retained Gram factors.}
The proof of (MA-SB.52) is valid for any locally
bounded positive weight with locally bounded inverse,
bounded $C_g$, actual closed supported invariant subspaces
$J_c$ satisfying $J_0\subset J_c$ for $c\le0$,
and the single actual hypothesis
$C_gP_0^Z\in\mathcal S_1$.
If their actual quotient trace is independent of
$c$, the same source-cutoff limit follows.
This gives an exact criterion for the character
families of MA-CGP: their supported ranges and
nesting are proved, and a positive-weight reference
lift with its boundary multiplicities is the only
additional hypothesis needed to apply this argument.
For $s=0$ that hypothesis is already MA-CGP-004,
so MA-SB-004 is recovered directly.
An application at $s>0$ uses a proved weighted
character lift; it does not treat a Blaschke
multiplier as unitary in the stronger norm.

For every fixed $X,Y$ the global trace in this
theorem is still the actual Gram-resolvent trace
of MA-GP-008 (or MA-CGP.18 in its proved character
scope).  The reference-level factorization estimates
that actual trace; it does not discard either
normal-equation factor or replace the projection
by an additive Fourier cutoff.
Its regularizer may first be removed at fixed
data, or selected along the path to give a trace
error tending to zero, exactly as in MA-SB-003.
No uniform regularizer rate is asserted.

Only fixed-operator trace-ideal continuity and
orthogonality in the actual Hilbert metric enter
the uniform tail estimate.  The all-weight
finite-length modules and their boundary truncations
enter through MA-GW's separately proved actual
trace.  This retains the cote~1 trace-framework
placement of BR-TR-01/02 without adding an
arithmetic positivity or zero-location premise.

\section{What the actual PW tail estimate proves for every compact test}
\label{sec:all-test-pw-tail}

\paragraph{MA-AT-001 (ambient trace class and a uniform tail, without a
negative-time divisor formula).}
Fix \(0<\alpha<1/2\) and \(s\ge0\), and use the one physical Hilbert space
\[
 H=L^2\!\left(\mathbb R,w(x)\,dx\right),\qquad
 w(x)=e^{2\alpha|x|}(1+x^2)^{s/2}.
 \tag{MA-AT.1}
\]
For \(c\le0\), let \(P_c=M_{\mathbf1_{[c,\infty)}}\), let \(J_c\) be the
actual closed global PW range with support edge \(c\), and let
\[
 Q_c=P_{J_c},\qquad P_{Z_c}=P_c-Q_c,\qquad
 P_+=P_0.
 \tag{MA-AT.2}
\]
Thus \(J_c\subset P_cH\), and \(J_0\subset J_c\).
The range assertion is the actual one of MA-GP-002 and MA-GW-001,
not a replacement of a source norm by an ordinary Hardy norm.
The same lemma applies to another actual nested PW family whenever
its causal ambient lift at each fixed \(c\) has been proved.

For every \(g\in C_c^\infty(\mathbb R)\), without a sign restriction
on its support,
\[
 C_gP_{Z_c}\in\mathcal S_1(H),\qquad
 \tau_c(g):=\operatorname{Tr}_H(C_gP_{Z_c})
 \quad(c\le0).
 \tag{MA-AT.3}
\]
This defines an actual ambient trace.  For a noncausal test it is
not being identified here with a convolution action on a quotient
that only carries positive translations.

For a nonzero test put
\(a=\min\operatorname{supp}g\) and
\(b=\max\operatorname{supp}g\).
For \(R>b\) there is the exact factorization
\[
 P_{[R,\infty)}C_gP_{Z_c}
   =P_{[R,\infty)}C_gP_{Z_0}P_+P_{Z_c}.
 \tag{MA-AT.4}
\]
Consequently the fixed tail function
\[
 \omega_g(R)=\|P_{[R,\infty)}C_gP_{Z_0}\|_1
       \longrightarrow0\quad(R\longrightarrow\infty)
 \tag{MA-AT.5}
\]
satisfies
\[
 \sup_{c\le0}\|P_{[R,\infty)}C_gP_{Z_c}\|_1
       \le\omega_g(R)\qquad(R>b).
 \tag{MA-AT.6}
\]
All projections and norms in these formulas use the same \(H\).
No equivalence constant depending on \(c\) is used.

Here is the precise consequence for the actual source cutoff.
Write \(Y=e^{-c}\ge1\), let \(0\le\rho_{X,\delta}\le1\) be a smooth
function supported on \([1/X,X]\), equal to one on
\([1/X+\delta,X-\delta]\), and assume
\[
 X\ge 2\max(1,e^{-a})Y,\qquad
 0<\delta\le(4X)^{-1},\qquad \log(X-\delta)>b.
 \tag{MA-AT.7}
\]
Then every trace in
\[
 \begin{split}
 \bigg|&
 \operatorname{Tr}_H(M_{\rho_{X,\delta}(e^\cdot)}C_gQ_c)
       -g(0)\int_c^\infty\rho_{X,\delta}(e^x)\,dx
       +\tau_c(g)\bigg|\\
 &\hspace{20mm}\le
       \omega_g\bigl(\log(X-\delta)\bigr)
 \end{split}
 \tag{MA-AT.8}
\]
is the trace of an actual trace-class operator.
In particular, along every source-thin path with \(Y/X\to0\),
the left expression tends to zero.  If \(g(0)=0\), the displayed
volume counterterm is absent.  The result reduces the all-test
joint-limit problem to the separate question of the limit of
the actual numbers \(\tau_{-\log Y}(g)\).

\begin{proof}
\emph{Translation gives the ambient lift for an arbitrary test.}
Let \(U_t f(x)=f(x-t)\) on the full space \(H\).
Moderateness of \(w\) gives, for every fixed \(t\),
\[
 \|U_t\|\le 2^{s/4}e^{\alpha|t|}(1+t^2)^{s/4}.
 \tag{MA-AT.9}
\]
Indeed \(e^{2\alpha|x+t|}\le
e^{2\alpha|t|}e^{2\alpha|x|}\) and
\(1+(x+t)^2\le2(1+t^2)(1+x^2)\).
Choose \(t_0>-a\) and set \(g_{t_0}(t)=g(t-t_0)\).
Then \(g_{t_0}\in C_c^\infty((0,\infty))\) and
\[
 C_{g_{t_0}}=U_{t_0}C_g,\qquad
 C_gP_{Z_c}=U_{-t_0}C_{g_{t_0}}P_{Z_c}.
 \tag{MA-AT.10}
\]
MA-GW-003 proves the second factor on the right is trace class
on this actual \(H\), for each fixed \(c\).
The first factor is bounded.  This proves (MA-AT.3), without
commuting a translation past \(P_{Z_c}\) and without constructing
a negative-time action on \(H/J_c\).

\emph{Nesting and a physical support calculation.}
If \(v\in H\) and \(j\in J_0\), then \(P_+j=j\) and \(j\in J_c\).
Orthogonality in \(H\) therefore gives
\[
 \langle P_+P_{Z_c}v,j\rangle_H
     =\langle P_{Z_c}v,j\rangle_H=0.
 \]
The vector \(P_+P_{Z_c}v\) is supported on \([0,\infty)\),
so it belongs to \(Z_0=P_+H\ominus J_0\).  Hence
\[
 P_+P_{Z_c}=P_{Z_0}P_+P_{Z_c},\qquad
 \|P_+P_{Z_c}\|\le1.
 \tag{MA-AT.11}
\]
There is no assertion that the two projections commute.
For \(x\ge R>b\), the integral
\(C_gf(x)=\int g(t)f(x-t)\,dt\) samples only \(x-t>0\).
Consequently \(P_{[R,\infty)}C_g
 =P_{[R,\infty)}C_gP_+\), first on compactly supported
vectors and then by boundedness on \(H\).
Combining this with (MA-AT.11) proves (MA-AT.4).
Since \(C_gP_{Z_0}\) is a fixed trace-class operator,
strong convergence \(P_{[R,\infty)}\to0\) implies (MA-AT.5).
The contraction in (MA-AT.11) proves the uniform estimate.

\emph{The actual cutoff trace and its diagonal term.}
For fixed compact position support, the operator
\(M_{\rho(e^\cdot)}C_gP_c\) is trace class.
To see this in the actual weighted space, conjugate by multiplication
by \(\sqrt w\).  Compact support of \(\rho\) and \(g\) confines both
kernel variables to a compact rectangle.  The conjugated kernel is
a smooth compact convolution kernel multiplied on the two sides by
bounded compactly supported functions, together with the bounded
input projection \(P_c\).  The smooth kernel is trace class, for
example by its rapidly summable Fourier coefficients on a containing
rectangle.  Thus the nonsmooth but bounded weight factors and the
sharp support projection cause no trace-class problem.
Its trace is
\[
 \operatorname{Tr}_H(M_{\rho(e^\cdot)}C_gP_c)
       =g(0)\int_c^\infty\rho(e^x)\,dx.
 \tag{MA-AT.12}
\]
For precision, cyclically move the bounded multiplication factors
against that fixed smooth trace-class kernel, so their product on
the diagonal is \(\rho(e^x)\mathbf1_{[c,\infty)}(x)\);
the two weight factors cancel.
The usual smooth-kernel diagonal formula then extends to this
bounded compact multiplier by bounded smooth approximation,
strong multiplier convergence and trace-ideal continuity.
This proves (MA-AT.12) also at the sharp endpoint \(c\).

The identity \(Q_c=P_c-P_{Z_c}\), (MA-AT.3), and the preceding
compact-kernel argument show that the first operator in (MA-AT.8)
is trace class and that the expression before its absolute value is
\[
 \operatorname{Tr}_H
     \bigl((1-M_{\rho_{X,\delta}(e^\cdot)})C_gP_{Z_c}\bigr).
 \tag{MA-AT.13}
\]
The image of \(C_gP_{Z_c}\) is supported on \([c+a,\infty)\).
By (MA-AT.7),
\[
 1/X+\delta\le5/(4X)<e^a/Y=e^{c+a},
 \]
so the entire lower transition is below that image.
The multiplier in (MA-AT.13) can therefore be restricted to
\([R,\infty)\), where \(R=\log(X-\delta)>b\).
It has norm at most one.  Apply (MA-AT.6) and
\(|\operatorname{Tr}T|\le\|T\|_1\) to obtain (MA-AT.8).
\end{proof}

\paragraph{Exact boundary of this result.}
For a causal test, MA-GW-002--003 identify \(\tau_c(g)\) with a
fixed half-plane divisor trace, independent of \(c\), and the
all-weight causal slow-band theorem follows.
For a test with negative support, (MA-AT.3)--(MA-AT.8) establish
actual trace class and a uniform cutoff error, but do not provide
a divisor formula for \(\tau_c(g)\) or its limit as \(c\to-\infty\).
In particular, a bounded full-space translation used in
(MA-AT.10) is not an inverse for the positive semigroup on the
PW quotient.  No cyclic trace step can replace that missing
identification.
The finite-length spectral modules used in MA-GW describe
positive convolution and its trace.  This lemma is an analytic
statement about their actual ambient completions and orthogonal
projections; their class in the trace framework does not supply
the unproved negative-time formula.

\section{The negative-test PW limit from an actual adjoint defect}
\label{sec:negative-pw-trace}

\paragraph{MA-AT-002 (the negative-support limit at polynomial weight zero).}
Retain the actual rational spaces and projections of MA-AT-001,
with \(s=0\) and \(0<\alpha<1/2\).  Let \(J\) be the completed
full theta ideal, let \(P_Z=1-P_J\), and put
\[
 d=\tfrac12-\alpha,\qquad u=\tfrac12+\alpha.
 \tag{MA-AT.20}
\]
For a test \(g\in C_c^\infty((-\infty,0))\), define the two causal
tests
\[
 \begin{aligned}
 h(t)&=e^{2\alpha t}\overline{g(-t)},\\
 \ell(t)&=(e^{-2\alpha t}-e^{2\alpha t})
                         \overline{g(-t)} \qquad(t>0),
 \end{aligned}
 \tag{MA-AT.21}
\]
extended by zero elsewhere.  Write
\[
 \begin{aligned}
 \tau_c(g)&=\operatorname{Tr}_H(C_gP_{Z_c}),\\
 \tau_{\rm strip}(g)&=\operatorname{Tr}_H(C_gP_Z),\\
 \tau_+(h)&=\operatorname{Tr}_H(C_hP_{Z_c}).
 \end{aligned}
 \tag{MA-AT.22}
\]
The first trace is the actual ambient trace of MA-AT-001.
The second is the actual MA-WP trace.  The last is independent
of \(c\) by MA-GW-002--003, since \(h\) is causal.
Then
\[
 \boxed{\displaystyle
 \lim_{c\to-\infty}\tau_c(g)
   =\tau_{\rm strip}(g)
        +\overline{\tau_+(h)-\tau_{\rm strip}(h)}.}
 \tag{MA-AT.23}
\]
This equality is first proved by convergence of an actual
trace-class adjoint defect.  Its established spectral evaluations
give the equivalent formula
\[
 \begin{split}
 \lim_{c\to-\infty}\tau_c(g)
   ={}&\sum_{d<\operatorname{Re}\rho<u}
          m_\rho G_0(\rho-\tfrac12)\\
     &+\sum_{\operatorname{Re}\rho\le d}
          m_\rho G_0\bigl(-2\alpha-\overline{\rho-\tfrac12}\bigr),
 \qquad
 G_0(z)=\int g(x)e^{zx}\,dx .
 \end{split}
 \tag{MA-AT.24}
\]
All zeros here and below are nontrivial completed zeros.
The lower boundary is included in the second sum.
Both series converge absolutely.
No conjugation symmetry of a zero divisor is needed in deriving
the reflected exponent.

\begin{proof}
\emph{The exact adjoint kernel.}
The physical weight is \(w(x)=e^{2\alpha|x|}\).
The weighted adjoint of full-space convolution is
\[
 (C_g^*f)(x)
   =\int_0^\infty
       \overline{g(-t)}\,\frac{w(x-t)}{w(x)}f(x-t)\,dt .
 \tag{MA-AT.25}
\]
This follows by changing variables in the weighted inner product,
initially for compact vectors.  The resulting operator is bounded,
as is \(C_g^*\), so the identity extends to \(H\).
Set
\[
                              E=C_g^*-C_h.
 \tag{MA-AT.26}
\]
Choose \(0<a<b\) with \(\operatorname{supp}h\subset[a,b]\).
If \(x\le0\), both \(x\) and \(x-t\) are nonpositive and
\(w(x-t)/w(x)=e^{2\alpha t}\); hence \(Ef(x)=0\).
If \(x>b\), both are positive and the ratio is
\(e^{-2\alpha t}\); hence \(Ef(x)=C_\ell f(x)\).
Thus, as actual operators on \(H\),
\[
 M_{\mathbf1_{(-\infty,0]}}E=0,\qquad
 P_RE=P_RC_\ell\quad(R>b).
 \tag{MA-AT.27}
\]
These are support and kernel identities, not equalities of
scalar trace distributions.

\emph{A fixed compact part and a uniform actual tail.}
For fixed \(R>b\), the operator
\[
                         A_R=M_{\mathbf1_{[0,R]}}E
 \tag{MA-AT.28}
\]
is trace class on the full \(H\).  Its input variable lies in
the compact interval \([-b,R-a]\).
After conjugating by \(\sqrt w\), each of its two terms is a
smooth compact convolution kernel multiplied on the two sides
by bounded compact factors.  This is trace class by the compact
kernel argument of MA-AT-001.  The sharp endpoints and the cusp
of \(w\) cause no difficulty, because those factors are bounded
multipliers of an already trace-class smooth kernel.

The actual projections satisfy
\[
                   P_{Z_c}\longrightarrow P_Z
                   \quad\hbox{strongly as }c\to-\infty.
 \tag{MA-AT.29}
\]
Indeed \(P_c\to1\) strongly, while the increasing actual PW ranges
\(J_c\) exhaust \(J\) after closure, by MA-GP-001--002.
Their orthogonal projections converge strongly to \(P_J\).
This is the actual source condensation in the same physical norm.
Because all these projections are self-adjoint, right
multiplication of the fixed trace-class \(A_R\) gives
\[
             \|A_R(P_{Z_c}-P_Z)\|_1\longrightarrow0.
 \tag{MA-AT.30}
\]

The inclusion \(J_0\subset J\), exactly as in MA-AT.11, also gives
\(P_+P_Z=P_{Z_0}P_+P_Z\), with \(\|P_+P_Z\|\le1\).
Apply the tail factorization to the causal test \(\ell\), both
for \(P_{Z_c}\) and for \(P_Z\).
The one fixed operator \(C_\ell P_{Z_0}\) is trace class by
MA-GW-003.  Therefore, writing
\(\Omega_\ell(R)=\|P_RC_\ell P_{Z_0}\|_1\), one has
\[
 \begin{aligned}
 \|P_REP_{Z_c}\|_1&\le\Omega_\ell(R),\\
 \|P_REP_Z\|_1&\le\Omega_\ell(R),\qquad
                  \Omega_\ell(R)\longrightarrow0.
 \end{aligned}
 \tag{MA-AT.31}
\]
In particular \(EP_{Z_c}\) and \(EP_Z\) are trace class, since
their remaining pieces are \(A_RP_{Z_c}\) and \(A_RP_Z\).
Combining (MA-AT.27)--(MA-AT.31) proves the quantitative estimate
\[
 \|EP_{Z_c}-EP_Z\|_1
   \le\|A_R(P_{Z_c}-P_Z)\|_1+2\Omega_\ell(R)
       \longrightarrow0
 \tag{MA-AT.32}
\]
by first taking \(c\to-\infty\) for fixed \(R\), then
\(R\to\infty\).
Only the defect converges in trace norm here; no such assertion
has been made for \(C_gP_{Z_c}\).

\emph{Adjunction and trace cyclicity after trace class is known.}
MA-AT-001 proves \(C_gP_{Z_c}\in\mathcal S_1\).
Moreover
\(C_g^*P_{Z_c}=C_hP_{Z_c}+EP_{Z_c}\) is trace class by the
causal lift and (MA-AT.31).
The adjoint \(P_{Z_c}C_g^*\) is also trace class.
Cyclicity for these actual products, or their decomposition
into the two diagonal corners of \(P_{Z_c}\), now gives
\[
 \overline{\tau_c(g)}
       =\operatorname{Tr}(C_g^*P_{Z_c})
       =\tau_+(h)+\operatorname{Tr}(EP_{Z_c}).
 \tag{MA-AT.33}
\]
For the actual full-strip projection, MA-WP supplies
\(C_gP_Z,C_hP_Z\in\mathcal S_1\), and (MA-AT.31) supplies
the same defect factor.  Thus the identical justified argument gives
\[
 \overline{\tau_{\rm strip}(g)}
       =\tau_{\rm strip}(h)+\operatorname{Tr}(EP_Z).
 \tag{MA-AT.34}
\]
Subtract (MA-AT.34) from (MA-AT.33), use (MA-AT.32), and
conjugate.  This proves (MA-AT.23) entirely through the actual
operators.

\emph{Evaluation of the already proved trace difference.}
At \(s=0\), MA-GW and MA-WP show that
\(\tau_+(h)-\tau_{\rm strip}(h)\) is the trace over
\(\operatorname{Re}\rho\le d\), with full multiplicities.
For \(\lambda=\rho-\tfrac12\), a direct change of variables gives
\[
 \overline{\int_0^\infty
       e^{2\alpha t}\overline{g(-t)}e^{\lambda t}\,dt}
      =G_0(-2\alpha-\overline{\lambda}).
 \tag{MA-AT.35}
\]
The compact smooth test has rapid vertical Mellin decay,
uniformly on each fixed real strip.  The same is true after
the displayed reflection, so the known unit-height zero count
justifies conjugation and evaluation term by term.
This proves (MA-AT.24).
\end{proof}

\paragraph{The resulting actual joint limit and its local comparison.}
For the source position profiles of MA-AT-001 and every path
\(X,Y\to\infty\), \(Y/X\to0\), \(\delta X\to0\),
the threshold (MA-AT.7) is eventually satisfied.
Since \(g(0)=0\), the actual global trace consequently has limit
\[
 \lim\operatorname{Tr}(M_{\rho_{X,\delta}}C_gQ_Y^{\rm glob})
   =-\tau_{\rm strip}(g)
        -\overline{\tau_+(h)-\tau_{\rm strip}(h)}.
 \tag{MA-AT.36}
\]
This also holds for the fixed annular profile
\(a_X(t)=A(t/X)-A(Xt)\) used in MA-LA-014.
Indeed the output of \(C_gP_{Z_c}\) starts at
\(c+\min\operatorname{supp}g\).  When \(X/Y\) is sufficiently
large, \(A(Xt)\) vanishes on that output, while \(A(t/X)=1\)
through \(x=\log X\).  The same fixed upper-tail estimate applies.
The compact diagonal term for \(P_c\) remains zero.

Choose the fixed finite set \(S\) to contain every prime-power
norm or reciprocal norm met by the test.  The actual local
quotient limit and Meyer's smooth comparison give
\(\operatorname{Tr}T^S_{X,Y,X^{-1}}-P(g)\to-\tau_{\rm all}(g)\).
Subtracting this already established local limit from (MA-AT.36)
gives
\[
 \begin{split}
 &\lim\left\{\operatorname{Tr}(M_{a_X}C_gQ_Y^{\rm glob})
       -\bigl(\operatorname{Tr}T^S_{X,Y,X^{-1}}-P(g)\bigr)\right\}\\
 &\quad=
   \sum_{\operatorname{Re}\rho\ge u}
       m_\rho G_0(\rho-\tfrac12)\\
 &\qquad\quad+
   \sum_{\operatorname{Re}\rho\le d}m_\rho
       \left\{G_0(\rho-\tfrac12)
         -G_0\bigl(-2\alpha-\overline{\rho-\tfrac12}\bigr)\right\}.
 \end{split}
 \tag{MA-AT.37}
\]
At the lower boundary the reflected exponent equals the original
one, so those terms in the second line's lower-divisor correction
cancel exactly.  No exterior zero is asserted to exist.

\paragraph{Scope, characters, and the trace framework.}
The operator proof used only the physical weight, actual nested
supported ranges, their strong condensation, and their proved
causal and full-strip trace-class lifts.
It therefore applies to the actual nontrivial character families
of MA-CGP-001--004 and MA-WP as well, at \(s=0\), with
\(Y=Q\Lambda\), their own divisor, and their actual CQ local
operator; the polar term is then zero.
The conjugate in (MA-AT.35) comes from Hilbert adjunction.
It is not a replacement of a character by its contragredient,
nor does it use a real-coefficient functional equation.

Adding the causal result and this negative result gives an
actual formula for any compact smooth test whose support
avoids a neighbourhood of zero.
This statement does not yet treat every test satisfying only
\(g(0)=0\), and (MA-AT.23)--(MA-AT.37) have not here been
extended to \(s>0\).  At positive polynomial weight the defect
does not vanish on the left, so (MA-AT.32) requires an additional
uniform estimate there.
The finite-length blocks and their lengths remain the spectral
inputs already proved in MA-GW and MA-WP.  The new conclusion is
the actual operator convergence (MA-AT.32), which justifies how
those traces enter the negative-test limit; no categorical trace
class or arithmetic positivity statement replaces that analysis.

\section{Uniform weighted smoothing and the reflected negative-test limit}
\label{sec:weighted-negative-pw-trace}

\paragraph{MA-AT-003 (every polynomial weight, with the actual small left
defect).}
Fix \(0<\alpha<1/2\), \(s\ge0\), and the actual rational PW spaces
of MA-GW and MA-AT-001.  All ambient projections in the statement use
the single physical space
\[
 H_s=L^2(\mathbb R,w_s(x)\,dx),\qquad
 w_s(x)=e^{2\alpha|x|}p_s(x),\qquad
 p_s(x)=(1+x^2)^{s/2}.
 \tag{MA-AT.40}
\]
There is a finite integer \(q\), depending only on the fixed
quotient and \(s\), with the following uniform smoothing property.
Fix \(0<a<b\).  Let \(\mathfrak a(x,t)\) be measurable in \(x\),
\(C^q\) in \(t\) after extension by zero, supported in
\(a\le t\le b\), and put
\[
 \begin{aligned}
 M_q(\mathfrak a)
   &=\max_{0\le j\le q}\ \mathop{\rm ess\,sup}_{x,t}
                         |\partial_t^j\mathfrak a(x,t)|,\\
 (\mathcal K_{\mathfrak a}f)(x)
   &=\int_a^b\mathfrak a(x,t)f(x-t)\,dt .
 \end{aligned}
 \tag{MA-AT.41}
\]
Then
\[
 \sup_{c\le0}
       \|\mathcal K_{\mathfrak a}P_{Z_c}\|_1
        \le C_{\alpha,s,a,b,q}\,M_q(\mathfrak a).
 \tag{MA-AT.42}
\]
No derivative in the output variable \(x\) is required.
In particular sharp output restrictions are allowed.
This is an estimate on the actual ambient operator in \(H_s\),
not on a uniformly equivalent surrogate norm.
No preservation of the theta ideal by \(\mathcal K_{\mathfrak a}\)
is assumed.

For \(g\in C_c^\infty((-\infty,0))\), retain
\(h(t)=e^{2\alpha t}\overline{g(-t)}\) and define the actual
weighted-adjoint defect
\[
                         E_s=C_g^{*,H_s}-C_h .
 \tag{MA-AT.43}
\]
One has the actual convergence
\[
                       E_sP_{Z_c}\longrightarrow E_sP_Z
                           \quad\hbox{in }\mathcal S_1(H_s)
                              \quad(c\to-\infty).
 \tag{MA-AT.44}
\]
Consequently, with \(r_\rho^+\) and
\(r_\rho^{\rm strip}\) exactly as in MA-GW, and
\(\delta_\rho^-=r_\rho^+-r_\rho^{\rm strip}\),
\[
 \begin{split}
 \lim_{c\to-\infty}\operatorname{Tr}(C_gP_{Z_c})
   &=
     \tau_{{\rm strip},s}(g)
       +\overline{\tau_{+,s}(h)-\tau_{{\rm strip},s}(h)}\\
   &=
     \tau_{{\rm strip},s}(g)
       +\sum_\rho\delta_\rho^-\,
          G_0\bigl(-2\alpha-\overline{\rho-\tfrac12}\bigr).
 \end{split}
 \tag{MA-AT.45}
\]
Here \(G_0(z)=\int g(x)e^{zx}\,dx\).
The lower deficiency is precisely
\[
 \delta_\rho^-=
 \begin{cases}
 m_\rho,&\operatorname{Re}\rho<d,\\
 m_\rho-\min(m_\rho,b_s),&\operatorname{Re}\rho=d,\\
 0,&\operatorname{Re}\rho>d,
 \end{cases}
 \qquad
 d=\tfrac12-\alpha,\quad
 b_s=\#\{j\ge0:2j<s-1\}.
 \tag{MA-AT.46}
\]
No generic edge, distance from the zeros, or zero-location premise
is assumed.

\begin{proof}
\emph{Uniform semigroup bounds in the actual translated metric.}
Use the exact unitary \(U_c:H_c\to H_{W_c}\) of MA-GW:
\[
 \begin{aligned}
 (U_cf)(y)&=e^{\alpha(c+y)}f(c+y),\\
 W_c(y)&=e^{2\alpha|c+y|-2\alpha(c+y)}
                        (1+(c+y)^2)^{s/2},\qquad y\ge0.
 \end{aligned}
 \tag{MA-AT.47}
\]
We do not compare this norm with \(\mathcal H_s\) using a
constant depending on \(c\).
For the actual normalized right shifts
\((S_t v)(y)=\mathbf1_{y\ge t}v(y-t)\), \(t\ge0\), one has
\[
 \frac{W_c(y+t)}{W_c(y)}
 \le 2^{s/2}(1+t^2)^{s/2},\qquad
 \|S_t\|_{H_{W_c}}\le2^{s/4}(1+t^2)^{s/4}.
 \tag{MA-AT.48}
\]
The exponential factor in the ratio is at most one by
\(|x+t|-|x|\le t\); the polynomial factor satisfies
\(1+(x+t)^2\le2(1+t^2)(1+x^2)\).
These bounds are uniform for every \(c\le0\).
They pass to the actual quotient
\(\mathscr Q_c=H_{W_c}/U_cJ_c\).

\emph{A uniform trace-class power of the actual quotient resolvent.}
The actual range identity of MA-GW-001 gives
\[
 U_cJ_c=\overline{\operatorname{ran}T_\phi}^{\,H_{W_c}},
 \qquad
 \Phi(z)=\frac{(z-u)\zeta(z+d)}{(z+2)^N},\qquad
 u=\tfrac12+\alpha.
 \tag{MA-AT.49}
\]
The same fixed scalar \(\Phi\) and its same exponentially
decreasing causal kernel \(\phi\) work for every \(c\).
Only this already proved range identity is used here;
no constant from a fixed-\(c\) norm equivalence enters
the following estimates.
Let \(\mathscr A_c\) be the quotient shift generator.
The actual resolvent formulas MA-GW.14--.18 are
\[
 \begin{aligned}
 \mathscr R_c(\lambda)
   &=\int_0^\infty e^{-\lambda t}S_t^{\mathscr Q_c}\,dt
                       &&(\operatorname{Re}\lambda>0),\\
 \mathscr R_c(\lambda)
   &=T_{\lambda,c}/\Phi(-\lambda),\\
 T_{\lambda,c}
   &=\int_0^\infty\psi_\lambda(r)S_r^{\mathscr Q_c}\,dr,\qquad
 \psi_\lambda(r)=\int_r^\infty\phi(t)e^{\lambda(t-r)}\,dt .
 \end{aligned}
 \tag{MA-AT.50}
\]
The second formula is used where its denominator is nonzero
in the exponential-kernel half-plane, with the same removable
singularities as in MA-GW.  The two-sided generator identities
proved there make these formulas actual resolvents.
By (MA-AT.48), for every fixed \(\epsilon<\eta\),
\[
 \sup_{\substack{c\le0\\\operatorname{Re}\lambda\le\epsilon}}
       \|T_{\lambda,c}\|
 \le C_s\int_0^\infty |\phi(t)|e^{\max(\epsilon,0)t}
                   \int_0^t(1+r)^{s/2}\,dr\,dt<\infty .
 \tag{MA-AT.51}
\]
Thus all numerator bounds in the complete MA-GW contour proof
are uniform in \(c\).

Fix \(0<\epsilon<\eta\), \(C_0>u+2\), and use the same good
height rectangles \(\Gamma_n\) in the fixed strip
\([-C_0,\epsilon]\).  The denominator is scalar and independent
of \(c\).  The MT-M18 minimum-modulus proof, as used in
MA-GW.24--.26, therefore gives fixed constants \(C,B\) such that
\[
 \begin{aligned}
 \sup_{c\le0}\sup_{\lambda\in\Gamma_n}
             \|\mathscr R_c(\lambda)\|
       &\le C(2+|n|)^B,\\
 \sup_{c\le0}\|\mathscr R_c(\lambda)\|
       &\le C(1+|\lambda|)^{N-1}
                         &&(\operatorname{Re}\lambda\le-C_0).
 \end{aligned}
 \tag{MA-AT.52}
\]
The right vertical boundary instead uses the first integral
in (MA-AT.50), so apparent scalar poles already removed there
are not reintroduced.
The rank of each contour projection is
\(O(\log(2+|n|))\), uniformly in \(c\).
Indeed MA-GW's actual dual ODE proof gives exactly the same
interior lengths and the same allowed upper-edge jets for
every \(c\); the multiplicity count is fixed.

Choose a real \(\nu>\epsilon\) and an integer
\[
                         q>\max\{B+2,N+2,2\}.
 \tag{MA-AT.53}
\]
Repeat the proved resolvent contour calculation with
\((\nu-\lambda)^{-q}\).
Its finite-rank block trace norms are bounded by
\(C\log(2+|n|)(2+|n|)^{B-q}\).
The right-line integral is the actual
\((\nu-\mathscr A_c)^{-q}\), by the semigroup Laplace
formula and Fourier inversion.
The fixed-strip horizontal tails vanish uniformly in \(c\)
by (MA-AT.52).
The remaining integral on the line \(-M\) is bounded,
uniformly in \(c\), by \(CM^{N-q}\), which tends to zero.
This repeats the actual identification in MA-GW.32 with
uniform constants, rather than identifying a formal
residue sum.  In particular
\[
              \sup_{c\le0}\|(\nu-\mathscr A_c)^{-q}\|_1<\infty .
 \tag{MA-AT.54}
\]
The central strip is fixed throughout; no estimate on an
expanding minimum-modulus strip is being used.

\emph{The actual adjoint domain and uniform graph estimates.}
Write \(W=W_c\), \(Z_W=U_cZ_c\), and
\[
 B_cv=W^{-1}(Wv)',\qquad
 \mathcal R_c=(\nu-B_c)^{-q}P_{Z_W}.
 \tag{MA-AT.55}
\]
The adjoint positive shifts leave \(Z_W\) invariant.
Orthogonal realization of the quotient and Hilbert adjunction
turn (MA-AT.54) into an actual ambient trace-norm bound for
\(\mathcal R_c\).
The full adjoint resolvent has a uniform operator norm by
(MA-AT.48).  Expanding
\(B_c(\nu-B_c)^{-1}=\nu(\nu-B_c)^{-1}-1\) gives
\[
 \sup_{c\le0}
     \left(\|\mathcal R_c\|_1+
                 \sum_{j=0}^q\|B_c^j\mathcal R_c\|\right)<\infty .
 \tag{MA-AT.56}
\]
All powers are actual closed-generator powers.
As in the domain proof of MA-GW.38, multiplication by \(W\)
identifies this adjoint generator with left translation on
\(L^2(W^{-1}dy)\), with no endpoint boundary condition.
The direct difference quotient and its distributional converse
give
\[
 v\in\mathcal D(B_c^q)
 \ \Longleftrightarrow\
 (Wv)^{(j)}\in L^2(W^{-1}dy)\quad(0\le j\le q).
 \tag{MA-AT.57}
\]
The translation majorants in that argument are uniform by
(MA-AT.48).  In particular \(b=Wv\) belongs locally to
\(H^q\), with common derivatives through order \(q-1\)
at the cusp \(L=-c\).

On either side of \(L\), all positive-order derivatives
of \(\log W\) are bounded uniformly in \(c\).
The first is bounded by \(4\alpha+s/2\); higher derivatives
of \(\tfrac s2\log(1+x^2)\) are bounded, and the
exponential contribution has only the fixed first-derivative
jump at zero.
Consequently, for each fixed \(j\),
\[
 |\partial_y^j(W^{-1})(y)|\le C_j W(y)^{-1}
       \quad(y\ne L),\qquad
 C_T^{-1}\le\frac{W(y+t)}{W(y)}\le C_T
       \quad(|t|\le T,\ y,y+t\ge0),
 \tag{MA-AT.58}
\]
uniformly in \(c\), with the corresponding one-sided values
at \(L\).

For \(y_0=0\), or \(y_0=L\) when \(L>0\), apply the elementary
Sobolev endpoint inequality on \([y_0,y_0+1]\) to \(b=Wv\).
Using (MA-AT.58) and (MA-AT.57) gives
\[
 |b^{(j)}(y_0)|
   \le C_q W(y_0)^{1/2}
              \sum_{i=0}^q\|B_c^iv\|_{H_W}
                    \qquad(0\le j<q).
 \tag{MA-AT.59}
\]
Using the right interval at the cusp avoids a shrinking
interval when \(L\to0\).
In particular \(W(L)=1\); at the endpoint the possibly
large value \(W(0)\) is explicitly retained.

\emph{Uniform integration by parts for a variable causal kernel.}
Under \(U_c\), the kernel in (MA-AT.41) becomes
\[
 F_c(x,t)=e^{\alpha t}\mathfrak a(c+x,t),\qquad
 (\widetilde{\mathcal K}_cv)(x)
     =\int_0^\infty F_c(x,x-y)v(y)\,dy .
 \tag{MA-AT.60}
\]
Its \(t\)-derivatives through \(q\) are bounded by a fixed
constant times \(M_q(\mathfrak a)\), independently of \(c\).
For \(v\in\mathcal D(B_c^q)\), the exact pre-integration
expression is
\[
 \widetilde{\mathcal K}_c(\nu-B_c)^qv(x)
  =\int_0^\infty
        \frac{F_c(x,x-y)}{W(y)}
                   (\nu-\partial_y)^q b(y)\,dy .
 \tag{MA-AT.61}
\]
Split the input integral at \(L\) when \(L>0\) and integrate
by parts \(q\) times, just as in MA-GW.40--.43.
Only \(t\)-derivatives of \(F_c(x,t)\) occur: \(x\) is held
fixed in this input integration.
Thus measurability, rather than differentiability, in the
output variable is sufficient.

The bulk is an integral operator with kernel supported in
\(a\le x-y\le b\); its coefficients are bounded by
\(C M_q(\mathfrak a)\) after multiplication by \(W(y)\).
The uniform weight ratios in (MA-AT.48), or the bounded
translation estimate on that fixed interval, give
\[
                         \|\mathcal B_c\|
                                \le C M_q(\mathfrak a).
 \tag{MA-AT.62}
\]
For example, after conjugation by \(\sqrt W\), its absolute
kernel is bounded by
\(C M_q(\mathfrak a)\mathbf1_{a\le x-y\le b}\);
the two Schur integrals are at most a fixed multiple of
\((b-a)M_q(\mathfrak a)\).

Each remaining boundary term is an output profile times
one of the common jets \(b^{(j)}(y_0)\), \(j<q\), with
\(y_0=0\) or \(L\).
The signed coefficient formula is exactly MA-GW.43 with
\(k(x-y)\) replaced by \(F_c(x,x-y)\).
Its profile is supported in
\(y_0+a\le x\le y_0+b\).
By (MA-AT.58) its actual norm satisfies
\[
 \|\text{profile}_{c,y_0,j}\|_{H_W}
                \le C M_q(\mathfrak a)\,W(y_0)^{-1/2}.
 \tag{MA-AT.63}
\]
Indeed its pointwise bound is \(C M_q/W(y_0)\), and
the weight throughout that fixed-length output interval
is comparable to \(W(y_0)\).
At the cusp take the difference of the two one-sided
coefficients; the same bound holds.  At \(L=0\) there
is only the endpoint term.

Now substitute \(v=\mathcal R_cf\) in (MA-AT.61).
The left side is
\(\widetilde{\mathcal K}_cP_{Z_W}f\).
The bulk is \(\mathcal B_c\mathcal R_c\), with trace norm
bounded by (MA-AT.56) and (MA-AT.62).
By (MA-AT.56) and (MA-AT.59), each boundary functional
composed with \(\mathcal R_c\) has norm at most
\(C W(y_0)^{1/2}\).
This cancels precisely the factor \(W(y_0)^{-1/2}\)
in (MA-AT.63).
There are only finitely many such rank-one terms,
with a number depending on \(q\).
This proves (MA-AT.42) on \(H_W\), and hence on the
actual physical \(H_s\) by the exact unitary \(U_c\).
Causality keeps its output inside \(H_c\), so extension
by zero adds no operator or trace norm.

\emph{The uniformly small left adjoint defect.}
The actual weighted-adjoint formula is still MA-AT.25,
now with \(w=w_s\).
For \(x<0\) its difference from \(C_h\) has amplitude
\[
 h(t)\left(\frac{p_s(x-t)}{p_s(x)}-1\right),
                       \qquad a\le t\le b.
 \tag{MA-AT.64}
\]
For every fixed \(j\le q\) and \(R>2b+2\),
\[
 \sup_{\substack{x<-R\\a\le t\le b}}
 \left|\partial_t^j
        \left(\frac{p_s(x-t)}{p_s(x)}-1\right)\right|
                              \le C_{j,s,b}/R .
 \tag{MA-AT.65}
\]
For \(j=0\), integrate
\((\log p_s)'(z)=sz/(1+z^2)\) over a bounded interval.
For \(j\ge1\), use
\(|p_s^{(j)}(z)|/p_s(z)\le C_j(1+|z|)^{-j}\)
there and the bounded ratio of \(p_s(x-t)\) and \(p_s(x)\).
Leibniz' rule gives the same \(O(R^{-1})\) bound for
all \(t\)-derivatives of (MA-AT.64).
Multiply this amplitude by the sharp output indicator
\(\mathbf1_{x<-R}\) and apply (MA-AT.42).
No derivative of that indicator is taken.  We obtain
\[
          \sup_{c\le0}
             \|M_{\mathbf1_{(-\infty,-R)}}E_sP_{Z_c}\|_1
                              \le C_g/R .
 \tag{MA-AT.66}
\]
The cusp input at physical zero cannot contribute to this
left output, since the kernel is causal.
The preceding uniform proof also explicitly controls the
possible endpoint at the moving \(c\).

\emph{Reference upper tails, the full projection, and trace-norm convergence.}
The single reference operator \(E_sP_{Z_0}\) is trace class.
To verify this directly with the proved estimate, on an input
supported in \([0,\infty)\) use the globally smooth-in-\(t\)
amplitude
\[
 \widetilde{\mathfrak a}(x,t)
  =\overline{g(-t)}
       \left(e^{-2\alpha t}\frac{p_s(x-t)}{p_s(x)}
                              -e^{2\alpha t}\right).
 \tag{MA-AT.67}
\]
Its required derivatives are uniformly bounded in \(x\).
On every contributing input and output, \(x-t\ge0\)
and \(x>0\), so it agrees with the actual defect.
Hence \(\mathcal K_{\widetilde{\mathfrak a}}P_{Z_0}
=E_sP_{Z_0}\), and (MA-AT.42) at \(c=0\) applies.
Finite propagation and same-H nesting now give, for \(R>b\),
\[
 \begin{aligned}
 P_RE_sP_{Z_c}
     &=P_RE_sP_{Z_0}P_+P_{Z_c},\\
 \|P_RE_sP_{Z_c}\|_1
     &\le\Omega_E(R):=\|P_RE_sP_{Z_0}\|_1
                                  \longrightarrow0 .
 \end{aligned}
 \tag{MA-AT.68}
\]
The identical identity holds with \(P_Z\) in place of \(P_{Z_c}\).

The actual strong condensation \(P_{Z_c}\to P_Z\) holds
in \(H_s\) by MA-GP-001 and MA-GW.
For each fixed \(R\), the bounded operator
\(M_{x<-R}E_sP_{Z_c}\) converges strongly to
\(M_{x<-R}E_sP_Z\).
The uniform trace-norm bound (MA-AT.66) passes to the limit:
\[
                  \|M_{x<-R}E_sP_Z\|_1\le C_g/R .
 \tag{MA-AT.69}
\]
For completeness, this is trace-ideal lower semicontinuity,
not continuity of the trace under a bare strong limit:
pair the operators with each finite-rank test operator,
use the strong limit and the uniform trace-norm bound,
and then use the finite-rank dual characterization of
the trace norm.  It gives membership in \(\mathcal S_1\)
and the bound, but does not yet identify their traces.

The middle operator \(A_R=M_{\mathbf1_{[-R,R]}}E_s\)
is a fixed compact-kernel trace-class operator on the full
physical space.  Its input is in a fixed compact interval,
and both weighted-adjoint and ordinary kernels factor as
smooth compact convolution kernels with bounded compact
weight multipliers, as in MA-AT-002.
Thus \(\|A_R(P_{Z_c}-P_Z)\|_1\to0\).
All three pieces are now trace class, and
\[
 \|E_sP_{Z_c}-E_sP_Z\|_1
   \le \|A_R(P_{Z_c}-P_Z)\|_1
                   +2C_g/R+2\Omega_E(R).
 \tag{MA-AT.70}
\]
First take \(c\to-\infty\) for fixed \(R\), then \(R\to\infty\).
This proves the actual convergence (MA-AT.44).

\emph{The trace limit and its reflected boundary deficiency.}
MA-AT-001 supplies \(C_gP_{Z_c}\in\mathcal S_1\);
MA-GW supplies \(C_hP_{Z_c}\in\mathcal S_1\).
Equation (MA-AT.70) and its proof show that
\(E_sP_{Z_c},E_sP_Z\) are trace class.
The actual identities \(C_g^*P_{Z_c}=C_hP_{Z_c}+E_sP_{Z_c}\)
and their full-strip versions therefore make every cyclic
trace product legitimate.
Repeat MA-AT.33--.35 with these operators to obtain the
first line of (MA-AT.45), before evaluating any divisor sum.
The proven actual GW and WP multiplicities give precisely
\(\delta_\rho^-\) in (MA-AT.46).
The change of variables (MA-AT.35) is unchanged.
Rapid vertical Mellin decay makes the resulting series
absolutely convergent, proving the second line of (MA-AT.45).
\end{proof}

\paragraph{Actual joint limit and local comparison at every \(s\ge0\).}
Combine (MA-AT.45) with the actual all-test upper-tail estimate
MA-AT-001.  For strictly negative \(g\), the source-thin
slow-band limit, and the common fixed annular-profile limit, are
\[
 \lim\operatorname{Tr}(M_\rho C_gQ_Y^{\rm glob})
    =-\tau_{{\rm strip},s}(g)
       -\sum_\rho\delta_\rho^-\,
            G_0\bigl(-2\alpha-\overline{\rho-\tfrac12}\bigr).
 \tag{MA-AT.71}
\]
The profiles and parameter order are exactly those in
MA-AT-001--002; \(\delta X\to0\), \(Y/X\to0\), and the
eventual lower-support threshold are retained.
The actual local quotient still uses its own proved Hilbert
norm.  For a fixed \(S\) containing every test prime contribution,
its smooth Meyer comparison is
\(\operatorname{Tr}T^S-P(g)\to-\tau_{\rm all}(g)\).
Put \(\delta_\rho^+=m_\rho-r_\rho^+\), namely full multiplicity
strictly above \(u\), the deficiency
\(m_\rho-\min(m_\rho,b_s)\) at \(u\), and zero below.
The actual global minus polar-corrected local limit is
\[
 \begin{split}
 &\sum_\rho\delta_\rho^+G_0(\rho-\tfrac12)\\
 &\quad+\sum_\rho\delta_\rho^-
       \left\{G_0(\rho-\tfrac12)
         -G_0\bigl(-2\alpha-\overline{\rho-\tfrac12}\bigr)\right\}.
 \end{split}
 \tag{MA-AT.72}
\]
The lower-edge deficiency cancels in the second line because
reflection fixes every exponent on that edge.
No upper or lower exterior zero, and no missing edge jet,
is asserted to exist.

\paragraph{Character corollary: both requested families at every weight.}
Equations (MA-AT.42), (MA-AT.44)--(MA-AT.46) and
(MA-AT.71)--(MA-AT.72) also hold for every nonprincipal primitive
Dirichlet character, of either parity, and for the nontrivial
finite-order characters of \(\mathbb Q(\zeta_3)\) covered by
MA-GW-004, including the non-induced \(\psi_{269,3}\).
Use their own primitive divisor and multiplicities, the actual
projections \(P_{\eta,c}^Z\), and
\[
 Q=|D_K|N\mathfrak f_\eta,\qquad
 Y=Q\Lambda,\qquad c=-\log Y,\qquad
 \Phi_\eta(z)=\frac{L_K(z+d,\bar\eta)}{(z+2)^N}.
 \tag{MA-AT.73}
\]
Here the actual range identity is MA-GW.55--.60.
Its causal scalar kernel and denominator are fixed independently
of \(c\).  The uniform physical shift estimate (MA-AT.48)
therefore makes the actual numerator estimate MA-GW.63 uniform
in \(c\), exactly as in (MA-AT.51).
The fixed-character contour and finite-rank estimates
MA-GW.71--.72 have the same uniform constants after this change;
the far-left exponent is \(N\), rather than \(N-1\).
Taking \(q>\max\{B+2,N+2,2\}\), the remainder is
\(O(M^{N+1-q})\) and the same actual resolvent contour argument
gives (MA-AT.54).
All subsequent weighted domain, endpoint, cusp and variable-kernel
estimates are identical.  The strong source condensation is
MA-CGP-003, and MA-GW.48/52 give exactly the boundary deficiencies
used in the character version of (MA-AT.45).
Thus these are discharged character applications of the
actual-operator proof, not merely a proposed interface.

The intermediate scalar \(\bar\eta\) in (MA-AT.73) has already
been converted by the completed functional equation to the
\(\eta\) divisor in MA-GW-004.  The Hilbert-adjoint reflection
in (MA-AT.45) does not change that character label.
In the paired local comparison use the actual CQ operator
with its primitive conductor and different factors; its polar
term is zero.  For \(\psi_{269,3}\), \(Q=3\cdot269^2\), with
the prime over \(3\) retaining its primitive Euler factor.

\paragraph{All-test boundedness and the remaining origin scope.}
Combining (MA-AT.42) for ordinary causal convolution with the
bounded full-space translation in MA-AT.10 also proves
\(\sup_{c\le0}\|C_gP_{Z_c}\|_1<\infty\) for each
\(g\in C_c^\infty(\mathbb R)\).
On any fixed compact test support this is bounded by a fixed
finite \(C^q\) seminorm.  This uniform boundedness alone does not
identify the limit on a nonsmooth one-sided split.

The theorem proves no negative-translation group on the finite
PW quotient.  It proves the limit by actual ambient adjunction
and trace-class convergence.  It covers strictly negative tests,
and, with MA-SB-005, tests supported away from zero.
A test satisfying only \(g(0)=0\) still requires the separate
near-zero argument; a high-order kernel estimate does not
silently justify a nonsmooth one-sided split.
The finite-length \(\mathbb C[z]\)-modules and their boundary
truncations are the already proved spectral trace inputs.
Their underlying finite-dimensional traces enter only after
the actual convergence (MA-AT.70); no positivity, Tate twist,
or zero-location conclusion follows from that placement.

\section{Uniform control at logarithmic zero}
\label{sec:near-zero-trace}

The causal and anticausal test classes do not give a smooth
partition of an arbitrary test through zero.  The following
argument controls the omitted neighborhood in the actual
trace norm.  It uses the proved weighted spectral synthesis,
but needs no uniform minimum-modulus estimate in a moving norm.

\paragraph{MA-NZ-001 (a uniform actual trace-ideal estimate near zero).}
Use the rational family of MA-GW-001--003, or either character
family of MA-GW-004.  Fix $0<\alpha<1/2$, $s\ge0$ and the
single physical space
\[
 H=L^2(\mathbb R,w(x)\,dx),\qquad
 w(x)=e^{2\alpha|x|}(1+x^2)^{s/2}.
 \tag{MA-NZ.1}
\]
Let $c\le0$, let $J_c\subset H_c=\mathbf1_{[c,\infty)}H$
be the actual supported PW range, and let $P_c^Z$ project
orthogonally onto $Z_c=H_c\ominus J_c$, extended by zero.
For $g\in C_c^\infty(\mathbb R)$ with $g(0)=0$, choose
$\chi\in C_c^\infty((-2,2))$ equal to one on $[-1,1]$, and set
\[
                 r_\varepsilon(t)=g(t)\chi(t/\varepsilon).
 \tag{MA-NZ.2}
\]
Then
\[
       \lim_{\varepsilon\downarrow0}
       \sup_{c\le0}\|C_{r_\varepsilon}P_c^Z\|_1=0.
 \tag{MA-NZ.3}
\]
All the operators in this assertion act on the actual $H$.
The result holds at every polynomial weight, with the exact
upper-edge truncations of MA-GW.

\begin{proof}
\emph{1. Uniform dissipativity and an actual Hilbert--Schmidt resolvent.}
On the physical supported half-line use the positive translations
$U_tf(x)=\mathbf1_{x\ge c+t}f(x-t)$, $t\ge0$.
The logarithm of $w$ is globally Lipschitz and satisfies
\[
 (\log w)'(x)\le2\alpha+s/2\quad\text{almost everywhere}.
\]
Consequently $\|U_t\|\le e^{\omega t}$ for
$\omega=\alpha+s/4$, independently of $c$.
The actual quotient semigroup has the same bound.
Its Hilbert adjoint on $Z_c$ also has the same bound.
The full half-line adjoint generator is
\[
 \begin{aligned}
 \mathcal L_cv&=w^{-1}(wv)',\\
 \mathcal D(\mathcal L_c^2)
   &=\{v:(wv)^{(j)}\in L^2((c,\infty),w^{-1}dx),\ 0\le j\le2\}.
 \end{aligned}
 \tag{MA-NZ.4}
\]
There is no boundary condition at $c$.
This is the graph domain proved in MA-GW.37--.38 or
MA-GW.78, expressed in physical coordinates.
Let $B_c$ be its generator restricted to the adjoint-invariant
space $Z_c$.  Choose a fixed real $\nu>\omega$, put
$a_0=\nu-\omega>0$, and extend
\[
                   R_c=(\nu-B_c)^{-1}P_c^Z
 \tag{MA-NZ.5}
\]
by zero on $Z_c^\perp$.
The actual vectorwise Laplace resolvent gives
$\|R_c\|\le a_0^{-1}$.

Differentiating the quasicontractive semigroup norm at zero
on generator-domain vectors gives
$\operatorname{Re}\langle B_cv,v\rangle\le\omega\|v\|^2$.
The same holds for the adjoint generator, since its semigroup
has the same norm bound.  Applying these two inequalities
to the two resolvents proves, on $H$,
\[
 \begin{gathered}
 0\le a_0R_c^*R_c\le\operatorname{Re}R_c,
 \qquad 0\le a_0R_cR_c^*\le\operatorname{Re}R_c,\\
 \operatorname{Re}R_c=(R_c+R_c^*)/2.
 \end{gathered}
 \tag{MA-NZ.6}
\]
For example, with $v=R_cf$ in $Z_c$,
$\operatorname{Re}\langle R_cf,f\rangle
 =\nu\|v\|^2-\operatorname{Re}\langle v,B_cv\rangle
 \ge a_0\|R_cf\|^2$.

Write $\lambda_\rho=\rho-1/2$ and let $r_\rho^+$ be the
actual retained half-line lengths from MA-GW, in the indicated
character sector.  The exact kernel assertion there implies
that $Z_c$ is the closed span of the Riesz representatives
\[
  \mathbf1_{[c,\infty)}(x)\,w(x)^{-1}x^j
             e^{\overline{\lambda_\rho}x},
                     \qquad 0\le j<r_\rho^+.
 \tag{MA-NZ.7}
\]
This proves the needed completeness of the adjoint blocks;
it is not inferred merely from completeness of an unrelated
primal spectral family.
Let $F_{T,c}$ be the span of these vectors with
$|\operatorname{Im}\rho|\le T$, and let $E_{T,c}$ be its
orthogonal projection in $H$.
These are finite-dimensional invariant spaces for $B_c$
and $R_c$, have dimensions independent of $c$, and increase
densely to $Z_c$.
Their eigenvalues for $B_c$ are
$\overline{\lambda_\rho}$ with the stated lengths.

Since the trace on a finite invariant space counts those
eigenvalues, even with nontrivial Jordan blocks,
\[
 \begin{aligned}
 \operatorname{Tr}(E_{T,c}\operatorname{Re}R_cE_{T,c})
   &=\sum_{|\operatorname{Im}\rho|\le T}
       r_\rho^+\operatorname{Re}(\nu-\lambda_\rho)^{-1},\\
 S_\nu&=\sum_\rho
       r_\rho^+\operatorname{Re}(\nu-\lambda_\rho)^{-1}<\infty.
 \end{aligned}
 \tag{MA-NZ.8}
\]
Every term is nonnegative.  The containing strip bounds
each resolvent factor by a constant times
$(1+|\operatorname{Im}\rho|^2)^{-1}$, and the
multiplicity-weighted $O(\log T)$ local count proves
summability.
Choose orthonormal bases successively along these nested
finite spaces.  Positivity and (MA-NZ.8) show that
$(\operatorname{Re}R_c)^{1/2}$ is Hilbert--Schmidt;
hence $\operatorname{Re}R_c$ is trace class with trace
$S_\nu$, independent of $c$.
Equation (MA-NZ.6) now gives $R_c\in\mathcal S_2$ and
$\|R_c\|_2^2\le S_\nu/a_0$.

Let $S_\nu(T)$ be the tail of the sum in (MA-NZ.8).
Compressing both positive inequalities in (MA-NZ.6)
to $1-E_{T,c}$ gives both uniform tails
\[
 \begin{aligned}
 \|R_c(1-E_{T,c})\|_2^2&\le S_\nu(T)/a_0,\\
 \|(1-E_{T,c})R_c\|_2^2&\le S_\nu(T)/a_0,\\
 \|R_c^2-R_cE_{T,c}R_c\|_1
       &\le a_0^{-1}\sqrt{S_\nu S_\nu(T)}.
 \end{aligned}
 \tag{MA-NZ.9}
\]
In particular $R_c^2$ is trace class with a uniform
finite-block approximation.  The full half-line resolvent
preserves $Z_c$, so
$R_c^2=(\nu-\mathcal L_c)^{-2}P_c^Z$ as an actual operator.
The identities
$\mathcal L_cR_c^2=\nu R_c^2-R_c$ and
$\mathcal L_c^2R_c^2=\nu^2R_c^2-2\nu R_c+P_c^Z$
also give a uniform bound into the second graph domain.

\emph{2. Finite-block regularity without a moving-norm estimate.}
For each fixed $T$, the norm of $B_c|_{F_{T,c}}$ is
bounded independently of $c$.
Here is an elementary finite-dimensional proof.
A complex matrix is unitarily upper triangularizable:
choose a unit eigenvector, decompose its orthogonal
complement, and apply induction to the compression on
that complement.  The resulting triangular matrix has
the eigenvalues on its diagonal.
In this basis quasidissipativity says that
$\omega I-\operatorname{Re}B_c$ is positive semidefinite.
Its two-by-two principal minors therefore give, for $i<j$,
\[
 |(B_c)_{ij}|
   \le2\sqrt{(\omega-\operatorname{Re}\lambda_i)
                     (\omega-\operatorname{Re}\lambda_j)}.
 \tag{MA-NZ.10}
\]
The real parts lie in the fixed containing strip, the
heights are bounded by $T$, and the dimension is fixed.
Bounding the norm by the diagonal norm plus the
Hilbert--Schmidt norm of the strict upper triangle
gives the claimed bound $C_T$.
This uses neither a basis condition number nor zero separation.

Put $\ell=w'/w$ on the two open half-lines.  It is bounded
by $2\alpha+s/2$ in absolute value.
For a vector $v\in F_{T,c}$, the identity
$v'=B_cv-\ell v$ bounds its weighted first derivative
by $C_T\|v\|$ on $(c,\infty)$.
The one-dimensional Sobolev trace inequality on $[c,c+1]$
also gives
\[
             \sqrt{w(c)}\,|v(c+)|\le C_T\|v\|_H,
 \tag{MA-NZ.11}
\]
uniformly in $c$, because the ratio $w(c+t)/w(c)$
is uniformly bounded above and below for $0\le t\le1$.
Extend $v$ by zero to the full line.
The fundamental theorem of calculus bounds the translated
difference away from the endpoint by $C_T|t|\|v\|$.
On the endpoint interval of length $|t|$, (MA-NZ.11)
and the same derivative bound give $C_T\sqrt{|t|}\|v\|$.
For negative translations apply the bounded inverse shift
to the positive-translation estimate.
Thus, for $|t|\le1$,
\[
      \|U_tv-v\|_H\le C_T\sqrt{|t|}\,\|v\|_H
                      \qquad(v\in F_{T,c},\ c\le0).
 \tag{MA-NZ.12}
\]
Here $U_t$ is the full-line translation; its small-shift
norm is uniformly bounded by the same Lipschitz weight bound.

\emph{3. The actual second-order ambient factorization.}
For any smooth compact $r$ write
$a=w^{-1}$ and $b=wv$ on $(c,\infty)$.
For $v\in\mathcal D(\mathcal L_c^2)$, two integrations
by parts in
$\int_c^\infty r(x-y)a(y)(\nu-\partial_y)^2b(y)dy$
give a bounded bulk operator
\[
 \mathcal B_r
   =C_{r''}-2C_{r'}M_{\nu-\ell}
          +C_rM_{(\nu-\ell)^2-\ell'_{\rm reg}}.
 \tag{MA-NZ.13}
\]
The bounded regular derivative is
$\ell'_{\rm reg}(x)=s(1-x^2)/(1+x^2)^2$.
Its omitted jump is handled explicitly at the cusp.
The input to each bulk convolution is extended by zero
below $c$, whereas the output is a vector in full $H$.

The lower endpoint contribution is exactly
\[
 \begin{split}
 \mathcal E_{c,r}v(x)
  ={}&\bigl[a(c)(2\nu r(x-c)-r'(x-c))
                        +a'(c+)r(x-c)\bigr]b(c)\\
     &-a(c)r(x-c)b'(c).
 \end{split}
 \tag{MA-NZ.14}
\]
When $c<0$, splitting the integral at zero adds
$-4\alpha r(x)b(0)$, since $a'(0+)-a'(0-)=-4\alpha$.
At $c=0$ there is only the right-sided endpoint term.
The graph domain in (MA-NZ.4) makes $b,b'$ continuous
across the interior cusp; hence there are no further
terms.  The compact kernel gives zero endpoint terms
at infinity.  These calculations are valid for Sobolev
domain vectors by interval integration by parts.

Apply this identity to $v=R_c^2f$.  Its left side is
$C_rP_c^Zf$.  Thus it is the actual operator equality
\[
 \begin{split}
 C_rP_c^Z={}&\mathcal B_rR_c^2+\mathcal E_{c,r}R_c^2\\
 &-\mathbf1_{c<0}\,4\alpha\,r\otimes
                              [\,f\mapsto(wR_c^2f)(0)\,].
 \end{split}
 \tag{MA-NZ.15}
\]
In particular the last two terms are actual finite-rank maps.

The normalized endpoint functionals
$b(c)/\sqrt{w(c)}$ and $b'(c)/\sqrt{w(c)}$
are bounded uniformly by the second graph norm.
To verify this, apply the ordinary $H^2$ trace inequality
on $[c,c+1]$ and use the uniform comparability of $w$
there with $w(c)$; the three weighted integrals are
exactly the norms of $v,\mathcal L_cv,\mathcal L_c^2v$.
The functional $b(0)$ has the same bound using $[0,1]$
for every $c\le0$.
The uniform graph bound on $R_c^2$ proved after
(MA-NZ.9) therefore makes these normalized composed
functionals uniformly bounded in $c$.

\emph{4. Shrinking the actual test.}
The condition $g(0)=0$ gives, for $r=r_\varepsilon$,
\[
 \begin{gathered}
 \|r''\|_1=O(1),\qquad \int r''=0,\qquad
 \|r'\|_1=O(\varepsilon),\qquad \|r\|_1=O(\varepsilon^2),\\
 \|r'\|_2=O(\varepsilon^{1/2}),\qquad
 \|r\|_2=O(\varepsilon^{3/2}).
 \end{gathered}
 \tag{MA-NZ.16}
\]
For example $g(t)=O(|t|)$ and bounded $g',g''$ on
$|t|\le2\varepsilon$ bound separately the three
terms in $(g\chi(t/\varepsilon))''$.
The cancellation of its integral follows because $r$
is a genuine smooth compact full-line test.
These estimates and the full-line small-shift bound
show $\|\mathcal B_r\|\le C$, uniformly in $c,\varepsilon$.
On a fixed $F_{T,c}$ the principal term satisfies
\[
 C_{r''}v=\int r''(t)(U_tv-v)dt,
 \qquad
 \|\mathcal B_r|_{F_{T,c}}\|
                       \le C_T\sqrt\varepsilon.
 \tag{MA-NZ.17}
\]
Indeed (MA-NZ.12) controls that integral and the other
bulk terms have norms $O(\varepsilon)$ and
$O(\varepsilon^2)$ by (MA-NZ.16).

The range of $R_cE_{T,c}R_c$ is in $F_{T,c}$ and its
trace norm is at most $S_\nu/a_0$.
Using (MA-NZ.9) and (MA-NZ.17) gives
\[
 \sup_{c\le0}\|\mathcal B_rR_c^2\|_1
  \le C a_0^{-1}\sqrt{S_\nu S_\nu(T)}
                   +C_T\sqrt\varepsilon\,S_\nu/a_0.
 \tag{MA-NZ.18}
\]
First let $\varepsilon\downarrow0$ and then $T\to\infty$.
The bulk tends to zero uniformly.

For the moving endpoint, the weighted norms of
$r(\,\cdot-c)$ and $r'(\,\cdot-c)$ are bounded by
$C\sqrt{w(c)}\|r\|_2$ and
$C\sqrt{w(c)}\|r'\|_2$.
Since $a(c)=1/w(c)$ and
$|a'(c+)|\le C/w(c)$, the factor $w(c)$ cancels
against the normalized graph functionals just proved.
Equation (MA-NZ.14) therefore has trace norm
$O(\sqrt\varepsilon)$ after composition with $R_c^2$,
uniformly in $c$.
The interior cusp term has trace norm
$O(\varepsilon^{3/2})$, using the bounded weight near
zero and the uniform functional at zero.
Together with (MA-NZ.15) and (MA-NZ.18), this proves
(MA-NZ.3).
\end{proof}

\paragraph{MA-NZ-002 (the full $g(0)=0$ class has no unaccounted contact term).}
Write $\tau_c(g)=\operatorname{Tr}(C_gP_c^Z)$, the actual
ambient trace of MA-AT-001.
Suppose its limit as $c\to-\infty$ has been proved for
every compact smooth test vanishing on a neighborhood
of zero.  Then it has a unique limit for every
$g\in C_c^\infty(\mathbb R)$ satisfying $g(0)=0$:
\[
 \lim_{c\to-\infty}\tau_c(g)
  =\lim_{\varepsilon\downarrow0}\,\lim_{c\to-\infty}
       \tau_c\bigl(g-r_\varepsilon\bigr).
 \tag{MA-NZ.19}
\]
This extends a proved positive/negative-support formula
without a nonsmooth splitting or a false $C^\infty$
approximation claim.

\begin{proof}
MA-NZ-001 bounds
$|\tau_c(g)-\tau_c(g-r_\varepsilon)|$
by a quantity tending to zero uniformly in $c$.
For every fixed $\varepsilon$ the latter test is
smooth and vanishes near zero, so its limit exists
by hypothesis.  The same uniform estimate for two
values of $\varepsilon$ makes those limits Cauchy.
The elementary three-term estimate then proves
(MA-NZ.19) and its uniqueness.

When the formula away from zero is a finite linear
combination of one-sided Mellin sums over divisors
in fixed vertical strips with $O(\log T)$ local
count, its extension is obtained by the same
one-sided integrals of $g$.
For either half-line, two integrations by parts give
\[
 \left|\int_{\mathbb R_\pm}r_\varepsilon(t)e^{(\sigma+i\gamma)t}dt\right|
   \le C\min\{\varepsilon^2,(1+\gamma^2)^{-1}\},
 \tag{MA-NZ.20}
\]
uniformly on a fixed real-part interval.
At zero the value is zero; the first-derivative
boundary term and the second-derivative integral
are uniformly bounded by (MA-NZ.16).
Summing (MA-NZ.20), split at $|\gamma|=1/\varepsilon$,
gives $O(\varepsilon\log(2/\varepsilon))$.
Thus those actual spectral descriptions pass to
the limit whenever their away-from-zero formula
has separately been established.
\end{proof}

The theorem supplies no negative-time divisor formula
on its own.  It removes the remaining endpoint issue
after such a formula is proved for tests supported
strictly on the negative side.
On the requested space $g(0)=0$, no additional
derivative of a delta distribution can remain.
If one elects to extend a resulting functional to
all tests, two extensions agreeing on this space
differ only by a multiple of evaluation at zero:
subtract $g(0)$ times one fixed test of value one
there to see this directly.  Existence of an
unrenormalized limit for tests with $g(0)\ne0$
is not asserted.

All arithmetic inputs are the actual retained jets,
their exact kernel description and the fixed-function
local count already proved in MA-GW.  The new uniform
argument is the dissipative resolvent factorization,
finite-dimensional triangular estimate and explicit
second-order endpoint calculation above.
It uses no $\mathsf A$, $M$, zero separation,
unproved weighted inner projection or extra
arithmetic positivity principle.

\section{What operator positivity the draft actually uses}
\label{sec:positivity}

\paragraph{MA-POS-001 (source audit and scope).}
A full-text search of the public draft for positivity, positive
operators, self-adjointness, and the associated projection and
singular-value arguments, followed by reading each relevant context,
gives the following inventory. The entries distinguish positivity
on an ambient Hilbert space from positivity compatible with the
arithmetic convolution action on the full zero-jet quotient.
The latter is the specific property addressed by GRH\_NIGHT
MT-M20 and MT-M23. None of the operator-positivity uses identified
in this inventory asserts or requires that property.

\begin{longtable}{p{0.25\linewidth}p{0.38\linewidth}p{0.28\linewidth}}
\hline
Exact draft locator & Operator and use & Does it require faithful
full-jet positivity?\\
\hline
\endhead
Remark 4.3.1.6, p.~176 & The bounded nonnegative Borel argument
$\varphi(A)$ of a unitary operator $A$ in the regular representation;
orthogonal spectral projections justify positivity and self-adjointness.
& No. This is the ordinary unitary regular representation.\\[3pt]
Corollary 4.3.3.4, pp.~184--185 & The compact-isotype averaging
operator is Haar volume times the orthogonal isotypic projector.
Peter--Weyl orthogonality proves positivity.
& No. It concerns the compact factor, whose action is unitary.\\[3pt]
Theorem 9.2.1.2(3), proof p.~405 & Three nonnegative indicator
multipliers are positive idempotents; hence their ranges give the
orthogonal splitting used in the Paley--Wiener diagram.
& No. Multiplication by an indicator is positive in an ordinary
weighted $L^2$ space.\\[3pt]
Remark 9.3.3.2, pp.~427--428 & A smooth real cutoff
$0\le\rho\le1$ gives a positive self-adjoint multiplier $M_\rho$.
& No. The inner product is the explicitly weighted integral.\\[3pt]
Remark 9.3.3.4, p.~428 & The absolute indicator cutoffs are positive
orthogonal projectors.
& No. The separate false injectivity assertion audited in MT-M34/35
does not make their positivity false.\\[3pt]
Definition 10.4.1.1, p.~456; Theorem 10.4.2.2, p.~461 & The
geometric and spectral orthogonal projections are positive by
definition. Self-adjointness of $M_\rho$ is used to pass a strong
cutoff limit through a trace-class operator.
& No. Orthogonality does not make the nonunitary translation
representation a star representation.\\[3pt]
Theorem 12.1.1.2, proof pp.~499--500 & An arbitrary nuclear
operator is decomposed into a linear combination of four positive
nuclear operators. The positive case is treated by spectral
expansion, then the general case by linearity.
& No. Positivity is an auxiliary reduction, not a claim that the
arithmetic convolution operator is positive.\\[3pt]
Lemma C.3.4.4(5), p.~699 & $|P_iF|$ is positive and finite rank;
its eigenvalues bound the trace norm of the finite-rank error.
& No. This is the absolute value of an arbitrary operator.\\[3pt]
Lemma C.3.5.3, pp.~699--700; Theorem C.3.6.1, pp.~703--704 &
The positive compact operator $|A|=(A^*A)^{1/2}$ supplies
singular values, the min--max formulas, and the singular-value
expansion.
& No. $|A|$ is positive regardless of whether $A$ is normal.\\[3pt]
Corollary C.3.5.5, pp.~702--703; Proposition C.3.6.5,
pp.~705--707 & Self-adjoint cutoff approximants convert strong
convergence into convergence in a Schatten ideal when composed
with an operator already in that ideal.
& No. The trace-class premise remains a separate premise.\\[3pt]
Section C.3.7, p.~707 & A kernel criterion is stated specifically
for positive operators; the Volterra example is used to explain
why the same criterion fails for arbitrary operators.
& No. The stated restriction is not replaced by a claim of
positivity for every geometric kernel.\\
\hline
\end{longtable}

\emph{Verification and limit of the audit.}
For every entry the surrounding definition, argument, and use of
the operator were read, rather than inferred from a keyword alone.
The multiplier entries follow directly from
$\langle M_\rho f,f\rangle=\int\rho|f|^2w\ge0$.
The compact-isotype entry acts by the positive Haar volume on its
isotype and by zero on the mutually orthogonal other isotypes.
The four-positive-operators reduction on p.~499 is the usual
decomposition of the real and imaginary self-adjoint parts into
positive and negative parts; all four remain trace class. Hence
it does not transfer positivity to the original arithmetic operator.
The absolute-value entries concern $A^*A$ and its nonnegative
square root, not the eigenvalues of $A$ itself. These observations
verify the stated scope of each use. They are not an independent
audit of every analytic theorem surrounding those uses.

Several other occurrences must not be added to this list as if
they were operator positivity. Sections 1.4.6 and 1.7.2,
pp.~30 and 35, discuss Weil's positive \emph{functional} and say
that the proof does not use it. The discussion on pp.~528 and
533 refers to an unfinished positivity criterion, including
``Theorem ??'' and the empty Theorem 13.3.1.2 on p.~522.
It supplies no additional positive arithmetic operator. The
positive Radon functional on p.~506 and the positive measure and
upper density on p.~635 are measure-theoretic positivity.
Likewise, ``positive cutoff'' in Definition 9.3.3.3 refers to
one side of a norm interval; operator positivity follows instead
from the nonnegative multiplier. We found no separate assertion
elsewhere in the checked positivity contexts that the faithful
full zero-jet action is a positive star representation. The audit
does not turn the absence of such an assertion into a proof of
the missing geometric trace identity.

\subsection{The weighted adjoint is different from arithmetic inversion}

\paragraph{MA-POS-002 (a direct weighted-space obstruction).}
Let $H=L^2(\mathbb R,w(x)\,dx)$ with the weight in (MA-TR-1),
where $\alpha>0$. For $U_yf(x)=f(x-y)$,
\[
 U_y^*g(x)=\frac{w(x+y)}{w(x)}g(x+y).
 \tag{MA-POS-1}
\]
In particular the adjoint is not $U_{-y}$ for $y\ne0$.
Even a compactly supported smooth convolution square need not
give a positive operator in this weighted space. More precisely,
if $0\ne f\in C_c^\infty(\mathbb R)$,
$f^\sharp(x)=\overline{f(-x)}$, and $g=f*f^\sharp$, then
$C_g$ is bounded on $H$ but is not self-adjoint and therefore
is not positive.

\begin{proof}
Substitute $x=u+y$ in
$\int f(x-y)\overline{g(x)}w(x)\,dx$ to obtain
(MA-POS-1). Boundedness of translations follows from the weight
estimate (M26.10), and integration over the compact support of
$g$ proves boundedness of $C_g$.
Conjugate to unweighted $L^2$ by the unitary map
$Vf=\sqrt w f$. The resulting integral kernel is
\[
 K(x,z)=g(x-z)\sqrt{\frac{w(x)}{w(z)}}.
\]
Since $g(-u)=\overline{g(u)}$, the adjoint kernel is
$g(x-z)\sqrt{w(z)/w(x)}$. Now
$g(0)=\int|f|^2>0$, so continuity makes $g(y)\ne0$ for all
sufficiently small $y$. For a positive such $y$ and $z>0$,
the weight is strictly increasing and $w(z+y)\ne w(z)$.
Thus the two continuous kernels differ on a nonempty open
rectangle away from the cusp at zero. Testing there against
two compactly supported smooth functions proves that the
bounded operators differ from their adjoints. Every positive
bounded operator is self-adjoint, proving the assertion.
\end{proof}

This calculation contains no zeta zeros. It explains precisely
why positivity of the source's multipliers and orthogonal
projections cannot be promoted to arithmetic star compatibility.
It is not an additional claim that the author made that promotion.
The source itself distinguishes the two projection orders in
(10.4.1.10)--(10.4.1.11), pp.~457--458.

\paragraph{MA-POS-003 (verification of the source's test-square observation, with its finite-sum form).}
Masullo already states the single-square observation in
\S1.7.2, p.~35. It is verified here and extended to finite
sums; it is not claimed as a newly discovered single-square
obstruction.
For an idele class group $C_F$ with Haar measure and involution
$f^\sharp(g)=\overline{f(g^{-1})}$, the test space
$\mathcal D^0(C_F)$ of MT-M30 contains no nonzero convolution
square $f*f^\sharp$ with $f\in\mathcal D(C_F)$.
\emph{Proof.} The idele class group is abelian, hence unimodular,
and its identity $e$ has norm one. Directly,
\[
 (f*f^\sharp)(e)=\int_{C_F}|f(g)|^2\,d^\times g.
\]
If the convolution square belongs to $\mathcal D^0$, the left
side is zero, so $f=0$ almost everywhere, and then identically
by continuity.\hfill$\square$
The same proof applies to a finite sum of such squares: its
value at $e$ is a sum of nonnegative squared norms. Thus an
identity restricted to norm-one-vanishing tests is not by itself
a positivity test on nonzero squares. MT-M30 provides an
ideal-valued correction that recovers equality of the relevant
linear functionals; it expressly does not say that this
correction preserves positivity.

\paragraph{MA-POS-004 (individual spectral projectors do not exhaust an operator).}
There is a separate literal convergence error in the positive-case
argument on p.~500. The paragraph immediately before (12.1.1.7)
defines $P_m$ as the projector onto
$\ker(T-\sigma_m(T)\operatorname{Id})$, with distinct singular
values listed once each, and then asserts
$P_mTP_m\to T$ in $\mathcal S_p$, $p=1,2$.
The paragraph preceding (12.1.1.12) similarly asserts convergence
of these individual traces to $\operatorname{Tr}T$.
Both assertions are false; cumulative projectors give the repair.

\emph{Proof and explicit counterexample.}
On $L^2(\mathbb R)$ choose an orthonormal Hermite basis
$(e_m)_{m\ge1}$ and define $Te_m=2^{-m}e_m$.
This is a positive trace-class operator of trace one with
distinct eigenvalues. Its individual eigenspace projector gives
\[
 \begin{aligned}
 \|T-P_mTP_m\|_1&=1-2^{-m}\longrightarrow1,\\
 \|T-P_mTP_m\|_2^2&=\tfrac13-4^{-m}\longrightarrow\tfrac13,\\
 \operatorname{Tr}(P_mTP_m)&=2^{-m}\longrightarrow0.
 \end{aligned}
\]
Thus none of the asserted individual limits holds.
Instead put $E_N=\sum_{m\le N}P_m$. For an arbitrary positive
trace-class $T$, its spectral expansion gives
\[
 \|T-E_NTE_N\|_1
   =\sum_{m>N}\sigma_m(T)\operatorname{rank}P_m\longrightarrow0,
\]
and hence convergence also in $\mathcal S_2$ and of traces.
The kernels $k_m$ of (12.1.1.7) remain the individual summands;
their partial sums are the kernels for $E_NTE_N$.
This repairs precisely the two uses of exhaustion.
\hfill$\square$
The source's later displayed partial sums show the intended
repair, but do not make its individual-projector sentences true.
This defect does not contradict the valid decomposition into
four positive operators, nor refute the corrected general
trace--kernel theorem. It is an additional source check, and
uses no zero-location premise.

\subsection{Faithfulness, multiplicity, and the determinant programme}

The exact known result remains MT-M20: positivity compatible with
the full convolution star action, in a faithful Hilbert completion
of the full zero-jet quotient, is equivalent to the corresponding
GRH assertion \emph{and simplicity}. MT-M23 applies
Hu--Kaneko--Martin--Schildkraut, Theorem 1.1, to nonabelian Galois
Dedekind zeta functions and rules out that faithful positive
completion there. These are carried as checked predecessor
results, not claimed as new theorems of this audit.

The adjective ``faithful'' must specify what structure it retains.
A single operator can be positive on a vector space that also
carries a nonunitary Jordan action: the identity operator is a
trivial example. The obstruction concerns simultaneous star
compatibility with the full arithmetic action. In a local jet
block this action contains a nonzero nilpotent extension; a
positive invariant Hermitian structure would make its centered
generator skew-adjoint, and hence diagonalizable, destroying
that extension. Positivity of one unrelated operator does not
have that implication.

For the actual determinant comparison, MT-M29 already supplies
the correct distinction. On the associated graded, repeated
zero ordinates may give repeated eigenvalues of a positive
trace-class inverse-square operator, with all multiplicities
retained. That construction discards the Jordan extensions. On
the positive-frequency part its Fredholm determinant is
\[
 \det(1-wB)=\prod_{\gamma>0}
              (1-w/\gamma^2)^{m_\gamma},
\]
with the central zero removed as specified in MT-M29.
This is compatible with multiple zeros; it is not a faithful
positive realization of the original full jet module.

In the inherited Grothendieck programme the precise archive
cross-references are cote~1, scans~079--080 (the determinant
coefficients, Corollaries 4--5, and the crossed-out product/trace
argument) and scans~111--112 (the genus-one discussion), as
audited in ALPOGE\_BRIDGES BR-TR-01/02 and F-SM-11/12.
Those scans were not newly supplied or freshly inspected here.
Passing from a finite filtration to its graded object preserves
its class in $K_0$, its traces, and the corresponding determinant
coefficients; this is why a $K$-theoretic or $\lambda$-ring
identity cannot on its own certify faithfulness of the jet action.
The algebraic formalism is SGA~6, Expos\'e~V, \S\S2 and 7,
especially (7.1.1), as placed in GRH\_NIGHT AG-LAMBDA.
The analytic positivity of a Hilbert completion is additional
structure, not a consequence of that formalism or of an
unconstructed cohomological realization over
$\operatorname{Spec}\mathbb Z$.

MA-POS-001 therefore clears a specific misdiagnosis: the known
full-jet obstruction is not a blanket refutation of every
positive cutoff used by the draft. MA-POS-002 identifies the
actual distinction of adjoints, and MA-POS-003 identifies the
restriction on the test cone. Neither proves the absent
global-to-local trace transport.

\IfFileExists{modules/characters.tex}{\section{The character extension and an explicit imaginary quadratic test}
\label{sec:characters}

\subsection{Conventions and the part already supplied by F-SM-13}

This section treats task~1(d).  The primitive Dirichlet numerical
formula is ALPOGE\_BRIDGES F-SM-13, whose proof is not repeated.
The new work is the precise imaginary quadratic specialization,
an ideal-indexed prime-carrying map with its norm obstruction, and
an admissible finite-$S$ example.  None of these arguments establishes
the missing zero-location-free narrow-strip transfer in Masullo
\S12.3, p.~514, or the invocation of that transfer in the proof of
Theorem~12.5.1.1, p.~520.

For $h\in C_c^\infty(\mathbb R_{>0})$ with $h(1)=0$, write
\[
 \widehat h(z)=\int_0^\infty h(t)t^z\,\frac{dt}{t},\qquad
 H(s)=\widehat h(s-\tfrac12),\qquad
 P(h)=\widehat h(\tfrac12)+\widehat h(-\tfrac12).
 \tag{MA-CH.1}
\]
For a primitive Dirichlet character $\chi$ of conductor $q$ and
parity $\chi(-1)=(-1)^a$, F-SM-13 gives
\[
 \begin{split}
 A_a(h)&=\int_0^\infty
   \frac{e^{-(a+1/2)x}}{1-e^{-2x}}
       \bigl(h(e^x)+h(e^{-x})\bigr)\,dx,\\
 B_\chi(h)&=\sum_{n\ge1}\frac{\Lambda(n)}{\sqrt n}
       \bigl(\chi(n)h(n)+\overline\chi(n)h(n^{-1})\bigr),\\
 A_a(h)+B_\chi(h)&=\boldsymbol1_{q=1}P(h)
    -\sum_{\rho}m_\chi(\rho)\widehat h(\rho-\tfrac12).
 \end{split}
 \tag{MA-CH.2}
\]
Here and below the full sum is over the nontrivial completed zeros;
trivial finite-$L$ zeros are not counted separately.  The last
identity with the full sum is unconditional.  The use of $A$ in
F-SM-13 is exactly the replacement of that sum by the sum in
$|\operatorname{Re}\rho-1/2|<\alpha$, for every $\alpha>0$.
The conductor and parity are primitive ones: an imprimitive
Dirichlet series with deleted Euler factors is a different
function.  The conjugate character on $h(n^{-1})$ is essential.
For $q>1$ there is no polar term.  These are the same conventions
as F-SM-13, equations preceding and including
\texttt{SM-numeric-strip}, and its Mellin-inversion calculation.

To fix the adelic orientation for both fields, use GRH\_NIGHT
MT-M22 and MT-M24.  If the primitive finite Euler factors are
$(1-c_{\mathfrak p}(N\mathfrak p)^{-s})^{-1}$, let $\beta$ be
Meyer's idele-class character specified by
\[
       \beta(p_v)^{-1}=c_{\mathfrak p_v}
       \quad(v\notin S),\qquad
       \omega_s=\beta|\cdot|^s .
 \tag{MA-CH.3}
\]
The inverse in this equation is Meyer's inverse in the Euler
multiplier, not a naming convention to be suppressed.  Put
$C=C_K$, $C^1=\ker|\cdot|$, and $\xi=\beta|_{C^1}$.
Choose a norm splitting $a:\mathbb R\to C$, $|a(x)|=e^x$.
Since $\beta$ has finite order, $\beta(a(x))=1$.
With probability Haar measure $dc$ on $C^1$, the actual projector
and sector test are
\[
 \begin{aligned}
 (e_\xi f)(g)&=\int_{C^1}\xi(c)f(cg)\,dc,
       &e_\xi^2&=e_\xi,\\
 \phi_{\xi,h}(g)&=\beta(g)^{-1}h(|g|),
       &\widehat\phi_{\xi,h}(\omega_s)&=\widehat h(s).
 \end{aligned}
 \tag{MA-CH.4}
\]
The range of $e_\xi$ consists of the functions satisfying
$f(cg)=\xi(c)^{-1}f(g)$, equivalently those on which left
translation by $c$ acts by $\xi(c)$.  Substitution $u=cd$ in the
integral proves this assertion and idempotence; character
orthogonality proves the Mellin formula and the vanishing of all
other compact components.  Thus (MA-CH.4) is an actual sector
projection, not merely a spectral correspondence.  It commutes
with the adelic summation and Fourier constructions in the sense
specified by MT-M22/M24; no assertion that Fourier transform
preserves one character rather than exchanges dual sectors is
needed.

Meyer's normalization is \S3.1, author p.~16, extended to $C_K$
in \S5.1, author p.~38.  His global trace and divisor statements
are Theorems~5.8 and~5.11, author pp.~43--44 and~47.
For the actual image criterion use author pp.~46--47, and for
its transfer to the specified Masullo denominator use MT-M22
(M22.7)--(M22.10) and MT-M24 (M24.6)--(M24.9).
These are published Meyer input and the already supplied
comparison results, not consequences of $A$ or $M$.

There is a numerical normalization to retain when returning to
the draft.  Masullo Definition~A.3.2.1 and (A.3.2.1), pp.~559--560,
give compact volume
\[
 v_K=\frac{2^{r_1}(2\pi)^{r_2}h_KR_K}
                  {w_K\sqrt{|D_K|}}.
 \tag{MA-CH.5}
\]
The measure in (MA-CH.4) is his measure divided by $v_K$.
Convolution using his measure and test $\phi/v_K$ is exactly
convolution using the normalized measure and test $\phi$.
With a fixed unrescaled test the former operator, and its trace
when defined, is $v_K$ times the latter.  MT-M24 explains why
this nonzero scalar does not change the theta-image subspace;
it must nevertheless be retained in a numerical trace
comparison.  The draft's local Theorem~12.4.1.1, p.~514,
likewise divides by the norm-one local volume.  All formulas
in this section use the normalized convention.

\subsection{The imaginary quadratic explicit formula and its exact gate}

\paragraph{MA-CH-001 (the finite-order imaginary quadratic specialization).}
Let $K$ be imaginary quadratic, and let $\eta$ be a nontrivial
primitive finite-order Hecke character of conductor $\mathfrak f$.
Write $c_{\mathfrak p}=\eta(\mathfrak p)$ when $\eta$ is
unramified at $\mathfrak p$ and $c_{\mathfrak p}=0$ otherwise.
This definition uses inertia over $K$; ramification of $K/\mathbb Q$
alone does not force $c_{\mathfrak p}=0$.
Put
\[
 \begin{aligned}
 Q_\eta&=|D_K|N\mathfrak f,\\
 \Lambda_\eta(s)&=Q_\eta^{s/2}\Gamma_{\mathbb C}(s)L_K(s,\eta),\\
 \Gamma_{\mathbb C}(s)&=2(2\pi)^{-s}\Gamma(s),\\
 A_{\mathbb C}(h)&=\int_0^\infty
       \frac{e^{-x/2}}{1-e^{-x}}
          \bigl(h(e^x)+h(e^{-x})\bigr)\,dx,\\
 B_\eta(h)&=\sum_{\mathfrak p}\sum_{k\ge1}
       \frac{\log N\mathfrak p}{(N\mathfrak p)^{k/2}}
       \left(c_{\mathfrak p}^{\,k}h((N\mathfrak p)^k)
          +\overline{c_{\mathfrak p}}^{\,k}
                       h((N\mathfrak p)^{-k})\right),\\
 \tau_\eta^{\rm full}(h)&=\sum_\rho
           m_\eta(\rho)\widehat h(\rho-\tfrac12).
 \end{aligned}
 \tag{MA-CH.6}
\]
Then
\[
              A_{\mathbb C}(h)+B_\eta(h)
                         =-\tau_\eta^{\rm full}(h).
 \tag{MA-CH.7}
\]
The prime sum is finite, the integral and zero sum converge
absolutely, and this is an unconditional equality of continuous
linear functionals on the indicated test space.  For the trivial
character the same statement has $P(h)$ on the right before
$-\tau^{\rm full}$.

For the open-strip distribution
\[
 \tau_{\eta,\alpha}(h)=
     \sum_{|\operatorname{Re}\rho-1/2|<\alpha}
            m_\eta(\rho)\widehat h(\rho-\tfrac12),
 \tag{MA-CH.8}
\]
the exact discrepancy is
\[
 A_{\mathbb C}(h)+B_\eta(h)+\tau_{\eta,\alpha}(h)
       =-\!\sum_{|\operatorname{Re}\rho-1/2|\ge\alpha}
            m_\eta(\rho)\widehat h(\rho-\tfrac12).
 \tag{MA-CH.9}
\]
Consequently (MA-CH.7) with every positive narrow strip follows
from RH for this one $L_K(s,\eta)$, hence from $M$.  The
corresponding Dirichlet statement follows from $A$.  The
imaginary quadratic assertion does not follow from the
Dirichlet scope of $A$.

\emph{Proof of the new specialization.}
A finite-order continuous character is trivial on the connected
complex archimedean group; thus its archimedean factor is the
single $\Gamma_{\mathbb C}(s)$ in (MA-CH.6).  The primitive
completion is entire for nontrivial $\eta$, and its functional
equation is
$\Lambda_\eta(s)=W_\eta\Lambda_{\bar\eta}(1-s)$.
The full explicit formula, equivalently Meyer's
Theorems~5.8/5.11 on the sector (MA-CH.4), is unconditional.
To evaluate its terms with the stated normalization, the same
contour pairing as in F-SM-13 is
\[
 \frac1{2\pi i}\int_{(c)}
 \left(\frac{\Lambda_\eta'}{\Lambda_\eta}(s)H(s)
       +\frac{\Lambda_{\bar\eta}'}{\Lambda_{\bar\eta}}(s)
                                                  H(1-s)\right)ds,
 \qquad c>1.
 \tag{MA-CH.10}
\]
It equals the full zero sum, with the two polar residues
subtracted only for the trivial character.  The standard
fixed-character polynomial bounds on vertical strips and the
order-one zero count, already used in MT-M18/M22 and F-SM-13,
justify the contour pairing and absolute convergence against
the rapidly decreasing Mellin transform.  Alternatively its
convergence is the summability in Meyer Theorem~5.8; no
critical-line assertion enters either justification.

The two constant logarithmic derivatives contribute
$(\log Q_\eta-2\log(2\pi))h(1)=0$.  On $\operatorname{Re}s=c$
the Euler logarithmic derivative is absolutely convergent,
and Mellin inversion gives, for $n>1$,
\[
 \frac1{2\pi i}\int_{(c)}H(s)n^{-s}ds=n^{-1/2}h(n),\qquad
 \frac1{2\pi i}\int_{(c)}H(1-s)n^{-s}ds=n^{-1/2}h(n^{-1}).
 \tag{MA-CH.11}
\]
Hence the finite terms in (MA-CH.10) are exactly $-B_\eta(h)$,
including the conjugate coefficient on the inverse argument.
For the gamma term use
\[
 \frac{\Gamma'}{\Gamma}(s)
    =-\gamma+\int_0^\infty
                    \frac{e^{-x}-e^{-sx}}{1-e^{-x}}\,dx
       \quad(\operatorname{Re}s>0).
\]
Truncate the integral, pair it with $H(s)+H(1-s)$, and apply
(MA-CH.11).  The constants pair with $h(1)=0$, while the
remaining term is exactly $-A_{\mathbb C}(h)$.  Near zero the
numerator $h(e^x)+h(e^{-x})$ is $O(x^2)$, and for large $x$
it vanishes, so the resulting integral is absolutely
convergent.  The truncated digamma expressions have a
uniform bound by a polynomial in $|\operatorname{Im}s|$ on
the fixed integration line, which permits passage to the
limit against $H$.  This proves (MA-CH.7).

The same archimedean coefficient can be checked directly from
the complex local Weil integral.  With $z=re^{i\theta}$ and
$t=|z|_{\mathbb C}=r^2$, the canonical multiplicative measure is
$dt/t\,d\theta/(2\pi)$.  The angular average
\[
 \int_0^{2\pi}\frac{d\theta}{2\pi|1-re^{i\theta}|^2}
                  =\frac1{|1-r^2|}
 \quad(r\ne1)
\]
turns that integral into
$\int_0^\infty h(t)\sqrt t\,|1-t|^{-1}dt/t$, which is
$A_{\mathbb C}(h)$ after splitting at $1$.
Thus no extra factor of $2$ is implicit in the notation
$\Gamma_{\mathbb C}$.

Equation (MA-CH.9) is subtraction of absolutely convergent
sums.  It also specifies the exact residual gate.  For a
fixed $\alpha$, its right side vanishes for every allowed
test if and only if there is no completed zero outside the
open strip.  Indeed both full and partial zero distributions
are continuous on the exponential Schwartz space and
annihilate the actual centered theta ideal.  MT-M30 allows
the norm-one vanishing restriction without loss of the
quotient, and density gives all such smooth tests.  The
zero-jet interpolation of MT-M18 then detects any one
excluded zero; the other jets can be prescribed zero.
This argument keeps multiplicities and also detects a zero
on the strip boundary.  Accordingly, equality for all
$\alpha>0$ is precisely RH for this single character, not
an additional proof of it.  For actual Hilbert operator
traces use MT-M26/MT-M27 on wide strips, or their
zero-location-dependent narrow-strip extensions.  The
displayed partial sum alone is not a claim of trace class
for a boundary-sensitive arbitrary quotient.\hfill$\square$

\subsection{An ideal-indexed transport and its escaping correction}

\paragraph{MA-CH-002 (a twisted summation embedding and a norm obstruction).}
Let $K=\mathbb Q$ with a nonprincipal primitive Dirichlet
character, or let $K$ be imaginary quadratic with the
nontrivial character $\eta$ above.  Extend its Euler
coefficient multiplicatively to integral ideals by
$c(\mathfrak a)=0$ when $\mathfrak a$ meets the conductor.
For $K=\mathbb Q$ the ideals are positive integers and
$c(n)=\chi(n)$.  If $F\in C_c^\infty(\mathbb R_{>0})$, define
\[
 (\mathcal Z_\eta F)(t)=
       \sum_{\mathfrak a\ne0}
            \frac{c(\mathfrak a)}{\sqrt{N\mathfrak a}}
                           F((N\mathfrak a)t).
 \tag{MA-CH.12}
\]
This is an injective, norm-translation-equivariant map into
the actual centered closed theta ideal identified in
MT-M22/MT-M24.  It is continuous from each fixed compact
support test space to the exponential Schwartz space.
Its Mellin transform and inverse on its image are
\[
 \begin{split}
 \widehat{\mathcal Z_\eta F}(z)
        &=L_K(z+\tfrac12,\eta)\widehat F(z),\\
 F(t)&=\sum_{\mathfrak a\ne0}
    \frac{\mu_K(\mathfrak a)c(\mathfrak a)}{\sqrt{N\mathfrak a}}
             (\mathcal Z_\eta F)((N\mathfrak a)t).
 \end{split}
 \tag{MA-CH.13}
\]
The inverse is an ideal-indexed twisted M\"obius sum.
No surjectivity onto the whole theta ideal is asserted.

Put $(U_y f)(t)=f(e^{-y}t)$ and, for a finite set $S$ of
prime ideals, put
\[
 P_{\eta,S}=\prod_{\mathfrak p\in S}
  \left(1-c_{\mathfrak p}(N\mathfrak p)^{-1/2}
                                      U_{-\log N\mathfrak p}\right).
 \tag{MA-CH.14}
\]
For $g=\mathcal Z_\eta F$ one has the exact finite-Euler formula
\[
       P_{\eta,S}g(t)=
       \sum_{\substack{\mathfrak a\ne0\\(\mathfrak a,S)=1}}
            \frac{c(\mathfrak a)}{\sqrt{N\mathfrak a}}
                                     F((N\mathfrak a)t).
 \tag{MA-CH.15}
\]
If $\operatorname{supp}F\subset[a,b]$, the correction
$P_{\eta,S}g-F$ is supported in $(0,b/q_S]$, where $q_S$
is the least norm of an omitted prime with nonzero Euler
coefficient.  Hence it eventually vanishes on each compact
subset of $\mathbb R_{>0}$ as $S$ increases through all
unramified primes.

Nevertheless fix any nontrivial completed zero $\rho$ with
$\lambda=\rho-1/2$ and $|\operatorname{Re}\lambda|<\alpha$.
For $s\ge0$ use the same centered norm as in the main section,
\[
 \|u\|_{\alpha,s}^2=
  \int_{\mathbb R}|u(e^x)|^2e^{2\alpha|x|}(1+x^2)^{s/2}\,dx.
 \tag{MA-CH.16}
\]
There is a compactly supported $F$ with $\widehat F(\lambda)=1$,
and for that $F$ every finite $S$ satisfies
\[
 \|P_{\eta,S}\mathcal Z_\eta F-F\|_{\alpha,s}
   \ge
   \left(\int_{\mathbb R}
        e^{2\operatorname{Re}\lambda x-2\alpha|x|}
                          (1+x^2)^{-s/2}\,dx\right)^{-1/2}>0.
 \tag{MA-CH.17}
\]
This excludes convergence in this norm, despite eventual
agreement on every compact radial set.  The statement is
unconditional whenever the strip contains the chosen zero;
wide strips contain every nontrivial zero.  Applying it on
every arbitrarily narrow strip uses $A$ in the Dirichlet case,
or the single-character RH assumption in the Hecke case.

\emph{Proof.}
For each $t>0$, only the finitely many ideals with norm at
most $b/t$ occur in (MA-CH.12), locally uniformly in $t$.
Termwise differentiation is therefore valid.  Mellin
integration is absolutely convergent for
$\operatorname{Re}z>1/2$ and gives the first equation in
(MA-CH.13).  The finite Hecke $L$-function is entire in
the nontrivial cases considered here.  On each fixed
vertical strip it has at most polynomial growth; multiplication
by the compact-test Mellin transform gives rapid decay of
every desired order.  Mellin inversion may consequently
be shifted to any line $\operatorname{Re}z=-N$.
It gives $D^j\mathcal Z_\eta F(t)=O_{N,j}(t^N)$ as
$t\downarrow0$, with $D=t\,d/dt$, while the function is
identically zero for $t>b$.  The same argument bounds every
exponential Schwartz seminorm by finitely many seminorms of
$F$ on its fixed support.  Thus $g(e^x)$ belongs to the
exponential Schwartz space and depends continuously on $F$.

At every completed nontrivial zero $\rho$, the first equation
in (MA-CH.13) vanishes to at least its full multiplicity at
$z=\rho-1/2$.  MT-M22 (M22.10), followed by the centering
in MT-M24 (M24.3)/(M24.9), identifies exactly these simultaneous
jet vanishings with the actual centered theta ideal.
This proves membership using the specified adelic quotient,
without replacing its denominator by a merely formal
Euler-product ideal.

The right side of the second equation in (MA-CH.13) is again
finite for each fixed $t$: $g(u)=0$ for $u>b$.
Substituting (MA-CH.12) gives a finite double sum.  At a
nonunit ideal $\mathfrak b$ its coefficient is
\[
   \frac{c(\mathfrak b)}{\sqrt{N\mathfrak b}}
                         \sum_{\mathfrak a\mid\mathfrak b}
                                           \mu_K(\mathfrak a)=0;
\]
at the unit ideal it is $1$.  This proves the inverse,
including ideals meeting the conductor, where the extended
coefficient is zero.  In particular the map is injective.
The radial scaling in every summand proves translation
equivariance directly.

The same finite cancellation with only primes in $S$ proves
(MA-CH.15).  A nonunit ideal with nonzero coefficient and
no prime factor in $S$ has norm at least $q_S$, which proves
the support assertion.  There are only finitely many primes
of bounded norm, so $q_S\to\infty$ along any exhaustion.
Every factor in (MA-CH.14) preserves the theta ideal, either
by translation invariance or by its entire Mellin multiplier.

The functional $u\mapsto\widehat u(\lambda)$ is continuous
for the norm (MA-CH.16), with dual norm equal to the square
root of the integral in (MA-CH.17), by Cauchy--Schwarz.
Strict interiority of $\lambda$ makes that integral finite.
It vanishes on the theta ideal and its Hilbert closure.
Choose a compact smooth $b_0$ on the log line with
$\int b_0(x)dx=1$, and set
$F(e^x)=e^{-\lambda x}b_0(x)$.  Then
$\widehat F(\lambda)=1$, whereas
$\widehat{P_{\eta,S}g}(\lambda)=0$ for every $S$.
Cauchy--Schwarz gives (MA-CH.17).  No assertion about a zero
on the boundary of the strip is used.\hfill$\square$

\paragraph{What changes in the geometric intertwiner.}
For $\mathbb Q$ with a nontrivial Dirichlet character the
compact factor in (MA-CH.4), the conductor-level finite
Schwartz data, parity at the real place, and the coefficients
$\chi(n)$ must all be included.  The trivial-character
integer theta formula is therefore not an unchanged formula
for this sector.  For imaginary quadratic $K$ its exact
adelic starting point is
\[
  e_\xi\Phi_K f(g)
      =\int_{C_K^1}\xi(c)
              \sum_{u\in K^\times} f(ucg)\,dc,
 \qquad f\in\mathcal S_0(\mathbb A_K).
 \tag{MA-CH.18}
\]
At a fixed finite level its archimedean sums are sums on
rank-two lattices coming from fractional ideals, with the
ideal-class sectors and the finite congruence data retained.
A general ray character requires its conductor data even
when the ordinary ideal class group is trivial.  To pass
to an $S$-idele presentation, $S$ must include archimedean
places, the character conductor, the differential/base-field
ramification data, and enough primes to make the relevant
$S$-ideal class group trivial; see Meyer \S5.1,
author pp.~37--38, and the hypotheses of Masullo
Theorem~12.4.1.1, p.~514.

Equations (MA-CH.12)--(MA-CH.15) construct a concrete
prime-carrying radial subspace of that actual ideal, with
ideal M\"obius inversion.  They do not assert that a single
integer theta sum parametrizes all the lattices in
(MA-CH.18), or that their inverse is bounded for a narrow
Hilbert quotient.  Equation (MA-CH.17) proves that the
finite-Euler inverses cannot supply such a norm convergence
on a strip containing a zero.  The missing step is therefore
unchanged in its analytic nature: a valid global-to-local
operator trace comparison must retain the omitted ideal-prime
terms and justify the required topology.  The smooth
identification MT-M24 alone supplies neither that limit nor
zero location.

\subsection{The cubic character, its primitive local factors, and a trace gap}

\paragraph{MA-CH-003 (the primitive conductor of the concrete cubic character).}
Take the character from MA-PF-003/MA-PF-005:
\[
 K=\mathbb Q(\zeta_3)=\mathbb Q(\sqrt{-3}),\qquad
 E=K(269^{1/3}),\qquad \eta=\psi_{269,3}.
 \tag{MA-CH.19}
\]
Either faithful character of $\operatorname{Gal}(E/K)=C_3$
has primitive conductor
\[
     \mathfrak f_\eta=(269),\qquad
     Q_\eta=3\cdot269^2.
 \tag{MA-CH.20}
\]
At the unique prime over $3$ the extension $E/K$ splits
completely, so the primitive Hecke factor is
$(1-3^{-s})^{-1}$.  At the prime over $269$ the factor is
$1$, since the character is nontrivial on inertia there.
At the unique prime $(2)$, whose norm is $4$, the character
value is $1$ and its factor is $(1-4^{-s})^{-1}$.

The non-descent and Artin induction statements are the
already proved MA-PF-003/MA-PF-005:
\[
 \eta^c=\eta^{-1}\ne\eta,
 \qquad
 L_K(s,\eta)=\frac{\zeta_{\mathbb Q(269^{1/3})}(s)}
                         {\zeta_{\mathbb Q}(s)}.
 \tag{MA-CH.21}
\]
In particular this is not the norm pullback of a Dirichlet
character.  Those background identifications are not being
reproved; the local conductor and retained Euler factors
in (MA-CH.20) are the new refinement.

\emph{Proof of the local refinement.}
At $X=-1$ the polynomial $f(X)=X^3-269$ satisfies
$v_3(f(-1))=v_3(-270)=3>2v_3(f'(-1))=2$.
Thus the strong Hensel criterion gives a root in $\mathbb Q_3$.
For an elementary verification, substitute $X=-1+9t$:
\[
       f(-1+9t)=27(27t^3-9t^2+t-10).
\]
The map $t\mapsto10+9t^2-27t^3$ is a contraction of
$\mathbb Z_3$ with Lipschitz constant at most $1/9$ and
has a fixed point.  Since $K_3$ contains $\zeta_3$, it
then contains all three roots.  The prime over $3$ splits
completely in $E/K$, proving the asserted factor and zero
conductor exponent there.

Since $269\equiv2\pmod3$, it is inert and unramified in
$K/\mathbb Q$.  Over its completion the same polynomial
is Eisenstein of degree $3$.  Its extension is totally
ramified and tame, because the residue characteristic is
$269$.  A faithful rank-one character has no inertia
invariants and has zero Swan conductor, hence Artin
conductor $1-0+0=1$ at that prime.  Outside $3\cdot269$
the polynomial discriminant $-27\cdot269^2$ is a unit,
so the splitting-field extension is unramified.  There is
no archimedean finite-order conductor at the complex place.
This proves (MA-CH.20) and the factor at $269$.

Finally $2$ is inert in $K$, so its residue field is
$\mathbb F_4$.  Modulo that prime $X^3-269$ is $X^3-1$,
whose three distinct roots are exactly $\mathbb F_4^\times$.
The prime is unramified and its Frobenius is the identity
on all three roots.  Thus $c_{(2)}=1$, as asserted.
\hfill$\square$

\paragraph{MA-CH-004 (an admissible finite-$S$ trace gap for this character).}
For the character (MA-CH.19), let $S_0$ consist of the complex
place and the primes above $3$ and $269$.
There is a real $h\in C_c^\infty(\mathbb R_{>0})$, vanishing
at $1$, such that
\[
  h(4)=1,\qquad A_{\mathbb C}(h)=0,\qquad
  B_{\eta,S_0}(h)=0,\qquad B_\eta(h)=\log2.
 \tag{MA-CH.22}
\]
Here $B_{\eta,S_0}$ restricts (MA-CH.6) to the finite
places in $S_0$.  Consequently the canonical finite-$S_0$
local Weil functional is $0$, whereas the full geometric
functional is $\log2$ and the full zero trace is $-\log2$.
The same zero trace is the narrow-strip numerical trace
under RH for this one $\eta$, hence under $M$.
The discrepancy is unconditional as a comparison with the
full trace; it does not assume $M$.

\emph{Proof.}
Choose a nonnegative smooth bump $v$ supported in
$(3.9,4.1)$ with $v(4)=1$, and a nonzero nonnegative smooth
bump $w$ supported in $(1.4,1.6)$.  Both archimedean
integrals are strictly positive.  Define
\[
             h=v-\frac{A_{\mathbb C}(v)}{A_{\mathbb C}(w)}w.
 \tag{MA-CH.23}
\]
Then $h(1)=0$, $h(4)=1$ and $A_{\mathbb C}(h)=0$.
Every integral ideal norm is an integer.  The only
prime-ideal power norm in the support of $v$ is $4$,
coming from the inert prime $(2)$ with exponent one;
there is no prime ideal of norm $2$, and no prime-power
norm in the support of $w$.  Both supports lie above $1$,
so the inverse-argument terms also vanish.
By MA-CH-003 this gives
$B_\eta(h)=\log4/\sqrt4=\log2$.
The primes in $S_0$ contribute nothing: $3^k$ is outside
the supports for every $k\ge1$, and the coefficient at
$269$ is zero.  This proves (MA-CH.22).

This $S_0$ meets the local setup's actual arithmetic
hypotheses.  It contains all field and character
ramification.  Moreover $\mathbb Z[\zeta_3]$ is
norm-Euclidean: rounding the two real coordinates in the
basis $1,\zeta_3$ gives an error $a+b\zeta_3$ with
$|a|,|b|\le1/2$ and norm
$a^2-ab+b^2\le3/4<1$.  Euclidean division follows, so
the ideal class group and its $S_0$ localization are
trivial.  The character is unramified outside $S_0$.
Thus this is not a comparison with an inadmissible
archimedean-only $S$.

The canonical complex local term is the angular integral
computed in MA-CH-001.  At an unramified finite place,
Meyer's local Weil distribution expands into the two
prime-power sums in (MA-CH.6); the normalization is
Definition~5.7/Theorem~5.8, author pp.~43--44, and
\S6.1, the displayed expansion immediately after
Definition~6.1.  At a ramified character the
remaining unit contribution is supported at norm $1$,
and the conductor/discriminant terms also multiply
$h(1)$, so they vanish here.  Accordingly its finite-$S_0$
value is $A_{\mathbb C}(h)+B_{\eta,S_0}(h)=0$.
Equation (MA-CH.7) supplies the full value and its sign.
Any transfer preserving this uncorrected finite-$S_0$
functional therefore cannot identify it with the full
global functional.  A transfer carrying the omitted
ideal-prime term is not excluded.\hfill$\square$

\subsection{Artin motives and the scope of the extension}

The character $\eta$ in (MA-CH.19) is a rank-one finite-monodromy
sheaf over $K$ away from its conductor.  Pushing it from $K$
to $\mathbb Q$ gives the rank-two Artin motive
$\operatorname{Ind}_{G_K}^{G_{\mathbb Q}}\eta$ with dihedral
image $S_3$.  Its $L$-function is (MA-CH.21), including the
local inertia-invariant factors checked in MA-CH-003.
It is irreducible, since its two characters on $G_K$ are
$\eta$ and $\eta^{-1}$, and it is odd, since complex
conjugation exchanges the two induced lines and has
eigenvalues $1,-1$.  Its analytic continuation is already
that of a Hecke $L$-function.  The unresolved narrow-strip
assertion for this motive is its own RH, not Dirichlet RH.

The precise geometric dictionary is SGA~1, Expos\'e~V,
\S7, for finite \'etale covers and finite fundamental-group
sets, followed by the idempotent-completed representation
category of Artin motives in Deligne (1979), Definition~6.1
and \S\S6.2--6.3.  Deligne (0.5.3) identifies restriction of
scalars with induction; the motives here have weight zero
and Hodge type $(0,0)$.  MA-CH-003 computes their local
conductor and inertia data; it does not assign geometric
cohomology to the analytic zero jets.  MA-CH-001 is the
adelic trace formula on a fixed character component of
$C_K$.  MA-CH-002 is a statement about the topology of its
summation ideal, and MA-CH-004 exhibits the Euler data
lost by an uncorrected local trace.  In the cote~1
trace/determinant programme, scans~079--080 and~111--112,
these are respectively a character specialization, a
failure of the proposed approximation in the required norm,
and an omitted trace term; the archive locators are carried
through ALPOGE\_BRIDGES BR-TR-01/BR-TR-02 and F-SM-11/F-SM-12,
not newly certified by a scan inspection in this section.

For the example, a field of definition for the faithful
rank-one character contains $\zeta_3$, while its induced
rank-two representation has a rational model, the standard
representation of $S_3$.  Thus the failure to descend as a
rank-one norm character is not a failure of motivic descent
of its induced object.  This distinction is exactly why
the rank-one Dirichlet scope of $A$ does not cover it.

\paragraph{Input and check ledger for this section.}
MA-CH-001 is unconditional for the full distribution; its
Dirichlet narrow specialization uses $A$, and its Hecke
narrow specialization uses the RH of that one character
(sufficiently $M$).  MA-CH-002 is unconditional with its
explicit interior-zero hypothesis and uses no outside
preprint.  MA-CH-003 and MA-CH-004 are unconditional;
using their full trace as every narrow trace again needs
the single-character RH assumption.  The Kummer example
is shared with the PF26 analysis, but its field, character,
conductor and trace-gap proof do not require the PF26
factoring preprint.  The event register records the actual
source checks, full new-proof reading, and independent
local-conductor check separately.

\paragraph{References specific to this section.}
Meyer, R. (2005). On a representation of the idele class
group related to primes and zeros of $L$-functions.
\emph{Duke Mathematical Journal, 127}, 519--595.
Author version arXiv:math/0311468v3, \S3.1, pp.~16--17;
\S5.1, pp.~37--38; Theorems~5.8 and~5.11, pp.~43--44
and~47; image criterion, pp.~46--47; \S6.1, the local
expansion immediately after Definition~6.1.
\url{https://arxiv.org/abs/math/0311468}.
Deligne, P. (1979). Valeurs de fonctions $L$ et p\'eriodes
d'int\'egrales. In \emph{Automorphic forms, representations
and $L$-functions} (Proceedings of Symposia in Pure
Mathematics, Vol.~33, Part~2, pp.~313--346).
American Mathematical Society. Equation~(0.5.3),
Definition~6.1 and \S\S6.2--6.3; the Ghate--Beineke
translation has the relevant passages on PDF pp.~4 and~20.
\url{https://mathweb.tifr.res.in/~eghate/Deligne.pdf}.
}{}
\section{Prime-carrying operators in the character sectors}
\label{sec:character-commutator}

The scalar character formula of MA-CH-001 has an actual operator
counterpart.  This section extends the one-sided construction
MA-MEY-001--003 to ideal Euler factors and checks the dual character
in the second diagonal.  The construction uses absolute Euler
convergence on its one-sided completions.  It does not use the
location of any completed zero.

\paragraph{MA-CC-001 (the character Euler commutator).}
Let $K$ be a number field and let $\eta$ be a primitive finite-order
Hecke character.  Choose a finite place set $S$ containing the
archimedean places, all ramification of $\eta$, and the finitely
many places needed for Meyer's structural local description and
the fixed additive character.  In the explicit imaginary quadratic
example of MA-CH-003, one may use $S_0$ from MA-CH-004.
Write $\mathscr I^S$ for the integral ideals prime to $S$,
$c_{\mathfrak a}=\eta(\mathfrak a)$, and $\mu_K$ for the ideal
M\"obius function.  For $b>1/2$, put
\[
\begin{aligned}
 \mathcal H_b&=L^2(\mathbb R,e^{2bx}\,dx),
       &T_\ell u(x)&=u(x+\ell),\\
 Z_{\eta,S}&=\sum_{\mathfrak a\in\mathscr I^S}
       c_{\mathfrak a}(N\mathfrak a)^{-1/2}
                         T_{\log N\mathfrak a},\\
 J_{\eta,S}&=\sum_{\mathfrak a\in\mathscr I^S}
       \mu_K(\mathfrak a)c_{\mathfrak a}
       (N\mathfrak a)^{-1/2}T_{\log N\mathfrak a}.
\end{aligned}
\tag{MA-CC.1}
\]
These are mutually inverse bounded operators, with absolutely
convergent operator series.  If $g\in C_c^\infty(\mathbb R)$ and
$\rho$ is smooth with compactly supported derivative and limits
$\rho_\pm$ at the two ends, the actual operator
\[
 K_{\eta,S,\rho}(g)
   =C_g Z_{\eta,S}[J_{\eta,S},M_\rho]
\tag{MA-CC.2}
\]
is trace class, and
\[
 \operatorname{Tr}K_{\eta,S,\rho}(g)
  =(\rho_- -\rho_+)
       \sum_{\substack{\mathfrak p\notin S\\k\ge1}}
       (\log N\mathfrak p)c_{\mathfrak p}^{\,k}
       (N\mathfrak p)^{-k/2}g(k\log N\mathfrak p).
\tag{MA-CC.3}
\]
The sum is finite.  Trace-norm estimates are uniform for a bounded
test family with fixed compact support.

Fix $\phi$ with $\phi_-=1$, $\phi_+=0$ and
$\phi(x)+\phi(-x)=1$.  In the actual compact-character sector
of Meyer's common space, the centered diagonal entries of
$P_+^S-P_-$ are
\[
 D_{\eta,S}=Z_{\eta,S}[J_{\eta,S},M_\phi],
 \qquad R D_{\bar\eta,S}R,
 \qquad Ru(x)=u(-x).
\tag{MA-CC.4}
\]
In particular the second term involves the dual Euler coefficients.
With $\beta(p_v)^{-1}=c_{\mathfrak p_v}$ as in MT-M22 and
probability Haar on $C_K^1$, set
\[
\begin{aligned}
 f_g(z)&=\beta(z)^{-1}|z|^{-1/2}g(\log|z|),\\
 B_{\eta,\overline S}(g)&=
 \sum_{\substack{\mathfrak p\notin S\\k\ge1}}
 \frac{\log N\mathfrak p}{(N\mathfrak p)^{k/2}}
 \left(c_{\mathfrak p}^{\,k}g(k\log N\mathfrak p)
 +\overline{c_{\mathfrak p}}^{\,k}g(-k\log N\mathfrak p)\right).
\end{aligned}
\tag{MA-CC.5}
\]
Then the common-space defect, computed on its actual nuclear
extension, satisfies
\[
 \mathcal D_S(f_g)
 =\operatorname{Tr}\lambda(f_g)(P_- -P_+^S)
 =-B_{\eta,\overline S}(g).
\tag{MA-CC.6}
\]
Consequently the operator identity in MA-HM-002 becomes
\[
 \operatorname{Tr}_{Z}C_g^Z
 =\chi_{\pi_+}(f_g)-\mathcal A_S(f_g)
                           -B_{\eta,\overline S}(g)
\tag{MA-CC.7}
\]
whenever its exact uniform-interior hypothesis holds.  The smooth
Meyer quotient identity before transport to $Z$ is unconditional.

\begin{proof}
Unique factorization of nonzero ideals replaces integer
factorization in MA-MEY-001.  Since $|c_{\mathfrak a}|=1$ outside
$S$ and $\|T_{\log N\mathfrak a}\|_{\mathcal H_b}
=(N\mathfrak a)^{-b}$, both norm series are dominated by
$\zeta_K(b+1/2)$.  Their product is absolutely convergent, and
the coefficient of the shift indexed by $\mathfrak c$ is
\[
 c_{\mathfrak c}(N\mathfrak c)^{-1/2}
                \sum_{\mathfrak d\mid\mathfrak c}\mu_K(\mathfrak d).
\]
This is zero unless $\mathfrak c=\mathcal O_K$, where it is one.
Distinct ideals may have the same norm; absolute convergence
allows this grouping by ideals before combining equal shifts.
Thus both operator products are the identity.

The compact-kernel estimate (MA-MEY.5), whose full proof is already
given, applies to each real shift.  It bounds the trace norm of
the term indexed by $(\mathfrak a,\mathfrak d)$ in (MA-CC.2) by
a constant times
\[
 (N\mathfrak a)^{-b-1/2}(N\mathfrak d)^{-b-1/2}
                                      (1+\log N\mathfrak d).
\]
The sum is finite.  Indeed for $\sigma=b+1/2>1$, choose
$0<\varepsilon<\sigma-1$ and use
$1+\log t\ll_\varepsilon t^\varepsilon$ together with the
absolute convergence of $\zeta_K(\sigma-\varepsilon)$.
The dependence on $g$ is through finitely many seminorms on its
fixed support, exactly as in (MA-MEY.5)--(MA-MEY.7).

Put $v=\log N\mathfrak a$ and $w=\log N\mathfrak d$.
The kernel of $C_gT_v[T_w,M_\rho]$ is
$g(x+v+w-z)(\rho(z)-\rho(z-w))$.  The weight conjugation is one
on the diagonal, so its trace is
$g(v+w)w(\rho_+-\rho_-)$ by (MA-MEY.9).
The absolute trace-norm convergence just proved permits summation
and regrouping.  The resulting coefficient at $\mathfrak c$ is
\[
 (\rho_+-\rho_-)
 c_{\mathfrak c}(N\mathfrak c)^{-1/2}g(\log N\mathfrak c)
 \sum_{\mathfrak d\mid\mathfrak c}
                          \mu_K(\mathfrak d)\log N\mathfrak d.
\]
The last sum is $-\Lambda_K(\mathfrak c)$, where
$\Lambda_K(\mathfrak p^k)=\log N\mathfrak p$ and it is zero
otherwise.  The proof is the same finite divisor expansion as
for integers: the coefficient of any one $\log N\mathfrak p$
is canceled by $\prod_{\mathfrak q\mid\mathfrak c,
\mathfrak q\ne\mathfrak p}(1-1)$ unless there is only one prime
divisor.  This proves (MA-CC.3).  There are only finitely many
integral ideals of bounded norm, proving finiteness of the trace sum.

It remains to check that these operators are the diagonal entries
of the actual adelic construction, including the dual orientation.
Write a raw vector in the $\beta$ sector as
\[
 v(ca(x))=\beta(c)^{-1}e^{-x/2}u(x),
 \qquad c\in C_K^1,\quad |a(x)|=e^x.
\tag{MA-CC.8}
\]
A finite-order character is trivial on this connected norm
splitting.  The raw translation $\lambda_{p_v}^{-1}$ becomes
$c_{\mathfrak p_v}(N\mathfrak p_v)^{-1/2}
T_{\log N\mathfrak p_v}$.  This follows either directly from
(MA-CC.8) or from (M22.4) with
$\beta(p_v)^{-1}=c_{\mathfrak p_v}$.
Thus the restriction of Meyer's $\mathcal L_S$ is exactly
$Z_{\eta,S}$, and its inverse is $J_{\eta,S}$.

The raw inversion $J_0v(z)=|z|^{-1}v(z^{-1})$ sends
(MA-CC.8) to
\[
 (J_0v)(ca(x))=\beta(c)e^{-x/2}u(-x).
\tag{MA-CC.9}
\]
It therefore exchanges the $\beta$ and $\beta^{-1}$ sectors and
becomes the linear reflection $R$ from $\mathcal H_{-b}$ to
$\mathcal H_b$ (or in the reverse direction).  It is not complex
conjugation of the vector.  The middle Euler multiplier, now in
the inverse character sector, has coefficients
$\overline{c_{\mathfrak p}}$ by unitarity.  Substitution in the
two explicit projection matrices (MA-MEY.13)--(MA-MEY.14)
proves (MA-CC.4), using $RM_{1-\phi}R=M_\phi$.

Meyer's Lemma~5.6 gives the nuclear extension of the complete
smoothed defect, including its off-diagonal entries, from the
paired weighted Hilbert spaces into the common smooth space.
Restriction to the fixed compact character and composition with
the continuous Hilbert inclusion give a trace-class operator.
The same nuclear rank-one representation computes its trace in
both spaces: each pairing is unchanged by that composition and
the scalar series is absolutely convergent.  Hence the trace is
the sum of the two Hilbert diagonal traces.  Since
$RC_gR=C_{g^\vee}$, (MA-CC.3) yields
\[
 \operatorname{Tr}\lambda(f_g)(P_+^S-P_-)
 =\operatorname{Tr}_{\mathcal H_b}C_gD_{\eta,S}
  +\operatorname{Tr}_{\mathcal H_b}C_{g^\vee}D_{\bar\eta,S}
 =B_{\eta,\overline S}(g).
\]
This proves (MA-CC.6).  Meyer's actual trace additivity
(MA-HM.23), and then the actual nuclear comparison (MA-HM.16),
give (MA-CC.7).  The first step uses no zero location; the
second uses precisely (MA-HM.2).
\end{proof}

\paragraph{MA-CC-002 (the non-Dirichlet cubic gap is an operator trace).}
For $K=\mathbb Q(\zeta_3)$, the character $\psi$ of
$K(269^{1/3})/K$ and the admissible set $S_0$ of MA-CH-003--004,
let $h$ be the test constructed in MA-CH-004 and put
$g(x)=h(e^x)$.  Then the actual one-sided prime-carrying operator
satisfies
\[
 \operatorname{Tr}K_{\psi,S_0,\phi}(g)=\log2,
 \qquad
 \operatorname{Tr}K_{\bar\psi,S_0,\phi}(g^\vee)=0,
 \qquad \mathcal D_{S_0}(f_g)=-\log2.
\tag{MA-CC.10}
\]
The genuine smooth quotient trace is consequently $-\log2$.
Under the same Hilbert interior-strip premise, the actual
orthogonal geometric cutoff limit is $+\log2$.

\begin{proof}
The previously checked construction has support strictly above
one, $h(4)=1$, and no other prime-ideal-power norm in its support.
The only contributing ideal is the unique prime above $2$,
of norm $4$, with $c_{\mathfrak p}=1$.
Its term in (MA-CC.3) is
$(\log4)/\sqrt4=\log2$; the reciprocal test has no positive
norm contribution.  This proves the three operator equalities.
MA-CH-004 also gives $A_{\mathbb C}(h)=0$ and $B_{\psi,S_0}(h)=0$.
The local Fourier-cutoff trace $\mathcal A_{S_0}(f_g)$ is
therefore zero by the local calculation, and the polar quotient
vanishes in this nontrivial character sector.  Trace additivity
gives the smooth quotient trace without comparing zero sums.
MA-HM-002 and MA-LIFT-003 give the last conclusion, with their
stated premise retained.  On wide strips it is unconditional;
on every positive narrow strip, RH for this particular character
is sufficient and is implied by $M$, not by the Dirichlet scope
of $A$.
\end{proof}

For primitive Dirichlet characters the same formulas use integer
ideals and the conjugate values $\overline\chi(p)^k$ on the
reciprocal terms.  The local real factor retains the parity in
(MA-CH.2).  The changes for an imaginary quadratic field are the
ideal norm, the complex local factor and conductor, and the dual
finite character in the reflected Euler multiplier.  The
operator proof itself needs no positivity and no simplicity of
zeros.  A compact cutoff has equal endpoint limits and hence
zero commutator trace in (MA-CC.3), so MA-MEY-002's obstruction
to replacing a one-sided cutoff by compact ones persists in
any character sector with a nonzero prime sum.

In the language of the requested geometric programme, the finite
character is a rank-one finite-monodromy local system on the
appropriate arithmetic open scheme, or its Artin motive over
$K$; duality replaces its Frobenius eigenvalues by their inverses.
These are the SGA~1, Expos\'e~V, \S7 and Deligne
\emph{Valeurs de fonctions $L$}, Definition~6.1 and
\S\S6.2--6.3 placements already used in MA-CH-003 and MA-S5.
The analytic commutator is an additional construction on the
idele-class representation, not an asserted proper algebraic
pushforward.  Its trace/extension mechanism has the inherited
cote~1, scans~079--080 and 111--112 comparison through
BR-TR-01/02 and F-SM-11/12.  None of these algebraic descriptions
removes the Hilbert strip-containment premise.

\paragraph{Exact source scope.}
The source maps are Meyer (2005), \path{arXiv:math/0311468v3},
Lemmas~5.3--5.6 and Theorem~5.8, author pp.~39--44; the
character and dual conventions are GRH\_NIGHT MT-M22,
(M22.1), (M22.4)--(M22.6).  The convergence and kernel estimates
used from this run are MA-MEY-001--003.  All ideal regroupings,
the reflected character computation and the operator specialization
are proved above.  The underlying algebraic character and its
local factors are the already proved MA-CH-003--004.

\section{A coherent additive character for the global Fourier step}
\label{sec:additive-normalization}

The local Fourier computations of MA-CQ use a specified character
and its self-dual measure. Their shell supports and absolute Gauss
normalizations hold for either additive sign. A global Poisson
identity has the extra requirement that the character annihilate
the diagonal arithmetic lattice. That requirement must be checked
before the character PW construction.

\paragraph{MA-NR-001 (the printed signs and their explicit repair).}
Read the finite character in Masullo Example~A.3.1.6, p.~557,
with the natural inclusion
$\mathbb Q_p/\mathbb Z_p=\mathbb Z[1/p]/\mathbb Z\to\mathbb R/\mathbb Z$
printed there. It is
$\psi_p^+(x)=\exp(2\pi i\{x\}_p)$, while
Example~A.3.1.5 on the same page gives
$\psi_\infty^+(x)=\exp(2\pi ix)$.
Their product is the character defined in Example~A.3.1.8,
p.~558. With these literal signs, the self-annihilator assertion
for $\mathcal O_{F,S}$ in Example~A.3.1.11, p.~559,
following (A.3.1.10), is false already for
$F=\mathbb Q$, $S=\{\infty,3\}$.
In fact, for $\Gamma=\mathbb Z[1/3]$ and
$x=1$, $y=1/3$, both diagonally embedded in $\Gamma$,
\[
 \psi_\infty^+(xy)\psi_3^+(xy)
                 =\exp(4\pi i/3)\ne1.
 \tag{MA-NR.1}
\]
Thus even the asserted inclusion
$\Gamma\subseteq\Gamma^\perp$ fails under the printed pairing.
This is a sign convention defect, not a defect of the printed
self-dual Haar measures.

A coherent choice, used for the global constructions here, is
\[
 \vartheta_\infty(x)=\exp(2\pi ix),\qquad
 \vartheta_p(x)=\exp(-2\pi i\{x\}_p),\qquad
 \widehat f(y)=\int f(x)\overline{\vartheta(xy)}\,d\nu(x).
 \tag{MA-NR.2}
\]
For number fields compose these characters with the local traces;
at a complex place the archimedean character is
$\exp(2\pi i\operatorname{Tr}_{\mathbb C/\mathbb R}z)$.
The finite Gauss constants in MA-CQ and MA-CGP are defined by
their displayed integrals with this choice. Changing a finite
additive sign reflects the corresponding Fourier coordinate.
It preserves the finite norm-support bound and the value at zero,
so the corrected source PW spaces are unchanged by this repair.

\begin{proof}
The displayed counterexample follows because the ordinary
$3$-adic fractional part of $1/3$ is $1/3$.
For the repaired global rational character, every rational $r$
satisfies
\[
                    r-\sum_p\{r\}_p\in\mathbb Z.
 \tag{MA-NR.3}
\]
Only primes in its denominator contribute. To prove the assertion,
subtract the indicated finite sum. At each finite prime $p$ the
$p$-th summand removes the negative-power part of $r$, while the
other summands are $p$-integral. The rational difference is
therefore integral at every prime, and hence an integer.
Equation (MA-NR.3) proves that the product in (MA-NR.2)
is trivial on the diagonal $\mathbb Q$.

The exact local repair can also be checked without invoking a
global duality theorem. For any prime $p$, use the pairing
\[
 (x_\infty,x_p),(y_\infty,y_p)\longmapsto
    \exp\bigl(2\pi i(x_\infty y_\infty-\{x_py_p\}_p)\bigr)
 \tag{MA-NR.4}
\]
on $\mathbb R\times\mathbb Q_p$.
The diagonal $\Gamma=\mathbb Z[1/p]$ annihilates itself,
by (MA-NR.3). Conversely suppose $(x_\infty,x_p)$
annihilates $\Gamma$. Testing against $1$ gives
$x_\infty-\{x_p\}_p\in\mathbb Z$, so
$x_\infty=r\in\Gamma$.
Subtract the diagonal point $(r,r)$, which already annihilates
$\Gamma$. Testing the remaining point $(0,x_p-r)$ against
$p^{-k}$ for every $k\ge0$ gives
$x_p-r\in p^k\mathbb Z_p$ for every $k$.
Their intersection is zero, so $x_p=r$.
This proves $\Gamma^\perp=\Gamma$ for the repaired pairing.
Equivalently, under the two positive signs the annihilator is
$\{(r,-r):r\in\mathbb Z[1/p]\}$, which explains the exact
coordinate reflection missing from the printed claim.

If a Fourier kernel $\overline{\psi(xy)}$ is replaced by
$\overline{\psi(-xy)}$, its transform is $\widehat f(-y)$.
Reflection preserves Haar measure, absolute norm, and the origin.
Thus the self-dual measure and all norm-support conditions remain
the same. The local annihilator of $\mathcal O_v$ is still the
inverse different, since it is closed under negation. The local
Gauss phase must instead be recomputed using the chosen kernel;
MA-CQ defines it by an integral, so no unverified phase convention
is imported. These observations prove all asserted repair claims.
\end{proof}

\paragraph{Scope and trace-framework placement.}
This finding uses neither $A$, $M$, nor zero location.
It does not refute the self-duality of the adelic measure in
Example~A.3.1.8: the product with the printed signs still gives
an isomorphism with the Pontryagin dual. What fails is the
identification of the diagonal lattice with its annihilator,
used in the arithmetic Poisson step. The corrected construction
uses a coherent character throughout; the choice is recorded
rather than absorbed into an unnamed root-number scalar.

The multiplicative finite-monodromy characters and their
inertia-invariant Euler polynomials, placed in SGA~1,
Expos\'e~V, \S7 and Deligne, Definition~6.1 and
\S\S6.2--6.3 in MA-CH, do not change under an additive
Fourier-sign repair. This is a normalization of the analytic
Fourier pairing and of the associated Gauss constants, not a
new motive or weight assertion. In the cote~1 trace comparison
(scans~079--080 and~111--112 via BR-TR-01/02), the paired
operators must use the same corrected Fourier identification
before trace cyclicity is applied.

\section{Actual character sections and their local Fourier cutoff limits}
\label{sec:character-quotient-repair}

The algebraic phases in MA-LA-013 do have an actual local
character-sector realization.  Ramified places require a
finite Fourier transform of the unit character, and the
resulting Gauss factor and conductor shift cannot be
replaced by an unramified Euler factor.  For odd Dirichlet
characters the real Fourier transform is not an involution;
the finite Gauss factors supply exactly the missing sign.
We construct these maps explicitly for $K=\mathbb Q$ and
for $K=\mathbb Q(\zeta_3)$, the field of MA-CH-003.
Both have class number one and finite unit group.

Let $S_f$ contain every ramified place of the finite-order
idele-class character $\beta$, and every place in the
different of $K/\mathbb Q$.  Choose a global generator
$\pi_v$ of each prime $v\in S_f$, positive when $K=\mathbb Q$.
Put
\[
\begin{aligned}
 q_v&=N\mathfrak p_v,&
 d_v&=\operatorname{ord}_v(\mathfrak d_{K_v/\mathbb Q_p}),&
 f_v&=\operatorname{ord}_v(\mathfrak f_\beta),\\
 b_v&=\beta_v(\pi_v),&
 \kappa_v&=\beta_v^{-1}|_{\mathcal O_v^\times},&
 \theta_{\beta,v}(x)&=\kappa_v(x)\mathbf1_{\mathcal O_v^\times}(x).
\end{aligned}
\tag{MA-CQ.1}
\]
Thus $|b_v|=1$.  The primitive Euler coefficient is
$c_v=b_v$ if $f_v=0$, and is zero if $f_v>0$.
Our idele convention is precisely (M22.1): the class of
the inverse prime idele has $\beta(p_v)^{-1}=c_v$ at an
unramified place.  For a primitive Dirichlet character
$\chi$, this means $c_p=\chi(p)$ and
$\beta_p|_{\mathbb Z_p^\times}=\chi_p^{-1}$ at its
conductor places.  The raw compact type is $\beta^{-1}$,
as in MA-CH-001; it is not silently inverted below.

Use the coherent local additive character $\vartheta_v$ of
MA-NR-001 (positive archimedean trace and negative finite
fractional trace) and its self-dual measure $\nu_v$.  We write the Fourier kernel as
$\overline{\vartheta_v(xy)}$, fixing every Gauss phase by this
choice.  The conventions are
\[
 \nu_{\mathbb R}=dx,\qquad
 \nu_{\mathbb C}=2\,dx\,dy,\qquad
 \nu_v(\mathcal O_v)=q_v^{-d_v/2}.
\tag{MA-CQ.2}
\]
At a finite place, the annihilator of $\mathcal O_v$ is
$\pi_v^{-d_v}\mathcal O_v$.
These measure and different normalizations are Masullo's
Definition A.3.1.2, printed p.~556, and Examples
A.3.1.5--A.3.1.10, pp.~557--558; MA-NR-001 supplies
the explicitly coherent additive signs.
The following calculation derives the phases from the
chosen Fourier integral instead of importing a root-number
formula with an unspecified sign.

\paragraph{MA-CQ-001 (the actual finite Fourier factors).}
If $f_v=0$, then
\[
 \mathcal F_v\theta_{\beta,v}(y)
   =q_v^{-d_v/2}
       \left(\mathbf1_{\pi_v^{-d_v}\mathcal O_v}(y)
         -q_v^{-1}
           \mathbf1_{\pi_v^{-(d_v+1)}\mathcal O_v}(y)\right).
\tag{MA-CQ.3}
\]
If $f_v>0$, put $m_v=f_v+d_v$ and
\[
 \tau_v(\kappa_v)
   =\int_{\mathcal O_v^\times}\kappa_v(u)
       \overline{\vartheta_v(\pi_v^{-m_v}u)}\,d\nu_v(u),
 \qquad
 \epsilon_v(\kappa_v)=q_v^{m_v/2}\tau_v(\kappa_v).
\tag{MA-CQ.4}
\]
Then
\[
\begin{aligned}
 \mathcal F_v\theta_{\beta,v}(y)
   &=\tau_v(\kappa_v)
        \kappa_v^{-1}(\pi_v^{m_v}y)
            \mathbf1_{\pi_v^{-m_v}\mathcal O_v^\times}(y),\\
 |\epsilon_v(\kappa_v)|&=1,\qquad
 \epsilon_v(\kappa_v)\epsilon_v(\kappa_v^{-1})
                                      =\kappa_v(-1).
\end{aligned}
\tag{MA-CQ.5}
\]
In particular the ramified Fourier transform is one
valuation shell, not the unramified all-pass series.

\begin{proof}
The Fourier transform of $\mathbf1_{\mathcal O_v}$ is
$q_v^{-d_v/2}\mathbf1_{\pi_v^{-d_v}\mathcal O_v}$.
Subtract the transform of $\mathbf1_{\pi_v\mathcal O_v}$,
using additive scaling by $\pi_v$, to obtain (MA-CQ.3).

For $f=f_v>0$, the function
$\kappa_v\mathbf1_{\mathcal O_v^\times}$ is invariant
under additive translations by $\pi_v^f\mathcal O_v$.
Its transform is therefore supported in the annihilator
$\pi_v^{-(f+d_v)}\mathcal O_v$.
If $\operatorname{ord}_v y>-(f+d_v)$, choose a unit
$u_0$ on which $\kappa_v$ is nontrivial but for which
$u_0-1\in\pi_v^{f-1}\mathcal O_v$.
For $f=1$ this means a suitable unit of $\mathcal O_v$;
the conductor definition supplies one.
Then $\vartheta_v(xy(u_0-1))=1$ for every unit $x$.
Changing $x$ to $u_0x$ in the Fourier integral
multiplies it by the nontrivial value $\kappa_v(u_0)$,
so that integral is zero.
Only the shell of valuation $-(f+d_v)$ remains.
For $y=\pi_v^{-(f+d_v)}v_0$ with $v_0$ a unit,
the substitution $u\mapsto v_0^{-1}u$ gives the
first formula in (MA-CQ.5).

Plancherel determines its absolute constant without
an additional Gauss-sum theorem.  The input squared
norm is $\nu_v(\mathcal O_v^\times)$.
The output squared norm is
$|\tau_v|^2q_v^{f+d_v}\nu_v(\mathcal O_v^\times)$.
Consequently $|\tau_v|=q_v^{-(f+d_v)/2}$.
For the last identity apply the Fourier transform twice.
Its square is reflection.  Scaling the shell in the
first line of (MA-CQ.5) gives
\[
 \mathcal F_v^2\theta_{\beta,v}
   =q_v^{f+d_v}
        \tau_v(\kappa_v)\tau_v(\kappa_v^{-1})
                            \theta_{\beta,v}
   =\kappa_v(-1)\theta_{\beta,v}.
\]
This proves both identities in its second line.
\end{proof}

\subsection{A section before Fourier transformation, and a Hilbert completion after it}

Let $\mu_K=\mathcal O_K^\times$, of order two or six, and
$\Gamma_S=\mu_K\prod_{v\in S_f}\pi_v^{\mathbb Z}$.
On the actual polar-free LF space
\[
 \mathcal V_S=
    \mathcal S_{\rm flat}(K_\infty)
       \mathbin{\otimes_\iota}
            \bigotimes_{v\in S_f} C_c^\infty(K_v^\times)
\tag{MA-CQ.6}
\]
let $L_\gamma f(x)=f(\gamma^{-1}x)$.
Flatness at the archimedean origin means that all real
partial derivatives vanish there.  Define
\[
 P_0=\prod_{v\in S_f}M_{\mathbf1_{\mathcal O_v^\times}},
 \qquad E_\mu=\frac1{|\mu_K|}\sum_{\zeta\in\mu_K}L_\zeta,
 \qquad
 \mathscr B_S=P_0
      \sum_{\mathbf n\in\mathbb Z^{S_f}}
             L_{\prod_v\pi_v^{n_v}}E_\mu .
\tag{MA-CQ.7}
\]
After $P_0$, the sum is finite on every LF stage.
The same proof as (MA-LA.46)--(MA-LA.47) shows that
$\mathscr B_S$ is a continuous projection with kernel
exactly the $\Gamma_S$-coinvariant relation space.
For clarity, decompose a vector into its finitely many
valuation tuples; moving each tuple to the unit tuple
changes it by a single unit relation, and averaging over
$\mu_K$ changes it by finitely many more.
The image is the unit-shell space invariant under
simultaneous multiplication by $\mu_K$.
Its kernel is closed, so this is a continuous section
of the separated coinvariant quotient.

Project further to the compact type $\beta^{-1}$ by
the actual compact-group average
\[
        \Pi_\beta f(x)=
            \int_{K_\infty^1\times\prod_{v\in S_f}\mathcal O_v^\times}
                            \beta(c)f(cx)\,dc ,
\tag{MA-CQ.8}
\]
with probability Haar measure.
This is an idempotent because the group is compact
and $\beta$ is a character.  It commutes with the
valuation section.  Its unit-shell vectors are exactly
\[
                u(x_\infty)\prod_{v\in S_f}\theta_{\beta,v}(x_v).
\tag{MA-CQ.9}
\]
For $K=\mathbb Q$, $u(-x)=(-1)^e u(x)$, where
$\chi(-1)=(-1)^e$.
For $K=\mathbb Q(\zeta_3)$ the finite-order character
is trivial at the connected complex place, so $u$
is radial.  The product in (MA-CQ.9) is invariant
under $\mu_K$, because $\beta$ is trivial on
principal ideles.  Thus this is a genuine compact
sector of the actual coinvariant quotient.

Complete this section using the archimedean norm
$L^2(K_\infty,\nu_\infty)$ on $u$, in its indicated
parity or radial subspace; call the resulting space
$\mathscr H_\beta$.  This is a specified quotient
completion, not a quotient norm induced from the
full additive $S$-adelic $L^2$ space.
The latter interpretation would contradict
(MA-LA.58).
The multiplicative norm-one compact quotient has
Haar mass one, as in MA-CH-001.
Equivalently, if a lift is integrated on the
multiplicative cover with counting measure on
$\Gamma_S$, give its compact fiber mass $|\mu_K|$;
the factor $1/|\mu_K|$ in a unit-shell lift then
gives exactly the normalized norm integral below.
This records the Haar choice independently of the
fixed self-dual additive measures (MA-CQ.2).

\paragraph{MA-CQ-002 (the actual dual-sector Fourier maps).}
On the parity or radial archimedean space define the
commuting unitaries
\[
 (V_vu)(x)=q_v^{1/2}u(\pi_v x),\qquad
 D_{v,c}=(1-\overline c\,q_v^{-1/2}V_v^*)
                              (1-cq_v^{-1/2}V_v)^{-1}.
\tag{MA-CQ.10}
\]
The archimedean rotation in $\pi_v$ acts trivially on
the complex radial sector.  Define
\[
 M_{\beta,v}=
 \begin{cases}
   b_v^{d_v}V_v^{-d_v}D_{v,\overline b_v},& f_v=0,\\
   \epsilon_v(\kappa_v)b_v^{f_v+d_v}
                                V_v^{-(f_v+d_v)},& f_v>0,
 \end{cases}
 \qquad M_{\beta,S}=\prod_{v\in S_f}M_{\beta,v}.
\tag{MA-CQ.11}
\]
Then actual Fourier transformation followed by
unit-shell periodization of (MA-CQ.9) is
\[
          H_\beta=M_{\beta,S}\mathcal F_\infty:
                        \mathscr H_\beta\longrightarrow
                                      \mathscr H_{\overline\beta}.
\tag{MA-CQ.12}
\]
It is unitary, and
\[
                         H_{\overline\beta}H_\beta=1.
\tag{MA-CQ.13}
\]
These formulas include odd parity; no assertion
$\mathcal F_{\rm odd}^2=1$ is made.
The folded formulas hold normally on compact
archimedean annuli and in the stated Hilbert norm.
The intermediate vectors need not be flat at zero.

\begin{proof}
All factors in (MA-CQ.10)--(MA-CQ.11) are unitary:
$V_v$ is normalized by the archimedean determinant
$|\pi_v|_\infty=q_v$, the numerator of $D_{v,c}$
is the adjoint of its commuting invertible
denominator, and (MA-CQ.5) gives the modulus of
each ramified scalar.

We verify that they are the actual periodization,
including the character orientation.
For $w\ne v$, $\pi_v$ is a unit at $w$.
The global character relation and the triviality
of a finite-order character on the positive
archimedean norm subgroup give
\[
                \prod_{\substack{w\in S_f\\w\ne v}}
                              \kappa_w(\pi_v)=b_v .
\tag{MA-CQ.14}
\]
After Fourier transformation the unit characters
are $\kappa_w^{-1}$.
In the term $L_{\prod_v\pi_v^{n_v}}$ of the
periodization, their angular evaluations contribute
exactly $\prod_v b_v^{n_v}$.
Indeed at $w$, removing the valuation factor
$\pi_w^{-n_w}$ leaves the unit
$\prod_{v\ne w}\pi_v^{-n_v}$, whose
$\kappa_w^{-1}$ value is
$\prod_{v\ne w}\kappa_w(\pi_v)^{n_v}$.
Multiplication over $w$ proves the assertion.

At an unramified $v$, (MA-CQ.3) permits
$n_v\le d_v$, with coefficient
$q_v^{-d_v/2}(1-q_v^{-1})$, and $n_v=d_v+1$,
with coefficient $-q_v^{-d_v/2-1}$.
Set $k=d_v-n_v\ge0$ in the first range.
The resulting archimedean operator is
\[
 b_v^{d_v}V_v^{-d_v}
 \left((1-q_v^{-1})
       \sum_{k\ge0}\overline b_v^{\,k}q_v^{-k/2}V_v^k
                   -b_v q_v^{-1/2}V_v^*\right)
       =b_v^{d_v}V_v^{-d_v}D_{v,\overline b_v}.
\tag{MA-CQ.15}
\]
The series converges absolutely in operator norm.
At a ramified $v$, the only term is
$n_v=f_v+d_v$.  Its scalar is
$\tau_v(\kappa_v)b_v^{f_v+d_v}$.
Converting the unnormalized archimedean dilation
to $V_v^{-(f_v+d_v)}$ supplies the factor
$q_v^{(f_v+d_v)/2}$, and hence gives the second
line of (MA-CQ.11).
Taking the finite product proves (MA-CQ.12).
Rapid decay of $\mathcal F_\infty u$ and its
derivatives proves normal convergence on annuli;
the geometric operator-norm bounds prove Hilbert
convergence.

This also verifies independence of a representative.
Every vector of the fixed compact type in
$\mathcal V_S$ is, modulo finitely many unit
relations, the unit-shell representative given
by (MA-CQ.7).
Additive Fourier scaling gives
$\mathcal F_S L_\gamma
       =L_{\gamma^{-1}}\mathcal F_S$
for $\gamma\in\Gamma_S$, since $|\gamma|_S=1$.
The just-established convergent folded series
can be reindexed under these translations and
kills their differences.
Thus it produces the same Hilbert vector from
every such representative.  This is a proved
extension through the chosen section, not an
application of $\mathscr B_S$ as a bounded map
on the entire additive Hilbert cover.

It remains to check Fourier square, especially
in odd parity.
On the real parity-$e$ space,
$\mathcal F_\infty^2=(-1)^e$; on the complex
radial space its square is one.
In both cases
$\mathcal F_\infty V_v\mathcal F_\infty^{-1}
        =V_v^*$.
For an unramified factor,
\[
 \mathcal F_\infty D_{v,\overline b_v}
                  \mathcal F_\infty^{-1}=D_{v,b_v}^{-1},
\]
so its factors in
$M_{\overline\beta,S}
 (\mathcal F_\infty M_{\beta,S}
                         \mathcal F_\infty^{-1})$
cancel.  At a ramified place the two scalars
multiply to $\kappa_v(-1)$ by (MA-CQ.5),
while both dilation and $b_v$ factors cancel.
Their total product is
$\prod_{v\in S_f}\kappa_v(-1)=\beta_\infty(-1)$
by the principal-idele relation.
This is $(-1)^e$ over $\mathbb Q$ and one in
the complex radial case.  Multiplication by
$\mathcal F_\infty^2$ therefore gives exactly
(MA-CQ.13).
\end{proof}

Set
\[
       D_{\beta,S}=M_{\overline\beta,S},
 \qquad Q_{\infty,Y}
    =\mathcal F_\infty^{-1}
          M_{B(|\xi|_\infty/Y)}\mathcal F_\infty .
\tag{MA-CQ.16}
\]
On the fixed archimedean parity or radial sector this also
equals $\mathcal F_\infty M_B\mathcal F_\infty^{-1}$,
because the square of $\mathcal F_\infty$ is the scalar
just computed.  The actual quotient Fourier cutoff is
\[
 H_{\overline\beta}M_{B(|\xi|_\infty/Y)}H_\beta
                  =D_{\beta,S}Q_{\infty,Y}D_{\beta,S}^* .
\tag{MA-CQ.17}
\]
This is an identity of actual bounded maps between the
specified character quotient completions.  In particular
the unramified factors on its right are $D_{v,c_v}$,
with the arithmetic coefficient $c_v$, not its
conjugate.  The intermediate Fourier map changes the
sector, which is essential to this orientation.

\subsection{The actual coupled cutoff limit, including parity and the complex place}

Let $d_\infty=1$ over $\mathbb R$ and $d_\infty=2$
over $\mathbb C$, so $|x|_\infty=|x|_{\rm Euclidean}^{d_\infty}$.
For $h\in C_c^\infty(\mathbb R_{>0})$ define
\[
 (A_hu)(x)
    =\int_0^\infty h(t)t^{-1/2}
                 u(t^{-1/d_\infty}x)\,\frac{dt}{t},
                  \qquad g(y)=h(e^y).
\tag{MA-CQ.18}
\]
This is the integrated normalized norm-dilation action
on the compact type, with norm-one Haar mass one.
The logarithmic unitary on that sector is
\[
 (\mathscr U_\infty u)(y)
   =c_\infty e^{y/2}u(e^{y/d_\infty}),
 \qquad c_{\mathbb R}=\sqrt2,\quad
                         c_{\mathbb C}=\sqrt{2\pi}.
\tag{MA-CQ.19}
\]
The complex constant uses $2\,dx\,dy$ in (MA-CQ.2).
Direct substitution gives
$\mathscr U_\infty A_h\mathscr U_\infty^{-1}=C_g$
and $\mathscr U_\infty V_v\mathscr U_\infty^{-1}
       =W_{\log q_v}$, where $(W_Rf)(y)=f(y+R)$.

Choose real profiles $A,B$ on $[0,\infty)$, smooth,
between zero and one, equal to one on $[0,1]$ and
zero on $[2,\infty)$.
For $X\ge2\epsilon>0$ and $Y>0$, put
\[
\begin{aligned}
 a_{\epsilon,X}(x)
       &=A(|x|_\infty/X)-A(|x|_\infty/\epsilon),\\
 T^{\beta,S}_{X,Y,\epsilon}(h)
       &=M_{a_{\epsilon,X}}A_h
          H_{\overline\beta}M_{B(|\xi|_\infty/Y)}H_\beta\\
       &=M_{a_{\epsilon,X}}A_h
                    D_{\beta,S}Q_{\infty,Y}D_{\beta,S}^* .
\end{aligned}
\tag{MA-CQ.20}
\]
This is an operator on the actual quotient completion
$\mathscr H_\beta$ just specified.
Its two Fourier maps act in the indicated order on
the two dual character sectors.

Write
\[
\begin{aligned}
 B_{\beta,S}(h)
   &=\sum_{\substack{v\in S_f\\ f_v=0}}\sum_{k\ge1}
      (\log q_v)q_v^{-k/2}
        \bigl(c_v^kh(q_v^k)+\overline c_v^{\,k}h(q_v^{-k})\bigr),\\
 A_e(h)&=\int_0^\infty
       \frac{e^{-(e+1/2)x}}{1-e^{-2x}}
                       \bigl(h(e^x)+h(e^{-x})\bigr)\,dx
                   \quad(K=\mathbb Q),\\
 A_{\mathbb C}(h)&=\int_0^\infty
       \frac{e^{-x/2}}{1-e^{-x}}
                       \bigl(h(e^x)+h(e^{-x})\bigr)\,dx
                   \quad(K=\mathbb Q(\zeta_3)).
\end{aligned}
\tag{MA-CQ.21}
\]
Use $A_\infty=A_e$ or $A_{\mathbb C}$ as appropriate.
For $h(1)=0$ these are ordinary convergent integrals.
The prime sums are finite.

\paragraph{MA-CQ-003 (the repaired character quotient limit).}
Every operator in (MA-CQ.20) is trace class and has a
nuclear lift into its flat unit-shell character section.
For every compact smooth $h$ with $h(1)=0$,
\[
\begin{aligned}
 \operatorname{Tr}T^{\beta,S}_{X,Y,\epsilon}(h)
       &\longrightarrow A_\infty(h)+B_{\beta,S}(h)
             &&(XY\to\infty,\ \epsilon Y\to0),\\
 \operatorname{Tr}T^{\beta,S}_{X,Y,\epsilon}(h)
       &\longrightarrow0
             &&(\epsilon Y\to\infty,\ X\ge2\epsilon).
\end{aligned}
\tag{MA-CQ.22}
\]
No A, M, or zero-location assumption is used.
For the finite-place product the actual unitary
factor has the commutator trace
\[
\begin{aligned}
 &\operatorname{Tr}\bigl(
       C_g(D_{\beta,S}^*M_\phi D_{\beta,S}-M_\phi)\bigr)\\
 &\hspace{12mm}
   =B_{\beta,S}(h)
       -\left(\sum_{v\in S_f}(f_v+d_v)\log q_v\right)g(0),
\end{aligned}
\tag{MA-CQ.23}
\]
where $\phi(-\infty)=1$, $\phi(+\infty)=0$ and
$\phi'\in C_c^\infty$.
The shift term is
$-\log(|D_K|N\mathfrak f_\beta)\,g(0)$ when $S_f$
contains all the places specified above.
It vanishes in (MA-CQ.22), but its origin in the
actual Fourier factors remains explicit.

\begin{proof}
\emph{Finite-cutoff maps and their nuclear section.}
In logarithmic coordinates $M_aA_h$ has the
smooth kernel
\[
 \bigl(A(e^y/X)-A(e^y/\epsilon)\bigr)h(e^{y-z}).
\tag{MA-CQ.24}
\]
It is supported in a compact rectangle.
Its Fourier-series rank-one expansion, with the
fixed smooth output cutoff specified in
MA-LA-011, is nuclear into the smooth compact
output interval with every derivative.
Returning by (MA-CQ.19) gives a nuclear map
into flat archimedean Schwartz functions of
the specified parity or radial type.
Multiplication by the fixed finite unit
characters places the output in the actual
section (MA-CQ.9).
The remaining factor in (MA-CQ.20) is bounded.
This proves finite-cutoff nuclearity and trace
class.  The continuous inclusion of the section
into its Hilbert completion and the same
rank-one series give the two trace pairings;
thus the quotient trace is justified by
actual maps, without extending the section
map to the full additive Hilbert cover.

\emph{The real parity and complex radial Fourier cutoffs.}
First set $D_{\beta,S}=1$ and compute the
archimedean trace directly.
Let
$K_Y=\mathcal F_\infty^{-1}(B(|\xi|_\infty/Y))$.
This is a radial Schwartz function (even in
the real case), with
\[
 K_Y(x)=Y K_1(Y^{1/d_\infty}x),
                         \qquad \int K_1\,d\nu_\infty=1.
\tag{MA-CQ.25}
\]
Extend the operator on the indicated compact
sector by zero on the other angular sectors.
The parity projector is
$\frac12(1+(-1)^e R)$; the radial projector
over $\mathbb C$ is the probability average
of rotations.
The resulting kernel is Schwartz in both
additive variables: its output variable has
compact annular support and its input
variable has uniform Schwartz decay, because
$h$ has compact norm support.
Its actual trace is therefore its diagonal
integral, justified by an absolutely
summable Fourier or Hermite rank-one
expansion as in MA-RS-002.

For the real parity-$e$ sector the trace is
\[
 \int_0^\infty h(t)t^{-1/2}
    \frac12\int_{\mathbb R}a_{\epsilon,X}(x)
       \left[
        K_Y((t^{-1}-1)x)
        +(-1)^eK_Y((t^{-1}+1)x)
       \right]dx\,\frac{dt}{t}.
\tag{MA-CQ.26}
\]
For the complex radial sector it is
\[
 \int_0^\infty h(t)t^{-1/2}
    \int_0^{2\pi}\int_{\mathbb C}
       a_{\epsilon,X}(x)
       K_Y((t^{-1/2}e^{i\theta}-1)x)
           \,d\nu_\infty(x)\,\frac{d\theta}{2\pi}\,\frac{dt}{t}.
\tag{MA-CQ.27}
\]
Both formulas are obtained from the actual
angular projector, not by assigning an
archimedean scalar distribution.

For every nonzero real or complex scalar $L$,
substitute $z=Y^{1/d_\infty}x$ in the inner
integral.  It becomes
\[
 \int\left[
        A(|z|_\infty/(XY))
         -A(|z|_\infty/(\epsilon Y))
      \right]K_1(Lz)\,d\nu_\infty(z).
\tag{MA-CQ.28}
\]
If $XY\to\infty$ and $\epsilon Y\to0$,
dominated convergence gives
$|\det_{\mathbb R}L|^{-1}$.
If $\epsilon Y\to\infty$, then also
$XY\to\infty$, and the same integral
tends to zero.
Its absolute value is bounded uniformly
by $\|K_1\|_1|\det_{\mathbb R}L|^{-1}$.
The annular multiplier is bounded by one
because $X\ge2\epsilon$ and
$a=A(|x|_\infty/X)(1-A(|x|_\infty/\epsilon))$.

For $t\ne1$ the real determinant bounds
give
\[
 \frac{t^{-1/2}}2
      \left(\frac1{|t^{-1}-1|}
                        +\frac1{t^{-1}+1}\right).
\tag{MA-CQ.29}
\]
In the complex case use the elementary
angular identity
\[
 \int_0^{2\pi}
   \frac{d\theta/(2\pi)}{|1-r e^{i\theta}|^2}
                       =\frac1{|1-r^2|}\quad(r>0,\ r\ne1).
\tag{MA-CQ.30}
\]
For $r<1$ it follows by expanding the two
geometric series and integrating their
equal powers; for $r>1$ factor out $r^2$
and apply the first case.
Thus the complex majorant after angular
integration is
$t^{-1/2}/|1-t^{-1}|$.
Multiplying either majorant by
$|h(t)|\,dt/t$ is integrable, since
$h(1)=0$ and $h$ has compact support
inside $(0,\infty)$.
The exceptional set $t=1$ has measure
zero.  Dominated convergence therefore
justifies both full trace limits.

For the first real limit the remaining
coefficient is
\[
 \frac{\sqrt t}{2}
           \left(\frac1{|1-t|}
                          +\frac{(-1)^e}{1+t}\right).
\]
Split the norm integral at $t=1$ and
substitute $t=e^x$ and $t=e^{-x}$.
The coefficient of $h(e^x)+h(e^{-x})$
is
\[
 \frac{e^{-x/2}}2
    \left(\frac1{1-e^{-x}}
                          +\frac{(-1)^e}{1+e^{-x}}\right)
       =\frac{e^{-(e+1/2)x}}{1-e^{-2x}}.
\tag{MA-CQ.31}
\]
This proves the parity formula $A_e$,
including the odd case.
The corresponding complex coefficient
is $\sqrt t/|1-t|$, which gives
$e^{-x/2}/(1-e^{-x})$ after the same
two substitutions.  This proves
$A_{\mathbb C}$ with coefficient one.
The second archimedean limit is zero.

\emph{Finite primes, conductor and different shifts.}
In logarithmic coordinates (MA-CQ.10)
has exactly the Laurent series of
MA-LA-013 with $p$ replaced by $q_v$
and phase $c_v$.  Its trace-class
piece estimate and coefficient
calculation therefore give the
$v$-term of $B_{\beta,S}$ when
$f_v=0$.
The additional factor $V_v^{-d_v}$
is a unitary shift.  Its commutator
with $M_\phi$ has multiplier
$\phi(y+d_v\log q_v)-\phi(y)$,
whose integral is $-d_v\log q_v$.
After convolution its trace is
$-d_v(\log q_v)g(0)$.
At a ramified place there is no
Euler all-pass factor: the scalar
in (MA-CQ.11) cancels from conjugation,
and its shift contributes
$-(f_v+d_v)(\log q_v)g(0)$.
All these commutators become trace
class after $C_g$, by the compact
cutoff-difference kernel estimate
of MA-LA-013.
Bounded cyclicity for the finite
commuting product gives their sum.
This proves (MA-CQ.23).
The equality
$\prod_v q_v^{f_v+d_v}
       =|D_K|N\mathfrak f_\beta$
is the norm formula for conductor
and different, with all such places
included in $S_f$.

For the exact coupled order put
$\Phi(y)=A(e^y)$ and define the fixed
trace-class operator
\[
 E_{\beta,S}(g)=
       (D_{\beta,S}^*M_\Phi D_{\beta,S}-M_\Phi)C_g.
\tag{MA-CQ.32}
\]
The order with convolution on the
right follows by taking the adjoint
of the left-order formula for
$g^*(y)=\overline{g(-y)}$.
The two directional character
coefficients in (MA-CQ.23) and
complex conjugation give the same
trace, including the real conductor
coefficient.  Hence
$\operatorname{Tr}E_{\beta,S}(g)
        =B_{\beta,S}(h)$ when $g(0)=0$.

All $D$ factors and $C_g$ commute
with $W_R$.  Positive archimedean
Fourier scaling gives on either
compact sector
\[
                W_RQ_{\infty,Y}W_{-R}
                         =Q_{\infty,e^R Y}.
\tag{MA-CQ.33}
\]
Repeating the actual bounded
cyclicity calculation (MA-LA.75),
with the right order in (MA-CQ.32),
gives
\[
\begin{aligned}
 \operatorname{Tr}T^{\beta,S}_{X,Y,\epsilon}(h)
  -\operatorname{Tr}\bigl(M_aA_hQ_{\infty,Y}\bigr)
   &=\operatorname{Tr}\bigl(E_{\beta,S}(g)Q_{\infty,XY}\bigr)\\
   &\quad-\operatorname{Tr}
           \bigl(E_{\beta,S}(g)Q_{\infty,\epsilon Y}\bigr).
\end{aligned}
\tag{MA-CQ.34}
\]
In this step $C_g$ is never commuted
past the Fourier cutoff.
The bounded multipliers
$Q_{\infty,L}$ converge strongly
to one as $L\to\infty$ and to zero
as $L\to0$.  Pairing these limits
with the fixed trace-class operator
$E_{\beta,S}(g)$ gives
$B_{\beta,S}(h)$ in the first
regime and zero in the second.
Combining with the independently
proved archimedean limits proves
(MA-CQ.22).
\end{proof}

\paragraph{MA-CQ-004 (the non-induced cubic character in an actual local quotient).}
Use the character $\psi=\psi_{269,3}$ of MA-CH-003.
Let
\[
 S_0=\{\infty,\mathfrak p_3,\mathfrak q_{269}\},
 \qquad S_1=S_0\cup\{\mathfrak p_2\},
 \qquad
 N\mathfrak p_3=3,\quad N\mathfrak q_{269}=269^2,\quad
 N\mathfrak p_2=4 .
\tag{MA-CQ.35}
\]
The section for $\beta=\psi$ consists of radial flat
complex Schwartz functions times
$\psi_{\mathfrak q_{269}}^{-1}$
on its unit shell, and the trivial unit characters
at the other selected finite places.
It is an actual section of the $S$-unit coinvariants.
Its Fourier cutoff factors can be written
\[
\begin{aligned}
 D_{\psi,S_0}
   &=\epsilon_\psi\,
       V_3^{-1}D_{3,1}\,V_{269^2}^{-1},\\
 D_{\psi,S_1}
   &=D_{\psi,S_0}D_{4,1},\qquad |\epsilon_\psi|=1 .
\end{aligned}
\tag{MA-CQ.36}
\]
Here the subscripts on $V$ name the norms of the
chosen prime generators, and
$\epsilon_\psi=\epsilon_{\mathfrak q_{269}}
       (\psi_{\mathfrak q_{269}}|_{\mathcal O^\times})$
is the explicit finite Fourier factor (MA-CQ.4).
The prime at $3$ retains its primitive Euler factor.
The prime at $269$ contributes a conductor shift
and Gauss factor, with no Euler all-pass factor.

For the test $h$ of MA-CH-004,
$h(1)=0$, $A_{\mathbb C}(h)=0$, and the actual
coupled local traces satisfy
\[
\begin{aligned}
 \lim\operatorname{Tr}T^{\psi,S_0}_{X,Y,\epsilon}(h)&=0,\\
 \lim\operatorname{Tr}T^{\psi,S_1}_{X,Y,\epsilon}(h)&=\log2
                   \qquad(XY\to\infty,\ \epsilon Y\to0).
\end{aligned}
\tag{MA-CQ.37}
\]
These are unconditional limits on the specified
local quotient completions.  They do not invoke
Hecke GRH or identify those completions with a
global narrow orthogonal projection.

\begin{proof}
The arithmetic input is the fully checked
MA-CH-003: $K=\mathbb Q(\zeta_3)$ has class
number one and unit group $\mu_6$;
$K(269^{1/3})/K$ is cyclic of degree three;
$\psi$ has conductor $(269)$.
The extension splits at the prime over $3$
and at the prime over $2$, so their primitive
coefficients are one.  The local conductor
at $\mathfrak q_{269}$ has exponent one
and its inertia invariants are zero.
The different has exponent one at
$\mathfrak p_3$ and zero at the other
two selected places.

Choose global generators $1-\zeta_3$, $269$
and $2$, respectively.  Their archimedean
norms are $3$, $269^2$ and $4$.
The uniformizer value $b_{269}$ is one:
the principal-idele relation for $269$
has trivial archimedean character and
all its other local components are
units at unramified characters.
The values $b_3=b_2=1$ are the already
verified split Frobenius values.
Substitution of these data into
(MA-CQ.11), with $\beta$ replaced by
$\overline\beta$ as in (MA-CQ.16),
gives exactly (MA-CQ.36).
The unit invariance of the section
follows from the global character
relation on $\mu_6$, and the actual
section and Fourier maps are those
proved in MA-CQ-002.
No exact numerical value of the
Gauss phase is needed; it is a
specified finite integral, with its
unit modulus proved in MA-CQ-001.

The test from MA-CH-004 has one
bump near norm four with value one
there, and a second bump in
$(1.4,1.6)$ chosen to cancel its
complex archimedean integral.
Its supports avoid the prime
powers at $3$, and there is no
Euler factor at the ramified
prime over $269$.
Thus $B_{\psi,S_0}(h)=0$.
Adding $\mathfrak p_2$ contributes
$(\log4)/\sqrt4=\log2$;
the reciprocal support contributes
zero.  The shift term in
(MA-CQ.23) is
$-\log(3\cdot269^2)h(1)=0$.
Equation (MA-CQ.22) now gives both
actual limits in (MA-CQ.37).
\end{proof}

\paragraph{Grothendieck placement and the exact global boundary.}
The finite-order character is a rank-one finite-monodromy
representation of the arithmetic fundamental group, hence
a rank-one finite-monodromy etale local system on its
unramified open arithmetic curve.
At an unramified prime its Euler polynomial is
$\det(1-\operatorname{Fr}_vT)=1-c_vT$.
At a ramified prime the determinant is taken on inertia
invariants; in the cubic example that space is zero
at $\mathfrak q_{269}$, so the primitive local factor
is one there.
These are the SGA~1, Expos\'e~V, \S7 monodromy
convention and Deligne, \emph{Valeurs de fonctions L},
Definition~6.1 and \S\S6.2--6.3, with the exact
source scope already audited in MA-CH.
For $\psi$, induction to $\mathbb Q$ is the verified
odd two-dimensional $S_3$ Artin representation of
MA-CH-003, and $\psi^c=\psi^{-1}\ne\psi$ proves that
it is not induced from a rational norm character.

The compact character projector, unit-shell section,
dual Fourier maps and Gauss factors constructed here
are analytic realizations of that character data.
The factor $q_v^{-1/2}$ and the norm-half centering
in (MA-CQ.18)--(MA-CQ.19) arise from normalized
unitary dilation and additive measure.
They are not an actual integral Tate twist of
the Artin motive.  No half-integral Tate object
or new arithmetic weight construction is being
asserted.  Finite monodromy gives the usual
weight-zero Frobenius data; it does not by itself
give a global zero-location bound for the
associated $L$-function.

All limits above concern the specifically constructed
local quotient completions.  The section is not
bounded on the full additive Hilbert cover, the
intermediate Fourier vectors need not be polar-free,
and no equality with the global weighted orthogonal
Paley--Wiener condensation is inferred.
The missing global-to-local factorization isolated
after MA-LA-013 remains a separate operator statement.
In particular neither A nor M is used here to
manufacture an independent narrow identity.
The global corrected cutoff and its discarded-divisor
criterion remain MA-WP, MA-EG and MA-PR.

\paragraph{Exact source and proof scope.}
The actual additive measures and Fourier conventions
were checked in the supplied Masullo text at Definition
A.3.1.2, p.~556, and Examples A.3.1.5--A.3.1.10,
pp.~557--558, including (A.3.1.6) and (A.3.1.8).
The different normalization is also stated at
Remark B.4.4.6 and its proof, p.~623.
The LF polar-free and finite-valuation conventions
are the explicit ones already checked in MA-LA-009--011,
with source Definition~5.3.4.1, p.~259,
pp.~260--261, and Proposition~2.1.4.1(1),
pp.~71--72.
The character convention is (M22.1)--(M22.4),
and the source compact Haar normalization is the
probability normalization of MA-CH-001.
The finite character-shell transforms, their
Gauss moduli, the parity sign, the actual
quotient Fourier maps, both archimedean
coupled cutoff limits and the finite-prime
transport are proved above.
The arithmetic of the cubic example is used
from MA-CH-003--004 with its previously
completed proof checks.
No source PDF or extraction is included in
the deliverable, and no new independent
global zero-location argument is claimed.


\section{General number-field annular transport on the original flat quotient}
\label{sec:general-annular-transport}

\paragraph{Status.}
The complete author reading and two independent
whole-proof readings have passed.  It addresses the pre-existing question
\texttt{Q-MA-LA-GENERALS}.  The field, source topology, Fourier
operator and two physical norm cutoffs below are part of the statement.
They are not replaced by a radial additive model.

Let $K$ be a fixed number field of signature $(r_1,r_2)$ and degree
$n=r_1+2r_2$.  Let $S$ be a finite structural set of places containing
all infinite places.  Here structural means Meyer's sufficiently
large condition
$C_K/\mathcal O_{\overline S}^{\times}\simeq C_S$; no prime in $S$
is required to be principal.  Set
\[
 G=\mathbb A_S^\times,\qquad
 \Gamma=\mathcal O_{K,S}^\times,\qquad
 G^1=\ker|\cdot|_S,\qquad C_S=G/\Gamma,
 \qquad C_S^1=G^1/\Gamma.
 \tag{MA-GA.1}
\]
The $S$-unit group is discrete and cocompact in $G^1$.
In particular $C_S^1$ is compact, including when the ordinary
unit rank is positive.  Use the positive homomorphic norm splitting
$\nu(e^x)$ whose infinite coordinates are $e^{x/n}$ and whose
finite coordinates are $1$.  The normalized absolute value at a
complex place is the square of the ordinary absolute value, so
$|\nu(e^x)|_S=e^x$.  Give $C_S^1$ probability Haar measure.

The original flat source, with its LF topology, is
\[
 V_S=\mathcal S_{00}(K_\infty)
       \mathbin{\widehat\otimes}_{\iota}
       \bigotimes_{v\in S_f}C_c^\infty(K_v^\times).
 \tag{MA-GA.2}
\]
Here every ordinary partial derivative of the infinite component
vanishes on each coordinate hyperplane $x_v=0$.  At a finite
place the factor means a locally constant, compactly supported
function whose support avoids zero, extended by zero to $K_v$.
It does not mean the larger multiplicative Schwartz space with
unbounded valuation support.  A finite-place LF stage fixes a
compact punctured support and an open invariance subgroup; its
finite-place factor is finite dimensional.  We specify
the topology by these strict LF stages, as in MA-LA.41;
no interchange of a projective tensor product with arbitrary
inductive limits is assumed.
Intrinsically this is the coordinate-flat closed subspace
of the additive Schwartz--Bruhat space
$\mathcal S(\mathbb A_S)$ with its inherited LF topology.
The inductive tensor notation in (MA-GA.2) records that
topology.  We do not assert that the printed projective
tensor symbol in Masullo's (5.3.4.1), taken as a separate
topological prescription, equals it.

Fix a smooth unitary character $\chi$ of $C_S^1$ and extend it
trivially along $\nu$.  Its raw compact type is $\chi^{-1}$.
In the centered norm coordinate the corresponding Hilbert space
is $H_\chi=L^2(\mathbb R,dx)$.  A scalar function $F$ represents
$\chi(\bar k)^{-1}F(x)$, where $k\in G^1$ and $x=\log|t|_S$.
Write
\[
 \mathcal E=\{F\in C^\infty(\mathbb R):
      \sup_x e^{N|x|}|F^{(j)}(x)|<\infty
          \text{ for every }N,j\geq0\}.
 \tag{MA-GA.3}
\]
All character and norm identifications below are made on this
actual quotient, with the above Haar normalization.

Choose $A,B\in C^\infty([0,\infty))$, between zero and one,
equal to one on $[0,1]$ and zero on $[2,\infty)$.  For
$X\geq2\epsilon>0$, $Y>0$, and $g\in C_c^\infty(\mathbb R)$,
let
\[
\begin{aligned}
 a_{\epsilon,X}(x)&=A(e^x/X)-A(e^x/\epsilon),\\
 (C_gF)(x)&=\int_{\mathbb R}g(y)F(x-y)\,dy,\\
 Q_{\chi,Y}&=\mathfrak S_{\bar\chi}
               M_{B(e^x/Y)}\mathfrak S_\chi,\\
 T^\chi_{X,Y,\epsilon}(g)&=M_{a_{\epsilon,X}}C_gQ_{\chi,Y}.
\end{aligned}
 \tag{MA-GA.4}
\]
The operators $\mathfrak S_\chi:H_\chi\to H_{\bar\chi}$ are
the actual descended additive Fourier transforms with the
coherent local additive characters and self-dual measures of
MA-NR-001/MA-CQ.  Thus $B(e^x/Y)$ cuts the physical total norm
on the Fourier side.  Neither this cutoff nor $A(e^x/X)$ is
a cutoff in an individual archimedean coordinate.

Let $RF(x)=F(-x)$, interchanging the two compact sectors, and put
\[
 U_\chi=\mathfrak S_{\bar\chi}R,
 \qquad \mathcal M(U_\chi F)(z)=m_\chi(z)\mathcal MF(z),
 \qquad \mathcal MF(z)=\int_{\mathbb R}F(x)e^{zx}\,dx.
 \tag{MA-GA.5}
\]
This defines $m_\chi$ by the actual local Fourier operator;
all its conductor, different and phase factors are retained.

\paragraph{MA-GA-001 (general structural-$S$ source, lift and annular rate).}
The separated $\Gamma$-coinvariants of (MA-GA.2) identify
topologically with the smooth functions on $C_S$ having every
exponential norm weight.  Their fixed $\chi^{-1}$ compact
sector is precisely (MA-GA.3).  Centered periodization has a
continuous section into the original source; its central
Hilbert completion is the stated $H_\chi$.

The operator $T^\chi_{X,Y,\epsilon}(g)$ in (MA-GA.4) is trace class on that completion
and has a continuous nuclear lift into (MA-GA.2).  The lift
and its quotient operator have the same trace, computed from
one nuclear series.  The intermediate Fourier vectors are
allowed in the actual completion and need not be flat.

For $g(0)=0$, define the actual local trace
\[
 \mathcal W_{S,\chi}(g)
    =-\frac1{2\pi}\int_{\mathbb R}
             \mathcal Mg(it)\frac{m_\chi'(it)}{m_\chi(it)}\,dt.
 \tag{MA-GA.6}
\]
For each $0<\delta<1/2$ there is a constant $C_\delta$,
depending on the fixed field, $S,\chi,A,B,g,\delta$ but not
on $X,Y,\epsilon$, such that, with $U=XY$ and $V=\epsilon Y$,
\[
\begin{aligned}
 \left|\operatorname{Tr}T^\chi_{X,Y,\epsilon}(g)
                    -\mathcal W_{S,\chi}(g)\right|
   &\leq C_\delta(U^{-\delta}+V^\delta)
          &&(U\geq1,\ V\leq1),\\
 \left|\operatorname{Tr}T^\chi_{X,Y,\epsilon}(g)\right|
   &\leq C_\delta V^{-\delta}
          &&(V\geq1).
\end{aligned}
 \tag{MA-GA.7}
\]
For tests in a fixed compact support the constants are
controlled by finitely many smooth seminorms of $g$.
Consequently the two coupled limits are, respectively,
$\mathcal W_{S,\chi}(g)$ when $XY\to\infty$ and
$\epsilon Y\to0$, and zero when $\epsilon Y\to\infty$.
These statements allow a nontrivial class group, every unit
rank and every archimedean signature.  They use no global
zero-location hypothesis.

\begin{proof}
\emph{1. The original flat topology and its periodization.}
For a real coordinate let $r_v$ be its ordinary absolute
value, and for a complex coordinate let it be its ordinary
modulus.  Restriction of the flat archimedean source to
$K_\infty^\times$ is characterized, for every $N\geq0$
and every invariant differential operator $D$, by
\[
 \sup_{t\in K_\infty^\times}
    \exp\!\left(N\sum_{v\mid\infty}|\log r_v(t)|\right)
       |D f(t)|<\infty.
 \tag{MA-GA.8}
\]
These seminorms and the original flat subspace topology are
equivalent.  Here is a direct proof of the estimates needed
in both directions.  Split the coordinates into $r_v\leq1$
and $r_v>1$.  In each small coordinate apply Taylor's formula
along the segment from zero to that coordinate, to any
prescribed order.  Every lower Taylor coefficient is zero
on the entire coordinate hyperplane, also after differentiating
in the other coordinates.  Iterating gives a product of
arbitrarily high powers of the small $r_v$.  A sufficiently
large ordinary Schwartz weight bounds the product of negative
powers of all large $r_v$.  Invariant derivatives are finite
sums of polynomial coordinate coefficients times ordinary
derivatives, so the same argument applies after such a
derivative.  Each seminorm in (MA-GA.8) is therefore bounded
by finitely many higher ordinary Schwartz seminorms.

Conversely ordinary derivatives on the punctured coordinates
are combinations of invariant derivatives with coefficients
bounded by powers of $r_v^{-1}$.  The estimates (MA-GA.8)
imply that every such derivative tends to zero faster than
every power at each coordinate hyperplane, uniformly when
the other coordinates range over a compact set.  Extending
by zero gives a smooth function there: in one real direction
use the fundamental theorem of calculus on the segment from
zero, then repeat for every derivative; in a complex coordinate
use the same argument in its two real directions.  Repeating
over the coordinates proves the asserted simultaneous flatness.
The large-coordinate estimates give all ordinary Schwartz
seminorms.  This proves the converse continuity as well.

It also proves density of $C_c^\infty(K_\infty^\times)$ in
the flat source.  Indeed multiply by smooth cutoffs that equal
one on $|\log r_v|\leq R$ and vanish on
$|\log r_v|\geq R+1$, for every $v$.  The error tends to zero
in each seminorm (MA-GA.8), using a strictly higher exponential
weight.  At finite places this procedure changes neither
the support nor the invariance level.  Thus
$C_c^\infty(G)$ is dense in $V_S$, with convergence in each
original finite-place LF stage.

Let $\mathcal E(C_S)$ mean the analogous all-exponential-weight
smooth space on $C_S\simeq C_S^1\times\mathbb R$, with the
usual LF union over finite compact invariance levels.
Define centered periodization by
\[
 (\mathcal P f)(\bar t)
      =|t|_S^{1/2}\sum_{\gamma\in\Gamma}f(\gamma t).
 \tag{MA-GA.9}
\]
The sum and all its differentiated, exponentially weighted
versions converge, and $\mathcal P:V_S\to\mathcal E(C_S)$
is continuous.  We verify the point that is absent from a
finite-unit-shell proof.  Cover the compact quotient $C_S^1$
by representatives in one compact subset of $G^1$.
On a fixed source LF stage and these representatives, the
finite valuations of a nonzero summand in (MA-GA.9) permit
only finitely many valuation vectors of $\gamma$.
Each nonempty such fiber is a coset of $\mathcal O_K^\times$.
The archimedean logarithms of ordinary units form a lattice
in the norm-zero real hyperplane, with finite kernel the
roots of unity.  This also follows by projecting their
discrete cocompact subgroup through the compact angular
kernel.  The lattice has polynomial counting growth.

For a representative $k$ in the fixed compact set and
$t=k\nu(e^x)$, the vector of archimedean logarithms of
$\gamma t$ is a fixed bounded vector plus its unit-lattice
vector plus $(x/n,\ldots,x/n)$.  Normalized complex logs
have the corresponding factors two.  Norm equivalence
between this direct decomposition and Euclidean coordinates
gives constants $c,C>0$ such that
\[
 \sum_{v\mid\infty}|\log r_v(\gamma k\nu(e^x))|
       \geq c\bigl(|x|+\|\ell(\gamma)\|\bigr)-C
 \tag{MA-GA.10}
\]
on each of the finitely many allowed valuation fibers.
The fiber representatives are absorbed in $C$.
An arbitrarily high weight in (MA-GA.8) now dominates the
factor $e^{x/2}$, every required norm weight and the lattice
sum.  Differentiation only replaces the source seminorm by
one of higher order.  This proves the assertion uniformly
on the compact quotient.  The finite-place invariance level
of the periodization is fixed by that of the source stage,
so this is continuity in the stated LF topology.

\emph{2. A section and the exact closed relation space.}
Because the action of the discrete subgroup $\Gamma$ on
$G^1$ is proper and its quotient is compact, choose
$b_0\in C_c^\infty(G^1)$, $b_0\geq0$, whose translates
cover the quotient with positive sum.  The sum is locally
finite.  The function
\[
 b(k)=\frac{b_0(k)}{\sum_{\gamma\in\Gamma}b_0(\gamma k)}
 \quad\text{satisfies}\quad
       \sum_{\gamma\in\Gamma}b(\gamma k)=1.
 \tag{MA-GA.11}
\]
It is smooth, compactly supported and has a fixed finite
invariance level.  The formula includes finite torsion;
no division by an unrecorded number of roots of unity occurs.
Define, for $\Psi\in\mathcal E(C_S)$,
\[
 (\sigma\Psi)(k\nu(e^x))=e^{-x/2}b(k)\Psi(\bar k,x).
 \tag{MA-GA.12}
\]
Then $\mathcal P\sigma=1$.  This section takes values in
the original $V_S$, continuously.  Indeed on its support
all archimedean moduli are comparable to $e^{x/n}$ and all
finite valuations are in a fixed finite set.  Every source
seminorm (MA-GA.8) is therefore controlled by finitely many
norm-exponential and compact-direction seminorms of $\Psi$.
The finite invariance level is the common refinement of
those of $b$ and $\Psi$.  The estimates just proved give
the flat zero extension, not merely a multiplicative
Schwartz function.  On compact norm support the image is
in $C_c^\infty(G)$ in a fixed original LF stage.

Let $L_\gamma f(t)=f(\gamma^{-1}t)$.  The kernel of
$\mathcal P$ is exactly
\[
 \ker\mathcal P
     =\overline{\operatorname{span}
          \{L_\gamma f-f:\gamma\in\Gamma,\ f\in V_S\}}^{\,V_S}.
 \tag{MA-GA.13}
\]
For a compactly supported $f$ put
$h_\gamma(t)=b(k)f(\gamma t)$, $t=k\nu(e^x)$.
Only finitely many $h_\gamma$ are nonzero: their possible
$t$ lie in a compact set, and properness of the action
then permits only finitely many $\gamma$.  Directly,
\[
 \sum_\gamma h_\gamma=\sigma\mathcal P f,
 \qquad
 \sum_\gamma L_\gamma h_\gamma=f,
 \qquad
 f-\sigma\mathcal P f
          =\sum_\gamma(L_\gamma-1)h_\gamma.
 \tag{MA-GA.14}
\]
Each term belongs to the original compact flat source.
For general $f$ use the compact approximations from step 1
and continuity of both maps.  If $\mathcal P f=0$ this
places $f$ in the closed relation space.  The converse
follows by reindexing the convergent periodization sum.
Equations (MA-GA.12--14) give the topological quotient
isomorphism, with continuous inverse induced by $\sigma$.

Compact-character projection on $C_S^1$, with its probability
Haar measure, is a continuous projector on this quotient.
Its $\chi^{-1}$ range is exactly $\mathcal E$ of (MA-GA.3).
In particular
\[
 (\sigma_\chi F)(k\nu(e^x))
       =e^{-x/2}b(k)\chi(\bar k)^{-1}F(x)
 \tag{MA-GA.15}
\]
is a continuous section of this sector.  Centering turns
$L^2(C_S,|t|_S\,d^\times t)$ on the sector into $L^2(dx)$.
Since $C_c^\infty(\mathbb R)\subset\mathcal E$, its completion
is exactly $H_\chi$.  This construction neither assumes a
trivial ideal class group nor chooses generators of primes.
Compact averaging occurs on $C_S^1$, not on the generally
noncompact group $G^1$.

The balanced splitting is a coordinate choice, not a change
of the physical operator.  If the original source uses
another positive homomorphic splitting, write
$\nu_0(e^x)=\kappa(x)\nu(e^x)$ and
$\chi(\kappa(x))=e^{i\lambda x}$.  The raw $\chi^{-1}$
convention gives $F_0(x)=e^{-i\lambda x}F(x)$.
With $D=M_{e^{-i\lambda x}}$ one has exactly
\[
 D C_gD^{-1}=C_{e^{-i\lambda\cdot}g},\qquad
 D M_{A(e^x/X)}D^{-1}=M_{A(e^x/X)},\qquad
 Q^{\nu_0}_{\chi,Y}=D Q^\nu_{\chi,Y}D^{-1}.
 \tag{MA-GA.15a}
\]
The Fourier-side identity follows by using the conjugate
coordinate change in the dual compact sector; the two
changes cancel around its physical norm multiplier.
Thus the theorem for every compact smooth $g$ applies
to the original splitting with the displayed test twist.
Its condition $g(0)=0$ is preserved.  Neither physical
norm cutoff has changed.

\emph{3. The actual local Fourier operator and its bounds.}
Use the additive Fourier transform on $\mathcal S(\mathbb A_S)$,
the chosen self-dual measures, and
$Jf(t)=|t|_S^{-1}f(t^{-1})$ on $G$.
Meyer's descent lemma applies to the equivariant operator
$\mathcal F J$: it is a continuous unitary multiplier on
$L^2(C_S,|t|_S\,d^\times t)$.  Thus the descended additive
Fourier operator interchanges the two compact sectors and
reverses the norm coordinate.  Since $-1\in\Gamma$, its
square on the quotient is one.  In the present notation,
\[
 \mathfrak S_{\bar\chi}\mathfrak S_\chi=1,
 \qquad U_\chi\text{ is unitary},\qquad
 \mathfrak S_{\bar\chi}=U_\chi R,
 \quad\mathfrak S_\chi=R U_\chi^{-1}.
 \tag{MA-GA.16}
\]
This is the descent of the actual transform on the source,
not a choice of an unrelated unitary with the same scalar
trace.  The flat quotient embeds into the full local
coinvariant space $\mathcal T(C_S)$ by (MA-GA.9).
Meyer's full local-coinvariant theorem identifies that space
with the separated coinvariants of $\mathcal S(\mathbb A_S)$.
Our additional argument (MA-GA.8--14) identifies the smaller
flat quotient with its own, generally stronger, topology.
No claim that Fourier preserves flatness is required.

For completeness the exact local functional equation can
define its multiplier without a Fourier-sign convention
ambiguity.  Define $\gamma_v(s,\omega_v,\psi_v)$ by
\[
 Z_v(\mathcal F_v f,\omega_v^{-1},1-s)
       =\gamma_v(s,\omega_v,\psi_v)Z_v(f,\omega_v,s).
\]
Then Mellin unfolding of (MA-GA.9), and $U_\chi=\mathcal FJ$
on its sector, give
\[
 m_\chi(z)=\prod_{v\in S}
    \gamma_v(1/2-z,\chi_v^{-1},\psi_v).
 \tag{MA-GA.17}
\]
The raw $\chi^{-1}$ convention is important here: its
coefficient is extracted by integrating against $\chi$, so
the unfolded Mellin transform is $Z(f,\chi,1/2+z)$.
For $s=1/2+z$, substitution $t\mapsto t^{-1}$ gives
$Z(Jf,\chi^{-1},1-s)=Z(f,\chi,s)$.  Applying the displayed
functional equation with the character $\chi^{-1}$ and
parameter $1-s$ therefore gives exactly (MA-GA.17), rather
than its reciprocal or its conjugate-character version.
Equivalently each factor is
\[
 \gamma_v(1/2-z,\chi_v^{-1},\psi_v)
   =\varepsilon_v(1/2-z,\chi_v^{-1},\psi_v)
       \frac{L_v(1/2+z,\chi_v)}
            {L_v(1/2-z,\chi_v^{-1})}.
 \tag{MA-GA.17a}
\]
The normalization of the quotient Haar measure cancels in
this ratio.  Here $\psi_v$ means the actual Fourier kernel,
including the conjugation specified in MA-CQ.

Explicitly at a real place write
$\chi_v(t)=\operatorname{sgn}(t)^{e_v}|t|_v^{i\tau_v}$.
At a complex place write
$\chi_v(t)=(t/|t|_{\mathrm{ordinary}})^{j_v}|t|_v^{i\tau_v}$.
Here $e_v\in\{0,1\}$, $j_v\in\mathbb Z$ and
$\tau_v\in\mathbb R$; $\chi_v^{-1}$ changes $\tau_v,j_v$
to their negatives and leaves $e_v$ unchanged.
At an unramified finite place let $\omega_v=\chi_v(\pi_v)$
for any local uniformizer.  Thus
$L_v(s,\chi_v)=(1-\omega_vq_v^{-s})^{-1}$ and
$\chi_v^{-1}(\pi_v)=\overline{\omega_v}$.
The real and complex local $L$ factors are respectively
$\Gamma_{\mathbb R}(s+i\tau_v+e_v)$ and
$\Gamma_{\mathbb C}(s+i\tau_v+|j_v|/2)$.

Tate's real and complex tables express the archimedean
factors, up to their actual exponential monomial and unit
phase, as
\[
 \frac{\Gamma_{\mathbb R}(1/2+z+i\tau_v+e_v)}
      {\Gamma_{\mathbb R}(1/2-z-i\tau_v+e_v)},
 \qquad
 \frac{\Gamma_{\mathbb C}(1/2+z+i\tau_v+|j_v|/2)}
      {\Gamma_{\mathbb C}(1/2-z-i\tau_v+|j_v|/2)}.
 \tag{MA-GA.18}
\]
Here $e_v\in\{0,1\}$, $j_v\in\mathbb Z$, and
$\tau_v\in\mathbb R$ are the fixed local parameters of
the unitary character.  The finite unramified factors have
the form
\[
 c_v e^{d_vz}
    \frac{1-\overline{\omega_v}q_v^{z-1/2}}
         {1-\omega_v q_v^{-z-1/2}},
 \qquad |c_v|=|\omega_v|=1,\quad d_v\in\mathbb R;
 \tag{MA-GA.19}
\]
the ramified factors are nonzero exponential monomials
with the actual Gauss phase.  Only their analytic bounds
are being read from these forms; (MA-GA.17) retains every
normalizing factor.  The numerator gamma poles lie on or
to the left of $\Re z=-1/2$; the finite denominator poles
lie on that line.  Denominator gamma poles give zeros,
not new poles.  Stirling's formula and its differentiated
forms therefore give
\[
 \begin{gathered}
 m_\chi\text{ holomorphic on }\Re z>-1/2,\qquad
 |m_\chi(it)|=1,\\
 m_\chi^{(j)}(c+it)\text{ has polynomial growth in }t
       \text{ on every fixed closed vertical substrip},\\
 (m_\chi^{-1})^{(j)}(c+it)\text{ has polynomial growth
       on every closed substrip of }-1/2<\Re z<1/2.
 \end{gathered}
 \tag{MA-GA.20}
\]
There is no uniformity claim as $K,S,\chi$ vary.  Polynomial
bounds for each derivative suffice below.

\emph{4. An exact kernel and its two weighted output estimates.}
Write $b_1(x)=B(e^x)$, $v(x)=1-b_1(-x)$ and
$T_v=M_v-U_\chi M_vU_\chi^{-1}$.  Thus $v=1$ far left,
$v=0$ far right, $v'$ is compactly supported, and
$\int v'=-1$.  The bounded identities (MA-GA.16) imply
\[
 Q_{\chi,1}=U_\chi M_{b_1(-x)}U_\chi^{-1}
                =M_{b_1(-x)}+T_v.
 \tag{MA-GA.21}
\]
Let $G(z)=\mathcal Mg(z)$ and
$P_v(z)=\int v'(x)e^{zx}\,dx$.  For $c>-1/2$ set
$H_c=L^2(\mathbb R,e^{2cx}dx)$ and use the unitary Fourier
map $F\mapsto(2\pi)^{-1/2}\mathcal M F(c+it)$ on $H_c$.
For $|c|<1/2$, the actual operator $K_g=C_gT_v$, viewed
from $H_0$ to $H_c$, has kernel
\[
 k_c(t,u)=\frac{G(z)}{2\pi}
       \frac{P_v(z-w)}{z-w}
       \left(\frac{m_\chi(z)}{m_\chi(w)}-1\right),
 \qquad z=c+it,\quad w=iu.
 \tag{MA-GA.22}
\]
We justify both the literal operator and the negative
output lines.  Start with input having smooth compact
Fourier support.  Its image under $U_\chi^{-1}$ is
Schwartz on the central line by (MA-GA.20).  If $c>0$,
multiplication by $v$ maps this vector into the Schwartz
space on line $c$.  The actual weighted Fourier extension
is compatible with the central one.  For compact inputs
this is the local functional equation on the compact
source section; for the present input, cut off in the
physical norm coordinate and pass to the limit in both
central and $c$-weighted Schwartz topologies.  The
polynomial multiplier bounds give convergence in the
weighted topology, and central unitarity identifies the
same local-distribution limit.  This is also the
compatibility in Meyer's local Fourier estimates.  Integration by parts gives
the convergent transform
\[
       \int_{\mathbb R}v(x)e^{(z-w)x}\,dx
                      =-\frac{P_v(z-w)}{z-w}
                     \qquad(c>0).
 \tag{MA-GA.23}
\]
Subtracting the two literal compositions in $T_v$ yields
(MA-GA.22), with the displayed sign.  There is no use of
the Fourier transform of a nonintegrable step as an
ordinary function on the central line.

The apparent singularity at $z=w$ in (MA-GA.22) is
removable: its quotient factor is
$m_\chi'(w)/m_\chi(w)$ there.  The bounds (MA-GA.20)
and the rapid vertical decay of $G(z)$ and $P_v(z-w)$
show that $k_c$ is Schwartz in $(t,u)$, uniformly for
$c$ in every compact subinterval of $(-1/2,1/2)$.
Near $z=w$ use the integral formula for the divided
difference; away from it use the polynomial bounds
directly.  The inequality
$1+|u|\leq(1+|t|)(1+|t-u|)$ shows that the two rapid
decays dominate every joint polynomial.  The same
argument applies after any $t,u$ derivative.

For the compact-Fourier-support inputs just used, the
right side of (MA-GA.22), integrated against the input,
is holomorphic in $z$ on this strip and rapidly decreasing
on every closed vertical substrip.  Shift the inverse
Mellin contour from a positive $c$ to any
$-1/2<c<1/2$.  The horizontal integrals tend to zero by
these bounds, and no pole is crossed.  Thus all these
kernels represent the same actual central output, not
different regularizations.  Density then identifies the
extended operator on all of $H_0$.

A Schwartz kernel on $\mathbb R^2$ is trace class.  One
direct proof expands it in the product Hermite basis:
integration by parts with both harmonic oscillators
gives coefficients decreasing faster than every power
in both indices.  Their absolute sum is finite, giving
a nuclear series of rank-one operators with summable
norms.  The same series permits the integral of the
diagonal, by the corresponding uniform Schwartz bounds.
Applied to (MA-GA.22), this proves
\[
 K_g:H_0\longrightarrow H_c\text{ is trace class}
                  \quad(-1/2<c<1/2).
 \tag{MA-GA.24}
\]
Its trace-class norm on any one line is bounded by
finitely many Schwartz kernel seminorms, hence by
finitely many seminorms of $g$ on a fixed compact
support.  On the central diagonal $P_v(0)=-1$, so
\[
       \operatorname{Tr}_{H_0}K_g
          =-\frac1{2\pi}\int_{\mathbb R}
                 G(it)\frac{m_\chi'(it)}{m_\chi(it)}\,dt
          =\mathcal W_{S,\chi}(g).
 \tag{MA-GA.25}
\]
This establishes the local trace from the actual operator
before any use of its scalar diagonal.  A constant phase
does not alter this trace; an actual exponential monomial
adds a constant multiple of $g(0)$, which is zero at the
asserted limit scope.  Neither is deleted from the operator.

\emph{5. The original finite-cutoff nuclear lift.}
For a fixed annulus $a=a_{\epsilon,X}$, the operator
$M_aC_g$ has the smooth compactly supported kernel
$a(x)g(x-y)$ on $\mathbb R^2$.  It maps $H_0$ into
$C_c^\infty(I)$ for one fixed compact output interval
$I$, continuously in every smooth seminorm.  It is
nuclear into that Fr\'echet space, not only trace class
into $H_0$.  To see the stronger assertion, choose an
interval containing the input support with an interior
margin, expand in its orthonormal Fourier basis in $y$,
and call the resulting compact-output coefficient
functions $v_j(x)$.  Repeated integration by parts in
$y$ gives
\[
       \|v_j\|_{C^k}\leq C_{k,N}(1+|j|)^{-N}
                         \quad\text{for every }k,N.
 \tag{MA-GA.26}
\]
There are positive numbers $d_j$ with $\sum_jd_j<\infty$
such that $\{v_j/d_j\}$ is bounded in every $C^k$ seminorm.
For example, choose successively increasing index blocks
on whose tails the first $r$ seminorms are at most
$2^{-r}(1+|j|)^{-r}$, and on the $r$th block use
$d_j=(1+|j|)^{-r/2}$, enlarging the finitely many initial
values; the block endpoints may be increased to make
the total sum finite.  Each fixed seminorm is then
bounded on all later blocks, as well as on the finitely
many earlier ones.

Thus $M_aC_g=\sum_jd_j(v_j/d_j)\otimes\ell_j$, where
the Fourier coefficient functionals $\ell_j$ have
uniformly bounded norm on $H_0$.  Composing them with
the bounded contraction $Q_{\chi,Y}$ gives one nuclear
series for $T^\chi_{X,Y,\epsilon}(g):H_0\to C_c^\infty(I)$.
The vector family is bounded in every output smooth
seminorm, and the functional bounds are uniform in $Y$
when the annulus and $g$ are fixed.

Compose its vector terms with (MA-GA.15).  They are now
bounded in one original compact, punctured LF stage of
$V_S$, flat at every coordinate hyperplane.  This is
the promised continuous nuclear return
\[
 \mathcal A_{X,Y,\epsilon}:H_\chi\longrightarrow V_S,
       \qquad \mathcal A_{X,Y,\epsilon}
                         =\sigma_\chi T^\chi_{X,Y,\epsilon}(g).
 \tag{MA-GA.27}
\]
The section alone is only asserted continuous; nuclearity
is asserted after this smoothing factor.  Let
$j_\chi:V_S\to H_\chi$ be centered periodization followed
by compact-character projection and the inclusion of
$\mathcal E$ in $H_\chi$.  Then
\[
 j_\chi\mathcal A_{X,Y,\epsilon}
            =T^\chi_{X,Y,\epsilon}(g),\qquad
 \mathcal A_{X,Y,\epsilon}j_\chi
          \text{ is a nuclear endomorphism of }V_S,
\]
\[
 \operatorname{Tr}_{V_S}(\mathcal A_{X,Y,\epsilon}j_\chi)
          =\operatorname{Tr}_{H_\chi}T^\chi_{X,Y,\epsilon}(g).
 \tag{MA-GA.28}
\]
Both traces use exactly the preceding summable series:
the scalar terms agree because $j_\chi\sigma_\chi=1$.
They are independent of its representation.  One can
verify the needed approximation property on the original
source directly: cut off each archimedean log coordinate
as in step 1, project onto a fixed finite-place stage,
and approximate on the resulting punctured compact
coordinate boxes by smooth finite Fourier sums with
a surrounding compact cutoff.  These finite-rank maps
converge to the identity uniformly on bounded sets in
each source seminorm; higher weights in (MA-GA.8)
control the first cutoff error.  This also supplies a
direct finite-rank justification of trace cyclicity.

\emph{6. Exact physical scaling and the uniform tails.}
Put $(W_RF)(x)=F(x+R)$.  Norm reversal by the actual
Fourier transform gives
\[
 W_RQ_{\chi,Y}W_{-R}=Q_{\chi,e^RY},\qquad
 W_{-\log Y}T^\chi_{X,Y,\epsilon}(g)W_{\log Y}
                       =T^\chi_{XY,1,\epsilon Y}(g).
 \tag{MA-GA.29}
\]
This is the same scaling of the two original physical
cutoffs as in MA-LA-014.  Use $U=XY$, $V=\epsilon Y$,
$a=A(e^x/U)-A(e^x/V)$, and (MA-GA.21).  The compact
kernel of $M_aC_gM_{b_1(-x)}$ has trace
$g(0)\int a(x)b_1(-x)\,dx$.  Therefore the exact trace
identity, including its subtraction at the identity, is
\[
 \operatorname{Tr}T^\chi_{X,Y,\epsilon}(g)
     =g(0)\int_{\mathbb R}a(x)b_1(-x)\,dx
                         +\operatorname{Tr}(M_aK_g).
 \tag{MA-GA.30}
\]
The first term vanishes at the stated $g(0)=0$ scope.
No assumption that the support of $g$ avoids zero is made.

Fix $0<\delta<1/2$.  The multiplication map
$1-A(e^x/U):H_\delta\to H_0$ has norm at most
$U^{-\delta}$, since its support is in $x\geq\log U$.
The map $A(e^x/V):H_{-\delta}\to H_0$ has norm at most
$(2V)^\delta$, since its support is in
$x\leq\log(2V)$.  Equations (MA-GA.24--25) and
$1-a=1-A(e^x/U)+A(e^x/V)$ yield
\[
 \bigl|\operatorname{Tr}(M_aK_g)-\operatorname{Tr}K_g\bigr|
       \leq U^{-\delta}\|K_g:H_0\to H_\delta\|_1
            +(2V)^\delta\|K_g:H_0\to H_{-\delta}\|_1.
 \tag{MA-GA.31}
\]
Finally $U\geq2V$ implies $0\leq a\leq1$ and
$\operatorname{supp}a\subset[\log V,\log(2U)]$.
Consequently
\[
      |\operatorname{Tr}(M_aK_g)|
           \leq V^{-\delta}\|K_g:H_0\to H_\delta\|_1.
 \tag{MA-GA.32}
\]
These inequalities prove all the asserted rates and limits.
They also specify exactly the uniformity that is required;
no limiting interchange over compact characters, unit
lattices, fields or finite sets of places has been used.
\end{proof}

\paragraph{Source locators and interpretation.}
Masullo's public draft, privately downloaded for this run,
defines the flat-hyperplane
space in Lemma 5.3.3.3(3), printed/PDF pp.~252--257,
especially (5.3.3.10), and the $S$-local source in
Definition 5.3.4.1, p.~259, (5.3.4.1), with its LF
properties discussed in Lemma 5.3.4.2, pp.~260--261.
The direct Taylor proof (MA-GA.8) is the proof used here
for the needed topology and density estimates.

Meyer (2005), \S\S3.1--3.2, supplies the $S$-unit quotient
and the actual unitary descent; exact supplied-source
locators are \texttt{eq:def\_Fourier}, \texttt{eq:def\_J},
\texttt{eq:def\_EFou}, and
\texttt{lem:Fourier\_unitary\_ICL}, including its full proof.
The weighted extension is Proposition 4.4,
\texttt{pro:Fourier\_estimate}.  The original full source
identification is the complete argument ending at
\texttt{cor:Sch\_local\_coinvariants}, preceded by
\texttt{lem:S\_local\_temperation} and
\texttt{the:Sch\_Adel\_temperation}; the structural condition
is the sufficiently-large-$S$ definition preceding the
global construction.  These locators do not identify
the flat quotient with a different bornology: that
additional point is proved in (MA-GA.8--15).

The analytic local forms are Tate (1950), Theorem 2.4.1,
thesis p.~(2.11), and \S2.5, real p.~(2.14), complex
p.~(2.16), nonarchimedean p.~(2.19), respectively PDF
pages 18, 21, 23 and 26 of the original-thesis scan.
Formula (MA-GA.17) fixes their use to the actual chosen
Fourier convention.  The scalar trace is also Meyer's
general-character local trace formula (20), p.~22;
it was derived here from the actual kernel before being
identified with that local distribution.  In the rational
even sector it gives $A_0(h)+B_S(h)$ of MA-LA-014 for
$g(x)=h(e^x)$; in the sectors of MA-CQ it gives precisely
their ramified and archimedean local terms.  The fixed
compact-character projector has trace one, so no extra
class-number or regulator multiplicity is inserted.

The nuclear comparison is an actual source-to-quotient
comparison with the common pairing (MA-GA.28).  It
does not assert the separate global-to-local identity
FT002, or any positivity or location of its discarded
zeros.  For finite-order characters it retains the same
finite-monodromy sector as MA-CQ; allowing a general
unitary compact character does not assign it an
unconstructed finite Artin motive.

\section{Two literal flat-source estimates and their direct repair}
\label{sec:general-flat-source-scope}

\paragraph{MA-GA-002 (counterexamples to two printed scalar inequalities).}
The literal scalar inequalities (5.3.3.21), printed/PDF
p.~256, and (5.3.3.23), p.~257, of Masullo's public draft
are false under their stated flat-source hypotheses.
This is a statement about those displayed estimates,
not a counterexample to the flat-space topological
identification in Lemma 5.3.3.3(3).  The latter estimates
needed in MA-GA-001 are proved directly by higher-order
Taylor bounds in (MA-GA.8).

\begin{proof}
In (5.3.3.21) take $\alpha=a=1$ and $\beta=0$.
The printed inequality specializes to
\[
 \sup_{x\ne0}(1+|\log|x||)|x|\,|f(x)|
       \leq \sup_x(1+|x|)|x|\,|f(x)|.
 \tag{MA-GA.33}
\]
Absolute values around the function do not weaken the
counterexample, since our functions are nonnegative.
Choose $\eta\in C_c^\infty((-1/4,1/4))$,
$0\leq\eta\leq1$, $\eta(0)=1$.  For $M\geq1$ put
\[
 f_M(x)=
 \begin{cases}
 x^{-1}\eta(\log x+M),&x>0,\\
 0,&x\leq0.
 \end{cases}
 \tag{MA-GA.34}
\]
Each $f_M$ is smooth, compactly supported away from zero,
and hence belongs to $\mathcal S_{00}(\mathbb R)$.
At $x=e^{-M}$ the left side of (MA-GA.33) is at least
$1+M$.  On the entire support, its right side is at most
$1+e^{-M+1/4}$.  Since $M>e^{-M+1/4}$ for $M\geq1$,
the inequality fails.  The ratio of the two sides is
unbounded as $M\to\infty$; thus merely inserting a
universal constant in that same one-seminorm inequality
would not repair it.  Using higher Schwartz derivatives,
as in (MA-GA.8), does repair the required continuity claim.
The failure does not depend on allowing derivative order
zero.  For $\alpha=a=\beta=1$ the same family satisfies
$(x\partial_x)(xf_M)=\eta'(\log x+M)$ and
$x^2f_M'=\eta'(\log x+M)-\eta(\log x+M)$.
At any fixed $t_0$ with $\eta'(t_0)\ne0$, the left
supremum is at least
$(1+M-|t_0|)|\eta'(t_0)|$, while the right one is at
most $(1+e^{-M+1/4})\|\eta'-\eta\|_\infty$.
This also fails for all sufficiently large $M$.

The first inequality in (5.3.3.23) asserts, for $|x|<A$,
\[
 \left|\frac{\varphi(x)}x\right|
       \leq\frac1{2A}\int_{-A}^{A}|\varphi'(y)|\,dy.
 \tag{MA-GA.35}
\]
Again the example is nonnegative at the evaluation points,
so it also disproves the version without the first absolute
value.  There is a single admissible flat function for which
this fails at arbitrarily small $A$.  Choose a nonnegative
smooth bump $\theta$, supported in $(-1/10,1/10)$,
increasing up to its value $\theta(0)=1$ and decreasing
afterward.  Its total variation is $2$.  Set
\[
 x_j=4^{-j},\qquad a_j=e^{-16^j},\qquad
 \varphi(x)=\sum_{j\geq1}a_j
                  \theta\!\left(\frac{x-x_j}{x_j}\right)
       \quad(x>0),\qquad \varphi(x)=0\quad(x\leq0).
 \tag{MA-GA.36}
\]
The supports $[0.9x_j,1.1x_j]$ are disjoint.  On the
$j$th support every derivative of order $k$ is bounded
by $C_k a_jx_j^{-k}$.  For every $N,k$ this is
$o(x_j^N)$ as $j\to\infty$.  The derivative functions
therefore extend continuously by zero at zero, with
all their difference quotients tending to zero.
Induction on $k$ proves that $\varphi$ is smooth and
flat at zero.  It is compactly supported and hence
belongs to $\mathcal S_{00}(\mathbb R)$ and, after
restriction, to the all-weight space
$\mathcal S_{(-\infty,\infty)}(\mathbb R^\times)$
used in that printed proof.

Take $A_j=2x_j$ and evaluate at $x=x_j$.
No support with index less than $j$ meets $[-A_j,A_j]$,
and every support with index at least $j$ is contained
there.  Consequently
\[
 \left|\frac{\varphi(x_j)}{x_j}\right|=\frac{a_j}{x_j},
 \qquad
 \frac1{2A_j}\int_{-A_j}^{A_j}|\varphi'(y)|\,dy
          =\frac1{2x_j}\sum_{m\geq j}a_m
          <\frac{a_j}{x_j}.
 \tag{MA-GA.37}
\]
For the strict inequality, $a_{m+1}/a_m\leq e^{-240}<1/2$
for $m\geq1$, so $\sum_{m>j}a_m<a_j$.
This disproves (MA-GA.35) along $A_j\downarrow0$ for
one fixed admissible function.

The correct estimate uses the segment from zero to the
actual point.  For $x>0$, flatness and the fundamental
theorem of calculus give
\[
 \left|\frac{\varphi(x)}x\right|
       \leq\frac1x\int_0^x|\varphi'(y)|\,dy
       \leq\frac{C_N}{N+1}x^N
       \quad\text{if }|\varphi'(y)|\leq C_Ny^N.
 \tag{MA-GA.38}
\]
The negative direction is identical after reflection.
Repeating this estimate for every derivative gives the
valid flat extension argument used in MA-GA-001.
\end{proof}

The two counterexamples are local source-audit corrections.
They do not alter the actual norm-cutoff theorem or supply
an additional arithmetic zero-exclusion condition.

\section{Character Paley--Wiener ranges and their actual Gram transport}
\label{sec:character-global-pw}

This section treats nonprincipal primitive Dirichlet characters and
$\psi=\psi_{269,3}$ over $\mathbb Q(\zeta_3)$.  It identifies an
actual finite global range, rather than only the local quotient of
MA-CQ.  The principal rational character is MA-GP-001--008.
The finite-vector construction below also works for every nontrivial
finite-order Hecke character of the latter class-number-one field.
All compact character measures have mass one, as in MA-CH;
Masullo's other Haar normalization requires the rescaling recorded
in (MA-CH.5).

Put $K=\mathbb Q$ or $\mathbb Q(\zeta_3)$, and let $\eta$ be
one of the stated nontrivial primitive characters.  Write
\[
 d=[K:\mathbb Q]\in\{1,2\},\quad
 w_K=|\mathcal O_K^\times|\in\{2,6\},\quad
 Q=|D_K|N\mathfrak f_\eta.
 \tag{MA-CGP.1}
\]
At infinity use parity $e$ when $K=\mathbb Q$, with
$\eta(-1)=(-1)^e$, and radial functions when $K$ is complex.
Let $\mathscr S_\eta$ denote that closed subspace of the
archimedean Schwartz space.  The real Fourier measure is $dx$
and the complex one is $2\,dx\,dy$.  For global Poisson we
explicitly choose the coherent adelic additive character:
at infinity it is $\exp(2\pi i\operatorname{Tr}_{K_\infty/\mathbb R}x)$,
and at finite places it is the negative fractional-trace
character.  Its product is trivial on $K$.  The Fourier
kernel is its complex conjugate.  All Gauss integrals in
the MA-CQ formulas used below are evaluated with this choice.
This does not presume coherence of independently printed
positive-sign local examples in the source.  Conjugating a
finite local additive character reflects that Fourier
coordinate; it preserves its norm support and its value at
zero.  The source's invariant PW support and polar conditions
are therefore unchanged by this explicit normalization.
A Fourier transform changes $\eta$ to
$\bar\eta$; their archimedean parity or radial type is the same.
For an integral ideal put $c(\mathfrak a)=\eta(\mathfrak a)$,
extended by zero on ideals meeting the conductor.

\paragraph{MA-CGP-001 (the primitive theta model and its exact reflected scale).}
For $f\in\mathscr S_\eta$ define
\[
 (\Theta_\eta f)(t)
    =w_K\sqrt t\sum_{\mathfrak a\ne0}
            c(\mathfrak a)f\bigl((N\mathfrak a\,t)^{1/d}\bigr),
                   \qquad t>0.
 \tag{MA-CGP.2}
\]
The values on the positive real ray determine the parity or
radial function.  This is the actual centered, fixed-character
adelic theta image, and
\[
       \Theta_\eta:\mathscr S_\eta
          \xrightarrow{\ \sim\ }\mathcal I_\eta^{\rm cen}
 \tag{MA-CGP.3}
\]
is a topological isomorphism onto the actual smooth ideal
of MT-M22/MT-M24, centered at $1/2$.
There is an explicitly specified phase $\epsilon_\eta^\theta$
of modulus one such that
\[
 \Theta_\eta f(t)
      =\epsilon_\eta^\theta\,
           \Theta_{\bar\eta}(\mathcal F_\infty f)(1/(Qt)).
 \tag{MA-CGP.4}
\]
It is a theta Fourier phase, not an unspecified completed
$L$-function root-number convention.

\begin{proof}
\emph{The actual finite vector and its Fourier transform.}
Use the idele character $\beta$ and local notation
$b_v,\kappa_v,f_v,d_v$ of MA-CQ, with
$\kappa_v=\beta_v^{-1}|_{\mathcal O_v^\times}$.
Take the primitive finite vector
\[
 \varphi_{\eta,v}=
   \begin{cases}
     \mathbf1_{\mathcal O_v},&f_v=0,\\
     \kappa_v\mathbf1_{\mathcal O_v^\times},&f_v>0.
   \end{cases}
 \tag{MA-CGP.5}
\]
A nontrivial character in either field has at least one ramified
finite place: a finite-order character unramified everywhere
would factor through the trivial ideal class group, and its
archimedean connected component is trivial.  At such a place
both the value and the integral of the finite vector vanish.
Consequently $f\otimes\varphi_{\eta,\rm fin}$ satisfies
both adelic polar conditions for every $f\in\mathscr S_\eta$;
no two archimedean moment conditions are imposed.

Choose the positive prime generators over $\mathbb Q$, or
arbitrary prime generators over $\mathbb Q(\zeta_3)$ as in
MA-CQ, and put
\[
 m_v=f_v+d_v,\qquad
 \delta=\prod_v\pi_v^{m_v},\qquad |\delta|_\infty=Q,
\]
\[
 \epsilon_\eta^\theta
      =\prod_{f_v>0}\epsilon_v(\kappa_v)
                                      \prod_v b_v^{m_v}.
 \tag{MA-CGP.6}
\]
The products are finite.  At an unramified place the Fourier
transform of $\mathbf1_{\mathcal O_v}$ is
$q_v^{-d_v/2}\mathbf1_{\pi_v^{-d_v}\mathcal O_v}$.
At a ramified place MA-CQ-001 gives the single shell with
constant $q_v^{-m_v/2}\epsilon_v(\kappa_v)$.
When the separate local dilations $\pi_v^{m_v}$ are replaced
by the one global dilation $\delta$, the other unit components
contribute $\prod_v b_v^{m_v}$, by the principal-idele relation
(MA-CQ.14).  Thus, as an equality of actual finite functions,
\[
 \widehat\varphi_{\eta,\rm fin}(y)
       =Q^{-1/2}\epsilon_\eta^\theta
                                \varphi_{\bar\eta,\rm fin}(\delta y).
 \tag{MA-CGP.7}
\]
The phases in (MA-CGP.6) have modulus one by MA-CQ-001.
Their dual product is $(-1)^e$ in the real case and one in
the complex case; this agrees with $\mathcal F_\infty^2$.
In particular the odd real case has not lost its Fourier sign.

The finite support in (MA-CGP.5) restricts the adelic sum to
integral elements prime to the conductor.  By the principal
character relation its value on such an element $a$ is
$\beta_\infty(a)c((a))$.  Thus it is $c((a))$ on the
positive generator over $\mathbb Q$ and on every generator
in the complex case.  There are $w_K$ generators of each integral ideal.
Their archimedean parity or radial factors compensate the unit
character, so their contributions agree.  This proves that
(MA-CGP.2) is the actual full adelic theta sum on the norm
splitting, with the compact character projected as in MA-CH.

Adelic Poisson summation, with both polar terms zero, gives
$t^{-1}$ times the sum of the Fourier transforms at archimedean
argument $t^{-1/d}a$.  Insert (MA-CGP.7) and replace
$a$ by $\delta^{-1}a$.  At infinity $\delta$ is positive
in the real case; its rotation is immaterial in the complex
radial case.  Multiplying by $\sqrt t$ gives precisely
(MA-CGP.4), including its argument $1/(Qt)$ and constant.
This also gives continuity into exponential Schwartz space:
for large $t$ use Schwartz decay in (MA-CGP.2), and for small
$t$ use (MA-CGP.4).  Ideal counting is $O(T)$ in these two
fields, by interval counting or the rank-two integral lattice;
these estimates also apply after every norm derivative.

\emph{Onto the specified smooth image.}
Here the class-number-one restriction is used essentially.
Every norm-one idele class has a representative in
$K_\infty^1\times\prod_v\mathcal O_v^\times$; its kernel
is the finite global unit group.  Compact character averaging
can therefore be performed on this compact group before
summation.  At infinity it gives exactly the stated parity or
radial Schwartz function.  At a ramified finite place it gives
a finite sum of valuation-shell dilates of
$\kappa_v\mathbf1_{\mathcal O_v^\times}$: compact support
bounds valuations below, and nontrivial unit type together
with local constancy forces vanishing near zero.
At an unramified place the averaged function is constant
on unit shells and constant sufficiently near zero.  It is
a finite linear combination of dilates of
$\mathbf1_{\mathcal O_v}$.  Only finitely many places vary.

A finite-place dilation by the global prime generator
$\pi_v^j$ is moved through the theta sum by reindexing its
principal element.  The other finite unit components supply
the scalar $b_v^j$, and the archimedean factor becomes
$f(\pi_v^j\,\cdot)$, still in $\mathscr S_\eta$.
Thus every projected adelic Schwartz image is a finite linear
combination of images (MA-CGP.2), and hence is one such image.
MT-M24 (M24.7)--(M24.9) identifies this actual image, already
closed, with $\mathcal I_\eta^{\rm cen}$.

Injectivity follows without a spectral assertion.  Put
$F_f(t)=w_K\sqrt t f(t^{1/d})$.
Then (MA-CGP.2) is exactly the ideal sum in (MA-CH.12).
On every compact positive interval its inverse is
\[
 F_f(t)=\sum_{\mathfrak a\ne0}
      \frac{\mu_K(\mathfrak a)c(\mathfrak a)}{\sqrt{N\mathfrak a}}
                   (\Theta_\eta f)((N\mathfrak a)t).
 \tag{MA-CGP.8}
\]
Schwartz decay at infinity and the ideal count make the
series and the substituted double series absolutely convergent,
with all derivatives, on such intervals.  The ideal divisor
sum $\sum_{\mathfrak b\mid\mathfrak a}\mu_K(\mathfrak b)$
is one only for the unit ideal, proving (MA-CGP.8).
Both spaces in (MA-CGP.3) are Fr\'echet, and the target is
closed by MT-M24.  The continuous bijection just proved is
therefore a topological isomorphism by the open mapping theorem.
\end{proof}

\paragraph{MA-CGP-002 (exact finite source support and its Hardy divisor).}
For $\Lambda>0$ let
\[
 \mathscr P_{\eta,\Lambda}
  =\{\Theta_\eta f:f\in\mathscr S_\eta,
          \ \operatorname{supp}\widehat f
                     \subset\{|\xi|_\infty\le\Lambda\}\}.
 \tag{MA-CGP.9}
\]
Then, as an equality of actual smooth spaces,
\[
 \mathscr P_{\eta,\Lambda}
   =\{g\in\mathcal I_\eta^{\rm cen}:
                              \operatorname{supp}g\subset[1/(Q\Lambda),\infty)\}.
 \tag{MA-CGP.10}
\]
This level has an exact source interpretation.  Let $S$ contain
infinity, the conductor and the different, and take $S'=S$
in Masullo Definition~7.3.2.3, pp.~354--355 (thus no additional
unit-shell Euler-deletion transition).  Extend its $S$-local
Schwartz data by $\mathbf1_{\mathcal O_v}$ outside $S$,
apply the actual global character theta map, and use the
invariant Fourier support bound $\prod_{v\in S}|\xi_v|_v\le Y$.
The resulting smooth global level is exactly
\[
                  \mathscr P_{\eta,Y/Q}.
 \tag{MA-CGP.11}
\]
Its Hilbert closure is therefore the actual corresponding
orthogonal range of Definition~11.2.1.3, pp.~488--489,
with this specified choice of finite-level data.

For $0<\alpha<1/2$, polynomial weight $s=0$, and
$c=-\log(Q\Lambda)$, use the equivalent half-line norm
$L^2([c,\infty),e^{2\alpha x}dx)$ and
$(U_cg)(y)=e^{\alpha(c+y)}g(e^{c+y})$.
The Laplace transform is unscaled, with boundary norm
$(2\pi)^{-1}\int|F(it)|^2dt$.
The Hilbert closure $J_{\eta,\Lambda}$ then satisfies
\[
 \begin{aligned}
 \mathcal L U_cJ_{\eta,\Lambda}
       &=B_{\eta,\alpha}H^2(\operatorname{Re}z>0),\\
 Z(B_{\eta,\alpha})&=
   \{\alpha+\tfrac12-\rho_\eta:
                         \operatorname{Re}\rho_\eta<\alpha+\tfrac12\},
 \end{aligned}
 \tag{MA-CGP.12}
\]
with full completed-zero multiplicities.  The inner function is
pure Blaschke: no singular inner factor or positive delay remains.
This does not assume RH for the character.

\begin{proof}
Equation (MA-CGP.4) proves the forward support inclusion in
(MA-CGP.10).  Conversely, apply it and then (MA-CGP.8) in
the dual sector.  If $g$ vanishes below $1/(Q\Lambda)$,
then $\Theta_{\bar\eta}\widehat f(t)=0$ for $t>\Lambda$;
every term of its M\"obius inverse also vanishes there.
Hence $\widehat f(t^{1/d})=0$ for $t>\Lambda$, and parity
or radiality gives the asserted full band support.
The band condition is closed in Schwartz space, so (MA-CGP.3)
also proves that this smooth range is closed.

For (MA-CGP.11), the finite Fourier support of the canonical
vector has maximum norm product $\prod_v q_v^{f_v+d_v}=Q$.
Thus the frequency band $Y/Q$ supplies actual $S$-local
Schwartz data satisfying the source's product support bound
$Y$.  This proves one inclusion, with the stated normalization.
Conversely use global Poisson on arbitrary such source data.
Outside $S$ the Fourier factors are still $\mathbf1_{\mathcal O_v}$,
since all different places are in $S$.
For a nonzero principal element $a$ contributing to the reflected
sum, $|a|_v\le1$ outside $S$, and the product formula gives
$|a|_S\ge1$.  At norm $t<1/Y$, its reflected $S$-argument
has norm at least $1/t>Y$, so every summand is zero.
Compact character averaging preserves this statement: the
class-number-one compact representatives above have norm one
in $S$ and unit components outside it.
The actual projected image therefore lies in the smooth ideal
and is supported in $[1/Y,\infty)$.  Equation (MA-CGP.10)
gives the reverse inclusion.  Smooth closure preserves the
same support condition, proving (MA-CGP.11) at the source's
closed-image level as well.

We give the Hardy normalization and common divisor explicitly.
Right translation on the support half-line preserves the smooth
ideal and (MA-CGP.10).  Under $U_c$ it is a scalar times the
unilateral translation.  The Laplace image of the closure is
therefore a nonzero invariant subspace $BH^2$, by the same
vectorwise resolvent average and Beurling theorem proved in
MA-GP-003.
For a frequency test $u=\widehat f$ supported in a positive
annulus inside the band, put $U(r)=u(r^{1/d})$ and
$s_z=z+1/2-\alpha$.  Direct Mellin integration, first where
it converges absolutely, gives
\[
 (\mathcal L U_c\Theta_\eta f)(z)
   =w_K\epsilon_\eta^\theta e^{zc}Q^{z-\alpha}
          L_K(s_z,\bar\eta)
                   \int_0^\Lambda U(r)r^{s_z-1}\,dr.
 \tag{MA-CGP.13}
\]
Analytic continuation gives the equality on the whole right
half-plane.  There $\operatorname{Re}s_z>1/2-\alpha>0$,
so the gamma factors have no poles, and the nontrivial finite
$L$-function is entire.  No polar cancellation or moment-zero
constraint is required.  At every prescribed $s_z$ the final
integral can be made nonzero by an annular test.  Thus the
common zeros are exactly those of $L_K(s_z,\bar\eta)$,
to their exact orders.  The completed functional equation
$\Lambda_\eta(s)=W_\eta\Lambda_{\bar\eta}(1-s)$ changes
their indexing to precisely (MA-CGP.12).  This explains why
the dual character in (MA-CGP.13) is essential.

Every core output is flat at $c$ and decreases faster than
all exponentials at $+\infty$.  Its Laplace transform is
entire and, by repeated integration by parts, decreases faster
than every power uniformly on the closed right half-plane.
The disk image used in MA-GP-003 is consequently bounded
and analytically extends across every boundary point other
than infinity.  The same elementary singular-measure argument
there excludes finite boundary atoms and all finite singular
support; only a delay $e^{-az}$ at infinity could remain.
To exclude it, choose $u$ positive on the positive ray arbitrarily
close to the band boundary, flat there and zero near the origin,
with the required parity or radial extension.  For $t$ just
above $1/(Q\Lambda)$ only the unit ideal contributes to the
reflected formula (MA-CGP.4).  Its value is a nonzero constant
phase times that positive test.  The range has no positive
delay.  This proves the pure Blaschke assertion, as well as
Blaschke summability through its actual Hardy realization.
\end{proof}

The restriction $S'=S$ in (MA-CGP.11) is substantive.  If
additional unit-shell transition factors are inserted at
unramified places, their finite Euler multipliers and changed
levels must be retained; this theorem does not identify those
Euler-deleted ranges with the primitive range.  The fixed-base
levels just proved already form a dense increasing band exhaustion of
the full primitive character image.  The absolute hypercube
levels contain every one of its canonical tests once their
level exceeds the finite coordinate bounds and the
archimedean band, so their union is dense as well.

\paragraph{MA-CGP-003 (all-weight condensation and a concrete character Gram kernel).}
For every $\alpha>0$ and $s\ge0$, the orthogonal projections
$Q_{\eta,\Lambda}$ onto the closures of (MA-CGP.9) increase
strongly to the actual geometric projection $P_{G,\eta}$.
For $A=M_\chi C_h$ with compact smooth norm tests,
\[
       \|AQ_{\eta,\Lambda}-AP_{G,\eta}\|_1\longrightarrow0.
 \tag{MA-CGP.14}
\]
The actual half-line Gram formula of MA-GP-004 applies to
these ranges at every polynomial weight; for $s=0$ its
inner multiplier is the explicitly identified $B_{\eta,\alpha}$.

There is also an explicit frequency Gram transport.  Write
\[
 a_n=\sum_{N\mathfrak a=n}\overline{c(\mathfrak a)},\qquad
 \gamma_K=\begin{cases}2,&K=\mathbb Q,\\2\pi,&K=\mathbb Q(\zeta_3),\end{cases}
\]
\[
 \mathcal P_{\eta,\Lambda}=L^2((0,\Lambda),\gamma_K\,dy).
 \tag{MA-CGP.15}
\]
This is precisely the auxiliary ordinary archimedean band
norm on $U(y)=u(y^{1/d})$; no weighted local norm is implicit.
On the dense smooth parity or radial band core, including
functions nonzero at the origin when the type permits them,
Poisson gives
\[
 (T_{\eta,\Lambda}U)(t)
   =\frac{w_K\epsilon_\eta^\theta}{\sqrt{Qt}}
                      \sum_{n\ge1}a_n U(n/(Qt)).
 \tag{MA-CGP.16}
\]
It is closable and injective after closure, and its range
closure is $J_{\eta,\Lambda}$.
For the norm (MA-GP.52) and $A=M_\chi C_h$, its core form
has the bounded trace-class extension
\[
 \begin{aligned}
 (K_{\eta,\chi,h}U)(y)
   ={}&\frac{w_K^2}{\gamma_K}
      \sum_{m,n\ge1}\frac{\overline{a_m}a_n}{m}
       \chi\!\left(\frac{m}{Qy}\right)
       w_{\alpha,s}\!\left(\log\frac{m}{Qy}\right)\\
      &\hspace{12mm}\times
       \int h(r)r^{1/2}U(nry/m)\,\frac{dr}{r}.
 \end{aligned}
 \tag{MA-CGP.17}
\]
For $\operatorname{supp}\chi\subset[a,b]$ and
$\operatorname{supp}h\subset[r_-,r_+]$, only
$m\le Q\Lambda b$ and $n\le Q\Lambda b/r_-$ contribute.
Let $G_\eta=\overline T_{\eta,\Lambda}^{\,*}\overline T_{\eta,\Lambda}$.
Then the actual fixed-level trace is
\[
 \operatorname{Tr}_H(AQ_{\eta,\Lambda})
    =\lim_{\varepsilon\downarrow0}
       \operatorname{Tr}_{\mathcal P_{\eta,\Lambda}}
       \bigl((G_\eta+\varepsilon)^{-1/2}
                K_{\eta,\chi,h}(G_\eta+\varepsilon)^{-1/2}\bigr),
 \tag{MA-CGP.18}
\]
with the same two actual trace-norm limits as (MA-GP.68).
It is a Hilbert trace transport; representation intertwining
or an LF lift of its Gram resolvents is not asserted.

\begin{proof}
Smooth frequency cutoffs preserving parity or radiality
approximate every $f\in\mathscr S_\eta$ in Schwartz topology.
Both adelic polar conditions remain automatic from the fixed
nontrivial finite vector.  Equations (MA-CGP.3) and (MA-CGP.9)
give density of the increasing union in the actual smooth
ideal and in its weighted closure.  Strong projection convergence
follows.  Since $A$ itself is trace class by MA-LIFT-003,
finite-rank approximation and self-adjoint strong convergence
prove (MA-CGP.14), for all the stated weights.
The equivalent half-line norm has the same ratio
$R(x)=e^{2\alpha|x|-2\alpha x}$ as (MA-GP.40)--(MA-GP.44);
its positive bounded inverse proof applies without change.

The source norm in (MA-CGP.15) follows directly from
$2\int_0^\infty|u(r)|^2dr$ on a real parity sector and
$4\pi\int_0^\infty|u(r)|^2r\,dr$ on the complex radial
sector, substituting $y=r^d$.
The dense core consists of actual smooth archimedean functions
with frequency support in the open band.  It need not vanish
near the origin.  Dilation into a slightly smaller band
shows that this core gives the same range closure as the
closed-band definition (MA-CGP.9).
Every compact output localization of (MA-CGP.16) has only
finitely many terms and is bounded on this auxiliary source,
by Cauchy--Schwarz and the substitutions $y=n/(Qt)$.
This proves closability exactly as in MA-GP-005.
The finite inverse coefficients
$b_n=\sum_{N\mathfrak a=n}\mu_K(\mathfrak a)
                              \overline{c(\mathfrak a)}$
satisfy $\sum_{d\mid n}a_db_{n/d}=\mathbf1_{n=1}$.
Applied to the locally finite Poisson sum on a graph-domain
vector, they prove injectivity of the closure, with the same
null-set justification as (MA-GP.72).

For the form calculation, the squared scalar in (MA-CGP.16)
is $w_K^2/Q$.  In the pairing with the $m$-term put
$y=m/(Qt)$; the measure $t^{-2}dt$ supplies $Q\,dy/m$.
Dividing by the source density $\gamma_K\,dy$ gives
exactly $w_K^2/(\gamma_K m)$ and the conjugate coefficient
$\overline{a_m}a_n$ in (MA-CGP.17).
After $z=nry/m$, each kernel is a bounded weight multiplier
times
\[
       \frac{w_K^2\overline{a_m}a_n}
                    {\gamma_K\sqrt{mn}\sqrt{yz}}
                              h\!\left(\frac{mz}{ny}\right).
 \tag{MA-CGP.19}
\]
Its support is away from both coordinate origins.  The finite
smooth-kernel trace-class argument of MA-GP-008 applies,
without a zero-mean projection: no such condition belongs
to this nontrivial finite-vector source.
Graph-core approximation extends the form identity to
$\operatorname{Dom}\overline T$ in both variables.
Spectral calculus then maps the whole source into that domain
under $(G_\eta+\varepsilon)^{-1/2}$ and gives a contraction
$F_\varepsilon=\overline T(G_\eta+\varepsilon)^{-1/2}$.
It and its adjoint converge strongly to the actual polar
isometry and its adjoint.  Applying the proved trace-ideal
continuity and cross-space cyclicity from MA-GP-008 gives
(MA-CGP.18), including both trace-norm limits.  These are
actual bounded maps, not formal products with an unspecified
adjoint domain.
\end{proof}

\paragraph{The two requested character families and the remaining asymptotic.}
For primitive Dirichlet $\chi$ the scale is $Q=q$, and the
Hardy zeros in (MA-CGP.12) are the zeros of its own primitive
completed function, with parity retained throughout Fourier
transformation.  For $\psi_{269,3}$ the scale is
$Q=3\cdot269^2$, with $\delta=(1-\zeta_3)269$ and
$b_3=b_{269}=1$.  Thus (MA-CGP.6) uses the Gauss phase of
$\psi^{-1}$ on the ramified unit shell; the quotient cutoff
in MA-CQ-004 uses the dual phase, as it should.
In both cases the source invariant level $Y$ has the same
support edge $-\log Y$, after the exact primitive band
conversion $\Lambda=Y/Q$.

The fixed-character global orthogonal range, its common Hardy
divisor, its bounded half-line Gram and its explicit local cyclic
operator are now specified.  This reduces task~1(d)'s remaining
global/local issue to an asymptotic comparison of that actual
Gram-corrected operator with the separately constructed local
character cutoff of MA-CQ.  Its Gauss and conductor factors
cannot be dropped.  Nothing here changes the all-weight
retained multiplicities of MA-WP or supplies an independent
narrow-strip return.  Dirichlet RH is the stated scope of $A$;
the cubic character is outside that scope and has the distinct
single-character Hecke RH question already identified in MA-CH.

\paragraph{Source and Grothendieck scope.}
MT-M22 (M22.1)--(M22.10) and MT-M24, especially
(M24.2), (M24.6)--(M24.9), were read completely for the actual
character, Poisson and closed-image identifications used here.
The mixed support definitions and proof were read at Masullo
pp.~350--357, especially Definition~7.3.2.3 and
Proposition~7.3.2.4; the Hilbert projection definition was
read at pp.~487--489.  The exact $S'=S$ specialization is
stated above instead of discarding the source's extra Euler
transition factors.  The finite Fourier calculations are the
independently checked MA-CQ-001--002.  The Hardy argument
uses the Beurling/Burnol source scope already given in
MA-GP-003, with every new character and conductor factor
derived here.

The rank-one finite-monodromy character and its induced Artin
object are those of MA-CH and MA-CQ.  The finite-length
$\mathbb C[z]$ zero modules retain the full common-zero
multiplicities before Hilbert truncation; a Gram normalization
changes the analytic metric, not that algebraic length.
The actual trace-class cyclic map is the ordinary Hilbert
counterpart of the inherited cote~1 trace/determinant programme,
BR-TR-01/02, scans~079--080 and~111--112.  Those archive
locators are inherited, not newly inspected here.  No
arithmetic positivity, half-integral Tate object, or new
$\ell$-adic weight realization is claimed.

\section{The actual character PW trace and its two orders of limits}
\label{sec:character-pw-trace}

The character range theorem MA-CGP-002 permits the finite-level
operator calculation of MA-GP-009 to be made on the two requested
character families.  The character and conductor enter through
an actual range and its orthogonal projection.  In particular,
substituting a character into a numerical zero sum is not the
argument for trace class or for the cutoff limits below.

\paragraph{MA-CGP-004 (the character finite-level trace and its order defect).}
Let $\eta$ be a nonprincipal primitive Dirichlet character, or
$\psi_{269,3}$ over $\mathbb Q(\zeta_3)$; the same proof applies
to the other nontrivial finite-order characters covered by
MA-CGP-001.  Fix $0<\alpha<1/2$, polynomial weight $s=0$,
and the normalized character-sector Hilbert space
\[
 H=L^2(\mathbb R,e^{2\alpha|x|}\,dx),\qquad
 Q=|D_K|N\mathfrak f_\eta.
 \tag{MA-CGP.20}
\]
For a primitive archimedean norm-band $\Lambda>0$ put
\[
 \begin{aligned}
 c&=-\log(Q\Lambda),& H_c&=\mathbf1_{[c,\infty)}H,\\
 \mathcal Q_{\eta,c}&=H_c/J_{\eta,\Lambda},&
 Z_{\eta,c}&=H_c\ominus J_{\eta,\Lambda}.
 \end{aligned}
 \tag{MA-CGP.21}
\]
The range $J_{\eta,\Lambda}$ is the actual one in MA-CGP-002,
and $Q_{\eta,\Lambda}$ is its orthogonal projection in $H$.
Let $P_{\eta,c}^Z$ project onto $Z_{\eta,c}$, extended by
zero on $H_c^\perp$.  For every additive
$g\in C_c^\infty((0,\infty))$, set
$G_0(\lambda)=\int g(t)e^{\lambda t}\,dt$.
Then
\[
 \begin{aligned}
 C_gP_{\eta,c}^Z&\in\mathcal S_1(H),\qquad
 C_g^{\mathcal Q_{\eta,c}}\in\mathcal S_1(\mathcal Q_{\eta,c}),\\
 \operatorname{Tr}_H(C_gP_{\eta,c}^Z)
  &=\operatorname{Tr}_{\mathcal Q_{\eta,c}}C_g^{\mathcal Q_{\eta,c}}
    =\sum_{\operatorname{Re}\rho_\eta<1/2+\alpha}
             m_\eta(\rho_\eta)G_0(\rho_\eta-1/2).
 \end{aligned}
 \tag{MA-CGP.22}
\]
This absolutely convergent trace is independent of $\Lambda$.

For the uniformly bounded compact smooth exhaustion
$\chi_R(x)=\chi_0(x/R)$, $\chi_0=1$ near zero, the actual
localized traces have the two iterated limits
\[
 \begin{aligned}
 \lim_{\Lambda\to\infty}\lim_{R\to\infty}
    \operatorname{Tr}(M_{\chi_R}C_gQ_{\eta,\Lambda})
   &=-\sum_{\operatorname{Re}\rho_\eta<1/2+\alpha}
                     m_\eta(\rho_\eta)G_0(\rho_\eta-1/2),\\
 \lim_{R\to\infty}\lim_{\Lambda\to\infty}
    \operatorname{Tr}(M_{\chi_R}C_gQ_{\eta,\Lambda})
   &=-\sum_{1/2-\alpha<\operatorname{Re}\rho_\eta<1/2+\alpha}
                     m_\eta(\rho_\eta)G_0(\rho_\eta-1/2).
 \end{aligned}
 \tag{MA-CGP.23}
\]
Their difference, in this order, is exactly
\[
       -\sum_{\operatorname{Re}\rho_\eta\le1/2-\alpha}
                    m_\eta(\rho_\eta)G_0(\rho_\eta-1/2).
 \tag{MA-CGP.24}
\]
The lower boundary is included.  No exterior or boundary zero
is asserted to exist, and no nonzero difference is asserted
for an arbitrary particular test.

\begin{proof}
\emph{The actual character quotient, including the dual indexing.}
Use
\[
 (U_cf)(y)=e^{\alpha(c+y)}f(c+y),\qquad
 R_c(y)=e^{2\alpha|c+y|-2\alpha(c+y)},\qquad y\ge0.
 \tag{MA-CGP.25}
\]
It is unitary from $H_c$ to $L^2((0,\infty),R_c(y)dy)$.
The finite-level norm ratio lies between one and
$e^{4\alpha\max(-c,0)}$.  MA-CGP-002 identifies the actual
closed ideal, in the equivalent ordinary half-line norm,
with the inverse Laplace image of
$B_{\eta,\alpha}H^2(\operatorname{Re}z>0)$.
Its model-space zeros are
\[
             z_j=\alpha+\tfrac12-\rho_{\eta,j}=a_j+ib_j,
                              \qquad a_j>0.
 \tag{MA-CGP.26}
\]
The precise source calculation behind this indexing is
(MA-CGP.13): the reflected Fourier data carries
$L_K(z+1/2-\alpha,\bar\eta)$.
The completed functional equation changes its zero $1-\rho_\eta$
to the displayed $z_j$.  Hence it is the divisor of $\eta$,
not an accidentally substituted dual divisor, that occurs
in (MA-CGP.22).
The scalar theta phase and the conductor factor in that
Laplace formula do not add zeros.

Right translations preserve the actual supported ideal.
On the equivalent half-line space they are scalar multiples
of the unilateral semigroup $S_t$ of MA-GP-009.
The actual quotient is therefore boundedly similar to the
Hardy model-space quotient, retaining the original weighted
orthogonal projection as a separate operator.
Put $k(t)=e^{\alpha t}g(t)$ and
\[
        \mathcal G(z)=\int_0^\infty k(t)e^{-zt}\,dt
                                      =G_0(\alpha-z).
 \tag{MA-CGP.27}
\]
Actual quotient convolution is similar to
$T_{\mathcal G}=P_{K_B}M_{\mathcal G}|_{K_B}$,
where $K_B=H^2\ominus B_{\eta,\alpha}H^2$.
The support of $g$ inside $(0,\infty)$ is what makes this
an actual causal action on $H_c$ and its ideal.

\emph{Trace class before evaluating the trace.}
The containing critical strip for this fixed Hecke function
and its unconditional zero count give
\[
       0<a_j\le\alpha+1/2<1,
                  \qquad\sum_j(1+b_j^2)^{-1}<\infty.
 \tag{MA-CGP.28}
\]
These are the fixed-character versions of the NT-Z10 input
already used in MT-M18/MT-M22; they are not RH statements.
No uniform separation from the upper strip edge is assumed.
The Malmquist basis, its completeness, and its boundary
Poisson norm calculation are exactly (MA-GP.89)--(MA-GP.92):
for any bounded analytic $F$ with
$|F(it)|\le C/(1+|t|)$,
\[
                  \|M_Fe_j\|^2\le C_F(1+b_j^2)^{-1}.
 \tag{MA-CGP.29}
\]
The proof uses only the upper bound on $a_j$ in (MA-CGP.28),
the normalized Poisson mass, and $\int|F(it)|^2dt<\infty$;
it does not use a symmetry of the zero set under conjugation.
Thus the multiplier and its compression are Hilbert--Schmidt.
Since $\mathcal G$ decreases faster than every power on the
closed right half-plane, both
$r(z)=(\kappa+z)^{-1}$ and
$F_2(z)=(\kappa+z)\mathcal G(z)$ satisfy those hypotheses.
The quotient multiplier homomorphism gives
$T_{\mathcal G}=T_rT_{F_2}\in\mathcal S_1$ and
$T_r^2\in\mathcal S_1$.
In this actual orthonormal basis, the analytic Poisson formula
then gives
\[
                \operatorname{Tr}T_{\mathcal G}
                   =\sum_j\mathcal G(z_j)
                   =\sum_jG_0(\rho_{\eta,j}-1/2).
 \tag{MA-CGP.30}
\]
The diagonal trace sum is absolute by trace class, and its
zero-series form is also absolute by the fixed-character
zero count and rapid Mellin decay.
Bounded similarity transfers this quotient trace to the
specified quotient norm.

\emph{The actual ambient lift.}
The quotient semigroup resolvent is $T_r$ in the Hardy model;
its square is trace class, as just proved.
The adjoint semigroup on the actual orthogonal representative
has generator
$\mathcal L_Ru=R^{-1}(Ru)'$, with its genuine half-line
graph domain and no artificial boundary condition at zero.
All the hypotheses of the actual half-line lift calculation
(MA-GP.96)--(MA-GP.105) are now verified for this quotient.
That calculation depends only on this resolvent square,
$R=R_c$, and the compact test $k$, not on the zeta divisor.
It therefore applies to $Z_{\eta,c}$ without changing its
endpoint or cusp terms.
Explicitly, if $a=R^{-1}$ and $v=Ru$, they are
\[
 \begin{aligned}
 &[a(0)(2\kappa k-k')+a'(0+)k]v(0)-a(0)k\,v'(0),\\
 &-4\alpha\,k(\,\cdot-L)v(L)\quad(c<0,\ L=-c).
 \end{aligned}
 \tag{MA-CGP.31}
\]
The second term is absent for $c\ge0$.
The bulk is a bounded operator times the trace-class
resolvent square; these displayed terms are finite rank,
with bounded graph-domain trace functionals.
Thus the actual $C_gP_{\eta,c}^Z$ is trace class on $H$.
Trace cyclicity with its actual orthogonal projection
identifies its trace with quotient convolution, proving
(MA-CGP.22).  This argument has not replaced the weighted
projection by an ordinary Hardy projection.

\emph{The two limits of actual cutoff operators.}
The exact orthogonal decomposition is
$Q_{\eta,\Lambda}=P_c-P_{\eta,c}^Z$, with
$P_c=\mathbf1_{[c,\infty)}$.
The compact kernel argument and bounded smooth approximation
of $P_c$ give
$\operatorname{Tr}(M_{\chi_R}C_gP_c)=0$, since $g(0)=0$.
For fixed $\Lambda$, strong convergence of $M_{\chi_R}$
paired with the trace-class operator $C_gP_{\eta,c}^Z$
therefore yields the first inner limit in (MA-CGP.23).
It is already independent of $\Lambda$ by (MA-CGP.22).

For the other order, MA-CGP-003 proves actual trace-norm
convergence of the finite PW condensation at fixed $R$.
The all-weight actual lift and cutoff theorem MA-WP then
give the trace on the completed full-strip quotient for
this primitive character.  At $s=0$ the edge condition
$2j<s-1=-1$ retains no boundary jet.  Precisely the strict
strip in the second line of (MA-CGP.23) remains.
Subtracting the two absolutely convergent sums proves
(MA-CGP.24), including its lower-edge multiplicities.
\end{proof}

For the source invariant product level $Y$, MA-CGP-002 gives
$Y=Q\Lambda$, so the support edge is $c=-\log Y$.
Thus the theorem applies to that actual source family, with
$S'=S$ and the character's conductor and different included.
Every finite localized trace used above also has the concrete
Gram-resolvent realization (MA-CGP.18).  Its regularizer is
taken to zero at those fixed data; it is not interchanged
with either position or band limit by this theorem.

For Dirichlet characters the conductor is $q$, and for the
non-induced cubic example it is the analytic scale
$3\cdot269^2$.  The phase and scale have been proved in
MA-CGP-001, while the dual-sector local cutoff is MA-CQ.
The calculation above is an actual global character extension
of MA-GP-009.  It neither identifies that global operator
with the local cutoff nor supplies a uniform joint-cutoff
estimate.  The growth of the fixed-level norm equivalence
still prevents such a conclusion from the present proof.
Under RH for the particular character the two displayed
sums agree, but that conditional agreement is not used
anywhere in the operator proof.

The finite-length zero modules and ordinary trace ideals
have the same Grothendieck placement as MA-GP-009 and
MA-CGP-003.  The proof adds no arithmetic positivity,
new motive, or archival source assertion.  Its source
inputs are the fully proved actual character range,
the fixed-character zero count and the already proved
half-line endpoint/cusp lift; all new quotient and
character normalizations are specified above.

\section{Polynomial-weighted causal synthesis for actual character PW levels}
\label{sec:character-weighted-pw}

The nontrivial character source has no archimedean polar moment
constraint.  Its reflected theta formula nevertheless has an exact
conductor scale and a dual character.  We retain both, identify the
actual weighted range by a bounded causal convolution, and prove
the spectral and ambient trace assertions in that quotient.
No weighted Beurling identification is assumed.

Let \(K=\mathbb Q\) or \(\mathbb Q(\zeta_3)\), and let \(\eta\)
be a nonprincipal primitive Dirichlet character in the first case,
or a nontrivial finite-order Hecke character covered by
MA-CGP-001 in the second.  This includes \(\psi_{269,3}\).
Write
\[
 \begin{gathered}
 D_\infty=[K:\mathbb Q],\quad
 w_K=|\mathcal O_K^\times|,\quad Q=|D_K|N\mathfrak f_\eta,\\
 d=\tfrac12-\alpha,\quad u=\tfrac12+\alpha,
 \qquad 0<\alpha<\tfrac12,\quad s\ge0 .
 \end{gathered}
 \tag{MA-GW.45}
\]
Here \(D_\infty=1,2\) and \(w_K=2,6\), respectively.
The degree \(D_\infty\) is distinct from the strip parameter \(d\).
For the primitive archimedean norm band \(\Lambda>0\), set
\[
 Y=Q\Lambda,\quad c=-\log Y,\quad
 H=L^2(\mathbb R,e^{2\alpha|x|}(1+x^2)^{s/2}dx),\quad
 H_c=\mathbf1_{[c,\infty)}H .
 \tag{MA-GW.46}
\]
Let \(J_{\eta,\Lambda}\) be the actual Hilbert closure of
\(\mathscr P_{\eta,\Lambda}\) in (MA-CGP.9), written in the
logarithmic variable.  Its source invariant level is exactly
\(Y\), with \(S'=S\), by (MA-CGP.11); no extra unit-shell
Euler deletion is inserted.  Let \(Q_{\eta,\Lambda}\) be its
orthogonal projection in \(H\), and put
\[
 \mathcal Q_{\eta,c}=H_c/J_{\eta,\Lambda},\qquad
 Z_{\eta,c}=H_c\ominus J_{\eta,\Lambda},\qquad
 P_{\eta,c}^Z=\text{the orthogonal projection onto }Z_{\eta,c}.
 \tag{MA-GW.47}
\]
The last projection is extended by zero on \(H_c^\perp\).

\paragraph{MA-GW-004 (the actual character theorem at every polynomial weight).}
Let \(\rho\) run through the zeros of the primitive entire
completion \(\Lambda_\eta\), with multiplicities \(m_\eta(\rho)\).
Define
\[
 b_s=\#\{j\in\mathbb Z_{\ge0}:2j<s-1\},\qquad
 r^+_{\eta,\rho}=
 \begin{cases}
 m_\eta(\rho),&\operatorname{Re}\rho<u,\\
 \min(m_\eta(\rho),b_s),&\operatorname{Re}\rho=u,\\
 0,&\operatorname{Re}\rho>u .
 \end{cases}
 \tag{MA-GW.48}
\]
The actual closed subspace \(J_{\eta,\Lambda}\subset H_c\)
is precisely the common kernel of the continuous Mellin jets
of orders \(0\le j<r^+_{\eta,\rho}\) at \(\rho-\tfrac12\).
The corresponding finite spectral blocks span a dense subspace
of the actual quotient.

For every additive \(g\in C_c^\infty((0,\infty))\), both
actual operators
\[
 C_g^{\mathcal Q_{\eta,c}}\in\mathcal S_1(\mathcal Q_{\eta,c}),
 \qquad
 C_gP_{\eta,c}^Z\in\mathcal S_1(H)
 \tag{MA-GW.49}
\]
are trace class, and
\[
 \operatorname{Tr}_{\mathcal Q_{\eta,c}}C_g^{\mathcal Q_{\eta,c}}
   =\operatorname{Tr}_H(C_gP_{\eta,c}^Z)
   =\sum_\rho r^+_{\eta,\rho}\,\mathcal Mg(\rho-\tfrac12),
 \qquad
 \mathcal Mg(z)=\int_0^\infty g(t)e^{zt}dt .
 \tag{MA-GW.50}
\]
The sum is absolutely convergent and independent of \(\Lambda\).
For the uniformly bounded smooth compact exhaustion
\(\chi_R(x)=\chi_0(x/R)\), with \(\chi_0=1\) near zero,
\[
 \begin{aligned}
 \lim_{\Lambda\to\infty}\lim_{R\to\infty}
    \operatorname{Tr}(M_{\chi_R}C_gQ_{\eta,\Lambda})
       &=-\sum_\rho r^+_{\eta,\rho}\mathcal Mg(\rho-\tfrac12),\\
 \lim_{R\to\infty}\lim_{\Lambda\to\infty}
    \operatorname{Tr}(M_{\chi_R}C_gQ_{\eta,\Lambda})
       &=-\sum_\rho r^{\rm strip}_{\eta,\rho}
                                      \mathcal Mg(\rho-\tfrac12),
 \end{aligned}
 \tag{MA-GW.51}
\]
where
\[
 r^{\rm strip}_{\eta,\rho}=
 \begin{cases}
 m_\eta(\rho),&d<\operatorname{Re}\rho<u,\\
 \min(m_\eta(\rho),b_s),&
            \operatorname{Re}\rho=d\text{ or }\operatorname{Re}\rho=u,\\
 0,&\text{otherwise}.
 \end{cases}
 \tag{MA-GW.52}
\]
The first limit minus the second equals
\[
 -\sum_{\operatorname{Re}\rho<d}
             m_\eta(\rho)\mathcal Mg(\rho-\tfrac12)
 -\sum_{\operatorname{Re}\rho=d}
        \bigl(m_\eta(\rho)-\min(m_\eta(\rho),b_s)\bigr)
                                  \mathcal Mg(\rho-\tfrac12).
 \tag{MA-GW.53}
\]
The theorem asserts neither the existence of such omitted zeros
or jets nor the nonvanishing of this expression for a given test.
It uses no \(\mathsf A\), \(M\), generic-edge assumption, or
separation between zeros.

\begin{proof}
\emph{1. The actual source and its bounded causal symbol.}
The map
\[
 (U_cf)(y)=e^{\alpha(c+y)}f(c+y),\qquad
 W_c(y)=e^{2\alpha|c+y|-2\alpha(c+y)}
                                      (1+(c+y)^2)^{s/2}
 \tag{MA-GW.54}
\]
is unitary from \(H_c\) onto \(H_{W_c}=L^2((0,\infty),W_cdy)\).
For this fixed \(c\), that norm is equivalent to
\(\mathcal H_s=L^2((0,\infty),(1+y^2)^{s/2}dy)\).
The polynomial weight has not been removed.

Choose an integer \(N\) sufficiently large and set
\[
              \Phi_\eta(z)=
                 \frac{L_K(z+d,\bar\eta)}{(z+2)^N}.
 \tag{MA-GW.55}
\]
There is no pole-canceling numerator factor: the nontrivial
finite Hecke function is entire.  Put \(\epsilon_0=d/2\).
The fixed-character vertical polynomial bound of NT-Z10,
and the absolutely convergent Euler product on a right
half-plane, permit \(N>3\) with
\[
 \begin{aligned}
 |\Phi_\eta(\sigma+it)|&\le C(1+|t|)^{-2}
                  &&(-\epsilon_0\le\sigma\le2),\\
 |\Phi_\eta(z)|&\le C|z|^{-N}
                  &&(\operatorname{Re}z\ge2).
 \end{aligned}
 \tag{MA-GW.56}
\]
There is no singularity on this closed half-plane.
The inverse Laplace integral may be moved to
\(\operatorname{Re}z=-\epsilon_0\) by finite rectangles
and the uniform vertical decay.  It gives a continuous
kernel \(\phi_\eta\) with
\(|\phi_\eta(t)|\le Ce^{-\epsilon_0t}\) for \(t\ge0\).
For \(t<0\), move instead to \(\operatorname{Re}z=M\).
The absolute integral is at most \(Ce^{Mt}M^{1-N}\);
letting \(M\to\infty\) proves that the kernel is causal.
Fourier inversion on a vertical line, followed by analytic
continuation, gives \(\mathcal L\phi_\eta=\Phi_\eta\) on
\(\operatorname{Re}z>-\epsilon_0\).

The unilateral shifts \(S_tf(y)=\mathbf1_{y\ge t}f(y-t)\)
on \(\mathcal H_s\) are strongly continuous and satisfy
\[
 \|S_t\|\le2^{s/4}(1+t^2)^{s/4},\qquad
 T_{\phi_\eta}f=\int_0^\infty\phi_\eta(t)S_tf\,dt
                                      \in\mathcal B(\mathcal H_s).
 \tag{MA-GW.57}
\]
The inequality follows from
\(1+(y+t)^2\le2(1+y^2)(1+t^2)\).
Continuity follows on compact smooth functions and then
by density.  The integral is vectorwise, with an integrable
operator-norm majorant.  Its Laplace action is multiplication
by \(\Phi_\eta\).

We now prove the exact range equality
\[
 U_cJ_{\eta,\Lambda}
        =\overline{\operatorname{ran}T_{\phi_\eta}}^{\,\mathcal H_s}.
 \tag{MA-GW.58}
\]
Let \(u_\infty=\widehat f\) be any actual smooth frequency
test of the required parity or radial type supported in
\(\{|\xi|_\infty\le\Lambda\}\), and put
\(U(r)=u_\infty(r^{1/D_\infty})\).
No value or moment condition at zero is imposed.
Indeed the nontrivial primitive finite vector in
(MA-CGP.5) already kills both adelic polar terms.
At least one finite place is ramified in these
class-number-one character families, as proved in
MA-CGP-001.

The actual reflected theta calculation (MA-CGP.13) gives
\[
 \begin{aligned}
 (\mathcal L U_c\Theta_\eta f)(z)
  &=w_K\epsilon_\eta^\theta e^{zc}Q^{z-\alpha}
       L_K(z+d,\bar\eta)\int_0^\Lambda U(r)r^{z+d-1}dr\\
  &=C_\eta(Y)\,L_K(z+d,\bar\eta)\,\mathcal La(z),\\
 a(q)&=e^{-dq}U(\Lambda e^{-q}),\qquad
 C_\eta(Y)=w_K\epsilon_\eta^\theta\,Y^dQ^{-1/2}.
 \end{aligned}
 \tag{MA-GW.59}
\]
For completeness the scalar simplification is
\(e^{zc}Q^{z-\alpha}\Lambda^{z+d}
=Q^{-\alpha}\Lambda^d=Y^dQ^{-1/2}\).
The first equality is also valid for a frequency test
nonzero at the origin: initially take \(\operatorname{Re}(z+d)>1\),
where the ideal sum and the integral converge absolutely.
Both sides then continue to \(\operatorname{Re}z>0\).
Thus no annular-density assertion is needed for this inclusion.

At \(q=0\), \(a\) and all its derivatives vanish because
the smooth compact frequency test is flat at the band edge.
At infinity its derivatives are \(O(e^{-dq})\).
For real parity this follows from the smooth positive-ray
function \(U\) and its logarithmic derivatives; for complex
radial smooth functions, \(u_\infty(\sqrt r)\) is smooth at
\(r=0\), and the same bound holds.  In particular it holds
even when \(U(0)\ne0\).
Consequently \((\partial+2)^Na\in\mathcal H_s\), all
integration-by-parts boundary terms at zero vanish, and
(MA-GW.59) becomes
\[
 U_c\Theta_\eta f
           =C_\eta(Y)\,T_{\phi_\eta}(\partial+2)^Na.
 \tag{MA-GW.60}
\]
Laplace uniqueness justifies this equality of actual
half-line functions.  This proves inclusion of every
actual source core vector in the bounded convolution range.

Conversely, any \(a\in C_c^\infty((0,\infty))\) gives an
actual frequency test by
\(U(\Lambda e^{-q})=e^{dq}a(q)\).
Its parity or radial extension is smooth, supported away
from the origin and the band edge, and imposes no moment
condition.  Its Fourier inverse belongs to the actual
archimedean Schwartz source.  Formula (MA-GW.60) applies.
For \(h\in C_c^\infty((0,\infty))\), solve
\((\partial+2)^Na=h\) with zero initial traces using the
causal kernel \(e^{-2q}q^{N-1}/(N-1)!\).
The solution is flat at zero and all its derivatives
decrease exponentially, up to a fixed polynomial factor.
Smooth compact truncations \(a_M\) satisfy
\((\partial+2)^Na_M\to h\) in \(\mathcal H_s\).
Boundedness of \(T_{\phi_\eta}\), the nonzero scalar
\(C_\eta(Y)\), and density of these \(h\) prove the reverse
inclusion in (MA-GW.58).  Equivalent fixed-\(c\) norms
preserve this same closed subspace.

\emph{2. The actual quotient resolvent and its character index.}
Set \(\mathcal J=\overline{\operatorname{ran}T_{\phi_\eta}}\)
and \(\mathcal Q=\mathcal H_s/\mathcal J\).
The shifts commute with \(T_{\phi_\eta}\), hence induce a
strongly continuous polynomially bounded semigroup on
\(\mathcal Q\).  Let \(A\) be its generator.
For \(\operatorname{Re}\lambda>0\) its actual resolvent is
\[
 R(\lambda)=\int_0^\infty e^{-\lambda t}S_t^{\mathcal Q}dt.
 \tag{MA-GW.61}
\]
The generator difference quotient, integrated against this
exponentially damped kernel, verifies both inverse identities
and the generator-domain inclusion.

For \(\operatorname{Re}\lambda<\epsilon_0\), define
\[
 \psi_\lambda(r)=\int_r^\infty
                  \phi_\eta(t)e^{\lambda(t-r)}dt,\qquad
 T_\lambda=\int_0^\infty\psi_\lambda(r)S_r^{\mathcal Q}dr.
 \tag{MA-GW.62}
\]
Differentiation under the exponentially dominated integral
makes this operator-norm holomorphic there.  Moreover
\[
 \sup_{\operatorname{Re}\lambda\le\epsilon}\|T_\lambda\|<\infty
                                      \qquad(\epsilon<\epsilon_0).
 \tag{MA-GW.63}
\]
An explicit bound is a constant times
\(\int_0^\infty|\phi_\eta(t)|e^{\max(\epsilon,0)t}
\int_0^t(1+r)^{s/2}dr\,dt\).
It is uniform in height and throughout that entire left
half-plane.

The identities
\(\psi_\lambda'=-\phi_\eta-\lambda\psi_\lambda\) and
\(\psi_\lambda(0)=\Phi_\eta(-\lambda)\) give
\[
 (\lambda-A)T_\lambda=\Phi_\eta(-\lambda)1,\qquad
 T_\lambda(\lambda-A)=\Phi_\eta(-\lambda)1
                                      \quad\text{on }\mathcal D(A).
 \tag{MA-GW.64}
\]
Before quotienting the right side is
\(\Phi_\eta(-\lambda)1-T_{\phi_\eta}\).
The first identity follows by a shift difference quotient
and integration by parts in \(r\); the weighted exponential
majorants above prove that the result lies in the generator
domain.  Applying the same calculation to a domain vector
proves the second identity.  All operators preserve
\(\mathcal J\), because their convolutions commute.
Thus, wherever the scalar is nonzero,
\[
                         R(\lambda)=T_\lambda/\Phi_\eta(-\lambda)
 \tag{MA-GW.65}
\]
is the actual two-sided inverse of the closed generator.
It agrees with (MA-GW.61) on their intersection.
These formulas give a meromorphic resolvent on the plane;
all apparent poles in the right half-plane are removable.

For \(\operatorname{Re}\lambda\le0\), the finite-function
argument \(d-\lambda\) has positive real part.
Its archimedean factors are nonzero and holomorphic:
they are \(\Gamma_{\mathbb R}(d-\lambda+e)\) in the real
case and \(\Gamma_{\mathbb C}(d-\lambda)\) in the complex
case, up to the nonzero conductor factor.
The entire completed functional equation
\(\Lambda_{\bar\eta}(w)=W_{\bar\eta}\Lambda_\eta(1-w)\)
therefore identifies every possible pole, including its
multiplicity, as
\[
 \lambda_\rho=\rho-u,\qquad \operatorname{Re}\rho\le u,
 \qquad\text{pole order at most }m_\eta(\rho).
 \tag{MA-GW.66}
\]
The divisor is that of \(\eta\), not \(\bar\eta\).
The classical containing strip places these possible
poles in \([-u,0]+i\mathbb R\).
No pole or artificial zero occurs at \(d-\lambda=1\);
the finite character function here is nontrivial and
entire.  Any zero there would be an actual completed
zero paired with the completed endpoint at zero.
Trivial finite-\(L\) zeros have \(d-\lambda\in\mathbb R_{\le0}\) in these
finite-order families and lie outside the left-half-plane
calculation.  No hidden gamma divisor is added.

\emph{3. Finite rank and the exact edge jets.}
At a possible pole \(\lambda_0\) of scalar multiplicity \(m\),
define \(E=(2\pi i)^{-1}\int_\gamma R(\lambda)d\lambda\)
on a small positive circle.  Integrating the resolvent
identity over nested circles gives \(E^2=E\) and commutation
with the resolvent.  Since \(AR(z)=zR(z)-1\), the contour
integrals also converge in each required graph norm and give
\[
 (A-\lambda_0)^mE=0,\qquad
 [(\lambda-\lambda_0)^{-j}]R(\lambda)
                                =(A-\lambda_0)^{j-1}E .
 \tag{MA-GW.67}
\]
These identities alone do not assert finite rank; we prove it.

The continuous complex-linear dual of \(\mathcal H_s\)
is represented by densities \(v\) through
\(\ell_v(f)=\int_0^\infty f(y)v(y)dy\), with
\[
                \int_0^\infty|v(y)|^2(1+y^2)^{-s/2}dy<\infty .
 \tag{MA-GW.68}
\]
The quotient dual is the annihilator of \(\mathcal J\).
Its transpose translations act by \(v(y+t)\).
Thus every vector in \(\operatorname{ran}E'\) satisfies
\((\partial-\lambda_0)^mv=0\) distributionally on the
connected half-line.  Multiplication by \(e^{-\lambda_0y}\)
and successive integration show that it has the form
\(e^{\lambda_0y}p(y)\), with \(\deg p<m\).
If \(\operatorname{Re}\lambda_0<0\), every such polynomial
satisfies (MA-GW.68).  If \(\operatorname{Re}\lambda_0=0\),
its highest degree \(j\) satisfies it exactly when
\(2j<s-1\).  Equality gives logarithmic divergence.
This bounds the dimension of \(\operatorname{ran}E'\)
by \(m\) or \(\min(m,b_s)\).  The range of a bounded
projection and the range of its transpose have the same
finite dimension once either is finite: restricting the
transpose identifies it with the full continuous dual of
the closed range.  Hence this is an actual finite-rank bound.

It is attained.  For each allowed \(0\le j<m\), the
continuous functional
\[
                    \ell_j(f)=\int_0^\infty
                                   y^j e^{\lambda_0y}f(y)dy
 \tag{MA-GW.69}
\]
annihilates \(T_{\phi_\eta}f\).
After substituting the causal convolution and expanding
\((y+t)^j\), each coefficient is a derivative of
\(\Phi_\eta(-\lambda)\) at \(\lambda_0\) of order at most
\(j<m\), and therefore vanishes.  The exponential kernel
bound and (MA-GW.68) justify all integrals.
These independent functionals form a finite triangular
chain under translations.  The Laplace resolvent
(MA-GW.61) on that chain is its ordinary finite-dimensional
resolvent; the small contour integral fixes each member.
They belong to \(\operatorname{ran}E'\), giving the
matching lower bound and its full dual range.
Equation (MA-GW.67) now proves finite rank of every
principal part.

Pulling (MA-GW.69) back by \(U_c\) gives, up to a nonzero
constant,
\[
 f\longmapsto\int_c^\infty
             (x-c)^j e^{(\lambda_0+\alpha)x}f(x)dx .
 \tag{MA-GW.70}
\]
Here \(\lambda_0+\alpha=\rho-\tfrac12\).
The invertible triangular change from \((x-c)^j\) to
\(x^j\) gives the stated Mellin jets and the lengths
(MA-GW.48).

\emph{4. Uniform contour bounds and actual spectral synthesis.}
Fix \(0<\epsilon<\epsilon_0\) and \(C_0>u+2\).
The right-line resolvent at \(\operatorname{Re}\lambda=\epsilon\)
is uniformly bounded by (MA-GW.57) and (MA-GW.61).
On the whole region \(\operatorname{Re}\lambda\le-C_0\),
the Euler product gives
\[
 |L_K(d-\lambda,\bar\eta)|
   \ge\prod_{\mathfrak p}(1-(N\mathfrak p)^{-(d+C_0)})
   =\zeta_K(d+C_0)^{-1}>0 .
\]
Together with (MA-GW.63) and (MA-GW.65), this yields
\[
                   \|R(\lambda)\|\le C(1+|\lambda|)^N
                         \qquad(\operatorname{Re}\lambda\le-C_0).
 \tag{MA-GW.71}
\]
The constant is uniform as the left line moves.
The exponent is \(N\), rather than \(N-1\) in the
pole-canceled zeta case.

The fixed-character good-height proof MT-M18,
(M18.5)--(M18.8), applies to \(L_K(\,\cdot\,,\bar\eta)\).
It uses the unconditional local count from NT-Z10,
the Euler lower bound at \(2+iT\), the logarithmic-distance
averaging on a half-unit interval, and the finite
disk-Blaschke/Harnack argument.  In particular it gives a
polynomial lower bound for the finite \(L\)-function along
one fixed horizontal strip segment, independently of
separations of zeros inside the blocks.
Reflecting the chosen heights under \(w=d-\lambda\)
gives \(t_n=n+O(1)\), with gaps between \(1/2\) and \(3/2\).
For
\(\Gamma_n=\partial([-C_0,\epsilon]+i[t_n,t_{n+1}])\),
all low contours adjusted individually if necessary,
put \(E_n=(2\pi i)^{-1}\int_{\Gamma_n}R(\lambda)d\lambda\).
We obtain
\[
 \sup_{\Gamma_n}\|R(\lambda)\|\le C(2+|n|)^B,\qquad
 \operatorname{rank}E_n=O(\log(2+|n|)).
 \tag{MA-GW.72}
\]
The right side uses the actual right resolvent if the
scalar denominator has a removable zero there.
Every horizontal side avoids all relevant zeros.
The rank bound follows from the finite-rank calculation
and the completed local zero count.
Only this one central strip uses the good-height estimate;
the expanding left region uses the uniform bound
(MA-GW.71).

For \(k\in C_c^\infty((0,\infty))\) put
\(H_k(\lambda)=\int k(t)e^{\lambda t}dt\) and
\[
 V_n(k)=\frac1{2\pi i}\int_{\Gamma_n}
                             H_k(\lambda)R(\lambda)d\lambda .
 \tag{MA-GW.73}
\]
Their ranges lie in \(\operatorname{ran}E_n\) by the
resolvent contour identities.  Repeated integration by
parts in \(H_k\) and (MA-GW.72) give
\[
 \sum_n\|V_n(k)\|_1
  \le C_M\sum_n\log(2+|n|)(2+|n|)^{B-M}<\infty
                                             \qquad(M>B+2).
 \tag{MA-GW.74}
\]
Trace class is therefore established before taking residues
or scalar traces.

On the line \(\operatorname{Re}\lambda=\epsilon\), inserting
the actual Laplace resolvent and applying Fourier inversion
gives
\[
 C_k^{\mathcal Q}
   =\frac1{2\pi i}\int_{\epsilon-i\infty}^{\epsilon+i\infty}
                                  H_k(\lambda)R(\lambda)d\lambda .
 \tag{MA-GW.75}
\]
Fubini is justified by the exponentially damped polynomial
semigroup bound and rapid vertical decay of \(H_k\).
Finite rectangles with good horizontal heights move this
integral to \(-C_0\) and produce precisely the block sum;
the horizontal integrals vanish by (MA-GW.72).
There are no poles farther left.  Move the remaining
line to \(-M\), using (MA-GW.71), and then let \(M\to\infty\).
If \(\operatorname{supp}k\subset[a,b]\), \(a>0\), then
\[
 |H_k(-M+it)|
   \le C_{\ell,k}e^{-Ma}(1+M)^\ell(1+|t|)^{-\ell}.
\]
For \(\ell>N+2\), (MA-GW.71) makes the remaining
operator-norm integral tend to zero.
Consequently the actual operator identity is
\[
                          C_k^{\mathcal Q}=\sum_nV_n(k)
                                      \quad\text{in }\mathcal S_1 .
 \tag{MA-GW.76}
\]
Choose a positive-time smooth approximate identity supported
in \([\varepsilon,2\varepsilon]\).
Its quotient convolutions tend strongly to the identity.
Equation (MA-GW.76) proves density of the finite blocks.
If every retained jet vanishes on a quotient vector, the
full dual ranges determined above show that every \(E_n\)
kills it; then every \(C_k^{\mathcal Q}\) kills it, and
the approximate identity makes it zero.
This proves the exact closed-kernel assertion, with no
remaining unsynthesized quotient.

On each finite block the trace is its length times
\(H_k(\lambda_\rho)\); all higher nilpotent terms have
trace zero.  Set \(k(t)=e^{\alpha t}g(t)\).
Then \(H_k(\lambda_\rho)=\mathcal Mg(\rho-\tfrac12)\).
The trace-norm series and bounded similarity to the
actual quotient norm prove its trace assertion in
(MA-GW.50).  Absolute convergence also follows from
NT-Z10 and rapid vertical Mellin decay on the fixed strip.

We also need an actual trace-class resolvent power.
Take \(\nu>\epsilon\) and an integer
\(q>\max\{B+2,N+2\}\).  Repeat the same contours with
\(H(\lambda)=(\nu-\lambda)^{-q}\).
On the right line its actual Laplace kernel is
\(t^{q-1}e^{-\nu t}/(q-1)!\), so Fourier inversion gives
\((\nu-A)^{-q}\).
The block series converges in trace norm by (MA-GW.72).
The far-left remainder is at most
\(CM^{N+1-q}\) by (MA-GW.71); the horizontal integrals
vanish by the same chosen inequalities.
Thus
\[
                            (\nu-A)^{-q}\in\mathcal S_1(\mathcal Q).
 \tag{MA-GW.77}
\]

\emph{5. The actual ambient lift, with its endpoint and cusp.}
Return by the bounded similarity to \(H_{W_c}\), and
write \(W=W_c\).  The adjoint shifts and their generator are
\[
 (S_t^{*,W}v)(y)=\frac{W(y+t)}{W(y)}v(y+t),\qquad
 L_Wv=W^{-1}(Wv)' .
\]
Multiplication by \(W\) is unitary onto \(L^2(W^{-1}dy)\),
where this is the left-translation generator with no
boundary condition at zero.  The fundamental theorem of
calculus gives its difference-quotient domain; the converse
follows by testing against compact interior functions.
Induction gives
\[
 v\in\mathcal D(L_W^q)
 \quad\Longleftrightarrow\quad
 (Wv)^{(j)}\in L^2(W^{-1}dy)\quad(0\le j\le q).
 \tag{MA-GW.78}
\]
The required translation majorants follow from (MA-GW.57)
and fixed-\(c\) norm equivalence.

The actual orthogonal complement of the invariant ideal
is invariant under these adjoints.
The full resolvent restricts to it by its Laplace integral.
Adjunction of (MA-GW.77) therefore gives
\[
 \mathcal R_q=(\nu-L_W)^{-q}P_{U_cZ_{\eta,c}}
                   \in\mathcal S_1(H_W),
 \tag{MA-GW.79}
\]
and it maps boundedly into the \(q\)th graph domain.
This last assertion follows from
\(L_W(\nu-L_W)^{-1}=\nu(\nu-L_W)^{-1}-1\).

To make the ambient factorization explicit, put \(a=W^{-1}\),
\(b=Wv\), and \(k=e^{\alpha\,\cdot}g\).  For a domain vector,
\[
 C_k^+(\nu-L_W)^qv(x)
       =\int_0^\infty k(x-y)a(y)(\nu-\partial_y)^qb(y)dy .
 \tag{MA-GW.80}
\]
Split at \(L=-c\) if \(c<0\).
On each smooth piece every \(W\partial^j(W^{-1})\)
is bounded, because all positive-order derivatives of
\[
 2\alpha|c+y|-2\alpha(c+y)+\tfrac s2\log(1+(c+y)^2)
\]
are bounded there and have finite one-sided cusp values.
After integration by parts the bulk is the bounded operator
\[
 \mathcal B_{k,q}v(x)
   =\int W(y)(\nu+\partial_y)^q
                         \bigl(k(x-y)/W(y)\bigr)v(y)dy .
 \tag{MA-GW.81}
\]
It is a finite sum of bounded input multipliers followed
by causal convolutions with derivatives of \(k\).
For each derivative order \(j\ge1\), the full boundary
expression on each interval is
\[
 \binom qj\nu^{q-j}
 \sum_{\ell=0}^{j-1}(-1)^{j+\ell}
   [\,\partial_y^\ell(k(x-y)a(y))\,
                          b^{(j-1-\ell)}(y)\,]_{\rm left}^{\rm right}.
 \tag{MA-GW.82}
\]
The endpoint at infinity vanishes.
Keep the lower-endpoint term at zero and the difference
of the two one-sided terms at the cusp.
By (MA-GW.78), \(b\in H^q_{\rm loc}\) and has common
derivative traces through order \(q-1\) there.
Each resulting term is a compact smooth output profile,
a linear combination of \(k^{(i)}(x)\) or \(k^{(i)}(x-L)\),
times a graph-bounded functional \(b^{(j)}(0)\) or
\(b^{(j)}(L)\), \(j<q\).
The one-dimensional Sobolev trace inequality supplies
this bound.  In particular the endpoint and cusp have
not been omitted.

Apply (MA-GW.80)--(MA-GW.82) to \(v=\mathcal R_qf\).
The result is
\(C_k^+P_{U_cZ_{\eta,c}}=\mathcal B_{k,q}\mathcal R_q\)
plus finitely many bounded rank-one maps.
This proves actual ambient trace class.
Conjugating by \(U_c\) and extending by zero proves
the second assertion of (MA-GW.49).
Trace cyclicity with the actual orthogonal projection
identifies this ambient trace with the quotient trace,
already established.  This proves all of (MA-GW.50).

\emph{6. The two orders of actual cutoff limits.}
The exact orthogonal decomposition in \(H\) is
\[
                 Q_{\eta,\Lambda}=P_c-P_{\eta,c}^Z,
                   \qquad P_c=\mathbf1_{[c,\infty)} .
\]
For a compact smooth position multiplier, \(M_\chi C_gP_c\)
is trace class with trace zero.  To verify this before
taking the diagonal, extend the compact-position
convolution kernel to a smooth compact rectangle.
Its rapidly convergent double Fourier expansion is nuclear.
The weight on each compact coordinate projection is bounded
and boundedly invertible, so it only inserts bounded
multipliers.  Restricting by \(P_c\) preserves trace class.
Bounded smooth approximation of that indicator and
trace-norm continuity identify the trace with
\(g(0)\int\chi=0\).
Consequently
\[
 \operatorname{Tr}(M_{\chi_R}C_gQ_{\eta,\Lambda})
     =-\operatorname{Tr}(M_{\chi_R}C_gP_{\eta,c}^Z).
 \tag{MA-GW.83}
\]
For fixed \(\Lambda\), strong convergence of the uniformly
bounded exhaustion paired with (MA-GW.49) gives its limit
in trace norm and hence the first trace in (MA-GW.51).
That trace is already independent of \(\Lambda\).

For the opposite order, the actual all-weight condensation
(MA-CGP.14) gives trace-norm convergence at each fixed
position cutoff.  Its proof uses smooth frequency cutoffs
in the actual character source, both automatic polar
conditions, and strong monotone convergence of the actual
orthogonal projections; it does not use the \(s=0\)
Hardy identification.
The completed global cutoff theorem MA-WP-003 then gives
the actual full-strip trace, with interior multiplicity
\(m_\eta(\rho)\) and exactly \(2j<s-1\) at each edge.
These are precisely (MA-GW.52).
This proves the second order in (MA-GW.51).
Subtracting the two absolutely convergent series yields
(MA-GW.53): the upper-edge truncations cancel and only
the omitted lower-edge lengths remain.
\end{proof}

\paragraph{Exact analytic inputs and scope.}
The source formula is (MA-CGP.13), with its preceding
actual finite-vector and Poisson proofs (MA-CGP.5)--(MA-CGP.7).
Those proofs retain the raw compact type, Gauss phase,
real parity, complex additive measure \(2\,dx\,dy\),
and exact conductor scale.
For Dirichlet characters \(Q=q\); for the non-induced
cubic character \(Q=3\cdot269^2\).
The constant in (MA-GW.59) is therefore fully specified,
rather than absorbed into a changed character divisor.
No archimedean moment or canceled-pole factor belongs
to the present nontrivial character range.

The unconditional fixed-character polynomial growth,
local count and Euler-center bound are NT-Z10,
(NT-J-upper)--(NT-J-lower), in the supplied number-theory
module.  The complete good-height and disk minimum-modulus
proof is MT-M18, (M18.5)--(M18.8).
These arguments were reread for the present extension;
they concern a fixed Hecke function and give no uniformity
in conductor or field.
The inherited primary completed-function locators are
Ramakrishnan--Valenza, \emph{Fourier Analysis on Number Fields},
\S7.4, Theorem~7-19, printed p.~279, and \S7.7,
Theorem~7-28, printed p.~289.
No GRH statement is used to obtain those analytic bounds.

The finite-length modules here are the proved finite-rank
principal parts of an actual weighted quotient resolvent.
Their trace-class smoothing, graph-domain lift and cutoff
limits precede evaluation of the zero sums.
For the character's rank-one finite-monodromy Artin data,
and the non-induced cubic representation, the exact
Grothendieck/Artin placement remains MA-CH and MA-CQ.
For the inherited trace framework the available locators
remain cote~1 scans~079--080 and~111--112,
ALPOGE\_BRIDGES BR-TR-01/02.
No new archive inspection, arithmetic positivity,
half-integral Tate object, or zero-location theorem is
asserted by this extension.

All norm comparisons are for a fixed finite band.
No uniform trace-norm estimate as \(\Lambda\to\infty\),
joint-limit interchange, or equality with the separately
constructed local character Fourier cutoff follows from
this theorem.

\section{The physical Fourier cutoff does not intertwine the global PW projection}
\label{sec:physical-projection-obstruction}

\paragraph{Status and exact question.}
The complete author and root independent whole-proof readings
are recorded. Further independent review has its own event scope.
It tests a literal operator identity left open in
\texttt{MA-Q1-CQ003} and \texttt{MA-Q1-RF003}.  It does not
identify the local and global Hilbert norms, or replace a
projection by its trace.  The counterexample is first constructed
on genuine Schwartz PW data and then transferred to the original
coordinate-flat LF core of MA-GA-001.

Let $\chi=\chi_4$ be the primitive odd character modulo $4$,
and put $S=\{\infty,2\}$.  Write
\[
 \chi(n)=0\ (2\mid n),\qquad
 \chi(4n+1)=1,\qquad\chi(4n+3)=-1,
 \qquad \beta(z)=L(z,\chi).
 \tag{MA-PI.1}
\]
For an odd archimedean Schwartz function $f$ define the actual
centered global theta image in logarithmic coordinates by
\[
 Y_f(x)=2e^{x/2}\sum_{n\geq1}\chi(n)f(ne^x).
 \tag{MA-PI.2}
\]
This is the primitive character model of MA-CGP-001, with
the actual finite vector
$\theta_2=\chi|_{\mathbb Z_2^\times}\mathbf1_{\mathbb Z_2^\times}$
at $2$ and $\mathbf1_{\mathbb Z_p}$ outside $S$.
Both adelic polar moments vanish at the ramified finite place.
The image $Y_f$ belongs to the all-exponential Schwartz space
$\mathcal E$ for every such $f$.

Fix
\[
 0<a<1/2,\qquad s\geq0,\qquad
 H_{a,s}=L^2\!\left(\mathbb R,
           e^{2a|x|}(1+x^2)^{s/2}\,dx\right).
 \tag{MA-PI.3}
\]
Let $Q^{\mathrm{PW}}_{S,Y,a,s}$ be the actual orthogonal
projection onto the closed invariant global PW image of
product-norm Fourier level $Y$, with $S'=S$, in this same
Hilbert space.  This is the corrected centered realization
of the global projection in the specified MA-GP/MA-CGP/MA-GW
norm.  Masullo Definition 11.2.1.3, (11.2.1.7--8),
printed/PDF pp.~488--489, defines the corresponding
$a$-twisted level inside its $W^{s,2}$ realization; here
the image and orthogonal projection are explicitly fixed
in (MA-PI.3).  No inference that equivalent norms preserve
orthogonal projectors is made.  This is not the auxiliary
ordinary-frequency projector used to parametrize the range
in MA-GP/MA-CGP.

For each $Y>0$ let $\rho_Y\in C^\infty([0,\infty))$ satisfy
\[
 0\leq\rho_Y\leq1,\qquad
 \rho_Y(r)=1\quad(0\leq r\leq Y/2),\qquad
 \rho_Y(r)=0\quad(r\geq2Y).
 \tag{MA-PI.4}
\]
Every source-thin smooth PW-ball profile eventually has
these two plateau properties.  No derivative bound on the
family is required.  Use exactly the original physical
Fourier cutoff on $S$-local data:
\[
 \mathfrak B_{S,Y}
       =\mathcal F_S^{-1}M_{\rho_Y(|\xi|_S)}\mathcal F_S.
 \tag{MA-PI.5}
\]
Its actual descended version acts on the local quotient
completion constructed in MA-CQ and MA-GA.  The vectors
below lie in a canonical source subspace on which all
operations in (MA-PI.5) are literal Schwartz operations.

\paragraph{MA-PI-001 (an original flat-source obstruction at large finite levels).}
For each fixed $a,s$ in (MA-PI.3), there is $Y_0>0$ such
that, for every $Y\geq Y_0$ and every smooth profile
(MA-PI.4), there is an odd, archimedean-flat Schwartz
function $f_Y$ for which
\[
 Q^{\mathrm{PW}}_{S,Y,a,s}Y_{f_Y}
              \ne Y_{\mathfrak b_Y f_Y},
 \qquad
 \mathfrak b_Y f
     :=\mathcal F_\infty^{-1}
          \bigl(\rho_Y(4|\xi|)\widehat f(\xi)\bigr).
 \tag{MA-PI.6}
\]
The actual $S$-source input is
$\varphi_Y=f_Y\otimes\theta_2\in V_S$ of MA-GA.2.
Moreover
\[
 \mathfrak B_{S,Y}\varphi_Y
           =(\mathfrak b_Y f_Y)\otimes\theta_2
                         \in\mathcal S_0(\mathbb A_S).
 \tag{MA-PI.7}
\]
Thus (MA-PI.6) is a failure on the original flat core of
the canonical intertwining
\[
 Q^{\mathrm{PW}}_{S,Y,a,s}\,\vartheta_S
                 \stackrel{?}{=}\vartheta_S\,B_{S,Y},
 \tag{MA-PI.8}
\]
where $\vartheta_S$ is the actual centered global theta
embedding of the local coinvariant source and $B_{S,Y}$
is the descended physical cutoff.  Both sides are defined
on the witness.  The assertion holds at every narrow
width and every polynomial weight, including the draft's
$s\geq3$ scope, at arbitrarily large physical cutoff levels.
No zero-location assumption is used.

\begin{proof}
\emph{1. The actual finite source and the cutoff scale.}
Use the coherent additive characters of MA-NR-001/MA-CQ:
the archimedean Fourier kernel is $e^{-2\pi i x\xi}$ and
the $2$-adic kernel is its coherent positive fractional-part
counterpart.  The self-dual measure has
$\operatorname{vol}(\mathbb Z_2)=1$.
The two residue classes $1+4\mathbb Z_2$ and
$3+4\mathbb Z_2$ give directly
\[
 \mathcal F_2\theta_2(\xi)
     =\frac{i}{2}\chi(4\xi)
                           \mathbf1_{\mathbb Z_2^\times}(4\xi).
 \tag{MA-PI.9}
\]
Indeed on $\xi=u/4$, $u$ odd, their two integrals are
$\tfrac14(e^{\pi iu/2}-e^{3\pi iu/2})=\tfrac i2\chi(u)$;
at the other points of $\tfrac14\mathbb Z_2$ the difference
is zero, and outside that ball each integral is zero.
The finite Fourier support is exactly the shell
$|\xi|_2=4$.  Consequently multiplication by the physical
product-norm cutoff is exactly multiplication by
$\rho_Y(4|\xi_\infty|)$ on the archimedean Fourier factor.
Inverting both transforms proves (MA-PI.7), with the
finite phase canceled by the actual inverse.

The outside-$S$ integer balls restrict the global theta
sum to integers.  The two signs contribute equally because
both $\chi$ and $f$ are odd, giving (MA-PI.2) with no
extra constant.  The finite vector and its Fourier
transform both vanish at zero, so all these adelic
vectors are polar free independently of archimedean
moments.  The product formula makes the global theta
map constant on the $S$-unit coinvariant relations.
Thus the calculation is valid after taking the actual
local quotient; it is not a raw-cover representative
substituted for that quotient.

If $\widehat f$ is smooth, compactly supported, and has
support in $|\xi_\infty|\leq Y/4$, then the whole finite
Fourier support in (MA-PI.9) has product norm at most
$Y$.  It is compact and satisfies both polar conditions.
Its global theta image therefore belongs to the literal
invariant PW source level, and hence to
$\operatorname{Ran}Q^{\mathrm{PW}}_{S,Y,a,s}$.
Only this direct membership is used below; no equality
between different candidate PW ranges is needed.

\emph{2. An exact one-sided theta Gram form.}
For the intermediate calculation let $-1/2<b<1/2$ and
put $\sigma=1/2+b\in(0,1)$.
For $u\in C_c^\infty(\mathbb R_{>0})$, extend $u$ oddly
to $\mathbb R$ and put $f=\mathcal F_\infty^{-1}u$.
Define its logarithmic frequency vector by
\[
 A_u(r)=e^{(1/2-b)r}u(e^r),\qquad
 \widehat A_u(t)=\int_{\mathbb R}A_u(r)e^{-itr}\,dr.
 \tag{MA-PI.10}
\]
Each $A_u$ is a compact smooth function, and conversely
every compact smooth $A$ is obtained by
$u(\xi)=\xi^{b-1/2}A(\log\xi)$ on $\xi>0$.

Write $M_f(z)=\int_0^\infty f(t)t^{z-1}\,dt$.
The actual theta Mellin transform is
\[
 \mathcal M Y_f(z)=2\beta(z+1/2)M_f(z+1/2).
 \tag{MA-PI.11}
\]
It follows first by absolutely convergent unfolding in
$\Re(z+1/2)>1$ and then by analytic continuation; the
theta image is in $\mathcal E$ and $f$ is odd Schwartz.
The actual inverse odd Fourier transform gives
\[
 M_f(z)=i\,\frac{\Gamma_{\mathbb R}(z+1)}
                       {\Gamma_{\mathbb R}(2-z)}M_u(1-z),
 \qquad 0<\Re z<1.
 \tag{MA-PI.12}
\]
For clarity, $f(x)=2i\int_0^\infty u(r)\sin(2\pi xr)\,dr$.
Insert $e^{-\varepsilon x}$ in its Mellin integral and
integrate first in $x$.  The resulting gamma expression
converges uniformly for $r$ in the compact positive
support of $u$.  Letting $\varepsilon\downarrow0$ gives
$2i(2\pi)^{-z}\Gamma(z)\sin(\pi z/2)M_u(1-z)$.
The gamma reflection and duplication identities give
exactly (MA-PI.12).  This also fixes its Fourier sign.

Mellin Plancherel on the line $b$ now yields the exact
sesquilinear form
\[
\begin{split}
 \int_{\mathbb R}Y_f(x)\overline{Y_h(x)}e^{2bx}\,dx
   &=\frac1{2\pi}\int_{\mathbb R}
             W_b(t)\widehat A_u(t)
                         \overline{\widehat A_v(t)}\,dt,\\
 W_b(t)&=4\left|
       \frac{\Gamma_{\mathbb R}(\sigma+1+it)}
            {\Gamma_{\mathbb R}(2-\sigma-it)}
                         \right|^2|\beta(\sigma+it)|^2,
\end{split}
 \tag{MA-PI.13}
\]
where $h=\mathcal F_\infty^{-1}v$.
This is an inner product of actual theta images in a
specified one-sided norm.  It is not asserted to equal
the symmetric norm (MA-PI.3) at this stage.

The multiplier $W_b$ is a real, even, nonnegative tempered
function.  To verify the needed growth without a
mean-square theorem, let $A(X)=\sum_{n\leq X}\chi(n)$;
then $|A(X)|\leq1$.  Partial summation gives, for
$\Re z=\sigma>0$,
\[
 \beta(z)=\sum_{n\leq N}\chi(n)n^{-z}
               -A(N)N^{-z}+z\int_N^\infty A(x)x^{-z-1}\,dx,
 \tag{MA-PI.14}
\]
and the last two terms have absolute value at most
$N^{-\sigma}(1+|z|/\sigma)$.  Taking
$N\asymp1+|\Im z|$, and bounding the finite sum by
$\sum_{n\leq N}n^{-\sigma}$, proves
$\beta(\sigma+it)=O_\sigma((1+|t|)^{1-\sigma})$
for $0<\sigma<1$.  Stirling gives
\[
 \left|
       \frac{\Gamma_{\mathbb R}(\sigma+1+it)}
            {\Gamma_{\mathbb R}(2-\sigma-it)}\right|^2
       \sim (|t|/(2\pi))^{2\sigma-1}.
 \tag{MA-PI.15}
\]
Thus $W_b(t)=O_b(1+|t|)$ for each $|b|<1/2$.

\emph{3. Its inverse-Fourier distribution is not compactly supported.}
The precise functional equation needed here follows
from the same actual theta model.  Formula (MA-PI.9)
gives the finite theta phase $i$ and conductor $4$.
For $f_0(t)=t e^{-\pi t^2}$ one has
$\widehat f_0=-if_0$.  The actual Poisson identity
of MA-CGP.4 therefore reads
$Y_{f_0}(\log t)=Y_{f_0}(\log(1/(4t)))$.
Taking its Mellin transform and using
$2M_{f_0}(z)=\Gamma_{\mathbb R}(z+1)$ proves
\[
 \Gamma_{\mathbb R}(z+1)\beta(z)
       =4^{1/2-z}\Gamma_{\mathbb R}(2-z)\beta(1-z).
 \tag{MA-PI.16}
\]
Both sides continue as the same entire completed function.
Its entireness can also be seen directly from the rapid
theta decay at both ends and the entire reciprocal gamma
factor.  This proof records the root and conductor constants
instead of importing an unspecified functional-equation sign.

Put $d=1/2-b$.  Equations (MA-PI.13) and (MA-PI.16)
show that the real-axis multiplier has the entire continuation
\[
 \mathscr W_b(z)=4^{1-2b}\beta(d+iz)\beta(d-iz).
 \tag{MA-PI.17}
\]
Let $\kappa_b=\mathcal F_-^{-1}W_b$ in the tempered
distribution sense, with inverse normalization $1/(2\pi)$.
Suppose it had compact support in $[-R,R]$.  On a fixed
compact neighborhood of that support, a distribution
has finite order: for some $C,N$, its value is bounded
by $C$ times the maximum of derivatives through order
$N$ of the test function.  Applying this to a compact
cutoff times $e^{-izx}$ shows that its Fourier transform
is entire and satisfies
\[
 |\mathscr W_b(z)|\leq C'(1+|z|)^N e^{(R+1)|\Im z|}.
 \tag{MA-PI.18}
\]
Differentiation in $z$ is allowed in the same finite-order
bound, and agreement on the real axis identifies this
entire transform with (MA-PI.17).

This exponential-type bound is impossible.  Iterating
the gamma recurrence in (MA-PI.16) gives the exact values
\[
 \beta(-2n)=(-1)^n
        \frac{2\,4^n(2n)!}{\pi^{2n+1}}\beta(2n+1)
                  \qquad(n\geq0).
 \tag{MA-PI.19}
\]
For example the $n=0$ identity gives $\beta(0)=1/2$.
For the recurrence calculation one uses
$\Gamma(1/2-n)=(-4)^n n!\sqrt\pi/(2n)!$.
On the positive real axis the alternating series gives
$1-3^{-t}\leq\beta(t)\leq1$ for $t>0$.
At $T_n=2n+d$, (MA-PI.17) and (MA-PI.19) therefore give,
for $n\geq1$,
\[
 |\mathscr W_b(iT_n)|
    =4^{1-2b}|\beta(-2n)|\beta(2n+2d)
       \geq c_b(4/\pi^2)^n(2n)!,\qquad c_b>0.
 \tag{MA-PI.20}
\]
Since $(2n)!\geq n^n$, this grows faster than every
polynomial times $e^{C T_n}$.  It contradicts
(MA-PI.18).  Thus $\kappa_b$ is not compactly supported.
It is even because $W_b$ is even.

\emph{4. Separated genuine source data with nonzero Gram pairing.}
For any prescribed real cut positions $A<B$, there are
$A_u\in C_c^\infty((B,\infty))$ and
$A_v\in C_c^\infty(( -\infty,A))$ for which the
right side of (MA-PI.13) is nonzero.  We give the
distribution argument, including its quantifier.
If every such pairing vanished, then
$\kappa_b$ would vanish on $(-\infty,A-B)$.
Indeed take any compact test $\psi$ supported there,
choose $r_0>B$ with $r_0+\operatorname{supp}\psi\subset
(-\infty,A)$, approximate $\delta_{r_0}$ by smooth
compact functions in $(B,\infty)$, and pair their
convolutions with the fixed translated test
$\overline{\psi(\,\cdot-r_0)}$.
The assumed zero pairings pass to the distributional
limit and give $\langle\kappa_b,\psi\rangle=0$.
Evenness gives the corresponding vanishing on
$(B-A,\infty)$, which would make $\kappa_b$ compactly
supported.  Step 3 excludes this.

This argument supplies ordinary smooth compact functions,
not distributional input vectors: choose a nonzero test
of $\kappa_b$ beyond the prescribed gap, then take a
sufficiently concentrated member of the indicated
approximate identity.  The limiting pairing is nonzero,
so an actual finite member already has nonzero pairing.
Equivalently one may search any fixed countable dense
families of compact smooth functions in the two open
half-lines; continuity of (MA-PI.13) guarantees a pair
with nonzero absolutely convergent integral.  No claim
is made for every arbitrarily preassigned pair of bounded
bands.  The assertion is existence of suitable compact
bands beyond each prescribed physical cut and gap.

Use this result with $b=-a$, $A=\log(1/8)$ and
$B=\log(1/2)$.  Choose such $u,v$, and let
$f=\mathcal F_\infty^{-1}u$,
$h=\mathcal F_\infty^{-1}v$.  They are odd Schwartz,
with smooth compact Fourier supports respectively in
$|\xi|>1/2$ and $0<|\xi|<1/8$, and
\[
 D:=\int_{\mathbb R}Y_f(x)\overline{Y_h(x)}e^{-2ax}\,dx
                           \ne0.
 \tag{MA-PI.21}
\]
All functions here are actual primitive theta source
vectors and their actual images.  There is no requirement
that $f$ or $h$ be flat at zero at this intermediate step.

\emph{5. Large physical levels in the specified symmetric weighted norm.}
For $L>0$ use the unitary archimedean frequency dilation
\[
 u_L(\xi)=e^{-L/2}u(e^{-L}\xi),\qquad
 v_L(\xi)=e^{-L/2}v(e^{-L}\xi),\qquad Y=e^L.
 \tag{MA-PI.22}
\]
Their inverse transforms are
$f_L(t)=e^{L/2}f(e^Lt)$ and
$h_L(t)=e^{L/2}h(e^Lt)$.
Directly from the actual series (MA-PI.2),
\[
       Y_{f_L}(x)=Y_f(x+L),\qquad
       Y_{h_L}(x)=Y_h(x+L).
 \tag{MA-PI.23}
\]
The supports chosen in step 4 and (MA-PI.4) imply
\[
 \mathfrak b_Yf_L=0,\qquad
 \mathfrak b_Yh_L=h_L,\qquad
 Y_{h_L}\in\operatorname{Ran}Q^{\mathrm{PW}}_{S,Y,a,s}.
 \tag{MA-PI.24}
\]
For the last statement the finite product Fourier norm
is at most $4(e^L/8)=Y/2<Y$, so step 1 applies.
These equalities use only the fixed and zero plateaus,
not the values or derivatives of the profile in transition.

The pairing in the specified global norm obeys
\[
 \begin{split}
 &e^{-2aL}L^{-s}
       \langle Y_{f_L},Y_{h_L}\rangle_{H_{a,s}}\\
 &\quad=\int_{\mathbb R}Y_f(x)\overline{Y_h(x)}
       e^{2a|x-L|-2aL}
       \frac{(1+(x-L)^2)^{s/2}}{L^s}\,dx
             \longrightarrow D\ne0.
 \end{split}
 \tag{MA-PI.25}
\]
For $s=0$ the factor $L^{-s}$ is one.  Pointwise, the
two weight factors tend to $e^{-2ax}$ and one.
For $L\geq1$, they are bounded by
$C_s e^{2a|x|}(2+|x|)^s$.  The two actual theta images
are in $\mathcal E$, so dominated convergence is valid.
Consequently the pairing in (MA-PI.25) is nonzero for
every sufficiently large $L$, independently of the
permitted profile.  Orthogonality and (MA-PI.24) give
\[
 \left\langle
 Q^{\mathrm{PW}}_{S,Y,a,s}Y_{f_L}
              -Y_{\mathfrak b_Yf_L},\,Y_{h_L}
                       \right\rangle_{H_{a,s}}
       =\langle Y_{f_L},Y_{h_L}\rangle_{H_{a,s}}\ne0.
 \tag{MA-PI.26}
\]
This is already an actual PW-source counterexample.
It is obtained for all sufficiently large $Y=e^L$, not
by asserting a failure of a limiting scalar trace.

\emph{6. Transfer to the original coordinate-flat LF core.}
Fix one such finite $Y$, its profile, and the corresponding
$f_L,h_L$.  Choose an even $\omega\in C_c^\infty(\mathbb R)$
equal to one near zero.  Put
\[
 f_{L,\varepsilon}(t)
        =(1-\omega(t/\varepsilon))f_L(t),\qquad
 r_\varepsilon(t)=\omega(t/\varepsilon)f_L(t).
 \tag{MA-PI.27}
\]
Then $f_{L,\varepsilon}$ is odd and flat in a neighborhood
of zero, and its tensor with $\theta_2$ is in the actual
source $V_S$ of MA-GA.2.  The approximation is not
asserted in the topology of flat functions at this step;
the precise global and cutoff-output estimates follow.

Because $f_L$ is odd, write
$r_\varepsilon(t)=\varepsilon g_\varepsilon(t/\varepsilon)$,
where
$g_\varepsilon(y)=\omega(y)f_L(\varepsilon y)/\varepsilon$.
For $0<\varepsilon\leq1$ these are a bounded family in
odd compactly supported Schwartz space.  Indeed the
zeroth derivative is bounded by the first derivative
of $f_L$ near zero; higher derivatives acquire factors
$\varepsilon^{j-1}$, and the fixed compact cutoff
handles every Schwartz seminorm.
The actual theta map is continuous from that Schwartz
space to $\mathcal E$, by the direct sum at infinity
and its Poisson-reflected sum at zero, as in MA-CGP-001.
Scaling (MA-PI.2) gives
\[
 Y_{r_\varepsilon}(x)
       =\varepsilon^{3/2}
          Y_{g_\varepsilon}(x-\log\varepsilon),\qquad
 \|Y_{r_\varepsilon}\|_{H_{a,s}}
       \leq C\varepsilon^{3/2-a}
                    (1+|\log\varepsilon|)^{s/2}
                         \longrightarrow0.
 \tag{MA-PI.28}
\]
The estimate follows from the same weight inequality
as (MA-PI.25).  Its exponent is positive throughout
the asserted narrow scope.

There is also convergence after the literal smooth
physical cutoff.  Differentiating the Fourier integral
of $r_\varepsilon$ gives
\[
 \sup_\xi|\partial_\xi^j\widehat r_\varepsilon(\xi)|
      \leq C_j\varepsilon^{j+2}
                         \qquad(j\geq0).
 \tag{MA-PI.29}
\]
For the fixed finite $Y$ and profile, multiplying by
$\rho_Y(4|\xi|)$ gives compactly supported smooth
functions tending to zero in every smooth seminorm.
The inverse Fourier transforms therefore tend to zero
in Schwartz space.  Theta continuity then yields
\[
 \|Y_{\mathfrak b_Yr_\varepsilon}\|_{H_{a,s}}
                   \longrightarrow0.
 \tag{MA-PI.30}
\]
No uniform derivative bound in the profile was used:
the flat approximation is selected after that profile
and finite level are fixed.

Since $\mathfrak b_Yf_L=0$, (MA-PI.28--30) and boundedness
of the actual orthogonal projection show that
\[
 \begin{split}
 &\left\langle
 Q^{\mathrm{PW}}_{S,Y,a,s}Y_{f_{L,\varepsilon}}
       -Y_{\mathfrak b_Yf_{L,\varepsilon}},\,Y_{h_L}
                  \right\rangle_{H_{a,s}}\\
 &\hspace{4em}\longrightarrow
                  \langle Y_{f_L},Y_{h_L}\rangle_{H_{a,s}}
                         \ne0.
 \end{split}
 \tag{MA-PI.31}
\]
Choose any sufficiently small positive $\varepsilon$.
This supplies the actual finite-$\varepsilon$ flat
witness $f_Y=f_{L,\varepsilon}$ in (MA-PI.6).
All original maps on it have the source and target
declared above; no distribution or nonflat PW vector
is substituted for the final source input.
\end{proof}

\paragraph{Exact source and conclusion scope.}
The invariant product-norm PW data and their polar
conditions are Masullo Definition 7.3.2.3,
(7.3.2.12--17), printed/PDF pp.~354--355, with $S'=S$.
The source global projection is Definition 11.2.1.3,
(11.2.1.7--8), pp.~488--489; the corrected centered
Hilbert realization used here is specified after
(MA-PI.3).  The physical local Fourier
cutoff is precisely the first bullet of Step II on
p.~516 and the composition on p.~517, in the proof of
Theorem 12.4.1.1.  These passages were freshly read for
the map tested here.  The coherent finite sign is the
explicit repair already fixed in MA-NR/MA-CQ; (MA-PI.9)
recomputes its required value directly.

The gamma form in (MA-PI.12) agrees with Tate (1950),
Theorem 2.4.1 and the real table on thesis p.~(2.14),
PDF p.~21.  The exact functional equation used to
prove the new obstruction is derived in (MA-PI.16),
and its required growth consequence is proved, not
imported from a mean-square or universality theorem.

The result refutes the canonical theta intertwining
(MA-PI.8), including on the actual flat core and with
source-thin smooth profiles at large finite levels.
It does not assert that the draft printed that exact
intertwining equation: its Step III kernel identification
does not itself supply such an equation or its proof.
The result explains why an actual quotient section and
equality of PW ranges cannot supply this literal map
identification by themselves.

The Gram-corrected global projection and the two actual
Gram factors in MA-GP-008/MA-CGP-003 remain valid.
The theorem neither rules out those different maps nor
changes the proved source-return operators.  It does
not refute an $A$-supported scalar trace identity,
remove the arithmetic FT002 divisor, or claim that
every conceivable factorization has failed.  Its
obstruction is a nonzero inner product of actual
source images in the specified global norm, witnessed
before any trace is taken.

\paragraph{Grothendieck placement.}
The character $\chi_4$ is the rank-one finite Artin
representation of $\operatorname{Gal}(\mathbb Q(i)/\mathbb Q)$
which is $-1$ on complex conjugation.  On
$\operatorname{Spec}\mathbb Z[1/2]$ it is the associated
finite-monodromy etale local system; its unramified Euler
polynomial is $1-\chi_4(p)T$, and its inertia-invariant
space at $2$ is zero.  The exact source conventions are
SGA~1, Expos\'e~V, \S7, and Deligne, \emph{Valeurs de
fonctions L}, Definition~6.1 and \S\S6.2--6.3, already
audited in MA-CH/MA-CQ and not claimed freshly reread here.
The norm and orthogonal projection in (MA-PI.3--8) are
analytic metric data; the obstruction does not alter
this Artin representation or assert a new weight or
functoriality statement.  The inherited cote~1
trace/determinant programme, BR-TR-01/02, scans~079--080
and~111--112, requires actual comparison maps before
trace cyclicity is applied.  Those archive locators are
inherited; the present proof determines one such proposed
map identity to be false, while retaining the actual Gram
and nuclear-return comparisons already constructed.

\section{The same obstruction in the printed twisted Hilbert realization}
\label{sec:printed-projection-obstruction}

\paragraph{Status and purpose.}
The complete author reading and two independent whole-proof
readings have passed; their actual check scopes are recorded separately.  MA-PI-001 fixes the corrected centered
symmetric norm.  Here the original printed Hilbert realization is
derived from its own definitions, and its orthogonal projector is
tested directly.  The change of realization below is an actual
isometry, not an appeal to equivalence of norms.

Retain $\chi=\chi_4$, $S=\{\infty,2\}$, the actual finite vector
$\theta_2$, the canonical adelic source $\varphi_f$, and the
physical cutoff $\mathfrak B_{S,Y}$ from (MA-PI.1--9).
The outside-$S$ factors of $\varphi_f$ are
$\mathbf1_{\mathbb Z_p}$.  For every odd Schwartz $f$, both
adelic polar conditions hold at the ramified finite place.
Let $A$ denote the source's twist parameter, to distinguish it
from the symmetric width $a$ in MA-PI-001.

\paragraph{The exact source coordinates and metric.}
Masullo Definition~2.2.3.1, equation~(2.2.3.1),
printed/PDF pp.~81--82, defines the \emph{uncentered} theta
function
\[
 \Phi_{\mathbb Q}(\varphi)(g)
       =\sum_{\xi\in\mathbb Q^\times}\varphi(\xi g).
 \tag{MA-PI.32}
\]
Definition~9.1.2.1, equation~(9.1.2.3), p.~399, defines
the norm on $W^{s,2}(C_{\mathbb Q})$ by
\[
 \|F\|_{s,2}^2
     =\int_{C_{\mathbb Q}}|F(g)|^2
              (1+|\log|g||^2)^{s/2}\,d\mu(g).
 \tag{MA-PI.33}
\]
The compact-character component is denoted
$\mathcal H_s^\chi=W^{s,2}(C_{\mathbb Q})(\chi|_{C_{\mathbb Q}^1})$.
Definition~10.1.1.1 and its preceding equations
(10.1.1.1--3), pp.~434--435, realize the twist by
the actual function $|g|^A\Phi_{\mathbb Q}(\varphi)(g)$.
The same twist appears in the invariant PW level of
Definition~11.2.1.3, equations~(11.2.1.7--8), pp.~488--489.
Let $Q^{\mathrm{src}}_{S,Y,A,s}$ be that printed
orthogonal projection in $\mathcal H_s^\chi$.

Choose the positive archimedean norm splitting $\nu$ and write
$g=\nu(e^x)k$, where $k\in C_{\mathbb Q}^1$.
As in the paragraph following (9.1.2.3), the Haar measure
has the exact product decomposition $dx\,d\mu_1(k)$.
Put $m=\mu_1(C_{\mathbb Q}^1)>0$; this retains the source's
Haar normalization rather than silently assigning it mass one.
For the actual canonical character source one has
\[
 \begin{split}
 (T_Af)(\nu(e^x)k)
   &:=|g|^A\Phi_{\mathbb Q}(\varphi_f)(g)\\
   &=\chi(k)^{-1}e^{b x}Y_f(x),
          \qquad b=A-\tfrac12,
 \end{split}
 \tag{MA-PI.34}
\]
where $Y_f$ is exactly (MA-PI.2).
Indeed the actual finite vector restricts the theta sum to
integers with coefficient $\chi_4(n)$, and the two signs
give the factor $2$ in (MA-PI.2).  Its factor $e^{x/2}$
is the centering, while the source twist is $e^{Ax}$.
The compact-character dependence is that of MA-CGP.2--7;
here $\chi^{-1}=\chi$, so no inverse-character ambiguity
affects (MA-PI.34).

Define the one-sided Hilbert space and an actual unitary map by
\[
 \begin{split}
 H^{\mathrm{one}}_{b,s}
    &=L^2\bigl(\mathbb R,e^{2bx}(1+x^2)^{s/2}\,dx\bigr),\\
 (U_bF)(\nu(e^x)k)
    &=m^{-1/2}\chi(k)^{-1}e^{bx}F(x).
 \end{split}
 \tag{MA-PI.35}
\]
Equations (MA-PI.33) and $|\chi(k)|=1$ prove that $U_b$
is isometric onto $\mathcal H_s^\chi$: the compact-character
projector shows that every function in that component has
the displayed form almost everywhere.  Its inverse is the
well-defined fiber coefficient
\[
 (U_b^{-1}H)(x)=m^{-1/2}e^{-bx}
        \int_{C_{\mathbb Q}^1}H(\nu(e^x)k)\chi(k)\,d\mu_1(k).
\]
This formula is understood for almost every $x$, using
Fubini and Cauchy--Schwarz on the compact fiber.  In particular
\[
 T_A f=\sqrt m\,U_bY_f,\qquad
 \langle T_Af,T_Ah\rangle_{s,2}
     =m\int_{\mathbb R}Y_f(x)\overline{Y_h(x)}
                           e^{2bx}(1+x^2)^{s/2}\,dx.
 \tag{MA-PI.36}
\]
Consequently $U_b^{-1}Q^{\mathrm{src}}_{S,Y,A,s}U_b$
is the orthogonal projector onto the actual inverse image
of the printed PW range.  This assertion follows from
unitarity; no assertion about orthogonal projectors under
a merely bounded norm equivalence is used.

\paragraph{MA-PI-002 (failure for the original printed projection).}
For every $0<A<1$ and every $s\geq0$, there is $Y_0>0$
such that, for every $Y\geq Y_0$ and every smooth physical
profile (MA-PI.4), there is an odd Schwartz function $f_Y$
flat in a neighborhood of zero for which
\[
 Q^{\mathrm{src}}_{S,Y,A,s}T_Af_Y
                     \ne T_A(\mathfrak b_Yf_Y).
 \tag{MA-PI.37}
\]
The input $f_Y\otimes\theta_2$ is in the original
coordinate-flat LF source of MA-GA.2.  Its physical
cutoff output is genuine $S$-local Schwartz data satisfying
both polar conditions, so the right side is the actual
printed global theta map on that output.
In particular (MA-PI.37) applies to $A=1/2$, the exact
twist specified in the proof of Theorem~12.4.1.1 on p.~515.
The original printed canonical intertwining is therefore
false, without a norm-realization qualification.
This does not assert the failure of every different
Gram-corrected factorization.

\begin{proof}
\emph{1. Original source membership and physical action.}
The finite Fourier calculation (MA-PI.9) shows that the
canonical data have shell $|\xi_2|_2=4$ and that
\[
 \mathfrak B_{S,Y}(f\otimes\theta_2)
       =(\mathfrak b_Yf)\otimes\theta_2,
 \quad
 \mathfrak b_Yf=\mathcal F_\infty^{-1}
                      \bigl(\rho_Y(4|\xi|)\widehat f(\xi)\bigr).
 \tag{MA-PI.38}
\]
This is a literal identity on the stated canonical Schwartz
data.  The norm multiplier is even, so the equivalent
expression $\mathcal F_\infty M_\rho\mathcal F_\infty^{-1}$
gives the same operator: $\mathcal F_\infty^2$ is reflection
and reflection commutes with $M_\rho$.  The same observation
applies to coordinatewise reflected finite Fourier conventions.
Thus the unchanged physical norm cutoff of the source is used.

If $\widehat h$ is compact smooth and supported in
$|\xi_\infty|<Y/4$, the entire $S$-Fourier support is
compact with product norm $<Y$.  Both polar conditions
hold at $2$.  Definition~7.3.2.3, (7.3.2.12--17),
pp.~354--355, with $S'=S$, therefore puts the actual
theta image in the source invariant PW level.  Twisting
by $|g|^A$ and taking the closure gives
$T_Ah\in\operatorname{Ran}Q^{\mathrm{src}}_{S,Y,A,s}$.
This is direct membership, not an assumed equality of
different possible PW ranges.

\emph{2. A nonzero cross pairing for this exact one-sided line.}
Put $b=A-1/2\in(-1/2,1/2)$.
Equations (MA-PI.10--20) already prove, for every such $b$,
the exact one-sided theta Gram identity and that its
inverse-Fourier distribution is not compactly supported.
Those formulas were derived for the whole open range of
$b$, before the specialization $b=-a$ in MA-PI-001.
Their distributional separation argument therefore
supplies odd Schwartz functions $f,h$ whose Fourier
transforms are compact smooth with respective supports
\[
 \begin{gathered}
 \operatorname{supp}\widehat f\subset\{|\xi|>1/2\},\qquad
 \operatorname{supp}\widehat h\subset\{0<|\xi|<1/8\},\\
 D_b:=\int_{\mathbb R}Y_f(x)\overline{Y_h(x)}e^{2bx}\,dx\ne0.
 \end{gathered}
 \tag{MA-PI.39}
\]
These are genuine Schwartz source vectors.  No zero-location
or mean-square hypothesis is used; the separation follows
from the proved factorial growth of the entire beta product.

\emph{3. Dilation in the printed polynomial norm.}
Use the actual dilations $f_L(t)=e^{L/2}f(e^Lt)$ and
$h_L(t)=e^{L/2}h(e^Lt)$, and set $Y=e^L$.
Their Fourier supports and the two profile plateaus give
\[
 \mathfrak b_Y f_L=0,\qquad
 \mathfrak b_Y h_L=h_L,\qquad
 T_Ah_L\in\operatorname{Ran}Q^{\mathrm{src}}_{S,Y,A,s}.
 \tag{MA-PI.40}
\]
Exact theta scaling gives $Y_{f_L}(x)=Y_f(x+L)$ and the
same formula for $h$.  Substituting into the actual
printed inner product (MA-PI.36), for $L\geq1$, yields
\[
 \begin{split}
 &m^{-1}e^{2bL}L^{-s}
        \langle T_A f_L,T_A h_L\rangle_{s,2}\\
 &\qquad=\int_{\mathbb R}Y_f(x)\overline{Y_h(x)}e^{2bx}
                  \frac{(1+(x-L)^2)^{s/2}}{L^s}\,dx
                         \longrightarrow D_b\ne0.
 \end{split}
 \tag{MA-PI.41}
\]
The polynomial ratio tends to one and is bounded by
$C_s(2+|x|)^s$.  Both theta images belong to $\mathcal E$,
so dominated convergence proves (MA-PI.41).  The pairing
is therefore nonzero for every sufficiently large $L$,
independently of the choice of transition profile.
Orthogonality of the actual printed projector and
(MA-PI.40) imply
\[
 \bigl\langle
 Q^{\mathrm{src}}_{S,Y,A,s}T_Af_L
       -T_A(\mathfrak b_Y f_L),\ T_Ah_L\bigr\rangle_{s,2}
        =\langle T_Af_L,T_Ah_L\rangle_{s,2}\ne0.
 \tag{MA-PI.42}
\]

\emph{4. A finite witness in the original flat source.}
Fix such $Y$ and its smooth profile, and choose the even
compact cutoff $\omega=1$ near zero as in (MA-PI.27).
Set $f_{L,\varepsilon}=f_L-r_\varepsilon$, where
$r_\varepsilon(t)=\omega(t/\varepsilon)f_L(t)
                  =\varepsilon g_\varepsilon(t/\varepsilon)$.
The $g_\varepsilon$ form a bounded family of odd compact
Schwartz functions.  Using (MA-PI.28) and the exact
one-sided norm in (MA-PI.36) gives
\[
 \|T_A r_\varepsilon\|_{s,2}
      \leq C\varepsilon^{3/2+b}
                    (1+|\log\varepsilon|)^{s/2}
       =C\varepsilon^{A+1}(1+|\log\varepsilon|)^{s/2}
                          \longrightarrow0.
 \tag{MA-PI.43}
\]
For completeness, the substitution $x=y+\log\varepsilon$
in the squared integral contributes $\varepsilon^3$
from theta scaling and $\varepsilon^{2b}$ from the
exponential weight.  The remaining polynomial ratio is
bounded by a constant times
$(1+|\log\varepsilon|)^s(1+|y|)^s$; uniform exponential
Schwartz bounds for $Y_{g_\varepsilon}$ finish the estimate.
The exponent $A+1$ is positive throughout the asserted range.

For the fixed profile, (MA-PI.29) proves that
$\mathfrak b_Y r_\varepsilon\to0$ in odd Schwartz space:
its Fourier transform is smooth with fixed compact support,
and every derivative tends uniformly to zero.  The actual
theta map is continuous into $\mathcal E$, and multiplication
by $e^{bx}$ is continuous from $\mathcal E$ into polynomial
$L^2$.  Hence
\[
       \|T_A(\mathfrak b_Y r_\varepsilon)\|_{s,2}\to0.
 \tag{MA-PI.44}
\]
The actual orthogonal projector is bounded.  Equations
(MA-PI.43--44) show that the pairing in (MA-PI.42), with
$f_L$ replaced by $f_{L,\varepsilon}$, converges to its
nonzero value.  Choose a sufficiently small positive
$\varepsilon$ and put $f_Y=f_{L,\varepsilon}$.
This function is Schwartz and flat near the archimedean
origin; its finite factor is supported away from zero.
It is therefore an actual vector of the original inherited
LF source, and it proves (MA-PI.37).  The right side
remains genuine Schwartz polar-free data by (MA-PI.38).
\end{proof}

\paragraph{Source and conclusion boundaries.}
The exact source definitions, compact Haar product decomposition,
twisting convention, invariant PW range and orthogonal projection
were read in full at the locators given above.  The physical
cutoff is the source's Step~II, p.~516, and composition on
p.~517, in Theorem~12.4.1.1.  Its proof explicitly uses
$W_{1/2}^{s,2}$ on p.~515.  Formula (MA-PI.34) therefore
identifies that actual case with $b=0$, without a centering
guess or an equivalent-norm substitution.

The finite Artin representation, inertia convention and
Grothendieck placement are exactly those recorded after
MA-PI-001; changing the analytic Hilbert realization does
not change that rank-one object.  The result determines
the specified canonical theta/physical-cutoff map identity
to be false in both the corrected and the literal printed
realizations.  It does not attribute an unprinted equation
to the source, and it does not rule out a different
Gram-corrected comparison, alter the valid nuclear returns,
or remove the arithmetic FT002 condition.

\section{A rational trivial-sector obstruction with both polar conditions}
\label{sec:rational-physical-projection-obstruction}

\paragraph{Proof-review status.}
The complete author and independent whole-proof readings have
passed. Their exact proof and source-interface scopes are recorded
in the separate check events.

We now use exactly the rational trivial sector, $S=\{\infty\}$
and $S'=S$, with no added unit-shell transition places,
and the actual theta map and PW source of (MA-GP.1)--(MA-GP.2):
\[
 \begin{gathered}
 \mathcal U_0=\{f\in\mathcal S(\mathbb R): f(-t)=f(t),\quad
                         f(0)=0,\ \int_{\mathbb R}f(t)dt=0\},\\
 (\Theta f)(x)=2e^{x/2}\sum_{n\ge1}f(ne^x),\qquad
 \mathcal U_{0,Y}=\{f\in\mathcal U_0:
                  \operatorname{supp}\widehat f\subseteq[-Y,Y]\}.
 \end{gathered}
 \tag{MA-PI.70}
\]
The Fourier kernel is $e^{-2\pi it\xi}$. The range
$\Theta\mathcal U_0$ is in the all-exponential Schwartz
space $\mathcal E$, and the map into that space is continuous.
This is the actual source identification of MA-TR-101 and
the Poisson identity (MA-GP.8), not a prescribed divisor model.

Two specified Hilbert metrics will be treated independently:
\[
 \begin{aligned}
 H^{\rm sym}_{a,s}
  &=L^2\bigl(\mathbb R,e^{2a|x|}(1+x^2)^{s/2}dx\bigr),
             &&0<a<1/2,\ s\ge0,\\
 H^{\rm one}_{b,s}
  &=L^2\bigl(\mathbb R,e^{2bx}(1+x^2)^{s/2}dx\bigr),
             &&|b|<1/2,\ s\ge0.
 \end{aligned}
 \tag{MA-PI.71}
\]
In either one of these fixed metrics, put
\[
 J_Y=\overline{\Theta\mathcal U_{0,Y}},\qquad Q_Y=P_{J_Y}.
 \tag{MA-PI.72}
\]
Thus $Q_Y$ is its actual orthogonal projection. In the symmetric
case it is exactly the corrected global projector of MA-GP.
The one-sided assertion is a separately proved statement in
its displayed metric; equivalent norms are not used to transfer
orthogonal projections.

Let $\rho_Y\in C^\infty([0,\infty))$ satisfy
\[
 0\le\rho_Y\le1,\quad \rho_Y(r)=1\ (r\le Y/2),\quad
 \rho_Y(r)=0\ (r\ge2Y),\qquad
 \mathfrak b_Y=\mathcal F^{-1}M_{\rho_Y(|\xi|)}\mathcal F.
 \tag{MA-PI.73}
\]
This is the original physical Fourier cutoff for this $S$.

\paragraph{MA-PI-003 (the principal rational canonical intertwining fails).}
For either fixed metric in (MA-PI.71), there is $Y_0$ such that,
for every $Y\ge Y_0$ and every profile (MA-PI.73), there exists
an even Schwartz function $f_Y$, identically zero near the
additive origin, satisfying
\[
 f_Y\in\mathcal U_0,\qquad \mathfrak b_Y f_Y\in\mathcal U_0,
 \qquad Q_Y\Theta f_Y\ne\Theta\mathfrak b_Y f_Y.
 \tag{MA-PI.74}
\]
Both actual theta images therefore belong to $\mathcal E$.
The failure is not caused by an undefined polar output.
The input lies in the original flat source, with its two
global polar conditions. No information about the locations
or multiplicities of zeta zeros is used.

\begin{proof}
\emph{1. Enforcing the two source moments before the Gram calculation.}
Fix $b\in(-1/2,1/2)$ for this calculation and set
$\sigma=1/2+b$, $d=1-\sigma$.
Given $B\in C_c^\infty(\mathbb R)$, put
\[
 A=(\partial_r+\sigma)B,\qquad
 u(\xi)=\xi^{\sigma-1}A(\log\xi)\quad(\xi>0),
 \qquad u(-\xi)=u(\xi),\qquad f=\mathcal F^{-1}u.
 \tag{MA-PI.75}
\]
The frequency vector is smooth with compact support away from
zero. Integration by parts gives
\[
 \int_0^\infty u(\xi)d\xi
       =\int_{\mathbb R}e^{\sigma r}
                  (B'(r)+\sigma B(r))dr=0.
 \tag{MA-PI.76}
\]
Hence $f(0)=2\int_0^\infty u=0$ and
$\int f=\widehat f(0)=u(0)=0$. Thus $f\in\mathcal U_0$.
The differential operation does not enlarge the support of
$B$ in logarithmic frequency.

Write $M_f(z)=\int_0^\infty f(t)t^{z-1}dt$ and
$\widehat B(t)=\int B(r)e^{-itr}dr$. Actual unfolding, first
in $\Re z>1$ and then by continuation, gives
\[
 \mathcal M(\Theta f)(z-1/2)=2\zeta(z)M_f(z).
 \tag{MA-PI.77}
\]
The even inverse Fourier integral, with an Abel factor
$e^{-\varepsilon t}$ before interchanging integrals, gives
\[
 M_f(z)=\frac{\Gamma_{\mathbb R}(z)}
                    {\Gamma_{\mathbb R}(1-z)}M_u(1-z),
                         \qquad0<\Re z<1.
 \tag{MA-PI.78}
\]
Indeed $f(t)=2\int_0^\infty u(r)\cos(2\pi tr)dr$.
The Abel integral tends to
$2(2\pi)^{-z}\Gamma(z)\cos(\pi z/2)M_u(1-z)$,
uniformly over the compact positive support of $u$;
reflection and duplication give precisely (MA-PI.78).
On the line $z=\sigma+it$, one has
$M_u(1-z)=(\sigma+it)\widehat B(t)$.

The actual Poisson identity $\Theta f(x)=\Theta\widehat f(-x)$
in (MA-GP.8), together with (MA-PI.77)--(MA-PI.78), fixes
the functional-equation normalization
\[
 \frac{\Gamma_{\mathbb R}(z)}{\Gamma_{\mathbb R}(1-z)}
                  \zeta(z)=\zeta(1-z).
 \tag{MA-PI.79}
\]
For example use a nonzero $f$ constructed in (MA-PI.75)
in the common strip, cancel its nonzero Mellin transform on an open set,
and use meromorphic continuation. This is only the usual
rational functional equation in the precise Fourier convention
needed here; it contains no zero-location assertion.

For $B,C$ and their corresponding $f,h$, Mellin Plancherel now gives
\[
 \begin{aligned}
 \int_{\mathbb R}(\Theta f)(x)
       \overline{(\Theta h)(x)}e^{2bx}dx
    &=\frac1{2\pi}\int_{\mathbb R}
             \mathscr V_b(t)\widehat B(t)
                               \overline{\widehat C(t)}dt,\\
 \mathscr V_b(t)&=4(t^2+\sigma^2)
                       \zeta(d+it)\zeta(d-it).
 \end{aligned}
 \tag{MA-PI.80}
\]
In particular this is a Gram form of genuine sources with
both polar conditions, not a formal positive symbol.

\emph{2. The constrained Gram distribution has unbounded support.}
The real-axis function $\mathscr V_b$ is even, nonnegative,
and polynomially bounded. To check the last property directly,
the usual first partial-summation identity
\[
 \zeta(z)=\frac{z}{z-1}
               -z\int_1^\infty \{x\}x^{-z-1}dx,
                   \qquad\Re z>0,\ z\ne1,
 \tag{MA-PI.81}
\]
gives $\zeta(d+it)=O_d(1+|t|)$ on the fixed line
$0<d<1$. The integral follows initially from summing
$\lfloor x\rfloor$ in $\Re z>1$ and continues there
absolutely for $\Re z>0$. Thus $\mathscr V_b$ defines
a tempered distribution.

The same expression in a complex variable $t$ is entire.
Its only possible poles are the simple zeta poles at
$t=\pm i\sigma$; the factor $t^2+\sigma^2$ cancels
both. Its growth on the imaginary axis is not of
exponential type. From (MA-PI.79) and the gamma recurrence,
\[
 \zeta(-2n-1)=(-1)^{n+1}
       \frac{2(2n+1)!}{(2\pi)^{2n+2}}\zeta(2n+2),
                         \qquad n\ge0.
 \tag{MA-PI.82}
\]
Put $T_n=2n+1+d$. Since both zeta values at the positive
arguments in the following display exceed one, for $n\ge1$
\[
 \begin{aligned}
 |\mathscr V_b(iT_n)|
  &=4(T_n^2-\sigma^2)
          |\zeta(-2n-1)|\zeta(2n+1+2d)\\
  &\ge \frac{8(T_n^2-\sigma^2)(2n+1)!}{(2\pi)^{2n+2}}.
 \end{aligned}
 \tag{MA-PI.83}
\]
This exceeds every bound $C(1+T_n)^N e^{RT_n}$ eventually.
If the inverse Fourier distribution $\kappa_b$ of
$\mathscr V_b$ had compact support, its finite order on a
fixed compact neighborhood, applied to a cutoff times
$e^{-izx}$, would give exactly such an entire exponential
bound, as proved in (MA-PI.18). Hence $\kappa_b$ is not
compactly supported. It is even.

Consequently, for every two cut positions $A_0<B_0$ there
are compact smooth $B,C$, supported respectively in
$(B_0,\infty)$ and $(-\infty,A_0)$, with the Gram pairing
(MA-PI.80) nonzero. Here is the exact support argument.
If all such pairings vanished, testing $B$ by a smooth
approximation to a delta at $r_0>B_0$ would show that
$\kappa_b$ vanishes on $(-\infty,A_0-B_0)$. Each compact
test in that interval permits an $r_0$ with its translate
below $A_0$. Evenness would give vanishing beyond the
opposite end too, contradicting unbounded support.
A nonzero limiting distribution pairing already has a
nonzero ordinary smooth approximate-delta member. Thus
the resulting $B,C$, and hence $u,v,f,h$, are actual
smooth functions throughout.

Choose $A_0=\log(1/2)$ and $B_0=\log2$. Equations
(MA-PI.75)--(MA-PI.76) supply $f,h\in\mathcal U_0$ with
\[
 \operatorname{supp}\widehat f\subset\{|\xi|>2\},\qquad
 \operatorname{supp}\widehat h\subset\{0<|\xi|<1/2\},
 \quad D_b:=\langle\Theta f,\Theta h\rangle_{L^2(e^{2bx}dx)}
                              \ne0.
 \tag{MA-PI.84}
\]
These supports are also compact. No arbitrary preassigned
pair of bounded bands is asserted to work.

\emph{3. Physical dilation in each specified metric.}
Put $Y=e^L$, $f_L(t)=e^{L/2}f(e^Lt)$ and
$h_L(t)=e^{L/2}h(e^Lt)$. The two moments stay zero,
and actual summation gives
\[
 \Theta f_L(x)=\Theta f(x+L),\qquad
 \Theta h_L(x)=\Theta h(x+L),\qquad
 \mathfrak b_Y f_L=0,\qquad h_L\in\mathcal U_{0,Y}.
 \tag{MA-PI.85}
\]
The last two assertions use only the plateaus in
(MA-PI.73), so they hold uniformly over their profiles.
For the symmetric metric choose $b=-a$ in step 1.
Then the same dominated-convergence calculation as
(MA-PI.25) gives
\[
 e^{-2aL}L^{-s}
       \langle\Theta f_L,\Theta h_L\rangle_{H^{\rm sym}_{a,s}}
                        \longrightarrow D_{-a}\ne0.
 \tag{MA-PI.86}
\]
For the one-sided metric use its own parameter $b$ in
step 1. Direct substitution $x\mapsto x-L$ gives
\[
 e^{2bL}L^{-s}
       \langle\Theta f_L,\Theta h_L\rangle_{H^{\rm one}_{b,s}}
 =\int\Theta f(x)\overline{\Theta h(x)}e^{2bx}
       \frac{(1+(x-L)^2)^{s/2}}{L^s}dx
                         \longrightarrow D_b\ne0.
 \tag{MA-PI.87}
\]
For $L\ge1$ the polynomial ratio is at most
$C_s(2+|x|)^s$ and tends to one. Both theta images are
in $\mathcal E$, which proves domination also for $s=0$.
There is therefore one $Y_0$ for the fixed metric such
that the pairing is nonzero for all $Y\ge Y_0$,
independently of the profile. Since $\Theta h_L\in J_Y$,
\[
 \langle Q_Y\Theta f_L-\Theta\mathfrak b_Y f_L,
                      \Theta h_L\rangle_H
                 =\langle\Theta f_L,\Theta h_L\rangle_H\ne0.
 \tag{MA-PI.88}
\]
The bandlimited input $f_L$ is not flat. The final step
produces a genuine flat input and retains both polar
conditions also after the cutoff.

\emph{4. Two fixed flat correctors for a fixed physical profile.}
Fix one $Y\ge Y_0$ and its profile. The inverse Fourier
transform $k_Y=\mathcal F^{-1}(\rho_Y(|\xi|))$ is a real
even Schwartz function, and is not constant. Indeed a
constant Schwartz function is zero, whereas the Fourier
multiplier equals one near zero. On even compact smooth
functions supported away from the origin the two
functionals
\[
 I(v)=\int_{\mathbb R}v(t)dt,\qquad
 \ell_Y(v)=(\mathfrak b_Yv)(0)
                         =\int_{\mathbb R}k_Y(t)v(t)dt
 \tag{MA-PI.89}
\]
are linearly independent. Otherwise testing arbitrarily
small even bumps away from zero would force $k_Y$ to be
constant on $\mathbb R\setminus\{0\}$, and hence everywhere
by continuity. Thus there exist fixed real even functions
$\psi_1,\psi_2\in C_c^\infty(\mathbb R\setminus\{0\})$
such that
\[
 T_Y=\begin{pmatrix}
 I(\psi_1)&I(\psi_2)\\
 \ell_Y(\psi_1)&\ell_Y(\psi_2)
 \end{pmatrix}
                         \quad\hbox{is invertible}.
 \tag{MA-PI.90}
\]
This is finite-dimensional linear algebra after the
functional independence; no uniform choice over profiles
is claimed or needed.

Choose a fixed even $\omega\in C_c^\infty(\mathbb R)$
equal to one near zero. For $0<\varepsilon\le1$, put
\[
 \begin{gathered}
 r_\varepsilon(t)=\omega(t/\varepsilon)f_L(t),\qquad
 \binom{c_1(\varepsilon)}{c_2(\varepsilon)}
       =T_Y^{-1}\binom{I(r_\varepsilon)}{\ell_Y(r_\varepsilon)},\\
 f_{L,\varepsilon}=f_L-r_\varepsilon
                    +c_1(\varepsilon)\psi_1
                    +c_2(\varepsilon)\psi_2,\qquad
 e_\varepsilon=f_L-f_{L,\varepsilon}.
 \end{gathered}
 \tag{MA-PI.91}
\]
The function $f_{L,\varepsilon}$ is even, Schwartz, and
zero near zero. Equations (MA-PI.85) and (MA-PI.90)--(MA-PI.91)
give $I(f_{L,\varepsilon})=0$ and
$\ell_Y(f_{L,\varepsilon})=0$. Since the multiplier
is one at zero, also
$I(\mathfrak b_Y f_{L,\varepsilon})=I(f_{L,\varepsilon})=0$.
Consequently
\[
 f_{L,\varepsilon},\quad e_\varepsilon,\quad
                 \mathfrak b_Y f_{L,\varepsilon},\quad
                 \mathfrak b_Y e_\varepsilon
                            \in\mathcal U_0.
 \tag{MA-PI.92}
\]
In particular every theta image subsequently used is
in its actual all-exponential source image.

The next estimates are necessary: cutting off near zero
does not tend to zero in every Schwartz seminorm.
Because $f_L$ is even and $f_L(0)=0$, both its value and
first derivative vanish at zero. Therefore
\[
 r_\varepsilon(t)=\varepsilon^2g_\varepsilon(t/\varepsilon),
 \quad g_\varepsilon(y)=\omega(y)
                    f_L(\varepsilon y)/\varepsilon^2,
 \quad |c_1(\varepsilon)|+|c_2(\varepsilon)|=O(\varepsilon^3).
 \tag{MA-PI.93}
\]
The $g_\varepsilon$ form a bounded family of even compact
smooth functions with $g_\varepsilon(0)=0$. For derivatives
of order zero and one this follows from Taylor's formula
using the two vanishings; higher derivatives of the
non-cutoff factor have coefficients $\varepsilon^{j-2}$,
$j\ge2$. The last bound in (MA-PI.93) follows from
$\|r_\varepsilon\|_{L^1}=O(\varepsilon^3)$,
boundedness of $k_Y$, and the fixed inverse $T_Y^{-1}$.

\emph{5. A uniform actual theta estimate for the flattened error.}
Fix a smooth real step $\phi$ equal to one far left and
zero far right, and put $P(x)=\phi(x)e^{-x/2}$.
For every even Schwartz $v$ with $v(0)=0$, actual
Poisson summation gives
\[
 R_v(x):=\Theta v(x)-I(v)P(x)\in\mathcal E,
 \tag{MA-PI.94}
\]
continuously as a function of $v$ on that closed Schwartz
subspace. At the left end the exact Poisson formula is
$\Theta v(x)=I(v)e^{-x/2}+\Theta\widehat v(-x)$,
whose last term decreases faster than every exponential;
at the right end use the direct theta sum. The same
estimates apply to every derivative. This proves the
claimed continuity, including for bounded Schwartz
families, without applying theta continuity on $\mathcal U_0$
to a vector outside that subspace.

Set $q_\varepsilon=I(g_\varepsilon)$, a bounded scalar.
By (MA-PI.90)--(MA-PI.93),
$\sum_j c_jI(\psi_j)=I(r_\varepsilon)=\varepsilon^3q_\varepsilon$.
Scaling the actual theta series and subtracting its
polar part gives the exact identity
\[
 \begin{split}
 \Theta e_\varepsilon(x)
   ={}&\varepsilon^{5/2}R_{g_\varepsilon}(x-\log\varepsilon)
                      -\sum_{j=1}^2c_j(\varepsilon)R_{\psi_j}(x)\\
     &+\varepsilon^3q_\varepsilon e^{-x/2}
            [\phi(x-\log\varepsilon)-\phi(x)].
 \end{split}
 \tag{MA-PI.95}
\]
All terms on the right are in either specified $H$.
Put $c=a$ for the symmetric metric and $c=|b|$ for
the one-sided metric; thus $0\le c<1/2$ and the
one-sided norm is bounded above by the symmetric
$c$ norm. Translation of a bounded $\mathcal E$ family
by $-\log\varepsilon$ has norm at most
$C\varepsilon^{-c}(1+|\log\varepsilon|)^{s/2}$.
For the last term the bracket is supported in an
interval contained in $[\log\varepsilon-C,C]$.
Its absolute value is bounded independently of
$\varepsilon$. Hence
\[
 \begin{split}
 &\bigl\|\varepsilon^3e^{-x/2}
        [\phi(x-\log\varepsilon)-\phi(x)]\bigr\|_H^2\\
 &\qquad\le C\varepsilon^6
      \int_{\log\varepsilon-C}^{C}
           e^{-x+2c|x|}(1+x^2)^{s/2}dx\\
 &\qquad\le C'\varepsilon^{5-2c}
                          (1+|\log\varepsilon|)^s.
 \end{split}
 \tag{MA-PI.96}
\]
The fixed bounded positive part of the integral is
absorbed in the last bound. Combining (MA-PI.93)--(MA-PI.96)
proves
\[
 \|\Theta e_\varepsilon\|_H
   \le C\bigl[\varepsilon^{5/2-c}
                 (1+|\log\varepsilon|)^{s/2}+\varepsilon^3\bigr]
                                      \longrightarrow0.
 \tag{MA-PI.97}
\]
This estimate has kept the canceling polar pieces
together; none is separately presumed to be a vector
of the narrow Hilbert space.

There is also convergence after the fixed physical
cutoff. Direct differentiation of the Fourier integral gives
\[
 \sup_\xi|\partial_\xi^j\widehat r_\varepsilon(\xi)|
                         \le C_j\varepsilon^{j+3}.
 \tag{MA-PI.98}
\]
After multiplication by the fixed compact smooth
multiplier $\rho_Y(|\xi|)$, these estimates and
$c_j=O(\varepsilon^3)$ show that
$\mathfrak b_Y e_\varepsilon\to0$ in Schwartz space,
with every seminorm $O(\varepsilon^3)$. It belongs
to $\mathcal U_0$ by (MA-PI.92). The genuine continuous
theta map on that subspace therefore gives
\[
                 \|\Theta\mathfrak b_Ye_\varepsilon\|_H
                                      \longrightarrow0.
 \tag{MA-PI.99}
\]
The derivative constants may depend on the already
fixed profile; the flat approximation is chosen
after that profile and finite level have been fixed.

Finally use the bounded orthogonal projection $Q_Y$
and test against the actual vector $\Theta h_L\in J_Y$.
Equations (MA-PI.88), (MA-PI.97) and (MA-PI.99) give
\[
 \begin{split}
 &\langle Q_Y\Theta f_{L,\varepsilon}
            -\Theta\mathfrak b_Y f_{L,\varepsilon},\Theta h_L\rangle_H\\
 &\qquad\longrightarrow
                       \langle\Theta f_L,\Theta h_L\rangle_H\ne0.
 \end{split}
 \tag{MA-PI.100}
\]
Some finite positive $\varepsilon$ therefore supplies
$f_Y=f_{L,\varepsilon}$ satisfying (MA-PI.74).
Every input and output asserted there is an actual
Schwartz source vector with the required moments.
\end{proof}

\paragraph{The exact printed principal realization.}
The source isometry in (MA-PI.32)--(MA-PI.36) also has
the following direct trivial-character specialization.
Put $m=\mu_1(C_{\mathbb Q}^1)$ and $b=A-1/2$. For
$g=\nu(e^x)k$ and the source $f\otimes\mathbf1_{\widehat{\mathbb Z}}$,
\[
 \begin{aligned}
 (T_A^0f)(g)
   &:=|g|^A\Phi_{\mathbb Q}
                   (f\otimes\mathbf1_{\widehat{\mathbb Z}})(g)
       =e^{bx}\Theta f(x),\\
 (U_b^0F)(\nu(e^x)k)&=m^{-1/2}e^{bx}F(x).
 \end{aligned}
 \tag{MA-PI.101}
\]
The finite integer balls are unchanged by finite units,
so the theta function is constant on the compact fiber.
The two signs give the factor two in (MA-PI.70).
The exact printed norm (MA-PI.33) makes $U_b^0$ unitary
from $H^{\rm one}_{b,s}$ onto the trivial compact-character
component of $W^{s,2}(C_{\mathbb Q})$, and
$T_A^0=\sqrt m\,U_b^0\Theta$.
For $S=\{\infty\}$ the printed invariant PW source is
exactly the compact Fourier support condition in
(MA-PI.70). An arbitrary source can be replaced by its
even part without changing its theta sum or either
polar condition, and this preserves the Fourier band.
Therefore its actual printed orthogonal projection obeys
\[
 Q^{\rm src}_{\infty,Y,A,s}
                       =U_b^0Q_Y^{\rm one}(U_b^0)^{-1}.
 \tag{MA-PI.102}
\]
This is a unitary identity derived from the source
definitions, not a change to an equivalent metric.
For every $0<A<1$, (MA-PI.74) in its one-sided metric
consequently gives genuine flat witnesses with
\[
 Q^{\rm src}_{\infty,Y,A,s}T_A^0f_Y
                       \ne T_A^0\mathfrak b_Yf_Y.
 \tag{MA-PI.103}
\]
In particular this proves the same principal rational
canonical obstruction at the exact twist $A=1/2$
used on p.~515, with both actual source and cutoff
output polar free. The compact Haar mass has been
retained throughout the identification.

\paragraph{Source attribution and the precise conclusion.}
The theta source and rational PW support are the actual
MA-GP.1--2 objects. Their source locators are Masullo,
Definition 7.3.2.1, (7.3.2.1--7), p.~351, and
Definition 7.3.2.3, (7.3.2.12--17), pp.~354--355;
the Hilbert orthogonal projections are Definitions
11.2.1.1/11.2.1.3, pp.~487--489, with the explicit
metric here retained. The physical cutoff is the
first Step~II bullet on p.~516 and (12.4.1.8--10),
p.~517. Its action is computed on the actual source.
The even gamma factor agrees with Tate (1950),
Theorem~2.4.1 and the real table, thesis p.~(2.14),
PDF p.~21; its needed normalization and the precise
factorial consequence are derived above. These
are source interfaces, not a fresh audit of Tate's
whole global proof.

This proves failure of the canonical theta-map
intertwining $Q_Y\Theta=\Theta\mathfrak b_Y$ in the
main rational trivial sector, on the corrected
symmetric metric and separately on each displayed
one-sided metric. It does not assert that the draft
prints that exact equation, or exclude every different
nonlocal or Gram-corrected comparison map. The actual
Gram factors of MA-GP remain valid. The scalar trace
identity supplied by $A$, the independent FT002
deficiency gate, and the distinction between a
restricted ideal return and the full quotient trace
are unchanged.

\paragraph{Grothendieck placement.}
The arithmetic sector is the trivial rank-one Artin
representation, the constant finite-monodromy local
system on $\operatorname{Spec}\mathbb Z$, with Euler
factor $(1-p^{-z})^{-1}$. Its placement is SGA~1,
Expos\'e~V, \S7, and Deligne, \emph{Valeurs de fonctions L},
Definition~6.1 and \S\S6.2--6.3, in the inherited
conventions of MA-CH and MA-CQ. The two source moments
are analytic polar conditions. The metric obstruction
does not change the Artin representation, a $K_0$
class, or a motivic weight. In cote~1, scans~079--080
and~111--112 through BR-TR-01/02, the trace/determinant
comparison still requires defined continuous/nuclear
maps. This proof refutes one exact proposed map while
leaving the previously constructed actual comparison
operators at their proved scopes.

\section{The non-induced imaginary-quadratic physical-projection obstruction}
\label{sec:imaginary-projection-obstruction}

\paragraph{Status and scope.}
The complete author reading and two independent whole-proof
readings have passed; their actual check scopes are recorded separately.  It uses the actual non-induced cubic character
already fixed in MA-CH-003.  The different prime is retained in $S$,
and the final input belongs to the original coordinate-flat source.
Only direct membership in the full source PW range is needed.
The modified finite test vector below is not identified with the
primitive canonical vector or with a new Artin character.

Keep the established arithmetic data
\[
 \begin{gathered}
 K=\mathbb Q(\zeta_3),\quad E=K(269^{1/3}),\quad
 \psi=\psi_{269,3},\quad
 \mathfrak p=(\sqrt{-3}),\quad \mathfrak q=(269),\\
 N\mathfrak p=3,\quad q=N\mathfrak q=269^2,\quad
 \mathfrak f_\psi=\mathfrak q,\quad Q=3q,\quad
 S=\{\infty,\mathfrak p,\mathfrak q\},\qquad S'=S.
 \end{gathered}
 \tag{MA-PI.120}
\]
Here $Q$ is the primitive analytic conductor.
The class number is one, $|\mathcal O_K^\times|=6$,
and $\psi_{\mathfrak p}$ is trivial: the prime over $3$
splits in $E/K$.  The local different exponent at
$\mathfrak p$ is one.  At $\mathfrak q$ the character
has conductor exponent one and no inertia invariants.
These facts, including $\psi^c=\psi^{-1}\ne\psi$ and
non-descent from a Dirichlet norm character, are the
already proved MA-CH-003/MA-PF-003/005 facts.

Use the coherent additive characters of MA-CGP-001.
At infinity the Fourier kernel is
$e^{-2\pi i\operatorname{Tr}_{\mathbb C/\mathbb R}(z\xi)}$,
with self-dual measure $2\,d(\operatorname{Re}z)\,d(\operatorname{Im}z)$.
The normalized complex absolute value is $|z|_{\mathbb C}=|z|^2$.
Put $\pi=\sqrt{-3}$, and take the finite functions
\[
 \begin{split}
 h_{\mathfrak p}
   &=\mathbf1_{\mathcal O_{\mathfrak p}^{\times}}
            -3\mathbf1_{\pi\mathcal O_{\mathfrak p}^{\times}}
     =\mathbf1_{\mathcal O_{\mathfrak p}}
            -4\mathbf1_{\pi\mathcal O_{\mathfrak p}}
            +3\mathbf1_{\pi^2\mathcal O_{\mathfrak p}},\\
 \kappa_{\mathfrak q}
   &=\psi_{\mathfrak q}^{-1}|_{\mathcal O_{\mathfrak q}^{\times}}\,
                             \mathbf1_{\mathcal O_{\mathfrak q}^{\times}},\\
 \varphi_f&=f\otimes h_{\mathfrak p}\otimes\kappa_{\mathfrak q}
                    \quad\text{on }\mathbb A_{K,S}.
 \end{split}
 \tag{MA-PI.121}
\]
The global extension has $\mathbf1_{\mathcal O_v}$ outside $S$.
Let $\mathcal T_S$ denote its actual centered, fixed-$\psi$
global theta map, and put
$Y_f=\mathcal T_S(\varphi_f)$ for radial $f\in\mathcal S(\mathbb C)$.
The finite ramified factor has value and integral zero,
so all these global data satisfy both polar conditions,
without conditions on $f(0)$ or its integral.

For $Y>0$ let $\rho_Y$ be a real smooth profile with
\[
 0\le\rho_Y\le1,\qquad
 \rho_Y(r)=1\ (r\le Y/2),\qquad
 \rho_Y(r)=0\ (r\ge2Y),\qquad
 \mathfrak B_{S,Y}
      =\mathcal F_S^{-1}M_{\rho_Y(|\xi|_S)}\mathcal F_S.
 \tag{MA-PI.122}
\]
For either of the independently specified metrics
\[
 \begin{aligned}
 H^{\rm sym}_{a,s}
    &=L^2(\mathbb R,e^{2a|x|}(1+x^2)^{s/2}dx),
                 &&0<a<1/2,\ s\ge0,\\
 H^{\rm one}_{b,s}
    &=L^2(\mathbb R,e^{2bx}(1+x^2)^{s/2}dx),
                 &&|b|<1/2,\ s\ge0,
 \end{aligned}
 \tag{MA-PI.123}
\]
write $Q^H_{S,Y}$ for the orthogonal projection onto the
closure of the full actual invariant global PW image,
with the source product Fourier bound $Y$ and $S'=S$.
This is the corrected centered realization in the first
metric.  Its exact printed realization in the second
metric is identified below.  The projection is onto the
full original source level, not only the image of the
particular family (MA-PI.121).

\paragraph{MA-PI-004 (a non-induced character witness in the original domains).}
For either fixed metric in (MA-PI.123), there is $Y_0>0$
such that, for every $Y\ge Y_0$ and every profile
(MA-PI.122), there is a radial Schwartz function $f_Y$
identically zero near the complex origin for which
\[
 \varphi_{f_Y}\in V_S,\qquad
 Q^H_{S,Y}Y_{f_Y}
       \ne\mathcal T_S(\mathfrak B_{S,Y}\varphi_{f_Y}).
 \tag{MA-PI.124}
\]
Here $V_S$ is the original inherited coordinate-flat LF
source of MA-GA.2.  The cutoff output is actual
$S$-local Schwartz data satisfying both polar conditions.
Thus $Y_{f_Y}$ and
$\mathcal T_S(\mathfrak B_{S,Y}\varphi_{f_Y})$
belong to $\mathcal E$.
The set $S$ includes the different and the conductor and
satisfies the arithmetic hypotheses of the original
source setup, including Theorem~12.4.1.1, p.~514.
No Hecke zero-location premise is used.

\begin{proof}
\emph{1. The exact finite Fourier shells and unchanged physical cutoff.}
The self-dual volume at $\mathfrak p$ is
$v=\operatorname{vol}(\mathcal O_{\mathfrak p})=3^{-1/2}$.
Annihilators of balls, with different exponent one, give
\[
 \widehat h_{\mathfrak p}(\xi)
   =v\left(\mathbf1_{\{|\xi|_{\mathfrak p}\le3\}}
       -\frac43\mathbf1_{\{|\xi|_{\mathfrak p}\le9\}}
       +\frac13\mathbf1_{\{|\xi|_{\mathfrak p}\le27\}}\right)
   =3^{-3/2}h_{\mathfrak p}(\pi^3\xi).
 \tag{MA-PI.125}
\]
Indeed $\widehat{\mathbf1_{\pi^j\mathcal O}}=
3^{-j}v\,\mathbf1_{\pi^{-j-1}\mathcal O}$, for $j=0,1,2$.
The values in (MA-PI.125) are $-v$ on the shell of norm
$9$, $v/3$ on the shell of norm $27$, and zero elsewhere.
The original $h_{\mathfrak p}$ is supported away from zero
and has integral $v-4v/3+3v/9=0$.

The primitive finite Gauss calculation of MA-CQ-001,
with conductor exponent one and different exponent zero,
gives a single Fourier shell
$|\xi_{\mathfrak q}|_{\mathfrak q}=q$ for
$\kappa_{\mathfrak q}$, with a nonzero Gauss factor of
modulus $q^{-1/2}$.  Its exact phase cancels when the
actual inverse Fourier transform is applied.
The finite product Fourier norms of (MA-PI.121) are
therefore exactly $9q$ and $27q$.

For later use split (MA-PI.125) into its two shells.
Their actual inverse Fourier transforms are
\[
 h_9=\mathbf1_{\mathcal O_{\mathfrak p}}
                  -3\mathbf1_{\pi\mathcal O_{\mathfrak p}},
 \qquad
 h_{27}=3\mathbf1_{\pi^2\mathcal O_{\mathfrak p}}
                  -\mathbf1_{\pi\mathcal O_{\mathfrak p}},
 \qquad h_9+h_{27}=h_{\mathfrak p}.
 \tag{MA-PI.126}
\]
For example
$\mathcal F_{\mathfrak p}^{-1}\mathbf1_{\{|\xi|\le3^k\}}
 =3^kv\,\mathbf1_{\pi^{k-1}\mathcal O_{\mathfrak p}}$.
Multiplying the shell differences by $-v$ and $v/3$
proves (MA-PI.126), including every constant.
Consequently the physical cutoff has the literal decomposition
\[
 \begin{split}
 \mathfrak B_{S,Y}\varphi_f
   &=\sum_{r\in\{9,27\}}
          f_{Y,r}\otimes h_r\otimes\kappa_{\mathfrak q},\\
 f_{Y,r}
   &=\mathcal F_{\mathbb C}^{-1}
           \bigl(\rho_Y(rq|z|_{\mathbb C})\widehat f(z)\bigr).
 \end{split}
 \tag{MA-PI.127}
\]
This computation retains the two finite Fourier bands;
it does not replace them by a single cutoff of a different
level.  Every summand is Schwartz and has both global
polar moments zero because of $\kappa_{\mathfrak q}$.
Conjugating coherent Fourier signs only reflects norm
coordinates, leaving these norm cutoffs unchanged.

\emph{2. Actual theta, its finite filter, and its effective reflection scale.}
Let $Y_f^{\rm prim}$ denote the primitive-vector theta
of MA-CGP-001.  With $f$ radial its exact formula is
\[
 Y_f^{\rm prim}(x)
    =6e^{x/2}\sum_{\mathfrak a\ne0}
             \psi(\mathfrak a)
                       f\bigl((N\mathfrak a\,e^x)^{1/2}\bigr).
 \tag{MA-PI.128}
\]
The ideal coefficient is extended by zero at $\mathfrak q$.
Class number one and the six units give the factor $6$.
The equality $\psi_{\mathfrak p}=1$ makes the actual
finite change in (MA-PI.121) exactly
\[
 Y_f(x)=Y_f^{\rm prim}(x)
          -\frac4{\sqrt3}Y_f^{\rm prim}(x+\log3)
          +Y_f^{\rm prim}(x+2\log3).
 \tag{MA-PI.129}
\]
There is no unrecorded uniformizer phase.  The principal-idele
relation on $\pi$ forces $\psi_{\mathfrak q}(\pi)=1$:
the archimedean component is trivial, the component at
$\mathfrak p$ is trivial, and all other finite components
are unramified on the unit $\pi$.  Any complex rotation
from multiplication by $\pi$ is invisible to radial $f$.
The factor $3^{-j/2}$ from the centered shift accounts for
the two coefficients in (MA-PI.129).

The actual Poisson formula for the primitive finite vector,
or (MA-PI.125) together with its ramified Gauss factor,
therefore yields
\[
 Y_f(x)=\epsilon_\psi\,Y^{\,\bar\psi}_{\widehat f}
                     (-x-\log Q_*),\qquad
 Q_*=27q=9Q,\qquad |\epsilon_\psi|=1.
 \tag{MA-PI.130}
\]
Here the superscript means that the same $h_{\mathfrak p}$
is used and the ramified character is conjugated.
For a direct check using (MA-PI.129), the primitive
reflection is at $-\log Q$.  Replacing $x$ by
$-x-\log(9Q)$ in the three dual shifted terms reverses
their order; their coefficients $1,-4/\sqrt3,1$ are
unchanged.  This proves (MA-PI.130), with the same phase.
Thus $Q_*$ is an effective scale of this test vector,
not a replacement for the primitive conductor $Q=3q$.

The direct ideal sum at $x\to+\infty$ and the reflected
sum at $x\to-\infty$ prove
$Y_f\in\mathcal E$, continuously for radial Schwartz $f$.
For completeness, a Schwartz bound of order $2M$ on the
archimedean variable bounds the direct sum by a constant
times $e^{-(M-1/2)x}\sum_{\mathfrak a}(N\mathfrak a)^{-M}$,
for $M>1$.  Every norm derivative satisfies the same kind
of estimate.  The reflected identity and continuity of
Fourier transform on Schwartz space give the other end.
This proves continuity by finitely many source seminorms
for every prescribed target seminorm.  It applies uniformly
to bounded Schwartz families.

Write $F(t)=f(\sqrt t)$ for the radial norm profile, and
$M_F(s)=\int_0^\infty F(t)t^{s-1}dt$.
Actual unfolding in $\Re s>1$ gives
\[
 \begin{split}
 \mathcal M Y_f(s-\tfrac12)&=6A_\psi(s)M_F(s),\\
 A_\psi(s)&=E_3(s)L_K(s,\psi),\\
 E_3(s)&=(1-3^{-s})(1-3^{1-s})
           =1-4\,3^{-s}+3\,3^{-2s}.
 \end{split}
 \tag{MA-PI.131}
\]
The factor replaces the primitive local geometric series
at $\mathfrak p$ by the weights $1,-3,0,\ldots$.
One has exactly
$E_3(s)=3^{1-2s}E_3(1-s)$.

The Gaussian $f_0(z)=e^{-2\pi|z|^2}$ is self-dual for the
declared complex Fourier measure, and
$M_{F_0}(s)=(2\pi)^{-s}\Gamma(s)=\Gamma_{\mathbb C}(s)/2$,
where $\Gamma_{\mathbb C}(s)=2(2\pi)^{-s}\Gamma(s)$.
Taking its Mellin transform in (MA-PI.130) proves
\[
 \Gamma_{\mathbb C}(s)A_\psi(s)
   =\epsilon_\psi Q_*^{1/2-s}
                \Gamma_{\mathbb C}(1-s)A_{\bar\psi}(1-s).
 \tag{MA-PI.132}
\]
Both completed expressions are entire, by the
all-exponential theta bounds.  Multiplication by the
entire reciprocal gamma function also proves that
$A_\psi$ and $A_{\bar\psi}$ are entire.  The effective
scale in (MA-PI.132) is derived from the actual source,
not assigned as an arithmetic conductor.
For arbitrary radial Schwartz $f$, $M_F(s)$ is holomorphic
in $\Re s>0$.  Analytic continuation therefore extends
(MA-PI.131) to that half-plane, which includes every
line used in the Gram calculation.

\emph{3. The exact radial Gram form and its temperateness.}
Fix $b\in(-1/2,1/2)$ and put
$\sigma=1/2+b$, $d=1-\sigma$.
Given $A_u\in C_c^\infty(\mathbb R)$, define a radial
frequency vector by
\[
 u(z)=U(|z|_{\mathbb C}),\qquad
 U(t)=t^{\sigma-1}A_u(\log t),\qquad
 f=\mathcal F_{\mathbb C}^{-1}u.
 \tag{MA-PI.133}
\]
It is smooth with compact support away from the complex
origin; $f$ is radial Schwartz.
The radial local Fourier equation is
\[
 M_F(s)=\frac{\Gamma_{\mathbb C}(s)}
                  {\Gamma_{\mathbb C}(1-s)}M_U(1-s),
                 \qquad0<\Re s<1.
 \tag{MA-PI.134}
\]
This is Tate's local equation, Theorem~2.4.1, with the
angular index zero in the complex table, thesis p.~(2.16),
PDF p.~23.  The table gives
$(2\pi)^{1-2s}\Gamma(s)/\Gamma(1-s)$.
Its angular multiplicative-measure factor $2\pi$ occurs
on both sides and cancels for our norm-profile Mellin
integrals.  The self-dual Gaussian above fixes the phase
as one.  On radial functions the inverse and forward
Fourier transforms agree, since their square is reflection.
These observations verify (MA-PI.134) in the declared
normalization rather than importing an unspecified gamma factor.

For $\widehat A_u(t)=\int A_u(r)e^{-itr}dr$ one has
$M_U(1-\sigma-it)=\widehat A_u(t)$.
Equations (MA-PI.131)--(MA-PI.134) and Mellin Plancherel
therefore give, for $f,h$ associated to $u,v$,
\[
 \begin{split}
 \int_{\mathbb R}Y_f(x)\overline{Y_h(x)}e^{2bx}dx
   &=\frac1{2\pi}\int_{\mathbb R}
         W_b(t)\widehat A_u(t)\overline{\widehat A_v(t)}dt,\\
 W_b(t)&=36Q_*^{-2b}
             A_\psi(d+it)A_{\bar\psi}(d-it).
 \end{split}
 \tag{MA-PI.135}
\]
For real $t$ this symbol is real and nonnegative, because
$A_{\bar\psi}(\bar s)=\overline{A_\psi(s)}$.
That identity follows first from the absolutely convergent
ideal series and the real-coefficient $E_3$, and then by
analytic continuation.

Here is an actual-source proof that $W_b(t)\,dt$ is tempered;
no additional vertical-growth theorem is needed.
Fix a nonnegative nonzero $A_0\in C_c^\infty(\mathbb R)$.
Its Fourier transform has modulus at least $c>0$ on
$[-\delta,\delta]$ for some $\delta>0$.
For each $T\in\mathbb R$ use
$A_T(r)=e^{iTr}A_0(r)$ in (MA-PI.133).
The resulting $u_T$ have support in one fixed compact
annulus.  The chain rule shows that each fixed Schwartz
seminorm of $u_T$ is at most $C(1+|T|)^N$: powers and
derivatives of $|z|_{\mathbb C}$ and its logarithm are
bounded on that annulus, and each derivative of the
modulation contributes at most a power of $T$.
Fourier continuity gives the same assertion for any
fixed seminorm of $f_T$.  The actual theta continuity
proved after (MA-PI.130), followed by the continuous
inclusion $\mathcal E\to L^2(e^{2bx}dx)$, consequently gives
$\|Y_{f_T}\|^2\le C'(1+|T|)^{N'}$ for fixed $C',N'$.
Applying the positive identity (MA-PI.135) and
$\widehat A_T(t)=\widehat A_0(t-T)$ yields
\[
       \int_{T-\delta}^{T+\delta}W_b(t)dt
                         \le C''(1+|T|)^{N'}.
 \tag{MA-PI.136}
\]
A covering by intervals of this fixed length now gives
$\int W_b(t)(1+|t|)^{-M}dt<\infty$ for sufficiently
large $M$.  Thus it is a positive tempered measure and
its inverse Fourier transform $\kappa_b$ is defined
as a tempered distribution.

\emph{4. Noncompact support, with Hermitian rather than assumed even symmetry.}
The symbol has the entire continuation
\[
       \mathscr W_b(z)
         =36Q_*^{-2b}A_\psi(d+iz)A_{\bar\psi}(d-iz).
 \tag{MA-PI.137}
\]
It cannot satisfy an exponential-type bound.  At
$s_n=-n-1/2$, (MA-PI.132) and the gamma reflection
identity give
\[
 |A_\psi(s_n)|
   =Q_*^{n+1}(2\pi)^{-2n-2}
       \frac{\Gamma(n+3/2)^2}{\pi}
                            |A_{\bar\psi}(n+3/2)|.
 \tag{MA-PI.138}
\]
Indeed
$\Gamma(-n-1/2)\Gamma(n+3/2)=(-1)^{n+1}\pi$.
Absolute convergence of the ideal series at $2$ and
dominated convergence show that $L_K(t,\bar\psi)\to1$
as the real variable $t\to+\infty$.  Also $E_3(t)\to1$.
Both required positive-axis $A$ factors therefore have
modulus at least $1/2$ for all sufficiently large $n$.
With $T_n=n+1/2+d$, (MA-PI.137)--(MA-PI.138) imply
\[
 |\mathscr W_b(iT_n)|
   \ge c_b\,Q_*^{n+1}(2\pi)^{-2n-2}\Gamma(n+3/2)^2.
 \tag{MA-PI.139}
\]
Since $\Gamma(n+3/2)\ge(\sqrt\pi/2)n!$, this grows
faster than every polynomial times $e^{R T_n}$.
If $\kappa_b$ had compact support, the finite-order
distribution estimate applied to a cutoff times
$e^{-izx}$ would give precisely such a bound for its
entire Fourier transform, as proved in (MA-PI.18).
Agreement on the real axis identifies that transform
with (MA-PI.137), a contradiction.

Although no evenness of $W_b$ is assumed, its reality
on the real axis gives the Hermitian distribution identity
$\kappa_b(-x)=\overline{\kappa_b(x)}$.
In particular its support is symmetric.  The separated
test argument from MA-PI-001 therefore applies unchanged:
if every compact smooth high/low log-frequency pair
beyond a fixed gap had zero Gram pairing, approximate
a delta in the high variable.  This would make
$\kappa_b$ vanish on the corresponding negative
half-line; Hermitian symmetry would give vanishing
on the positive half-line too.  That would make its
support compact.  A nonzero limiting distribution
pairing has a nonzero ordinary smooth approximate-delta
member, so actual compact test functions suffice.

Choose the gap to obtain radial Schwartz $f,h$ with
compact smooth Fourier supports and
\[
 \begin{gathered}
 \operatorname{supp}\widehat f
       \subset\{|z|_{\mathbb C}>2/(9q)\},\qquad
 \operatorname{supp}\widehat h
       \subset\{0<|z|_{\mathbb C}<1/(54q)\},\\
 D_b:=\int Y_f(x)\overline{Y_h(x)}e^{2bx}dx\ne0.
 \end{gathered}
 \tag{MA-PI.140}
\]
No assertion is made for every preassigned pair of
bounded annuli.

\emph{5. Direct PW membership and actual dilation in both metrics.}
For $L>0$ put $Y=e^L$ and use the complex unitary
frequency dilation
\[
 u_L(z)=e^{-L/2}u(e^{-L/2}z),\qquad
 f_L(z)=e^{L/2}f(e^{L/2}z),
 \tag{MA-PI.141}
\]
and similarly for $v,h$.  Complex norm frequencies
are multiplied by $e^L$.  Exact theta summation gives
$Y_{f_L}(x)=Y_f(x+L)$ and the same formula for $h$.
The smallest finite Fourier norm is $9q$ and the
largest is $27q$.  Equations (MA-PI.122) and
(MA-PI.140) therefore imply
\[
 \mathfrak B_{S,Y}\varphi_{f_L}=0,\qquad
 \mathfrak B_{S,Y}\varphi_{h_L}=\varphi_{h_L},\qquad
 Y_{h_L}\in\operatorname{Ran}Q^H_{S,Y}.
 \tag{MA-PI.142}
\]
For the last inclusion the entire $S$-Fourier support
is compact with product norm $<Y/2<Y$.
Both polar conditions hold at $\mathfrak q$.
This is direct membership in Masullo's original
Definition~7.3.2.3 with $S'=S$, followed by its
actual global image and Hilbert closure.  It does not
require the modified family to exhaust that range.

For the symmetric metric choose $b=-a$ when constructing
(MA-PI.140); for the one-sided metric use its own $b$.
The exact two changes of variable give respectively
\[
 \begin{split}
 e^{-2aL}L^{-s}
     \langle Y_{f_L},Y_{h_L}\rangle_{H^{\rm sym}_{a,s}}
                    &\longrightarrow D_{-a}\ne0,\\
 e^{2bL}L^{-s}
     \langle Y_{f_L},Y_{h_L}\rangle_{H^{\rm one}_{b,s}}
                    &\longrightarrow D_b\ne0.
 \end{split}
 \tag{MA-PI.143}
\]
The polynomial ratio is
$(1+(x-L)^2)^{s/2}/L^s$, bounded for $L\ge1$ by
$C_s(2+|x|)^s$ and tending to one.
In the first line the normalized exponential factor
tends to $e^{-2ax}$ and is bounded by $e^{2a|x|}$;
in the second it is exactly $e^{2bx}$.
The actual $\mathcal E$ bounds prove domination.
Thus the pairings are nonzero for all sufficiently
large $Y$, uniformly over the permitted profiles.
Orthogonality and (MA-PI.142) give a nonzero pairing
of $Q^H_{S,Y}Y_{f_L}
-\mathcal T_S(\mathfrak B_{S,Y}\varphi_{f_L})$
against the actual range vector $Y_{h_L}$.

\emph{6. A finite flat witness, including the two-shell cutoff output.}
Fix this finite $Y$ and profile.  Choose a radial
$\omega\in C_c^\infty(\mathbb C)$ equal to one near zero,
and put
\[
 f_{L,\varepsilon}(z)
       =(1-\omega(z/\varepsilon))f_L(z),\qquad
 r_\varepsilon(z)=\omega(z/\varepsilon)f_L(z)
                        =g_\varepsilon(z/\varepsilon).
 \tag{MA-PI.144}
\]
Here $g_\varepsilon(w)=\omega(w)f_L(\varepsilon w)$
is a bounded family of radial compact smooth functions.
No value or derivative vanishing of $f_L$ is needed:
the finite ramified factor already removes both polar
terms.  The original input
$\varphi_{f_{L,\varepsilon}}$ is coordinate-flat at
every place of $S$.

The actual theta series gives exactly
$Y_{r_\varepsilon}(x)=
\varepsilon Y_{g_\varepsilon}(x-2\log\varepsilon)$.
Uniform $\mathcal E$ bounds then yield
\[
 \begin{split}
 \|Y_{r_\varepsilon}\|_{H^{\rm sym}_{a,s}}
    &\le C\varepsilon^{1-2a}
                    (1+|\log\varepsilon|)^{s/2}\longrightarrow0,\\
 \|Y_{r_\varepsilon}\|_{H^{\rm one}_{b,s}}
    &\le C\varepsilon^{1+2b}
                    (1+|\log\varepsilon|)^{s/2}\longrightarrow0.
 \end{split}
 \tag{MA-PI.145}
\]
For the second line substitute $x=y+2\log\varepsilon$
in the squared norm; it gives the exact factor
$\varepsilon^{2+4b}$.  For the first line use
$e^{2a|y+2\log\varepsilon|}
 \le\varepsilon^{-4a}e^{2a|y|}$.
All asserted exponents are positive.

For every real-coordinate multi-index $\alpha$,
differentiation of the complex Fourier integral gives
\[
 \sup_\xi|\partial_\xi^\alpha\widehat r_\varepsilon(\xi)|
                         \le C_\alpha\varepsilon^{|\alpha|+2}.
 \tag{MA-PI.146}
\]
The power $2$ is the real dimension of $\mathbb C$;
the self-dual factor $2$ and Fourier constants only
change $C_\alpha$.  For each of the two fixed shells,
multiplication by $\rho_Y(rq|z|_{\mathbb C})$ gives
compact smooth Fourier functions tending to zero
in every smooth seminorm.  Their inverse transforms
therefore tend to zero in radial Schwartz space.

Apply the exact decomposition (MA-PI.127) to this error.
For either fixed finite factor $h_9$ or $h_{27}$ its
actual theta map is continuous from radial Schwartz
space to $\mathcal E$.  One can see this directly
from (MA-PI.126): its image is respectively
$Y_f^{\rm prim}-\sqrt3Y_f^{\rm prim}(\,\cdot+\log3)$
or
$Y_f^{\rm prim}(\,\cdot+2\log3)
 -3^{-1/2}Y_f^{\rm prim}(\,\cdot+\log3)$.
The ramified factor remains unchanged and removes
both polar moments in each summand.  Consequently
\[
 \|\mathcal T_S(\mathfrak B_{S,Y}\varphi_{r_\varepsilon})\|_H
                                      \longrightarrow0.
 \tag{MA-PI.147}
\]
The profile and its derivative constants were fixed
before choosing $\varepsilon$; no uniform family
derivative bound has been assumed.

Use boundedness of the actual orthogonal projector and
test the difference for $f_{L,\varepsilon}$ against
$Y_{h_L}$.  By (MA-PI.145)--(MA-PI.147), the pairing
converges to the nonzero pairing obtained after
(MA-PI.143).  A sufficiently small finite positive
$\varepsilon$ supplies $f_Y=f_{L,\varepsilon}$ and
proves (MA-PI.124).  All asserted inputs and outputs
are actual source vectors; no distributional or
nonflat limiting vector replaces this final witness.
\end{proof}

\paragraph{The exact printed Hilbert realization.}
Let $m=\mu_1(C_K^1)>0$ be the actual compact Haar mass.
With the positive complex norm splitting
$g=\nu(e^x)k$, the source definition gives
\[
 \begin{split}
 (T_A\varphi)(\nu(e^x)k)
     &:=|g|^A\Phi_K(\varphi\otimes\mathbf1_{\mathcal O^{S}})(g)
       =\psi(k)^{-1}e^{bx}\mathcal T_S(\varphi)(x),\\
 (U_bF)(\nu(e^x)k)
     &=m^{-1/2}\psi(k)^{-1}e^{bx}F(x),
                    \qquad b=A-\tfrac12 .
 \end{split}
 \tag{MA-PI.148}
\]
These formulas apply to the source vectors and cutoff
outputs used above, which have the same compact type.
The printed polynomial norm makes $U_b$ a unitary
from $H^{\rm one}_{b,s}$ onto that compact-character
component.  Its inverse is compact-fiber integration
against $\psi$, multiplied by $m^{-1/2}e^{-bx}$,
as in (MA-PI.35).  Thus the actual printed orthogonal
projector is conjugate to the one-sided projector
used above, by equality of the full source ranges
under this actual isometry.  Only direct membership
of the test vector in that range was used in the proof.
Multiplying (MA-PI.124) by this unitary gives failure
of the printed canonical intertwining for every
$0<A<1$, in particular the original twist $A=1/2$.
The flat-source error exponent in that realization
is $1+2b=2A$, exactly as in (MA-PI.145).

\paragraph{Sources and Grothendieck placement.}
The original flat source is MA-GA.2 with its inherited
LF topology.  The invariant PW domain is Masullo
Definition~7.3.2.3, (7.3.2.12--17), pp.~354--355,
with $S'=S$ and no extra unit-shell transition set.
Its printed global orthogonal projector is Definition
11.2.1.3, (11.2.1.7--8), pp.~488--489.
The uncentered theta, printed polynomial norm and
actual twist are the source definitions checked in
(MA-PI.32)--(MA-PI.36); their general-number-field
form is unchanged here.  The physical cutoff is
the first Step~II bullet on p.~516 and composition
on p.~517 in Theorem~12.4.1.1.  Our $S$ contains
the field-ramified place $3$, as its p.~514
hypothesis requires.  The complex local gamma
normalization was freshly checked in Tate's actual
Theorem~2.4.1 and complex table, thesis pp.~(2.11)
and~(2.16), PDF pp.~18 and~23.  The new filter,
reflection scale, Gram temperateness and all error
estimates are proved above.

The arithmetic object remains the rank-one finite
Artin character of $\operatorname{Gal}(E/K)=C_3$,
with the established nontrivial two-dimensional
$S_3$ induction to $\mathbb Q$.  At $\mathfrak q$
its inertia invariants vanish; at $\mathfrak p$
it is unramified and has primitive Euler factor
$(1-3^{-s})^{-1}$.  The factor $E_3$ is imposed by
the explicit local test vector and is not a new
arithmetic conductor or a change of motive.
The finite-monodromy convention is SGA~1,
Expos\'e~V, \S7, and Deligne, \emph{Valeurs de fonctions L},
Definition~6.1 and \S\S6.2--6.3, with the inherited
MA-CH/MA-CQ audit scope.  The cote~1 comparison
programme, BR-TR-01/02, scans~079--080 and~111--112,
still requires actual comparison maps before trace
cyclicity; those archive locators are inherited.
This result refutes the specified canonical physical
cutoff/orthogonal projection map in a genuinely
non-induced sector.  It does not assert that the source
literally prints that proposed intertwining equation.
It does not refute a different
Gram-corrected transport, alter the valid nuclear
returns, or remove the arithmetic FT002 condition.

\section{Documentary dependency audit of the proposed global-to-local calculation}
\label{sec:source-dependency-audit}

This section records a fresh complete reading of Masullo~(2026),
\S\S12.3--12.5, printed pp.~514--520.  It compares the exact
source assertions with results already proved in this report.
It introduces no new mathematical theorem or arithmetic conclusion.

\paragraph{The absent comparison.}
Section~12.3, p.~514, has its heading and the footnote
``Include section'', but no argument.  On p.~516, Step~III
explicitly requires an identification of the local truncated
kernel with the kernel denoted by $k^{\mathrm{PW}}$.
On p.~517, the proof computes the kernel of a local physical
Fourier-cutoff composition and assigns it that notation.
This computation does not identify that kernel with the one
defined using the weighted global orthogonal PW projection.

MA-PI-001--003 refute a precise canonical theta intertwining
between the two constructions on actual source vectors.
MA-PI-002 uses the literal printed Hilbert realization,
including the norm-half twist specified on p.~515;
MA-PI-003 also retains both polar conditions in the principal
sector.  MA-PI-004 supplies the non-induced cubic character
over $\mathbb Q(\zeta_3)$, retaining the different prime over
$3$, the conductor prime over $269$, both physical Fourier
shells and genuine flat source data.
The draft does not print the exact canonical equation
tested by these results.  The documentary conclusion is therefore
that the required identification is absent and this particular
candidate fails.  No claim about every different factorization,
or every scalar trace equality, follows.

\paragraph{Hypotheses and cutoffs on p.~514.}
Theorem~12.4.1.1 retains class-group triviality after localization,
all archimedean places, all places ramified over $\mathbb Q$,
and the indicated outside-$S$ unit invariance of the test.
Field ramification cannot be replaced merely by character
ramification.  Footnote~2 flags unfinished work on a one-sided
support convention, an independent lower cutoff, and the resulting
limit.  The actual repaired local quotient operators are supplied
in MA-LA-014--015, MA-CQ-003--004 and MA-GA-001 at their stated
scopes.  The frequency operator in Step~II, p.~516, is explicitly
a product-norm ball.  The annular-frequency alternative of
MA-LA-016 is a different operator.

\paragraph{The source and trace factorization on p.~515.}
The proof fixes the norm-half twisted Sobolev realization,
asserts nuclearity on a quotient PW level, and invokes a
factorization through a local theta image before cyclic trace
transfer.  Its topological-copy assertion cites
Proposition~5.3.4.4 and is accompanied by a footnote asking for
that proposition to be updated.  The cited transform/image
criterion does not itself specify the continuous inverse,
nuclear lift, or commuting operator diagram needed here.
Range inclusion alone does not justify cyclic transfer.
The actual local sections and returns in MA-LA-011, MA-CQ,
MA-GA and MA-MO-017/028 have their proved domains and metrics;
they do not by themselves identify the global orthogonal
projection with the physical Fourier cutoff.

\paragraph{Literal assertions decided by the earlier source audit.}
MA-LA-008 identifies the output norm-half factor present in the
p.~515 operator but absent from the raw $K$ in (12.4.1.7),
p.~516.  The exact repair is $A=UK$; a change of coordinates
also requires consistent conjugation of the other operators.
MA-LA-009 refutes the raw tensor containment following that
display: equivariance does not make the raw output constant on
input $S$-unit coinvariant classes.  An output quotient or a
specified section is required.
MA-LA-010 refutes the asserted additive-Schwartz domain/codomain
of mixed-place product-norm cutoffs, even after an annular
position repair.  They remain bounded additive $L^2$ multipliers;
that weaker fact does not establish the printed source mapping.

\paragraph{Local kernels and nuclearity on pp.~516--518.}
Equations~(12.4.1.8)--(12.4.1.10) compute the local physical
composition.  Its kernel and the ensuing local orbital-volume
calculation still require the actual quotient/source domains.
MA-LA-012 rules out interpreting a nonzero unit-equivariant
operator of this form as compact on the full additive $L^2$
cover.  This is an attribution boundary: the p.~515 nuclearity
claim concerns a quotient PW level, and p.~517 uses a fundamental
domain and unit sum.  We do not attribute an explicit full-cover
trace-class theorem to the draft.  Actual quotient trace-class
operators exist and have been constructed above.
The p.~518 logarithmic-volume calculation supplies neither
the missing global map nor the lower-cutoff/rate analysis
flagged in footnote~2.

\paragraph{Rates and the final reuse on pp.~519--520.}
Page~519 requires both smoothing-width products to vanish,
but does not impose the additional slow-band relation between
the large radii.  MA-LA-015 proves profile comparison for the
actual repaired local operators; its paired global comparison
retains the stated slow-band hypothesis.
Theorem~12.5.1.1 assumes vanishing on the entire norm-one
subgroup.  The p.~520 choice of $S$ retains field and character
ramification, structural/class-group conditions, test invariance,
and all relevant small primes.

The p.~520 argument first decomposes the global kernel into
its $S$-part and outside-$S$ parts, then removes the latter
using the cited shell theorem.  It next invokes
Theorem~12.4.1.1 to identify the global $S$-kernel trace with
the local distributions minus the two polar terms, and finally
uses the spectral construction to obtain the relative-strip
zero sum.  The third step is the unsupplied comparison.
The local calculation alone gives neither equality of these
kernels, compatibility of their Hilbert metrics, control of
the limiting remainder, nor the polar subtraction.
The missing \S12.3 is the intended location of such a bridge.

\paragraph{Scope of the resulting conclusion.}
The source contains literal repairable errors, assertions
refuted on their stated domains, a local calculation for which
actual quotient repairs are now proved, and a separate missing
global comparison.  These categories must not be conflated.
MA-FT-001/002 and the Gram/resolvent results retain their exact
operator limits and arithmetic deficiency condition.
The canonical-map counterexamples do not eliminate that
condition or rule out all other comparison maps.
This documentary pass uses the previously recorded whole-proof
reviews of those results; it does not claim a fresh complete
proof review of every cited module.  Its row-level evidence is
in \texttt{source\_dependency\_audit.tsv}.

\section{Actual slow-band transport in the two character families}
\label{sec:character-slow-band}

\paragraph{MA-SB-004 (the character slow-band limit with its exact remainder).}
Let $\eta$ be a nonprincipal primitive Dirichlet character, or
the non-$\mathbb Q$-induced cubic character $\psi_{269,3}$ of
MA-CH-003.  More generally the statement applies to the
nontrivial finite-order characters of $\mathbb Q(\zeta_3)$
covered by MA-CGP-001--004.  Use probability Haar measure
on the norm-one compact group, the coherent additive
character of MA-NR-001, and the primitive finite factors
with $c_v=\beta(p_v)^{-1}$ as in MA-CH.3.
Fix
\[
 \begin{gathered}
 0<\alpha<\tfrac12,\qquad s=0,\qquad
 g\in C_c^\infty((0,\infty)),\qquad h(t)=g(\log t),\\
 G_0(z)=\int g(x)e^{zx}\,dx,\qquad
 Q=|D_K|N\mathfrak f_\eta,\qquad Y=Q\Lambda\ge1.
\end{gathered}
 \tag{MA-SB.30}
\]
Thus $h(1)=0$.  Let $Q_{\eta,Y}^{\rm glob}$ be the actual
orthogonal PW projection at invariant source product level
$Y$, with $S'=S$ and $S$ containing the conductor and
different.  Equivalently it is the primitive archimedean
norm-band projection of MA-CGP-002 at $\Lambda=Y/Q$.
No Euler factor has been removed by this equality.

Let $\rho_{X,\delta}$ be a source cutoff supported on
$[1/X,X]$, equal to one on
$[1/X+\delta,X-\delta]$, and between zero and one.
There are two fixed functions
$\omega_{\eta,g},\nu_{\eta,g}:[0,\infty)\to[0,\infty)$,
both tending to zero at infinity, such that whenever
$X\ge2Y\ge2$ and $0<\delta\le1/(4X)$,
\[
 \begin{split}
 &\left|\operatorname{Tr}_H
      (M_{\rho_{X,\delta}}C_gQ_{\eta,Y}^{\rm glob})
                +\tau_{\eta,+}(g)\right|\\
 &\hspace{15mm}\le
    \omega_{\eta,g}(\log(X/2))
        +\nu_{\eta,g}(\log(X/2)),\\
 &\tau_{\eta,+}(g)
     =\sum_{\operatorname{Re}\rho_\eta<1/2+\alpha}
                 m_\eta(\rho_\eta)G_0(\rho_\eta-1/2).
\end{split}
 \tag{MA-SB.31}
\]
The trace is on the actual character-sector space
$H=L^2(\mathbb R,e^{2\alpha|x|}\,dx)$.
In particular its limit along every source-thin path
$X,Y\to\infty$, $Y/X\to0$, is $-\tau_{\eta,+}(g)$.

For a comparison with the actual local quotient, use the
common fixed profiles from MA-CQ-003:
\[
 \begin{aligned}
 a_X(t)&=A(t/X)-A(Xt),\\
 T^{\beta,S}_{X,Y,X^{-1}}(h)
  &=M_{a_X}A_hD_{\beta,S}Q_{\infty,Y}D_{\beta,S}^*,\\
 Q_{\infty,Y}
  &=\mathcal F_\infty^{-1}
       M_{B(|\xi|_\infty/Y)}\mathcal F_\infty.
\end{aligned}
 \tag{MA-SB.32}
\]
Here $A,B$ equal one on $[0,1]$ and zero on $[2,\infty)$,
and $M_{a_X}$ acts through the archimedean norm in the
local operator and through $t=e^x$ in the global operator.
The finite unitary is the actual
$D_{\beta,S}=M_{\bar\beta,S}$ of MA-CQ.16--.17;
the intermediate Fourier map exchanges the dual sectors.
Choose the fixed $S$ also to include every prime ideal
whose positive or inverse prime-power norm meets
$\operatorname{supp}h$.
Then the two actual limits exist and their difference is
\[
 \begin{split}
 &\lim_{\substack{X,Y\to\infty\\Y/X\to0}}
 \left\{\operatorname{Tr}_H(M_{a_X}C_gQ_{\eta,Y}^{\rm glob})
            -\operatorname{Tr}T^{\beta,S}_{X,Y,X^{-1}}(h)\right\}\\
 &\hspace{12mm}=
    \sum_{\operatorname{Re}\rho_\eta\ge1/2+\alpha}
                  m_\eta(\rho_\eta)G_0(\rho_\eta-1/2).
\end{split}
 \tag{MA-SB.33}
\]
There is no polar term in these nontrivial primitive
character sectors.  The upper boundary is included in
the remainder.  Neither its nonemptiness nor its value
on an arbitrary particular test is asserted.

\begin{proof}
\emph{The actual range and its normalization.}
By MA-CGP-002 the source range is the actual centered
character ideal supported on
$x\ge-\log(Q\Lambda)=-\log Y$.
Put $L=\log Y$ and $r=x+L$.
The unitary $f\mapsto e^{\alpha x}f(x)$, followed by
this translation, gives exactly the metric
\[
 R_L(r)=
 \begin{cases}e^{4\alpha(L-r)},&0\le r\le L,\\
 1,&r\ge L.
 \end{cases}
 \tag{MA-SB.34}
\]
In the equivalent ordinary half-line norm the range is
$V_\eta L^2(0,\infty)$, where
\[
 V_\eta=\mathcal L^{-1}M_{B_{\eta,\alpha}}\mathcal L,
 \qquad P_{K_\eta}=1-V_\eta V_\eta^*.
 \tag{MA-SB.35}
\]
The actual pure Blaschke divisor consists of
$z=\alpha+1/2-\rho_\eta$ with positive real part,
including the full multiplicities.  This is the
$\eta$ divisor: the $\bar\eta$ in the reflected
Fourier formula MA-CGP.13 has already been changed
by the completed functional equation.  Neither
complex conjugation of the final divisor nor a
further conductor dilation is allowed here.

Put $k(t)=e^{\alpha t}g(t)$.
MA-CGP-004, at $c=0$ where the metric is one, proves
the actual ambient lift
$C_k^+P_{K_\eta}\in\mathcal S_1$ and identifies its
trace as $\tau_{\eta,+}(g)$.
This verifies the operator hypothesis of MA-SB-001;
a summable formal spectrum would not suffice.
The proof there supplies the fixed trace-class
Hankel operator $H_{B_\eta k}$ and hence the tails
\[
 \begin{aligned}
 \omega_{\eta,g}(r)&=
       \|P_{[r,\infty)}C_k^+P_{K_\eta}\|_1,\\
 \nu_{\eta,g}(r)&=
       \|P_{[r,\infty)}H_{B_\eta k}\|_1.
\end{aligned}
 \tag{MA-SB.36}
\]
Both vanish at infinity with no separation of the
zeros from the half-plane boundary.

\emph{The actual weighted Gram trace.}
The precise weighted projector is
\[
 V_\eta\mathfrak G_{\eta,L}^{-1}V_\eta^*M_{R_L},
 \qquad
 \mathfrak G_{\eta,L}=V_\eta^*M_{R_L}V_\eta .
 \tag{MA-SB.37}
\]
Causality and $R_L-1$ supported on $[0,L]$ give
$D_{\eta,L}=\mathfrak G_{\eta,L}^{-1}-1
 =P_{[0,L]}D_{\eta,L}P_{[0,L]}$ and
$-1\le D_{\eta,L}\le0$.
Every hypothesis used in MA-SB.18--.22 is now
verified for this fixed $V_\eta$ and $k$.
In particular those identities retain the right
metric factor, establish trace class before
cyclicity, and insert the input support projection
only inside an already defined trace-class trace.
They prove the same error bound
$\omega_{\eta,g}(R)+\nu_{\eta,g}(R-L)$.

For the source cutoff, $X\ge2Y$ and the stated
$\delta$ remove its entire lower transition from
the support half-line.  Its half-line profile is
one up to $R=L+\log(X-\delta)$.
This proves (MA-SB.31), including uniformity in
$1\le Y\le X/2$.
For $a_X$, the lower summand vanishes on that
half-line and the upper profile is one through
$R=L+\log X$.  Thus the same actual global limit
holds for the profile in (MA-SB.32).

\emph{The actual local limit and its trace comparison.}
MA-CQ-003 applies with $\epsilon=X^{-1}$:
$XY\to\infty$ and $\epsilon Y=Y/X\to0$.
It gives, on its actual local quotient completion,
\[
 \operatorname{Tr}T^{\beta,S}_{X,Y,X^{-1}}(h)
       \longrightarrow A_\infty(h)+B_{\beta,S}(h).
 \tag{MA-SB.38}
\]
The parity or complex archimedean coefficient and
the primitive dual finite factors are precisely
those of MA-CQ.21.  The conductor/different
dilations in $D_{\beta,S}$ remain in this operator;
their scalar trace contribution is proportional
to $g(0)=0$, as proved in MA-CQ.23.
The choice of $S$ gives $B_{\beta,S}=B_\eta$.

Meyer's actual smooth quotient trace and its
character-sector comparison MA-CC-001 identify
$A_\infty+B_\eta$ with
$-\tau_\eta^{\rm full}(g)$.  The polar quotient
is zero because the character is nontrivial.
This uses the unconditional smooth comparison,
not its separate uniform-interior Hilbert version.
Subtracting the two already proved operator
limits gives $\tau_\eta^{\rm full}-\tau_{\eta,+}$;
their established divisor descriptions then give
(MA-SB.33), with its upper-edge convention.
\end{proof}

The comparison with the actual completed full-strip
trace of MA-WP is also exact:
\[
 \begin{split}
 &\lim_{\substack{X,Y\to\infty\\Y/X\to0}}
    \operatorname{Tr}(M_{a_X}C_gQ_{\eta,Y}^{\rm glob})
       -\left(-\tau_{\eta,\alpha}(g)\right)\\
 &\hspace{10mm}=-\sum_{\operatorname{Re}\rho_\eta\le1/2-\alpha}
                  m_\eta(\rho_\eta)G_0(\rho_\eta-1/2).
 \end{split}
 \tag{MA-SB.39}
\]
At $s=0$ no boundary jet is continuous, which explains
the included lower edge on the right.
Every finite global trace here has the concrete
character Gram-resolvent realization MA-CGP.18.
At fixed $X,Y$, take that proved regularizer limit
first, or choose a regularizer whose fixed-data
trace error tends to zero along the path.
No uniform regularizer rate, LF lift of its inverse,
or omission of either Gram factor is asserted.

For $\psi_{269,3}$ the scale in (MA-SB.30) is
$Q=3\cdot269^2$, and its actual local finite
unitary is the one in MA-CQ-004.  In particular the
prime over $3$ retains its Euler factor, whereas
the ramified prime over $269$ contributes its
conductor dilation.  This distinction is already
part of the proved operator, not a modification
of the primitive $L$-function.
The result uses neither $\mathsf A$ nor $M$.
Its scope remains $s=0$, the stated positive
logarithmic support, and the actual source level
$S'=S$.  The paired local comparison uses the
specified fixed profiles, without an unproved
uniform claim for arbitrary varying local smoothings.

\section{The all-weight slow-band limit in the character sectors}
\label{sec:weighted-character-slow-band}

\paragraph{MA-SB-006 (the positive-weight character hypothesis is discharged).}
Take either primitive character family of MA-GW-004,
including $\psi_{269,3}$, and fix
$0<\alpha<1/2$, $s\ge0$, and
$g\in C_c^\infty((0,\infty))$.
Use the same primitive and coherent additive conventions
as MA-SB-004.  Put
\[
 \begin{gathered}
 Q=|D_K|N\mathfrak f_\eta,\qquad Y=Q\Lambda,
 \qquad c=-\log Y,\\
 r_{\eta,\rho}^+=
 \begin{cases}
 m_\eta(\rho),&\operatorname{Re}\rho<1/2+\alpha,\\
 \min(m_\eta(\rho),b_s),&\operatorname{Re}\rho=1/2+\alpha,\\
 0,&\operatorname{Re}\rho>1/2+\alpha,
 \end{cases}\\
 \tau_{\eta,+,s}(g)=
              \sum_\rho r_{\eta,\rho}^+G_0(\rho-1/2),
 \qquad b_s=\#\{j\ge0:2j<s-1\}.
 \end{gathered}
 \tag{MA-SB.65}
\]
Here $j$ is an integer and $G_0(z)=\int g(x)e^{zx}dx$.
Let $Q_{\eta,Y}^{\rm glob}$ be the actual invariant
source PW projection, with $S'=S$.
For the source cutoff in MA-SB-005 the uniform estimate
is now unconditional in this character sector:
\[
 \begin{split}
 &\left|\operatorname{Tr}_H
    (M_{\rho_{X,\delta}}C_gQ_{\eta,Y}^{\rm glob})
                         +\tau_{\eta,+,s}(g)\right|\\
 &\hspace{20mm}\le
     \Omega_{\eta,g,s}(\log(X-\delta))
                       \le\Omega_{\eta,g,s}(\log(X/2)),\\
 &\Omega_{\eta,g,s}(R)=
       \|P_RC_gP_{\eta,0}^Z\|_1\longrightarrow0.
 \end{split}
 \tag{MA-SB.66}
\]
The conditions are $Y\ge1$, $X\ge2Y$,
$0<\delta\le1/(4X)$ and
$\log(X-\delta)>\max\operatorname{supp}g$.
The reference projection $P_{\eta,0}^Z$ is at source
level $Y=1$, hence at primitive norm band $\Lambda=1/Q$.
It is orthogonal in the actual physical space
$L^2(\mathbb R,e^{2\alpha|x|}(1+x^2)^{s/2}dx)$.

In particular every source-thin slow path gives
$-\tau_{\eta,+,s}(g)$ as its actual global trace limit.
For the common fixed annular profile and the actual
local character quotient of MA-CQ-003, with $S$
containing all test prime contributions, the paired
limit is
\[
 \begin{split}
 &\lim_{\substack{X,Y\to\infty\\Y/X\to0}}
 \left\{\operatorname{Tr}(M_{a_X}C_gQ_{\eta,Y}^{\rm glob})
             -\operatorname{Tr}T^{\beta,S}_{X,Y,X^{-1}}(h)\right\}\\
 &\quad=\sum_{\operatorname{Re}\rho>1/2+\alpha}
                      m_\eta(\rho)G_0(\rho-1/2)\\
 &\qquad+\sum_{\operatorname{Re}\rho=1/2+\alpha}
       \bigl(m_\eta(\rho)-\min(m_\eta(\rho),b_s)\bigr)
                                      G_0(\rho-1/2).
 \end{split}
 \tag{MA-SB.67}
\]
There is no polar term for these nontrivial characters.

\begin{proof}
MA-CGP-002 gives the actual supported source range
at $c=-\log(Q\Lambda)$ and its nesting in this same
physical norm.  MA-GW-004, including its exact
primitive finite-vector calculation, proves the
actual reference lift $C_gP_{\eta,0}^Z\in\mathcal S_1$.
It also proves the $c$-independent actual quotient
trace with exactly the lengths in (MA-SB.65).
Thus every hypothesis of the abstract reference-level
argument in MA-SB-005 is satisfied.
In particular
\[
 \begin{aligned}
 P_0P_{\eta,c}^Z&=P_{\eta,0}^ZP_0P_{\eta,c}^Z,\\
 P_RC_gP_{\eta,c}^Z
      &=P_RC_gP_{\eta,0}^ZP_0P_{\eta,c}^Z.
 \end{aligned}
 \tag{MA-SB.68}
\]
for $c\le0$ and $R>\max\operatorname{supp}g$.
All factors act in the same actual weighted Hilbert
space and the rightmost factor $P_0P_{\eta,c}^Z$
has norm at most one.
The fixed trace-class tail and the identical
lower-transition exclusion give (MA-SB.66).

For the annular profile, the global error is bounded
by $\Omega_{\eta,g,s}(\log X)$.
MA-CQ-003 gives the actual local limit on its
specified quotient completion, retaining the
finite conductor/different dilations and the
dual-sector Fourier order.  The local completion
is not given the global polynomial weight.
The unconditional smooth character trace comparison
MA-CC-001 gives that local limit as
$-\tau_\eta^{\rm full}(g)$, with no polar quotient.
Subtracting the actual limits therefore gives
$\tau_\eta^{\rm full}-\tau_{\eta,+,s}$, which is
exactly (MA-SB.67).  This verifies the boundary
deficiency as well as the strict upper-exterior part.
\end{proof}

Relative to the actual completed full-strip trace,
the difference is the negative lower-exterior
trace and its unretained lower-edge length:
\[
 \begin{split}
 &-\sum_{\operatorname{Re}\rho<1/2-\alpha}
                        m_\eta(\rho)G_0(\rho-1/2)\\
 &\quad-\sum_{\operatorname{Re}\rho=1/2-\alpha}
       \bigl(m_\eta(\rho)-\min(m_\eta(\rho),b_s)\bigr)
                                       G_0(\rho-1/2).
 \end{split}
 \tag{MA-SB.69}
\]
This is also the difference of the two iterated
actual limits in MA-GW.51--.53.
No exterior zero or deficient edge jet is asserted
to exist.  The positive-length deficiency divisor
is detected by MA-JP-004 if it is nonempty.
All finite global traces retain the actual two
Gram resolvents of MA-CGP.18, in the fixed-data
regularizer sense specified in MA-SB-005.
The local comparison retains its fixed-profile
scope.  No $\mathsf A$, $M$, generic strip edge,
or uniform zero-separation hypothesis is used.

\section{Causal tests separate a nonempty zero divisor}
\label{sec:causal-trace-separation}

\paragraph{MA-JP-004 (resolvent-power separation, including two polar moments).}
Let
\[
 D=\sum_{\lambda\in\mathcal Z}d_\lambda[\lambda],
 \qquad d_\lambda\in\mathbb Z_{>0},\qquad
 u\le\operatorname{Re}\lambda\le v,
 \tag{MA-JP.40}
\]
be a locally finite divisor in a fixed vertical strip.
Assume its count includes multiplicities and satisfies
\[
 \sum_{\substack{\lambda\in\mathcal Z\\
                   T\le\operatorname{Im}\lambda<T+1}}
     d_\lambda\le C\log(2+|T|)\qquad(T\in\mathbb R).
 \tag{MA-JP.41}
\]
For $g\in\mathcal D_+=C_c^\infty((0,\infty))$, put
\[
 E_a(g)=\int_0^\infty g(t)e^{at}\,dt,\qquad
 \tau_D(g)=\sum_{\lambda\in\mathcal Z}d_\lambda E_\lambda(g).
 \tag{MA-JP.42}
\]
This sum is absolutely convergent and defines a continuous
linear functional on $\mathcal D_+$.
It is identically zero if and only if $D$ is empty.

More precisely, fix two distinct prescribed complex
exponents $a_1,a_2$.  Then
\[
 \begin{split}
 \tau_D|_{\ker E_{a_1}\cap\ker E_{a_2}}=0
 &\quad\Longleftrightarrow\quad
 \tau_D\in\operatorname{span}\{E_{a_1},E_{a_2}\}\\
 &\quad\Longleftrightarrow\quad
 \mathcal Z\subset\{a_1,a_2\}.
 \end{split}
 \tag{MA-JP.43}
\]
If there is a divisor point outside those two exponents,
therefore, there exists $g\in\mathcal D_+$ with
\[
       E_{a_1}(g)=E_{a_2}(g)=0,
                         \qquad \tau_D(g)=1.
 \tag{MA-JP.44}
\]
The scalar normalization may be complex.
The support is strictly positive; no nonnegativity of
the values of $g$ is asserted.

\begin{proof}
\emph{The functional and its convergence.}
On a fixed compact interval inside $(0,\infty)$,
integration by parts $N$ times gives
\[
 |E_{\sigma+i\gamma}(g)|
  \le C_{N,u,v,\operatorname{supp}g}
        (1+|\gamma|)^{-N}
                \max_{0\le j\le N}\|g^{(j)}\|_\infty
                  \quad(u\le\sigma\le v).
 \tag{MA-JP.45}
\]
Equation (MA-JP.41) makes the sum absolutely convergent
already with $N=2$.  The same estimate proves continuity
on each fixed-support test space and hence on its
usual inductive limit $\mathcal D_+$.

\emph{A compact-test approximation to a resolvent power.}
Choose an integer $k\ge4$ and a real
$\kappa>\max(v,\operatorname{Re}a_1,
\operatorname{Re}a_2)+1$.
Let smooth real profiles $\psi,\chi$ on $[0,\infty)$
have values in $[0,1]$, with
\[
 \begin{aligned}
 &\psi=0\text{ on }[0,1],\quad \psi=1\text{ on }[2,\infty),\\
 &\chi=1\text{ on }[0,1],\quad \chi=0\text{ on }[2,\infty).
 \end{aligned}
\]
For $0<\epsilon<1/4$ and $R\ge1$ set
\[
 g_{\epsilon,R}(t)=t^{k-1}e^{-\kappa t}
                    \psi(t/\epsilon)\chi(t/R).
 \tag{MA-JP.46}
\]
These are actual members of $\mathcal D_+$.
If $\sigma$ lies in any fixed interval containing
$[u,v]$ and $\operatorname{Re}a_1,
\operatorname{Re}a_2$ below $\kappa$, both the
$L^1$ norm and the second-derivative $L^1$ norm of
$e^{\sigma t}g_{\epsilon,R}(t)$ are bounded independently
of $\epsilon,R$.
Indeed the differentiated lower-transition terms are
supported on $[\epsilon,2\epsilon]$; a total of $j\le2$
derivatives there gives an integral $O(\epsilon^{k-j})$.
The terms without a differentiated lower profile have
the integrable powers $t^{k-1-j}$ near zero.
On $[R,2R]$ all differentiated upper profiles are
controlled by a polynomial in $R$ times a fixed
decreasing exponential.  The remaining integrals
are bounded by the same exponential moments.
Thus twice integrating by parts, with zero endpoint
terms for each compact test, proves the uniform bound
\[
 |E_{\sigma+i\gamma}(g_{\epsilon,R})|
                  \le C_\kappa(1+\gamma^2)^{-1}.
 \tag{MA-JP.47}
\]
Pointwise dominated convergence in the integral gives
\[
 E_\lambda(g_{\epsilon,R})\longrightarrow
       \int_0^\infty t^{k-1}e^{-(\kappa-\lambda)t}\,dt
       =\frac{\Gamma(k)}{(\kappa-\lambda)^k}.
 \tag{MA-JP.48}
\]
Equations (MA-JP.41) and (MA-JP.47) justify this
limit in the entire divisor sum, with both endpoint
cutoffs removed.  This justification is needed:
the uncut resolvent kernel itself is not a compact
smooth test supported away from zero.

\emph{Meromorphic pole separation.}
Suppose first that
$\tau_D=A_1E_{a_1}+A_2E_{a_2}$ on $\mathcal D_+$.
Apply that equality to (MA-JP.46), take the justified
limit, and cancel $\Gamma(k)$.  For every sufficiently
large real $\kappa$ this gives
\[
 F_D(\kappa):=
    \sum_{\lambda\in\mathcal Z}
                     \frac{d_\lambda}{(\kappa-\lambda)^k}
      =\frac{A_1}{(\kappa-a_1)^k}
             +\frac{A_2}{(\kappa-a_2)^k}.
 \tag{MA-JP.49}
\]
The series on the left converges normally on every
compact subset avoiding $\mathcal Z$.
For its tail, $\kappa$ in a fixed compact set gives
$|\kappa-\lambda|\ge |\operatorname{Im}\lambda|/2$
once the latter is sufficiently large; (MA-JP.41)
then proves normal convergence.  The remaining
finitely many terms are holomorphic off their poles.
Thus $F_D$ is meromorphic on $\mathbb C$, with
principal part
$d_{\lambda_0}(\kappa-\lambda_0)^{-k}$ at every
$\lambda_0\in\mathcal Z$.
There are no other singular terms there and this
coefficient is nonzero.

The identity theorem first applies on the right
half-plane and then on the complement of the locally
finite discrete union
$\mathcal Z\cup\{a_1,a_2\}$, which is connected.
Consequently (MA-JP.49) is a meromorphic identity.
Any divisor point outside $\{a_1,a_2\}$ would give
a noncancelling pole on the left and a holomorphic
right side there, a contradiction.
If the divisor is supported in those two points,
its defining formula is plainly in their span.
The same argument with a zero right side proves
$\tau_D=0$ on $\mathcal D_+$ if and only if the
divisor is empty.  In fact the coefficients $A_i$
in an identity are exactly the divisor lengths
at $a_i$, with zero when that point is absent.

\emph{The exact linear-algebra step imposing both moments.}
The map
$E=(E_{a_1},E_{a_2}):\mathcal D_+\to\mathbb C^2$
is onto.  To check this concretely, choose a smooth
bump $b$ supported in a sufficiently small interval
about a positive point so that both $E_{a_i}(b)$
are nonzero.  This is possible by approximating
a positive point mass and using continuity of
the two exponential functions.
Choose a positive shift $\Delta$ for which
$e^{(a_1-a_2)\Delta}\ne1$.
The bump and its right translate $b(t-\Delta)$
both belong to $\mathcal D_+$, and their moment
vectors have determinant
\[
 E_{a_1}(b)E_{a_2}(b)
          \bigl(e^{a_2\Delta}-e^{a_1\Delta}\bigr)\ne0.
 \tag{MA-JP.50}
\]
If $\tau_D$ vanishes on $\ker E$, it descends to a
linear functional on the actual two-dimensional
quotient identified by $E$.  It is therefore a
linear combination of $E_{a_1},E_{a_2}$.
The converse is immediate.  This proves the first
equivalence of (MA-JP.43); the pole argument proves
the second.
When the support is not contained in those two
points, a vector in $\ker E$ consequently has
nonzero $\tau_D$ value.  Dividing it by that value
gives (MA-JP.44).
\end{proof}

\paragraph{Application to the actual trace and the fast counterpath.}
Take the actual retained divisor from MA-WP,
$\lambda=\rho-1/2$ with positive length $r_\rho$,
and discard terms with $r_\rho=0$.
It satisfies (MA-JP.40)--(MA-JP.41), since
$0<r_\rho\le m_\rho$ and the fixed-function local
zero count is the one used in MT-M18 and MA-WP.
It is contained in a fixed vertical strip and
has no point at the polar exponents $\pm1/2$.
For these compact tests MA-WP and its actual
ambient lift already prove
\[
 \tau_D(g)=\operatorname{Tr}_{H/J}C_g^{H/J}
                    =\operatorname{Tr}_H(C_gP_Z).
 \tag{MA-JP.51}
\]
Thus the functional separated above is the trace
of an actual trace-class operator, not a scalar
substitute for an unconstructed quotient action.
If the retained divisor is nonempty, (MA-JP.44)
with $a_1=1/2,a_2=-1/2$ supplies the stronger
version of MA-JP-002:
\[
 \begin{gathered}
 g\in C_c^\infty((0,\infty)),\qquad g(0)=0,\\
 G(1/2)=G(-1/2)=0,\qquad
 \operatorname{Tr}_{H/J}C_g^{H/J}=1.
 \end{gathered}
 \tag{MA-JP.52}
\]
The actual counterpath construction MA-JP-003
therefore applies with a causal test, without
any other change to its trace-ideal argument.
For the rational trivial sector, Hardy's established
existence of a critical-line zero (1914, printed
p.~1012; full reference and source scope in MA-JP-003) supplies a
retained point for every $\alpha>0$ and $s\ge0$.
This is an explicit unconditional zero-location
\emph{existence} input, already identified in
MA-JP-003; it is neither RH nor an assumption of
an exterior zero.  For a different function the
nonemptiness hypothesis is retained unless its
own suitable zero has separately been supplied.

\paragraph{Application to the exact slow-band correction.}
Apply the same lemma to the upper divisor
\[
 D_{\eta,+}^{\rm out}
  =\sum_{\operatorname{Re}\rho_\eta\ge1/2+\alpha}
          m_\eta(\rho_\eta)[\rho_\eta-1/2].
 \tag{MA-JP.53}
\]
It satisfies the same count and strip assumptions.
The actual correction in MA-SB-003, or MA-SB-004
for the two character families, vanishes for every
causal compact test if and only if this divisor
is empty.  The equivalence remains true when
both polar moments are required to vanish,
since nontrivial completed zeros do not give
the exponents $\pm1/2$.
If this divisor is nonempty, there is a normalized
causal test for which the actual correction is one.
This last statement is conditional on nonemptiness;
the lemma proves no exterior-zero existence or
zero-location theorem.

The separation proof uses finite-length multiplicities
through actual traces and a normally convergent
meromorphic resolvent-power function.  It neither
requires simple zeros nor assigns a sign to values
of the complex test.  In the trace-framework placement
of BR-TR-01/02, the retained finite-length blocks
contribute their lengths to the trace; the analytic
separation of those contributions is the additional
argument proved here, not an arithmetic positivity
consequence of that categorical description.

\section{The actual slow-band formula for the complete source test class}
\label{sec:full-test-slow-band}

This section joins the separately proved causal, negative-support,
and norm-one estimates.  Its formula concerns the original weighted
orthogonal PW projections and actual trace-class cutoff operators.
It is not an assignment of negative translations to a one-sided
quotient.  In particular it records the difference between the
slow-band joint limit and the completed full-strip limit.

\paragraph{Notation and actual source families.}
Fix $0<\alpha<1/2$ and $s\ge0$, and put
\[
 \begin{gathered}
 H=L^2(\mathbb R,w(x)\,dx),\qquad
 w(x)=e^{2\alpha|x|}(1+x^2)^{s/2},\\
 d=\tfrac12-\alpha,\qquad u=\tfrac12+\alpha,
 \qquad b_s=\#\{j\in\mathbb Z_{\ge0}:2j<s-1\}.
 \end{gathered}
 \tag{MA-FT.1}
\]
The sector is either the rational trivial sector, a nonprincipal
primitive Dirichlet character, or a nontrivial finite-order
character over $\mathbb Q(\zeta_3)$ in the actual class-number-one
construction of MA-CGP-001--003.  The latter includes the explicit
non-induced character $\psi_{269,3}$ of MA-CH-003.
All zeros below are those of this sector's own primitive completed
$L$-function, with their full multiplicities $m_\rho$; trivial zeros
cancelled by the archimedean completion are not in the divisor.
Write $\lambda_\rho=\rho-1/2$ and
\[
 \begin{aligned}
 r_\rho^+&=
 \begin{cases}
 m_\rho,&\operatorname{Re}\rho<u,\\
 \min(m_\rho,b_s),&\operatorname{Re}\rho=u,\\
 0,&\operatorname{Re}\rho>u,
 \end{cases}\\
 r_\rho^{\rm st}&=
 \begin{cases}
 m_\rho,&d<\operatorname{Re}\rho<u,\\
 \min(m_\rho,b_s),&\operatorname{Re}\rho\in\{d,u\},\\
 0,&\operatorname{Re}\rho\notin[d,u],
 \end{cases}\\
 \delta_\rho^-&=r_\rho^+-r_\rho^{\rm st},\qquad
 \delta_\rho^+=m_\rho-r_\rho^+,\qquad
 R_\alpha(z)=-2\alpha-\overline z.
 \end{aligned}
 \tag{MA-FT.2}
\]
Thus $\delta^-$ consists of full lengths strictly below $d$ and
the missing length at $d$; $\delta^+$ consists of full lengths
strictly above $u$ and the missing length at $u$.
These are nonnegative integers.  They do not assert the existence
of any exterior zero or deficient boundary block.

Let $J_c\subset H_c=\mathbf1_{[c,\infty)}H$ be the actual PW
theta-image closure at support edge $c\le0$, and let
\[
 P_c=M_{\mathbf1_{[c,\infty)}},\qquad
 Q_c=P_{J_c},\qquad P_c^Z=P_c-Q_c,\qquad
 \tau_c(g)=\operatorname{Tr}_H(C_gP_c^Z).
 \tag{MA-FT.3}
\]
These projections are orthogonal in the single space $H$.
For the rational family $Y=e^{-c}$ is the source PW level.
For a character with conductor $\mathfrak f$,
$Q=|D_K|N\mathfrak f$ and $Y=Q\Lambda=e^{-c}$, where
$\Lambda$ is its primitive archimedean band.
MA-GP-010 and MA-CGP-002 identify these with the actual invariant
finite-$S$ global PW ranges when $S'=S$.  No extra unit-shell
transition is silently erased.  For fixed finite data the
Gram-resolvent realization remains MA-GP-008 or MA-CGP.18, with
both of its normal-equation factors.

For $g\in C_c^\infty(\mathbb R)$ with $g(0)=0$, define
\[
 \begin{aligned}
 G_+(z)&=\int_0^\infty g(t)e^{zt}\,dt,&
 G_-(z)&=\int_{-\infty}^0g(t)e^{zt}\,dt,\\
 G(z)&=G_+(z)+G_-(z),&
 \tau_{{\rm st},s}(g)&=\sum_\rho r_\rho^{\rm st}G(\lambda_\rho).
 \end{aligned}
 \tag{MA-FT.4}
\]
The two half-line integrals are scalars; no nonsmooth split of
$g$ is being inserted into an unproved operator functional.

\paragraph{MA-FT-001 (full-test actual limit, with both deficiencies).}
For the actual operators (MA-FT.3),
\[
 \begin{split}
 L_{\alpha,s}(g)
   &:=\lim_{c\to-\infty}\tau_c(g)\\
   &=\tau_{{\rm st},s}(g)
     +\sum_\rho\delta_\rho^-
       \left\{G_+(\lambda_\rho)
                      +G_-(R_\alpha\lambda_\rho)\right\}.
 \end{split}
 \tag{MA-FT.5}
\]
Every series displayed here and below is absolutely convergent.
For any source-thin cutoff path $X,Y\to\infty$, $Y/X\to0$,
and $\delta X\to0$, using the source profiles of MA-AT-001,
\[
 \lim\operatorname{Tr}_H
       (M_{\rho_{X,\delta}(e^\cdot)}C_gQ_{-\log Y})
                    =-L_{\alpha,s}(g).
 \tag{MA-FT.6}
\]
The same formula holds for the common fixed annular profile
$a_X(t)=A(t/X)-A(Xt)$ of MA-LA-014 and MA-CQ-003--004.

Choose a fixed finite set $S$ containing every prime-power
contribution, including the reciprocal contributions, met by
$h(t)=g(\log t)$, and also every conductor or different prime
in the character sectors.  Let $T^S_{X,Y,X^{-1}}(h)$ denote the actual
repaired local quotient operator with that common annular profile;
in the character sector this is the actual operator of MA-CQ,
with the conductor and different normalizations there.
Put $P_\eta(g)=G(1/2)+G(-1/2)$ in the rational trivial sector
and $P_\eta(g)=0$ in the stated nontrivial sectors.  Then
\[
 \begin{split}
 \Delta_{\alpha,s}(g)
 &:=\lim_{\substack{X,Y\to\infty\\Y/X\to0}}
  \left\{\operatorname{Tr}_H(M_{a_X(e^\cdot)}C_gQ_{-\log Y})
       -\left(\operatorname{Tr}T^S_{X,Y,X^{-1}}(h)
                                     -P_\eta(g)\right)\right\}\\
 &=\sum_\rho\delta_\rho^+G(\lambda_\rho)
    +\sum_\rho\delta_\rho^-
       \left\{G_-(\lambda_\rho)
                         -G_-(R_\alpha\lambda_\rho)\right\}.
 \end{split}
 \tag{MA-FT.7}
\]
At the lower edge $\operatorname{Re}\rho=d$ one has
$R_\alpha\lambda_\rho=\lambda_\rho$, so its contribution
to the second sum in (MA-FT.7) cancels exactly.  The lower-edge
length has not been lost from (MA-FT.5), and is not replaced
by a full multiplicity before that cancellation.

\begin{proof}
\emph{Actual operator inputs before a scalar evaluation.}
MA-GW-001--004 give the actual source range, its positive
semigroup, the retained half-plane and strip lengths, and the
ambient trace-class lift.  MA-AT-001 then proves
$C_gP_c^Z\in\mathcal S_1(H)$ for every smooth compact $g$,
without a support-sign condition.  These statements are in
the specified weighted space, not just in an equivalent
unweighted model at a fixed level.

First suppose $g$ vanishes on a neighbourhood of zero.
Its positive and negative restrictions extend by zero to
smooth compact tests $g_+$ and $g_-$ with strictly positive
and strictly negative support.  The actual causal trace of
MA-SB-005/006 and MA-GW gives
\[
 \tau_c(g_+)=\sum_\rho r_\rho^+G_+(\lambda_\rho).
 \tag{MA-FT.8}
\]
For the negative test, MA-AT-003, including its proved
character corollary (MA-AT.73), gives
\[
 \lim_{c\to-\infty}\tau_c(g_-)
  =\sum_\rho r_\rho^{\rm st}G_-(\lambda_\rho)
      +\sum_\rho\delta_\rho^-G_-(R_\alpha\lambda_\rho).
 \tag{MA-FT.9}
\]
The reflection in this formula comes from the actual
weighted Hilbert adjoint and the trace-norm limit of its
defect.  It does not replace the primitive sector by its
contragredient.  Add the two actual trace identities, use
$r^+=r^{\rm st}+\delta^-$, and obtain (MA-FT.5) in this case.

\emph{Removal of the norm-one neighbourhood.}
Choose $\chi\in C_c^\infty((-2,2))$ equal to one on
$[-1,1]$, with $0\le\chi\le1$, and put
\[
 r_\epsilon(t)=g(t)\chi(t/\epsilon),\qquad
 g_\epsilon=g-r_\epsilon.
 \tag{MA-FT.10}
\]
Each $g_\epsilon$ is an actual smooth compact test vanishing
near zero.  MA-NZ-001 proves the operator estimate
\[
 \eta_g(\epsilon):=
     \sup_{c\le0}\|C_{r_\epsilon}P_c^Z\|_1
                   \longrightarrow0.
 \tag{MA-FT.11}
\]
Thus $|\tau_c(g)-\tau_c(g_\epsilon)|\le\eta_g(\epsilon)$
uniformly in $c$.  For any $\epsilon,\epsilon'>0$, taking
the already established limits for $g_\epsilon$ and
$g_{\epsilon'}$ gives
\[
 |L_{\alpha,s}(g_\epsilon)-L_{\alpha,s}(g_{\epsilon'})|
              \le\eta_g(\epsilon)+\eta_g(\epsilon').
 \tag{MA-FT.12}
\]
These limits are Cauchy as $\epsilon\downarrow0$.
The same uniform estimate shows $\tau_c(g)$ tends to their
limit: choose $\epsilon$ first to make the two approximation
errors small, and then take $c$ sufficiently negative for
the fixed test $g_\epsilon$.  This is the operator argument
of MA-NZ-002, not convergence of a nonsmooth split in the
smooth test-function topology.

We verify that its scalar value is precisely (MA-FT.5).
For $z=\sigma+i\gamma$ in a fixed vertical strip,
the half-line transforms of $r_\epsilon$ satisfy
\[
 \left|\int_{\mathbb R_\pm}
                r_\epsilon(t)e^{zt}\,dt\right|
       \le C\min\{\epsilon^2,(1+\gamma^2)^{-1}\}.
 \tag{MA-FT.13}
\]
Here $r_\epsilon(0)=0$, while its first derivative at zero
need not vanish.  The $O(\epsilon^2)$ bound follows from
$g(t)=O(t)$ on the shrinking support.  Two integrations
by parts give the other bound with that derivative's
boundary term retained, since
$\|r_\epsilon''\|_1=O(1)$ and
$\|r_\epsilon'\|_1=O(\epsilon)$.
The multiplicity-weighted zero count on each unit-height
interval is $O(\log(2+|\gamma|))$ in these fixed sectors.
Summing (MA-FT.13), first over $|\gamma|\le1/\epsilon$
and then over the remaining intervals, gives
$O(\epsilon\log(2/\epsilon))$.
The reflected exponents have the same imaginary parts
and real parts in another fixed bounded interval.
The same bound therefore applies to every correction
sum.  Each is absolutely convergent for a fixed $g$,
again by two integrations by parts and $g(0)=0$.
All scalar expressions for $g_\epsilon$ consequently
tend to those in (MA-FT.5) for $g$.

\emph{The original cutoff operators.}
MA-AT.8 bounds the difference between the trace in
(MA-FT.6) and $-\tau_{-\log Y}(g)$ by
$\omega_g(\log(X-\delta))\to0$.  Its lower-support
condition is eventually satisfied because $Y/X\to0$,
the support of $g$ is fixed, and $\delta X\to0$.
Its diagonal volume term is zero because $g(0)=0$.
Combining it with (MA-FT.5) proves (MA-FT.6).

For the fixed annular profile, let
$a=\min\operatorname{supp}g$ when $g\ne0$.
The image of $C_gP_c^Z$ is supported in $[c+a,\infty)$.
Eventually $Xe^{c+a}>2$, so the lower term $A(Xe^x)$
vanishes on that image.  The upper term $A(e^x/X)$
is one for $x\le\log X$.  The exact tail factorization
MA-AT.4 therefore bounds the remaining trace error
by $\omega_g(\log X)\to0$.
The compact-position identity term again has trace zero.
This proves the annular version of (MA-FT.6).

The near-zero passage can also be made directly in these
cutoff traces.  For any bounded compact multiplier $M_\rho$
with norm at most one, the actual compact-kernel diagonal
formula and $r_\epsilon(0)=0$ give
\[
 \begin{split}
 \operatorname{Tr}(M_\rho C_{r_\epsilon}Q_c)
     &=-\operatorname{Tr}(M_\rho C_{r_\epsilon}P_c^Z),\\
 |\operatorname{Tr}(M_\rho C_{r_\epsilon}Q_c)|
     &\le\eta_g(\epsilon).
 \end{split}
 \tag{MA-FT.14}
\]
The bound is uniform in both source parameters.  Thus
the use of MA-NZ introduces no additional hidden
order-of-limits assumption at the original operator level.

\emph{The paired local trace.}
The actual local cutoff theorem MA-LA-014, or MA-CQ-003--004
in the character sector, applies to the same fixed
annular profile: $XY\to\infty$ and $Y/X\to0$ are its
two scale conditions.  The geometric limit includes
every prime term by the chosen finite $S$.
The already established full smooth explicit formula,
Meyer's Theorems~5.8 and~5.11 in the rational normalization
of MA-TR/MA-EG, and MA-CH-001 with MA-CQ in the stated
character normalization, give
\[
 \lim\left(\operatorname{Tr}T^S_{X,Y,X^{-1}}(h)
                                -P_\eta(g)\right)
           =-\sum_\rho m_\rho G(\lambda_\rho).
 \tag{MA-FT.15}
\]
This use of the explicit formula needs no information
about which side of $d$ or $u$ a zero lies on.
For characters the conductor term and different factors
are those of MA-CQ; the contact conductor term vanishes
on this class because $g(0)=0$, rather than being omitted
from the local construction.
Subtract (MA-FT.15) from (MA-FT.6), insert (MA-FT.5),
and use $m-r^{\rm st}=\delta^++\delta^-$.
This gives (MA-FT.7).  On the lower edge
$\lambda=-\alpha+i\gamma$, whence
$-2\alpha-\overline\lambda=-\alpha+i\gamma=\lambda$,
proving the final cancellation.
\end{proof}

\paragraph{MA-FT-002 (the full-test arithmetic gate is exact).}
For any one of the stated sectors at fixed $\alpha,s$,
the following conditions are equivalent:
\begin{enumerate}
\item The paired actual local and global limits in
(MA-FT.7) agree for every compact smooth $g$ with $g(0)=0$.
\item They agree for every $g\in C_c^\infty((0,\infty))$.
\item $\delta_\rho^+=0$ for every zero $\rho$.
\item $r_\rho^{\rm st}=m_\rho$ for every zero $\rho$.
\end{enumerate}
One may additionally require both polar moments
$G(1/2)=G(-1/2)=0$ in conditions~1 and~2 without
changing this equivalence.  Requiring these conditions
for every $0<\alpha<1/2$, at any fixed $s\ge0$,
is exactly the Riemann hypothesis for the stated sector.

\begin{proof}
Condition~1 implies~2.  For a strictly positive test,
$G_-=0$, and (MA-FT.7) becomes the trace functional of
the nonnegative integer divisor
\[
                   D^+=\sum_\rho\delta_\rho^+[\lambda_\rho].
 \tag{MA-FT.16}
\]
Its support is locally finite in a fixed vertical
strip and its count is bounded by the sector's full
multiplicity count.  MA-JP-004 therefore proves that
this functional vanishes on all such tests if and
only if $D^+$ is empty.  This proves~2 implies~3.
The same lemma gives a witness with both prescribed
polar moments zero whenever the divisor is nonempty:
nontrivial completed zeros in these sectors do not
have $\lambda=\pm1/2$.  This uses the classical
nonvanishing at the boundary points~0 and~1, already
part of the completed-sector setup, and does not
locate any nontrivial zero within the critical strip.

For completeness, the symmetry needed for~3 implies~4
does not require the character to be self-dual.
The unitary completed functional equation relates
$\Lambda_\eta(z)$ to $\Lambda_{\bar\eta}(1-z)$,
while conjugation relates
$\Lambda_{\bar\eta}(z)$ to
$\overline{\Lambda_\eta(\overline z)}$.
Their combination sends each zero in the same sector
to $1-\overline\rho$, preserving multiplicity.
The rational trivial sector has the same symmetry.
Thus any full or missing length below $d$ is reflected
to exactly the same full or missing length above $u$;
at the two edges the same number $b_s$ is used.
Condition~3 hence also gives $\delta^-=0$, which is~4.
Condition~4 makes both sums in (MA-FT.7) vanish, so it
implies~1.  This also proves the equivalence after
adding the two polar-moment conditions.

If all zeros have real part $1/2$, they lie strictly
inside every strip of positive width, so~4 holds for
every such width, independently of their multiplicity.
Conversely, a zero with real part different from$1/2$
is strictly exterior to a sufficiently narrow positive
strip, where~4 fails.  This proves the last assertion.
\end{proof}

\paragraph{Limit order, source locator, and input boundary.}
The actual completed full-strip trace is
$-\tau_{{\rm st},s}(g)$, by MA-WP-003.
Subtracting it from (MA-FT.6) leaves
\[
 -\sum_\rho\delta_\rho^-
       \{G_+(\lambda_\rho)+G_-(R_\alpha\lambda_\rho)\}.
 \tag{MA-FT.17}
\]
This is the exact order-of-limits correction; it is
not silently identified with zero.  Likewise the
fixed local smoothing profile in this theorem is
not an arbitrary path of thinning local profiles.

The source gap remains at Masullo's blank \S12.3,
p.~514, the asserted comparison in Theorem~12.4.1.1,
pp.~514--519, and the omitted global-to-local bridge
used on p.~520 in Theorem~12.5.1.1.  The printed cutoff
is Definition~9.3.3.1, p.~427; the thin-smoothing
conditions on p.~519 alone do not specify $Y/X\to0$.
MA-JP-001--003 already refute a rate-independent
comparison under only those conditions.  Here the
slow-band path is explicitly supplied and its entire
actual operator limit is proved.  Its agreement with
the unaltered full geometry has precisely the arithmetic
condition MA-FT-002.  No whole-plane zero-location
theorem is used in MA-FT-001 or its equivalence proof.
Using the stipulated theorem $A$ to satisfy that
condition over $\mathbb Q$ or in a Dirichlet sector is
an explicit zero-location input, hence circular for
the requested independent-method test.  The specified
non-induced cubic Hecke sector is not covered by $A$;
its Riemann hypothesis is still the condition, and
$M$ is only a sufficient conditional input.

\paragraph{Trace and finite-length placement.}
The trace-class operators in this theorem use the
same analytic trace framework as Grothendieck's
cote~1, scans~079--080 and~111--112, through the checked
predecessor placements ALPOGE\_BRIDGES BR-TR-01/02.
The integers $m_\rho,r_\rho^{\rm st},\delta_\rho^\pm$
are lengths of finite jet modules supported at
$\lambda_\rho$ in the affine spectral line.
For a finite block, additivity of trace is additivity
on the corresponding finite-length
$\mathbb C[z]$-modules and their $K_0$ classes;
the absolute convergence proved above is what permits
their analytic sum.  The $\lambda$-operations and
Fredholm determinant of the retained actual quotient
are already identified in MA-EG-002, with the exact
SGA~6 placement there.  This theorem does not promote
the reflected half-line expression to a new algebraic
endomorphism, construct an arithmetic weight theorem,
or make the positive source projections a faithful
positive action on the full jet space.  Those stronger
claims do not follow from a trace-class limit.

\section{Arbitrary thin smoothings on the actual local quotient}
\label{sec:local-variable-profiles}

The local theorems MA-LA-014 and MA-CQ-003--004 were first
proved with fixed profiles.  Their actual Fourier covariance
and trace-class compact-input and compact-output kernels
remove that restriction.  No bounds on derivatives of the
varying profiles will be imposed.

Use one of the actual rational or character local completions
\(\mathscr H_\beta\) constructed in MA-LA-011 or MA-CQ-002,
with the specified finite set \(S\), coherent additive character,
primitive finite factors, and probability compact Haar measure.
Let \(H_\beta:\mathscr H_\beta\to\mathscr H_{\bar\beta}\)
be its actual unitary Fourier map.  Thus
\(H_{\bar\beta}H_\beta=1\), including the finite Gauss phases
that cancel the odd real Fourier square.
For norm profiles \(\rho,b\), put
\[
 T(\rho,b;h)
    =M_{\rho(|x|_\infty)}A_h
                      H_{\bar\beta}M_{b(|\xi|_\infty)}H_\beta,
 \qquad g(y)=h(e^y).
 \tag{MA-LA.80}
\]
The rational trivial case is included.  These are the actual
quotient operators; the finite conductor and different
dilations are part of \(H_\beta\) throughout.

\paragraph{MA-LA-015 (profile comparison and the full local source limit).}
For a bounded position profile supported in a compact norm
annulus and a bounded frequency profile, (MA-LA.80) is trace
class.  When the position profile is smooth, it has the same
actual nuclear lift into the flat unit-shell section as
MA-LA-014 or MA-CQ-003, for every fixed profile.

The relevant exact covariance is
\[
 H_\beta A_h=A_{h^\vee}H_\beta,\qquad
 A_hH_{\bar\beta}=H_{\bar\beta}A_{h^\vee},\qquad
 h^\vee(t)=h(t^{-1}),\quad g^\vee(y)=g(-y).
 \tag{MA-LA.81}
\]
There are fixed nonnegative functions
\(\omega_g^{\rm out}(r),\omega_g^{\rm in}(r)\to0\)
as \(r\downarrow0\), depending only on the test in its
logarithmic local Hilbert realization, with the following
property.  If a bounded log multiplier \(a\), of norm at most
one, is supported in any interval of radius \(r\le1\), then
\[
 \|M_aC_g\|_1\le\omega_g^{\rm out}(r),\qquad
 \|C_gM_a\|_1\le\omega_g^{\rm in}(r).
 \tag{MA-LA.82}
\]
The bounds are independent of the centre of the interval.
For supports in \(k\) such windows, the corresponding bound
is \(k\) times the right side.
More quantitatively, for \(\operatorname{supp}g\subset[-m,m]\),
\[
 \omega_g^{\rm out}(r)+\omega_g^{\rm in}(r)
       \le C_m\|g\|_{C^4}\sqrt r,\qquad
 \|g\|_{C^4}=\max_{0\le j\le4}\|g^{(j)}\|_\infty ,
 \quad0<r\le1 .
 \tag{MA-LA.86a}
\]
The constant depends on the fixed compact auxiliary cutoff
and \(m\), not on a moving profile or its derivatives.

Here is a general moving-profile version of the local limit.
Let \(\ell_n<u_n\), \(v_n\in\mathbb R\), and
\(\epsilon_n\to0\).  Let \(0\le\rho_n,b_n\le1\)
have the following logarithmic support and plateau properties:
\[
 \begin{aligned}
 \rho_n(e^y)&=0&&
       (y\notin[\ell_n-\epsilon_n,u_n+\epsilon_n]),\\
 \rho_n(e^y)&=1&&
       (\ell_n+\epsilon_n\le y\le u_n-\epsilon_n),\\
 b_n(e^y)&=1&&(y\le v_n-\epsilon_n),\\
 b_n(e^y)&=0&&(y\ge v_n+\epsilon_n).
 \end{aligned}
 \tag{MA-LA.83}
\]
The position profiles are smooth; the usual smooth ball
frequency profiles may be used, and the trace comparison
also permits bounded frequency multipliers.
No monotonicity is required within a transition window.
Suppose
\[
                u_n+v_n\longrightarrow+\infty,\qquad
                \ell_n+v_n\longrightarrow-\infty .
 \tag{MA-LA.84}
\]
For every compact smooth \(h\) with \(h(1)=0\), the actual
trace satisfies
\[
       \operatorname{Tr}T(\rho_n,b_n;h)
            \longrightarrow A_\infty(h)+B_{\beta,S}(h).
 \tag{MA-LA.85}
\]
Here the archimedean term, primitive directional prime
coefficients and conductor convention are exactly
MA-CQ.21--.23, or MA-LA.69--.70 in the rational sector.
If instead \(\ell_n+v_n\to+\infty\) and
\(u_n-\ell_n\to+\infty\), the limit is zero.

\begin{proof}
\emph{Actual finite-profile trace class and the small-window bounds.}
Use the exact logarithmic local unitary MA-CQ.19, or its
rational version.  It identifies \(A_h\) with \(C_g\) on
ordinary \(L^2(\mathbb R,dy)\), and norm multipliers with
their functions of \(e^y\).
For a smooth compact position multiplier,
the kernel of \(M_\rho C_g\) is smooth and supported in a
compact rectangle.  A double Fourier-series expansion on a
containing rectangle has rapidly summable coefficients.
It therefore gives an actual nuclear map into the smooth
compact output interval, with every derivative seminorm.
The finite unit-shell section puts that output into the
same flat Schwartz source as in MA-CQ.24.
Composition with the remaining bounded Fourier-cutoff
operator preserves this nuclear lift and its Hilbert trace.
No uniform bound for these LF derivative seminorms is
claimed when the profile varies.
For a merely bounded compact position multiplier, factor it
through a fixed smooth compact multiplier equal to one on
its support.  This still proves Hilbert trace class, but
does not assert a lift into a smooth section for a
discontinuous output profile.

Fix \(\chi\in C_c^\infty(\mathbb R)\) equal to one on
\([-2,2]\).  Both
\[
 \begin{aligned}
 K_g^{\rm out}&=M_\chi C_g\in\mathcal S_1,\\
 K_g^{\rm in}&=C_gM_\chi\in\mathcal S_1,\qquad
 P_r=M_{\mathbf1_{[-r,r]}},\\
 \omega_g^{\rm out}(r)&=\|P_rK_g^{\rm out}\|_1,\qquad
 \omega_g^{\rm in}(r)=\|K_g^{\rm in}P_r\|_1 .
 \end{aligned}
 \tag{MA-LA.86}
\]
Trace class follows from the same compact-rectangle
calculation.  The projections \(P_r\) converge strongly
to zero as \(r\downarrow0\).
Finite-rank approximation of a trace-class operator
therefore shows that both displayed functions tend to zero.
For a multiplier supported in \([-r,r]\), factor through
\(P_r\) and \(\chi\) to obtain (MA-LA.82).
Unitary log translation moves any centre to zero and
commutes with \(C_g\), proving that the bound is independent
of the centre.  Splitting a multiplier supported on finitely
many windows into bounded pieces proves the last assertion.

For (MA-LA.86a), choose containing rectangles for the two
smooth kernels with side lengths depending only on \(m\)
and the fixed \(\chi\).
Their normalized trigonometric bases \(e_j,f_k\) have
supremum norms bounded independently of \(j,k\).
Two integrations by parts in each variable give Fourier
coefficients
\[
 |a_{jk}|
     \le C_m\|g\|_{C^4}(1+j^2)^{-1}(1+k^2)^{-1}.
\]
There are no boundary terms because the kernels vanish
near the edges of the rectangles.
Only derivatives of \(g\) of total order at most four
occur.  Hence the coefficient matrix is absolutely
summable with sum at most \(C_m\|g\|_{C^4}\).
The rank-one Fourier expansion for \(K_g^{\rm out}\)
now gives
\(\|P_rK_g^{\rm out}\|_1
 \le\sum_{j,k}|a_{jk}|\|P_re_j\|_2\|f_k\|_2
 \le C_m\|g\|_{C^4}\sqrt r\).
For \(K_g^{\rm in}P_r\), restrict the input basis instead.
This proves the stated rate for both moduli.

\emph{Fourier covariance with the correct dual sector.}
For normalized norm dilation
\(\mathcal D_tu(x)=t^{-1/2}u(t^{-1/d_\infty}x)\),
additive Fourier scaling gives
\(\mathcal F_\infty\mathcal D_t
 =\mathcal D_{t^{-1}}\mathcal F_\infty\).
The exponent \(1/2\) is the unitary norm normalization;
the real or complex additive measure scales by \(t\).
Every actual finite Fourier factor in MA-CQ.10--.11
is a function of commuting norm dilations, including
its conductor/different shift and Gauss scalar.
It consequently commutes with every \(A_h\).
Integrating the dilation identity against \(h(t)\,dt/t\)
proves (MA-LA.81).  The second identity can also be
obtained from the first using \(H_{\bar\beta}H_\beta=1\).
This proof keeps the intermediate \(\bar\beta\) sector;
it does not replace the odd Fourier map by an involution
before its finite phases have been included.

For two frequency profiles the actual operator difference
thus factors as
\[
 \begin{split}
 T(\rho,b_1;h)-T(\rho,b_2;h)
   &=M_\rho H_{\bar\beta}
           A_{h^\vee}M_{b_1-b_2}H_\beta,\\
 \|T(\rho,b_1;h)-T(\rho,b_2;h)\|_1
   &\le\|C_{g^\vee}M_{(b_1-b_2)(e^\cdot)}\|_1 .
 \end{split}
 \tag{MA-LA.87}
\]
For two position profiles the corresponding bound is
\[
 \|T(\rho_1,b;h)-T(\rho_2,b;h)\|_1
       \le\|M_{(\rho_1-\rho_2)(e^\cdot)}C_g\|_1 .
 \tag{MA-LA.88}
\]
The multipliers and Fourier maps outside the indicated
trace-class factors have norm at most one.
These are operator trace-norm bounds, not a formal
cyclic rearrangement of an unconstructed trace.
In particular no derivatives of either varying profile
occur, and no convolution has been commuted past a
Fourier cutoff.

\emph{Comparison with fixed smooth profiles.}
Fix \(0<\eta<\min(1/4,\log2)\).
Choose smooth nonincreasing \(A_\eta,B_\eta\) on
\([0,\infty)\), between zero and one, equal to one on
\([0,1]\) and zero on \([e^\eta,\infty)\).
Define
\[
 \begin{aligned}
 \rho_n^\eta(t)
   &=A_\eta(t/e^{u_n})-A_\eta(t/e^{\ell_n}),\\
 b_n^\eta(t)&=B_\eta(t/e^{v_n}),\qquad
 T_n^\eta=T(\rho_n^\eta,b_n^\eta;h).
 \end{aligned}
 \tag{MA-LA.89}
\]
These are the fixed-profile operators of MA-LA-014 or
MA-CQ-003, with upper position radius \(e^{u_n}\),
lower radius \(e^{\ell_n}\), and frequency radius
\(e^{v_n}\).  The necessary separation of the two
position endpoints holds eventually, because
(MA-LA.84) implies \(u_n-\ell_n\to\infty\).

The difference \(\rho_n-\rho_n^\eta\) is supported
in two log windows about \(\ell_n,u_n\), each of
radius at most \(\epsilon_n+\eta\).
The difference \(b_n-b_n^\eta\) is supported in one
log window about \(v_n\) with the same radius.
All these differences have absolute value at most one.
Equations (MA-LA.82), (MA-LA.87)--(MA-LA.88) give
\[
 \|T(\rho_n,b_n;h)-T_n^\eta\|_1
   \le2\omega_g^{\rm out}(\epsilon_n+\eta)
               +\omega_{g^\vee}^{\rm in}(\epsilon_n+\eta).
 \tag{MA-LA.90}
\]
For \(\operatorname{supp}g\subset[-m,m]\), (MA-LA.86a)
bounds this error by
\(C_m\|g\|_{C^4}\sqrt{\epsilon_n+\eta}\).
This is a quantitative comparison between actual profile
realizations.  It does not assert a rate for the remaining
fixed-profile geometric limit.
For each fixed \(\eta\), the already proved fixed-profile
theorem gives \(\operatorname{Tr}T_n^\eta\to
A_\infty+B_{\beta,S}\), by (MA-LA.84).
The limit is independent of \(\eta\).
Taking the upper limit in \(n\) in (MA-LA.90), then
\(\eta\downarrow0\), proves (MA-LA.85).
The fixed-profile zero limit under the second stated
regime gives the other conclusion in exactly the same
way.  Only differences of profile realizations are
controlled in trace norm here; the varying operators
\(T(\rho_n,b_n;h)\) themselves are not asserted to
converge in \(\mathcal S_1\).
\end{proof}

\paragraph{The printed source smoothing conditions and the full global comparison.}
For the position cutoff in Masullo, Definition~9.3.3.1,
printed p.~427, take a smooth \(0\le\rho_n\le1\) supported
in \([1/X_n,X_n]\), equal to one on
\([1/X_n+\delta_n,X_n-\delta_n]\), with its stipulated
multiplicative symmetry if desired.
For the actual PW frequency ball, take \(0\le b_n\le1\)
equal to one up to \(Y_n-\delta_n\), and zero from
\(Y_n\) onward.  An outward thin smoothing is covered
as well.  Inward smoothing satisfies (MA-LA.83) with
\[
 \begin{gathered}
 \ell_n=-\log X_n,\quad u_n=\log X_n,\quad v_n=\log Y_n,\\
 \epsilon_n=
 \max\{\log(1+\delta_nX_n),
             -\log(1-\delta_n/X_n),
             -\log(1-\delta_n/Y_n)\}\longrightarrow0 .
 \end{gathered}
 \tag{MA-LA.91}
\]
For positive \(\delta_n\) all denominators here are valid
eventually.  The printed conditions
\(\delta_nX_n\to0\), \(\delta_nY_n\to0\) on p.~519,
together with \(X_n,Y_n\to\infty\), imply this conclusion.
Independent position and frequency transition widths
are equally allowed when their corresponding logarithmic
widths tend to zero.
The bounds permit arbitrarily rapid smooth oscillations
inside each transition interval.

The support convention at frequency is not inferred
from the annular definition on p.~427.
Masullo's Step~II on printed p.~516 explicitly specifies
the Fourier subset \(\{|x|_{F,S}\le\lambda'\}\);
the displayed composition on pp.~516--517 is its
Fourier-conjugated ball multiplier.
The footnote on p.~514 separately discusses the
one-sided position intention and an independently
moving lower endpoint.  The present theorem extends
the actual repaired annular position operator and
that specified frequency ball, retaining a smooth
flat section at each finite annular cutoff.

If also \(Y_n/X_n\to0\), then (MA-LA.84) holds and
the local limit is (MA-LA.85), for every such choice
of smooth source profiles.  The global source-profile
limit has already been proved in MA-FT-001 at every
polynomial weight and for every compact smooth
\(g\) with \(g(0)=0\).
Thus its paired local correction (MA-FT.7), its
full-strip comparison (MA-FT.17), and its exact
arithmetic gate MA-FT-002 remain valid with these
arbitrary thin source local smoothings, in place
of the initially fixed local profiles.
For a character the level remains \(Y=Q\Lambda\),
and \(S\) includes the conductor/different places
and all prime-power contributions to the test.
No finite Euler factor, Gauss phase or conductor
dilation is removed in making this extension.

\paragraph{MA-LA-016 (an annular frequency cutoff is a different operator).}
Keep the same source position cutoffs, but now replace
the PW frequency ball by a smooth annular frequency
profile \(\widetilde b_n\), bounded between zero and one,
supported in \([1/Y_n,Y_n]\), and equal to one on
\([1/Y_n+\delta_n,Y_n-\delta_n]\).
Assume
\[
 X_n,Y_n\to\infty,\qquad Y_n/X_n\to0,\qquad
 \delta_nX_n\to0,\qquad\delta_nY_n\to0 .
 \tag{MA-LA.92}
\]
Then its actual local operator is trace class, with
a nuclear flat-section lift for every smooth
position profile, and
\[
             \operatorname{Tr}T(\rho_n,\widetilde b_n;h)
                         \longrightarrow0
                         \qquad(h(1)=0).
 \tag{MA-LA.93}
\]
This is not the frequency operator specified on
Masullo p.~516.

\begin{proof}
For fixed \(0<\eta<\min(1/4,\log2)\), use the same
position reference as before and put
\[
 \widetilde b_n^\eta(t)
        =B_\eta(t/Y_n)-B_\eta(Y_nt).
 \tag{MA-LA.94}
\]
For large \(n\) this is bounded between zero and one.
Its difference from \(\widetilde b_n\) is supported
in two shrinking log windows about
\(-\log Y_n,\log Y_n\), up to the fixed reference
width \(\eta\).
The lower relative source width is now
\(\log(1+\delta_nY_n)\to0\).
The same proof of (MA-LA.90), with two frequency
windows instead of one, makes the trace error
arbitrarily small as first \(n\to\infty\), then
\(\eta\to0\).

Linearity of the actual bounded frequency multiplier
gives the exact fixed-reference decomposition
\[
 \begin{split}
 T(\rho_n^\eta,\widetilde b_n^\eta;h)
    ={}&T(\rho_n^\eta,B_\eta(\,\cdot/Y_n);h)\\
       &-T(\rho_n^\eta,B_\eta(Y_n\,\cdot);h).
 \end{split}
 \tag{MA-LA.95}
\]
Both terms are the actual ball operators already
constructed, with frequency radii \(Y_n\) and
\(1/Y_n\), respectively.
For the first, the two coupled products are
\(X_nY_n\to\infty\) and \(Y_n/X_n\to0\).
For the second, they are
\(X_n/Y_n\to\infty\) and \(1/(X_nY_n)\to0\).
The fixed-profile local theorem therefore gives
the same limit \(A_\infty(h)+B_{\beta,S}(h)\)
for both terms.
Their difference tends to zero, and the profile
comparison proves (MA-LA.93).
\end{proof}

This alternative changes an actually detectable trace.
For example, choose a nonzero nonnegative
\[
                   h\in C_c^\infty((1,\sqrt2)).
 \tag{MA-LA.96}
\]
Then \(h(1)=0\), every prime-power and reciprocal
prime-power contribution is zero, and the real
parity or complex archimedean integral in MA-CQ.21
is strictly positive.
The actual ball operator therefore has positive
limit \(A_\infty(h)\), whereas (MA-LA.93) gives
zero for the annular frequency alternative.
This test uses no zero-location information.
It is a scope counterexample to changing the
frequency support convention, not an attribution
of that changed convention to the printed PW map.

The trace framework is the one already justified
by the actual nuclear section and common Hilbert
trace in MA-LA-014 and MA-CQ-003.
The new ingredient is the derivative-free
trace-norm comparison (MA-LA.87)--(MA-LA.90).
It enlarges the proved cutoff class while leaving
the finite-length arithmetic deficiency of
MA-FT-002 explicit; it provides no further
zero-location or arithmetic positivity theorem.

\section{One Gram regularization for every compact smooth test}
\label{sec:universal-gram-regularizer}

The finite cyclic construction in MA-GP-008 removes its
regularizer at fixed data.  Its use along a moving source path
need not choose a different regularizer for each test.
The following trace-ideal argument fixes that quantifier without
asserting an effective rate for a small singular value of a theta map.

\paragraph{MA-GP-011 (a test-independent regularizer on a prescribed path).}
Fix an actual rational or character source family for which
MA-GP-008 or MA-CGP-003 constructs the closed theta operator.
The physical space is the same fixed
$H=L^2(\mathbb R,w(x)\,dx)$ as above; $w$ and $w^{-1}$
are bounded on compact intervals.  Let any prescribed sequence
of source data be indexed by $n\ge1$.  At its $n$th level write
\[
 \begin{gathered}
 T_n=V_nG_n^{1/2},\qquad G_n=T_n^*T_n,\qquad Q_n=V_nV_n^*,\\
 F_{n,\varepsilon}=V_n\sqrt{G_n/(G_n+\varepsilon)},
 \qquad Q_{n,\varepsilon}=F_{n,\varepsilon}F_{n,\varepsilon}^*,\\
 A_{n,g}=M_{\chi_n(e^\cdot)}C_g.
 \end{gathered}
 \tag{MA-GP.150}
\]
Here $T_n$ is closed and densely defined on its specified
auxiliary Hilbert source, $V_n$ is its polar partial isometry,
and $\chi_n$ is any of the smooth compact source position
cutoffs, with $|\chi_n|\le1$.  The source may vary with $n$.

There is a single decreasing dyadic sequence
$\varepsilon_n\downarrow0$, depending on the source path
and its cutoffs but independent of $g$, such that every
$g\in C_c^\infty(\mathbb R)$ satisfies
\[
 \begin{split}
 &\|A_{n,g}(Q_{n,\varepsilon_n}-Q_n)\|_1\\
 &\quad+
 \|F_{n,\varepsilon_n}^*A_{n,g}F_{n,\varepsilon_n}
                             -V_n^*A_{n,g}V_n\|_1
                    \longrightarrow0.
 \end{split}
 \tag{MA-GP.151}
\]
More precisely, for all $n$ with
$\operatorname{supp}g\subset[-n,n]$, the left side is at most
$2^{-n}\|g\|_{C^5}$.
The convergence is therefore uniform on each bounded
$C^5$ family of tests supported in a fixed compact interval.
No requirement $g(0)=0$ is needed for this regularizer estimate.

For each fixed $n$ the actual domain-correct cyclic identity is
\[
 \begin{split}
 \operatorname{Tr}_H(A_{n,g}Q_{n,\varepsilon_n})
 &=\operatorname{Tr}
       (F_{n,\varepsilon_n}^*A_{n,g}F_{n,\varepsilon_n})\\
 &=\operatorname{Tr}
       \left((G_n+\varepsilon_n)^{-1/2}
              K_{n,g}(G_n+\varepsilon_n)^{-1/2}\right),
 \end{split}
 \tag{MA-GP.152}
\]
with the actual finite rational-dilation kernel of MA-GP-008
or its character kernel of MA-CGP-003.  Thus along every
prescribed source path whose levels, cutoff profiles, and relative
rates satisfy MA-FT-001, these same regularized cyclic
operators have that full-test limit simultaneously
for every compact smooth $g$ with $g(0)=0$.

\begin{proof}
\emph{A fixed-order compact-kernel estimate.}
For a fixed $n$ and a fixed test-support interval $[-m,m]$,
choose a compact interval $I$ containing the logarithmic
support of $\chi_n$, and a compact interval $J$ containing
$I+[-m,m]$.  Choose smooth compact functions $\psi_I,\psi_J$
equal to one on these respective intervals.  On unweighted
$L^2(\mathbb R)$ the smooth compact kernel
\[
             K_g(x,y)=\psi_I(x)g(x-y)\psi_J(y)
 \tag{MA-GP.153}
\]
defines a trace-class operator with
\[
                       \|K_g\|_1\le C_{n,m}\|g\|_{C^4}.
 \tag{MA-GP.154}
\]
Here is a finite-order proof of the bound.  Enclose its support
in a rectangle whose sides are strictly outside that support,
and periodically extend the kernel on that rectangle.
Apply $(1-\partial_x^2)(1-\partial_y^2)$ before taking its
double Fourier coefficients.  Their absolute values are bounded
by a constant times
$\|g\|_{C^4}(1+j^2)^{-1}(1+k^2)^{-1}$.
The corresponding rank-one Fourier kernels have uniformly
bounded trace norms, and the displayed double sequence is
summable.  Restricting their absolutely convergent nuclear
series to the original rectangle proves (MA-GP.154).
The cutoff derivatives in this argument are fixed with $n,m$;
no derivative of the varying test of order greater than four
is used.

Conjugation of $A_{n,g}$ by multiplication by $\sqrt w$
expresses it as $K_g$ multiplied on its two sides by
bounded compactly supported factors.  Specifically the
output factor is $\chi_n(e^x)\sqrt{w(x)}$ on $I$, and
the input factor is $w(y)^{-1/2}$ on $J$.
The two auxiliary smooth cutoffs equal one wherever
$A_{n,g}$ can have a nonzero kernel.  Consequently
\[
                       \|A_{n,g}\|_1\le C'_{n,m}\|g\|_{C^4}.
 \tag{MA-GP.155}
\]
Local boundedness of $w,w^{-1}$ suffices; its cusp at zero
is not differentiated.  The constant may grow with $n$.
No uniform-in-$n$ bound is claimed or needed here.

\emph{Compactness of a bounded higher-derivative test family.}
Let
\[
 \mathcal B_m=\{g\in C_c^\infty(\mathbb R):
       \operatorname{supp}g\subset[-m,m],\ \|g\|_{C^5}\le1\}.
 \tag{MA-GP.156}
\]
Its closure is compact in the $C^4$ norm: the derivatives
through order four are uniformly bounded, their highest
derivative is uniformly Lipschitz by the fifth-derivative
bound, and the lower derivatives are equicontinuous as well.
The Arzel\`a--Ascoli theorem, successively for those finitely
many derivatives, gives compactness with consistent limiting
derivatives.  The support remains in the same interval.
Equation (MA-GP.155) therefore makes
$\{A_{n,g}:g\in\mathcal B_m\}$ relatively compact in
$\mathcal S_1(H)$ for each fixed $n,m$.

\emph{Uniform removal of the regularizer on that compact set.}
For fixed $n$, spectral calculus gives
$F_{n,\varepsilon}\to V_n$ and
$F_{n,\varepsilon}^*\to V_n^*$ strongly as
$\varepsilon\downarrow0$, with all these operators bounded
by one.  The same is true of $Q_{n,\varepsilon}\to Q_n$,
including their adjoints.  These are the actual polar
operators of the closed source, so no unbounded adjoint
product is inserted into a trace.

For every fixed $A\in\mathcal S_1(H)$, it follows that
\[
 \begin{split}
 \mathcal E_{n,\varepsilon}(A)
 &:=\|A(Q_{n,\varepsilon}-Q_n)\|_1\\
 &\quad+\|F_{n,\varepsilon}^*AF_{n,\varepsilon}
                              -V_n^*AV_n\|_1
                 \longrightarrow0.
 \end{split}
 \tag{MA-GP.157}
\]
Indeed the assertion follows first for rank-one $A$ from
strong convergence on both sides and then for arbitrary
$A$ from finite-rank approximation and the uniform bounds.
The difference of the errors for two inputs has absolute
value at most $4\|A-A'\|_1$, by the trace-ideal inequality.
A finite net in the compact closure from (MA-GP.156)
therefore gives
\[
 \sup_{g\in\mathcal B_m}
             \mathcal E_{n,\varepsilon}(A_{n,g})
                  \longrightarrow0
             \quad(n,m\text{ fixed}).
 \tag{MA-GP.158}
\]

At stage $n$, take $m=n$.  Among sufficiently small dyadic
numbers choose the first one, in decreasing order, below
$\min(2^{-n},\varepsilon_{n-1}/2)$ for which the supremum
in (MA-GP.158) is at most $2^{-n}$; at $n=1$ omit the
previous-stage bound.  The convergence in (MA-GP.158)
ensures that the set is nonempty.  This defines
$\varepsilon_n$ using the entire fixed unit test family,
so it is independent of the particular $g$ later tested.
For any fixed nonzero $g$ and all sufficiently large $n$,
the function $g/\|g\|_{C^5}$ belongs to $\mathcal B_n$.
Homogeneity proves the quantitative bound in (MA-GP.151);
the zero test is immediate.

Finally MA-GP-008 and MA-CGP-003 give (MA-GP.152) at
each fixed level, with their already proved graph-domain
and trace-class identities.  Equation (MA-GP.151) changes
its global trace by a quantity tending to zero.
Every independently established actual source limit,
in particular MA-FT-001 on its prescribed paths, is
therefore the limit of these same cyclic regularized
operators for every test in the stated class.
\end{proof}

\paragraph{What the construction does and does not add.}
The choice is mathematically specified by compact operator
norms.  It is not an algorithm for computing a numerical
regularizer from $X,Y$ and the conductor, and it asserts
no polynomial or other effective rate.
It removes a test-dependent choice from the actual
finite Gram transport without changing its arithmetic
remainder.  The two Gram factors and every primitive
character normalization remain present.  Its trace and
Fredholm identities have the same cote~1 placement as
MA-GP-008; no extra positivity, full-jet faithfulness,
or zero-location premise is used.

\section{The exact trace-norm obstruction in the PW exhaustion}
\label{sec:pw-trace-norm-convergence}

Strong convergence of the actual finite PW projections was proved
in MA-GP-001 and MA-CGP-003.  It does not by itself imply trace-norm
convergence of an uncut smoothed projection.  The following result
identifies that additional condition exactly, and proves its
sufficiency on the original weighted Hilbert space.

Use any of the actual rational or character families of MA-FT-001,
with \(0<\alpha<1/2\), \(s\ge0\), and
\[
 H=L^2(\mathbb R,w(x)\,dx),\qquad
 w(x)=e^{2\alpha|x|}(1+x^2)^{s/2}.
 \tag{MA-TN.1}
\]
Retain \(J_c,P_c,Q_c,P_c^Z\) from (MA-FT.3), and put
\[
 J=\overline{T_{1/2}\mathcal I}^{\,H},\qquad
 P_G=P_J,\qquad Z=J^\perp,\qquad P^Z=1-P_G.
 \tag{MA-TN.2}
\]
The full-strip lengths \(r_\rho^{\rm st}\), half-line lengths
\(r_\rho^+\), and nonnegative difference
\(\delta_\rho^-=r_\rho^+-r_\rho^{\rm st}\) are exactly (MA-FT.2).
In particular every length here is a length of an actual
finite primary quotient, with the boundary-jet thresholds
already proved in MA-WP and MA-GW.

Let \(U_tf(x)=f(x-t)\) on \(H\).  Write
\[
 \omega=\alpha+s/4,\qquad \nu>\omega,\qquad
 a_0=\nu-\omega,\qquad
 \ell(x)=(\log w)'(x)\quad(x\ne0).
 \tag{MA-TN.3}
\]
Thus \(|\ell|\le2\omega\).  On \(Z_c=P_c^ZH\) let \(B_c\)
be the adjoint quotient generator of MA-NZ, and on \(Z\)
let \(B\) be the adjoint of the actual full quotient
translation generator.  Extend their resolvents by zero:
\[
 R_c=(\nu-B_c)^{-1}P_c^Z,\qquad
 R=(\nu-B)^{-1}P^Z,\qquad
 A_c=\operatorname{Re}R_c,\quad A=\operatorname{Re}R.
 \tag{MA-TN.4}
\]
These are actual bounded operators on the single space \(H\).

\paragraph{MA-TN-001 (an exact positive trace defect).}
As \(c\to-\infty\), both \(R_c\to R\) and \(R_c^*\to R^*\)
strongly.  The operators \(A_c,A\) are positive and trace class,
and
\[
 \begin{aligned}
 \operatorname{Tr}A_c
   &=\sum_\rho r_\rho^+
                 \operatorname{Re}(\nu-\lambda_\rho)^{-1},\\
 \operatorname{Tr}A
   &=\sum_\rho r_\rho^{\rm st}
                 \operatorname{Re}(\nu-\lambda_\rho)^{-1},\\
 \lim_{c\to-\infty}\|A_c-A\|_1
   &=D_\nu:=
       \sum_\rho\delta_\rho^-
                 \operatorname{Re}(\nu-\lambda_\rho)^{-1},
       \qquad \lambda_\rho=\rho-\tfrac12.
 \end{aligned}
 \tag{MA-TN.5}
\]
All three sums converge absolutely.  In particular \(D_\nu=0\)
if and only if every \(\delta_\rho^-\) is zero.  If this happens,
\[
              \|R_c-R\|_2\longrightarrow0,\qquad
              \|R_c^2-R^2\|_1\longrightarrow0.
 \tag{MA-TN.6}
\]
There is no assertion that this arithmetic condition holds
without an additional argument.

\begin{proof}
\emph{Actual strong limits.}
The dense increasing PW exhaustion of MA-GP-001 or MA-CGP-003
gives \(Q_c\to P_G\) strongly, while \(P_c\to1\) strongly.
Consequently \(P_c^Z=P_c-Q_c\to P^Z\) strongly.  On the
supported half-line the adjoint of positive translation is
\(P_cU_t^*P_c\).  Since the actual \(J_c\) is invariant under
positive translation, its orthogonal complement in \(H_c\)
is invariant under this adjoint.  The adjoint semigroup,
extended by zero, is therefore \(P_cU_t^*P_c^Z\).
The full \(J\) is translation invariant, so the corresponding
full-space expression is \(U_t^*P^Z\).  The physical bound
\(\|U_t\|\le e^{\omega t}\), \(t\ge0\), gives the integrable
majorant in the vectorwise identities
\[
 R_c=\int_0^\infty e^{-\nu t}P_cU_t^*P_c^Z\,dt,\qquad
 R=\int_0^\infty e^{-\nu t}U_t^*P^Z\,dt .
 \tag{MA-TN.7}
\]
Strong convergence follows by dominated convergence for
each vector.  Taking adjoints in these integrals gives
integrands \(P_c^ZU_tP_c\) and \(P^ZU_t\), so the same
argument proves adjoint strong convergence.  All resolvent
norms are at most \(a_0^{-1}\).

\emph{Positive traces from the actual complete jet spaces.}
MA-NZ.6--.8 gives
\[
 a_0R_c^*R_c\le A_c,\qquad
 a_0R_cR_c^*\le A_c
 \tag{MA-TN.8}
\]
and the first trace in (MA-TN.5).  The identical energy
argument applies to \(R\).  For completeness, MA-WP.4 says
that \(Z\) is the closed span of the continuous retained
jet representatives
\[
       w(x)^{-1}x^j e^{\overline{\lambda_\rho}x},
                  \qquad 0\le j<r_\rho^{\rm st}.
 \tag{MA-TN.9}
\]
They are in the graph domain of the full adjoint generator,
whose expression is \(w^{-1}(w\,\cdot)'\).
Finite height spans are invariant and have eigenvalues
\(\overline{\lambda_\rho}\) with these exact lengths.
Their orthogonal projections increase strongly to \(P^Z\).
The trace of the compression of \(A\) to such a finite span
is the corresponding sum of
\(r_\rho^{\rm st}\operatorname{Re}(\nu-\lambda_\rho)^{-1}\).
Positivity and an orthonormal exhaustion, exactly as in
MA-NZ.8, prove that \(A\) is trace class and give its
asserted trace.

For every retained point \(\operatorname{Re}\lambda_\rho\le\alpha\).
Thus the summands are positive, and the real part of each
resolvent factor is \(O((1+|\operatorname{Im}\rho|^2)^{-1})\).
The fixed-function multiplicity-weighted local zero count
proves absolute convergence.  It also proves the assertion
about vanishing of \(D_\nu\), without assuming that any
discarded point exists.

\emph{The positive-operator trace-norm lemma.}
We prove the precise form needed here.  Suppose \(T_j,T\ge0\)
are trace class, \(T_j\to T\) strongly, and
\(\operatorname{Tr}T_j=S\) is constant, with
\(\operatorname{Tr}T\le S\).  Then
\[
             \|T_j-T\|_1\longrightarrow S-\operatorname{Tr}T.
 \tag{MA-TN.10}
\]
Choose a finite-rank orthogonal projection \(F\), and put
\(F^\perp=1-F\).  Strong convergence on the finite-dimensional
range of \(F\), and self-adjointness, show that the three
blocks with at least one \(F\) in the decomposition of
\(T_j-T\) tend to zero in trace norm.  For example
\((T_j-T)F\) has uniformly bounded rank and its
Hilbert--Schmidt norm tends to zero by summing over an
orthonormal basis of \(FH\).
The remaining block has the upper bound
\[
 \|F^\perp(T_j-T)F^\perp\|_1
 \le \operatorname{Tr}T_j-\operatorname{Tr}(FT_jF)
       +\operatorname{Tr}T-\operatorname{Tr}(FTF).
 \tag{MA-TN.11}
\]
First let \(j\) tend to the limit.  Then let \(F\) range
through a finite-rank exhaustion for the positive trace-class
operator \(T\).  The upper bound tends to
\(S-\operatorname{Tr}T\).  The reverse bound is
\(\|T_j-T\|_1\ge|\operatorname{Tr}(T_j-T)|\).
This proves (MA-TN.10); the proof applies to the parameter
\(c\to-\infty\), without discretizing its values.
Apply it to \(A_c,A\) to prove the last line of (MA-TN.5).

\emph{Hilbert--Schmidt convergence when the defect vanishes.}
Suppose \(D_\nu=0\), so \(A_c\to A\) in trace norm.
For a finite-rank orthogonal \(F\), (MA-TN.8) gives
\[
 \|R_c(1-F)\|_2^2
       \le a_0^{-1}\operatorname{Tr}((1-F)A_c).
 \tag{MA-TN.12}
\]
The analogous estimate holds for \(R\).
Choose \(F\) so that the right side for \(A\) is small;
trace-norm convergence makes it uniformly small for all
sufficiently negative \(c\).  On \(FH\), strong convergence
implies \(\|(R_c-R)F\|_2\to0\).  Splitting the right input
by \(F+(1-F)\) proves \(\|R_c-R\|_2\to0\).
The two Hilbert--Schmidt norms are bounded, so
\[
 \|R_c^2-R^2\|_1
   \le \|R_c\|_2\|R_c-R\|_2
             +\|R_c-R\|_2\|R\|_2\longrightarrow0 .
 \tag{MA-TN.13}
\]
This proves (MA-TN.6).
\end{proof}

\paragraph{MA-TN-002 (the actual ambient convergence criterion).}
The following conditions are equivalent:
\begin{enumerate}
 \item \(\delta_\rho^-=0\) for every \(\rho\);
 \item for every \(g\in C_c^\infty(\mathbb R)\),
 \(\|C_gP_c^Z-C_gP^Z\|_1\to0\);
 \item the convergence in condition 2 holds for every
 \(g\in C_c^\infty((0,\infty))\).
\end{enumerate}
These are actual ambient operators, not only quotient
compressions or equalities of scalar traces.
Strong convergence \(C_gP_c^Z\to C_gP^Z\) holds for every
fixed compact smooth \(g\) without condition 1.
If condition 1 fails, a causal test can be chosen for which
\[
 \liminf_{c\to-\infty}
       \|C_gP_c^Z-C_gP^Z\|_1
 \ge\left|\sum_\rho\delta_\rho^-G(\lambda_\rho)\right|>0,
 \qquad G(z)=\int g(t)e^{zt}\,dt .
 \tag{MA-TN.14}
\]
The test may also satisfy \(G(1/2)=G(-1/2)=0\).

\begin{proof}
Only sufficiency requires a new endpoint argument.
All operators whose norms are asserted to converge are
trace class by MA-AT-001 and MA-WP-003.
Assume condition 1.  Set \(k(x)=\nu-\ell(x)/2\);
it is a fixed bounded positive multiplier with \(k\ge a_0\).
For \(v\) in the first graph domain of the full half-line
adjoint generator, integration by parts gives
\[
 \operatorname{Re}\langle v,(\nu-\mathcal L_c)v\rangle_H
       =\int_c^\infty k(x)|v(x)|^2w(x)\,dx
                     +\tfrac12 w(c)|v(c+)|^2 .
 \tag{MA-TN.15}
\]
There is no cusp contribution in this first-order identity:
\(w\) is locally absolutely continuous.  The boundary
term at infinity vanishes.  Indeed
\(\sqrt w\,v\) belongs to ordinary \(H^1(c,\infty)\),
because \(\ell\) is bounded, and an \(H^1\cap L^2\)
function tends to zero at infinity.
The full-line version has neither endpoint term, by
the same argument at both ends.

Define the bounded scalar functional
\[
       e_c(f)=\sqrt{w(c)}(R_cf)(c+).
 \tag{MA-TN.16}
\]
Its boundedness follows from (MA-TN.15), applied to
\(v=R_cf\).  Polarization gives the exact operator identities
\[
 A_c=R_c^*M_kR_c+\tfrac12e_c^*e_c,\qquad
 A=R^*M_kR .
 \tag{MA-TN.17}
\]
MA-TN-001 gives convergence of the left sides in trace norm
and of \(R_c\) in Hilbert--Schmidt norm.  Consequently
\(R_c^*M_kR_c\to R^*M_kR\) in trace norm.  Subtracting
in (MA-TN.17) gives
\[
                         \|e_c\|^2\longrightarrow0 .
 \tag{MA-TN.18}
\]
This controls a moving boundary functional in its actual
norm, rather than by a nonuniform equivalence of weights.

Use the second-order bulk operator of MA-NZ.13:
\[
 \mathcal B_g=C_{g''}-2C_{g'}M_{\nu-\ell}
             +C_gM_{(\nu-\ell)^2-\ell'_{\rm reg}},
 \quad
 \ell'_{\rm reg}(x)=s(1-x^2)/(1+x^2)^2.
 \tag{MA-TN.19}
\]
It is bounded on \(H\) for every fixed compact smooth \(g\).
The exact second-order integration by parts in MA-NZ.14--.15,
now with \(r=g\), gives for \(c<0\)
\[
 C_gP_c^Z=\mathcal B_gR_c^2+\mathcal E_{c,g}R_c^2
                              -4\alpha\,g\otimes\kappa_c,
 \qquad
 \kappa_c(f)=(wR_c^2f)(0).
 \tag{MA-TN.20}
\]
The same calculation on the full line gives
\[
             C_gP^Z=\mathcal B_gR^2-4\alpha\,g\otimes\kappa,
 \qquad \kappa(f)=(wR^2f)(0).
 \tag{MA-TN.21}
\]
All graph domains and the cusp contribution are exactly
those checked in MA-NZ; in particular the jump is
\(\ell(0+)-\ell(0-)=4\alpha\).

Here is the moving endpoint in a form controlled by
(MA-TN.18).  Put \(g_c(x)=g(x-c)\).  Since
\(\mathcal L_cR_c^2=\nu R_c^2-R_c\), the endpoint
formula MA-NZ.14 becomes
\[
 \mathcal E_{c,g}R_c^2
  =\frac{(\nu-\ell(c))g_c-g_c'}{\sqrt{w(c)}}
                                  \otimes(e_cR_c)
       +\frac{g_c}{\sqrt{w(c)}}\otimes e_c .
 \tag{MA-TN.22}
\]
The sign and the coefficient follow by replacing
\((wR_c^2f)'(c)\) by
\(\sqrt{w(c)}(\nu e_cR_cf-e_cf)\) in that formula.
The norms of the two target vectors in (MA-TN.22)
are bounded uniformly in \(c<0\): on the fixed translated
support, \(w(c+t)/w(c)\le e^{2\omega|t|}\).
Together with \(\|R_c\|\le a_0^{-1}\), (MA-TN.18) proves
\(\|\mathcal E_{c,g}R_c^2\|_1\to0\).

The cusp functional has a fixed-vector description.
The actual half-line differential resolvent, squared, is
\[
 (\nu-\mathcal L_c)^{-2}h(x)
   =w(x)^{-1}\int_x^\infty
                  (y-x)e^{-\nu(y-x)}w(y)h(y)\,dy ,
       \qquad x\ge c .
 \tag{MA-TN.23}
\]
This follows either by two vectorwise semigroup integrals
or by composing their absolutely convergent kernels.
For \(c<0\), evaluation at \(x=0\) therefore says
\(\kappa_c(f)=\langle P_c^Zf,k_0\rangle_H\), where
\(k_0(y)=\mathbf1_{y\ge0}y e^{-\nu y}\), with the
inner product linear in its first argument.
The same formula holds with \(P^Z\) for \(\kappa\).
The vector \(k_0\) belongs to \(H\), since \(\nu>\alpha\).
Strong convergence of the orthogonal projections gives
\(\|\kappa_c-\kappa\|\to0\).

Now (MA-TN.6) controls the bulk difference in
(MA-TN.20)--(MA-TN.21), (MA-TN.22) controls the endpoint,
and the fixed-vector formula controls the cusp.
Their sum converges in trace norm.  This proves
condition 1 implies condition 2 for every compact smooth
test, without a condition at zero.

Condition 2 implies condition 3 by restriction.
For a causal test, the actual finite quotient trace
of MA-GW and the actual full-strip trace of MA-WP give
\[
 \operatorname{Tr}(C_gP_c^Z-C_gP^Z)
                     =\sum_\rho\delta_\rho^-G(\lambda_\rho).
 \tag{MA-TN.24}
\]
It is independent of \(c\).  The multiplicity-weighted
local count makes the series absolutely convergent.
If condition 3 holds, this sum vanishes for every
compact causal test.  MA-JP-004, applied to the
nonnegative integer divisor with coefficients
\(\delta_\rho^-\), then makes that divisor empty.
The completed nontrivial divisor in these primitive sectors
has no point at \(\rho=0\) or \(1\), as in MA-FT-002.
Thus any nonempty deficiency contains a point away from
the removed exponents \(\lambda=\pm1/2\).
The version of MA-JP-004 with those two polar moments
therefore gives the final assertion and (MA-TN.14).
Strong convergence, independently of
all these conditions, follows from \(P_c^Z\to P^Z\)
and boundedness of \(C_g\).
\end{proof}

\paragraph{Actual cutoff consequence and arithmetic scope.}
Under the equivalent conditions of MA-TN-002, let
\(M_n=M_{\chi_n}\) be any uniformly bounded compact
position multipliers converging strongly to one,
and let \(c_n\to-\infty\).  There is no required relative
rate.  Then the renormalized actual operators satisfy
\[
 \begin{aligned}
 M_nC_gQ_{c_n}-M_nC_gP_{c_n}
      &=-M_nC_gP_{c_n}^Z
                  \longrightarrow-C_gP^Z
                              \quad\hbox{in }\mathcal S_1(H).
 \end{aligned}
 \tag{MA-TN.25}
\]
Indeed split the difference into
\(-M_n(C_gP_{c_n}^Z-C_gP^Z)\) and
\(-(M_n-1)C_gP^Z\).  The first tends to zero by
MA-TN-002; the second does so by finite-rank approximation
of the fixed trace-class operator.  When \(g(0)=0\),
the subtracted compact-input term has trace zero
by MA-AT.12.  It is still an operator term and is not
discarded in the asserted trace-norm identity.
This does not assert trace-norm convergence of the
unrenormalized global operator, nor does it change the
different local geometric limits on different joint paths.
In particular the fast-path counterexample MA-JP-003
is unaffected.

The own-sector functional equation used in MA-FT-002
pairs the upper and lower deficiencies, including their
boundary lengths.  Thus the condition in MA-TN-002 is
equivalent to the exact upper-deficiency condition of
MA-FT-002.  For all positive widths it is the same sector's
RH statement.  The stipulated \(\mathsf A\) supplies it
for the rational and Dirichlet sectors explicitly as
zero-location input; \(M\) suffices for the specified
non-induced cubic Hecke sector.  The equivalence and
the actual conditional operator convergence above use
neither assertion in their proofs.

\paragraph{Trace-framework placement and what has changed.}
The coefficients \(\delta_\rho^-\) remain differences of
finite lengths in the spectral primary modules of MT-M18
and MA-WP.  Their additive trace contribution is the
finite-length \(K_0\) interpretation already used in
MA-EG and MA-FT, and the analytic traces belong to the
nuclear/Fredholm framework of Grothendieck, cote~1,
scans 79--80 and 111--112, as checked in
ALPOGE\_BRIDGES BR-TR-01/02.
The ordinary positive operator \(A_c\) is the real part
of an actual dissipative resolvent.  It is not a
faithful arithmetic star-action on the full jet module,
and it yields no weight, purity, or motivic positivity
theorem.  The Artin-motive interpretation of the character
sector remains that of MA-CH and MA-CQ, with SGA~1,
Expos\'e~V, \S7 and Deligne (1979), \S6.1--6.3 as the
already recorded exact placement.
This result identifies an operator topology consequence
of the existing arithmetic gate; it is not counted as a
strict reduction of that gate.

\section{The five-dimensional Artin gate}
\label{sec:ma-s5}

\paragraph{Status after the Fourier argument.}
MA-FQ-004 below proves $\beta=b+n$ under the exact hypothesis
$(\mathcal R_B)$ stated in this section. The bounds, positive-excess
signatures, and formal order witness developed here remain valid
intermediate deductions; their arithmetic positive-excess branch
is excluded by that later theorem. The hypothesis
$(\mathcal R_B)$ is not proved for every extension in the family.
The input $M$ supplies it and also supplies the additional Hecke
zero-location hypotheses needed for GRH of all irreducible factors.

\paragraph{Notation and the audited input.}
We retain precisely Calegari's family: the Galois group of $L/\mathbb Q$
is $S_5$, complex conjugation is of type $(12)(34)$, and $5$ is
unramified with Frobenius of the same type. Write
$E=L^{A_5}$, $F=L^{F_{20}}$, $H=EF$, and let $\eta$ be the character
of $E/\mathbb Q$. In the character-table notation of
GRH\_NIGHT MT-S506, put
\[
 V=\rho_4,\quad W=\rho_5,\quad X=\rho_6,\qquad
 B=\zeta_H/\zeta_F=L_F(s,\eta_F),\quad C=L(s,\eta),
 \quad A=L_f(s,\varpi).
\]
Here $\varpi$ is Calegari's specified cuspidal comparison representation,
not an arbitrary degree-five automorphic representation. The notation
$A$ in this section is an $L$-function; the Dirichlet GRH theorem is
denoted $\mathsf A$ when it is needed. All finite $L$-functions include
the ramified factors. Set
\[
 \begin{gathered}
 \beta=\operatorname{ord}_{1/2} B,\quad b=\operatorname{ord}_{1/2} C,
 \quad n=\operatorname{ord}_{1/2} A,\quad
 k=\beta-b-n,\\
 (\mathcal R_B):\ B(\sigma)\ne0\quad
 (0<\sigma<1,\ \sigma\ne1/2,\ \sigma\in\mathbb R).
 \end{gathered}
\]
We use the two all-prime identities
\[
 B=C L_W,\qquad \zeta_F=\zeta_{\mathbb Q}L_{W\eta}.
 \tag{MA-S5.1}
\]
They are GRH\_NIGHT MT-S506, equation (S506.6), and MT-S508,
equation (S508.2); their character calculations and induction conventions
were rechecked in this run.

\subsection*{MA-S5-001. What the earlier comparison actually proves}

\paragraph{Audited finding.}
The assertions MT-S504, MT-S505, MT-S507 and MT-S508 pass the checks
specified below, with their stated hypotheses. In particular, under
$(\mathcal R_B)$,
\[
 k\ge0,\qquad
 k=0\ \Longleftrightarrow\ \Omega\text{ is finite}
 \ \Longleftrightarrow\
 \varpi\text{ and }W\text{ have matching factors at every place}.
 \tag{MA-S5.2}
\]
The last condition also gives the all-place match for the quadratic
twist. It proves automorphy and holomorphy of the nontrivial irreducible
Artin factors in this family. GRH for all these factors requires the
additional zero-location input $M$, or separate corresponding Hecke
GRH hypotheses. The condition $(\mathcal R_B)$ alone is not such an input.

Calegari's initial comparison is equality on a set of primes of
Dirichlet density one. It is not equality at all but finitely many
places. Replacing the first assertion by the second would assume
the finiteness of $\Omega$ that is to be proved.

\paragraph{Complete check of the deductions used.}
The primary locators are Calegari (2013), Theorem 1.1 (author-hosted
publisher PDF p.~2), Theorem 4.7 (p.~12), Lemma 4.8 (pp.~13--14),
Landau Lemma 4.9 (p.~14), Lemma 4.10 (pp.~14--15), and the final
proof of Theorem 1.2 (p.~15). These are publisher-PDF page numbers,
not the arXiv-v1 page numbers used in MT-S504. The complete proofs
of the comparison lemmas, their intervening Landau argument, and
the final deduction were read. The construction of the cuspidal
representations remains Calegari's cited theorem.

There are also harmless sign/asymptotic slips in the printed proofs
of Lemma 4.6 and Theorem 4.7 (publisher-PDF p.~12). In the first,
the density-weighted coefficient is $-1/3$, so the logarithm of
the comparison quotient is
$-\frac13\log(1/(s-1))+O(1)$ as $s\downarrow1$.
The quotient is not just holomorphic but nonzero at $1$, since
each of its four factors has an exact simple pole there. Its
bounded logarithm supplies the contradiction with the corrected
asymptotic. In the second, the corrected identity is
\[
 \log\left|\frac{L(s,\Pi\times\Pi^\vee)}
                       {L(s,\operatorname{End}\xi)}\right|
       =8\sum_{p\in\Omega}p^{-s}+O(1).
\]
The left side is bounded for the same simple-pole reason.
Monotone convergence therefore gives
$\sum_{p\in\Omega}p^{-1}<\infty$, which is stronger than density
zero. The printed $o(|\log(s-1)|^{-1})$ does not follow from
boundedness; the valid density conclusion uses
$O(1)=o(\log(1/(s-1)))$. These corrections do not change any
comparison theorem used here.

For clarity, with $G$ denoting the archimedean factors, the quotient
that is actually used is
\[
 R(s)=\frac{G_{\varpi}(s)A(s)}{G_W(s)L_W(s)}
     =R_{\mathcal B}(s)D_\Omega(s),\qquad
 D_\Omega(s)=\prod_{p\in\Omega}
       \left(\frac{1+p^{-s}}{1-p^{-s}}\right)^2.
 \tag{MA-S5.3}
\]
For the convention $L_p=P_p(p^{-s})^{-1}$, the finite factor in
$R_{\mathcal B}$ is $P_{W,p}/P_{\varpi,p}$. The reversed displayed
finite-polynomial ratio in Calegari's Lemma 4.8 is the normalization
slip already corrected in MT-S508 (S508.9); it does not change the
real-interval nonvanishing of the finite factor. At $p\in\Omega$,
\[
 P_{\varpi,p}(T)=(1+T^2)(1-T)^3,\qquad
 P_{W,p}(T)=(1+T^2)(1-T)(1+T)^2.
\]
The finite and archimedean factor $R_{\mathcal B}$ is holomorphic
and nonzero at real $0<s<1$. If $\Omega$ is infinite, its positive
Euler series has an abscissa $\gamma\in(0,1)$ and a pole at
$\gamma$. The strict upper bound uses convergence at $1$, which
follows from the proof of Theorem 4.7, rather than merely from
density zero. The exclusion of $\gamma=0$ is the arbitrary-finite-subset
argument on publisher p.~14. Thus Lemmas 4.8--4.9 supply an actual
real pole, and Lemma 4.10 supplies exact comparison if $\Omega$
is finite.

Substitution of (MA-S5.1) gives
\[
 R(s)=\frac{G_{\varpi}(s)}{G_W(s)}\frac{A(s)C(s)}{B(s)},
 \qquad \operatorname{ord}_{1/2}R=-k.
\]
Under $(\mathcal R_B)$ every possible real pole is at $1/2$.
Infinite $\Omega$ therefore implies $k>0$; finite $\Omega$ gives
$R=1$ and $k=0$. These two exhaustive cases prove the central
equivalence and its inequality. This is a check of the proof of
MT-S508, not a new proof of its automorphic construction.

MT-S504 assumes the stronger input $M$ and $\operatorname{ord}_{1/2}L_W\le0$.
Its signed Brauer argument makes $L_W$ nonzero and holomorphic on
the other real points; a real pole of $R$ would then require a
positive central order of $L_W$. Consequently its implication is
valid. MT-S505 correctly uses that implication to exclude a negative
central order, then uses $\zeta_{\mathbb Q}(\sigma)<0$ on $(0,1)$
in the second identity (MA-S5.1). It does not infer nonvanishing
of $L_W$ from its holomorphy. The parity part concerns meromorphic
orders, for which the Laurent-series argument is valid.

For MT-S507 the whole displayed cone proof and its sharpness proof
were checked. Its premise $c,d\ge0$ is an extra holomorphy input.
The exhaustive finite calculation was re-executed after removing
the exploratory optimizer: $156$ subgroups, $19$ conjugacy classes,
and $43$ induced-linear multiplicity vectors were recovered. Every
vector was reconstructed as a nonnegative integer combination of
the twelve displayed generators. The four sharpness witnesses
were checked in integer arithmetic. The files in
\texttt{checks/s5/} contain that computation. These witnesses are
cone points, not number fields.

Finally, the half-plane and one-$B$ consequences in MT-S508 were
checked with their precise scope. Full GRH for $B$ gives GRH for
$A$ and $C$ and critical-line divisor support for $L_W$. It does
not by itself give GRH for $L_V,L_{V\eta},L_X,L_{W\eta}$.
This last distinction is needed when interpreting task 2's
phrase ``holomorphy and GRH for every irreducible factor.''

\subsection*{MA-S5-002. The residual excess is even}

\paragraph{Statement.}
In this family the three orders $\beta,b,n$ are even, independently
of $(\mathcal R_B)$. Consequently, under $(\mathcal R_B)$,
\[
        k\in 2\mathbb Z_{\ge0}.
\]
At this stage the hypothetical positive branch of MT-S508 is therefore
\[
             k\ge2\quad\text{and }k\text{ even},
 \tag{MA-S5.4}
\]
not the potentially odd excess allowed by its stated bound $k\ge1$.
The orders $n_\eta=\operatorname{ord}_{1/2}L_f(s,\varpi\otimes\eta)$
and $f=\operatorname{ord}_{1/2}\zeta_F$ are even as well.

\begin{proof}
The completed Dedekind zeta functions have functional equation with
sign $+1$. The completed quotients for $H/F$ and $E/\mathbb Q$
therefore give sign $+1$ for the quadratic Hecke functions $B$ and
$C$. Their gamma and conductor factors are holomorphic and nonzero
at $1/2$. If a nonzero meromorphic function satisfies
$\Lambda(1/2+t)=\Lambda(1/2-t)$, its leading Laurent term
$a t^r$ satisfies $a=(-1)^r a$, so $r$ is even. Applying this to
$B,C$, and $\zeta_F$ gives the first two parities and the assertion
about $f$.

At every unramified prime, the Satake multiset of $\varpi$ is
invariant under inversion. Away from $\Omega$ it is the multiset
of the real orthogonal representation $W$; on $\Omega$ it is
$\{1,1,1,i,-i\}$. Thus $\varpi$ and $\varpi^\vee$ agree at all
but finitely many places. Ordinary strong multiplicity one applies
to these two already automorphic representations, and gives
$\varpi\simeq\varpi^\vee$. Notice that this use of multiplicity
one is different from the unavailable all-but-finite comparison
of $\varpi$ with $W$.

A self-dual cuspidal representation of odd degree is of orthogonal
type. Arthur (2013), Theorems 1.4.1 and 1.5.3(a), gives this
orthogonal/symplectic classification, and Theorem 1.5.3(b), applied
also to the trivial representation of $\mathrm{GL}_1$, gives
$\epsilon(1/2,\varpi)=+1$. The exact epsilon-factor statement is
also stated by Arthur as Theorem 3(b) in his author-hosted
\emph{The Endoscopic Classification of Representations}, p.~8;
the specialization is explicitly recorded as Theorem 3.3.6 in
Dalal--Gerbelli-Gauthier, \texttt{arXiv:2410.01976v2}. The input
here is Arthur's classical theorem, not a new theorem from that
preprint. Since the completed standard $L$-function is holomorphic
at the centre and its gamma factor is nonzero there, the same
Laurent argument gives $n\in2\mathbb Z_{\ge0}$.

The quadratic twist is again self-dual cuspidal of odd degree, so
the same argument gives even $n_\eta$. Finally $k=\beta-b-n$ is
even, and MT-S508 supplies $k\ge0$ under $(\mathcal R_B)$.
No zero-location statement was used in any parity argument.
\end{proof}

\paragraph{Why MT-ART10--MT-ART12 do not close this gate.}
Those statements were read with their hypotheses: the integral
class and hermitian invariants in MT-ART11--MT-ART12 give
functional-equation parities, particularly on the symplectic
character group, and MT-ART10's forced zero is a quaternionic
example. The present $S_5$ irreducibles are real orthogonal and
have root number $+1$ (Calegari, p.~3). No negative sign is available
to force or rule out a positive even excess. The additional
automorphic root-number input above does eliminate the odd values
of $k$, but parity alone does not imply $k=0$.

\subsection*{MA-S5-003. A Rankin--Selberg bound for the excess}

\paragraph{Statement.}
Let $\Pi$ and $\xi$ be Calegari's four-dimensional automorphic
representation and finite-image representation of the central
extension of $S_5$ used in his section 4. Put
\[
 c=\operatorname{ord}_{1/2}L_V,\quad
 d=\operatorname{ord}_{1/2}L_{V\eta},\quad
 e=\operatorname{ord}_{1/2}L_X,
\]
and let $r_4,r_5$ be the central orders of the completed
Rankin--Selberg functions of $\Pi\times\Pi^\vee$ and
$\varpi\times\varpi^\vee$, respectively. Then
\[
 r_4=n+d+e-k,\qquad
 r_5=r_4+c+f.
 \tag{MA-S5.5}
\]
In particular, under $(\mathcal R_B)$,
\[
     0\le k\le n+d+e.
 \tag{MA-S5.6}
\]
If $A(1/2)L_{V\eta}(1/2)L_X(1/2)\ne0$, the comparison closes.
Combining with MA-S5-002, failure to close requires
\[
 2\le k\le n+d+e,\qquad k\in2\mathbb Z,
 \qquad n,d,e\in2\mathbb Z_{\ge0}.
 \tag{MA-S5.7}
\]
Thus it requires a central zero in at least one of these three
known entire automorphic functions. This is a necessary condition;
it is not an asserted nonvanishing theorem for them.

\begin{proof}
We first make the finite representation calculation explicit.
Scalar central elements act trivially on $\operatorname{End}\xi$,
so it is a representation of $S_5$. In the order of classes
$1,2A,2B,3A,4A,5A,6A$, its character is
\[
                    (16,0,0,1,0,1,3).
\]
For the odd classes, this is the squared absolute trace of the
projective matrices in Calegari's section 2, publisher-PDF p.~6:
the traces are $0,0,\sqrt3$ up to scalar. On $A_5$, it follows
from $\xi|_{N.A_5}=\operatorname{Sym}^3\varrho$. Indeed the
symmetric-cube traces at projective orders $2,3,5$ have absolute
squares $0,1,1$, respectively. For example, if the two eigenvalues
of a determinant-one lift are $e^{i\theta},e^{-i\theta}$, the
symmetric-cube trace is
$2\cos3\theta+2\cos\theta$; taking $\theta=\pi/2,\pi/3,\pi/5$
or $2\pi/5$ gives these squares. Comparing every entry with
the $S_5$ table proves
\[
          \operatorname{End}\xi=\mathbf1\oplus V\eta\oplus W\oplus X.
 \tag{MA-S5.8}
\]
An entrywise product and inner product in the same character table
gives
\[
 W\otimes W=\mathbf1\oplus V\oplus V\eta\oplus W\oplus W\eta\oplus X.
 \tag{MA-S5.9}
\]
Both identities, and the Euler-polynomial calculation below, have
an exact integer/cyclotomic certificate in \texttt{checks/s5/}.

Enlarge the finite bad-prime set $\mathcal B$ if needed. At a good
prime outside $\Omega$, the Satake multiset of $\Pi$ is the
multiset of $\xi$ up to a common sign; their endomorphism
multisets therefore agree. At $p\in\Omega$, up to a scalar root
of unity and replacing $i$ by $-i$, the two multisets are
\[
          S_\Pi=\{1,i,1,i\},\qquad S_\xi=\{1,i,-1,-i\}.
\]
These are the second and fourth exceptional rows of Calegari's
$4A$ table, publisher p.~11. For every positive integer $j$,
\[
 \left|\sum_{\alpha\in S_\Pi}\alpha^j\right|^2
 -\left|\sum_{\alpha\in S_\xi}\alpha^j\right|^2
 =\begin{cases}8,&j\text{ odd},\\0,&j\text{ even}.\end{cases}
 \tag{MA-S5.10}
\]
For odd $j$ the first squared trace is $8$ and the second is
$0$; for $j\equiv2\pmod4$ both are $0$; for $4\mid j$ both
are $16$. The scalar root of unity cancels in each endomorphism
trace. Expanding Euler logarithms in $\Re s>1$ gives
\[
 \log\frac{L_p(s,\Pi\times\Pi^\vee)}{L_p(s,\operatorname{End}\xi)}
 =8\sum_{j\text{ odd}}\frac{p^{-js}}j
 =2\log\left(\frac{1+p^{-s}}{1-p^{-s}}\right)^2.
\]
Multiplying at good primes, then completing, yields a meromorphic
identity
\[
 \frac{\Lambda(s,\Pi\times\Pi^\vee)}
      {\Lambda(s,\operatorname{End}\xi)}
          =U_4(s)D_\Omega(s)^2,
 \tag{MA-S5.11}
\]
where $U_4$ consists of finitely many local factors, gamma factors,
and conductor powers and is holomorphic and nonzero at $1/2$.

Here is the check of that last, potentially important assertion.
The base change of $\Pi$ to $E$ has the finite-image local
parameters of $\operatorname{Sym}^3\varrho$. Restriction of a
Weil--Deligne parameter to the finite-index local Weil group
over $E$ has zero monodromy and finite image; consequently the
original parameter has zero monodromy and finite image as well.
Its image is a union of finitely many translates of the finite
restricted image. Thus its endomorphism parameter has eigenvalues
that are roots of unity, including at bad primes. No polynomial
$\det(1-p^{-s}\operatorname{Frob})$ can vanish at $s=1/2$.
The Artin parameter has the same property. Since $E$ is real,
base change is split at infinity, where the parameters are also
finite-image parameters; the gamma factors are products of
$\Gamma_{\mathbb R}(s)$ and $\Gamma_{\mathbb R}(s+1)$ and are
nonzero and finite at the centre. Conductor powers never vanish.
This proves the assertion about $U_4$ without an all-place
matching assumption.

By (MA-S5.3), $\operatorname{ord}_{1/2}D_\Omega=-k$.
Artin multiplicativity in (MA-S5.8), together with
$\zeta_{\mathbb Q}(1/2)\ne0$, gives
\[
 \operatorname{ord}_{1/2}\Lambda(s,\operatorname{End}\xi)=d+x+e,
 \qquad x=\operatorname{ord}_{1/2}L_W=\beta-b=n+k.
\]
Taking orders in (MA-S5.11) therefore proves
$r_4-(d+x+e)=-2k$, or the first identity in (MA-S5.5).

For the degree-five Rankin--Selberg comparison, the odd-power
traces at an exceptional prime are $3$ and $-1$, and the even-power
traces agree. Their squared-trace difference is again $8$ on odd
powers and $0$ on even powers. The same calculation, now with
(MA-S5.9), gives
\[
 r_5=c+d+x+y+e-2k,
 \qquad y=\operatorname{ord}_{1/2}L_{W\eta}=f.
\]
Subtracting the formula for $r_4$ proves $r_5-r_4=c+f$.
The bad-factor argument is identical, using the finite-image
base change of $\varpi$.

The completed Rankin--Selberg function of two unitary cuspidal
representations is holomorphic at $1/2$: its possible poles lie
on $\Re s=0$ or $1$. An exact usable locator is Cogdell (2004),
\emph{Lectures on $L$-functions, Converse Theorems, and Functoriality
for $\mathrm{GL}_n$}, Theorem 9.2(i), printed p.~74 (PDF p.~78),
which states the Jacquet--Shalika analytic result. Thus $r_4\ge0$.
Calegari's Theorem 1.1 gives $d,e\ge0$, and $n\ge0$ because
$A$ is entire. Under $(\mathcal R_B)$, MT-S508 gives $k\ge0$.
These prove (MA-S5.6). If the stated three central values are
nonzero, the upper bound is zero. Finally $d,e$ are even by the
root-number-$+1$ functional equations of the real $S_5$
representations, and MA-S5-002 proves the remaining parities.
\end{proof}

\subsection*{MA-S5-006. Symmetric-square holomorphy removes the comparison order}

\paragraph{Statement.}
With the same notation, the sharper bound is
\[
                 0\le k\le d+e
       \qquad\text{under }(\mathcal R_B).
 \tag{MA-S5.16}
\]
Consequently $L_{V\eta}(1/2)L_X(1/2)\ne0$ suffices to close the
comparison, regardless of the central order of $A$. A hypothetical
positive branch must satisfy
\[
       2\le k\le d+e,\qquad k,d,e\in2\mathbb Z_{\ge0}.
 \tag{MA-S5.17}
\]

\begin{proof}
Let $\nu=\psi^{-1}$ be Calegari's finite-order similitude character.
The exact inputs are his Lemma 4.2 and the following displayed
decomposition (publisher-PDF pp.~9--10), and Lemma 4.6
(publisher-PDF p.~12):
\[
 \Pi^\vee=\Pi\otimes\nu^{-1},\qquad
       (\wedge^2\Pi)\otimes\nu^{-1}=\mathbf1\boxplus\varpi.
\]
Choose a finite set $S$ containing the archimedean places, the
places above $2$, and all ramified places in these identities.
At every unramified place, the decomposition of the tensor square
into symmetric and exterior squares gives
\[
 L^S(s,\Pi\times\Pi^\vee)
  =L^S(s,\Pi,\operatorname{Sym}^2\otimes\nu^{-1})
       \zeta_{\mathbb Q}^S(s)A^S(s).
\]
This identity first holds in an absolute-convergence half-plane
and then meromorphically. The omitted factors are holomorphic
and nonzero at $1/2$, by the finite-image parameter check in
MA-S5-003. The character $\nu$ has finite order, and the central
character of $\Pi$ is unitary: its base change has finite-order
central character, so a power of it has finite order already
over $\mathbb Q$. Thus Takeda (2015), Theorem 7.1, applies to
the twisted symmetric-square factor. It is holomorphic at $1/2$.

Write $q$ for its central order. Since $\zeta_{\mathbb Q}(1/2)\ne0$,
the displayed identity gives $r_4=q+n$. MA-S5-003 gives
$r_4=n+d+e-k$. Hence
\[
                         q=d+e-k\ge0.
\]
Combining this with MT-S508 proves (MA-S5.16); MA-S5-002 and the
Artin parities already checked give (MA-S5.17). No nonvanishing
of the symmetric-square factor has been asserted or used.
\end{proof}

\paragraph{Source-check boundary.}
Takeda's result is the published theorem in
\emph{Mathematische Zeitschrift} 281 (2015), 103--157, not an
additional conjectural preprint assumption. The exact primary
locator read here is Theorem 7.1 and its complete terminal proof,
\texttt{arXiv:1506.04791v1}, pp.~49--50; that proof reduces the
claim to its Theorem 6.2 and the earlier integral representation.
The latter classical upstream results were not independently
re-proved here. The new deduction above and its local
normalizations were checked in full. The representation-ring
form of the same deduction is
$(\operatorname{Sym}^2\xi)\otimes\nu^{-1}=V\eta\oplus X$,
obtained by subtracting
$(\wedge^2\xi)\otimes\nu^{-1}=\mathbf1\oplus W$
from (MA-S5.8). Thus this is a symmetric-power statement in the
same Artin-motive $\lambda$-ring as MA-S5-003.

\subsection*{MA-S5-004. A stronger exceptional-prime signature}

\paragraph{Statement.}
Assume $(\mathcal R_B)$ and $k>0$. Besides MT-S508's Abelian law,
one has the ordinary weighted asymptotic
\[
  \sum_{\substack{p\le X\\p\in\Omega}}
           \frac{\log p}{\sqrt p}
            \sim \frac{k}{4}\log X,
 \tag{MA-S5.12}
\]
and, writing $\pi_\Omega(X)=\#\{p\le X:p\in\Omega\}$,
\[
                 \lim_{X\to\infty}
                  \frac{\log \pi_\Omega(X)}{\log X}=\frac12.
 \tag{MA-S5.13}
\]
This is an exponent statement for the unweighted counting
function, not an asymptotic $\pi_\Omega(X)\sim c\sqrt X/\log X$.
In particular any upper bound $\pi_\Omega(X)=O(X^{1/2-\delta})$
with fixed $\delta>0$ would close the comparison.

\begin{proof}
The proof of MT-S508 gives absolute convergence of the Euler
logarithm in $\Re s>1/2$ and
\[
 \log D_\Omega(s)=4\sum_{p\in\Omega}p^{-s}+E(s),
\]
where $E$ is holomorphic in $\Re s>1/3$; its proof is the
locally uniform comparison with $\sum_p p^{-3\Re s}$. Since
$D_\Omega$ has a pole of order $k$ at $1/2$, logarithmic
differentiation along $s=1/2+\epsilon$ gives
\[
 \sum_{p\in\Omega}\frac{\log p}{p^{1/2+\epsilon}}
                  =\frac{k}{4\epsilon}+O(1).
 \tag{MA-S5.14}
\]
We give the positive-measure Tauberian step explicitly, rather
than use an unstated strengthening of the Abelian law.
Let $\nu$ be the positive measure on $[0,\infty)$ with mass
$(\log p)/\sqrt p$ at $\log p$ for $p\in\Omega$, and put
$a=k/4$. Its Laplace transform $F(t)$ satisfies
$F(t)\sim a/t$ as $t\downarrow0$. Define positive measures
$\nu_T$ by $\nu_T(J)=T^{-1}\nu(TJ)$. Then for every $v>0$,
\[
       \int_0^\infty e^{-v x}\,d\nu_T(x)
            =T^{-1}F(v/T)\longrightarrow a/v.
\]
The finite measures $d\lambda_T(x)=e^{-x}d\nu_T(x)$ have
total mass tending to $a$. Their tails satisfy, for $R>0$,
\[
 \lambda_T([R,\infty))
   \le e^{-R/2}\int e^{-x/2}\,d\nu_T(x)
   =e^{-R/2}T^{-1}F(1/(2T)),
\]
which is uniformly small for large $T$ as $R\to\infty$.
Thus they are tight. Every weakly convergent subsequence has
Laplace transforms $a/(v+1)$, $v\ge0$, and hence equals
$a e^{-x}dx$. For completeness, uniqueness follows by pushing
the measures forward under $x\mapsto e^{-x}$ to $[0,1]$:
equality of the Laplace transforms at nonnegative integers is
equality of all polynomial moments, and polynomials are dense
in the continuous functions on $[0,1]$. Therefore all of
$\lambda_T$ converge weakly to $a e^{-x}dx$.
Multiplying by $e^x$ on the compact interval $[0,1]$ and
approximating its indicator from above and below by continuous
functions gives
\[
           T^{-1}\nu([0,T])=\nu_T([0,1])\longrightarrow a.
\]
There are no atoms at the endpoints in the limiting measure,
so this last approximation introduces no boundary term. With
$T=\log X$, it proves (MA-S5.12).

For each fixed $0<u<1$, subtraction of (MA-S5.12) at $X$ and
$X^u$ gives
\[
 \sum_{\substack{X^u<p\le X\\p\in\Omega}}
       \frac{\log p}{\sqrt p}
              \sim a(1-u)\log X.
\]
For large $X$, the decreasing function $(\log t)/\sqrt t$
is at most $u(\log X)X^{-u/2}$ on this interval. Consequently
\[
 \pi_\Omega(X)\ge
     \left(\frac{a(1-u)}u+o(1)\right)X^{u/2}.
\]
This proves that the lower limit in (MA-S5.13) is at least
$u/2$, and then at least $1/2$ by letting $u\uparrow1$.
Conversely, absolute convergence of
$\sum_{p\in\Omega}p^{-\sigma}$ for every $\sigma>1/2$
gives $\pi_\Omega(X)\le X^\sigma\sum_{p\in\Omega}p^{-\sigma}$.
The upper limit is at most $\sigma$, hence at most $1/2$.
This proves the limit and the final criterion.
\end{proof}

\subsection*{MA-S5-005. The checked central constraints still permit positive excess}

\paragraph{Statement and its scope.}
The parity, factorization, complete monomial-cone inequalities,
and new Rankin--Selberg and symmetric-square central inequalities
do not, as a system
of constraints on integer orders, force $k=0$. This statement
does not construct a number field and does not establish that
$k>0$ occurs arithmetically.

\begin{proof}
For the coordinates of MT-S506 take
\[
 (a,b,c,d,x,y,e)=(0,0,0,2,2,0,0),
 \qquad n=0,\quad n_\eta=2,\quad\beta=2,\quad f=0.
 \tag{MA-S5.15}
\]
Then $k=\beta-b-n=2=x-n=n_\eta-f$. Every coordinate that is
required to be nonnegative is nonnegative and even. The seven
nontrivial generators in MT-S507 are
\[
 (x+y,b+x,a+y,c+x+e,d+y+e,a+c+x,b+d+y)
                       =(2,2,0,2,2,2,2).
\]
Hence all $43$ monomial tests are nonnegative. The new formulas
give $r_4=n+d+e-k=0$ and $r_5=r_4+c+f=0$. All relevant
factorization orders agree, and the symmetric-square order is
$q=d+e-k=0$. This verifies the claimed
compatibility in integer arithmetic.

It can also be realized by even meromorphic germs, solely to
check the local functional-equation and multiplicative relations.
With $t=s-1/2$, set $L_W=t^2$, $L_{V\eta}=t^2$,
$A_\eta=t^2$, $B=t^2$, and set $A,C,L_V,L_X,L_{W\eta}$ and
the zeta-base germ equal to $1$. Put $D=t^{-2}$ and both
Rankin--Selberg germs equal to $1$. The two tensor-product
identities then have the required orders and every germ is even.
These germs have no prescribed arithmetic Euler factors. They
are a witness to the insufficiency of the listed central
constraints, not a counterexample to the Artin conjecture or to
the requested arithmetic equality.
\end{proof}

\subsection*{MA-S5-007. The additional input needed by the multiplicity-one route}

\paragraph{A sufficient extra hypothesis, not an input used above.}
Suppose that there exists an isobaric sum $\Theta$ of unitary cuspidal
automorphic representations of $\mathrm{GL}_{10}/\mathbb Q$ whose
unramified local factors agree, outside a finite set, with
$L(s,\Pi,\operatorname{Sym}^2\otimes\nu^{-1})$.
Then $\Omega$ is finite and $k=0$. This is a conditional reduction;
Takeda's holomorphy theorem does not assert this automorphic
realization, and no such realization is supplied in this run.

\begin{proof}
Let $\Sigma=\pi_{V\eta}\boxplus\pi_X$, using the two cuspidal
transfers already supplied by Calegari. The identity at the end of
MA-S5-006 gives matching local factors for $\Theta$ and $\Sigma$
outside $\Omega$ and a finite set $S$. Their Satake parameters at
these good primes have modulus one. Write their cuspidal
decompositions with respect to their combined set of distinct
unitary cuspidal constituents $\tau_i$ as multiplicities
$m_{1,i},m_{2,i}\ge0$. Consider the quotient of partial functions
\[
 Q(s)=\frac{L^S(s,\Theta\times\Theta^\vee)
                 L^S(s,\Sigma\times\Sigma^\vee)}
                {L^S(s,\Theta\times\Sigma^\vee)
                 L^S(s,\Sigma\times\Theta^\vee)}.
\]
At a good prime its Euler logarithm is
$\sum_{j\ge1}|\operatorname{tr}T_{\Theta,p}^j-
\operatorname{tr}T_{\Sigma,p}^j|^2p^{-js}/j$.
It is zero outside $\Omega$. On $\Omega$ the traces are bounded
in absolute value by $10$, so $\sum_{p\in\Omega}p^{-1}<\infty$
from MA-S5-001 proves that $\log Q(\sigma)$ has a finite limit
as $\sigma\downarrow1$. Thus $Q$ has neither a zero nor a pole
at $1$.

For each pair of cuspidal constituents the partial
Rankin--Selberg function is nonzero at $1$ and has a simple pole
there precisely when the two constituents agree. Here the
simple-pole assertion is Cogdell (2004), Theorems 9.1--9.2 and the
partial-factor check on printed pp.~74--75; the nonvanishing is
Shahidi (1980), the theorem on p.~462. Multiplicativity therefore
gives
\[
 -\operatorname{ord}_{s=1}Q(s)
              =\sum_i(m_{1,i}-m_{2,i})^2.
\]
The left side is zero, so all multiplicities agree and
$\Theta=\Sigma$ as isobaric representations. At an exceptional
good prime their proposed local factors would have quotient
$((1+p^{-s})/(1-p^{-s}))^2$, which is not identically one.
Consequently all exceptional primes lie in the finite omitted
set. Calegari's Lemma 4.10 gives $A=L_W$, and (MA-S5.1)
then gives $\beta=b+n$, so $k=0$.
\end{proof}

\paragraph{Why a recent analytic multiplicity-one theorem does not provide it.}
Let $u_p,v_p$ be the prime coefficients of the twisted
symmetric-square function and $L_{V\eta}L_X$, respectively.
The exact local comparison gives $u_p-v_p=4$ on $\Omega$ and
$0$ outside it, apart from finitely many primes. Farmer, Pitale,
Ryan, and Schmidt (2026), Theorem 2.3, requires the distance
condition (2.20) in their revised primary preprint, pp.~11--12:
in particular, the series
\[
       \sum_p\frac{p|u_p-v_p|^2\log p}{p^\sigma}
 \tag{MA-S5.18}
\]
must converge for every $\sigma>1$. If $(\mathcal R_B)$ and
$k>0$ hold, MA-S5-004 gives
$\sum_{p\in\Omega,p\le X}(\log p)/\sqrt p
\sim(k/4)\log X$. Hence (MA-S5.18) diverges already at
$\sigma=3/2$, and by positivity also at every
$1<\sigma\le3/2$. Thus the theorem's hypothesis actually fails
under that hypothetical branch. Its full proof is not needed for this
hypothesis check and is not claimed independently audited.
This excludes that particular proposed application; it proves
neither the impossibility of a different analytic proof nor the
arithmetic existence of positive excess.

\paragraph{The successive reductions and the completed central argument.}
The Fourier argument in MA-FQ-001--004 proves $k=0$ under
$(\mathcal R_B)$. Before that argument, the checked necessary
conditions for a positive excess were (MA-S5.17), together with
(MA-S5.12)--(MA-S5.13).
The root-number route eliminates odd excess, the factorization
route isolates $k$ exactly, and the Rankin--Selberg route forces
additional central zeros; the symmetric-square route improves
the gate further to (MA-S5.17). Ordinary strong multiplicity one
cannot start from density-one comparison. The finite monomial
cone cannot supply the missing noncoincidence or nonvanishing:
(MA-S5.15) satisfies it even after the new inequalities.
Treating the meromorphic degree-zero quotient as a member of
the ordinary Selberg class would also assume a missing
holomorphy/finitely-many-poles condition. Finally, producing a
global finite-image Galois representation from the bounded local
parameters of the nonregular automorphic $\varpi$ is a separate
unproved realization step in this argument; the local parameters
alone do not furnish it. Nonvanishing of $L_{V\eta}L_X$, an
exceptional-prime estimate stronger than exponent $1/2$, and
the extra functoriality in MA-S5-007 were possible sufficient
inputs at that stage; none is needed in the completed proof.
That proof uses the exact equality of the gamma kernels,
finite-total-variation removal of the bad Euler terms, atomic
positive definiteness of the integer divisor, and the coset-ring
theorem. It excludes an infinite exceptional-prime set without
proving any of those additional central-value or functoriality
hypotheses. The ordinary real-axis hypothesis $(\mathcal R_B)$
remains explicit.

\paragraph{Grothendieck placement.}
Let $S$ contain the ramified primes and work over
$U=\operatorname{Spec}\mathbb Z[1/S]$. The $S_5$-cover gives
finite-monodromy lisse sheaves and their rational Artin motives.
SGA~1, Expos\'e~V, \S7 identifies finite \'etale covers with finite
continuous fundamental-group sets. Deligne (1979),
\emph{Valeurs de fonctions $L$ et p\'eriodes d'int\'egrales},
Definition 6.1 and \S\S6.2--6.3, identifies their Karoubian
rational linearization with Artin motives, whose realizations
have Hodge type $(0,0)$ and weight zero; equation (0.5.3) identifies
restriction of scalars with induction. Thus (MA-S5.1) is the
finite-pushforward decomposition $\operatorname{Ind}_{F_{20}}^{S_5}
\eta_F=\eta\oplus W$, while (MA-S5.8)--(MA-S5.9) are tensor
identities in the representation $K_0$ ring. Their exterior and
symmetric powers belong to the representation $\lambda$-ring
formalism, as placed in GRH\_NIGHT MT-ART11 and AG-LAMBDA
(SGA~6, Expos\'e~V, \S\S2,7, especially (7.1.1)).

More explicitly, MA-S5-001 concerns the induced Artin-motive
decomposition and the effectivity of its analytic divisor;
MA-S5-002 concerns the parity of that weight-zero motive's
functional equation together with the automorphic comparison.
MA-S5-003 concerns tensor products and the descended
endomorphism motive, and MA-S5-006 concerns its symmetric-power
summand. MA-S5-004 concerns the Frobenius places at which the
automorphic Satake data fail to equal the motive's Euler data;
its weighted counting measure is analytic and is not asserted
to be the dimension function of a cohomology group.
MA-S5-005 is only a compatible assignment of analytic orders
to the representation-ring identities, with no realization by
an Artin motive or finite \'etale cover. MA-S5-007 separates the
known symmetric-power Euler data of this Artin motive from the
extra automorphic realization needed by its conditional route.
These distinctions use
the same exact SGA~1 and Deligne locators just given.

The parity assertion is about the functional equation of a
weight-zero Artin motive and its known automorphic comparison;
the Rankin--Selberg bound is an inequality between analytic
central multiplicities. None of these constructions identifies
those multiplicities with dimensions of a positive faithful
arithmetic cohomology. They refine the Artin-divisor effectivity
gate of GRH\_NIGHT Q-ARTIN-POS and MT-S508. In the inherited
cote~1 trace/determinant programme (ALPOGE\_BRIDGES BR-TR-01,
BR-TR-02; sheets 79--80 and 111--112), they constrain a candidate
spectral divisor but do not supply the absent positivity theorem.
The latter archive locators are programme cross-references from
the supplied runs, not a claim that those scans were independently
re-read in this S5 audit.

\paragraph{Bibliographic locators used in this section.}
Calegari, F. (2013). \emph{The Artin conjecture for some
$S_5$-extensions}. \emph{Mathematische Annalen, 356}, 191--207.
\url{https://doi.org/10.1007/s00208-012-0839-4}.
Author-hosted publisher PDF:
\url{https://www.math.uchicago.edu/~fcale/papers/S5.pdf}.
The character table is on PDF p.~4; the projective $\xi$ table
is on p.~6; the exceptional $4A$ Satake table is on p.~11.

Arthur, J. (2013). \emph{The endoscopic classification of
representations: Orthogonal and symplectic groups} (Colloquium
Publications, Vol.~61). American Mathematical Society.
Theorems 1.4.1 and 1.5.3. The author's equivalent root-number
statement is Theorem 3(b), p.~8, at
\url{https://www.claymath.org/library/cw/arthur/pdf/arthur-endoscopic-tifr.pdf}.
The older overview's provisional account of the then unfinished
twisted trace formula is not used to replace the published theorem.

Cogdell, J. W. (2004). \emph{Lectures on $L$-functions, converse
theorems, and functoriality for $\mathrm{GL}_n$}. In
\emph{Lectures on automorphic $L$-functions} (Fields Institute
Monographs, Vol.~20). American Mathematical Society.
Author-hosted notes, Theorems 9.2 and 9.3, printed pp.~74--75:
\url{https://people.math.osu.edu/cogdell.1/fields-www.pdf}.

Takeda, S. (2015). On a certain metaplectic Eisenstein series and
the twisted symmetric square $L$-function.
\emph{Mathematische Zeitschrift, 281}, 103--157.
\url{https://doi.org/10.1007/s00209-015-1476-x}.
Theorem 7.1, author preprint pp.~49--50:
\url{https://arxiv.org/abs/1506.04791}.

Shahidi, F. (1980). On nonvanishing of $L$-functions.
\emph{Bulletin of the American Mathematical Society, 2}(3),
462--464. \url{https://doi.org/10.1090/S0273-0979-1980-14769-2}.
The theorem on p.~462 and its full printed outline on
pp.~462--463 were read; the outline's upstream Eisenstein-series
and Whittaker results remain classical inputs.

Farmer, D. W., Pitale, A., Ryan, N. C., \& Schmidt, R. (2026).
Multiplicity one for analytic $L$-functions and applications.
\emph{Mathematical Proceedings of the Cambridge Philosophical
Society, 181}(1), 885--911.
\url{https://doi.org/10.1017/S0305004126102011}.
The exact hypothesis checked is (2.20) and Theorem 2.3,
\texttt{arXiv:1305.3972v2}, pp.~11--12:
\url{https://arxiv.org/abs/1305.3972}.

Deligne, P. (1979). Valeurs de fonctions $L$ et p\'eriodes
d'int\'egrales. In \emph{Automorphic forms, representations and
$L$-functions} (Proceedings of Symposia in Pure Mathematics,
Vol.~33, Part~2, pp.~313--346). American Mathematical Society.
\S0.5.3 and \S\S6.1--6.3; Ghate--Beineke translation, PDF
pp.~4 and 20:
\url{https://mathweb.tifr.res.in/~eghate/Deligne.pdf}.

\section{A repaired base-change step in the symplectic construction}
\label{sec:s5-symplectic-source-repair}

\paragraph{Literal source finding.}
Calegari's Lemma~4.2, author-hosted publisher PDF pp.~9--10,
contains on p.~9 the unnumbered equality
$\wedge^2\Pi_E=\mu\boxplus\mu$ in its hypothetical
cuspidal-induction branch, followed by a multiplicity-two
argument.  The same passage is in the supplied
\texttt{S5ArtinWeb.tex}, lines~563--572, under
\texttt{lemma:satake}.  The published page was checked
both as text and as a rendered image in the present audit.
The correct quadratic base change is
$\mu\boxplus\mu^\sigma$, where $\sigma$ is the nontrivial
element of $\operatorname{Gal}(E/\mathbb Q)$.
The identical-summand equality and its stated local-multiplicity
justification are not valid.  The intended conclusion has the
following complete repair using the same classical inputs.

\paragraph{MA-S5-008 (corrected summands and a constituent-degree repair).}
In Calegari's specified family, let $\pi_E$ be the established
cuspidal representation of $\operatorname{GL}_2(\mathbb A_E)$,
put $\omega=\omega_{\pi_E}$, and retain the established cuspidal
descent
\[
 \operatorname{BC}_{E/\mathbb Q}\Pi
        =\operatorname{Sym}^3\pi_E.
 \tag{MA-S5.40}
\]
Its exterior-square transfer $\Sigma=\wedge^2\Pi$ is
noncuspidal.  In fact it has the exact isobaric type
\[
 \Sigma=\chi\boxplus\Sigma_5,
 \qquad \Sigma_5\text{ cuspidal of degree five},
 \tag{MA-S5.41}
\]
with
\[
 \operatorname{BC}\chi=\omega^3,
 \qquad
 \operatorname{BC}\Sigma_5=\operatorname{Sym}^4\pi_E\otimes\omega.
 \tag{MA-S5.42}
\]
The representation $\Pi$ is of symplectic type, with similitude
character equal to the character $\nu$ fixed immediately before
the source's Lemma~4.2 by $\Pi^\vee\simeq\Pi\otimes\nu^{-1}$.
Consequently its unramified Satake parameters have the asserted
form
\[
                         \{a,b,\nu(p)/a,\nu(p)/b\}.
 \tag{MA-S5.43}
\]
This repair uses no comparison with $W$ at the exceptional
primes and no zero-location or real nonvanishing hypothesis.
In particular it does not use the later conclusion
$\nu=\psi^{-1}$ of Calegari's Lemma~4.6.

\begin{proof}
\emph{The actual base change has constituent degrees one and five.}
For a two-dimensional vector space representation $\tau$,
the polynomial identity is
\[
 \wedge^2\operatorname{Sym}^3\tau
   =(\operatorname{Sym}^4\tau\otimes\det\tau)
                                      \oplus(\det\tau)^3.
 \tag{MA-S5.44}
\]
For example, if $a,b$ are the two eigenvalues, the six
exterior products of $a^3,a^2b,ab^2,b^3$ are
\[
                 a^5b,\ a^4b^2,\ a^3b^3,\ a^3b^3,
                                  \ a^2b^4,\ ab^5.
 \tag{MA-S5.45}
\]
Five are the weights of $(ab)\operatorname{Sym}^4\tau$;
the remaining one is $(ab)^3$.  Equality on the diagonal
torus proves (MA-S5.44) as a semisimple polynomial
representation identity.

Kim's exterior-square transfer and the already used
symmetric-fourth transfer exist as isobaric automorphic
representations.  In this icosahedral case
$\operatorname{Sym}^4\pi_E$ is cuspidal, as in the opening
of Calegari's \S4.  Cyclic base change and exterior powers
commute at every common unramified place: both apply the
same exterior-power map to the local parameter restricted
to $E$.  Thus isobaric strong multiplicity one and
(MA-S5.40)--(MA-S5.45) give
\[
 \operatorname{BC}\Sigma
    =\omega^3\boxplus
                       (\operatorname{Sym}^4\pi_E\otimes\omega).
 \tag{MA-S5.46}
\]
This is an isobaric decomposition into a character and a
cuspidal degree-five representation.  It is not a claim
about double eigenvalue multiplicities at each prime.

Suppose first that $\Sigma$ were cuspidal of degree six.
Its noncuspidal quadratic base change forces
\[
 \Sigma=\operatorname{AI}_{E/\mathbb Q}\mu,
 \qquad
 \operatorname{BC}\Sigma=\mu\boxplus\mu^\sigma,
 \tag{MA-S5.47}
\]
where $\mu$ is cuspidal of degree three.  The second
identity is the restriction-of-induction formula at
unramified parameters, promoted by isobaric multiplicity
one; the two cosets give $\mu$ and its Galois conjugate.
Both constituents have degree three.  They cannot equal
the degree-one and degree-five cuspidal constituents in
(MA-S5.46), by uniqueness of isobaric decomposition.
Hence $\Sigma$ is noncuspidal.

Moreover $\mu\not\simeq\mu^\sigma$ in the hypothetical
cuspidal induction in (MA-S5.47): a Galois-invariant
cuspidal representation descends, and its induction is
the sum of the two quadratic twists of a descent.
Thus in precisely that branch the source cannot replace
the conjugate summand by an identical one.  This explains
the literal error independently of the repaired
constituent-degree contradiction.

To prove the stronger assertion (MA-S5.41), decompose
$\Sigma$ into its cuspidal isobaric constituents over
$\mathbb Q$.  Under quadratic base change a cuspidal
constituent of degree $r$ either remains cuspidal of
degree $r$, or splits into two cuspidal constituents of
degree $r/2$.  In order to contribute the degree-five
constituent of (MA-S5.46), a constituent over $\mathbb Q$
must have degree five or ten.  The total degree is six,
so it has degree five.  The remaining degree is one.
This proves (MA-S5.41), and uniqueness of the two
constituents after base change gives (MA-S5.42).

\emph{The classification now selects symplectic type explicitly.}
Apply the same published cuspidality classification that
the source invokes: Asgari--Raghuram, Theorem~1.1(ii).
We check its alternatives using its proved exterior-square
decompositions, rather than assuming that every noncuspidal
exterior square is symplectic.

In its tensor-product case (a), Proposition~3.1(a) gives
two isobaric representations of degree three.  None of
their cuspidal constituents has degree five.
In its dihedral Asai case (b), Proposition~3.2 gives
the automorphic induction of a twisted symmetric square
of a dihedral degree-two representation.  That symmetric
square is an isobaric sum of degrees one and two, possibly
further split.  Quadratic automorphic induction therefore
has constituents of degree at most four.
Finally, in its quadratic-induction case (d), Proposition~3.5
gives an Asai representation of degree four together with
an induced character of degree two.  Again every cuspidal
constituent has degree at most four.
All three alternatives contradict the cuspidal degree-five
constituent in (MA-S5.41).

The remaining alternative (c) says that $\Pi$ is a
transfer from a globally generic cuspidal representation
of $\operatorname{GSp}_4$.  If its similitude character
is $\widetilde\nu$, Proposition~3.3 gives
$\wedge^2\Pi=\widetilde\nu\boxplus\widetilde r_5$.
The unique degree-one constituent of (MA-S5.41) is
therefore $\widetilde\nu=\chi$; the remaining constituent
is $\Sigma_5$.  In particular
\[
                 \Pi^\vee\simeq\Pi\otimes\chi^{-1},
                 \qquad \omega_\Pi=\chi^2.
 \tag{MA-S5.48}
\]

\emph{Identification of the previously fixed similitude.}
The same constituent argument excludes quadratic
automorphic induction of $\Pi$ from \emph{any} quadratic
field, since Proposition~3.5 would always give only
constituents of degree at most four in $\wedge^2\Pi$.
Suppose $\Pi\simeq\Pi\otimes\xi$ for a nontrivial
Hecke character $\xi$.  Comparing central characters
gives $\xi^4=1$.  Either $\xi$ is quadratic, or
$\xi^2$ is a nontrivial quadratic self-twist.
The cyclic self-twist criterion, also proved as
Asgari--Raghuram Lemma~5.2, would then make $\Pi$ a
quadratic automorphic induction, a contradiction.
Thus $\Pi$ has no nontrivial character self-twist.

Compare (MA-S5.48) with the already established
$\Pi^\vee\simeq\Pi\otimes\nu^{-1}$.
It gives $\Pi\simeq\Pi\otimes\chi\nu^{-1}$, so
the absence of self-twists proves $\chi=\nu$.
This supplies the last step in the original lemma
without its unexplained quadratic-induction shortcut.

For clarity, the asserted parameter form holds at every
unramified place of $\Pi$, not only at places used in a
weak transfer.  Local self-duality with multiplier $\nu$
first makes $\nu$ unramified there and makes the Satake
multiset invariant under $z\mapsto\nu(p)/z$.
Its product is $\nu(p)^2$ by (MA-S5.48).
Nonfixed orbits already form pairs with product $\nu(p)$.
The remaining values are $\sqrt{\nu(p)}$ and
$-\sqrt{\nu(p)}$.  The product condition makes the
multiplicity of the negative value even; since the
remaining cardinality is even, the positive multiplicity
is even too.  Pair equal fixed values.  This proves
(MA-S5.43) in all cases, including repeated eigenvalues.
\end{proof}

\paragraph{Role, status, and exact source scope.}
The corrected conclusion is the symplectic parameter input
used in MA-S5-003/006 and MA-RB-005.  It restores that
input while leaving the exceptional-prime comparison and
the shared-real-zero gate unchanged.  No claim is made
that the earlier source readings detected this error:
the literal summand mistake was identified in the new
read for MA-RB-005 and audited separately here.
The outcome is a refuted printed intermediate equality
with a proved refinement, not a refuted construction theorem.

The exact primary Calegari locator is Lemma~4.2, publisher
PDF pp.~9--10, especially the unnumbered base-change
equality at the foot of p.~9.  The required symmetric-cube
descent is on p.~9 before Lemma~4.1, and symmetric-fourth
cuspidality is stated at the opening of \S4, p.~8.
The supplied source's plethysm also occurs in the proof
of Lemma~2.5; (MA-S5.44)--(MA-S5.45) reprove it here.
The author-hosted publisher version is
\url{https://math.uchicago.edu/~fcale/papers/S5.pdf}.
The source's cyclic base-change input remains
Arthur--Clozel, Theorem~4.2, printed p.~202;
its full trace-formula proof is not newly audited.

For the completion of the repair, the primary
Asgari--Raghuram author version
\url{https://math.okstate.edu/people/asgari/Files/cuspidality.pdf}
was freshly read at Theorem~1.1, p.~2; isobaric
multiplicity one, Theorem~2.1, p.~4; the complete proofs
of Propositions~3.1--3.3 and~3.5, pp.~8--10; and the
complete self-twist proof, Lemma~5.2, p.~11.
The main cuspidality classification itself, Kim's
transfer, symmetric-power cuspidality, and cyclic base
change are the same carried published inputs as in the
source construction.  They are not presented as new
unconditional functoriality proved by this audit.

Calegari, F. (2013). The Artin conjecture for some
$S_5$-extensions. \emph{Mathematische Annalen, 356},
191--207. \url{https://doi.org/10.1007/s00208-012-0839-4}.

Asgari, M., \& Raghuram, A. (2011).
A cuspidality criterion for the exterior square transfer
of cusp forms on $\mathrm{GL}(4)$.
In \emph{On certain $L$-functions} (Clay Mathematics
Proceedings, Vol.~13, pp.~33--53).
American Mathematical Society.

Arthur, J., \& Clozel, L. (1989).
\emph{Simple algebras, base change, and the advanced
theory of the trace formula}.
Annals of Mathematics Studies, Vol.~120.
Princeton University Press.

\section{The degree-zero distribution route and its exact remaining gate}
\label{sec:booker-route}

This section examines a further route to the $S_5$ comparison, using
Booker's published distributional converse theorem. It yields a new
necessary signature for the positive-excess branch, but does not
eliminate that branch.

\paragraph{The primary input and its normalization.}
For an $L$-datum $\mathcal F=(f,K,m)$, Booker's normalization is
\[
 L_{\mathcal F}(s)
 =\exp\left(\sum_{n\ge2}\frac{f(n)}{\log n}
                              n^{1/2-s}\right)
 =\sum_{n\ge1}a_{\mathcal F}(n)n^{-s}.
 \tag{MA-LD.1}
\]
Its degree is $2\lim_{x\downarrow0}xK(x)$. Positivity means that
$m(z)<0$ at only finitely many points. It is a condition on the
signed divisor, not on $f(n)$. The exact locators are Definition~1.3,
author p.~3; Examples~1.4(2)--(3), p.~4; Theorem~1.6, p.~5;
and Theorem~1.7, p.~6, in \texttt{arXiv:1308.3067v2}.
The multiplicity-one criteria used below are
\[
 \sum_{n\ge2}\frac{|f(n)|}{\log n}<\infty,
 \qquad
 \sum_{n\ge1}\frac{|a_{\mathcal F}(n)|}{\sqrt n}<\infty,
 \qquad
 \sum_{|\operatorname{Re}z|\le T}|m(z)|=o(T).
 \tag{MA-LD.2}
\]
Each is equivalent to $\mathcal F=0$. In degree zero, the additional
hypothesis of positivity also forces $\mathcal F=0$.
In particular, the second expression in (MA-LD.2) uses the
Dirichlet coefficients $a_{\mathcal F}$, not the logarithmic
coefficients $f$.

\paragraph{MA-LD-001 (an infinite signed discrepancy is necessary).}
Retain all of the family and comparison hypotheses in MA-S5-001,
including $(\mathcal R_B)$. Form the difference of the actual
automorphic and Artin data
\[
 \mathcal F_R=\mathcal F_{\varpi}-\mathcal F_W\in\mathcal L_0,
 \qquad L_{\mathcal F_R}(s)=\frac{A(s)}{L_W(s)}.
 \tag{MA-LD.3}
\]
Let $\mathscr R$ be its completed meromorphic function. It differs
from the quotient $R$ in (MA-S5.3) only by the nonvanishing
conductor exponential and a nonzero normalization constant.
Put
\[
 m_R(z)=\operatorname{ord}_{s=1/2+iz}\mathscr R(s),\qquad
 N_R(T)=\sum_{|\operatorname{Re}z|\le T}|m_R(z)|.
 \tag{MA-LD.4}
\]
If $k>0$, then $\mathscr R$ has infinitely many poles and infinitely
many zeros, all on $\operatorname{Re}s=1/2$, and
\[
                       N_R(T)\ne o(T).
 \tag{MA-LD.5}
\]
The critical-line assertion here concerns the signed quotient
divisor alone and uses no zero-location hypothesis beyond
$(\mathcal R_B)$. It is not a statement about the zeros of the
individual factors.

Consequently, under $(\mathcal R_B)$, each of the following is
equivalent to $k=0$: the quotient has only finitely many poles;
it has only finitely many zeros; or $N_R(T)=o(T)$.

\begin{proof}
Booker's Examples~1.4(2)--(3) attach actual $L$-data to the
cuspidal automorphic representation $\varpi$ and the Artin
representation $W$, respectively. Artin holomorphy is not
needed for the second construction. The class of $L$-data is
an additive group. Both representations have degree five,
so their difference has degree zero. The completed divisor is
the difference of their completed divisors, which proves the
normalizations in (MA-LD.3)--(MA-LD.4). In particular
$m_R(0)=-k$.

For a good prime, logarithmic differentiation of (MA-S5.3)
gives the exact coefficients
\[
 f_R(p^j)=
 \begin{cases}
 4(\log p)p^{-j/2},&p\in\Omega,\ j\text{ odd},\\
 0,&\text{otherwise}.
 \end{cases}
 \tag{MA-LD.6}
\]
The remaining coefficients are supported on powers of finitely
many bad primes. MA-S5-004 shows that
$\pi_\Omega(X)=X^{1/2+o(1)}$. Hence the Euler logarithm in
(MA-LD.1) converges absolutely and locally uniformly on
$\operatorname{Re}s>1/2$: its good first-prime terms are
$4\sum_{p\in\Omega}p^{-s}$, and its higher odd-power terms
are bounded on compact subsets by a constant times
$\sum_p p^{-3\operatorname{Re}s}$.
For each bad prime, axiom~(A1) of Definition~1.3 gives, for
every fixed $a>0$,
$f_R(p^j)(j\log p)^a=O_a(1)$; this also gives the required
local convergence. Exponentiation therefore makes
$L_{\mathcal F_R}$ holomorphic and nonvanishing on that
half-plane.

The divisor $m_R$ is integer-valued everywhere, since it is the
difference of two actual meromorphic divisors. Consequently the
continuation domain in Proposition~2.1(iii) is all of $\mathbb C$;
none of the branch cuts allowed for general $L$-data occurs here.
Proposition~2.1(i), author p.~9, supplies a canonical gamma
factor whose logarithm is holomorphic on
$\operatorname{Re}s>1/2$ and extends continuously to the
boundary. Thus $\mathscr R$ is holomorphic and nonvanishing
on the open right half-plane. Its functional equation,
Proposition~2.1(iv), reflects this conclusion to
$\operatorname{Re}s<1/2$. Its complete signed divisor is
therefore supported on the critical line. Thus their completed
orders cancel off that line. The argument supplies no equality
or inequality between their orders on the critical line.

If $\mathscr R$ had finitely many poles, then
$\mathcal F_R$ would be positive in Booker's precise sense.
The degree-zero case of Theorem~1.7, followed by
Theorem~1.6, would give $\mathcal F_R=0$, contradicting
$m_R(0)=-k<0$. If it had finitely many zeros, apply the same
argument to $-\mathcal F_R$. This proves the two infinitude
assertions. Theorem~1.6(v) gives (MA-LD.5).
Conversely, if $k=0$, MA-S5-001 gives the all-place comparison,
so the difference datum vanishes and all three finiteness
conditions hold.\end{proof}

\paragraph{MA-LD-002 (the summability criteria fail at the exact threshold).}
Under the positive-excess hypotheses of MA-LD-001,
\[
 \sum_{\substack{p\le X\\p\in\Omega}}p^{-1/2}
                  \sim\frac{k}{4}\log\log X,
 \qquad
 \sum_{2\le n\le X}\frac{|f_R(n)|}{\log n}
                  \sim k\log\log X.
 \tag{MA-LD.7}
\]
Also
\[
 \sum_{n\ge1}\frac{|a_{\mathcal F_R}(n)|}{\sqrt n}=\infty,
 \qquad
 \sum_{n\ge2}\frac{|f_R(n)|}{\sqrt n}<\infty.
 \tag{MA-LD.8}
\]
Thus replacing $a_{\mathcal F_R}$ by $f_R$ in Booker's criterion
would produce a false shortcut precisely in this branch.

\begin{proof}
Let
$U(X)=\sum_{p\le X,\ p\in\Omega}(\log p)/\sqrt p$.
MA-S5-004 proves $U(X)\sim(k/4)\log X$. Partial summation gives
\[
 \sum_{\substack{p\le X\\p\in\Omega}}p^{-1/2}
 =\frac{U(X)}{\log X}
       +\int_2^X\frac{U(t)}{t(\log t)^2}\,dt
 \sim\frac{k}{4}\log\log X.
 \tag{MA-LD.9}
\]
The first asymptotic follows by bounding the remainder
$U(t)-(k/4)\log t$ by an arbitrary fixed multiple of
$\log t$ beyond a sufficiently large point. The contribution
of odd powers $j\ge3$ to the second sum in (MA-LD.7) is
bounded by
$4\sum_p\sum_{j\ge3}p^{-j/2}/j<\infty$.
Each bad-prime contribution is finite by (A1), for example
using $|f_R(p^j)|\ll(j\log p)^{-2}$.
The prime term in (MA-LD.6) therefore proves that second
asymptotic.

At each $p\in\Omega$, the Dirichlet coefficient of
$L_{\mathcal F_R}$ is $a_{\mathcal F_R}(p)=4$.
The first assertion of (MA-LD.8) follows already from the
prime subsequence and (MA-LD.7). For the last assertion,
the good-prime contribution is
\[
 4\sum_{p\in\Omega}\sum_{j\text{ odd}}\frac{\log p}{p^j}.
\]
Its higher-power part converges absolutely. For its first
part, $\pi_\Omega(t)\ll t^{3/4}$ eventually, and integration
by parts gives
$\sum_{p\in\Omega}(\log p)/p<\infty$, since
$\int_2^\infty (\log t)t^{-5/4}dt<\infty$.
The finitely many bad-prime geometric tails also converge,
again by (A1).\end{proof}

\paragraph{MA-LD-003 (Euler positivity does not give datum positivity).}
There is an elementary genuine degree-zero example illustrating
the distinction. Let $\chi_5$ be the even primitive quadratic
character modulo $5$, and consider
\[
 L_*(s)=\frac{\zeta(s)}{L(s,\chi_5)},\qquad
 \mathcal F_* =\mathcal F_\zeta-\mathcal F_{\chi_5}.
 \tag{MA-LD.10}
\]
The two archimedean gamma kernels coincide, so $K_*=0$, and
\[
       f_*(n)=(1-\chi_5(n))\frac{\Lambda(n)}{\sqrt n}\ge0
       \quad(n>1).
 \tag{MA-LD.11}
\]
Nevertheless $\mathcal F_*$ is not positive: its completed
quotient has infinitely many poles and infinitely many zeros.

\begin{proof}
The data exist by the usual degree-one examples, and their
difference has degree zero. The logarithmic derivative gives
(MA-LD.11), including the ramified prime $5$, where
$\chi_5(5)=0$. Since $\chi_5(2)=-1$, its coefficient at $2$
is nonzero. If the completed quotient had finitely many
poles, Theorem~1.7 would force this nonzero datum to vanish.
Applying it to the negative datum gives the assertion about
zeros. The quotient has the usual completed functional
equation; the conductor ratio is absorbed in its exponential
factor.\end{proof}

This example is not a number field in the stipulated $S_5$
family and does not have the exceptional-prime sparsity above;
it has a pole at $s=1$. It isolates exactly the invalid general
inference from Euler positivity to positivity of the divisor.

\paragraph{Remaining gate and source-check scope.}
Even holomorphy of both individual factors would not make their
quotient have finitely many poles: a denominator zero creates a
pole whenever its order exceeds that of the numerator. The
degree-zero theorem thus supplies an equivalent cancellation
gate, not a proof of that cancellation. The input $M$ is a
zero-location assertion and supplies no such order inequality.
The ordinary Selberg-class theorem is inapplicable for the same
reason; its analytic-continuation axiom permits a pole only at
$1$, whereas the unresolved branch already has a pole at $1/2$.

The full proofs read for this check were Booker's Lemmas~2.3--2.4,
Proposition~2.1 and its proof, pp.~9--13; the complete
Theorem~1.6 proof in \S2.3, pp.~13--15; and the degree-below-one
part of Theorem~1.7, including Lemma~3.1 and its preceding
contour deformation, pp.~15--17. The exact use of positivity is
on p.~16: it permits a height above all poles before the contour
deformation. Positivity of the coefficients cannot justify that
step. This source check does not claim a new proof of all
automorphic analytic properties used to construct the example
data, nor an audit of the unused degree-one and higher-degree
classification.

Booker, A. R. (2015). $L$-functions as distributions.
\emph{Mathematische Annalen, 363}(1--2), 423--454.
\url{https://doi.org/10.1007/s00208-015-1178-z}.
Author manuscript: \url{https://arxiv.org/abs/1308.3067v2}.
This is a published input, separate from the user-specified
theorem $A$, the conditional claim $M$, and the outside
factoring preprint PF26.

\section{A common-zero signature for the remaining positive excess}
\label{sec:booker-rankin}

The signed-discrepancy conclusion in MA-LD-001 can be combined
with symmetric-square holomorphy on the whole critical line.
This gives a condition on common zeros of specified functions,
beyond the inequality at their central values.

Fix the finite set $S$ used in MA-S5-006 and put
\[
 Q_{\rm sym}(s)=L^S(s,\Pi,\operatorname{Sym}^2\otimes\nu^{-1}),
 \qquad D(s)=L_{V\eta}(s)L_X(s),\qquad s_t=\tfrac12+it.
 \tag{MA-LD.12}
\]
For real $t$ define the finite meromorphic orders
\[
 \begin{aligned}
 \beta(t)&=\operatorname{ord}_{s_t}B,&
 b(t)&=\operatorname{ord}_{s_t}C,&
 n(t)&=\operatorname{ord}_{s_t}A,\\
 d(t)&=\operatorname{ord}_{s_t}L_{V\eta},&
 e(t)&=\operatorname{ord}_{s_t}L_X,&
 q(t)&=\operatorname{ord}_{s_t}Q_{\rm sym},\\
 \kappa(t)&=\operatorname{ord}_{s_t}L_W-n(t)
             =\beta(t)-b(t)-n(t).
 \end{aligned}
 \tag{MA-LD.13}
\]
All of $\beta,b,n,d,e,q$ are nonnegative. Here the functions
$B,C$ are nonprincipal Hecke functions, $A,L_{V\eta},L_X$
are the known entire automorphic functions already used in
MA-S5-001/003, and the assertion for $Q_{\rm sym}$ is the
holomorphy part of Takeda's theorem. None of these assertions
requires a zero-location hypothesis. At $t=0$,
$\kappa(0)=k$.

\paragraph{MA-LD-004 (two unavoidable common-zero families).}
At every real $t$ one has
\[
 \begin{gathered}
 \boxed{q(t)=d(t)+e(t)-\kappa(t)},\\
 -(b(t)+n(t))\le\kappa(t)\le d(t)+e(t),
 \end{gathered}
 \tag{MA-LD.14}
\]
and therefore
\[
                 \beta(t)+q(t)=b(t)+n(t)+d(t)+e(t).
 \tag{MA-LD.15}
\]
Assume now $(\mathcal R_B)$ and $k>0$. Then both of the following
sets of heights are infinite and unbounded in absolute value:
\[
 \{t\in\mathbb R:B(s_t)=D(s_t)=0\},\qquad
 \{t\in\mathbb R:Q_{\rm sym}(s_t)=A(s_t)C(s_t)=0\}.
 \tag{MA-LD.16}
\]
This assertion is about literal coincident zeros at the same
height, with their actual orders. Define their overlap counts
with multiplicity by
\[
 \begin{aligned}
 N_{B,D}(T)&=\sum_{|t|\le T}\min\{\beta(t),d(t)+e(t)\},\\
 N_{Q,AC}(T)&=\sum_{|t|\le T}\min\{q(t),b(t)+n(t)\}.
 \end{aligned}
 \tag{MA-LD.17}
\]
The sums are over the discrete union of the relevant zero
heights. Then
\[
       N_{B,D}(T)+N_{Q,AC}(T)\ge N_R(T)\ne o(T).
 \tag{MA-LD.18}
\]
Consequently, finiteness of either common-zero set in
(MA-LD.16), or an $o(T)$ bound for the sum in (MA-LD.18),
would imply $k=0$ under $(\mathcal R_B)$.

\begin{proof}
First extend the local-factor check of MA-S5-003 from the
central point to the entire line $\operatorname{Re}s=1/2$.
Every finite correction factor there has roots that are
roots of unity in the variable $p^{-s}$, so its zeros and
poles have real part zero. The archimedean corrections are
ratios of products of $\Gamma_{\mathbb R}(s)$ and
$\Gamma_{\mathbb R}(s+1)$, because the relevant parameters
have finite image. Their poles are at nonpositive real
integers. Thus these factors, as well as conductor
exponentials, are holomorphic and nonvanishing at every
$s_t$. Deleting the finitely many additional factors in
the chosen $S$ has the same property.

Write $z(t)=\operatorname{ord}_{s_t}\zeta_{\mathbb Q}$ and
$r_4(t)=\operatorname{ord}_{s_t}
\Lambda(s,\Pi\times\Pi^\vee)$.
The meromorphic identity (MA-S5.11) and the character
decomposition (MA-S5.8) therefore give
\[
 \begin{split}
 r_4(t)
 &=z(t)+d(t)+\operatorname{ord}_{s_t}L_W+e(t)
                                       -2\kappa(t)\\
 &=z(t)+n(t)+d(t)+e(t)-\kappa(t).
 \end{split}
\]
The symmetric/exterior-square identity used in MA-S5-006,
at this same point, gives
\[
                         r_4(t)=z(t)+n(t)+q(t).
\]
The orders $z(t)$ need not vanish away from the centre,
but cancel exactly between these two formulas. This proves
the identity in (MA-LD.14). Takeda's Theorem~7.1 states
holomorphy of the indicated incomplete symmetric-square
function away from the possible points $0$ and $1$;
in particular $q(t)\ge0$ for every real $t$. This proves
the upper bound. The lower bound follows from
$\beta(t)=b(t)+n(t)+\kappa(t)\ge0$, and (MA-LD.15) is
the same pair of identities combined.

Under the positive-excess hypotheses, MA-LD-001 proves
infinitely many negative and positive orders for the
completed comparison quotient, all on the critical line.
Its order at $s_t$ is $-\kappa(t)$. Hence there are
infinitely many heights with $\kappa(t)>0$, and infinitely
many with $\kappa(t)<0$. At a height of the first kind,
\[
              \beta(t)\ge\kappa(t)>0,
              \qquad d(t)+e(t)\ge\kappa(t)>0.
\]
At a height of the second kind,
\[
              q(t)\ge-\kappa(t)>0,
              \qquad b(t)+n(t)\ge-\kappa(t)>0.
\]
These prove (MA-LD.16). The functions involved are nonzero
meromorphic functions, so each bounded segment contains
only finitely many zeros; the two sets are therefore
unbounded. The same inequalities give, at every height,
\[
 \min\{\beta(t),d(t)+e(t)\}
     +\min\{q(t),b(t)+n(t)\}\ge|\kappa(t)|.
\]
By critical-line support, $N_R(T)=\sum_{|t|\le T}|\kappa(t)|$.
Summation and (MA-LD.5) prove (MA-LD.18). The final sufficient
criteria follow by contradiction.
\end{proof}

This deduction gives no smaller numerical upper bound than
$k\le d(0)+e(0)$ at the centre. Its new requirement is an
infinite set of coincidences away from the centre whose
combined multiplicity count is not $o(T)$. No theorem excluding that
pattern is asserted. The zero-location input $M$ alone
does not bound these coincidences or their multiplicities.
In particular the argument does not apply a low-degree
converse theorem to the symmetric-square function, whose
degree is ten, and it does not assume that it is an
automorphic standard function on $\mathrm{GL}_{10}$.

\paragraph{Primary input and check scope.}
The new use of Takeda is his holomorphy on the whole line,
not merely at the central point. The exact primary source
read is Theorem~7.1 and its complete terminal proof,
\texttt{arXiv:1506.04791v1}, pp.~49--50. Its permitted $S$
is exactly the archimedean places, places dividing $2$,
and the ramified places of the representation and twist.
The finite-order twist and unitary central character
were checked in MA-S5-006. The proof invokes Takeda's
Theorem~6.2 and earlier integral representation; these
upstream representation-theoretic arguments remain
published inputs rather than newly reproved results.

Takeda, S. (2015). On a certain metaplectic Eisenstein
series and the twisted symmetric square $L$-function.
\emph{Mathematische Zeitschrift, 281}, 103--157.
\url{https://doi.org/10.1007/s00209-015-1476-x}.
\url{https://arxiv.org/abs/1506.04791v1}.

\section{The signed divisor as an integer positive-definite function}
\label{sec:s5-crystal}

The sign condition at the good primes survives the finitely many bad
Euler factors in a useful precise form. The correction for those
factors is continuous on the zero side, and therefore disappears
when individual zero atoms are extracted. The topology in the
conclusion is essential: the resulting function is positive definite
on the additive group $\mathbb R_d$ with its discrete topology.

\paragraph{MA-FQ-001 (finite Euler corrections preserve atomic positivity).}
Retain the family, the specified comparison representation $\varpi$,
and $(\mathcal R_B)$ of MA-S5-001. Let
$\mathcal F_R=(f_R,K_R,m_R)$ be the difference of the actual
automorphic and Artin data in MA-LD-001, with its actual completed
divisor. Then
\[
 K_R=0,\qquad
 \mu=\sum_{t\in\mathbb R}m_R(t)\delta_t
       \ \text{is a tempered, locally finite signed measure},
 \qquad m_R(0)=-k.
 \tag{MA-FQ.1a}
\]
All its masses are integers. Define $a(t)=-m_R(t)$, with value zero
away from the divisor. The function $a$ is positive definite on
$\mathbb R_d$, and
\[
 a(0)=k,\qquad |m_R(t)|\le k\quad(t\in\mathbb R),\qquad
 a\in B(\mathbb R_d),\qquad \|a\|_{B(\mathbb R_d)}=k.
 \tag{MA-FQ.1b}
\]
Here $B(\mathbb R_d)$ is the Fourier--Stieltjes algebra: the
transforms of finite complex measures on the compact character
group $\widehat{\mathbb R_d}$. In particular $a$ has finite integer
range and locally finite support in the usual real topology.
No positivity of the signed divisor as a measure is asserted.

\begin{proof}
\emph{The exact archimedean parameters.}
The quadratic field $E=L^{A_5}$ is totally real. The finite-image
two-dimensional representation $\varrho$ used by Calegari is odd
at each real place of $E$, so its local parameter is
$\mathbf1\oplus\operatorname{sgn}$. Calegari's construction gives
\[
 \operatorname{BC}_{E/\mathbb Q}(\varpi)
  =\operatorname{Sym}^4(\pi_E)\otimes\det(\pi_E)^{-2},
 \tag{MA-FQ.1c}
\]
where $\pi_E$ is the automorphic representation corresponding to
$\varrho$. At a real place the right side has parameter
$\mathbf1^{\oplus3}\oplus\operatorname{sgn}^{\oplus2}$:
the five symmetric-power monomials have signs $+,-,+,-,+$,
and the determinant twist has sign $\operatorname{sgn}^{-2}=1$.
Every real place of $\mathbb Q$ splits in the real quadratic
extension. Local base change there is the identity on the real
parameter, so the same parameter belongs to $\varpi_\infty$.
There is no descent ambiguity at infinity; in particular
$\eta_\infty=1$.

For $W$, complex conjugation has order two and trace one in
dimension five. Its $+1$ and $-1$ multiplicities are therefore
three and two. Thus the two actual gamma factors are exactly
\[
 G_{\varpi}(s)=G_W(s)
  =\Gamma_{\mathbb R}(s)^3\Gamma_{\mathbb R}(s+1)^2
  =\Gamma_{\mathbb R}(s)\Gamma_{\mathbb C}(s)^2.
 \tag{MA-FQ.1d}
\]
Booker's formula for the kernel in Example~1.4(2), also used for
Artin representations in Example~1.4(3), consequently gives
$K_R=0$. This proves equality of the kernels, not merely equality
of their residues at zero. The sources are Calegari's character
table, author-hosted publisher PDF p.~4; Lemma~3.2 and its
preceding paragraph, p.~8; and the opening of \S4, pp.~8--9.
The invoked modularity, symmetric-power transfer and cyclic
base-change results remain the classical inputs used in
Calegari's construction. This deduction does not assume the
desired all-place comparison over $\mathbb Q$.

\emph{The divisor and the precise Fourier identity.}
MA-LD-001 proves that the complete signed divisor of this actual
quotient is supported on $\operatorname{Re}s=1/2$, using
$(\mathcal R_B)$ and MA-S5-004; if $k=0$ the assertion follows
already from MA-S5-001. Multiplication by the nonzero conductor
exponential changes no order. Hence $m_R$ is integer valued,
real supported, and has $m_R(0)=-k$. Axiom~(A3) of Booker's
Definition~1.3 gives constants $C,A$ such that
\[
 \sum_{|t|\le T}|m_R(t)|\le C(1+T)^A.
 \tag{MA-FQ.1e}
\]
Integrality implies that there are only finitely many nonzero
masses in each compact interval. Summation over dyadic annuli
shows that $\int(1+|t|)^{-N}\,d|\mu|(t)<\infty$ for $N>A$.
Thus $\mu$ defines a tempered distribution as well as a locally
finite signed measure.

Use the following Fourier convention throughout this argument:
\[
 \widehat u(x)=\int_{\mathbb R}e^{itx}u(t)\,dt,
 \qquad
 u(t)=\frac1{2\pi}\int_{\mathbb R}e^{-itx}\widehat u(x)\,dx,
 \qquad
 \langle\widehat T,g\rangle=\langle T,\widehat g\rangle.
 \tag{MA-FQ.1f}
\]
Let $\mathcal B$ be the finite bad-prime set, disjoint from
$\Omega$. Define the positive locally finite measure
\[
 M_\Omega=\sum_{\substack{p\in\Omega\\j\ge1\ \mathrm{odd}}}
 \frac{4\log p}{p^{j/2}}
   \bigl(\delta_{j\log p}+\delta_{-j\log p}\bigr)
 \tag{MA-FQ.1g}
\]
and the possibly signed or complex measure
\[
 B_0=2f_R(1)\delta_0+
 \sum_{p\in\mathcal B}\sum_{j\ge1}
 \left(f_R(p^j)\delta_{j\log p}
       +\overline{f_R(p^j)}\delta_{-j\log p}\right).
 \tag{MA-FQ.1h}
\]
Its total variation is finite. Indeed, axiom~(A1) with exponent
two gives
\[
 \sum_{j\ge1}|f_R(p^j)|
 \le C_2\sum_{j\ge1}(j\log p)^{-2}<\infty
 \quad\text{for each fixed }p,
 \tag{MA-FQ.1i}
\]
and $\mathcal B$ is finite. This argument uses no sign assumption
on any bad Euler coefficient. The stronger root-of-unity
description of Calegari's bad local polynomials in Lemma~4.8
is consistent with this bound but is not needed here.

The measure $M_\Omega$ is tempered. For $k>0$, MA-S5-004 gives
$\sum_{p\le e^T,\ p\in\Omega}(\log p)/\sqrt p=O(1+T)$.
The total mass contributed by $j\ge3$ is bounded by
\[
 8\sum_p(\log p)\sum_{j\ge3}p^{-j/2}<\infty.
\]
Consequently $M_\Omega([-T,T])=O(1+T)$. If $k=0$, then
$\Omega$ is finite by MA-S5-001, and the full geometric tails
have finite mass, which gives the same conclusion.

For a compact smooth test $g$ satisfying
$g(-x)=\overline{g(x)}$, its Fourier transform in (MA-FQ.1f)
is real on $\mathbb R$. Booker's (A4), with $K_R=0$, reads
\[
 \langle\mu,\widehat g\rangle
 =-2\operatorname{Re}\sum_{n\ge1}f_R(n)g(\log n).
 \tag{MA-FQ.1j}
\]
The summation includes $n=1$ under the minus sign. Since
$f_R(1)$ and $g(0)$ are real, the term at zero is
$-2f_R(1)g(0)$. For $n\ge2$, the two atoms in (MA-FQ.1h)
give exactly twice the real part of $f_R(n)g(\log n)$.
The coefficients at good primes are the real nonnegative
coefficients in (MA-LD.6). All other coefficients with $n>1$
vanish unless $n$ is a prime power. It follows that
\[
                  \widehat{(-\mu)}=M_\Omega+B_0.
 \tag{MA-FQ.1k}
\]
To see that this is a complex distribution identity, decompose
an arbitrary compact smooth $g$ as $g=g_1+ig_2$, where
\[
 g_1(x)=\frac{g(x)+\overline{g(-x)}}2,
 \qquad
 g_2(x)=\frac{g(x)-\overline{g(-x)}}{2i}.
\]
Both $g_1$ and $g_2$ have the required Hermitian symmetry, and
both sides of (MA-FQ.1k) are complex linear in the test.
The identity on compact smooth tests extends to Schwartz tests
because all three distributions are tempered and compact
smooth functions are dense in the Schwartz space.

\emph{Cancellation by a continuous density.}
Define
\[
 q(t)=-\frac1{2\pi}\int_{\mathbb R}e^{-itx}\,dB_0(x),
 \qquad \nu=-\mu+q(t)\,dt.
 \tag{MA-FQ.1l}
\]
Finite total variation makes $q$ bounded and continuous, with
$\|q\|_\infty\le\|B_0\|_{\mathrm{TV}}/(2\pi)$.
Fourier inversion and Fubini on Schwartz tests give
$\widehat{q(t)\,dt}=-B_0$, so
$\widehat\nu=M_\Omega$. In particular $\nu$ is a distribution
of positive type. More explicitly, with
$\widetilde\psi(t)=\overline{\psi(-t)}$,
\[
 \langle\nu,\psi*\widetilde\psi\rangle
   =\frac1{2\pi}\int_{\mathbb R}
                         |\widehat\psi(-x)|^2\,dM_\Omega(x)
   \ge0\qquad(\psi\in C_c^\infty(\mathbb R)).
 \tag{MA-FQ.1m}
\]
The reflection $-x$ and the factor $(2\pi)^{-1}$ follow from
the Fourier convention (MA-FQ.1f). They do not affect positivity.

\emph{Extraction of the atomic coefficient function.}
Choose a nonzero real even $\phi\in C_c^\infty(\mathbb R)$.
Put $\phi_\epsilon(t)=\epsilon^{-1/2}\phi(t/\epsilon)$ and
$\Psi=\phi*\widetilde\phi$. Then
$\phi_\epsilon*\widetilde\phi_\epsilon(t)=\Psi(t/\epsilon)$,
where $\Psi$ is compactly supported and
$\Psi(0)=\|\phi\|_2^2>0$. For arbitrary finite choices
$x_1,\ldots,x_r\in\mathbb R$ and $c_1,\ldots,c_r\in\mathbb C$,
apply (MA-FQ.1m) to
$\psi_\epsilon=\sum_i c_i\phi_\epsilon(\,\cdot-x_i)$.
This gives
\[
 0\le\sum_{i,j}c_i\overline{c_j}\,
 \left\langle\nu,
  \Psi\left(\frac{\,\cdot-(x_i-x_j)}\epsilon\right)
 \right\rangle.
 \tag{MA-FQ.1n}
\]
There are only finitely many differences $x_i-x_j$.
Local finiteness of the support of $\mu$ therefore allows a
single sufficiently small $\epsilon$ for which the atomic
contribution at each such difference is precisely
$a(x_i-x_j)\Psi(0)$; there are no other atoms in the
corresponding shrinking support. The continuous contribution
has absolute value at most
$\epsilon\|q\|_\infty\|\Psi\|_1$ for each pair $i,j$.
Letting $\epsilon$ decrease to zero in (MA-FQ.1n) and dividing
by $\Psi(0)$ gives
\[
       \sum_{i,j}c_i\overline{c_j}\,a(x_i-x_j)\ge0.
 \tag{MA-FQ.1o}
\]
This is exactly positive definiteness on $\mathbb R_d$.
The two-by-two principal matrices give
$a(-t)=\overline{a(t)}$ and $|a(t)|\le a(0)=k$.

For completeness the standard Bochner implication here requires
no continuity in the usual real topology. Form the pre-Hilbert
space with formal vectors $e_x$, inner products
$\langle e_x,e_y\rangle=a(x-y)$, and quotient its nullspace.
Its completion carries commuting unitary translations
$U_z e_x=e_{x+z}$ and a vector $v=e_0$ of squared norm $k$.
If $k=0$ all the coefficients vanish. Otherwise the unital
commutative $C^*$-algebra generated by the $U_z$ is $C(X)$
for its compact spectrum $X$. A point $\omega\in X$ determines
the character $z\mapsto\omega(U_z)$ of $\mathbb R_d$.
The positive functional $T\mapsto\langle Tv,v\rangle$ is,
by the Riesz representation theorem, integration against a
positive measure of mass $k$ on $X$. Push it forward to
$\widehat{\mathbb R_d}$ to obtain a positive finite measure
$\sigma$ with
\[
        a(z)=\int_{\widehat{\mathbb R_d}}\chi(z)\,d\sigma(\chi),
        \qquad \sigma(\widehat{\mathbb R_d})=k.
 \tag{MA-FQ.1p}
\]
The map to the character group is continuous in its product
topology, so the pushforward is well defined even though the
compact character group need not be metrizable. This proves
$a\in B(\mathbb R_d)$ with norm at most $k$. Every representing
measure has total variation at least $|a(0)|=k$, proving the
norm equality. Finally $a$ is integer valued, so the bound
$|a|\le k$ gives finite range. Its support was already shown
to be locally finite. This proves every assertion.
\end{proof}

\paragraph{Proof scope and cited input.}
The argument above proves the atomic positivity assertion even
with arbitrary signs in the finitely many bad-prime progressions.
It does not require that $-\mu$ itself be of positive type;
the continuous correction in (MA-FQ.1l) is part of the proof.
No theorem about Fourier quasicrystals with uniformly discrete
support is invoked: uniform discreteness has not been assumed.
The next coset-ring step uses the locally compact group
$\mathbb R_d$, while local finiteness always uses the usual
real topology.

The exact Booker locators are Definition~1.3(A1), (A3), (A4),
author p.~3, and Examples~1.4(2)--(3), p.~4, in
\path{arXiv:1308.3067v2}. The actual critical-line-support
argument and its complete source-proof audit are MA-LD-001.
The compact-group Fourier--Stieltjes convention is the one used
in Sanders (2020), \S12, before Theorem~12.1; replacing each
character by its inverse changes neither the algebra nor its norm.
The last paragraph of the proof supplies the required special
case of the usual Bochner argument from the commutative
$C^*$-algebra and Riesz representation theorems, which are
classical background inputs.

Calegari, F. (2013). The Artin conjecture for some $S_5$-extensions.
\emph{Mathematische Annalen, 356}, 191--207.
\url{https://doi.org/10.1007/s00208-012-0839-4}.

Booker, A. R. (2015). $L$-functions as distributions.
\emph{Mathematische Annalen, 363}, 423--454.
\url{https://doi.org/10.1007/s00208-015-1178-z}.

\section{The discrete coset-ring step}
\label{sec:s5-coset}

Write $\mathbb R_d$ for the additive real group with the discrete
topology. Its coset ring is the Boolean algebra generated by the
cosets of all its abstract subgroups. Local finiteness below always
refers to the usual topology of $\mathbb R$: every compact interval
meets the set in finitely many points. This distinction matters
because an abstract subgroup used in a coset-ring representation
can be dense in the usual topology.

\paragraph{MA-FQ-002 (local finiteness removes dense cosets).}
If $E$ belongs to the coset ring of $\mathbb R_d$ and is locally
finite in $\mathbb R$, there are a finite set $F$ and finitely many
arithmetic progressions $\beta_j+\alpha_j\mathbb Z$, with
$\alpha_j>0$, such that
\[
 E\mathbin{\triangle}\bigcup_{j=1}^{J}
                  (\beta_j+\alpha_j\mathbb Z)=F.
 \tag{MA-FQ.2}
\]
The empty union is allowed. Consequently, if an integer-valued
function $a\in B(\mathbb R_d)$ has locally finite support in
$\mathbb R$, then there are integers $c_j,d_\ell$ and real numbers
$\alpha_j>0$, $\beta_j$, and $\gamma_\ell$, with finite index sets,
such that
\[
 a(t)=\sum_{j=1}^{J}c_j
                 1_{\beta_j+\alpha_j\mathbb Z}(t)
       +\sum_{\ell=1}^{L}d_\ell 1_{\{\gamma_\ell\}}(t).
 \tag{MA-FQ.3}
\]
In particular its atomic measure is the corresponding finite
integer linear combination of translated lattice combs and point
masses. No uniform separation of distinct points is assumed.

\begin{proof}
We first prove the assertion about sets, including the issue of
Boolean differences of dense subgroups. Every element of the
coset ring is a finite union of atoms of the form
\[
 A\setminus\bigcup_{i=1}^{n}B_i,
 \tag{MA-FQ.4}
\]
where $A$ is a coset of a subgroup, and each $B_i$ is a coset
contained in $A$. Indeed, take the finite Boolean algebra
generated by the cosets in one expression for $E$, intersect the
positive cosets, and intersect each excluded coset with that
intersection. If there are no positive cosets, use
$A=\mathbb R$. Empty atoms may be discarded.

Suppose an atom in (MA-FQ.4) is nonempty and that the subgroup
underlying $A$ is dense. Translate a surviving point to $0$.
We then have $A=H$, with $H$ a dense subgroup of $\mathbb R$,
and $B_i=x_i+H_i\subset H$, with $0\notin B_i$.
Consider the finite family
\[
 \left\{H\cap\bigcap_{i\in I}H_i:
                         I\subseteq\{1,\ldots,n\}\right\}.
\]
Choose an inclusion-minimal member $K$ that is dense in
$\mathbb R$; such a member exists because $H$ is in the family.
For each $i$, if $K\cap H_i$ is dense, minimality gives
$K\cap H_i=K$. Thus $K\subset H_i$, and $K\cap B_i$ is
empty: an intersection would imply $B_i=H_i$, contrary to
$0\notin B_i$.

In the other case, $K\cap H_i$ has proper closed subgroup
closure in $\mathbb R$. Such a closure is either $\{0\}$ or a
lattice, and is nowhere dense. The intersection $K\cap B_i$
is either empty or a coset of $K\cap H_i$, so its closure is
also nowhere dense. A finite union of these closures cannot
fill any open interval. Every open interval therefore contains
a nonempty open subinterval avoiding them, and the density of
$K$ supplies a point of $K$ in that subinterval. It follows that
\[
 K\setminus\bigcup_{i=1}^{n}B_i
\]
is dense in $\mathbb R$. This set lies in the atom, so a
nonempty atom with dense ambient subgroup cannot be a subset
of a locally finite set. Every nonempty atom in the expression
for $E$ must therefore have discrete ambient subgroup.

A discrete subgroup of $\mathbb R$ is either $\{0\}$ or
$\alpha\mathbb Z$ for some $\alpha>0$. The first case gives a
singleton atom. In the second case identify $A$ with
$\mathbb Z$. Each excluded coset $B_i$ is then either a
singleton or a residue class modulo a positive integer $n_i$.
Let $N$ be the least common multiple of the finitely many
$n_i$, with $N=1$ when there are none. Removing the residue
classes leaves a union of residue classes modulo $N$; removing
the singleton exclusions changes it at only finitely many
points. Hence every atom is a finite union of arithmetic
progressions up to a finite symmetric difference. Taking the
finite union proves (MA-FQ.2).

For the function assertion, Sanders's Theorem~12.3 applies to
the locally compact Abelian group $\mathbb R_d$. It gives
finitely many subgroups $H_i$ and integer-valued functions
$z_i$ on $\mathbb R_d/H_i$ with finite $\ell^1$ norms, such
that
\[
 a=\sum_i\sum_{W\in\mathbb R_d/H_i}z_i(W)1_W.
 \tag{MA-FQ.5}
\]
An integer-valued $\ell^1$ function has finite support.
Thus (MA-FQ.5) uses only finitely many cosets. The function
$a$ is bounded, hence has finite integer range, and each level
set $E_r=\{t:a(t)=r\}$ lies in the finite Boolean algebra
generated by those cosets. For $r\ne0$, the set $E_r$ is a
subset of the locally finite support of $a$; apply (MA-FQ.2)
to it.

It remains to pass from a union of progressions to an integer
linear combination of their indicators. Inclusion--exclusion
does this: every nonempty intersection of finitely many
progressions is either another progression or a singleton,
because it is a coset of the intersection of their discrete
subgroups. The finite symmetric differences add only finitely
many singleton indicators. Sum the resulting expressions for
$r1_{E_r}$ over the finite nonzero range of $a$. This proves
(MA-FQ.3), including the assertion about the associated atomic
measure.
\end{proof}

\paragraph{Primary sources and exact scope.}
The set classification is Schreiber's Example~1.6, p.~808;
its dense-coset input is Lemma~1.4 and Corollary~1.5, p.~807.
The proof above supplies the needed real-line argument directly,
choosing an actual surviving point before taking intersections.
Kolountzakis gives a second general proof in
\S3.2, Lemma~2 and Theorem~3, author pp.~7--8 of
\path{arXiv:math/9904065v1}; the real-line classification is
display (9), p.~7, attributed there to Rosenthal. The
Fourier--Stieltjes input is Sanders's Theorem~12.3 and definition
(12.2); only its finite qualitative representation is used.
The present argument does not require its quantitative bounds
and does not reprove Cohen's idempotent theorem.

Schreiber, B. M. (1970). On the coset ring and strong Ditkin
sets. \emph{Pacific Journal of Mathematics, 32}(3), 805--812.
\url{https://doi.org/10.2140/pjm.1970.32.805}.

Kolountzakis, M. N. (2000). On the structure of multiple
translational tilings by polygonal regions.
\emph{Discrete \& Computational Geometry, 23}(4), 537--553.
\url{https://doi.org/10.1007/s004540010014}.

Sanders, T. (2020). Bounds in Cohen's idempotent theorem.
\emph{Journal of Fourier Analysis and Applications, 26},
Article 25. \url{https://doi.org/10.1007/s00041-020-09732-y}.

\section{Fourier rigidity and the central excess}
\label{sec:s5-crystal-rigidity}

The argument here uses the actual signed divisor and its ordinary
distributional Fourier transform.  It does not assume that the
divisor is positive, that zeros are simple, or that the degree-ten
symmetric-square function is automorphic.  The analytic input is
the positive-definite atomic coefficient function constructed in
MA-FQ-001.  The finite-coset reduction is MA-FQ-002.

\paragraph{MA-FQ-003 (integer atoms and a prime-frequency obstruction).}
Let $a:\mathbb R\to\mathbb Z$ be positive definite when $\mathbb R$
is regarded as a discrete additive group.  Suppose its nonzero
support is locally finite in the ordinary real topology, and
the atomic measure
\[
                 \nu=\sum_{t\in\mathbb R}a(t)\delta_t
\tag{MA-FQ.20}
\]
is tempered.  Suppose also that its ordinary Fourier transform,
with convention
$\widehat\nu(x)=\int e^{itx}\,\nu(dt)$, is a purely atomic
locally finite complex measure.  Then there are finitely many
numbers $h_1,\ldots,h_N>0$ such that
\[
                 \operatorname{supp}\widehat\nu
                        \subseteq\bigcup_{j=1}^{N}h_j\mathbb Z.
\tag{MA-FQ.21}
\]
The empty union is allowed when $\nu=0$.
In particular, for any fixed $c>0$, only finitely many distinct
rational primes $p$ can have
\[
                         \widehat\nu(\{c\log p\})\ne0.
\tag{MA-FQ.22}
\]

\begin{proof}
The two-by-two positive-definite inequality gives
$|a(t)|\le a(0)$ for every $t$.  If $a(0)=0$, then $a=0$
and the conclusions follow.  Otherwise put $k=a(0)\in\mathbb N$.
The function has finite integer range contained in
$\{-k,-k+1,\ldots,k\}$.
By Bochner's theorem on the locally compact discrete abelian
group $\mathbb R_d$, it belongs to $B(\mathbb R_d)$, with
$\|a\|_{B(\mathbb R_d)}=k$.  Its representing measure lives on
the compact Pontryagin dual of $\mathbb R_d$; it is not being
identified with the ordinary Fourier transform in (MA-FQ.21).

For every nonzero value $r$ in the range, polynomial
interpolation on this finite range gives a polynomial $P_r$
with $P_r(r)=1$ and $P_r(q)=0$ at every other range value.
Thus $P_r(a)=1_{\{a=r\}}\in B(\mathbb R_d)$, since the
Fourier--Stieltjes algebra is closed under multiplication
and contains constants.  Cohen's idempotent theorem places
each such level set in the coset ring of $\mathbb R_d$.
These level sets are locally finite.  MA-FQ-002 therefore
writes their indicators, and hence $a$, as a finite integer
linear combination of indicators of ordinary discrete
cosets in $\mathbb R$.

Every nonzero discrete subgroup of $\mathbb R$ is of the
form $\alpha\mathbb Z$ with $\alpha>0$; the zero subgroup
gives a singleton.  Combining repeated terms, we obtain
an equality of tempered measures
\[
 \nu=\sum_{j=1}^{J}c_j\sum_{n\in\mathbb Z}
                         \delta_{\alpha_j n+\beta_j}
                  +\sum_{\ell=1}^{L}d_\ell\delta_{\tau_\ell},
 \qquad c_j,d_\ell\in\mathbb Z,\quad \alpha_j>0 .
\tag{MA-FQ.23}
\]
No positivity of the individual coefficients in this
representation is asserted or needed.  In particular,
the two Fourier transforms used in the preceding paragraph
and the next paragraph are kept distinct.

The ordinary Poisson summation identity for a lattice gives
\[
 \widehat{\sum_{n\in\mathbb Z}\delta_{\alpha n+\beta}}
    =\frac{2\pi}{\alpha}\sum_{m\in\mathbb Z}
        e^{2\pi i m\beta/\alpha}\delta_{2\pi m/\alpha}.
\tag{MA-FQ.24}
\]
This is an identity of tempered distributions, obtained by
pairing with a Schwartz test and applying the classical
one-dimensional Poisson formula.  A singleton transforms
to the ordinary continuous function $e^{i\tau x}$, interpreted
as its density against $dx$.  Thus (MA-FQ.23) transforms to
a measure supported on finitely many unshifted frequency
lattices, plus the absolutely continuous measure with density
\[
                         \sum_{\ell=1}^{L}d_\ell e^{i\tau_\ell x}.
\]
By hypothesis the complete transform is purely atomic.
Uniqueness of the local Lebesgue decomposition forces this
continuous density to vanish almost everywhere, hence
everywhere by continuity.  The remaining support is
contained in the finite union of lattices with
$h_j=2\pi/\alpha_j$, proving (MA-FQ.21).
Any cancellation among the lattice contributions can only
shrink their union of supports.

For a fixed $h>0$, suppose $c\log p=nh$ and
$c\log q=mh$ for distinct primes $p,q$ and positive
integers $n,m$.  Then $m\log p=n\log q$, so
$p^m=q^n$, contradicting unique factorization in
$\mathbb Z$.  Each lattice in (MA-FQ.21) therefore
contains $c\log p$ for at most one distinct prime.
This proves (MA-FQ.22).
\end{proof}

\paragraph{MA-FQ-004 (central equality in the stipulated family under
$(\mathcal R_B)$).}
Retain the exact $S_5$ family and the hypothesis
$(\mathcal R_B)$ of MA-S5-001.  Then
\[
                              \boxed{\beta=b+n.}
\tag{MA-FQ.25}
\]
Equivalently, the exceptional set $\Omega$ is finite and
the two rank-five comparison factors match at every place,
in the precise sense of MT-S508 and MA-S5-001.

\begin{proof}
MA-S5-001 first gives $k=\beta-b-n\ge0$.
Suppose for contradiction that $k>0$.
Use the actual difference datum
$\mathcal F_R=\mathcal F_{\varpi}-\mathcal F_W$ from
MA-LD-001 and its signed completed divisor
\[
              \mu=\sum_{t\in\mathbb R}m_R(t)\delta_t.
\]
Its support is on the critical line, written here in
height coordinates, by MA-LD-001.  That result uses
$(\mathcal R_B)$ and the positive-excess exceptional-prime
asymptotic, not the input $M$.
In particular $m_R(0)=-k$.

MA-FQ-001 proves that
$a(t)=-m_R(t)$ is an integer-valued positive-definite
function on $\mathbb R_d$.  It has locally finite support,
and $-\mu$ is tempered.  The same proof verifies that its
ordinary Fourier transform is a locally finite purely
atomic measure:
\[
                       \widehat{-\mu}=M_{\rm good}+B.
\tag{MA-FQ.26}
\]
Here the possible conductor atom and every bad-prime
correction are contained in the finite-total-variation
measure $B$.  The good-prime part has, for each
$p\in\Omega$, the nonzero atom
\[
              M_{\rm good}(\{\log p\})
                        =4(\log p)p^{-1/2}>0.
\tag{MA-FQ.27}
\]
The set $\Omega$ contains only good primes.  No power
of a different rational prime has norm $p$, so neither
another prime-power term nor a bad-prime correction
cancels this atom.  The conductor atom is at zero.
Consequently $\widehat{-\mu}(\{\log p\})\ne0$ for every
$p\in\Omega$.

Apply MA-FQ-003 to $\nu=-\mu$ with $c=1$.
It follows that $\Omega$ is finite.
MA-S5-001, the fully checked form of MT-S508,
then gives $k=0$, contradicting $k>0$.
Thus $k=0$ and (MA-FQ.25) follows.  The equivalent
all-place comparisons are the conclusion of that same
checked source criterion.
\end{proof}

\paragraph{Scope of the conclusion and the other inputs.}
The hypothesis $(\mathcal R_B)$ remains explicit.
This theorem does not prove it for every extension in the
family.  It excludes an additional central excess under
exactly that hypothesis, so it closes the central branch
that remained in MT-S508.  The all-place comparison gives
the Artin automorphy and holomorphy consequences recorded
in MA-S5-001.

GRH for all irreducible factors still requires the
appropriate Hecke zero-location input, sufficiently $M$.
Since $B$ is a Hecke function, $M$ also supplies
$(\mathcal R_B)$ in this family.  The Dirichlet theorem
$A$ alone does not replace that general Hecke premise.
No field with $k>0$ is constructed, and no unrestricted
$S_5$ Artin conjecture is claimed.

MA-LD-001--004 and the formal order assignment MA-S5-005
remain valid conditional deductions.  Their positive-excess
antecedent is now excluded by MA-FQ-004.  The formal
assignment was never an arithmetic realization; the
Fourier rigidity condition is the additional structure
that the central-order inequalities did not encode.
Similarly, the separate sufficient functoriality criterion
MA-S5-007 remains correct but is no longer needed to
eliminate this central branch.

For Artin motives, the conclusion identifies the actual
rank-five motive's $L$-function with its automorphic
comparison under the stipulated family hypotheses.
The finite-monodromy interpretation is the one in
SGA~1, Expos\'e~V, \S7 and Deligne,
\emph{Valeurs de fonctions $L$}, Definition~6.1,
\S\S6.2--6.3.  The proof uses harmonic analysis of the
signed divisor and prime frequencies in addition to
representation-ring identities.  It does not construct
a positive faithful Hilbert action on all zero jets,
and therefore does not evade or contradict MT-M20/M23.

\paragraph{Primary sources and proof-check boundary.}
The only deep harmonic-analysis input to MA-FQ-003 is
the classical idempotent theorem and its discrete-coset
consequence, with the precise versions checked in
MA-FQ-002.  An exact modern primary statement is
Sanders (2020), \S12, Theorems~12.1 and~12.3,
\emph{Bounds in Cohen's idempotent theorem},
\emph{Journal of Fourier Analysis and Applications, 26},
Article~25,
\url{https://doi.org/10.1007/s00041-020-09732-y}.
The passage from ordinary discrete cosets to lattice
combs and then to prime-frequency finiteness is proved
in full above.

The discrete-coset theorem is also stated and proved in
Kolountzakis (2000), \S3.2, Lemma~2 and Theorem~3,
\emph{On the structure of multiple translational tilings
by polygonal regions},
\emph{Discrete \& Computational Geometry, 23}, 537--553;
author version \url{https://arxiv.org/abs/math/9904065},
pp.~7--8.  The present argument uses the theorem in one
real dimension, with the ordinary topology used for
local finiteness and the discrete topology used for
the Fourier--Stieltjes algebra.

\section{Scope tests for the Fourier rigidity argument}
\label{sec:s5-scope-tests}

The central equality in MA-FQ-004 uses the actual completed divisor,
exact cancellation of the archimedean kernel, and the sign and
support of the remaining prime terms.  The following tests make
these hypotheses explicit.  They are not additional hypotheses in
that theorem and are not claims of first priority for the known
$L$-function or crystalline examples.

\paragraph{MA-FQ-005 (standard examples do not contradict the combined rigidity premises).}
The following statements hold.
\begin{enumerate}
\item The genuine degree-zero datum for $\zeta(s)/L(s,\chi_5)$,
where $\chi_5$ is the primitive even quadratic character modulo
five, has $K=0$ and nonnegative coefficients $f(n)$ for $n>1$.
Its full completed divisor is not supported on the critical line,
because it has poles at $s=0,1$.
\item The genuine degree-zero datum for
$L(s,\chi_5)/L(s,\chi_{13})$, with the two primitive even
quadratic characters, has infinitely many negative coefficients
at good primes.  Removing finitely many bad-prime progressions
does not give the sign hypothesis of MA-FQ-001.
\item The explicit integer-mass Fourier quasicrystal below is a
positive measure, but its atomic coefficient function is not
positive definite on $\mathbb R_d$.  Its ordinary Fourier
transform cannot be a nonnegative tempered measure plus a
finite-total-variation measure.  Its spectrum also contains a
frequency which is not the logarithm of a prime power.
\item The elementary one-prime function
$(1-p^{1/2-s})^{-1}$ has a central vertical pole lattice but
violates Booker's axiom (A1).  It is not an allowed bad-prime
correction in MA-FQ-001.
\end{enumerate}

\begin{proof}
\emph{The positive Euler ratio.}
Booker's actual Artin data for the trivial character and
$\chi_5$ have the same real gamma factor, so their difference
has $K=0$.  Logarithmic differentiation gives
\[
 f(n)=(1-\chi_5(n))\frac{\Lambda(n)}{\sqrt n}\ge0
                         \qquad(n>1).
 \tag{MA-FQ.30}
\]
The completed zeta function, before multiplying by any
pole-removing polynomial, has simple poles at $0$ and $1$.
The completed nontrivial Dirichlet function is nonzero at
both points: use its usual nonvanishing at $1$ and its
functional equation.  The completed quotient therefore
has these same two poles.  In height coordinates
$s=1/2+iz$, they are at $z=\pm i/2$, so the required real
support already fails, independently of the location of
any nontrivial zero.

Deleting those poles would change the explicit formula.
Indeed, for $h(z)=\int e^{ixz}g(x)\,dx$ their contribution is
\[
 h(i/2)+h(-i/2)=
                    \int_{\mathbb R}2\cosh(x/2)g(x)\,dx.
 \tag{MA-FQ.31}
\]
If the remaining signed divisor were real supported,
the formula for the Fourier transform of its negative
would contain the additional density $-2\cosh(x/2)$.
This is not the identity used in MA-FQ-001.
Its positive prime measure is not tempered either:
polynomial mass growth would imply convergence of every
positive Laplace transform, hence local convergence of
the Euler logarithm throughout $\operatorname{Re}s>1/2$.
Exponentiation would then remove the actual pole at $1$,
a contradiction.  Thus one cannot discard the endpoint
poles and retain the tempered positive-prime identity.

\emph{The ratio of two entire Dirichlet completions.}
Both characters are even, so the gamma kernels agree, and
their nontrivial completions are entire.  At an unramified
prime the difference datum has coefficient
\[
 f(p)=(\chi_5(p)-\chi_{13}(p))\frac{\log p}{\sqrt p}.
 \tag{MA-FQ.32}
\]
If $p\equiv3\pmod{65}$, the two character values are
$-1,+1$, respectively.  If $p\equiv11\pmod{65}$, they
are $+1,-1$.  These values follow directly from the
quadratic residues modulo five and thirteen.  Each of
the two reduced residue classes contains infinitely many
primes by the classical Dirichlet theorem.  In particular
no finite set of bad primes removes all the negative
coefficients.  This sign failure is independent of any
zero-location hypothesis.

\emph{The explicit crystalline test.}
Use Kurasov--Sarnak's polynomial
\[
 P(z_1,z_2)=1-\frac{z_1}{3}+\frac{z_2^2}{3}-z_1z_2^2,
 \qquad F(w)=P(2^{-w},3^{-w}).
 \tag{MA-FQ.33}
\]
Their Theorem~1 and equations (39), (41)--(43) construct
the positive crystalline zero measure of $F$.  Here is
a direct test of its relevant coefficient function.
Put
\[
 \xi_1=\log2,\quad \xi_2=\log3,\quad
 A=\xi_1/2+\xi_2,\quad D=\xi_1/2-\xi_2.
\]
Its real zero ordinates satisfy
\[
                 S(t)=3\sin(At)+\sin(Dt)=0.
 \tag{MA-FQ.34}
\]
One has $|D|<A$ and $A/D\notin\mathbb Q$, since a
rational ratio would make $\log2/\log3$ rational.
At a root, $|\sin(At)|\le1/3$, and hence
$|3A\cos(At)|\ge\sqrt8 A>|D|$.  Thus every root is
simple, including zero.  There is a nonzero root
$\lambda$ between $\pi/(2A)$ and $3\pi/(2A)$, because
$S$ is positive at the first endpoint and negative at
the second.

If $2\lambda$ were also a root, the double-angle formula
and $S(\lambda)=0$ would give
\[
           0=6\sin(A\lambda)
                     \{\cos(A\lambda)-\cos(D\lambda)\}.
 \tag{MA-FQ.35}
\]
If the sine vanishes, both $A\lambda$ and $D\lambda$
are integer multiples of $\pi$, contradicting
$A/D\notin\mathbb Q$.  Otherwise the cosines agree,
so $A\lambda\equiv D\lambda$ or $-D\lambda$ modulo
$2\pi$.  The corresponding equality or opposite
equality of the sines, together with $S(\lambda)=0$,
again makes both sines zero and gives the same
contradiction.  Therefore $2\lambda$ is not a root.

The atomic mass function $a$ of this measure consequently
has $a(0)=a(\lambda)=a(-\lambda)=1$ and
$a(2\lambda)=a(-2\lambda)=0$.  Its matrix on
$0,\lambda,2\lambda$ is not positive semidefinite:
\[
 (1,-1,1)
 \begin{pmatrix}1&1&0\\1&1&1\\0&1&1\end{pmatrix}
 (1,-1,1)^{\mathsf T}=-1.
 \tag{MA-FQ.36}
\]
Thus $a$ is not positive definite on $\mathbb R_d$.
If the ordinary Fourier transform of this atomic
measure were a nonnegative tempered measure plus a
finite measure, subtracting the latter's inverse
transform would be a bounded continuous correction.
The shrinking-bump proof in MA-FQ-001 would then make
its atomic coefficient function positive definite,
contradicting (MA-FQ.36).

There is also an independent Euler-support obstruction.
Set $v=z_2^2$.  Since
\[
 P(z_1,z_2)=(1+v/3)
       \left(1-z_1\frac{1/3+v}{1+v/3}\right),\qquad
 \frac{1/3+v}{1+v/3}
     =\frac13+\frac89v-\frac8{27}v^2+O(v^3),
\]
the coefficient of $z_1z_2^4$ in the actual Taylor
series for $\log P$ is $8/27$.
The Fourier summation formula therefore gives a
nonzero atom at $\xi_1+4\xi_2=\log162$.
The integer $162=2\cdot3^4$ is not a prime power.
This term cannot occur in the Euler logarithmic
derivative of the actual Artin/automorphic quotient
used in MA-FQ-001.

\emph{The central one-prime lattice.}
For the last example,
\[
 \log(1-p^{1/2-s})^{-1}
         =\sum_{j\ge1}\frac{p^{j(1/2-s)}}j.
\]
Comparison with Booker's normalization
$\sum_{n\ge2}(f(n)/\log n)n^{1/2-s}$
gives $f(p^j)=\log p$ for every $j$.
For any positive $N$, the sequence
$f(p^j)(j\log p)^N$ is unbounded.
This violates (A1).  Equivalently its reciprocal
frequency comb has nondecaying masses, rather than
the summable bad-prime tails of MA-FQ.1i.
This proves all four statements.
\end{proof}

\paragraph{Source locators and proof-check boundary.}
The Dirichlet examples use Booker's Definition~1.3,
author p.~3, and Examples~1.4(2)--(3), p.~4, in
\path{arXiv:1308.3067v2}; the first example is also
MA-LD-003.  Classical nonvanishing at one and
Dirichlet's theorem on primes in arithmetic progressions
are carried background inputs.  The exact character
values, endpoint correction, and A1 obstruction above
are checked directly.

Kurasov, P., \& Sarnak, P. (2020).
\emph{Stable polynomials and crystalline measures}.
Journal of Mathematical Physics, 61, 083501.
Author version \url{https://arxiv.org/abs/2004.05678}:
the summation formula and Theorem~1 with proof,
\S2, pp.~3--8; equations (39), (41)--(43),
\S\S2--3, pp.~8--10.  The complete invoked summation
argument and Theorem~1 proof were read in this run.
The matrix calculation above does not use the paper's
deeper additive-structure classification.

\paragraph{A checked family-hypothesis remark.}
Calegari's \S3 opening, publisher PDF p.~7, says that
three and five both split in the real quadratic field,
although the family hypotheses in Theorem~1.1 impose
only the condition at five.  This does not hide an
additional condition needed here: Sasaki's published
Theorem~1 requires complete splitting at five and
projective decomposition image of order two above
five, with no condition at three.  The stipulated
double-transposition Frobenius supplies those
requirements.  This is a statement/application check,
not a new proof of the modularity-lifting theorem:
Sasaki, S. (2013), \emph{On Artin representations and
nearly ordinary Hecke algebras over totally real fields},
Documenta Mathematica, 18, 997--1038,
Theorem~1, pp.~998--999,
\url{https://doi.org/10.4171/DM/419}.

\section{An abstract vanishing criterion for real-supported Euler data}
\label{sec:s5-crystal-criterion}

The positive-prime measure in the preceding argument does not need
an independent Tauberian growth estimate once the signed divisor is
known to be supported on the real height axis. The explicit-formula
axiom itself gives its polynomial growth. This yields the following
criterion for an actual $L$-datum; it is not restricted to the
$S_5$ comparison.

\paragraph{MA-FQ-006 (real integral divisor and eventual Euler positivity).}
Let $\mathcal F=(f,K,m)$ be an $L$-datum in the precise sense of
Booker, Definition~1.3. Suppose that
\[
 \begin{gathered}
 K(x)=0\quad(x>0),\qquad
 m(z)=0\quad(z\notin\mathbb R),\qquad
 m(t)\in\mathbb Z\quad(t\in\mathbb R),\\
 f(n)=0\quad\text{if }n\ge2\text{ is not a prime power},\\
 \text{there is a finite set of primes }\mathcal B
 \text{ such that }
 f(p^j)\in[0,\infty)
       \quad(p\notin\mathcal B,\ j\ge1).
 \end{gathered}
 \tag{MA-FQ.40}
\]
Then
\[
                  \mathcal F=(0,0,0).
 \tag{MA-FQ.41}
\]
In particular its associated $L$-function is identically one.
The hypothesis $K\equiv0$ is equality of the full kernel,
which is stronger than degree zero alone. No sign condition
on $m$, no assumption on its central value, and no separate
growth estimate for the good-prime measure are required.

\begin{proof}
\emph{The atomic divisor and the finite bad-prime correction.}
Use the ordinary Fourier convention (MA-FQ.1f):
\[
 \widehat u(x)=\int e^{itx}u(t)\,dt,\qquad
 \langle\widehat T,g\rangle=\langle T,\widehat g\rangle,
 \qquad
 u(t)=\frac1{2\pi}\int e^{-itx}\widehat u(x)\,dx.
 \tag{MA-FQ.42}
\]
By axiom (A3) there are constants $C,A$ with
\[
 \sum_{|t|\le T}|m(t)|\le C(1+T)^A.
 \tag{MA-FQ.43}
\]
Integrality makes the nonzero support locally finite in
the ordinary real topology. The signed atomic measure
$\mu=\sum_t m(t)\delta_t$ is tempered: summing (MA-FQ.43)
over dyadic annuli shows that its total variation is
integrable against $(1+|t|)^{-N}$ for $N>A$.

Define the positive locally finite measure
\[
 M_{\rm good}
   =\sum_{\substack{p\notin\mathcal B\\j\ge1}}
        f(p^j)\bigl(\delta_{j\log p}+\delta_{-j\log p}\bigr)
 \tag{MA-FQ.44}
\]
and the complex measure
\[
 B_0=2f(1)\delta_0+
    \sum_{\substack{p\in\mathcal B\\j\ge1}}
       \bigl(f(p^j)\delta_{j\log p}
                    +\overline{f(p^j)}\delta_{-j\log p}\bigr).
 \tag{MA-FQ.45}
\]
Local finiteness in (MA-FQ.44) follows because a bounded
interval contains logarithms of only finitely many
positive integers. Axiom (A1), with exponent two, gives
for every fixed prime $p$
\[
 \sum_{j\ge1}|f(p^j)|
       \le C_2\sum_{j\ge1}(j\log p)^{-2}<\infty.
 \tag{MA-FQ.46}
\]
Thus $B_0$ has finite total variation. This uses no sign
restriction at the bad primes. Also $f(1)\in\mathbb R$,
as required by (A1).

\emph{The explicit formula supplies the missing growth bound.}
For $g\in C_c^\infty$ with $g(-x)=\overline{g(x)}$,
its Fourier transform is real on $\mathbb R$.
Axiom (A4), with $K=0$, gives
\[
 \langle\mu,\widehat g\rangle
            =-2\operatorname{Re}
                    \sum_{n\ge1}f(n)g(\log n).
 \tag{MA-FQ.47}
\]
The sum includes $n=1$ under the minus sign. The
prime-power support and the signs at the good primes
give the identity
\[
                 \widehat{(-\mu)}=M_{\rm good}+B_0
                    \quad\text{on }C_c^\infty(\mathbb R).
 \tag{MA-FQ.48}
\]
This is a complex distribution identity, not merely an
identity on real tests. Indeed every $g$ decomposes as
$g=g_1+ig_2$, where
\[
 g_1(x)=\tfrac12\bigl(g(x)+\overline{g(-x)}\bigr),\qquad
 g_2(x)=\tfrac1{2i}\bigl(g(x)-\overline{g(-x)}\bigr).
 \tag{MA-FQ.49}
\]
Both have the requisite Hermitian symmetry, and both
sides of (MA-FQ.48) are complex linear in the test.

The left side of (MA-FQ.48) and $B_0$ are already
tempered. Hence the locally defined positive measure
$M_{\rm good}$ agrees on compact tests with the tempered
distribution $T=\widehat{(-\mu)}-B_0$.
Choose $\theta\in C_c^\infty(\mathbb R)$, nonnegative
and at least one on $[-1,1]$, and put
$\theta_R(x)=\theta(x/R)$ for $R\ge1$.
Continuity of $T$ on the Schwartz space bounds its
action by finitely many Schwartz seminorms. For some
integer $D$ and constant $C'$ this gives
\[
 \begin{aligned}
 M_{\rm good}([-R,R])
   &\le\langle M_{\rm good},\theta_R\rangle
      =\langle T,\theta_R\rangle\\
   &\le |\langle T,\theta_R\rangle|
      \le C'(1+R)^D.
 \end{aligned}
 \tag{MA-FQ.50}
\]
For clarity, every seminorm
$\sup_x|x^r\partial_x^\ell\theta_R(x)|$ is
$R^{r-\ell}\sup_y|y^r\theta^{(\ell)}(y)|$;
the finitely many such powers are bounded by one
fixed power of $1+R$. Positivity gives the first
inequality, and the middle pairing is real and
nonnegative. Dyadic summation now proves that
$M_{\rm good}$ is a tempered measure. Density of
$C_c^\infty$ in the Schwartz space extends
(MA-FQ.48) to an equality of tempered distributions.
No Tauberian theorem has entered this step.

\emph{Atomic positivity.}
Define the bounded continuous function and the tempered
distribution
\[
 q(t)=-\frac1{2\pi}\int e^{-itx}\,dB_0(x),
                 \qquad \nu=-\mu+q(t)\,dt.
 \tag{MA-FQ.51}
\]
Finite total variation of $B_0$ and Fourier inversion
give $\widehat{q(t)\,dt}=-B_0$, hence
$\widehat\nu=M_{\rm good}$.
For $\widetilde\psi(t)=\overline{\psi(-t)}$,
Fourier inversion and Fubini yield
\[
 \langle\nu,\psi*\widetilde\psi\rangle
      =\frac1{2\pi}\int
                    |\widehat\psi(-x)|^2\,dM_{\rm good}(x)
      \ge0\qquad(\psi\in C_c^\infty).
 \tag{MA-FQ.52}
\]
The integral is finite by the polynomial growth just proved.

We recall the precise atomic extraction so that it does not
depend on the $S_5$ specialization of MA-FQ-001.
Choose a real even $\phi\in C_c^\infty$ of $L^2$ norm one.
Let $\Psi=\phi*\widetilde\phi$, so $\Psi(0)=1$, and put
$\phi_\epsilon(t)=\epsilon^{-1/2}\phi(t/\epsilon)$.
Then $\phi_\epsilon*\widetilde\phi_\epsilon(t)
=\Psi(t/\epsilon)$.
For any finite $x_i\in\mathbb R$ and $c_i\in\mathbb C$,
substitute
$\psi_\epsilon=\sum_i c_i\phi_\epsilon(\,\cdot-x_i)$
in (MA-FQ.52).
Local finiteness of the atomic support isolates the atom
at each of the finitely many differences $x_i-x_j$
when $\epsilon$ is small. Each continuous-density term
has modulus at most
$\epsilon\|q\|_\infty\|\Psi\|_1$.
Letting $\epsilon$ tend to zero consequently gives
\[
 \sum_{i,j}c_i\overline{c_j}\,a(x_i-x_j)\ge0,
                     \qquad a(t)=-m(t).
 \tag{MA-FQ.53}
\]
Thus $a$ is positive definite on $\mathbb R_d$,
the additive real group with the discrete topology.
The two-by-two inequality gives
$|a(t)|\le a(0)=-m(0)$; this conclusion has not
assumed a sign or a prescribed value for $m(0)$.
If $a(0)=0$, the same inequality gives $a=0$.

\emph{Only finitely many primes can occur.}
The measure $\sum_t a(t)\delta_t=-\mu$ is tempered,
its coefficient function is integer valued and
positive definite on $\mathbb R_d$, and its
ordinary support is locally finite. Its ordinary
Fourier transform in (MA-FQ.48) is a locally finite
purely atomic measure. These are exactly the
hypotheses of MA-FQ-003. That theorem, proved from
Cohen's idempotent theorem, the real-line coset
classification and ordinary Poisson summation,
therefore gives a finite set of $h_\ell>0$ such that
\[
 \operatorname{supp}\widehat{(-\mu)}
                    \subseteq\bigcup_{\ell=1}^L h_\ell\mathbb Z.
 \tag{MA-FQ.54}
\]
The compact-dual representing measure on
$\widehat{\mathbb R_d}$ used in that proof is not
identified with this ordinary Fourier transform.
The empty union is permitted.

Call a good prime $p$ active when $f(p^j)>0$ for
at least one $j\ge1$. At $x=j\log p$ such a
coefficient is an atom of $\widehat{(-\mu)}$.
No other prime power has the same integer value
$p^j$, no bad-prime progression meets it, and the
conductor atom is at zero. Thus the atom cannot
cancel. Choose one such exponent $j_p$ for each
active prime. By (MA-FQ.54), $j_p\log p$ lies in
one of the finitely many frequency lattices.

Each one of those lattices can contain such
frequencies from at most one distinct prime,
even though the exponents $j_p$ may vary.
Indeed, if
$j_p\log p=nh$ and $j_q\log q=mh$ with
$m,n,j_p,j_q$ positive integers, then
$p^{mj_p}=q^{nj_q}$. Unique factorization forces
$p=q$. There are therefore only finitely many
active good primes.

\emph{Finite prime support forces the whole datum to vanish.}
Adjoin the finitely many bad primes to the active good
primes. The resulting finite set $\mathcal P$ contains
every prime whose powers support a nonzero coefficient
$f(n)$ for $n\ge2$. By (MA-FQ.46),
\[
                \sum_{n\ge2}|f(n)|
                =\sum_{p\in\mathcal P}\sum_{j\ge1}|f(p^j)|
                <\infty.
 \tag{MA-FQ.55}
\]
Consequently the complete Fourier measure
\[
 D=2f(1)\delta_0+
       \sum_{n\ge2}\bigl(f(n)\delta_{\log n}
                    +\overline{f(n)}\delta_{-\log n}\bigr)
          =\widehat{(-\mu)}
 \tag{MA-FQ.56}
\]
has finite total variation. Its inverse Fourier
transform is the bounded continuous density
\[
               r(t)=\frac1{2\pi}\int e^{-itx}\,dD(x).
 \tag{MA-FQ.57}
\]
Fourier inversion on Schwartz tests gives the
equality of tempered distributions $-\mu=r(t)\,dt$.
Both sides are locally finite measures.
The left side is purely atomic and the right
side absolutely continuous. Thus every atomic
mass is zero, so $\mu=0$ and $m=0$.
Equivalently uniqueness of the local Lebesgue
decomposition gives the same conclusion.
Taking Fourier transforms gives $D=0$.
The atom at zero is $2f(1)$, and the atom at
$\log n>0$ is exactly $f(n)$; these locations
are distinct. Hence $f(1)=0$ and $f(n)=0$
for every $n\ge2$. With $K=0$ this proves
(MA-FQ.41). The defining exponential series for
the associated $L$-function is then identically
one in $\operatorname{Re}s>1$.
\end{proof}

\paragraph{The finite nonintegral exceptions allowed by (A3).}
The same vanishing conclusion holds if the extra requirement
$m(t)\in\mathbb Z$ in (MA-FQ.40) is omitted and one retains
exactly Booker's axiom (A3), which allows only finitely many
points with nonintegral mass. Here is the complete modification.

First, the integral nonzero masses have absolute value at least
one, so (MA-FQ.43) bounds their number in every compact interval.
Adjoining the finitely many exceptional points preserves local
finiteness. The measure is tempered by the same dyadic estimate.
Every step through atomic positivity (MA-FQ.53) therefore applies
without change. Its two-by-two inequality bounds every value
of $a=-m$ by the finite real number $a(0)$. The integral values
in this bounded interval form a finite set, and the finitely
many exceptional points add only finitely many other values.
Thus $a$ has finite real range.

Bochner's theorem on $\mathbb R_d$, in the explicit
GNS--Riesz form proved in MA-FQ-001, gives
$a\in B(\mathbb R_d)$; that proof only uses positive
definiteness and the finite mass $a(0)$, not its integrality.
For each nonzero value $r$ in the finite range, choose a
polynomial taking value one at $r$ and zero at every other
range value. Applied to $a$, it gives
$1_{\{a=r\}}\in B(\mathbb R_d)$.
Cohen's idempotent theorem places this level set in the
coset ring. It is locally finite because the support of
$a$ is. The set assertion of MA-FQ-002, followed by
inclusion--exclusion as in that proof, expresses its
indicator as a finite integer linear combination of
ordinary lattice-coset and singleton indicators.
Multiplication by $r$ and summation over the finite
nonzero range expresses the atomic measure of $a$
as a finite \emph{real} linear combination of translated
lattice combs and singleton masses.

The ordinary Poisson formula (MA-FQ.24) applies to those
real coefficients just as it does to integer coefficients.
Its singleton terms give a continuous trigonometric-polynomial
density. Since the actual Fourier transform (MA-FQ.48) is
purely atomic, uniqueness of the local Lebesgue decomposition
eliminates that density. Its remaining support lies in
finitely many unshifted frequency lattices, exactly as in
(MA-FQ.54). The varying-prime-power argument, finite total
variation, inverse continuous density and atomwise vanishing
in the last two parts of the main proof now apply verbatim.
They give $\mathcal F=0$. This proves the extension using
precisely the finite-exception clause in (A3); the integral
case remains the stated main theorem.

\paragraph{Relation to the $S_5$ gate and to the scope tests.}
For the actual difference datum of MA-LD-001,
the gamma calculation in MA-FQ-001 gives $K_R=0$,
the signed orders are integral, and its good
prime-power coefficients are nonnegative.
The required real support has a shorter direct
proof from MT-S508, without MA-S5-004.
In its paragraph \emph{The positive Euler quotient
controls a whole half-plane}, the proof following
(S508.13) uses (S508.9), the positive Dirichlet
series and Landau's theorem.  Only the sufficient
conclusion $\gamma\le1/2$ is needed here.
Indeed the source first gives $\gamma<1$.
If $1/2<\gamma<1$, Landau's singularity at the real
abscissa would be a pole of the meromorphic quotient
at a point where $(\mathcal R_B)$ makes its
denominator nonzero, a contradiction.
Any case $\gamma<1/2$, including $\gamma\le0$,
already gives the required half-plane convergence;
no equality assertion about that abscissa is needed.
A finite $\Omega$ also causes no convergence problem.
For every $\sigma>1/2$, the convergent positive
Dirichlet series has prime coefficients equal to
four, so $\sum_{p\in\Omega}p^{-\sigma}<\infty$.
Its Euler logarithm therefore converges absolutely
and locally uniformly in $\operatorname{Re}s>1/2$.
The finitely many bad-prime logarithms converge
there by (A1), as in MA-LD-001. Thus the actual
finite quotient $A/L_W$ is holomorphic and
nonvanishing throughout that half-plane.

The exactly equal gamma factors cancel, and
the conductor exponential is nowhere zero.
Consequently the completed quotient also has
no divisor on $\operatorname{Re}s>1/2$,
including at $s=1$. Its completed functional
equation reflects the full divisor across
$\operatorname{Re}s=1/2$, so there is no divisor
on the open left half-plane either, including
at $s=0$. This verifies real height support
with no omitted endpoint poles. MA-FQ-006 now
applies directly. Its vanishing includes every
bad-prime logarithmic coefficient and the
conductor atom; equality of the local logarithmic
series implies equality of the normalized
standard Euler factors at every finite prime.
The archimedean factors were already equal.
It also gives $m_R(0)=-k=0$.

Thus MA-S5-004 and the original MA-LD-001
proof remain valid historical deductions, but
their separate Tauberian exceptional-prime
estimate is not needed for this final closure.
Temperedness of the positive prime measure
comes directly from (MA-FQ.50).
The hypothesis $(\mathcal R_B)$ remains an
explicit premise and is not proved here.

The exact hypotheses explain the comparisons
in MA-FQ-005. The ratio $\zeta/L(\,\cdot,\chi_5)$
has completed poles off the critical line;
the equal-parity Dirichlet ratio there has
infinitely many negative good-prime coefficients;
the single-prime central Euler pole violates
(A1); and the stated nonperiodic positive
Fourier crystal has mixed logarithmic
frequencies outside the prime-power support.
Those checked examples therefore do not meet
(MA-FQ.40).

\paragraph{Primary inputs and attribution.}
The exact axioms used are Booker (2015),
\emph{L-functions as distributions},
Definition~1.3(A1), (A3), (A4), author p.~3
of \path{arXiv:1308.3067v2}.
The finite-prime-support conclusion also
satisfies the hypothesis of Booker's
Theorem~1.6(ii), author p.~5, because
(MA-FQ.55) implies
$\sum_{n\ge2}|f(n)|/\log n<\infty$.
The last part of the proof above supplies
the needed vanishing directly in the
real-supported, zero-kernel case.
The deep harmonic-analysis input remains
Cohen's idempotent theorem in the precise
form checked in MA-FQ-002--003:
Sanders (2020), \S12, Theorems~12.1 and~12.3,
and the real-line locally finite coset
argument supplied in MA-FQ-002, with
Schreiber (1970), Lemma~1.4 and
Corollary~1.5, p.~807, Example~1.6,
p.~808, as its exact primary locator.
The criterion is a consequence of those
inputs and the explicit-formula axioms;
no claim of priority over known rigidity
results is made.

Booker, A. R. (2015). $L$-functions as distributions.
\emph{Mathematische Annalen, 363}, 423--454.
\url{https://doi.org/10.1007/s00208-015-1178-z}.

Sanders, T. (2020). Bounds in Cohen's idempotent theorem.
\emph{Journal of Fourier Analysis and Applications, 26},
Article~25.
\url{https://doi.org/10.1007/s00041-020-09732-y}.

Schreiber, B. M. (1970). On the coset ring and strong Ditkin sets.
\emph{Pacific Journal of Mathematics, 32}(3), 805--812.
\url{https://doi.org/10.2140/pjm.1970.32.805}.

\section{A discriminant-dependent part of the real Hecke gate}
\label{sec:s5-real-zero-free}

This section proves a part of $(\mathcal R_B)$ using the specified
Dirichlet theorem $\mathsf A$, without the general Hecke hypothesis
$M$.  It does not prove $(\mathcal R_B)$ on the remaining central
interval.  The notation is that of MA-S5:
$L/\mathbb Q$ is the specified $S_5$ extension,
$E=L^{A_5}$, $F=L^{F_{20}}$, $H=EF$,
$C=L(s,\eta)$, and $B=\zeta_H/\zeta_F=C L_W$.

\paragraph{MA-RB-001 (an explicit endpoint exclusion).}
Put
\[
 \Delta_L=|\operatorname{disc}L|,\qquad
 \mathcal H_L=1+\tfrac12\log\Delta_L,\qquad
 \kappa_5=(\sqrt5-1)^2,\qquad
 \delta_L=\min\{\tfrac12,\kappa_5/\mathcal H_L\}.
 \tag{MA-RB.1}
\]
Under $\mathsf A$, the quadratic Hecke function $B$ has no real zero in
\[
              (0,\delta_L)\ \cup\ (1-\delta_L,1).
 \tag{MA-RB.2}
\]
More precisely, if $1/2<\sigma_0<1$ and $B(\sigma_0)=0$, then
\[
        1-\sigma_0\ge\kappa_5/\mathcal H_L.
 \tag{MA-RB.3}
\]
The endpoints in (MA-RB.2) are open.  No assertion about a possible
central zero, or about any remaining real point, is included.

\begin{proof}
\emph{The signed Artin orders at the particular point.}
Fix a real $\sigma_0\in(0,1)\setminus\{1/2\}$ and suppose that
$r=\operatorname{ord}_{\sigma_0}B\ge1$.
The theorem $\mathsf A$ gives $C(\sigma_0)\ne0$.
Consequently the all-prime identity $B=C L_W$ gives
\[
                 \operatorname{ord}_{\sigma_0}L_W=r.
 \tag{MA-RB.4}
\]
This is a pointwise deduction from a quotient of holomorphic
functions with nonzero denominator, and does not assume global
holomorphy of the unknown five-dimensional Artin factor.

The other five-dimensional factor is also holomorphic at this
particular real point:
\[
                   L_{W\eta}=\zeta_F/\zeta_{\mathbb Q}.
 \tag{MA-RB.5}
\]
Indeed $\zeta_F$ is holomorphic on $(0,1)$, and
$\zeta_{\mathbb Q}(\sigma)<0$ there.  For the latter sign, the
convergent alternating Dirichlet series is strictly positive
for real $\sigma>0$, while its factor
$1-2^{1-\sigma}$ is negative when $0<\sigma<1$.
Calegari's Theorem~1.1 supplies holomorphy of the irreducible
four- and six-dimensional factors $L_V,L_{V\eta},L_X$.
The regular-representation factorization, including every
ramified Euler factor, is
\[
 \zeta_L
   =\zeta_{\mathbb Q}\,C\,
       L_V^4L_{V\eta}^4L_W^5L_{W\eta}^5L_X^6.
 \tag{MA-RB.6}
\]
Taking signed orders at $\sigma_0$, the first two orders are
zero and every other order on the right is nonnegative.
It follows that
\[
                 \operatorname{ord}_{\sigma_0}\zeta_L
                         \ge5r\ge5.
 \tag{MA-RB.7}
\]
In particular no unproved holomorphy of either
five-dimensional factor has entered this inequality.
Equations (MA-RB.4)--(MA-RB.6) use the already checked
MT-S506 (S506.6), MT-S508 (S508.2), and the regular-character
identity in MT-S507 (S507.3).

\emph{Completion, signature, and the exact Hadamard identity.}
Because $L/\mathbb Q$ is Galois and its complex conjugation is the
nonidentity element of type $(12)(34)$, no embedding of $L$ is
real.  Its degree is $|S_5|=120$, so its signature is
$(r_1,r_2)=(0,60)$.  With
\[
 \Gamma_{\mathbb R}(s)=\pi^{-s/2}\Gamma(s/2),
 \qquad
 \Gamma_{\mathbb C}(s)=2(2\pi)^{-s}\Gamma(s),
\]
take the completed entire function
\[
 \xi_L(s)
   =s(s-1)\Delta_L^{s/2}\Gamma_{\mathbb C}(s)^{60}\zeta_L(s).
 \tag{MA-RB.8}
\]
The multiplicative constant in this normalization is immaterial.
The usual Dedekind continuation and functional equation give
$\xi_L(s)=\xi_L(1-s)$ and
$\overline{\xi_L(\overline s)}=\xi_L(s)$.
The simple pole and nonzero residue of $\zeta_L$ at $1$ imply
$\xi_L(1)\ne0$; the functional equation gives $\xi_L(0)\ne0$.
The zeros of $\xi_L$ are the nontrivial zeros, with
$0\le\operatorname{Re}\rho\le1$, including their multiplicities.
Only this containing strip is needed.

Here are the needed growth and normalization details of the
Hadamard formula.  On $\operatorname{Re}s\ge2$ the Euler product
bounds $\zeta_L(s)$, and the gamma estimate gives
$|\xi_L(s)|\le\exp(C|s|\log(|s|+3))$.
The functional equation gives the same bound on
$\operatorname{Re}s\le-1$.
On the remaining fixed strip, the unconditional polynomial
bound for the finite Dedekind function, with its pole removed,
and Stirling's formula give the same bound.
These are the fixed-strip inputs explicitly established in
NT-Z10 and recalled in MT-M18.
Thus $\xi_L$ is entire of order at most one.
The same local zero count gives
$\sum_{|\rho|>1}|\rho|^{-2}<\infty$.
Hadamard factorization therefore has the form
\[
 \xi_L(s)=\xi_L(0)e^{a_Ls}
        \prod_\rho(1-s/\rho)e^{s/\rho},
 \qquad
 \frac{\xi_L'}{\xi_L}(s)
   =a_L+\sum_\rho\left(\frac1{s-\rho}+\frac1\rho\right).
 \tag{MA-RB.9}
\]
The series in the second formula converges locally normally
off the zeros.  The constant
$a_L=\xi_L'(0)/\xi_L(0)$ is real.
Moreover
$S_L=\sum_\rho\operatorname{Re}(1/\rho)$ converges absolutely:
for large $|\rho|$ each term has absolute value at most
$|\rho|^{-2}$, by the containing strip.
Differentiating the functional equation at $0$ and $1$ gives
$\xi_L'(1)/\xi_L(1)=-a_L$.
Taking real parts of (MA-RB.9) at $1$ and using the
multiset symmetry $\rho\mapsto1-\rho$ gives
$-a_L=a_L+2S_L$.
Hence $a_L=-S_L$.  Taking real parts at a real $\sigma>1$
now proves, with its constant fixed rather than suppressed,
\[
 \frac{\xi_L'}{\xi_L}(\sigma)
             =\sum_\rho\operatorname{Re}\frac1{\sigma-\rho}.
 \tag{MA-RB.10}
\]
All summands are nonnegative.  The displayed real series
converges absolutely, again by the containing strip and the
zero count.  This proves the particular Hadamard identity
used here from the standard factorization theorem.

\emph{The archimedean estimate with all signs retained.}
Logarithmically differentiating (MA-RB.8) on the real
half-line $\sigma>1$ gives the exact formula
\[
 \sum_\rho\operatorname{Re}\frac1{\sigma-\rho}
 =\frac1{\sigma}+\frac1{\sigma-1}
       +\frac12\log\Delta_L
       +60\bigl(\psi(\sigma)-\log(2\pi)\bigr)
       +\frac{\zeta_L'}{\zeta_L}(\sigma),
 \tag{MA-RB.11}
\]
where $\psi=\Gamma'/\Gamma$.
The last term is negative, since the absolutely convergent
Dedekind Euler product gives
\[
 -\frac{\zeta_L'}{\zeta_L}(\sigma)
    =\sum_{\mathfrak p}\sum_{j\ge1}
          \frac{\log N\mathfrak p}{(N\mathfrak p)^{j\sigma}}>0.
 \tag{MA-RB.12}
\]
For completeness, $\psi(x)\le\log x$ for $x>0$ follows
directly from the gamma integral.  The density
$u^{x-1}e^{-u}/\Gamma(x)$ on $(0,\infty)$ has expectation
of $u$ equal to $x$, and expectation of $\log u$ equal
to $\psi(x)$; Jensen's inequality for the concave logarithm
gives the claim.  Differentiation of the gamma integral is
justified on compact positive $x$ intervals by integrable
majorants at zero and infinity.
Therefore, when $1<\sigma\le2$,
\[
                \psi(\sigma)-\log(2\pi)
                    \le\log\bigl(\sigma/(2\pi)\bigr)<0.
\]
Writing $\sigma=1+t$, $0<t\le1$, and dropping only
nonpositive terms in (MA-RB.11), we obtain
\[
        \sum_\rho\operatorname{Re}\frac1{1+t-\rho}
                     \le \frac1t+\mathcal H_L.
 \tag{MA-RB.13}
\]

\emph{Optimization and reflection.}
Suppose now that $1/2<\sigma_0<1$ is a zero of $B$, and put
$d=1-\sigma_0$.  Its at least five occurrences in the
zero divisor of $\xi_L$ contribute at least $5/(t+d)$
to (MA-RB.13).  Consequently
\[
                  \frac5{t+d}\le\frac1t+\mathcal H_L
                      \qquad(0<t\le1).
 \tag{MA-RB.14}
\]
Here $\Delta_L>1$ follows even from the displayed formulas:
at $\sigma=2$, (MA-RB.11)--(MA-RB.12) and the gamma bound give
$0\le3/2+\frac12\log\Delta_L-60\log\pi$, which excludes
$\Delta_L=1$.  Since the discriminant is an integer,
$\mathcal H_L\ge1+\frac12\log2>\sqrt5-1$.
Thus $t=(\sqrt5-1)/\mathcal H_L$ is permitted.
Substitution in (MA-RB.14) gives
$d\ge(\sqrt5-1)^2/\mathcal H_L$, proving (MA-RB.3).

To reflect the result, the quadratic Hecke completion of $B$
has root number $+1$, as the quotient of the completed
Dedekind equations for $H$ and $F$.
Its conductor factor is
$(|\operatorname{disc}H|/|\operatorname{disc}F|)^{s/2}$,
and its archimedean factor is
\[
                \Gamma_{\mathbb R}(s)^4
                    \Gamma_{\mathbb R}(s+1)^2.
 \tag{MA-RB.15}
\]
Indeed the settled permutation character
$\operatorname{Ind}_{F_{20}}^{S_5}\mathbf1=\mathbf1+W\eta$
has value two at complex conjugation, giving signature
$(2,2)$ for $F$; since $E$ is real, the signature of $H$
is $(4,4)$.  Taking the ratio of their gamma factors yields
(MA-RB.15), equivalently $\Gamma_{\mathbb R}(s)^2
\Gamma_{\mathbb C}(s)^2$.
All these factors are finite and nonzero on $(0,1)$.
The completed functional equation therefore pairs a real
zero $\sigma_0$ with $1-\sigma_0$, of the same order.
Reflecting the upper exclusion and using the definition
of $\delta_L$ proves (MA-RB.2).
\end{proof}

\paragraph{The exact remaining obstruction.}
The proof uses $\mathsf A$ only to make the quadratic Dirichlet
factor $C$ nonzero at the specified off-center point.
It supplies no nonvanishing assertion for the quadratic Hecke
function $B$ over the nonabelian sextic field $F$ in
$[\delta_L,1-\delta_L]\setminus\{1/2\}$.
The fact that its character is pulled back from $\eta$ does
not turn its $L$-function over $F$ into a Dirichlet
$L$-function over $\mathbb Q$: induction gives
$\operatorname{Ind}_{G_F}^{G_{\mathbb Q}}\eta_F=\eta\oplus W$,
which is precisely the unresolved nonabelian factorization.

The simple-real-zero descent theorem does not fill that interval.
Heilbronn (1973), Theorem~1, printed p.~870, says that a simple
real zero of the Dedekind zeta function of a normal extension
descends to the compositum of its quadratic subfields.
For $L/\mathbb Q$ this compositum is $E$, and $\mathsf A$
excludes such a zero of $\zeta_E$ off the center.
But a zero of $B$ in the present argument already has
multiplicity at least five in $\zeta_L$, by (MA-RB.7);
it is outside that theorem's simplicity hypothesis.

There is no contradiction from the signed order identities alone.
For example, at a hypothetical reflected pair of real points one
may assign
\[
 (\operatorname{ord}\zeta_{\mathbb Q},
   \operatorname{ord}C,\operatorname{ord}L_V,
   \operatorname{ord}L_{V\eta},\operatorname{ord}L_W,
   \operatorname{ord}L_{W\eta},\operatorname{ord}L_X)
                       =(0,0,0,0,1,0,0).
 \tag{MA-RB.16}
\]
Every induced-character order is then nonnegative, because
it is a nonnegative multiplicity times the sole nonzero
irreducible order.  The regular order is five and the
quadratic-subfield order is zero.  Assigning the same data
at the reflected point is compatible with functional-equation
orders.  This is only a formal compatibility test; it neither
constructs an arithmetic zero nor proves that further analytic
arguments cannot exclude it.

\paragraph{Sources, input class, and motive interpretation.}
The modular holomorphy input is Calegari (2013), Theorem~1.1,
publisher-PDF p.~2, and the exact all-prime induction identities
are the previously checked MT-S506--MT-S508 locators above.
The classical completed Hecke continuation and functional
equation are Ramakrishnan and Valenza (1999), \S7.4,
Theorem~7-19, printed p.~279, and \S7.7, Theorem~7-28,
printed p.~289, as recorded in MT-M18; NT-Z10 supplies the
fixed-strip growth and zero count.  The necessary Hadamard
normalization, archimedean estimate, and optimization are
proved in full here.  No new zero-free-region theorem is an
input, and no priority claim is made for the classical
logarithmic-derivative method.  The input class of MA-RB-001
is $\mathsf A$, not $M$.

The regular-character identity in (MA-RB.6) is the
$L$-function identity of the regular Artin motive attached
to the finite \'etale cover $\operatorname{Spec}L\to
\operatorname{Spec}\mathbb Q$.  Its decomposition with
multiplicities equal to irreducible dimensions explains
the coefficient five in (MA-RB.7).
The exact available motive locators remain Deligne (1979),
Definition~6.1 and \S\S6.2--6.3, and SGA~1,
Expos\'e~V, \S7, already checked in MA-S5.
This interpretation supplies the factorization, not the
analytic inequality or a number-field purity assertion.

Heilbronn, H. (1973). On real zeros of Dedekind $\zeta$-functions.
\emph{Canadian Journal of Mathematics, 25}(4), 870--873.
\url{https://doi.org/10.4153/CJM-1973-090-3}.

\section{The remaining real comparison obstruction before \texorpdfstring{\((\mathcal R_B)\)}{the real-zero hypothesis}}
\label{sec:s5-abscissa-reduction}

The Fourier-rigidity theorem can be used before assuming
\((\mathcal R_B)\) to exclude an entire convergence branch.
It does not, by itself, put the divisor of a remaining comparison
on the critical line.  We prove a sharper unconditional reduction
of that remaining branch and identify precisely what its boundary
positivity supplies.

Retain the specified Calegari family and the notation
\(A=L_f(\,\cdot\,,\varpi)\), \(C=L(\,\cdot\,,\eta)\),
\(B=C L_W\), \(L_V,L_{V\eta},L_X\) from MA-S5.
The gamma factors of \(A\) and \(L_W\) are exactly equal,
as proved in MA-FQ-001.  Write
\[
 R(s)=\frac{A(s)}{L_W(s)}
       =U_{\mathcal B}(s)D_\Omega(s),\qquad
 D_\Omega(s)=\prod_{p\in\Omega}
                  \left(\frac{1+p^{-s}}{1-p^{-s}}\right)^2 .
 \tag{MA-RB.17}
\]
The finite factor \(U_{\mathcal B}\) is holomorphic and nonzero
on \(\operatorname{Re}s>0\).  Indeed every finite local parameter
in this comparison has finite image and zero monodromy by
the checked finite-base-change argument in MA-S5-003;
its local roots have modulus at most one.
The omitted archimedean ratio is one.
Conductor powers in the fully completed quotient are
nowhere-zero exponentials, so they change none of its orders.

\paragraph{MA-RB-002 (a shared real zero and multiplicity at least nine).}
If the specified all-place comparison does not hold, then
\(\Omega\) is infinite and the nonnegative Dirichlet series of
\(D_\Omega\) has a finite abscissa \(\gamma\) with
\[
                 \tfrac12<\gamma<1.
 \tag{MA-RB.18}
\]
It has a real pole there.  Put
\[
 \begin{gathered}
 k=-\operatorname{ord}_\gamma R
     =-\operatorname{ord}_\gamma D_\Omega\in\mathbb Z_{\ge1},
 \qquad n=\operatorname{ord}_\gamma A,\quad
 b=\operatorname{ord}_\gamma C,\\
 v_0=\operatorname{ord}_\gamma L_V,\quad
 v_1=\operatorname{ord}_\gamma L_{V\eta},\quad
 e=\operatorname{ord}_\gamma L_X .
 \end{gathered}
 \tag{MA-RB.19}
\]
Then
\[
 \begin{aligned}
 \operatorname{ord}_\gamma L_W&=n+k,\\
 \operatorname{ord}_\gamma B&=b+n+k,\\
 v_1+e&\ge k,\\
 \operatorname{ord}_\gamma\zeta_L&\ge b+5n+9k\ge9k .
 \end{aligned}
 \tag{MA-RB.20}
\]
Thus \(B\) and \(L_{V\eta}L_X\) have a common real zero at
\(\gamma\).  With \(\mathcal H_L=1+\frac12\log\Delta_L\)
as in MA-RB-001, one has the sharper bound
\[
       1-\gamma\ \ge\
                \frac{(3\sqrt{k}-1)^2}{\mathcal H_L}
                       \ \ge\ \frac4{\mathcal H_L}.
 \tag{MA-RB.21}
\]
These conclusions use neither \(\mathsf A\), \(M\), nor
\((\mathcal R_B)\).

In particular, absence of a common real zero of
\(B\) and \(L_{V\eta}L_X\) in
\[
                 \left(\tfrac12,\,1-\frac4{\mathcal H_L}\right]
 \tag{MA-RB.22}
\]
suffices for the exact comparison.
An empty interval satisfies that sufficient condition.
This is a comparison criterion; it is not a proof of
\((\mathcal R_B)\) or a zero-free assertion about every
zero of \(B\).

\begin{proof}
\emph{The convergence branch excluded by Fourier rigidity.}
The classical finite-\(\Omega\) comparison, Calegari's
Lemma~4.10, gives exact equality at all places.
Consequently failure implies \(\Omega\) is infinite.
Each Euler factor in (MA-RB.17) has nonnegative Taylor
coefficients, and its coefficient of \(p^{-s}\) is four.
Their product therefore gives a nonnegative Dirichlet
series in the absolute-convergence half-plane.

The previously checked proof of Calegari's Theorem~4.7
gives \(\sum_{p\in\Omega}p^{-1}<\infty\).
The Euler product consequently has a finite positive
limit at \(s=1\) from the right.
Its meromorphic continuation is (MA-RB.17).
If its abscissa were one, Landau's theorem would force
a singularity at one; meromorphy and that finite positive
limit make the point regular instead.
Thus its abscissa is less than one.
It is finite below: any one prime in \(\Omega\) already
has a nonnegative local series diverging on the real
half-line \(s\le0\).
Hence \(\gamma\ge0\).

Suppose \(\gamma\le1/2\).
For each real \(\sigma>1/2\), convergence of the positive
Dirichlet series and its prime coefficient four give
\(\sum_{p\in\Omega}p^{-\sigma}<\infty\).
The Euler logarithm then converges absolutely and locally
uniformly on \(\operatorname{Re}s>1/2\).
Together with the finite local unit this makes \(R\)
holomorphic and nonzero throughout that half-plane.

The genuine comparison \(L\)-datum, its integral signed
orders, its exact zero gamma kernel, its prime-power
support and eventual nonnegative good-prime coefficients
were constructed in MA-FQ-001/006 without using
\((\mathcal R_B)\) for those properties.
That hypothesis was used only to obtain real height
support.  Here the same support follows from the just-proved
half-plane nonvanishing and the completed functional
equation.  It reflects the divisor to the other open
half-plane; the conductor exponential has no divisor.
This includes both completed endpoints.
Thus MA-FQ-006 applies and the entire comparison datum
is zero.  Its good-prime coefficient
\(4(\log p)p^{-1/2}\) at every \(p\in\Omega\) then forces
\(\Omega\) to be empty, a contradiction.
We have proved (MA-RB.18), without a real-zero premise.

Landau's theorem at the now finite abscissa \(\gamma\)
forces a singularity of \(D_\Omega\).
Its meromorphic continuation makes that singularity
a pole.  Its positive integer order is \(k\) in
(MA-RB.19).  The finite local unit in (MA-RB.17) does
not change that order.

\emph{The symmetric-square constraint at this actual point.}
Every finite local unit and gamma unit in the proofs of
MA-S5-003 and MA-S5-006 remains holomorphic and nonzero
at any real point of \((0,1)\), not only at \(1/2\):
\(|p^{-\sigma}|<1\) excludes a finite local zero, while
the real gamma arguments there are positive.
The already proved meromorphic tensor identities therefore
give their pointwise order identities at \(\gamma\).

For clarity let \(\Pi\) and \(\nu\) be the four-dimensional
cuspidal representation and finite-order similitude
character in MA-S5-006.  If \(r_4\) is the order at \(\gamma\)
of its completed Rankin--Selberg function, then
\[
 \operatorname{End}\xi=\mathbf1\oplus V\eta\oplus W\oplus X,
 \qquad
 r_4=v_1+\operatorname{ord}_\gamma L_W+e-2k .
 \tag{MA-RB.23}
\]
The real function \(\zeta_{\mathbb Q}\) is nonzero on
\((0,1)\), by the alternating-series argument already
given in MA-RB-001.
Equation (MA-RB.17) gives
\(\operatorname{ord}_\gamma L_W=n+k\).
The exterior-square identity and tensor decomposition give
\[
 L^S(s,\Pi\times\Pi^\vee)
  =L^S(s,\Pi,\operatorname{Sym}^2\otimes\nu^{-1})
                               \zeta_{\mathbb Q}^S(s)A^S(s).
 \tag{MA-RB.24}
\]
The finite deleted factors are units at \(\gamma\).
The central character of \(\Pi\) is unitary and
\(\nu\) is finite order, as checked in MA-S5-006.

Takeda's published Theorem~7.1 states that this incomplete
twisted symmetric-square function is holomorphic everywhere
except for possible simple poles at zero and one.
In particular, it is holomorphic at the present
\(\gamma\in(1/2,1)\).
If \(q_\gamma\ge0\) is its order there, (MA-RB.23)--(MA-RB.24)
give
\[
       r_4=n+v_1+e-k=q_\gamma+n,\qquad
                         q_\gamma=v_1+e-k\ge0.
 \tag{MA-RB.25}
\]
This is an extension of the order deduction to a new real
point, not a new assertion of symmetric-square nonvanishing.
It proves the third line of (MA-RB.20).

The first two lines of (MA-RB.20) follow from \(R=A/L_W\)
and \(B=C L_W\), with signed orders retained.
Here \(n,b,v_0,v_1,e\ge0\) by standard Hecke holomorphy
and Calegari's established four- and six-dimensional
automorphy.  Also
\[
 L_{W\eta}=\zeta_F/\zeta_{\mathbb Q}
\]
is holomorphic at this particular real point; denote its
nonnegative order by \(x_1\).
Thus the actual regular-character factorization gives
\[
 \begin{aligned}
 \operatorname{ord}_\gamma\zeta_L
    &=b+4v_0+4v_1+5(n+k)+5x_1+6e\\
    &\ge b+5n+5k+4(v_1+e)
     \ge b+5n+9k .
 \end{aligned}
 \tag{MA-RB.26}
\]
No general holomorphy of the unknown \(L_W\) is being
assumed: its positive order \(n+k\) at this point is
already forced by the actual comparison pole.
Nor was \(C(\gamma)\ne0\) needed.  Its possibly positive
order \(b\) only strengthens (MA-RB.26).

\emph{The discriminant bound.}
The completed Dedekind logarithmic-derivative identity
and estimate (MA-RB.8)--(MA-RB.13) were proved from
unconditional analytic inputs; only the earlier deduction
of multiplicity from an arbitrary \(B\)-zero used
\(\mathsf A\).
They therefore apply here without that theorem.
Set \(\delta=1-\gamma\), so \(0<\delta<1/2\).
The at least \(9k\) copies of this real zero give
\[
                 \frac{9k}{t+\delta}
                     \le\frac1t+\mathcal H_L
                                  \qquad(0<t\le1).
 \tag{MA-RB.27}
\]
The optimizing value \(t=(3\sqrt{k}-1)/\mathcal H_L\)
is in the permitted interval.
Indeed \(t=1\) in (MA-RB.27) gives
\[
        9k\le(1+\delta)(1+\mathcal H_L)
                    <\tfrac32(1+\mathcal H_L).
\]
Since \(\mathcal H_L\ge1\), this implies
\(3\sqrt{k}-1<\mathcal H_L\).
At that optimizing value, solving (MA-RB.27) for
\(\delta\) yields exactly the first inequality in
(MA-RB.21).  Since \(k\ge1\), its numerator is at least
four.  This proves the remaining assertions, including
the shared-zero criterion (MA-RB.22).
\end{proof}

\paragraph{MA-RB-003 (what positivity on the rightmost line actually proves).}
In the nontrivial branch of MA-RB-002 define
\[
                    a_\gamma(t)=
                       -\operatorname{ord}_{\gamma+it}D_\Omega,
                              \qquad t\in\mathbb R .
 \tag{MA-RB.28}
\]
Then \(a_\gamma\) is an integer-valued positive-definite
function on the additive group \(\mathbb R\) equipped with
the discrete topology.  It satisfies
\[
        a_\gamma(0)=k,\qquad |a_\gamma(t)|\le k,
                         \qquad a_\gamma\in B(\mathbb R_{\rm d}).
 \tag{MA-RB.29}
\]
Its ordinary real support is locally finite.
Consequently MA-FQ-002 expresses it as a finite integer
linear combination of indicators of translated lattices
and singletons.

This is a boundary constraint.  It does not give the
real-supported explicit-formula hypothesis of MA-FQ-006.
More precisely, the actual good-prime measure at the
center \(1/2\),
\[
 \mathcal M_{\rm good}
   =4\sum_{p\in\Omega}\sum_{\substack{j\ge1\\j\ {\rm odd}}}
       (\log p)p^{-j/2}
             \bigl(\delta_{j\log p}+\delta_{-j\log p}\bigr),
 \tag{MA-RB.30}
\]
is not tempered when \(\gamma>1/2\).
The positive finite measures producing (MA-RB.29) lose
their mass at infinity in the ordinary real topology.
There is no ordinary finite-measure weak limit to which
one can transfer their exact prime-frequency support.

\begin{proof}
The logarithmic derivative in the convergence half-plane is
\[
 F(s)=-\frac{D_\Omega'(s)}{D_\Omega(s)}
       =4\sum_{p\in\Omega}\sum_{\substack{j\ge1\\j\ {\rm odd}}}
                              (\log p)p^{-js}.
 \tag{MA-RB.31}
\]
It converges absolutely for \(\operatorname{Re}s>\gamma\).
For example choose an intermediate exponent between
\(\gamma\) and \(\operatorname{Re}s\); the extra power
dominates the logarithm, and the remaining prime powers
are bounded geometric series.
For \(\varepsilon>0\), define the positive finite measure
on the ordinary real line
\[
 \nu_\varepsilon
   =4\varepsilon\sum_{p\in\Omega}
          \sum_{\substack{j\ge1\\j\ {\rm odd}}}
              (\log p)p^{-j(\gamma+\varepsilon)}
                                      \delta_{j\log p}.
 \tag{MA-RB.32}
\]
Its Fourier transform, with the unscaled minus convention,
is
\[
 \int e^{-itx}\,d\nu_\varepsilon(x)
                         =\varepsilon F(\gamma+\varepsilon+it).
 \tag{MA-RB.33}
\]
At every point \(\gamma+it\), meromorphy of \(D_\Omega\)
makes the right side tend to \(a_\gamma(t)\).
Indeed a local order \(m\) gives
\(-D_\Omega'/D_\Omega=-m/(s-\gamma-it)+O(1)\).
The same conclusion gives zero at every regular nonzero
point.

For finite \(t_1,\ldots,t_r\) and complex \(z_1,\ldots,z_r\),
\[
 \sum_{\ell,m}z_\ell\overline{z_m}
   \varepsilon F(\gamma+\varepsilon+i(t_\ell-t_m))
    =\int\left|\sum_\ell z_\ell e^{-it_\ell x}\right|^2
                                               d\nu_\varepsilon(x)\ge0 .
 \tag{MA-RB.34}
\]
Passing to the finite pointwise limit proves positive
definiteness.  Its two-by-two matrices give
\(|a_\gamma(t)|\le a_\gamma(0)=k\).
The usual GNS construction for the discrete group
represents \(a_\gamma\) as a matrix coefficient of a
unitary representation, proving membership in
\(B(\mathbb R_{\rm d})\), exactly as in MA-FQ-001.
Meromorphic local finiteness of the divisor gives local
finiteness of its support on the line
\(\operatorname{Re}s=\gamma\).
The values belong to the finite integer set
\(\{-k,\ldots,k\}\).
The already proved Cohen/coset result MA-FQ-002 now
gives the asserted finite lattice/singleton description.
No ordinary continuity of \(a_\gamma\) is asserted.

To check temperedness directly, suppose contrary to
(MA-RB.30)'s claim that its positive measure were tempered.
A fixed nonnegative compact smooth function, at least one
on \([-1,1]\), scaled to \([-T,T]\), gives
\(\mathcal M_{\rm good}([-T,T])\le C(1+T)^N\):
each relevant Schwartz seminorm of the scaled test is
bounded by a fixed power of \(1+T\), and positivity
bounds the interval mass by that pairing.
For every \(\sigma>1/2\), this polynomial mass bound implies
\[
 \int_0^\infty e^{-(\sigma-1/2)x}\,d\mathcal M_{\rm good}(x)
       =4\sum_{p\in\Omega}\sum_{j\ {\rm odd}}
                               (\log p)p^{-j\sigma}<\infty .
 \tag{MA-RB.35}
\]
This follows by summing over unit intervals, with an
exponentially decreasing majorant times a polynomial.
Since \(j\log p\ge\log2\), convergence of this positive
series also gives convergence of
\(\log D_\Omega(\sigma)
 =4\sum_{p,j\ {\rm odd}}p^{-j\sigma}/j\).
Monotone convergence of the finite Euler products then
gives convergence of its nonnegative Dirichlet series.
It follows that \(\gamma\le1/2\), contrary to MA-RB-002.
Thus \(\mathcal M_{\rm good}\) is not tempered.
Adding the known finite-total-variation bad-prime and
conductor correction cannot repair that failure.

Finally,
\[
        \nu_\varepsilon(\mathbb R)
             =\varepsilon F(\gamma+\varepsilon)\longrightarrow k>0 .
 \tag{MA-RB.36}
\]
On every fixed compact interval only finitely many
prime-power logarithms occur in (MA-RB.32), and each
coefficient tends to zero.  Thus
\(\nu_\varepsilon\) converges vaguely to zero while
its total mass tends to \(k\); it is not tight.
It has no narrow weak limit as finite measures on the ordinary
real line (that is, against bounded continuous test functions).
Equivalently its pointwise Fourier limit
\(a_\gamma\) is not ordinarily continuous at zero:
local finiteness gives a punctured neighborhood on which
it vanishes, whereas \(a_\gamma(0)=k\).
This is consistent with its positive definiteness on
\(\mathbb R_{\rm d}\), but supplies no ordinary
pure-atomic Fourier identity with the original prime
measure.  That is the precise identity used in
MA-FQ-003/006 and absent from this boundary argument.
\end{proof}

\paragraph{Tests of the remaining hypotheses.}
The simple profile \(a(t)=k\mathbf1_{\{0\}}(t)\) satisfies
every conclusion (MA-RB.29) and has locally finite support:
its finite positive-definiteness form is
\(k\sum_{\text{equal }t_\ell}|\sum z_\ell|^2\ge0\).
As an ordinary atomic measure it is \(k\delta_0\),
whose ordinary Fourier transform has a constant
continuous density.  The boundary construction has no
pure-atomic Fourier premise that would exclude this
possibility.  This is a test of the stated boundary
conditions, not an arithmetic example or a genuine
nonzero comparison \(L\)-datum.

The actual pole at \(s=\gamma>1/2\) itself violates
the real-height support premise: in the convention
\(s=1/2+iz\), it is the nonreal height
\(z=-i(\gamma-1/2)\).
The new positivity argument therefore does not extend
MA-FQ-006 by simply forgetting that premise.
Its output is (MA-RB.29), not the stronger ordinary
Fourier identity of that theorem.

The sharpened order constraints also remain formally
compatible.  At a hypothetical reflected pair
\(\gamma,1-\gamma\), assign
\[
 (b,n,v_0,v_1,e,\operatorname{ord}L_W,
                      \operatorname{ord}L_{W\eta},k)
                         =(0,0,0,1,0,1,0,1).
 \tag{MA-RB.37}
\]
Then \(q_\gamma=0\), \(r_4=0\), and
\(r_5=r_4+v_0+\operatorname{ord}L_{W\eta}=0\);
the regular order is nine.
All listed holomorphy and induced-character order
conditions are nonnegative.  Putting the same orders
at the reflected point is compatible with the
functional equations.  Unlike at the center, a
root-number-\(+1\) equation does not force the order
at one member of an off-center reflected pair to be
even.  This is only a formal order test, not a field
or a claim that these orders are arithmetically realized.

Heilbronn's simple-real-zero descent theorem still does
not apply, since the actual normal-closure multiplicity
would be at least nine.  Takeda's theorem supplies
holomorphy, not the nonvanishing of \(L_{V\eta}L_X\)
on the residual interval.  Neither a shared-zero
exclusion there nor \((\mathcal R_B)\) has been proved.
Moreover even an exact comparison would give
\(B=A C\); it would not by itself rule out real
off-center zeros of the automorphic numerator.

\paragraph{Exact sources and proof scope.}
The classical positive-Euler comparison is Calegari
(2013), Theorem~4.7, article-PDF p.~12, Lemmas~4.8--4.10,
article-PDF pp.~13--15.  These are the PDF page numbers
used in MT-S508, not the journal pagination.
The finite-factor normalization and
Landau application were checked in MT-S508 and MA-S5-001.
The actual gamma equality and genuine comparison
\(L\)-datum inputs are MA-FQ-001/006.
The finite representation identities and bad-local
unit assertions are the already checked MA-S5-003/006;
their certificates are not rerun here.

The new pointwise symmetric-square application uses
Takeda (2015), Theorem~7.1 and its entire terminal proof,
primary preprint \texttt{arXiv:1506.04791v1},
printed pp.~49--50.
Those two pages were freshly read to verify the
full holomorphy range used in (MA-RB.25).
Its earlier Eisenstein-series and integral-representation
theorems remain the same declared classical inputs
as in MA-S5-006; no new audit of all their proofs is
claimed.
The required Hadamard constant, gamma signs and
Dedekind estimate are proved in MA-RB-001,
(MA-RB.8)--(MA-RB.13), independently of its separate
use of \(\mathsf A\) to convert an arbitrary \(B\)-zero.
The boundary Fourier-series and non-temperedness
arguments are proved in full above.

The coefficient nine in (MA-RB.26) combines the
regular Artin-motive multiplicity five with the
symmetric-power identity
\((\operatorname{Sym}^2\xi)\otimes\nu^{-1}
=V\eta\oplus X\) already placed in the representation
and motive framework of MA-S5.
The exact inherited locators remain Deligne (1979),
Definition~6.1 and \S\S6.2--6.3, and SGA~1,
Expos\'e~V, \S7.
This interpretation accounts for the algebraic factors;
it does not supply the missing real nonvanishing.

Takeda, S. (2015). On a certain metaplectic Eisenstein series
and the twisted symmetric square \(L\)-function.
\emph{Mathematische Zeitschrift, 281}, 103--157.
\url{https://doi.org/10.1007/s00209-015-1476-x}.

\section{Two precise limits of the remaining positivity arguments}
\label{sec:s5-scaled-datum}

The following two results test proposed extensions of the comparison
argument.  The first is an actual nonzero $L$-datum.  The second is
a finite representation-ring certificate, without a claimed arithmetic
realization.  Neither supplies a smaller unresolved arithmetic gate
than MA-RB-002.

\paragraph{MA-FQ-007 (a genuine rescaled datum with abscissa below one).}
There is a nonzero $L$-datum $\mathcal F_2=(f_2,0,m_2)$ in the
exact sense of Booker, Definition~1.3, such that:
\begin{enumerate}
\item its full archimedean kernel is zero, its spectral
      multiplicities are integers, and its coefficients $f_2(N)$
      are nonnegative for every $N>1$ and vanish off prime powers;
\item its associated Dirichlet series has nonnegative coefficients
      and abscissa of convergence $3/4$;
\item its completed function has simple real poles at $1/4$ and
      $3/4$, and its spectral support is contained in
      $|\operatorname{Im}z|\le1/4$ but is not contained in
      $\mathbb R$;
\item it is \emph{not} positive in Booker's separate sense of
      having only finitely many negative spectral multiplicities.
\end{enumerate}
In particular, the real-height support assumption of MA-FQ-006
cannot be replaced by the assertion that the abscissa is less
than one, even while retaining all the other hypotheses of that
theorem.  This example is not an $S_5$ comparison datum: its
coefficients violate both the finite-image bound and the specified
odd-prime-power pattern of that comparison.

\begin{proof}
Let $\chi_5$ be the primitive even quadratic character modulo five.
Use the genuine difference datum already constructed in MA-LD-003:
\[
 \mathcal F_0=\mathcal F_\zeta-\mathcal F_{\chi_5}
             =(f_0,0,m_0),\qquad
 L_{\mathcal F_0}(w)=R_0(w)=\frac{\zeta(w)}{L(w,\chi_5)}.
 \tag{MA-FQ.60}
\]
The real gamma factors agree.  Booker's conductor term in
Examples~1.4 gives the exact value at one, while logarithmic
differentiation gives the remaining coefficients:
\[
 f_0(1)=\tfrac12\log5,\qquad
 f_0(n)=(1-\chi_5(n))\frac{\Lambda(n)}{\sqrt n}
                                      \quad(n>1).
 \tag{MA-FQ.61}
\]
Define
\[
 \begin{gathered}
 m_2(z)=m_0(2z),\qquad K_2=0,\\
 f_2(n^2)=2f_0(n)\quad(n\ge1),\qquad
 f_2(N)=0\quad\text{if $N$ is not a square}.
 \end{gathered}
 \tag{MA-FQ.62}
\]
We verify all four axioms, including the normalization at one.

For (A1), $f_2(1)=\log5$ is real.  For $N>1$,
\[
 \begin{gathered}
 |f_2(N)|\ll N^{-1/4}\log N,\\
 \sum_{N\le X}|f_2(N)|^2
     =4\sum_{n\le\sqrt X}|f_0(n)|^2
     \ll 1+(\log X)^3 .
 \end{gathered}
 \tag{MA-FQ.63}
\]
The first estimate implies every required bound
$f_2(N)\log^jN\ll_j1$.  For the second estimate use
$|f_0(n)|\le2(\log n)n^{-1/2}$ for $n>1$ and compare
the resulting harmonic sum with its integral.  Thus the
mean square is $O_\varepsilon(X^\varepsilon)$ for every
$\varepsilon>0$.  Axiom (A2) holds with the identically
zero kernel.

For (A3), the actual integral completed divisor of (MA-FQ.60)
is discrete and has the usual polynomial counting bound.
Scaling its ordinates by $1/2$ preserves both properties and
every integer multiplicity.  More explicitly, the total variation
of $m_2$ for $|\operatorname{Re}z|\le T$ equals that of $m_0$
for $|\operatorname{Re}z|\le2T$.  The old support lies in
$|\operatorname{Im}z|\le1/2$, including the two completed
endpoints, and the new support lies in
$|\operatorname{Im}z|\le1/4$.

For (A4), take its actual test $g\in C_c^\infty(\mathbb R)$
and $h(z)=\int_{\mathbb R}g(x)e^{ixz}\,dx$ with
$h(\mathbb R)\subset\mathbb R$.  Put
$\widetilde g(x)=2g(2x)$; its transform is
$\widetilde h(z)=h(z/2)$ and is again real on the real axis.
The original explicit formula, with $K_0=0$, gives
\[
 \begin{aligned}
 \sum_zm_2(z)h(z)
   &=\sum_zm_0(z)h(z/2)\\
   &=-2\operatorname{Re}\sum_{n\ge1}2f_0(n)g(2\log n)\\
   &=-2\operatorname{Re}\sum_{N\ge1}f_2(N)g(\log N).
 \end{aligned}
 \tag{MA-FQ.64}
\]
The sums include $n=N=1$.  Consequently the scaling also
has the correct conductor term, rather than only a formal
Euler identity.  There is no Jacobian in the multiplicities:
the first equality merely reindexes the discrete divisor.
This proves that (MA-FQ.62) is a genuine $L$-datum.

Booker's normalization of its associated function now gives,
initially in an absolute-convergence half-plane,
\[
 \begin{aligned}
 \log L_{\mathcal F_2}(s)
   &=\sum_{n\ge2}\frac{2f_0(n)}{\log n^2}
                              (n^2)^{1/2-s}\\
   &=\sum_{n\ge2}\frac{f_0(n)}{\log n}n^{1-2s}
     =\log R_0(2s-\tfrac12).
 \end{aligned}
 \tag{MA-FQ.65}
\]
It follows that the exact finite function and analytic conductor
are
\[
 L_{\mathcal F_2}(s)
      =\frac{\zeta(2s-1/2)}{L(2s-1/2,\chi_5)},\qquad
 Q_{\mathcal F_2}=e^{-2f_2(1)}=\frac1{25}.
 \tag{MA-FQ.66}
\]
An actual completed function with divisor $m_2$ is
\[
 \mathscr R_2(s)=5^{-s+1/4}
                    \frac{\zeta(2s-1/2)}{L(2s-1/2,\chi_5)}.
 \tag{MA-FQ.67}
\]
Indeed the original completed ratio is $5^{-w/2}R_0(w)$;
the common real gamma factor cancels exactly.  Under
$s\mapsto1-s$, its variable $w=2s-1/2$ changes to $1-w$,
so it has the original completed functional equation centered
at $s=1/2$.  At $s=1/2+iz$ its old height is $2z$, as
required by (MA-FQ.62).

For completeness, the nonzero value of $L(1,\chi_5)$ needed
at the endpoints can be checked without a new zero-location
input.  Its convergent, residue-class-grouped series has terms
\[
 \begin{aligned}
 L(1,\chi_5)
 &=\sum_{a\in\{1,6,11,\ldots\}}
       \left(\frac1a-\frac1{a+1}-\frac1{a+2}
                                      +\frac1{a+3}\right)\\
 &=\sum_{a\in\{1,6,11,\ldots\}}
       \left(\frac1{a(a+1)}-
                          \frac1{(a+2)(a+3)}\right)>0.
 \end{aligned}
 \tag{MA-FQ.68}
\]
The completed nontrivial Dirichlet function is therefore
nonzero at one and, by its functional equation, at zero.
The completed zeta function has a simple pole at each endpoint.
Thus the original completed ratio has simple poles at $w=0,1$,
and (MA-FQ.67) has simple poles at $s=1/4,3/4$.
Their heights are $z=i/4,-i/4$, respectively.  In particular
the real-height support hypothesis fails unconditionally.

The local Euler factors of $R_0$ are
\[
 \begin{cases}
  1,&\chi_5(p)=1,\\
  (1+p^{-w})/(1-p^{-w})
          =1+2\sum_{j\ge1}p^{-jw},&\chi_5(p)=-1,\\
  (1-5^{-w})^{-1},&p=5.
 \end{cases}
 \tag{MA-FQ.69}
\]
They have nonnegative coefficients.  If
$R_0(w)=\sum a_0(n)n^{-w}$, then
\[
 a_2(n^2)=a_0(n)n^{1/2}\ge0,
 \qquad a_2(N)=0\quad(N\text{ not a square}).
 \tag{MA-FQ.70}
\]
This series converges absolutely for $\operatorname{Re}s>3/4$.
It cannot have a smaller abscissa, because its sum has the
actual pole at $3/4$.  Nonnegativity identifies its ordinary
and absolute abscissae.  Hence its abscissa is exactly $3/4$.

By (MA-FQ.61)--(MA-FQ.62), all the coefficients $f_2(N)$
for $N>1$ are nonnegative and are supported on prime powers.
Nevertheless this is not a positive datum in Booker's
terminology.  MA-LD-003 proves, using the degree-zero case
of Booker's Theorem~1.7, that $m_0$ has infinitely many
negative and infinitely many positive values.  The bijective
scaling of the divisor preserves that property for $m_2$.
Thus Euler coefficient positivity has not been substituted
for the separate spectral finiteness hypothesis.

Finally, if $\chi_5(p)=-1$, then
\[
                    f_2(p^2)=\frac{4\log p}{\sqrt p}.
 \tag{MA-FQ.71}
\]
There are arbitrarily large such primes.  An elementary
argument suffices here: if their set were finite, choose
a positive integer $N$ with $N\equiv2\pmod5$ and
$N\equiv1\pmod p$ for each prime in that finite set.
The Chinese remainder theorem permits these congruences.
No prime in the list or equal to five divides $N$, but
$\chi_5(N)=-1$, contradicting the multiplicativity of
$\chi_5$ on the prime factorization of $N$.
The finite-image comparison has a uniform bound
$|f(p^j)|\le C(\log p)p^{-j/2}$ outside a finite set;
at $j=2$, (MA-FQ.71) violates it by an unbounded factor
$4\sqrt p/C$.  No finite choice of bad primes repairs
that defect.  Moreover $f_2(p)=0$ for every prime,
and, at $\chi_5(p)=-1$, the nonzero exponents are
$2,6,10,\ldots$, whereas the exceptional $S_5$ pattern
has nonzero coefficients at the odd exponents.  This
establishes the asserted distinction from an actual
comparison in the stipulated family.
\end{proof}

\paragraph{MA-RB-004 (a genuine-character certificate for all induced order inequalities).}
Let $G=S_5$ and retain its seven irreducible characters
$\mathbf1,\eta,V,V\eta,W,W\eta,X$, of respective dimensions
$1,1,4,4,5,5,6$.  For each integer $k\ge1$, the formal
off-center order assignment
\[
 \begin{gathered}
 o(V\eta)=o(W)=k,\\
 o(\mathbf1)=o(\eta)=o(V)=o(W\eta)=o(X)=0
 \end{gathered}
 \tag{MA-RB.40}
\]
has the genuine order character
\[
                    \Theta_\gamma=k(V\eta+W).
 \tag{MA-RB.41}
\]
For every subgroup $H\le G$ and every genuine character
$\chi$ of $H$, it satisfies
\[
 \begin{aligned}
 o(\operatorname{Ind}_H^G\chi)
   &=\langle\operatorname{Ind}_H^G\chi,\Theta_\gamma\rangle_G\\
   &=\langle\chi,\operatorname{Res}_H^G\Theta_\gamma\rangle_H
      \in\mathbb Z_{\ge0}.
 \end{aligned}
 \tag{MA-RB.42}
\]
Together with the formal automorphic numerator order $n=0$,
it attains equality in both $v_1+e\ge k$ and the regular
order bound $\operatorname{ord}_\gamma\zeta_L\ge9k$ of
MA-RB-002.  Thus adding only nonnegative induced-character
order constraints cannot strengthen either bound or exclude
this formal order pattern.

\begin{proof}
Extend $o$ additively to the representation ring.  The
irreducible characters are orthonormal, so (MA-RB.40) gives
$o(\psi)=\langle\psi,\Theta_\gamma\rangle_G$ for every
character $\psi$.  To verify the reciprocity equality directly,
use the finite induction formula and the fact that
$\Theta_\gamma$ is a class function:
\[
 \begin{aligned}
 \langle\operatorname{Ind}_H^G\chi,\Theta_\gamma\rangle_G
 &=\frac1{|G||H|}\sum_{x\in G}\sum_{h\in H}
                   \chi(h)\overline{\Theta_\gamma(xhx^{-1})}\\
 &=\frac1{|H|}\sum_{h\in H}
                   \chi(h)\overline{\Theta_\gamma(h)}
  =\langle\chi,\operatorname{Res}_H^G\Theta_\gamma\rangle_H .
 \end{aligned}
 \tag{MA-RB.43}
\]
Both characters in the last inner product are genuine.
Decomposing their finite-dimensional representations into
irreducibles makes that inner product a sum of products
of nonnegative integer multiplicities.  This proves
(MA-RB.42), without restricting $H$ to the previously
listed subgroups or $\chi$ to one-dimensional characters.

For the pointwise analytic order identities already proved
in MA-RB-002, the assignment is
\[
 (b,n,v_0,v_1,e,o(W),o(W\eta))=(0,0,0,k,0,k,0).
 \tag{MA-RB.44}
\]
It gives $o(A/L_W)=-k$, $o(B)=k$, and
$v_1+e=k$.  The previously checked symmetric-square
and Rankin--Selberg identities give
$q_\gamma=v_1+e-k=0$, $r_4=q_\gamma+n=0$, and
$r_5=r_4+v_0+o(W\eta)=0$.  Their holomorphy constraints
are therefore satisfied.  The regular-character order is
\[
 o(\operatorname{Reg}_G)
    =\Theta_\gamma(1)=4k+5k=9k.
 \tag{MA-RB.45}
\]
Equivalently, it is the regular Artin factorization with
irreducible multiplicities equal to their dimensions.

Under exact Artin induction, (MA-RB.42) is precisely the
formal order of $L_{L^H}(s,\chi)$ at the hypothetical
point.  In particular it satisfies every order
nonnegativity supplied by the holomorphy of an abelian
relative Hecke function, as well as every further
holomorphic induced genuine character.  Assign the same
orders at the distinct reflected point $1-\gamma$.
This is compatible with the order relation in the functional
equations; an off-center pair imposes no central parity
condition on the individual orders.  Consequently no
contradiction, strict improvement of either displayed bound,
or exclusion of the common-zero pattern follows from adding
only these nonnegative induced-order constraints.
\end{proof}

The certificate in MA-RB-004 does not construct a number
field, automorphic representation, or $L$-function realizing
these orders.  Additional analytic information, including
a suitable nonvanishing or a stronger independently proved
functorial constraint, could still rule them out.  The lemma
audits the scope of this induction/holomorphy argument; it
does not reduce the shared-real-zero gate of MA-RB-002.
Likewise the genuine $L$-datum of MA-FQ-007 is an exact
counterexample to the stated abstract extension, not to an
$S_5$ theorem or to the retained real-support criterion.

\paragraph{Primary locators and proof scope.}
The defining axioms, including the separate meaning of
\emph{positive}, are Booker, Definition~1.3, author p.~3;
the degree-one Artin/automorphic constructions and exact
conductor term are Examples~1.4(2)--(3), p.~4, in
\texttt{arXiv:1308.3067v2}.  Those two primary pages were
freshly read for the rescaling.  The group property of
$L$-data is stated on p.~4.  The earlier complete audit
of the degree-zero part of Theorem~1.7, pp.~15--17, is
MA-LD-003; that theorem is used here only for the final
spectral-positivity distinction, not to establish the four
axioms or the abscissa.  No new audit of the classical
degree-one continuation and functional equation is claimed.

The $S_5$ character notation, exact Artin induction, and
tensor identities in MA-RB-004 are the already checked
MA-S5-003/006 and MA-RB-002.  Formula (MA-RB.43) proves
the needed reciprocity calculation here.  The interpretation
as a genuine finite representation concerns nonnegative
character multiplicities; it supplies no arithmetic weight
or analytic nonvanishing assertion.

Booker, A. R. (2015). $L$-functions as distributions.
\emph{Mathematische Annalen, 363}(1--2), 423--454.
\url{https://doi.org/10.1007/s00208-015-1178-z}.
Author manuscript: \url{https://arxiv.org/abs/1308.3067v2}.

\section{Exact quadratic Schur tests of the residual comparison pole}
\label{sec:s5-quadratic-schur}

Retain the actual nontrivial branch of MA-RB-002 at its real
rightmost pole \(\gamma\in(1/2,1)\), with pole order \(k\ge1\).
Write
\[
 \begin{gathered}
 b=\operatorname{ord}_\gamma C,\quad
 n=\operatorname{ord}_\gamma A,\quad
 v_0=\operatorname{ord}_\gamma L_V,\quad
 v_1=\operatorname{ord}_\gamma L_{V\eta},\quad
 e=\operatorname{ord}_\gamma L_X,\\
 x=\operatorname{ord}_\gamma L_W=n+k,\qquad
 x_1=\operatorname{ord}_\gamma L_{W\eta}.
 \end{gathered}
 \tag{MA-RB.46}
\]
All these orders are nonnegative. The justification for \(x\)
is the actual comparison pole, not a general assumption of
Artin holomorphy. At this real point
\(L_{W\eta}=\zeta_F/\zeta_{\mathbb Q}\) is holomorphic, since
\(\zeta_{\mathbb Q}(\gamma)\ne0\).

Let \(\Pi,\nu,\varpi\) be Calegari's four-dimensional
representation, its finite-order similitude character, and its
specified five-dimensional comparison. A twist written outside
\(\operatorname{Sym}^2\) means a twist of that \(L\)-function.
In particular \(\operatorname{Sym}^2(\varpi)\otimes\eta\)
denotes a different operation from
\(\operatorname{Sym}^2(\varpi\otimes\eta)\).

\paragraph{MA-RB-005 (the quadratic Schur holomorphy constraints).}
Choose a finite set \(S\) containing infinity, \(2\), and every
ramified or omitted place in the comparisons. Define
\[
 \begin{aligned}
 Q_4(s)&=L^S(s,\Pi,\operatorname{Sym}^2\otimes\nu^{-1}),&
 Q_{4,\eta}(s)&=L^S(s,\Pi,\operatorname{Sym}^2
                                      \otimes\nu^{-1}\eta),\\
 Q_5(s)&=L^S(s,\varpi,\operatorname{Sym}^2),&
 Q_{5,\eta}(s)&=L^S(s,\varpi,\operatorname{Sym}^2\otimes\eta).
 \end{aligned}
 \tag{MA-RB.47}
\]
All four functions are holomorphic at \(\gamma\), and
\[
 \begin{array}{c|c}
 \text{function}&\text{order at }\gamma\\ \hline
 Q_4&v_1+e-k\\
 Q_{4,\eta}&v_0+e+k\\
 Q_5&v_0+n+x_1\\
 Q_{5,\eta}&b+v_1+n+x_1+2k .
 \end{array}
 \tag{MA-RB.48}
\]
The partial Euler products for \(\wedge^2\varpi\)
and \((\wedge^2\varpi)\otimes\eta\) are \(Q_4\) and
\(Q_{4,\eta}\). This supplies their continuation and
holomorphy without positing a separate automorphic
exterior-square lift of \(\varpi\).

Among the four order nonnegativities, only \(v_1+e\ge k\)
is not automatic from (MA-RB.46); that is already MA-RB-002.
Thus these quadratic constraints do not increase the
coefficient nine in its normal-closure multiplicity bound
or exclude the formal genuine-character assignment of
MA-RB-004. They do force the additional actual zeros
\[
 \operatorname{ord}_\gamma Q_{4,\eta}\ge k,\qquad
 \operatorname{ord}_\gamma Q_{5,\eta}\ge2k .
 \tag{MA-RB.49}
\]
No nonvanishing of either function is asserted.

\begin{proof}
\emph{Finite representation identities.}
Use the already checked \(S_5\) character table, with columns
\(\mathbf1,\eta,V,V\eta,W,W\eta,X\). The relevant new
calculations, in the specified class order, are
\[
 \begin{array}{c|rrrrrrr}
 &1&2A&2B&3A&6A&4A&5A\\ \hline
 \chi_W&5&1&1&-1&1&-1&0\\
 \chi_W(g^2)&5&5&5&-1&-1&1&0\\
 \chi_{\operatorname{Sym}^2W}&15&3&3&0&0&1&0\\
 \chi_{\wedge^2W}&10&-2&-2&1&1&0&0 .
 \end{array}
 \tag{MA-RB.50}
\]
For eigenvalues \(\lambda_j\), the sums of
\(\lambda_j\lambda_l\) over \(j\le l\) and \(j<l\)
are \((\chi_W(g)^2+\chi_W(g^2))/2\) and
\((\chi_W(g)^2-\chi_W(g^2))/2\), respectively. They
give every entry displayed. Adding columns of the
checked character table now proves
\[
 \operatorname{Sym}^2W
   =\mathbf1\oplus V\oplus W\oplus W\eta,\qquad
 \wedge^2W=V\eta\oplus X .
 \tag{MA-RB.51}
\]
Equality of characters proves the representation identities.
The table also gives \(X\eta=X\), so their external
\(\eta\)-twists have constituents
\(\eta,V\eta,W,W\eta\) and \(V,X\), respectively.
This reuses the existing table without reopening the
earlier finite representation certificates.

\emph{The general symplectic parameter identity.}
Calegari's Lemma~4.2 and the subsequent parameter formula
give, at every unramified place,
\[
 S_4=\{a,b,\nu/a,\nu/b\},\qquad
 T_5=\{1,ab/\nu,\nu/(ab),a/b,b/a\}.
 \tag{MA-RB.52}
\]
The eigenvalues of
\(\nu^{-1}\operatorname{Sym}^2 S_4\) are
\[
 \{a^2/\nu,b^2/\nu,\nu/a^2,\nu/b^2,
                   ab/\nu,a/b,b/a,\nu/(ab),1,1\}.
 \tag{MA-RB.53}
\]
The products of two distinct entries of \(T_5\) give
the same multiset: the pairs containing \(1\) give the
four middle entries, the two inverse pairs give \(1,1\),
and the remaining four pairs give the first four entries.
Hence
\[
                  \wedge^2T_5=\nu^{-1}\operatorname{Sym}^2S_4 .
 \tag{MA-RB.54}
\]
Multiplication by \(\eta(p)\) gives the external twist.
Taking reciprocal Euler polynomials outside \(S\) proves
the asserted identities with \(Q_4,Q_{4,\eta}\) initially
in an absolute-convergence half-plane. Equality with
these known functions then supplies the meromorphic
continuations of the exterior-square products. No
additional functorial lift is assumed.

\emph{Exceptional Euler factors.}
At \(p\in\Omega\) the class is \(4A\), and the two
five-dimensional multisets are
\[
        T_A=\{1,1,1,i,-i\},\qquad
        T_W=\{1,-1,-1,i,-i\}.
 \tag{MA-RB.55}
\]
For \(P_{\mathcal E}(T)=\prod_{\lambda\in\mathcal E}
(1-\lambda T)\), direct multiplication gives
\[
 \begin{array}{c|c|c}
 &\text{automorphic polynomial}&\text{Artin polynomial}\\ \hline
 \operatorname{Sym}^2&
 (1-T)^7(1+T)^2(1+T^2)^3&
 (1-T)^5(1+T)^4(1+T^2)^3\\
 \wedge^2&
 (1-T)^4(1+T^2)^3&
 (1-T)^2(1+T)^2(1+T^2)^3 .
 \end{array}
 \tag{MA-RB.56}
\]
Indeed, the multiplicities of \(1,-1,i,-i\) in the
two symmetric squares are \(7,2,3,3\) and \(5,4,3,3\).
In the exterior squares they are \(4,0,3,3\) and
\(2,2,3,3\). These counts prove all signs in (MA-RB.56).
The \(L\)-factor is \(P^{-1}\), so its automorphic to
Artin ratio in either row is
\[
       \Delta(T)=\frac{P_{\rm Artin}(T)}{P_{\rm automorphic}(T)}
                   =\left(\frac{1+T}{1-T}\right)^2 .
 \tag{MA-RB.57}
\]
Since \(\eta(p)=-1\) on \(4A\), the external twist
replaces \(T\) by \(-T\) and the ratio by \(\Delta^{-1}\).
At all other unramified places the original parameters
agree; so do their Schur functors and external twists.

Put \(D^S(s)=\prod_{p\in\Omega\setminus S}\Delta(p^{-s})\).
Removing finitely many primes of \(\Omega\) changes
\(D_\Omega\) by a holomorphic nonzero factor at
\(\gamma>0\); thus \(\operatorname{ord}_\gamma D^S=-k\).
Combining (MA-RB.51)--(MA-RB.57) gives
\[
 \begin{aligned}
 Q_4&=D^S L_{V\eta}^S L_X^S,&
 Q_{4,\eta}&=(D^S)^{-1}L_V^S L_X^S,\\
 Q_5&=D^S\zeta_{\mathbb Q}^S L_V^S L_W^S L_{W\eta}^S,&
 Q_{5,\eta}&=(D^S)^{-1}C^S L_{V\eta}^S L_W^S L_{W\eta}^S .
 \end{aligned}
 \tag{MA-RB.58}
\]
These are first identities of absolutely convergent Euler
products and then meromorphic identities. All removed
finite Artin factors are units at \(\gamma\): their
inverse Euler polynomials have roots of modulus at least
one, whereas \(p^{-\gamma}<1\). We use the four
\emph{partial} symmetric-square functions; no assertion
about unspecified ramified factors of a new lift is needed.
Deletion of \(S\) therefore changes none of the Artin orders
on the right. Also \(\operatorname{ord}_\gamma
\zeta_{\mathbb Q}^S=0\).

\emph{Known holomorphy and the exact orders.}
Takeda's Theorem~7.1 applies to both cuspidal representations
\(\Pi,\varpi\), with the four unitary twists in (MA-RB.47).
Their central characters are unitary and \(\nu,\eta\)
have finite order. These are the already checked
Calegari inputs: the finite-image base changes have
finite-order central character, hence the descended
central characters are unitary. Enlarge \(S\) as in
Takeda's theorem. Its only possible poles occur at
zero and one, so all four functions are holomorphic at
\(\gamma\in(1/2,1)\).
Taking orders in (MA-RB.58) gives
\[
 \begin{split}
 \operatorname{ord}_\gamma Q_4&=v_1+e-k,\\
 \operatorname{ord}_\gamma Q_{4,\eta}&=v_0+e+k,\\
 \operatorname{ord}_\gamma Q_5&=v_0+x+x_1-k=v_0+n+x_1,\\
 \operatorname{ord}_\gamma Q_{5,\eta}
             &=b+v_1+x+x_1+k=b+v_1+n+x_1+2k .
 \end{split}
 \tag{MA-RB.59}
\]
This proves (MA-RB.48)--(MA-RB.49). All functions used
here are genuine meromorphic \(L\)-functions or exact
Euler products identified with them; these are actual
pointwise identities, not freely assigned analytic orders.

\emph{Scope of these inequalities.}
The last three expressions in (MA-RB.59) are
automatically nonnegative under (MA-RB.46). The first
gives exactly \(v_1+e\ge k\). For every positive integer
\(k\), the formal assignment of MA-RB-004 is
\[
       (b,n,v_0,v_1,e,x,x_1)=(0,0,0,k,0,k,0).
\]
It gives
\[
 (\operatorname{ord}Q_4,\operatorname{ord}Q_{4,\eta},
       \operatorname{ord}Q_5,\operatorname{ord}Q_{5,\eta})
      =(0,k,0,3k),\qquad
 b+4v_0+4v_1+5x+5x_1+6e=9k .
 \tag{MA-RB.60}
\]
Thus it satisfies all the new inequalities and attains
the old regular-order lower bound. This proves the
specified failure to obtain a stricter bound or a
contradiction from this family of nonnegativities.
It does not assert arithmetic realization of that
formal assignment.
\end{proof}

\paragraph{Higher exterior powers of the same parameters.}
The multiset \(T_5\) in (MA-RB.52) is invariant under
inversion and has product one. Complementing a subset
therefore gives \(\wedge^{5-r}T_5=\wedge^rT_5\).
The same holds for \(W\): it is self-dual and its
determinant is one, since a transposition has three
eigenvalues \(1\) and two eigenvalues \(-1\), and
transpositions generate \(S_5\). Thus its third exterior
power repeats the second, its fourth repeats the
standard representation, and its fifth is trivial.
For \(S_4\), its product is \(\nu^2\) and its inverse
multiset is \(\nu^{-1}S_4\), so
\(\wedge^3S_4=\nu S_4\) and \(\wedge^4S_4=\nu^2\).
These introduce no further standard exterior-power
factor beyond the preceding ones and the original
standard functions. No theorem about higher symmetric
powers is inferred from Takeda's quadratic theorem.

\paragraph{Primary inputs and actual verification scope.}
The new character calculation uses Calegari's character
table in \S2. His Lemma~4.2 and the following displayed
decomposition and parameter formula supply (MA-RB.52);
Lemma~4.6 gives \(\nu=\psi^{-1}\); Theorem~4.7 and the
local calculation in Lemma~4.8 give (MA-RB.55).
The supplied primary LaTeX locators are
\texttt{lemma:satake}, \texttt{lemma:sim}, the theorem
immediately following the latter, and \texttt{lemma:aa}.
The exact source spans read for this deduction are
lines 554--602 and 735--807 of \texttt{S5ArtinWeb.tex}.
These locators do not conflate article-PDF and journal pages.
The literal base-change summand error newly detected in
the proof of Lemma~4.2 is recorded and repaired in
MA-S5-008. That complete repair preserves the symplectic
parameter input used here; the printed identical-summand
argument is not being endorsed.

The analytic input is Takeda (2015), Theorem~7.1,
\texttt{arXiv:1506.04791v1}, printed pp.~49--50.
Its entire terminal proof was read in the MA-RB-002 audit.
Here its hypotheses are checked for the additional
five-dimensional representation and the external twists.
The upstream Eisenstein-series and integral-representation
results remain the declared classical inputs; no new
full audit of those results is claimed.

The exact new calculations are reproducible in
\texttt{checks/s5/s5\_quadratic\_schur.py} and its JSON
output. They check the two character decompositions,
the symbolic \(\operatorname{GSp}_4\) weight identity,
the four exceptional ratios and the order expressions,
using the previously checked finite table.
The written proof is independent of floating-point
calculations. In the Artin-motive language fixed in
MA-S5, (MA-RB.51) is an identity of symmetric and exterior
operations in the representation \(\lambda\)-ring.
The inherited exact locators remain Deligne (1979),
Definition~6.1 and \S\S6.2--6.3, and SGA~1, Expos\'e~V,
\S7. This interpretation supplies no missing
nonvanishing at \(\gamma\).

Calegari, F. (2013). The Artin conjecture for some
\(S_5\)-extensions. \emph{Mathematische Annalen, 356},
191--207. \url{https://doi.org/10.1007/s00208-012-0839-4}.

Takeda, S. (2015). On a certain metaplectic Eisenstein series
and the twisted symmetric square \(L\)-function.
\emph{Mathematische Zeitschrift, 281}, 103--157.
\url{https://doi.org/10.1007/s00209-015-1476-x}.

\section{All Rankin cross products among the known comparison constituents}
\label{sec:s5-rankin-cross}

The following calculation tests every pair among the specified
five-dimensional comparison and its sign twist and the known
four- and six-dimensional cuspidal factors. It uses the
Rankin--Selberg functions of those actual cuspidal representations;
it assumes no automorphic tensor-product lift.

Retain the actual point \(\gamma\), the positive integer \(k\),
and the nonnegative orders of (MA-RB.46). Set
\[
 \begin{gathered}
 (\pi_1,\pi_2,\pi_3,\pi_4,\pi_5)
   =(\varpi,\varpi\otimes\eta,\pi_V,\pi_{V\eta},\pi_X),\\
 (U_1,U_2,U_3,U_4,U_5)=(W,W\eta,V,V\eta,X),\\
 \Theta=b\eta+v_0V+v_1V\eta+(n+k)W+x_1W\eta+eX,\qquad
 q=v_1+e-k\ge0.
 \end{gathered}
 \tag{MA-RB.61}
\]
Here \(\pi_V,\pi_{V\eta},\pi_X\) are the established
automorphic realizations of those Artin representations.
In particular \(\pi_V\) is distinct notation from the
four-dimensional symplectic comparison \(\Pi\) in MA-RB-005.
All five \(\pi_i\) are unitary cuspidal and self-dual.
For \(\varpi\) this was proved in MA-S5-002; sign twisting
preserves it. For the other three it follows from their
actual Artin correspondence and the real self-dual
representations, or from their local parameters and strong
multiplicity one.

\paragraph{MA-RB-006 (the complete fifteen-pair order constraints).}
Choose \(S\) to contain infinity and all ramified or omitted
places in this finite family, and retain \(D^S\) from
MA-RB-005. Put
\[
 \mathcal L_{ij}(s)=L^S(s,\pi_i\times\pi_j^\vee),\qquad
 E=\begin{pmatrix}
 2&-2&0&0&0\\
 -2&2&0&0&0\\
 0&0&0&0&0\\
 0&0&0&0&0\\
 0&0&0&0&0
 \end{pmatrix}.
 \tag{MA-RB.62}
\]
The exact meromorphic identities are
\[
          \mathcal L_{ij}(s)
             =(D^S(s))^{E_{ij}}L^S(s,U_i\otimes U_j)
                     \qquad(1\le i,j\le5).
 \tag{MA-RB.63}
\]
Each \(\mathcal L_{ij}\) is holomorphic at \(\gamma\), and
its order is
\[
       r_{ij}
          =\langle U_i\otimes U_j,\Theta\rangle_{S_5}
                       -E_{ij}k .
 \tag{MA-RB.64}
\]
Let
\[
 M=\begin{pmatrix}
 0&4&1&2&2\\
 4&0&2&1&2\\
 1&2&1&1&2\\
 2&1&1&1&2\\
 2&2&2&2&3
 \end{pmatrix}.
 \tag{MA-RB.65}
\]
Then \(r_{ij}\ge kM_{ij}\) entry by entry. Every one of
the fifteen holomorphy inequalities \(r_{ij}\ge0\) is
already implied by the nonnegative orders in (MA-RB.61)
and the earlier inequality \(q\ge0\). The formal
assignment \(\Theta=k(V\eta+W)\), \(n=0\), attains
\((r_{ij})=kM\) and regular order \(9k\).
Thus this complete family of Rankin holomorphy constraints
neither excludes that formal assignment nor improves the
previous normal-closure multiplicity bound.

\begin{proof}
\emph{The local comparison for all fifteen pairs.}
At an exceptional prime the class is \(4A\).
The five-dimensional parameter multisets \(T_A,T_W\)
are (MA-RB.55). For their \(j\)th power traces \(a_j,w_j\),
\[
 a_j-w_j=
  \begin{cases}4,&j\text{ odd},\\0,&j\text{ even},\end{cases}
 \quad
 a_j^2-w_j^2=
  \begin{cases}8,&j\text{ odd},\\0,&j\text{ even}.\end{cases}
 \tag{MA-RB.66}
\]
Indeed \(a_j=3,w_j=-1\) for odd \(j\), while for even
\(j\) both are \(3+i^j+(-i)^j\).
The sign character equals \(-1\) here, so the traces
for \(\varpi\otimes\eta,W\eta\) are
\((-1)^j a_j,(-1)^jw_j\).

The local standard parameter for \(V\) at \(4A\) is
\(\{1,-1,i,-i\}\). Its sign twist has the same multiset,
and the parameter of \(X=\wedge^2V\) is
\(\{1,-1,i,i,-i,-i\}\). Consequently the \(j\)th
power trace of each of \(V,V\eta,X\) is zero whenever
\(j\) is odd. This can also be read directly from the
character table: every odd power of a \(4A\) element
remains in \(4A\).

For local multisets \(P,Q\), the logarithm of the
Rankin Euler factor is
\(\sum_{j\ge1}\operatorname{tr}(P^j)
\operatorname{tr}(Q^j)T^j/j\).
All our local multisets are self-dual, so the use of
\(\pi_j^\vee\) changes none of these calculations.
In a mixed comparison with \(V,V\eta\), or \(X\),
the trace difference vanishes for even \(j\) by
(MA-RB.66), and for odd \(j\) by the zero trace of
the second factor. Its complete Euler ratio is
therefore one. The known three representations
already match at all unramified primes, so a pair
entirely among them has ratio one as well.

For the two diagonal five-dimensional pairs the
logarithmic ratio is
\[
       8\sum_{\substack{j\ge1\\j\ {\rm odd}}}\frac{T^j}{j}
                     =2\log\Delta(T).
 \tag{MA-RB.67}
\]
For the cross pair \((1,2)\), its extra factor
\((-1)^j\) negates this expression, giving
\(\Delta(T)^{-2}\). Exponentiating the absolutely
convergent logarithms near \(T=0\) proves the exact
rational local identities. Outside \(\Omega\) all
five original unramified parameters already match
their corresponding Artin ones.
Multiplying the local identities proves (MA-RB.63)
in an absolute-convergence half-plane, and hence
meromorphically. This proof determines the complete
local polynomials, not only their first coefficients.

\emph{The complete finite tensor calculation.}
For clarity the following table gives the multiplicities
of every irreducible in \(U_i\otimes U_j\); symmetric
entries have been listed once.
\[
 \begin{array}{c|rrrrrrr}
  ij&\mathbf1&\eta&V&V\eta&W&W\eta&X\\ \hline
  11&1&0&1&1&1&1&1\\
  12&0&1&1&1&1&1&1\\
  13&0&0&1&0&1&1&1\\
  14&0&0&0&1&1&1&1\\
  15&0&0&1&1&1&1&2\\
  22&1&0&1&1&1&1&1\\
  23&0&0&0&1&1&1&1\\
  24&0&0&1&0&1&1&1\\
  25&0&0&1&1&1&1&2\\
  33&1&0&1&0&1&0&1\\
  34&0&1&0&1&0&1&1\\
  35&0&0&1&1&1&1&1\\
  44&1&0&1&0&1&0&1\\
  45&0&0&1&1&1&1&1\\
  55&1&1&1&1&2&2&1
 \end{array}
 \tag{MA-RB.68}
\]
Here is an exact method proving every entry, rather
than an assumed tensor-product rule. In class order
\(1,2A,2B,3A,6A,4A,5A\), take the already checked
character-table columns and class sizes
\((1,10,15,20,20,30,24)\). Multiply the two selected
columns entry by entry. Its coefficient in column
\(\chi\) is
\[
 \frac1{120}\sum_{\mathcal C}|\mathcal C|\,
      \chi_{U_i}(\mathcal C)\chi_{U_j}(\mathcal C)
                                      \chi(\mathcal C).
\]
Substitution gives (MA-RB.68), and multiplying the
resulting integer rows back by the character table
recovers each product column. Equality of characters
then proves every displayed decomposition.
The old \(W\otimes W\) identity is included for the
complete matrix, using the same previously checked table.
All remaining entries also follow from the displayed
products and \(X\eta=X\), but no sign-twist convention
has been left implicit in (MA-RB.68).

\emph{Actual holomorphy and pointwise orders.}
The Rankin--Selberg theorem for two unitary cuspidal
representations is applicable to every \(\pi_i,\pi_j\).
For unequal dimensions the completed function is entire.
For equal dimensions its possible poles lie on
\(\operatorname{Re}s=0\) or \(1\), and thus away
from the present \(\gamma\). Deleting finite Euler
and archimedean factors introduces no pole: their
reciprocals are entire. In this finite-image situation
they are in fact units at a real point of \((0,1)\),
by the same local-parameter argument as MA-S5-003.
Hence every \(\mathcal L_{ij}\) is holomorphic at
\(\gamma\), as asserted.

On the Artin side all deleted finite factors are
units there. The order of the trivial constituent
is zero because \(\zeta_{\mathbb Q}(\gamma)\ne0\).
Artin multiplicativity and (MA-RB.61) therefore give
the inner product in (MA-RB.64). Also
\(\operatorname{ord}_\gamma D^S=-k\). Taking orders
in the genuine meromorphic identity (MA-RB.63)
proves (MA-RB.64) exactly, including its sign.
No general holomorphy of the unknown \(L_W\) was
assumed: its positive order \(n+k\) at this particular
point is the consequence already proved in MA-RB-002.

\emph{All inequalities and the formal witness.}
The two nonzero positive exponents in \(E\) give
\[
 r_{11}=r_{22}
     =v_0+v_1+n+k+x_1+e-2k
     =q+n+v_0+x_1\ge0 .
 \tag{MA-RB.69}
\]
The only negative exponents are \(E_{12}=E_{21}=-2\);
they add \(2k\) to an already nonnegative genuine-character
pairing. All other exponents vanish, so their
nonnegativity is immediate from the genuine tensor
characters and the nonnegative orders in \(\Theta\).
Thus every holomorphy inequality follows from the
prior order constraints.

More explicitly, subtracting \(kM_{ij}\) from
(MA-RB.64) gives the following exhaustive table:
\[
 \begin{array}{c|l}
 ij&r_{ij}-kM_{ij}\\ \hline
 11,22&q+n+v_0+x_1\\
 12&b+q+n+v_0+x_1\\
 13,24&e+n+v_0+x_1\\
 14,23&q+n+x_1\\
 15,25&e+q+n+v_0+x_1\\
 33,44&e+n+v_0\\
 34&b+q+x_1\\
 35,45&q+n+v_0+x_1\\
 55&b+q+2n+v_0+2x_1 .
 \end{array}
 \tag{MA-RB.70}
\]
Every entry is nonnegative, proving the stated
entrywise matrix bound. Under the formal assignment
\[
        (b,n,v_0,v_1,e,x_1)=(0,0,0,k,0,0)
\]
one has \(q=0\), so every entry in (MA-RB.70)
vanishes. Hence \(r=kM\), while the regular-character
order remains \(9k\), as in MA-RB-004.
The same tuple also satisfies all four quadratic
Schur inequalities of MA-RB-005. It follows that
adding all fifteen Rankin order nonnegativities
cannot by itself shrink that formal feasible set
or strengthen the bound it attains.
\end{proof}

The conclusion concerns exactly the inequalities supplied
by these known Rankin holomorphy theorems. It does not
assert that arbitrary formal order matrices are
arithmetically realizable, or rule out further analytic
restrictions on their simultaneous zeros. In particular
the numerous actual zeros forced by (MA-RB.65) could be
excluded by a suitable additional nonvanishing theorem;
the holomorphy theorem used here supplies none.
The residual shared-real-zero criterion of MA-RB-002
therefore remains open.

\paragraph{Primary inputs and verification scope.}
The unramified exceptional parameters and the real
\(S_5\) table are the same Calegari \S2 and
Theorem~4.7 inputs used in MA-RB-005. Its four- and
six-dimensional automorphy is Theorem~1.1, already
checked in MA-S5. The argument reuses those
correspondences and does not assume the missing
five-dimensional comparison.

For the analytic input, Cogdell (2004), Lecture~9,
Theorems~9.1(i) and 9.2(i), printed p.~74
(PDF p.~78), state the precise unequal- and
equal-dimension conclusions. The complete terminal
continuation and pole argument in Lecture~9,
\S\S2--3, printed p.~71 (PDF p.~75), was freshly
read for this application. It expresses the global
function through the global integrals and entire
local ratios, then chooses the local ratios nonzero
at the point in question. The underlying global
integral and local Whittaker theorems remain declared
classical inputs; their full earlier proofs are not
newly audited here.

The complete exact calculation is reproducible in
\texttt{checks/s5/s5\_rankin\_cross.py} and its JSON
output. It checks all fifteen product-character
decompositions, the full cyclotomic local ratios,
and the general-\(k\) order matrix. It reuses the
existing checked table and makes no arithmetic
zero-location computation.

Cogdell, J. W. (2004). Lectures on \(L\)-functions,
converse theorems, and functoriality for
\(\mathrm{GL}_n\). In \emph{Lectures on automorphic
\(L\)-functions} (Fields Institute Monographs,
Vol.~20). American Mathematical Society.
\url{https://people.math.osu.edu/cogdell.1/fields-www.pdf}.

\section{A variational bound retaining both real zeros and the known prime}
\label{sec:s5-variational-endpoint}

This refines the residual interval in MA-RB-002 without a
zero-location hypothesis.  It uses the exact signature and the
specified Frobenius at five, rather than replacing all their
contributions by a constant.  Throughout, \(L/\mathbb Q\) is
the actual Galois closure in Calegari's family and
\(\Delta=\Delta_L\) is its absolute discriminant.
Thus \([L:\mathbb Q]=120\), its signature is \((0,60)\),
and five is unramified with Frobenius of cycle type \(2B\).

\paragraph{MA-RB-007 (the exact one-variable interval bound).}
For \(\sigma>1\), define
\[
 \begin{split}
 H_\Delta(\sigma)={}&\frac1\sigma+\frac1{\sigma-1}
       +\frac12\log\Delta
       +60\bigl(\psi(\sigma)-\log(2\pi)\bigr)\\
       &-\frac{120\log5}{5^{2\sigma}-1},
       \qquad \psi=\Gamma'/\Gamma .
 \end{split}
 \tag{MA-RB.71}
\]
For every actual field in this family, \(H_\Delta(\sigma)>0\)
on the entire indicated half-line.  For an integer \(k\ge1\), put
\[
 \begin{aligned}
 F_{\Delta,k}(\sigma)
    &=(\sigma-\tfrac12)^2
          -\frac{9k(2\sigma-1)}{H_\Delta(\sigma)},\\
 J_{\Delta,k}&=\inf_{\sigma>1}F_{\Delta,k}(\sigma).
 \end{aligned}
 \tag{MA-RB.72}
\]
This is a finite real number, strictly less than \(1/4\),
and its infimum is attained at a point of \((1,\infty)\).

Suppose the remaining comparison-failure branch of
MA-RB-002 occurs.  Let \(\gamma\in(1/2,1)\) be its
rightmost real abscissa and \(k\ge1\) the order of the
comparison pole there.  Then
\[
 \begin{gathered}
 9k\left(\frac1{\sigma-\gamma}
               +\frac1{\sigma-1+\gamma}\right)
                  \le H_\Delta(\sigma)\qquad(\sigma>1),\\
             (\gamma-\tfrac12)^2\le J_{\Delta,k}.
 \end{gathered}
 \tag{MA-RB.73}
\]
Consequently \(J_{\Delta,k}\le0\) excludes that branch;
otherwise its exact bound from these retained terms is
\[
             \frac12<\gamma
                  \le\frac12+\sqrt{J_{\Delta,k}}<1.
 \tag{MA-RB.74}
\]
Since \(J_{\Delta,k}\le J_{\Delta,1}\), the criterion
\(J_{\Delta,1}\le0\) excludes the comparison-failure
branch without knowing its pole order.  If
\(J_{\Delta,1}>0\), the residual shared-zero criterion
in MA-RB-002 can be confined to
\[
 \begin{gathered}
 B(\gamma)=0,\qquad L(\gamma,V\eta)L(\gamma,X)=0,\\
 \frac12<\gamma\le\frac12+\sqrt{J_{\Delta,1}}.
 \end{gathered}
 \tag{MA-RB.75}
\]
These statements do not assert that either the numerical
criterion or the remaining nonvanishing condition holds
for every discriminant in the family.

\begin{proof}
\emph{The known finite-prime term and positivity of the denominator.}
The order of a \(2B\) element is two.  Since five is
unramified and \(L/\mathbb Q\) is Galois of degree 120,
every prime above five has residue degree two and there
are exactly sixty such primes.  Its complete contribution
to the positive Euler logarithmic derivative is therefore
\[
 E_5(\sigma)
    =60\sum_{j\ge1}\frac{\log25}{25^{j\sigma}}
    =\frac{120\log5}{5^{2\sigma}-1}.
 \tag{MA-RB.76}
\]
In particular this retains all powers of five, not only
the first local coefficient.

The completed Dedekind function and the constant in its
Hadamard logarithmic derivative were fixed and proved in
MA-RB.8--.12.  They give the exact, absolutely convergent
identity
\[
 \begin{split}
 H_\Delta(\sigma)
  ={}&\sum_\rho\operatorname{Re}\frac1{\sigma-\rho}\\
     &+\left(-\frac{\zeta_L'}{\zeta_L}(\sigma)
                                      -E_5(\sigma)\right).
 \end{split}
 \tag{MA-RB.77}
\]
The sum includes multiplicities of all nontrivial completed
zeros.  The unconditional containing strip
\(0\le\operatorname{Re}\rho\le1\) makes every summand
nonnegative for \(\sigma>1\).  The parenthesized term
is the sum of the strictly positive Euler terms at all
prime ideals not above five.  It is strictly positive,
for example because there is a prime ideal above two.
This proves \(H_\Delta>0\) for an actual field, independently
of the proposed comparison failure or the existence of
any specified zero.  Only the usual containing strip is
used; no critical-line location is used.

\emph{The two retained real zeros.}
MA-RB-002 proves that \(\zeta_L\) has a zero of
multiplicity at least \(9k\) at \(\gamma\).
The completed functional equation
\(\xi_L(s)=\xi_L(1-s)\) gives the same multiplicity
at \(1-\gamma\).  These points are distinct, lie in
\((0,1)\), and none of the completion factors has a
zero or pole there.  Thus both occurrences belong to
the nontrivial completed zero divisor in (MA-RB.77).
Retaining them and dropping only nonnegative other
terms gives the first assertion of (MA-RB.73).

Write \(a=\gamma-1/2\) and \(x=\sigma-1/2\).
Then \(0<a<1/2<x\), and their contribution is
\[
 9k\left(\frac1{x-a}+\frac1{x+a}\right)
          =\frac{18kx}{x^2-a^2}.
 \tag{MA-RB.78}
\]
All denominators are positive.  Multiplying the first
inequality of (MA-RB.73) by them gives
\(a^2\le x^2-18kx/H_\Delta(\sigma)\).
Taking the infimum over \(\sigma>1\) gives exactly
the second assertion.  The conclusions for nonpositive
and positive \(J\) follow, and monotonicity in \(k\)
follows directly from the positive denominator in
(MA-RB.72).  MA-RB-002 already identifies the two
shared vanishing factors at this actual \(\gamma\);
only their permitted interval is changed in (MA-RB.75).

\emph{Existence and range of the variational minimum.}
All terms in \(H_\Delta(1+t)-1/t\) are smooth at
\(t=0\).  Hence
\[
 \begin{aligned}
 H_\Delta(1+t)&=t^{-1}+O_\Delta(1),\\
 F_{\Delta,k}(1+t)
       &=\frac14+(1-9k)t+O_{\Delta,k}(t^2)
                         \qquad(t\downarrow0).
 \end{aligned}
 \tag{MA-RB.79}
\]
The coefficient of \(t\) is negative.  Thus the
continuous extension of \(F\) at \(\sigma=1\) has
value \(1/4\), and \(F\) is strictly below that
value at some interior point.
For completeness, take a gamma random variable \(T\)
of shape \(\sigma\) and rate one.  As in MA-RB-001,
\(\psi(\sigma)=\mathbb E\log T\) and
\(\mathbb ET=\sigma\).  When \(\sigma>1\), the gamma
recurrence also gives \(\mathbb E(T^{-1})=(\sigma-1)^{-1}\).
Jensen's inequality applied to \(T\) and \(T^{-1}\)
therefore gives
\(\log(\sigma-1)\le\psi(\sigma)\le\log\sigma\).
In particular \(\psi(\sigma)=\log\sigma+O(1/\sigma)\),
which implies
\[
 H_\Delta(\sigma)=60\log\sigma+O_\Delta(1),\qquad
 F_{\Delta,k}(\sigma)\longrightarrow+\infty
                       \quad(\sigma\longrightarrow\infty).
 \tag{MA-RB.80}
\]
On every compact subinterval of \((1,\infty)\),
\(H\) is positive and continuous, so \(F\) is finite
and continuous.  These facts confine the infimum to
a compact subinterval strictly inside \((1,\infty)\).
It is attained there, is finite, and is strictly below
\(1/4\).  This proves all assertions about the minimum.
\end{proof}

\paragraph{Strict improvement over the previous endpoint.}
The variational restriction is strictly stronger than the
one-sided bound of MA-RB.21 whenever the hypothetical
branch is feasible.  Here is a direct comparison that
does not require numerically finding the minimizer.
Put
\[
 \begin{gathered}
 \mathcal H_L=1+\tfrac12\log\Delta,\\
 t_0=\frac{3\sqrt{k}-1}{\mathcal H_L},\qquad
 \delta_0=\frac{(3\sqrt{k}-1)^2}{\mathcal H_L}.
 \end{gathered}
 \tag{MA-RB.81}
\]
Under that branch, MA-RB-002 proves \(0<t_0<1\)
and \(0<\delta_0\le1-\gamma<1/2\).
The exact old optimization gives
\[
              \frac{9k}{t_0+\delta_0}
                            =\frac1{t_0}+\mathcal H_L.
 \tag{MA-RB.82}
\]
For \(1<\sigma\le2\), the gamma-integral Jensen
bound proved in MA-RB-001 gives
\(\psi(\sigma)\le\log\sigma<\log(2\pi)\).
Also \(1/\sigma<1\), and \(E_5(\sigma)>0\).
Consequently
\[
 \begin{split}
 H_\Delta(1+t_0)&<\frac1{t_0}+\mathcal H_L\\
       &<9k\left(\frac1{t_0+\delta_0}
                         +\frac1{1+t_0-\delta_0}\right).
 \end{split}
 \tag{MA-RB.83}
\]
The last expression is (MA-RB.78) evaluated at
\(x=1/2+t_0\) and \(a_0=1/2-\delta_0>0\).
The same positive-denominator algebra now gives
\[
 J_{\Delta,k}\le F_{\Delta,k}(1+t_0)
                         <(\tfrac12-\delta_0)^2.
 \tag{MA-RB.84}
\]
If the branch survives, \(J>0\), so its new upper
endpoint is strictly smaller than \(1-\delta_0\).
If \(J\le0\), it has instead been excluded.
For a pole order not yet known, the choice \(k=1\)
therefore gives a strict refinement of the universal
endpoint \(1-4/\mathcal H_L\) whenever that branch
had remained feasible.

\paragraph{Source and method scope.}
The prime-five hypothesis is Calegari's Theorem~1.1,
assumption~(2), in the supplied \texttt{S5ArtinWeb.tex}
lines 123--136; assumption~(1) supplies the totally
imaginary signature.  The class \((12)(34)\) has
order two; its label \(2B\) follows the checked table.
The completed Hadamard identity, including its additive
constant and all sixty complex gamma factors, is the
already proved MA-RB.8--.12.  The multiplicity \(9k\)
and the shared-zero reduction are MA-RB-002.
The present optimization is a deduction from those
exact inputs, not a numerical test of a field or an
additional automorphy or zero-location theorem.
It does not use \(\mathsf A\), \(M\), or
\((\mathcal R_B)\).

\section{A universal Euler contribution at every rational prime}
\label{sec:s5-universal-euler}

The preceding variational argument can retain a single
lower bound for all of the remaining Euler terms,
using the discriminant to distinguish ramified from
unramified primes. This is one uniform estimate from
the finite \(S_5\)-cover; it is not a sequence of
separate retained-prime estimates.
Keep the actual normal closure \(L/\mathbb Q\),
\(\Delta=\Delta_L\), \(H_\Delta\), \(F_{\Delta,k}\),
and \(J_{\Delta,k}\) of MA-RB-007.

\paragraph{MA-RB-008 (the universal all-prime variational bound).}
For \(\sigma>1\), define
\[
 \begin{split}
 \mathcal B_\Delta(\sigma)
   &=120\sum_{\substack{p\ {\rm prime}\\p\ne5}}
                   \frac{\log p}{p^{6\sigma}-1}
        -114\sum_{p\mid\Delta}\frac{\log p}{p^{6\sigma}-1}\\
   &=-120\,\frac{\zeta_{\mathbb Q}'}{\zeta_{\mathbb Q}}(6\sigma)
       -\frac{120\log5}{5^{6\sigma}-1}
        -114\sum_{p\mid\Delta}\frac{\log p}{p^{6\sigma}-1},\\
 \mathcal H_\Delta(\sigma)
   &=H_\Delta(\sigma)-\mathcal B_\Delta(\sigma).
 \end{split}
 \tag{MA-RB.85}
\]
The prime on \(\zeta_{\mathbb Q}\) means differentiation
with respect to its argument. The hypothesis at five
gives \(5\nmid\Delta\). Thus \(\mathcal B_\Delta\)
retains coefficient \(120\) at every prime not dividing
\(5\Delta\), and coefficient \(6\) at every prime
dividing \(\Delta\). The prime five is excluded, since
its exact full Euler contribution is already retained
in \(H_\Delta\).

For every actual field in the specified family,
\[
             0<\mathcal H_\Delta(\sigma)<H_\Delta(\sigma)
                                  \qquad(\sigma>1).
 \tag{MA-RB.86}
\]
For each integer \(k\ge1\), put
\[
 \begin{aligned}
 \mathcal F_{\Delta,k}(\sigma)
   &=(\sigma-\tfrac12)^2
           -\frac{9k(2\sigma-1)}{\mathcal H_\Delta(\sigma)},\\
 \mathcal J_{\Delta,k}
   &=\inf_{\sigma>1}\mathcal F_{\Delta,k}(\sigma).
 \end{aligned}
 \tag{MA-RB.87}
\]
This infimum is finite and attained in \((1,\infty)\), and
\[
                  \mathcal J_{\Delta,k}<J_{\Delta,k}<\tfrac14.
 \tag{MA-RB.88}
\]
If the comparison-failure branch of MA-RB-002 occurs
with its actual pole order \(k\) at \(\gamma\), then
\[
 \begin{gathered}
 9k\left(\frac1{\sigma-\gamma}
             +\frac1{\sigma-1+\gamma}\right)
                 \le\mathcal H_\Delta(\sigma)
                                  \qquad(\sigma>1),\\
                  (\gamma-\tfrac12)^2
                            \le\mathcal J_{\Delta,k}.
 \end{gathered}
 \tag{MA-RB.89}
\]
In particular \(\mathcal J_{\Delta,1}\le0\) excludes
that branch. Otherwise the current shared-real-zero
criterion can be confined to
\[
 \begin{gathered}
 B(\gamma)=0,\qquad
       L(\gamma,V\eta)L(\gamma,X)=0,\\
 \tfrac12<\gamma
       \le\tfrac12+\sqrt{\mathcal J_{\Delta,1}}<1 .
 \end{gathered}
 \tag{MA-RB.90}
\]
If the previous interval in MA-RB-007 remained nonempty,
this either excludes its branch or strictly decreases
its upper endpoint. Neither the numerical exclusion
criterion nor nonvanishing on the remaining interval
is asserted for every discriminant in the family.

\begin{proof}
\emph{Residue degrees at ramified as well as unramified primes.}
Fix a rational prime \(p\) and a prime \(\mathfrak p\)
of \(L\) above it. Let \(D_{\mathfrak p}\) and
\(I_{\mathfrak p}\) be its decomposition and inertia
groups in \(G=\operatorname{Gal}(L/\mathbb Q)=S_5\).
Their residue-field exact sequence is
\[
 1\longrightarrow I_{\mathfrak p}
  \longrightarrow D_{\mathfrak p}
  \longrightarrow
       \operatorname{Gal}(\mathbb F_{p^{f_p}}/\mathbb F_p)
  \longrightarrow1 .
 \tag{MA-RB.91}
\]
Here \(f_p\) is the residue degree. The last group
is generated by the arithmetic residue Frobenius
\(u\mapsto u^p\) and has order \(f_p\). Choose any
lift \(d_p\in D_{\mathfrak p}\) of that generator.
The order of its image divides the order of \(d_p\).
The possible element orders in \(S_5\), by the cycle
types \(1,2A,2B,3A,6A,4A,5A\), are respectively
\(1,2,2,3,6,4,5\). Therefore
\[
                     1\le f_p\le6
                         \quad\hbox{for every rational prime }p.
 \tag{MA-RB.92}
\]
No assumption that \(I_{\mathfrak p}\) is trivial
has entered this argument. In particular it applies
to primes dividing the discriminant.

Because \(L/\mathbb Q\) is Galois, all primes over
\(p\) have this residue degree, a common ramification
index \(e_p\), and a common decomposition-group order.
Their number is \(g_p=120/(e_pf_p)\ge1\).
The local factor of the Dedekind zeta function is
\((1-p^{-f_ps})^{-g_p}\); its contribution to the
positive logarithmic derivative on the real
half-line is
\[
 E_p(\sigma)
     =\frac{g_pf_p\log p}{p^{f_p\sigma}-1}.
 \tag{MA-RB.93}
\]
This includes every power of every prime ideal
over \(p\), also when it is ramified.

For a fixed \(x>1\), the function
\(u\mapsto u/(x^u-1)\) is strictly decreasing on
\(u>0\). Indeed, after multiplying its derivative
by the positive denominator, its sign is that of
\(e^v(1-v)-1\), where \(v=u\log x>0\).
This last expression is zero at \(v=0\) and has
derivative \(-ve^v<0\).
Use \(x=p^\sigma\), \(f_p\le6\), and \(g_p\ge1\)
in (MA-RB.93). It follows that
\(E_p(\sigma)\ge6\log p/(p^{6\sigma}-1)\) at every
prime. At an unramified prime there is the stronger
identity \(g_pf_p=120\). Since \(f_p\le6\),
\[
 E_p(\sigma)=\frac{120\log p}{p^{f_p\sigma}-1}
                  \ge\frac{120\log p}{p^{6\sigma}-1}.
\]
Combining the two cases gives the uniform bound
\[
 E_p(\sigma)\ge c_p\frac{\log p}{p^{6\sigma}-1},
 \qquad
 c_p=
 \begin{cases}
 120,&p\nmid\Delta,\\
 6,&p\mid\Delta,
 \end{cases}
 \qquad \sigma>1.
 \tag{MA-RB.94}
\]
This is a bound for the full local Euler contribution,
not a claim that a ramified Frobenius itself has a
specified conjugacy class in \(S_5\). The comparison
is pointwise on the real half-line. It is not
coefficientwise positivity of the difference of
Dirichlet series, nor an inclusion of Artin motives.

\emph{Strict positivity after the universal subtraction.}
All the positive Euler sums just used converge
absolutely on \(\sigma>1\); the comparison series
in (MA-RB.85) converges already when \(6\sigma>1\).
The ordinary rational zeta Euler product gives
the logarithmic-derivative expression there, with
exactly the stated coefficient \(120\), the
exclusion of five, and the finite ramified correction
of \(120-6=114\).
Subtracting \(\mathcal B_\Delta\) from the exact identity
(MA-RB.77) gives
\[
 \begin{split}
 \mathcal H_\Delta(\sigma)
   ={}&\sum_\rho\operatorname{Re}\frac1{\sigma-\rho}\\
      &+\sum_{p\ne5}
          \left(E_p(\sigma)
                  -c_p\frac{\log p}{p^{6\sigma}-1}\right).
 \end{split}
 \tag{MA-RB.95}
\]
The zero sum includes all nontrivial completed
zeros with multiplicity and is absolutely convergent
in this real-part form, as proved in MA-RB-001.
Its terms are strictly positive for each zero in
the usual containing strip
\(0\le\operatorname{Re}\rho\le1\).
Every Euler remainder is nonnegative by
(MA-RB.94).

We verify that the completed divisor is nonempty,
without invoking any particular zero or zero location.
If it were empty, the same exact Hadamard identity
of MA-RB.8--.12 would give, for every \(\sigma>1\),
\[
 0=\frac1\sigma+\frac1{\sigma-1}
       +\frac12\log\Delta
       +60\bigl(\psi(\sigma)-\log(2\pi)\bigr)
       +\frac{\zeta_L'}{\zeta_L}(\sigma).
\]
The Euler series implies
\(\zeta_L'/\zeta_L(\sigma)\to0\) as \(\sigma\to\infty\):
each positive term tends to zero, and for
\(\sigma\ge2\) the series at \(2\) is an absolutely
summable majorant. The Jensen estimate already proved
in MA-RB-007 gives
\(\psi(\sigma)=\log\sigma+O(1/\sigma)\).
The right side would therefore be
\(60\log\sigma+O_\Delta(1)\), a contradiction.
The completed divisor is consequently nonempty.
Its sum in (MA-RB.95) is strictly positive at every
\(\sigma>1\). This proves
\(\mathcal H_\Delta(\sigma)>0\) even if every retained
Euler remainder happens to vanish.

Also \(\mathcal B_\Delta(\sigma)>0\), since its
positive-coefficient form is
\[
 120\sum_{p\nmid5\Delta}\frac{\log p}{p^{6\sigma}-1}
       +6\sum_{p\mid\Delta}\frac{\log p}{p^{6\sigma}-1}.
\]
This proves the other strict inequality in (MA-RB.86).
Discarding the discriminant information would instead
retain only \(6\sum_{p\ne5}\log p/(p^{6\sigma}-1)\).
The present single bound is strictly larger, because
there is an unramified prime different from five.
Only the discriminant-aware bound is used in the
minimum (MA-RB.87).

\emph{The retained real pair and the minimum.}
Under the hypothetical comparison failure, the
distinct real zeros \(\gamma,1-\gamma\) each
have multiplicity at least \(9k\), as proved
in MA-RB-002/007. Retain them in (MA-RB.95)
and discard only nonnegative terms. This gives
the first line of (MA-RB.89).
With \(a=\gamma-1/2\) and \(x=\sigma-1/2\),
their contribution is \(18kx/(x^2-a^2)\).
All denominators are positive. The same algebra
as (MA-RB.78) gives
\(a^2\le\mathcal F_{\Delta,k}(\sigma)\), and
taking the infimum proves the second line.

The function \(\mathcal B_\Delta\) is smooth on a
neighborhood of \(\sigma=1\) and tends to zero
at infinity. Indeed its differentiated prime
series converge uniformly on every closed
subinterval of \(6\sigma>1\).
Consequently the exact endpoint estimates of
MA-RB-007 become
\[
 \begin{aligned}
 \mathcal H_\Delta(1+t)&=t^{-1}+O_\Delta(1),\\
 \mathcal F_{\Delta,k}(1+t)
       &=\tfrac14+(1-9k)t+O_{\Delta,k}(t^2),\\
 \mathcal H_\Delta(\sigma)
       &=60\log\sigma+O_\Delta(1),\qquad
 \mathcal F_{\Delta,k}(\sigma)\longrightarrow+\infty .
 \end{aligned}
 \tag{MA-RB.96}
\]
The first two assertions concern \(t\downarrow0\);
the last line concerns \(\sigma\to\infty\).
The digamma estimate needed there is already
proved by the two Jensen inequalities in
MA-RB-007, so no new gamma asymptotic is assumed.
Continuity and strict positivity of
\(\mathcal H_\Delta\), the negative linear
coefficient \(1-9k\), and (MA-RB.96) imply
that the infimum in (MA-RB.87) is finite
and attained on a compact subinterval of
\((1,\infty)\), with value less than \(1/4\).

To obtain the strict improvement, let
\(\sigma_*\in(1,\infty)\) attain the old
minimum \(J_{\Delta,k}\), whose existence
is part of MA-RB-007. The strict positive
denominator comparison (MA-RB.86) gives
\[
 \mathcal J_{\Delta,k}
   \le\mathcal F_{\Delta,k}(\sigma_*)
   <F_{\Delta,k}(\sigma_*)=J_{\Delta,k}.
\]
This proves (MA-RB.88), including strictness,
without a numerical optimization or a
compactness assumption about a limiting
minimizer.

Finally, \(\mathcal F_{\Delta,k}\) decreases
pointwise with \(k\), so
\(\mathcal J_{\Delta,k}\le\mathcal J_{\Delta,1}\).
If the latter is nonpositive, (MA-RB.89)
contradicts \(\gamma>1/2\). Otherwise
(MA-RB.89) gives the interval in (MA-RB.90),
and MA-RB-002 supplies its two vanishing
factors. This proves all stated conclusions.
\end{proof}

\paragraph{Finite-cover and Artin-motive meaning of the local bound.}
The representation underlying \(\zeta_L\) is the regular
permutation representation \(\mathbb Q[G]\), the
Artin motive of the finite cover \(\operatorname{Spec}L\)
over \(\operatorname{Spec}\mathbb Q\).
At a possibly ramified prime, take
\(\mathbb Q[G]^{I_{\mathfrak p}}\).
For the left regular action, an invariant basis is
indexed by the \(I_{\mathfrak p}\)-orbits \(I_{\mathfrak p}h\).
There are \(120/e_p=g_pf_p\) of them.
The chosen lift \(d_p\) normalizes inertia and
permutes these orbits. Its \(r\)th power fixes
one if and only if \(d_p^r\in I_{\mathfrak p}\),
which is equivalent to \(f_p\mid r\).
Every cycle therefore has length \(f_p\),
and there are exactly \(g_p\) cycles. Thus
\[
 \det(1-Td_p\mid\mathbb Q[G]^{I_{\mathfrak p}})
                         =(1-T^{f_p})^{g_p}.
\]
This proves directly the Artin local-factor
interpretation of (MA-RB.93).
We used arithmetic residue Frobenius. Replacing
it by its inverse gives the same determinant
for these complete cyclic permutation blocks.
The argument concerns the finite cover and
its local inertia invariants; it does not
assert an effective or positive realization
of an unknown analytic zero divisor.

\paragraph{Established locators and scope.}
Calegari's Theorem~1.1, assumptions (1)--(2),
supplies the actual \(S_5\) normal closure,
signature and exact prime-five datum; the
supplied primary LaTeX locator is
\texttt{S5ArtinWeb.tex}, lines 123--136.
The standard Galois decomposition-group
sequence in (MA-RB.91) fixes the residue-field
convention, and the local regular-representation
calculation above explains every ramified
Euler coefficient used here.

For the finite-cover interpretation, the exact
inherited locators are SGA~1, Expos\'e~V, \S7,
and Deligne (1979), \emph{Valeurs de fonctions
\(L\) et p\'eriodes d'int\'egrales},
Definition~6.1 and \S\S6.2--6.3; equation
(0.5.3) supplies the previously fixed
restriction-of-scalars/induction convention.
Those categorical locators are the established
ones in MA-S5; no new claim that they prove
the analytic inequality is made.
The exact Hadamard identity and constant are
MA-RB.8--.12, the shared-zero multiplicity is
MA-RB-002, and the preceding variational
construction is MA-RB-007. The new proof uses
no further functoriality, critical-line
location, \(\mathsf A\), \(M\), or
\((\mathcal R_B)\).

\section{A genuine family on which the variational exclusion is nondecisive}
\label{sec:s5-endpoint-scope}

The exact all-prime inequality in MA-RB-008 narrows an actual
comparison obstruction. Its numerical exclusion criterion
cannot hold uniformly on Calegari's family. We prove this
using genuine fields satisfying both specified local
hypotheses. No zero of their \(L\)-functions is constructed.

\paragraph{MA-RB-009 (scope of the universal numerical criterion).}
Retain the final discriminant-aware functions
\(\mathcal H_\Delta,\mathcal F_{\Delta,k},\mathcal J_{\Delta,k}\)
of MA-RB-008, and put
\[
 C_0=60(1+\log(2\pi))+5\log5+\frac{40}{7},
 \qquad C_\Delta=\frac12\log\Delta-C_0 .
 \tag{MA-RB.97}
\]
For every actual field in the specified family and every
integer \(k\ge1\), if \(C_\Delta\ge18k\), then
\[
 \frac14-\frac{9k}{C_\Delta}
       \le\mathcal J_{\Delta,k}<\frac14.
 \qquad
 C_\Delta\ge40k\ \Longrightarrow\
                      \mathcal J_{\Delta,k}\ge\frac1{40}>0 .
 \tag{MA-RB.98}
\]

There are genuine fields in this family with arbitrarily
large normal-closure discriminant. An explicit construction
is as follows. For each rational prime \(q>5\), choose
the least positive solution \(a_q\) to the two congruences
\[
 qa_q\equiv2\pmod3,\qquad qa_q\equiv1\pmod5,
 \qquad
 f_q(X)=X^5+qa_qX+15q .
 \tag{MA-RB.99}
\]
Let \(K_q=\mathbb Q[X]/(f_q)\) and let \(L_q\) be its
normal closure. Then
\[
 \begin{gathered}
 \operatorname{Gal}(L_q/\mathbb Q)=S_5,\qquad
 [L_q:\mathbb Q]=120,\qquad
 (r_1(L_q),r_2(L_q))=(0,60),\\
 \text{complex conjugation has type }(12)(34),\\
 5\text{ is unramified and }
       \operatorname{Frob}_5\text{ has type }(12)(34),\qquad
 q\mid\Delta_{L_q}.
 \end{gathered}
 \tag{MA-RB.100}
\]
Consequently, for every fixed integer \(k\ge1\),
\[
                \mathcal J_{\Delta_{L_q},k}
                        \longrightarrow\frac14
                   \qquad(q\longrightarrow\infty
                         \text{ through primes}).
 \tag{MA-RB.101}
\]
In particular the criterion
\(\mathcal J_{\Delta,1}\le0\) does not exclude comparison
failure uniformly for this actual family.
This is a limitation of that criterion, not a construction
of comparison failure, a shared real zero, or a
\(\beta\)-excess example. It does not reduce the remaining
shared-zero gate.

\begin{proof}
\emph{Construction, irreducibility, and real signature.}
Because \(q\) is prime to \(15\), the Chinese remainder
theorem supplies \(a_q\), with \(1\le a_q\le15\).
In particular \(qa_q>0\).
The polynomial \(f_q\) is Eisenstein at \(q\):
every nonleading coefficient is divisible by \(q\),
and its constant coefficient has \(q\)-valuation one.
Here the irreducibility argument is elementary.
A factorization into two monic integer polynomials
of positive degree, which a rational factorization of
a monic polynomial would supply, reduces modulo \(q\)
to a factorization of \(X^5\). Both constant terms
would be divisible by \(q\), contrary to
\(q^2\nmid15q\).
Thus \(f_q\) is irreducible over \(\mathbb Q\), and
its Galois group acts transitively on its five roots.

On the real line,
\[
                         f_q'(x)=5x^4+qa_q>0.
\]
The polynomial tends to opposite infinities at the
two ends of that line. Hence it has exactly one
real root. Characteristic zero and irreducibility
give separability; the other four roots therefore
form two nonreal conjugate pairs. Complex conjugation
on the five roots has cycle type \(2B=(12)(34)\).

\emph{The two finite Frobenius calculations.}
The defining congruences give the exact factorizations
\[
 \begin{aligned}
 f_q(X)\bmod3
   &=X^5-X=X(X-1)(X+1)(X^2+1),\\
 f_q(X)\bmod5
   &=X^5+X=X(X^2+2)(X^2-2).
 \end{aligned}
 \tag{MA-RB.102}
\]
Over \(\mathbb F_3\), the first three linear factors
are distinct, and \(X^2+1\) is irreducible, since
\(-1\) is not a square. All four factors are distinct.
Over \(\mathbb F_5\), neither \(2\) nor \(-2=3\)
is a square, since the nonzero squares are \(1,4\).
The two quadratic factors and \(X\) are distinct.
Both reductions are therefore squarefree.

We spell out the local interpretation being used.
The roots of either reduction belong to a finite
residue-field extension and lift uniquely, by the
simple-root Hensel property, in the corresponding
unramified local extension. All characteristic-zero
roots of \(f_q\) lie there: there are five distinct
lifts and the polynomial has degree five.
Thus the local splitting extension is unramified
at \(3\) and at \(5\). The arithmetic residue
Frobenius permutes the reduced roots in cycles whose
lengths are their irreducible factor degrees.
Uniqueness of the lifts gives the same permutation
on the characteristic-zero roots.
This proves that \(3\) supplies a transposition,
while \(5\) supplies the specified \(2B\) element.
It also proves the required unramified assertion
in the normal closure, rather than merely in
the degree-five subfield.

\emph{Why the transitive group is all of \(S_5\).}
Let \(G\) be that transitive group, and let
\(\tau\in G\) be the transposition just obtained.
On the five roots, form a graph whose edges are
the two-element supports of all \(G\)-conjugates
of \(\tau\). The graph has an edge and is
\(G\)-invariant. Transitivity implies that all
its connected components have the same size.
That size divides the prime number \(5\), and
it is not one because the graph has an edge.
The graph is therefore connected.
The transpositions associated with the edges
of a connected graph generate its full symmetric
group: along a path, conjugating an edge transposition
by the preceding edge transpositions produces the
transposition of the first and last vertices.
All those edge transpositions lie in \(G\).
It follows that \(G=S_5\).

In particular the normal closure has degree \(120\).
Its complex conjugation is nontrivial. A Galois
number field with a real embedding would, after
that embedding, have every conjugate in \(\mathbb R\)
and hence trivial complex conjugation.
Thus \(L_q\) is totally imaginary, giving signature
\((0,60)\).

\emph{Ramification at the varying prime and unbounded discriminants.}
Extend the \(q\)-adic valuation to \(L_q\), normalized
by \(v(q)=1\), and let \(\theta\) be any root of \(f_q\).
If \(v(\theta)<0\), then \(\theta^5\) has valuation
strictly smaller than both other terms of the
polynomial equation, which is impossible for a sum
equal to zero. The same contradiction holds if
\(v(\theta)=0\).
If \(v(\theta)>0\), the linear term has valuation
at least \(1+v(\theta)>1\), while the constant
term has valuation one. Thus the minimum valuation
can occur twice only when
\[
                   5v(\theta)=1.
 \qquad
 v(\theta)=\frac15,\qquad
                    5\mid e_q(L_q/\mathbb Q).
 \tag{MA-RB.103}
\]
The last conclusion follows from the value group
\(e_q^{-1}\mathbb Z\) of this extended valuation.
Therefore \(q\) is ramified and \(q\mid\Delta_{L_q}\).
Here the standard discriminant criterion has its
literal finite-algebra meaning: ramification gives
a nonzero nilpotent in \(\mathcal O_{L_q}/q\mathcal O_{L_q}\);
multiplication by its product with any element is
nilpotent and has trace zero. The reduced trace pairing
is degenerate, so its discriminant vanishes modulo \(q\).
Consequently \(\Delta_{L_q}\ge q\), and these
discriminants tend to infinity with \(q\).
There are also infinitely many distinct normal
closures: finitely many number fields have only
finitely many ramified primes in total.
This completes all assertions in (MA-RB.100).

\emph{Uniform control of the retained Euler and gamma terms.}
The coefficients in the positive-prime form of
\(\mathcal B_\Delta\) are at most \(120\).
For all \(\sigma>1\),
\[
 \begin{split}
 0<\mathcal B_\Delta(\sigma)
 &\le120\sum_{n=2}^{\infty}\frac{\log n}{n^6-1}\\
 &\le120\,\frac{64}{63}\sum_{n=2}^{\infty}n^{-5}\\
 &\le120\,\frac{64}{63}
             \left(\frac1{32}+\int_2^\infty x^{-5}\,dx\right)
   =\frac{40}{7}.
 \end{split}
 \tag{MA-RB.104}
\]
The second line uses \(\log n\le n\) and
\(n^6-1\ge(63/64)n^6\).
This bound is uniform in the discriminant and its
ramified prime support. Also the exact known
prime-five term satisfies
\[
          E_5(\sigma)=\frac{120\log5}{5^{2\sigma}-1}
                                    \le5\log5.
\]
The gamma recurrence
\(\psi(\sigma+1)=\psi(\sigma)+1/\sigma\),
and the already proved Jensen lower bound at
\(\sigma+1\), give
\[
 \psi(\sigma)
     \ge\log\sigma-\frac1\sigma\ge-1
                                    \qquad(\sigma>1).
 \tag{MA-RB.105}
\]
Thus the exact denominator in MA-RB-008 obeys
\[
 \begin{aligned}
 \mathcal H_\Delta(\sigma)&\ge C_\Delta,\\
 \mathcal F_{\Delta,k}(\sigma)
   &=x^2-\frac{18kx}{\mathcal H_\Delta(\sigma)}
      \ge x^2-\frac{18kx}{C_\Delta},
          \qquad x=\sigma-\tfrac12>\tfrac12,
 \end{aligned}
 \tag{MA-RB.106}
\]
where the second line is used when \(C_\Delta>0\).
The first line drops only the positive rational
terms \(1/\sigma+1/(\sigma-1)\) and uses
(MA-RB.104)--(MA-RB.105).
If \(C_\Delta\ge18k\), the polynomial on the right
of the second line is increasing for \(x\ge1/2\):
its derivative is \(2x-18k/C_\Delta\ge0\).
Its value at \(x=1/2\) is \(1/4-9k/C_\Delta\).
Taking the infimum and using the already proved
strict upper bound \(\mathcal J_{\Delta,k}<1/4\)
gives (MA-RB.98). At \(C_\Delta\ge40k\) its lower
bound is at least \(1/4-9/40=1/40\).

Finally \(\Delta_{L_q}\to\infty\) implies
\(C_{\Delta_{L_q}}\to\infty\).
For each fixed \(k\), the two bounds in
(MA-RB.98) squeeze its variational minimum
to \(1/4\). This proves (MA-RB.101).
\end{proof}

\paragraph{Actual arithmetic content and limitation.}
The polynomial construction imposes the complex-conjugation
and prime-five hypotheses of Calegari's Theorem~1.1,
assumptions (1)--(2), in the supplied
\texttt{S5ArtinWeb.tex}, lines 123--136.
Calegari's theorem is not being cited as the construction
of this family; that construction and its Galois proof
are given above.
The residue-Frobenius convention is the same arithmetic
convention as (MA-RB.91) and the local regular Artin
factor in MA-RB-008.
The finite-cover placement is the established SGA~1,
Expos\'e~V, \S7 correspondence between finite \'etale
covers and finite continuous fundamental-group sets,
with the Artin-motive interpretation of Deligne (1979),
Definition~6.1 and \S\S6.2--6.3.
The elementary local argument above uses the standard
simple-root Hensel property and the normalized local
valuation; no new theorem locator is assigned to those
classical field-theory facts.

The uniform bounds use the exact Euler expression of
MA-RB-008 and the gamma recurrence and Jensen argument
already proved in MA-RB-007. They require neither
critical-line location, \(\mathsf A\), \(M\), nor
\((\mathcal R_B)\).
The genuine fields demonstrate that one cannot force
the displayed numerical minimum to be nonpositive
for every member merely from these family hypotheses.
They provide no example of the hypothetical comparison
pole or of either shared vanishing factor.
The remaining problem is still an actual arithmetic
shared-zero exclusion on the shorter interval of
MA-RB-008. No new reduction of that interval or
exclusion of its zeros is claimed here.

\section{The precise residual Galois realization gate}
\label{sec:s5-residual-realization}

The nonregular Galois-representation route has an exact elementary
scope test.  It supplies no new exclusion of the shared real zero
in MA-RB-002/008.
Let \(A=L_f(\,\cdot\,,\varpi)\), and let \(W\) be the specified
five-dimensional Artin representation.  Fix a finite set
\(S\) containing all bad primes of both objects.  At \(p\notin S\),
write
\[
 H_{A,p}(T)=\det(1-T\operatorname{Sat}_p\varpi),
 \qquad H_{W,p}(T)=\det(1-TW(\operatorname{Frob}_p)).
 \tag{MA-RB.107}
\]
These are the characteristic polynomials whose reciprocals
are the local Euler factors.

For a rational prime \(\ell\), let \((G_\ell)\) assert the
existence of a continuous semisimple representation
\[
 r_\ell:G_{\mathbb Q}\longrightarrow
                \operatorname{GL}_5(\overline{\mathbb F}_\ell)
 \tag{MA-RB.108}
\]
with the discrete topology on the target, unramified outside
some finite set \(S_\ell\supset S\cup\{\ell\}\), and satisfying
\[
 \det(1-Tr_\ell(\operatorname{Frob}_p))
                 =H_{A,p}(T)\pmod\ell
                 \quad(p\notin S_\ell).
 \tag{MA-RB.109}
\]
The polynomial on the right is integral, as verified below.
The finite set in this definition may depend on \(\ell\).

\paragraph{MA-RB-010 (odd residual realization is an equivalent gate).}
For every fixed odd prime \(\ell\), the following conditions
are equivalent:
\begin{enumerate}
\item the specified all-place comparison \(A=L_W\) holds;
\item \(\Omega\) is finite, or equivalently the good-prime
      polynomials in (MA-RB.107) agree outside a finite set;
\item \((G_\ell)\) holds.
\end{enumerate}
By contrast, \((G_2)\) holds without the comparison hypothesis.
Thus constructing a compatible representation in any one odd
residual characteristic would settle the comparison, but its
existence is an equivalent formulation of the remaining gate.
The result uses neither \(\mathsf A\), \(M\), nor
\((\mathcal R_B)\), and does not prove \((\mathcal R_B)\).

\begin{proof}
The exact good-prime alternative is the already audited
Calegari theorem following his similitude calculation:
\(H_{A,p}=H_{W,p}\) outside the Dirichlet-density-zero set
\(\Omega\), whereas at \(p\in\Omega\) the eigenvalue
multisets are respectively
\(\{1,i,-i,1,1\}\) and \(\{1,i,-i,-1,-1\}\).
Consequently
\[
 \begin{aligned}
 H_{A,p}(T)&=(1-T)^3(1+T^2),\\
 H_{W,p}(T)&=(1-T)(1+T)^2(1+T^2),\\
 H_{A,p}(T)-H_{W,p}(T)&=-4T(1-T)(1+T^2)
                         \quad(p\in\Omega).
 \end{aligned}
 \tag{MA-RB.110}
\]
In particular the two polynomials are distinct in every odd
characteristic and equal in characteristic two.
At other good primes the Artin polynomial is integral: its
roots are roots of unity, while the rational character of
\(W\), evaluated on all powers of Frobenius, makes its
coefficients rational by Newton identities.  Rational
algebraic integers are integers.  Formula (MA-RB.110) proves
integrality at the remaining good primes as well.

Fix a number-field realization of the finite-image
representation \(W\), and a prime above \(\ell\).
A stable integral lattice is obtained by summing the finitely
many translates of an initial lattice and localizing at that
prime.  Reduction gives a continuous finite-image
representation \(\overline W_\ell\); semisimplification does
not change any characteristic polynomial.  Indeed a
composition series makes every group element block upper
triangular, so its characteristic polynomial is the product
of those on the successive quotients.

Suppose \((G_\ell)\) holds.  Both \(r_\ell\) and
\(\overline W_\ell\) have finite image: a continuous image
of a profinite compact group in a discrete space is finite.
Their product therefore factors through a common finite
Galois quotient \(G\) of \(G_{\mathbb Q}\).
Let \(D\subseteq G\) be the set on which their
characteristic polynomials differ.  It is a union of
conjugacy classes.  If \(D\ne\varnothing\), Chebotarev's
theorem gives the unramified rational primes whose Frobenius
lies in \(D\) the positive Dirichlet density
\(\lvert D\rvert/\lvert G\rvert\).
Outside the finite excluded set, however, each such prime
belongs to \(\Omega\): at every other prime both residual
polynomials equal the reduction of \(H_{W,p}\).
This contradicts the zero Dirichlet density of \(\Omega\).
Hence \(D=\varnothing\).

It follows that the two polynomials in (MA-RB.107) are
congruent at every common good prime.  When \(\ell\) is odd,
(MA-RB.110) excludes every prime of \(\Omega\) outside the
finite excluded set.  Therefore \(\Omega\) is finite.
Conversely, if \(\Omega\) is finite, the semisimplification
of \(\overline W_\ell\) satisfies (MA-RB.109) after adding
\(\Omega\) and its finite ramification set to \(S_\ell\).
This proves the equivalence of the second and third
conditions.  The equivalence of the first and second is the
previously checked finite-exception comparison, Calegari's
Lemma~4.10, already used at the start of MA-RB-002; no
new all-place matching premise is being added here.

Finally (MA-RB.110) vanishes modulo two.  Thus
\(\overline W_2^{\mathrm{ss}}\) satisfies (MA-RB.109)
for all common good primes whether or not \(\Omega\)
is finite.  This proves the assertion about \((G_2)\).
\end{proof}

\paragraph{What the primary construction theorems supply.}
The source comparison is exact enough to check the numerical
Hecke hypotheses without building a new Galois representation.
Delete the fixed finite set \(S\), and let \(N^*\) be an integer
whose prime divisors contain \(S\).  For this retained
Hecke-polynomial collection the coefficient field is
\(\mathbb Q\), the coefficients are integral, and every
field-conjugate collection is the original one.
If \(a_j(p)\) is its coefficient of degree \(j\), every
Satake root has modulus one, so
\[
 |a_j(p)|\le \binom{5}{j},\qquad
 \sum_{p\nmid N^*}|a_j(p)|^2p^{-\sigma}
 \le \binom{5}{j}^{\!2}
                   \log\frac1{\sigma-1}+O(1)
                         \quad(\sigma\downarrow1).
 \tag{MA-RB.111}
\]
For the last bound, subtract \(\sum_p p^{-\sigma}\)
from the Euler logarithm of \(\zeta(\sigma)\).
The sum over powers at least two is bounded uniformly
for \(\sigma>1\), and
\(\log\zeta(\sigma)=\log(1/(\sigma-1))+O(1)\).
Removing finitely many primes changes only the bounded term.

These are the finite-set versions of assumptions (1), (2),
(3), and (5) in Kim and Yamauchi's \(\operatorname{GL}_n\)
construction.  Their Conjecture~5.1 and assumption (4) require
actual compatible semisimple residual Galois representations.
Section~6 reconstructs Artin representations after that input;
it does not construct the input.  The finite-set formulation
is explicit here: no rationality assertion at an omitted place
is hidden in (MA-RB.111), and no application of their complete
theorem without its Galois hypothesis is claimed.
MA-RB-010 shows that even one odd-characteristic instance of
the needed compatibility already contains the comparison
obstruction in this family.

The earlier Kim--Yamauchi \(\operatorname{GSp}_4\) theorem
also assumes residual representations, together with transfer,
integrality, and rationality of the real-analytic Hecke module.
Its stated infinity type is not a holomorphic limit of
discrete series.  Calegari--Geraghty's displayed nonregular
modularity-lifting theorems instead concern holomorphic
weight \((j,2)\), with \(j\ge4\), and Hodge--Tate weights
\([0,0,j-1,j-1]\), besides ordinary and residual-image
conditions.  These are not the all-zero weights of the
finite-image spin candidate.  Thus these precise inspected
theorems do not supply the missing representation here.
This is a hypothesis check, not a claim that no future or
other construction can do so.

\paragraph{Finite-cover interpretation and exact locators.}
The known Artin object \(W\) comes from a representation
of the finite Galois cover after the ramified primes are
removed.  Its integral lattice and residual realizations
are available independently of \(A\).  The proof uses the
common finite quotient of two such continuous residual
representations and its Frobenius classes.  It does not
assign an Artin motive to \(A\) from its finite list of
Hecke polynomials.  The inherited categorical locators are
SGA~1, Expos\'e~V, \S7, and Deligne (1979),
\emph{Valeurs de fonctions \(L\) et p\'eriodes d'int\'egrales},
Definition~6.1 and \S\S6.2--6.3; they are not being cited
as a construction of the missing compatible realization.

Calegari (2013), Theorem~4.7 and Lemma~4.10, supplies the
local alternative, density, and finite-exception closure;
the exact supplied LaTeX spans are
\texttt{S5ArtinWeb.tex}, lines 736--765 and 823--854.
Kim and Yamauchi (2016),
\emph{A conditional construction of Artin representations
for real analytic Siegel cusp forms of weight \((2,1)\)},
Theorem~1.1 and \S3.4, are in arXiv:1307.7619v5,
pp.~2--4 and 14--15,
\url{https://arxiv.org/pdf/1307.7619}; chapter DOI
\url{https://doi.org/10.1090/conm/664/13061}.
Their (2017) \emph{A uniform structure on subgroups of
\(\operatorname{GL}_n(\mathbb F_q)\) and its application to a
conditional construction of Artin representations of
\(\operatorname{GL}_n\)}, \emph{Journal of the Ramanujan
Mathematical Society}, 32(1), 75--99, is arXiv:1403.5535v3:
Theorem~1.1, pp.~2--3; Conjecture~5.1 and \S6,
pp.~13--16; and \S\S7--8, pp.~16--18,
\url{https://arxiv.org/pdf/1403.5535}.
Calegari and Geraghty (2020), \emph{Minimal modularity
lifting for nonregular symplectic representations},
\emph{Duke Mathematical Journal, 169}(5), 801--896,
Theorems~1.1--1.2, published pp.~802--803
(author-hosted advance PDF pp.~2--3),
\url{https://www.math.uchicago.edu/~fcale/papers/Siegel.pdf},
DOI \url{https://doi.org/10.1215/00127094-2019-0044},
supplies the different-weight hypothesis comparison.
The advance copy carries a 2019 copyright notice; the
published article is dated 2020.
The scope statements above do not certify complete independent
proof checks of those construction or modularity theorems.

\section{The precise scope of a positive Rankin comparison}
\label{sec:s5-rankin-scope}

This section tests a proposed replacement of the regular
\(S_5\)-representation by the known positive Rankin functions.
It proves the finite local comparisons, including at ramified
places, and an exact limitation of the resulting leading
near-one score.  It does not compare every possible full
Euler--gamma variational expression, and does not reduce
the arithmetic gate in MA-RB-008.

\paragraph{Review status.}
The complete post-writing author reading and two independent
whole-proof readings have passed.  The source-hypothesis
checks below have their own recorded scope.

Retain the actual extension \(L/\mathbb Q\), its discriminant
\(\Delta\), and the representations and actual order matrix
of MA-RB-006.  Write \(\operatorname{Reg}=\mathbb C[S_5]\).
For a finite Artin representation \(T\), let \(N(T)\)
denote its finite Artin conductor.  At a prime \(p\), use
its actual inertia group \(I_p\) and a lift \(\phi_p\)
of arithmetic residue Frobenius, and set
\[
 t_{T,p}(j)=
       \operatorname{tr}\bigl(\phi_p^j\mid T^{I_p}\bigr),
 \qquad
 E_T(\sigma)=\sum_p\sum_{j\ge1}
             t_{T,p}(j)(\log p)p^{-j\sigma}
              \quad(\sigma>1).
 \tag{MA-RB.150}
\]
For \(T=\operatorname{Reg}\), this is
\(-\zeta_L'/\zeta_L(\sigma)\).

\paragraph{MA-RB-011 (local positivity and the five-factor score).}
Let \(U\) be any finite-dimensional complex representation
of \(S_5\), of dimension \(d\).  Then
\[
 \begin{gathered}
 \log N(\operatorname{End}U)
                       \le \frac{d^2}{120}\log\Delta,\\
 t_{\operatorname{End}U,p}(j)
                  \ge\frac{d^2}{120}\,t_{\operatorname{Reg},p}(j)
                           \qquad(p,\ j\ge1).
 \end{gathered}
 \tag{MA-RB.151}
\]
The second inequality includes ramified primes.
For the known six-dimensional \(X\), its archimedean
parameter in \(\operatorname{End}X\) has dimensions
\(20\) and \(16\) for the trivial and sign characters,
respectively.  At the specified prime five,
\[
 E_{\operatorname{End}X,5}(\sigma)
       -\frac3{10}E_{\operatorname{Reg},5}(\sigma)
       =\frac{4(\log5)5^{-\sigma}}{1-5^{-2\sigma}}.
 \tag{MA-RB.152}
\]

For the separate score assertion, let
\[
 \begin{gathered}
 (U_1,\ldots,U_5)=(W,W\eta,V,V\eta,X),\\
 a_i\ge0,\quad a\ne0,\qquad
 p=\sum_i a_i^2,\quad
 d=5a_1+5a_2+4a_3+4a_4+6a_5,\quad
 m=a^{\mathsf T}Ma ,
 \end{gathered}
 \tag{MA-RB.153}
\]
where \(M\) is exactly (MA-RB.65).  Define the positive-part
score by
\[
 {\cal S}(a)=
       \frac{\max\{\sqrt m-\sqrt p,\,0\}^{\,2}}{d^2}.
 \tag{MA-RB.154}
\]
Then every such choice satisfies the strict bound
\[
                         {\cal S}(a)<\frac4{120}.
 \tag{MA-RB.155}
\]
This is the score obtained from the guaranteed multiplicity
when the unknown comparison order is bounded only by \(k\ge1\).
For the honest Artin subfamily \((V,V\eta,X)\), the conductor
bound in (MA-RB.151) applies to every integer Artin direct
sum, whose unramified factors match the corresponding
automorphic isobaric sum.
For weights involving the automorphic comparison \(A\),
no corresponding Artin-conductor bound is asserted.

\begin{proof}
\emph{The finite inertia and conductor comparisons.}
Average a positive definite Hermitian form over \(S_5\),
so that \(U\) is unitary.  The character of its endomorphism
representation is therefore
\[
             \chi_{\operatorname{End}U}(g)=|\chi_U(g)|^2.
 \tag{MA-RB.156}
\]
For any subgroup \(I\subset S_5\), character averaging gives
\[
 \dim(\operatorname{End}U)^I
   =\frac1{|I|}\sum_{i\in I}|\chi_U(i)|^2
   \ge\frac{d^2}{|I|}.
\]
The regular representation has \(I\)-invariant dimension
\(120/|I|\).  Subtracting both invariant dimensions from
the respective total dimensions yields
\[
 \operatorname{codim}(\operatorname{End}U)^I
      \le\frac{d^2}{120}\operatorname{codim}
                                      (\operatorname{Reg})^I.
 \tag{MA-RB.157}
\]
The finite Artin-conductor exponent at \(p\), in lower
ramification numbering, is the positive weighted sum
\[
 a_p(T)=\sum_{r\ge0}
          \frac{|I_{p,r}|}{|I_{p,0}|}
                        \operatorname{codim}T^{I_{p,r}},
       \qquad I_{p,0}=I_p.
 \tag{MA-RB.158}
\]
Only finitely many summands are nonzero.  Applying
(MA-RB.157) to each actual ramification group proves
\(a_p(\operatorname{End}U)\le(d^2/120)a_p(\operatorname{Reg})\).
The regular conductor--discriminant identity is
\(N(\operatorname{Reg})=\Delta\).  Multiplication over
primes, or summation after logarithms, proves the first
line of (MA-RB.151).  This uses no classification of
inertia subgroups.

\emph{The actual ramified Frobenius trace.}
The averaging projector
\(|I_p|^{-1}\sum_{i\in I_p}T(i)\) commutes with
\(T(\phi_p^j)\), since \(\phi_p\) normalizes inertia.
Taking the trace of their product gives
\[
 t_{\operatorname{End}U,p}(j)
   =\frac1{|I_p|}\sum_{i\in I_p}
                        |\chi_U(i\phi_p^j)|^2 .
 \tag{MA-RB.159}
\]
If \(\phi_p^j\notin I_p\), the coset contains no
identity element, so the analogous regular trace is
zero and the displayed sum is nonnegative.
If \(\phi_p^j\in I_p\), the coset is \(I_p\) and
contains the identity.  The sum is then at least
\(d^2/|I_p|\), while the regular trace is \(120/|I_p|\).
This proves the second line of (MA-RB.151).
It also proves that every summand defining
\(E_{\operatorname{End}U}\) is nonnegative.
Absolute convergence for \(\sigma>1\) follows from
the fixed-dimensional finite-image bound, so the
termwise inequalities may be summed there.

For \(X\), the already checked character table gives
\(\chi_X(2B)=-2\).  Consequently
\(\chi_{\operatorname{End}X}(2B)=4\), and solving
\(d_++d_-=36\), \(d_+-d_-=4\) gives \(d_+=20,d_-=16\).
At five inertia is trivial and Frobenius has order two.
Its odd powers have endomorphism trace \(4\) and regular
trace \(0\); its even powers have the respective traces
\(36\) and \(120\).  Summing the odd powers proves
(MA-RB.152).

\emph{The exact character bound for the complete weight cone.}
The known class order is
\((1,2A,2B,3A,6A,4A,5A)\).  The genuine character
\(\Theta=V\eta+W\) has values
\[
                   (9,-1,1,0,2,-1,-1).
 \tag{MA-RB.160}
\]
These are read from the previously checked table, not
from a new representation certificate.  In particular
\(\Theta(g)\le2\) for every nonidentity element.
For real weights in (MA-RB.153), use the class function
\(\chi_U=\sum_i a_i\chi_{U_i}\).
Orthogonality of the five distinct irreducible
characters gives
\(\langle|\chi_U|^2,\mathbf1\rangle=p\), and
\(\chi_U(1)=d\).
The exact discrepancy matrix in (MA-RB.62)--(MA-RB.64)
gives
\[
 \begin{split}
 m&=\langle|\chi_U|^2,\Theta\rangle
                                -2(a_1-a_2)^2\\
  &\le\frac{9d^2}{120}
                    +2\left(p-\frac{d^2}{120}\right)
   =2p+\frac7{120}d^2 .
 \end{split}
 \tag{MA-RB.161}
\]
The inequality is valid because \(|\chi_U(g)|^2\ge0\)
on every class; no positivity of the character
\(\Theta\) itself is assumed.

Put \(x=p/d^2\).  Cauchy--Schwarz gives
\[
 d^2\le
  (5^2+5^2+4^2+4^2+6^2)\sum_i a_i^2=118p.
\]
Nonnegativity of the weights gives
\(d\ge4\sum_i a_i\) and
\((\sum_i a_i)^2\ge p\).  Thus
\[
                         \frac1{118}\le x\le\frac1{16}.
 \tag{MA-RB.162}
\]
If \(m\le p\), the score is zero and (MA-RB.155)
follows.  Otherwise (MA-RB.161) gives
\[
 {\cal S}(a)
   \le\bigl(\sqrt{2x+7/120}-\sqrt x\bigr)^2
                     <\frac1{30}.
 \tag{MA-RB.163}
\]
To verify the strict inequality without numerical
optimization, its positive square-root comparison is
\(\sqrt{2x+7/120}<\sqrt x+1/\sqrt{30}\).
After squaring and moving terms it is
\(x+1/40<2\sqrt{x/30}\).  Both sides are positive.
A further square makes it exactly
\[
                 (x-\tfrac1{120})(x-\tfrac3{40})<0.
\]
The whole interval in (MA-RB.162) lies strictly between
these roots.  This proves (MA-RB.155).
If only \(V,V\eta,X\) occur, Cauchy--Schwarz gives the
stronger lower bound \(x\ge1/68\), so the same strict
inequality applies.
\end{proof}

\paragraph{Why the positive part and the pole term matter.}
The elementary near-one model underlying the proposed
score is
\[
       \frac{m}{t+\delta}\le\frac p t+h,
          \qquad t>0,\quad \delta>0,\quad h>0,\quad p>0.
 \tag{MA-RB.164}
\]
Here \(p\) is the actual pole order, rather than a
dimension-scaled fraction of it.  Rearranging gives
\(\delta\ge mt/(p+ht)-t\).
The derivative of the right side is
\(mp/(p+ht)^2-1\).  If \(m\le p\), its supremum
over \(t>0\) is zero, approached as \(t\downarrow0\).
If \(m>p\), its maximum occurs at
\(t=(\sqrt{mp}-p)/h\) and equals
\[
                  \frac{(\sqrt m-\sqrt p)^2}{h}.
 \tag{MA-RB.165}
\]
This proves the positive-part convention in (MA-RB.154).
For integer isobaric multiplicities among the known
five distinct cuspidals, the self-Rankin pole order
is exactly \(\sum_i a_i^2=p\).  The established order
matrix guarantees multiplicity at least \(km\) at the
actual residual zero.  The certificate above concerns
the unconditional lower bound \(k\ge1\), not an
assertion that all actual orders equal this bound.

Under a degree-proportional conductor bound
\(\log N\le(d^2/120)\log\Delta\), the leading constant
in this model is \(240{\cal S}(a)/\log\Delta\).
The regular representation has the corresponding
multiplicity \(9\), pole order \(1\), and degree \(120\),
giving \(8/\log\Delta\).
For \(\operatorname{End}X\) alone they give
\[
 \frac{(\sqrt3-1)^2}{0.15\log\Delta}
                         <\frac8{\log\Delta}.
 \tag{MA-RB.166}
\]
This comparison is a statement about the displayed
model and its leading conductor score.  It is not
a substitution of that model for an exact zero sum.

\paragraph{The exact archimedean and Euler comparison.}
The additional terms can be kept visibly separate.
Let
\[
 \begin{gathered}
 P(\sigma)=\frac1\sigma+\frac1{\sigma-1},
 \quad
 g_0(\sigma)=\frac{\Gamma_{\mathbb R}'}{\Gamma_{\mathbb R}}
                                                   (\sigma),
 \quad
 g_1(\sigma)=\frac{\Gamma_{\mathbb R}'}{\Gamma_{\mathbb R}}
                                                   (\sigma+1),\\
 H_X^{\rm ex}(\sigma)
   =P(\sigma)+\tfrac12\log N(\operatorname{End}X)
                         +20g_0(\sigma)+16g_1(\sigma)
                         -E_{\operatorname{End}X}(\sigma),\\
 H_{\rm reg}^{\rm ex}(\sigma)
   =P(\sigma)+\tfrac12\log\Delta
                         +60g_0(\sigma)+60g_1(\sigma)
                         -E_{\operatorname{Reg}}(\sigma).
 \end{gathered}
 \tag{MA-RB.167}
\]
Here the \(X\) expression uses the actual Artin
function \(L(s,\operatorname{End}X)\), with all its
inertia-invariant factors.  Its indicated normalized
completion is
\(s(s-1)N(\operatorname{End}X)^{s/2}
\Gamma_{\mathbb R}(s)^{20}\Gamma_{\mathbb R}(s+1)^{16}
L(s,\operatorname{End}X)\).
For the regular representation replace the conductor
by \(\Delta\), both exponents by \(60\), and the finite
function by \(\zeta_L\).  Logarithmic differentiation
proves that (MA-RB.167) gives their logarithmic
derivatives on \(\sigma>1\).
The already checked partial Rankin identity agrees
with this Artin function away from the finite omitted
set.  No identification of ramified tensor-product
factors is inferred here from equality of scalar
standard \(L\)-factors alone.  Such an identification
is unnecessary for the finite lemma or the score.
Their exact difference is
\[
 \begin{split}
 H_X^{\rm ex}-\tfrac3{10}H_{\rm reg}^{\rm ex}
   ={}&\tfrac7{10}P
       +\tfrac12\bigl(\log N(\operatorname{End}X)
                                      -\tfrac3{10}\log\Delta\bigr)\\
     &+\psi(\sigma/2)-\psi((\sigma+1)/2)\\
     &-\bigl(E_{\operatorname{End}X}
                                  -\tfrac3{10}E_{\operatorname{Reg}}\bigr).
 \end{split}
 \tag{MA-RB.168}
\]
Indeed
\(\Gamma_{\mathbb R}(s)=\pi^{-s/2}\Gamma(s/2)\),
so \(2(g_0-g_1)=\psi(\sigma/2)-\psi((\sigma+1)/2)\).
The latter difference is strictly negative:
differentiating the gamma integral gives
\(\psi'(x)=\operatorname{Var}(\log T)>0\) for a
gamma-distributed random variable of shape \(x>0\).
The conductor and Euler differences in (MA-RB.168)
have the favorable signs established above.
The positive term \(7P/10\) is the unscaled pole cost.
No sign for the whole right side, or universal ordering
of the corresponding exact variational minima, follows
merely from these separate signs.

\paragraph{Source scope, motive meaning and remaining gate.}
Calegari's Theorem~1.1, in the supplied primary
\texttt{S5ArtinWeb.tex}, lines123--136, supplies the
modularity of \(V,V\eta,X\).  The already checked
Rankin theorem and MA-RB-006 supply the actual
isobaric pole and residual-order assertions.
The ramification-filtration expression (MA-RB.158)
is the classical definition of the Artin conductor;
the regular conductor--discriminant identity is the
inherited input recorded with Neukirch, VII
(11.7)--(11.9), printed pp.~533--534/PDF pp.~541--542,
in GRH\_NIGHT NT-C06 and its source ledger.
That locator is not claimed as a new reading of
the original monograph pages in this section.

In the finite-cover language already fixed by SGA~1,
Expos\'e~V, \S7, and Deligne's Artin-motive locators
recorded in MA-S5, \(\operatorname{End}U=U\otimes U^\vee\)
is an actual finite Artin object.  Its local factor
uses the actual inertia invariants and residue
Frobenius in (MA-RB.159).  The inequality against
a rational multiple of the regular character is
a scalar inequality; it is not a submotive assertion
or positivity of a space of analytic zeros.

The newer exceptional-zero statements checked for
this route are Harcos--Thorner (2025),
\emph{A new zero-free region for Rankin--Selberg
\(L\)-functions}, arXiv:2303.16889v6, Theorem~1.1,
printed p.~3, and Harcos--Thorner (2026),
\emph{Tatuzawa's theorem for
Rankin--Selberg \(L\)-functions},
arXiv:2508.10844v3, Theorem~1.1/Corollary~1.4,
p.~3, and Theorem~1.6, pp.~4--5.
Their exact hypotheses and conductor-dependent
regions were inspected.  No proof of coverage of
the residual interval, nor a fresh full audit of
those source proofs, is claimed.
In particular, a possible quadratic twist satisfies
the latter theorem's condition
\(\pi\otimes\chi^2=\pi\); odd degree does not
contradict that condition.

Thus the local lemma and the cone score are checked
scope results.  They do not exclude the actual
shared zero in (MA-RB.90), prove \((\mathcal R_B)\),
or construct a positive-excess field.

\section{A relative-zeta envelope for the remaining real zero}
\label{sec:s5-relative-variational}

The positive completed divisor of a relative Dedekind zeta
function removes the base field's entire zero contribution.
This gives one finite-family refinement of MA-RB-008.
It does not assert that the remaining interval is empty,
and does not assume noncoincidence of zeros of different
Artin factors.

\paragraph{Review status.}
The complete post-writing author reading and two independent
whole-proof readings have passed. The primary-source checks
have their own exact recorded scope.

Keep the actual Calegari extension
\(L/\mathbb Q\), \(G=\operatorname{Gal}(L/\mathbb Q)=S_5\),
and the irreducible representations
\(\mathbf1,\eta,V,V\eta,W,W\eta,X\) of MA-RB-002.
For every nontrivial subgroup \(H\le G\), set
\[
 \begin{gathered}
 K_H=L^H,\qquad n_H=[K_H:\mathbb Q]=[G:H],\qquad
 d_H=120-n_H,\\
 h_U(H)=\dim U^H,\qquad
 m_H=5-h_W(H)
        +\min\{4-h_{V\eta}(H),\,6-h_X(H)\}.
 \end{gathered}
 \tag{MA-RB.170}
\]
Conjugate subgroups give the same data, so one may take
one representative of each conjugacy class. No subgroup
enumeration is needed below. Write \((r_{1,H},r_{2,H})\)
for the signature of \(K_H\), and put
\[
 \begin{gathered}
 a_H=60-r_{1,H}-r_{2,H},\qquad b_H=60-r_{2,H},\\
 D_H(s)=\frac{\zeta_L(s)}{\zeta_{K_H}(s)},\qquad
 \Xi_H(s)=
 \left(\frac{\Delta_L}{\Delta_{K_H}}\right)^{s/2}
       \Gamma_{\mathbb R}(s)^{a_H}
       \Gamma_{\mathbb R}(s+1)^{b_H}D_H(s).
 \end{gathered}
 \tag{MA-RB.171}
\]
Here \(\Gamma_{\mathbb R}(s)=\pi^{-s/2}\Gamma(s/2)\),
and \(\psi_{\mathbb R}=\Gamma_{\mathbb R}'/
\Gamma_{\mathbb R}\). Since \(H\ne\{1\}\),
\(n_H\le60\); in particular \(a_H,b_H\ge0\)
and \(a_H+b_H=d_H>0\).

\paragraph{MA-RB-012 (a strictly smaller relative-zeta gate).}
Every \(D_H\) is entire. Its completion \(\Xi_H\)
is entire of order at most one, has nonzero values at
zero and one, and satisfies
\(\Xi_H(s)=\Xi_H(1-s)\) and real conjugation symmetry.
Its zero divisor is nonempty and is contained, with
nonnegative integer multiplicities, in the nontrivial
completed divisor of \(\zeta_L\).

For real \(\sigma>1\), define the exact expression
\[
 \begin{split}
 \mathscr H_H(\sigma)
   ={}&\frac12\log\frac{\Delta_L}{\Delta_{K_H}}
       +a_H\psi_{\mathbb R}(\sigma)
       +b_H\psi_{\mathbb R}(\sigma+1)\\
      &+\frac{\zeta_L'}{\zeta_L}(\sigma)
       -\frac{\zeta_{K_H}'}{\zeta_{K_H}}(\sigma).
 \end{split}
 \tag{MA-RB.172}
\]
There are no pole terms in this formula. It extends
continuously to \(\sigma=1\) and is strictly positive
on \([1,\infty)\). For \(k\ge1\), let
\[
 \begin{split}
 F_{H,k}(\sigma)
    &=(\sigma-\tfrac12)^2
        -\frac{m_Hk(2\sigma-1)}{\mathscr H_H(\sigma)},\\
 J_{H,k}&=\inf_{\sigma>1}F_{H,k}(\sigma),\qquad
 \mathfrak J_{L,k}
      =\min_{\{1\}\ne H\le G}J_{H,k}.
 \end{split}
 \tag{MA-RB.173}
\]
Each infimum is finite and is a minimum after adjoining
the endpoint \(\sigma=1\). It need not be attained at
an interior point. One has, unconditionally,
\[
          \mathfrak J_{L,k}
                 <\mathcal J_{\Delta_L,k}<\frac14,
 \tag{MA-RB.174}
\]
where \(\mathcal J_{\Delta_L,k}\) is the previously
checked discriminant-aware bound of MA-RB-008.

If the actual comparison-failure branch of MA-RB-002
occurs at its real abscissa \(\gamma\), with pole order
\(k\ge1\), then
\[
 \begin{gathered}
 \operatorname{ord}_{\gamma}\Xi_H
   =\operatorname{ord}_{1-\gamma}\Xi_H\ge m_Hk
                                      \quad(H\ne\{1\}),\\
             (\gamma-\tfrac12)^2\le\mathfrak J_{L,k}
                                      \le\mathfrak J_{L,1}.
 \end{gathered}
 \tag{MA-RB.175}
\]
Consequently \(\mathfrak J_{L,1}\le0\) excludes that
branch. Otherwise its sufficient shared-zero gate
is confined to
\[
 \begin{gathered}
 B(\gamma)=0,\qquad
        L(\gamma,V\eta)L(\gamma,X)=0,\\
 \frac12<\gamma
       \le\frac12+\sqrt{\mathfrak J_{L,1}}<1.
 \end{gathered}
 \tag{MA-RB.176}
\]
The new interval is strictly smaller than the nonempty
interval in MA-RB.90. These assertions use neither
\(\mathsf A\), \(M\), nor \((\mathcal R_B)\).

\begin{proof}
\emph{The classical divisibility input in its exact established form.}
Aramata's main theorem states that, for a finite normal
extension of number fields \(L/K\), the finite quotient
\(\zeta_L/\zeta_K\) is entire. The exact primary locator
is Aramata (1933), p.~31, the main statement preceding
Satz~1; its proof is completed on p.~34 following
Satz~6. All four pages of that proof were freshly
read in the source audit. This is a cited classical
input, not a new holomorphy theorem.
Here \(L/K_H\) is normal with Galois group \(H\),
even when \(K_H/\mathbb Q\) is not normal. Therefore
\[
             D_H=\zeta_L/\zeta_{K_H}\quad\hbox{is entire}.
 \tag{MA-RB.177}
\]
The statement concerns the full Dedekind functions,
including every ramified factor. It does not require
Artin holomorphy for individual irreducible factors.
The Artin induction and nontrivial Hecke holomorphy
interface is also explicitly recorded in the supplied
Hu--Kaneko--Martin--Schildkraut source,
\texttt{lem:artin-L}(a)--(d) and
\texttt{cor:hol-quotient}, including the latter's
complete proof. Tags MA-RB.178--.180 are reserved
and unused after this citation-only consolidation.

\emph{Completion and the genuine excess divisor.}
For a number field \(K\), use the same normalization
as MA-RB-001:
\[
 \xi_K(s)=s(s-1)\Delta_K^{s/2}
      \Gamma_{\mathbb R}(s)^{r_1(K)}
      \Gamma_{\mathbb C}(s)^{r_2(K)}\zeta_K(s),
 \qquad
 \Gamma_{\mathbb C}(s)
       =\Gamma_{\mathbb R}(s)\Gamma_{\mathbb R}(s+1).
 \tag{MA-RB.181}
\]
The standard Dedekind continuation, functional
equation and fixed-strip bounds used in MA-RB.8--.10
give an entire function of order at most one,
nonzero at zero and one, with real conjugation
symmetry and \(\xi_K(s)=\xi_K(1-s)\).
Its zeros lie in \(0\le\operatorname{Re}s\le1\);
their count in \(|s|\le T\) is \(O_K(T\log(T+3))\).
The proof of MA-RB.8--.10 uses these properties
for a general Dedekind completion; the particular
signature there enters only its explicit gamma
derivative.

Complex conjugation in \(G\) has type \(2B\), so
the signature of \(L\) is \((0,60)\).
Taking the quotient of (MA-RB.181) for \(L\) and
\(K_H\) cancels both factors \(s(s-1)\) and gives
exactly \(\Xi_H\) in (MA-RB.171).
At a nontrivial completed zero of \(\xi_{K_H}\)
the archimedean gamma factors are holomorphic and
nonzero: the only possible gamma endpoint in
the containing strip is zero, and \(\xi_{K_H}(0)\ne0\).
Entireness of \(D_H\) consequently says that
the zero occurs in \(\xi_L\) with at least the
same order. There are no other zeros of
\(\xi_{K_H}\). The quotient
\[
                   \Xi_H=\xi_L/\xi_{K_H}
 \tag{MA-RB.182}
\]
is therefore entire, with a nonnegative excess
divisor \(\mathcal Z_H\), and is nonzero at both
endpoints. This also verifies cancellation of
every possible trivial-zero or gamma singularity
in the formula (MA-RB.171); no hidden endpoint
pole or zero is retained.

For completeness its order can be checked without
an estimate for a reciprocal near zeros.
Write the genus-one Hadamard products for the
two functions \(\xi_L,\xi_{K_H}\), with
\(E_1(z)=(1-z)e^z\).
The factors belonging to \(\xi_{K_H}\) cancel
exactly, because the products converge locally
normally and \(\sum_\rho|\rho|^{-2}<\infty\).
Thus
\[
 \Xi_H(s)=\Xi_H(0)e^{\alpha_Hs}
                     \prod_{\rho\in\mathcal Z_H}E_1(s/\rho),
                       \qquad \alpha_H\in\mathbb R.
 \tag{MA-RB.183}
\]
The excess count is bounded by that of \(\xi_L\).
For \(|s|\le R\), the factors with \(|\rho|\le2R\)
have total log-modulus at most \(O_L(R\log^2(R+3))\):
use
\(\log|E_1(s/\rho)|
 \le\log(1+R/|\rho|)+R/|\rho|\)
and partial summation of the zero count.
For \(|\rho|>2R\), use
\(\log|E_1(s/\rho)|\le C R^2/|\rho|^2\);
their total is \(O_L(R\log(R+3))\).
The finite minimum nonzero zero modulus is harmless.
This proves order at most one.
Reflection and conjugation follow directly from
(MA-RB.182).

\emph{The exact positive logarithmic derivative.}
In (MA-RB.183),
\(\sum_{\rho\in\mathcal Z_H}\operatorname{Re}(1/\rho)\)
converges absolutely by the containing strip.
Call this sum \(S_H\). Logarithmic differentiation
at zero and one and the reflection
\(\rho\mapsto1-\rho\) give
\(-\alpha_H=\alpha_H+2S_H\), hence
\(\alpha_H=-S_H\).
For every real \(\sigma>1\) this fixes the constant
in the Hadamard formula and proves
\[
 \frac{\Xi_H'}{\Xi_H}(\sigma)
      =\sum_{\rho\in\mathcal Z_H}
                    \operatorname{Re}\frac1{\sigma-\rho}
      =\mathscr H_H(\sigma).
 \tag{MA-RB.184}
\]
The displayed real-part series is absolutely
convergent. The second equality is the direct
logarithmic derivative of (MA-RB.171), so in
particular no \(1/\sigma\) or \(1/(\sigma-1)\)
term is present.

The divisor \(\mathcal Z_H\) cannot be empty.
The gamma estimate already proved from the
gamma integral in MA-RB-007 gives
\(\psi_{\mathbb R}(\sigma+e)
 =\tfrac12\log\sigma-\tfrac12\log(2\pi)
     +O(1/\sigma)\) for \(e=0,1\).
Both finite Dedekind logarithmic derivatives
tend to zero on the positive real axis, by
dominated convergence in their absolutely
convergent Euler series at \(\sigma\ge2\).
Consequently
\[
          \mathscr H_H(\sigma)
               =\frac{d_H}{2}\log\sigma+O_L(1)
                                      \quad(\sigma\longrightarrow\infty).
 \tag{MA-RB.185}
\]
An empty divisor would make the first expression
in (MA-RB.184) identically zero, contradicting
\(d_H>0\). Every term in (MA-RB.184) is positive
when \(\sigma>1\), which proves strict positivity.

The function \(\Xi_H\) is nonzero at one, so its
logarithmic derivative extends continuously there.
The real-part zero series also extends to one:
its tail is locally bounded by a constant times
\(\sum_{|\rho|>2}|\rho|^{-2}\), and the remaining
finite terms have no singularity at one.
At one all terms are nonnegative. The reflection
symmetry ensures that a nonempty divisor contains
at least one zero with real part strictly less
than one, even if some zeros lie on a boundary
line. That zero contributes strictly positively
at one. Therefore \(\mathscr H_H(1)>0\).

\emph{The residual multiplicity without an extra zero-location premise.}
The permutation representation of the embeddings of
\(K_H\) is \(\operatorname{Ind}_H^G\mathbf1\).
Frobenius reciprocity gives multiplicity \(h_U(H)\)
for each irreducible \(U\). Its all-prime
factorization and the regular factorization imply
\[
      D_H(s)=
           \prod_U L(s,U)^{\,\dim U-h_U(H)}.
 \tag{MA-RB.186}
\]
All exponents are nonnegative integers.
At the particular residual point \(\gamma\),
MA-RB.19--.26 proves the following pointwise
orders without assuming general holomorphy of
the unknown five-dimensional factors:
\[
 \begin{gathered}
 x_{\mathbf1}=0,\quad x_\eta=b\ge0,\quad
 x_V=v_0\ge0,\quad x_{V\eta}=v_1\ge0,\quad
 x_X=e\ge0,\\
 x_W=n+k\ge k,\qquad x_{W\eta}\ge0,\qquad
                         v_1+e\ge k.
 \end{gathered}
 \tag{MA-RB.187}
\]
In particular \(x_{W\eta}\ge0\) follows there
from \(L_{W\eta}=\zeta_F/\zeta_{\mathbb Q}\)
at the real point \(\gamma\in(1/2,1)\), where
\(\zeta_{\mathbb Q}\ne0\).
Taking orders in (MA-RB.186), discarding only
nonnegative summands, and using \(v_1+e\ge k\)
gives
\[
 \begin{split}
 \operatorname{ord}_\gamma D_H
   &\ge (5-h_W)(n+k)
         +(4-h_{V\eta})v_1+(6-h_X)e\\
   &\ge m_H k .
 \end{split}
 \tag{MA-RB.188}
\]
The gamma factors are units at \(\gamma\), so
this is also the order of \(\Xi_H\).
Its functional equation gives the same order
at \(1-\gamma\). The two real points are distinct.
Moreover \(m_H\ge1\): the character row of \(W\)
has value strictly less than five at every
nonidentity element, so \(W\) is faithful.
For nontrivial \(H\), \(h_W(H)<5\); the other
two terms defining the minimum are nonnegative.

Retain these two real zeros in the positive sum
(MA-RB.184). With \(u=\gamma-1/2\) and
\(x=\sigma-1/2>1/2\), one obtains
\[
      \frac{2m_H kx}{x^2-u^2}
                      \le\mathscr H_H(\sigma),
       \qquad u^2\le F_{H,k}(\sigma).
 \tag{MA-RB.189}
\]
All denominators are positive because
\(0<u<1/2<x\).
Taking the infimum and then the finite minimum
proves (MA-RB.175), apart from its final
monotonicity assertion.

\emph{The minima and the strict comparison with the old gate.}
Strict positivity on \([1,\infty)\) makes every
\(F_{H,k}\) continuous there. Formula (MA-RB.185)
gives \(F_{H,k}(\sigma)\to+\infty\) as
\(\sigma\to\infty\), while
\[
             F_{H,k}(1)
                    =\frac14-\frac{m_Hk}{\mathscr H_H(1)}
 \tag{MA-RB.190}
\]
is finite. Hence the infimum over the open
half-line equals a minimum on \([1,\infty)\).
This proves the asserted finiteness and the
endpoint qualification.

Take \(H=A_5\), so \(K_H=E\) is the real quadratic
subfield. Its reality follows because complex
conjugation \(2B\) is even. The only irreducible
representations with \(A_5\)-invariants are
\(\mathbf1,\eta\): an invariant subspace for a
normal subgroup is \(G\)-stable, and an irreducible
representation trivial on \(A_5\) factors through
\(G/A_5=C_2\). Thus \(m_{A_5}=9\).

Write \(\mathscr H_K^{\rm full}=\xi_K'/\xi_K\)
on \(\sigma>1\). The same exact Hadamard
argument as (MA-RB.184) shows
\(\mathscr H_E^{\rm full}(\sigma)>0\).
Its divisor is nonempty: otherwise its zero
sum would vanish identically, whereas the
degree-two gamma derivative gives
\(\mathscr H_E^{\rm full}(\sigma)
 =\log\sigma+O_E(1)\) at infinity.
Now (MA-RB.182), followed by the exact positive
Euler-remainder identity MA-RB.95, gives
\[
 0<\mathscr H_{A_5}(\sigma)
    =\mathscr H_L^{\rm full}(\sigma)
                    -\mathscr H_E^{\rm full}(\sigma)
    <\mathscr H_L^{\rm full}(\sigma)
    \le\mathcal H_{\Delta_L}(\sigma).
 \tag{MA-RB.191}
\]
Let \(\sigma_*>1\) be an attained minimizer for
\(\mathcal J_{\Delta_L,k}\), whose existence
was proved in MA-RB-008. Since the numerator
\(9k(2\sigma_*-1)\) is strictly positive,
(MA-RB.191) implies
\[
 \mathfrak J_{L,k}\le J_{A_5,k}
    \le F_{A_5,k}(\sigma_*)
    <\mathcal F_{\Delta_L,k}(\sigma_*)
     =\mathcal J_{\Delta_L,k}.
 \tag{MA-RB.192}
\]
This proves (MA-RB.174). The strict step is
available at an old interior minimizer, so it
does not require an interior minimizer for
the new function.

There is no need for \(\mathsf A\) even in this
quadratic case. Indeed \(\zeta_E=\zeta_{\mathbb Q}C\)
has order \(b\) at \(\gamma\), and MA-RB.26
already gives
\[
       \operatorname{ord}_\gamma D_{A_5}
          =\operatorname{ord}_\gamma\zeta_L-b
                              \ge5n+9k.
 \tag{MA-RB.193}
\]
The possible quadratic zero is subtracted
exactly, including its multiplicity.
Finally each \(F_{H,k}\) decreases with \(k\),
as \(m_H>0\); the infima and finite minimum
do also. Together with MA-RB-002's two vanishing
factors, this proves (MA-RB.175)--(MA-RB.176)
and the nonpositive-endpoint exclusion.
\end{proof}

\paragraph{Three directly verified subgroup choices.}
The following examples are part of the same finite
envelope, not successive new reductions:
\[
\begin{array}{c|r|rr|rrr|r|rr}
 H&n_H&r_{1,H}&r_{2,H}&h_W&h_{V\eta}&h_X&m_H&a_H&b_H\\ \hline
 A_5&2&2&0&0&0&0&9&58&60\\
 \langle(12)\rangle&60&0&30&3&1&3&5&30&30\\
 \langle(12345)\rangle&24&0&12&1&0&2&8&48&48
\end{array}
 \tag{MA-RB.194}
\]
For the transposition, the invariant dimensions
are the averages of the identity and \(2A\)
character values. The relevant \(2A\) values
are \(1,-2,0\). For a five-cycle subgroup,
its four nonidentity elements all have type \(5A\);
the relevant values are \(0,-1,1\).
These give exactly the displayed averages.
In the last two cases the subgroup contains
no element of type \(2B\). Complex conjugation
therefore fixes no coset in \(G/H\), because
a fixed coset would put a conjugate of \(2B\)
inside \(H\). Thus these fields are totally
imaginary, giving the displayed signatures.
This uses the existing checked character rows;
it is not a new finite-representation certificate.

\paragraph{A finite Euler evaluation with an explicit error.}
For clarity this gate uses genuine field data,
not an assigned value for an unknown zero sum.
For a number field \(K\), define its rational
prime-power coefficients by
\[
 -\frac{\zeta_K'}{\zeta_K}(s)
       =\sum_{n\ge2}\frac{\Lambda_K(n)}{n^s},
 \qquad
 \Lambda_K(p^j)
       =(\log p)\sum_{\substack{\mathfrak p\mid p\\
                               f_{\mathfrak p}\mid j}}
                                     f_{\mathfrak p}.
 \tag{MA-RB.195}
\]
The factorization of every rational prime gives
these coefficients, including ramified primes.
They are zero at integers that are not prime powers.
The displayed sum is at most \([K:\mathbb Q]\log p\),
since \(\sum_{\mathfrak p\mid p}e_{\mathfrak p}
 f_{\mathfrak p}=[K:\mathbb Q]\).
Put \(c_H(n)=\Lambda_L(n)-\Lambda_{K_H}(n)\).
One always has the sufficient elementary bound
\[
              |c_H(n)|\le(120+n_H)\log n .
 \tag{MA-RB.196}
\]
In particular its constant for \(H=A_5\) is
\(122\). No sign of \(c_H(n)\) is assumed.
For an integer \(N\ge3\) and \(\sigma>1\), set
\[
 \begin{split}
 \mathscr H_{H,N}(\sigma)
    ={}&\frac12\log\frac{\Delta_L}{\Delta_{K_H}}
       +a_H\psi_{\mathbb R}(\sigma)
       +b_H\psi_{\mathbb R}(\sigma+1)
       -\sum_{2\le n\le N}c_H(n)n^{-\sigma},\\
 T_{H,N}(\sigma)
    ={}&(120+n_H)N^{1-\sigma}
       \left(\frac{\log N}{\sigma-1}
                         +\frac1{(\sigma-1)^2}\right).
 \end{split}
 \tag{MA-RB.197}
\]
The function \((\log x)x^{-\sigma}\) is decreasing
for \(x\ge3\). Comparison of the absolute tail
with its integral on \([N,\infty)\) proves
\[
 \left|\mathscr H_H(\sigma)-\mathscr H_{H,N}(\sigma)\right|
             \le T_{H,N}(\sigma),\qquad
 U_{H,N}(\sigma):=\mathscr H_{H,N}(\sigma)+T_{H,N}(\sigma)
                    \ge\mathscr H_H(\sigma)>0.
 \tag{MA-RB.198}
\]
Therefore a comparison failure would satisfy
the fully finite upper bound
\[
 (\gamma-\tfrac12)^2
      \le(\sigma-\tfrac12)^2
               -\frac{m_Hk(2\sigma-1)}{U_{H,N}(\sigma)}.
 \tag{MA-RB.199}
\]
For every fixed \(\sigma>1\), its right side
converges to \(F_{H,k}(\sigma)\) as \(N\to\infty\).
Taking rational \(\sigma>1\), finite \(N\), and
the finite subgroup family consequently recovers
\(\mathfrak J_{L,k}\) as the infimum of these
certified upper bounds. To see equality, the
upper bounds are all at least the true infimum;
conversely choose a real \(\sigma>1\) arbitrarily
close to an infimum, approximate it by a rational
one using continuity, and then take \(N\) large.
This argument also covers a minimum at one by
approaching from the right. It does not use an
Euler series at \(\sigma=1\), or assert a uniform
tail bound there.

\paragraph{Scope of positive products of these relative factors.}
The finite envelope already contains the bound from
any positive integral product of the same \(D_H\)
when only the guaranteed multiplicities \(m_Hk\)
are retained. Indeed, for nonnegative coefficients
\(u_H\), not all zero, and each \(\sigma>1\),
\[
 \frac{\sum_Hu_Hm_H}{\sum_Hu_H\mathscr H_H(\sigma)}
 =\sum_H
   \frac{u_H\mathscr H_H(\sigma)}
        {\sum_Ju_J\mathscr H_J(\sigma)}
             \frac{m_H}{\mathscr H_H(\sigma)}
 \le\max_H\frac{m_H}{\mathscr H_H(\sigma)}.
\]
For integer \(u_H\), the numerator is the guaranteed
multiplicity divided by \(k\), and the denominator is
the logarithmic derivative of the completed product.
Its corresponding two-real-zero upper bound is
therefore at least \(\min_H F_{H,k}(\sigma)\).
Taking an infimum in \(\sigma\) cannot lower it
below \(\mathfrak J_{L,k}\), since this family is
finite. The pointwise algebra also holds for real
nonnegative coefficients, without defining fractional
holomorphic powers. Higher actual multiplicities or
additional arithmetic information are outside this
limited optimality assertion.

\paragraph{Finite-cover and source scope.}
There is a genuine \(G\)-equivariant inclusion
\(\mathbb Q[G/H]\hookrightarrow\mathbb Q[G]\),
given by \(gH\mapsto\sum_{h\in H}gh\).
Right averaging
\(|H|^{-1}\sum_{h\in H}R_h\) is an idempotent
onto this image and commutes with the left \(G\)-action.
Its kernel is therefore a rational complement, an actual
finite Artin representation with character
\(\operatorname{Reg}_G-\operatorname{Ind}_H^G\mathbf1\);
equivalently it is the induction from \(H\) of
the augmentation representation of \(H\).
This describes the degree-zero-cohomology
Artin-motive quotient for the finite cover
\(\operatorname{Spec}L\to\operatorname{Spec}K_H\),
viewed over \(\mathbb Q\). Inertia invariants are
exact in characteristic zero, so the local
quotient includes every ramified Euler factor.
The conductor is \(\Delta_L/\Delta_{K_H}\) by
the same conductor--discriminant convention as
in the exact completion. Its Frobenius traces
need not be nonnegative. The positive measure
in the proof is the completed excess zero
divisor supplied by holomorphic divisibility,
not coefficientwise Euler positivity.

The primary divisibility source is Aramata
(1933), \emph{Über die Teilbarkeit der Dedekindschen
Zetafunktionen}, \emph{Proceedings of the Imperial
Academy}, \(9(2)\), 31--34,
doi:10.3792/pia/1195580866:
the principal statement is on p.~31, its
finite-character construction is in Sätze~1--6,
and the final analytic proof is on p.~34.
All four primary pages were freshly read in
their rendered form. The paper invokes this
known theorem in its exact normal-extension form,
as (MA-RB.177), without reproducing its proof.
The supplied Hu--Kaneko--Martin--Schildkraut
primary LaTeX, \texttt{HigherOrderZeros.tex},
\texttt{lem:artin-L}(a)--(d),
\texttt{lem:aramata-brauer}, and
\texttt{cor:hol-quotient}, was also read for
its exact all-prime induction/holomorphy scope.
No general Artin holomorphy theorem is substituted
for the relative divisibility statement.

The Dedekind continuation/growth inputs and
fixed Hadamard constant are the established
NT-Z10 and MA-RB.8--.12; the actual residual
orders are MA-RB.19--.26. The strict prior
denominator comparison uses MA-RB.95 and the
attained old minimum in MA-RB.87--.88.
For the categorical interpretation the inherited
exact locators remain SGA~1, Exposé~V, \S7,
and Deligne (1979), \emph{Valeurs de fonctions
\(L\) et périodes d'intégrales},
Definition~6.1, \S\S6.2--6.3 and equation~(0.5.3).
The conductor--discriminant locator remains
Neukirch, VII(11.7)--(11.9), printed pp.~533--534.
These previously checked categorical and
conductor locators are not claimed to have
received a fresh complete-source audit here.

This is a strictly smaller necessary interval
for the actual failure branch and hence a
weaker sufficient shared-zero exclusion
condition. It is not a proof that every
\(\mathfrak J_{L,1}\) is nonpositive, a proof
of \((\mathcal R_B)\), or a construction of
an arithmetic positive-excess example.

\section{PF26: the algebraic input and its smaller analytic premise}
\label{sec:pf26-audit}

\subsection{Source and scope of the check}

The source checked here is the supplied 48-page PDF of OpenAI,
\emph{Deterministic Polynomial Factorization over Prime Fields}, dated
4 October 2026.  Printed and PDF page numbers agree. Its SHA-256 is
\texttt{ed88aa5ba802311a21c4f265bcef43515a16019f4b5afdd7f52ecbd36ce0204d}.
The complete proofs in Sections 2--10, pp.~6--44, were read, including all
the inputs to Theorem~1.2, Proposition~8.1 and Proposition~10.1. The
conclusions below are mathematical proof-reading conclusions, not a
formal verification or a claim that an implementation was run. PF26
remains an explicitly named outside preprint in the input register.

\paragraph{MA-PF-001 (the three interfaces are distinct).}
Put
\[
 B=20+(d+1)(\lceil\log_2p\rceil+1).
\]
The algebraic reduction of PF26, Theorem~1.2, p.~3, assumes an auxiliary
prime $\ell_r$ for each prime $r\le d$ only in the branch
$p>B^{200000}$. Its conditions, equation~(1.2), are
\[
 \ell_r\notin\{2,3,p,r\},\qquad
 \ell_r\equiv1\pmod{12r},\qquad
 p^{(\ell_r-1)/r}\not\equiv1\pmod{\ell_r}.
\]
The running time is polynomial in $B$ and the \emph{numerical values}
of these primes. Proposition~10.1, pp.~41--44, gives the explicit
sufficient bound
\[
 O(B^{500000}+B^{100}E^{20}),\qquad E=\max(2,\max_r\ell_r).
 \tag{MA-PF-1}
\]
The more flexible algebraic interface is Proposition~8.1, pp.~35--37:
for a completely split square-free polynomial of odd degree $N$, choose
an odd prime $r\mid N$; supply a field $K/\mathbb F_p$ of degree at
most $r-1$, a primitive $r$th root of unity in $K$, and a generator of
the full $r$-primary subgroup of $K^\times$. These data suffice for a
proper factor in a fixed polynomial number of bit operations. An
auxiliary prime is one way to construct these data; it is not part of
this latter interface.

The analytic premise used for the unconditional statement of PF26 is
\emph{not} the Riemann hypothesis. It is the theorem quoted as
Theorem~11.1, pp.~44--45, from the companion
\emph{Primitive roots for every admissible integer base}, Theorem~1.2:
for every cyclotomic field $K$ containing $\mu_{12}$ and every
finite-order Hecke character $\chi$ of $K$,
\[
 L_K(s,\chi)\ne0\quad\text{if }\Re s>1-10^{-6}.
 \tag{MA-PF-2}
\]
The principal character is included, its pole at $1$ is permitted, and
there is no conductor or imaginary-height restriction. PF26 does not
prove that companion theorem. Its own proof of the implication from
this premise is Lemma~11.2 and Proposition~11.3, pp.~45--46. The
algebraic reduction in Theorem~1.2 uses neither this premise nor zero
location. GRH is a sufficient substitute, as PF26 explicitly observes
on p.~47. In the present paper $A$ is the Dirichlet theorem under the
source convention fixed in the introduction, whereas $M$ is conditional.

\emph{Verification.} The stated distinctions follow directly from the
three source interfaces and the proofs: Sections~2--10 terminate after
constructing the primary table; Section~11 alone bounds the numerical
auxiliary-prime search. In particular, accepting the algebraic PF26
input does not accept its companion's analytic assertion. The proof
reading below found no gap in these algebraic interfaces.\hfill$\square$

\subsection{The completed check of CX-ALG-028}

\paragraph{MA-PF-002 (the finite residual family and an elementary test).}
For the large-characteristic branch $p>B^{200000}$, define
\[
 \mathcal B(p,m,d)=
 \{r\le d:r\text{ prime},\ r\nmid(p-1)m,
              \ p^{r-1}\equiv1\pmod{r^2}\}.
 \tag{MA-PF-3}
\]
Every member is odd and different from $p$. This is exactly the set in
GRH\_NIGHT CX-ALG-028. Under the algebraic PF26 input and $A$, adding
RH for $\zeta_{\mathbb Q(p^{1/r})}$ for $r\in\mathcal B(p,m,d)$
suffices to factor degree-$d$ polynomials over an explicitly supplied
$\mathbb F_{p^m}$ in deterministic polynomial time in
$(d+1)m\log p$. Thus the additional family is empty precisely when
\[
 \forall r\le d\text{ prime},\quad
 r\mid(p-1)m\quad\text{or}\quad p^{r-1}\not\equiv1\pmod{r^2}
 \tag{MA-PF-4}
\]
in that branch. In the small-characteristic branch no additional
analytic input is needed, irrespective of this arithmetic test.
Only primes actually required by the primary-table construction need
be included. This is a sufficient family, not a necessity theorem.

\emph{Proof and independent check of the transfer.}
The large-characteristic inequality gives $p>d$ and $p>3$. For an odd
prime $r\le d$ put $s=\operatorname{ord}_r(p)$ and
$a=v_r(p^s-1)$. Since $s\mid r-1$, writing $r-1=st$ gives
$r\nmid t$. The binomial identity, or the odd-prime lifting-the-exponent
identity, gives
\[
 v_r(p^{r-1}-1)=v_r((p^s)^t-1)=v_r(p^s-1)=a.
\]
Consequently the condition $a>1$ in CX-ALG-028 is precisely the
congruence in (MA-PF-3). The prime $2$ divides $p-1$, so it cannot
belong to the displayed set.

The quadratic primary entry can be obtained from a nonsquare $c$ of
$\mathbb F_p$, by raising $c$ to the odd part of $p-1$. If
$v_2(p-1)=1$, use $c=-1$. If $2\mid m$, compute the quadratic
subfield of the supplied field by solving $z^{p^2}=z$ over
$\mathbb F_p$, and choose $0\ne v$ with $v^p=-v$ there. Then
$c=v^2\in\mathbb F_p$ satisfies $c^{(p-1)/2}=-1$. Otherwise
$A$, with the classical least nonresidue consequence recorded in
GRH\_NIGHT CX-ALG-023, supplies $c$ by a polynomially bounded search.
This checks all quadratic cases, including the supplied-subfield case.

Proceed in increasing order through odd primes $r\le d$. Factoring
$\Phi_r$ has degree $r-1$, so uses only earlier table entries. An
irreducible factor gives
$K_r=\mathbb F_p(\zeta_r)$ of degree
$s=\operatorname{ord}_r(p)\mid r-1$. Put $Q=p^s$ and $R=r^a$.
If $a=1$, the element $\zeta_r$ is already the full primary
generator. If $r\mid p-1$, a small $r$th nonresidue furnished by
$A$ in $\mathbb F_p=K_r$, raised to $(p-1)/R$, is a full primary
generator. Testing whether a candidate is an $r$th power uses the
single modular exponent $(p-1)/r$; no factorization of $p-1$ is needed.

For every remaining case it suffices to construct an explicit
$U_r\simeq\mathbb F_{p^r}$. If $r\mid m$, solve
$z^{p^r}=z$ in the supplied field. Its kernel has dimension $r$ and
its inherited multiplication gives $U_r$. If $r\nmid m$, the
unhandled conditions are exactly $r\in\mathcal B(p,m,d)$.
GRH\_NIGHT CX-ALG-024 supplies $U_r$ in polynomial time from $A$
and RH for the indicated pure field. Its proof was reread in this
check: the factorization of zeta functions, the discriminant bound and
the Bach--Sorenson Theorem~5.1 input agree. One can use the invariant
algebra of PF26 Lemma~3.1, p.~11, to implement its final field
construction. Its auxiliary prime needs congruence modulo $r$ only;
the extra factor $12$ in PF26 equation~(1.2) is unnecessary at this
interface.

Here is the remaining algebraic transfer, with its orders checked.
Because $(r,s)=1$,
\[
 V_r=U_r\otimes_{\mathbb F_p}K_r\simeq\mathbb F_{p^{rs}}
\]
is a field of degree $r$ over $K_r$. Solve the $K_r$-linear equation
$u^Q=\zeta_r u$ for $u\ne0$. Such a solution exists since
$\zeta_r$ has norm one and the map $v\mapsto v^{Q-1}$ has precisely
the norm-one subgroup as its image. As
$v_r(Q^r-1)=a+1$ and $u^{Q-1}$ has order $r$, the $r$-part of
$\operatorname{ord}(u)$ is exactly $r^{a+1}=rR$. Hence
\[
 w=u^{(Q^r-1)/(rR)},\qquad \omega_r=w^r
\]
have exact orders $rR$ and $R$. The subgroup of $V_r^\times$ of
order $R$ is contained in $K_r^\times$, so $\omega_r$ belongs to
$K_r$; its coordinates are recovered by the known inclusion and a
linear solve. All dimensions are at most $r(r-1)$ and all exponent
lengths are polynomial in $r\log p$.

The induction is well founded. PF26 Proposition~2.3, pp.~9--10,
uses the quadratic entry to split a completely split polynomial of
even degree. Proposition~8.1, pp.~35--37, uses precisely the odd
entry just supplied. Its proof receives a field containing a full
primary generator; it does not inspect how that generator was
constructed. PF26 Lemma~9.1 and Proposition~9.2, pp.~38--40,
therefore factor every input below the next table prime with earlier
entries. The table is constructed once; it is not recursively rebuilt.
There are at most $d$ entries and polynomially many factoring calls.
The full reading of the implementation and size proof, reported below,
checks that the exponent is independent of $r$ and of the genus.

For completeness, over the supplied extension field take a square-free
$h$ and form
\[
 C=\mathbb F_{p^m}[X]/(h),\qquad
 C_0=\ker(z\mapsto z^p-z).
\]
If $h$ has $t$ irreducible factors, then $C_0\simeq\mathbb F_p^t$.
The kernel is obtained by $\mathbb F_p$-linear algebra on dimension
$m\deg h$. If $t>1$, a basis supplies a nonscalar $b$. Its minimal
polynomial over $\mathbb F_p$ is a square-free product of distinct
linear factors, of degree between $2$ and $t$. Factor it by the
prime-field procedure. For each root $a$, the gcd
$\gcd(h,b-a)$ selects the irreducible factors on which $b$ has value
$a$, producing a proper split. Degrees strictly decrease under
recursion. Derivative gcds handle multiplicities; if the derivative
vanishes, take coefficient $p$th roots by inverse Frobenius
$c\mapsto c^{p^{m-1}}$ and divide all positive exponents by $p$.
In the branch $p\le B^{200000}$, enumeration of the $p$ scalars in
the same fixed algebra is already polynomial in $B$. These arguments
complete the transfer and the empty-set criterion.\hfill$\square$

\paragraph{Unconditional primary-table promise.}
The stronger conclusion in the last sentence of CX-ALG-028 is also
valid: if every prime $r\le d$ either divides $m$ or has
$v_r(p^{\operatorname{ord}_r(p)}-1)=1$, the displayed construction
uses neither $A$ nor an additional number-field zero-location input.
For $r=2$ use the two elementary quadratic branches above; for odd
$r$ use $\zeta_r$ in the valuation-one branch and the supplied
subfield otherwise. The claim is conditional only on the algebraic
PF26 interface.

\subsection{One Kummer Hecke function for each residual field}

\paragraph{MA-PF-003 (the single-character identity).}
Let $r$ be an odd prime, $p\ne r$ a prime, and put
\[
 F_r=\mathbb Q(p^{1/r}),\quad K_r=\mathbb Q(\zeta_r),\quad
 E_r=K_r(p^{1/r}).
\]
Let $\psi_{p,r}$ be any faithful character of
$\operatorname{Gal}(E_r/K_r)\simeq C_r$, interpreted by class field
theory as its primitive finite-order Hecke character of $K_r$.
Choose the convention in which its unramified value is the Kummer
symbol
\[
 \psi_{p,r}(\mathfrak l)
   =\frac{\operatorname{Frob}_{\mathfrak l}(p^{1/r})}{p^{1/r}}
   =\left(\frac{p}{\mathfrak l}\right)_r
   \quad (\mathfrak l\nmid pr),
 \tag{MA-PF-5}
\]
with a fixed embedding $\mu_r\hookrightarrow\mathbb C$. Then,
including every ramified local factor,
\[
 \boxed{\displaystyle
   \frac{\zeta_{F_r}(s)}{\zeta_{\mathbb Q}(s)}
     =L_{K_r}(s,\psi_{p,r}) .}
 \tag{MA-PF-6}
\]
Each faithful choice gives the same induced $L$-function. In particular,
under $A$, RH for $\zeta_{F_r}$ is equivalent to RH for this single
Hecke $L$-function. The residual family in (MA-PF-3) is therefore an
explicit finite list of finite-order Hecke functions, one for each
$r\in\mathcal B(p,m,d)$.

\emph{Proof.} The Eisenstein argument in CX-ALG-024 gives
\[
 G=\operatorname{Gal}(E_r/\mathbb Q)
       \simeq\mathbb F_r\rtimes\mathbb F_r^\times,
\]
acting on the $r$ roots by affine transformations. Let $P$ denote the
permutation representation on $\mathbb F_r$ and $N$ its translation
subgroup. The character of $P-\boldsymbol1$ is $r-1$ at the
identity, $-1$ at each nonzero translation and $0$ at every element
whose multiplicative part differs from $1$. On the other hand,
$\operatorname{Ind}_N^G\psi_{p,r}$ is zero outside $N$ and its
restriction to $N$ is the sum of all $r-1$ nontrivial characters.
At a nonidentity element of $N$ their sum is $-1$ and at the identity
it is $r-1$. Thus
\[
   P\simeq\boldsymbol1\oplus
            \operatorname{Ind}_N^G\psi_{p,r}.
 \tag{MA-PF-7}
\]
The two sides are actual semisimple representations. Artin induction
formalism gives (MA-PF-6); using inertia-invariant spaces in local
Euler factors proves the equality at the ramified primes as well.
Equivalently, one must use the primitive Hecke character attached to
$E_r/K_r$, not arbitrarily remove every prime above $pr$ from its
$L$-function. Since the character is nontrivial, its Hecke function is
entire. In the open critical strip off the middle line $A$ makes
$\zeta_{\mathbb Q}$ nonzero, so (MA-PF-6) gives the RH equivalence.
This also recovers the already established identity of CX-ALG-024,
\[
 \zeta_{E_r}(s)=\zeta_{K_r}(s)
                     L_{K_r}(s,\psi_{p,r})^{r-1}.
 \tag{MA-PF-8}
\]
No general Artin-holomorphy conjecture is used.\hfill$\square$

\paragraph{Geometric formulation.}
Equation (MA-PF-7) is a decomposition of the Artin motive
$h^0(\operatorname{Spec}F_r)$ into its unit summand and the induction
of the rank-one Kummer motive over $K_r$. On
$\operatorname{Spec}\mathcal O_{K_r}[1/(pr)]$, the latter is a
finite-monodromy rank-one $\ell$-adic sheaf, for $\ell\nmid pr$,
with coefficient field containing $\mu_r$. Its good Frobenius
eigenvalues are roots of unity, hence have weight zero. The local
Euler polynomial is the exterior-power identity
\[
 \det(1-T\operatorname{Frob}\mid V^{I})
  =\sum_j(-T)^j\operatorname{Tr}
           (\operatorname{Frob}\mid\bigwedge^j V^{I}).
\]
This places the identity in the representation-ring and $\lambda$-ring
language; it does not deduce complex zero location from finite-field
weights. For finite-monodromy objects the precise foundational
dictionary is SGA~1, Expos\'e~V, \S7, opening paragraphs, original
pp.~140--141, recomposed p.~115 (arXiv:math/0206203v2, PDF p.~131).
The representation-ring use is the same one recorded in GRH\_NIGHT
CX-ALG-024 and its SGA~6 Expos\'e~V framework reference.

\subsection{A fixed strip is sufficient at the residual primes}

\paragraph{MA-PF-004 (a smaller sufficient analytic condition).}
Fix a number $\delta$ with $0<\delta\le1/2$, independently of
$p,m,d$. Assume the algebraic PF26 input and $A$. In the
large-characteristic branch, it suffices to assume
\[
 L_{K_r}(s,\psi_{p,r})\ne0
       \quad\text{for }\Re s>1-\delta,
       \qquad r\in\mathcal B(p,m,d),
 \tag{MA-PF-9}
\]
to obtain deterministic polynomial-time factorization, with an exponent
allowed to depend on the fixed $\delta$. Thus a proof of full RH for
the residual pure fields is more than this algorithm requires.
The assertion does not permit a different uncontrolled width
$\delta(p,m,d)$ for each input.

\emph{Proof.} For one residual $r$, use $K=\mathbb Q(\zeta_r)$ and
$E=K(p^{1/r})$. Equations (MA-PF-8)--(MA-PF-9) and $A$ place
all nontrivial zeros of both $\zeta_K$ and $\zeta_E$ in
$\Re s\le1-\delta$. The standard smoothed formula used and proved
in PF26 Lemma~11.2, p.~45, works for any fixed such $\delta$.
Indeed, for a fixed nonzero nonnegative
$\varphi\in C_c^\infty((1,2))$, write
$\Phi(s)=\int_0^\infty\varphi(t)t^{s-1}\,dt$ and
\[
 \Psi_D(x)=\sum_{\mathfrak p}\sum_{j\ge1}
       \log N\mathfrak p\,
       \varphi\bigl((N\mathfrak p)^j/x\bigr),
\]
where $\mathfrak p$ runs through the prime ideals of $D$ and
$\Delta_D$ is its absolute discriminant. Mellin inversion and the
smoothed explicit formula give
\[
 \Psi_D(x)=x\Phi(1)-\sum_\rho x^\rho\Phi(\rho)
                       +O_\varphi(\log\Delta_D+[D:\mathbb Q]).
\]
The zeros are counted with multiplicities. Rapid vertical decay of
$\Phi$ and the uniform local zero count give
$\sum_\rho|\Phi(\rho)|\ll_\varphi\log\Delta_D+[D:\mathbb Q]$.
Consequently
\[
 \Psi_D(x)=x\Phi(1)+
 O_\varphi\bigl(x^{1-\delta}
               (\log\Delta_D+[D:\mathbb Q])\bigr)
 \tag{MA-PF-10}
\]
for $D=K,E$. This argument uses the fixed strip, not the middle line.

The discriminant estimate in CX-ALG-024 gives
\[
 [E:\mathbb Q]=r(r-1),\qquad
 \log|D_E|\le r(2r-3)\log r+(r-1)^2\log p.
\]
Thus the degree-plus-log-discriminant quantity for either field is
$O(B^3)$, uniformly. Subtracting $r^{-1}$ times the formula for $E$
from the one for $K$ yields
\[
 \Psi_K(x)-r^{-1}\Psi_E(x)
  =(1-r^{-1})x\Phi(1)+O_\varphi(x^{1-\delta}B^3).
 \tag{MA-PF-11}
\]
Take $x=T_\delta B^{20/\delta}$ with a sufficiently large fixed
$T_\delta$. The error-to-main-term ratio is
$O_\varphi(T_\delta^{-\delta}B^{-17})$.

As in PF26 Proposition~11.3, p.~46, terms of prime-power exponent at
least two or residue degree at least two on the $K$ side contribute
at most
\[
 O_\varphi\bigl([K:\mathbb Q]\sqrt{x}\log^2(2x)\bigr).
\]
The degree-one terms above $p$ or $r$ contribute at most
$O_\varphi([K:\mathbb Q]\log(2x))$. A sufficiently large fixed
$T_\delta$ makes their sum smaller than half the positive main term
in (MA-PF-11). Every remaining degree-one prime ideal of $K$ lies
above a rational prime $\ell\equiv1\pmod r$ with
$\ell\notin\{p,r\}$. If $p$ is an $r$th power modulo $\ell$,
that prime ideal splits completely in $E/K$, and its first-power
contribution cancels exactly in (MA-PF-11). Discarding the remaining
negative terms gives an upper bound. Therefore a noncancelling
degree-one prime must occur, furnishing
\[
 \ell\ll_\delta B^{20/\delta},\qquad
 \ell\equiv1\pmod r,\qquad
 p^{(\ell-1)/r}\not\equiv1\pmod\ell.
 \tag{MA-PF-12}
\]
Search by trial division and modular powering. The invariant algebra
of PF26 Lemma~3.1 constructs $U_r\simeq\mathbb F_{p^r}$ using a
number of operations polynomial in $\ell$ and $\log p$. Apply the
transfer in MA-PF-002. Since $\delta$ is fixed, (MA-PF-12) is a
fixed polynomial bound, and all subsequent costs are polynomial.
No twist by $\mathbb Q(\zeta_{12})$ is needed.\hfill$\square$

The abstract motive and finite-\'etale dictionary is unchanged from
MA-PF-003. The difference is analytic: the prime-detection trace uses
an error $O(x^{1-\delta})$, not necessarily the square-root error.
The finite list (MA-PF-9) remains outside what $A$ alone establishes.
The quoted companion theorem (MA-PF-2) has a different, globally
quantified statement and has not been audited here.

\subsection{The first residual character is already imaginary quadratic}

\paragraph{MA-PF-005 (the cubic instance and the cutoff distinction).}
For $d=3$ in the large-characteristic branch,
\[
 \mathcal B(p,m,3)=\{3\}
 \quad\Longleftrightarrow\quad
       p\equiv-1\pmod9\ \text{ and }\ 3\nmid m.
 \tag{MA-PF-13}
\]
For these instances the only extra Hecke function in MA-PF-003 is
the order-three Kummer character for
\[
 \mathbb Q(\zeta_3,p^{1/3})/\mathbb Q(\zeta_3),\qquad
 \mathbb Q(\zeta_3)=\mathbb Q(\sqrt{-3}).
\]
It is not induced by a character over $\mathbb Q$. Infinitely many
primes meet this congruence and the exact large-characteristic cutoff
for fixed $m$ with $3\nmid m$ and $d=3$.

The concrete prime $p=269$ gives a convenient explicit character:
\[
 269^2-1=72360=3^3\cdot2680,
 \qquad \zeta_{\mathbb Q(269^{1/3})}/\zeta_{\mathbb Q}
       =L_{\mathbb Q(\sqrt{-3})}(s,\psi_{269,3}).
 \tag{MA-PF-14}
\]
It is an example of the arithmetic condition, but is \emph{not} a
case where PF26 needs an analytic assumption: for $d=3$ it has
$B=60$ and $269<60^{200000}$, hence falls in the unconditional
small-characteristic branch.

\emph{Proof.} In the large branch $p>3$. The prime $2$ is excluded.
For $r=3$, exclusion of $3\mid p-1$ leaves $p\equiv-1\pmod3$,
whose order modulo $3$ is $2$. Then
$v_3(p^2-1)>1$ is equivalent to $p\equiv-1\pmod9$.
The other exclusion is exactly $3\nmid m$. The polynomial
$X^3-p$ is Eisenstein at $p$ and has nonsquare discriminant
$-27p^2$, so its splitting field has group $S_3$ and its quadratic
subfield is $\mathbb Q(\sqrt{-3})$. Conjugation by the nontrivial
element of $\operatorname{Gal}(\mathbb Q(\sqrt{-3})/\mathbb Q)$
sends the order-three Kummer character to its inverse. These two
characters differ. A character restricted from $G_{\mathbb Q}$
would be invariant under this conjugation; equivalently a Hecke
character pulled back by the norm would be invariant. Thus this
character does not descend from $\mathbb Q$.

Dirichlet's theorem supplies arbitrarily large primes $p\equiv8$
modulo $9$. For fixed $d=3$, the function
$20+4(\lceil\log_2p\rceil+1)$ grows logarithmically; hence
$p>B^{200000}$ for all sufficiently large $p$. This proves the
infinite-family assertion. Trial division by the primes at most
$\sqrt{269}<17$ verifies primality of $269$, and the displayed
arithmetic verifies the example and its branch.\hfill$\square$

This is the concrete connection to the first non-Dirichlet imaginary
quadratic Hecke case in task~1(d). In the Artin-motive formulation it
is the two-dimensional irreducible summand of the degree-three
permutation motive of $X^3-p$. Holomorphy is already supplied by
Hecke induction; the analytic issue is its zeros. Neither the finite
monodromy nor the weight-zero statement supplies those zeros' location.

\subsection{Record of the full PF26 proof check}

The following checks record what was actually inspected. Every listed
proof was read in its entirety from the supplied PDF's extracted
text. The key interface and analytic-premise pages 35, 44 and 45
were additionally rendered and visually inspected. The source's
classical structure theorems remain cited theorems; their entire
literature was not re-refereed.

\begin{enumerate}
\item \textbf{Lemma~2.1, p.~7.} A nontrivial gcd produced over the
working extension is the product of a subset of $\mathbb F_p$-rational
linear factors, so is over $\mathbb F_p$. Polynomial degrees and
pivots reduce to scalar zero tests. A field inverse is attempted only
after every component is nonzero.
\item \textbf{Lemma~2.2, pp.~8--9.} Distinct-degree decomposition
needs no complete factorization. If $r\mid e$, the Frobenius
eigenspace $v^{Q^{e/r}}=\zeta_rv$ projects nontrivially to every
field component. The moment-curve grid has fewer than $h^2$ forbidden
values; the strict characteristic inequality supplies enough distinct
scalars. The valuation of the order of a nonzero solution is exactly
the full $r$-valuation of $Q^e-1$. Primary logarithms are digit loops,
not loops over the primary subgroup. Promise of an $r$th power makes
each logarithm divisible by $r$.
\item \textbf{Proposition~2.3, pp.~9--10.} Negating a difference
changes its primary logarithm by half the two-primary order. Each
unordered pair therefore receives one orientation. The row counts
cannot all be $(N-1)/2$ when $N$ is even, and $p>N$ preserves the
distinction after reduction modulo $p$.
\item \textbf{Lemma~3.1 and Proposition~3.2, pp.~11--13.} Fixed
functions under the index-$r$ cyclotomic subgroup form an algebra
whose geometric factors are its cosets. The verified residue test
makes Frobenius a transitive $r$-cycle, hence the algebra a field.
The tensor-product field and Frobenius eigenvector have the asserted
primary orders. The only factoring call has degree $r-1$, so the
table/factorization induction has no circular dependency.
\item \textbf{Lemmas~4.1--4.2, pp.~15--16.} Injectivity uses
injectivity of $1-\sigma$ on the infinity lattice, not on its image
in the Jacobian. In the label calculation the function action is
$h(X,Y)\mapsto h(X,\zeta^{-1}Y)$, which gives the sign in
equation~(4.7). Hilbert~90 removes the norm-one function; invariant
divisors decompose into reduced ramification coefficients and a
pullback divisor. Equal labels kill the class because the sum of
the ramification generators is zero.
\item \textbf{Lemma~4.3, Lemma~4.4 and Corollary~4.5,
pp.~17--18.} Frobenius minus one lowers the $\lambda$-filtration
by one; multiplication by $r$ lowers it by $r-1$. The binomial
induction gives the precise extension degree $r^{v_r(N)}$.
The quotient relation's second coordinate shows that fixed pairs
give genuinely fixed Jacobian points. Finite-field norm surjectivity
removes the constant obstruction to descending a divisor class and
its principal function.
\item \textbf{Proposition~4.6, pp.~18--19.} The relation needed
for the next label test is proved only after the previous labels
coincide. The last test would force
$Ne\in\lambda^{1+(r-1)v_r(N)}\Lambda$, contrary to the exact
lattice valuation. No unjustified simultaneous identities at all
levels are used.
\item \textbf{Proposition~5.1 and Lemma~5.2, pp.~20--23.}
Fractional ideals use finite quotients with explicit bounds. Their
inversion, intersections and products reduce to linear conditions.
Characteristic polynomials of multiplication by $X$ compress the
bound and compute pushforward divisors, including multiplicities.
The infinity Vandermonde argument supplies the degree bound for
the complete Riemann--Roch ansatz, and reduction gives absolute
degree at most $2g$.
\item \textbf{Lemma~5.3, pp.~23--24.} At each place a grid
contains a parameter attaining every component's least valuation.
The parameter may depend on the place, which is sufficient to
generate the product ideal; no one parameter must work everywhere.
Determinants compute the component norm without splitting the
coefficient algebra. The grid length and generator lists are
polynomial.
\item \textbf{Lemma~5.4, p.~24.} Binary addition and doubling
of reduced divisor classes stores correcting functions by circuits.
Valuations and nonzero leading coefficients propagate without
expanding the rational function. Since the final divisor has zero
coefficient at a ramification point, the final leading coefficient
is its value there.
\item \textbf{Theorem~6.1 and Lemma~6.2, pp.~25--27.} The norm
promise splits the cyclic algebra only for the existence argument.
The algorithm never uses the unknown splitting or norm solution.
The ramified and unramified local matrix models satisfy the
commutation relation with the stated sign. The local residue-root
promises follow from the norm valuation/residue identities.
Solving $Ve=he$ in a rank-one left ideal gives the norm equation
after $r$ iterations.
\item \textbf{Lemma~6.3 and Lemma~6.4, pp.~28--29.} Completed
orders descend by finite congruence data, glue to
$\operatorname{End}\mathcal V$ and have global sections as claimed.
The splitting proof is over the actual finite field. A maximal
line subbundle has quotient summand degrees no larger than its
degree, making the extension groups vanish. Evaluation gives a
full block triangular algebra. Its trace-pairing kernel is the
strictly triangular radical, whose common column kernel is the
first block; this yields the rank-one fiber idempotent by linear
algebra without diagonalizing a general algebra.
\item \textbf{Section~6.3 and Lemma~6.5, pp.~29--32.} A lift
$b$ of the fiber idempotent has constant characteristic polynomial
$T^{r-1}(T-1)$, so $b^r$ is a rank-one idempotent. Trace-dual basis
bounds control every local coefficient; the finite ansatz contains
all global sections. Precision $3M+3$ suffices after multiplication
by the prescribed powers of local uniformizers. The rational
linear solve has degrees bounded by minors, avoiding expression
tree growth.
\item \textbf{Proposition~7.1, pp.~33--34.} Normalizing the
pushforward function to residue one at a specified ramification
point changes a known norm only by an $r$th power in the constant
field. The norm-solver promise is therefore established. For the
norm-zero divisor $D-\operatorname{div}(a)$, the maximum of cyclic
partial sums solves the divisor equation exactly. Subtracting its
degree times the fixed ramification point gives degree zero and
does not change its $1-\sigma$ image.
\item \textbf{Proposition~8.1, pp.~35--37.} The binomial constant
extension has degree $r^{v_r(N)}\le N$ and supplies the entire
primary group. Every chosen division, rather than merely one
special sequence, is rational over that extension. The simultaneous
sum meets its scalar-grid bound. The reached-test invariant is
exactly Proposition~4.6's invariant. Binary infinity coefficients
are never materialized as large divisors.
\item \textbf{Lemma~9.1 and Proposition~9.2, pp.~38--40.}
The moment-curve separator has fewer than $t^3$ bad parameters
for $t$ components. Characteristic-polynomial square-freeness
detects a separating element, and substitution of a proper factor
selects a proper group of components. Recursive degrees decrease;
derivative-zero and repeated-factor branches preserve multiplicities.
The primary table is built only once.
\item \textbf{Proposition~10.1, pp.~41--44.} The displayed
norm, divisor and circuit dimensions are fixed powers of $B$;
the degree of the curve is not held constant. Moment-curve and
interpolation grids are smaller than $p>B^{200000}$. The
auxiliary algebra is charged by its numerical dimension $\ell-1$,
using elimination without an $\ell$-dependent interpolation grid.
The indexed-array estimate
$O(B^{201000}+B^{41}E^{10})$ squares to a bound covered by
(MA-PF-1). The exponent and the branch cutoffs were checked,
including their dependence on values rather than bit lengths of
auxiliary primes.
\end{enumerate}

\paragraph{Grothendieck's geometric contribution to the algorithm.}
For square-free $h$, $\operatorname{Spec}(\mathbb F_{p^m}[X]/h)$
is a finite \mbox{\'etale} scheme; irreducible factors are its connected
components and primitive idempotents are their indicator functions.
SGA~1, Expos\'e~I, Corollary~7.4, recomposed p.~9 (PDF p.~25),
gives the separable-polynomial criterion; Expos\'e~V, \S7,
recomposed p.~115, gives the finite fundamental-group-set dictionary.
PF26's norm algorithm also makes essential use of the projective-line
splitting theorem, not merely a general effectiveness slogan:
Grothendieck (1957), Th\'eor\`eme~2.1, p.~126, is its published
holomorphic source, and p.~127 discusses the algebraic version over
an algebraically closed field. The form over the original finite
field is proved explicitly in PF26 Lemma~6.3, p.~28. The global
section algebra's parabolic fiber and its rank-one idempotent are
the computational bridge. No archive cote for this new algorithm
has been identified, and none is asserted. The statements do not
turn a general existence theorem in EGA, FGA or SGA into a polynomial
algorithm without the size analysis just checked.

\section{The four-way input audit}
\label{sec:input-audit}

\paragraph{Scope and meaning of the deliverable.}
\texttt{INPUT\_CLASSIFICATION.tsv} contains one primary class for each
of the 362 rows of the supplied GRH\_NIGHT
\texttt{CONSEQUENCES.tsv}, 162 ALPOGE\_BRIDGES findings, and all 39
CX\_GRH findings described by the attached intake.  The ALPOGE set
includes the intake's full numeric ranges and contextual shorthand,
together with an explicitly marked conservative supporting block.
Question/current-gate rows are recorded separately in
\texttt{INPUT\_UNCLASSIFIED\_GATES.tsv}; a missing question statement is
not assigned a fictitious class \emph{none}.  The supplied CX intake
reports earlier proof checks, but the original CX proof package was not
available or resolved in this audit.  Those 39 classifications are
therefore explicitly \emph{intake-only}.

This is a dependency and provenance audit.  A row's original statement,
input, status, proof locator and Grothendieck locator are preserved
verbatim for the GRH catalogue.  Targeted source readings resolve the
ambiguous or additional inputs discussed below.  Other inherited proof
claims are carried as background with that exact check status; the table
does not certify a new complete reading of 563 mathematical proofs.

\paragraph{Conventions.}
The four values are exactly \texttt{A}, \texttt{M},
\texttt{outside preprint}, and \texttt{none}.  A denotes Alp\"oge's
Dirichlet-GRH theorem, used as a theorem as instructed.  M denotes the
conditional claim for every unitary Hecke character over every number
field in Masullo, \S1.2--\S1.2.3, pp.~12--16.  Its finite-order part
suffices for Brauer/Artin and Dedekind applications; potentially CM
positive-weight factors generally require infinite-order unitary Hecke
characters after the weight twist.  Theorem 13.3.1.2, p.~522, is not used
as a theorem statement.

The single-column precedence is
\[
 \texttt{M}>\texttt{outside preprint}>\texttt{A}>\texttt{none}.
\]
It is an administrative convention, not an ordering of mathematical
implication or necessity.  In particular, M does not prove PF26's
algebraic theorem.  The outside-input names and all additional premises
remain in separate columns.  A smaller exact analytic family is retained
even when M is its sufficient source.  The class \texttt{none} means no
use of A, M, or an identified substantive outside-preprint theorem in the
assertion being classified.  It does not mean no literature, no
hypotheses, or an unconditional solution of every surrounding problem.
For example NT-L15 and NT-Z19 retain the stronger weak mean-square
hypothesis, and NT-Z07 retains its independent non-CM $\mathrm{GL}_2$
GRH hypothesis even though its primary class is A.

\subsection{MA-IN-001: the three misleading leading A labels}

\paragraph{Finding.}
The full asserted scope of GRH\_NIGHT NT-Z13, NT-Z15 and NT-Z16 has
primary class M.  Their Dirichlet specializations have class A.  The
source catalogue's opening A label is weaker than the scope of the
complete row; the remainder of each original input field already
acknowledges the conditional Hecke branch.

\begin{proof}
NT-Z13 applies its argument both to zeta/Dirichlet functions and to
finite-order Hecke characters over general number fields.  Its stated
analytic hypotheses include the real or conjugate-paired gamma
parameters; the arbitrary infinite-order case is explicitly excluded.
NT-Z15 applies the central-multiplicity bound to the same Hecke scope.
NT-Z16 invokes the NT-Z13 analytic hypotheses for its first
\emph{nonnegative} ordinate bound.  In all three, the zero-location
input for the general Hecke branch is M's finite-order assertion.
A's exported Dirichlet assertion does not itself supply that branch.
Thus the union of actually asserted cases has primary class M under the
declared rule.  Direct inspection of the complete source input fields
shows that M is mentioned after the opening A text, so the correction
is to a misleading leading label rather than an allegation that the
source completely omitted M.  NT-Z14 inherits the same distinction and
is also classified M, without the misleading-leading-label flag.
\end{proof}

In geometric language, a finite-order Hecke character gives a rank-one
local system over its own number field.  Induction to $\mathbb Q$ need
not split into rank-one Dirichlet systems.  This is the Artin-motive
scope distinction used in GRH\_NIGHT AG-ARTIN and MT-ART01, and in
ALPOGE\_BRIDGES CC-F03; the finite-\'etale representation dictionary is
SGA~1, Expos\'e~V, \S7.  It is not a logical independence assertion
between the two GRH statements.

\subsection{MA-IN-002: the PF26 input in AG-CYCLIC-DIV}

\paragraph{Finding.}
GRH\_NIGHT AG-CYCLIC-DIV has class \texttt{outside preprint}, with
PF26 named.  Its original word ``neither'' correctly excludes A and M
but is insufficient as a four-way input label.  The full original input
field already names PF26; no undisclosed-dependency allegation is made.

\begin{proof}
The supplied geometry proof invokes PF26's odd-splitter construction in
\S4: Lemma~4.1, p.~15; Lemma~4.2, pp.~15--16; Lemmas~4.3--4.4,
p.~17; Corollary~4.5, p.~18; and Proposition~4.6, pp.~18--19.
These are affirmative algebraic inputs to the proof, not papers merely
named as targets of criticism.  No A/M zero-location input is needed for
this supplied algebraic step.  Consequently the required class is
\texttt{outside preprint}.  The separate full PF26 reading in MA-PF-001
checks this source chain; this documentary argument neither duplicates
that proof nor turns PF26 into a formally certified theorem.
\end{proof}

The geometric content is computation of a nontrivial idempotent in a
finite \'etale algebra, hence a decomposition of a finite \'etale
scheme (SGA~1, Expos\'e~V, \S7).  GRH\_NIGHT CX-ALG-027 and
CX-ALG-028 need both the algebraic PF26 input and an analytic family.
They have primary class M at full generality, with PF26 and the precise
A/pure-field premises explicitly retained.  When the bad-prime family
of MA-PF-002 is empty, the latter specialization uses PF26+A; the table
does not conceal this specialization under universal M.

\subsection{MA-IN-003: abelian motives and the named 2026 inputs}

\paragraph{Finding.}
The cp03 ALPOGE intake's ``none''/``unconditional'' description of the
full abelian-sector conclusions CC-AG01--03 and CC-XC01--04 must retain
the outside-preprint input van Hoften--Sempliner,
Hexagon~2610.00259v2.  The later intake's explicit preprint qualification
is the correct provenance.  CC-TATE03 has a second distinct input,
Hexagon~2610.00226v1.

\begin{proof}
The supplied module \texttt{consequences\_criteria\_Tate2026.tex}
states the chain explicitly.  CC-TATE01 imports van Hoften and
Sempliner's \emph{The Hodge Conjecture for CM abelian varieties},
Theorem~I, p.~1, version~2 of 18 October 2026.  CC-TATE02 applies
Milne's published 1999 Theorem~7.1, journal p.~72, and Propositions
7.3--7.4, pp.~73--75, together with his published 2002 Theorem~3.3
and Corollary~3.4, final author-manuscript pp.~8--9.  Those classical
implications assume precisely the all-CM Hodge input.  The abelian
numerical-motive classifications and arbitrary-abelian forward-cycle
targets in the flagged rows use this conclusion.  Neither A nor M
supplies the algebraic cycle theorem.

CC-TATE03 separately imports \emph{Milne's rationality conjecture},
Hexagon~2610.00226v1, contributed by Keerthi Madapusi and bearing the
PDF byline ``ChatGPT and Claude'': Theorem~1.1, p.~1;
Theorem~5.1, pp.~11--12; Proposition~A.1, pp.~13--14; and
Corollary~6.2, p.~13.  The rationality assertion uses this manuscript;
the simultaneous algebraic cycle on the special fibre additionally
uses the all-CM theorem.  Both are retained.  Milne's
\emph{On the Hodge, Tate, and standard conjectures for abelian
varieties}, Hexagon~2611.00244v1, Theorems~2.1, 3.1 and 3.5,
pp.~4, 6 and 8, is recorded as the cited 2026 corroborating note.
It is not counted as an indispensable third theorem where the proof
already cites the original 1999/2002 implications.  This establishes
the precise dependency correction.
\end{proof}

The relevant Grothendieck question is the passage from Frobenius
eigenvalues to actual motivic objects and correspondences in cote~15,
scans~100, 102, 104 and 106, as located in ALPOGE\_BRIDGES
CC-TATE04.  A spectral coincidence is not itself an algebraic
correspondence.  The table also preserves the important narrower
exceptions: CC-XC20--22 use their explicit Tate/supersingular Chow
decompositions and published Fu--Li results without the all-abelian
preprint.  BR-15-CLASSBOUND and BR-15-SLOPEBOUND use classical
Honda--Tate and ideal theory.  They are not assigned the broader input
merely because they occur in the same programme.

\subsection{MA-IN-004: a single class preserves a conjunction only with its notes}

\paragraph{Finding.}
The table's single primary label is compatible with all required
conjunctions and conditional subcases precisely because the secondary
inputs and assertion scope are retained.  A formal equivalence and an
application asserting its positive or zero-location side receive
different descriptions.

\begin{proof}
For a row let $S$ be its set of identified A/M/outside-theorem inputs
and $H$ its remaining hypotheses.  The primary label is the maximum of
$S$ under the declared administrative precedence, or \texttt{none}
if $S$ is empty.  The stored original input, outside-name list and
residual notes retain $S$ and $H$; thus the projection is never used to
replace the conjunction by its maximum.  In particular the projection
does not imply $\mathrm{M}\Rightarrow\mathrm{PF26}$.

For the formal criteria MT-M19, MT-M20 and MT-M27, the proof establishes
an equivalence without asserting that its GRH/positivity side holds.
These rows have class \texttt{none}; the A/M application is stated
separately, and simplicity remains additional in MT-M20.  In contrast,
AG-DET, AG-MOMENT, AG-HANKEL and AG-SIGNED explicitly include the
identified arithmetic Hecke/CM determinant or quotient application.
Their positive real squared-zero support uses M at full scope, or A
in the Dirichlet specialization.  Their elementary entire-function,
moment or Gram-matrix cores can be used with independently supplied
positive support, but that observation does not erase the arithmetic
input of the full row.  Likewise MT-M26 retains M for its asserted
every-positive-width narrow-strip application, whereas its wide-strip
core is unconditional.  These examples establish that the convention
does not convert an arithmetic zero-location premise into ``none''
merely by calling it a hypothesis.
\end{proof}

For cote~1, scans~79--80 and 111--112, the distinction is between the
general Fredholm/trace formalism and positivity of the specific
arithmetic operator $B$.  The source framework alone does not supply
its positive spectrum.  ALPOGE\_BRIDGES BR-TR-01--02 and
GRH\_NIGHT MT-M20/M23 preserve the further distinction between positive
traces and faithful positive full-jet operators.

\subsection{MA-IN-005: the coverage certificate}

\paragraph{Finding.}
There are 563 distinct classified run/ID pairs.  All 362 supplied
GRH catalogue rows occur exactly once, with their five original source
fields unchanged.  Every primary class is one of the required four
values.  The independent range/context audit passes.

\begin{proof}
The deterministic builder uses explicit reviewed ID assignments, not
an A/M keyword rule.  It refuses duplicate assignments or a mismatch
between the source catalogue and the assignment set.  The independent
script \texttt{checks/audit\_input\_classification.py} reads the emitted
TSV and original source files without importing the builder's decisions.
It verifies equality of the 362-ID sets, uniqueness of run/ID pairs,
exact preservation of statement/input/status/proof/Grothendieck fields,
the permitted class values, nonempty evidence/status/residual fields,
existing evidence files and named outside inputs.  It separately
expands the intake's numeric ranges (126 finding IDs), checks 22
contextual nonnumeric IDs, verifies the 39-entry CX finding set, and
requires all six CX gate IDs in the separate gate register.
The emitted certificate \texttt{checks/input\_independent\_audit.json}
records PASS and the table's SHA-256.  This proves the finite data
claims just stated; it does not prove the inherited mathematics.
\end{proof}

\begin{center}
\begin{tabular}{lrrrrr}
Run & A & M & outside preprint & none & total\\
\hline
GRH\_NIGHT &66&140&3&153&362\\
ALPOGE\_BRIDGES &25&0&50&87&162\\
CX\_GRH &8&14&5&12&39\\
\hline
Total &99&154&58&252&563
\end{tabular}
\end{center}

\subsection{MA-IN-006: other preprints, published sources and refutation targets}

\paragraph{Finding.}
The outside-preprint class must retain substantive imported manuscripts
beyond the four names listed in the prompt.  Publication status and
logical use both matter: an arXiv URL alone does not determine the class,
and refuting a manuscript does not import its theorem as a premise.

\begin{proof}
Direct dependency readings identify Engel--de Gaay Fortman--Schreieder,
arXiv:2507.15704v3, Theorem~1.1 and Corollary~1.2, p.~2, in AG-DP-EGS
and the very-general even-degree branch of AG-FOURIER-ISOGENY.
Their institution lists the work as a preprint.  CC-XC23 imports
Fu--Huang--Li, arXiv:2609.38801v2, Theorem~1.9, p.~5, with its
excellent-lifting, characteristic and Artin-invariant restrictions.
CC-WM01 imports Ito--Zou, arXiv:2610.07890v1, Theorem~B, p.~2,
and Wear, arXiv:2303.05610v1, Theorem~1.1.2, p.~3, in two of its
branches.  Wear's author page still lists the paper as submitted.
CC-WM02--03 import Petrov--Zavyalov,
arXiv:2607.14517v2, Theorem~1.2, p.~2, equivalently
Theorem~5.1 and Lemma~5.2, p.~18.  F-SM-32--33 use
Ramzi--Sosnilo--Winges, arXiv:2503.11338v1,
Definition~2.5 and Theorem~2.6(1),(4), pp.~5--6; the author lists
acceptance in JEMS, but an issued publication was not established in
this audit.  The table retains it as an accepted preprint with that
qualification.  CX-E01--05 use the separate OpenAI result E,
\emph{The Exponential Time Hypothesis is False}, dated 25 October 2026,
Theorem~1.1 and Corollary~1.4 as reported by the intake.  The official
repository locates it among preprints.  These are substantive premises
and are therefore named.

Conversely, primary publisher records establish that Deninger's
arXiv:2504.15767 result is published in \emph{Journal of Number Theory}
286 (2026), 47--66; Haine--Porta--Teyssier's exodromy paper is published
in \emph{Journal of Topology} 19(3) (2026); and the partial-regularity
theorem used by F-SM-15 is published as Chen, Lau, Lee, Ono and Zhang,
\emph{Archiv der Mathematik} 127 (2026), 165--178.
These are retained as published inputs.  Finally BKNO
arXiv:2508.17776, Corollary~3.2 and Proposition~4.17, is the target
of the explicit counterexample in AG-LOCAL-EXCEPTION; Thakur's
Proposition~2.9 is the target in BR-15-CLASSBOUND.  The counterexamples
do not assume the assertions they refute.  This establishes all the
claimed distinctions; exact URLs and dated evidence appear in
\texttt{modules/inputs\_source\_status.tsv}.
\end{proof}

\subsection{MA-IN-007: the exact Alp\"oge locator and its inspection boundary}

\paragraph{Finding.}
The supplied runs support the following precise citation for A:
Alp\"oge, \emph{Remark}, public four-page announcement,
\url{https://alpo.ge/remark.pdf}, p.~4, step~20; the reported interface
is \texttt{ALPOGE\_Main.lean}, with definitions in
\texttt{ALPOGE\_Statements.lean}.  This audit uses A as a theorem and
makes no fresh Lean-compilation claim.

\begin{proof}
The public PDF's p.~4 identifies step~20 as GRH.  The predecessor's
\path{research/source_and_masullo.md}, \S1.1, records its complete
reading of the following exported declarations:
\begin{verbatim}
riemannHypothesis : RiemannHypothesis := rhOfAll
riemannZeta_eq_zero_re : RiemannZetaZerosOnLine := rhOfAll
LFunction_ne_zero_of_half_lt_re
  : DirichletLNonvanishingRightOfHalf := allChi
LFunction_eq_zero_re : DirichletLZerosOnLine := claimAll.2.1
LFunction_eq_zero_re_of_isPrimitive
  : PrimitiveDirichletLZerosOnLine := critAll
\end{verbatim}
The same record identifies the pin
\texttt{0df444a360eaa60ab8c11dca51a86af692955474}
and the import \texttt{Supplements.Kernels.MainTheoremsAll}.
Its definitions establish the rational Dirichlet scope; no exported
arbitrary-number-field Hecke assertion is recorded there.
The present attachment contains the predecessor's inspection record,
not those Lean files or their full import tree.  Thus the exact filenames
and declarations can be cited as predecessor-inspected interface
evidence, while the public step locator was checked directly.  No new
formal-certificate event follows from either fact.
\end{proof}


\part{Unproved assertions and remaining research gates}
\section{Status boundaries}
The following are open assertions, not unchecked proofs presented
as findings. The exact evolving rows and tried approaches are
maintained in \texttt{QUESTIONS.tsv} and \texttt{LEDGER.md}.

\paragraph{Q-MA-GEOMETRIC.}
The literal rate-independent source comparison is refuted by
MA-JP-001--003: an allowed joint path has a nonzero actual Gram
remainder, with both polar moments removed. The separate canonical
bounded inverse is refuted by MA-MO-001. MA-MO-002--003 also
refute a bounded raw Euler-conjugated cutoff on every narrow
one-sided zeta or Dirichlet Hilbert line, including the actual
rational smooth theta image for its specified multiplier
continuation. MA-MO-004 gives the exact original/continued polar correction.
MA-MO-005 then proves unboundedness of the original raw first
diagonal on genuine smooth theta inputs with both polar moments
zero and smooth final outputs. A bounded extension of the full
Hilbert direct-sum matrix is therefore impossible on that core.
MA-MO-007 refutes invariance of the entire original smaller
common-sum smooth core, and MA-MO-008 proves failure of a
second raw iteration even after both first-output poles are
killed. MA-MO-010 constructs the actual once-defined smoothed
map on that pole-free domain in the fixed-gap smooth and
unconditional generic-width Gevrey regimes. MA-MO-011 identifies
its closure as a codimension-two subspace of the theta ideal;
every nonzero admissible smoothing has a surjective outgoing
two-moment map. Its minimal rank-two orthogonal correction is
a changed action. MA-MO-013 shows that this closed domain has
no nonzero translation-invariant subspace, and hence cannot
receive a nonzero translation-equivariant intertwiner.
These are exact domain obstructions, not a claim about every
alternative nonlocal map. MA-MO-009 separately extends the
raw first-diagonal obstruction to nontrivial imaginary-quadratic
Hecke characters, including the non-induced cubic example,
using literal smooth ideal sums and no polar subtraction.
MA-MO-012 proves an unconditional inverse-domain lemma for
every smooth ideal vector of every nontrivial finite-order
character. The complete original common-sum smoothing is then
an actual trace-class endomorphism of the closed theta ideal
in the generic Gevrey or fixed-gap smooth regimes. This is
a genuine ideal-domain construction; its ambient trace contains
the orthogonal ideal projection. MA-MO-014 compares its actual
polynomial-weight realizations and gives finite theta Gram
evaluation with an a posteriori certificate; MA-MO-015 gives
its explicit Blaschke/logistic kernel correction.
MA-MO-017 now supplies a genuine nuclear return to the
original local source and evaluates that projected ideal
trace by the local Fourier multiplier. MA-MO-019 evaluates both
one-sided traces at that same value and thereby proves cancellation
of the explicit Blaschke correction. The projection remains in the
argument: removing it would double a generally nonzero trace.
MA-MO-016 separately constructs a rational
polar-retaining common-space completion and equivariant exact
sequences. Its actual positive source column is trace class
into the common completion, with an explicit rank-at-most-two
outgoing quotient defect. MA-MO-017 also evaluates the rational
local source endomorphism on the positive completion; this
is not the full raw matrix defect on its polar columns.
MA-MO-018 extends the full original rational matrix to a trace-class
endomorphism of the augmented common completion under an explicitly
used zero-free right half-plane with a margin. Both original
inverses are identified on the positive narrow line and both polar
coordinates are retained. MA-MO-020 proves that its outgoing
positive-source quotient defect has exact rank two for every
nonzero compact test, even when both polar Mellin moments vanish.
MA-MO-021 constructs the actual invariant relative operator
\(C_g(1-P_+)\), evaluates its trace along the positive-source
exact sequence, and proves its Fredholm determinant factorization.
The full raw matrix trace is exactly the omitted finite-prime
sum for value-zero tests under the stated zero-free premise.
This does not define an induced raw-matrix quotient endomorphism.
MA-MO-023 strengthens the polar obstruction without a containing
half-plane premise: the known critical-line zeros alone force
outgoing rank two at every positive narrow width, including
the unconditional generic Gevrey source-column regime.
MA-MO-022 supplies the full nontrivial-character common matrix
under the explicit dual half-plane premise. Its actual ideal
restriction is the local source endomorphism, its quotient
action is minus convolution, and its trace is the outside-$S$
prime distribution for value-zero tests. The structural-$S$
condition handles general class groups and unit rank. A supplies
this premise for Dirichlet characters; the non-induced example
retains its own Hecke RH condition or M.
MA-MO-024 proves the actual converse obstruction: a strictly
exterior character zero prevents bounded full-source smoothing
for every nonzero compact test, even one vanishing at that
zero in Mellin transform. The every-width bounded-extension
criterion is equivalent to own-character RH; it is not a new
arithmetic exclusion. MA-MO-025 includes every deficient
boundary multiplicity, with the threshold equality proved by
weak convergence of actual orbit integrals. MA-MO-026 proves
the same exterior and boundary necessity on the rational
polar completion. Full retained multiplicity is necessary;
sufficiency without an interior margin is not inferred.
MA-MO-027 proves full-matrix sufficiency when there are finitely
many boundary zeros, every boundary multiplicity is fully retained,
and all other zeros lie a fixed distance inside the strip. Its
boundary inverse is the actual right-hand distributional limit of
the original wide inverse, with the exact Sobolev threshold.
MA-MO-028 separately removes the earlier generic/Gevrey and
fixed-gap/smooth restrictions from the actual completed-ideal
nuclear return and its local trace: these hold for all compact
smooth tests at every positive narrow width, including widths
meeting zeros. It also gives the rational bounded cutoff moments,
codimension-two closure and actual once-defined map into the ideal
at all these widths. The earlier singular boundary-line Riesz
formula is not extended. Neither unrestricted one-sided
boundedness nor the full no-margin matrix theorem follows.
MA-MO-030 proves the full matrix when the total multiplicity
in fixed-size near-edge rectangles is uniformly bounded by
$M$ and $s>2M-1$. Infinitely many boundary or approaching
zeros are allowed under that explicit condition; growing
uncontrolled near-edge counts remain outside the theorem.
MA-MO-031 proves that polynomial reciprocal growth is
equivalent to unrestricted one-sided bounded smoothing for
the stated universal class of small real cutoff perturbations.
MA-MO-032 proves the corresponding criterion with the physical
cutoff fixed and all tests varying on the single support $[1,2]$.
The conditional failure branch has genuine normalized core
countervectors and unbounded finite-operator norms despite
convergent traces. $A$ supplies the rational positive branch;
no zero-location-free arithmetic bound is inferred.
MA-MO-033--034 give a boundary-uniform reciprocal necessity for
the actual full common source, with the rational negative column
and both polar coordinates treated explicitly. A hypothetical
very tight interior cluster would violate it, producing genuine
compact-source witnesses and residual bad tests; no such
arithmetic configuration is asserted. MA-MO-035 proves full
trace-class sufficiency under a polynomial cluster-depth budget,
with a strict total-boundary-multiplicity threshold and the
critical square-root logarithm kept separately. These results
do not decide the general arithmetic no-margin question.
MA-MO-036 gives the exact analytic equivalence: under the
closed-strip premise, strong localized $H^{-s/2}$ right-boundary
limits with polynomial window norms are equivalent to universal
boundedness on the specified compact test supports, and to the
actual full trace-class matrix for every compact test. Both dual
sectors are retained; in the rational case necessity already
follows from the negative columns into the polar completion.
Full individual boundary retention follows. The arithmetic
polynomial norm estimate and FT002 remain unresolved.
MA-GA-001 separately completes the actual
general-field annular-rate and nuclear-source-lift theorem on
the inherited flat LF space, with unchanged physical norm
cutoffs. That local theorem does not supply the global FT002
deficiency condition or every original physical projection
factorization.
MA-PI-001--004 independently decide the canonical theta-map
identification: the original physical Fourier cutoff does
not intertwine the global orthogonal PW projection. The
proved witnesses include both $\chi_4$ and the principal
rational sector, in the corrected symmetric realization
and the literal printed twisted norm. In the principal
case the final flat input and its cutoff output both
satisfy the two polar conditions. The non-induced cubic
character over $\mathbb Q(\zeta_3)$ is also covered,
with the different place retained and both finite Fourier
shells kept in the physical cutoff. The source failure is
an actual nonzero Gram pairing, before taking any trace,
at every sufficiently large finite level. This refutes
that precise canonical map, not all possible intertwiners;
the previously constructed Gram corrections remain valid.
Conditionally on a zero on the chosen line, MA-MO-006 proves
that even nonzero compact left smoothing cannot make the
specified sequential multiplier composition closable on its
zero-avoiding compact-frequency core. This is a domain theorem
under that explicit hypothesis, not evidence that an off-critical
zero exists or a statement about every geometric domain.
These precise failures do
not assert that every possible prime-bearing map is impossible.
MA-MEY constructs the actual common-space prime correction;
MA-PR supplies actual partial nuclear returns and the exact
boundary-block return obstruction. MA-EG makes the unchanged
full narrow identity equivalent to retaining every zero
multiplicity in the actual quotient.

MA-PR-006--007 close the regularized return problem even at
exceptional widths for every locally eligible test, using the
actual whole-quotient block costs. MA-PR-009 gives the exact
finite-derivative-seminorm criterion for all compact eligible
tests. When there is no proper primary truncation, MA-PR-010
reduces unrestricted all-test returns to polynomial growth of
one canonical sequence of actual block inclusion norms.
MA-PR-011 adds the necessary absence of proper truncation to give
the unrestricted all-test characterization at every width.
MA-PR-012 bounds the actual whole-quotient cost by a complex-distance
edge potential. Its explicit local count conditions suffice, but
the ordinary count gives only a logarithm-squared loss.
MA-PR-013 supplies a lower bound from a whole tight cluster,
using the original block norm, repeated nodes and every oblique
projection cost. Neither adverse arithmetic clusters nor a
polynomial arithmetic potential budget are asserted.
These arithmetic norm estimates remain unproved at a general
width; regularized returns do not supply the original
common-space off-diagonal domains.

The finite global PW range, its two actual Gram factors, and its
regularized cyclic operator are now constructed by MA-GP and
MA-CGP. MA-GW-001--004 prove the actual finite quotients at every
polynomial weight, all edge ranks, ambient trace class, and both
iterated limits. MA-SB-005--006 prove the causal slow-band joint
limit in the original physical Hilbert metric. The exact paired
local/global correction is the upper exterior divisor plus the
unretained upper-edge lengths. MA-JP-004 proves that vanishing
on every causal compact test, even after both polar moments are
removed, is equivalent to emptiness of this deficiency divisor.
Thus this route no longer has an unproved causal cutoff estimate.
Its remaining full-identity gate is the independent absence of
that divisor, not the existence of a trace-class transport.

MA-AT-001 proves ambient trace class and a uniform cutoff error
for every compact smooth test. MA-AT-002 computes the negative-time
finite-level limit at $s=0$ by actual trace-norm convergence of
an adjoint defect; MA-AT-003 extends that convergence to every
polynomial weight using a uniform actual smoothing estimate.
Its lost lower zeros have reflected exponents. MA-NZ-001--002
now remove the norm-one neighbourhood in uniform actual trace norm,
so these limits extend to the complete requested class $g(0)=0$.
MA-FT-001 combines these results into the full actual slow-band
formula, including the reflected negative-half-line lower correction.
MA-FT-002 proves that its agreement with the local geometry for all
requested tests is exactly the same upper-deficiency emptiness criterion.
MA-LA-015 removes the formerly fixed-local-profile restriction
using actual trace-norm comparison of arbitrary thin transitions.
MA-GP-011 chooses one Gram regularizer along a prescribed path
for every compact smooth test. Neither changes the arithmetic gate.
MA-TN-001--002 identify convergence of the actual uncut ambient
convolutions in trace norm with the same full-retention condition.
Their exact positive-resolvent trace-norm defect is proved,
including the finite endpoint and weight-cusp contributions.
This equivalence adds a topology criterion; it is not an
independent proof of arithmetic divisor vanishing.
A one-sided quotient is not silently given a negative-time group action.
The numerical rational and Dirichlet identity is known from $A$;
using it to close the independent-method gate would use zero location.

\paragraph{The $S_5$ branch is closed under its stated hypothesis.}
Q-MA-S5 is proved by MA-FQ-004 under $(\mathcal R_B)$.
The earlier central-order and common-zero signatures remain
valid conditional deductions, but their positive-excess
antecedent is excluded by the later Fourier-rigidity argument.
MA-RB-001 excludes explicit discriminant-dependent endpoint
intervals for real zeros of $B$ under $\mathsf A$.
MA-RB-002 gives a sharper criterion for the comparison itself
without $\mathsf A$, $M$, or $(\mathcal R_B)$: failure forces a
common real zero of $B$ and $L_{V\eta}L_X$ in
$(1/2,1-4/\mathcal H_L]$, with multiplicity at least nine in
$\zeta_L$. MA-RB-003 proves boundary positive definiteness, but
also proves that the actual central prime measure in this branch
is not tempered; the ordinary Fourier identity of MA-FQ-006 is
therefore unavailable. This shared-zero exclusion remains open.
MA-RB-007 retains the reflected real zero, all sixty complex
gamma terms and the known prime-five contribution. Its attained
one-variable minimum $J_{\Delta_L,1}$ confines the residual
point to $(1/2,1/2+\sqrt{J_{\Delta_L,1}}]$ when that minimum
is positive; a nonpositive minimum excludes comparison failure.
MA-RB-008 additionally retains a uniform lower bound for every
other prime, including ramified primes, and proves that its
discriminant-aware minimum $\mathcal J_{\Delta_L,1}$ is strictly
smaller than $J_{\Delta_L,1}$. Its residual interval is
$(1/2,1/2+\sqrt{\mathcal J_{\Delta_L,1}}]$ when the minimum
is positive. MA-RB-009 constructs genuine fields in the specified
family with discriminants tending to infinity and proves that
this earlier discriminant-only minimum tends to $1/4$ on them.
That scope result supplies no positive-excess example and does
not by itself analyze the new relative envelope below.
MA-RB-012 now takes a finite envelope over the actual entire
relative functions $\zeta_L/\zeta_{L^H}$, for nontrivial
subgroups $H\le S_5$. It proves
$\mathfrak J_{L,k}<\mathcal J_{\Delta_L,k}$ without
$\mathsf A$, $M$ or $(\mathcal R_B)$. The current residual
interval is $(1/2,1/2+\sqrt{\mathfrak J_{L,1}}]$ if this
minimum is positive; a nonpositive minimum excludes failure.
This is a strictly smaller gate, not a general numerical
exclusion. The proof retains all ramified factors, handles a
possible minimum at one, and gives finite Euler error bounds.
The quadratic contribution is subtracted with its actual order,
so no noncoincidence of zeros is required.
MA-RB-010 identifies a compatible residual realization in
any one odd characteristic as an equivalent comparison gate;
characteristic two already has such a realization without
comparison. The inspected nonregular construction theorems
assume the missing odd-characteristic compatibility. This
source-scope result is not a smaller arithmetic gate.
MA-RB-011 proves ramified finite-Artin conductor and Euler
comparisons and a strict limitation of the five-factor leading
near-one Rankin score. It does not order all exact Euler--gamma
variational bounds or shrink the residual arithmetic interval.
The stronger hypothesis $(\mathcal R_B)$ remains unproved on
the remaining central interval for every extension in the family.
The input $M$ supplies it and also the stated all-factor GRH
consequence. No unrestricted $S_5$ assertion is inferred.

\paragraph{Q-MA-HECKE.}
The full imaginary-quadratic character explicit formula and
the stated local and ideal-summation deductions are established
in the character section. MA-CQ-001--004 construct the actual
Dirichlet and cubic character local quotient, dual Fourier map,
finite-cutoff nuclear lift and coupled trace limit, with
conductor and different factors retained. MA-CGP-001--003
identify the actual global primitive character PW ranges,
conductor-scaled support, Hardy divisor and Gram trace transport.
MA-GW-004 and MA-SB-006 now give the actual all-weight finite
character trace and causal slow-band comparison with its exact
deficiency divisor. MA-NR-001 records the coherent additive-sign
repair needed for global Poisson. Replacing its full zero trace by every
narrow trace independently of zero location remains open.
$M$ is a sufficient conditional input; $A$ supplies no general
replacement for this non-Dirichlet Hecke premise.

\paragraph{No unfinished proof is silently included as checked.}
Part~I includes proved conditional implications as such. Their
unproved antecedents are not counted as established results.
Any subsequent draft or approach not passing a complete proof
reading will be placed in this part or in the explicitly marked
research notes, with its precise gap.

\section{Primary source references not repeated in the modules}
Alp\"oge, L. (2026). \emph{Remark} [Public announcement].
Step~20, p.~4. \url{https://alpo.ge/remark.pdf}.
Lean-interface provenance: MA-IN-007.

Masullo. (2026). \emph{The music of the primes} [Unrefereed draft].
The source attribution follows the supplied request; this audit
does not infer a fuller author name from an unverified source.
\url{https://math216.wordpress.com/wp-content/uploads/2026/11/the-music-of-the-primes.pdf}.
Exact statements and equations are cited throughout.

Meyer, R. (2005). On a representation of the id\`ele class group
related to primes and zeros of $L$-functions.
\emph{Duke Mathematical Journal, 127}(3), 519--595.
\url{https://doi.org/10.1215/S0012-7094-04-12734-4}.
Author manuscript: \url{https://arxiv.org/abs/math/0311468}.
The summation and common-space comparison are Lemmas 5.3--5.4
and equation (31), author pp.~39--41, as checked in MT-M22/24.

Grothendieck, A. (with Raynaud, M.). (1971).
\emph{Rev\^etements \'etales et groupe fondamental (SGA~1)}
(Lecture Notes in Mathematics, Vol.~224). Springer.
Expos\'e~I, Corollary~7.4; Expos\'e~V, \S7.
\url{https://arxiv.org/abs/math/0206203}.

Berthelot, P., Grothendieck, A., \& Illusie, L. (Eds.). (1971).
\emph{Th\'eorie des intersections et th\'eor\`eme de Riemann--Roch
(SGA~6)} (Lecture Notes in Mathematics, Vol.~225). Springer.
Expos\'e~V, \S\S2, 7, especially (7.1.1), via the exact
predecessor placement AG-LAMBDA and MT-ART11.

\end{document}
